\documentclass[11pt]{article}
\usepackage[a4paper,margin=1.45cm]{geometry}
\usepackage[T1]{fontenc}
\usepackage{lmodern}
\usepackage{fontspec}
\usepackage[french]{babel}
\usepackage{amsmath,amssymb,mathtools,mathrsfs,array,booktabs,pdflscape,graphicx,adjustbox,diagbox,longtable}
\DeclareFontFamily{U}{rsfs}{\skewchar\font127}
\DeclareFontShape{U}{rsfs}{m}{n}{<-6.5>rsfs5<6.5-8>rsfs7<8->rsfs10}{}
\usepackage[protrusion=true,expansion=false]{microtype}
\setlength{\parindent}{1.25em}
\setlength{\parskip}{0.10em}
\allowdisplaybreaks
\setlength{\emergencystretch}{12em}
\sloppy
\hbadness=10000
\vbadness=10000
\hfuzz=7pt
\vfuzz=10pt
\raggedbottom
\newcommand{\widebar}[1]{\overline{#1}}
\newcommand{\A}{\Delta}
\newcommand{\D}{\Delta}
\newcommand{\R}{\varrho}
\newcommand{\hata}{\widehat{a}}
\newcommand{\hatv}{\widehat{v}}
\newcommand{\modu}[1]{\;\operatorname{mod} #1}
\newcommand{\mcell}[1]{\ensuremath{\begin{gathered}#1\end{gathered}}}
\newcommand{\pair}[2]{(#1\,|\,#2)}
\newcommand{\eps}{\varepsilon}
\newcommand{\rhoidx}{\varrho}
\providecommand{\mat}[1]{\begin{pmatrix}#1\end{pmatrix}}
\providecommand{\eqzero}[1]{\equiv 0\,(#1)}
\providecommand{\Norm}{\operatorname{N}}
\providecommand{\Gg}{\mathfrak G}
\providecommand{\Hh}{\mathfrak H}
\providecommand{\Ii}{\mathfrak I}
\providecommand{\PP}{\mathfrak P}
\providecommand{\oo}{\mathfrak o}
\providecommand{\Sp}{\operatorname{Sp}}
\makeatletter
\renewcommand\footnotesize{\@setfontsize\footnotesize{7.5}{8.7}\selectfont}
\makeatother

\providecommand{\Kfield}{\mathfrak{K}}
\providecommand{\Ssys}{\mathfrak{S}}
\providecommand{\Lsys}{\mathfrak{L}}
\providecommand{\Jdom}{\mathfrak{J}}
\providecommand{\Gdom}{\mathfrak{G}}
\providecommand{\Hgrp}{\mathfrak{H}}
\providecommand{\GaloisRes}{\Phi}
\providecommand{\pol}[2]{\left(#1\frac{\partial}{\partial #2}\right)}
\providecommand{\Deltaop}{\Delta\!\left(\frac{\partial}{\partial x},\frac{\partial}{\partial y},\frac{\partial}{\partial z}\right)}
\providecommand{\Omegaop}{\Omega}
\providecommand{\Afield}{\mathfrak{A}}
\providecommand{\Bfield}{\mathfrak{B}}
\providecommand{\Cfield}{\mathfrak{C}}
\providecommand{\Rfield}{\mathfrak{R}}
\providecommand{\Sdomain}{\mathfrak{S}}
\providecommand{\Xreal}{\mathfrak{X}}
\providecommand{\Yreal}{\mathfrak{Y}}
\providecommand{\bdelta}{\bar{\delta}}
\RequirePackage{tikz}
\newcommand{\NoetherRosetteGlyph}{%
  \begin{tikzpicture}[x=0.1363584905660377pt,y=-0.1363584905660377pt,baseline=(current bounding box.south)]
    \path[use as bounding box] (0,0) rectangle (108,108);
    \fill[even odd rule]
    (33.5,84) -- (32.5,83) -- (30.5,83) -- (29,81.5) -- (29,76.5) -- (30.5,75) -- (31.5,75) -- (32.5,74) -- (34.5,74) -- (36,72.5) -- (36,71.5) -- (35.5,71) -- (33.5,71) -- (33,70.5) -- (33,69.5) -- (30.5,67) -- (28.5,69) -- (26.5,69) -- (24.5,71) -- (23.5,71) -- (22.5,72) -- (21.5,72) -- (20.5,73) -- (18.5,73) -- (17.5,72) -- (16.5,73) -- (16,72.5) -- (16,68.5) -- (18.5,66) -- (19.5,66) -- (20,65.5) -- (20,64.5) -- (26,58.5) -- (26,56.5) -- (24,54.5) -- (24,53.5) -- (17,46.5) -- (17,45.5) -- (16.5,45) -- (15.5,45) -- (15,44.5) -- (15,42.5) -- (15.5,42) -- (17.5,42) -- (19.5,44) -- (20.5,44) -- (21,44.5) -- (20,45.5) -- (21,46.5) -- (21,47.5) -- (22,48.5) -- (22,49.5) -- (29,56.5) -- (29,57.5) -- (27,59.5) -- (27,60.5) -- (26,61.5) -- (26,62.5) -- (24,64.5) -- (24,65.5) -- (22,67.5) -- (22.5,68) -- (24.5,68) -- (25.5,67) -- (28.5,67) -- (29.5,66) -- (31.5,66) -- (32.5,65) -- (33.5,65) -- (34.5,64) -- (36.5,64) -- (37.5,63) -- (38.5,63) -- (39.5,62) -- (41.5,62) -- (42.5,61) -- (44.5,61) -- (45.5,60) -- (46.5,60) -- (47.5,59) -- (48.5,59) -- (49,58.5) -- (49,57.5) -- (49.5,57) -- (51.5,57) -- (52,56.5) -- (50.5,55) -- (49.5,55) -- (48.5,54) -- (47.5,54) -- (46.5,53) -- (45.5,53) -- (43,50.5) -- (43,49.5) -- (42,48.5) -- (42,47.5) -- (41.5,47) -- (40.5,47) -- (39.5,48) -- (33.5,48) -- (32,46.5) -- (32,44.5) -- (34,42.5) -- (34,40.5) -- (33.5,40) -- (32.5,40) -- (31.5,39) -- (26.5,39) -- (26,38.5) -- (26,35.5) -- (27.5,34) -- (31.5,34) -- (33,35.5) -- (33,36.5) -- (34.5,38) -- (37.5,38) -- (38,37.5) -- (38,36.5) -- (41,33.5) -- (41,31.5) -- (40,30.5) -- (40,27.5) -- (39,26.5) -- (39,25.5) -- (38,24.5) -- (38,23.5) -- (37,22.5) -- (37,19.5) -- (37.5,19) -- (38.5,20) -- (39.5,20) -- (41,21.5) -- (41,24.5) -- (42,25.5) -- (42,27.5) -- (43,28.5) -- (43,30.5) -- (44,31.5) -- (44,35.5) -- (45,36.5) -- (45,37.5) -- (45.5,38) -- (46.5,38) -- (47,38.5) -- (47,40.5) -- (48,41.5) -- (48,42.5) -- (49,43.5) -- (49,44.5) -- (50,45.5) -- (50,48.5) -- (53,51.5) -- (52,52.5) -- (52,55.5) -- (53.5,57) -- (56.5,57) -- (58.5,59) -- (59.5,59) -- (60.5,60) -- (61.5,60) -- (62,60.5) -- (62,61.5) -- (63,62.5) -- (63,64.5) -- (63.5,65) -- (64.5,65) -- (65.5,64) -- (68.5,64) -- (69.5,65) -- (73.5,65) -- (74,65.5) -- (74,66.5) -- (75,67.5) -- (75,68.5) -- (75.5,69) -- (76,68.5) -- (76,66.5) -- (76.5,66) -- (78.5,68) -- (80.5,68) -- (81.5,69) -- (84.5,69) -- (85,68.5) -- (84,67.5) -- (84,65.5) -- (83.5,65) -- (82.5,65) -- (82,64.5) -- (82.5,64) -- (84.5,64) -- (86,65.5) -- (86,66.5) -- (88.5,69) -- (89.5,69) -- (90,69.5) -- (90,70.5) -- (91,71.5) -- (91,72.5) -- (93,74.5) -- (91.5,76) -- (90.5,76) -- (88.5,74) -- (87.5,74) -- (86.5,73) -- (85.5,73) -- (83.5,71) -- (82.5,71) -- (81.5,70) -- (80.5,70) -- (79.5,69) -- (77.5,69) -- (74.5,72) -- (73.5,72) -- (73,72.5) -- (73,74.5) -- (73.5,75) -- (78.5,75) -- (79.5,74) -- (80.5,74) -- (81,74.5) -- (81,78.5) -- (80,79.5) -- (80,80.5) -- (78.5,82) -- (73.5,82) -- (73,81.5) -- (73,79.5) -- (72,78.5) -- (72,77.5) -- (71.5,77) -- (70.5,77) -- (69.5,76) -- (68.5,76) -- (68,76.5) -- (68,77.5) -- (67,78.5) -- (67,79.5) -- (66.5,80) -- (65.5,80) -- (65,80.5) -- (66,81.5) -- (64.5,83) -- (63,81.5) -- (63,79.5) -- (62,78.5) -- (62,77.5) -- (61,76.5) -- (61,74.5) -- (60,73.5) -- (60,72.5) -- (59,71.5) -- (59,68.5) -- (58,67.5) -- (58,64.5) -- (56,62.5) -- (56,61.5) -- (55.5,61) -- (54.5,62) -- (53.5,62) -- (52,63.5) -- (52,64.5) -- (51,65.5) -- (51,66.5) -- (50,67.5) -- (50,68.5) -- (49,69.5) -- (49,70.5) -- (48,71.5) -- (48,75.5) -- (46,77.5) -- (46,79.5) -- (44.5,81) -- (42.5,81) -- (42,80.5) -- (44,78.5) -- (44,77.5) -- (43.5,77) -- (41.5,77) -- (39.5,75) -- (38.5,76) -- (37.5,76) -- (37,76.5) -- (37,78.5) -- (36,79.5) -- (38,81.5) -- (38,82.5) -- (37.5,83) -- (34.5,83) -- cycle
    (65.5,22) -- (63.5,22) -- (63,21.5) -- (63.5,21) -- (65.5,21) -- (66,21.5) -- cycle
    (55.5,28) -- (53.5,28) -- (53,27.5) -- (53.5,27) -- (54.5,27) -- (57.5,24) -- (59.5,24) -- (60,24.5) -- (57.5,27) -- (56.5,27) -- cycle
    (64.5,26) -- (64,25.5) -- (64.5,25) -- (65,25.5) -- cycle
    (63.5,29) -- (63,28.5) -- (63,27.5) -- (63.5,27) -- (64.5,27) -- (65,27.5) -- cycle
    (58.5,55) -- (57,53.5) -- (57,52.5) -- (56,51.5) -- (56,49.5) -- (57,48.5) -- (57,46.5) -- (60,43.5) -- (60,42.5) -- (59.5,42) -- (58.5,42) -- (58,41.5) -- (59,40.5) -- (59,39.5) -- (61.5,37) -- (62.5,37) -- (63.5,36) -- (65.5,36) -- (66.5,37) -- (68.5,37) -- (70,35.5) -- (69,34.5) -- (69,32.5) -- (68,31.5) -- (68.5,31) -- (74.5,31) -- (75,31.5) -- (75,32.5) -- (76,33.5) -- (76,35.5) -- (75,36.5) -- (75,37.5) -- (72,40.5) -- (72,41.5) -- (71,42.5) -- (72,43.5) -- (72,44.5) -- (72.5,45) -- (73.5,45) -- (74.5,46) -- (75.5,46) -- (76.5,45) -- (78.5,45) -- (79.5,44) -- (80,44.5) -- (80,47.5) -- (79.5,48) -- (77.5,48) -- (76.5,49) -- (74.5,49) -- (73.5,50) -- (71.5,50) -- (70.5,51) -- (68.5,51) -- (67.5,52) -- (65.5,52) -- (64.5,53) -- (63.5,53) -- (62.5,54) -- (59.5,54) -- cycle
    (86.5,50) -- (85.5,50) -- (85,49.5) -- (86,48.5) -- (86,46.5) -- (85.5,46) -- (82.5,46) -- (81.5,47) -- (81,46.5) -- (82,45.5) -- (82,44.5) -- (82.5,44) -- (84.5,44) -- (86.5,42) -- (88.5,42) -- (89.5,41) -- (92.5,41) -- (93,41.5) -- (93,42.5) -- (94,43.5) -- (94,44.5) -- (93.5,45) -- (92.5,45) -- (88.5,49) -- (87.5,49) -- cycle
    (44,42.5) -- (43.5,42) -- (43,42.5) -- (43.5,43) -- cycle
    (45,43.5) -- (44.5,43) -- (44,43.5) -- (44.5,44) -- cycle
    (80.5,60) -- (79.5,60) -- (79,59.5) -- (79,56.5) -- (81,54.5) -- (81,53.5) -- (82.5,52) -- (83,52.5) -- (83,53.5) -- (82,54.5) -- (82,55.5) -- (81.5,56) -- (80.5,56) -- (80,56.5) -- (80,57.5) -- (81,58.5) -- (81,59.5) -- cycle
    (55.5,54) -- (55,53.5) -- (55.5,53) -- (56,53.5) -- cycle
    (43,65.5) -- (42.5,65) -- (42,65.5) -- (42.5,66) -- cycle
    (61,67.5) -- (60.5,67) -- (60,67.5) -- (60.5,68) -- cycle
    (47,69.5) -- (47,68.5) -- (46.5,68) -- (46,68.5) -- (46,69.5) -- (46.5,70) -- cycle
    (62,68.5) -- (61.5,68) -- (61,68.5) -- (61.5,69) -- cycle
    (67.5,97) -- (65.5,97) -- (64,95.5) -- (64,94.5) -- (61.5,92) -- (60.5,92) -- (54.5,86) -- (53.5,86) -- (49.5,90) -- (48.5,90) -- (47.5,91) -- (46.5,91) -- (45.5,92) -- (45,91.5) -- (45,90.5) -- (45.5,90) -- (46.5,90) -- (48.5,88) -- (49.5,88) -- (54.5,83) -- (55.5,83) -- (60.5,88) -- (61.5,88) -- (62.5,89) -- (63.5,89) -- (64,89.5) -- (64,90.5) -- (64.5,91) -- (66.5,91) -- (67,90.5) -- (66,89.5) -- (67,88.5) -- (65,86.5) -- (65,85.5) -- (65.5,85) -- (66.5,85) -- (67,85.5) -- (67,86.5) -- (68,87.5) -- (68,90.5) -- (69,91.5) -- (69,94.5) -- (70,95.5) -- (69.5,96) -- (68.5,96) -- cycle
    (40.5,96) -- (36.5,96) -- (36,95.5) -- (37,94.5) -- (37,91.5) -- (40.5,88) -- (41.5,88) -- (42,88.5) -- (41,89.5) -- (41,92.5) -- (42,93.5) -- (42,94.5) -- cycle;
  \end{tikzpicture}%
}
\providecommand{\noethpIIrosette}{\makebox[9.44444pt][c]{\raisebox{-3.85pt}[0pt][0pt]{\NoetherRosetteGlyph}}}
\providecommand{\NoetherRosetteInline}{\makebox[17.72671698pt][c]{\raisebox{-3.85pt}[10.87671698pt][3.85pt]{\NoetherRosetteGlyph}}}

\providecommand{\dd}{\mathrm d}
\providecommand{\Div}{\operatorname{Div}}
\providecommand{\G}{\mathfrak G}
\providecommand{\bd}{\bar{\delta}}
\providecommand{\p}{\partial}

\providecommand{\frS}{\mathfrak S}
\providecommand{\frp}{\mathfrak p}
\providecommand{\fra}{\mathfrak a}
\providecommand{\frb}{\mathfrak b}
\providecommand{\frc}{\mathfrak c}
\providecommand{\frf}{\mathfrak f}
\providecommand{\frD}{\mathfrak D}
\providecommand{\frr}{\mathfrak r}
\providecommand{\frm}{\mathfrak m}
\providecommand{\frn}{\mathfrak n}
\providecommand{\calA}{\mathcal A}
\providecommand{\calB}{\mathcal B}
\providecommand{\calN}{\mathcal N}



\providecommand{\frH}{\mathfrak H}
\providecommand{\frK}{\mathfrak K}
\providecommand{\frM}{\mathfrak M}
\providecommand{\frN}{\mathfrak N}
\providecommand{\frG}{\mathfrak G}
\providecommand{\frS}{\mathfrak S}
\providecommand{\tmod}[1]{\;(\operatorname{mod} #1)}


\providecommand{\mM}{\mathfrak{M}}
\providecommand{\mN}{\mathfrak{N}}
\providecommand{\mL}{\mathfrak{L}}
\providecommand{\mG}{\mathfrak{G}}
\providecommand{\mA}{\mathfrak{A}}
\providecommand{\mB}{\mathfrak{B}}
\providecommand{\mx}{\mathfrak{x}}
\providecommand{\my}{\mathfrak{y}}
\providecommand{\mz}{\mathfrak{z}}
\providecommand{\ma}{\mathfrak{a}}
\providecommand{\mb}{\mathfrak{b}}
\providecommand{\me}{\mathfrak{e}}
\providecommand{\mbarA}{\overline{\mathfrak{A}}}
\providecommand{\mbarB}{\overline{\mathfrak{B}}}
\providecommand{\bara}{\overline{\mathfrak{a}}}


\providecommand{\widebar}[1]{\overline{#1}}
\providecommand{\mM}{\mathfrak{M}}
\providecommand{\mN}{\mathfrak{N}}
\providecommand{\mL}{\mathfrak{L}}
\providecommand{\mG}{\mathfrak{G}}
\providecommand{\mA}{\mathfrak{A}}
\providecommand{\mB}{\mathfrak{B}}
\providecommand{\mX}{\mathfrak{X}}
\providecommand{\mY}{\mathfrak{Y}}
\providecommand{\mU}{\mathfrak{U}}
\providecommand{\mP}{\mathfrak{P}}
\providecommand{\mQ}{\mathfrak{Q}}
\providecommand{\mS}{\mathfrak{S}}
\providecommand{\mD}{\mathfrak{D}}
\providecommand{\mR}{\mathfrak{R}}
\providecommand{\mV}{\mathfrak{V}}
\providecommand{\ma}{\mathfrak{a}}
\providecommand{\mb}{\mathfrak{b}}
\providecommand{\mx}{\mathfrak{x}}
\providecommand{\my}{\mathfrak{y}}
\providecommand{\mz}{\mathfrak{z}}
\providecommand{\me}{\mathfrak{e}}
\providecommand{\mbarA}{\overline{\mathfrak{A}}}
\providecommand{\mbarB}{\overline{\mathfrak{B}}}
\providecommand{\bara}{\overline{\mathfrak{a}}}
\providecommand{\modu}[1]{\;\operatorname{mod} #1}

\providecommand{\ideal}[1]{\mathfrak{#1}}
\providecommand{\srcfn}[2]{%
  \begingroup
  \renewcommand{\thefootnote}{#1}%
  \footnote{#2}%
  \addtocounter{footnote}{-1}%
  \endgroup}
\providecommand{\srcfnmark}[1]{%
  \begingroup
  \renewcommand{\thefootnote}{#1}%
  \footnotemark%
  \addtocounter{footnote}{-1}%
  \endgroup}
\providecommand{\srcfntext}[2]{%
  \begingroup
  \renewcommand{\thefootnote}{#1}%
  \footnotetext{#2}%
  \endgroup}
\usepackage[unicode,hidelinks,hypertexnames=false,bookmarksopen=true]{hyperref}
\hypersetup{pdftitle={Emmy Noether — Écrits mathématiques en français},pdfauthor={Emmy Noether}}
\newcommand{\NoetherNavAnchor}[2]{\phantomsection\label{nav:#1}\pdfbookmark[1]{#2}{noether.#1}}
\begin{document}
\thispagestyle{empty}
\begin{center}
{\huge\bfseries Emmy Noether}\par
\vspace{0.8em}
{\Large Écrits mathématiques}\par
\vspace{0.5em}
{\large Traduction française --- texte de base}\par
\end{center}
\vspace{2.5em}

Ce recueil réunit les articles 1 à 43 d'Emmy Noether, le cours
\emph{Algèbre des grandeurs hypercomplexes} du semestre d'hiver 1929/30,
rédigé par M.\ Deuring, ainsi que l'article de Kapferer et Noether et son
complément. Il comprend également la bibliographie et les listes de notices,
comptes rendus et livres.

Les articles portent notamment sur la théorie des invariants, les idéaux,
les corps, les algèbres et leurs représentations. Le cours comporte une
introduction et trente et une sections, réparties en six chapitres.
Les titres, signatures et références bibliographiques propres à chaque texte
figurent avec celui-ci.

\vspace{1.5em}
\noindent\textbf{Lecture et sources.}
Le volume est accompagné d'un fichier \LaTeX{} cumulatif et d'une archive
contenant les fichiers de composition et les instructions de compilation.

\vfill
\begingroup\small
\noindent\textbf{Auteurs et traduction.}
Les œuvres sont d'Emmy Noether et des auteurs nommés dans chaque texte.
Le cours de 1929/30 a été rédigé par M.\ Deuring.
La traduction française est produite par IA. Le modèle employé pour sa
production initiale n'est pas établi par les éléments de provenance disponibles.
Corrections et composition : OpenAI Codex --- GPT-6 Astra, effort Ultra ; Claude Opus 5.5 (Anthropic), 26 octobre 2026.
Aucune traduction humaine ni relecture scientifique humaine n'est attribuée
à l'éditeur.

\medskip
\noindent\textbf{Sources.}
Édition française ED0014/v051 :
\href{https://doi.org/10.5281/zenodo.21921884}{10.5281/zenodo.21921884}.
Édition allemande ED0014/v051 :
\href{https://doi.org/10.5281/zenodo.21908301}{10.5281/zenodo.21908301}.
Les références des publications originales et les signatures restent dans les textes.

\medskip
\noindent\textbf{Droits.}
Les droits sur les œuvres originales et les contributions dérivées protégées
restent à leurs titulaires respectifs. Cette édition est proposée gratuitement
à des fins non commerciales et purement éducatives, dans la présomption de
l'éditeur que les titulaires ne s'y opposeraient pas. Il s'agit d'une présomption,
non d'une autorisation établie : le libre accès ne vaut pas autorisation de
publication dérivée. Les titulaires peuvent
\href{https://github.com/KokunoYumeto/emmy-noether-fr/issues/new}{signaler une objection sur GitHub};
l'éditeur s'engage à retirer les éléments concernés à leur demande.
Aucune nouvelle licence ni mise dans le domaine public n'est accordée.
\endgroup
\clearpage
\pdfbookmark[0]{Table des matières}{noether.contents}
\section*{Table des matières}
\begingroup
\small
\setlength{\parindent}{0pt}
\setlength{\tabcolsep}{0pt}
\begin{longtable}{@{}p{0.91\linewidth}@{\hspace{1em}}r@{}}
\multicolumn{2}{@{}l}{\textbf{Articles 1–43}}\\*[0.5em]
\hyperref[nav:P01]{1. Sur la formation du système de formes de la forme biquadratique ternaire} & \pageref{nav:P01}\\[0.35em]
\hyperref[nav:P02]{2. Sur la formation du système de formes de la forme biquadratique ternaire} & \pageref{nav:P02}\\[0.35em]
\hyperref[nav:P03]{3. Sur la théorie des invariants des formes à n variables} & \pageref{nav:P03}\\[0.35em]
\hyperref[nav:P04]{4. Sur la théorie des invariants des formes à n variables} & \pageref{nav:P04}\\[0.35em]
\hyperref[nav:P05]{5. Corps de fonctions rationnelles} & \pageref{nav:P05}\\[0.35em]
\hyperref[nav:P06]{6. Corps et systèmes de fonctions rationnelles} & \pageref{nav:P06}\\[0.35em]
\hyperref[nav:P07]{7. Le théorème de finitude pour les invariants des groupes finis} & \pageref{nav:P07}\\[0.35em]
\hyperref[nav:P08]{8. Sur la représentation entière rationnelle des invariants d’un système d’un nombre arbitraire de formes fondamentales} & \pageref{nav:P08}\\[0.35em]
\hyperref[nav:P09]{9. Les domaines les plus généraux d’entiers transcendants} & \pageref{nav:P09}\\[0.35em]
\hyperref[nav:P10]{10. Les équations fonctionnelles de l’application isomorphe} & \pageref{nav:P10}\\[0.35em]
\hyperref[nav:P11]{11. Équations à groupe prescrit} & \pageref{nav:P11}\\[0.35em]
\hyperref[nav:P12]{12. Invariants d’expressions différentielles arbitraires} & \pageref{nav:P12}\\[0.35em]
\hyperref[nav:P13]{13. Problèmes variationnels invariants} & \pageref{nav:P13}\\[0.35em]
\hyperref[nav:P14]{14. La théorie arithmétique des fonctions algébriques d’une variable dans son rapport avec les autres théories et avec la théorie des corps de nombres} & \pageref{nav:P14}\\[0.35em]
\hyperref[nav:P15]{15. La finitude du système des invariants entiers des formes binaires} & \pageref{nav:P15}\\[0.35em]
\hyperref[nav:P16]{16. Sur le développement en série dans la théorie des formes} & \pageref{nav:P16}\\[0.35em]
\hyperref[nav:P17]{17. Modules dans des domaines non commutatifs, provenant notamment d’expressions différentielles et d’expressions aux différences} & \pageref{nav:P17}\\[0.35em]
\hyperref[nav:P18]{18. Sur un travail de K. Hentzelt, tombé à la guerre, sur la théorie de l’élimination} & \pageref{nav:P18}\\[0.35em]
\hyperref[nav:P19]{19. Théorie des idéaux dans les anneaux} & \pageref{nav:P19}\\[0.35em]
\hyperref[nav:P20]{20. Un critère algébrique d’irréductibilité absolue} & \pageref{nav:P20}\\[0.35em]
\hyperref[nav:P21]{21. Calcul des variations formel et invariants différentiels} & \pageref{nav:P21}\\[0.35em]
\hyperref[nav:P22]{22. Sur la théorie des idéaux polynomiaux et des résultantes} & \pageref{nav:P22}\\[0.35em]
\hyperref[nav:P23]{23. Invariants algébriques et différentiels} & \pageref{nav:P23}\\[0.35em]
\hyperref[nav:P24]{24. Théorie de l’élimination et théorie générale des idéaux} & \pageref{nav:P24}\\[0.35em]
\hyperref[nav:P25]{25. Théorie de l’élimination et théorie des idéaux} & \pageref{nav:P25}\\[0.35em]
\hyperref[nav:P26]{26. Construction abstraite de la théorie des idéaux dans un corps de nombres algébriques} & \pageref{nav:P26}\\[0.35em]
\hyperref[nav:P27]{27. Nombres de Hilbert dans la théorie des idéaux} & \pageref{nav:P27}\\[0.35em]
\hyperref[nav:P28]{28. Caractères de groupes et théorie des idéaux} & \pageref{nav:P28}\\[0.35em]
\hyperref[nav:P29]{29. Le théorème de finitude des invariants des groupes linéaires finis de caractéristique p} & \pageref{nav:P29}\\[0.35em]
\hyperref[nav:P30]{30. Construction abstraite de la théorie des idéaux dans les corps de nombres et de fonctions algébriques} & \pageref{nav:P30}\\[0.35em]
\hyperref[nav:P31]{31. Le théorème du discriminant pour les ordres d’un corps algébrique de nombres ou de fonctions} & \pageref{nav:P31}\\[0.35em]
\hyperref[nav:P32]{32. Avec R. Brauer : sur les corps de décomposition minimaux des représentations irréductibles} & \pageref{nav:P32}\\[0.35em]
\hyperref[nav:P33]{33. Grandeurs hypercomplexes et théorie des représentations dans une conception arithmétique} & \pageref{nav:P33}\\[0.35em]
\hyperref[nav:P34]{34. Grandeurs hypercomplexes et théorie des représentations} & \pageref{nav:P34}\\[0.35em]
\hyperref[nav:P35]{35. Sur les domaines maximaux de fonctions à coefficients entiers} & \pageref{nav:P35}\\[0.35em]
\hyperref[nav:P36]{36. Différentiation des idéaux et différente} & \pageref{nav:P36}\\[0.35em]
\hyperref[nav:P37]{37. Base normale dans les corps sans ramification supérieure} & \pageref{nav:P37}\\[0.35em]
\hyperref[nav:P38]{38. Avec R. Brauer et H. Hasse : preuve d’un théorème principal dans la théorie des algèbres} & \pageref{nav:P38}\\[0.35em]
\hyperref[nav:P39]{39. Systèmes hypercomplexes dans leurs rapports avec l’algèbre commutative et la théorie des nombres} & \pageref{nav:P39}\\[0.35em]
\hyperref[nav:P40]{40. Algèbre non commutative} & \pageref{nav:P40}\\[0.35em]
\hyperref[nav:P41]{41. Le théorème du genre principal pour les corps de nombres relativement galoisiens} & \pageref{nav:P41}\\[0.35em]
\hyperref[nav:P42]{42. Produits croisés décomposés et leurs ordres maximaux} & \pageref{nav:P42}\\[0.35em]
\hyperref[nav:P43]{43. Différentiation des idéaux et différente} & \pageref{nav:P43}\\[0.35em]
\multicolumn{2}{@{}l}{\textbf{Algèbre des grandeurs hypercomplexes}}\\*[0.5em]
\hyperref[nav:BOOK]{Algèbre des grandeurs hypercomplexes — cours de la professeure E. Noether, semestre d’hiver 1929/30, rédigé par le professeur M. Deuring} & \pageref{nav:BOOK}\\[0.35em]
\hyperref[nav:INTRO]{Introduction} & \pageref{nav:INTRO}\\[0.35em]
\hyperref[nav:S01]{§ 1. Définition de la représentation directe — définition de la représentation réciproque} & \pageref{nav:S01}\\[0.35em]
\hyperref[nav:S02]{§ 2. Classes de représentations} & \pageref{nav:S02}\\[0.35em]
\hyperref[nav:S03]{§ 3. Modules de représentation} & \pageref{nav:S03}\\[0.35em]
\hyperref[nav:S04]{§ 4. Le rapport entre modules de représentation et représentations} & \pageref{nav:S04}\\[0.35em]
\hyperref[nav:S05]{§ 5. Extension du corps des coefficients des systèmes hypercomplexes} & \pageref{nav:S05}\\[0.35em]
\hyperref[nav:S06]{§ 6. Les représentations irréductibles des systèmes commutatifs} & \pageref{nav:S06}\\[0.35em]
\hyperref[nav:S07]{§ 7. Les isomorphismes d’un corps} & \pageref{nav:S07}\\[0.35em]
\hyperref[nav:S08]{§ 8. Le corps de décomposition} & \pageref{nav:S08}\\[0.35em]
\hyperref[nav:S09]{§ 9. Les isomorphismes de $Z_1$ sur les $Z_i$} & \pageref{nav:S09}\\[0.35em]
\hyperref[nav:S10]{§ 10. Le groupe de Galois} & \pageref{nav:S10}\\[0.35em]
\hyperref[nav:S11]{§ 11. Le théorème fondamental de la théorie de Galois} & \pageref{nav:S11}\\[0.35em]
\hyperref[nav:S12]{§ 12. La signification formelle des composantes $e_i$} & \pageref{nav:S12}\\[0.35em]
\hyperref[nav:S13]{§ 13. Le théorème général de prolongement} & \pageref{nav:S13}\\[0.35em]
\hyperref[nav:S14]{§ 14. L’anneau de groupe} & \pageref{nav:S14}\\[0.35em]
\hyperref[nav:S15]{§ 15. Les anneaux de groupe des groupes abéliens} & \pageref{nav:S15}\\[0.35em]
\hyperref[nav:S16]{§ 16. Les relations entre caractères} & \pageref{nav:S16}\\[0.35em]
\hyperref[nav:S17]{§ 17. La théorie de Galois des groupes abéliens} & \pageref{nav:S17}\\[0.35em]
\hyperref[nav:S18]{§ 18. Un lemme} & \pageref{nav:S18}\\[0.35em]
\hyperref[nav:S19]{§ 19. Représentations des systèmes hypercomplexes simples bilatères dans des corps non commutatifs étendant leurs corps de coefficients} & \pageref{nav:S19}\\[0.35em]
\hyperref[nav:S20]{§ 20. Corps non commutatifs} & \pageref{nav:S20}\\[0.35em]
\hyperref[nav:S21]{§ 21. La théorie de Galois des corps non commutatifs} & \pageref{nav:S21}\\[0.35em]
\hyperref[nav:S22]{§ 22. Le groupe des corps gauches de centre donné} & \pageref{nav:S22}\\[0.35em]
\hyperref[nav:S23]{§ 23. Systèmes de facteurs} & \pageref{nav:S23}\\[0.35em]
\hyperref[nav:S24]{§ 24. Multiplication des systèmes de facteurs} & \pageref{nav:S24}\\[0.35em]
\hyperref[nav:S25]{§ 25. Représentation normale de $\mathfrak K_r$ avec sous-corps commutatif maximal galoisien} & \pageref{nav:S25}\\[0.35em]
\hyperref[nav:S26]{§ 26. Multiplication des représentations croisées} & \pageref{nav:S26}\\[0.35em]
\hyperref[nav:S27]{§ 27. Représentation des $\mathfrak K_r$ comme produits croisés} & \pageref{nav:S27}\\[0.35em]
\hyperref[nav:S28]{§ 28. Théorème du produit pour les systèmes de facteurs} & \pageref{nav:S28}\\[0.35em]
\hyperref[nav:S29]{§ 29. Théorème du genre principal dans le cas minimal} & \pageref{nav:S29}\\[0.35em]
\hyperref[nav:S30]{§ 30. Spécialisation aux corps de décomposition cycliques} & \pageref{nav:S30}\\[0.35em]
\hyperref[nav:S31]{§ 31. Applications du cas particulier cyclique} & \pageref{nav:S31}\\[0.35em]
\multicolumn{2}{@{}l}{\textbf{Complément, bibliographie et listes}}\\*[0.5em]
\hyperref[nav:KAPFERER]{Conditions de multiplicité nécessaires et suffisantes pour le théorème fondamental de Noether sur les fonctions algébriques — par H. Kapferer} & \pageref{nav:KAPFERER}\\[0.35em]
\hyperref[nav:SUPPLEMENT]{Complément, en commun avec E. Noether} & \pageref{nav:SUPPLEMENT}\\[0.35em]
\hyperref[nav:BIBLIOGRAPHIE]{Bibliographie} & \pageref{nav:BIBLIOGRAPHIE}\\[0.35em]
\hyperref[nav:NOTICES]{Communications brèves parues dans les Jahresberichte der Deutschen Mathematiker-Vereinigung} & \pageref{nav:NOTICES}\\[0.35em]
\hyperref[nav:REVIEWS]{Comptes rendus parus dans les Jahresberichte der Deutschen Mathematiker-Vereinigung} & \pageref{nav:REVIEWS}\\[0.35em]
\hyperref[nav:LIVRES]{Livres réalisés avec la participation d’Emmy Noether} & \pageref{nav:LIVRES}\\[0.35em]
\end{longtable}
\endgroup
\clearpage
\providecommand{\mat}[1]{\begin{pmatrix}#1\end{pmatrix}}

\begin{center}
{\Large\bfseries \pdfbookmark[0]{Articles 1–43}{noether.group.P01}\NoetherNavAnchor{P01}{1. Sur la formation du système de formes de la forme biquadratique ternaire}1. Sur la formation du système de formes de la forme\\ biquadratique ternaire}\par
\vspace{1.2em}
Par Emmy Noether.\par
\vspace{0.8em}
\emph{Sitz. Ber. d. Physikal.-mediz. Sozietät in Erlangen} 39 (1907), p. 176--179
\end{center}

\vspace{4em}
\begin{center}
(Extrait de la dissertation de l'autrice.)
\end{center}

Des travaux de Gordan, Maisano et Pascal\footnote{P. Gordan, \emph{Über das volle Formensystem der ternären biquadratischen Form} $f=x_1^3x_2+x_2^3x_3+x_3^3x_1$ (\emph{Math. Annalen} t. XVII, p. 217 à 233, 1880); G. Maisano, 1. \emph{Sistemi completi dei primi cinque gradi della forma ternaria biquadratica e degl' invarianti, covarianti e contravarianti di sesto grado} (\emph{Giorn. di Battaglini} XIX). 2. \emph{Sui covarianti indipendenti di $6^\circ$ grado nei coefficienti della forma biquadratica ternaria} (\emph{Rend. Circ. Mat. di Palermo} I, 1887); E. Pascal, \emph{Contributo alla teoria della forma ternaria biquadratica e delle sue varie decomposizioni in fattori} (mémoire couronné par la R. Accademia delle scienze fisiche e matematiche di Napoli, 1905).} portent sur le système de formes de la forme biquadratique ternaire. M. Gordan établit le système de formes complet, composé de 54 formations, de la forme automorphe spéciale
\[
f=x_1^3x_2+x_2^3x_3+x_3^3x_1,
\]
en prenant pour base des principes analogues à ceux qu'il avait donnés pour les systèmes de formes dans le domaine binaire.

Chez M. Maisano, pour la forme biquadratique générale, les formes jusqu'au 5\ieme{} ordre\footnote{Par \og ordre \fg{} on entendra la dimension en les coefficients, par \og degré \fg{} celle en les variables.} inclusivement sont établies, ainsi que quelques invariants, covariants et contravariants d'ordre supérieur, suivant la méthode employée par M. Gordan dans le tome I des \emph{Math. Annalen} pour la forme cubique ternaire. M. Pascal s'occupe, en utilisant les résultats de Maisano, principalement de la question de la décomposition de la forme biquadratique en facteurs\footnote{Dans la seconde note brève, M. Maisano tente de démontrer une dépendance linéaire des 3 covariants de 6\ieme{} ordre et de 6\ieme{} degré. Contrairement à cette démonstration incomplète, M. Pascal croit avoir démontré l'indépendance linéaire des 3 covariants de 6\ieme{} ordre, de même que celle des 3 invariants de 9\ieme{} ordre. Mais qu'il existe en fait une relation linéaire entre les 3 covariants, d'une part, et les 3 invariants, d'autre part, est apparu au cours de mes calculs; les formules explicites, dans la notation de Pascal, sont:
\[
0=20\Omega_1+6\Omega_2-3\Omega_3+4f\cdot C_1-12f\cdot C_2-4A\cdot\Delta.
\]
\[
0=4A^3-15A\cdot B+30C-90D+9E.
\]}. 

\textbf{Le but de mes recherches est l'établissement du système de formes pour la forme biquadratique ternaire générale; pour commencer, seules les bases principales sont données, et l'on établit un système dit \og relativement complet \fg{}\footnote{Pour la dénomination, voir Gordan--Kerschensteiner, \emph{Vorlesungen über Invariantentheorie}, t. II, p. 227.}.}

L'idée fondamentale de la formation de systèmes de formes ternaires est la même que dans le domaine binaire. En partant d'un premier système relativement complet --- le système des formes reprises du binaire ---, on parvient, suivant une loi déterminée, à des systèmes de module toujours plus élevé, jusqu'à ce que le système d'un module --- le module étant pris comme forme fondamentale --- devienne fini et connu, ou encore jusqu'à ce qu'un module puisse être ramené à des formes ayant des invariants pour facteur. Par transvection du système relativement complet avec le système du module, on obtient dans le premier cas le système absolument complet, tandis que, dans le second cas, le système relativement complet et le système absolument complet deviennent identiques. En raison de la finitude des systèmes de formes démontrée en général par M. Hilbert, ce procédé doit nécessairement arriver à un terme.

Dans notre cas, le module $(abc)$ du premier système relativement complet se ramène aux modules
\[
\Delta=(abc)^2a_x^2b_x^2c_x^2\quad\text{et}\quad \nu=(abu)^4;
\]
il en résulte aisément le système relativement complet mod $\nu$. Comme suite de modules, nous choisissons maintenant les formes
\[
\nu;\qquad \nu^{(\nu)}=(\nu\nu_1x)^4=s_x^4,
\]
\[
\nu(s)=(ss'u)^4\ldots
\]
et nous adjoignons comme autres modules deux formes quadratiques qui apparaissent dans la formation du système relativement complet mod $s$, $u_\varrho^2$ et $t_x^2$. On peut alors montrer que le module qui suit $s$, $(ss'u)^4$, est réductible au module $(\varrho,t)$. Mais comme le système simultané de deux formes quadratiques est fini et connu, le terme caractérisé ci-dessus est ainsi atteint. Comme système relativement complet mod $(\varrho,t)$, on obtient 331 formations. ---

J'ajoute encore quelques indications concernant la méthode employée.

Le procédé fondamental pour engendrer des formes, le procédé de contraction, peut se définir dans le domaine ternaire comme suit:

Si, dans un produit symbolique
\[
s_x^m t_x^n u_\sigma^\mu u_\tau^\nu
\]
on remplace les paires de facteurs
\[
\begin{array}{c|c|c|c|c}
 & s_xt_x & u_\sigma u_\tau & s_xu_\sigma\ \text{ou}\ s_xu_\tau & t_xu_\tau\ \text{ou}\ t_xu_\sigma\\
\text{resp. par} & (stu) & (\sigma\tau x) & s_\sigma\ \text{ou}\ s_\tau & t_\tau\ \text{ou}\ t_\sigma\\
\text{contraction} & \mathrm{I} & \mathrm{II} & \mathrm{III} & \mathrm{IV}
\end{array}
\]
alors les formes qui en résultent sont issues de la forme primitive par contraction\footnote{Gordan, \emph{Math. Annalen}, t. XVII, p. 219.}.

On peut encore toujours poser ici $s_\sigma=0;\ t_\tau=0$. Pour le rapport entre les contractions individuelles vaut alors le théorème:

\textbf{Les contractions I et II sont des contractions fondamentales, à partir desquelles, indépendamment de l'ordre de composition, les contractions III et IV peuvent être composées. Autrement dit: pour former toutes les formes issues d'une expression donnée par contraction, il suffit d'appliquer les contractions I et II}\footnote{Le théorème souffre une exception dans le cas où $n$ et $\mu$, resp. $m$ et $\nu$, s'annulent simultanément; la \og série de formes \fg{} se réduit alors au terme diagonal.}.

Nous définissons comme \og série de formes \fg{} une forme initiale avec toutes les formes qui en proviennent par contraction en elle-même. D'après le théorème, une série de formes $s_x^m t_x^n u_\sigma^\mu u_\tau^\nu$ peut être disposée en un schéma rectangulaire. En progressant, on obtient:
\begin{enumerate}
\item d'une colonne vers la droite, toutes les formes issues par contraction I des formes voisines;
\item d'une ligne vers le bas, toutes les formes issues par contraction II des formes placées au-dessus, et par conséquent toutes les formes issues par contraction.
\end{enumerate}

Sous ces hypothèses, les théorèmes sur les réducteurs peuvent être énoncés sous leur forme la plus générale et, avec la définition
\begin{center}
\og Un réducteur est une série de formes réductible \fg{}
\end{center}
de la manière suivante:

\begin{center}
\textbf{Si la forme initiale d'une série de formes est réductible parce que l'un de ses membres est issu par contraction avec un réducteur, et si la forme finale de la série de formes est née de ce même membre par contraction, alors la série de formes tout entière est réductible.}
\end{center}

Comme autres méthodes de réduction, il faut mentionner:
\begin{enumerate}
\item la \og double réduction \fg{}, c'est-à-dire la réduction d'une forme de deux manières afin d'obtenir une relation entre les formes supérieures;
\item la \og contraction avec des formes décomposées \fg{}, qui présente en partie le caractère de la réduction par réducteurs, en partie celui de la \og double réduction \fg{}.
\end{enumerate}

\clearpage
\setcounter{footnote}{0}

\begin{center}
{\Large\bfseries \NoetherNavAnchor{P02}{2. Sur la formation du système de formes de la forme biquadratique ternaire}2. Sur la formation du système de formes de la forme\\ biquadratique ternaire\footnote{Un bref extrait de l'introduction et du chapitre I a paru dans les \emph{Sitzungsberichte der phys.-med. Sozietät Erlangen 1907}, p. 176 à 179.}}\par
\vspace{0.8em}
Par Mlle Emmy Noether à Erlangen.\par
\vspace{1.2em}
Journal f. d. reine u. angew. Math. 134 (1908), p. 23--90 et deux tableaux
\end{center}

\vspace{3em}

\begin{center}
\textbf{Introduction.}
\end{center}

Des travaux de Gordan, Maisano et Pascal\footnote{P. Gordan, sur le système complet de formes de la forme biquadratique ternaire
\[
f=x_1^3x_2+x_2^3x_3+x_3^3x_1.
\]
(\emph{Math. Annalen}, t. XVII (1880), p. 217--233.)
G. Maisano, 1) \emph{Sistemi completi dei primi cinque gradi della forma ternaria biquadratica e degl' invarianti, covarianti e contravarianti di sesto grado}. (Giorn. di Battaglini XIX.)
2) \emph{Sui covarianti indipendenti di $6^\circ$ grado nei coefficienti della forma biquadratica ternaria}. (Rend. Circ. Mat. di Palermo I, 1887.)
E. Pascal, \emph{Contributo alla teoria della forma ternaria biquadratica e delle sue varie decomposizioni in fattori}. (Memoria premiata dalla R. Accademia delle scienze fisiche e matematiche di Napoli. 1905.)} portent sur le système de formes de la forme biquadratique ternaire. M. Gordan établit le système complet de formes, composé de 54 formations, de la forme automorphe spéciale :
\[
f=x_1^3x_2+x_2^3x_3+x_3^3x_1
\]
en se fondant sur des principes analogues à ceux qu'il avait donnés pour les systèmes de formes dans le domaine binaire.

Chez M. Maisano, les formes jusqu'au cinquième ordre\footnote{Par \og{} ordre \fg{} on entendra la dimension dans les coefficients, par \og{} degré \fg{} celle dans les variables.} inclusivement sont établies pour la forme biquadratique générale, ainsi que quelques invariants, covariants et contravariants d'ordre supérieur, suivant la méthode appliquée par M. Gordan, dans le tome I des \emph{Math. Annalen}, à la forme cubique ternaire. M. Pascal, utilisant les résultats de Maisano, s'occupe principalement de la question de la décomposition de la forme biquadratique en facteurs\footnote{Dans la seconde courte note, M. Maisano tente de démontrer une dépendance linéaire des trois covariants du 6e ordre et du 6e degré. Contrairement à cette preuve incomplète, M. Pascal croit avoir démontré l'indépendance linéaire des trois covariants du 6e ordre, ainsi que celle des trois invariants du 9e ordre. Mais qu'il existe en fait une relation linéaire entre les trois covariants d'une part, et les trois invariants d'autre part, les formules (13), \S{} 11, et (22), note du \S{} 17, du présent travail le montrent, ces relations y étant données explicitement. D'après une communication de M. Pascal, la contradiction s'explique par une erreur numérique dans l'expression de son contravariant $p$ (nommé ici $\varrho$), p. 46 de son mémoire.}.

\emph{Le but des recherches qui suivent est l'établissement du système de formes de la forme biquadratique ternaire générale ; plus précisément, on ne donne dans ce travail que les fondements principaux, en établissant un \og{} système relativement complet \fg{}\footnote{Pour la terminologie, voir \S{} 3a.}.}

Le travail se rattache étroitement à celui de Gordan ; toutefois, les principes qui n'y étaient donnés que de manière tout à fait générale devaient d'abord être élaborés en détail. À l'aide d'un théorème sur la relation entre les contractions et par l'introduction de la \og{} série de formes \fg{} (\S{} 1), les théorèmes de réduction sont formulés avec précision et complétés (\S{} 3), tandis que l'établissement récurrent de développements particuliers en séries (\S{} 2) fournit le moyen calculatoire permettant d'effectuer réellement les réductions.

L'idée fondamentale de la formation des systèmes de formes ternaires est la même que dans le domaine binaire. Partant d'un premier système relativement complet --- le système des formes reprises du binaire ---, on parvient, suivant une loi déterminée, à des systèmes de module toujours plus élevé, jusqu'à ce que le système d'un module --- le module étant pris comme forme fondamentale --- devienne fini et connu, ou encore jusqu'à ce qu'un module puisse être ramené à des formes ayant des invariants pour facteurs. Par transvection du système relativement complet avec le système du module, il naît dans le premier cas le système absolument complet, tandis que dans le second cas le système relativement complet et le système absolument complet deviennent identiques. En raison de la finitude des systèmes de formes démontrée en général par M. Hilbert, ce procédé doit nécessairement aboutir.

Dans notre cas, le module $(abc)$ du premier système relativement complet (\S{} 4) est ramené aux modules
\[
\Delta=(abc)^2a_x^2b_x^2c_x^2\quad\text{et}\quad \nu=(abu)^4
\]
(\S{} 5), puis l'on trouve le système relativement complet modulo $\nu$ (\S{} 6). Ces résultats sont déjà donnés dans le travail de Gordan pour la forme biquadratique générale ; notre mode de déduction se transpose toutefois plus aisément aux formes de degré supérieur.

Comme série des modules, nous choisissons maintenant les formes (\S{} 7) :
\[
\nu,\qquad \nu(\nu)=(\nu\nu_1x)^4=s_x^4,\qquad \nu(s)=(ss'u)^4,\ldots .
\]

Lors de la formation du système relativement complet modulo $s$, deux formes quadratiques apparaissant dans le système, $u_\varrho^2$ et $t_x^2$, sont adjointes comme modules (\S{} 9 et 10), et l'on forme par conséquent le système relativement complet modulo $(s,\varrho,t)$ ; on montre ensuite (\S{} 17) que le module $(ss'u)^4$ est réductible aux modules $\varrho$ et $t$. Le système immédiatement supérieur passe ainsi en un système relativement complet modulo $(\varrho,t)$. Mais comme le système simultané de deux formes quadratiques est fini et connu, et se compose dans le cas général de 20 formations\footnote{Voir Clebsch--Lindemann, \emph{Vorlesungen über Geometrie}. (t. I, p. 288 sq.) Système de deux formes cogrédientes. Gordan, \emph{Über Büschel von Kegelschnitten}. (\emph{Math. Ann.} t. XIX, p. 530.) Système de deux formes contragrédientes.}, l'établissement du système relativement complet modulo $(\varrho,t)$ atteint l'achèvement indiqué ci-dessus. La transvection de ce système avec le système de $(\varrho,t)$, en vue de former le système absolument complet, est réservée.

On trouve, comme système relativement complet modulo $(\varrho,t)$, 331 formations, ordonnées dans le tableau joint selon leur degré dans les variables $x$ et $u$.

Pour finir, nous donnons encore un aperçu des symboles introduits :
\begin{align*}
f&=a_x^4,\\
\theta&=\theta_x^4u_\vartheta^2=(abu)^2a_x^2b_x^2,\\
K&=K_x^6u_k^3=(a\theta u)u_\vartheta^2a_x^3\theta_x^3,\\
j&=u_j^6=(a\theta u)^4u_\vartheta^2,\\
\Delta&=\Delta_x^6=a_\vartheta^2a_x^2\theta_x^4=(abc)^2a_x^2b_x^2c_x^2,\\
N&=N_x^8u_n=(a\Delta u)a_x^3\Delta_x^5,\\
\nu&=u_\nu^4=(abu)^4,\\
H&=H_x^2u_\eta^4=(\nu\nu_1x)^2u_\nu^2u_{\nu_1}^2,\\
L&=L_x^3u_l^6=(\nu\eta x)H_x^2u_\nu^3u_\eta^3,\\
g&=g_x^6=(\nu\eta x)^4H_x^2,\\
\sigma&=u_\sigma^6=H_\nu^2u_\nu^2u_\eta^4,\\
s&=s_x^4=(\nu\nu_1x)^4,\\
Z&=Z_x^4u_\zeta^2=(ss'u)^2s_x^2s_x'^2,\\
\varrho&=u_\varrho^2=\theta_\nu^4u_\vartheta^2,\\
t&=t_x^2=a_\eta^4H_x^2,\\
i&=a_\nu^4,\\
J&=s_\nu^4.
\end{align*}

\begin{center}
{\Large\bfseries Chapitre I. \quad Théorèmes généraux sur les formes ternaires.}
\end{center}

\begin{center}
\textbf{\S{} 1. \quad Procédé de contraction. \quad Séries de formes.}
\end{center}

Le procédé fondamental pour produire des formes est le procédé de contraction, qui se définit dans le domaine ternaire de la manière suivante. Soit donné un produit symbolique :
\[
s_x^m t_x^n u_\sigma^\mu u_\tau^\nu .
\]
Si l'on remplace les paires de facteurs
\[
s_xt_x\qquad u_\sigma u_\tau\qquad s_xu_\sigma\ \text{ou}\ s_xu_\tau
\qquad t_xu_\tau\ \text{ou}\ t_xu_\sigma
\]
respectivement par
\[
(stu)\qquad (\sigma\tau x)\qquad s_\sigma\ \text{ou}\ s_\tau
\qquad t_\tau\ \text{ou}\ t_\sigma,
\]
contraction
\[
\text{I}\qquad\text{II}\qquad\text{III}\qquad\text{IV},
\]
les formes obtenues sont issues de la forme primitive par contraction\footnote{Gordan, \emph{Math. Annalen} t. XVII p. 219.}.

L'expression $s_x^m t_x^n u_\sigma^\mu u_\tau^\nu$ peut alors être comprise soit comme le produit effectif de deux formes $S\cdot T$, soit comme une seule forme avec plusieurs séries de symboles :
\[
s_x^m t_x^n u_\sigma^\mu u_\tau^\nu=A_x^{m+n}u_\lambda^{\mu+\nu}.
\]
Dans le premier cas nous parlons de la contraction de la forme $S$ avec $T$, dans le second de la \og{} contraction de la forme $A$ en elle-même \fg{}. Pour abréger, nous supposerons toujours dans la suite (ce qui est effectivement le cas pour toutes les formes considérées plus loin)
\[
\tag{a}
s_\sigma^\lambda=0,\qquad t_\tau^\varkappa=0
\]
pour tous $\lambda$ et $\varkappa$ non nuls, et conjointement avec toutes les autres contractions. Cela signifie que la forme $s_x^m t_y^n u_\sigma^\mu v_\tau^\nu$ est une dite \og{} forme normale \fg{}, c'est-à-dire qu'elle s'annule sous l'application du procédé $\Omega$, aussi bien relativement aux variables $x$ et $u$ que relativement aux variables $y$ et $v$.

La condition (a) ne constitue donc aucune restriction, mais peut toujours être atteinte\footnote{Voir Gordan, \emph{Math. Annalen} t. V p. 104.}.

La relation entre les différentes contractions est donnée par le théorème suivant :

\emph{Théorème I. Les contractions I et II sont des contractions fondamentales à partir desquelles, indépendamment de l'ordre de composition, les contractions III et IV peuvent être composées. Autrement dit : pour former toutes les formes issues par contraction d'une expression donnée, il suffit d'appliquer les contractions I et II.}

Preuve : d'après le théorème d'identité, respectivement le théorème du produit pour les matrices, il vient, en tenant compte de (a) :
\[
(\widehat{st}\sigma x)=-t_\sigma,\qquad
(\widehat{\sigma\tau}su)=-s_\tau,
\]
\[
(\widehat{st}\tau x)=s_\tau,\qquad
(\widehat{\sigma\tau}tu)=t_\sigma,
\]
\[
(stu)(\sigma\tau x)=s_\tau+t_\sigma-u_x\cdot s_\tau t_\sigma .
\]
Par composition des contractions I et II naissent donc les contractions III et IV, aussi bien lorsqu'on applique la contraction II à la contraction I, c'est-à-dire aux facteurs $(stu)u_\sigma u_\tau$, que lorsqu'on applique la contraction I à la contraction II, aux facteurs $(\sigma\tau x)s_xt_x$.

Nous définissons comme \emph{série de formes}\footnote{Voir Clebsch, \emph{Über eine Fundamentalaufgabe der Invariantentheorie} (Abh. der Gött. Ges. d. Wiss. t. XVII) : \S{} 17 et fin du \S{} 18. La \og{} série de formes \fg{} se distingue du \og{} système équivalent proprement réduit \fg{} qui y est introduit par le principe de l'ordre selon les formes supérieures.} une forme initiale avec toutes les formes qui en sont issues par contraction en elle-même, et nous désignons la série de formes par la forme initiale. On appellera ici \og{} forme supérieure \fg{} toute forme ayant davantage de contractions en elle-même, tandis que parmi les contractions de même rang $s_\tau$ et $t_\sigma$, l'un doit être normalisé comme supérieur. De la loi de formation de la série de formes, il résulte :

1) Toute forme de la série de formes est linéaire dans les symboles de la forme initiale.

2) La série de formes est déterminée univoquement pour des formes initiales ayant chacune deux séries de symboles contragrédients ; pour des formes initiales ayant davantage de séries de symboles, elle doit être normalisée univoquement par distinction de contractions particulières parmi celles que relie le théorème d'identité.

D'après le théorème I, une série de formes $s_x^m t_x^n u_\sigma^\mu u_\tau^\nu$ ($m>n;\ \mu>\nu$) peut être rangée selon le schéma rectangulaire suivant :
\begingroup\small
\[
\begin{array}{c|c|c|c|c}
s_x^m t_x^n u_\sigma^\mu u_\tau^\nu
& (stu)
& (stu)^2
& \cdots
& \noethpIIrosette\\[0.3em]
(\sigma\tau x)
& t_\sigma,\ s_\tau
& t_\sigma(stu),\ s_\tau(stu)
& \cdots
& \\[0.3em]
(\sigma\tau x)^2
& t_\sigma(\sigma\tau x),\ s_\tau(\sigma\tau x)
& t_\sigma^2,\ t_\sigma s_\tau,\ s_\tau^2
& \cdots
& \vdots\\[0.3em]
\vdots&\vdots&\vdots&\ddots&\vdots\\[0.3em]
(\sigma\tau x)^\nu
& t_\sigma(\sigma\tau x)^{\nu-1},\ s_\tau(\sigma\tau x)^{\nu-1}
& t_\sigma^2(\sigma\tau x)^{\nu-2},\ t_\sigma s_\tau(\sigma\tau x)^{\nu-2},\ s_\tau^2(\sigma\tau x)^{\nu-2}
& \cdots
& \vdots\\[0.3em]
\vdots
& t_\sigma(\sigma\tau x)^\nu
& t_\sigma^2(\sigma\tau x)^{\nu-1},\ t_\sigma s_\tau(\sigma\tau x)^{\nu-1}
& \cdots
& \vdots\\[0.3em]
\vdots
& \vdots
& t_\sigma^2(\sigma\tau x)^\nu
& \cdots
& \vdots
\end{array}
\]
\endgroup
\begingroup\small
\[
\begin{array}{r@{\quad}c}
\noethpIIrosette\srcfnmark{*)}
&
\begin{array}{c|c|c}
(stu)^n & \cdots & \cdots\\[0.3em]
t_\sigma(stu)^{n-1},\ s_\tau(stu)^{n-1}
& s_\tau(stu)^n & \cdots\\[0.3em]
t_\sigma^2(stu)^{n-2},\ t_\sigma s_\tau(stu)^{n-2},\ s_\tau^2(stu)^{n-2}
& t_\sigma s_\tau(stu)^{n-1},\ s_\tau^2(stu)^{n-1}
& s_\tau^2(stu)^n .
\end{array}
\end{array}
\]
\endgroup
\srcfntext{*)}{Ici, et de façon analogue dans la suite, la partie inférieure indiquée par \NoetherRosetteInline{} doit être raccordée à l'endroit du schéma situé en haut à droite et désigné de même.}

On obtient ici, lorsque l'on

1) avance d'une colonne vers la droite, toutes les formes qui naissent des formes adjacentes par contraction I,

2) avance d'une ligne vers le bas, toutes les formes qui naissent des formes situées au-dessus par contraction II,

et donc toutes les formes qui naissent par contraction.

Une forme quelconque de la série totale : $t_\sigma^\kappa s_\tau^\lambda(stu)^\mu$ ou $t_\sigma^\kappa s_\tau^\lambda(\sigma\tau x)^\nu$ constitue le début d'une nouvelle série de formes, composée de toutes les formes du schéma rectangulaire ayant la forme $t_\sigma^\kappa s_\tau^\lambda(stu)^\mu$ ou $t_\sigma^\kappa s_\tau^\lambda(\sigma\tau x)^\nu$ comme coin supérieur gauche et possédant le facteur $t_\sigma^\kappa s_\tau^\lambda$.

Remarque finale. Le théorème I subit une exception dans le cas où les nombres $n$ et $\mu$ (respectivement $m$ et $\nu$) deviennent nuls, de sorte que les contractions I, II et IV n'existent plus, tandis que la contraction III (respectivement IV) existe encore. Dans ce cas, nous appelons série de formes le terme diagonal du schéma :
\[
\begin{array}{ccccc}
s_x^m u_\tau^\nu &&&& t_x^n u_\sigma^\mu\\
& s_\tau && t_\sigma &\\
& s_\tau^2 && t_\sigma^2 &\\
&&\ddots&&\\
&&s_\tau^\nu && t_\sigma^\mu .
\end{array}
\]
Conformément à l'absence des contractions I et II, le développement en série suivant les polaires de la série de formes se réduit alors toujours au terme initial ; les théorèmes sur les réduisants restent toutefois valables.

\begin{center}
\textbf{\S{} 2. \quad Développements en série suivant les polaires de la série de formes.}
\end{center}

Pour les formes à $n$ variables, deux espèces de développements en série ont été explicitement établies\footnote{Voir Gordan, \emph{Über Kombinanten}. \S{} 2 et 5 (\emph{Math. Ann.} t. V).} :

1) pour les formes ayant chacune une série de variables contragrédientes : $s_x^m u_\tau^\nu$, une série procédant suivant les puissances de $u_x$, dont les coefficients sont des \og{} formes normales \fg{} ;

2) pour les formes à deux séries de variables cogrédientes $s_x^m t_x^n$, un développement suivant les polaires des covariants élémentaires (polaires de la série de formes $s_x^m t_x^n$), qui dans le cas ternaire s'écrit comme suit\footnote{Voir Study, \emph{Methoden zur Theorie der ternären Formen} (Teubner 1889) II. \S{} 7. Contrairement à celle donnée là, notre déduction fournit la méthode pour déterminer rapidement par le calcul les coefficients numériques $C_{pr\varrho}$. Le développement de Clebsch--Study, sous la spécialisation a), a'), serait :
\[
(stu)^\mu u_\tau^{\nu-\varkappa}v_\tau^\varkappa s_x^m t_x^{\,n-\lambda}(tuv)^\lambda
=\sum_{p,r,\varrho} C'_{pr\varrho}
\bigl[s_\tau^\varrho(stu)^{\mu+r-\varrho}\bigr]_{\widehat{uv}^{\lambda-r+\varrho}v^{\varkappa-\varrho},\,w^p,\,(w=x\widehat{uv})},
\]
et ne permet donc pas la simplification qui intervient déjà, dans notre développement, au cours du calcul, lorsque l'on sait que tous les termes contenant le facteur $u_x$ doivent être omis.} :
\[
\tag{I}
s_x^m t_x^{n-\lambda}t_y^\lambda
=\sum_{\varrho}
\frac{\binom{m}{\varrho}\binom{\lambda}{\varrho}}
{\binom{m+n-\varrho+1}{\varrho}}\cdot
\bigl[(st\widehat{xy})^\varrho s_x^{m-\varrho}t_x^{n-\varrho}\bigr]_{y^{\lambda-\varrho}}.
\]

Nous combinons les deux développements en établissant, pour les formes ayant chacune deux séries de variables contragrédientes :
\[
s_x^m t_x^{n-\lambda}t_y^\lambda u_\sigma^\mu u_\tau^{\nu-\varkappa}v_\tau^\varkappa
\]
des développements suivant les puissances de $u_x$, dont les coefficients sont des polaires composées de la série de formes $s_x^m t_x^n u_\sigma^\mu u_\tau^\nu$, prises selon les variables $x$ et $u$.

Il suffit d'établir la série pour deux séries contragrédientes de symboles (respectivement de variables), car, en réunissant deux à deux les séries de symboles (de variables) en une seule nouvelle série, les formes ayant davantage de séries de symboles (de variables) se ramènent à ce cas.

Afin que toutes les formes de la série de formes ne contiennent que deux séries de symboles, respectivement de variables, nous spécialisons :
\[
\begin{array}{ll}
\text{a) } \sigma=\widehat{st}, & \text{a') } y=\widehat{uv},\\[0.3em]
\text{b) } s=\widehat{\sigma\tau}, & \text{b') } v=\widehat{xy}
\end{array}
\]
et nous effectuons le développement pour la spécialisation a), a').

Par transfert direct du développement I, on obtient des polaires composées pour les formes $(stu)^\mu s_x^m t_x^{n-\lambda}t_y^\lambda$, tandis que pour les formes $(stu)^\mu u_\tau^{\nu-\varkappa}v_\tau^\varkappa s_x^m t_x^{n-\lambda}t_y^\lambda$, au lieu de polaires composées, apparaissent des termes dont l'évaluation rend nécessaire l'introduction de variables auxiliaires toujours nouvelles, à n'identifier qu'à la fin, ce qui complique inutilement le calcul.

Nous suivons donc une voie indirecte, en représentant les polaires composées de toutes les formes de la série de formes, c'est-à-dire les expressions
\[
\bigl[s_\tau^\varrho(stu)^{\mu-\varrho+r}\bigr]_{y^\lambda}^{v^\varkappa}
\]
par les termes initiaux de ces polaires, puis en obtenant le développement voulu suivant les polaires par inversion du système récurrent d'équations qui en résulte.

D'après le principe de transfert, le développement suivant vaut pour la $\lambda$-ième polaire d'une forme $s_x^m t_x^n:[s_x^m t_x^n]_{y^\lambda}$ ($\lambda\le n$) (le cas $\lambda>n$ se ramène au premier par le développement de $[s_y^m t_y^n]_{x^{m+n-\lambda}}$)\footnote{Gordan, \emph{Die Resultante binärer Formen} (chap. I, \S{} 5). Rend. Circ. Mat. di Palermo XXII. 1906.} :
\[
\bigl[s_x^m t_x^n\bigr]_{y^\lambda}
=\sum_r (-1)^r\,
\frac{\binom{m}{r}\binom{\lambda}{r}}{\binom{m+n}{r}}\cdot
(st\widehat{xy})^r s_x^{m-r}t_x^{n-\lambda}t_y^{\lambda-r},
\]
et de même pour la $\varkappa$-ième polaire de $u_\sigma^\mu u_\tau^\nu:[u_\sigma^\mu u_\tau^\nu]^{v^\varkappa}$
\[
\bigl[u_\sigma^\mu u_\tau^\nu\bigr]^{v^\varkappa}
=\sum_\varrho (-1)^\varrho\,
\frac{\binom{\mu}{\varrho}\binom{\varkappa}{\varrho}}{\binom{\mu+\nu}{\varrho}}\,
(\sigma\tau\widehat{uv})^\varrho u_\sigma^{\mu-\varrho}u_\tau^{\nu-\varkappa}v_\tau^{\varkappa-\varrho}.
\]

Par multiplication, il s'ensuit, en introduisant les spécialisations a) et a'),
\[
\begin{aligned}
\bigl[s_x^m t_x^n(stu)^\mu u_\tau^\nu\bigr]_{\widehat{uv}^{\lambda}}^{v^\varkappa}
={}&
\sum_{r,\varrho}c_{r\varrho}\cdot
(st\widehat{x\widehat{uv}})^r(\widehat{st}\tau\widehat{uv})^\varrho
s_x^{m-r}t_x^{n-\lambda}(tuv)^{\lambda-r}(stu)^{\mu-\varrho}
u_\tau^{\nu-\varkappa}v_\tau^{\varkappa-\varrho},
\end{aligned}
\]
et de là, par développement suivant les puissances de $u_x$, en distinguant le premier terme ($\sum'_{r,\varrho}$ signifie que le terme $r=0$, $\varrho=0$ est à omettre) et en tenant compte de (a) p. 26,
\[
\tag{II}
\begin{aligned}
&s_x^m t_x^{n-\lambda}(tuv)^\lambda(stu)^\mu u_\tau^{\nu-\varkappa}v_\tau^\varkappa\\
&\quad =
\bigl[s_x^m t_x^n(stu)^\mu u_\tau^\nu\bigr]_{\widehat{uv}^{\lambda}}^{v^\varkappa}\\
&\qquad
-\sum_{r,\varrho}' c_{r\varrho}\cdot s_\tau^\varrho(stu)^{\mu-\varrho+r}
u_\tau^{\nu-\varkappa}v_\tau^{\varkappa-\varrho}
s_x^{m-r}t_x^{n-\lambda}(tuv)^{\lambda-r+\varrho}\cdot v_x^r\\
&\qquad
+u_x\cdot\sum_{\varrho}\cdot\sum_{r=1}^{\lambda} c_{r\varrho}\cdot
s_\tau^\varrho(stu)^{\mu-\varrho+r-1}(stv)
u_\tau^{\nu-\varkappa}v_\tau^{\varkappa-\varrho}
s_x^{m-r}t_x^{n-\lambda}(tuv)^{\lambda-r+\varrho}\cdot v_x^{r-1}\\
&\qquad
\vdots\\
&\qquad
-(-1)^\lambda u_x^\lambda\cdot\sum_\varrho c_{\lambda\varrho}\cdot
s_\tau^\varrho(stu)^{\mu-\varrho}(stv)^\lambda
u_\tau^{\nu-\varkappa}v_\tau^{\varkappa-\varrho}
s_x^{m-\lambda}t_x^{n-\lambda}(tuv)^\varrho,
\end{aligned}
\]
où
\[
c_{r\varrho}=(-1)^{r+\varrho}\cdot
\frac{\binom{m}{r}\binom{\lambda}{r}}{\binom{m+n}{r}}\cdot
\frac{\binom{\mu}{\varrho}\binom{\varkappa}{\varrho}}{\binom{\mu+\nu}{\varrho}}
\qquad
\begin{array}{l}
\text{pour }\lambda\le n,\\
\varkappa\le\nu.
\end{array}
\]
et de même pour les autres combinaisons des nombres
\[\lambda,n;\varkappa,\nu.\]
L'expression
\[s_x^m t_x^{n-\lambda}(tuv)^\lambda(stu)^\mu u_\tau^{\nu-\varkappa}v_\tau^\varkappa\]
est donc ramenée à une polaire composée et, par application de l'identité
\[
(stv)u_\tau=(stu)v_\tau-s_\tau(tuv),
\]
à une somme d'expressions formées de manière analogue, qui correspondent à des formes supérieures de la série de formes.

Par itération on parvient au développement :
\[
\tag{III}
s_x^m t_x^{n-\lambda}(tuv)^\lambda(stu)^\mu u_\tau^{\nu-\varkappa}v_\tau^\varkappa
=
\sum_{p,r,\varrho}u_x^p\cdot C_{pr\varrho}\cdot
\bigl[s_\tau^\varrho(stu)^{\mu-\varrho+r}\bigr]_{\widehat{uv}^{\lambda-r+\varrho}v^{\varkappa-\varrho+p}\cdot v_x^{r-p}}.
\]
Les coefficients numériques $C_{pr\varrho}$ se calculent univoquement et récursivement à partir des coefficients $c_{r\varrho}$, $c'_{r\varrho}$ du système d'équations obtenu par itération de (II.) (voir l'exemple p. 42). Des développements analogues s'obtiennent pour les autres spécialisations.

\begin{center}
\textbf{\S{} 3. \quad Théorèmes de réduction.}
\end{center}

Les théorèmes connus sur la réduction des formes et des systèmes de formes doivent maintenant être mis en rapport avec les paragraphes 1 et 2 ; quelques définitions empruntées à la théorie des formes binaires sont nécessaires à cet effet.

a) Nous appelons \emph{système relativement complet}, modulo une série de formes donnée, un système de formes ayant la propriété que toutes les formations qui naissent par contraction d'un produit quelconque de ces formes peuvent s'exprimer comme fonctions rationnelles entières des formes du système, à un terme additif près d'expressions qui sont nées par contraction des formes du système avec le système du module\footnote{Gordan--Kerschensteiner, \emph{Vorlesungen über Invariantentheorie}. t. II, p. 227.}.

b) Nous appelons une forme \emph{réductible} lorsqu'elle peut s'exprimer par des formes qui ont des invariants pour facteurs ou par des \og{} formes supérieures \fg{} ; c'est-à-dire par des formes supérieures de la série totale de formes à laquelle elle appartient, ou par des formes qui contiennent les symboles du module à un ordre supérieur.

Les théorèmes sur les réduisants\footnote{Gordan, \emph{Math. Annalen} t. XVII, p. 222.} peuvent alors s'énoncer sous leur forme la plus générale de la façon suivante :

Définition : \emph{Un réduisant est une série de formes réductible.}

\emph{Théorème II. Si la forme initiale d'une série de formes est réductible du fait que l'un de ses termes est issu par contraction avec un réduisant (a un réduisant pour facteur), et si la forme finale de la série de formes est née de ce même terme par contraction, alors toute la série de formes est réductible}\footnote{La simplification apportée par le théorème II se voit dans l'exemple suivant : pour la forme spéciale $f=x_1^3x_2+x_2^3x_3+x_3^3x_1$, $a_\nu^2a_x^2u_\nu^2$ est un réduisant. Il en résulte, d'après le théorème II, la réduction de la série de formes
\[
\theta_\nu(\vartheta\nu x)\theta_x^3u_\vartheta u_\nu^2
=
a_\nu^2(abu)-\underbrace{a_\nu b_\nu(abu)}_{0},
\]
puisque $\theta_\nu^4=a_\nu^2b_\nu^2(abu)^2$ est issue du terme $a_\nu^2(abu)$ par contraction. Chez Gordan, loc. cit. p. 231, la réduction des formes
\[
\theta_\nu(\vartheta\nu x),\quad \theta_\nu(\vartheta\nu x)^2,\quad
\theta_\nu^2,\quad \theta_\nu^2(\vartheta\nu x),\quad
\theta_\nu^2(\vartheta\nu x)^2
\]
est en revanche effectuée séparément.}.

Preuve : La série de formes naît du terme initial par contraction en lui-même, donc d'une part par contraction supérieure avec le réduisant, et d'autre part par contraction du réduisant en lui-même. Dans les deux cas, d'après \S{} 2, toutes les formes de la série de formes peuvent être représentées comme polaires (transvections) sur la série de formes du réduisant.

Le passage de la forme initiale à toute la série de formes n'a plus lieu pour les autres méthodes de réduction. Celles-ci sont :

1) la dite \og{} réduction double \fg{}.

Nous entendons par là la réduction, de deux manières, d'une expression qui a deux réduisants indépendants l'un de l'autre pour facteurs, ce qui fait naître des relations entre les formes supérieures auxquelles on a réduit.

Par l'application systématique de la \og{} réduction double \fg{}, on doit obtenir d'une part toutes les relations\footnote{Ainsi ont été trouvées par exemple, au moyen de la \og{} réduction double \fg{} : la réduction de la forme $a_\eta^3$ (\S{} 10), la seule forme que Maisano ait introduite en trop dans le système ; puis les relations mentionnées dans la note à l'introduction entre les 3 covariants et les 3 invariants.}, et d'autre part, lorsque la \og{} réduction double \fg{}, appliquée à diverses expressions, ne fournit pas de nouvelles relations, on dispose à la fois d'un critère pour l'indépendance des formes supérieures et d'un contrôle du calcul.

2) \emph{Contraction avec des formes décomposables.}

Par \og{} formes décomposables \fg{} on entendra --- avec une légère généralisation de la notion usuelle --- des formes qui peuvent s'exprimer par des produits de formes de degré inférieur et par des formes \og{} supérieures \fg{}. Les produits de formes passent, lors d'une contraction simple, en produits ; de même, les produits de formes non linéaires lors d'une contraction double dans les spécialisations a') et b') (\S{} 2), d'après le théorème du produit :
\[
a_yb_y a_xb_x=\frac12a_y^2\cdot b_x^2+\frac12b_y^2\cdot a_x^2-\frac12(ab\widehat{yx})^2.
\]
Les premières polaires (transvections) et certaines secondes polaires de formes décomposables sont donc à nouveau des formes décomposables. Il s'ensuit :

Un terme d'une transvection simple (respectivement double) sur une forme décomposable, né par contraction simple (respectivement double) avec une forme décomposable, devient lui-même une forme décomposable :

a) si les formes immédiatement supérieures dans la série de formes de la forme décomposable se décomposent ou deviennent réductibles, ou encore si, lors d'une contraction simple, les spécialisations $y=\widehat{uv}$ ou $v=\widehat{xy}$ se présentent ;

b) si les autres contractions simples conduisent à des formes supérieures (voir \S{} 12, B. $H_k$ et $H_k(k\eta x)$).

Un tel terme d'une transvection peut aussi se décomposer :

c) si les autres contractions simples conduisent à des formes inférieures qui se décomposent indépendamment de la contraction avec la forme décomposable, c'est-à-dire peuvent s'exprimer par des produits et par la forme à réduire, mais où l'on doit chaque fois s'assurer par le calcul que le coefficient numérique du terme considéré ne s'annule pas identiquement. Ce n'est rien d'autre que la \og{} réduction double \fg{} de la forme inférieure (voir \S{} 23, C. $H_\vartheta(\theta H u)(\theta s u)$).

Des formes décomposables naissent par toutes les contractions simples avec des déterminants fonctionnels\footnote{voir Gordan, loc. cit. \S{} 3.}, et en outre par certaines contractions supérieures. On a
\[
(st\widehat{yx})s_y=\frac12\,t\cdot s_y^2-\frac12\,s\cdot t_y^2+\frac12(st\widehat{yx})^2,
\]
\[
(st\widehat{yx})s_y^2=\frac13\,t\cdot s_y^3-\frac13\,s\cdot t_y^3+s_y(st\widehat{yx})^2-\frac13(st\widehat{yx})^3.
\]
Des formules analogues valent pour les formes dualistiques
\[
(\sigma\tau\widehat{vu})v_\sigma
\quad\text{et}\quad
(\sigma\tau\widehat{vu})v_\sigma^2.
\]

D'après les théorèmes de ce paragraphe, il résulte, pour l'établissement de toutes les formes irréductibles qui naissent de la contraction de $S$ avec $T$, la règle suivante :

On procède, en commençant par les contractions les plus basses, suivant une voie quelconque fixée à l'avance (par exemple ligne par ligne ou symétriquement au terme diagonal), dans la formation de la série de formes $ST$, jusqu'à parvenir à une forme ayant un réduisant pour facteur. La série de formes définie par cette forme doit être omise dès que la forme finale satisfait à la condition énoncée dans le théorème II. Après élimination de toutes les séries de formes réductibles, il faut encore enlever, par application de la réduction double, toutes les autres formes réductibles ainsi que toutes les formes décomposables. Les emplacements de la série totale de formes qui sont recouverts deux fois par des séries de formes réductibles indépendamment les unes des autres donnent lieu à la réduction de formes supérieures (voir \S{} 11, D).

\begin{center}
\textbf{Chapitre II. \quad Systèmes fondamentaux relativement complets.}

\textbf{\S{} 4. \quad Premier système relativement complet.}
\end{center}

Dans la suite, nous nous limitons à la forme biquadratique, à l'exception du premier théorème. Toutefois, moyennant un choix approprié des modules, les résultats des \S\S{} 5 et 6 se transposent aisément aux formes de degré supérieur, tout comme ils valent pour la forme cubique\footnote{Pour la forme cubique, les formules correspondantes, déduites de manière analogue, sont :
\[
(abc)a_x^2b_y^2c_z^2
=\frac13(sxy)(sxz)(syz)+\frac23\Delta_x\Delta_y\Delta_z\cdot(xyz),
\qquad (\mathrm{IV}.)
\]
\[
\begin{gathered}
s=(abc)(abu)(acu)(bcu),\qquad \Delta=(abc)^2a_xb_xc_x,\\
t=\Delta_\vartheta(\Delta\theta u)^2u_\vartheta
=\theta_s^2u_su_\vartheta^2,
\qquad T=a_t^3,
\end{gathered}
\]
\[
\left.
\begin{aligned}
(a\theta u)^2&=\frac13u_xs,\\
a_\vartheta&=\frac13u_x\cdot\Delta,
&a_\vartheta(a\theta u)&=0,
&a_\vartheta(a\theta u)^2&=s,\\
a_\vartheta^2&=\Delta,
&a_\vartheta^2(a\theta u)&=0.
\end{aligned}
\right\}\tag{1}
\]
}.

On obtient un premier système relativement complet pour les formes de degré quelconque d'après le théorème connu suivant :


\clearpage
\noindent\textbf{Théorème III.}
\emph{Les formes reprises du domaine binaire, c'est-à-dire les formes qui naissent des formes du système de la forme binaire correspondante par bordage selon le principe de transfert de Clebsch, forment un système relativement complet mod $(abc)$.}

\noindent\textbf{Preuve :}
À une relation binaire exprimant que les formes $Q'_1,\ldots,Q'_n$ forment un système complet d'une forme fondamentale binaire :
\[
Q'-F(Q')=0=\sum_\lambda\bigl\{(ab)c_x+(bc)a_x+(ca)b_x\bigr\}\varphi_\lambda\srcfnmark{*)}
\]
\srcfntext{*)}{Gordan--Kerschensteiner, loc. cit., p.~134.}
correspond, par bordage, la relation ternaire :
\[
Q-F(Q_i)=\sum_\lambda u_x\cdot(abc)\varphi_\lambda .
\]
Or c'est l'équation de définition d'un système relativement complet mod $(abc)$. Pour la forme biquadratique, le système des formes reprises se compose des formes suivantes :
\[
\begin{gathered}
f=a_x^4,\qquad
\theta=\theta_x^4u_\vartheta^2=(abu)^2a_x^2b_x^2,\qquad
K=K_x^6u_k^3=(a\theta u)u_\vartheta^2a_x^3\theta_x^3,\\
j=u_j^6=(a\theta u)^4u_\vartheta^2,\qquad
\nu=u_\nu^4=(abu)^4,
\end{gathered}
\]
parmi lesquelles les quatre premières forment un système relativement complet mod $((abc),\nu)$.

\begin{center}
\textbf{\S{} 5. \quad Réduction du module $(abc)$ aux modules $\A$ et $\nu$.}
\end{center}

Le module $(abc)$, donné seulement par un facteur entre crochets, doit, conformément à la définition (\S{} 3, a), être ramené à des formes du système. Autrement dit : la série de formes $(abc)$ doit être normalisée univoquement (voir \S{} 1).

\[
a_s^2a_xu_s=\frac13a_s^3\cdot u_x=\frac13S\cdot u_x. \tag{2}
\]
\[
\left.
\begin{aligned}
(a\A u)^2&=a_s-\frac{S}{6}u_x^2,\\
(a\A u)^3&=0.
\end{aligned}
\right\} \tag{3}
\]
\[
\left.
\begin{aligned}
\theta(s)=(ss_1x)^2&=2a_t+\frac13S\cdot\theta-\frac13Tu_x^2,\\
\theta(t)=(tt_1x)^2&=-\frac13S\cdot a_t+\frac23T\cdot a_s+\frac1{18}S^2\cdot\theta-\frac1{18}u_x^2\cdot S\cdot T.
\end{aligned}
\right\} \tag{4}
\]

Série des modules : $(abc)$ ; $\A$ ; $s$ ; $t$ (le système du module $t$ se compose de la seule forme $t$).

Nous montrerons que les formes
\[
\A=(abc)^2a_x^2b_x^2c_x^2,\qquad \nu=(abu)^4
\]
suffisent comme module, c'est-à-dire que les formes
\[
\begin{gathered}
\A=(abc)^2a_x^2b_x^2c_x^2,\qquad
a_\nu=(abc)(bcu)^3a_x^3,\\
a_\nu^2=(abc)^2(bcu)^2a_x^2,\qquad
i=a_\nu^4=(abc)^4
\end{gathered}
\]
définissent la série de formes $(abc)$. Dans la série de formes $a_\nu$, la forme $a_\nu^3$ manque d'après l'identité :
\[
a_\nu^3=(abc)^3a_x(bcu)=\frac13 u_x\cdot(abc)^4=\frac13 i\cdot u_x.
\]

Pour ramener un module à un module supérieur, nous devons, d'après la définition, ramener les contractions les plus générales, ou encore toutes les contractions particulières entrant en considération avec le module, à des contractions avec le module supérieur.

Les contractions les plus générales naissent, dans notre cas, des expressions :
\[
(abc)a_x^3b_y^3c_z^3,\qquad
(abc)a_x^2b_y^2c_z^2a_xb_xc_x
\]
(reprises des formes cubiques),
\[
(abc)a_xb_yc_za_x^2b_x^2c_x^2
\]
(reprise des formes quadratiques), et par polarisation de ces expressions.

Pour calculer ces expressions par les polaires de $\A$ et de la série de formes $a_\nu$, respectivement par leurs termes initiaux, nous suivons, comme au \S{} 2, la voie indirecte. Nous développons les termes polaires
\[
(abc)^2a_x^2b_y^2c_z^2,\qquad
a_\nu(\nu xy)(\nu yz)(\nu zx)a_xa_ya_z
=(abc)(bc\widehat{xy})(bc\widehat{yz})(bc\widehat{zx})a_xa_ya_z\quad\text{etc.}
\]
comme agrégats linéaires d'expressions ayant le facteur symbolique $(abc)$, et obtenons par inversion du système d'équations les relations cherchées, ainsi qu'une relation entre les polaires prises selon différentes combinaisons des variables. Pour l'évaluation servent le théorème d'identité et le théorème du produit pour les déterminants.

Le système d'équations, pour $(abc)a_x^3b_y^3c^3$, est :
\begingroup\small
\[
\begin{aligned}
1)\quad
\sum_3 a_\nu(\nu yz)^3a_x^3
&=6(abc)a_x^3b_y^3c_z^3
-6\sum_3(abc)a_x^3b_y^2b_zc_z^2c_y\footnotemark,\\[0.4em]
2)\quad
3(abc)^2a_x^2b_y^2c_z^2\cdot(xyz)
-\frac12\sum_3a_\nu^2(\nu yz)^2a_x^2\cdot(xyz)
&=2\sum_3(abc)a_x^3b_y^2b_zc_z^2c_y\\
&\quad-4\sum_3(abc)a_xa_ya_zb_y^2b_xc_z^2c_x,\\[0.4em]
3)\quad
a_\nu(\nu xy)(\nu yz)(\nu zx)a_xa_ya_z
&=2\sum_3(abc)a_xa_ya_zb_y^2b_xc_z^2c_x,\\[0.4em]
4)\quad
\sum_3a_\nu(\nu yz)^3a_x^3
-\frac32\sum_3a_\nu^2(\nu yz)^2a_x^2\cdot(xyz)
+\frac16 i\cdot(xyz)^3
&=6\sum_3(abc)a_xa_ya_zb_y^2b_xc_z^2c_x.
\end{aligned}
\]
\endgroup
\footnotetext{$\sum_3$ désigne la somme de trois termes obtenue à partir du terme initial par permutation cyclique des variables. Les équations valent aussi séparément pour chacun des termes de la somme.}

Il en résulte pour
\[
(abc)a_x^3b_y^3c_z^3,
\]
et, par un calcul analogue, pour
\[
(abc)a_x^2b_y^2c_z^2a_xb_xc_x,\qquad
(abc)a_xb_yc_za_x^2b_x^2c_x^2:
\]
\[
\begin{aligned}
(abc)a_x^3b_y^3c_z^3
&=\frac32(abc)^2a_x^2b_y^2c_z^2\cdot(xyz)
+\frac32 a_\nu(\nu xy)(\nu yz)(\nu zx)a_xa_ya_z\\
&\quad-\frac1{12}i\cdot(xyz)^3,\\[0.4em]
(\text{IV.}.)\quad
(abc)a_x^2b_y^2c_z^2a_xb_xc_x
&=\frac23(abc)^2a_x^2b_yb_xc_zc_x\cdot(xyz)
\\
&\quad+\frac13a_\nu(\nu xy)(\nu yz)(\nu zx)a_x^3
-\frac16 a_\nu^2(\nu xy)(\nu zx)a_x^2\cdot(xyz),\\[0.4em]
(abc)a_xb_yc_za_x^2b_x^2c_x^2
&=\frac16(abc)^2a_x^2b_x^2c_x^2\cdot(xyz).
\end{aligned}
\]

Nous avons ainsi obtenu un système relativement complet mod $(\A,\nu)$, composé des quatre formes : $f,\theta,K,j$. Il en résulte que toutes les transvections de $f$ sur $\theta$ qui ne sont pas reprises du domaine binaire, de même que la transvection $(a\theta u)^2$, réductible dans le binaire à $(ab)^4$, peuvent s'exprimer par les symboles $\A$ et $\nu$.

Nous obtenons les formules de réduction fondamentales pour la suite, par spécialisation du développement IV ou, plus brièvement, par calcul direct (permutation des symboles), pour $a_\vartheta(a\theta u)^2$, $a_\vartheta^2(a\theta u)^2$, $(a\theta u)^2$ selon l'ansatz suivant :
\[
\begin{aligned}
a_\vartheta(a\theta u)^2
&=(abc)(bcu)(abu)^2-\frac13(abc)(bcu)\{(bcu)a_x-u_x\cdot(abc)\}^2,\\
a_\vartheta^2(a\theta u)^2
&=(abc)^2(abu)^2-\frac13(abc)^2\{(bcu)a_x-u_x\cdot(abc)\}^2
\end{aligned}
\]
pour $(a\theta u)^2$\srcfn{*)}{Gordan--Kerschensteiner, loc. cit., p.~92.} :
\[
\begin{aligned}
(abu)a_y^3b_x^3&=\frac32u_\vartheta(\vartheta yx)\theta_y^2+\frac14(\nu yx)^3,\\
(abu)(acu)^3b_x^3&=\frac32u_\vartheta(\vartheta\widehat{au}x)(\theta au)^2+\frac14(\nu\widehat{au}x)^3:
\end{aligned}
\]
\[
\left.
\begin{aligned}
(a\theta u)^2&=\frac16\nu\cdot f-\frac23u_x\cdot a_\nu
+\frac23u_x^2\cdot a_\nu^2-\frac{i}{18}u_x^4,\\
\underline{a_\vartheta}&=\frac13u_x\cdot\A,\\
a_\vartheta(a\theta u)&=0,\\
a_\vartheta^2&=\A,\\
(a\theta u)^3&=0,\\
a_\vartheta(a\theta u)^2&=-\frac56a_\nu+\frac76u_x\cdot a_\nu^2-\frac{i}{9}u_x^3,\\
a_\vartheta^2(a\theta u)&=0,\\
(a\theta u)^4&=j,\\
a_\vartheta(a\theta u)^3&=0,\\
a_\vartheta^2(a\theta u)^2&=\frac23a_\nu^2-\frac{i}{9}u_x^2.
\end{aligned}
\right\} \tag{1}
\]

Nous y joignons la formule, issue elle aussi de la réduction du module $(abc)$ :
\[
a_\nu^3a_xu_\nu=\frac13 i\cdot u_x. \tag{2}
\]

\emph{Des formules (1.) et (2.) et des formules construites de façon analogue pour la forme fondamentale $\nu$, toutes les formules de réduction ultérieures se déduisent par polarisation}, à l'exception des relations pour les formes décomposables. Des formules (1.) nous voyons :

La série de formes :
\[
\begin{array}{rcl}
a_\vartheta(a\theta u)^2 &\text{conduit à}& \text{des symboles }\nu\text{ seuls},\\[0.2em]
a_\vartheta &\text{conduit à}& \text{des symboles }\A\text{ et }\nu,\\[0.2em]
(a\theta u)^2 &\text{conduit à}& \text{des symboles }j,\A\text{ et }\nu,\\[0.2em]
(a\theta u)^2\text{ sous la contraction I} &\text{conduit à}& \text{des symboles }j\text{ et }\nu,\\[0.2em]
(a\theta u)^2\text{ sous la contraction II} &\text{conduit à}& \text{des symboles }\A\text{ et }\nu.
\end{array}
\]

Nous pouvons donc remplacer :
\begin{enumerate}
\item mod $(\A,\nu)$ : les contractions avec $j$ par les contractions correspondantes avec $(a\theta u)^2$ ou $(a\theta u)^3$ ;
\item mod $(\nu)$ : les contractions avec $\A$ par les contractions correspondantes avec $a_\vartheta$ ou $a_\vartheta(a\theta u)$, ou encore par des contractions particulières avec $(a\theta u)^2$ (qui ne conduisent qu'à une contraction I simple).
\end{enumerate}
En particulier, on a :
\[
\begin{aligned}
(a\theta u)^2a_y^2\theta_x^2&=\frac13(jyx)^2\pmod{\nu},&
(a\theta u)^2a_y\theta_y&=-\frac16(jyx)^2\pmod{\nu},\\[0.4em]
(a\theta u)^2v_\vartheta^2&=\frac13(\A uv)^2\pmod{\nu},&
(a\theta u)(a\theta v)u_\vartheta v_\vartheta&=-\frac16(\A uv)^2\pmod{\nu},\\[0.4em]
(a\theta u)^2a_yv_\vartheta^2&=\frac13(\A uv)^2\A_y\pmod{\nu}.
\end{aligned}
\]

\begin{center}
\textbf{\S{} 6. \quad Réduction du module $(\A,\nu)$ au module $(\nu)$.}
\end{center}

Les formules de réduction du paragraphe précédent nous donnent le moyen d'établir directement le système relativement complet mod $\nu$ ; nous verrons que nous n'avons à adjoindre au système des formes reprises que les formes $\A$ et $(a\A u)=N$ (de manière analogue au cas de la forme cubique, et plus généralement valable).

Pour réduire le module $\A$, nous considérons toutes les contractions particulières entrant en ligne de compte, c'est-à-dire les contractions du système relativement complet mod $(\A,\nu)$ avec le système de $\A$, et nous partons des formes de plus bas ordre dans les coefficients, les contractions de $f$ avec $\A$.

Pour réduire les formes $(a\A u)^2$, $(a\A u)^3$, $(a\A u)^4$, nous remplaçons, d'après \S{} 5, les symboles $\A$ par des formes inférieures de la série de formes, c'est-à-dire que nous développons les expressions
\[
a_\vartheta b_\vartheta(b\theta u)^2,\qquad
a_\vartheta(a\theta u)b_\vartheta(b\theta u)^2,\qquad
a_\vartheta(a\theta u)b_\vartheta(b\theta u)^3
\]
1) suivant les polaires de la série de formes $a_\vartheta$ (symboles $\A$ et $\nu$),

2) suivant les polaires de la série de formes $b_\vartheta(b\theta u)^2$ (symboles $\nu$)

et nous obtenons les formules de réduction (calcul voir ci-dessous) :
\begin{align}
\text{a)}\quad (a\A u)^2
&=\frac12(\vartheta\nu x)^2-\frac25 f\cdot a_\nu^2
-\frac45u_x\cdot\theta_\nu(\vartheta\nu x)^2
+u_x^2\left\{\frac16i\cdot f-\frac1{40}s\right\},\notag\\
\text{b)}\quad (a\A u)^3
&=-\frac3{10}\theta_\nu(\vartheta\nu x),\tag{3}\\
\text{c)}\quad \underline{(a\Delta u)^4}
&=\frac{21}{10}\theta_\nu^2-\frac3{10}H
-\frac65u_x\cdot\theta_\nu^3+\frac1{10}u_x^2\cdot\varrho.\notag
\end{align}

(Les mêmes résultats sont aussi obtenus par la réduction double des expressions $a_\vartheta^2(b\theta u)^2$, $a_\vartheta^2(b\theta u)^3$, $a_\vartheta^2(a\theta u)^2(b\theta u)^2$ et $a_\vartheta^2(a\theta u)(b\theta u)^3$, après élimination de $a_j^2$.)

Des formules (3.) il résulte que $(a\A u)^2$ est un réduisant relativement au module $\nu$. On en déduit :

1) La réduction du système de $\A$ : la série de formes $(\A\A' u)^2\equiv a_\vartheta^2(a\A u)^2$ a le réduisant $(a\A u)^2$ pour facteur.

2) La réduction du système né par contraction avec le module $\A$ :

La série de formes $(\theta\A u)^2$ a le réduisant $(a\A u)^2$ pour facteur.

Pour les formes $(\theta\A u)$ et $\A_\vartheta$, on obtient :
\[
(a\A u)^2(abu)=(\theta\A u)-u_x\cdot\A_\vartheta(\theta\A u),
\]
\[
(a\A u)(a\A b)=-\frac12\A_\vartheta+\frac12u_x\cdot\A_\vartheta^2.
\]

Mais du réduisant $(\theta\A u)$ résulte la réduction du système né par contraction avec le module $\A$.

\bigskip


Le système relativement complet mod $\nu$ se compose donc des six formes :
\[
f,\ \theta,\ K,\ j,\ \A,\ N.
\]

Comme exemple du développement en série du \S{} 2, nous donnons en détail la déduction de la formule (3.)a ; dans les calculs ultérieurs, on écrira directement la formule finale III. (Les polaires des formes nulles ont déjà été omises pendant le calcul ; la première ligne donne les valeurs inscrites d'après le développement II, la seconde ligne les valeurs finales trouvées récursivement en commençant par la dernière équation du système.) D'après le développement II, on obtient :


\begin{align*}
\text{1)}\quad&m=3,\quad n=4,\quad \lambda=2,\quad \mu=0,\quad \nu=\varkappa=1,\\
a_\vartheta b_\vartheta(b\theta u)^2
&=[a_\vartheta]_{\widehat{ub^2}}b+\frac67\{a_\vartheta b_\vartheta(a\theta u)(b\theta u)-u_x\cdot a_\vartheta b_\vartheta(a\theta b)(b\theta u)\}\\
&\quad-\frac17\{a_\vartheta(a\theta u)^2b_\vartheta-2u_x\cdot a_\vartheta(a\theta u)(a\theta b)b_\vartheta+u_x^2\cdot a_\vartheta(a\theta b)^2b_\vartheta\}\\
&=\frac23(a\Delta u)^2+\frac15[a_\vartheta(a\theta u)^2]b+\frac1{30}f\cdot[a_\vartheta^2(a\theta u)^2]-\frac25u_x\cdot[a_\vartheta(a\theta u)^2]b^2\\
&\quad-\frac1{15}u_x\cdot[a_\vartheta^2(a\theta u)^2]b+\frac15u_x^2\cdot[a_\vartheta(a\theta u)^2]b^3+\frac1{30}u_x^2\cdot[a_\vartheta^2(a\theta u)^2]b^2;\\[.6em]
\text{2)}\quad&m=2,\quad n=3,\quad \lambda=1,\quad \mu=1,\quad \nu=\varkappa=1,\\
\underline{a_\vartheta(a\theta u)b_\vartheta(b\theta u)}
&=\frac25\{a_\vartheta(a\theta u)^2b_\vartheta-u_x\cdot a_\vartheta(a\theta u)(a\theta b)b_\vartheta\}+\frac12a_\vartheta^2(b\theta u)^2\\
&\quad-\frac15\{a_\vartheta^2(b\theta u)(a\theta u)-u_x\cdot a_\vartheta^2(a\theta b)(b\theta u)\}\\
&=\frac12(a\Delta u)^2+\frac25[a_\vartheta(a\theta u)^2]b+\frac1{15}[a_\vartheta^2(a\theta u)^2]\cdot f\\
&\quad-\frac25u_x\cdot[a_\vartheta(a\theta u)^2]b^2-\frac1{10}u_x\cdot[a_\vartheta^2(a\theta u)^2]b+\frac1{30}u_x^2\cdot[a_\vartheta^2(a\theta u)^2]b^2;\\[.6em]
\text{3)}\quad&m=2,\quad n=3,\quad \lambda=1,\quad \mu=1,\quad \nu=1,\quad \varkappa=2,\\
a_\vartheta(a\theta b)b_\vartheta(b\theta u)
&=\frac25\{a_\vartheta(a\theta b)(a\theta u)b_\vartheta-u_x\cdot a_\vartheta(a\theta b)^2b_\vartheta\}\\
&=\frac25[a_\vartheta(a\theta u)^2]b^2+\frac1{30}[a_\vartheta^2(a\theta u)^2]b-\frac25u_x\cdot[a_\vartheta(a\theta u)^2]b^3-\frac1{30}u_x\cdot[a_\vartheta^2(a\theta u)^2]b^2;\\[.6em]
\text{4)}\quad&m=1,\quad n=2,\quad \lambda=0,\quad \mu=2,\quad \nu=\varkappa=1,\\
a_\vartheta(a\theta u)^2b_\vartheta
&=[a_\vartheta(a\theta u)^2]b+\frac23a_\vartheta^2(a\theta u)(b\theta u)\\
&=[a_\vartheta(a\theta u)^2]b+\frac16[a_\vartheta^2(a\theta u)^2]\cdot f-\frac16u_x\cdot[a_\vartheta^2(a\theta u)^2]b;\\[.6em]
\text{5)}\quad&m=1,\quad n=2,\quad \lambda=0,\quad \mu=2,\quad \nu=1,\quad \varkappa=2,\quad(\varkappa>\nu),\\
\underline{a_\vartheta(a\theta u)(a\theta b)b_\vartheta}
&=[a_\vartheta(a\theta u)^2]b^2+\frac13a_\vartheta^2(a\theta b)(b\theta u)\\
&=[a_\vartheta(a\theta u)^2]b^2+\frac1{12}[a_\vartheta^2(a\theta u)^2]b-\frac1{12}u_x\cdot[a_\vartheta^2(a\theta u)^2]b^2;\\[.6em]
\text{6)}\quad&m=1,\quad n=2,\quad \lambda=0,\quad \mu=2,\quad \nu=1,\quad \varkappa=3,\\
a_\vartheta(a\theta b)^2b_\vartheta&=[a_\vartheta(a\theta u)^2]b^3;\\[.6em]
\text{7)}\quad&m=2,\quad n=4,\quad \lambda=2,\quad \mu=\nu=\varkappa=0,\quad (\text{d'après le développement I}),\\
a_\vartheta^2(b\theta u)^2&=[a_\vartheta^2]_{\widehat{ub^2}}+\frac1{10}\{[a_\vartheta^2(a\theta u)^2]\cdot f-2u_x\cdot[a_\vartheta^2(a\theta u)^2]b+u_x^2\cdot[a_\vartheta^2(a\theta u)^2]b^2\};\\[.6em]
\text{8)}\quad&m=1,\quad n=3,\quad \lambda=1,\quad \mu=1,\quad \nu=\varkappa=0,\quad (\text{développement I}),\\
a_\vartheta^2(a\theta u)(b\theta u)&=\frac14[a_\vartheta^2(a\theta u)^2]\cdot f-\frac14u_x\cdot[a_\vartheta^2(a\theta u)^2]b;\\[.6em]
\text{9)}\quad&m=1,\quad n=3,\quad \lambda=1,\quad \mu=\varkappa=1,\quad \nu=0,\quad (\text{développement I}),\\
a_\vartheta^2(a\theta b)(b\theta u)&=\frac14[a_\vartheta^2(a\theta u)^2]b-\frac14u_x\cdot[a_\vartheta^2(a\theta u)^2]b^2.
\end{align*}

De 1) et 4) il résulte :
\[
\begin{aligned}
(a\Delta u)^2
&=\frac65[a_\vartheta(a\theta u)^2]b+\frac15f\cdot[a_\vartheta^2(a\theta u)^2]+\frac35u_x\cdot[a_\vartheta(a\theta u)^2]b^2-\frac3{20}u_x\cdot[a_\vartheta^2(a\theta u)^2]b\\
&\quad-\frac3{10}u_x^2\cdot[a_\vartheta(a\theta u)^2]b^3-\frac1{20}u_x^2\cdot[a_\vartheta^2(a\theta u)^2]b^2,\\
(a\Delta u)^2
&=-a_\nu b_\nu+\frac35f\cdot a_\nu^2+\frac45u_x\cdot a_\nu^2b_\nu-\frac3{20}u_x^2\cdot a_\nu^2b_\nu^2-\frac1{12}u_x^2\cdot i\cdot f.
\end{aligned}
\]

(Pour la conversion en symboles $\theta$, voir \S\ 8 et \S\ 9.) La formule (3.)a pourrait encore se calculer plus brièvement par transfert direct du développement I, en introduisant la variable auxiliaire $w=x\widehat{ub}$, tandis que, dès les formules (3.)b et (3.)c, la voie que nous avons suivie devient plus simple ; pour (3.)b, on sait d'avance, d'après \S\ 9, que tous les termes contenant le facteur $u_x$ doivent être omis. On obtient, pour calculer (3.)b et (3.)c :
\[
\begin{aligned}
(3.)\mathrm{b)}\quad 1)&\quad a_\vartheta(a\theta u)b_\vartheta(b\theta u)^2=\frac12(a\Delta u)^3+\frac45[a_\vartheta(a\theta u)^2]_{\widehat{ub}}b+\frac{77}{360}[a_\vartheta^2(a\theta u)^2]_{\widehat{ub}},\\
2)&\quad a_\vartheta(a\theta u)b_\vartheta(b\theta u)^2=[a_\vartheta(a\theta u)^2]_{\widehat{ub}}b+\frac{13}{36}[a_\vartheta^2(a\theta u)^2]_{\widehat{ub}};\\[.5em]
(3.)\mathrm{c)}\quad 1)&\quad a_\vartheta(a\theta u)b_\vartheta(b\theta u)^3=\frac12(a\Delta u)^4+\frac65[a_\vartheta(a\theta u)^2]_{\widehat{ub^2}}b+\frac{53}{120}[a_\vartheta^2(a\theta u)^2]_{\widehat{ub^2}}\\
&\qquad-\frac65u_x\cdot[a_\vartheta(a\theta u)^2]_{\widehat{ub^2}}b^2-\frac35u_x\cdot[a_\vartheta^2(a\theta u)^2]_{\widehat{ub^2}}b+\frac{19}{120}u_x^2\cdot[a_\vartheta^2(a\theta u)^2]_{\widehat{ub^3}}b^2,\\
2)&\quad a_\vartheta(a\theta u)b_\vartheta(b\theta u)^3=\frac34[a_\vartheta^2(a\theta u)^2]_{\widehat{ub^2}}.
\end{aligned}
\]

Remarque finale : les transvections de $f$ sur $j$ qui sont réductibles d'après \S\ 5 se calculent de manière analogue. On a
\[
\tag{4}
\begin{aligned}
a_j&=-\theta_\nu+u_x\cdot\left\{\frac92\theta_\nu^2-\frac12H\right\}-3u_x^2\cdot\theta_\nu^3+\frac12u_x^3\cdot\varrho,\\
a_j^2&=\frac{21}{10}\theta_\nu^2-\frac3{10}H-2u_x\cdot\theta_\nu^3+\frac12u_x^2\cdot\varrho,\\
a_j^3&=-\frac7{10}\theta_\nu^3+\frac12u_x\cdot\varrho,\\
a_j^4&=\frac35\varrho.
\end{aligned}
\]

De (3.) et (4.) il résulte, par transfert du domaine binaire,
\[
(\theta\theta'u)^2=\frac13\,j\cdot f\pmod{\nu}.\footnotemark
\]
Les autres formes de la série de formes $(\theta\theta'u)^2$ et la forme $\theta'_\vartheta$ sont réductibles aux symboles $\nu$. Nous pouvons donc remplacer :
\[
\modu{\nu}:\quad \text{les contractions avec } f\cdot j
\text{ par les contractions correspondantes avec }(\theta\theta'u)^2.
\]
On a en outre :
\[
(\vartheta\vartheta'x)^2=\frac43 f\cdot\A,
\qquad
\theta_{\vartheta'}(\vartheta\vartheta'x)=-N.
\]
\footnotetext{Gordan--Kerschensteiner, loc. cit. p. 181.}

\clearpage

\begin{center}
{\Large\bfseries Chapitre III.\quad Le système relativement complet modulo $(s,\varrho,t)$.}
\end{center}

\begin{center}
\textbf{\S\ 7.\quad Aperçu de la formation du système.}
\end{center}

Pour passer du système relativement complet modulo $\nu$ au système relativement complet relativement au module immédiatement supérieur, selon la définition de \S\ 3, il faut former le système de $\nu$ par rapport à ce module supérieur et le contracter avec les formes du système relativement complet modulo $\nu$. Mais $\nu=u_\nu^4$, en tant que contravariant du quatrième degré de la forme $f$, est dans une relation dualiste avec elle. Si donc on considère $\nu$ comme forme fondamentale, on obtient un système relativement complet de six formes modulo $\nu(\nu)=(\nu\nu_1x)^4=s_x^4$.

Ainsi, pour former le système relativement complet modulo $s$ (système III), il faut former les transvections de
\[
\begin{array}{ll}
\text{Système I :} & f,\ \theta,\ K,\ j,\ \A,\ N\quad \text{ avec}\vphantom{\dfrac11}\\[0.3em]
\text{Système II :} & \nu,\ H=(\nu\nu_1x)^2u_\nu^2u_{\nu_1}^2,\ L=(\nu\eta x)H_x^2u_\nu^3u_\eta^3,\ g=(\nu\eta x)^4H_x^2,\\
& \sigma=H_\nu^2u_\nu^2u_\eta^4,\ (\nu\sigma x).
\end{array}
\]
On verra (formule (13)) que la forme $g$ est réductible, de sorte que le système II n'est constitué que de cinq formes.

Au cours du calcul, les formes quadratiques sont attachées sous forme de modules
\[
\varrho=u_\varrho^2=\theta_\nu^4u_\vartheta^2,
\qquad
 t=t_x^2=a_\eta^4H_x^2,
\]
de sorte qu'un système relativement complet est obtenu modulo $(s,\varrho,t)$ ; et le module $(\varrho,t)$ doit être considéré comme un module supérieur au module $s$.

L'ordre des formes individuelles doit suivre l'ordre des coefficients. Nous normalisons la contraction
\[
\theta_\eta(K_\eta,\theta_\iota,\ldots)
\]
comme supérieure à la contraction
\[
H_\vartheta(H_\kappa,L_\vartheta,\ldots).
\]
Ainsi, d'après \S\ 1 et \S\ 3b, la succession des formes individuelles du même ordre est déterminée de manière univoque.

La formation des formes d'ordre égal s'effectuera sous forme de tableau :

I. Indication des formes des systèmes I et II qui doivent être réunies pour atteindre l'ordre déterminé dans les coefficients.

II. Disposition des formes irréductibles selon le schéma de la succession des formes.

III. Réduction des séries de formes ou des formes réductibles ou décomposables, c'est-à-dire des points où la succession totale des formes est interrompue ; cela montre que toutes les formations irréductibles sont épuisées par celles indiquées.

IV. Conséquences : 1) indication des réduisants nouvellement apparus ; 2) indication des formes qui n'entrent pas, dans les produits avec des formes du même système, en contraction avec des formes de l'autre système.

On verra que seuls interviennent :

1) les contractions d'une forme du système I avec une forme du système II ;

2) la contraction de puissances de $\theta$, respectivement de $H$, ou de deux formes de l'un des systèmes avec une forme du second système.

Nous obtiendrons le schéma de contraction suivant :
\[
\begin{array}{r@{\ }c@{\ }l@{\qquad}c@{\qquad}r@{\ }c@{\ }l}
\multicolumn{3}{c}{\text{Système I}} & \text{contracté avec} & \multicolumn{3}{c}{\text{Système II}}\\[0.25em]
1. & \text{ordre} & f & & 2. & \text{ordre} & \nu\\
2. & ,, & \theta & & 4. & ,, & H\\
3. & ,, & K,j,\A & & 6. & ,, & L,\sigma\\
4. & ,, & N,\theta^2,f\cdot j & & 8. & ,, & (\nu\sigma x),H^2\\
5. & ,, & \theta\cdot K & & 10. & ,, & H\cdot L\\
6. & ,, & \theta^3,\A\cdot j & & 12. & ,, & H^3.
\end{array}
\]

En plus de la \og{} double réduction \fg{}, les deux formes spéciales servent de contrôles de calcul :

\[
\begin{array}{ll}
\text{I.}&\begin{array}{rlrlrl}
f&=\displaystyle\sum_3x_1^3x_2,& a_\nu^2&=\frac{i}{6}u_x^2,& s&=\frac{i}{3}f,\\
a_\eta^2&=-\frac19 i\theta,& \A_\nu^2&=\frac{i}{6}\theta,& \sigma&=-\frac{i}{6}j,\\
g&=-\frac23i\Delta,& \varrho&=0,& t&=0,
\end{array}\\[1.2em]
\text{II.}&\begin{array}{rlrlrl}
f&=\displaystyle\sum_3x_1^4,& s&=\frac43if,& a_\eta^2&=\frac29i\theta,\\
\Delta_\nu^2&=\frac1{15}i\theta,& \sigma&=\frac43ij,& g&=\frac43i\Delta,\\
\varrho&=0,& t&=0.
\end{array}
\end{array}
\]

Observation finale : les formules suivantes correspondent aux formules (1.), (2.), (3.), (4.) pour la forme fondamentale $\nu$ :
\[
\tag{5.}
\begin{array}{@{}c|c|c@{}}
\cdot & \underline{H_\nu}=\dfrac13u_x\cdot\sigma & \cdot\\[0.45em]
\underline{(\nu\eta x)^2}=A & \underline{H_\nu(\nu\eta x)}=0 & \underline{H_\nu^2}=\sigma\\[0.45em]
\underline{(\nu\eta x)^3}=0 & \underline{H_\nu(\nu\eta x)^2}=B & \underline{H_\nu^2(\nu\eta x)}=0\\[0.45em]
\underline{(\nu\eta x)^4}=g & \underline{H_\nu(\nu\eta x)^3}=0 & \underline{H_\nu^2(\nu\eta x)^2}=C,
\end{array}
\]
\[
\begin{aligned}
A&=\frac16s\cdot\nu-\frac23u_x\cdot s_\nu+\frac23u_x^2\cdot s_\nu^2-\frac{J}{18}u_x^4,\\
B&=-\frac56s_\nu+\frac76u_x\cdot s_\nu^2-\frac{J}{9}u_x^3,\\
C&=\frac23s_\nu^2-\frac{J}{9}u_x^2.
\end{aligned}
\]
\[
\tag{6.}
s_\nu^3u_\nu s_x=\frac13u_x\cdot s_\nu^4=\frac13J\cdot u_x.
\]

\[
\tag{7}
\begin{aligned}
\text{a)}\quad (\nu\sigma x)^2
&=\frac12(Hsu)^2-\frac25\nu\cdot s_\nu^2-
\frac45u_x\cdot s_\eta(Hsu)^2
+u_x^2\left\{\frac16J\cdot\nu-\frac1{40}(ss'u)^4\right\},\\
\text{b)}\quad (\nu\sigma x)^3&=-\frac3{10}s_\eta(Hsu),\\
\text{c)}\quad (\nu\sigma x)^4&=\frac{21}{10}s_\eta^2-\frac3{10}Z-
\frac65u_x\cdot s_\eta^3+\frac1{10}u_x^2\cdot s_\eta^4.
\end{aligned}
\]
\[
\tag{8}
\begin{aligned}
\text{a)}\quad g_\nu&=-s_\eta+u_x\left\{\frac92s_\eta^2-\frac12Z\right\}-3u_x^2\cdot s_\eta^3+\frac12u_x^3\cdot s_\eta^4,\\
\text{b)}\quad g_\nu^2&=\frac{21}{10}s_\eta^2-\frac3{10}Z-2u_x\cdot s_\eta^3+\frac12u_x^2\cdot s_\eta^4,\\
\text{c)}\quad g_\nu^3&=-\frac7{10}s_\eta^3+\frac12u_x\cdot s_\eta^4,\\
\text{d)}\quad g_\nu^4&=\frac35s_\eta^4.
\end{aligned}
\]

\begin{center}
\textbf{\S\ 8.\quad Formes de troisième ordre (système III).}
\end{center}

\noindent\textbf{Contraction de $f$ avec $\nu$.}

\noindent\emph{Formes irréductibles :}
\[
\begin{array}{cccc}
a_\nu&\cdot&\cdot&\cdot\\
\cdot&a_\nu^2&\cdot&\cdot\\
\cdot&\cdot&\cdot&\cdot\\
\cdot&\cdot&\cdot&a_\nu^4=i.
\end{array}
\]

\noindent\emph{Réduction :}
\[
a_\nu^3=\frac13 i u_x\qquad\text{(formule (2)).}
\]
\noindent\emph{Conséquences :} 1) $a_\nu^3$ est réduisant. 2) $f$ n'entre en contraction avec $\nu$ que dans les produits avec des formes contragrédientes $(j)$. Il en va de même pour $\nu$ par rapport à $f$.

\begin{center}
\textbf{\S\ 9.\quad Formes de quatrième ordre (système III).}
\end{center}

\noindent\textbf{Contraction de $\theta$ avec $\nu$.}

\noindent\emph{Formes irréductibles :}
\[
\begin{array}{cccc}
(\vartheta\nu x)&\theta_\nu&\cdot&\cdot\\
(\vartheta\nu x)^2&\theta_\nu(\vartheta\nu x)&\theta_\nu^2&\cdot\\
\cdot&\theta_\nu(\vartheta\nu x)^2&\cdot&\theta_\nu^3.
\end{array}
\]

\noindent\emph{Réductions :} À partir du réduisant $a_\nu^3$, par double réduction des expressions $a_\nu^3(abu)$, $a_\nu^3b_\nu$, $a_\nu^3b_\nu(abu)$ suivant le schéma :

\[
\begin{aligned}
\theta_\nu^2(\vartheta\nu x)
&=a_\nu^2(\widehat{ab}\nu x)(abu)-\frac13(\nu\nu_1x)^3=a_\nu b_\nu(\widehat{ab}\nu x)(abu)+\frac16(\nu\nu_1x)^3\\
&=a_\nu^3(abu)-a_\nu^2b_\nu(abu)=2a_\nu^2b_\nu(abu)=\frac23a_\nu^3(abu),\\[.4em]
\theta_\nu^2(\vartheta\nu x)^2
&=a_\nu^2(\widehat{ab}\nu x)^2-\frac13(\nu\nu_1x)^4=a_\nu b_\nu(\widehat{ab}\nu x)^2+\frac16(\nu\nu_1x)^4\\
&=i\cdot f-2a_\nu^3b_\nu+a_\nu^2b_\nu^2-\frac13s=2a_\nu^3b_\nu-2a_\nu^2b_\nu^2+\frac16s=\frac23i\cdot f-\frac23a_\nu^3b_\nu-\frac16s,\\[.4em]
\theta_\nu^3(\vartheta\nu x)&=a_\nu^2b_\nu(\widehat{ab}\nu x)(abu)=a_\nu^3b_\nu(abu).
\end{aligned}
\]

Les formules sont obtenues :
\[
\tag{9}
\begin{array}{c|c|c}
\underline{\theta_\nu^2(\vartheta\nu x)}=0
& \underline{\theta_\nu^3(\vartheta\nu x)}=0
& \underline{\theta_\nu^4u_\vartheta^2}=u_\varrho^2=\varrho\quad(\text{module adjoint, \S\ 7}),\\[0.6em]
\underline{\theta_\nu^2(\vartheta\nu x)^2}=\dfrac49 i f-\dfrac16s
& &
\end{array}
\]

\noindent\emph{Conséquences :} 1) $\theta_\nu^2(\vartheta\nu x)^2$ est réduisant. 2) $\nu$ n'entre en contraction avec $\theta$ que dans les produits à formes contragrédientes $(g)$.

\begin{center}
\textbf{\S\ 10.\quad Formes de cinquième ordre (système III).}
\end{center}

\[
\text{Contraction de }\left\{\begin{array}{l}
K,j,\A \text{ avec }\nu,\\
f\text{ avec }H.
\end{array}\right.
\]

\noindent\emph{Formes irréductibles :}
\[
\begin{array}{c@{\qquad}c@{\qquad}c@{\qquad}c}
\text{A)} & \begin{array}{cccc}
\cdot&K_\nu&\cdot&\cdot\\
(k\nu x)^2&\cdot&K_\nu^2&\cdot\\
(k\nu x)^3&K_\nu(k\nu x)^2&\cdot&\cdot\\
\cdot&K_\nu(k\nu x)^3&\cdot&\cdot,
\end{array}
&
\text{B)}\ \begin{array}{c}(j\nu x)\\(j\nu x)^2\\(j\nu x)^3\\ \cdot\end{array}
&
\text{C)}\ \begin{array}{c}\A_\nu\\\cdot\\\cdot\\\cdot\end{array}
\end{array}
\]
\[
\text{D)}\qquad
\begin{array}{cccc}
(aHu)&(aHu)^2&\cdot&\\
a_\eta&a_\eta(aHu)&a_\eta(aHu)^2&\\
\cdot&a_\eta^2&\cdot&\cdot
\end{array}
\]

\noindent\emph{Réductions :}

\noindent A) Séquence de formes $K_\nu^3$ : réduisant $a_\nu^3$.

Forme $K_\nu^2(k\nu x)^2$ : réduisant $\theta_\nu^2(\vartheta\nu x)^2$.
\[
(k\nu x),\quad K_\nu(k\nu x),\quad K_\nu^2(k\nu x)
\]
Ce sont des formes décomposables selon \S\ 3.

\noindent B) Forme
\[
\begin{aligned}
(j\nu x)^4&\equiv(\widehat{a\theta}\nu x)^4
=a_\nu^4\cdot\theta+\theta_\nu^4\cdot f-4a_\nu^3\theta_\nu
-4a_\nu\{a_\nu\cdot\theta_x-(\widehat{a\theta}\nu x)\}^3\\
&\quad+6a_\nu^2\{a_\nu\theta_x-(\widehat{a\theta}\nu x)\}^2
=\R\cdot f+\frac53 i\theta-6a_\nu^2(\widehat{a\theta}\nu x)^2
+4a_\nu(\widehat{a\theta}\nu x)^3,
\end{aligned}
\]
et donc
\[
(j\nu x)^4=\R\cdot f+\frac53 i\theta\cdot\modu{(\A,\nu)},\qquad (\text{cf. fin de \S\ 5}).
\]

\noindent C) Forme $\A_\nu^4$ : réduisant $\theta_\nu^4$.

Formes $\A_\nu^2$ et $\A_\nu^3$ : réduisant $a_\nu^3$ ou $\theta_\nu^4$ par substitution de la forme $\A$ par les formes inférieures de la succession de formes (fin de \S\ 5), c'est-à-dire par double réduction des expressions
\[
a_\vartheta a_\nu^3,
\qquad a_\vartheta a_\nu^3\theta_\nu,
\qquad a_\vartheta\theta_\nu^4.
\]
Le double calcul de $\A_\nu^3$ donne lieu à la réduction de la forme supérieure $a_\eta^3$.

\noindent\emph{Formules de réduction :}
\[
\tag{10}
\begin{aligned}
\text{a)}\quad \A_\nu^2
&=\frac35a_\eta^2-\frac3{10}(asu)^2+\frac13 i\theta
+\frac15u_x^2\cdot(2a_\R^2-t),\\
\text{b)}\quad \A_\nu^3
&=-\frac3{10}a_\R+\frac35u_x\left\{\frac{13}{10}a_\R^2-\frac25t\right\},\\
\text{c)}\quad \A_\nu^4&=a_\R^2-\frac25t.
\end{aligned}
\]

\noindent\emph{Dérivation des formules :}
\[
\begin{aligned}
\text{a)}\quad
a_\vartheta a_\nu^3
&=\frac13 i\theta
=[a_\vartheta]_{\nu^3}+\frac67[a_\vartheta^2]_{\nu^2}
+\frac65[a_\vartheta(a\theta u)^2]_{\nu^2}\,
\widehat{\nu}x^2
+\frac{11}{30}[a_\vartheta^2(a\theta u)^2]\,\widehat{\nu x}^2\\
&\quad-\frac67u_x\cdot[a_\vartheta^2]_\nu^3
-\frac16u_x\cdot[a_\vartheta^2(a\theta u)^2]_{\nu}\,\widehat{\nu x}^2,\\
\frac13 i\theta
&=\A_\nu^2-\frac23u_x\cdot\A_\nu^3-a_\nu a_{\nu_1}(\nu\nu_1x)^2
+\frac25a_\nu^2(\nu\nu_1x)^2
+\frac15u_x\cdot a_\nu^2a_{\nu_1}(\nu\nu_1x)^2;
\end{aligned}
\]
\[
\begin{aligned}
\text{b)}\quad
 a_\vartheta a_\nu^3\theta_\nu
&=0=[a_\vartheta]_{\nu^4}+\frac9{14}[a_\vartheta^2]_{\nu^3}
+\frac35[a_\vartheta(a\theta u)^2]_{\nu^2}\,\widehat{\nu x}^2
+\frac{17}{120}[a_\vartheta^2(a\theta u)^2]_{\nu}\,\widehat{\nu x}^2\\
&\quad-\frac9{14}u_x\cdot[a_\vartheta^2]_{\nu^4}
-\frac1{24}u_x\cdot[a_\vartheta^2(a\theta u)^2]_{\nu^2}\,\widehat{\nu x}^2,\\
0&=\frac56\A_\nu^3-\frac12u_x\cdot\A_\nu^4-\frac14a_\eta^3+\frac1{20}u_x\cdot a_\eta^4;
\end{aligned}
\]
\[
\begin{aligned}
\text{b')}\quad
a_\vartheta\theta_\nu^4
&=a_\R=a_\vartheta\theta_\nu\{a_\nu-(a\theta\widehat{\nu x})\}^3\\
&=-\frac32[a_\vartheta^2]_{\nu^3}+\frac35[a_\vartheta(a\theta u)^2]_{\nu^2}\,\widehat{\nu}x^2
-\frac{37}{120}[a_\vartheta^2(a\theta u)^2]_{\nu}\,\widehat{\nu x}^2\\
&\quad+\frac32u_x\cdot[a_\vartheta^2]_{\nu^4}
+\frac{49}{120}u_x\cdot[a_\vartheta^2(a\theta u)^2]_{\nu^2}\,\widehat{\nu x}^2,\\
a_\R&=-\frac32\A_\nu^3+\frac32u_x\cdot\A_\nu^4-\frac{11}{20}a_\eta^3+\frac7{20}u_x\cdot a_\eta^4;
\end{aligned}
\]
\[
\begin{aligned}
\text{c)}\quad
 a_\vartheta^2\theta_\nu^4
&=a_\R^2=[a_\vartheta^2]_{\nu^4}+\frac35[a_\vartheta^2(a\theta u)^2]_{\nu^2}\,\widehat{\nu x}^2,\\
a_\R^2&=\A_\nu^4+\frac25a_\eta^4.
\end{aligned}
\]

\noindent D) Réduction de $a_\eta^3$ par double réduction de $\A_\nu^3$ (C).

Les autres réductions sont obtenues par le calcul dual correspondant à celui de \S\ 9.

\noindent\emph{Formules de réduction :}
\[
\tag{11}
\begin{array}{rl|rl}
\underline{a_\eta^2(aHu)}&=-\dfrac16(asu)^3,&
\underline{a_\eta^2(aHu)^2}&=\dfrac49 i\nu-\dfrac16(asu)^4,\\[0.7em]
\underline{a_\eta^3}&=-a_\R+\dfrac35u_x\cdot a_\R^2+\dfrac15u_x\cdot t,
& \underline{a_\eta^3(aHu)}&=0,\\[0.7em]
&&\underline{a_\eta^4}&=t\quad(\text{adjoint comme module}).
\end{array}
\]

\noindent\emph{Conséquences :}

1) $(j\nu x)^4$, $\A_\nu^2$, $a_\eta^2(aHu)$ sont des réduisants.

2) $\nu$ n'entre que dans les produits aux formes contragrédientes en contraction avec $K,j,\A$ ; il en va de même pour $f$ par rapport à $H$ et pour $j$ par rapport à $\nu$, puisque $\theta_\nu(\vartheta\cdot j,\nu,x)^3$ devient réductible selon \S\ 11 B.

\begin{center}
\textbf{\S\ 11.\quad Formes de sixième ordre (système III).}
\end{center}
\[
\text{Contraction de }\left\{\begin{array}{l}
\theta^2,
f\cdot j,
N \text{ avec }\nu,\\
\theta\text{ avec }H.
\end{array}\right.
\]

\noindent\emph{Formes irréductibles :}
\[
\renewcommand{\arraystretch}{1.25}
\begin{array}{@{}r@{\quad}c@{\qquad\qquad}r@{\quad}c@{\qquad\qquad}r@{\quad}c@{}}
\text{A)}&(\vartheta^2\nu x)^3&\text{B)}&a_\nu(j\nu x)&\text{C)}&N_\nu,\\[0.45em]
&\cdot&&\cdot&&\\[-0.1em]
&\theta_\nu(\vartheta^2\nu x)^3,&&a_\nu^2(j\nu x),&&
\end{array}
\]
\[
\renewcommand{\arraystretch}{1.25}
\begin{array}{@{}r@{\quad}c|c|c|c@{}}
&\cdot&(\theta Hu)&(\theta Hu)^2&\cdot\\
\text{D)}&(\vartheta\eta x)&\theta_\eta&H_\vartheta(\theta Hu),\ \theta_\eta(\theta Hu)&\theta_\eta(\theta Hu)^2\\
&(\vartheta\eta x)^2&H_\vartheta(\vartheta\eta x),\ \theta_\eta(\vartheta\eta x)&\theta_\eta^2&\theta_\eta^2(\theta Hu)\\
&\cdot&\theta_\eta(\vartheta\eta x)^2&\cdot&\cdot
\end{array}
\]

\noindent\emph{Réductions :}

\noindent A) $(\vartheta^2\nu x)^4$ est décomposé selon la formule (3)a :
\[
(a\A\widehat{\vartheta}x)^2=\frac12(\vartheta^2\nu x)^4-\frac35 f\cdot a_\nu^2(\nu\vartheta x)^2
=\frac13\A^2+f\cdot\A_\vartheta^2.
\]

\noindent B) $a_\nu(j\nu x)^2$ se décompose selon la formule (9)a, en remplaçant, selon la fin de \S\ 6, $f\cdot j$ par $(\theta\theta'u)^2$ et en tenant compte de la réductibilité de $\theta_\vartheta'$ :
\[
\frac13a_\nu(j\nu x)^2=(\widehat{\theta\theta'}\nu x)^2\theta_\nu
=\theta_\nu^3\cdot\theta-\theta_\nu^2\theta_\nu'
=\theta_\nu^3\cdot\theta+\theta_\nu^2(\widehat{\vartheta\nu}\theta'u)
=\theta_\nu^3\cdot\theta-\frac23\theta_\nu^3\cdot\theta.
\]
Formes $a_\nu(j\nu x)^3$ et $a_\nu^2(j\nu x)^2$ : réduisants $a_j$ et $(j\nu x)^4$ :
\[
a_\nu(j\nu x)^3=a_j(j\nu x)^3-(j\nu x)^3(j\nu\widehat{au}),
\]
\[
a_\nu^2(j\nu x)^2=a_j^2(j\nu x)^2-2a_j(j\nu x)^2(j\nu\widehat{au})+(j\nu x)^2(j\nu\widehat{au})^2.
\]

\noindent C) Séquence de formes $N_\nu^2$ ; réduisant $\A_\nu^2$. $(n\nu x)$ et $N_\nu(n\nu x)$ se décomposent en transvections sur des déterminants fonctionnels selon \S\ 3.

\noindent D) \emph{Aperçu des réductions :}

a) Séquence de formes $\theta_\eta^2H_\vartheta$ : réduisant $a_\eta^2(aHu)$. (Conduit à des formes avec des symboles supérieurs.)

b) Suite de formes $\theta_\eta H_\vartheta$ : réduisant $a_\eta^2(aHu)$ par double réduction. (Cela conduit à des formes avec des symboles supérieurs et à des formes supérieures de la séquence totale des formes, c'est-à-dire à la séquence de formes $\theta_\eta^2$.) Au lieu de $\theta_\eta H_\vartheta$, nous considérons la séquence de formes
\[
\theta_\eta(\vartheta\eta x)(\theta Hu)=\theta_\eta H_\vartheta+\theta_\eta^2,
\]
et on y remplace les symboles $\theta$ par la forme inférieure $(abu)$, arrivant ainsi à la double réduction de la succession des formes
\[
a_\eta^2(aHu)(abu)=\theta_\eta H_\vartheta+\frac32\theta_\eta^2\modu{s}
\]
jusqu'à la forme définitive
\[
a_\eta^3b_\eta(bHu)(abu)=\theta_\eta^3H_\vartheta+\frac12\theta_\eta^4.
\]

c) Séquence de formes $\theta_\eta^3$ : par double réduction des formes qui appartiennent simultanément aux séquences $\theta_\eta H_\vartheta$ et $\theta_\eta^2H_\vartheta$, c'est-à-dire des formes de la séquence $\theta_\eta^2H_\vartheta$, vers les formes avec des symboles supérieurs d'une part, et vers les formes avec des symboles supérieurs et la succession supérieure de formes $\theta_\eta^3$ d'autre part.

d) Formes $\theta_\eta^2(\vartheta\eta x)$, $\theta_\eta^2(\vartheta\eta x)^2$, $\theta_\eta^3(\vartheta\eta x)$ : par double réduction à l'aide du réduisant $a$, de manière analogue au calcul de \S\ 9.

e) Forme $\theta_\eta^2(\theta Hu)^2$ et $\theta_\eta^2(\vartheta\eta x)^2$ : par double réduction des expressions $\A_\nu^2(a\A u)^4$ et $(a\A\widehat{\nu}x)^4$ en utilisant les réduisants $\A_\nu^2$ et $(a\A u)^4$.

f) Formes $H_\vartheta$ et $H_\vartheta^2$ : formes décomposables. Pour $H_\vartheta^2$ cela résulte de la double réduction de l'expression $\A_\nu^2(a\A u)^2$, ou encore du calcul direct des produits $(a_\nu^2)^2$, $a_\nu^2\cdot b_{\nu_1}$ par formes du système. (Le calcul analogue de $(a_\nu)^2$ donne une relation entre les formes décomposables.)

g) Réduction des formes supérieures car les formes de la séquence $\theta\cdot H$ admettent deux possibilités de réduction (elles appartiennent à deux séquences de formes réductibles), à savoir :
\begin{quote}\small
$g$ : par double réduction de $\theta_\eta^2(\vartheta\eta x)^2$ ; admet les réductions d) et e).\par
$s_\vartheta^2(\theta su)$ : par double réduction de $\theta_\eta^3(\vartheta\eta x)$ ; admet les réductions c) et d).\par
$s_\vartheta^2(\theta su)^2$ : par double réduction de $\theta_\eta^2H_\vartheta^2$ ; admet la réduction a) et possède comme facteur le réduisant $\theta_\nu^2(\vartheta\nu x)^2$.\par
$s_\nu^2$ : par double réduction de $\theta_\eta^3H_\vartheta$ ; admet les réductions a) et c).
\end{quote}
La démonstration de l'indépendance des formes restantes résulte d'une double réduction de la séquence de formes $\A_\nu^2(a\A u)^2\equiv\theta_\eta^2$.

\noindent\emph{Exécution des réductions :}

L'existence des réductions a), b), c), d) est claire a priori, puisque selon la loi de formation indiquée, les relations qui naissent dans la double réduction sont indépendantes les unes des autres. Pour les réductions e), f), g), il faut vérifier par calcul qu'aucune identité ne se produit.

Nous donnons aussi, symétriquement par rapport au terme diagonal dans la suite des formes $\theta\cdot H$ et en avançant vers la forme $\theta_\eta$, une partie des formules obtenues par les réductions a), b), c), d), avec indication du calcul lorsqu'elle est utile ; elles seront utilisées tantôt dans les réductions e), f), g), tantôt plus loin.

\noindent 1) $H_\vartheta$ et $H_\vartheta^2$ selon f). Avec l'approche\footnote{Le signe $\sim$ au lieu du signe d'égalité signifie que les produits de $u_x$ de formes supérieures à celles issues des contractions III et IV sont omis.}
\[
\begin{aligned}
a_\nu\cdot b_{\nu_1}^2
&=\nu\cdot a_\nu b_\nu^2-(aHu)(bHu)b_\eta
\sim \nu\cdot a_\nu b_\nu^2-f\cdot a_\eta(aHu)^2-(abu)(aHu)b_\eta+(abu)^2b_\eta,\\
a_\nu^2\cdot b_{\nu_1}^2
&=b_{\nu_1}^2\{a_{\nu_1}u_\nu-(a\widehat{\nu\nu_1}u)\}^2
\sim \nu\cdot a_\nu^2b_\nu^2-2a_\eta b_\eta(aHu)(bHu)-\frac16(asu)^2(bsu)^2\\
&\sim \nu\cdot a_\nu^2b_\nu^2-2a_\eta^2(aHu)(abu)-\frac16(asu)^2(bsu)^2+(\vartheta\eta x)^2(\theta Hu)^2
\end{aligned}
\]
En remplaçant la valeur de $\theta_\eta H_\vartheta$, les formules apparaissent :
\[
\tag{12}
\begin{aligned}
\text{a)}\quad H_\vartheta&\sim 2a_\nu\cdot a_\nu^2+2\nu\cdot\theta_\nu(\vartheta\nu x)^2+2f\cdot a_\eta(aHu)^2-2\theta_\eta,\\
\text{b)}\quad H_\vartheta^2&\sim (a_\nu^2)^2+2\theta_\eta^2+\frac23(\theta su)^2+\frac1{12}s\cdot\nu-\frac19 i f\cdot\nu-\frac16f\cdot(asu)^4.
\end{aligned}
\]

\noindent 2) $\theta_\eta H_\vartheta$ selon b) :
\[
\begin{aligned}
a_\eta^2(aHu)(abu)&\sim [a_\eta^2(aHu)]_{\widehat{bu}}+\frac12 f\cdot[a_\eta^2(aHu)^2]
=a_\eta b_\eta(aHu)(abu)\\
&\quad+a_\eta(\widehat{ab}\eta x)(abu)(aHu)\sim\theta_\eta(\vartheta\eta x)(\theta Hu)+\frac12\theta_\eta^2+\frac14(\nu\eta x)^2.
\end{aligned}
\]
Des formules (5) et (11) il résulte :
\[
\theta_\eta H_\vartheta\sim -\frac32\theta_\eta^2-\frac14(\theta su)^2-\frac1{12}s\cdot\nu+\frac29i f\cdot\nu.
\]

\noindent 3) $\theta_\eta H_\vartheta(\theta Hu)$ est calculé selon b) de manière analogue à $\theta_\eta H_\vartheta$ :
\[
\theta_\eta H_\vartheta(\theta Hu)\sim-\theta_\eta^2(\theta Hu)-\frac16(\theta su)^3.
\]

\noindent 4) $\theta_\eta H_\vartheta(\vartheta\eta x)$ selon b); $\theta_\eta^2(\vartheta\eta x)$ selon d) (cf. \S\ 9) :
\[
\theta_\eta H_\vartheta(\vartheta\eta x)\sim-\theta_\eta^2(\vartheta\eta x)-\frac16s_\vartheta(\theta su)
\sim\frac13(\vartheta\R x)-\frac16s_\vartheta(\theta su),
\]
\[
\theta_\eta^2(\vartheta\eta x)\sim\frac23a_\eta^3(abu)\sim-\frac13(\vartheta\R x).
\]

\noindent 5)
\[
\begin{array}{ccc}
\theta_\eta H_\vartheta^2\text{ selon b)},&
\theta_\eta^2H_\vartheta\text{ selon a)},&
\theta_\eta^2H_\vartheta\text{ selon b), et par là }\theta_\eta^3\text{ selon c)},\\
\text{de }a_\eta^2(aHb)(bHu);&
\text{de }a_\eta^2(Hab)(abu);&
\text{de }a_\eta^3(bHu)(abu).
\end{array}
\]
\[
\theta_\eta H_\vartheta^2\sim-\frac32\theta_\eta^2H_\vartheta-\frac14s_\vartheta(\theta su)^2+\frac16s_\nu-\frac49ia_\nu
\sim-\theta_\R-\frac16s_\nu-\frac19ia_\nu,
\]
\[
\theta_\eta^2H_\vartheta\sim\frac23\theta_\R-\frac16s_\vartheta(\theta su)^2+\frac29s_\nu-\frac29ia_\nu,
\]
\[
\theta_\eta^2H_\vartheta\sim\frac23\theta_\R-\frac23\theta_\eta^3+\frac5{36}s_\nu,
\qquad
\theta_\eta^3\sim\frac14s_\vartheta(\theta su)^2-\frac18s_\nu+\frac13i\cdot a_\nu.
\]

\noindent 6) $\theta_\eta^2(\theta Hu)^2$ selon e) :
\[
\A_\nu^2(a\A u)^4=[(a\A u)^4]_{\nu^2}
=-\frac{12}{5}\theta_\eta^2(\theta Hu)^2-\frac3{10}\sigma-\frac1{20}(\theta su)+\nu\cdot\R
\quad(\text{selon la formule (3)c}),
\]
\[
\A_\nu^2(a\A u)^4=[\A_\nu^2]^4_{\widehat{u a}}=\frac35\theta_\eta^2(\theta Hu)^2+\frac15\sigma-\frac3{10}(\theta su)^4+\frac13ij
\quad(\text{selon la formule (10)a}),
\]
\[
\theta_\eta^2(\theta Hu)^2=-\frac16\sigma+\frac1{12}(\theta su)^4-\frac19ij+\frac13\nu\cdot\R.
\]

\noindent 7) $\theta_\eta^2(\vartheta\eta x)^2$ selon d) et e). De là résulte la réduction de la forme supérieure $g$.

Par un calcul analogue à celui de \S\ 9 il résulte :
\[
\begin{aligned}
\theta_\eta^2(\vartheta\eta x)^2
&=\frac12t\cdot f-\frac12a_\eta^2b_\eta-\frac1{12}g
=\frac23t\cdot f-\frac23a_\eta^3b_\eta-\frac16g\\
&=-\frac16g-\frac13(\vartheta\R x)^2+\frac4{15}f\cdot a_\R^2+\frac8{15}f\cdot t
\quad(\text{selon la formule (11)}).
\end{aligned}
\]
Un calcul analogue à celui qui précède donne :
\[
(a\A\widehat{\nu}x)^4=[(a\A u)^4]_{\widehat{\nu}x^4}
=\frac{21}{10}\theta_\eta^2(\vartheta\eta x)^2+\frac7{10}s_\vartheta^2-\frac3{10}g
\quad(\text{selon la formule (3)c}),
\]
\[
\begin{aligned}
(a\A\widehat{\nu x})^4&=-\frac13i\A+6[\A_\nu^2]a^2-4[\A_\nu^3]a+\A_\nu^4\cdot f\\
&= -\frac{36}{5}\theta_\eta^2(\vartheta\eta x)^2-\frac95s_\vartheta^2-\frac35g-\frac35(\vartheta\R x)^2\\
&\quad+\frac53i\A+\frac{37}{25}f\cdot a_\varrho^2+\frac{74}{25}f\cdot t,\quad(\text{selon la formule (10)}).
\end{aligned}
\]
Donc
\[
\begin{aligned}
\frac{93}{10}\theta_\eta^2(\vartheta\eta x)^2
&=-\frac3{10}g-\frac52s_\vartheta^2-\frac35(\vartheta\R x)^2\\
&\quad+\frac53i\A+\frac{37}{25}f\cdot a_\R^2+\frac{74}{25}f\cdot t,
\end{aligned}
\]
et en comparant les deux valeurs de $\theta_\eta^2(\vartheta\eta x)^2$ selon g) :
\[
\tag{13.)*}
0=g-2s_\vartheta^2+2(\vartheta\R x)^2+\frac43i\A-\frac45f\cdot a_\R^2-\frac85f\cdot t.
\]
\srcfntext{*)}{Il s'agit de la relation mentionnée dans la note de l'introduction entre les trois covariants du sixième ordre et du sixième degré.}

\noindent 8) $\theta_\eta^2H_\vartheta(\theta Hu)$ selon a) et b); donc, $\theta_\eta^3(\theta Hu)$ selon c) (les formes avec le facteur $u_x$ disparaissent) :
\[
\theta_\eta^2H_\vartheta(\theta Hu)=-\frac16s_\vartheta(\theta su)^3-\frac13(\nu\R x),
\]
\[
\theta_\eta^2H_\vartheta(\theta Hu)=-\frac13\theta_\eta^3(\theta Hu)-\frac13(\nu\R x),
\]
\[
\theta_\eta^3(\theta Hu)=\frac12s_\vartheta(\theta su)^3.
\]

\noindent 9) $\theta_\eta^2H_\vartheta(\vartheta\eta x)$ selon a) et b); donc, $\theta_\eta^3(\vartheta\eta x)$ selon c). $\theta_\eta^3(\vartheta\eta x)$ selon d); d'où la réduction de la forme supérieure $s_\vartheta^2(\theta su)$ selon g) (les formes avec le facteur $u_x$ disparaissent) :
\[
\theta_\eta^2H_\vartheta(\vartheta\eta x)=\frac13(atu),
\]
\[
\theta_\eta^2H_\vartheta(\vartheta\eta x)=\frac13(atu)+\frac13\theta_\eta^3(\vartheta\eta x)-\frac16s_\vartheta^2(\theta su),
\]
\[
\theta_\eta^3(\vartheta\eta x)=\frac12s_\vartheta^2(\theta su),
\]
\[
\theta_\eta^3(\vartheta\eta x)=\frac25\theta_\R(\vartheta\R x)+\frac15(atu),
\]
\[
\tag{14}
s_\vartheta^2(\theta su)=\frac45\theta_\R(\vartheta\R x)+\frac25(atu).
\]

\noindent 10) $\theta_\eta^2H_\vartheta^2$ selon a) et via le réduisant $\theta_\nu^2(\vartheta\nu x)^2$; $\theta_\eta^3H_\vartheta$ selon a) et c); $\theta_\eta^4$ selon c); d'où la réduction des formes supérieures $s_\vartheta^2(\theta su)^2$ et $s_\nu^2$ selon g) :
\[
\begin{aligned}
\text{D'après a)}\quad
\theta_n^2H_\vartheta^2&=[a_\eta^2(aHu)^2]b^2+\frac13[a_\eta^4]_{\widehat{bu^2}}-\frac13H_\nu^2(\nu\eta x)^2\\
&=\frac49i\cdot a_\nu^2-\frac16s_\vartheta^2(\theta su)^2-\frac5{18}s_\nu^2+\frac13(atu)^2+\frac1{54}J\cdot u_x^2.
\end{aligned}
\]
\[
\begin{aligned}
\text{Par }\theta_\nu^2(\vartheta\nu x)^2 :
\quad\theta_\eta^2H_\vartheta^2
&=[\theta_\nu^2(\vartheta\nu x)^2]_{\nu_1^2}+\frac13[\theta_\nu^4]\widehat{\nu_1}x^2-\frac13s_\vartheta^2(\theta su)^2\\
&=\frac49i\cdot a_\nu^2-\frac16s_\nu^2-\frac13s_\vartheta^2(\theta su)^2+\frac13(\nu\R x)^2.
\end{aligned}
\]
\[
\begin{aligned}
\text{D'après a)}\quad
\theta_\eta^3H_\vartheta&=-[a_\eta^2(aHu)]_{\widehat{bu}}b^2-\frac12[a_\eta^2(aHu)^2]b^2-\frac13[a_\eta^3]_{\widehat{bu^2}}b+\frac1{18}[a_\eta^4]_{\widehat{bu^2}}\\
&\quad+\frac12u_x\cdot[a_\eta^2(aHu)^2]b^3\\
&=-\frac29ia_\nu^2+\frac1{12}s_\nu^2+\frac4{15}\theta_\R^2-\frac7{90}(\nu\R x)^2+\frac1{30}(atu)^2+\frac2{27}i^2\cdot u_x^2-\frac1{36}J\cdot u_x^2.
\end{aligned}
\]
\[
\begin{aligned}
\text{D'après c)}\quad
\theta_\eta^3H_\vartheta :
a_\eta^3(Hab)(bHu)&=\frac12(atu)^2
=\theta_\eta^2H_\vartheta(\theta Hu)(\vartheta\eta x)+\frac14H_\nu^2(\nu\eta x)^2\\
&\quad+\frac12\theta_\eta^2H_\vartheta^2
=\frac32\theta_\eta^2H_\vartheta^2+\theta_\eta^3H_\vartheta-u_x\cdot[a_\eta^2(aHu)^2]b^3\\
&\quad+\frac14H_\nu^2(\nu\eta x)^2,
\end{aligned}
\]
et donc
\[
\begin{aligned}
\theta_\eta^3H_\vartheta&=-\frac23i\cdot a_\nu^2+\frac5{12}s_\nu^2+\frac12(\nu\R x)^2-\frac12(atu)^2\\
&\quad+\frac4{27}i^2\cdot u_x^2-\frac1{12}J\cdot u_x^2.
\end{aligned}
\]
D'après c),
\[
\begin{aligned}
\theta_\eta^4&=\frac49ia_\nu^2-\frac16s_\nu^2+\frac{16}{15}\theta_\R^2\\
&\quad-\frac{14}{45}(\nu\R x)^2+\frac2{15}(atu)^2.
\end{aligned}
\]
D'après g), à partir des deux valeurs pour $\theta_\eta^2H_\vartheta^2$ :
\[
\tag{15}
\begin{aligned}
s_\vartheta^2(\theta su)^2
&=\frac23s_\nu^2-2(atu)^2+2(\nu\R x)^2-\frac19J\cdot u_x^2\\
&=\frac89i\cdot a_\nu^2+\frac8{15}\theta_\R^2-\frac{14}{15}(atu)^2+\frac{38}{45}(\nu\R x)^2-\frac4{27}i^2\cdot u_x^2.
\end{aligned}
\]
D'après g), à partir des deux valeurs pour $\theta_\eta^3H_\vartheta$ :
\[
\tag{16}
s_\nu^2=\frac43i\cdot a_\nu^2+\frac45\theta_\R^2+\frac85(atu)^2-\frac{26}{15}(\nu\R x)^2+\frac16J\cdot u_x^2-\frac29i^2\cdot u_x^2.
\]
La formule (16) fournit la base de la réduction du module $(ss'u)^4$, et la formule (13) pour la réduction du système de $s$ (\S\ 17).

\noindent\emph{Conséquences :}

1) $a_\nu(j\nu x)^3$, $\theta_\eta H_\vartheta$, $\theta_\eta^2(\vartheta\eta x)$, $\theta_\eta^2(\theta Hu)^2$ sont des réduisants ; $a_\nu(j\nu x)^2$, $H_\vartheta$ et $H_\vartheta^2$ sont des formes décomposables.

2) La forme du système II, $j(\nu)=g=(\nu\eta x)^4H_x^2$, est réductible.

3) Comme la seule forme contragrédiente $g$ devient réductible, $\nu$ n'entre pas du tout, ni en puissances ni en produits avec des formes du même système, dans la contraction avec le système I.

$\theta$ entre en contraction uniquement considéré comme contravariant en produits ou puissances, et par conséquent $H$ uniquement considéré comme covariant (cf. \S\ 13 B).

\clearpage

\begin{center}
\textbf{\S\ 12.\quad Formes de septième ordre.}
\end{center}
\[
\text{Contraction de }\left\{\begin{array}{l}
\theta\cdot K\text{ avec }\nu,\\
K,j,\A\text{ avec }H,\\
f\text{ avec }L,\sigma.
\end{array}\right.
\]

\noindent\emph{Formes irréductibles :}
\[
\renewcommand{\arraystretch}{1.35}
\begin{array}{@{}r@{\quad}c|c|c|c@{}}
&\cdot&\cdot&(KHu)^2&\cdot\\
&\cdot&K_\eta&\cdot&K_\eta(KHu)^2\\
\text{B)}&(k\eta x)^2&\cdot&K_\eta^2&\cdot\\
&(k\eta x)^3&H_k(k\eta x)^2,\ K_\eta(k\eta x)^2&\cdot&\cdot\\
&\cdot&K_\eta(k\eta x)^3&\cdot&\cdot
\end{array}\, ,
\]
\[
\renewcommand{\arraystretch}{1.35}
\begin{array}{@{}r@{\quad}c@{\qquad}c@{\qquad}c@{\qquad\qquad}r@{\quad}c@{}}
\text{C)}&(j\eta x)&H_j&\cdot&\text{D)}&(\A Hu)\\
&\cdot&H_j(j\eta x)&H_j^2,&&\A_\eta,
\end{array}
\]
\[
\renewcommand{\arraystretch}{1.35}
\begin{array}{@{}r@{\quad}c@{\qquad}c@{\qquad}c@{\qquad}c@{}}
&\cdot&(aLu)^2&(aLu)^3&\cdot\\
\text{E)}&a_l&\cdot&a_l(aLu)^2&a_l(aLu)^3\\
&\cdot&a_l^2&\cdot&\cdot
\end{array}\, .
\]

\noindent\emph{Réductions :}

\noindent A) $(\vartheta\cdot k,\nu,x)^4$ se décompose comme une transvection simple sur la forme décomposable $(\vartheta^2\nu x)^4$.

\noindent B) Série de formes $H_kK_\eta$ : réduisant $H_\vartheta\theta_\eta$.

Série de formes $K_\eta^2(KHu)$ : réduisant $a_\eta^2(aHu)$.

Série de formes $K_\eta^2(k\eta x)$ : réduisant $\theta_\eta^2(\vartheta\eta x)$.

La double réduction qui en résulte pour la série de formes $K_\eta^3$ donne lieu à la réduction de la série supérieure $(j\eta x)^3$.

$(KHu)$, $(k\eta x)$, $K_\eta(k\eta x)$ sont des formes décomposables selon \S\ 3.

$H_k(KHu)$, $K_\eta(KHu)$ se décomposent comme provenant d'une contraction simple avec la forme décomposable $K_\nu(k\nu x)$ et le déterminant fonctionnel
\[
K_\nu^2\equiv\theta_\nu^2(a\theta u)\equiv a_\nu^2(a\theta u)\modu{\nu^2}.
\]
La double représentation de $K_\nu^2$ correspond à une relation entre formes décomposables. On obtient :
\[
\begin{aligned}
H_k(KHu)
&=2K_{\nu_1}(\nu\nu_1k)=2\theta_{\nu_1}(\nu\nu_1\widehat{a\theta})\\
&=2\theta_\nu^2\cdot a_\nu-a_\nu^2\cdot\theta_\nu+f\cdot\theta_\eta(\theta Hu)^2-f\cdot \nu\cdot\theta_\nu^3\\
&\qquad\quad(\text{théorème du produit})\modu{(\nu^3)},
\end{aligned}
\]
\[
\begin{aligned}
K_\eta(KHu)&=-K_\nu^2(\nu\nu_1x)=-a_\nu^2(a\theta u)(\nu\nu_1x)=-\theta_\nu^2(a\theta u)(\nu\nu_1x)\\
&=a_\eta(aHu)^2\cdot\theta+a_\nu^2\cdot\theta_\nu
=-\theta_\eta(\theta Hu)^2\cdot f-\theta_\nu^2\cdot a_\nu:
\end{aligned}
\]
Par conséquent
\par\smallskip\noindent
\begin{minipage}[t]{3.4em}\vspace*{0pt}(17.)\srcfnmark{*)}\end{minipage}%
\begin{minipage}[t]{\dimexpr\linewidth-3.4em\relax}\vspace*{0pt}\centering
\(\displaystyle
0=a_\nu^2\cdot\theta_{\nu_1}+\theta_\nu^2\cdot a_{\nu_1}
+f\cdot\theta_\eta(\theta Hu)^2+\theta\cdot a_\eta(aHu)^2\modu{(\nu^3)}.
\)
\end{minipage}\par\smallskip
\srcfntext{*)}{Dans la formule (17), les produits de formes parmi lesquelles figure une contravariante, comme dans ce cas
$f\cdot\theta_\nu^3\cdot \nu$, sont omis dans toutes les relations entre formes décomposables, car ils
n'ont aucune importance lors de l'application des formules au système $s_{IV}$.}
Les formes $H_k$, $H_k(k\eta x)$, $H_k^2$, $H_k^2(k\eta x)$ se décomposent comme provenant d'une contraction simple avec les formes décomposables $H_\vartheta$ et $H_\vartheta^2$ (cf. \S\ 3. 2), a) et b)).

On a
\[
H_k=H_\vartheta(a\theta u)\sim[H_\vartheta]_{\widehat{ua}}-f\cdot H_\vartheta(\theta Hu)
\qquad(\text{spécialisation }y=\widehat{uv}\text{ de 2) a}).
\]
\[
H_k(k\eta x)=H_\vartheta a_\eta-H_\vartheta\theta_\eta\cdot f,
\qquad
H_\vartheta a_\eta=\frac54[H_\vartheta]a-\frac14H_\vartheta a_\vartheta
\quad(\S\ 3. 2)b),
\]
\[
H_\vartheta^2(a\theta u)=[H_\vartheta^2]_{\widehat{ua}},
\qquad
H_\vartheta^2a_\eta=[H_\vartheta^2]a=H_k^2(k\eta x)+H_\vartheta^2\theta_\eta\cdot f.
\]

\noindent C) Série de formes $(j\eta x)^3$ : réductible par double réduction de la série de formes $K_\eta^3$ selon B), suivant le schéma :
\[
\begin{aligned}
0&=a_\eta^3(a\theta u)=\theta_\eta^3(a\theta u)
=\{\theta_\eta+(\widehat{a\theta}\eta x)\}^3(a\theta u)-\theta_\eta^3(a\theta u)\\
&=\frac12(j\eta x)^3\modu{(\nu^3,s,\R,t,i,J)}
\qquad(\text{cf. fin du \S\ 5}).
\end{aligned}
\]
Le schéma analogue vaut jusqu'à la forme finale de la série :
\[
0=H_\vartheta^2(\widehat{a\theta}\eta x)a_\eta^3
=H_\vartheta^2(\widehat{a\theta}\eta x)\theta_\eta^3=\frac12H_j^2(j\eta x)^4\modu{(\nu^3,s,\R,t,i,J)}.
\]
Forme $(j\eta x)^2$ : elle se décompose 1) par double polarisation de la relation pour la forme décomposable $H_\vartheta^2$ ((12.) b) ;

2) par calcul direct du produit $a_\nu\cdot\theta_\nu^3$ ; d'où la décomposition de la forme supérieure $(\A Hu)^2$. On obtient :
\[
\tag{18}
\displaystyle
\left.\begin{array}{ll}
\text{a)}&\dfrac13(\A Hu)^2+\dfrac13(j\eta x)^2=\theta_\nu^2\cdot a_{\nu_1}^2,\\[0.6em]
\text{b)}&(j\eta x)^2=\dfrac43\theta_\nu^3\cdot a_{\nu_1},
\end{array}\right\}\modu{(\nu^3,s,\R,t,i,J)}
\]
suivant le schéma :
\[
\begin{aligned}
H_\vartheta^2(a\theta u)^2=\frac13(\A Hu)^2
&=2\theta_\eta^2(a\theta u)(aHu)+\theta_\nu^2\cdot a_{\nu_1}^2
\quad(\text{théorème du produit})\\
&=(a\theta u)(aHu)\{a_\eta-(\widehat{a\theta}\eta x)\}^2+\theta_\nu^2\cdot a_{\nu_1}^2,
\end{aligned}
\]
\[
\begin{aligned}
a_{\nu_1}\cdot\theta_\nu^3
&=-(a\widehat{\nu\nu_1}u)\theta_\nu^2(\theta\widehat{\nu\nu_1}u)-\frac12(a\widehat{\nu\nu_1}u)\theta_\nu\theta_{\nu_1}(\theta\widehat{\nu\nu_1}u)\\
&=-\frac32\theta_\eta^2(\theta Hu)(aHu)=-\frac32\theta_\eta^2(\theta Hu)(a\theta u)\\
&=-\frac32(a_\eta-(\widehat{a\theta}\eta x))^2\{(aHu)-(a\theta u)\}(a\theta u).
\end{aligned}
\]

\noindent D) Série de formes $\A_\eta(\A Hu)$ : réduisant $\A_\nu^2$.

$(\A Hu)^2$ est une forme décomposable selon C).

\noindent E) Série de formes $a_l^2(aLu)$ : réduisant $a_\eta^2(aHu)$. Les formes $(aLu)$ et $a_l(aLu)$ se décomposent en transvections sur des déterminants fonctionnels selon \S\ 3.

\noindent F) Série de formes $a_\sigma^3$ : réduisant $a_\eta^3$.

Formes $a_\sigma^2$ et $a_\sigma^3$ : réduisants $a_\nu^3$ et $a_\eta^3$ par double réduction des expressions correspondantes
\[
H_\nu a_\nu^3\quad\text{et}\quad H_\nu a_\eta^3,
\qquad
H_\nu a_\nu^3a_\eta\quad\text{et}\quad H_\nu a_\eta^4
\]
(cf. \S\ 10. C), réduction de $\A_\nu^2$ et $\A_\nu^3$).

D'où la réduction pour les formes supérieures $a_\nu s_\nu(asu)^2$, $a_\nu^2s_\nu(asu)^2$ et pour les formes dans les symboles $\R,t(a_\nu,a_\varrho^2)$, en tenant compte de (16.) :
\[
\tag{19}
\displaystyle
\left.\begin{array}{rcl}
a_\nu s_\nu(asu)^2&=&\dfrac13 i\theta_\nu^2-\dfrac1{18}iH,\\[0.5em]
a_\nu^2s_\nu(asu)^2&=&0,
\end{array}\right\}\modu{(\R,t)}.
\]
La forme $a_\sigma$ se décompose par polarisation de (12.)b ou, plus brièvement, par développement direct du produit $a_\nu^2\cdot\theta_\nu^3$ suivant le schéma :
\[
a_\nu^2\cdot\theta_{\nu_1}^3=a_\nu^2\theta_{\nu_1}\{a_{\nu_1}-(\widehat{a\theta}\nu_1x)\}^2
=-2a_\eta(aHu)(a\theta H)(\widehat{a\theta}\eta x)+(\widehat{a\theta}\nu_1x)^2a_\nu^2\theta_{\nu_1}
\]
et en tenant compte de la réduction de $Hj(j\eta x)^2$ (C) :
\[
\tag{20}\displaystyle a_\sigma=2a_\nu^2\cdot\theta_{\nu_1}^3\modu{(s,\R,t,i)}.
\]

\noindent\emph{Conséquences :}

1) $(j\eta x)^3$, $\A_\eta(\A Hu)$, $a_\sigma^2$ sont des réduisants ; $(j\eta x)^2$, $(\A Hu)^2$, $a_\sigma$ sont des formes décomposables.

2) $f$ n'entre en contraction avec $L$ que dans les produits avec des formes contragrédientes ; $f$, et donc aussi $K$ et $N$, n'entrent pas du tout en contraction avec $\sigma$ et, par conséquent, pas non plus avec $(\nu\sigma x)$. $j$ n'entre que dans des produits avec des formes contragrédientes ; $\A$ n'entre pas du tout, en produits, dans des contractions avec $H$ et $L$ (cela découle de $\A_\eta(\A Hu)$ et $(\A Hu)^2$). De la même manière $\sigma$, et donc $(\nu\sigma x)$, n'entre dans les produits dans des contractions avec aucune forme du système I. Dans les produits du système II, compte tenu des conséquences de \S\ 11, il ne reste que des puissances de $H$ et des produits de $H$ avec $L$.

\begin{center}
\textbf{\S\ 13.\quad Formes de huitième ordre (système III).}
\end{center}
\[
\text{Contraction de }\left\{\begin{array}{l}
\A\cdot j\text{ avec }\nu,\\
\theta^2,
 f\cdot j,
 N\text{ avec }H,\\
\theta\text{ avec }L,\sigma.
\end{array}\right.
\]

\noindent\emph{Formes irréductibles :}
\[
\renewcommand{\arraystretch}{1.28}
\begin{array}{@{}r@{\quad}c@{\qquad\qquad}r@{\quad}c@{\qquad}c@{}}
\text{A)}&\A_\nu(j\nu x),&\text{B)}&(\vartheta^2\eta x)^3&\cdot\\
&&&(\vartheta^2\eta x)^4&H_\vartheta(\vartheta^2\eta x)^3,\ \theta_\eta(\vartheta^2\eta x)^3,
\end{array}
\]
\[
\renewcommand{\arraystretch}{1.28}
\begin{array}{@{}r@{\quad}c@{\qquad}c@{\qquad\qquad}r@{\quad}c@{}}
\text{C)}&(aHu)H_j&\cdot&\text{D)}&H_n;\ N_\eta,\\
&\cdot&a_\eta(aHu)H_j,&&
\end{array}
\]
\[
\renewcommand{\arraystretch}{1.30}
\begin{array}{@{}r@{\quad}c@{\qquad}c@{\qquad}c@{\qquad}c@{\qquad}c@{}}
&\cdot&\cdot&(\theta Lu)^2&(\theta Lu)^3&\cdot\\
&\cdot&\theta_l&\cdot&L_\vartheta(\theta Lu)^2,\ \theta_l(\theta Lu)^2&\theta_l(\theta Lu)^3\\
\text{E)}&(\vartheta lx)^2&\cdot&\theta_l^2&\cdot&\cdot\\
&\cdot&\theta_l(\vartheta lx)^2&\cdot&\cdot&\cdot
\end{array}
\]

\noindent\emph{Réductions :}

\noindent A) $\A_\nu(j\nu x)^3$ : réduisant $a_\nu(j\nu x)^3$.

$\A_\nu(j\nu x)^2$ est réductible par la forme décomposable $a_\nu(j\nu x)^2$ selon \S\ 3. 2)b, en tenant compte du réduisant $\theta_{\vartheta'}$. En effet, d'après \S\ 11 B) on obtient :
\[
\frac3{10}\A_\nu(j\nu x)^2+\frac35a_\nu(j\nu\vartheta)a_\vartheta(j\nu x)
+\frac1{10}a_\nu(j\nu\vartheta)^2
=\frac35\theta_\nu^3\cdot\theta_{\vartheta''}^{\prime2}+\frac25\theta_\nu^3\theta_{\vartheta''}\theta_{\vartheta''}'.
\]
(réduisants $a_j$ et $\theta_{\vartheta'}$).

\noindent B) Les formes $H_\vartheta(\vartheta'\eta x)^2$, $H_\vartheta^2(\vartheta'\eta x)$, $H_\vartheta^2(\vartheta'\eta x)^2$ se décomposent comme résultant de contractions simples successives avec les formes décomposables $H_\vartheta$ et $H_\vartheta^2$, respectivement $H_\vartheta a_\eta$ et $H_\vartheta^2a_\eta$, selon \S\ 3. 2)b) et selon la relation :
\[
H_\vartheta(\vartheta'\eta x)^2=2H_\vartheta a_\eta^2\cdot f-2H_\vartheta a_\eta b_\eta.
\]

\noindent C) Série de formes $a_\eta(j\eta x)^2$ : réduisants $a_j$ et $(j\eta x)^3$.

La forme $a_\eta(j\eta x)$ se décompose : réduisant $a_j$ et contraction simple avec la forme décomposable $(j\eta x)^2$ selon \S\ 3, 2)a (spécialisation $y=uv$).

La forme $a_\eta H_j$ se décompose : 1) par contraction simple avec la forme décomposable $a_\nu(j\nu x)^2$ ;

2) par double contraction spéciale avec la forme décomposable $(j\eta x)^2$ ; d'où une relation entre formes décomposables.

À partir de $a_\nu(j\nu x)^2$ :
\[
0=\theta_\nu^3\cdot\theta_\nu+2a_\eta H_j\modu{(s,\R,t,i,J)}.
\]
De $(j\eta x)^2$ :
\[
0=\theta_\nu^3\cdot\theta_\nu-\frac32a_\eta H_j\modu{(s,\R,t,i,J)}
\]
(les produits avec contravariants sont omis), suivant le schéma :
\[
0=\frac12a_\nu(abu)^2\theta_\nu^3
+\frac12a_\nu(abu)\theta_{\nu_1}^3(\theta bu)
-\frac34(\widehat{j\eta}bu)^2+\frac34H_j(\widehat{j\eta}bu)
-\frac18f\cdot H_j^2
\]
et en permutant les symboles en $a_\nu(abu)(\theta bu)\theta_{\nu_1}^3$.

\noindent D) Série de formes $N_\eta(NHu)$ : réduisant $\A_\nu^2$.

$(NHu)$, $(n\eta x)$, $N_\eta(n\eta x)$, $H_n(NHu)$ sont des formes décomposables (cf. \S\ 12 B) ; $(NHu)^2$ se décompose par contraction simple avec la forme $(\A Hu)^2$.

\noindent E) Séries de formes $\theta_lL_\vartheta$, $\theta_l^2(\theta Lu)^2$, $\theta_l^2(\vartheta lx)$ :

Réduisants : $\theta_\eta H_\vartheta$, $\theta_\eta^2(\theta Hu)^2$, $\theta_\eta^2(\vartheta\eta x)$.

Les formes $(\theta Lu)$, $(\vartheta lx)$, $L_\vartheta$, $L_\vartheta(\theta Lu)$, $\theta_l(\theta Lu)$, $L_\vartheta(\vartheta lx)$, $\theta_l(\vartheta lx)$, $L_\vartheta^2$, $L_\vartheta^2(\theta Lu)$, $\theta_l^2(\theta Lu)$ se décomposent. (Par dualité avec les formes décomposables aux formes décomposables de \S\ 12 B).)

\noindent F) Série de formes $\theta_\sigma(\vartheta\sigma x)$ : réduisant $a_\sigma^2$.

Les formes $(\vartheta\sigma x)$ et $(\vartheta\sigma x)^2$ se décomposent, comme provenant d'une contraction simple avec la forme décomposable $a_\sigma$ ; $\theta_\sigma$, issue de la double contraction, se décompose, puisque la série de formes
\[
a_\nu^2(a\theta u)\theta_{\nu_1}^3\equiv H_j^2(j\eta x)\equiv0\;\text{modulo }(s,\R,t,i,J):
\]
\[
\theta_\sigma=a_\sigma(abu)^2=a_\nu^2(abu)^2\cdot\theta_\nu^3.
\]

\noindent\emph{Conséquences :}

$\theta^3$ ou $\theta^2K$ ne se contracte pas avec $H$ ; $\theta$, considéré uniquement comme un contravariant, entre dans les puissances ou produits en contraction avec $L$, mais n'entre pas du tout en contraction avec $\sigma$ et $(\nu\sigma x)$.

$N$ n'entre en produits dans des contractions avec aucune forme du système II, ni avec des puissances ou des produits de formes du système II. Dans les produits du système I, compte tenu des conséquences des derniers paragraphes, ne subsistent que les produits $f\cdot j$, $\A\cdot j$, les puissances de $\theta$ et les produits $\theta\cdot K$.

\begin{center}
\textbf{\S\ 14.\quad Formes de neuvième ordre (système III).}
\end{center}
\[
\text{Contraction de }\left\{\begin{array}{l}
\theta\cdot K\text{ avec }H,\\
K,j,\A\text{ avec }L,\\
j,\A\text{ avec }\sigma,\\
f\text{ avec }H^2.
\end{array}\right.
\]

\noindent\emph{Formes irréductibles :}
\[
\renewcommand{\arraystretch}{1.35}
\begin{array}{@{}r@{\quad}c@{\qquad}c@{\qquad\qquad}r@{\quad}c@{\qquad}c@{\qquad}c@{\qquad}c@{\qquad}c@{}}
&&&&\cdot&\cdot&\cdot&(KLu)^3&\cdot\\
\text{A)}&(\vartheta\cdot k,\eta,x)^4&\cdot&\text{B)}&\cdot&\cdot&\cdot&\cdot&K_l(KLu)^3\\
&\cdot&H_k(\vartheta\cdot k,\eta,x)^4,&&\cdot&\cdot&K_l^2&\cdot&\cdot\\
&&&& (klx)^3&\cdot&\cdot&\cdot&\cdot\\
&&&&\cdot&K_l(Klx)^3&\cdot&\cdot&\cdot,
\end{array}
\]
\[
\renewcommand{\arraystretch}{1.35}
\begin{array}{@{}r@{\quad}c@{\qquad}c@{\qquad}r@{\quad}c@{\qquad}r@{\quad}c@{\qquad}r@{\quad}c@{\qquad}c@{}}
\text{C)}&L_j&\cdot&\text{D)}&\A_l,&\text{E)}&(j\sigma x),&\text{G)}&(aH^2u)^3&(aH^2u)^4\\
&\cdot&L_j^2,&&&&&&\cdot&a_\eta(aH^2u)^3.
\end{array}
\]

\noindent\emph{Réductions :}

A) Les formes $H_\vartheta(k\eta x)^3$, $H_\vartheta^2(k\eta x)^2$, $H_\vartheta^2(k\eta x)^3$ se décomposent comme résultant de contractions simples successives avec des formes décomposables.

B) Séries de formes $K_lL_k$, $K_l^2(KLu)$, $K_l^2(klx)$ :

Réduisants : $\theta_\eta H_\vartheta$, $a_\eta^2(aHu)$, $\theta_\eta^2(\vartheta\eta x)$.

Les formes $K_l$, $K_l(KLu)$, $K_l(klx)$, $(KLu)^2$, $(klx)^2$ se décomposent en transvections sur des déterminants fonctionnels (\S\ 3).

Les formes $L_k$, $L_k(KLu)$, $L_k(KLu)^2$, $L_k(klx)$, $L_k(klx)^2$, $L_k^2$, $L_k^2(KLu)$, $L_k^2(klx)$, $L_k^3$, $K_l(KLu)^2$, $K_l(klx)^2$ se décomposent comme provenant d'une contraction simple avec les formes décomposables :
\[
L_\vartheta,
\ H_k(KHu),
\ L_\vartheta(\vartheta lx),
\ H_k^2,
\ L_\vartheta^2,
\ L_\vartheta^2(\theta Lu),
\ K_\eta(KHu),
\ \theta_l(\vartheta lx)
\]
selon \S\ 3. 2), a) et b).

C) Série de formes $(jlx)^3$ : réduisant $(j\eta x)^3$.

Les formes $(jlx)^2$ et $L_j(jlx)$ apparaissent par contraction simple avec la forme décomposable $(j\eta x)^2$.

D) Série de formes $\A_l(\A Lu)$ : réduisant $\A_\nu^2$.

Formes $(\A Lu)^2$, $(\A Lu)^3$ : contraction simple avec la forme décomposable $(\A Hu)^2$.

E) Série de formes $(j\sigma x)^2$ : réduisant $a_\sigma^2$ par substitution de $j$ par des formes inférieures de la série de formes (\S\ 5) :
\[
a_\sigma^2(a\theta u)^2=\frac13(j\sigma x)^2\ldots,
\qquad
 a_\sigma^2(a\theta\sigma x)^2=\frac13(j\sigma x)^4.
\]
F) Série de formes $\A_\sigma^2$ : réduisant $a_\sigma^2$.

Forme $\A_\sigma$ : réduisant $a_\sigma^2$ par substitution de $\A$ par les formes inférieures de la série :
\[
a_\vartheta a_\sigma^2=\frac23\A_\sigma\modu{\nu}.
\]

\noindent\emph{Conséquences :} Seule la forme $j$ entre en contraction avec $\sigma$ et $(\nu\sigma x)$.

$N$ n'entre pas en contraction avec $L$, puisque la forme correspondant à $\A_l$, à savoir $N_l$, est décomposable.

\begin{center}
\textbf{\S\ 15.\quad Formes de dixième ordre (système III).}
\end{center}
\[
\text{Contraction de }\left\{\begin{array}{l}
\theta^2,
\ f\cdot j\text{ avec }L,\\
\theta\text{ avec }H^2.
\end{array}\right.
\]

\noindent\emph{Formes irréductibles :}
\[
\begin{array}{c@{\quad}l@{\qquad}c@{\quad}l}
\text{A)}&\begin{array}{cc}(\vartheta^2lx)^4&\cdot\\ \cdot&
\theta_l(\vartheta^2lx)^4,\end{array}&
\text{B)}&\begin{array}{cc}(aLu)^2L_j&\cdot\\ \cdot&
a_l(aLu)^2L_j,\end{array}\\[0.8em]
\text{C)}&\begin{array}{ccc}(\theta H^2u)^3
&(\theta H^2u)^4&\cdot\\ \cdot
&H_\vartheta(\theta H^2u)^3,
&\theta_\eta(\theta H^2u)^3.\end{array}
\end{array}
\]

\noindent\emph{Réductions :}

A) et B) : décomposition des formes restantes par contraction simple avec des formes décomposables.

C) Décomposition des formes restantes selon \S\ 13 B), par le calcul dual correspondant.

\noindent\emph{Conséquences :}

$\theta^2\cdot K$ n'entre pas en contraction avec $L$ ; en conséquence, $H^2L$ n'entre pas en contraction avec $K$. $H^3$ et $H^2L$ n'entrent pas dans les contractions avec $\theta$. Avec cela, le schéma de contraction de \S\ 7 est démontré.

\begin{center}
\textbf{\S\ 16.\quad Formes d'ordre 11, 12, 13 et 15 (système III).}
\end{center}

\noindent\emph{Ordre 11.}
\[
\text{Contraction de }\left\{\begin{array}{l}
\theta\cdot K\text{ avec }L,\\
K,j\text{ avec }H^2,\\
j\text{ avec }(\nu\sigma x),\\
f\text{ avec }H\cdot L.
\end{array}\right.
\]
\noindent\emph{Formes irréductibles :}
\[
\begin{array}{c@{\quad}l@{\qquad}c@{\quad}l}
\text{A)}&\begin{array}{cc}(\vartheta\cdot k,l,x)^5&\cdot\\ \cdot&
\theta_l(\vartheta\cdot k,l,x)^5,\end{array}&
\text{B)}&\begin{array}{cc}(KH^2u)^4&\cdot\\ \cdot&
K_\eta(KH^2u)^4,\end{array}\\[0.8em]
\text{C)}&(\nu\sigma j),&
\text{D)}&(a,H\cdot L,u)^4.
\end{array}
\]

\noindent\emph{Ordre 12.}
\[
\text{Contraction de }\left\{\begin{array}{l}
\theta^3\text{ avec }L,\\
\theta\text{ avec }H\cdot L.
\end{array}\right.
\]
\noindent\emph{Formes irréductibles :}
\[
\text{E)}\ (\vartheta^3lx)^6,
\qquad
\text{F)}\ (\theta,H\cdot L,u)^4,
\qquad L_\vartheta(\theta,H\cdot L,u^4).
\]

\noindent\emph{Ordre 13.}
\[
\text{Contraction de }K\text{ avec }H\cdot L.
\]
\noindent\emph{Formes irréductibles :}
\[
\text{G)}\ (K,H\cdot L,u)^5,
\qquad K_\eta(K,H\cdot L,u)^5.
\]

\noindent\emph{Ordre 15.}
\[
\text{Contraction de }K\text{ avec }H_3.
\]
\noindent\emph{Forme irréductible :}
\[
\text{H)}\ (KH^3u)^6.
\]

\noindent\emph{Réductions :}

Les formes $H_{j^3}^2$, $H_{j^4}^2$ admettent le réduisant $L_j^3$, de la relation :
\[
(\nu lx)L_x^3=-\frac12H^2\modu{(\sigma,s)}.
\]
La décomposition ou la réductibilité des formes restantes découle des considérations des paragraphes précédents.

\noindent\emph{Conséquences :}

D'après le schéma de contraction de \S\ 7 et les conséquences de \S\S\ 8--15, les formes irréductibles du système III sont ainsi épuisées (117 formes).

\clearpage

\begin{center}
{\Large\bfseries Chapitre IV.\quad Le système relativement complet modulo $(\varrho,t)$.}
\end{center}

\begin{center}
\textbf{\S\ 17.\quad Réduction du module $(ss'u)^4$ et du système de $s$.}
\end{center}

Pour former le système relativement complet immédiatement ci-dessus, par analogie avec \S\ 7 il faut former le système de $s$ par rapport au module suivant
\[
\nu(s)=(ss'u)^4
\]
et le contracter avec les formes du système relativement complet modulo $s$. On montrera qu'en introduisant le module $(\varrho,t)$ comme module supérieur,
\[
\begin{array}{ll}
\text{A)}&\text{le module }(ss'u)^4\text{ devient réductible,}\vphantom{\dfrac11}\\
\text{B)}&\text{le système de }s\text{ se compose de la seule forme }s.
\end{array}
\]

A) La réduction du module $(ss'u)^4$ s'obtient immédiatement à partir du réduisant $s_\nu^2$ trouvé par la formule (16.), \S\ 11. De
\[
 s_\nu^2=\frac43 i a_\nu^2+\frac45\theta_\varrho^2+\frac85(atu)^2-
 \frac{26}{15}(\nu\varrho x)^2+\frac16J\cdot u_x^2-
 \frac29i^2\cdot u_x^2
\]
Il s'ensuit, par polarisation de $x$ vers $\nu_1$ :
\par\smallskip\noindent
\begin{minipage}[t]{3.4em}\vspace*{0pt}(21.)\srcfnmark{*)}\end{minipage}%
\begin{minipage}[t]{\dimexpr\linewidth-3.4em\relax}\vspace*{0pt}\centering
\(\displaystyle\begin{aligned}
(ss'u)^4&=-\frac23J\cdot \nu+6s_\nu^2s_{\nu_1}^2\\
&=\frac{24}{5}\theta_\nu^2\theta_\varrho^2+
\frac{48}{5}a_\nu^2(atu)^2-
\frac{52}{5}H_\varrho^2+
\frac43i\cdot(asu)^4+\frac13J\cdot \nu-
\frac49i^2\cdot \nu,
\end{aligned}\)
\end{minipage}\par\smallskip
et avec elle la réduction du module $(ss'u)^4$.

\srcfntext{*)}{De la formule (21.) découle la relation, mentionnée dans la note de l'introduction,
entre les trois invariants du neuvième ordre, par application de la formule (10.)c :
\begin{equation*}
\tag{22}
\begin{aligned}
(ass')^4&=\frac{24}{5}\theta_\nu^2\theta_\varrho^2a_\nu^2a_\vartheta^2+
\frac{48}{5}a_\nu^2b_\nu^2(tab)^2-
\frac{52}{5}H_\varrho^2a_\eta^4+
\frac53J\cdot i-\frac49i^3,\\
(ass')^4&=\frac{24}{5}a_\varrho^2a_{\varrho'}^2-
\frac{12}{5}t_\varrho^2+\frac53J\cdot i-\frac49i^3.
\end{aligned}
\end{equation*}}

B) La réduction de
\[
Z=(ss'u)^2=\theta(s)
\]
et à partir de là la réduction du système de $s$ s'obtient en remplaçant la forme $Z$ par la forme inférieure $g_\nu^2$ selon la formule (8.)b, \S\ 7, en tenant compte de la relation
\[
 s_\eta^2=s_\nu^2(\nu\nu_1x)^2-\frac13Z.
\]
La forme $g_\nu^2$ a, selon la formule (13.), \S\ 11, le réduisant $s_\nu^2$ comme facteur. On obtient, par les formules (8.) et (16.) :
\[
 g_\nu^2=-Z+\frac{14}{5}ia_\eta^2+\frac{14}{15}i(asu)^2
 \modu{(\varrho,t)},
\]
et, par les formules (13.), (14.), (15.), (16.) :
\[
\begin{aligned}
g_\nu^2&=2[s_\vartheta^2]_{\nu^2}-\frac43 i\A_\nu^2
      =2s_\vartheta^2s_\nu^2-\frac65[s_\vartheta^2(\theta su)^2]\widehat{\nu x}^{2}
        -\frac43i\A_\nu^2\\
&=\frac45i\cdot a_\eta^2-\frac25i\cdot(asu)^2+\frac13J\cdot\theta
\modu{(\varrho,t)}:
\end{aligned}
\]
De là, il s'ensuit
\[
\tag{23}
Z=2i\cdot a_\eta^2+\frac43i\cdot(asu)^2-\frac13J\cdot\theta
 \modu{(\varrho,t)}.
\]

Le système relativement complet modulo $((ss'u)^4,\varrho,t)$ passe donc à un système relativement complet modulo $(\varrho,t)$ ; et de là naît, par transvection sur le système connu de deux formes quadratiques, le système absolument complet. Pour former le système relativement complet modulo $(\varrho,t)$, puisque selon la formule (23.) le système de $s$ se réduit à la forme $s$, il faut former les transvections du système relativement complet modulo $(s,\varrho,t)$ avec $s$ et avec les puissances de $s$. On verra que les puissances de $s$ n'entrent en contraction qu'avec le système I.

Nous définissons comme \og{} formes supérieures \fg{}, parmi les formes du même ordre et du même degré, celles qui sont issues des formes supérieures du système relativement complet modulo $(s,\varrho,t)$ par contraction avec $s$, en considérant la contraction avec $s$ uniquement comme une polarisation des formes originales. Ainsi, l'ordre des formes individuelles est déterminé de manière univoque (cf. \S\ 1, séries de formes 2). Par exemple :
\[
(\theta_\eta(\theta Hu),s,u)^2\text{ est supérieure à }s_\eta((\theta Hu)^2,s,u)\text{\srcfnmark{*)}},
\]
\[
(\theta_\eta(\theta Hu)^2,s,u)\text{ est supérieure à }(\theta_\eta(\theta Hu),s,u)^2.
\]

Nous divisons le système résultant en quatre sous-systèmes :
\[
\begin{array}{ll}
\text{Système }s_I:&\text{issu de la contraction de }s\text{ et de puissances de }s\text{ avec le système I;}\\
\text{Système }s_{II}:&\text{issu de la contraction de }s\text{ avec le système II;}\\
\text{Système }s_{III}:&\text{issu de la contraction de }s\text{ avec les formes individuelles}\\
&\text{(sans produits) du système III et avec des produits d'une}\\
&\text{forme de chacun des systèmes I et II;}\\
\text{Système }s_{IV}:&\text{issu de la contraction de }s\text{ avec des produits d'une forme de chacun}\\
&\text{des systèmes I et III, II et III, III et III.}
\end{array}
\]
Remarque : Désormais, toutes les formules doivent être comprises modulo $(\varrho,t)$ ; à partir de la formule (27.), modulo $(\varrho,t,i,J,\text{ formes supérieures})$, sans que cela soit expressément indiqué à chaque fois.

Pour la réduction on regroupe les formules obtenues précédemment :
\[
\tag{24}
\begin{aligned}
s_\nu^2&=\frac43i\,a_\nu^2+\frac16J\cdot u_x^2-\frac29i^2\cdot u_x^2,\\
s_\nu^3&=\frac13J\cdot u_x,
& s_\nu^4&=J,\\
g&=2s_\vartheta^2-\frac43i\Delta,\\
s_\vartheta^2(\theta su)&=0,\\
s_\vartheta^2(\theta su)^2&=\frac89i\cdot a_\nu^2-\frac4{27}i^2\cdot u_x^2,\\
Z&=2i\cdot a_\eta^2+\frac43i(asu)^2-\frac13J\cdot\theta,\\
g_\nu&=-s_\eta+u_x\left\{2ia_\eta^2-\frac23i(asu)^2+\frac23J\cdot\theta\right\},\\
g_\nu^2&=\frac45ia_\eta^2-\frac25i(asu)^2+\frac13J\cdot\theta,
&g_\nu^3&=0,\quad g_\nu^4=0.
\end{aligned}
\]
\srcfntext{*)}{Pour abréger l'écriture, nous introduisons dans le système le terme de la
transvection à la place de la transvection complète $[(\theta Hu)^2]_{\widehat{su}}s$.
Lors de l'établissement des formules, toutes les transvections sont remplacées par leurs termes
initiaux ; par exemple $(\theta_\eta(\theta Hu),s,u)^2$ par
$\theta_\eta(\theta Hu)(\theta su)^2$.}

\noindent\emph{Conséquences :}
\[
 s_\nu^2,\quad s_\vartheta^2(\theta su),\quad s_\vartheta^2s_\nu,\quad s_\vartheta^2\theta_\nu
\]
sont des réduisants.

\clearpage

\begin{center}
\textbf{\S\ 18.\quad Système $s_I$.}
\end{center}

\[
\text{Contraction de}\left\{\begin{array}{l}
s\text{ avec } f;\ \theta;\ K,j,\Delta;\ f\cdot j,N;\ \Delta\cdot j,\\
s^2\text{ avec }\theta,K.
\end{array}\right.
\]

\noindent\emph{Formes irréductibles.}

\noindent Ordre 5 :
\[
(asu);\qquad (asu)^2;\qquad (asu)^3;\qquad (asu)^4.
\]

\noindent Ordre 6 :
\[
\begin{array}{cccc}
(\theta su)&(\theta su)^2&(\theta su)^3&(\theta su)^4\\
s_\vartheta&s_\vartheta(\theta su)&s_\vartheta(\theta su)^2&s_\vartheta(\theta su)^3\\
\cdot&s_\vartheta^2&\cdot&\cdot
\end{array}
\]

\noindent Ordre 7 :
\[
\begin{array}{c@{\quad}c@{\quad}c@{\quad}c@{\qquad}l}
\cdot&(Ksu)^2&(Ksu)^3&(Ksu)^4&\\
\text{A)}\ s_k&s_k(Ksu)&s_k(Ksu)^2&s_k(Ksu)^3&
\text{B)}\ s_j,\\
\cdot&s_k^2&\cdot&\cdot,&
\text{C)}\ (\Delta su),\ (\Delta su)^2.
\end{array}
\]

\noindent Ordre 8 :
\[
\begin{array}{ll}
\text{A)}\ s_j(asu),\ s_j(asu)^2,\ s_j(asu)^3,
&
\text{B)}\ \begin{array}[t]{cc}
\cdot&(Nsu)^2\\
s_n&s_n(Nsu).
\end{array}
\end{array}
\]

\noindent Ordre 10 :
\[
\text{A)}\ s_j(\Delta su),\qquad \text{B)}\ s_\vartheta(\theta s'u)^4.
\]

\noindent Ordre 11 :
\[
\begin{array}{cc}
(Ks^2u)^6&\cdot\\
s_k(Ks^2u)^5&s_k(Ks^2u)^6.
\end{array}
\]

\noindent\emph{Réductions :}

Les séries de formes
\[
 s_\vartheta^2(\theta su),\quad s_k^2(Ksu),\quad s_j^2,\quad (\Delta su)^3,\quad (Nsu)^3
\]
ont le réduisant $s_\vartheta^2(\theta su)$ pour facteur ; $s_j^2$ et $(\Delta su)^3$ s'obtiennent en remontant de $j$ et $\Delta$ à des formes inférieures de la série de formes :
\[
\begin{aligned}
s_\vartheta^2(\theta su)(a\theta u)^3&=-\frac12s_j^2,\\
s_\vartheta^2(\theta su)(a\theta u)^2&=\frac13(\Delta su)^3,\\
s_\vartheta^2(\theta su)^2(a\theta u)^2&=\frac13(\Delta su)^4+\frac13s_j^2.
\end{aligned}
\]
Formes $s_\vartheta^2s_{\vartheta'}$ et $s_\vartheta^2s_{\vartheta'}^2$ : réduisants $s_\vartheta^2(\theta su)$ et $\theta_{\vartheta'}$ :
\[
\begin{aligned}
s_\vartheta^2s_{\vartheta'}&=s_\vartheta^2\theta_{\vartheta'}-s_\vartheta^2(\theta s\widehat{\vartheta' x}),\\
s_\vartheta^2s_{\vartheta'}^2&=s_\vartheta^2s_{\vartheta'}\theta_{\vartheta'}-s_\vartheta^2s_{\vartheta'}(\theta s\widehat{\vartheta' x}).
\end{aligned}
\]
Série de formes $s_j(\Delta su)^2$ : réduisants $\Delta_j$ et $(\Delta su)^3$.

\noindent\emph{Conséquences :}

$s$ n'entre pas en puissances dans la contraction avec $f,j,\Delta,N$. Les formes $f,j,\Delta,N$ n'entrent dans la contraction que dans des produits avec des formes contragrédientes ; toutefois, dans les produits avec $f$, ces formes peuvent encore être quadratiques en $x$, et dans les produits avec $\Delta$ elles peuvent encore être linéaires en $x$. Il faut donc considérer les produits de formes suivants :
\[
 f\cdot(0,\lambda),\quad f\cdot(1,\lambda),\quad f\cdot(2,\lambda),\quad
 j\cdot(\mu,0),\quad N,\quad
 \Delta\cdot(0,\lambda)\quad \Delta\cdot(1,\lambda)\quad(\lambda>0),
\]
où $(\mu,\lambda)$ désigne une forme de degré $\mu$ en $x$ et de degré $\lambda$ en $u$.

Ni $\theta$ ni $K$ n'entrent, dans la contraction avec $s$ ou $s^2$, en puissances ou en produits avec des formes arbitraires. Pour les produits avec le déterminant fonctionnel $K$, $K\cdot\varphi_x$ ($\varphi\ne f$ ou $\theta$), la relation entre formes décomposables est
\[
K\cdot\varphi_x=(a\theta u)\varphi_x=f\cdot(\varphi\theta u)-\theta\cdot(\varphi a u).
\]
Le schéma ci-dessus pour la contraction du système $s_I$ est ainsi démontré.

\begin{center}
\textbf{\S\ 19.\quad Système $s_{II}$.}
\end{center}

Contraction de $s$ avec $\nu;\ H;\ L;\ H^2$.

\noindent\emph{Formes irréductibles :}
\[
\begin{array}{c@{\qquad}c@{\qquad}c@{\qquad}c}
\text{Ordre 6} & \text{Ordre 8} & \text{Ordre 10} & \text{Ordre 12}\\[0.3em]
s_\nu & (Hsu),\ (Hsu)^2 & (Lsu)^2,\ (Lsu)^3 & (H^2su)^3\\
\cdot & s_\eta & s_l & \cdot\\
\cdot & \cdot & \cdot & \cdot\\
s_\nu^4=J & \cdot & \cdot & \cdot
\end{array}
\]

\noindent\emph{Réductions :}

Séries de formes $s_\eta(Hsu)$ et $s_l(Lsu)$ : réduisant $s_\nu^2$.

Série de formes $s_\sigma$ : réduisant $s_\nu^2$ issu de $H$, $s_\nu^2=\dfrac23s_\sigma$.

$(H^2su)^4$ se décompose par la formule (7.)a (cf. \S\ 11 A). Par conséquent $(H\cdot L,s,u)^4$ se décompose également.

\noindent\emph{Conséquences :}

$s^2$ n'entre pas dans la contraction avec le système II. La forme $\nu$ n'entre que dans des produits avec des formes contragrédientes. La forme $H$, considérée comme covariant, entre dans la contraction dans des produits avec des formes $(1,\lambda)$, $(2,\lambda)$, $(3,\lambda)$, pour $\lambda$ arbitraire. Les formes $\sigma$ et $(\nu\sigma x)$ n'entrent pas du tout, et $L$, comme déterminant fonctionnel, n'entre pas dans des produits pour la contraction ; les transvections complètes sur les covariants du système III deviennent réductibles. Les produits des systèmes I et II qui restent sont
\[
f\cdot \nu,\qquad \Delta\cdot \nu.
\]

\begin{center}
\textbf{\S\ 20.\quad Formes du 7e ordre (Système $s_{III}$).}
\end{center}

Contraction de $s$ avec $f\cdot \nu$, $a_\nu$, $a_\nu^2$.

\noindent\emph{Formes irréductibles :}
\[
\begin{gathered}
s_\nu(asu),\quad a_\nu(asu),\quad s_\nu(asu)^2,\quad a_\nu(asu)^2,\quad
s_\nu(asu)^3,\quad a_\nu(asu)^3,\\
a_\nu s_\nu,\qquad a_\nu s_\nu(asu),\qquad a_\nu^2(asu),\qquad a_\nu^2(asu)^2.
\end{gathered}
\]

\noindent\emph{Réductions :}

Les réductions évidentes qui proviennent de la contraction directe avec les réduisants $s_\nu^2$ et $s_\eta(Hsu)$ ne seront pas mentionnées, pas plus que la décomposition des transvections avec des déterminants fonctionnels.

Pour les formes $a_\nu s_\nu(asu)^2$ et $a_\nu^2s_\nu(asu)^2$, voir la formule (19.), \S\ 12 ; en outre,
\[
a_\nu s_\nu(asu)^3=a_\nu(a\widehat{\nu_1\nu_2}u)^3(\nu_1\nu_2\nu)=0.
\]
Série de formes $a_\nu^2s_\nu$ : réduisant $s_\vartheta^2(\theta su)$, obtenu en remontant de $a_\nu^2$ à des formes inférieures, c'est-à-dire par double réduction des expressions
\[
s_\vartheta^2(a\theta s)(a\theta u),\qquad
s_\vartheta^2(a\theta s)(a\theta u)^2
\]
d'après l'ansatz
\[
\begin{aligned}
s_\vartheta^2(a\theta s)(a\theta u)
&=[(a\theta u)^2]s^3+\frac15[a_\vartheta(a\theta u)^2]s^2+\frac1{60}[a_\vartheta^2(a\theta u)^2]s\\
&=\frac12a_\nu^2s_\nu-\frac23a_\nu s_\nu^2,\\
s_\vartheta^2(a\theta s)(a\theta u)^2
&=\frac45[a_\vartheta(a\theta u)^2]_{\widehat{us}}s^2+\frac{23}{180}[a_\vartheta^2(a\theta u)^2]_{\widehat{us}}s\\
&=-\frac12a_\nu^2s_\nu(asu).
\end{aligned}
\]
Ici $(a_\nu^2b_\nu(abu)=0)$.

Les formules obtenues sont
\[
\tag{25}
\begin{array}{rcl@{\qquad}rcl}
\underline{a_\nu^2s_\nu}&=&\dfrac43i\cdot a_\nu^2b_\nu,
& \underline{a_\nu s_\nu(asu)^2}&=&\dfrac13i\cdot\theta_\nu^2-\dfrac1{18}iH,\\[0.7em]
\underline{a_\nu s_\nu(asu)^3}&=&0,
& \underline{a_\nu^2s_\nu(asu)}&=&0,\\[0.7em]
\underline{a_\nu^2s_\nu(asu)^2}&=&0.
\end{array}
\]

\noindent\emph{Conséquences :}

1) $a_\nu s_\nu(asu)^2$ et $a_\nu^2s_\nu$ sont des réduisants.

2) On obtient ainsi les valeurs des formes $(\Delta su)^3$, $(\Delta su)^4$, déjà reconnues réductibles en \S\ 19 ; de même pour la série de formes correspondante $(agu)^2$, qui, par contraction avec $\nu$, donne lieu à une double réduction au moyen du réduisant $g_\nu$ :
\[
\tag{26}
\begin{aligned}
\text{a)}\quad (\Delta su)^3&=-\frac65a_\nu^2(asu)+3a_\nu s_\nu(asu)+2i\cdot\theta_\nu(\vartheta\nu x),\\
\text{b)}\quad (\Delta su)^4&=-\frac25a_\nu^2(asu)^2+\frac43i\theta_\nu^2+\frac19iH.
\end{aligned}
\]
Des formules (24.) et (26.), modulo $(i,J)$ (cf. p. 71),
\[
\tag{27}
\begin{aligned}
\text{a)}\quad (agu)^2&=\frac23(\Delta su)^2-\frac43a_\nu s_\nu+\frac35s\cdot a_\nu^2,\\
\text{b)}\quad (agu)^3&=2a_\nu s_\nu(asu)-a_\nu^2(asu)^2,\\
\text{c)}\quad (agu)^4&=-a_\nu^2(asu)^2.
\end{aligned}
\]


\begin{center}
\textbf{\S\ 21.\quad Formes du 8e ordre (Système $s_{III}$).}
\end{center}

Contraction de $s$ avec les formes du 4e ordre (Système III). (\S\ 9.)

\noindent\emph{Formes irréductibles :}
\[
\begin{array}{c|c|c}
\cdot
%
& ((\vartheta \nu x),s,u)^2,\ (\theta_\nu su)
& \bigstar\\[0.45em]
(\vartheta \nu x)s_\nu,\ ((\vartheta \nu x^2),s,u)
& ((\vartheta \nu x)^2,s,u)^2,\ \theta s_\nu
& \\[0.45em]
(\vartheta \nu x)^2s_\nu & ((\vartheta \nu x)^2,s,u)s_\nu,\ (\theta_\nu(\vartheta \nu x)^2,s,u)
&
\end{array}
\]
\[
\begin{array}{c|c|c|c}
\bigstar
& ((\vartheta \nu x),s,u)^3,\ (\theta_\nu su)^2
& ((\vartheta \nu x),s,u)^4,\ (\theta_\nu su)^3
& \cdot\\[0.45em]
& ((\vartheta \nu x)^2,s,u)^3,\ (\theta_\nu(\vartheta \nu x),s,u)^2,\ (\theta_\nu^2su)
& (\theta_\nu(\vartheta \nu x),s,u)^3,\ (\theta_\nu^2su)^2
& (\theta_\nu(\vartheta\nu x),s,u)^4.\\[0.45em]
& \cdot
& (\theta_\nu^3su)
& \cdot
\end{array}
\]

\noindent\emph{Réductions :}

$((\vartheta \nu x),s,u)s_\nu$ se décompose comme transvection sur des déterminants fonctionnels, d'après \S\ 3.

Série de formes $((\vartheta \nu x),s,u)^2s_\nu$ : réduisants $s_\nu^2$ et $s_\vartheta^2\theta_\nu$ :
\[
(\widehat{\vartheta \nu}su)s_\nu(\theta su)=s_\nu s_\vartheta(\theta su)=-\frac12(\widehat{\vartheta \nu}su)^2(\theta su),
\]
\[
\theta_\nu(\widehat{\vartheta \nu}su)s_\nu=\theta_\nu s_\nu s_\vartheta=-\frac12\theta_\nu(\widehat{\vartheta \nu}su)^2.
\]
Série de formes $\theta_\nu s_\nu(\theta su)$ : réduisant $a_\nu s_\nu(asu)^2$ par double réduction des expressions
\[
a_\nu s_\nu(asu)^2(abu)\quad\text{et}\quad a_\nu s_\nu(asu)^2(abu)(sbu);
\]
$\theta_\nu s_\nu(\theta su)^3$ par double réduction de $(agu)^3(abu)(bgu)^3$.

Série de formes $\theta_\nu(\vartheta \nu x)s_\nu$ : réduisant $a_\nu^2s_\nu$ issu de $\theta_\nu(\vartheta \nu x)s_\nu=a_\nu^2(abu)s_\nu$.

Série de formes $(\theta_\nu(\vartheta \nu x)^2,s,u)$ : double réduction des formes de la série $\theta_\nu(\widehat{\vartheta \nu}su)s_\nu$, qui appartiennent en même temps aux séries de formes $((\vartheta \nu x)su)^2s_\nu$ et $\theta_\nu(\vartheta \nu x)s_\nu$.

La forme $(\theta_\nu(\vartheta \nu x),s,u)$ se décompose comme transvection simple sur le déterminant fonctionnel $\theta_\nu(\vartheta \nu x)=a_\nu^2(abu)$.

La forme $(\vartheta \nu x)^2,s,u)^4$ : double réduction de l'expression $(a\A u)^2(\A su)^4$, ou encore de $(\theta gu)^4$, à partir des formules (27.) et de façon analogue à la réduction de $(j\eta x)^4$ (\S\ 12 C)).

La double réduction donne les formules :
\[
\tag{28}
\begin{aligned}
0&=\theta_\nu s_\nu(\theta su)-\frac16(\widehat{\vartheta \nu}su)^2(\theta su)-\frac1{12}(Hsu),\\
0&=\theta_\nu s(\theta su)^2-\frac14(\widehat{\vartheta \nu}su)^2(\theta su)^2-\frac1{24}(Hsu)^2,\\
0&=(\widehat{\vartheta \nu}su)^2(\theta su)^2+2\theta_\nu(\widehat{\vartheta \nu}su)(\theta su)^2
  +3\theta_\nu^2(\theta su)^2-\frac13H(su)^2,\\
0&=\theta_\nu s_\nu(\theta su)^3+\frac12\theta_\nu(\widehat{\vartheta \nu}su)(\theta su)^3,\\
0&=\theta_\nu s_\nu s_\vartheta-\frac14s_\eta,
\quad 0=\theta_\nu s_\nu s_\vartheta(\theta su),
\quad 0=\theta_\nu s_\nu s_\vartheta(\theta su)^2.
\end{aligned}
\]
Le dernier groupe correspond à $(\theta_\nu(\vartheta \nu x)^2,s,u)^2\ldots$.

\noindent\emph{Conséquences :}

1) $\theta(\vartheta \nu x)s_\nu$, $(\widehat{\vartheta \nu}su)(\theta su)s_\nu$, $\theta_\nu(\widehat{\vartheta \nu}su)^2$, $(\widehat{\vartheta \nu}su)^2(\theta su)^2$ sont des réduisants.

2) Les formes de 4e ordre $(5,\lambda)$, $(6,\lambda)$, etc. n'entrent pas dans la contraction avec $s^2$.

3) Pour la série de formes $(\theta gu)^2$, les formules (27.), compte tenu des réductions de ce paragraphe, donnent les formules suivantes :
\[
\tag{29}
\begin{aligned}
\text{a)}\quad (\theta gu)^2&=\frac43\theta_\nu s_\nu+\frac23s_\nu s_\vartheta-\frac13s\cdot\theta_\nu^2-\frac19s\cdot H-\frac13f\cdot a_\nu^2(asu)^2+\frac13a_\nu^2\cdot(asu)^2,\\
\text{b)}\quad (\theta gu)^3&=-\frac13(\widehat{\vartheta \nu}su)^2(\theta su)-\theta_\nu(\widehat{\vartheta \nu}su)(\theta su)-\theta_\nu^2(\theta su),\\
\text{c)}\quad (\theta gu)^4&=2\theta_\nu^2(\theta su)^2-\frac13(Hsu)^2,\\
g\vartheta(\theta gu)&=(\widehat{\vartheta \nu}su)(\vartheta \nu x)s_\nu-\theta_\nu(\vartheta \nu x)(\widehat{\vartheta \nu}su),\\
g\vartheta(\theta gu)^2&=\frac23s\cdot\theta_\nu^3,
\quad g\vartheta(\theta gu)^3=\frac13\theta_\nu^3(\theta su),
\quad g\vartheta(\theta gu)^4=0.
\end{aligned}
\]

\begin{center}
\textbf{\S\ 22.\quad Formes du 9e ordre (Système $s_{III}$).}
\end{center}

Contraction de $s$ avec les formes du 5e ordre (Système III) et le produit $\A\cdot \nu$ (\S\ 10).

\noindent\emph{Formes irréductibles :}
\[
\begin{array}{c@{\quad}c|c|c|c}
\text{A)}& \cdot & \cdot & (K_\nu su)^2 & \bigstar\\[0.45em]
& \cdot & ((k\nu x)^2,s,u)^2,\ K_\nu s_\nu & ((k\nu x)^2,s,u)^3 & \\[0.45em]
& (k\nu x)^2s_\nu,\ ((k\nu x)^3,s,u) & ((k\nu x)^3,s,u)^2 & \cdot & \cdot\\[0.45em]
& (k\nu x)^3s_\nu & ((k\nu x)^3,s,u)s_\nu,\ (K_\nu(k\nu x)^3,s,u) & \cdot & \cdot\\[0.45em]
\bigstar & (K_\nu su)^3 & (K_\nu su)^4 & \cdot & \\[0.45em] & \cdot & \cdot & (K_\nu^2su)^4, &
\end{array}
\]
\[
\begin{array}{c@{\quad}c|c@{\qquad}c@{\quad}c}
\text{B)}& (j\nu x)s_\nu,\ ((j\nu x)^2,s,u) & ((j\nu x)^2,s,u)^2 && \\
& (j\nu x)^3,s,u) & \cdot, && \text{C)}\quad s_\nu(\A su),\ (\A_\nu su),
\end{array}
\]
\[
\begin{array}{c@{\quad}c|c|c}
\text{D)}& \cdot & ((aHu),s,u)^2,\ ((aHu)^2,s,u) & \bigstar\\[0.45em]
& (aHu)s_\eta,\ (a_\eta su) & (aHu)^2s_\eta,\ (a_\eta su)^2 & \\[0.45em]
& \cdot & (a_\eta^2su) & \\
\bigstar& ((aHu),s,u)^3,\ (aHu)^2,s,u)^2 & ((aHu),s,u)^4 & \\
& (a_\eta su)^3,\ (a_\eta(aHu),s,u)^2,\ (a_\eta(aHu)^2,s,u) & (a_\eta(aHu),s,u)^3. &
\end{array}
\]

\noindent\emph{Réductions :}

A) Les réductions résultent directement des réductions du \S\ 21 par contraction simple. La série de formes $K_\nu s_\nu(Ksu)^2$ a pour facteurs les réduisants $a_\nu s_\nu(asu)^2$ et $\theta_\nu s_\nu(\theta su)^2$. Par conséquent, les formes supérieures $K_\nu^2(Ksu)^2$, $K_\nu^2(Ksu)^3$ se décomposent d'après l'ansatz :
\[
0=a_\nu s_\nu(asu)^2(a\theta u)=K_\nu s_\nu(Ksu)^2=\theta_\nu s_\nu(\theta su)^2(a\theta u).
\]

B) Série de formes $(j\nu x)^2s_\nu$ : réduisant $a_\nu^2s_\nu$ issu de
\[
(j\nu x)^2s_\nu=3a_\nu^2s_\nu(a\theta u)^2.
\]
$((j\nu x)^3,su)^2$ : réduisant $a_\nu^2s_\nu$ issu de
\[
((j\nu x)^3,s,u)^2=-6a_\nu^2s_\nu(a\theta u)^2(\theta su).
\]

C) Formes $\A_\nu(\A su)^2$ et $s_\nu(\A su)^2$ : réduisant $g_\nu$ issu de la formule (27.)a.

Forme $\A_\nu s_\nu$:
1) réduisant $a_\nu^2s_\nu$ issu de $a_\vartheta a_\nu^2s_\nu=\frac23\A_\nu s_\nu$ ;
2) se décompose par la formule (27.)a au moyen de la double réduction de l'expression $(\A s\widehat{\nu x})^2$.

Par conséquent, la forme supérieure $a_\eta s_\eta$ se décompose.

D) Série de formes $(aHu)^2s_\eta(asu)$ : réduisant $a_\nu s_\nu(asu)^2$ par double réduction de la série de formes $[a_\nu s_\nu(asu)^2]\widehat{\nu_1x}$ ; réduction de $(aHu)^2s_\eta(asu)^2$ à partir de $a_\nu s_\nu(asu)^2$ et $a_\nu s_\nu(asu)^3$. D'où la réduction de $(a_\eta,s,u)^4$.

Série de formes $a_\eta(aHu)s_\eta$ : réduisants $a_\nu^2s_\nu$, $a_\eta^2s_\eta$ par double réduction de l'expression $(\A s\widehat{\nu x})^3$, puisque $a_\eta^2s_\eta$ ne provient pas, par contraction, du terme réductible $a_\nu^2s_\nu(\nu\nu_1x)$ de la forme $a_\eta(aHu)s_\eta$.

Série de formes $(a_\eta^2su)^2$ : réduisant $a_\nu^2s_\nu$ issu de
\[
a_\nu^2s_\nu s_{\nu_1}=-\frac12a_\eta^2(Hsu)^2.
\]
La forme $(a_\eta(aHu),s,u)$ se décompose comme transvection sur un déterminant fonctionnel.

Les formules de réduction obtenues sont :
\[
\tag{30}
\begin{aligned}
\text{a)}\quad 0&=(aHu)^2(asu)s_\eta-a_\eta(asu)(Hsu)^2-a_\eta(aHu)(asu)(Hsu),\\
\text{b)}\quad 0&=(aHu)^2(asu)^2s_\eta-a_\eta(aHu)(asu)^2(Hsu),\\
\text{c)}\quad 0&=a_\eta(asu)^2(Hsu)^2,\\
0&=a_\eta s_\eta+\frac14a_\nu^2\cdot g-\frac14s\cdot a_\eta^2,\\
0&=a_\eta(aHu)s_\eta-\frac12a_\eta^2(Hsu),
\quad 0=a_\eta^2(Hsu)^2,\ldots,
\quad 0=a_\eta^2s_\eta.
\end{aligned}
\]

\noindent\emph{Conséquences :}

1) $(j\nu x)^2s_\nu$, $\A_\nu s_\nu$, $\A_\nu(\A su)^2$, et par conséquent $N_\nu(Nsu)^2$, $(aHu)^2(asu)s_\eta$, $a_\eta(aHu)s_\eta$, $a_\eta^2(Hsu)^2$, sont des réduisants ; $a_\eta s_\eta$ est une forme décomposable qui passe, sous toutes les contractions simples et doubles, en formes décomposables.

2) Les formes de 5e ordre $(5,\lambda)$, $(6,\lambda)$, etc. n'entrent pas dans la contraction avec $s^2$.

\begin{center}\textbf{\S\ 23.\quad Formes du dixième ordre (Système $s_{III}$).}\end{center}
Contraction de $s$ avec les formes du sixième ordre (Système III) (\S\ 11).

\noindent\emph{Formes irréductibles :}
\[
\begin{array}{r@{\quad}c}
\text{A)}&
\begin{array}{c|c}
(\vartheta^2\nu x)^3s_\nu&\cdot\\[0.45em]
\cdot&(\vartheta^2\nu x)^3s_\nu s_\vartheta,\ \theta_\nu(\vartheta^2\nu x)^3s_\vartheta,
\end{array}
\end{array}
\]
\[
\begin{array}{r@{\quad}c}
\text{B)}&
\begin{array}{c|ccc}
\cdot&(a_\nu(j\nu x),s,u)^2&(a_\nu(j\nu x),s,u)^3&(a_\nu(j\nu x),s,u)^4\\[0.45em]
a_\nu(j\nu x)s_\nu&\cdot&(a_\nu^2(j\nu x),s,u)^2&(a_\nu^2(j\nu x),s,u)^3,
\end{array}
\end{array}
\]
\[
\begin{array}{r@{\quad}c}
\text{C)}&
\begin{array}{c|c|c}
\cdot&\cdot&((\theta Hu),s,u)^2,\ ((\theta Hu)^2,s,u)\\[0.35em]
\cdot&((\vartheta\eta x),s,u)^2,\ (\theta Hu)s_\eta,\ (\theta_\eta su)&((\vartheta\eta x),s,u)^3,\ (\theta_\eta su)^2\\[0.35em]
((\vartheta\eta x)^2,s,u)&\cdot&(\theta_\eta^2su)\\[0.35em]
\cdot&(\theta_\eta(\vartheta\eta x)^2,s,u)&\cdot
\end{array}
\end{array}
\]
\[
\begin{array}{c|c}
((\theta Hu),s,u)^3,\ ((\theta Hu)^2,s,u)^2&((\theta Hu),s,u)^4\\[0.45em]
(H_\vartheta(\theta Hu),s,u)^2,\ (\theta_\eta(\theta Hu),s,u)^2,\ (\theta_\eta(\theta Hu)^2,s,u)&(\theta_\eta su)^4,\ (\theta_\eta(\theta Hu),s,u)^3\\[0.45em]
\cdot&(\theta_\eta^2(\theta Hu),s,u)^2
\end{array}
\]
\noindent\emph{Réductions :}
A) Les formes $((\vartheta^2\nu x)^3,s,u)^2$, $((\vartheta^2\nu x)^3,s,u)^3$ se décomposent :
\[
(\vartheta \nu x)(\widehat{\vartheta\nu}su)(\widehat{\vartheta'}\nu su)
=(\vartheta \nu x)s_\vartheta s_{\vartheta'}-(\vartheta\nu x)s_\nu s_{\vartheta'}
=-2\theta\cdot(\vartheta\nu x)s_\nu s_\vartheta+(\vartheta\nu x)\cdot s_\vartheta^2.
\]
Les formes restantes ont soit des réduisants pour facteurs, soit sont des contractions simples avec des formes décomposées.

B) réduisants $a_\nu^2s_\nu$ et $(\widehat{j\nu}su)s_\nu$.

C) 1. La série de formes $(\theta Hu)^2s_\eta$ : a) réduisant $(aHu)^2(asu)s_\eta$ ; b) double réduction des séries de formes $(\theta gu)^3g_\nu$ et $g_\nu(agu)^3(bgu)^2(abu)$. D'où la réduction des formes supérieures $(\theta_\eta su)^3$, $(H_\vartheta(\theta Hu),s,u)^3$ et $(a_\nu\cdot b_{\nu_1}^2,s,u)^4$ [cette dernière forme remplaçant $(H_\vartheta su)^4$].

2. $((\vartheta\eta x),s,u)^4$: réduisant $(a_\eta su)^4$.

3. $((\vartheta\eta x)^2,s,u)^3$ : réduisant $s_\vartheta^2s_\nu$ issu de $(\widehat{\vartheta\eta}su)^2(Hsu)$. Les formes $(\vartheta\eta x)s_\eta$, $(\vartheta\eta x)^2s_\eta$, $((\vartheta\eta x)^2,s,u)^2$, $\theta_\eta s_\eta$ se décomposent par contraction avec la forme décomposée $a_\eta s_\eta$.

4. Séries de formes $H_\vartheta(\theta Hu)s_\eta$, $\theta_\eta(\theta Hu)s_\eta$, $H_\vartheta(\vartheta\eta x)s_\eta$, $\theta_\eta(\vartheta\eta x)s_\eta$ : réduisants $a_\eta(aHu)s_\eta$ et $a_\eta^2s_\eta$.

La série de formes $(\theta_\eta(\vartheta\eta x),s,u)^2$ : réduisant $a_\eta^2(Hsu)^2$ [à l'exception de $(\theta_\eta^2(\theta Hu),s,u)^2$, qui provient par contraction du terme $a_\eta^2(bsu)(Hsu)$].

5. $(H_\vartheta(\vartheta\eta x),s,u)$, $(H_\vartheta(\vartheta\eta x),s,u)^2$ : double réduction de $a_\eta(aHu)s_\eta(abu)$ et $g_\vartheta(\theta gu)g_\nu$.

La série de formes $(H_\vartheta(\vartheta\eta x),s,u)^3$ : réduisant $((\vartheta \nu x)^2,s,u)^2s_\nu$.

6. $(H_\vartheta(\theta Hu),s,u)$ se décompose par contraction simple avec la forme décomposée $(\theta_\nu(\vartheta \nu x),s,u)$ d'après \S\ 3. 2), c), c'est-à-dire par double réduction de la forme décomposée inférieure $(H_\vartheta su)^2$ :\footnote{Le signe $\equiv$ signifie que des produits de formes de degré inférieur ont été omis.}
\begin{align*}
\begin{aligned}
0&\equiv \theta_\nu(\vartheta \nu\nu_1)(\theta su)+\frac23\theta_\nu^2(\vartheta \nu_1 x)(\theta su)+\theta_\nu(\vartheta \nu x)(\theta su)s_{\nu_1}\\
&\equiv -\frac12H_\vartheta(\theta Hu)(\theta su)+\theta_\eta(\theta su)(Hsu)+\frac12H_\vartheta(\theta su)(Hsu)\\
&\equiv -\frac12H_\vartheta(\theta Hu)(\theta su)+\theta_\eta(\theta su)(Hsu)+\frac12[H_\vartheta]_{\widehat{su}^2}-\frac12\cdot\frac35H_\vartheta(\theta Hu)(\theta su)\\
&\equiv -\frac35H_\vartheta(\theta Hu)(\theta su).
\end{aligned}
\end{align*}
Par double réduction, d'après 1. et 5., on obtient les formules suivantes :
\[
\begin{array}{@{}r@{\quad}l@{}}
&\displaystyle
\begin{aligned}
0={}&\theta_\eta(\theta su)(Hsu)^2+\frac14H_\vartheta(\theta Hu)(\theta su)^2-\frac34\theta_\eta(\theta Hu)(\theta su)(Hsu)\\
&-\frac54\theta_\eta(\theta Hu)^2(\theta su)-(asu)^3\cdot b_\eta(bHu)^2-a_\nu(asu)^3\cdot b_\nu^2,
\end{aligned}\\[0.55em]
(31.)&\displaystyle 0=H_\vartheta(\theta Hu)(\theta su)^3-7\theta_\eta(\theta Hu)(\theta su)^2(Hsu),\\[0.55em]
&\displaystyle 0=a_\nu(asu)^2b_{\nu_1}^{2}(bsu)^2-\theta_\eta(\theta su)^2(Hsu)^2+6\theta_\eta(\theta Hu)(\theta su)^2(Hsu),\\[0.55em]
&\displaystyle 0\equiv (H_\vartheta(\vartheta\eta x),s,u),\qquad
0\equiv H_\vartheta(\widehat{\vartheta\eta}su)(Hsu)-4\theta_\eta^2(Hsu).
\end{array}
\]
\noindent\emph{Conséquences :}
\begin{enumerate}
\item $(\theta Hu)^2s_\eta$, $H_\vartheta(\theta Hu)s_\eta$, $\theta_\eta(\theta Hu)s_\eta$, $H_\vartheta(\vartheta\eta x)s_\eta$, $\theta_\eta(\vartheta\eta x)s_\eta$, $(H_\vartheta(\vartheta\eta x),s,u)$, $(\theta_\eta(\vartheta\eta x),s,u)^2$ sont des réduisants.
\item $s^2$ n'entre pas dans la contraction avec les formes $(5,\lambda)$, etc., d'ordre 6.
\end{enumerate}

\begin{center}\textbf{\S\ 24.\quad Formes du 11e ordre (Système $s_{III}$).}\end{center}
Contraction de $s$ avec les formes du 7e ordre (Système III) (\S\ 12).

\noindent\emph{Formes irréductibles :}
\[
\text{A)}\quad
\begin{array}{c|c|c|c}
\cdot&\cdot&\cdot&((KHu)^2su)^2\\[0.35em]
\cdot&\cdot&(K_\eta su)^2&\cdot\\[0.35em]
\cdot\cdot&\cdot&\cdot&\cdot\\[0.35em]
((k\eta x)^3,s,u)\cdot&\cdot&\cdot&\cdot\\[0.35em]
\cdot&(K_\eta(k\eta x)^3,s,u)&\cdot&\cdot
\end{array}
\]
\[
\begin{array}{c|c}
((KHu)^2,s,u)^3&((KHu)^2,s,u)^4\\[0.45em]
\cdot&(K_\eta(KHu)^2,s,u)^3
\end{array}
\]
\[
\text{B)}\quad ((j\eta x),s,u)^2,\ (H_jsu)\ \big|\ ((j\eta x),s,u)^3,
\qquad \text{C)}\quad (\D_\eta su),
\]
\[
\text{D)}\quad
\begin{array}{c|c|c}
\cdot&((aLu)^2,s,u)^2,\ ((aLu)^3,s,u)&((aLu)^2,s,u)^3\\[0.5em]
(aLu)^2s_l,\ (a_lsu)^2&(a_lsu)^3&(a_l(aLu)^2,s,u)^2.
\end{array}
\]
\noindent\emph{Réductions :}

A) Réduction ou décomposition par contraction simple avec les formes correspondantes dans les symboles $\theta\cdot H$.

$(K_\eta(KHu)^2,s,u)$ et $(K_\eta(KHu)^2,s,u)^2$ : décomposition d'après \S\ 3. 2), a), par contraction avec $K_\nu^2(Ksu)$, $K_\nu^2(Ksu)^2$.

$(K_\eta(KHu),s,u)^4\equiv(a_\nu^2\cdot\theta_{\nu_1},s,u)^4$ : décomposition d'après \S\ 3. 2), b), à partir de la forme décomposée $K_\nu^2(Ksu)^3$.

B) Série de formes $(j\eta x)s_\eta$ : réduisant $(j\nu x)^2s_\nu$.

Forme $H_j(\widehat{j\eta}su)(Hsu)$: réduisant $(j\nu x)(\widehat{j\nu}su)^2$.

Forme $H_j(j\eta x)(Hsu)$ : se décompose en réduisant la forme décomposée $((j\eta x)^2,s,u)^2$ au moyen du réduisant $(j\nu x)^2s_\nu$ (formule (18.)a) :
\[
0=-2(j\nu x)^2s_\nu s_{\nu_1}=[(j\eta x)^2]_{\widehat{su}^2}+H_j(j\eta x)(Hsu).
\]
C) réduisants $\D_\nu s_\nu$ et $\D_\nu(\D su)^2$.

D) Réduction et décomposition par contraction simple avec les formes correspondantes dans les symboles $f\cdot H$.

E) De (30.)c on obtient la réduction de la série de formes décomposée $(a_\sigma su)^2$ :
\begin{align*}
\begin{aligned}
0&=a_\eta(Hsu)^2(as\widehat{\nu x})^2
=a_\eta(Hsu)^2(a\widehat{\nu\eta}u)^2+2a_\eta(a\widehat{\nu\eta}u)(\widehat{\nu\eta}su)(Hsu)^2
=\frac13a_\sigma(asu)^2,\\
0&=a_\eta a_\nu(Hsu)^2(as\widehat{\nu x})(asu)
=a_\eta(a\widehat{\nu\eta}u)^2(Hsu)^2(asu)+a_\eta(a\widehat{\nu\eta}u)(\widehat{\nu\eta}su)(Hsu)^2(asu)
=\frac13a_\sigma(asu)^3.
\end{aligned}
\end{align*}
\noindent\emph{Conséquence :} $s^2$ n'entre pas dans la contraction avec les formes d'ordre 7.

\begin{center}\textbf{\S\ 25.\quad Formes des 12e, 13e, 14e et 15e ordres (Système $s_{III}$).}\end{center}
Contraction de $s$ avec les formes des ordres 8, 9, 10, 11 (Système III) (\S\S\ 13--16).

\noindent\emph{Formes irréductibles :}

\noindent 12e ordre.
\[
\begin{adjustbox}{max width=0.98\linewidth}
$\displaystyle
\begin{array}{c@{\qquad}c@{\qquad}c}
\text{A)}&
\begin{array}{c|c}
((\vartheta^2\eta x)^4,s,u)&\cdot\\[0.35em]
\cdot&\theta_\eta(\vartheta^2\eta x)^3s_\vartheta,
\end{array}
&
\text{B)}\quad ((aHu)H_j,s,u)^2\;\big|\;((aHu)H_j,s,u)^3,
\\[1.0em]
\text{C)}&
\begin{array}{c|c|c}
\cdot&((\theta Lu)^2,s,u)^2,\ ((\theta Lu)^3,s,u)&((\theta Lu)^2,s,u)^3\\[0.35em]
(\theta_lsu)^2&\cdot&(\theta_l(\theta Lu)^2,s,u)^2.
\end{array}
&
\end{array}$
\end{adjustbox}
\]
\noindent\emph{Réductions :} Les réductions provenant de la contraction directe avec des réduisants ou des formes décomposées ne seront plus mentionnées.

B) décomposition de $((aHu)H_j,s,u)$ et $(a_\eta(aHu)H_j,s,u)$ comme transvections sur des déterminants fonctionnels.

décomposition de $(a_\eta(aHu)H_j,s,u)^2$ : double réduction de la forme décomposée $[\theta_\nu\cdot\theta_{\nu_1}^{\prime 3}]_{\widehat{su}^3}$ (\S\ 13. C) :
\begin{align*}
0&\equiv[\theta_{\nu_1}^{\prime 3}\cdot\theta_\nu]_{\widehat{su}^3}
\equiv-\frac34\theta_\nu(\theta su)^2(\theta\theta'u)\theta_{\nu_1}^{\prime 3}\\
&\equiv-\frac98(\widehat{\theta\theta'}\eta\widehat{su})(\theta\theta'u)(\theta Hu)(\theta'Hu)\theta_\eta'(\theta su)
\equiv-\frac38a_\eta(aHu)H_j(asu)^2.
\end{align*}

C) Des formules (31.) résulte la réduction des formes décomposées
\[
 [L_\vartheta(\theta Lu)]_{\widehat{su}^3}\quad\hbox{et}\quad [\theta_l(\theta Lu)]_{\widehat{su}^3},
\]
c'est-à-dire des produits $[\theta_\nu^2\cdot H]_{\widehat{su}^3}$ et $[a_\nu^2\cdot(bHu)^2]_{\widehat{su}^3}$ :
\[
\begin{array}{@{}r@{\quad}l@{}}
&\displaystyle 0\equiv\theta_\nu^2(\theta su)^2(Hsu),\qquad
0\equiv [L_\vartheta(\theta Lu)]_{\widehat{su}^4}-7[\theta_l(\theta Lu)]_{\widehat{su}^4},\\[0.55em]
&\displaystyle 0\equiv a_\nu^2(asu)^2(bHu)^2(bsu)+6\theta_l(\theta Lu)^2(\theta su)(Lsu),\\[0.55em]
(32.)&\displaystyle 0\equiv\theta_\nu^3(\theta su)(Hsu)^2
\quad\hbox{à partir de }0\equiv\theta_l^2(\theta Lu)(\theta su)(Lsu)^2\\[-0.05em]
&\displaystyle\hfill(\hbox{réduisant }a_\eta^2(Hsu)^2).
\end{array}
\]

\noindent 13e ordre.
\[
\begin{array}{rl}
\text{A)}&((KLu)^3,s,u)^2;\ ((KLu)^3,s,u)^3,\\[0.25em]
\text{B)}&(L_jsu)^2,\\[0.25em]
\text{C)}&((aH^2u)^3,s,u)^2.
\end{array}
\]
\noindent 14e ordre.
\[
\begin{array}{rl}
\text{A)}&((aLu)^2L_j,s,u)^2,\\[0.25em]
\text{B)}&((\theta H^2u)^3,s,u)^2.
\end{array}
\]
\noindent 15e ordre.
\[
 ((KH^2u)^4,s,u)^2.
\]
\noindent\emph{Réductions :} décomposition de $(K_l(KLu)^3,s,u)^2$, $(H_\vartheta(\theta H^2u)^3,s,u)$, $((K,H\cdot L,u)^5,s,u)$ d'après \S\ 3. 2), c) ; c'est-à-dire en tenant compte simultanément des formes décomposées $(K_l(KLu)^2,s,u)^3$, $(H_\vartheta(\theta H'u)^2,s,u)^2$, $((K,H\cdot L,u)^4,s,u)^2$.

\noindent\emph{Conséquence :} $s^2$ n'entre pas dans la contraction avec les formes individuelles du système III.

\begin{center}\textbf{\S\ 26.\quad Système $s_{IV}$.}\end{center}
\noindent 1) \emph{Contraction de $s$ avec les systèmes I$\cdot$III.}

\noindent\emph{Réductions :} d'après les réductions des derniers paragraphes $(a_\nu^2s_\nu,\ (aHu)^2(asu)s_\eta$, etc.), les formes $(0,\lambda)$, $(1,\lambda)$, $(2,\lambda)$ ne permettent pas simultanément les contractions I et III ; par conséquent, d'après \S\ 18, les formes $f,\D,N$ n'entrent pas dans des produits pour la contraction vers le système de formes $s_{IV}$.

La forme $j$ n'entre pas dans la contraction puisque, par les mêmes réductions, les contractions contragrédientes à $j$,
\[
(K_\nu(k\nu x)^3,s,u),\qquad ((\vartheta^2\eta x)^4,s,u)\quad\hbox{etc.}
\]
correspondent aux contractions cogrédientes à $j$ :
\[
K_\nu(k\nu x)^2s_k,
\qquad (\vartheta^2\eta x)^3s_\vartheta\quad\hbox{etc.}
\]

\noindent 2) \emph{Contraction de $s$ avec les systèmes II$\cdot$III}, c'est-à-dire de $s$ avec $H\cdot(1,\lambda)$, $H\cdot(2,\lambda)$, $H\cdot(3,\lambda)$.

\noindent\emph{Formes irréductibles :}
\[
\begin{array}{rl}
\text{A)}&(a_\nu\cdot H,s,u)^4,
((aHu)^2\cdot H,s,u)^3,
((aHu)^2\cdot H,s,u)^4,\\[0.25em]
&((aLu)^2\cdot H,s,u)^4,
((aLu)^3\cdot H,s,u)^3,
((aH^2u)^3\cdot H,s,u)^4.\\[0.45em]
\text{B)}&(a_\nu^2\cdot H,s,u)^3,
(a_\eta(aHu)^2\cdot H,s,u)^3,
(a_l(aLu)^2\cdot H,s,u)^4.\\[0.45em]
\text{C)}&(\theta_\nu\cdot H,s,u)^4,
((\theta Hu)^2\cdot H,s,u)^3,
((\theta Hu)^2\cdot H,s,u)^4,\\[0.25em]
&((\theta Lu)^2\cdot H,s,u)^4,
((\theta Lu)^3\cdot H,s,u)^3,
((\theta H^2u)^3\cdot H,s,u)^4,\\[0.45em]
\text{D)}&(\theta_\eta(\theta Hu)^2\cdot H,s,u)^3,
(\theta_l(\theta Lu)^2\cdot H,s,u)^4,\\[0.45em]
\text{E)}&((KLu)^3\cdot H,s,u)^4,
((KH^2u)^4\cdot H,s,u)^4,\\[0.45em]
\text{F)}&((j\nu x)^2\cdot H,s,u)^3,
((j\nu x)^2\cdot H,s,u)^4,
((j\eta x)\cdot H,s,u)^4,\\[0.25em]
&(H_j\cdot H,s,u)^3,
(L_j\cdot H,s,u)^4.
\end{array}
\]
\noindent\emph{Réductions :} $\nu$ n'entre pas dans la contraction, de façon analogue à $j$.

B) et D). $(a_\nu^2\cdot H,s,u)^4$ et $(\theta_\nu^2\cdot H,s,u)^4$ se décomposent par les formules (27.)c) et (29.)c), en tenant compte de $(gHu)^2=-\frac13s\cdot\sigma$ :
\[
(a_\nu^2\cdot H,s,u)^4\equiv\frac13(asu)^4\cdot\sigma,
\qquad
(\theta_\nu^2\cdot H,s,u)^4\equiv-\frac16(\theta su)^4\cdot\sigma;
\]
d'où la décomposition de $(a_\eta(aHu)\cdot H,s,u)^4$, $(\theta_\eta(\theta Hu)\cdot H,s,u)^4$.

$(\theta_\nu^2\cdot H,s,u)^3$, $(\theta_\nu^3\cdot H,s,u)^3$, et par conséquent $(\theta_\eta^2(\theta Hu)\cdot H,s,u)^4$, se décomposent d'après \S\ 25, formes d'ordre 12, C).

\noindent 3) \emph{Contraction de $s$ avec les systèmes III$\cdot$III}, c'est-à-dire de $s$ avec les formes du système III :
\[
(3,\lambda)\cdot(3,\lambda),\quad (3,\lambda)\cdot(2,\lambda),\quad (3,\lambda)\cdot(1,\lambda),\quad
(2,\lambda)\cdot(2,\lambda),\quad (2,\lambda)\cdot(1,\lambda),\quad
(1,\lambda)\cdot(1,\lambda).
\]
\noindent\emph{Formes irréductibles :}
\[
\begin{array}{rl}
\text{A)}&(a_\nu^2\cdot a_\nu^2,s,u)^3,
(a_\nu^2\cdot a_\nu^2,s,u)^4,
(a_\nu^2\cdot a_\eta(aHu)^2,s,u)^3,\\[0.45em]
\text{B)}&(a_\nu\cdot H_j,s,u)^4,
((aHu)^2\cdot H_j,s,u)^3,
((aLu)^2\cdot H_j,s,u)^4,\\[0.25em]
&((aLu)^3\cdot H_j,s,u)^2,
((aH^2u)^3\cdot H_j,s,u)^3.
\end{array}
\]
\noindent\emph{Réductions :}

les produits II$\cdot$III, pour le même ordre dans les symboles $\nu$ et $f$, doivent être considérés comme supérieurs aux produits III$\cdot$III.

Les réductions proviennent en partie de relations entre produits (syzygies), en partie du remplacement des produits par les formes décomposées correspondantes et de la réduction de la contraction de ces formes avec $s$. Les produits contenant des contravariants sont toujours omis, puisqu'ils passent dans des produits.

A) $f$ quadratique, symboles $\nu$ d'ordre arbitraire. D'après \S\ 11, les formules (12.), et par un calcul analogue, on a :
\begin{align*}
1)\quad 0&=a_\nu\cdot b_{\nu_1}+\frac12f\cdot(aHu)^2-\frac14\theta\cdot H+\frac12u_xH_\vartheta-\frac14u_x^2\cdot H_\vartheta^2,\\
2)\quad 0&=a_\nu\cdot b_{\nu_1}^2+f\cdot a_\eta(aHu)^2-\theta_\eta-\frac12H_\vartheta,\\
3)\quad 0&=a_\nu^2\cdot b_{\nu_1}^2-H_\vartheta^2+2\theta_\eta^2.
\end{align*}
Il en résulte que
\begin{align*}
4)\quad 0&=a_\nu^2\cdot b_{\nu_1}^2(j\nu_1x)\quad(\text{d'après }3)),\\
5)\quad L_\vartheta(\theta Lu)&=-a_\nu\cdot b_\eta(bHu)^2+\frac34a_\nu^2\cdot(bHu)^2+\frac34\theta_\nu^2\cdot H\quad(\text{d'après }2)),\\
6)\quad L_\vartheta(\theta Lu)&=-6a_\nu\cdot b_\eta(bHu)^2+2a_\nu^2\cdot(bHu)^2-\frac12\theta_\nu^2\cdot H\quad(\text{d'après }1)),\\
7)\quad 0&=a_\nu\cdot(bHu)^2+f\cdot(aLu)^3+\theta_\nu\cdot H\quad(\text{d'après }1)),
\end{align*}
\[
\left.
\begin{aligned}
8)\quad 0&=a_\nu\cdot(bLu)^3+\frac34(\theta Hu)^2\cdot H\\
9)\quad 0&=(aHu)^2\cdot(bHu)^2+2(\theta Hu)^2\cdot H
\end{aligned}
\right\}\quad\text{(d'après 7)),}
\]
\begin{align*}
10)\quad 0&=a_\eta(aHu)^2\cdot b_\eta(bHu)^2\quad(\text{d'après }5)\text{ et }6)),\\
11)\quad 0&= [a_\nu^2\cdot b_\eta(bHu)]_{\widehat{su}^4}\quad(\text{d'après la formule (30.)b)}).
\end{align*}
Les réductions de
\[
(H_\vartheta su)^4,
\qquad (L_\vartheta(\theta Lu),s,u)^3,
\qquad (L_\vartheta(\theta Lu),s,u)^4
\]
(formules (31.) et (32.)), avec les formules 1) à 11), donnent la réduction de toutes les contractions de $s$ avec des produits quadratiques dans les symboles $f$ et arbitraires en $\nu$.

B) produits de deux des symboles $f,\theta,j$ ; $\nu$ d'ordre arbitraire. D'après les formules (17.), (18.), (20.) et par calcul direct, les relations sont :
\[
\left.
\begin{aligned}
1)\quad 0&=a_\nu\cdot\theta_{\nu_1}+\frac14\theta\cdot(aHu)^2+\frac14f\cdot(\theta Hu)^2,\\
2)\quad 0&=\theta_\nu\cdot\theta_{\nu_1}+\frac12\theta\cdot(\theta Hu)^2
\end{aligned}
\right\}\quad\text{(calcul direct analogue à A 1)).}
\]
Donc :
\[
\left.
\begin{aligned}
3)\quad 0&=3a_\nu\cdot(\theta Hu)^2+\theta\cdot(aLu)^3+2f\cdot(\theta Lu)^3\\
4)\quad 0&=3\theta_\nu\cdot(aHu)^2+f\cdot(\theta Lu)^3+2\theta\cdot(aLu)^3
\end{aligned}
\right\}\quad\text{(d'après 1) et A 6)),}
\]
\begin{align*}
5)\quad 0&=a_\nu\cdot(\theta Lu)^3,\quad 0=\theta_\nu\cdot(aLu)^3,\quad 0=(aHu)^2\cdot(\theta Hu)^2
\quad\text{(d'après 3), 4) et A 8)),}\\
6)\quad 0&=a_\nu\cdot\theta_\nu^2+\theta_\nu\cdot a_\nu^2+f\cdot\theta_\eta(\theta Hu)^2+\theta\cdot a_\eta(aHu)^2\quad(\text{formule (17.)}),\\
7)\quad 0&=\theta_\nu\cdot\theta_\nu^2+\theta\cdot\theta_\eta(\theta Hu)^2+\frac16f\cdot H_j\\
&\qquad [\text{dérivation analogue à celle de }6)\text{ à partir de }\theta_\nu^2(\theta\theta'\widehat{\nu}_1x)],\\
8)\quad 0&=a_\nu\cdot(j\nu x)^2+f\cdot H_j\quad(\text{d'après }6)),\\
9)\quad 0&=(aHu)^2\cdot(j\nu x)^2-2a_\nu\cdot H_j\quad(\text{d'après }8)),\\
10)\quad 0&=\theta_\nu\cdot(j\nu x)^2,\quad 0=\theta\cdot H_j\quad(\text{d'après }7)\text{ et }8)),
\end{align*}
\[
\left.
\begin{aligned}
11)\quad 0&=a_\nu\cdot\theta_\nu^3-\frac34(j\eta x)^2,\\
12)\quad 0&=a_\nu^2\cdot\theta_\nu^2-\frac13(j\eta x)^2-\frac13(\D Hu)^2
\end{aligned}
\right\}\quad\text{(formule (18.)),}
\]
\begin{align*}
13)\quad 0&=a_\nu^2\cdot\theta_\nu^3-\frac12a_\sigma\quad(\text{formule (20.)}),\\
14)\quad 0&=\theta_\nu\cdot\theta_\nu^3,
&0&=a_\eta H_j\quad(\S\ 13. C)).
\end{align*}
Les formules ont été poussées assez loin pour que les produits deviennent polarisables : c'est-à-dire que, par contraction, un seul nouveau produit apparaît, parce que a) la contraction du produit en lui-même devient réductible, et b) les autres produits issus de la contraction deviennent réductibles ou conduisent à des contravariants (cf. \S\ 3. 2), a) et b)).

Les formules 1) à 5) donnent la réduction de toutes les contractions de $s$ avec les produits issus de $a_\nu\cdot\theta_\nu$ par contraction. Les formules 6) à 10) donnent la réduction des produits issus de $a_\nu\cdot\theta_\nu^2$ par contraction. D'après \S\ 24 A), en tenant compte des formules 6) à 10), les contractions de $s$ avec les produits issus de $a_\nu^2\cdot\theta_\nu$ se décomposent.

On a :
\begin{align*}
0&\equiv a_\nu^2(asu)^2\theta_{\nu_1}(\theta su)^2,
&0&\equiv a_\nu^2(asu)\theta_{\nu_1}(\theta su)^3-K_\eta(KHu)^2(Ksu)^3,\\
0&\equiv a_\nu^2(asu)^2\theta_{\nu_1}(\theta su),
&0&\equiv a_\nu^2(asu)\theta_{\nu_1}(\theta su)^2,\\
0&\equiv a_\nu^2(asu)\theta_{\nu_1}(\theta su).
\end{align*}
Les réduisants
\[
 ((j\eta x)^2,s,u)^3
 \ [\hbox{à partir de }(\vartheta\eta\widehat{su})^2(Hsu).\ \S\ 23.\ C\ 3)],
\]
\[
(H_j(j\eta x),s,u)^2\ [\S\ 24,\ B],
\quad ((\D Hu)^2,s,u)^2\ [\S\ 24,\ C],
\]
\[
(a_\sigma su)^2\ [\S\ 24,\ E]
\]
avec les formules 11) à 14), donnent la réduction de tous les produits issus de $a_\nu\cdot\theta_{\nu_1}^3$, $a_\nu^2\cdot\theta_{\nu_1}^2$, $a_\nu^2\cdot\theta_{\nu_1}^3$.

Les calculs précédents se résument dans la \emph{conclusion} suivante :

Le système relativement complet modulo $\modu{(\R,t)}$ se compose des 331 formes et sous-systèmes suivants, reproduits également ci-dessous dans l'ordre tabulaire.

\[
\left.
\begin{array}{lr}
u_x&1\ \text{forme}\\
\text{Système I}&6\ \text{formes}\\
\text{Système II}&5\ \text{formes}\\
\text{Système III}&117\ \text{formes}
\end{array}
\right\}
\quad
\begin{gathered}
\text{Système relativement complet mod }(s,\varrho,t)\\
[\text{cf. Tableau I}].
\end{gathered}
\]
\[
\left.
\begin{array}{lr}
s&1\ \text{forme}\\
\text{Système }s_I&35\ \text{formes}\\
\text{Système }s_{II}&9\ \text{formes}\\
\text{Système }s_{III}&125\ \text{formes}\\
\text{Système }s_{IV}&32\ \text{formes}
\end{array}
\right\}
\quad [\text{cf. Tableau II}].
\]

\clearpage
\begin{landscape}
\thispagestyle{plain}
\begin{center}\bfseries Tableau I (cf. p. 90)\end{center}
\vspace{-0.6em}
\begingroup
\setlength{\tabcolsep}{2pt}
\renewcommand{\arraystretch}{1.35}
\tiny
\begin{center}
\begin{adjustbox}{max width=0.98\linewidth,max totalheight=0.78\textheight,center}
\begin{tabular}{|c|c|c|c|c|c|c|c|c|c|c|c|c|c|c|c|c|c|c|c|c|c|c|}
\hline
\diagbox{$u$}{$x$} & 0 & 1 & 2 & 3 & 4 & 5 & 6 & 7 & 8 & 9 & 10 & 11 & 12 & 13 & 14 & 15 & 16 & 17 & 18 & 19 & 20 & 21 \\ \hline
0 & \mcell{i} &  &  &  & \mcell{f} &  & \mcell{\A} &  & \mcell{K_{\nu}(k\nu x)^3} &  & \mcell{K_{\eta}(k\eta x)^3} &  &  &  & \mcell{(\vartheta^2\eta x)^4} & \mcell{H_k(k\cdot\vartheta,\eta x)^4} &  & \mcell{\theta_l(\vartheta\cdot k,lx)^5} &  &  &  & \mcell{(\vartheta^3lx)^6} \\ \hline
1 &  & \mcell{u_x} &  &  &  & \mcell{\theta_\nu(\vartheta\nu x)^2} &  & \mcell{\theta_\eta(\vartheta\eta x)^2} & \mcell{N} & \mcell{(k\nu x)^3} & \mcell{\theta_\nu(\vartheta^2\nu x)^3} & \mcell{(k\eta x)^3} & \mcell{H_{\vartheta}(\vartheta^2\eta x)^3\\ \theta_\eta(\vartheta^2\eta x)^3} &  & \mcell{\theta_l(\vartheta^2lx)^4} &  & \mcell{(\vartheta k,\eta,x)^4} &  & \mcell{(\vartheta k,l,x)^5} &  &  &  \\ \hline
2 &  &  & \mcell{a_\nu^2} &  & \mcell{\theta\\ a_\eta^2} &  & \mcell{(\vartheta\nu x)^2} & \mcell{K_\nu(k\nu x)^2} & \mcell{(\vartheta\eta x)^2} & \mcell{H_k(k\eta x)^2\\K_\eta(k\eta x)^2} &  & \mcell{(\vartheta^2\nu x)^3\\ K_l(klx)^3} &  & \mcell{(\vartheta^2\eta x)^3} &  & \mcell{(\vartheta^2lx)^4} &  &  &  &  &  &  \\ \hline
3 &  & \mcell{\theta_\nu^3} &  & \mcell{a_\nu} & \mcell{\theta_\nu(\vartheta\nu x)} & \mcell{\A_\nu\\a_\eta} & \mcell{K\\H_{\vartheta}(\vartheta\eta x)\\\theta_\eta(\vartheta\eta x)} & \mcell{\A_\eta} & \mcell{(k\nu x)^2\\\theta_l(\vartheta lx)^2} &  & \mcell{(k\eta x)^2} &  & \mcell{(klx)^3} &  &  &  &  &  &  &  &  &  \\ \hline
4 & \mcell{\nu} &  & \mcell{H\\\theta_\nu^2} & \mcell{(j\nu x)^3\\a_\eta(aHu)} & \mcell{\theta_\eta^2} & \mcell{(\vartheta\nu x)\\a_l^2} &  & \mcell{N_\nu\\(\vartheta\eta x)} &  & \mcell{H_n\\N_\eta\\(\vartheta lx)^2} &  &  &  &  &  &  &  &  &  &  &  &  \\ \hline
5 &  & \mcell{a_\eta(aHu)^2} & \mcell{\theta_\eta^2(\theta Hu)} & \mcell{\theta_\nu} & \mcell{K_\nu^2\\(aHu)} & \mcell{\theta_\eta} & \mcell{K_\eta^2\\(\A Hu)\\a_l} &  & \mcell{\A_l} &  &  &  &  &  &  &  &  &  &  &  &  &  \\ \hline
6 & \mcell{j\\\sigma} &  & \mcell{(j\nu x)^2\\(aHu)^2} & \mcell{L\\a_\nu^2(j\nu x)\\H_{\vartheta}(\theta Hu)\\\theta_\eta(\theta Hu)} &  & \mcell{K_\nu\\\theta_l^2} &  & \mcell{K_\eta} &  &  &  &  &  &  &  &  &  &  &  &  &  &  \\ \hline
7 &  & \mcell{\theta_\eta(\theta Hu)^2} & \mcell{H_j(j\eta x)\\a_l(aLu)^2} &  & \mcell{a_\nu(j\nu x)\\(\theta Hu)} &  & \mcell{\A_\nu(j\nu x)\\\theta_l} & \mcell{K_l^2} &  &  &  &  &  &  &  &  &  &  &  &  &  &  \\ \hline
8 & \mcell{H_j^2\\a_l(aLu)^3} & \mcell{(\nu\sigma x)\\(j\nu x)} & \mcell{(\theta Hu)^2} & \mcell{K_\eta(KHu)^2\\(j\eta x)\\(aLu)^2} &  &  &  &  &  &  &  &  &  &  &  &  &  &  &  &  &  &  \\ \hline
9 &  & \mcell{H_j\\(aLu)^3} & \mcell{a_\eta(aHu)H_j\\L_{\vartheta}(\theta Lu)^2\\\theta_l(\theta Lu)^2} &  & \mcell{(KHu)} &  &  &  &  &  &  &  &  &  &  &  &  &  &  &  &  &  \\ \hline
10 & \mcell{\theta_l(\theta Lu)^3} & \mcell{L_j^2\\(j\sigma x)\\a_\eta(aH^2u)^3} &  & \mcell{H_j(aHu)\\(\theta Lu)^2} &  &  &  &  &  &  &  &  &  &  &  &  &  &  &  &  &  &  \\ \hline
11 &  & \mcell{(\theta Lu)^3} & \mcell{K_l(KLu)^3\\L_j\\(aH^2u)^3} &  &  &  &  &  &  &  &  &  &  &  &  &  &  &  &  &  &  &  \\ \hline
12 & \mcell{(aH^2u)^4} & \mcell{a_l(aLu)^2L_j\\H_{\vartheta}(\theta H^2u)^3\\\theta_\eta(\theta H^2u)^3} &  & \mcell{(KLu)^3} &  &  &  &  &  &  &  &  &  &  &  &  &  &  &  &  &  &  \\ \hline
13 & \mcell{(\nu\sigma j)} &  & \mcell{(aLu)^2L_j\\(\theta H^2u)^3} &  &  &  &  &  &  &  &  &  &  &  &  &  &  &  &  &  &  &  \\ \hline
14 & \mcell{(\theta H^2u)^4} & \mcell{K_\eta(KH^2u)^4\\(a,H\!\cdot L,u)^4} &  &  &  &  &  &  &  &  &  &  &  &  &  &  &  &  &  &  &  &  \\ \hline
15 & \mcell{L_\vartheta(\theta,H\!\cdot L,u)^4} &  & \mcell{(KH^2u)^4} &  &  &  &  &  &  &  &  &  &  &  &  &  &  &  &  &  &  &  \\ \hline
16 &  & \mcell{(\theta,H\!\cdot L,u)^4} &  &  &  &  &  &  &  &  &  &  &  &  &  &  &  &  &  &  &  &  \\ \hline
17 & \mcell{K_\eta(K,H\!\cdot L,u)^5} &  &  &  &  &  &  &  &  &  &  &  &  &  &  &  &  &  &  &  &  &  \\ \hline
18 &  & \mcell{(K,H\!\cdot L,u)^5} &  &  &  &  &  &  &  &  &  &  &  &  &  &  &  &  &  &  &  &  \\ \hline
19 &  &  &  &  &  &  &  &  &  &  &  &  &  &  &  &  &  &  &  &  &  &  \\ \hline
20 &  &  &  &  &  &  &  &  &  &  &  &  &  &  &  &  &  &  &  &  &  &  \\ \hline
21 & \mcell{(KH^3u)^6} &  &  &  &  &  &  &  &  &  &  &  &  &  &  &  &  &  &  &  &  &  \\ \hline
\end{tabular}
\end{adjustbox}
\end{center}
\vspace{-0.2em}
\begin{center}\footnotesize (Les nombres de la ligne horizontale supérieure indiquent le degré en $x$; les nombres de la première colonne verticale indiquent le degré en $u$.)\end{center}
\endgroup
\end{landscape}
\clearpage
\begin{landscape}
\thispagestyle{plain}
\begin{center}\bfseries Tableau II (cf. p. 90)\end{center}
\vspace{-0.6em}
\begingroup
\setlength{\tabcolsep}{2pt}
\renewcommand{\arraystretch}{1.35}
\tiny
\begin{center}
\begin{adjustbox}{max width=0.98\linewidth,max totalheight=0.78\textheight,center}
\begin{tabular}{|c|c|c|c|c|c|c|c|c|c|c|c|c|c|c|c|c|c|}
\hline
\diagbox{$u$}{$x$} & 0 & 1 & 2 & 3 & 4 & 5 & 6 & 7 & 8 & 9 & 10 & 11 & 12 & 13 & 14 & 15 & 16 \\ \hline
0 & \mcell{J} &  &  &  & \mcell{s} &  & \mcell{s_\vartheta^2} &  &  &  &  & \mcell{s_n} & \mcell{(k\nu x)^3s_\nu} & \mcell{(\vartheta^2\nu x)^3s_\nu s_\vartheta\\\theta_\nu(\vartheta^2\nu x)^3s_\vartheta} &  & \mcell{\theta_\eta(\vartheta^2\eta x)^3s_\vartheta} &  \\ \hline
1 &  &  &  &  &  &  & \mcell{(asu)} & \mcell{s_\vartheta} & \mcell{s_k^2\\(\Delta su)} & \mcell{s_n(Nsu)\\(\vartheta\nu x)^2s_\nu} & \mcell{((k\nu x)^3,s,u)s_\nu\\(K_\nu(k\nu x)^3,s,u)} &  & \mcell{(K_\eta(k\eta x)^3,s,u)} &  & \mcell{(\vartheta^2\nu x)^3s_\nu} &  & \mcell{((\vartheta^2\eta x)^4,s,u)} \\ \hline
2 &  &  &  &  & \mcell{(asu)^2} & \mcell{s_\vartheta(\theta su)} & \mcell{(\Delta su)^2\\a_\nu s_\nu} & \mcell{((\vartheta\nu x)^2,s,u)s_\nu\\(\theta_\nu(\vartheta\nu x)^2,s,u)} &  & \mcell{s_k\\(\theta_\eta(\vartheta\eta x)^2,s,u)} &  & \mcell{(k\nu x)^2s_\nu\\((k\nu x)^3,s,u)} &  & \mcell{((k\eta x)^3,s,u)} &  &  &  \\ \hline
3 &  &  & \mcell{(asu)^3} & \mcell{s_\vartheta(\theta su)^2\\s_\nu} & \mcell{a_\nu s_\nu(asu)\\a_\nu^2(asu)} & \mcell{s_\eta} & \mcell{(\theta su)\\(a_\eta^2su)} & \mcell{s_k(Ksu)} & \mcell{(Nsu)^2\\(\vartheta\nu x)s_\nu\\((\vartheta\nu x)^2,s,u)} & \mcell{((k\nu x)^3,s,u)^2} & \mcell{((\vartheta\eta x)^2,s,u)} &  &  &  &  &  &  \\ \hline
4 & \mcell{(asu)^4} & \mcell{s_\vartheta(\theta su)^3} & \mcell{a_\nu^2(asu)^2} & \mcell{(\theta_\nu^3su)} & \mcell{(\theta su)^2} & \mcell{s_k(Ksu)^2\\a_\nu(asu)\\s_\nu(asu)} & \mcell{((\vartheta\nu x)^2,s,u)^2\\\theta_\nu s_\nu} & \mcell{s_\nu(\Delta su)\\(\Delta_\nu su)\\(aHu)s_\eta\\(a_\eta su)} &  & \mcell{(\Delta_\eta su)} &  &  &  &  &  &  &  \\ \hline
5 &  &  & \mcell{(\theta su)^3} & \mcell{s_k(Ksu)^3,s_j,\\s_\vartheta(\theta s^\prime u)^4\\s_\nu(asu)^2\\a_\nu(asu)^2} & \mcell{(Hsu)\\((\vartheta\nu x)^2,s,u)^3\\(\theta_\nu(\vartheta\nu x),s,u)^2\\(\theta_\nu^2su)} & \mcell{((j\nu x)^3,s,u)\\(aHu)^2s_\eta\\(a_\eta su)^2} & \mcell{(Ksu)^2\\s_l\\(\theta_\eta^2su)} &  & \mcell{((k\nu x)^2,s,u)^2\\K_\nu s_\nu} &  &  &  &  &  &  &  &  \\ \hline
6 & \mcell{(\theta su)^4} & \mcell{s_\nu(asu)^3\\a_\nu(asu)^3} & \mcell{(Hsu)^2\\(\theta_\nu(\vartheta\nu x),s,u)^3\\(\theta_\nu^2su)^2} & \mcell{(a_\eta su)^3\\(a_\eta(aHu),s,u)^2\\(a_\eta(aHu)^2,s,u)} & \mcell{(Ksu)^3} & \mcell{s_j(asu)\\((\vartheta\nu x),s,u)^2\\(\theta_\nu su)} & \mcell{((k\nu x)^2,s,u)^3} & \mcell{s_j(\Delta su)\\a_\nu(j\nu x)s_\nu\\((\vartheta\eta x),s,u)^2\\(\theta Hu)s_\eta\\(\theta_\eta su)} &  &  &  &  &  &  &  &  &  \\ \hline
7 & \mcell{(\theta_\nu(\vartheta\nu x),s,u)^4} & \mcell{(a_\eta(aHu),s,u)^3} & \mcell{(Ksu)^4\\(\theta_\eta^2(\theta Hu),s,u)^2\\(a_\nu^2\cdot a_\nu^2,s,u)^3} & \mcell{s_j(asu)^2\\s_k(Ks^2u)^5\\((\vartheta\nu x),s,u)^3\\(\theta_\nu su)^2} & \mcell{(j\nu x)s_\nu\\((j\nu x)^2,s,u)\\((aHu),s,u)^2\\((aHu)^2,s,u)} & \mcell{((\vartheta\eta x),s,u)^3\\(\theta_\eta su)^2} & \mcell{(aLu)^2s_l\\(a_lsu)^2} &  &  &  &  &  &  &  &  &  &  \\ \hline

8 & \mcell{(a_\nu^2\cdot a_\nu^2,s,u)^4} & \mcell{s_j(asu)^3\\s_k(Ks^2u)^6\\((\vartheta\nu x),s,u)^4\\(\theta_\nu su)^3} & \mcell{((j\nu x)^2,s,u)^2\\((aHu),s,u)^3\\((aHu)^2,s,u)^2} & \mcell{(Lsu)^2\\(a_\nu^2(j\nu x),s,u)^2\\(H_\vartheta(\theta Hu),s,u)^2\\(\theta_\eta(\theta Hu),s,u)^2\\(\theta_\eta(\theta Hu)^2,s,u)} & \mcell{(a_lsu)^3} & \mcell{(K_\nu su)^2} &  & \mcell{(K_\eta su)^2} &  &  &  &  &  &  &  &  &  \\ \hline
9 & \mcell{(K_\nu^2su)^4\\((aHu),s,u)^4} & \mcell{(Lsu)^3\\(a_\nu^2(j\nu x),s,u)^3\\(\theta_\eta su)^4\\(\theta_\eta(\theta Hu),s,u)^3} & \mcell{(Ks^2u)^6\\(a_l(aLu)^2,s,u)^2\\(a_\nu^2\cdot H,s,u)^3} & \mcell{(K_\nu su)^3} & \mcell{(a_\nu(j\nu x),s,u)^2\\((\theta Hu),s,u)^2\\((\theta Hu)^2,s,u)} &  & \mcell{(\theta_lsu)^2} &  &  &  &  &  &  &  &  &  &  \\ \hline
10 &  & \mcell{(K_\nu su)^4\\(a_\nu^2\cdot a_\eta(aHu)^2,s,u)^3} & \mcell{(a_\nu(j\nu x),s,u)^3\\((\theta Hu),s,u)^3\\((\theta Hu),s,u)^2} & \mcell{((j\eta x),s,u)^2\\(H_jsu)\\((aLu)^2,s,u)^2\\((aLu)^3,s,u)} &  &  &  &  &  &  &  &  &  &  &  &  &  \\ \hline
11 & \mcell{(a_\nu(j\nu x),s,u)^4\\((\theta Hu),s,u)^4} & \mcell{(K_\eta(KHu)^2,s,u)^3\\((j\eta x),s,u)^3\\((aLu)^2,s,u)^3\\(a_\nu\cdot H,s,u)^4} & \mcell{(H^2su)^3\\(\theta_l(\theta Lu)^2,s,u)^2} &  & \mcell{((KHu)^2,s,u)^2} &  &  &  &  &  &  &  &  &  &  &  &  \\ \hline
12 &  & \mcell{(a_\eta(aHu)^2\cdot H,s,u)^3} & \mcell{((KHu)^2,s,u)^3} & \mcell{((aHu)H_j,s,u)^2\\((\theta Lu)^2,s,u)^2\\((\theta Lu)^3,s,u)} &  &  &  &  &  &  &  &  &  &  &  &  &  \\ \hline
13 & \mcell{((KHu)^2,s,u)^4} & \mcell{((aHu)H_j,s,u)^3\\((\theta Lu)^2,s,u)^3\\(\theta_\nu\cdot H,s,u)^4} & \mcell{(L_jsu)^2\\((aH^2u)^3,s,u)^2\\((j\nu x)^2\cdot H,s,u)^3\\((aHu)^2\cdot H,s,u)^3} &  &  &  &  &  &  &  &  &  &  &  &  &  &  \\ \hline
14 & \mcell{((j\nu x)^2\cdot H,s,u)^4\\((aHu)^2\cdot H,s,u)^4} & \mcell{(\theta_\eta(\theta Hu)^2\cdot H,s,u)^3} &  & \mcell{((KLu)^3,s,u)^2} &  &  &  &  &  &  &  &  &  &  &  &  &  \\ \hline
15 & \mcell{(a_l(aLu)^2\cdot H,s,u)^4} & \mcell{((KLu)^3,s,u)^3} & \mcell{((aLu)^2L_j,s,u)^2\\((\theta H^2u)^3,s,u)^2\\((\theta Hu)^2\cdot H,s,u)^3} &  &  &  &  &  &  &  &  &  &  &  &  &  &  \\ \hline
16 & \mcell{((\theta Hu)^2\cdot H,s,u)^4\\(a_\nu\cdot H_j,s,u)^4} & \mcell{((j\eta x)\cdot H,s,u)^4\\(H_j\cdot H,s,u)^3\\((aLu)^2\cdot H,s,u)^4\\((aLu)^3\cdot H,s,u)^3} &  &  &  &  &  &  &  &  &  &  &  &  &  &  &  \\ \hline
17 & \mcell{(\theta_l(\theta Lu)^2\cdot H,s,u)^4} &  & \mcell{((KH^2u)^4,s,u)^2} &  &  &  &  &  &  &  &  &  &  &  &  &  &  \\ \hline
18 &  & \mcell{((\theta Lu)^2\cdot H,s,u)^4\\((\theta Lu)^3\cdot H,s,u)^3\\((aHu)^2\cdot H_j,s,u)^3} &  &  &  &  &  &  &  &  &  &  &  &  &  &  &  \\ \hline
19 & \mcell{((aH^2u)^3\cdot H,s,u)^4\\(L_j\cdot H,s,u)^4} &  &  &  &  &  &  &  &  &  &  &  &  &  &  &  &  \\ \hline
20 &  & \mcell{((KLu)^3\cdot H,s,u)^4} & \mcell{((aLu)^3\cdot H_j,s,u)^2} &  &  &  &  &  &  &  &  &  &  &  &  &  &  \\ \hline
21 & \mcell{((\theta H^2u)^3\cdot H,s,u)^4\\((aLu)^2\cdot H_j,s,u)^4} &  &  &  &  &  &  &  &  &  &  &  &  &  &  &  &  \\ \hline
22 &  &  &  &  &  &  &  &  &  &  &  &  &  &  &  &  &  \\ \hline
23 & \mcell{((KH^2u)^4\cdot H,s,u)^4} & \mcell{((aH^2u)^3\cdot H,s,u)^3} &  &  &  &  &  &  &  &  &  &  &  &  &  &  &  \\ \hline
\end{tabular}
\end{adjustbox}
\end{center}
\vspace{-0.2em}
\begin{center}\footnotesize (Les nombres de la ligne horizontale supérieure indiquent le degré en $x$; les nombres de la première colonne verticale indiquent le degré en $u$.)\end{center}
\endgroup
\end{landscape}

\clearpage
\setcounter{footnote}{0}

\begin{center}
{\Large\bfseries \NoetherNavAnchor{P03}{3. Sur la théorie des invariants des formes à n variables}3. Sur la théorie des invariants des formes à $n$ variables\footnote{Conférence donnée à la réunion des naturalistes de Salzbourg, 1909.}}\par
\vspace{1.2em}
Par \textsc{Emmy Noether}, à Erlangen.\par
\medskip
Jahresbericht der DMV 19 (1910), p. 101--104
\end{center}

\vspace{1em}
Dans la théorie projective des invariants des formes à $n$ variables, le problème principal est résolu : la finitude du système de formes est démontrée. Une série d'autres questions n'ont cependant pas été traitées, à savoir celles qui concernent les méthodes d'étude des relations entre formes. Je voudrais communiquer brièvement les résultats que j'ai obtenus à propos de telles questions, mais, pour plus de clarté, esquisser d'abord les problèmes dans le domaine ternaire, où ils sont connus.

Je pense d'abord aux deux théorèmes fondamentaux de la méthode symbolique : le théorème selon lequel toutes les constructions invariantes peuvent être représentées symboliquement, et le second, selon lequel on obtient toutes les relations existant entre invariants par application successive d'un petit nombre d'identités symboliques ; ainsi la méthode symbolique donne toutes les relations sans sortir du domaine des invariants.\footnote{Study : \emph{Methoden zur Theorie der ternären Formen}. II. \S\ 6. (Leipzig, Teubner, 1889.)} Sur cette base viennent ensuite les procédés de génération des formes, le procédé symbolique de contraction et les procédés non symboliques de différentiation, la théorie des développements en séries, etc. À ce sujet je voudrais encore mentionner un théorème démontré dans ma dissertation,\footnote{\emph{Über die Bildung des Formensystems der ternären biquadratischen Form}. \S\ 1. (Crelle's Journal Bd. 134. 1908.)} selon lequel toutes les contractions peuvent être engendrées par application successive des deux contractions \og{} cogrédientes \fg{} possibles ; par là j'entends, pour un produit symbolique $s_x t_x u_\sigma u_\tau$, les deux contractions $(stu)$ et $(\sigma\tau x)$. Sur ce théorème on peut fonder la théorie des réduisants et de la réduction des systèmes de formes, de manière analogue au domaine binaire, comme je l'ai montré dans ma dissertation.

Si nous passons maintenant au domaine $n$-aire, le premier théorème de la méthode symbolique lui-même ne vaut déjà que dans une mesure conditionnelle. Comme on le sait, dans les formes quaternaires apparaissent déjà, outre les coordonnées de points et de plans $x$ et $u$, des coordonnées de droites $p_{ik}$, les sous-déterminants formés à partir de deux rangées de coordonnées de points. De façon analogue, dans le domaine $n$-aire, nous avons des rangées de variables $p_\varrho$, dont les éléments individuels sont les sous-déterminants formés à partir de $\varrho$ rangées $x$, et des rangées de symboles $S^\varrho$, dont les éléments se comportent comme des sous-déterminants formés à partir de $\varrho$ rangées $s$ cogrédientes aux $u$. La représentation symbolique par déterminants ne vaut, comme on le sait, que lorsque les rangées de symboles $S^\varrho$ sont résolues en rangées $s$ et que celles-ci sont réparties entre différents déterminants ; seuls ont alors une signification réelle les agrégats de déterminants qui dépendent des $S^\varrho$. Mais on ne peut embrasser d'un coup d'\oe{}il quels agrégats de déterminants réalisent cela ; de ce fait, tous les autres théorèmes mentionnés dans le domaine ternaire disparaissent.

Je voudrais maintenant indiquer un principe simple permettant d'obtenir une représentation symbolique explicite dans les rangées $S^\varrho$, et par là de retrouver aussi les autres théorèmes. C'est une application du théorème connu sur les matrices correspondantes : \og{} À chaque rangée de variables $p_\varrho$ on peut associer de manière biunivoque une autre rangée $q_{n-\varrho}$, dont les éléments individuels sont des sous-déterminants formés à partir de $n-\varrho$ rangées $u$ liées aux $\varrho$ rangées $x$ par les équations $(u\mid x)=0$. \fg{} De même, à chaque rangée de symboles $S^\varrho$ est associée une autre rangée $S_{n-\varrho}$, qui se comporte comme si elle était composée de $n-\varrho$ rangées cogrédientes aux $x$.

Le théorème admet encore une seconde formulation, appropriée aux applications : \og{} À tout produit matriciel $(v^{(1)}\ldots v^{(\varrho)}\mid y^{(1)}\ldots y^{(\varrho)})$,\footnote{C'est-à-dire la somme des produits des déterminants correspondants, formés à partir de la matrice des $\varrho$ rangées $v$ et de celle des $\varrho$ rangées $y$.} où les rangées de la seconde matrice sont contragrédientes à celles de la première, est associé un déterminant ; et, réciproquement, tout déterminant peut être transformé en un produit matriciel ayant un nombre prescrit $\varrho$ de rangées. \fg{} Déterminant et produit matriciel apparaissent ainsi comme entièrement équivalents, et le premier théorème de la méthode symbolique prend la forme suivante : \og{} Toutes les constructions invariantes peuvent être représentées symboliquement par des produits matriciels. \fg{} À partir de deux rangées de symboles $S^\sigma$, $T^\tau$, on peut maintenant, à l'aide de rangées de variables, former très facilement un produit matriciel qui contient explicitement ces rangées de symboles et représente donc une construction invariante de la forme $(S^\sigma\mid p_\sigma)(T^\tau\mid p_\tau)$, à savoir le suivant :
\begin{equation}
 \pair{q_{n-\sigma-\lambda}S^\sigma}{T_{n-\tau}p_{\tau-\lambda}},
 \qquad (\lambda=1,2\ldots\tau\ \text{ou}\ n-\sigma).
\tag{1}
\end{equation}

Plus importante est la réciproque : toutes les constructions invariantes peuvent être représentées explicitement par des produits matriciels, de sorte que le procédé ci-dessus constitue l'extension du procédé de contraction binaire et ternaire. Pour démontrer cette réciproque, il est nécessaire de transférer au domaine $n$-aire le second théorème de la méthode symbolique, à savoir d'établir la totalité des identités indépendantes indécomposables dans lesquelles toutes les rangées $S^\varrho$, $p_\varrho$ sont considérées comme indécomposables. Ces identités résultent des identités connues, qui ne contiennent que des rangées $x$ et $u$,\footnote{E. Pascal dans \emph{Memorie d. R. Acc. d. Lincei}. (4) 5 (1888), p. 375.} en remplaçant le théorème de produit pour les déterminants par un système de théorèmes de produit pour les matrices, et en contractant différents produits matriciels en un seul au moyen précisément de ces théorèmes de produit. On parvient ainsi à deux formules dualistiques, dont une seule sera indiquée :
\begin{equation}
\begin{aligned}
 \pair{S^{\varrho_1}S^{\varrho_2}\ldots S^{\varrho_k}}{x p_{\varrho-1}}
 &=\sum_{i=1}^{k}\eps\,\pair{S_1^{\rho_1}\cdots S_{i-1}^{\rho_{i-1}}S_{i+1}^{\rho_{i+1}}\cdots S_k^{\rho_k}q_{n-\rho+1}}{S_{n-\rho_i}x},\\
 \varrho&=\sum_{i=1}^{k}\varrho_i\leqq n;\qquad \eps=\pm1.
\end{aligned}
\tag{2}
\end{equation}

La seule spécialisation contenue dans ces formules, à savoir l'intervention d'une rangée $x$, ne peut être évitée ; avec des rangées de symboles et de variables tout à fait générales, nous arrivons à une \og{} identité de décomposition \fg{}, une identité dans laquelle une rangée $p_\varrho$ est décomposée en les rangées individuelles $x$. Le cas particulier le plus simple de l'identité de décomposition a déjà été utilisé par Clebsch, dans sa \og{} tâche fondamentale de la théorie des invariants \fg{}, pour déduire les covariants élémentaires dans les développements en séries. À l'aide de ces identités, qui assurent le passage des expressions déterminantales non explicites au produit matriciel, on peut alors effectivement montrer que la représentation matricielle fournit la représentation symbolique explicite souhaitée.

Pour le procédé de contraction, en revanche, vaut un théorème tout à fait analogue à celui du domaine ternaire, sur lequel on peut également fonder de manière analogue les théorèmes de réduction. Si l'on désigne dans (1) le nombre $\lambda$ comme le défaut de la contraction, et si l'on appelle cogrédientes les contractions pour lesquelles $\sigma=\tau$, alors toutes les contractions peuvent être engendrées par application successive des contractions cogrédientes de défaut 1. La preuve de ce théorème repose essentiellement sur le fait que les formes les plus générales peuvent être remplacées par des \og{} formes normales \fg{}, c'est-à-dire par des formes qui s'annulent identiquement sous l'application de toutes les opérations différentielles invariantes. Dans le domaine quaternaire, ce dernier fait a été démontré par M. Mertens ;\footnote{Über invariante Gebilde quaternärer Formen. Wiener Berichte. Bd. 98. 1889.} notre connaissance des opérations différentielles les plus générales -- qui, comme on le sait, proviennent des procédés de contraction par remplacement des rangées de symboles par des symboles différentiels -- nous permet, en utilisant l'identité de décomposition, de transférer presque mot pour mot la preuve de Mertens au domaine $n$-aire.

\clearpage
\setcounter{footnote}{0}

\begin{center}
{\Large\bfseries \NoetherNavAnchor{P04}{4. Sur la théorie des invariants des formes à n variables}4. Sur la théorie des invariants des formes à $n$ variables\footnote{Voir une brève communication à la réunion des naturalistes de Salzbourg. Jahresbericht d. deutsch. Math.-Vereinigung. Bd.\ 19, p.\ 101.}}\par
\vspace{0.5em}
Par Mlle\ \emph{Emmy Noether} à Erlangen.\par
\medskip
\rule{0.22\linewidth}{0.4pt}\par
\vspace{0.7em}
Journal f. d. reine u. angew. Math. 139 (1911), p. 118--154
\end{center}

\section*{Introduction.}
Dans la théorie projective des invariants des formes à $n$ variables, les questions qui se rattachent au problème principal -- la démonstration de la finitude du système de formes -- ont à peine été traitées dès que l'on dépasse les domaines binaire et ternaire. Ce sont les questions relatives à la dépendance mutuelle des formes et aux méthodes permettant de maîtriser cet enchaînement des formes. Ces méthodes reposent essentiellement -- comme cela a été complètement démontré dans les domaines binaire et ternaire -- sur deux théorèmes que M. Study appelle les \og théorèmes fondamentaux de la méthode symbolique \fg{}, à savoir le théorème de la représentabilité symbolique des formes et le théorème des identités symboliques.\footnote{E. Study, \emph{Methoden zur Theorie der ternären Formen} (Leipzig, Teubner, 1889). II. \S\ 6. -- Sur l'apparition de ces théorèmes fondamentaux également dans la théorie des invariants de groupes particuliers de transformations, voir Study, \emph{Das Apollonische Problem} (Math. Ann. Bd. 49 [1897]) et \emph{Zur Differentialgeometrie der analytischen Kurven} (Transactions of the American Math. Society, Bd. 1 [1909]).} Ce dernier théorème affirme que toutes les relations existant entre des formations invariantes entières s'obtiennent par application successive d'un petit nombre d'identités symboliques; autrement dit, la méthode symbolique, et elle seule, donne la connaissance de l'enchaînement des formes sans sortir du domaine des invariants.

M. Gordan a déjà indiqué, dans son programme d'Erlangen,\footnote{P. Gordan, \emph{Über das Formensystem binärer Formen}. S. 46. (Leipzig, Teubner, 1875).} la cause de la difficulté que présente le passage à $n$ variables: elle tient à ce que le premier théorème de la méthode symbolique n'est plus valable dans sa forme générale. Dans le domaine $n$-aire apparaissent en effet des séries de variables d'ordre supérieur $p_\rho$\footnote{Pour la notation, voir \S\ 1.} et, de manière analogue, des séries symboliques d'ordre supérieur $S^\rho$; on arrive à ces séries aussi bien lorsque, comme Clebsch l'a fait à l'origine,\footnote{A. Clebsch, über eine Fundamentalaufgabe der Invariantentheorie. (Abh. der Ges. d. Wiss. zu Göttingen, Bd. XVII, 1872).} on part de formes contenant arbitrairement beaucoup de séries $p_\rho$, que lorsque, suivant F. Deruyts et A. Capelli,\footnote{Voir A. Capelli, \emph{Lezioni sulla teoria delle forme algebriche} (Napoli 1902), \S\ 22--24, et la littérature qui y est citée.} on part de formes ne contenant qu'arbitrairement beaucoup de séries $x$ cogrédientes. Or il n'existe pas de représentation symbolique qui contiendrait explicitement les séries $p_\rho$, $S^\rho$;\footnote{Définition plus précise de la \og représentation explicite \fg{} au \S\ 2.} il faut décomposer les $p_\rho$, $S^\rho$ en les séries individuelles $x$ et les séries $s$ contragrédientes correspondantes, puis distribuer celles-ci entre différents déterminants. Sans doute peut-on, au moyen de conditions différentielles (voir la formule (19.)), vérifier qu'un agrégat donné de déterminants possède la propriété de ne dépendre que des séries $p_\rho$, $S^\rho$; mais on n'obtient pas ainsi de vue d'ensemble sur la totalité des formations invariantes ni sur les procédés qui les engendrent.\footnote{Un perfectionnement du symbolisme, dû à M. Study et très précieux pour le calcul (communiqué par F. Engel, Ber. d. K. Sächs. Ges. d. Wiss., math.-phys. Kl., 1900, p. 126, et, pour le domaine quaternaire, par E. Study, \emph{Geometrie der Dynamen}, p. 135), n'est pas non plus une représentation explicite dans notre sens. Il ne conduirait guère, par exemple, aux procédés de différentiation non symboliques servant à engendrer les formes.}

Nous montrons maintenant, aux \S\S\ 2 et 6, comment, en appliquant le \og théorème des matrices correspondantes \fg{},\footnote{Si l'on excepte l'\emph{Ausdehnungslehre} de Grassmann de 1862, no 112, ce résultat se trouve d'abord chez A. Brill, über zwei Berührungsprobleme, \S\ 2 (Math. Ann. Bd. 4. 1871).} on retrouve le premier théorème fondamental dans toute sa généralité, en remplaçant la représentation par déterminants par une représentation au moyen de \og produits matriciels \fg{}. Au \S\ 2, on montre que tout produit matriciel contenant explicitement les séries $S^\rho,p_\rho$ est une formation invariante de la forme fondamentale; la réciproque donnée au \S\ 6, à savoir que toute formation invariante peut se représenter explicitement par des produits matriciels, ne réussit qu'après le transfert du second théorème fondamental (\S\S\ 3--5).

Ce théorème est connu pour les formes qui ne contiennent que des séries de variables $x$.\footnote{E. Pascal, Giorn. di mat. 26 (1888), p. 33, 102; Rendic. d. Accad. d. Linc. (4) 4 (1888), p. 119; ib. Mem. (4) 5 (1888), p. 375.} Le problème général consistant à établir la totalité des identités indécomposables dans lesquelles toutes les séries $S^\rho,p_\rho$ sont considérées comme indécomposables a été formulé par M. Study;\footnote{\og Methoden \dots \fg{}, p. 204, note 17).} il voit en même temps dans le nombre des identités qui apparaissent une mesure de la difficulté des méthodes dans le domaine $n$-aire. Puisqu'il est possible de condenser la totalité de ces identités en une seule formule (36.), cette difficulté est du moins réduite au minimum. Dans les applications, il apparaîtra cependant souvent certains cas particuliers distingués de (36.), tels que (20.), (21.), caractérisés par le fait qu'ils admettent la représentation explicite sans substitution de nouvelles séries; ou encore la formule (31.), appelée \og identité de décomposition \fg{}, puisqu'une série $p_\rho$ y apparaît comme décomposée, en apparence, en les séries $y$.

Sur ces deux théorèmes fondamentaux, la suite de la théorie se construit alors facilement. Le premier théorème donne directement les procédés d'engendrement des formes, le procédé symbolique de contraction et les procédés de différentiation non symboliques,\footnote{Dans le travail \og Invariante Gebilde von Nullsystemen \fg{} (Sitzungsber. d. Ak. d. Wiss. Wien, Bd. 97, Abt. IIa, 1888), M. Mertens est arrivé, par une autre voie qui ne se généralise pas directement, au procédé de contraction quaternaire. L'invariant alterné de 6 séries $S^2$ qui y apparaît diffère de celui qui naît chez nous, par une seule application de la substitution (11.), d'un agrégat d'invariants pairs. Pour un cas particulier, on trouve aussi le procédé de contraction quaternaire chez M. P. Gordan, \emph{Das simultane System von zwei quadratischen quaternären Formen} (Math. Ann. Bd. 56. 1903).} tandis que l'identité de décomposition élimine, parmi ces procédés, tous ceux qui sont équivalents (\S\ 6). La connaissance des procédés de différentiation donne ensuite le transfert de la notion de \og forme normale \fg{} au domaine $n$-aire et, au moyen d'une méthode de preuve indiquée par M. Mertens\footnote{F. Mertens, Über invariante Gebilde quaternärer Formen. (Sitzber. d. Ak. d. Wiss. Wien. Bd. 98. Abt. IIa. 1889.)} dans le domaine quaternaire, le théorème selon lequel un système de formations invariantes peut être remplacé par un système équivalent de formes normales (\S\ 7). Sous cette dernière hypothèse, j'ai montré dans ma dissertation\footnote{Über die Bildung des Formensystems der ternären biquadratischen Form. (Ce journal, Bd. 134. 1908.)} que le procédé de contraction ternaire peut être engendré par application successive de deux contractions fondamentales. Sur la base de ce théorème et de la théorie du développement en séries, j'ai ensuite développé, en introduisant la notion de \og série de formes \fg{}, les principaux théorèmes de réduction pour les systèmes de formes. De façon tout à fait analogue, dans le domaine $n$-aire, le procédé de contraction peut être engendré à partir de $n-1$ contractions fondamentales (\S\ 8), et les théorèmes de réduction peuvent être établis (\S\ 9).

\paragraph{Note complémentaire.} À l'occasion de la communication de Salzbourg, M. le professeur E. Müller, de Vienne, a attiré mon attention sur le fait que le calcul par produits matriciels qui sous-tend ce travail est, en principe, identique au calcul de l'\emph{Ausdehnungslehre} de Grassmann.\footnote{Hermann Grassmanns gesammelte mathematische und physikalische Werke. Ersten Bandes zweiter Teil: Die Ausdehnungslehre von 1862. In Gemeinschaft mit H. Grassmann d. Jüngeren herausgegeben v. Fr. Engel (Leipzig, Teubner 1896).} En effet, le \og produit combinatoire \fg{} $[S^{\rho_1}S^{\rho_2}\cdots S^{\rho_i}]$ de Grassmann peut être conçu comme une matrice donnée par ces séries (voir \S\ 2): le \og produit progressif \fg{} comme la matrice des séries elles-mêmes, le \og régressif \fg{} comme la matrice des séries associées $S_{n-\rho_1},S_{n-\rho_2},\ldots S_{n-\rho_i}$; quant à la matrice correspondant au \og produit mixte \fg{}, elle naît lorsque plusieurs séries $S$ sont réunies, par la substitution (9.), en de nouvelles séries $P,Q$. Notre \og série associée \fg{} correspond alors au \og complément \fg{} de Grassmann, et notre \og produit matriciel \fg{} à un \og produit mixte de degré d'ordre 0 \fg{}.

Sous cette correspondance, le développement donné dans l'\emph{Ausdehnungslehre} de 1862, no 113, devient identique à une formule duale de notre formule (32.). Dans le cas particulier $\rho=1$, (32.) se transforme en (17.), et la formule duale en (16.). C'est de ces cas particuliers du développement de Grassmann que M. Müller\footnote{E. Müller, Beiträge zur Grassmannschen Ausdehnungslehre. 1. Mitteilung: Einige allgemeine Sätze. (Sitzungsber. d. Ak. d. Wiss. Wien; Bd. 118. Abt. IIa. Juli 1909), Satz 16 und 18.} a récemment tiré les identités (20.), (21.) mentionnées plus haut.

Notre dérivation donne, à la différence de celle de Grassmann--Müller, en même temps la preuve de la complétude des identités, ce qui apparaît comme l'essentiel pour les questions considérées ici; et elle va de (20.), (21.) jusqu'aux identités indécomposables les plus générales (36.), qui ne semblent pas avoir été remarquées jusqu'ici.

\clearpage
\section*{\S\ 1. Notations et définitions. Résumé de résultats connus.}
Nous désignons, comme d'usage, par série de variables $x$ la série des éléments $x_1,x_2,\ldots,x_n$, et de même par série de variables $p_\rho$ ($\rho=1,2,\ldots,n-1$) la série des $\binom n\rho$ éléments $p_{i_1i_2\ldots i_\rho}$, où $p_{i_1i_2\ldots i_\rho}$ désigne les $\binom n\rho$ déterminants formés à partir de la matrice des $\rho$ séries $x^{(1)},x^{(2)},\ldots,x^{(\rho)}$,
\[
\left|\begin{array}{cccc}
 x_{i_1}^{(1)}&x_{i_2}^{(1)}&\cdots&x_{i_\rho}^{(1)}\\
 x_{i_1}^{(2)}&x_{i_2}^{(2)}&\cdots&x_{i_\rho}^{(2)}\\
 \vdots&\vdots&&\vdots\\
 x_{i_1}^{(\rho)}&x_{i_2}^{(\rho)}&\cdots&x_{i_\rho}^{(\rho)}
\end{array}\right|,
\qquad (i_1<i_2<\cdots<i_\rho=1,2,\ldots,n)
\]
Ces déterminants sont donc les éléments de la série $p_\rho$. D'après le théorème des matrices correspondantes,\footnote{Voir par exemple Gordan--Kerschensteiner, \emph{Vorlesungen über Invariantentheorie}, Bd. I, S. 99.} à toute série $p_\rho$ est associée biunivoquement une seconde série $q_{n-\rho}$ par la relation, valable pour chaque élément de la série,
\begin{equation}
 p_{i_1i_2\ldots i_\rho}=(-1)^{\frac{\rho(\rho+1)}{2}+i_1+i_2+\cdots+i_\rho}q_{j_1j_2\ldots j_{n-\rho}},
\tag{1}
\end{equation}
où les $i$ et les $j$ sont rangés dans l'ordre et forment ensemble une permutation complète des nombres de 1 à $n$. La série $q_{n-\rho}$ se compose des $\binom n\rho$ déterminants de degré $n-\rho$, $q_{j_1j_2\ldots j_{n-\rho}}$, formés à partir de la matrice de $n-\rho$ séries $u^{(1)},u^{(2)},\ldots,u^{(n-\rho)}$ contragrédientes à $x$. Ces séries satisfont les $\rho\cdot(n-\rho)$ équations de condition:
\begin{equation}
 (u^{(k)}\mid x^{(l)})=\sum_{\alpha=1}^{n}u_\alpha^{(k)}x_\alpha^{(l)}=0;
 \qquad (k=1,2,\ldots,n-\rho;\ l=1,2,\ldots,\rho).
\tag{2}
\end{equation}
Nous les supposerons normalisées (avec Clebsch\footnote{a. a. O. \S\ 5.}) de telle sorte que le facteur de proportionnalité qui interviendrait autrement dans (1.) soit posé égal à un.

D'après Clebsch\footnote{a. a. O. \S\ 12.} (et de même d'après les recherches ultérieures de F. Deruyts et A. Capelli, voir l'Introduction), les formes les plus générales peuvent être ramenées à celles qui contiennent au plus une fois chacune des $n-1$ séries $p_\rho$, et qui peuvent donc se représenter symboliquement de la manière suivante:
\begin{equation}
 (S^1\mid p_1)^{m_1}(S^2\mid p_2)^{m_2}\cdots(S^{n-1}\mid p_{n-1})^{m_{n-1}},
\tag{3}
\end{equation}
où nous avons posé, par analogie avec (2.),
\[
 (S^\rho\mid p_\rho)=\sum S^{i_1i_2\ldots i_\rho}p_{i_1i_2\ldots i_\rho}.
\]
En raison des relations identiques non linéaires qui existent entre les éléments d'une série $p_\rho$, les séries symboliques $S^\rho$ peuvent être normalisées de façon à satisfaire à ces mêmes identités, c'est-à-dire à se comporter comme les déterminants d'une matrice à $\rho$ séries dont les séries $s^{(1)},s^{(2)},\ldots s^{(\rho)}$ sont contragrédientes à $x$.\footnote{Clebsch, a. a. O. \S\ 4.} On peut donc associer biunivoquement à toute série $S^\rho$ une autre série $S_{n-\rho}$ par la relation:
\begin{equation}
 S_{i_1i_2\ldots i_{n-\rho}}=(-1)^{\frac{(n-\rho)(n-\rho+1)}2+i_1+i_2+\cdots+i_{n-\rho}}S^{j_1j_2\ldots j_\rho},
\tag{4}
\end{equation}
où, de nouveau, les $i$ et les $j$ sont rangés dans l'ordre et épuisent ensemble les nombres de 1 à $n$. Les éléments $S_{i_1i_2\ldots i_{n-\rho}}$ de la série $S_{n-\rho}$ se comportent comme les déterminants formés à partir de $(n-\rho)$ séries $\sigma^{(1)},\sigma^{(2)},\ldots\sigma^{(n-\rho)}$ -- cogrédientes à $x$ --, et les séries $s$ et $\sigma$ sont liées par les $\rho\cdot(n-\rho)$ conditions:
\begin{equation}
 (s^{(k)}\mid \sigma^{(l)})=0.\qquad(k=1,2,\ldots\rho;\ l=1,2,\ldots n-\rho).
\tag{5}
\end{equation}
En vertu des relations (1.) et (4.), chaque facteur de (3.) admet la quadruple représentation symbolique équivalente:
\begin{equation}
 (S^\rho\mid p_\rho)=(S^\rho q_{n-\rho})=(S_{n-\rho}p_\rho)=(-1)^{\rho(n-\rho)}(q_{n-\rho}\mid S_{n-\rho}\footnotemark),
\tag{6}
\end{equation}
\footnotetext{Voir Clebsch, a. a. O. \S\ 5.}
où, comme toujours dans la suite, $(S^\rho q_{n-\rho})$ et $(S_{n-\rho}p_\rho)$ sont des abréviations pour les déterminants $(s^{(1)}\cdots s^{(\rho)}u^{(1)}\cdots u^{(n-\rho)})$ et $(\sigma^{(1)}\cdots\sigma^{(n-\rho)}x^{(1)}\cdots x^{(\rho)})$; le développement de Laplace montre que ceux-ci ne dépendent que des séries $S^\rho,q_{n-\rho}$, respectivement $S_{n-\rho},p_\rho$, ce que la notation ci-dessus doit indiquer. Les expressions $(S^\rho\mid p_\rho)$ et $(q_{n-\rho}\mid S_{n-\rho})$ désignent, d'après ce qui précède, les produits des matrices des séries $s$ et $x$, respectivement des séries $u$ et $\sigma$; par produit matriciel, on entend ici la somme des produits de déterminants correspondants, c'est-à-dire la valeur du déterminant de la matrice produit.

Le nombre caractéristique $\rho\cdot(n-\rho)$ appartenant à une série symbolique (respectivement à une série de variables) sera appelé, d'après Clebsch, le poids de la série;\footnote{a. a. O. \S\ 11.} le nombre $\rho$, qui indique le nombre des séries $s$, sera appelé indice de poids supérieur; le nombre $n-\rho$, qui indique le nombre des séries $\sigma$, indice de poids inférieur. Il résulte des relations (4.) qu'une série symbolique peut intervenir dans une seule et même relation tantôt avec un indice de poids supérieur, tantôt avec un indice de poids inférieur; il en va de même pour l'intervention des séries correspondantes $p_\rho$ et $q_{n-\rho}$. Remarquons encore que nous désignerons généralement: les séries de variables de degré supérieur dont la matrice associée est composée de séries $x$ par $p$, celles dont la matrice est composée de séries $u$ par $q$; les séries symboliques cogrédientes à $x$ par de petites lettres grecques: $\sigma,\tau,\alpha,\beta$; les séries symboliques cogrédientes à $u$ par de petites lettres latines: $s,t,a,b$; toutes les autres séries symboliques par de grandes lettres latines munies d'un indice de poids: $S^\rho$, $T^\tau$, $A^\rho$ ou bien $S_{n-\sigma}$, $T_{n-\tau}$, $A_{n-\rho}$. Selon que l'indice de poids est supérieur ou inférieur, nous écrivons aussi les différents éléments d'une série avec des indices supérieurs ou inférieurs, comme par exemple dans (4.).

\clearpage
\section*{\S\ 2. Représentation par produits matriciels.}
Dans ce qui suit, nous entendrons toujours par \og produit matriciel \fg{} la somme des produits de déterminants correspondants de deux matrices ayant le même nombre de lignes, les lignes de la seconde matrice étant contragrédientes à celles de la première. Les lignes de chaque matrice peuvent alors être: 1. données séparément; 2. réunies, comme en (6.), en une série $S^\rho$; 3. réunies en différentes séries $S^{\rho_1},S^{\rho_2}\ldots$ Le produit matriciel sera noté par le symbole $(\ \mid\ )$ employé dans (6.). La représentation (6.) est une application particulière d'une seconde formulation du théorème des matrices correspondantes: \og À tout produit matriciel à $\rho$ lignes $(v^{(1)}v^{(2)}\cdots v^{(\rho)}\mid y^{(1)}\cdots y^{(\rho)})$ est associé un déterminant; et réciproquement, à tout déterminant on peut associer un produit matriciel ayant un nombre prescrit $\rho$ de lignes ($\rho=1,2,\ldots,n-1$). \fg{}\footnote{Pour $\rho>n$, le produit matriciel s'annule, comme on sait; pour $\rho=n$, il se transforme en produit de deux déterminants.} En effet, il suffit de réunir les lignes $y$ en une série $p_\rho$ (respectivement les lignes $v$ en une série $q_\rho$) et d'effectuer les substitutions données par (1.) pour passer d'un produit matriciel donné au déterminant, et réciproquement. La valeur du déterminant n'est toutefois pas donnée par ses éléments individuels; elle ne résulte qu'après développement de Laplace. Puisque déterminant et produit matriciel se révèlent ainsi équivalents, le théorème de la représentabilité symbolique des formes (premier théorème fondamental) s'exprime aussi de la manière suivante:

\emph{Théorème I.} \og Toutes les formations invariantes peuvent être représentées symboliquement par des produits matriciels; réciproquement, tous les produits matriciels sont des formations invariantes. \fg{}\footnote{Les formations invariantes de formes fondamentales données ne sont naturellement que les agrégats de produits matriciels qui dépendent des séries symboliques $S^\rho\ldots$.}

La représentation par produits matriciels permet de surmonter la difficulté essentielle de la représentation non explicite imposée par la représentation au moyen de déterminants.\footnote{Voir l'Introduction.} Nous entendons par \og représentation explicite \fg{} une représentation dans laquelle n'entrent que les séries $S^\sigma,T^\tau,\ldots$ elles-mêmes, ou des séries $R^\rho$ qui s'en déduisent univoquement, mais dans laquelle les séries partielles individuelles $s,t,\ldots$ n'apparaissent plus. Nous obtiendrons au \S\ 6, pour toutes les formations invariantes qui ne dépendent que des séries symboliques $S^\sigma,T^\tau,\ldots$, une représentation explicite en ces séries symboliques, et par là le transfert au domaine $n$-aire des procédés générateurs: procédé de contraction et procédés de différentiation. La possibilité de la représentation explicite ressort de ce que seuls les produits matriciels les plus simples peuvent correspondre exactement aux déterminants qui apparaissent dans la représentation déterminantielle; tous les autres produits matriciels sont donc des regroupements de plusieurs déterminants en une seule expression. Car, si l'on décompose les séries symboliques et les séries de variables en les séries individuelles $s,u,x$, il doit nécessairement apparaître, d'après le premier théorème fondamental dans sa forme ordinaire, un agrégat de déterminants. La preuve réciproque, à savoir que tous les agrégats de déterminants représentant des formations invariantes de formes fondamentales données peuvent se regrouper en produits matriciels explicites, s'obtient en utilisant le second théorème fondamental (\S\ 6).

Pour obtenir maintenant, à partir de deux séries symboliques $S^\sigma,T^\tau$ et à l'aide de séries de variables, une formation invariante de la forme fondamentale $(S^\sigma\mid p_\sigma)(T^\tau\mid p_\tau)$ contenant explicitement toutes les séries (procédé de contraction), il suffit de former un produit matriciel du troisième type, conformément à la définition donnée au début du paragraphe:
\begin{equation}
 \pair{q_{n-\sigma-\lambda}S^\sigma}{T_{n-\tau}p_{\tau-\lambda}},
 \qquad (\lambda=1,2,\ldots,\tau,n-\sigma\footnotemark).
\tag{7}
\end{equation}
\footnotetext{Le nombre $\lambda$ va toujours jusqu'au plus petit des deux derniers nombres indiqués.}
Développé, (7.) donne:
\begingroup
\setlength{\arraycolsep}{2pt}
\begin{equation}
\begin{aligned}
&\pair{v^{(1)}\cdots v^{(n-\sigma-\lambda)}s^{(1)}\cdots s^{(\sigma)}}
       {\alpha^{(1)}\cdots\alpha^{(n-\tau)}y^{(1)}\cdots y^{(\tau-\lambda)}} \\
&=\sum
\left|\begin{array}{ccc}
 v_{i_1}^{(1)}&\cdots&v_{i_{n-\lambda}}^{(1)}\\
 \vdots&&\vdots\\
 v_{i_1}^{(n-\sigma-\lambda)}&\cdots&v_{i_{n-\lambda}}^{(n-\sigma-\lambda)}\\
 s_{i_1}^{(1)}&\cdots&s_{i_{n-\lambda}}^{(1)}\\
 \vdots&&\vdots\\
 s_{i_1}^{(\sigma)}&\cdots&s_{i_{n-\lambda}}^{(\sigma)}
\end{array}\right|
\left|\begin{array}{ccc}
 \alpha_{i_1}^{(1)}&\cdots&\alpha_{i_{n-\lambda}}^{(1)}\\
 \vdots&&\vdots\\
 \alpha_{i_1}^{(n-\tau)}&\cdots&\alpha_{i_{n-\lambda}}^{(n-\tau)}\\
 y_{i_1}^{(1)}&\cdots&y_{i_{n-\lambda}}^{(1)}\\
 \vdots&&\vdots\\
 y_{i_1}^{(\tau-\lambda)}&\cdots&y_{i_{n-\lambda}}^{(\tau-\lambda)}
\end{array}\right|\\
&=\sum_{i,k,l}\eps\,q_{k_1k_2\ldots k_{n-\sigma-\lambda}}
 S^{k_{n-\sigma-\lambda+1}\ldots k_{n-\lambda}}
 T_{l_1\ldots l_{n-\tau}}
 p_{l_{n-\tau+1}\ldots l_{n-\lambda}},
\end{aligned}
\tag{8}
\end{equation}
\endgroup
c'est-à-dire une expression qui ne dépend que des séries $S,T,q,p$. Ici $\eps=\pm1$\footnote{Cette notation sera conservée dans toute la suite.} et la somme doit être comprise de telle sorte que les indices de chaque série $S\ldots$ soient rangés dans l'ordre, et que, pour chaque combinaison $i_1i_2\ldots i_{n-\lambda}$, les $k$ aussi bien que les $l$ parcourent séparément toutes les permutations encore possibles des nombres $i$. Le nombre $\lambda$ qui intervient dans (7.) sera appelé le défaut de la contraction dans le cas $\sigma\geqq\tau$; la raison de cette dernière restriction apparaîtra au \S\ 6.

Les déterminants des deux matrices qui interviennent dans (7.) peuvent être considérés comme les éléments de nouvelles séries $Q^{n-\lambda},P_{n-\lambda}$, dont les séries associées seront $Q_\lambda,P^\lambda$. D'après (8.), ces séries $Q,P$ peuvent s'exprimer au moyen des séries originelles $S$ et $q$, respectivement $T$ et $p$. Nous l'indiquons par les équations:
\begin{equation}
\begin{aligned}
 q_{n-\sigma-\lambda}S^\sigma&=Q^{n-\lambda}\sim Q_\lambda, &\quad\text{ou encore}\quad Q_\lambda&=[q_{n-\sigma-\lambda}S^\sigma]_\lambda,\\
 T_{n-\tau}p_{\tau-\lambda}&=P_{n-\lambda}\sim P^\lambda, &\quad\text{ou encore}\quad P^\lambda&=[T_{n-\tau}p_{\tau-\lambda}]^\lambda.
\end{aligned}
\tag{9}
\end{equation}
En tenant compte de (4.), on obtient alors, comme en (6.), pour le produit matriciel (7.) la quadruple représentation symbolique équivalente:
\begin{equation}
\pair{q_{n-\sigma-\lambda}S^\sigma}{T_{n-\tau}p_{\tau-\lambda}}
=(q_{n-\sigma-\lambda}S^\sigma P^\lambda)=(Q_\lambda T_{n-\tau}p_{\tau-\lambda})=(-1)^{\lambda(n-\lambda)}\pair{P^\lambda}{Q_\lambda}.
\tag{10}
\end{equation}
On peut effectuer sur (7.) une autre substitution de nouvelles séries en posant:
\begin{equation}
 \pair{q_{n-\sigma-\lambda}S^\sigma}{T_{n-\tau}p_{\tau-\lambda}}
 =\pair{R^{\sigma+\lambda}}{p_{\sigma+\lambda}}(R^{\tau-\lambda}p_{\tau-\lambda}).
\tag{11}
\end{equation}
Ici encore, (8.) montre que les produits des éléments des séries $R$ s'expriment univoquement par les séries $S$ et $T$. Les substitutions (9.) et (11.) seront notamment utilisées lorsque l'on appliquera à (7.) une nouvelle contraction avec d'autres séries.

La représentation matricielle nous met en mesure d'énoncer de façon invariante les relations quadratiques entre les éléments d'une série $p_\rho$ (dont toutes les autres découlent, comme on sait), c'est-à-dire, d'après le \S\ 1, les relations linéaires dans les coefficients de la forme fondamentale auxquelles satisfont les séries $S^\rho$. Ce sont les identités:
\begin{equation}
 \pair{q_{n-\rho-\lambda}S^\rho}{S_{n-\rho}p_{\rho-\lambda}}=0,
 \qquad (\lambda=1,2,\ldots,\rho,n-\rho),
\tag{12}
\end{equation}
où les $p$ et les $q$ doivent être considérés comme des grandeurs arbitraires. Nous verrons au \S\ 5 que les identités dont le défaut $\lambda$ est impair sont des conséquences de celles de défaut plus élevé. Les relations (12.) sont une conséquence directe des relations (5.); car on a:
\[
 \pair{q_{n-\rho-\lambda}S^\rho}{S_{n-\rho}p_{\rho-\lambda}}
 =(v^{(1)}\cdots v^{(n-\rho-\lambda)}s^{(1)}\cdots s^{(\rho)}\mid \sigma^{(1)}\cdots\sigma^{(n-\rho)}y^{(1)}\cdots y^{(\rho-\lambda)}),
\]
et l'application du théorème du produit donne, en vertu de (5.), un déterminant à $(n-\lambda)$ lignes qui s'annule. Les identités (12.) disent que, pour rechercher tous les types de contractions, il suffit de contracter le produit de formes linéaires dans chaque série de variables; ou encore, d'après le \S\ 6, que la forme normalisée $(S^\rho\mid p_\rho)^m$ ne possède aucune formation invariante linéaire dans les coefficients.\footnote{Voir un appendice de Study dans les oeuvres de Grassmann, premier volume, seconde partie, p. 510 (condition pour qu'une grandeur $S^\rho$ soit simple).}

Les conditions (12.) suffisent aussi à caractériser les séries $p_\rho$. En effet, la propriété selon laquelle les éléments individuels de $p_\rho$ sont des déterminants d'une matrice est une propriété invariante; elle doit donc pouvoir s'exprimer invariantement. Or, d'après le théorème VI, les conditions (12.) donnent la restriction invariante la plus grande possible pour les $p_\rho$, à savoir que la forme linéaire construite avec la série $p_\rho$ comme coefficients ne possède absolument aucune formation invariante.

\clearpage
\section*{\S\ 3. Les identités symboliques.}
Nous devons maintenant transférer le second théorème fondamental de la méthode symbolique, c'est-à-dire établir la totalité des identités irréductibles indécomposables dans lesquelles toutes les séries $S^\rho,p_\rho$ sont regardées comme indécomposables. Nous fixons encore la convention suivante: conformément à la signification des séries symboliques et des séries de variables, on considérera comme composées les identités qui naissent lorsque l'on spécialise la série de variables $p_\rho$ en $p_{\rho-1}y$ et que l'on applique aux expressions obtenues l'une des identités du système; en revanche, les identités qui contiennent $k$ séries symboliques devront compter comme indécomposables même lorsqu'elles proviennent, par le procédé ci-dessus, d'identités ayant moins de séries symboliques. Il n'est pas nécessaire que les séries $S^\rho,p_\rho$ entrent explicitement dans les identités; pour les concevoir comme grandeurs indécomposables, il suffit de montrer que les produits matriciels qui apparaissent ne dépendent que de ces séries. Nous montrerons cependant que, dans ce cas, en appliquant en partie la substitution (11.), on peut toujours obtenir une représentation explicite. Les identités générales s'obtiennent, au moyen du principe de la représentation matricielle, à partir des identités particulières données par M. Pascal, qui ne contiennent que des séries $x$ et $u$ et qui constituent une transposition directe de celles données par M. Study pour le domaine ternaire:\footnote{Pour la formulation du problème et la littérature, voir l'Introduction.}
\begin{align}
\sum_{i=1}^n(-1)^{i+1}\pair{a^{(i)}}{x}(a^{(1)}\cdots a^{(i-1)}a^{(i+1)}\cdots a^{(n)}u)&=(-1)^{n+1}\pair{u}{x}(a^{(1)}\cdots a^{(n)}),\tag{13}\\
\sum\pm\pair{a^{(1)}}{x^{(1)}}\pair{a^{(2)}}{x^{(2)}}\cdots\pair{a^{(n)}}{x^{(n)}}&=(a^{(1)}\cdots a^{(n)})(x^{(1)}\cdots x^{(n)}),\tag{14}\\
\sum_{i=1}^n(-1)^{i+1}\pair{u}{\alpha^{(i)}}(\alpha^{(1)}\cdots\alpha^{(i-1)}\alpha^{(i+1)}\cdots \alpha^{n}x)&=(-1)^{n+1}\pair{u}{x}(\alpha^{(1)}\cdots\alpha^{(n)}).\tag{15}
\end{align}
Deux autres identités proviennent de (13.) et de (15.) par les substitutions $(a^{(i)}\mid x)=(a^{(i)}q_{n-1})$, respectivement $(u\mid \alpha^{(i)})=(p_{n-1}\alpha^{(i)})$; selon notre point de vue, elles ne doivent donc pas être considérées comme essentiellement différentes. (14.) représente le théorème du produit pour les déterminants; nous voulons montrer comment on peut remplacer (13.) et (14.) (et, dualement, (15.) et (14.)) par le système des théorèmes du produit formés de manière appropriée pour les matrices. Le théorème du produit, formé une fois pour des matrices à $\rho$ lignes et une fois pour des matrices à $(\rho-1)$ lignes, s'écrit:
\[
\pair{a^{(1)}a^{(2)}\cdots a^{(\rho)}}{xp_{\rho-1}}
=(a^{(1)}a^{(2)}\cdots a^{(\rho)}\mid xy^{(2)}\cdots y^{(\rho)})
=\sum\pm\pair{a^{(1)}}{x}\pair{a^{(2)}}{y^{(2)}}\cdots\pair{a^{(\rho)}}{y^{(\rho)}},
\]
\[
\sum\pm\pair{a^{(2)}}{y^{(2)}}\cdots\pair{a^{(\rho)}}{y^{(\rho)}}
=(a^{(2)}\cdots a^{(\rho)}\mid y^{(2)}\cdots y^{(\rho)})
=(a^{(2)}\cdots a^{(\rho)}\mid p_{\rho-1})=(a^{(2)}\cdots a^{(\rho)}q_{n-\rho-1});
\]
il en résulte le système de théorèmes du produit, où $\rho=2,3,\ldots n$:
\begin{equation}
\begin{aligned}
\pair{a^{(1)}a^{(2)}\cdots a^{(\rho)}}{xp_{\rho-1}}
 &=\sum_{i=1}^{\rho}(-1)^{i+1}\pair{a^{(i)}}{x}(a^{(1)}\cdots a^{(i-1)}a^{(i+1)}\cdots a^{(\rho)}\mid p_{\rho-1})\\
 &=\sum_{i=1}^{\rho}(-1)^{i+1}\pair{a^{(i)}}{x}(a^{(1)}\cdots a^{(i-1)}a^{(i+1)}\cdots a^{(\rho)}q_{n-\rho+1}).\footnotemark
\end{aligned}
\tag{16}
\end{equation}
\footnotetext{Les identités (16.) peuvent aussi être conçues comme les identités (13.) reprises du domaine de $\rho$ variables.}
Pour $\rho=n$, (16.) se transforme en (13.); si nous spécialisons encore $p_{n-1}=y^{(2)}p_{n-2}$, $p_{n-2}=y^{(3)}p_{n-3}$, etc., et si nous appliquons les identités (16.) aux expressions $(a^{(1)}\cdots a^{(i-1)}a^{(i+1)}\cdots a^{(n)}\mid y^{(2)}p_{n-2})$, $(a^{(1)}\cdots a^{(i-1)}a^{(i+1)}\cdots a^{(k-1)}a^{(k+1)}\cdots a^{(n)}\mid y^{(3)}p_{n-3})$, etc., nous passons de (13.) à (14.). L'identité (14.) est donc, selon notre définition, composée à partir des identités (16.); autrement dit, le système (16.) est équivalent aux identités (13.) et (14.). Le système (16.) satisfait à l'une des conditions exigées pour les identités générales: il contient la série $p_{\rho-1}$ comme grandeur indécomposable et, en même temps, sous forme explicite. C'est le seul système qui satisfasse à cette condition et à cette condition seulement. En effet, on ne peut former à partir des séries $a,x,p$ aucun autre déterminant ni produit matriciel; et il ne peut exister entre ceux-ci aucune relation indépendante de (16.), car autrement, par élimination de $(a^{(1)}\cdots a^{(\rho)}\mid xp_{\rho-1})$, il naîtrait une relation entre $\rho$ ($\rho\leqq n$) expressions $(a^{(i)}\mid x)(a^{(1)}\cdots a^{(i-1)}a^{(i+1)}a^{(\rho)}q_{n-\rho+1})$, alors que, d'après (13.), au moins $n+1$ expressions de ce type doivent entrer dans une relation. Des considérations analogues donnent le système équivalent à (15.) et (14.):
\begin{equation}
\begin{aligned}
\pair{u q_{\rho-1}}{\alpha^{(1)}\cdots \alpha^{(\rho)}}
 &=\sum_{i=1}^{\rho}(-1)^{i+1}\pair{u}{\alpha^{(i)}}(q_{\rho-1}\mid \alpha^{(1)}\cdots \alpha^{(i-1)}\alpha^{(i+1)}\cdots \alpha^{(\rho)})\\
 &=\sum_{i=1}^{\rho}(-1)^{i+1}\pair{u}{\alpha^{(i)}}(p_{n-\rho+1}\alpha^{(1)}\cdots \alpha^{(i-1)}\alpha^{(i+1)}\cdots \alpha^{(\rho)}).
\end{aligned}
\tag{17}
\end{equation}
Pour passer de (16.) à des identités entre des séries arbitraires $A^{\rho_1},B^{\rho_2},\ldots,x,p$, nous observons que ces identités doivent devenir identiques à (16.) dès que nous décomposons:
$A^{\rho_1}=a^{(1)}a^{(2)}\cdots a^{(\rho_1)}$, $B^{\rho_2}=b^{(1)}b^{(2)}\cdots b^{(\rho_2)}$, etc.; et qu'elles ne peuvent se regrouper en expressions finales que par les opérations déterminées par (16.), dès que les expressions individuelles sont du type de celles qui interviennent dans (16.). Nous obtenons ainsi, si nous posons encore
\begin{equation}
 \rho=\rho_1+\rho_2+\cdots\rho_k\leqq n,
\tag{18}
\end{equation}
\begin{align*}
&\pair{a^{(1)}\cdots a^{(\rho_1)}b^{(1)}\cdots b^{(\rho_2)}\cdots k^{(1)}\cdots k^{(\rho_k)}}{xp_{\rho-1}}
  =\pair{A^{\rho_1}B^{\rho_2}\cdots K^{\rho_k}}{xp_{\rho-1}} \\
&\quad=\sum_{i=1}^{\rho_1}(-1)^{i+1}\pair{a^{(i)}}{x}(a^{(1)}\cdots a^{(i-1)}a^{(i+1)}\cdots a^{(\rho_1)}B^{\rho_2}\cdots K^{\rho_k}q_{n-\rho+1})\\
&\qquad+(-1)^{\rho_1\rho_2}\sum_{i=1}^{\rho_2}(-1)^{i+1}\pair{b^{(i)}}{x}(b^{(1)}\cdots b^{(i-1)}b^{(i+1)}b^{(\rho_2)}A^{\rho_1}\cdots K^{\rho_k}q_{n-\rho+1})\\
&\qquad\vdots\\
&\qquad+(-1)^{(\rho_1+\cdots\rho_{k-1})\rho_k}\sum_{i=1}^{\rho_k}(-1)^{i+1}\pair{k^{(i)}}{x}(k^{(1)}\cdots k^{(i-1)}k^{(i+1)}\cdots k^{(\rho_k)}A^{\rho_1}\cdots (K-1)^{\rho_{k-1}}q_{n-\rho+1}).
\end{align*}
Chaque somme prise séparément (mais non pas déjà une partie de la somme) contient alors effectivement les séries $A^{\rho_1},B^{\rho_2}\ldots$ comme grandeurs indécomposables; car elle satisfait aux conditions nécessaires et suffisantes à cet effet:\footnote{Capelli, a. a. O. \S\ 23, donne des conditions suffisantes: $\bigl(a^{(k)}\mid \frac{\partial}{\partial a^{(i)}}\bigr)=0$, $(k>i,\ i=1,2,\ldots,\rho_1-1)$. On en déduit le plus simplement les conditions nécessaires et suffisantes par application du procédé des crochets de Poisson d'après Wellstein, Von den Differentialgleichungen der projektiven Invarianten, Math. Ann. 67 (1909) \S\ 3.}
\begin{equation}
 \left(a^{(i+1)}\middle|\frac{\partial}{\partial a^{(i)}}\right)=0;\ \ldots\ ;
 \left(k^{(j+1)}\middle|\frac{\partial}{\partial k^{(j)}}\right)=0.
 \quad (i=1,2,\ldots,\rho_1-1;\ j=1,2,\ldots,\rho_k-1)
\tag{19}
\end{equation}
Nous montrons qu'elle contient aussi toutes les séries de façon explicite. En effet, si nous posons $B^{\rho_2}\cdots K^{\rho_k}q_{n-\rho+1}=q_{n-\rho_1+1}$, alors, d'après (16.) et (10.),
\begin{align*}
\sum(-1)^{i+1}\pair{a^{(i)}}{x}(a^{(1)}\cdots a^{(i-1)}a^{(i+1)}\cdots a^{(\rho_1)}q_{n-\rho_1+1})
&=\pair{A^{\rho_1}}{xp_{\rho_1-1}}\\
&=(A_{n-\rho_1}xp_{\rho_1-1})=(-1)^{(\rho_1+1)(n+1)}\pair{q_{n-\rho_1+1}}{A_{n-\rho_1}x}\\
&=(-1)^{(\rho_1+1)(n+1)}\pair{B^{\rho_2}\cdots K^{\rho_k}q_{n-\rho+1}}{A_{n-\rho_1}x}.
\end{align*}
Si nous calculons de manière analogue les autres sommes et si, ensuite, puisque les séries sont déjà caractérisées comme distinctes par leurs indices de poids, nous posons $B^{\rho_2}=A^{\rho_2}$, $C^{\rho_3}=A^{\rho_3},\ldots$, nous obtenons les identités cherchées:
\begin{equation}
\begin{aligned}
\pair{A^{\rho_1}A^{\rho_2}\cdots A^{\rho_k}}{xp_{\rho-1}}
&=\sum_{i=1}^{k}(-1)^{\delta_i}\pair{A^{\rho_1}\cdots A^{\rho_{i-1}}A^{\rho_{i+1}}\cdots A^{\rho_k}q_{n-\rho+1}}{A_{n-\rho_i}x},\\
\delta_i&=(\rho_1+\rho_2+\cdots+\rho_{i-1})\rho_i+(\rho_i+1)(n+1),
\end{aligned}
\tag{20}
\end{equation}
et, à partir de (17.), les identités duales:
\begin{equation}
\begin{aligned}
\pair{u q_{\sigma-1}}{A_{\sigma_1}A_{\sigma_2}\cdots A_{\sigma_k}}
&=\sum_{i=1}^{k}(-1)^{\eps_i}\pair{u A^{n-\sigma_i}}{p_{n-\sigma+1}A_{\sigma_1}\cdots A_{\sigma_{i-1}}A_{\sigma_{i+1}}\cdots A_{\sigma_k}},\\
\eps_i&=(\sigma_1+\sigma_2+\cdots+\sigma_{i-1})\sigma_i+(\sigma_i+1)(n+\sigma).
\end{aligned}
\tag{21}
\end{equation}
\emph{Théorème II.} Avec (20.) et (21.) est épuisée la totalité des identités indécomposables entre les séries $A^{\rho_i},x,p$, respectivement $A_{\sigma_i},u,q$, et en même temps la totalité des identités indécomposables qui admettent une représentation explicite sans effectuer la substitution (11.).

La première partie du théorème résulte des considérations qui ont conduit à (20.) et (21.); la seconde partie résulte de l'impossibilité d'une identité entre deux séries symboliques:
\[
\pair{q_\rho A^\sigma}{B_{\rho+\sigma-\tau}p_\tau}
=\eps_1\pair{q_\rho B^{n-\rho-\sigma+\tau}}{A_{n-\sigma}p_\tau}
+\eps_2\pair{q_\rho q_{n-\tau}}{B_{\rho+\sigma-\tau}A_{n-\sigma}},
\]
où $\rho\leqq\tau\leqq\sigma$ et où aucune des séries n'est de poids $1\cdot(n-1)$. (D'autres hypothèses sur $\rho,\tau,\sigma$ se ramènent à celle-ci en permutant les désignations des séries.) Cette impossibilité suit par exemple du développement (8.) dans la spécialisation $q_\rho=lM^{\rho-1}$, $q_{n-\tau}=lN^{n-\tau-1}$, si l'on pose $l_1=1$, $M^{2,3,\ldots\rho}=1$, $A^{\rho+1,\ldots\rho+\sigma}=1$, tous les autres éléments des séries $l,M,A$ étant égaux à zéro, tandis que $N$ et $B$ sont arbitraires. Puisque les identités entre séries arbitraires, lorsque l'on décompose une série $p_\tau$ en $y^{(1)}y^{(2)}\cdots y^{(\tau)}$, se transforment en identités composées à partir de (20.), la seconde partie de l'assertion est démontrée dès que (20.) peut être engendrée à partir d'identités ne comportant que deux séries symboliques. On obtient en effet de
\begin{equation}
\pair{A^{\rho_1}A^{\rho_2}}{xp_{\rho-1}}
=\eps\pair{A^{\rho_1}q_{n-\rho+1}}{A_{n-\rho_2}x}
+\pair{A^{\rho_1}}{xp_{\rho_1-1}},
\quad \text{où }p_{\rho_1-1}\sim A^{\rho_2}q_{n-\rho+1},
\tag{20a}
\end{equation}
par la spécialisation $A^{\rho_1}=B^{\sigma_1}B^{\sigma_2}$, c'est-à-dire par
\[
\pair{B^{\sigma_1}B^{\sigma_2}}{xp_{\rho_1-1}}
=\eps\pair{B^{\sigma_1}q_{n-\rho_1+1}}{B_{n-\sigma_2}x}
+\pair{B^{\sigma_1}}{xp_{\sigma_1-1}},
\quad \text{où } p_{\sigma_1-1}\sim B^{\sigma_2}q_{n-\rho_1+1},
\]
une identité (20.) à trois séries symboliques; puis, par d'autres spécialisations de $B^{\sigma_1}=C^{\tau_1}C^{\tau_2}$, etc., les identités générales (20.). L'inexistence d'une représentation explicite d'une identité entre deux séries symboliques et deux séries de variables arbitraires permet donc de conclure à l'inexistence dans le cas d'un nombre arbitrairement grand de séries symboliques, lorsqu'aucune série $x$ ou $u$ ne figure parmi les séries de variables.

\section*{\S\ 4. Les identités de décomposition.}
Nous considérons maintenant le cas le plus général des identités entre des séries symboliques et des séries de variables arbitraires, dans les expressions finales desquelles, sans effectuer la substitution (11.), une série $p_\tau$ apparaît décomposée en ses séries individuelles $y^{(1)}\cdots y^{(\tau)}$; nous désignons ces identités sous le nom d'\og{} identités de décomposition \fg{}. À la suite de la remarque faite à la fin du paragraphe précédent, nous les déterminons d'abord dans le cas de deux séries symboliques seulement; nous développons donc l'expression $\pair{q_\rho A^\sigma}{B_{\rho+\sigma-\tau}p_\tau}$. Nous posons alors:
\begin{equation}
\begin{aligned}
q_{\rho-\alpha}^{(i_1\ldots i_\alpha)}&\sim p_{n-\rho}\,y^{(i_1)}\cdots y^{(i_\alpha)},\qquad
p_{\tau-\alpha}^{(i_1\ldots i_\alpha)}=y^{(i_{\alpha+1})}\cdots y^{(i_\tau)},\qquad
p_\tau=y^{(1)}\cdots y^{(\tau)}.
\end{aligned}
\tag{22}
\end{equation}
Par rapport aux nombres $\rho,\sigma,\tau$, il faut distinguer quatre cas:
\[
\begin{array}{llll}
1)&\rho+\sigma\leqq n,& \tau\leqq\rho,& \tau\leqq\sigma;\\[2pt]
2)&\rho+\sigma\leqq n,& \tau\geqq\rho,& \tau\leqq\sigma;\\[2pt]
3)&\rho+\sigma\leqq n,& \tau\leqq\rho,& \tau\geqq\sigma;\\[2pt]
4)&\rho+\sigma\leqq n,& \tau\geqq\rho,& \tau\geqq\sigma,
\end{array}
\]
dont nous voulons retenir le premier, afin de caractériser ensuite brièvement les autres. D'après (20.), (10.) et (9.), lorsque nous décomposons la série $y^{(1)}y^{(2)}\cdots y^{(\tau)}B_{\rho+\sigma-\tau}$ en $y^{(1)}$ et $y^{(2)}\cdots y^{(\tau)}B_{\rho+\sigma-\tau}$, on obtient:
\begin{equation}
\begin{aligned}
\pair{q_\rho A^\sigma}{y^{(1)}y^{(2)}\cdots y^{(\tau)}B_{\rho+\sigma-\tau}}
&=\pair{q_{\rho-1}^{(1)}A^\sigma}{y^{(2)}\cdots y^{(\tau)}B_{\rho+\sigma-\tau}}\\
&\quad+(-1)^\rho\pair{q_\rho[A_{n-\sigma}y^{(1)}]^{\sigma-1}}{y^{(2)}\cdots y^{(\tau)}B_{\rho+\sigma-\tau}}.
\end{aligned}
\tag{23}
\end{equation}
Le dernier terme de (23.) a maintenant exactement la forme de l'expression initiale: seul $\sigma$ y est remplacé par $\sigma-1$, et $\tau$ par $\tau-1$; il admet donc de nouveau l'équation (23.). Nous pouvons donc, puisque $\tau\leqq\sigma$, appliquer l'opération (23.) $\tau$ fois et, en tenant compte de (10.), parvenir à la relation:
\begin{equation}
\begin{aligned}
\pair{q_\rho A^\sigma}{p_\tau B_{\rho+\sigma-\tau}}
&=(-1)^{\rho\sigma+(\sigma+\tau)(n+1)}
  \pair{q_\rho B^{n-\rho-\sigma+\tau}}{A_{n-\sigma}p_\tau}\\
&\quad+\sum_{i_1=1}^{\tau}(-1)^{\rho(i_1-1)}
 \pair{q_{\rho-1}^{(i_1)}[A_{n-\sigma}y^{(1)}y^{(2)}\cdots y^{(i_1-1)}]^{\sigma-i_1+1}}
 {y^{(i_1+1)}\cdots y^{(\tau)}B_{\rho+\sigma-\tau}}.
\end{aligned}
\tag{24}
\end{equation}
Chaque terme placé sous le signe de sommation est à son tour du type de l'expression initiale; $\rho$ y est remplacé par $\rho-1$, $\sigma$ par $\sigma-i_1+1$, et $\tau$ par $\tau-i_1$, de sorte que la condition 1) est satisfaite plus fortement encore. Nous pouvons donc, puisque $\tau\leqq\rho$, appliquer l'opération (24.) $\tau$ fois et obtenir l'identité de décomposition cherchée:
\begin{equation}
\begin{aligned}
\pair{q_\rho A^\sigma}{p_\tau B_{\rho+\sigma-\tau}}
&=\eps_{i_0}\pair{q_\rho B^{n-\rho-\sigma+\tau}}{A_{n-\sigma}p_\tau}\\
&\quad+\sum_{i_1=1}^{\tau}\eps_{i_1}
 \pair{q_{\rho-1}^{(i_1)}B^{n-\rho-\sigma+\tau}}{A_{n-\sigma}p_{\tau-1}^{(i_1)}}\\
&\quad+\sum_{i_2=i_1+1}^{\tau}\eps_{i_2}
 \pair{q_{\rho-2}^{(i_1i_2)}B^{n-\rho-\sigma+\tau}}{A_{n-\sigma}p_{\tau-2}^{(i_1i_2)}}\\ &\quad\vdots\\
&\quad+\eps_{i_\tau}\pair{q_{\rho-\tau}^{(i_1i_2\ldots i_\tau)}B^{n-\rho-\sigma+\tau}}{A_{n-\sigma}}\\
&=\sum_{\alpha=0}^{\tau}\sum_{i_\alpha=i_{\alpha-1}+1}^{\tau}\eps_{i_\alpha}
 \pair{q_{\rho-\alpha}^{(i_1\ldots i_\alpha)}B^{n-\rho-\sigma+\tau}}{A_{n-\sigma}p_{\tau-\alpha}^{(i_1\ldots i_\alpha)}}.
\end{aligned}
\tag{25}
\end{equation}
Pour le facteur $\eps$, il résulte de (24.):
\begin{equation}
\begin{aligned}
\eps_{i_\alpha}&=(-1)^{\delta_{i_\alpha}},\qquad
\delta_{i_\alpha}=\sum_{\beta=1}^{\alpha}(\rho-\beta+1)(i_{\beta-1}+i_\beta+1)\\
&\quad +(\rho-\alpha)(\sigma+i_\alpha+\alpha)+(\sigma+\tau+\alpha)(n+1),
\end{aligned}
\tag{26}
\end{equation}
où il faut poser $i_0=0$. Le signe de sommation dans (25.) doit être compris ainsi: à chaque combinaison $i_1i_2\ldots i_{\alpha-1}$, où $i_1<i_2\cdots<i_{\alpha-1}$, on adjoint toutes les valeurs $i_\alpha$ pour lesquelles $i_{\alpha-1}<i_\alpha\leqq\tau$. Les sommes individuelles dans (25.) satisfont, comme on s'en convainc facilement en tenant compte de (26.), aux équations différentielles analogues à (19.):
\begin{equation}
 \left(\frac{\partial}{\partial y^{(i)}}\middle|y^{(i+1)}\right)=0,
 \qquad (i=1,2,\ldots\tau-1),
\tag{27}
\end{equation}
et ne dépendent donc que de la série $p_\tau$. L'identité de décomposition est ainsi une identité indécomposable entre les séries $A,B,q,p$; nous reviendrons au paragraphe suivant sur sa représentation explicite.

Dans le cas 2), l'opération (23.) peut également être appliquée $\tau$ fois, mais l'opération (24.) s'interrompt après la $\rho$-ième fois. Dans l'expression finale de (25.), il faut donc remplacer le premier signe de sommation par $\sum_{\alpha=0}^{\rho}$. Dans les cas 3) et 4), l'opération (23.) ne peut être exécutée que $\sigma$ fois; à la place de (24.) intervient alors la relation:
\begin{equation}
\begin{aligned}
\pair{q_\rho A^\sigma}{p_\tau B_{\rho+\sigma-\tau}}
&=(-1)^{\rho\sigma}
 (q_{\rho-(\tau-\sigma)}^{(\sigma+1\ldots\tau)}B^{n-\rho+\tau-\sigma})(A_{n-\sigma}p_{\tau-(\tau-\sigma)}^{(\sigma+1\ldots\tau)})\\
&\quad+\sum_{i_1=1}^{\sigma}(-1)^{\rho(i_1-1)}
 \pair{q_{\rho-1}^{(i_1)}[A_{n-\sigma}y^{(1)}\cdots y^{(i_1-1)}]^{\sigma-i_1+1}}
 {y^{(i_1+1)}\cdots y^{(\tau)}B_{\rho+\sigma-\tau}};
\end{aligned}
\tag{28}
\end{equation}
la répétition $(\tau-\sigma)$ fois de (28.), où le signe de sommation devient $\sum_{i_\beta=i_{\beta-1}+1}^{\sigma+\beta-1}$ -- comme le montre l'application $(\sigma-i_1+1)$ fois de (23.) aux termes sous le signe de sommation dans (28.) -- donne tous les termes de la somme individuelle correspondant à l'indice $\alpha=\tau-\sigma$ dans (25.), avec le signe correspondant. Mais alors les cas 3) et 4) passent aux cas 1) et 2), car les nombres correspondant à $\sigma,\tau$ deviennent: $\sigma'=\sigma-i_{\tau-\sigma}+\tau-\sigma$, $\tau'=\tau-i_{\tau-\sigma}=\sigma'$; et si l'on poursuit le procédé, on a $\tau'<\sigma'$. Il faut donc, dans l'expression finale de (25.), remplacer le premier signe de sommation par $\sum_{\alpha=\tau-\sigma}^{\tau,\rho}$.

De manière analogue, (21.) donne l'identité de décomposition dualiste, où une série $q_\rho$ est décomposée en les séries individuelles $v^{(1)},\ldots,v^{(\rho)}$. Posons
\begin{equation}
 \dot p_{\tau-\beta}^{(i_1\ldots i_\beta)}\sim v^{(i_1)}\cdots v^{(i_\beta)}q_{n-\tau};\qquad
 \dot q_{\rho-\beta}^{(i_1\ldots i_\beta)}=v^{(i_{\beta+1})}\cdots v^{(i_\rho)};\qquad
 q_\rho=v^{(1)}\cdots v^{(\rho)},
\tag{29}
\end{equation}
il vient alors:
\begin{equation}
 \pair{q_\rho A^\sigma}{p_\tau B_{\rho+\sigma-\tau}}
 =\sum_{\beta=\max(0,\tau-\sigma)}^{\min(\tau,\rho)}
   \sum_{i_\beta=i_{\beta-1}+1}^{\rho}\eps
\pair{\dot q_{\rho-\beta}^{(i_1\ldots i_\beta)}B^{n-\rho-\sigma+\tau}}{A_{n-\sigma}\dot p_{\tau-\beta}^{(i_1\ldots i_\beta)}} ,
\tag{30}
\end{equation}
où l'indice de sommation $\beta$ parcourt les valeurs du plus grand des nombres inférieurs jusqu'au plus petit des nombres supérieurs. Tous les termes d'une somme individuelle dans (25.) et dans (30.) sont les polaires d'une seule forme, obtenues par les procédés polaires
\[
\left(\frac{\partial}{\partial p_{\tau-\alpha}}\middle|p_{\tau-\alpha}^{(i_1\ldots i_\alpha)}\right),\qquad
\left(q_{\rho-\alpha}^{(i_1\ldots i_\alpha)}\middle|\frac{\partial}{\partial q_{\rho-\alpha}}\right).
\]
Pour exprimer cela, nous désignons par $\Pi$ un agrégat de telles opérations polaires, multiplié par des constantes, et nous obtenons alors, à la place de (25.) et (30.), la relation:
\begin{equation}
 \pair{q_\rho A^\sigma}{p_\tau B_{\rho+\sigma-\tau}}
 =\sum_{\alpha=\max(0,\tau-\sigma)}^{\min(\tau,\rho)}
 \Pi\pair{q_{\rho-\alpha}B^{n-\rho-\sigma+\tau}}{A_{n-\sigma}p_{\tau-\alpha}}.
\tag{31}
\end{equation}

\section*{\S\ 5. Les identités de décomposition sous forme explicite.}
Si, dans (25.), pour le cas 4), nous spécialisons $\tau$ en $\sigma+\rho$, il vient:
\begin{equation}
 \pair{q_\rho A^\sigma}{y^{(1)}\cdots y^{(\rho+\sigma)}}
 =\sum_{1\le i_1<\cdots<i_\rho\le \rho+\sigma}
 \eps_{i_\rho}\pair{q_\rho}{y^{(i_1)}\cdots y^{(i_\rho)}}
 \pair{A^\sigma}{y^{(i_{\rho+1})}\cdots y^{(i_{\rho+\sigma})}},
\tag{32}
\end{equation}
où $\eps_{i_\rho}=(-1)^{\delta_{i_\rho}}$ et où $\delta_{i_\rho}=\rho(\rho+1)/2+i_1+\cdots+i_\rho$. C'est une forme du théorème du produit pour les matrices plus générale que (16.) et (17.); elle résulte du développement de Laplace du déterminant de la matrice des produits
\[
 \sum\pm(v^{(1)}y^{(1)})\cdots(v^{(\rho)}y^{(\rho)})(a^{(1)}y^{(\rho+1)})\cdots(a^{(\sigma)}y^{(\rho+\sigma)}).
\]
Le théorème du produit plus général est donc composé à partir des identités (20.), respectivement (21.), et le choix de la forme spéciale (16.), (17.) s'explique ainsi.

La relation (32.) donne maintenant le moyen de représenter explicitement les identités de décomposition. Pour cela, nous considérons les formes qui apparaissent dans (31.) comme des formes fondamentales; c'est-à-dire que nous effectuons la substitution (11.), qui ne change rien à la représentation explicite dans les séries $A,B$, puisque, d'après (8.), les séries nouvellement introduites $R$ peuvent s'exprimer par les séries $A,B$ elles-mêmes et non par leurs séries partielles $a^{(1)}a^{(2)}\ldots$, $b^{(1)}b^{(2)}\ldots$. Soit donc
\begin{equation}
 \pair{q_{\rho-\alpha}B^{n-\rho-\sigma+\tau}}{A_{n-\sigma}p_{\tau-\alpha}}
 =\bigl(R_{\rho-\alpha}p_{n-\rho+\alpha}\bigr)
  \pair{R^{\tau-\alpha}}{p_{\tau-\alpha}},
\tag{33}
\end{equation}
il vient alors, pour la somme individuelle de (25.) correspondant à l'indice $\alpha$:
\begin{equation}
\begin{aligned}
&\sum \eps_{i_\alpha}\bigl(R_{\rho-\alpha}p_{n-\rho}y^{(i_1)}\cdots y^{(i_\alpha)}\bigr)
 \pair{R^{\tau-\alpha}}{y^{(i_{\alpha+1})}\cdots y^{(i_\tau)}}\\
&\quad=\sum \eps_{i_\alpha}\bigl([R_{\rho-\alpha}p_{n-\rho}]^\alpha y^{(i_1)}\cdots y^{(i_\alpha)}\bigr)
 \pair{R^{\tau-\alpha}}{y^{(i_{\alpha+1})}\cdots y^{(i_\tau)}}\\
&\quad=\eps\pair{[R_{\rho-\alpha}p_{n-\rho}]^\alpha R^{\tau-\alpha}}{p_\tau}
 =\eps\pair{R^{\tau-\alpha}q_{n-\tau}}{R_{\rho-\alpha}p_{n-\rho}}.
\end{aligned}
\tag{34}
\end{equation}
On a ainsi effectivement obtenu une expression finale explicite; et la forme symétrique dans laquelle $q_\rho$ et $p_\tau$ interviennent montre que (30.) conduit à la même expression finale, donc indépendante de la décomposition initiale. Une différence entre (25.) et (30.) ne subsiste que tant que les séries $y$ et $v$ ont elles-mêmes une signification autonome, c'est-à-dire tant que l'identité de décomposition devient, selon notre définition, une identité décomposable. Pour parvenir à des identités contenant davantage de séries symboliques, il suffit, d'après la fin du \S\ 3, de décomposer une série $q_\rho$ en $q_{\rho_1}C^{\sigma_1}$, respectivement $p_\tau$ en $p_{\tau_1}C_{\tau-\tau_1}$. Les expressions qui en résultent:
\begin{equation}
 \pair{q_{\rho_1}C^{\sigma_1}}{[R^{\tau-\alpha}q_{n-\tau}]_\alpha R_{\rho_1+\sigma_1-\alpha}},
 \qquad
 \pair{[R_{\rho-\alpha}p_{n-\rho}]^\alpha R^{\tau-\alpha}}{p_{\tau_1}C_{\tau-\tau_1}},
\tag{35}
\end{equation}
sont alors exactement du type de celles qui apparaissent pour deux séries symboliques; elles admettent donc de nouveau, après une substitution correspondant à (33.), une représentation explicite, de sorte qu'une répétition donne la représentation explicite pour un nombre arbitraire de séries. Remarquons encore que la première et la dernière somme, respectivement dans (25.) et (30.), donnent la représentation explicite sans application de (33.), par la seule formule (32.), ou bien se réduisent tout simplement à un seul terme contenant explicitement toutes les séries. La relation qui naît de (32.) par décomposition de $q_\rho$ en $q_{\rho_1}C^{\sigma_1}$, etc., peut être comprise comme le théorème du produit sous sa forme la plus générale.

En résumé, nous montrerons:

\emph{Théorème III:} Avec les identités
\begin{equation}
 \pair{q_\rho A^\sigma}{p_\tau B_{\rho+\sigma-\tau}}
 =\sum_{\alpha=\max(0,\tau-\sigma)}^{\min(\tau,\rho)}
 \eps\pair{R^{\tau-\alpha}q_{n-\tau}}{R_{\rho-\alpha}p_{n-\rho}},
\tag{36}
\end{equation}
où les séries $R$ sont déterminées par (33.), et avec celles qui en résultent par décomposition de $q_\rho,p_\tau$ en $q_{\rho_1}C^{\sigma_1}$, etc., la totalité des identités indécomposables est épuisée.

En effet, les identités qui en naissent sont les plus générales de leur espèce, puisque les séries individuelles ne sont plus soumises à aucune restriction quant au poids ou au nombre, contrairement à ce qui était encore le cas dans (20.), (21.). Mais il n'existe pas non plus d'autres identités entre ces séries; car, lorsqu'on décompose une série $p_\tau$ en $y^{(1)}\cdots y^{(\tau)}$, les identités les plus générales sont composées à partir de (20.), donc, d'après la fin du \S\ 3, à partir de (20a.); et la composition discutée à la suite de (16.) est précisément celle qui a été effectuée au \S\ 4. Le passage à un nombre arbitraire de séries est fourni, comme le montre la discussion de (20a.), par l'application toujours répétée de la même opération aux expressions qui résultent de la décomposition de $q_\rho$ en $q_{\rho_1}C^{\sigma_1}$, etc. Il en résulte aussi que le passage direct à partir de (20.), au lieu de (20a.), ne fournit rien de nouveau; cela se vérifie facilement par le calcul, en spécialisant une fois une série $q_\rho$ dans (25.), et en effectuant une autre fois, par la méthode du \S\ 4, le passage à partir de (20.). Dans les deux cas apparaissent les mêmes sommes, qui, par application du théorème général du produit (voir plus haut), passent aux identités qui découlent de (36.). Les identités (20.), (21.) correspondent aux cas particuliers $\tau=1$, $\rho=1$; il n'apparaît alors que deux sommes individuelles dans chaque cas, donc, conformément à une remarque précédente, une représentation explicite sans la substitution (33.), ce qui concorde avec ce qui précède.

Pour finir, mentionnons brièvement deux cas particuliers de l'identité de décomposition. Posons dans (31.): $\tau=\sigma$, $\rho=n-\sigma$, et désignons la série $p_\tau$ par $p'_\sigma\sim q'_{n-\sigma}$; il vient:
\begin{equation}
\begin{aligned}
\pair{A^\sigma}{p_\sigma}\pair{B^\sigma}{p'_\sigma}
-\pair{A^\sigma}{p'_\sigma}\pair{B^\sigma}{p_\sigma}
=\sum_{\alpha=1}^{\sigma,n-\sigma}
\Pi\pair{q_{n-\sigma-\alpha}B^\sigma}{A_{n-\sigma}p_{\sigma-\alpha}}.
\end{aligned}
\tag{37}
\end{equation}
C'est la formule de Clebsch\footnote{Clebsch, a. a. O., p. 25--32. Clebsch démontre que les termes individuels des seconds membres ne dépendent que des séries $A,B$; la représentation explicite s'obtiendrait par application de (32.) à ses expressions $\Phi_h$.} pour développer une forme à deux séries cogrédientes $p_\sigma,p'_\sigma$ suivant des covariants élémentaires. Les covariants élémentaires appartenant à une forme $\pair{A^\sigma}{p_\sigma}^m\pair{B^\sigma}{p'_\sigma}^n$ sont alors les suivants:
\begin{equation}
 \left\{\sum_{\alpha=1}^{\sigma,n-\sigma}
 \pair{q_{n-\sigma-\alpha}B^\sigma}{A_{n-\sigma}p_{\sigma-\alpha}}
 \right\}^{\rho}
 \pair{A^\sigma}{p_\sigma}^{m-\rho}\pair{B^\sigma}{p_\sigma}^{n-\rho},
 \qquad (\rho=0,1,\ldots m,n).
\tag{38}
\end{equation}
Ils proviennent, pour le dire brièvement, de la forme fondamentale par \og{} contraction cogrédiente de défaut $\alpha$ \fg{} (cf. \S\ 2).

Posons deuxièmement: $\tau=\sigma-\lambda$, $\rho=n-\sigma-\lambda$; $B^\sigma=A^\sigma$. Il vient alors:
\begin{equation}
\begin{aligned}
\pair{q_{n-\sigma-\lambda}A^\sigma}{A_{n-\sigma}p_{\sigma-\lambda}}
&=(-1)^\lambda\pair{q_{n-\sigma-\lambda}A^\sigma}{A_{n-\sigma}p_{\sigma-\lambda}}\\
&\quad+(-1)^{(n-\sigma)(\sigma-\lambda)}
 \sum_{\alpha=1}^{\sigma-\lambda,n-\sigma-\lambda}
 \Pi\pair{q_{n-\sigma-\lambda-\alpha}A^\sigma}{A_{n-\sigma}p_{\sigma-\lambda}}.
\end{aligned}
\tag{39}
\end{equation}
Ainsi, comme on l'a remarqué au \S\ 2, les conditions (12.) qui caractérisent les séries symboliques sont, pour un défaut impair $\lambda$, une conséquence des conditions de défaut plus élevé. Pour une forme linéaire $\pair{A^\sigma}{p_\sigma}\equiv\pair{B^\sigma}{p_\sigma}$, il résulte de (39.) que ses formations invariantes irréductibles qui sont quadratiques dans les coefficients se réduisent à des formations de défaut pair; cela concorde avec le fait connu qu'une forme linéaire, dans les domaines binaire et ternaire, ne possède aucune formation invariante.

Remarquons que les identités générales ont été originairement déduites sous l'hypothèse que les séries individuelles satisfont aux conditions (12.). Toutefois, comme ces séries individuelles n'entrent que linéairement dans les identités, ces dernières valent aussi pour les séries plus générales qui ne satisfont pas aux conditions (12.).

\section*{\S\ 6. Représentation explicite des formations invariantes.\\Procédés de contraction et de différentiation.}
Les résultats des paragraphes précédents nous permettent de mener à son terme le théorème, énoncé au \S\ 2, de la représentabilité symbolique (premier théorème fondamental), en y ajoutant la représentabilité explicite; nous montrons en effet:

\emph{Théorème IV:} \og{} Toutes les formations invariantes de formes fondamentales arbitraires peuvent être représentées explicitement par des produits matriciels; c'est-à-dire que, dans les expressions finales, n'entrent, outre les séries de variables $p_\rho$, que les séries $S^\sigma,\ldots T^\tau$ des formes initiales ou les séries $P,Q,R$ qui en sont dérivées par la substitution (9.) ou (11.). \fg{}

Pour la preuve, supposons qu'une formation invariante dépende, outre les séries de variables, des séries $S^\sigma,\ldots T^\tau$; et que ces séries soient suffisamment décomposées en séries individuelles pour permettre une représentation par déterminants:
\[
S^\sigma=S^{\sigma_1}\cdots S^{\sigma_i},\qquad
\ldots,\qquad
T^\tau=T^{\tau_1}\cdots T^{\tau_k}.
\]
Ces séries sont rangées dans un ordre arbitraire en soi, mais fixé une fois pour toutes, $S^\sigma,\ldots T^\tau$. Nous choisissons alors un tel agrégat de déterminants qui dépende de la première série totale $S^\sigma=S^{\sigma_1}\cdots S^{\sigma_i}$, et non seulement de certaines séries partielles. Or cette dépendance est, d'après les \S\S\ 3--5, une conséquence des identités générales. Si nous décomposons maintenant une série de variables $p_\rho$ en les séries $y,v$, ou l'une des séries ultérieures $T$ en les séries individuelles $t$, respectivement $\tau$, les identités les plus générales sont composées à partir de (20.) et (21.). Mais ces dernières identités conduisent toujours à une expression contenant explicitement la série $S^\sigma=S^{\sigma_1}\cdots S^{\sigma_i}$, à savoir le membre de gauche de (20.), (21.); on s'en convainc en effet par (19.), car seule la somme complète du membre de droite -- et non une partie de cette somme correspondant à un terme de (20a.) -- possède la propriété de dépendre uniquement de $S^\sigma$. Si l'on poursuit le procédé, toutes les séries qui sont une fois entrées explicitement demeurent aussi explicitement présentes; l'application des identités peut seulement exiger, le cas échéant, la réunion de plusieurs séries en une seule, c'est-à-dire la substitution (9.). S'il ne reste finalement qu'une seule série $T^\tau$ à mettre sous forme explicite, deux cas sont possibles. Il peut encore apparaître une série de variables libre, c'est-à-dire non liée à d'autres séries par la substitution (9.); sa décomposition en séries partielles $y$ ou $v$ permet d'achever le procédé; il apparaît alors, comme dans (31.), des polaires d'une ou plusieurs formes contenant explicitement toutes les séries. S'il n'apparaît plus aucune série de variables libre, nous reconnaissons, d'après la manière dont elles sont nées, dans la totalité des expressions qui contiennent les séries individuelles $t$ ou $\tau$ tout en dépendant de la série $T^\tau$, les termes d'une identité de décomposition; or celles-ci admettent toujours, d'après le \S\ 5 et par application de la substitution (11.), la représentation explicite dans toutes les séries. Le théorème ci-dessus est ainsi démontré en toute généralité. Remarquons encore que, lorsque seules deux séries symboliques apparaissent, la substitution (11.) n'intervient jamais; en effet, à cause de $\sigma,n-\sigma,\tau,n-\tau<n$, il ne peut n'entrer aucune série de variables que si les deux séries symboliques sont réunies dans un seul déterminant, donc apparaissent explicitement; mais dès que des séries de variables entrent en jeu, (34.) et (33.) montrent que la substitution (11.) devient superflue.

Il résulte de ce qui vient d'être dit que, pour engendrer toutes les formations invariantes, il suffit de réunir les séries symboliques, en leur adjoignant des séries de variables, en tous les produits matriciels possibles. Le produit matriciel le plus général est alors de la forme:
\begin{equation}
 \pair{P^{\mu_1}\cdots P^{\mu_i}}{Q_{\nu_1}\cdots Q_{\nu_k}},\qquad \sum\mu=\sum\nu,
\tag{40}
\end{equation}
où les $P$ et les $Q$ peuvent être des séries des formes initiales, ou bien être issues des séries initiales par une suite de substitutions (9.) ou (11.), ou enfin être des séries de variables. Mais les expressions (40.) peuvent être engendrées par réunion successive, deux à deux, de séries symboliques en un produit matriciel, à l'aide de séries de variables; car on obtient de
\[
 \pair{q_{n-\mu-\lambda}P^\mu}{Q_{\nu_1}p_{n-\nu_1-\lambda}}=\bigl(R_\lambda Q_{\nu_1}p_{n-\nu_1-\lambda}\bigr)
\]
par réunion avec la série $Q_{\nu_2}$ une expression
\[
 \pair{[R_\lambda Q_{\nu_1}]^{n-\lambda-\nu_1}}{Q_{\nu_2}p_{n-\nu_1-\nu_2-\lambda}}
 =\pair{q_{n-\mu-\lambda}P^\mu}{Q_{\nu_1}Q_{\nu_2}p_{n-\nu_1-\nu_2-\lambda}},
\]
donc, en poursuivant le procédé, une expression (40.). Si, en outre, $P$ et $Q$ ne sont pas les séries des formes initiales, (9.) et (11.), combinés avec ce qui vient d'être montré, prouvent qu'ils sont eux aussi nés d'une suite de réunions de deux séries symboliques à la fois. Il en résulte:

\emph{Théorème V:} Le procédé qui engendre toutes les formations invariantes, le procédé de contraction, consiste à réunir successivement, deux à deux, des séries symboliques en un produit matriciel, à l'aide de séries de variables.

À partir de deux séries symboliques $S^\sigma,T^\tau$, où $\sigma\geqq\tau$, on peut maintenant former les produits matriciels suivants:
\begin{align}
&\pair{q_{n-\sigma-\lambda}S^\sigma}{T_{n-\tau}p_{\tau-\lambda}}; &&(\lambda=1,2,\ldots\tau,n-\sigma)\tag{41}\\
&\pair{q_{n-\tau-\mu}T^\tau}{S_{n-\sigma}p_{\sigma-\mu}}; &&(\mu=\sigma-\tau+\lambda)\tag{42}\\
&\pair{q_{n-\tau-\nu}T^\tau}{S_{n-\sigma}p_{\sigma-\nu}}. &&(\nu=1,2,\ldots\sigma-\tau)
\tag{43}
\end{align}
Mais (31.) donne pour (42.) et (43.) les valeurs:
\begin{align}
&\pair{q_{n-\sigma-\lambda}S^\sigma}{T_{n-\tau}p_{\tau-\lambda}}
 +\sum_{\alpha}\Pi\pair{q_{n-\sigma-\lambda-\alpha}S^\sigma}{T_{n-\tau}p_{\tau-\lambda-\alpha}};\tag{42a}\\
&\Pi(q_{n-\sigma}S^\sigma)(T_{n-\tau}p_\tau)
 +\sum_{\beta}\Pi\pair{q_{n-\sigma-\beta}S^\sigma}{T_{n-\tau}p_{\tau-\beta}}.
\tag{43a}
\end{align}
Les formes (42.) peuvent donc s'exprimer linéairement par (41.) et par des polaires de (41.), tandis que les formes (43.) s'expriment par des polaires de la forme fondamentale et des formes (41.). De même, (31.) donne la dépendance linéaire des formes (41.) par rapport à (42.) et aux polaires de (42.). D'après la définition de Clebsch,\footnote{Clebsch, a. a. O. \S\ 7.} les formes (42.) sont donc équivalentes aux formes (41.), tandis que (43.) ramène à la forme fondamentale. Le procédé engendrant est donc donné à droits égaux par (41.) ou par (42.); nous voulons définir comme procédé de contraction le procédé (41.), plus simple relativement au nombre $\lambda$. Le nombre $\lambda$, le défaut de la contraction (cf. \S\ 2), donne le nombre des séries manquantes pour le produit déterminantal, et cela dans le produit matriciel qui, pour une contraction donnée, contient le nombre maximal de séries. Puisque $\lambda$ ne court que jusqu'au plus petit des deux nombres $\tau,n-\sigma$, où $\sigma\geqq\tau$, on a:
\begin{equation}
 \lambda\leqq\frac n2,\ n\text{ pair};\qquad
 \lambda\leqq\frac{n-1}{2},\ n\text{ impair}.
\tag{44}
\end{equation}
La réduction de (42.) et (43.) à (41.) demeure valable même lorsque, par une contraction ultérieure, les séries $p,q$ sont devenues des séries symboliques; car, d'après (36.), il vient:
\[
 \pair{A^{n-\tau-\varkappa}T^\tau}{S_{n-\sigma}B_{\sigma-\varkappa}}
 =\sum_{\alpha=0,\varkappa-\sigma+\tau}^{\tau,n-\sigma}
 \eps\pair{R^{\tau-\alpha}B^{n-\sigma+\varkappa}}{A_{\tau+\varkappa}R_{n-\sigma-\alpha}},
\]
où les séries $R$ sont issues de $S$ et $T$ d'après (33.). On obtient donc ces formes ne contenant que des séries symboliques par contraction de $\pair{q_\tau A^{n-\tau-\varkappa}}{B_{\sigma-\varkappa}p_{n-\sigma}}$, respectivement de $\pair{q_{\tau-\alpha}B^{n-\sigma+\varkappa}}{A_{\tau+\varkappa}p_{n-\sigma-\alpha}}$ -- selon que $\sigma-\tau\gtreqless 2\varkappa$ -- avec la forme fondamentale et les formes (41.).

Du procédé symbolique de contraction, on tire de la manière connue les procédés invariants de différentiation, pourvu que l'on remplace seulement les séries symboliques $S^\sigma,T_{n-\tau}$ par les séries de symboles de différentiation $\frac{\partial}{\partial p_\sigma}$, $\frac{\partial}{\partial q_{n-\tau}}$; comme on le sait, à chaque produit $\frac{\partial}{\partial p_{i_1\ldots i_\sigma}}\cdot\frac{\partial}{\partial q_{j_1\ldots j_{n-\tau}}}$ revient alors la signification réelle $\frac{\partial^2}{\partial p_{i_1\ldots i_\sigma}\partial q_{j_1\ldots j_{n-\tau}}}$. Puisque les séries $\frac{\partial}{\partial p_\sigma}$, $\frac{\partial}{\partial q_{n-\tau}}$ sont cogrédientes aux séries $S^\sigma,T_{n-\tau}$, elles satisfont aux mêmes identités que celles-ci, donc spécialement aussi à celles qui existent entre (42.) et (42a.), (43.) et (43a.). En résumé, nous obtenons:

\emph{Théorème VI:} Toutes les formations invariantes de formes fondamentales arbitraires naissent de celles-ci par application répétée du procédé symbolique de contraction (41.) ou aussi des procédés invariants de différentiation:
\begin{equation}
 \left(q_{n-\sigma-\lambda}\frac{\partial}{\partial p_\sigma}\middle|\frac{\partial}{\partial q_{n-\tau}}p_{\tau-\lambda}\right),
 \qquad(\sigma\geqq\tau,\ \lambda=1,2,\ldots\sigma,n-\tau).
\tag{45}
\end{equation}
Il faut encore remarquer qu'à chaque contraction (41.) le poids total de la forme, c'est-à-dire la somme des poids des séries de variables individuelles (cf. \S\ 1), diminue du nombre positif $2(\lambda^2+\lambda(\sigma-\tau))$. Il en résulte, indépendamment de la démonstration générale de finitude, la finitude de la série de formes, c'est-à-dire de la totalité des formations invariantes d'une forme fondamentale donnée qui sont linéaires dans les coefficients.\footnote{Cf. Clebsch, a. a. O. \S\S\ 11 et 12. Finitude du système des covariants élémentaires.}

Remarquons encore que le théorème VI vaut aussi pour les formes qui ne satisfont pas aux conditions (12.). En effet, chaque facteur de (3.), $\pair{S^\rho}{p_\rho}^{m_\rho}$, peut être remplacé par un produit de $m_\rho$ facteurs linéaires $\pair{A^\rho}{p_\rho}\pair{B^\rho}{p_\rho}\cdots\pair{M^\rho}{p_\rho}$; les formations invariantes passent ainsi à celles d'un système de formes linéaires, qui n'ont besoin d'être identifiées entre elles qu'après coup. Pour chaque forme linéaire individuelle, cependant, les conditions (12.), qui, lors de l'introduction de symboles différentiels, ne contiennent que des quotients différentiels du second ordre, sont toujours satisfaites.

\section*{\S\ 7. Développement des formes générales suivant des formes normales.}
Nous déduirons plus loin (\S\ 8) une propriété principale du procédé de contraction (41.), à savoir sa composition à partir de contractions fondamentales; mais il nous faut pour cela transférer au domaine $n$-aire un théorème important donné par M. Mertens\footnote{F. Mertens, \emph{Über invariante Gebilde quaternärer Formen} (voir l'Introduction), cité comme \og{} Mertens \fg{}.} dans le domaine quaternaire. Ce théorème affirme qu'un système fini arbitraire de formes peut être remplacé par un système équivalent de formes normales. Par \og{} forme normale \fg{}, nous entendons ici -- en généralisant la définition binaire et ternaire\footnote{Study, a. a. O. p. 54.} -- une forme dont toutes les formations invariantes linéaires dans les coefficients s'annulent identiquement, à l'exception de la forme fondamentale elle-même. D'après le théorème VI (\S\ 6), une forme normale $\varphi$ satisfait donc au système d'équations différentielles:
\begin{equation}
 \left(q_{n-\sigma-\lambda}\frac{\partial}{\partial p_\sigma}\middle|\frac{\partial}{q_{n-\tau}}p_{\tau-\lambda}\right)=0,
 \qquad(\sigma\geqq\tau;\ \lambda=1,2,\ldots\tau,n-\sigma),
\tag{46}
\end{equation}
où les séries $q_{n-\sigma-\lambda},p_{\tau-\lambda}$ doivent être considérées comme des grandeurs arbitraires. Dans le domaine quaternaire -- en tenant compte de (39.) pour $\sigma=\tau=2$ -- ces équations différentielles coïncident entièrement avec les 10 équations différentielles données par M. Mertens sans introduction des séries arbitraires $p,q$ (a. a. O., formule 28)); leur connaissance, jointe au théorème VI, nous permet de transférer presque mot pour mot sa méthode de preuve à $n$ variables. La seule chose à ajouter est la vérification, obtenue au moyen de l'identité de décomposition (30.), que les formes construites sont des formes normales; dans le domaine quaternaire, cette vérification s'obtient encore sans transformation supplémentaire. Il suffit donc de donner une brève exposition de la preuve, et cela, par simplicité, dans le seul cas utile pour la suite: celui d'une unique forme fondamentale $f$ et de sa série de formes (cf. fin du \S\ 6), car, pour les formations invariantes de degré quelconque dans les coefficients d'un nombre quelconque de formes fondamentales, la preuve reste la même.

L'idée fondamentale de la preuve est de remplacer, par l'introduction des coefficients transformés, une forme à $n-1$ séries $p_\rho$ par des formes à $n$ séries $\xi$ cogrédientes aux $x$, les séries des grandeurs de transformation. De ces formes se déduisent directement, par des procédés invariants de différentiation, les formes normales cherchées. Le développement des formes générales suivant des formes normales s'effectue au moyen d'un développement en série pour les formes à $\rho$ ($\rho\leqq n$) séries cogrédientes $\xi$; des procédés invariants de différentiation ramènent ensuite aux formes dans les séries initiales $p_\rho$.

Soient $\xi^{(1)},\ldots\xi^{(n)}$ les séries des grandeurs de transformation, de sorte que l'on ait:
\begin{equation}
\begin{aligned}
x_i&=\xi_i^{(1)}x'_1+\xi_i^{(2)}x'_2+\cdots+\xi_i^{(n)}x'_n,\\
p_{i_1\ldots i_\rho}&=\sum_j(\xi_{i_1}^{(j_1)}\cdots \xi_{i_\rho}^{(j_\rho)})p'_{j_1\ldots j_\rho}.
\end{aligned}
\tag{47}
\end{equation}
Les coefficients transformés des formes $f,\varphi,\ldots$ seront notés $(f),(\varphi),\ldots$, et soit
\[
 f\equiv\pair{S^1}{p_1}^{m_1}\pair{S^2}{p_2}^{m_2}\cdots\pair{S^{n-1}}{p_{n-1}}^{m_{n-1}}.
\]
On a alors:
\begin{equation}
\begin{aligned}
(f)&=\pair{S^1}{\xi^{(1)}}^{m_{11}}\cdots\pair{S^1}{\xi^{(n)}}^{m_{1n}}
\pair{S^2}{\xi^{(1)}\xi^{(2)}}^{m_{21}}\cdots
\pair{S^{n-1}}{\xi^{(2)}\cdots\xi^{(n)}}^{m_{n-1,n}}\\
&=\pair{s^{(1)}}{\xi^{(1)}}^{k_1}\pair{s^{(2)}}{\xi^{(2)}}^{k_2}\cdots\pair{s^{(n)}}{\xi^{(n)}}^{k_n},
\end{aligned}
\tag{48}
\end{equation}
où $m_{11}+\cdots+m_{1n}=m_1$, etc. De chaque coefficient $(f)$ pour lequel
\begin{equation}
 k_1\geqq k_2\geqq k_3\cdots\geqq k_n,
\tag{49}
\end{equation}
résulte, par le procédé de différentiation qui épuise exactement les séries $\xi$,
\begin{equation}
\left(\frac{\partial}{\partial\xi^{(1)}}\middle|p_1\right)^{k_1-k_2}
\left(\frac{\partial}{\partial\xi^{(1)}}\frac{\partial}{\partial\xi^{(2)}}\middle|p_2\right)^{k_2-k_3}
\cdots
\left(\frac{\partial}{\partial\xi^{(1)}}\cdots\frac{\partial}{\partial\xi^{(n-1)}}\middle|p_{n-1}\right)^{k_{n-1}-k_n}
\left(\frac{\partial}{\partial\xi^{(1)}}\cdots\frac{\partial}{\partial\xi^{(n)}}\right)^{k_n}
\tag{50}
\end{equation}
la forme
\begin{equation}
 \varphi=\bigl(s^{(1)}\mid p_1\bigr)^{k_1-k_2}
 \bigl(s^{(1)}s^{(2)}\mid p_2\bigr)^{k_2-k_3}\cdots
 \bigl(s^{(1)}\cdots s^{(n-1)}\mid p_{n-1}\bigr)^{k_{n-1}-k_n}
 \bigl(s^{(1)}s^{(2)}\cdots s^{(n)}\bigr)^{k_n}.
\tag{51}
\end{equation}
Les formes $\varphi$ sont les formes normales cherchées. Elles possèdent en effet les propriétés suivantes, qui définissent les formes normales:

1. Ce sont des formations invariantes de la forme fondamentale $f$, linéaires dans les coefficients. Cela résulte du fait qu'elles proviennent de $f$ par les procédés invariants de différentiation $\left(\frac{\partial}{\partial p_\rho}\middle|\xi^{(i_1)}\xi^{(i_2)}\cdots\xi^{(i_\rho)}\right)$ et (50.). Elles peuvent donc s'exprimer linéairement (d'après la fin du \S\ 6) par des polaires de la série finie de formes de $f$. On obtient ces polaires en considérant les variables $p_\rho$ qui entrent dans $\varphi$ comme les coefficients d'une forme qui se décompose en formes linéaires et qui a pour séries de variables les $\mathfrak{p}_\sigma$:
\begin{equation}
 H=(q_{n-1}\mid \mathfrak{p}_{n-1})^{k_1-k_2}(q_{n-2}\mid \mathfrak{p}_{n-2})^{k_2-k_3}\cdots(q_1\mid \mathfrak{p}_1)^{k_{n-1}-k_n}.
\tag{52}
\end{equation}
Les invariants simultanés des séries de formes $f$ et $H$ donnent toutes les polaires à considérer, puisque ces dernières doivent avoir, par rapport aux variables de chaque série $p_\rho$, le même degré que $\varphi$ elle-même.

2. Elles satisfont aux équations différentielles (46.). Si l'on applique en effet à $\varphi$ le procédé différentiel (45.), il vient:
\[
 \varphi'=\bigl(q_{n-\sigma-\lambda}s^{(1)}\cdots s^{(\sigma)}\mid [s^{(1)}\cdots s^{(\tau)}]_{n-\tau}p_{\tau-\lambda}\bigr).\qquad (\sigma\ge\tau)
\]
De l'identité de décomposition (30.) il résulte, si l'on identifie les séries $q_\sigma=s^{(1)}\cdots s^{(\sigma)}$, $p_{n-\tau}=[s^{(1)}\cdots s^{(\tau)}]_{n-\tau}$ formées à partir des séries $s$ avec les $q_\rho,p_\tau$ qui y apparaissent:
\[
 \varphi'=\sum_{\beta=\sigma-\tau+\lambda}^{n-\tau,\sigma}\sum_{i_\beta}\eps\,
 \pair{\dot{q}_{\sigma-\beta}^{(i_1\ldots i_\beta)}q_{n-\tau+\lambda}}{p_{\sigma+\lambda}\dot{p}_{n-\tau-\beta}^{(i_1\ldots i_\beta)}}.
\]
On a alors:
\[
 \dot{p}_{n-\tau-\beta}^{(i_1\ldots i_\beta)}\sim s^{(1)}\cdots s^{(\tau)}s^{(i_1)}\cdots s^{(i_\beta)},
\]
où $i_1,\ldots i_\beta$ désignent $\beta$ nombres quelconques, distincts les uns des autres, choisis parmi les nombres 1 à $\sigma$. Puisque $\beta>\sigma-\tau$, il s'ensuit que, parmi ces nombres $i_1,\ldots i_\beta$, il s'en trouve nécessairement entre 1 et $\tau$; mais alors $\dot{p}_{n-\tau-\beta}^{(i_1\ldots i_\beta)}$ s'annule identiquement, et par conséquent chaque produit matriciel individuel, donc aussi $\varphi'$, s'annule.

Pour établir l'équivalence du système des formes normales avec la série de formes $f$, il reste seulement à montrer que toutes les formes de la série de formes se développent suivant des polaires des formes normales, au sens fixé sous 1. Soit $\Pi$ un agrégat d'opérations polaires
\begin{equation}
 \left(\frac{\partial}{\partial\xi^{(i)}}\middle|\xi^{(k)}\right),
 \qquad (i\ne k;\ i=1,2,\ldots,\rho;\ k=1,2,\ldots,\rho)
\tag{53}
\end{equation}
et soit $\Pi^\sigma$ un agrégat des opérations spéciales (53.) pour lesquelles $i=\sigma$ et $k<\sigma$. Pour une forme $F$ des $\rho$ séries $\xi^{(1)},\ldots,\xi^{(\rho)}$, qui soit du degré $k_\rho$ dans la série $\xi^{(\rho)}$, on a le développement\footnote{F. Mertens, \emph{Über eine Formel der Determinantentheorie}. Sitzungsber. d. Ak. d. Wiss., Vienne, math.-naturw. Kl., Bd. 91. 2. Abt. (1885). Formule 20).}:
\begin{equation}
 F=\sum_{\alpha=0}^{k_\rho}\Pi F_\alpha,
 \qquad(\rho\leqq n),
\tag{54}
\end{equation}
où $F_\alpha$ est du $\alpha$-ième degré dans la dernière série $\xi^{(\rho)}$ et provient de $F$ par les opérations différentielles invariantes
\begin{equation}
 \left(\frac{\partial}{\partial\xi^{(1)}}\cdots\frac{\partial}{\partial\xi^{(\rho)}}\middle|\xi^{(1)}\cdots\xi^{(\rho)}\right)^\alpha
\tag{55}
\end{equation}
et $\Pi^\rho$. Si, dans la production de $F_\alpha$, on remplace le procédé (55.) par
\begin{equation}
 \left(\frac{\partial}{\partial\xi^{(1)}}\cdots\frac{\partial}{\partial\xi^{(\rho)}}\middle|p_\rho\right)^\alpha,
\tag{56}
\end{equation}
il naît une forme $F_\alpha(p_\rho)$ qui ne contient plus que $\rho-1$ séries $\xi^{(1)},\ldots,\xi^{(\rho-1)}$, à laquelle (54.) peut de nouveau être appliquée. La continuation du procédé conduit à des formes $G(p_\rho,p_{\rho-1},\ldots,p_1)$. Pour obtenir à partir d'elles le développement de $F$ selon (54.), il faut remplacer les séries individuelles $p_\sigma$ par $\xi^{(1)},\ldots,\xi^{(\sigma)}$ et appliquer aux formes qui en résultent le procédé polaire $\Pi$ (53.); cela donne (cf. (48.)) une somme de coefficients transformés $(G)$. Pour $\rho=n$, le procédé (55.) est conservé au premier pas. Si $c$ désigne des constantes numériques et si, pour abréger, on pose encore
\begin{equation}
 (\xi^{(1)}\xi^{(2)}\cdots\xi^{(n)})=\Delta;
 \qquad
 \left(\frac{\partial}{\partial\xi^{(1)}}\frac{\partial}{\partial\xi^{(2)}}\cdots\frac{\partial}{\partial\xi^{(n)}}\right)=\nabla,
\tag{57}
\end{equation}
il vient, dans le cas $\rho=n$:
\begin{equation}
 F=\sum c\Delta^\alpha(G),
\tag{58}
\end{equation}
tandis que, pour $\rho<n$, seul $\alpha=0$, donc $\sum c(G)$, apparaît au membre de droite.

Appliquons (58.) à un coefficient transformé $(f)$ (48.). En raison de la permutabilité des procédés polaires $\Pi^\sigma$ avec toutes les opérations (56.) pour lesquelles $\rho\ge\sigma$, et puisque les polaires de coefficients transformés donnent de nouveau des coefficients transformés, on obtient:
\[
 G=\left(\frac{\partial}{\partial\xi^{(1)}}\middle|p_1\right)^{\beta_1}
 \left(\frac{\partial}{\partial\xi^{(1)}}\frac{\partial}{\partial\xi^{(2)}}\middle|p_2\right)^{\beta_2}
 \cdots
 \left(\frac{\partial}{\partial\xi^{(1)}}\cdots\frac{\partial}{\partial\xi^{(n-1)}}\middle|p_{n-1}\right)^{\beta_{n-1}}
 \nabla^{\beta_n}\Pi(f)=\sum c\varphi,
\]
et de (58.) il résulte:
\begin{equation}
 (f)=\sum c\Delta^\alpha(\varphi)\qquad\text{(Mertens, formule 23)).}
\tag{59}
\end{equation}
Pour passer de (59.) au développement de $f$ lui-même, considérons de nouveau les variables $p_\rho$ comme les coefficients d'une forme $H$ (52.) dont les degrés $m_1,\ldots m_{n-1}$ correspondent à ceux de $f$, et soit $m=m_1+\cdots+m_{n-1}$. On a alors, d'après la définition des invariants:
\[
 f\cdot\Delta^m=\sum(f)(H)=\sum c\Delta^\alpha(\varphi)(H),
\]
et l'application du procédé $\nabla^m$ qui épuise exactement les séries $\xi$ donne:
\begin{equation}
 f=\sum c\Phi,
\tag{60}
\end{equation}
où les formes $\Phi$ sont dérivées de $\varphi$ et $H$ par des procédés invariants de différentiation par rapport aux variables, donc sont des invariants simultanés des séries de formes $\varphi$ et $H$.\footnote{Au lieu de cette conclusion directe, Mertens introduit au \S\ 7 une série d'autres procédés différentiels qui ont essentiellement pour but de former la série de formes $H$; c'est ce que réalise notre théorème VI.} Mais, puisque $\varphi$ est une forme normale, la série de formes $\varphi$ se réduit à la forme $\varphi$ elle-même; la série de formes $H$ contient toutes les combinaisons invariantes de variables possibles. Les invariants simultanés de $\varphi$ et de $H$ sont donc les polaires de $\varphi$ au sens indiqué sous 1.; autrement dit, (60.) représente le développement de $f$ suivant des polaires des formes normales. Il nous reste encore à démontrer le même développement pour une forme arbitraire $\theta$ de la série de formes $f$. Soient $\Gamma$ les formes normales déduites de $(\theta)$ par (50.), de sorte que l'on ait:
\begin{equation}
 (\theta)=\sum c\Delta^\beta(\Gamma).
\tag{61}
\end{equation}
Les formes $\theta$ sont caractérisées par la relation valable pour chaque coefficient transformé:
\[
 (\theta)\Delta^\sigma=\sum c(f)\qquad\text{(Mertens, formule 9)).}
\]
Mais comme toute forme $F$ des $n$ séries $\xi$ qui contient le facteur $\Delta^\sigma$ contient aussi le second $\nabla^\sigma\Pi F$ (Mertens, formule 20)), il s'ensuit:
\[
 (\theta)=\sum c\nabla^\sigma(f),\qquad\text{donc}\qquad \Gamma=\sum c\varphi,
\]
et, de (61.), la relation correspondant à (59.):
\[
 (\theta)=\sum c\Delta^\beta(\varphi),
\]
qui entraîne pour $\theta$ le développement (60.) suivant des polaires des formes normales $\varphi$. En résumé, nous avons:

\emph{Théorème VII.} Toute série de formes $f$ peut être remplacée par un système équivalent de formes normales.

\section*{\S\ 8. Réduction du procédé de contraction aux contractions fondamentales.}
À la suite du théorème VII, nous effectuons maintenant la réduction du procédé de contraction aux contractions fondamentales, mentionnée au début du \S\ 7. Nous appelons (cf. \S\ 5) toute contraction de deux séries cogrédientes $S^\rho,T^\rho$ une \og{} contraction cogrédiente \fg{}, et plus spécialement la contraction cogrédiente de défaut 1 de deux séries $S^\rho,T^\rho$ une \og{} contraction fondamentale du type $\rho$ \fg{}. Nous prouvons alors, en complétant le théorème VI, le

\emph{Théorème VIII:} Le procédé général de contraction (41.) peut se composer à partir des $(n-1)$ contractions fondamentales de type $\rho$, où $\rho=1,2,\ldots n-1$. En d'autres termes: pour engendrer toutes les formations invariantes, il suffit d'appliquer successivement les $n-1$ contractions fondamentales.

Dans le domaine binaire, le procédé de contraction (41.) coïncide directement avec l'unique contraction fondamentale possible; dans le domaine ternaire, j'ai démontré le théorème VIII au \S\ 1 de ma dissertation (cf. Introduction). Là, à cause de (44.), les contractions cogrédientes générales sont encore identiques aux contractions fondamentales, et il est remarquable que le théorème VIII, pour $n$ variables, réalise la réduction aux contractions fondamentales, non pas seulement la réduction beaucoup plus faible aux contractions cogrédientes. La preuve est, dans son principe, calquée sur la preuve ternaire, en construisant les contractions les plus générales à partir des contractions fondamentales, au moyen des identités symboliques. Nous engendrerons:
\begin{enumerate}
\item des contractions arbitraires de défaut 1 par des contractions cogrédientes de défaut 1, c'est-à-dire par des contractions fondamentales (analogue du cas ternaire);
\item des contractions arbitraires de défaut $\lambda$ par des contractions cogrédientes de défaut $\lambda$ et par des contractions arbitraires de défaut 1;
\item des contractions cogrédientes de défaut $\lambda$ par des contractions cogrédientes de défaut $\lambda-1$ et par des contractions arbitraires de défaut 1.
\end{enumerate}
Ce qui est essentiel, c'est que nous supposons, d'après le théorème VII, les deux formes à contracter $S$ et $T$ données comme formes normales. D'après (46.), les relations suivantes valent donc:
\begin{equation}
\begin{aligned}
\pair{q_{n-\sigma_1-\lambda}S^{\sigma_1}}{S_{n-\sigma_2}p_{\sigma_2-\lambda}}&=0,
&& (\sigma_1\ge\sigma_2,\ \lambda=1,2,\ldots,\sigma_2,n-\sigma_1),\\
\pair{q_{n-\tau_1-\lambda}T^{\tau_1}}{T_{n-\tau_2}p_{\tau_2-\lambda}}&=0,
&& (\tau_1\ge\tau_2,\ \lambda=1,2,\ldots,\tau_2,n-\tau_1).
\end{aligned}
\tag{62}
\end{equation}
Il suffit en outre, d'après la fin du \S\ 2, de supposer les formes $S$ et $T$ linéaires dans toutes les séries de variables; c'est-à-dire que les formes construites appartiennent à la série de formes
\[
 f=\pair{S^1}{p_1}\pair{S^2}{p_2}\cdots\pair{S^{n-1}}{p_{n-1}}
   \pair{T^1}{p_1}\pair{T^2}{p_2}\cdots\pair{T^{n-1}}{p_{n-1}}.
\]
1) Engendrement de $\pair{q_{n-\sigma-1}S^\sigma}{T_{n-\tau}p_{\tau-1}}$, où $\sigma>\tau$, à partir des contractions fondamentales des types $\tau+\alpha$, $\alpha=0,1,\ldots\sigma-\tau$. De la contraction fondamentale du type $\tau$
\[
 \pair{q_{n-\tau-1}S^\tau}{T_{n-\tau}p_{\tau-1}}
\]
nous obtenons, par la substitution $w\sim T_{n-\tau}p_{\tau-1}$, une forme de la série de formes $f$, avec trois séries d'indice supérieur de poids (cf. \S\ 1) égal à $\tau+1$:
\begin{equation}
 \pair{wS^\tau}{p_{\tau+1}}\pair{S^{\tau+1}}{p_{\tau+1}}\pair{T^{\tau+1}}{p_{\tau+1}}.
\tag{63}
\end{equation}
Une contraction fondamentale du type $\tau+1$, appliquée à (63.), donne:
\[
 \pair{q_{n-\tau-2}wS^\tau}{S_{n-\tau-1}p_\tau};
\]
il en résulte, par (31.) et par la substitution $q_{n-\tau-2}w=q'_{n-\tau-1}$, la valeur:
\[
 \eps\bigl(q_{n-\tau-2}wS^{\tau+1}\bigr)\pair{S^\tau}{p_\tau}
 +\sum_{\alpha=1}^{\tau,n-\tau-1}\Pi\pair{q'_{n-\tau-1-\alpha}S^{\tau+1}}{S_{n-\tau}p_{\tau-\alpha}}.
\]
L'expression sous le signe de sommation s'annule identiquement d'après (62.); nous sommes donc parvenus à une forme
\[
 \pair{wS^{\tau+1}}{p_{\tau+2}}\pair{S^{\tau+2}}{p_{\tau+2}}\pair{T^{\tau+2}}{p_{\tau+2}},
\]
qui provient de (63.) par transformation de $\tau$ en $\tau+1$. La répétition $(\sigma-\tau)$ fois du procédé conduit donc à la contraction cherchée:
\[
 \pair{wS^\sigma}{p_{\sigma+1}}=\eps\pair{q_{n-\sigma-1}S^\sigma}{T_{n-\tau}p_{\tau-1}}.
\]
Nous indiquons brièvement le procédé effectué par la formule:
\begin{equation}
 T^\tau S^\tau,\quad S^{\tau+1},\quad S^{\tau+2},\ldots S^\sigma;
\tag{64}
\end{equation}
nous voyons que seule la propriété de forme normale de $S$, et non celle de $T$, est utilisée. Un procédé dualiste de (64.) peut être effectué, en n'utilisant que la propriété de forme normale de $T$, dans l'ordre inverse, à savoir:
\begin{equation}
 S^\sigma T^\sigma,
\quad T^{\sigma-1},\quad T^{\sigma-2},\ldots,T^\tau,
\tag{65}
\end{equation}
d'après les relations:
\begin{align*}
\pair{q_{n-\sigma-1}S^\sigma}{T_{n-\sigma}p_{\sigma-1}}&=\eps\bigl(T_{n-\sigma}y p_{\sigma-1}\bigr),\quad y\sim q_{n-\sigma-1}S^\sigma,\\
\pair{q_{n-\sigma}T^{\sigma-1}}{T_{n-\sigma}yp_{\sigma-2}}&=\eps\bigl(T_{n-\sigma+1}y p_{\sigma-2}\bigr),\\
&\vdots\\
\pair{q_{n-\tau-1}T^\tau}{T_{n-\tau-1}yp_{\tau-1}}&=\eps\pair{q_{n-\sigma-1}S^\sigma}{T_{n-\tau}p_{\tau-1}}.
\end{align*}
Signalons encore le lien avec la diminution du poids total indiquée à la fin du \S\ 6: pour les contractions de défaut 1, cette diminution a la valeur $2(1+(\sigma-\tau))$, et, pour les contractions fondamentales, la valeur $2\cdot1$. Donc, si la composition à partir de contractions fondamentales est seulement possible, il faut nécessairement $\sigma-\tau+1$ contractions de ce type, comme on vient de le réaliser.

2) Engendrement des contractions générales à partir des contractions cogrédientes et fondamentales. La diminution de poids vaut, pour une contraction générale (41), $2(\lambda^2+\lambda(\sigma-\tau))=2[\lambda^2+(\sigma-\tau)(1+(\lambda-1))]$; elle donne donc la possibilité d'une décomposition en une contraction cogrédiente de défaut $\lambda$ (diminution de poids $2\lambda^2$) et $\sigma-\tau$ contractions de défaut 1 avec différence des indices de poids $\lambda-1$ (diminution de poids $2\{1+(\lambda-1)\}$). Dans le cas $\lambda=1$, cette décomposition coïncide avec celle donnée sous 1); nous montrons qu'elle se réalise effectivement.

De la contraction cogrédiente de défaut $\lambda$
\[
 \pair{q_{n-\tau-\lambda}S^\tau}{T_{n-\tau}p_{\tau-\lambda}}
\]
nous obtenons, par la substitution $R^\lambda\sim T_{n-\tau}p_{\tau-\lambda}$, une forme avec les deux séries d'indices supérieurs de poids $\tau+\lambda$ et $\tau+1$:
\begin{equation}
 \pair{R^\lambda S^\tau}{p_{\tau+\lambda}}\pair{S^{\tau+1}}{p_{\tau+1}}.
\tag{66}
\end{equation}
La série dont l'indice de poids est le plus bas, $\tau+1$, appartient à une forme normale; (66.) admet donc une contraction de défaut 1 composée de $\lambda$ contractions fondamentales, dans l'ordre de (65.):
\[
 (R^\lambda S^\tau)S^{\tau+\lambda},\quad S^{\tau+\lambda-1},\quad S^{\tau+\lambda-2},\ldots S^{\tau+1}.
\]
Cela donne, en tenant compte de (31.) et de (62.):
\[
 \pair{q_{n-\tau-\lambda-1}R^\lambda S^\tau}{S_{n-\tau-1}p_\tau}=\eps\pair{R^\lambda S^{\tau+1}}{p_{\tau+\lambda+1}},
\]
donc une forme qui provient de (66.) par transformation de $\tau$ en $\tau+1$, exactement comme sous 1) avec (63.). La répétition $(\sigma-\tau)$ fois du procédé conduit à la contraction cherchée:
\[
 \pair{R^\lambda S^\sigma}{p_{\sigma+\lambda}}=\eps\pair{q_{n-\sigma-\lambda}S^\sigma}{T_{n-\tau}p_{\tau-\lambda}}.
\]
Dans tout le procédé, seule la propriété de forme normale de $S$ a été utilisée; considérer $T$ comme forme normale conduit au procédé dualiste, avec renversement complet de l'ordre, d'après les formules:
\begin{align*}
\pair{q_{n-\sigma-\lambda}S^\sigma}{T_{n-\sigma}p_{\sigma-\lambda}}&=\eps\bigl(T_{n-\sigma}R_\lambda p_{\sigma-\lambda}\bigr),\quad R_\lambda\sim q_{n-\sigma-\lambda}S^\sigma,\\
\pair{q_{n-\sigma}T^{\sigma-1}}{T_{n-\sigma}R_\lambda p_{\sigma-\lambda-1}}&=\eps\bigl(T_{n-\sigma+1}R_\lambda p_{\sigma-\lambda-1}\bigr),\\
&\vdots\\
\pair{q_{n-\tau-1}T^\tau}{T_{n-\tau-1}R_\lambda p_{\tau-\lambda}}&=\eps\pair{q_{n-\sigma-\lambda}S^\sigma}{T_{n-\tau}p_{\tau-\lambda}}.
\end{align*}
3) Engendrement des contractions cogrédientes de défaut $\lambda$ à partir des contractions cogrédientes de défaut $\lambda-1$ et des contractions fondamentales. La diminution de poids pour les contractions cogrédientes de défaut $\lambda$ vaut $2\lambda^2=2((\lambda-1)^2+2\lambda-1)$, et laisse donc entrevoir une décomposition en une contraction cogrédiente de défaut $\lambda-1$ et $2\lambda-1$ contractions fondamentales. Nous la trouvons le plus simplement en posant d'abord spécialement $n=2\lambda$. L'unique contraction cogrédiente possible de défaut $\lambda$ devient $\pair{S^\lambda}{T_\lambda}=\pair{S^\lambda}{T_{n/2}}$; elle contient une série $u$ et une série $x$ de moins que la contraction de défaut $\lambda-1$ des mêmes séries, $\pair{uS^\lambda}{T_\lambda x}$. Réciproquement, nous composons la contraction de défaut $\lambda$ à partir de contractions fondamentales qui provoquent la disparition d'une série $u$ et d'une série $x$ (diminution de poids $2(n-1)=2(2\lambda-1)$), qui engendrent en même temps une nouvelle série symbolique $R^\lambda$, et d'une contraction cogrédiente de défaut $\lambda-1$. La contraction la plus simple qui réalise cela est la suivante, de défaut 1:
\[
 \pair{sT^\lambda}{\sigma p_\lambda}=\eps\pair{S^{2\lambda-1}}{R_{\lambda-1}p_\lambda}=\eps\pair{R^\lambda}{p_\lambda},
\]
où l'on a posé:
\[
 R^{\lambda+1}=sT^\lambda,
 \quad s=S^1,
 \quad \sigma=S_1,
 \quad R^\lambda\sim\sigma R_{\lambda-1}.
\]
D'après (65.) et (64.), elle est composée des $2\lambda-1$ contractions fondamentales:
\begin{equation}
 T^\lambda S^\lambda,S^{\lambda-1},\ldots S^1;
 \qquad R^{\lambda+1}S^{\lambda+1},S^{\lambda+2},\ldots S^{2\lambda-1}.
\tag{67}
\end{equation}
Une contraction de défaut $\lambda-1$ donne, en tenant compte de (21.) et de (62.):
\[
 \pair{uS^\lambda}{\sigma R_{\lambda-1}x}=\eps\pair{R^{\lambda+1}}{S_\lambda x}
 =\eps\pair{sT^\lambda}{S^\lambda x}=\eps\pair{S^\lambda}{T_\lambda}.
\]
Le cas général, pour deux séries $S^\rho,T^\rho$ et un $n$ quelconque, se transfère directement à partir du cas spécial. À la place de (67.) interviennent les $2\lambda-1$ contractions fondamentales suivantes:
\begin{equation}
 T^\rho S^\rho,S^{\rho-1},\ldots S^{\rho-\lambda+1};
 \qquad R^{\rho+1}S^{\rho+1},S^{\rho+2},\ldots S^{\rho+\lambda-1}.
\tag{68}
\end{equation}
La première moitié des contractions donne la forme:
\[
 \pair{q_{n-\rho-1}T^\rho}{S_{n-\rho+\lambda-1}p_{\rho-\lambda}},
\]
d'où naît, par les substitutions
\[
 S_{n-\rho+\lambda-1}p_{\rho-\lambda}\sim w,
 \qquad T^\rho w=R^{\rho+1},
\]
et par application de la seconde moitié des contractions:
\[
 \pair{q_{n-\rho-\lambda}S^{\rho+\lambda-1}}{R_{n-\rho-1}p_\rho}.
\]
Si l'on substitue:
\[
 q_{n-\rho-\lambda}S^{\rho+\lambda-1}\sim y,
\]
il vient, par une contraction cogrédiente de défaut $\lambda-1$:
\begin{equation}
 \pair{q_{n-\rho-\lambda+1}S^\rho}{yR_{n-\rho-1}p_{\rho-\lambda+1}}.
\tag{69}
\end{equation}
Nous montrons que c'est la contraction cherchée de défaut $\lambda$. Si l'on substitue
\[
 R_{n-\rho-1}p_{\rho-\lambda+1}=R_{n-\lambda}
\]
et si l'on remarque qu'en décomposant $y$, il vient
\[
 \pair{q_{n-\rho-\lambda+1}R^\lambda}{S_{n-\rho}y}
 =\eps\pair{q_{n-\rho-\lambda}S^{\rho+\lambda-1}}{S_{n-\rho}p'_{\rho-1}}=0,
\]
alors (20.) donne pour (69.) la valeur:
\[
 \eps\pair{S^\rho R^\lambda}{p_{\rho+\lambda-1}y}
 =\eps\pair{q_{n-\rho-\lambda}S^{\rho+\lambda-1}}{[S^\rho R^\lambda]_{n-\rho-\lambda}p_{\rho+\lambda-1}}.
\]
Mais cette contraction est du type (43.), donc équivalente à
\[
 \pair{q_{n-\rho-\lambda}S^\rho}{R_{n-\rho-1}p_{\rho-\lambda+1}}
 =\eps\pair{wT^\rho}{R_\lambda p_{\rho-\lambda+1}},
 \qquad R_\lambda\sim q_{n-\rho-\lambda}S^\rho.
\]
Ici $R_\lambda$ est déjà de la forme finale désirée; l'expression correspond à la forme $\pair{sT^\lambda}{S_\lambda x}$ dans le cas spécial. Si l'on remarque maintenant que la décomposition de $w$ et de $R_\lambda$ donne:
\[
 \pair{wR^{n-\lambda}}{T_{n-\rho}p_{\rho-\lambda+1}}
 =\eps\pair{q'_{\lambda-1}R^{n-\lambda}}{S_{n-\rho+\lambda-1}p_{\rho-\lambda}}
 =\eps\pair{q''_{n-\rho-1}S^\rho}{S_{n-\rho+\lambda-1}p_{\rho-\lambda}}=0,
\]
on obtient par (21.) la valeur:
\[
 \pair{wq_{n-\rho+\lambda-1}}{T_{n-\rho}R_\lambda}
 =\eps\pair{q_{n-\rho+\lambda-1}[T_{n-\rho}R_\lambda]^{\rho-\lambda}}{S_{n-\rho+\lambda-1}p_{\rho-\lambda}}
 \quad\text{équivalente à}\quad
 \pair{q_{n-\rho-\lambda}S^\rho}{T_{n-\rho}p_{\rho-\lambda}}.
\]
Dans tout le procédé, seule la propriété de forme normale de $S$, et non celle de $T$, a été utilisée; la contraction cogrédiente de défaut $\lambda-1$ servant à engendrer (69.) peut donc être remplacée par $2\lambda-3$ contractions fondamentales et par une contraction cogrédiente de défaut $\lambda-2$, qui admet à son tour une décomposition en contractions fondamentales, etc. De manière analogue, on obtient l'engendrement en utilisant la propriété de forme normale de $T$ et en appliquant le procédé dualiste, c'est-à-dire les contractions fondamentales:
\[
 S^\rho T^\rho,\quad T^{\rho+1},\ldots T^{\rho+\lambda-1};
 \qquad R^{\rho-1}T^{\rho-1},\quad T^{\rho-2},\ldots T^{\rho-\lambda+1},
\]
et une contraction dualiste de (69.). Il naît alors la contraction équivalente à la précédente
\[
 \pair{q_{n-\rho-\lambda}T^\rho}{S_{n-\rho}p_{\rho-\lambda}}.
\]

Si l'on ajoute la décomposition obtenue sous 2), on a bien obtenu un engendrement des contractions les plus générales à partir des contractions fondamentales. Il reste encore à montrer que cet engendrement est unique, abstraction faite de la permutation de l'ordre, c'est-à-dire qu'à chaque contraction correspond un et un seul nombre $\alpha_\rho$ de contractions fondamentales du type $\rho$, où $\rho=1,2,\ldots n-1$.

Soient $m_\rho$ les degrés dans les séries de variables $p_\rho$ de la forme initiale $f$, et $l_\rho$ les degrés de la forme née par contraction. Chaque contraction fondamentale du type $\rho$ transforme respectivement les degrés $m_{\rho-1},m_\rho,m_{\rho+1}$ en $m_{\rho-1}+1,m_\rho-2,m_{\rho+1}+1$. Les $\alpha$ satisfont donc au système d'équations:
\begin{equation}
 m_\rho-l_\rho=-\alpha_{\rho-1}+2\alpha_\rho-\alpha_{\rho+1},
 \qquad (\rho=1,2,\ldots n-1),
\tag{70}
\end{equation}
où il faut poser $\alpha_0$ et $\alpha_n=0$, et dont le déterminant a la valeur $n$.\footnote{Si $\Delta_n$ désigne le déterminant à $n$ rangées, on a la formule de récurrence: $\Delta_n=2\Delta_{n-1}-\Delta_{n-2}$; c'est-à-dire $\Delta_n-\Delta_{n-1}=\Delta_{n-1}-\Delta_{n-2}=\cdots=\Delta_2-\Delta_1=1$, puisque $\Delta_2=3$, $\Delta_1=2$.} Les $\alpha$ sont donc ainsi déterminés de façon unique, et même, d'après le théorème VIII, comme entiers positifs; cela donne des conditions pour les différences $m_\rho-l_\rho$.

Les résultats de ce paragraphe valent en vertu des relations (62.); mais il ne faut pas croire que, dans le cas où les relations (62.) ne valent pas, les expressions finales soient équivalentes aux expressions initiales. La définition de l'équivalence donnée au \S\ 6 admet en effet aussi la formulation suivante:\footnote{Comparer la réflexion à la fin du paragraphe suivant.} \og{} Deux formes sont équivalentes si et seulement si elles diffèrent par des formes de contraction plus élevée, c'est-à-dire par des formes qui présentent une diminution de poids plus élevée. \fg{} Cette diminution de poids plus élevée a toujours lieu lorsque les deux séries $q_{n-\sigma-\lambda},p_{\tau-\lambda}$ qui entrent dans la contraction sont réellement des séries de variables, comme c'est le cas dans (41.), (42.), dans les contractions qui apparaissent comme équivalentes dans ce paragraphe, ou encore dans les formes équivalentes du théorème VII. Mais elle n'a pas nécessairement lieu lorsque ces séries sont dérivées de séries symboliques par substitution, comme c'était le cas pour les formes de ce paragraphe qui étaient exprimables les unes par les autres en vertu de (62.). On voit facilement que les formes qui s'annulent présentent même en partie un poids total plus élevé.

\section*{\S\ 9. Séries de formes. Théorèmes de réduction\srcfn{*)}{Voir les considérations correspondantes dans le domaine ternaire. Diss. \S\S\ 1--3.}.}
Nous définissons maintenant de façon un peu plus précise la série de formes introduite à la fin du \S\ 6 comme \og{} l'ensemble des formations invariantes, linéaires dans les coefficients, d'une forme initiale donnée; c'est-à-dire l'ensemble des formes issues de la forme initiale par contraction sur elle-même, ordonnées suivant les formes supérieures \fg{}.

Pour expliquer les formes supérieures, supposons que la forme initiale soit donnée, d'après le théorème VII, comme produit de deux formes normales $S\cdot T$; nous appelons alors formes supérieures des formes obtenues par une contraction supérieure sur elles-mêmes, et, parmi les formes de même contraction sur elles-mêmes, celles qui sont spécialement normalisées. Plus précisément: soit la forme $A$ issue de $\alpha_\rho$ contractions fondamentales du type $\rho$, et la forme $B$ de $\beta_\rho$ contractions fondamentales du type $\rho$. Alors $B$ est supérieure à $A$:
\begin{enumerate}
\item si, dans le système d'inégalités $\beta_\rho\geqq\alpha_\rho$, $\rho=1,2,\ldots n-1$, le signe d'inégalité stricte vaut pour au moins une valeur de $\rho$;\footnote{Si, dans le système d'inégalités, le signe $>$ vaut en partie et le signe $<$ en partie, alors les formes ne sont pas comparables entre elles, puisqu'à une forme issue d'une autre par contraction correspond toujours un système de valeurs positif de contractions fondamentales, donc, dans notre cas, des valeurs positives ou nulles $\beta_\rho-\alpha_\rho$.}
\item si le signe d'égalité vaut pour toutes les valeurs de $\rho$, mais qu'un plus petit nombre de séries symboliques sont réunies en un seul produit matriciel. Cette normalisation se révèle nécessaire à cause des réductions du \S\ 8. Ainsi, par exemple, $\pair{S^{n-1}}{T_{n-1}}$ est supérieure à $\pair{S^{n-1}}{[S^1T^\rho]_{n-\rho-1}p_\rho}$.
\item Si le signe d'égalité vaut pour toutes les valeurs de $\rho$ et si le nombre des séries symboliques réunies en un produit matriciel est le même, $B$ est rendue supérieure à $A$ par une normalisation arbitraire en soi, mais fixée une fois pour toutes. Cette normalisation est toujours possible en raison du nombre fini de formes appartenant à un système de valeurs $\alpha$; elle peut, par exemple, être obtenue en distinguant l'une des formes, $S$ (plus grand nombre de séries symboliques $S$, plus grande somme des indices supérieurs de poids des séries $S$, etc.).
\end{enumerate}
Il est commode de se représenter la série de formes comme une variété de points de réseau de dimension $(n-1)$, les $n-1$ directions correspondant aux $n-1$ contractions fondamentales. À chaque forme correspond alors un et un seul système de valeurs $\alpha$, c'est-à-dire un point de réseau à coordonnées entières positives, tandis que les différentes formes appartenant à un système de valeurs $\alpha$ sont déterminées de façon unique, dans leur ordre, par la normalisation ci-dessus. Une forme quelconque de la série de formes fournit la forme initiale d'une nouvelle série de formes, contenue comme partie dans la série totale de formes; plus précisément, elle est contenue dans la partie formée par les formes que l'on obtient, à partir de la nouvelle forme initiale, en avançant dans les $n-1$ directions positives.

En vertu du théorème VIII et de la théorie générale des développements en séries, les théorèmes de réduction des domaines binaire et ternaire valent exactement pour les séries de formes (cf. ma dissertation, \S\ 3). Nous allons encore déduire brièvement le théorème principal. Nous faisons précéder la définition:

\og{} Dans un système donné de formes, soit un certain nombre fini de formes réuni en un module; nous appelons réductible une forme qui peut s'exprimer au moyen de formes ayant des invariants pour facteurs, ou au moyen de formes supérieures; c'est-à-dire au moyen de formes supérieures de la série totale de formes à laquelle la forme appartient, ou au moyen de formes qui contiennent les symboles du module dans un ordre supérieur. Nous appelons réducteur une série de formes réductible. \fg{} On a alors:

\emph{Théorème IX:} Si la forme initiale d'une série de formes est réductible parce qu'elle est issue d'une contraction avec un réducteur, alors toute la série de formes issue de cette forme initiale est réductible.

Pour la preuve, remarquons qu'une forme $h$ issue de la contraction d'une forme $f$ avec une seconde forme $g$ peut, à cause de (45.), être considérée comme une forme $f$ contenant plusieurs séries de variables cogrédientes $p_\rho,p'_\rho$. Or de telles formes peuvent, d'après la théorie générale du développement en séries, être représentées comme des polaires des covariants élémentaires de $f$, en appliquant successivement le développement en séries à chaque paire de séries de variables cogrédientes $p_\rho,p'_\rho$. D'après (38.), les covariants élémentaires qui apparaissent sont les formes engendrées par contraction cogrédiente. Mais, comme l'ordre d'application du développement en séries est encore arbitraire, nous pouvons en particulier choisir un ordre qui corresponde à l'engendrement des contractions générales à partir des contractions fondamentales. Le système de covariants élémentaires qui apparaît alors devient identique à la série de formes $f$; les formes issues d'une contraction cogrédiente de défaut plus élevé se rangent en effet parmi celles qui sont issues de contractions fondamentales. On parvient toutefois aussi aux polaires de la série de formes $f$ pour un ordre arbitraire du développement en séries, auquel peut correspondre un second système de covariants élémentaires. Comme la série de formes épuise l'ensemble des formations invariantes, toute forme du second système peut s'exprimer linéairement par des polaires de la série de formes, et plus précisément par des polaires des formes qui présentent la même diminution de poids ou une diminution plus élevée. Ce dernier point résulte du fait que (cf. \S\ 7) les polaires peuvent être considérées comme des invariants simultanés d'une série de formes $H$ (52.), avec les degrés correspondant à la forme $h$, et de la série de formes $f$. À chaque contraction de $H$ sur elle-même correspond une diminution de poids qui, après l'exécution de la substitution (11.), se transfère aux séries symboliques et donc aux formes de $f$.\footnote{Cette considération est aussi à la base de la seconde formulation de la notion d'équivalence (fin du \S\ 8). On trouvera d'autres développements sur les différents systèmes de covariants élémentaires dans le domaine ternaire chez Study, loc. cit., II, \S\ 7, théorèmes 7) et 8). En vertu de notre théorème VII, les mêmes résultats valent aussi pour $n$ variables.}

Une forme quelconque de la série de formes $h$ est issue d'une contraction de $h$ sur elle-même; c'est-à-dire, d'une part, d'une contraction supérieure de $g$ avec $f$, et, d'autre part, d'une contraction de $f$ ou de $g$ sur elle-même. Dans les deux cas, le développement en séries conduit, exactement comme on l'a montré ci-dessus pour la forme initiale $h$, à des polaires de la série de formes $f$. Si, en particulier, $f$ est un réducteur, il en résulte un développement suivant les polaires du réducteur; la série de formes $h$ devient donc réductible.

Les applications du théorème à la réduction des systèmes complets de formes s'obtiennent de façon analogue à celles des domaines binaire et ternaire.

\clearpage
\setcounter{footnote}{0}

\begin{center}
{\Large\bfseries \NoetherNavAnchor{P05}{5. Corps de fonctions rationnelles}5. Corps de fonctions rationnelles\footnote{Conférence prononcée à la \emph{Naturforscherversammlung}, Vienne, 1913.}}\par
\vspace{1.2em}
Par \textsc{Emmy Noether}, à Erlangen.\par
\medskip
Jahresbericht der DMV 22 (1913), p. 316--319
\end{center}

\vspace{2em}

Les questions suivantes remontent à l'origine à des conversations avec E. Fischer. Quelques-unes de ces questions ont d'ailleurs déjà été posées et résolues par E. Steinitz pour le cas particulier d'une seule indéterminée, et même sous des hypothèses plus générales sur le corps des coefficients.\footnote{E. Steinitz: \emph{Algebraische Theorie der Körper}, en particulier \S\ 24. Journal de Crelle 137. 1910.}

Par \og corps de fonctions rationnelles \fg{}, j'entends un corps dont les éléments sont des fonctions rationnelles de $n$ indéterminées, à coefficients dans un corps de nombres donné, qui peut en particulier comprendre aussi tous les nombres complexes. Des exemples de tels corps sont le corps des fonctions symétriques de $n$ quantités, ou, plus généralement, l'ensemble des fonctions rationnelles de $n$ indéterminées qui admettent les permutations d'un certain groupe (les \emph{Gattungsbereiche} de Lagrange); de même le corps des invariants.

Entre les deux premiers corps mentionnés et ce dernier, il existe une différence caractéristique: les premiers contiennent $n$ fonctions algébriquement indépendantes, les fonctions symétriques élémentaires; tandis que le nombre des invariants algébriquement indépendants est toujours plus petit que le nombre des indéterminées. Il est donc important que, de façon générale, on puisse associer biunivoquement de tels corps de la seconde espèce à des corps de la première espèce, en remplaçant quelques-unes des indéterminées par des nombres. Dans ce qui suit, je me limiterai donc aux corps de la première espèce, c'est-à-dire aux corps possédant $n$ fonctions algébriquement indépendantes; en vertu de l'association, les résultats valent alors en général.

Les questions se groupent autour de trois notions de base: base rationnelle, base minimale et base d'intégrité.

1. Par \og base rationnelle \fg{}, j'entends un nombre fini de fonctions du corps telles que toute fonction du corps puisse se représenter comme combinaison rationnelle de ce nombre fini de fonctions, à coefficients dans le domaine de coefficients prescrit. On montre par des considérations simples -- qui, dans le cas d'une seule indéterminée, se trouvent aussi chez E. Steinitz, loc. cit. -- l'existence de la base rationnelle pour tout corps de fonctions rationnelles. Par exemple, d'après le théorème de Lagrange, les fonctions symétriques élémentaires et une fonction appartenant au groupe forment une base rationnelle pour les \emph{Gattungsbereiche} de Lagrange mentionnés plus haut.

2. Toute base rationnelle doit contenir, pour les corps de la première espèce, au moins $n$ fonctions; si elle contient exactement $n$ fonctions, qui doivent alors être algébriquement indépendantes, je l'appelle \og base minimale \fg{}. La question de l'existence de la base minimale peut être résolue en partie par des théorèmes de Lüroth, Castelnuovo et Enriques.\footnote{J. Lüroth: \emph{Beweis eines Satzes über rationale Kurven}. Math. Ann. 9. 1875. -- G. Castelnuovo: \emph{Sulla razionalità delle involuzioni piane}. Math. Ann. 44. 1893. -- F. Enriques: \emph{Sopra una involuzione non razionale dello spazio}. Rend. Acc. Linc. vol. 21. 21 janvier 1912.} D'après ceux-ci, il existe toujours, pour les corps d'une seule indéterminée, une base minimale: la fonction de Lüroth du corps, déterminée de façon unique à une transformation homographique près (cf. Steinitz \S\ 24). Pour les corps de deux indéterminées, une base minimale existe toujours, et en général seulement, si l'on admet une extension algébrique du domaine des coefficients; tandis que, pour trois indéterminées et davantage, une base minimale n'existe généralement pas du tout. On n'a pas encore tenté de caractériser les corps particuliers qui possèdent une base minimale.

J'ai poursuivi la question de la base minimale pour les \emph{Gattungsbereiche} de Lagrange. Elle prend ici la signification suivante: \og Si le \emph{Gattungsbereich} de Lagrange appartenant au groupe $G$ possède une base minimale, alors on peut construire rationnellement, par représentation paramétrique, l'ensemble des équations ayant le groupe prescrit $G$; et cela pour tout corps de nombres qui contient le domaine des coefficients de la base minimale -- donc, en particulier, pour un corps de nombres quelconque si ce domaine des coefficients est le corps des nombres rationnels. \fg{}

L'exemple le plus simple est fourni par le groupe symétrique; ici les fonctions symétriques élémentaires forment une base minimale à coefficients rationnels. Il en résulte, comme représentation paramétrique, simplement l'équation à coefficients indéterminés. Le théorème d'irréductibilité de Hilbert montre comment passer de celle-ci à des équations sans affect, en nombre arbitraire, sur tout corps de nombres.

Or l'existence de la base minimale pour tous les groupes qui apparaissent dans les équations de degré 3 et 4 résulte directement des théorèmes de Lüroth et de Castelnuovo; il apparaît en outre que cette base minimale possède des coefficients rationnels. On peut donc, pour tout corps de nombres quelconque, construire rationnellement l'ensemble des équations de degré 3 et 4 ayant un groupe prescrit. Pour les groupes cycliques et diédraux, il existe une base minimale dont le domaine des coefficients est le corps des racines $n$-ièmes de l'unité; il en va de même pour quelques autres groupes métacycliques. La question de savoir si la base minimale est même caractéristique de certains groupes doit rester indécise.

3. Pour arriver à la dernière notion de base, je considère l'ensemble des polynômes contenus dans le corps. Par \og base d'intégrité \fg{}, j'entends un nombre fini de ces polynômes tel que tout polynôme du corps puisse se représenter comme combinaison rationnelle entière de ce nombre fini de polynômes, à coefficients dans le domaine de coefficients prescrit. En termes modernes, il s'agit d'une famille finie qui engendre l'anneau comme algèbre sur le domaine de coefficients prescrit, et non d'une base linéaire ni d'une base de module. La question de l'existence de la base d'intégrité -- qui contient par exemple, comme cas particulier, la finitude du système d'invariants -- a été posée par Hilbert, sous une forme un peu différente, dans ses \og Problèmes mathématiques \fg{},\footnote{D. Hilbert: \emph{Mathematische Probleme}. Conférence de Paris 1900. Problème 14. Göttinger Nachr. 1900.} comme le problème des fonctions relativement entières.

Pour l'instant, je ne peux indiquer qu'une classe de corps pour lesquels une base d'intégrité existe. Si, en effet, le corps contient un système de $n$ polynômes tel que la résultante homogène des termes de plus haute dimension de ces polynômes soit non nulle, alors il existe une base d'intégrité; elle est donnée par les coefficients de tous les $u$ de l'équation, irréductible dans le corps, de la forme linéaire:
\[
  u_0=u_1x_1+u_2x_2+\cdots+u_nx_n.\footnotemark
\]
\footnotetext{Cette équation irréductible se trouve, dans le cas d'une seule indéterminée, chez E. Steinitz, dans la preuve du théorème de Lüroth.}
Si, en particulier, la valeur de la résultante est égale à une unité du domaine des coefficients, alors l'intégralité de la représentation est elle aussi assurée. Un exemple de cette classe de corps est fourni par les \emph{Gattungsbereiche} de Lagrange, puisque la résultante des fonctions symétriques élémentaires a la valeur $\pm 1$. Un autre exemple, sans intégralité, est donné par tous les corps qui ne contiennent qu'un seul polynôme algébriquement indépendant; ici la base d'intégrité est identique à la fonction de Lüroth du plus petit sous-corps contenant tous les polynômes. Ainsi, comme on sait, tous les invariants projectifs rationnels entiers d'une forme quadratique à $n$ variables sont des fonctions entières du discriminant.

La condition indiquée est seulement suffisante, nullement nécessaire, pour l'existence de la base d'intégrité. On peut construire facilement des corps dans lesquels la résultante de $n$ polynômes quelconques s'annule, et qui possèdent néanmoins une base d'intégrité donnée par les coefficients de cette équation irréductible. Il se pose alors la question de savoir si, peut-être, dans le cas le plus général lui aussi, les coefficients de cette équation fournissent la base d'intégrité.

\clearpage
\setcounter{footnote}{0}
\providecommand{\Kfield}{\mathfrak{K}}
\providecommand{\Ssys}{\mathfrak{S}}
\providecommand{\Mfield}{\mathfrak{M}}
\providecommand{\Lsys}{\mathfrak{L}}
\providecommand{\Jdom}{\mathfrak{J}}

\begin{center}
{\Large\bfseries \NoetherNavAnchor{P06}{6. Corps et systèmes de fonctions rationnelles}6. Corps et systèmes de fonctions rationnelles}\par
\vspace{0.6em}
Emmy Noether (Erlangen).\par
\vspace{1.2em}
\emph{Math. Ann.} 76 (1915), p. 161--196
\end{center}

\vspace{1.5em}

Le présent travail traite des questions de base pour des systèmes arbitraires de fonctions rationnelles et de fonctions rationnelles entières; et, à vrai dire, les méthodes employées de la théorie des corps font apparaître comme essentiel le règlement de ces questions pour les corps formés de fonctions rationnelles --- les corps de fonctions rationnelles\footnote{Cf. une communication préliminaire à l'occasion de la réunion des naturalistes à Vienne: \emph{Rationale Funktionenkörper}, \emph{Jahresber. d. D. Math.-Ver.} 22 (1913), où est donné un aperçu des questions et des résultats concernant les corps de fonctions.} ---, tandis que la généralisation des résultats à des systèmes arbitraires se présente comme une conséquence.

Parmi les questions fondamentales relatives aux systèmes généraux, seule était connue jusqu'ici l'existence d'un système fini de générateurs, garantie par le théorème de la base de Hilbert (\emph{Math. Ann.} 36). Dans ce qui suit, la question de la représentabilité rationnelle est entièrement résolue par l'existence d'une base rationnelle pour tout système arbitraire (\S\ 7); la base rationnelle des corps est déjà obtenue au \S\ 4. Cette existence de la base rationnelle permet partout de partir du corps ou du système défini abstraitement, et d'éviter par là les difficultés qui ne proviennent que du choix spécial de la base rationnelle, et non du système lui-même, comme l'apparition de dénominateurs particuliers ou de points fondamentaux des fonctions de base.

Pour les corps est également traitée la question de la base minimale, c'est-à-dire d'une base rationnelle composée de fonctions algébriquement indépendantes (\S\ 6). Les méthodes de la théorie des corps conduisent en outre à une base rationnelle distinguée, la base d'involution\footnote{Le nom a été choisi parce que, dans son interprétation géométrique, le corps de fonctions rationnelles représente l'involution la plus générale dans un espace linéaire.} (\S\ 5), qui devient particulièrement importante pour les domaines d’intégrité constitués de polynômes (\S\ 8). Elle y donne une représentation avec un dénominateur fixe (une puissance d'une fonction déterminée par le domaine d’intégrité), comme on le connaît par exemple dans le cas spécial de la représentation typique des invariants.

De la représentabilité rationnelle on tire alors les conséquences les plus étendues possible pour la représentabilité rationnelle entière, c'est-à-dire pour la question de la finitude au sens étroit.

Les théorèmes de finitude connus jusqu'à présent reposent tous sur le fait que le théorème de la base de Hilbert garantit la finitude pour tous les systèmes dans lesquels une représentation
\[
  F=A_1f_1+A_2f_2+\cdots+A_kf_k
\]
entraîne une seconde représentation
\[
  F=B_1f_1+B_2f_2+\cdots+B_kf_k,
\]
où les $B_i$ appartiennent au système et sont de degré inférieur à celui de $F$.

Au contraire, le théorème sur la base rationnelle et les méthodes de la théorie des corps permettent de conclure à la finitude sous des hypothèses d'une nature tout autre.

Ainsi, comme réponse partielle à un problème de Hilbert (\emph{Problèmes mathématiques}, 14), on obtient une classe de fonctions relativement entières (domaines relativement entiers de première espèce) qui sont des domaines d’intégrité finis et dont la base d'intégrité est donnée par la base d'involution mentionnée ci-dessus (\S\ 10). Par analogie avec la preuve par Hilbert du théorème de Kronecker sur le système fondamental des entiers algébriques (\emph{Systèmes complets d'invariants}, \S\ 2; \emph{Math. Ann.} 42), on peut en outre montrer que tous les systèmes réguliers de polynômes --- c'est-à-dire tous les systèmes qui peuvent être envoyés sur des systèmes sans points fondamentaux --- possèdent une base d'intégrité (\S\ 12), tandis que des systèmes non réguliers sans base d'intégrité peuvent être indiqués facilement. Comme exemple simple de ce fait on peut mentionner le théorème: \og{}Dans tout système arbitraire de polynômes en nombre infini d'une seule indéterminée, il existe un nombre fini de ces polynômes tel que chaque polynôme du système puisse être représenté comme combinaison rationnelle entière de ce nombre fini.\fg{} D'autres exemples se trouvent au \S\ 13.

Enfin, les derniers paragraphes (\S\S\ 14 et 15) donnent un raffinement des théorèmes de base relativement à l'intégralité.

Un outil essentiel de l'étude est le principe de transfert du \S\ 3, selon lequel il suffit de considérer toutes les questions de base soulevées ici pour les systèmes dans lesquels le nombre des indéterminées coïncide avec le nombre des fonctions algébriquement indépendantes du système.

L'impulsion du présent travail est venue de conversations avec M. E. Fischer, en particulier d'une question de lui concernant la base minimale des \emph{Gattungsbereiche} de Lagrange. Les recherches qui s'y rattachent, et qui conduisirent à la construction d'équations à groupe prescrit (cf. la communication citée sur les \og{}Corps de fonctions rationnelles\fg{}, deuxième partie), sont réservées à une publication ultérieure.

L'extension essentielle du présent travail par rapport à la communication originelle consiste en ceci que les théorèmes de base qui y étaient énoncés seulement pour des corps ou des domaines d’intégrité spéciaux sont ici transférés à des systèmes arbitraires. Ce transfert réussit d'une manière simple au moyen de la notion de plus petit corps contenant un système arbitraire, comme généralisation du corps des fractions d'un domaine d’intégrité (\S\ 7). J'ai été conduite à former cette notion par le fait que M. K. Hentzelt, après avoir appris l'existence de la base rationnelle des corps, put démontrer l'existence de la base rationnelle pour des systèmes arbitraires par un détour passant par le théorème de la base de Hilbert. C'est aussi l'observation de K. Hentzelt selon laquelle les raisonnements que j'avais appliqués aux domaines d’intégrité valaient pour des systèmes arbitraires soumis aux mêmes hypothèses qui m'a conduite au théorème sur la base d'intégrité des systèmes réguliers; de même, je dois à M. Hentzelt quelques remarques particulières de moindre étendue.

Mentionnons enfin que, dans le cas d'une seule indéterminée, la base rationnelle et la base minimale des corps se trouvent dans l'\emph{Algebraische Theorie der Körper} de E. Steinitz, \S\ 24 (\emph{Journal de Crelle} 137); le corps des coefficients y est pris comme un corps au sens le plus général. La question de savoir jusqu'à quel point les théorèmes donnés ici restent valables sous ces hypothèses plus générales n'est pas abordée dans ce qui suit; en tout cas les démonstrations exigent une modification, puisque, par exemple, les moyens tirés de la théorie des déterminants fonctionnels font défaut pour les corps de caractéristique $p$.

\section*{\S\ 1. Corps $\Kfield_{n\rho}$ et systèmes $\Ssys_{n\rho}$.}

Dans ce qui suit, il s'agit d'étudier des systèmes de fonctions rationnelles de $n$ indéterminées, en particulier des corps de fonctions rationnelles.

Par $f(x_1\cdots x_n)$, $g(x_1\cdots x_n)$, \ldots, ou en abrégé $f(x)$, $g(x)$, \ldots, on entendra donc toujours des fonctions rationnelles de $x_1\cdots x_n$; elles seront supposées données sous forme réduite, c'est-à-dire avec numérateur et dénominateur premiers entre eux. Les fonctions rationnelles entières (polynômes) seront désignées par des lettres majuscules, $F(x)$, $G(x)$, \ldots. Le corps des coefficients $\Omega$ est supposé être un corps de nombres quelconque; il peut en particulier contenir aussi tous les nombres complexes. Le corps $\Omega$ peut en outre contenir un nombre fini de paramètres.

Les grandeurs de $\Omega$, donc toutes les fonctions de degré zéro, sont toujours comptées comme appartenant au système.

Avec ces conventions, le corps constitué de fonctions rationnelles --- le corps de fonctions rationnelles --- peut être défini abstraitement.

\emph{Définition I: Un système de fonctions rationnelles est appelé corps s'il satisfait aux conditions suivantes:}
\begin{enumerate}
\item \emph{Avec $f(x)$, et pour toute grandeur $c$ de $\Omega$, il contient aussi $c\cdot f(x)$.}
\item \emph{Avec $f(x)$ et $g(x)$, il contient toujours aussi $f(x)+g(x)$, $f(x)\cdot g(x)$ et --- pour $g(x)\ne0$ --- le quotient $f(x):g(x)$.}\footnote{Ici $f(x)$ et $g(x)$ peuvent aussi être de degré zéro, c'est-à-dire des grandeurs de $\Omega$.}
\end{enumerate}

Sont des corps spéciaux ceux qui naissent par adjonction à $\Omega$ d'un nombre fini de fonctions rationnelles
\[
  f_1(x),\ f_2(x),\ \cdots,\ f_k(x);
\]
ils seront désignés par
\[
  \Omega\bigl(f_1(x),\ldots,f_k(x)\bigr)
  \quad\text{ou}\quad
  \Omega(f_1,\ldots,f_k).
\]
\footnote{Que ces \og{}corps spéciaux\fg{} épuisent déjà la totalité de tous les corps résulte du \S\ 4.}
Parmi ces corps spéciaux nous mentionnons aussi le corps obtenu en adjoignant $x_1\cdots x_n$ à $\Omega$,
\[
  \Omega(x_1,\ldots,x_n),
\]
qui est identique à la totalité des fonctions rationnelles de $x_1\cdots x_n$ à coefficients dans $\Omega$.

Nous introduisons encore (à la suite de E. Steinitz) la notion de corps intermédiaire:
\begin{quote}
\og{}Un corps $\Mfield$ est situé entre $\mathfrak{L}_1$ et $\mathfrak{L}_2$ s'il contient $\mathfrak{L}_1$ et est contenu dans $\mathfrak{L}_2$;\fg{}
\end{quote}
la Définition I admet alors aussi la formulation suivante:

Le corps des fonctions rationnelles de $n$ indéterminées est un corps intermédiaire entre $\Omega$ et $\Omega(x_1\cdots x_n)$ qui contient effectivement les $n$ indéterminées dans au moins une fonction.

À côté du nombre des indéterminées apparaît, pour les systèmes de fonctions rationnelles, un autre nombre caractéristique, le nombre des fonctions algébriquement indépendantes, ou rang algébrique, suivant la définition suivante.

\emph{Définition II: Un système de fonctions rationnelles a le rang algébrique $\rho$ si l'on peut indiquer $\rho$ fonctions du système qui sont algébriquement indépendantes, mais si tout ensemble de $(\rho+1)$ fonctions du système est algébriquement dépendant.}\footnote{On appelle $\rho$ fonctions $f_1(x)\cdots f_\rho(x)$ algébriquement indépendantes si, de toute relation
\[
F\bigl(f_1(x)\cdots f_\rho(x)\bigr)=0\quad\text{identiquement en }x_1\cdots x_n,
\]
il résulte que
\[
F(\lambda_1\cdots \lambda_\rho)=0\quad\text{identiquement en }\lambda_1\cdots \lambda_\rho.
\]
Dans le cas opposé, elles sont dites algébriquement dépendantes.}

Les systèmes de fonctions rationnelles (en nombre fini ou infini) de $n$ indéterminées et de rang algébrique $\rho$ seront désignés par le double indice
\[
  \Ssys_{n\rho}.
\]
En particulier, par
\[
  \Kfield_{n\rho}
\]
il faudra entendre un corps de $n$ indéterminées et de rang algébrique $\rho$. On a l'inégalité:
\[
  1\leq \rho\leq n.
\]

Nous donnons encore quelques exemples de corps $\Kfield_{nn}$ et $\Kfield_{n\rho}$:

1. Des corps $\Kfield_{nn}$ sont les \emph{Gattungsbereiche} de Lagrange, que l'on peut définir comme
\begin{quote}
\og{}la totalité des fonctions rationnelles (et rationnelles entières) de $x_1\cdots x_n$ qui admettent les permutations d'un groupe de permutations $G$ de $x_1\cdots x_n$.\fg{}
\end{quote}
Ils contiennent toujours le corps des fonctions symétriques, donc aussi les $n$ fonctions symétriques élémentaires algébriquement indépendantes.

2. Des corps $\Kfield_{n\rho}$ sont les corps d'invariants projectifs, formés de la totalité des invariants rationnels (et rationnels entiers) d'une ou de plusieurs formes fondamentales.\footnote{Il est opportun de ne considérer que des transformations de déterminant $1$; alors toute combinaison rationnelle d'un nombre arbitraire d'invariants admet encore la transformation, et l'on obtient effectivement un corps. Les invariants homogènes et isobares pris seuls ne forment pas un corps, mais ils forment un système $\Ssys_{n\rho}$.} Ici $\rho<n$, puisque le nombre des invariants algébriquement indépendants est toujours plus petit que le nombre des coefficients.

L'énoncé correspondant vaut pour les corps d'invariants relatifs aux sous-groupes du groupe projectif.

\section*{\S\ 2. Lemme sur des fonctions algébriquement dépendantes.}

Le lemme de ce paragraphe montrera au \S\ 3 que le rang algébrique d'un système est son nombre véritablement caractéristique. Le lemme montre en effet qu'en spécialisant numériquement, successivement, quelques indéterminées, de telle sorte que $\rho$ fonctions algébriquement indépendantes restent algébriquement indépendantes, une fonction arbitraire qui dépend algébriquement de ces $\rho$ fonctions ne peut ni s'annuler identiquement ni devenir indéterminée.

Soit donc le système
\begin{equation}
  g_1(x),\ g_2(x),\ldots,g_\rho(x)
\tag{1}
\end{equation}
de rang algébrique $\rho$, où $\rho<n$, et soit $h(x)$ une fonction rationnelle algébriquement dépendante des fonctions (1). D'après des théorèmes connus, le rang de la matrice fonctionnelle de (1) est égal à $\rho$. Les indéterminées peuvent donc être nommées, par exemple, de sorte que
\begin{equation}
 \frac{\partial(g_1,g_2,\ldots,g_\rho)}{\partial(x_1,x_2,\ldots,x_\rho)}
 =\frac{\Gamma(x)}{G(x)^{\rho+1}}\ne0,
\tag{2}
\end{equation}
où $G(x)$ désigne, comme d'habitude, le dénominateur commun des fonctions (1).

\emph{Choisissons maintenant le nombre $a$ de telle sorte que le produit $G(x)\cdot\Gamma(x)$ ne soit pas divisible par $(x_i-a)$}, où $i$ est entendu comme l'un des nombres $\rho+1,\rho+2,\ldots,n$. Ainsi, pour $x_i=a$, le dénominateur commun des fonctions (1) ne peut pas s'annuler, et les $\rho$ fonctions
\begin{equation}
  [g_1(x)]_{x_i=a},\ldots,[g_\rho(x)]_{x_i=a}
\tag{3}
\end{equation}
sont également algébriquement indépendantes. Comme nous le dirons, le rang algébrique du système (1) est conservé sous la substitution $x_i=a$.

Soit donnée l'équation irréductible satisfaite par $h(x)$ relativement aux fonctions (1):
\begin{equation}
\begin{aligned}
&A_0(g_1,\ldots,g_\rho)\cdot\,h(x)^\tau
 +A_1(g_1,\ldots,g_\rho)\cdot\,h(x)^{\tau-1}
 +\cdots \\
&\hspace{6em}+A_\tau(g_1,\ldots,g_\rho)=0,
\end{aligned}
\tag{4}
\end{equation}
où $A_0(g_1,\ldots,g_\rho)$ et $A_\tau(g_1,\ldots,g_\rho)$ ne s'annulent pas identiquement. Pour passer de là à une identité entre polynômes, nous posons
\begin{equation}
  g_i(x)=G_i(x):G(x);\qquad h(x)=H_1(x):H_2(x),
\tag{5}
\end{equation}
et nous remarquons que --- du fait de la forme réduite supposée pour $h(x)$ --- $H_1(x)$ et $H_2(x)$ sont premiers entre eux. En vertu de (5), (4) devient
\begin{equation}
\begin{aligned}
&B_0(G,G_1,\ldots,G_\rho)\cdot\,G(x)^{\sigma_0\cdot}H_1(x)^\tau \\
&\quad+B_1(G,G_1,\ldots,G_\rho)\cdot\,G(x)^{\sigma_1\cdot}H_1(x)^{\tau-1}H_2(x)+\cdots \\
&\quad+B_\tau(G,G_1,\ldots,G_\rho)\cdot\,G(x)^{\sigma_\tau\cdot}H_2(x)^\tau=0,
\end{aligned}
\tag{6}
\end{equation}
où les $B_k$ sont des formes homogènes en $G,G_i$, et où au moins l'un des exposants non négatifs $\sigma_k$ doit être nul.

Supposons maintenant que $H_1(x)$ soit divisible par $(x_i-a)$. Alors, d'après (6),
\[
  B_\tau(G\,G_1\cdots G_\rho)\cdot G(x)^{\sigma_\tau}\cdot H_2(x)^\tau
\]
serait également divisible par $(x_i-a)$. Cela est exclu par hypothèse pour $G(x)$; donc, de la divisibilité de $B_\tau$ résulterait aussi la divisibilité de
\[
  A_\tau=B_\tau:G^\lambda,
\]
et il existerait entre les fonctions (3) la relation $A_\tau=0$, contrairement à l'indépendance algébrique supposée. Enfin $H_2(x)$, étant premier avec $H_1(x)$, ne peut pas non plus être divisible par $(x_i-a)$;\footnote{Cette dernière conclusion, reposant sur la forme réduite, ne serait plus possible pour une substitution $x_i=a$, $x_k=b$ tandis que les autres hypothèses sont conservées. Alors $H_1(x)$ et $H_2(x)$ pourraient s'annuler simultanément, et le quotient deviendrait indéterminé sous la substitution, égal à $=\frac00$; par exemple
\[
 h(x)=\frac{x_3(\alpha x_1+\beta x_2)+x_4(\gamma x_1+\delta x_2)}
 {x_3(\alpha' x_1+\beta' x_2)+x_4(\gamma' x_1+\delta' x_2)}
\]
sous la substitution $x_3=0$, $x_4=0$.} ainsi l'hypothèse selon laquelle $H_1(x)$ serait divisible par $(x_i-a)$ conduit à une contradiction. De la même façon on montre que $H_2(x)$ ne peut pas non plus être divisible par $(x_i-a)$. La fonction $[h(x)]_{x_i=a}=k(x)$ ne peut donc ni s'annuler identiquement ni devenir indéterminée.

La conclusion précédente peut maintenant être appliquée de nouveau au système (3) et à la fonction $k(x)$, de nouveau supposée sous forme réduite; la poursuite du procédé conduit finalement au lemme suivant.

\emph{Lemme: Les deux systèmes}
\[
  g_1(x),\ g_2(x),\ldots,g_\rho(x),
\]
\[
  g_1(x),\ g_2(x),\ldots,g_\rho(x),\ h(x)
\]
\emph{doivent avoir chacun le rang algébrique $\rho$, où $\rho<n$, et supposons, par exemple, que}
\[
 \frac{\partial(g_1,\ldots,g_\rho)}{\partial(x_1,\ldots,x_\rho)}
 =\frac{\Gamma(x)}{G(x)^{\rho+1}}\ne0.
\]
\emph{Les nombres}
\[
  a_n,\ a_{n-1},\ldots,a_{\rho+1}
\]
\emph{sont à choisir de telle sorte que, sous la substitution}
\[
  x_n=a_n,\quad x_{n-1}=a_{n-1},\ldots,\quad x_{\rho+1}=a_{\rho+1},
\]
\emph{le produit $G(x)\cdot\Gamma(x)$ ne s'annule pas, c'est-à-dire que le rang algébrique du système (1) soit conservé. Dans les substitutions successives}
\[
\begin{aligned}
 [h(x)]_{x_n=a_n}&=h^{(n-1)}(x),\qquad
 [h^{(n-1)}(x)]_{x_{n-1}=a_{n-1}}=h^{(n-2)}(x),\ldots,\\
 [h^{(\rho+1)}(x)]_{x_{\rho+1}=a_{\rho+1}}&=h^*(x_1,x_2,\ldots,x_\rho),
\end{aligned}
\]
\emph{que les fonctions $h(x)$, $h^{(n-1)}(x)$, \ldots, $h^{(\rho+1)}(x)$ soient chaque fois mises sous forme réduite. Sous ces hypothèses, dans la fonction $h^*(x_1,\ldots,x_\rho)$ ni le numérateur ni le dénominateur ne peuvent s'annuler identiquement, et la fonction ne peut pas non plus devenir $\frac{0}{0}$.}\footnote{Par exemple, dans la fonction $h(x)$ de la note précédente, la substitution successive donne
\[
 [h(x)]_{x_4=0}=h^{(3)}(x)=\frac{(\alpha x_1+\beta x_2)}{(\alpha' x_1+\beta' x_2)}.\qquad
 h^*(x_1 x_2)=\frac{\alpha x_1+\beta x_2}{\alpha' x_1+\beta' x_2}.
\]}

\section*{\S\ 3. Application biunivoque des systèmes $\Ssys_{n\rho}$ sur des systèmes $\Ssys^*_{\rho\rho}$.}

Du lemme nous tirerons immédiatement une application des systèmes $\Ssys_{n\rho}$ sur des systèmes $\Ssys^*_{\rho\rho}$, en remplaçant quelques indéterminées par des nombres; le rang algébrique $\rho$ se présente ainsi comme le nombre véritablement caractéristique.

Puisque $\Ssys_{n\rho}$ a le rang algébrique $\rho$, il existe $\rho$ fonctions de $\Ssys_{n\rho}$,
\begin{equation}
  g_1(x),\ldots,g_\rho(x),
\tag{1}
\end{equation}
qui sont algébriquement indépendantes. Si, en outre, $f(x)$ désigne une fonction arbitraire de $\Ssys_{n\rho}$, alors le système
\begin{equation}
  g_1(x),\ldots,g_\rho(x),\ f(x)
\tag{2}
\end{equation}
a également le rang algébrique $\rho$, et les systèmes (1) et (2) satisfont aux conditions du lemme.\footnote{Au besoin, les indéterminées sont à renuméroter.}

Si donc on détermine les nombres
\[
  a_n,\ a_{n-1},\ldots,a_{\rho+1}
\]
d'après le lemme, de telle sorte que le rang algébrique $\rho$ du système (1) soit conservé par la substitution --- ce qui est possible d'un nombre arbitrairement grand de façons ---, alors, dans la fonction définie par
\[
\begin{aligned}
 [f(x)]_{(x_n=a_n)}&=f^{(n-1)}(x),\qquad
 [f^{(n-1)}(x)]_{(x_{n-1}=a_{n-1})}=f^{(n-2)}(x),\\
 [f^{(\rho+1)}(x)]_{(x_{\rho+1}=a_{\rho+1})}&=f^*(x_1,\ldots,x_\rho),
\end{aligned}
\]
à savoir
\begin{equation}
  f^*(x_1,\ldots,x_\rho),
\tag{3}
\end{equation}
ni le numérateur ni le dénominateur ne peuvent s'annuler identiquement.

Faisons maintenant parcourir à $f(x)$ toutes les fonctions (en nombre fini ou infini) de $\Ssys_{n\rho}$. Nous considérons le système formé par la totalité de toutes les fonctions (3); il a les propriétés suivantes:

1. Il a le rang algébrique $\rho$; car il contient les fonctions provenant de (1),
\[
  g_1^*(x_1\cdots x_\rho)\cdots g_\rho^*(x_1\cdots x_\rho),
\]
qui, à cause du choix de $a_n,a_{n-1},\ldots,a_{\rho+1}$, sont algébriquement indépendantes. Le système doit donc être désigné par
\[
  \Ssys^*_{\rho\rho}.
\]
\footnote{De même, par $f^*(x)$, $g^*(x)$, \ldots, il faut toujours entendre des fonctions d'application, c'est-à-dire des fonctions de $x_1\cdots x_\rho$.}

2. La correspondance des fonctions $f(x)$ et $f^*(x)$ est sans exception biunivoque.

En effet, si $f_1(x)$ et $f_2(x)$ appartiennent à $\Ssys_{n\rho}$, leur différence
\[
  d(x)=f_1(x)-f_2(x)
\]
dépend elle aussi algébriquement du système (1); de plus
\[
  d^*(x)=f_1^*(x)-f_2^*(x).
\]
Puisque $d(x)$, avec le système (1), satisfait aux conditions du lemme, il résulte nécessairement de
\[
  d(x)\ne0
\]
que
\[
  d^*(x)\ne0;
\]
autrement dit, à des fonctions différentes de $\Ssys_{n\rho}$ correspondent des fonctions différentes de $\Ssys^*_{\rho\rho}$.

3. De toute relation entre fonctions de $\Ssys^*_{\rho\rho}$,
\[
  F\bigl(f_1^*(x),f_2^*(x),\ldots,f_k^*(x)\bigr)=0
  \quad\text{identiquement en }x_1\cdots x_\rho,
\]
il résulte, pour les fonctions de $\Ssys_{n\rho}$ qui leur correspondent de manière unique, que
\[
  F\bigl(f_1(x),f_2(x),\ldots,f_k(x)\bigr)=0
  \quad\text{identiquement en }x_1\cdots x_n.
\]
Car avec $f_1(x),f_2(x),\ldots,f_k(x)$, toute combinaison rationnelle de ces fonctions dépend elle aussi algébriquement du système (1).

Cela résulte en outre de
\begin{equation}
\begin{aligned}
 h(x)&=F\bigl(f_1(x),\ldots,f_k(x)\bigr):\\
 h^*(x)&=F\bigl(f_1^*(x),\ldots,f_k^*(x)\bigr).
\end{aligned}
\tag{4}
\end{equation}
D'après le lemme, donc,
\[
  h(x)\ne0\quad\text{implique}\quad h^*(x)\ne0;
\]
ce qu'il fallait démontrer.

En résumé, nous obtenons

\emph{Théorème I: À chaque système $\Ssys_{n\rho}$ de fonctions rationnelles (en nombre fini ou infini) peut être adjoint --- d'un nombre arbitrairement grand de façons --- un système biunivoque $\Ssys^*_{\rho\rho}$, de telle sorte que toutes les relations existant entre les fonctions de $\Ssys^*_{\rho\rho}$ soient également conservées dans $\Ssys_{n\rho}$ et réciproquement. En particulier, d'après (4), à un corps $\Kfield_{n\rho}$ correspond de nouveau un corps $\Kfield^*_{\rho\rho}$.}

Du Théorème I nous tirons la conséquence qui simplifie toutes les démonstrations ultérieures: \emph{les propriétés des systèmes $\Ssys_{n\rho}$ qui sont caractérisées par un nombre fini ou infini de relations entre fonctions du système valent pour les systèmes $\Ssys_{n\rho}$ dès qu'elles valent pour un système d'application $\Ssys^*_{\rho\rho}$.}

Nous appellerons brièvement cette conséquence le \og{}principe de transfert\fg{}.

L'exemple le plus simple de cette application est fourni par la théorie des invariants algébriques. En effet, par une transformation linéaire on peut toujours fixer numériquement certains coefficients de la forme fondamentale de manière que toutes les relations valables pour la forme spéciale restent aussi valables en général.

\section*{\S\ 4. Base rationnelle des corps $\Kfield_{n\rho}$.}

Dans les recherches suivantes, groupées autour des notions de base, nous utiliserons constamment le \og{}principe de transfert\fg{} du \S\ 3. Nous démontrons d'abord l'existence de la base pour les corps, mais nous énonçons les définitions de façon générale:

\emph{Définition III: Un nombre fini de fonctions de $\Ssys_{n\rho}$ est appelé base rationnelle si toute fonction de $\Ssys_{n\rho}$ peut être représentée comme combinaison rationnelle de ce nombre fini, à coefficients dans $\Omega$.}

La base rationnelle doit contenir $\rho$ fonctions algébriquement indépendantes, puisque le rang algébrique d'un système dérivé d'une base rationnelle est égal au rang algébrique de la base rationnelle.

Nous montrons d'abord que la démonstration d'existence de la base rationnelle admet le principe de transfert. En effet, la Définition III peut être formulée ainsi:

Les fonctions de $\Ssys_{n\rho}$
\[
  g_1(x),\ g_2(x),\ldots, g_k(x)
\]
forment une base rationnelle si et seulement si sont satisfaites les relations, en nombre infini,
\begin{equation}
  f(x)=\gamma\bigl(g_1(x),\ldots,g_k(x)\bigr),
\tag{1}
\end{equation}
où $f(x)$ parcourt toutes les fonctions de $\Ssys_{n\rho}$ et où $\gamma$ désigne une fonction rationnelle de $k$ arguments, variant avec $f$.

Ainsi, si dans le système d'application $\Ssys^*_{\rho\rho}$ une base rationnelle est donnée par
\[
  g_1^*(x),\ g_2^*(x),\ldots,g_k^*(x),
\]
ce qui entraîne les relations
\[
  f^*(x)=\gamma\bigl(g_1^*(x)\cdots g_k^*(x)\bigr),
\]
alors, par le Théorème I, les relations (1) sont satisfaites pour $\Ssys_{n\rho}$. Le \emph{principe de transfert pour la base rationnelle} dit donc:

\emph{De l'existence d'une base rationnelle pour un système d'application $\Ssys^*_{\rho\rho}$ résulte la base rationnelle du système originel $\Ssys_{n\rho}$.}\footnote{Le principe de transfert se formule de façon correspondante pour chaque question de base traitée ici.}

Pour démontrer l'existence de la base rationnelle de tous les corps $\Kfield_{n\rho}$, nous pouvons donc nous restreindre aux corps $\Kfield^*_{\rho\rho}$; ici une généralisation directe d'un raisonnement employé par E. Steinitz conduit au but.\footnote{E. Steinitz, \emph{Algebraische Theorie der Körper}, \S\ 24, \emph{Journal de Crelle} 137 (1909).}

Soit
\begin{equation}
  g_1^*(x_1,\ldots,x_\rho),\ldots,
  g_\rho^*(x_1,\ldots,x_\rho)
\tag{2}
\end{equation}
un système de $\rho$ fonctions algébriquement indépendantes de $\Kfield^*_{\rho\rho}$. Par élimination de $x_1\cdots x_\rho$, on obtient $\rho$ relations:
\begin{equation}
\begin{aligned}
  F_1\bigl(g_1^*(x),\ldots,g_\rho^*(x),x_1\bigr)&=0,\ldots,\\
  F_\rho\bigl(g_1^*(x),\ldots,g_\rho^*(x),x_\rho\bigr)&=0,
\end{aligned}
\tag{3}
\end{equation}
où, dans chaque relation $F_i=0$, l'indéterminée $x_i$ intervient effectivement; sinon, contrairement à l'hypothèse, il y aurait une dépendance algébrique entre les fonctions (2). Les relations (3) disent que
\[
  x_1,x_2,\ldots,x_\rho
\]
sont des éléments algébriques relativement à
\[
  \Omega\bigl(g_1^*(x),\ldots,g_\rho^*(x)\bigr).
\]
Le corps obtenu par adjonction de $x_1,\ldots,x_\rho$,
\[
  \Omega\bigl(g_1^*,\ldots,g_\rho^*,x_1,
  \ldots,x_\rho\bigr)=\Omega(x_1,
  \ldots,x_\rho),
\]
est donc une extension algébrique finie de $\Omega(g_1^*,\ldots,g_\rho^*)$. D'après un théorème connu,\footnote{Comparer par exemple Weber, \emph{Lehrbuch d. Algebra} 1, \S\ 144 (1re éd.) ou \S\ 55 (petite édition).} $\Kfield^*_{\rho\rho}$, comme corps intermédiaire entre $\Omega(g_1^*\cdots g_\rho^*)$ et $\Omega(x_1\cdots x_\rho)$, est lui-même une extension algébrique de $\Omega(g_1^*\cdots g_\rho^*)$, tandis que $\Omega(x_1\cdots x_\rho)$ est une extension algébrique de $\Kfield^*_{\rho\rho}$. Ainsi $\Kfield^*_{\rho\rho}$ naît de $\Omega(g_1^*\cdots g_\rho^*)$ par adjonction d'un nombre fini d'éléments de $\Kfield^*_{\rho\rho}$, ou encore par l'adjonction d'une fonction convenablement choisie; c'est-à-dire que l'on a
\[
  \Kfield^*_{\rho\rho}
  =\Omega(g_1^*\cdots g_\rho^*g_{\rho+1}^*\cdots g_k^*)
  =\Omega(g_1^*\cdots g_\rho^*g_0^*).
\]

Par le principe de transfert nous avons donc

\emph{Théorème II: Tout corps $\Kfield_{n\rho}$ possède une base rationnelle de rang algébrique $\rho$; la base rationnelle peut en particulier être choisie de telle sorte qu'elle contienne au plus $\rho+1$ fonctions.}

Soit donné un système arbitraire de $\rho$ fonctions algébriquement indépendantes de $\Kfield_{n\rho}$, $h_1(x),\ldots,h_\rho(x)$, et supposons qu'on puisse l'étendre par la méthode précédente à une base rationnelle
\[
  h_1(x),\ldots,h_\rho(x);
  h_{\rho+1}(x),\ldots,h_\sigma(x).
\]
Par la propriété définissante, il existe des relations:
\[
\begin{aligned}
  h_i(x)&=\chi_i\bigl(g_1(x),\ldots,g_k(x)\bigr)
    &&(i=1,2,\ldots,\sigma),\\
  g_j(x)&=\psi_j\bigl(h_1(x),\ldots,h_\sigma(x)\bigr)
    &&(j=1,2,\ldots,k),
\end{aligned}
\]
où $\chi_i$ et $\psi_j$ désignent des fonctions rationnelles de leurs arguments. Ainsi:

À partir de n'importe quelle base rationnelle on obtient toute autre base rationnelle par une transformation birationnelle; et:

Si
\[
  h_1(x)\cdots h_\rho(x)
\]
est un système quelconque de $\rho$ fonctions algébriquement indépendantes de $\Kfield_{n\rho}$, alors $\Kfield_{n\rho}$ est un corps algébrique fini sur
\[
  \Omega\bigl(h_1(x)\cdots h_\rho(x)\bigr).
\]

\section*{\S\ 5. Forme d'involution et base d'involution des corps $\Kfield_{n\rho}$.}

Alors que toute base rationnelle construite au \S\ 4 avait l'inconvénient de dépendre des fonctions initiales choisies arbitrairement, nous montrons maintenant l'existence d'une \emph{base rationnelle déterminée uniquement par le corps}.

D'après le \S\ 4, $\Omega(x_1\cdots x_\rho)$ est une extension algébrique --- disons de degré $i$ --- de $\Kfield^*_{\rho\rho}$ et, comme
\[
  \Kfield^*_{\rho\rho}(x_1,\ldots,x_\rho)=\Omega(x_1,\ldots,x_\rho),
\]
elle naît par adjonction des éléments $x_1,\ldots,x_\rho$, qui sont algébriques relativement à $\Kfield^*_{\rho\rho}$. Nous considérons donc la fonction entière $\Phi^*(z,u)$ irréductible relativement à $\Kfield^*_{\rho\rho}$, ayant pour zéro
\[
  z=u_1x_1+u_2x_2+\cdots+u_\rho x_\rho=u(x),
\]
et qui est de degré $i$ en $z$. De plus, comme produit des grandeurs conjuguées à $[z-u(x)]$, $\Phi^*(z,u)$ est une forme homogène de dimension $i$ en $z,u_1,\ldots,u_\rho$; elle contient donc, comme coefficient de $z^i$, une fonction de $\Kfield^*_{\rho\rho}$ indépendante de $u$. Si nous rendons ce coefficient le plus élevé égal à $1$ par division, alors la fonction irréductible $\Phi^*(z,u)$ est déterminée de manière unique. La forme homogène ainsi définie,
\begin{equation}
\begin{aligned}
  \Phi^*(z,u)
  &=z^i+
    \sum{}' g^*_{\alpha\alpha_1\ldots\alpha_\rho}(x_1,
    \ldots,x_\rho)\cdot
    z^\alpha u_1^{\alpha_1}\cdots u_\rho^{\alpha_\rho},\\
  &\hspace{8em}(\alpha+\alpha_1+\cdots+\alpha_\rho=i),
\end{aligned}
\tag{1}
\end{equation}
doit être appelée la \emph{forme d'involution}\footnote{La raison de ce nom deviendra claire à la fin du paragraphe; $\sum'$ signifie que la sommation doit s'étendre à tous les termes sauf $z^i$.} de $\Kfield^*_{\rho\rho}$; de (1) on obtient
\begin{equation}
\begin{aligned}
  \Phi(z,u)
  &=z^i+
    \sum{}' g_{\alpha\alpha_1\ldots\alpha_\rho}(x)\cdot
    z^\alpha u_1^{\alpha_1}\cdots u_\rho^{\alpha_\rho},\\
  &\hspace{8em}(\alpha+\alpha_1+\cdots+\alpha_\rho=i),
\end{aligned}
\tag{2}
\end{equation}
comme \emph{forme d'involution} de $\Kfield_{n\rho}$.

Les coefficients de la forme d'involution $\Phi^*(z,u)$ se révéleront être une base rationnelle de $\Kfield^*_{\rho\rho}$; c'est-à-dire que vaut la relation
\begin{equation}
  \Omega\bigl(g^*_{\alpha\alpha_1\ldots\alpha_\rho}(x_1,
  \ldots,x_\rho)\bigr)=\Kfield^*_{\rho\rho}
  \quad(\alpha+\alpha_1+\cdots+\alpha_\rho=i).
\tag{3}
\end{equation}
Pour la preuve nous observons que $\Phi^*(z,u)$ a des coefficients dans $\Omega(g^*_{\alpha\alpha_1\ldots\alpha_\rho}(x))$ et, à cause de son irréductibilité relativement à $\Kfield^*_{\rho\rho}$, l'est à plus forte raison relativement à ce sous-corps. Donc $\Phi^*(z,u)$ donne l'équation irréductible de $u(x)$ relativement à $\Omega(g^*_{\alpha\alpha_1\ldots\alpha_\rho}(x))$. Il en résulte que $\Omega(x_1\cdots x_\rho)$ est une extension algébrique --- de degré $i$ --- de $\Omega(g^*_{\alpha\alpha_1\ldots\alpha_\rho}(x))$. Comme $\Kfield^*_{\rho\rho}$ se trouve entre $\Omega(g^*_{\alpha\alpha_1\ldots\alpha_\rho}(x))$ et $\Omega(x_1\cdots x_\rho)$, l'identité des degrés entraîne, comme on sait, l'identité des corps, donc la relation (3). Par le principe de transfert, les coefficients de $\Phi(z,u)$ fournissent de même une base rationnelle de $\Kfield_{n\rho}$; ainsi nous avons

\emph{Théorème III: Les coefficients de la forme d'involution constituent une base rationnelle déterminée par le corps, la base d'involution; pour les corps $\Kfield_{\rho\rho}$ cette base est déterminée de manière unique.}\footnote{Pour les corps $\Kfield_{n\rho}$, la forme d'involution peut dépendre de l'application; l'unicité peut être obtenue par une application à coefficients indéterminés.}

Nous ajoutons une interprétation irrationnelle de ce fait, qui établit le lien avec la théorie géométrique des involutions.

La décomposition de $\Phi^*(z,u)$ en facteurs linéaires est donnée, dans un corps d'extension, par
\begin{equation}
\begin{aligned}
 \Phi^*(z,u)&=(z-u_1x_1-\cdots-u_\rho x_\rho)
 (z-u_1x_1^{(1)}-\cdots-u_\rho x_\rho^{(1)})\cdots\\
 &\qquad\cdot
 (z-u_1x_1^{(i-1)}-\cdots-u_\rho x_\rho^{(i-1)}).
\end{aligned}
\tag{4}
\end{equation}
La comparaison avec (1) montre que les fonctions
\[
  g^*_{\alpha\alpha_1\ldots\alpha_\rho}(x_1,
  \ldots,x_\rho)\quad(\alpha+\alpha_1+\cdots+\alpha_\rho=i)
\]
sont identiques aux fonctions symétriques élémentaires des lignes de grandeurs
\begin{equation}
\begin{array}{cccc}
  x_1 & x_2 & \cdots & x_\rho\\
  x_1^{(1)} & x_2^{(1)} & \cdots & x_\rho^{(1)}\\
  \vdots & \vdots & & \vdots\\
  x_1^{(i-1)} & x_2^{(i-1)} & \cdots & x_\rho^{(i-1)}.
\end{array}
\tag{5}
\end{equation}
Ainsi toute fonction de $\Kfield^*_{\rho\rho}$ est une combinaison rationnelle de ces fonctions élémentaires, symétriques dans les lignes (5), et réciproquement toute fonction symétrique des lignes (5) appartient à $\Kfield^*_{\rho\rho}$.

Ainsi, lorsque
\[
  f^*(x)=\frac{F^*(x)}{H^*(x)}
\]
parcourt toutes les fonctions de $\Kfield^*_{\rho\rho}$, il existe un système d'une infinité de relations:
\begin{equation}
  \frac{F^*(x)}{H^*(x)}=
  \frac{F^*(x^{(1)})}{H^*(x^{(1)})}=\cdots=
  \frac{F^*(x^{(i-1)})}{H^*(x^{(i-1)})},
\tag{6}
\end{equation}
ou, sous forme condensée,
\[
  F^*(x^{(k)})H^*(x^{(l)})-
  F^*(x^{(l)})H^*(x^{(k)})=0
  \quad(k,l=0,1,\cdots,i-1),
\]
où ni $F^*(x^{(k)})$ ni $H^*(x^{(k)})$ ne s'annule identiquement, puisqu'ils sont conjugués de $F^*(x)$ et $H^*(x)$ et formellement symétriques.

Réciproquement, pour toute ligne de grandeurs $\xi$ qui satisfait aux relations, en nombre infini,
\begin{equation}
  F^*(x)H^*(\xi)-H^*(x)F^*(\xi)=0,
  \qquad H^*(\xi)\ne0,
\tag{7}
\end{equation}
on obtient son appartenance aux lignes de grandeurs (5).

En effet, si $G^*(x)$ désigne le dénominateur commun des coefficients de $\Phi^*(z,u)$, alors, par l'hypothèse (7), la fonction
\[
  G^*(\xi)\cdot\Phi^*(z,u;\xi),
\]
qui est également entière en $\xi$, ne s'annule pas identiquement --- c'est-à-dire que les coefficients de $\Phi^*(z,u;\xi)$ ne deviennent pas indéterminés --- et elle a pour zéro
\[
  z=u_1\xi_1+u_2\xi_2+\cdots+u_\rho\xi_\rho.
\]
À cause de (7), $\Phi^*(z,u;x)$ a donc aussi le zéro $z=u(\xi)$; c'est-à-dire que $\xi$ appartient aux lignes de grandeurs (5). L'assertion correspondante vaut, d'après (6), si $\xi$ satisfait à un système d'équations relativement à une ligne de grandeurs conjuguée:
\[
  F^*(x^{(k)})H^*(\xi)-H^*(x^{(k)})F^*(\xi)=0;
  \qquad H(\xi)\ne0.
\]
Ainsi les lignes de grandeurs (5) sont définies comme les solutions $\xi$ du système d'équations
\[
  F^*(x)H^*(\xi)-H^*(x)F^*(\xi)=0;
  \qquad H(\xi)\ne0,
\]
où $F^*(x):H^*(x)$ parcourt toutes les fonctions de $\Kfield^*_{\rho\rho}$; ici la ligne $x$ peut aussi être remplacée par n'importe quelle ligne conjuguée.

Géométriquement, cela signifie ce qui suit, si l'on interprète chaque ligne $x$ comme un point: les solutions de (7) représentent, dans un espace linéaire de dimension $\rho$, $\infty^\rho$ groupes de $i$ points chacun, de telle sorte que chaque point du groupe détermine le groupe individuel de manière unique. Un tel système de groupes de points est appelé une \og{}\emph{involution dans un espace linéaire}\fg{}.\footnote{Les solutions exclues de (7), pour lesquelles $F^*(\xi)=0$, $H^*(\xi)=0$, ainsi que les solutions exclues du système d'équations correspondant à la base d'involution, s'appellent points fondamentaux de l'involution; on y compte aussi ceux qui proviennent de l'homogénéisation, les points \og{}situés à l'infini\fg{}.}

De façon correspondante, un corps $\Kfield_{n\rho}$ caractérise une involution, constituée de courbes, de surfaces, etc., dans un espace linéaire.

D'après ce qui précède, le degré $i$ de $\Omega(x_1\cdots x_\rho)$ relativement à $\Kfield^*_{\rho\rho}$ est égal au nombre des systèmes de solutions de (7) satisfaisant à la condition auxiliaire $H^*(\xi)\ne0$ --- que nous pouvons appeler les systèmes de solutions de $\Kfield^*_{\rho\rho}$. Par conséquent, le raisonnement utilisé pour démontrer l'existence de la base d'involution permet aussi la formulation irrationnelle suivante:

\og{}Si $\mathfrak{L}^*$ est un diviseur de $\Kfield^*_{\rho\rho}$ et si chaque système de solutions du corps $\mathfrak{L}^*$ est aussi un système de solutions de $\Kfield^*_{\rho\rho}$, alors $\mathfrak{L}^*$ et $\Kfield^*_{\rho\rho}$ sont identiques.\fg{}

L'énoncé correspondant vaut, d'après le principe de transfert, pour les diviseurs $\mathfrak{L}$ de $\Kfield_{n\rho}$.

\section*{\S\ 6. Base minimale des corps $\Kfield_{n\rho}$.}

Ce paragraphe donne une vue d'ensemble de théorèmes connus, mais sous une forme élargie rendue possible par le principe de transfert et par l'existence de la base rationnelle.

\emph{Définition IV: Une base rationnelle de $\Kfield_{n\rho}$ est appelée base minimale si elle ne contient que le nombre minimal $\rho$ de fonctions, qui doivent alors être algébriquement indépendantes.}

Des théorèmes de Lüroth, Castelnuovo et Enriques\footnote{J. Lüroth: \emph{Beweis eines Satzes über rationale Kurven}, \emph{Math. Ann.} 9 (1875). --- G. Castelnuovo: \emph{Sulla razionalità delle involuzioni piane}, \emph{Math. Ann.} 44 (1893). --- F. Enriques: \emph{Sopra una involuzione non razionale dello spazio}, \emph{Rend. Acc. Linc.}, Vol. 21, 21 janv. 1912.} résultent les faits suivants:

I. Tout corps $\Kfield_{n1}$ possède une base minimale, la fonction de Lüroth du corps, déterminée de manière unique à une transformation linéaire fractionnaire près.

II. Pour les corps $\Kfield_{n2}$, une base minimale existe toujours, et en général seulement, si l'on permet une extension algébrique du corps des coefficients $\Omega$.

III. Pour les corps $\Kfield_{n\rho}$ ($\rho\ge3$), il n'existe en général aucune base minimale. Une caractérisation des corps spéciaux $\Kfield_{n\rho}$ possédant une base minimale n'a pas encore été tentée.

Nous indiquons brièvement la démonstration du théorème de Lüroth pour les corps $\Kfield^*_{11}$ donnée par E. Steinitz:\footnote{Loc. cit., \S\ 24.}

Soit $\Phi^*(z,u)$ la forme d'involution de $\Kfield^*_{11}$, et soit $\psi^*(x)$ un coefficient quelconque de $\Phi^*(z,u)$ qui contient effectivement l'indéterminée $x$. Il apparaît que le degré de $\Omega(x)$ relativement à $\Omega(\psi^*(x))$ est égal au degré de $\Omega(x)$ relativement à $\Kfield^*_{11}$, d'où résulte l'identité des corps. En outre, si $\Omega(\chi^*(x))$ est égal à $\Omega(\psi^*(x))$, alors la dépendance rationnelle réciproque de $\chi^*(x)$ et $\psi^*(x)$ implique la dépendance linéaire fractionnaire des deux fonctions.

Par le principe de transfert, le théorème de Lüroth admet donc aussi la formulation:

\emph{La base minimale des corps $\Kfield_{n1}$ consiste en une fonction quelconque de la base d'involution; et deux fonctions quelconques de la base d'involution de $\Kfield_{n1}$ dépendent l'une de l'autre par une transformation linéaire fractionnaire.}

G. Castelnuovo démontre le fait II, au moyen de la théorie géométrique de l'involution, pour les corps $\Kfield^*_{22}$ dérivés d'une base rationnelle de trois fonctions. Comme le principe de transfert reste aussi valable sous une extension algébrique finie du corps des coefficients $\Omega$, l'existence de la base rationnelle donne ainsi la preuve pour tous les corps $\Kfield_{n2}$.

À propos de III, il faut remarquer qu'Enriques construit une involution non rationnelle dans l'espace tridimensionnel; c'est-à-dire un corps $\Kfield_{33}$ qui ne possède pas de base minimale.

\clearpage
\section*{\S\ 7. Système arbitraire $\Ssys_{n\rho}$, famille linéaire $\Lsys_{n\rho}$ et domaine d’intégrité $\Jdom_{n\rho}$ de fonctions rationnelles. Existence de la base rationnelle.}

De l'existence de la base rationnelle des corps $\Kfield_{n\rho}$ nous déduirons maintenant la base rationnelle pour des systèmes arbitraires de fonctions rationnelles et de fonctions rationnelles entières. Nous rangeons ces systèmes dans l'ordre selon lequel ils naissent successivement les uns des autres par les opérations élémentaires d'addition, de multiplication et de division.

I. Comme toujours, $\Ssys_{n\rho}$ désigne un système arbitraire de fonctions rationnelles (ou rationnelles entières) de $n$ indéterminées, de rang algébrique $\rho$. Les grandeurs provenant de $\Omega$ sont comptées comme appartenant au système.

II. Le système $\Lsys_{n\rho}$ s'appelle une \emph{famille linéaire} s'il satisfait aux conditions suivantes:

1. avec $f(x)$, $\Lsys_{n\rho}$ contient toujours aussi $c\cdot f(x)$, où $c$ est une grandeur arbitraire de $\Omega$;

2. avec $f(x)$ et $g(x)$, $\Lsys_{n\rho}$ contient toujours aussi $f(x)+g(x)$.

III. La famille linéaire $\Jdom_{n\rho}$ s'appelle un \emph{domaine d’intégrité} sous la condition supplémentaire:

3. avec $f(x)$ et $g(x)$, $\Jdom_{n\rho}$ contient toujours aussi $f(x)\cdot g(x)$.

IV. Le domaine d’intégrité $\Kfield_{n\rho}$ s'appelle un \emph{corps} sous la condition supplémentaire:

4. avec $f(x)$ et $g(x)$ --- où $g(x)\ne0$ ---, $\Kfield_{n\rho}$ contient toujours aussi le quotient $f(x):g(x)$.

Les conditions 1 à 4 sont celles indiquées au \S\ 1 pour le corps.\footnote{$f(x)$ et $g(x)$ peuvent aussi être de degré zéro, c'est-à-dire être des grandeurs de $\Omega$.}

Ces domaines doivent maintenant être déduits successivement les uns des autres.

a) Un système arbitraire $\Ssys_{n\rho}$ est élargi en une famille linéaire $\Lsys_{n\rho}$ par l'exigence suivante:

$\Lsys_{n\rho}$ contient, en plus de $\Ssys_{n\rho}$, toutes les fonctions $h(x)$ et seulement celles qui peuvent être représentées linéairement, avec des coefficients de $\Omega$, au moyen d'un nombre fini de fonctions de $\Ssys_{n\rho}$:
\[
  h(x)=c_1f_1(x)+c_2f_2(x)+\cdots+c_kf_k(x),
\]
où $f_1(x)\cdots f_k(x)$ parcourent toutes les fonctions de $\Ssys_{n\rho}$, et $c_1,
\ldots,c_k$ toutes les grandeurs de $\Omega$.

En effet, l'adjonction de combinaisons rationnelles des fonctions du système conserve le rang algébrique d'un système. $\Lsys_{n\rho}$ satisfait aux conditions 1 et 2, et est donc une famille linéaire. En outre, toute famille linéaire contenant $\Ssys_{n\rho}$ contient aussi $\Lsys_{n\rho}$ d'après 1 et 2; ainsi $\Lsys_{n\rho}$ est la plus petite famille linéaire contenant $\Ssys_{n\rho}$.

b) De même, d'une famille linéaire $\Lsys_{n\rho}$ naît le plus petit domaine d’intégrité $\Jdom_{n\rho}$ qui la contient, par l'exigence suivante:

$\Jdom_{n\rho}$ contient, en plus de $\Lsys_{n\rho}$, toutes les fonctions $h(x)$ et seulement celles qui peuvent être représentées comme \emph{combinaison rationnelle entière} --- avec des coefficients de $\Omega$ --- d'un nombre fini de fonctions de $\Lsys_{n\rho}$, $f_1(x)\cdots f_k(x)$, où $f_1(x)\cdots f_k(x)$ parcourent toutes les fonctions de $\Lsys_{n\rho}$.

c) Enfin, d'un domaine d’intégrité $\Jdom_{n\rho}$ naît le plus petit corps $\Kfield_{n\rho}$ qui le contient, par l'exigence suivante: $\Kfield_{n\rho}$ contient, en plus de $\Jdom_{n\rho}$, toutes les fonctions $h(x)$ et seulement celles qui peuvent être représentées comme quotients de deux fonctions de $\Jdom_{n\rho}$.

Si l'on réunit les trois élargissements en une seule étape, on obtient:

\emph{Tout système $\Ssys_{n\rho}$ peut être élargi au plus petit corps $\Kfield_{n\rho}$ qui le contient par l'exigence suivante:}

\emph{$\Kfield_{n\rho}$ contient, en plus de $\Ssys_{n\rho}$, toutes les fonctions $h(x)$ et seulement celles qui peuvent être représentées comme combinaisons rationnelles --- avec des coefficients de $\Omega$ --- d'un nombre fini de fonctions de $\Ssys_{n\rho}$, $f_1(x)\cdots f_k(x)$, où $f_1(x)\cdots f_k(x)$ parcourent toutes les fonctions de $\Ssys_{n\rho}$.}

Ce dernier fait implique immédiatement l'existence de la base rationnelle pour des domaines arbitraires:

Soit $\Kfield_{n\rho}$ le plus petit corps contenant $\Ssys_{n\rho}$, et soit une base rationnelle de $\Kfield_{n\rho}$ donnée par
\begin{equation}
  h_1(x),\ h_2(x),\ldots,h_\sigma(x).
\tag{1}
\end{equation}
D'après la définition de $\Kfield_{n\rho}$, il existe des relations:
\begin{equation}
\begin{aligned}
 h_1(x)&=\gamma_1\bigl(g_1(x),\ldots,g_\lambda(x)\bigr);\ldots;\\
 h_\sigma(x)&=\gamma_\sigma\bigl(g_\mu(x),\ldots,g_\nu(x)\bigr),
\end{aligned}
\tag{2}
\end{equation}
où $r_1,
\ldots,r_\sigma$ sont des fonctions rationnelles de leurs arguments, et
\begin{equation}
  g_1(x),\ldots,g_\lambda(x),\ldots,
  g_\mu(x),\ldots,g_\nu(x)
\tag{3}
\end{equation}
appartiennent toutes à $\Ssys_{n\rho}$. En vertu des relations (2), (3) représente, en plus de (1), également une base rationnelle de $\Kfield_{n\rho}$; puisque les fonctions (3) appartiennent à $\Ssys_{n\rho}$, c'est en même temps une base rationnelle de $\Ssys_{n\rho}$. Nous avons donc

\emph{Théorème IV: Un système arbitraire $\Ssys_{n\rho}$ de fonctions rationnelles (ou rationnelles entières) possède une base rationnelle, constituée d'un nombre fini de fonctions du système, de rang algébrique $\rho$.}

Soit maintenant en particulier $\Lsys_{n\rho}$ une famille linéaire, et soit
\begin{equation}
  g_1(x),\ldots,g_\rho(x),\ g_{\rho+1}(x),\ldots,g_\nu(x)
\tag{4}
\end{equation}
une base rationnelle de $\Lsys_{n\rho}$. Supposons, par exemple, que
\[
  g_1(x)\cdots g_\rho(x)
\]
soient algébriquement indépendantes. Si l'on pose alors
\[
  g(x)=c_1g_{\rho+1}(x)+c_2g_{\rho+2}(x)+\cdots+c_{\nu-\rho}g_\nu(x),
\]
on peut choisir les $c_i$ comme nombres de $\Omega$ de telle sorte que
\begin{equation}
  g_1(x),\ldots,g_\rho(x),\ g(x)
\tag{5}
\end{equation}
devienne une base rationnelle du plus petit corps contenant $\Lsys_{n\rho}$, à savoir $\Kfield_{n\rho}$. Mais puisque $\Lsys_{n\rho}$ est une famille linéaire, $g(x)$ appartient à $\Lsys_{n\rho}$, et (5) représente une base rationnelle de $\Lsys_{n\rho}$; c'est-à-dire:

\emph{Théorème V: La base rationnelle de toute famille linéaire $\Lsys_{n\rho}$ de fonctions rationnelles (ou rationnelles entières) --- donc en particulier la base rationnelle de tout domaine d’intégrité $\Jdom_{n\rho}$ et de tout corps $\Kfield_{n\rho}$ --- peut notamment être choisie de façon qu'elle contienne au plus $\rho+1$ fonctions (et au moins $\rho$).}

La borne trouvée au \S\ 4 pour le nombre des fonctions de la base rationnelle des corps repose donc sur la propriété qu'a le corps d'être une famille linéaire, tandis que la restriction à $\rho$ fonctions dans le cas de l'existence de la base minimale utilise toutes les hypothèses du corps.

Il reste à remarquer que, d'après \S\ 3 (4), toutes les propriétés caractéristiques des systèmes et des plus petits systèmes qui les contiennent se conservent par l'application.

\section*{\S\ 8. Base d'involution des domaines d’intégrité $\Jdom_{n\rho}$ formés de polynômes.}

Dans les paragraphes suivants, il s'agit de systèmes $\Ssys_{n\rho}$ dont les éléments sont des fonctions rationnelles entières (polynômes); d'abord de domaines d’intégrité $\Jdom_{n\rho}$ formés de polynômes. Pour ces domaines $\Jdom_{n\rho}$, il faut fixer d'abord les notions de divisibilité.

Soient $F(x),G(x)$ deux polynômes de $\Jdom_{n\rho}$, et soit $L(x)$ un diviseur commun de ces polynômes, de sorte que l'on ait
\[
  F(x)=L(x)\cdot F_1(x);\qquad G(x)=L(x)\cdot G_1(x).
\]
$L(x)$ s'appelle un \emph{diviseur commun dans $\Jdom_{n\rho}$} si et seulement si les trois polynômes $L(x),F_1(x),G_1(x)$ appartiennent à $\Jdom_{n\rho}$.

Deux polynômes de $\Jdom_{n\rho}$ sont dits \emph{premiers entre eux dans $\Jdom_{n\rho}$} s'ils ne possèdent aucun diviseur commun dans $\Jdom_{n\rho}$.

Soit $\Kfield_{n\rho}$ le plus petit corps contenant $\Jdom_{n\rho}$. Alors, d'après la définition, toute fonction de $\Kfield_{n\rho}$ possède une représentation
\[
  f(x)=F_1(x):F_2(x),
\]
où $F_1(x),F_2(x)$ appartiennent à $\Jdom_{n\rho}$ et sont premiers entre eux dans $\Jdom_{n\rho}$. De même, plusieurs fonctions de $\Kfield_{n\rho}$ peuvent être réduites à un dénominateur commun dans $\Jdom_{n\rho}$:
\[
  f_1(x)=\frac{F_1(x)}{F(x)},\ldots,
  f_k(x)=\frac{F_k(x)}{F(x)}.
\]
$F(x)$ s'appelle \emph{un plus petit dénominateur commun dans $\Jdom_{n\rho}$} si et seulement si les polynômes de $\Jdom_{n\rho}$,
\[
  F(x),F_1(x),\ldots,F_k(x),
\]
sont premiers entre eux dans $\Jdom_{n\rho}$.

Une telle représentation peut être possible de plusieurs manières, car les lois de décomposition unique en fonctions irréductibles ne valent pas en général dans $\Jdom_{n\rho}$.

Nous examinons maintenant la signification de la base d'involution de $\Kfield_{n\rho}$ pour $\Jdom_{n\rho}$, en transférant aux domaines d’intégrité les notions de forme d'involution et de base d'involution. Soit
\[
  \Phi(z,u)=z^i+\sum{}' g_{\alpha\alpha_1\ldots\alpha_\rho}(x)
  z^\alpha u_1^{\alpha_1}\cdots u_\rho^{\alpha_\rho},
  \qquad
  (\alpha+\alpha_1+\cdots+\alpha_\rho=i)
\]
une forme d'involution de $\Kfield_{n\rho}$, et soit $G(x)$ un plus petit dénominateur commun dans $\Jdom_{n\rho}$ des coefficients de $\Phi(z,u)$. Par multiplication par $G(x)$, $\Phi(z,u)$ passe en une forme $\Psi(z,u)$ dont les coefficients sont des polynômes de $\Jdom_{n\rho}$ et sont premiers entre eux dans $\Jdom_{n\rho}$:
\[
  G(x)\cdot\Phi(z,u)=\Psi(z,u)
  =G(x)z^i+
  \sum{}' G_{\alpha\alpha_1\ldots\alpha_\rho}(x)
  z^\alpha u_1^{\alpha_1}\cdots u_\rho^{\alpha_\rho},
  \qquad
  (\alpha+\alpha_1+\cdots+\alpha_o=i).
\]

Toute forme $\Psi(z,u)$ qui provient d'une forme d'involution $\Phi(z,u)$ de $\Kfield_{n\rho}$ par multiplication par un polynôme de $\Jdom_{n\rho}$, de telle sorte que les coefficients de $\Psi(z,u)$ soient des polynômes de $\Jdom_{n\rho}$ premiers entre eux dans $\Jdom_{n\rho}$, sera appelée une \emph{forme d'involution} de $\Jdom_{n\rho}$; et les coefficients de $\Psi(z,u)$ seront appelés une \emph{base d'involution} correspondante.

Il est d'abord clair, par la signification de la base d'involution de $\Kfield_{n\rho}$, que toute base d'involution de $\Jdom_{n\rho}$ est en même temps une base rationnelle. Pour la suite de l'étude, considérons un domaine image $\Jdom^*_{\rho\rho}$ dont le plus petit corps contenant $\Kfield^*_{\rho\rho}$ possède comme équation irréductible de zéro
\[
  z=u_1x_1+\cdots+u_\rho x_\rho
\]
l'équation $\Phi^*(z,u)=0$.

D'après \S\ 5, les coefficients de
\[
  \Phi^*(z,u)=z^i+\sum{}' g^*_{\alpha\alpha_1\ldots\alpha_\rho}(x)
  z^\alpha u_1^{\alpha_1}\cdots u_\rho^{\alpha_\rho},
  \qquad
  (\alpha+\alpha_1+\cdots+\alpha_\rho=i)
\]
représentent les fonctions symétriques élémentaires des lignes de grandeurs conjuguées relativement à $\Kfield^*_{\rho\rho}$,
\[
  x,\ x^{(1)}\cdots x^{(i-1)}.
\]
Soit $F^*(x)$ un polynôme quelconque de $\Kfield^*_{\rho\rho}$. Comme
\[
  F^*(x)=\frac1i\Bigl(F^*(x)+F^*(x^{(1)})+
  \cdots+F^*(x^{(i-1)})\Bigr),
\]
il s'ensuit que $F^*(x)$ est une fonction symétrique entière de ces lignes de grandeurs; par conséquent il peut s'exprimer comme combinaison rationnelle entière des fonctions symétriques élémentaires, c'est-à-dire des fonctions
\[
  g^*_{\alpha\alpha_1\ldots\alpha_\rho}(x).
\]
Si maintenant
\[
  G^*(x)\cdot\Phi^*(z,u)=\Psi^*(z,u)
  =G^*(x)z^i+
  \sum{}'G^*_{\alpha\alpha_1\ldots\alpha_\rho}(x)
  z^\alpha u_1^{\alpha_1}\cdots u_\rho^{\alpha_\rho},
  \qquad
  (\alpha+\alpha_1+\cdots+\alpha_\rho=i)
\]
est une forme d'involution de $\Jdom^*_{\rho\rho}$, alors tout polynôme de $\Kfield^*_{\rho\rho}$, et en particulier tout polynôme de $\Jdom^*_{\rho\rho}$, est une combinaison rationnelle entière des quotients
\[
  G^*_{\alpha\alpha_1\ldots\alpha_\rho}(x):G^*(x).
\]

En tenant compte du principe de transfert, on obtient:

\emph{Théorème VI: À tout domaine d’intégrité de polynômes $\Jdom_{n\rho}$ il existe au moins une base rationnelle distinguée, la base d'involution, telle que dans la représentation rationnelle de tous les polynômes de $\Jdom_{n\rho}$ par les fonctions de la base d'involution n'apparaisse au dénominateur qu'une puissance d'une fonction fixe (le coefficient dominant de la forme d'involution correspondante).}

\section*{\S\ 9. Domaines relativement entiers $\Gdom_{n\rho}$ formés de polynômes.}

Dans ce qui suit, nous considérons des domaines d’intégrité spéciaux formés de polynômes, qui constituent un stade intermédiaire entre les domaines d’intégrité et les corps.

Un domaine d’intégrité de polynômes $\Gdom_{n\rho}$ s'appelle un \og{} domaine relativement entier \fg{} sous la condition:

\emph{4a) $\Gdom_{n\rho}$ contient, avec $F(x)$ et $G(x)$ --- où $G(x)\ne0$ ---, le quotient $F(x):G(x)$ toujours et seulement lorsque ce quotient est entier dans les indéterminées, c'est-à-dire lorsqu'il est un polynôme.}

Il faut d'abord démontrer l'identité des domaines relativement entiers avec la totalité des polynômes d'un corps $\Kfield_{n\sigma}$ ($\sigma\ge \rho$) de fonctions rationnelles.

Soit $\Kfield_{n\rho}$ le plus petit corps contenant $\Gdom_{n\rho}$. Pour toute fonction $f(x)$ de $\Kfield_{n\rho}$ il existe une représentation comme quotient de deux polynômes de $\Gdom_{n\rho}$:
\[
  f(x)=F_1(x):F_2(x),
\]
où les polynômes de degré zéro sont également admis. Dès que $f(x)$ devient un polynôme, il appartient à $\Gdom_{n\rho}$ d'après la condition 4a); et puisque $\Gdom_{n\rho}$ ne peut contenir aucun autre polynôme, \emph{tout domaine relativement entier $\Gdom_{n\rho}$ est identique à la totalité des polynômes du plus petit corps qui le contient.}

Réciproquement, soit $\Kfield_{n\sigma}$ un corps quelconque de fonctions rationnelles (donc pas nécessairement le plus petit corps contenant un système donné), et soit $\Jdom$ la totalité des polynômes contenus dans $\Kfield_{n\rho}$. $\Jdom$ satisfait aux conditions 1, 2, 3 d'un domaine d’intégrité et à la condition 4a), et il est donc un domaine relativement entier; c'est-à-dire que \emph{la totalité des polynômes contenus dans un corps $\Kfield_{n\sigma}$ forme un domaine relativement entier $\Gdom_{n\rho}$ ($\rho\le\sigma$).}\footnote{Il peut aussi arriver que $\rho=0$; c'est-à-dire que la totalité des polynômes de $\Kfield_{n\sigma}$ consiste en éléments du domaine des coefficients $\Omega$.}

Nous montrons encore l'identité des domaines relativement entiers avec les \og{} domaines d’intégrité de fonctions relativement entières \fg{} de Hilbert.\footnote{D. Hilbert: \emph{Mathematische Probleme}. Conférence de Paris 1900. Problème 14 (\emph{Göttinger Nachrichten} 1900).}

Soit
\begin{equation}
  G_1(x),\ldots,G_k(x)
\tag{1}
\end{equation}
une base rationnelle de $\Gdom_{n\rho}$. D'après les conditions 1, 2, 3, 4a, toute combinaison rationnelle des polynômes (1) qui est entière dans les indéterminées, c'est-à-dire qui devient un polynôme, appartient à $\Gdom_{n\rho}$. Réciproquement, tout polynôme de $\Gdom_{n\rho}$ peut, d'après la notion de base rationnelle, s'exprimer rationnellement au moyen des polynômes (1); ainsi $\Gdom_{n\rho}$ est \emph{identique à la totalité de ces combinaisons rationnelles des polynômes (1) qui sont entières dans les indéterminées, c'est-à-dire qui sont des polynômes.} Nous avons donc la définition de Hilbert à partir de la base rationnelle spéciale, tandis que notre définition n'emploie que les propriétés abstraites du domaine.

Pour les domaines relativement entiers, la définition d'un diviseur commun donnée au \S\ 8 pour les domaines d’intégrité généraux se simplifie:

Soient $F(x),G(x)$ deux polynômes de $\Gdom_{n\rho}$, et soit $L(x)$ un diviseur commun de ces polynômes. Alors $L(x)$ s'appelle un \emph{diviseur commun dans $\Gdom_{n\rho}$} si et seulement si $L(x)$ appartient à $\Gdom_{n\rho}$.

Car l'appartenance des quotients $F(x):L(x)$ et $G(x):L(x)$ résulte alors de la condition 4a).

Une autre propriété essentielle des domaines relativement entiers est que la notion d'être \og{} algébriquement entier relativement à $\Gdom_{n\rho}$ \fg{} peut être définie au moyen de n'importe quelle équation exactement comme on le sait pour les corps de nombres algébriques ou pour le corps $\Omega(x_1,
\ldots,x_n)$.

\emph{Définition V: Une grandeur $\xi$ s'appelle algébriquement entière relativement à $\Gdom_{n\rho}$ --- ou encore relativement algébriquement entière --- si elle satisfait à une équation dont le coefficient dominant est l'unité et dont les autres polynômes appartiennent à $\Gdom_{n\rho}$.}

Supposons en effet que la grandeur relativement algébriquement entière $\xi$ satisfasse par exemple à l'équation
\[
  L(z)=z^\lambda+G_1(x)z^{\lambda-1}+\cdots+G_\lambda(x)=0.
\]
$L(z)$ est divisible par la fonction entière irréductible relativement au plus petit corps contenant $\Kfield_{n\rho}$,
\[
  M(z)=z^\mu+H_1(x)z^{\mu-1}+\cdots+H_\mu(x).
\]
Mais de cette divisibilité il résulte, par l'extension connue du théorème de Gauss\footnote{Comparer par exemple J. König: \emph{Einleitung in die allgemeine Theorie der algebraischen Größen} (Leipzig, B.~G. Teubner, 1903), chap. IX, \S\ 2, ou Weber (petite édition), \S\ 20.}, que les fonctions rationnelles $H_1(x),\ldots,H_\mu(x)$ sont des polynômes de $\Gdom_{n\rho}$, et par conséquent appartiennent à $\Gdom_{n\rho}$. Ainsi la définition de la grandeur relativement algébriquement entière s'obtient à partir de l'équation irréductible. En particulier, tout polynôme $F(x)$ de $\Gdom_{n\rho}$ est aussi algébriquement entier relativement à $\Gdom_{n\rho}$, puisqu'il satisfait à l'équation $z-F(x)=0$.

Il faut encore mentionner que, pour un domaine d’intégrité $\Jdom_{n\rho}$ de polynômes, de façon analogue au plus petit corps contenant, le plus petit domaine relativement entier $\Gdom_{n\rho}$ qui le contient est défini par l'exigence suivante:

$\Gdom_{n\rho}$ contient, avec $\Jdom_{n\rho}$, tous les polynômes $H(x)$, et seulement ceux-là, qui peuvent être représentés comme quotients de deux polynômes de $\Jdom_{n\rho}$.

De cette définition il résulte encore: De toute forme d'involution et de toute base d'involution de $\Jdom_{n\rho}$, il naît une forme et une base correspondantes de $\Gdom_{n\rho}$ en séparant tout au plus un diviseur commun dans $\Gdom_{n\rho}$ de tous les polynômes de la base.

En effet, $\Jdom_{n\rho}$ et $\Gdom_{n\rho}$ possèdent le même plus petit corps contenant $\Kfield_{n\rho}$; et toute forme d'involution de $\Jdom_{n\rho}$ provient d'une forme d'involution de $\Kfield_{n\rho}$ par multiplication par un polynôme de $\Jdom_{n\rho}$, et possède donc des coefficients dans $\Gdom_{n\rho}$, lesquels peuvent toutefois avoir un diviseur commun dans $\Gdom_{n\rho}$.\footnote{Par exemple, pour le domaine d’intégrité $\Jdom_{22}$ formé de toutes les combinaisons rationnelles entières de $x_1^2$ et $x_1x_2$, on obtient comme forme d'involution:
\[
\Psi(z,u)=x_1^2\cdot z^2-x_1^4u_1^2-2x_1^3x_2u_1u_2-x_1^2x_2^2u_2^2.
\]
Comme $x_1^2$ appartient à $\Jdom_{22}$, et donc à plus forte raison au plus petit domaine relativement entier $\Gdom_{22}$ qui le contient, et qu'il est par conséquent un diviseur commun dans $\Gdom_{22}$, il en résulte comme forme d'involution de $\Gdom_{22}$:
\[
\Psi(z,u)=z^2-x_1^2u_1^2-2x_1x_2u_1u_2-x_2^2u_2^2.
\]}

Enfin il faut noter que le \emph{domaine image} d'un domaine relativement entier $\Gdom_{n\rho}$, qui d'après \S\ 7 est de nouveau un domaine d’intégrité $\Jdom^*_{\rho\rho}$ de polynômes, n'a pas besoin d'être lui-même un domaine relativement entier. Car si $F(x)$ et $G(x)$ sont deux polynômes de $\Gdom_{n\rho}$ dont le quotient n'est pas un polynôme et par conséquent n'appartient pas à $\Gdom_{n\rho}$, alors le quotient $F^*(x):G^*(x)$ n'appartient pas non plus à $\Jdom^*_{\rho\rho}$. Mais ce quotient, provenant de la spécialisation de quelques indéterminées, peut très bien être un polynôme, et pour $\Jdom^*_{\rho\rho}$ la condition 4a) n'est alors pas remplie.\footnote{Que $\Gdom_{32}$ consiste par exemple en tous les polynômes qui sont des combinaisons rationnelles de $F=x_1^2x_2$ et $G=x_1^2(1+x_3)$. La spécialisation admissible $x_3=0$ donne $F^*:G^*=x_2$, tandis que $F:G$ n'est pas un polynôme. $\Jdom^*_{22}$ n'est donc pas un domaine relativement entier.}
\clearpage
\section*{\S\ 10. Domaines relativement entiers de première espèce et leur base d'intégrité.}

La notion introduite au \S\ 9 (définition V), \og{} relativement algébriquement entière \fg{}, permet d'indiquer une classe de domaines relativement entiers qui sont des domaines d’intégrité finis, et donc de répondre affirmativement, au moins pour cette classe, à une question posée par D. Hilbert (\emph{Problèmes mathématiques}, problème 14). Nous donnons la définition suivante.

\emph{Définition VI: Un nombre fini de polynômes de $\Gdom_{n\rho}$ s'appelle une base d'intégrité --- au sens historique d'un système fini de générateurs algébriques de l'anneau, et non d'une base linéaire ou modulaire --- si tout polynôme de $\Gdom_{n\rho}$ peut être représenté comme combinaison rationnelle entière de ce nombre fini, avec des coefficients de $\Omega$. Un domaine d’intégrité de polynômes qui possède une base d'intégrité s'appelle un domaine d’intégrité fini.}

La classe de domaines relativement entiers à considérer, pour laquelle la base d'involution se révélera être une base d'intégrité, nous la désignons comme domaines relativement entiers de première espèce selon la définition suivante.

\emph{Définition VII: Un domaine relativement entier $\Gdom_{n\rho}$ s'appelle de première espèce s'il possède, comme domaine image, un domaine relativement entier $\Gdom^*_{\rho\rho}$ tel que tout polynôme arbitraire en $x_1,\ldots,x_\rho$ (à coefficients dans $\Omega$) soit algébriquement entier relativement à $\Gdom^*_{\rho\rho}$.}

D'après cette définition, les indéterminées $x_1,\ldots,x_\rho$, et par conséquent aussi la forme linéaire $u(x)=u_1x_1+\cdots+u_\rho x_\rho$, sont algébriquement entières relativement à $\Gdom^*_{\rho\rho}$. La fonction entière irréductible ayant pour zéro $z=u(x)$ --- la forme d'involution du plus petit corps contenant $\Kfield^*_{\rho\rho}$ --- possède donc l'unité comme coefficient dominant et, comme autres coefficients, des polynômes de $\Gdom^*_{\rho\rho}$; elle est donc en même temps forme d'involution de $\Gdom^*_{\rho\rho}$, et aucune autre forme d'involution ne peut exister dans $\Gdom^*_{\rho\rho}$. Par le principe de transfert, $\Gdom_{n\rho}$ possède donc aussi une forme d'involution dont le coefficient dominant est l'unité; et comme, d'après le théorème VI, dans la représentation rationnelle de tous les polynômes de $\Gdom_{n\rho}$ par la base d'involution correspondante seul ce coefficient dominant entre au dénominateur, la base d'involution devient une base d'intégrité. Ainsi:

\emph{Théorème VII. Tout domaine relativement entier de première espèce est un domaine d’intégrité fini; la base d'intégrité est donnée par la base d'involution obtenue à partir du domaine image qui le caractérise. En particulier, la base d'intégrité des domaines relativement entiers de première espèce $\Gdom_{\rho\rho}$ est donnée par la base d'involution, définie de manière unique, de $\Gdom_{\rho\rho}$.}

Nous ajoutons un critère pour les domaines relativement entiers de première espèce. Soit $\Gdom_{n\rho}$ de première espèce et soit une base d'intégrité de $\Gdom_{n\rho}$ donnée par
\begin{equation}
  F_1(x),F_2(x),\ldots,F_k(x).
\tag{1}
\end{equation}
Alors, par un lemme auxiliaire de Hilbert,\footnote{D. Hilbert: \emph{Über die vollen Invariantensysteme}, Math. Ann. 42 (1893), \S\ 1. Le lemme auxiliaire qui y est énoncé seulement pour des polynômes homogènes vaut de même pour les polynômes non homogènes.} il existe $\rho$ polynômes algébriquement indépendants de $\Gdom_{n\rho}$,
\begin{equation}
  G_1(x),G_2(x),\ldots,G_\rho(x),
\tag{2}
\end{equation}
tels que les polynômes (1), et par suite tout polynôme de $\Gdom_{n\rho}$, soient algébriquement entiers relativement aux polynômes (2). Il en va de même de tout polynôme du domaine image $\Gdom^*_{\rho\rho}$ relativement aux polynômes correspondants
\begin{equation}
  G_1^*(x),G_2^*(x),\ldots,G_\rho^*(x).
\tag{3}
\end{equation}

Par la définition VII, tout polynôme arbitraire en $x_1,\ldots,x_\rho$ est donc lui aussi algébriquement entier relativement aux polynômes (3).

Réciproquement, supposons qu'un domaine image $\Jdom^*_{\rho\rho}$ d'un domaine relativement entier quelconque $\Gdom_{n\rho}$ contienne $\rho$ polynômes (3) relativement auxquels tout polynôme en $x_1,\ldots,x_\rho$ soit algébriquement entier. Nous montrons alors que $\Jdom^*_{\rho\rho}$ devient par là un domaine relativement entier $\Gdom^*_{\rho\rho}$, et ensuite que $\Gdom^*_{\rho\rho}$, et par conséquent $\Gdom_{n\rho}$, sont de première espèce.

En effet, soit $F^*(x)$ un polynôme du plus petit corps $\Kfield^*_{\rho\rho}$ contenant $\Jdom^*_{\rho\rho}$. Par hypothèse, il est algébriquement entier relativement aux polynômes (3). Comme la fonction correspondante dans $\Kfield_{n\rho}$ est alors algébriquement entière relativement à $\rho$ polynômes de $\Gdom_{n\rho}$, elle est un polynôme $F(x)$ et appartient donc à $\Gdom_{n\rho}$; par conséquent $F^*(x)$ appartient à $\Jdom^*_{\rho\rho}$. Le domaine image $\Jdom^*_{\rho\rho}$ contient donc tous les polynômes de $\Kfield^*_{\rho\rho}$ et est un domaine relativement entier $\Gdom^*_{\rho\rho}$. Pour $\Gdom^*_{\rho\rho}$, et donc aussi pour $\Gdom_{n\rho}$, il résulte immédiatement des définitions V et VII et de l'hypothèse qu'ils sont de première espèce; c'est-à-dire:

\emph{Les domaines relativement entiers de première espèce $\Gdom_{n\rho}$ sont caractérisés par le fait que, dans un domaine image $\Jdom^*_{\rho\rho}$, il existe $\rho$ polynômes relativement auxquels tout polynôme arbitraire en $x_1,\ldots,x_\rho$ est algébriquement entier. De cette hypothèse il résulte que le domaine image lui-même est un domaine relativement entier.}\footnote{Par exemple, dans le domaine $\Gdom_{32}$ mentionné dans la remarque finale du \S\ 9, la spécialisation admissible $x_1=1$ donne les deux polynômes $F^*=x_2$, $G^*=1+x_3$, relativement auxquels tout polynôme arbitraire en $x_2,x_3$ est algébriquement entier. $\Gdom_{32}$ est donc de première espèce, et même avec $F(x)$ et $G(x)$ comme base d'intégrité.}

\section*{\S\ 11. Lemme auxiliaire sur la dépendance algébriquement entière.}

Nous donnons maintenant un critère pour que tout polynôme arbitraire en $x_1,\ldots,x_\rho$ soit algébriquement entier relativement à $\rho$ polynômes en $x_1,\ldots,x_\rho$,
\begin{equation}
  G_1^*(x),\ldots,G_\rho^*(x).
\tag{1}
\end{equation}

Si les polynômes (1) sont spécialement des formes homogènes, alors, pour tout système particulier de valeurs qui fait s'annuler (1), toutes les formes qui en dépendent algébriquement entière s'annulent simultanément; donc en particulier aussi $x_1,\ldots,x_\rho$.

Les formes (1) ne peuvent donc s'annuler pour aucun système de valeurs autre que
\[
  x_1=0,
  \ldots,
  x_\rho=0;
\]
leur résultante doit être différente de zéro.

Réciproquement, cette condition est aussi suffisante et admet une extension aux polynômes non homogènes selon le lemme auxiliaire suivant.

\emph{Lemme auxiliaire: Une condition suffisante pour que tout polynôme arbitraire en $x_1,\ldots,x_\rho$ soit algébriquement entier relativement à $\rho$ polynômes en $x_1,\ldots,x_\rho$}
\[
  \tag{1} G_1^*(x),\ldots,G_\rho^*(x)
\]
\emph{est le non-évanouissement de la résultante, prise par rapport à $x_1,\ldots,x_\rho$, des termes de plus haute dimension}
\[
  H_1^*(x),\ldots,H_\rho^*(x)
\]
\emph{des polynômes (1):}
\begin{equation}
  R_0=\left[
  \begin{array}{ccc}
    H_1^*&\cdots&H_\rho^*\\
    x_1&\cdots&x_\rho
  \end{array}
  \right]\ne0.
\tag{2}
\end{equation}
\emph{Pour les formes homogènes (1), la condition (2) est nécessaire.}\footnote{Pour la théorie des résultantes, comparer F. Mertens, \emph{Zur Theorie der Elimination} (I. et II. partie), Sitzungsber. d. k. Ak. d. Wiss. Wien, Math.-naturw. Kl. Abt. IIa, Bd. 108 (1899). Reproduit partiellement dans J. König, \emph{Theorie d. algebr. Größen}, chap. VI.}

Pour la démonstration, nous formons les polynômes
\begin{equation}
\begin{gathered}
  \bar G_1^*(x)=G_1^*(x)-\lambda_1,
  \ldots,
  \bar G_\rho^*(x)=G_\rho^*(x)-\lambda_\rho,\\
  \bar\Phi^*(x)=\Phi^*(x)-\mu,
\end{gathered}
\tag{3}
\end{equation}
où $\lambda_1,\ldots,\lambda_\rho,\mu$ désignent des indéterminées et où $\Phi^*(x)$ désigne un polynôme arbitraire en $x_1,\ldots,x_\rho$. Soit la résultante des polynômes (3) par rapport à $x_1,\ldots,x_\rho,1$\footnote{C'est-à-dire la résultante homogène par rapport à $x_1,\ldots,x_\rho,x_0$, lorsque les polynômes (3) sont homogénéisés au moyen d'une indéterminée supplémentaire $x_0$ de telle sorte que le degré en les $x$ soit conservé.} donnée par
\begin{equation}
  R(\lambda,\mu)=\left[
  \begin{array}{cccc}
   \bar G_1^*&\cdots&\bar G_\rho^*&\bar\Phi^*\\
   x_1&\cdots&x_\rho&1
  \end{array}
  \right];
\tag{4}
\end{equation}
c'est une fonction rationnelle entière de $\lambda_1,\ldots,\lambda_\rho,\mu$. D'après F. Mertens, dans le développement de (4) selon les puissances de $\mu$, le coefficient dominant peut être indiqué explicitement;\footnote{A. a. O. \S\ 3 (preuve du théorème I), et explicitement \S\ 6 (détermination de l'ensemble de tous les termes de la résultante $R$ qui contiennent comme facteur un produit de puissances donné $a_{\lambda\lambda}^{p_\lambda}a_{\mu\mu}^{p_\mu}\cdots a_{\sigma\sigma}^{p_\sigma}$). Dans notre cas, on pose $p_\lambda=0$, $p_\mu=0,\ldots,p_\sigma=\tau$, $a_{\sigma\sigma}=(-\mu)$. Reproduit chez König, chap. VI, \S\ 9.} on obtient
\begin{equation}
  (-1)^\tau\cdot R(\lambda,\mu)
  =R_0^\sigma\cdot\mu^\tau+A_1(\lambda)\mu^{\tau-1}+\cdots+A_\tau(\lambda),
\tag{5}
\end{equation}
où $A_1(\lambda),\ldots,A_\tau(\lambda)$ désignent des fonctions entières à valeurs entières de $\lambda_1,\ldots,\lambda_\rho$ et des coefficients des polynômes (3). Ce développement montre d'abord que, sous l'hypothèse (2), $R_0\ne0$, $R(\lambda,\mu)$ non plus ne peut pas s'annuler identiquement. L'identité en $x_1,\ldots,x_\rho,\lambda_1,\ldots,\lambda_\rho,\mu$,
\[
  R(\lambda,
\mu)\equiv 0\pmod{\bar G_1^*(x),\ldots,
\bar G_\rho^*(x),\bar\Phi^*(x)},
\]
passe, par la substitution
\[
  \lambda_i=G_i^*(x),\qquad \mu=\Phi^*(x),
\]
en
\begin{equation}
\begin{aligned}
  R(G_i^*(x),\Phi^*(x))
  &=R_0^\sigma\cdot\Phi^{*\tau}(x)
    +A_1(G_i^*(x))\Phi^{*\tau-1}(x)+\cdots \\
  &\qquad +A_\tau(G_i(x))=0,
\end{aligned}
\tag{6}
\end{equation}
ce qui prouve le lemme auxiliaire, puisque $R_0$ est, par hypothèse, un nombre de $\Omega$.

\section*{\S\ 12. La base d'intégrité des systèmes réguliers $\Ssys_{n\rho}$ formés de polynômes.}

Le lemme auxiliaire du \S\ 11, avec la conclusion du \S\ 10, donne immédiatement le fait suivant:

\emph{S'il existe, dans un domaine image $\Jdom^*_{\rho\rho}$ de $\Gdom_{n\rho}$, $\rho$ polynômes tels que la résultante homogène de leurs termes de plus haute dimension ne s'annule pas identiquement --- $R_0\ne0$ ---, alors $\Gdom_{n\rho}$ est de première espèce et par conséquent est un domaine d’intégrité fini, avec la base d'involution comme base d'intégrité.}

La condition $R_0\ne0$ permet maintenant, au moyen du lemme auxiliaire, une seconde preuve d'existence de la base d'intégrité, due essentiellement à D. Hilbert;\footnote{\emph{Über die vollen Invariantensysteme}, \S\ 2. Dans le cas de Hilbert, sous l'hypothèse de l'existence obtenue autrement de la base d'intégrité du corps des invariants, il ne s'agit que d'une interprétation de cette base au moyen du système fondamental de Kronecker des grandeurs algébriquement entières d'un corps.} cette preuve ne fait certes pas reconnaître la base d'involution comme telle, mais elle vaut uniformément pour tous les systèmes $\Ssys_{n\rho}$ ayant la propriété indiquée.

Soient en effet les $\rho$ polynômes de $\Jdom^*_{\rho\rho}$ pour lesquels $R_0\ne0$, et qui par conséquent sont algébriquement indépendants, donnés par
\begin{equation}
  G_1^*(x),\ldots,G_\rho^*(x).
\tag{1}
\end{equation}
D'après le lemme auxiliaire, tout polynôme arbitraire en $x_1,\ldots,x_\rho$, donc en particulier tout polynôme du plus petit corps $\Kfield^*_{\rho\rho}$ contenant $\Ssys^*_{\rho\rho}$, est algébriquement entier relativement aux polynômes (1). Si les polynômes correspondants dans $\Ssys_{n\rho}$ sont donc
\begin{equation}
  G_1(x),\ldots,G_\rho(x),
\tag{2}
\end{equation}
alors, par le principe de transfert, tout polynôme du plus petit corps $\Kfield_{n\rho}$ contenant $\Ssys_{n\rho}$, et donc en particulier tout polynôme de $\Ssys_{n\rho}$, est algébriquement entier relativement aux polynômes (2). Réciproquement, toute fonction de $\Kfield_{n\rho}$ qui est algébriquement entière relativement aux polynômes (2) est aussi un polynôme. Ainsi, si l'on considère $\Kfield_{n\rho}$ (d'après le \S\ 4) comme un corps algébrique sur $\Omega(G_1,\ldots,G_\rho)$, la totalité des polynômes de $\Kfield_{n\rho}$ --- le domaine relativement entier $\Gdom_{n\rho}$ --- est identique aux grandeurs algébriquement entières de ce corps. Soit $G(x)$ un polynôme de $\Kfield_{n\rho}$ qui complète (2) en une base rationnelle, et soit $D(x)$ le discriminant de l'équation irréductible de degré $k$ à laquelle $G(x)$ satisfait relativement à $\Omega(G_1,\ldots,G_\rho)$. Alors tout polynôme de $\Kfield_{n\rho}$, et en particulier tout polynôme $F(x)$ de $\Ssys_{n\rho}$, comme grandeur algébriquement entière, admet la représentation
\begin{equation}
  F(x)=\frac{\Gamma_1(G_i)\cdot G(x)^{k-1}+\Gamma_2(G_i)\cdot G(x)^{k-2}+\cdots+\Gamma_k(G_i)}{D(x)}.
\tag{3}
\end{equation}

Faisons maintenant parcourir à $F(x)$ tous les polynômes de $\Ssys_{n\rho}$. L'application du théorème I de la base de Hilbert (Math. Ann. 36) au système formé par les fonctions correspondantes
\[
  v_1\Gamma_1(G_i)+v_2\Gamma_2(G_i)+\cdots+v_k\Gamma_k(G_i)
\]
montre alors l'existence d'un nombre fini de polynômes de $\Ssys_{n\rho}$,
\begin{equation}
  F_1(x),F_2(x),\ldots,F_\lambda(x),
\tag{4}
\end{equation}
tels que tout polynôme de $\Ssys_{n\rho}$ admette une représentation
\begin{equation}
  F(x)=A_1(G_i)F_1(x)+A_2(G_i)\cdot F_2(x)+\cdots+A_\lambda(G_i)F_\lambda(x),
\tag{5}
\end{equation}
où $A_1,\ldots,A_\lambda$ désignent des polynômes en les $G_i(x)$. Les polynômes (2) et (4) forment donc une base d'intégrité de $\Ssys_{n\rho}$; c'est-à-dire:

\emph{Théorème VIII: S'il existe, dans un système image $\Ssys^*_{\rho\rho}$ du système de polynômes $\Ssys_{n\rho}$, $\rho$ polynômes tels que la résultante des termes de plus haute dimension ne s'annule pas identiquement --- $R_0\ne0$ ---, alors $\Ssys_{n\rho}$ possède une base d'intégrité.}

Ce théorème admet encore une légère généralisation. Il suffit que les $\rho$ polynômes (1), pour lesquels $R_0\ne0$, appartiennent au domaine d’intégrité $\Jdom^*_{\rho\rho}$ qui est le plus petit domaine contenant $\Ssys^*_{\rho\rho}$. En effet, les raisonnements conduisant au théorème VIII montrent qu'alors encore la représentation (5) vaut. Mais, d'après la définition, tout polynôme $L(x)$ du plus petit domaine d’intégrité contenant $\Jdom_{n\rho}$ peut être représenté de manière entière et rationnelle par un nombre fini --- dépendant de $L(x)$ --- de polynômes de $\Ssys_{n\rho}$; il existe donc $\sigma$ polynômes de $\Ssys_{n\rho}$,
\begin{equation}
  K_1(x),K_2(x),\ldots,K_\sigma(x),
\tag{6}
\end{equation}
tels que les polynômes $G_i(x)$ deviennent des combinaisons rationnelles entières de (6). Ainsi (6) et (4) donnent une base d'intégrité. En introduisant la terminologie suivante:

\emph{Définition VIII: Un système $\Ssys_{n\rho}$ de polynômes est appelé régulier si, dans un domaine image $\Jdom^*_{\rho\rho}$ du plus petit domaine d’intégrité qui le contient, il existe $\rho$ polynômes tels que la résultante de leurs termes de plus haute dimension ne s'annule pas identiquement --- $R_0\ne0$ ---}

nous avons donc:

\emph{Théorème IX: Tout système régulier $\Ssys_{n\rho}$ de polynômes possède une base d'intégrité.}\footnote{Que des systèmes non réguliers sans base d'intégrité existent, Hilbert l'a montré par un exemple (Gött. Nachr. 1891, p. 233). L'exemple suivant est encore plus simple:
\[
  \Ssys_{22}=x_1,
  x_1x_2,
  x_1x_2^2,
  \ldots,
  x_1x_2^\nu,
  \ldots\quad\text{à l'infini.}
\]
L'existence de la base d'intégrité serait ici équivalente à l'existence d'un entier positif $\nu$ tel que les identités
\[
  x_1x_2^{\nu+\sigma}=\sum c\cdot x_1^{i_0}(x_1x_2)^{i_1}(x_1x_2)^{i_2}\cdots(x_1x_2^\nu)^{i_\nu}
  \quad(\sigma=1,2,\ldots)
\]
soient satisfaites. La comparaison des dimensions en $x_1,x_2$, déjà pour $\sigma=1$, conduit aux deux équations conditionnelles pour les exposants
\[
  1=i_0+i_1+i_2+\cdots+i_\nu,
  \qquad
  \nu+1=i_1+2i_2+\cdots+\nu\cdot i_\nu,
\]
qui manifestement n'ont aucune solution en entiers non négatifs. Ainsi $\Ssys_{22}$ ne possède aucune base d'intégrité.}

Nous donnons encore, de façon analogue au \S\ 5 pour l'involution, une interprétation géométrique des systèmes réguliers. Pour cela, nous considérons, comme là en (7), le système d'équations
\begin{equation}
  L^*(x)-L^*(\xi)=0,
\tag{7}
\end{equation}
où $L^*(x)$ parcourt tous les polynômes de $\Jdom^*_{\rho\rho}$. En homogénéisant au moyen de $x_0,\xi_0$, supposons que (7) passe au système d'équations entre formes homogènes
\begin{equation}
  \xi_0^{m\cdot}M^*(x)-x_0^{m\cdot}M^*(\xi)=0,
\tag{8}
\end{equation}
où, si $H^*(x)$ désigne les termes de plus haute dimension de $L^*(x)$, on a la relation
\begin{equation}
  H^*(x)\equiv M^*(x)\pmod{x_0}.
\tag{9}
\end{equation}

Comme \emph{points fondamentaux} de $\Jdom^*_{\rho\rho}$ (comparer la note au \S\ 5), nous désignons les solutions de (8) qui sont indépendantes de $x$ --- donc différentes des lignes conjuguées à $x$. Les points fondamentaux sont donc les systèmes de valeurs différents de $\xi_1=0,\ldots,\xi_\rho=0$ pour lesquels on a simultanément
\[
  \xi_0^m=0,
  \qquad M^*(\xi)=0,
\]
ou, d'après (9),
\[
  \xi_0=0,
  \qquad H^*(\xi)=0.
\]

S'il existe maintenant $\rho$ formes $H^*(\xi)$ pour lesquelles $R_0\ne0$, donc sans système commun de solutions, alors $\Jdom^*_{\rho\rho}$ ne peut posséder aucun point fondamental. Si, en particulier, $\Ssys^*_{\rho\rho}$ consiste en formes homogènes, alors pour $R_0\ne0$ le système $\Ssys^*_{\rho\rho}$ lui-même ne peut posséder aucun point fondamental, puisqu'un tel point serait en même temps un point fondamental de $\Jdom^*_{\rho\rho}$. Mais la réciproque vaut aussi. Pour cela, considérons le système $\Jdom(H^*)$ formé de la totalité des formes $H^*(x)$; c'est de nouveau un domaine d’intégrité.

Du théorème I de la base de Hilbert et du lemme auxiliaire de Hilbert (comparer la citation du \S\ 10), il résulte l'existence de $\rho$ formes de $\Jdom(H^*)$ dont l'annulation entraîne l'annulation de toutes les formes de $\Jdom(H^*)$. Si maintenant $\Jdom^*_{\rho\rho}$ ne possède aucun point fondamental, ces $\rho$ formes ne peuvent pas s'annuler simultanément; leur résultante $R_0\ne0$.

Ainsi:

\emph{Les systèmes réguliers $\Ssys_{n\rho}$ de polynômes sont caractérisés par le fait qu'un domaine image $\Jdom^*_{\rho\rho}$ du plus petit domaine d’intégrité qui les contient ne possède aucun point fondamental. Pour les systèmes réguliers $\Ssys^*_{\rho\rho}$ de formes homogènes, le système lui-même ne possède déjà aucun point fondamental.}\footnote{Le système non régulier $\Ssys_{22}$ donné dans la note précédente possède le point fondamental
\[
  \xi_1:\xi_2:\xi_0=0:1:0.
\]}

\section*{\S\ 13. Exemples de domaines relativement entiers de première espèce et de systèmes réguliers.}

1. Comme premier exemple d'un domaine relativement entier de première espèce $\Gdom_{nn}$, nous considérons la \emph{totalité des polynômes en $x_1,\ldots,x_n$ qui admettent un groupe de permutations donné $G$} --- donc la totalité des polynômes d'un domaine de genre de Lagrange. Puisque $x_1,\ldots,x_n$, en vertu de l'identité
\begin{equation}
  x_k^n-\sigma_1(x)x_k^{n-1}+\sigma_2(x)x_k^{n-2}+\cdots\pm\sigma_n(x)=0
  \qquad(k=1,2,
\cdots,n),
\tag{1}
\end{equation}
sont algébriquement entiers relativement aux fonctions symétriques élémentaires appartenant à $\Gdom_{nn}$, et donc, d'après la définition V, algébriquement entiers relativement à $\Gdom_{nn}$, $\Gdom_{nn}$ est de première espèce. De plus $\Gdom_{nn}$, comme domaine de première espèce, puisque ses éléments sont des formes homogènes et que $n=\rho$, forme un système régulier (par la conclusion du \S\ 10). En effet, à cause de (1), la résultante $R_0$ des fonctions symétriques élémentaires ne peut pas s'annuler identiquement; d'après les règles de calcul de la théorie des résultantes, $R_0$ se calcule comme étant $\pm1$. La \emph{forme d'involution} de $\Gdom_{nn}$ est déterminée de manière unique --- puisque $G_{nn}$ est de première espèce et que $n=\rho$ --- comme la forme d'involution du domaine de genre de Lagrange correspondant. Comme les grandeurs conjuguées à $x_1,\ldots,x_n$ relativement à ce domaine naissent des permutations du groupe $G$, cette forme d'involution est donc donnée par
\begin{equation}
\begin{aligned}
  \Psi(z,u)&=\prod_G
  (z-u_1x_{i_1}-u_2x_{i_2}-\cdots-u_nx_{i_n})\\
  &=z^i+\sum{}'G_{\alpha\alpha_1\ldots\alpha_n}(x)
  z^\alpha u_1^{\alpha_1}\cdots u_n^{\alpha_n},
\end{aligned}
\tag{2}
\end{equation}
où le produit doit être étendu à toutes les permutations de $G$. La formule (2) représente la \og{} forme de Galois du groupe $G$ \fg{}, et il vaut donc le fait suivant:

\emph{Tous les polynômes en $x_1,\ldots,x_n$ qui admettent un groupe de permutations $G$ sont des combinaisons rationnelles entières des coefficients $G_{\alpha\alpha_1\ldots\alpha_n}(x)$ de la forme de Galois du groupe.}

2. Comme exemple de \emph{systèmes réguliers}, considérons les systèmes $\Ssys_{n1}$ de polynômes, c'est-à-dire les systèmes qui ne contiennent qu'un seul polynôme algébriquement indépendant.

En effet, soit $\Ssys^*_{11}$ un système image de $\Ssys_{n1}$ et soit
\[
  F^*(x)=a_0x^k+a_1x^{k-1}+\cdots+a_k\qquad(a_0\ne0)
\]
un polynôme de $\Ssys^*_{11}$. Alors
\begin{equation}
  R_0=a_0\ne0\srcfnmark{*)}.
\tag{3}
\end{equation}
\srcfntext{*)}{Mertens, loc. cit., \S\ 5.}
Ainsi $\Ssys_{n1}$ est régulier d'après la définition VIII, et par conséquent:

\emph{Tout système $\Ssys_{n1}$ de polynômes, en particulier tout système d'une infinité de polynômes d'une indéterminée, possède une base d'intégrité.}\footnote{Les systèmes non réguliers les plus simples possibles sans base d'intégrité sont donc des systèmes $\Ssys_{22}$. Que de tels systèmes existent effectivement est montré par l'exemple de Hilbert et par l'exemple donné dans la note du \S\ 12.}

3. Nous spécialisons maintenant les systèmes $\Ssys_{n1}$ en domaines relativement entiers $\Gdom_{n1}$, qui, comme systèmes réguliers d'après le \S\ 12, sont des domaines relativement entiers de première espèce. Leur base d'intégrité est donc donnée par une base d'involution, qui est en même temps la base d'involution du plus petit corps contenant $\Kfield_{n1}$. Mais, d'après le \S\ 6, une fonction quelconque de la base d'involution de $\Kfield_{n1}$ est en même temps une base minimale --- nous l'avons appelée la fonction de Lüroth du corps --- et toute fonction de la base d'involution dépend de cette fonction de Lüroth par une transformation linéaire fractionnaire. Dans notre cas, la fonction de Lüroth, puisqu'elle appartient à $\Gdom_{n1}$, est un polynôme; tout autre polynôme de la base d'involution dépend donc linéairement fractionnairement de cette fonction de Lüroth $L(x)$ et, à cause de (3), algébriquement entière, donc entière et linéairement. Celle-ci devient ainsi la base d'intégrité. La forme d'involution de $\Gdom^*_{11}$, en posant aussi $(u_1=1)$, est
\[
  \Psi^*(z)=L^*(z)-L^*(x),
\]
d'où $L(x)$ apparaît de nouveau comme base d'intégrité. Ainsi:

\emph{Les domaines relativement entiers $\Gdom_{n1}$ sont de première espèce; leur base d'intégrité consiste en un polynôme, la fonction de Lüroth du plus petit corps contenant.}

4. Comme exemple d'un domaine relativement entier $\Gdom_{n1}$, mentionnons encore spécialement les invariants projectifs rationnels entiers d'une forme quadratique de $s$ variables. Conformément à ce qui précède, on sait qu'ils peuvent être représentés comme combinaisons rationnelles entières du déterminant $|a_{ik}|$ de la forme; et ce déterminant est en même temps la fonction de Lüroth du corps d'invariants correspondant.

\section*{\S\ 14. Systèmes formés de polynômes entiers.}

Les théorèmes de base obtenus doivent maintenant être renforcés au regard de l'intégralité arithmétique. Le domaine des coefficients $\Omega$ sera donc désormais supposé être un corps de nombres algébriques (fini); et $[\Omega]$ désignera la totalité des nombres algébriquement entiers de $\Omega$, que nous appellerons brièvement nombres entiers. Par polynôme entier, on entendra un polynôme à coefficients dans $[\Omega]$; $F(x),G(x),\ldots$ ne désigneront désormais que des polynômes entiers.

Nous considérons, en analogie avec le \S\ 7, les différents systèmes de polynômes entiers:

Un système quelconque de polynômes entiers sera désigné comme système entier $\Ssys_{n\rho}$.

La famille linéaire entière $\Lsys_{n\rho}$ est un système entier qui, avec $F(x)$ et $F_1(x)$, contient toujours aussi $c_1F(x)+c_2F_1(x)$, où $c_1,c_2$ sont des nombres entiers quelconques de $[\Omega]$.

Le domaine d’intégrité entier $\Jdom_{n\rho}$ est une famille linéaire entière qui, avec $F(x)$ et $G(x)$, contient toujours aussi $F(x)\cdot G(x)$.

Le domaine relativement entier entier $\Gdom_{n\rho}$ est un domaine d’intégrité qui, avec $F(x)$ et $G(x)$ --- où $G(x)\ne0$ ---, contient le quotient $F(x):G(x)$ toujours et seulement lorsque celui-ci est un polynôme entier.

Le plus petit domaine d’intégrité entier $\Jdom_{n\rho}$ contenant le système $\Ssys_{n\rho}$ consiste en tous les polynômes $H(x)$ qui sont des combinaisons rationnelles entières et entières de coefficients d'un nombre fini de polynômes de $\Ssys_{n\rho}$;

le plus petit domaine relativement entier entier $\Gdom_{n\rho}$ qui le contient consiste en tous les polynômes entiers $H(x)$ qui sont des combinaisons rationnelles d'un nombre fini de polynômes de $\Ssys_{n\rho}$.

Passons maintenant au transfert des théorèmes de base. En ce qui concerne la base rationnelle, il faut remarquer que les notions de corps, de plus petit corps contenant et de base rationnelle, qui reposent seulement sur la \og{} rationalité \fg{}, ne subissent aucune modification. En outre, le système image $\Ssys^*_{\rho\rho}$ d'un système entier $\Ssys_{n\rho}$ peut être choisi, de façons arbitrairement nombreuses, encore comme système entier. Comme, pour la famille linéaire entière aussi, les raisonnements conduisant à la borne du nombre de générateurs restent valables, on obtient:

\emph{Théorème V$'$: Tout système entier $\Ssys_{n\rho}$ possède une base rationnelle; la base rationnelle d'une famille linéaire entière $\Lsys_{n\rho}$ peut en particulier être choisie de telle sorte qu'elle contienne au plus $\rho+1$ polynômes.}

Nous examinons maintenant l'analogue de la base d'involution des domaines d’intégrité de polynômes (\S\ 8). À cette fin, le théorème sur les \emph{fonctions symétriques de lignes de grandeurs} doit être étendu relativement à \emph{l'intégralité sur les entiers} par le lemme auxiliaire suivant.

\emph{Lemme auxiliaire. Les fonctions symétriques entières et à coefficients entiers de $n$ lignes de grandeurs $(yz\cdots t)$:}
\begin{equation}
\begin{array}{cccc}
  y_1&z_1&\cdots&t_1\\
  y_2&z_2&\cdots&t_2\\
  \vdots&\vdots&&\vdots\\
  y_n&z_n&\cdots&t_n
\end{array}
\tag{1}
\end{equation}
\emph{sont des fonctions entières à coefficients entiers d'un nombre fini d'entre elles, à savoir des fonctions symétriques élémentaires}
\begin{equation}
\begin{gathered}
  \sigma_1(y),\sigma_2(y),\ldots,\sigma_n(y);\quad
  \sigma_1(z),\sigma_2(z),\ldots,\sigma_n(z);\quad\ldots;\\
  \sigma_1(t),\ldots,\sigma_n(t);
\end{gathered}
\tag{2}
\end{equation}
\emph{et des fonctions}
\begin{equation}
  \chi_1(y,z,\ldots,t),\ldots,\chi_k(y,z,\ldots,t),
\tag{3}
\end{equation}
\emph{qui sont algébriquement entières et entières relativement aux fonctions (2).}

En effet, les $n$ grandeurs $y_1,\ldots,y_n$ dépendent algébriquement entièrement et intégralement, sur les entiers, de $\sigma_1(y),\ldots,\sigma_n(y)$; l'assertion correspondante vaut pour $z_1,\ldots,z_n$ relativement à $\sigma_1(z),\ldots,\sigma_n(z)$, et ainsi de suite. Le corps des fonctions symétriques des lignes de grandeurs (1) forme donc un corps algébrique sur le corps déterminé par les fonctions (2), de telle manière que les grandeurs algébriquement entières et entières de ce corps sont identiques aux fonctions symétriques entières et à coefficients entiers des lignes de grandeurs (1). Si l'on prend donc (3) comme système fondamental de Kronecker, le lemme auxiliaire est démontré.

Soit maintenant $\Jdom_{n\rho}$ un domaine d’intégrité entier, $\Jdom^*_{\rho\rho}$ un domaine image entier, et
\[
  \Psi^*(z,u)=G^*(x)\cdot z^i+
  \sum{}'G^*_{\alpha\alpha_1\ldots\alpha_\rho}(x)\cdot
  z^\alpha u_1^{\alpha_1}\cdots u_\rho^{\alpha_\rho}
\]
une forme d'involution de $\Jdom^*_{\rho\rho}$; ici $G^*(x)$ peut aussi avoir dimension zéro, c'est-à-dire être un nombre de $[\Omega]$. Les fonctions symétriques élémentaires correspondant aux fonctions (2) sont alors données par certaines parmi les quotients
\begin{equation}
  G^*_{\alpha\alpha_1\ldots\alpha_\rho}(x):G^*(x).
\tag{4}
\end{equation}
Comme, d'après le lemme auxiliaire, les fonctions (3) dépendent algébriquement entièrement et intégralement des entiers de (2), les fonctions de $\Kfield^*_{\rho\rho}$ correspondant à (3) sont données --- par l'extension du théorème de Gauss aux polynômes à coefficients algébriquement entiers --- par
\begin{equation}
 K_\lambda^*(x):(G^*(x))^\sigma,
\tag{5}
\end{equation}
où les $K_\lambda^*(x)$ sont des polynômes entiers de $\Kfield^*_{\rho\rho}$. Ainsi les $K_\lambda^*(x)$ appartiennent au domaine lorsqu'il s'agit de domaines relativement entiers; en général, ce sont des quotients de polynômes de $\Jdom^*_{\rho\rho}$.

Si maintenant $F^*(x)$ est un polynôme entier arbitraire de $\Jdom^*_{\rho\rho}$, on a
\[
  F^*(x)=\frac1i\Bigl(F^*(x)+F^*(x^{(1)})+\cdots+F^*(x^{(i-1)})\Bigr),
\]
et donc $i\cdot F^*(x)$ devient une fonction symétrique entière et à coefficients entiers des lignes de grandeurs $x,x^{(1)},\ldots,x^{(i-1)}$.

Par le lemme auxiliaire et le principe de transfert, il en résulte donc:

\emph{Théorème VI$'$: Pour les domaines d’intégrité entiers $\Jdom_{n\rho}$, chaque forme d'involution détermine une base rationnelle distinguée de telle sorte que tout polynôme $F(x)$ de $\Jdom_{n\rho}$ admette une représentation}
\[
  F(x)=\frac{\Gamma(H(x),H_1(x),\ldots,H_\sigma(x))}{i\cdot H(x)^k},
\]
\emph{où $\Gamma$ désigne une fonction entière à coefficients entiers. Pour les domaines relativement entiers entiers $\Gdom_{n\rho}$, avec comme domaine image un domaine relativement entier entier $\Gdom^*_{\rho\rho}$, le dénominateur $H(x)$ est donné par le coefficient dominant de la forme d'involution correspondante.}

\section*{\S\ 15. Domaines relativement entiers entiers de première espèce et systèmes réguliers.}

Pour les systèmes entiers, la base d'intégrité entière prend la place de la base d'intégrité; c'est-à-dire un nombre fini de polynômes entiers du système tels que tout polynôme entier du système puisse être représenté comme combinaison entière à coefficients entiers de ce nombre fini.

Les théorèmes sur la base d'intégrité entière sont essentiellement parallèles aux précédents, et seules les preuves demandent de légères modifications.

D'abord, la définition V (\S\ 9) se transfère directement.

\emph{Définition V$'$: Une grandeur $\xi$ s'appelle algébriquement entière et entière relativement au domaine relativement entier entier $\Gdom_{n\rho}$ si elle satisfait à une équation dont le coefficient dominant est une unité de $[\Omega]$ et dont les autres polynômes appartiennent à $\Gdom_{n\rho}$.}

L'extension du théorème de Gauss aux polynômes à coefficients algébriquement entiers montre, comme au \S\ 9, que la définition à l'équation irréductible peut être obtenue à partir de là.

De la définition V$'$ on obtient ensuite le transfert de la définition VII.

\emph{Définition VII$'$: Un domaine relativement entier entier $\Gdom_{n\rho}$ s'appelle de première espèce si, comme domaine image, il possède un domaine relativement entier $\Gdom^*_{\rho\rho}$ tel que tout polynôme entier arbitraire en $x_1,\ldots,x_\rho$ soit algébriquement entier et entier relativement à $\Gdom^*_{\rho\rho}$.}

Pour montrer la finitude de ces domaines de première espèce, nous remarquons d'abord que, comme au \S\ 10, il résulte de la définition que le coefficient dominant de la forme d'involution est une unité de $[\Omega]$.

D'après le théorème VI$'$ (\S\ 14), tout polynôme $F(x)$ de $\Gdom_{n\rho}$ admet donc une représentation
\[
  F(x)=\frac{\Gamma(H_1(x),\ldots,H_\sigma(x))}{i}.
\]
En appliquant ici au système formé des polynômes entiers $\Gamma(H_1,\ldots,H_\sigma)$ le théorème II de la base de Hilbert à coefficients entiers (Math. Ann. 36), énoncé originairement pour des coefficients numériques rationnels entiers, dans son extension aux polynômes algébriquement entiers,\footnote{Si $w_1,\ldots,w_k$ désigne un système fondamental des entiers algébriques de $\Omega$, et si l'on pose
\[
  \Gamma(H_1,\ldots,H_\sigma)=\Gamma_1(H)w_1+\cdots+\Gamma_k(H)w_k,
\]
alors $\Gamma_1(H),\ldots,\Gamma_k(H)$ ont des coefficients numériques rationnels entiers. L'application du théorème II au système formé des formes $v_1\Gamma_1(H)+\cdots+v_k\Gamma_k(H)$ montre directement la validité de l'extension.} il s'ensuit l'existence de la base d'intégrité entière; c'est-à-dire:

\emph{Théorème VII$'$: Tout domaine relativement entier entier de première espèce possède une base d'intégrité entière.}

Nous transférons encore la définition des systèmes réguliers:

\emph{Définition VIII$'$: Un système entier $\Ssys_{n\rho}$ s'appelle entier-régulier si, dans un domaine image $\Jdom^*_{\rho\rho}$ du plus petit domaine d’intégrité entier qui le contient, il existe $\rho$ polynômes tels que la résultante des termes de plus haute dimension soit une unité de $[\Omega]$ --- $R_0=\eps$.}

Pour transférer la preuve d'existence de la base d'intégrité, nous remarquons que le lemme auxiliaire du \S\ 11 admet la formulation plus précise suivante, qui résulte des équations (5) et (6) du \S\ 11 et du fait que $A_1(\lambda),\ldots,A_\tau(\lambda)$ sont des fonctions entières à coefficients entiers de $\lambda_1,\ldots,\lambda_\rho$:

\emph{Une condition suffisante pour que tout polynôme entier arbitraire en $x_1,\ldots,x_\rho$ dépende algébriquement entièrement et intégralement de $\rho$ polynômes entiers est que la résultante des termes de plus haute dimension de ces polynômes soit une unité de $[\Omega]$ --- $R_0=\eps$.}

De ce lemme auxiliaire résulte exactement comme au \S\ 12, pour tout polynôme $F(x)$ d'un système entier-régulier, la représentation (3), \S\ 12, où maintenant $\Gamma_1(G_i),\ldots,\Gamma_k(G_i)$ sont des polynômes entiers en les $G_i$. Et de plus, comme au \S\ 12, l'application du théorème II de Hilbert (Math. Ann. 36), dans son extension aux polynômes algébriquement entiers, jointe à la définition du plus petit domaine d’intégrité entier contenant, donne l'existence de la base d'intégrité entière; c'est-à-dire:

\emph{Théorème IX$'$: Tout système entier-régulier possède une base d'intégrité entière.}

Comme exemple d'un domaine relativement entier entier de première espèce, en analogie avec l'exemple 1 du \S\ 13, renvoyons à la totalité $\Gdom_{nn}$ des polynômes entiers d'un domaine de genre de Lagrange. Que $\Gdom_{nn}$ soit entier de première espèce résulte de l'identité
\[
  x_k^n-\sigma_1(x)x_k^{n-1}+\cdots\pm\sigma_n(x)=0
  \qquad(k=1,2,
\cdots,n);
\]
puisque la résultante des fonctions symétriques élémentaires a la valeur $\pm1$, $\Gdom_{nn}$ est aussi un système entier régulier.

Un autre exemple est donné, d'après le lemme auxiliaire du \S\ 14, par la totalité des fonctions symétriques entières et à coefficients entiers de $n$ lignes de grandeurs.

Erlangen, mai 1914.
\clearpage
\setcounter{footnote}{0}
\section*{\NoetherNavAnchor{P07}{7. Le théorème de finitude pour les invariants des groupes finis}7. Le théorème de finitude pour les invariants des groupes finis}
\begin{center}
Par Emmy Noether, à Erlangen.\\[0.75em]
Math. Ann. 77 (1916), pp. 89--92
\end{center}

Dans ce qui suit on donnera une démonstration de finitude entièrement élémentaire des invariants des groupes finis --- ne reposant que sur la théorie des fonctions symétriques --- qui fournit en même temps une indication effective du système complet; tandis que la démonstration ordinaire, fondée sur le théorème de la base de Hilbert (Ann. 36), n'est qu'une démonstration d'existence.\footnote{Cf. par exemple Weber, \emph{Lehrbuch der Algebra} (2e éd.), vol. 2, \S\ 57.}

Soit le groupe fini $\Hgrp$ composé des $h$ transformations linéaires (de déterminant non nul) $A_1\cdots A_h$, où $A_k$ désigne la transformation linéaire
\[
  x_1^{(k)}=\sum_{\nu=1}^{n}a_{1\nu}^{(k)}x_\nu,\cdots,
  x_n^{(k)}=\sum_{\nu=1}^{n}a_{n\nu}^{(k)}x_\nu,
\]
ou, sous forme abrégée, $(x^{(k)})=A_k(x)$. Le groupe $\Hgrp$ transforme donc la rangée $(x)$ d'éléments $x_1\cdots x_n$ en les rangées $(x^{(k)})$ d'éléments $x_1^{(k)}\cdots x_n^{(k)}$. Puisque l'identité doit être contenue parmi $A_1\cdots A_h$, la rangée $(x)$ elle-même est aussi contenue parmi les rangées $(x^{(k)})$. Par invariant entier rationnel (absolu) du groupe on entend une fonction entière rationnelle de $x_1\cdots x_n$ qui reste identiquement inchangée par l'application de $A_1\cdots A_h$; pour un tel invariant $f(x)$ on a donc
\begin{equation}
  f(x)=f(x^{(1)})=\cdots=f(x^{(k)})=
  \frac1h\cdot\sum_{k=1}^{h}f(x^{(k)}).
\tag{1}
\end{equation}

\subsection*{1.}
La formule (1) exprime que $f(x)$ est une fonction symétrique entière rationnelle des rangées de grandeurs $(x^{(1)})\cdots(x^{(h)})$ --- et il s'agit, puisque chaque terme $f(x^{(k)})$ ne contient que la seule rangée $(x^{(k)})$, du cas le plus simple, ordinairement appelé cas uniforme. D'après le théorème connu sur les fonctions symétriques de rangées de grandeurs,\footnote{Cf. la note à 2.} $f(x)$ est donc représentable entière et rationnellement par les fonctions symétriques élémentaires de ces rangées, c'est-à-dire par les coefficients $G_{\alpha\alpha_1\cdots\alpha_n}(x)$ de la \og{} résolvante de Galois \fg{}
\[
\begin{aligned}
  \Phi(z,u)
  &=\prod_{k=1}^{h}\bigl(z+u_1x_1^{(k)}+\cdots+u_nx_n^{(k)}\bigr) \\
  &=z^h+
    \sum G_{\alpha\alpha_1\cdots\alpha_n}(x)
    z^{\alpha}u_1^{\alpha_1}\cdots u_n^{\alpha_n},
\end{aligned}
\qquad
\left(\begin{array}{c}
(\alpha+\alpha_1+
\cdots+
\alpha_n=h)\\
\alpha\ne h
\end{array}\right),
\]
où les $G_{\alpha\alpha_1\cdots\alpha_n}(x)$ sont des invariants de degré $\alpha_1+\cdots+\alpha_n$ en les $x$. Ainsi est démontré:

\emph{Les coefficients $G_{\alpha\alpha_1\cdots\alpha_n}(x)$ de la résolvante de Galois forment un système complet d'invariants du groupe, de telle sorte que tout invariant se représente entière et rationnellement par ces invariants en nombre fini.}

\subsection*{2.}
En partant de (1), on peut encore faire la considération suivante, plus élémentaire encore,\footnote{Cela revient à une démonstration du théorème sur les fonctions symétriques de rangées de grandeurs dans le cas uniforme mentionné ci-dessus.} qui conduit en même temps à un second système complet. Posons
\[
  f(x)=a+b x_1^{\mu_1}\cdots x_n^{\mu_n}
       +\cdots+c x_1^{\nu_1}\cdots x_n^{\nu_n},
\]
où $a,b,\cdots,c$ désignent des constantes. Alors, d'après (1), on a
\[
\begin{aligned}
 h\cdot f(x)
 &=h\cdot a+b\cdot\sum_{k=1}^{h}x_1^{(k)^{\mu_1}}\cdots x_n^{(k)^{\mu_n}}
     +\cdots
     +c\cdot\sum_{k=1}^{h}x_1^{(k)^{\nu_1}}\cdots x_n^{(k)^{\nu_n}} \\
 &=h\cdot a+b\cdot J_{\mu_1\cdots\mu_n}+\cdots+c\cdot J_{\nu_1\cdots\nu_n}.
\end{aligned}
\]
Tout invariant se compose donc entière et linéairement à partir des invariants spéciaux
\[
  J_{\mu_1\cdots\mu_n}
  =\sum_{k=1}^{h}x_1^{(k)^{\mu_1}}\cdots x_n^{(k)^{\mu_n}},
\]
et il suffit de prouver l'assertion pour ceux-ci. Mais, à un facteur numérique près, $J_{\mu_1\cdots\mu_n}$ est le coefficient de $u_1^{\mu_1}\cdots u_n^{\mu_n}$ dans l'expression
\[
  S_\mu=\sum_{k=1}^{h}
  \bigl(u_1x_1^{(k)}+\cdots+u_nx_n^{(k)}\bigr)^{\mu},
\]
où $\mu=\mu_1+\cdots+\mu_n$; cette expression est la $\mu$-ième somme de puissances des $h$ formes linéaires
\[
  \xi_1=u_1x_1^{(1)}+\cdots+u_nx_n^{(1)},\cdots,
  \xi_h=u_1x_1^{(h)}+\cdots+u_nx_n^{(h)}.
\]
On sait que les sommes de puissances infiniment nombreuses $S_\mu$ sont des combinaisons entières rationnelles de
\[
  S_1,S_2,\cdots,S_h,
\]
dont les coefficients sont donnés par les invariants
\[
  J_{\mu_1\cdots\mu_n}\qquad(\mu_1+\cdots+\mu_n\leqq h).
\]
On a ainsi obtenu un second système complet:

\emph{Un système complet d'invariants du groupe est donné par tous les invariants $J_{\mu_1\cdots\mu_n}$ pour lesquels $\mu_1+\cdots+\mu_n\leqq h$, où $h$ désigne l'ordre du groupe.}

Le lien avec 1. est établi par la remarque que les sommes de puissances $S_1\cdots S_h$ peuvent être remplacées par les fonctions symétriques élémentaires de $\xi_1\cdots\xi_h$, ce qui conduit au système complet des $G_{\alpha\alpha_1\cdots\alpha_n}(x)$ donné là. Les deux résultats montrent que tous les invariants s'expriment entière et rationnellement par ceux dont le degré en les $x$ ne dépasse pas l'ordre $h$ du groupe.

\subsection*{3.}
Du résultat qui vient d'être démontré on peut tirer des conséquences pour la représentation rationnelle. Tout invariant rationnel absolu se représente --- comme on le voit immédiatement en multipliant le numérateur et le dénominateur par les conjugués du dénominateur relativement au groupe --- comme quotient de deux invariants absolus entiers rationnels, non nécessairement premiers entre eux. Il s'ensuit que tout invariant rationnel absolu s'exprime rationnellement par les coefficients $G_{\alpha\alpha_1\cdots\alpha_n}(x)$ de la résolvante de Galois $\Phi(z,u)$. Ce théorème peut aussi être prouvé facilement, de diverses manières, sans passer par le théorème de finitude; il se trouve déjà formulé dans Weber II, \S\ 58.\footnote{La démonstration qui y est donnée n'est cependant pas concluante. On montre seulement que la fonction $\Psi(t)$ de la formule (7) a des invariants pour coefficients, mais non que ces invariants soient exprimables rationnellement par les coefficients de $\Phi(t)$. Cette lacune est évitée si, pour représenter les $x_i^{(k)}$ par $\xi_k$, on emploie à la place de la formule d'interpolation de Lagrange un procédé de différentiation connu. Il s'agit de la relation obtenue en différentiant l'identité $\Phi(-\xi_k,u)=0$ par rapport à tous les $u$:
\[
 \left[\frac{\partial\Phi}{\partial u_i}
      -x_i^{(k)}\cdot\frac{\partial\Phi}{\partial z}\right]_{z=-\xi_k}=0.
\]
Par conséquent, pour toute fonction rationnelle $\omega(x)$, les formules (7) et (8) de Weber doivent être remplacées par la formule suivante:
\[
 \omega(x_1^{(k)}\cdots x_n^{(k)})=
 \left.
 \omega\!\left(
   \frac{\frac{\partial\Phi}{\partial u_1}}{\frac{\partial\Phi}{\partial z}}\cdots,
   \frac{\frac{\partial\Phi}{\partial u_n}}{\frac{\partial\Phi}{\partial z}}
 \right)\right|_{z=-\xi_k},
\]
qui ne contient plus que des coefficients de $\Phi(z,u)$, et dont la sommation sur tous les $k$, pour des invariants $\omega(x)$, donne la représentation cherchée.}

\subsection*{4.}
Enfin, mentionnons que les résultats déduits ici sont implicitement contenus dans mon travail \og{} Corps et systèmes de fonctions rationnelles \fg{};\footnote{Math. Ann. 76, p. 161 (1915).} plus précisément, le système complet donné en 1. se trouve dans les théorèmes VI et VII, et une démonstration de la représentation rationnelle indépendante de ce théorème de finitude se trouve dans le théorème III.\footnote{Pour les invariants relatifs, les théorèmes VIII et IX contiennent une démonstration de finitude qui repose également sur un autre fondement que la démonstration ordinaire, mais qui n'est elle aussi qu'une démonstration d'existence; cela tient au fait que les invariants relatifs ne forment pas un corps.} Mais, dans le cas spécial présent, la marche de la preuve et le résultat se simplifient considérablement par rapport à la théorie générale. C'est par des conversations avec E. Fischer que j'ai été amenée à appliquer les recherches générales aux invariants des groupes finis; c'est également de lui que proviennent, quant au fond, la démonstration donnée sous 2 --- dans le cas général le système complet qui y apparaît n'est pas définissable rationnellement --- et la note à 3.

Erlangen, mai 1915.

\clearpage
\setcounter{footnote}{0}
\section*{\NoetherNavAnchor{P08}{8. Sur la représentation entière rationnelle des invariants d’un système d’un nombre arbitraire de formes fondamentales}8. Sur la représentation entière rationnelle des invariants d'un système d'un nombre arbitraire de formes fondamentales}
\begin{center}
Par Emmy Noether, à Erlangen.\\[0.75em]
Math. Ann. 77 (1916), pp. 93--102
\end{center}

À la fin de sa contribution au volume commémoratif de Schwarz, \og{} Sur les invariants d'un système d'un nombre arbitraire de formes fondamentales \fg{}, D. Hilbert énonce la conjecture que ces invariants se représentent entière et rationnellement par les invariants en nombre fini du système $(J,PJ)$, où $J$ désigne le système complet d'invariants des formes fondamentales linéairement indépendantes, et où les invariants $PJ$ proviennent de $J$ par des procédés polaires.

Je montre dans ce qui suit, en I, que cette conjecture est effectivement vraie, et cela comme une conséquence simple du théorème de réduction d'après lequel toute forme d'un nombre arbitraire de rangées de $n$ variables chacune se représente comme somme de polaires de formes qui ne contiennent que $n$ rangées fixes et qui sont elles-mêmes déduites de la forme initiale par des procédés polaires. Ce théorème de réduction est un cas spécial, formulé d'abord par F. Mertens, de la généralisation, donnée par A. Capelli, F. Mertens et J. Deruyts, du développement en séries de Clebsch--Gordan aux formes d'un nombre arbitraire de rangées de $n$ variables chacune; cette généralisation est démontrée par transformation formelle des procédés de différentiation.\footnote{F. Mertens: \og{} Über eine Formel der Determinantentheorie. \fg{} (formule 5), Sitzb. d. Ak. d. Wiss. Wien, vol. 91, II. Abt. (1885), et plus complètement (avec applications) dans: \og{} Über ganze Funktionen von $m$ Systemen von je $n$ Unbestimmten \fg{}. Monatsh. f. Math. u. Phys. 4e année (1893). Pour quelques-unes des démonstrations de Capelli (depuis 1882), cf. par exemple Math. Ann. 37, ou les \og{} Lezioni sulla teoria delle forme algebriche \fg{} autographiées (1902). J. Deruyts: \emph{Essai d'une théorie générale des formes algébriques}, Mém. d. l. Soc. Roy. des Sciences de Liège (2) 17 (1892).}

En II, je donne encore une démonstration plus conceptuelle du théorème de réduction, en concevant la question comme un problème d'équivalence de faisceaux linéaires de formes. Cette conception, ainsi qu'une partie essentielle des raisonnements employés, je les emprunte à un travail \og{} Sur les systèmes de modules algébriques \fg{}\footnote{E. Fischer: Über algebraische Modulsysteme und lineare homogene partielle Differentialgleichungen mit konstanten Koeffizienten, \S\ 1--4, J. f. M. 140 (1911).} et à des théorèmes inédits d'E. Fischer qui s'y rattachent, dont celui-ci m'a aimablement permis l'utilisation.

Enfin, je montre encore en III comment les raisonnements correspondants conduisent aussi aux développements en séries les plus généraux.

\subsection*{I.}
Soit $\theta$ une forme homogène dans chacune des $N>n$ rangées
\[
  A_1,A_2,\ldots,A_N,
\]
où, en général, la rangée $A_k$ se compose des $n$ éléments
\[
  A_k^{(1)},A_k^{(2)},\ldots,A_k^{(n)}.
\]
Les coefficients de $\theta$ peuvent être des indéterminées ou des grandeurs d'un domaine de rationalité donné. Désignons en outre par $P$ une fonction entière rationnelle --- à coefficients entiers rationnels --- des procédés polaires
\[
  P_{hk}=\left(A_h\frac{\partial}{\partial A_k}\right)
  =A_h^{(1)}\frac{\partial}{\partial A_k^{(1)}}+
   A_h^{(2)}\frac{\partial}{\partial A_k^{(2)}}+
   \cdots+
   A_h^{(n)}\frac{\partial}{\partial A_k^{(n)}},
\]
où le produit $P_{hk}P_{ij}\theta$ doit être compris comme l'application de $P_{hk}$ à $P_{ij}\theta$.

Le théorème de réduction mentionné plus haut est alors donné par l'identité
\begin{equation}
  \theta=\sum PZ,
\tag{1}
\end{equation}
où les formes $Z$ ne contiennent que les rangées
\[
  A_1,A_2,\ldots,A_n
\]
et sont déduites de $\theta$ par des procédés polaires $P$.

Considérons maintenant le système de formes fondamentales $(F)$, composé des $N$ formes du même ordre et du même nombre de variables
\[
  F_1,F_2,\ldots,F_N,
\]
dont les coefficients sont pris respectivement comme les $N$ rangées
\[
  A_1,A_2,\ldots,A_N.
\]
Soit $(J)$ le système complet --- composé d'un nombre fini d'invariants --- de $F_1,\ldots,F_n$, de sorte que tout invariant de $F_1,\ldots,F_n$ s'exprime entière et rationnellement par les invariants $(J)$. La conjecture de Hilbert à démontrer est alors que, pour les invariants simultanés $S$ de $F_1,\ldots,F_N$, un système complet est donné par les invariants en nombre fini $(J,PJ)$, où $PJ$ désigne tous les invariants déductibles de $J$ par des procédés polaires $P$.

En effet, tout invariant simultané $S$, à coefficients entiers rationnels, de $(F)$ est une forme homogène dans chacune des $N$ rangées $A_1,\ldots,A_N$ --- à coefficients entiers rationnels --- et l'identité (1) peut donc lui être appliquée. Comme l'application des procédés polaires $P$ à $S$ donne de nouveau, comme on le sait, des invariants, les formes $Z$ sont également des invariants, à savoir des invariants simultanés de $F_1,\ldots,F_n$, puisqu'elles ne contiennent que les rangées $A_1,\ldots,A_n$; par hypothèse, elles sont donc des fonctions entières rationnelles des invariants $(J)$. L'identité (1) devient ainsi
\[
  S=\sum PY,
\]
où les $Y$ désignent des produits de puissances des invariants $(J)$. Mais, par la définition de $P$, l'exécution effective des procédés $P$ sur $Y$ consiste en l'exécution successive, répétée un nombre fini de fois, des opérations polaires simples $P_{hk}$; d'après la règle de différentiation des produits, cela conduit à des sommes de produits d'invariants $(J)$ et $(PJ)$. La même chose vaut pour l'application de $P_{hk}$ à un produit tiré de $(J)$ et $(PJ)$; donc $PY$ aussi ne donne que des combinaisons entières rationnelles de $(J,PJ)$. \emph{La représentabilité entière rationnelle de tous les invariants $S$ par $(J,PJ)$ est ainsi démontrée.}

\subsection*{II.}
Avant de démontrer le théorème de réduction, rappelons quelques faits sur les faisceaux linéaires de formes. Par \emph{faisceau linéaire de formes sur $K$}, nous entendons un système de formes de même dimension, complet en ce sens qu'avec $\theta_1$ et $\theta_2$ il contient toujours $c_1\theta_1+c_2\theta_2$, où $c_1,c_2$ sont des grandeurs d'un domaine de rationalité prescrit $K$, et où $K$ contient le domaine des coefficients des $\theta$. Parmi ces formes, il existe toujours un nombre fini --- soit $\rho$ --- de formes linéairement indépendantes sur $K$, tandis que $\rho+1$ formes quelconques satisfont à une relation linéaire à coefficients dans $K$; le faisceau a le \emph{rang $\rho$} relativement à $K$.\footnote{On obtient des faisceaux linéaires de formes plus généraux --- dont le rang n'a pas besoin d'être fini --- en supprimant la restriction de même dimension, ou encore la restriction que $K$ contienne le domaine des coefficients des $\theta$.} Nous utiliserons le fait évident suivant: si $\mathfrak L$ et $\mathfrak T$ sont deux faisceaux linéaires sur $K$ de même rang $\rho$, et si $\mathfrak T$ est un sous-faisceau de $\mathfrak L$, alors $\mathfrak L$ et $\mathfrak T$ sont identiques (équivalents).

Nous passons maintenant au théorème de réduction. L'identité de I,
\[
  \theta=\sum PZ,
\]
passe en une identité analogue pour $P\theta$ lorsqu'on applique le même procédé polaire $P$ aux deux membres; le théorème de réduction donne donc en même temps une représentation de toutes les polaires de $\theta$ par des sommes des polaires spéciales $PZ$. Comme, réciproquement, ces polaires spéciales sont sûrement contenues parmi toutes les polaires, le théorème de réduction affirme l'équivalence de ces deux faisceaux linéaires de formes; en particulier, il affirme aussi l'équivalence des sous-faisceaux dont les formes ont, dans chaque rangée individuelle, la même dimension que $\theta$. Pour prouver cette équivalence, nous insérons un faisceau intermédiaire déterminé par les formes $Z$. Pour simplifier, nous donnons d'abord la preuve dans le cas de trois rangées binaires $(N=3,n=2)$, puis nous indiquerons brièvement la preuve générale, entièrement parallèle.

Que $\theta(x,y,z)$ ait donc les dimensions $\alpha,\beta,p$ respectivement dans les rangées
\[
  x_1x_2;\quad y_1y_2;\quad z_1z_2;
\]
et que les coefficients de $\theta$ soient d'abord des indéterminées $a$ (ou des nombres rationnels). Nous considérons les trois faisceaux linéaires de formes suivants.

1. Le faisceau $\mathfrak L$, formé de toutes les formes $f(x,y,z)$ qui proviennent de $\theta$ par des procédés polaires $P$ et qui, dans les rangées individuelles $x,y,z$, ont la même dimension que $\theta$; $\mathfrak L$ contient en particulier $\theta$. Si $f(x,y,z)$ est regardée comme une forme en $x,y,z$ et dans les indéterminées $a$, le corps des coefficients est le corps $R$ des nombres rationnels; pour $K$ aussi il faut prendre $R$, car l'expression $c_1P_1+c_2P_2$ ne redevient un procédé $P$ que lorsque $c_1,c_2$ sont des entiers rationnels; soit $\mathfrak L$ de rang $\rho$ relativement à $R$.

2. Le faisceau $\mathfrak S$, formé de toutes les formes $g(\lambda;x,y)$ définies par
\[
  g(\lambda;x,y)=f(x,y,\lambda_1x+\lambda_2y),
\]
ou, exprimé par des procédés polaires,
\[
\begin{aligned}
  g(\lambda;xy)
  &=\frac1{p!}\left[(\lambda_1x+\lambda_2y)\frac{\partial}{\partial z}\right]^p f(x,y,z) \\
  &=g_0(xy)\lambda_1^p+g_1(x,y)\lambda_1^{p-1}\lambda_2+\cdots+g_p(xy)\lambda_2^p,
\end{aligned}
\]
où
\[
  i!(p-i)!\,g_i(xy)=
  \left(y\frac{\partial}{\partial z}\right)^i
  \left(x\frac{\partial}{\partial z}\right)^{p-i}f(xyz).
\]
\footnote{Comme en I,
\[
 \left(t\frac{\partial}{\partial z}\right)=t_1\frac{\partial}{\partial z_1}+t_2\frac{\partial}{\partial z_2},
\]
et donc en particulier
\[
 \left[(\lambda_1x+\lambda_2y)\frac{\partial}{\partial z}\right]
 = (\lambda_1x_1+\lambda_2y_1)\frac{\partial}{\partial z_1}
 + (\lambda_1x_2+\lambda_2y_2)\frac{\partial}{\partial z_2},
\]
\[
 \left[z\left(\frac{\partial}{\partial\lambda_1}\frac{\partial}{\partial x}
 +\frac{\partial}{\partial\lambda_2}\frac{\partial}{\partial y}\right)\right]
 =z_1\left(\frac{\partial}{\partial\lambda_1}\frac{\partial}{\partial x_1}
 +\frac{\partial}{\partial\lambda_2}\frac{\partial}{\partial y_1}\right)
 +z_2\left(\frac{\partial}{\partial\lambda_1}\frac{\partial}{\partial x_2}
 +\frac{\partial}{\partial\lambda_2}\frac{\partial}{\partial y_2}\right).
\]}
Les $g_i(xy)$ donnent ainsi des formes $Z$ de l'identité (1). Les $g$ sont des formes à coefficients rationnels entiers en $x,y,\lambda$ et dans les indéterminées $a$; soit $\mathfrak S$ de rang $\sigma$ relativement à $R$.

3. Le faisceau $\mathfrak T$, obtenu de $\mathfrak S$ par le procédé polaire inverse, comme $\mathfrak S$ est obtenu de $\mathfrak L$; les formes $h$ de $\mathfrak T$ sont définies par
\[
\begin{aligned}
 h(x,y,z)
 &=\frac1{p!}\left[z\left(\frac{\partial}{\partial\lambda_1}\frac{\partial}{\partial x}
 +\frac{\partial}{\partial\lambda_2}\frac{\partial}{\partial y}\right)\right]^p g(\lambda;xy)\\
 &=h_0(xyz)+h_1(xyz)+\cdots+h_p(xyz),
\end{aligned}
\]
où, d'après 2.,
\[
  h_i(xyz)=
  \left(z\frac{\partial}{\partial x}\right)^{p-i}
  \left(z\frac{\partial}{\partial y}\right)^i g_i(xy).
\]
Les $h_i$ pris séparément, et par suite aussi $h$, sont donc des formes $PZ$ et des sommes de telles formes; soit $\mathfrak T$ de rang $\tau$ relativement à $R$.

Comme le montre la construction des formes $h(x,y,z)$ de $\mathfrak T$, ces formes proviennent de $\theta$ par des procédés polaires $P$, et dans les rangées $x,y,z$ elles ont séparément la même dimension que $\theta$; le faisceau $\mathfrak T$ est donc un sous-faisceau de $\mathfrak L$. D'après le fait rappelé ci-dessus, il ne reste qu'à démontrer l'égalité des rangs. Dans cette démonstration, on n'utilise plus la structure spéciale des $f(xyz)$ --- à savoir qu'elles sont les polaires d'une seule et même forme $\theta$.

\emph{a) Pour comparer les rangs de $\mathfrak L$ et de $\mathfrak S$, nous employons le lemme a): de $f(x,y,\lambda_1x+\lambda_2y)=0$, il s'ensuit nécessairement que $f(x,y,z)=0$.} Cela résulte immédiatement de la relation identique entre trois rangées binaires
\[
  (xy)z+(yz)x+(zx)y=0,
  \qquad (xy)=(x_1y_2-x_2y_1)\cdots,
\]
qui implique
\[
  f[x,y,(zy)x+(xz)y]=(xy)^p f(x,y,z).
\]
Ainsi, de $f(x,y,\lambda_1x+\lambda_2y)=0$ (identiquement en $\lambda$), il résulte, comme le montre la substitution $\lambda_1=(zy)$, $\lambda_2=(xz)$, et puisque $(xy)\ne0$, nécessairement $f(x,y,z)=0$.

Par ce lemme il existe une correspondance biunivoque entre les formes de $\mathfrak S$ et celles de $\mathfrak L$. En effet, de
\[
  f_1(x,y,\lambda_1x+\lambda_2y)=g(\lambda;xy),\qquad
  f_2(x,y,\lambda_1x+\lambda_2y)=g(\lambda;xy)
\]
il s'ensuit nécessairement que
\[
  f_1(xyz)=f_2(xyz).
\]
De plus, toute relation dans $\mathfrak S$ est conservée entre les formes correspondantes uniques dans $\mathfrak L$ (et réciproquement). Car de
\[
  c_1g^{(1)}(\lambda;xy)+c_2g^{(2)}(\lambda;xy)+\cdots+c_\varkappa g^{(\varkappa)}(\lambda;xy)=0
\]
il s'ensuit nécessairement, par le lemme,
\[
  c_1f^{(1)}(x,y,z)+c_2f^{(2)}(x,y,z)+\cdots+c_\varkappa f^{(\varkappa)}(x,y,z)=0.
\]
\emph{Ainsi est démontrée l'égalité des rangs de $\mathfrak L$ et $\mathfrak S$; $\rho=\sigma$.}

\emph{b) Pour comparer les rangs de $\mathfrak S$ et de $\mathfrak T$, nous employons le lemme correspondant b): de $h(x,y,z)=0$, il s'ensuit nécessairement que $g(\lambda;xy)=0$.} Pour la preuve, remarquons que, d'après 3., $h(x,y,z)=0$ équivaut à l'annulation des produits de puissances individuels
\[
 \left(\frac{\partial}{\partial\lambda_1}\frac{\partial}{\partial x_1}
 +\frac{\partial}{\partial\lambda_2}\frac{\partial}{\partial y_1}\right)^{p_1}
 \left(\frac{\partial}{\partial\lambda_1}\frac{\partial}{\partial x_2}
 +\frac{\partial}{\partial\lambda_2}\frac{\partial}{\partial y_2}\right)^{p_2}g(\lambda;x,y),
 \qquad (p_1+p_2=p).
\]
Par différentiation ultérieure, on obtient aussi l'annulation des produits de puissances
\[
\begin{aligned}
&\left(\frac{\partial}{\partial x_1}\right)^{\alpha_1}
 \left(\frac{\partial}{\partial x_2}\right)^{\alpha_2}
 \left(\frac{\partial}{\partial y_1}\right)^{\beta_1}
 \left(\frac{\partial}{\partial y_2}\right)^{\beta_2} \\
&\quad\cdot
 \left(\frac{\partial}{\partial\lambda_1}\frac{\partial}{\partial x_1}
 +\cdots\right)^{p_1}
 \left(\frac{\partial}{\partial\lambda_1}\frac{\partial}{\partial x_2}
 +\cdots\right)^{p_2}
 g(\lambda,xy)
\end{aligned}
\]
Par conséquent, toute combinaison linéaire de ces produits de puissances s'annule, et donc toute forme
\[
  f\left(\frac{\partial}{\partial x},\frac{\partial}{\partial y},
  \frac{\partial}{\partial\lambda_1}\frac{\partial}{\partial x}
  +\frac{\partial}{\partial\lambda_2}\frac{\partial}{\partial y}\right)g(\lambda,xy).
\]
En choisissant $f$ spécialement de sorte que
\[
  f(x,y,\lambda_1x+\lambda_2y)=g(\lambda;xy),
\]
on obtient donc aussi
\[
  g\left(\frac{\partial}{\partial\lambda};
    \frac{\partial}{\partial x},
    \frac{\partial}{\partial y}\right)g(\lambda,x,y)=0;
\]
ou bien, si
\[
  g(\lambda,xy)=\sum G_i\lambda_1^{\varkappa_1}\lambda_2^{\varkappa_2}
  x_1^{\varkappa_3}x_2^{\varkappa_4}y_1^{\varkappa_5}y_2^{\varkappa_6},
\]
où les $G_i$ sont des formes linéaires à coefficients entiers rationnels dans les indéterminées $a$, alors
\[
  \sum \varkappa_1!\varkappa_2!\varkappa_3!\varkappa_4!\varkappa_5!\varkappa_6!\,G_i^2=0,
\]
et par suite $g(\lambda;xy)=0$.

Du lemme b), il s'ensuit exactement comme en a) que \emph{les rangs de $\mathfrak S$ et $\mathfrak T$ sont égaux; $\sigma=\tau$.}

D'après a) et b), on a donc $\rho=\tau$; par conséquent --- puisque $\mathfrak T$ est un sous-faisceau de $\mathfrak L$ --- l'équivalence de ces deux faisceaux linéaires est démontrée. Toute forme de $\mathfrak L$, en particulier aussi $\theta$, se représente donc comme combinaison linéaire à coefficients entiers rationnels de $\rho$ formes linéairement indépendantes $h(x,y,z)$:
\[
  \theta=\sum c_i h^{(i)}(x,y,z)=\sum PZ.
\]
Cette identité, obtenue pour les coefficients indéterminés $a$ de $\theta$, demeure valable lorsqu'on remplace les indéterminées $a$ par des grandeurs d'un domaine de rationalité quelconque; ainsi \emph{le théorème de réduction est démontré en général dans le cas de trois rangées binaires.}

La preuve générale, pour $N$ rangées de $n$ variables chacune, est entièrement parallèle. Que $\theta(A)$ ait la dimension $\alpha_i$ dans la rangée $A_i$. Nous considérons de nouveau les trois faisceaux linéaires de formes:

1. le faisceau $\mathfrak L$, formé de toutes les formes $f(A)$ qui proviennent de $\theta(A)$ par des procédés polaires $P$ et qui, dans chaque rangée individuelle $A_i$, ont la même dimension que $\theta(A)$;

2. le faisceau $\mathfrak S$, formé de toutes les formes $g(\lambda;A_1\cdots A_n)$ qui proviennent de $f(A)$ par les substitutions
\[
\begin{aligned}
 A_{n+1}&=\lambda_{n+1}^{(1)}A_1+\lambda_{n+1}^{(2)}A_2+
        \cdots+\lambda_{n+1}^{(n)}A_n\\
 &\vdotswithin{=}\\
 A_N&=\lambda_N^{(1)}A_1+\lambda_N^{(2)}A_2+
        \cdots+\lambda_N^{(n)}A_n
\end{aligned}
\]
à cette occasion, les coefficients des produits de puissances individuels des $\lambda$ dans $g(\lambda;A_1\cdots A_n)$ donnent des formes $Z$ de l'identité (1);

3. le faisceau $\mathfrak T$, formé de toutes les formes $h(A)$ définies par
\[
\begin{aligned}
 h(A)=&\left[A_{n+1}\left(
 \frac{\partial}{\partial\lambda_{n+1}^{(1)}}\frac{\partial}{\partial A_1}
 +\cdots+
 \frac{\partial}{\partial\lambda_{n+1}^{(n)}}\frac{\partial}{\partial A_n}\right)\right]^{\alpha_{n+1}} \\
&\cdots
\left[A_N\left(
 \frac{\partial}{\partial\lambda_N^{(1)}}\frac{\partial}{\partial A_1}
 +\cdots+
 \frac{\partial}{\partial\lambda_N^{(n)}}\frac{\partial}{\partial A_n}\right)\right]^{\alpha_N}
 g(\lambda;A_1\cdots A_n).
\end{aligned}
\]
Ici $h(A)$ est une somme de formes $PZ$.

À cause de la relation identique entre $n+1$ rangées quelconques, le lemme a) subsiste; de même le lemme b) reste valable, puisque le nombre de variables n'est pas intervenu dans sa preuve. Les faisceaux $\mathfrak L$ et $\mathfrak T$ ont donc le même rang, et --- puisque $\mathfrak T$ est un sous-faisceau de $\mathfrak L$ --- ils sont identiques. Par conséquent l'identité
\[
  \theta=\sum PZ
\]
vaut pour des coefficients indéterminés, et par suite pour des coefficients qui sont des grandeurs d'un domaine de rationalité quelconque; \emph{le théorème de réduction est ainsi démontré en général.}

\subsection*{III.}
Nous indiquons encore brièvement comment des raisonnements analogues donnent aussi le \emph{développement en séries général} (pour $N\le n$). Soit $Z$ une forme séparément homogène dans les $n$ rangées $A_1,\ldots,A_n$ de $n$ grandeurs chacune; et soit $\Delta$ le déterminant de ces rangées. Nous prouvons d'abord l'identité correspondant à (1):
\begin{equation}
  Z\equiv\sum PH \mod{\Delta},
\tag{2}
\end{equation}
où les formes $H$ ne contiennent plus que les rangées $A_1\cdots A_{n-1}$ et sont déduites de $Z$ par des procédés $P$.

La preuve consiste à transformer les relations déduites en II en congruences modulo $\Delta$. Pour simplifier, prenons pour base trois rangées ternaires $x,y,z$; les coefficients sont supposés indéterminés.

Les trois faisceaux linéaires $\mathfrak L,\mathfrak S,\mathfrak T$ sont définis comme en II dans le cas des rangées binaires; soient leurs rangs $\rho,\sigma,\tau$, tandis que $\rho^*$ et $\tau^*$ désignent les rangs de $\mathfrak L$ et $\mathfrak T$ modulo $\Delta$ (c'est-à-dire qu'il y a $\rho^*$, mais non $\rho^*+1$, formes de $\mathfrak L$ linéairement indépendantes mod $\Delta$). Alors, comme en II, on a

\emph{Lemme a): de $f(x,y,\lambda_1x+\lambda_2y)=0$, il s'ensuit nécessairement que $f(xyz)\equiv0\pmod{\Delta}$ (et réciproquement).} En effet, par la relation entre quatre rangées ternaires,
\[
  \Delta_1x-\Delta_2y+\Delta_3z=\Delta t,
  \qquad \Delta_1=(yzt)\cdots,
\]
il existe toujours entre trois rangées ternaires une relation linéaire modulo $\Delta$:
\[
  \Delta_1x-\Delta_2y+\Delta_3z\equiv0\mod{\Delta},
\]
et l'on a par conséquent, comme en II,
\[
  f(x,y,-\Delta_1x+\Delta_2y)\equiv\Delta_3^p f(x,y,z)\pmod{\Delta}.
\]
De $f(x,y,\lambda_1x+\lambda_2y)=0$ (identiquement en $\lambda$), il résulte donc --- puisque $\Delta$ est irréductible et que $\Delta_3$ n'est pas divisible par $\Delta$ --- nécessairement $f(x,y,z)\equiv0\pmod{\Delta}$. La réciproque est claire, car $\Delta(x,y,\lambda_1x+\lambda_2y)=0$. On conclut alors comme en II que $\rho^*=\sigma$.

Le lemme b) subsiste avec sa preuve, et l'on a donc $\sigma=\tau$.

Pour démontrer $\tau=\tau^*$, on emploie le

\emph{Lemme c): la congruence $h(x,y,z)\equiv0\pmod{\Delta}$ entraîne nécessairement $h(x,y,z)=0$.} En effet, comme polaire, $h(x,y,z)$ s'annule sous l'application du procédé $\Omega$ connu; cela s'exprime par l'équation différentielle
\[
  \Delta\!\left(\frac{\partial}{\partial x},\frac{\partial}{\partial y},\frac{\partial}{\partial z}\right)h(xyz)
  =\left(\frac{\partial}{\partial x}\frac{\partial}{\partial y}\frac{\partial}{\partial z}\right)h(x,y,z)=0.\srcfnmark{*)}
\]
\srcfntext{*)}{D'après (3), on a
\[
c\cdot\Omega(h)=\left(
\frac{\partial}{\partial \xi}\frac{\partial}{\partial \eta}
\left(\frac{\partial}{\partial\lambda_1}\frac{\partial}{\partial \xi}
+\frac{\partial}{\partial\lambda_2}\frac{\partial}{\partial \eta}\right)
\right)
\left(x\frac{\partial}{\partial \xi}\right)^{\alpha-1}
\left(y\frac{\partial}{\partial \eta}\right)^{\beta-1}
\left[
z\left(\frac{\partial}{\partial\lambda_1}\frac{\partial}{\partial \xi}
+\frac{\partial}{\partial\lambda_2}\frac{\partial}{\partial \eta}\right)
\right]^{p-1}\!\cdot g(\lambda;\xi,\eta),
\]
et $\Omega(h)$ s'annule donc en raison de l'annulation du facteur déterminant.}
Si maintenant $h(x,y,z)=\Delta(x,y,z)\cdot k(x,y,z)$, alors, par différentiation ultérieure, on a aussi l'annulation de
\[
  k\!\left(\frac{\partial}{\partial x},\frac{\partial}{\partial y},\frac{\partial}{\partial z}\right)
  \Delta\!\left(\frac{\partial}{\partial x},\frac{\partial}{\partial y},\frac{\partial}{\partial z}\right)
  =
  h\!\left(\frac{\partial}{\partial x},\frac{\partial}{\partial y},\frac{\partial}{\partial z}\right)h(x,y,z),
\]
et par conséquent celle de $h(x,y,z)$. On a donc $\rho^*=\sigma=\tau=\tau^*$; comme $\mathfrak T$ est un sous-faisceau de $\mathfrak L$, l'identité (2) est démontrée. Le même raisonnement donne la preuve pour $n$ variables.

En appliquant l'identité (2) au multiplicateur de $\Delta$, on obtient un développement
\[
  Z=\varphi_0+\Delta\varphi_1+\Delta^2\varphi_2+\cdots+\Delta^\mu\varphi_\mu,
\]
où les $\varphi_i$ s'annulent chacun sous l'application du procédé $\Omega$. La répétition $i$ fois du procédé donne, puisque l'on a en général $\Omega(\Delta\cdot\psi)=c_1\psi+c_2\Delta\cdot\Omega(\psi)$,
\[
  c\cdot \Omega^i(Z)\equiv\varphi_i\pmod{\Delta};
\]
d'autre part, d'après (2), on a
\[
  c\cdot\Omega^i(Z)\equiv\sum PH_i\pmod{\Delta},
\]
où les $H_i$ sont déduites de la forme $\Omega^i(Z)$ de la même façon que $H$ est déduite de $Z$. D'après le lemme c) (étendu à $n$ variables), on obtient donc $\varphi_i=\sum PH_i$, et par conséquent
\begin{equation}
  Z=\sum PH+\Delta\cdot\sum PH_1+\Delta^2\cdot\sum PH_2+
  \cdots\Delta^\mu\cdot\sum PH_\mu
\tag{3}
\end{equation}
\emph{comme le développement en séries le plus général d'une forme de $n$ rangées de $n$ variables chacune, suivant les polaires de formes de seulement $n-1$ rangées.}

Le développement en séries pour les formes de $N$ rangées de $n$ variables $(N<n)$ s'obtient alors, comme on sait, par une substitution simple. Nous posons (pour $N=3$, $n>3$)
\[
  u=\xi_1x+\xi_2y+\xi_3z,
  \quad v=\eta_1x+\eta_2y+\eta_3z,
  \quad w=\zeta_1x+\zeta_2y+\zeta_3z,
\]
et développons $H(u,v,w)=Z(\xi,\eta,\zeta)$ d'après (3). Comme $\xi,\eta,\zeta$ n'interviennent alors que dans les combinaisons $u,v,w$, on a
\[
  \left(\eta\frac{\partial}{\partial \xi}\right)=
  \left(v\frac{\partial}{\partial u}\right)\cdots,
  \qquad
  \Omega=\sum\left(\frac{\partial}{\partial u}
  \frac{\partial}{\partial v}
  \frac{\partial}{\partial w}\right)_{ikl}(xyz)_{ikl}
  =\nabla_{uvw}(xyz).
\]
Par la spécialisation $\xi_1=\eta_2=\zeta_3=1$; $\xi_2=\xi_3=\cdots=\zeta_2=0$, $H(u,v,w)$ devient $H(x,y,z)$, et $\nabla_{uvw}(xyz)$ devient, à un facteur numérique près, $\nabla_{xyz}(xyz)=\nabla_{xyz}$, car l'application de $(y\partial/\partial x)$, etc., aux déterminants $(xyz)_{ikl}$ les fait s'annuler. Comme l'assertion correspondante vaut aussi pour $N$ rangées, il résulte de (3) un développement
\begin{equation}
  H=\sum P\Phi+\sum P\Phi_1+\cdots+\sum P\Phi_\mu,
\tag{4}
\end{equation}
où les $\Phi_i$ proviennent de $\nabla^i_{A_1\ldots A_N}H$ par des procédés $P$ et ne contiennent plus que les rangées $A_1,\ldots,A_{N-1}$, abstraction faite des combinaisons de déterminants qui apparaissent dans $\nabla^i$ et qui, comme on l'a mentionné, n'entrent pas en considération lors de la polarisation.

Si l'on remplace alors ces déterminants dans $\nabla^i$ par des indéterminées $p$, les formes qui en résultent peuvent de nouveau être développées d'après (4); on obtient ainsi finalement un développement
\begin{equation}
  H=\left[\sum P\psi(p)\right]_{p=\text{déterminants de }A_1\cdots A_N},
\tag{5}
\end{equation}
où les $\psi$ proviennent de la forme initiale par les procédés $P$ et $\nabla_{A_1\ldots A_N}(p)$, $\nabla_{A_1\ldots A_{N-1}}(p)$, etc. Avec ces développements et leurs combinaisons --- par exemple en appliquant (3), puis (4) et (5), aux formes $Z$ qui apparaissent dans (1) --- tous les développements connus pour les formes d'un nombre arbitraire de rangées de $n$ variables cogrédientes chacune sont épuisés.

Erlangen, 5 janvier 1915.
\clearpage
\setcounter{footnote}{0}
\begin{center}
{\Large\bfseries \NoetherNavAnchor{P09}{9. Les domaines les plus généraux d’entiers transcendants}9. Les domaines les plus généraux d'entiers transcendants}\par
\vspace{0.75em}
Par Emmy Noether, à Erlangen.\par
\vspace{1.2em}
\emph{Math. Ann.} 77 (1916), p. 103--128
\end{center}

\vspace{1.5em}

Au moyen du théorème du bon ordre, E. Zermelo a construit un domaine d'\og{} entiers transcendants \fg{}, c'est-à-dire un domaine numérique $\mathfrak G$ auquel appartiennent les propriétés abstraites suivantes :
\begin{enumerate}
\item[I.] La somme, la différence et le produit de deux nombres de $\mathfrak G$ sont encore des nombres de $\mathfrak G$.
\item[II.] Tout nombre réel ou complexe est quotient de deux nombres de $\mathfrak G$.
\item[III.] Tout nombre rationnel entier, respectivement algébrique entier, est un élément de $\mathfrak G$.
\item[IV.] Aucun nombre rationnel, respectivement algébrique, non entier n'appartient à $\mathfrak G$.\footnote{E. Zermelo, \emph{Über ganze transzendente Zahlen}, Math. Ann. 75, p. 434 (1914).}
\end{enumerate}

La construction repose sur le fait que le théorème du bon ordre permet de reconnaître l'existence d'une \emph{base algébrique} de tous les nombres, c'est-à-dire d'un \emph{système} $\mathsf{H}$ \emph{de nombres} $\eta$ entre lesquels il n'existe aucune relation algébrique, mais au moyen desquels tous les autres nombres peuvent s'exprimer algébriquement. En raison de leur indépendance algébrique, ces nombres de base $\eta$ peuvent être regardés comme des indéterminées; le domaine recherché devient donc un domaine entier de fonctions rationnelles et algébriques des indéterminées $\eta$,\footnote{C'est pourquoi les considérations du présent article sont en partie parallèles à celles de l'article antérieur de l'autrice : \emph{Körper und Systeme rationaler Funktionen}, Math. Ann. 76, p. 161 (1915).} dont les coefficients appartiennent au corps $\mathsf{K}$ de tous les nombres algébriques. Le domaine spécial $\mathfrak G_\eta$ construit par Zermelo est alors caractérisé par les propriétés de base suivantes :

$\mathfrak G_\eta$ contient :
\begin{enumerate}
\item[1)] tous les nombres de base $\eta$;
\item[2)] tous les polynômes en les $\eta$ à coefficients entiers;\footnote{Dans tout ce qui suit, \og{} entier \fg{} signifie que les coefficients appartiennent au domaine entier $[\mathsf{K}]$ de tous les entiers algébriques. Pour la définition de $\mathfrak G_\eta$, il suffirait aussi, dans 1)--5), de ne considérer que des coefficients entiers rationnels; pour des domaines plus généraux, les propriétés de base deviendraient cependant moins simples.}
\item[3)] toutes les fonctions algébriques entières des $\eta$ à coefficients entiers.
\end{enumerate}
$\mathfrak G_\eta$ exclut :
\begin{enumerate}
\item[4)] toutes les fonctions rationnellement fractionnaires des $\eta$ à coefficients entiers;
\item[5)] toutes les fonctions algébriquement fractionnaires des $\eta$ à coefficients entiers.\footnote{Zermelo, loc. cit., \S{} 3.}
\end{enumerate}

Il est maintenant facile de voir, en comparant les exemples des \S{}\S{} 2 et 3, qu'il existe des domaines $\mathfrak G$ essentiellement différents du domaine de Zermelo $\mathfrak G_\eta$ : c'est-à-dire des domaines $\mathfrak G$ pour lesquels, sous aucun bon ordre, toutes les propriétés de base 1)--5) ne sont satisfaites, et qui par conséquent ne peuvent être produits par la construction de Zermelo sous aucun bon ordre. Autrement dit, il existe des domaines $\mathfrak G$ qui ne peuvent pas être appliqués biunivoquement et isomorphiquement sur $\mathfrak G_\eta$. Ainsi se pose la question des domaines les plus généraux d'entiers transcendants; elle sera complètement résolue dans ce qui suit.

À cette fin il faut d'abord déterminer lesquelles des propriétés de base sont des conséquences des propriétés abstraites de définition et lesquelles proviennent de la construction spéciale. On montrera (\S{} 1) que, pour tout domaine $\mathfrak G$ prescrit, il existe un bon ordre tel que les propriétés de base 1) et 2) soient satisfaites; ainsi tout domaine $\mathfrak G$ peut être construit au moyen d'une base algébrique de tous les nombres. Aucune des autres propriétés de base n'a besoin d'être satisfaite (\S{}\S{} 2 et 3). En particulier il s'ensuit qu'une autre propriété abstraite du domaine de Zermelo n'appartient pas à tous les domaines $\mathfrak G$ : la clôture algébrique-entière. Elle peut se formuler ainsi :
\begin{enumerate}
\item[V.] Tout nombre qui est entier sur le sous-anneau engendré par un nombre fini de nombres de $\mathfrak G$ appartient à $\mathfrak G$.
\end{enumerate}

Les domaines spéciaux pour lesquels V vaut en plus de I--IV seront appelés domaines $\mathfrak H$ d'entiers transcendants algébriquement-entiers. Pour les domaines les plus généraux $\mathfrak H$, on obtient les résultats simples suivants (\S{} 4) :
\begin{enumerate}
\item[a)] L'intersection de tous les domaines $\mathfrak H$ constructibles au moyen d'une même base $H$ est donnée par toutes les fonctions algébriques entières de la base, c'est-à-dire par le domaine de Zermelo $\mathfrak G_\eta$, qui est lui-même un domaine $\mathfrak H$. Cela donne aussi une caractérisation abstraite du domaine de Zermelo.
\item[b)] Tout domaine $\mathfrak H$ résulte de l'adjonction algébriquement entière\footnote{Explication de l'expression au \S{} 4.} d'un système arbitraire $\mathfrak S(\eta)$ à $\mathfrak G_\eta$, où $\mathfrak S(\eta)$ ne doit satisfaire qu'à la seule condition qu'aucun nombre algébriquement fractionnaire n'entre dans le domaine par cette adjonction.
\end{enumerate}

Les résultats relatifs aux domaines les plus généraux $\mathfrak G$, lorsqu'on prend pour fondement une base algébrique, sont moins simples. Ici l'intersection de tous les domaines $\mathfrak G$ n'est elle-même pas un domaine $\mathfrak G$, et la construction des domaines les plus généraux devient par conséquent plus difficile à embrasser (\S{} 5).

Afin d'obtenir des résultats analogues à ceux des domaines $\mathfrak H$, la base algébrique est remplacée par une \emph{base rationnelle} $\Theta$, c'est-à-dire par un \emph{système} $\Theta$ \emph{de nombres} $\vartheta$ qui permet d'exprimer tous les nombres rationnellement, tandis que, sous au moins un bon ordre, aucun nombre de base ne peut être exprimé rationnellement par un nombre fini de précédents. Cette base rationnelle $\Theta$ doit encore être soumise à la condition accessoire que le domaine entier $[\Theta]$, formé de tous les polynômes entiers en les $\vartheta$, ne contienne aucun nombre algébriquement fractionnaire. La construction d'une telle base avec condition accessoire est donnée au \S{} 7. Alors, pour les domaines les plus généraux $\mathfrak G$ --- que l'on pourrait aussi appeler domaines d'entiers transcendants rationnellement entiers ---, on obtient à nouveau les résultats simples suivants (\S{} 6) :
\begin{enumerate}
\item[a)] L'intersection de tous les domaines $\mathfrak G$ constructibles au moyen de la même base rationnelle $\Theta$ est donnée par toutes les fonctions rationnelles entières de la base, c'est-à-dire par le domaine $[\Theta]$, qui est lui-même un domaine $\mathfrak G$.
\item[b)] Tout domaine $\mathfrak G$ résulte de l'adjonction rationnellement entière d'un système arbitraire $\mathfrak S(\vartheta)$ à $[\Theta]$, où $\mathfrak S(\vartheta)$ ne doit satisfaire qu'à la seule condition qu'aucun nombre algébriquement fractionnaire n'entre dans le domaine par cette adjonction.
\end{enumerate}

Pour avoir une analogie complète avec les domaines $\mathfrak H$, il faudrait retrancher les nombres algébriques des propriétés de définition III et IV des domaines $\mathfrak G$. Avec cette modification de III et IV, la construction des éléments entiers au moyen d'une base rationnelle peut être transférée à un corps arbitraire (\S{} 10), tandis que la construction de Zermelo présuppose que le corps soit algébriquement clos.

Lorsque $\mathsf{H}$ parcourt toutes les bases algébriques, respectivement lorsque $\Theta$ parcourt toutes les bases rationnelles satisfaisant à la condition accessoire, les résultats a) et b) donnent la totalité de tous les domaines $\mathfrak H$ et $\mathfrak G$. Si l'on réunit en une classe tous les domaines isomorphes, on peut alors extraire de cette totalité un sous-ensemble qui représente la totalité de toutes les classes --- au moins dénombrablement infinies d'après le \S{} 9 (\S{} 8).

\subsection*{\S{} 1. Démonstration des propriétés de base 1) et 2) pour tous les domaines $\mathfrak G$.}

Soit $\mathfrak G$ un domaine arbitrairement prescrit d'entiers transcendants, de sorte qu'il possède les propriétés abstraites I--IV. En suivant le procédé appliqué par E. Zermelo au corps de tous les nombres complexes, choisissons dans $\mathfrak G$ une \emph{base algébrique} $\mathsf{H}$ \emph{de tous les nombres de $\mathfrak G$}. Comme, d'après II, tout nombre réel ou complexe peut être représenté comme quotient de deux nombres de $\mathfrak G$, $\mathsf{H}$ est en même temps une base algébrique de \emph{tous} les nombres. Ainsi, pour tout domaine $\mathfrak G$ prescrit, la base algébrique de tous les nombres peut être choisie de façon à appartenir à $\mathfrak G$; la propriété de base 1) est donc satisfaite. De plus, par III, $\mathfrak G$ contient le domaine entier $[\mathsf{K}]$ de tous les entiers algébriques, et par conséquent, par I, aussi tous les polynômes en les éléments de base $\eta$ à coefficients dans $[\mathsf{K}]$, c'est-à-dire tous les polynômes entiers en les $\eta$, qui forment le domaine entier $[\mathsf{H}]$. La propriété de base 2) est donc elle aussi toujours satisfaite; autrement dit, tout domaine $\mathfrak G$ peut être construit au moyen d'une base algébrique $\mathsf{H}$ de tous les nombres de telle manière que $\mathfrak G$ contienne le domaine entier $[\mathsf{H}]$ comme diviseur.

\textbf{Remarque.} Néanmoins, en partant d'une base $\mathsf{H}$, on peut naturellement construire un domaine $\mathfrak G$ pour lequel les propriétés de base 1) et 2) ne valent pas relativement à $\mathsf{H}$ elle-même, mais valent relativement à une autre base $\mathsf{H}'$. Par exemple, que $\mathfrak G'$ provienne du domaine de Zermelo $\mathfrak G_\eta$ en remplaçant, dans toute la construction, un certain élément de base $\eta$ par un polynôme $f(\eta)$. Alors $\mathfrak G'$ ne contient ni $\eta$ ni $[\mathsf{H}]$; mais si, dans la base $\mathsf{H}$, on remplace l'élément de base $\eta$ par $\zeta=f(\eta)$ --- ce par quoi $\mathsf{H}$ passe de nouveau à une base algébrique $\mathsf{H}'$ ---, les propriétés de base 1) et 2) sont satisfaites pour $\mathfrak G'$ relativement à $\mathsf{H}'$.

\subsection*{\S{} 2. Exclusion des propriétés de base 4) et 5).}

Nous montrons maintenant que les domaines $\mathfrak G$ n'ont à satisfaire aucune des autres propriétés de base. D'abord 4) et 5) sont exclues en remplaçant, dans la construction de Zermelo, le domaine entier $[\mathsf{H}]$ par un domaine de \og{} fonctionnelles \fg{} rationnelles entières.

D'après Weber,\footnote{H. Weber, \emph{Lehrbuch der Algebra}. Petite édition, \S{}\S{} 96 et 97.} on appelle \emph{fonctionnelle rationnelle} toute fonction rationnelle à coefficients rationnels écrite sous la représentation univoquement déterminée suivante :
\[
 A(\eta_1\cdots \eta_t)
   = a\cdot\frac{E_1(\eta_1\cdots \eta_t)}{E_2(\eta_1\cdots \eta_t)},
\]
où $E_1,E_2$ sont des polynômes primitifs premiers entre eux à coefficients entiers rationnels, et où $a$ est un nombre rationnel positif, la valeur absolue de $A$. Si, en particulier, $a$ est un \emph{entier rationnel}, alors $A$ s'appelle une \emph{fonctionnelle rationnelle entière}. Comme cas particuliers, les fonctionnelles rationnelles entières comprennent les entiers rationnels et les polynômes à coefficients entiers rationnels; tandis que les nombres rationnellement fractionnaires et les polynômes à coefficients rationnellement fractionnaires doivent être rangés parmi les fonctionnelles rationnelles fractionnaires.

Soit maintenant le domaine $\mathfrak G$ constitué de la totalité $\mathfrak J$ de toutes les fonctionnelles rationnelles entières d'un nombre fini quelconque d'éléments de base $\eta$, et de toutes les fonctions qui sont entières sur $\mathfrak J$, c'est-à-dire qui satisfont à une équation dont le coefficient dominant est l'unité et dont les autres coefficients sont des fonctionnelles rationnelles entières.

La démonstration des propriétés I--IV pour $\mathfrak G$ est exactement parallèle à la démonstration correspondante chez Zermelo et ne sera qu'indiquée brièvement. D'abord II et III sont claires, puisque $\mathfrak G$ contient le domaine de Zermelo $\mathfrak G_\eta$ comme diviseur. Par le théorème de Gauss sur les polynômes à coefficients entiers rationnels, les sommes, différences et produits de fonctionnelles rationnelles entières sont encore des fonctionnelles rationnelles entières; donc $\mathfrak J$ forme un domaine entier, et par les arguments connus\footnote{Comparer par exemple Weber, \S{} 97, 8.} $\mathfrak G$ devient lui aussi un domaine entier; ainsi I est satisfaite. En outre, du théorème de Gauss étendu\footnote{Weber, \S{} 20, 5.} résultent les deux propositions auxiliaires suivantes : \og{} Si $z$ est entier sur $\mathfrak J$ en vertu d'une \emph{certaine} équation, alors il l'est aussi en vertu de l'équation \emph{irréductible} que $z$ satisfait relativement au corps de toutes les fonctionnelles rationnelles \fg{},\footnote{Comparer Weber, \S{} 97, 4.} et \og{} l'équation irréductible sur le corps de tous les nombres rationnels et satisfaite par un nombre algébrique reste irréductible sur le corps de toutes les fonctionnelles rationnelles \fg{}. Ainsi $\mathfrak G$ ne peut contenir aucun nombre algébrique fractionnaire; IV, et donc toutes les propriétés abstraites I--IV, sont vérifiées.

D'autre part, sous aucun bon ordre on ne peut choisir dans $\mathfrak G$ une base telle que $\mathfrak G$ exclue toutes les fonctions rationnellement et algébriquement fractionnaires des éléments de base. Soit en effet $z(\eta)$ une fonction quelconque du domaine qui doit être choisie comme élément de base, donc définie par une équation
\[
 z^k+A_1(\eta)z^{k-1}+\cdots+A_k(\eta)=0,
\]
où
\[
 A_i(\eta)=a_i\cdot\frac{E_1(\eta)}{E_2(\eta)}
\]
sont des fonctionnelles rationnelles entières. Alors la fonction $\frac{a_k}{z}$ appartient à $\mathfrak G$, de même que $\frac{a_k}{\sqrt z}$, $\frac{a_k}{\sqrt[3]z}$, etc.\footnote{Tout à fait de même, toute fonction algébrique des $\eta$ peut être transformée en élément du domaine fonctionnel $\mathfrak G$ par multiplication par un entier rationnel. Ce domaine possède ainsi la propriété que tout nombre complexe peut être représenté comme quotient d'un élément de $\mathfrak G$ et d'un entier rationnel, en analogie avec les nombres algébriques.} Les propriétés de base 4) et 5) sont donc exclues pour la totalité des domaines $\mathfrak G$.

\textbf{Remarque.} Un exemple au \S{} 3 montrera que $\mathfrak G$ peut, sous tout bon ordre, contenir des fonctionnelles rationnellement fractionnaires sans contenir un nombre rationnellement ou algébriquement fractionnaire.

\subsection*{\S{} 3. Exclusion de la propriété de base 3).}

Pour exclure la propriété de base 3), nous utilisons un \emph{domaine de module} généralisé dans lequel les arguments des polynômes du domaine ne sont plus les indéterminées, mais toutes les fonctions algébriques entières des indéterminées,\footnote{Par \og{} fonctions algébriques entières \fg{} on entend toujours des fonctions à coefficients \emph{entiers}; c'est-à-dire des fonctions satisfaisant à une équation dont le coefficient dominant est l'unité et dont les autres coefficients sont des polynômes en les $\eta$ à coefficients entiers algébriques, donc des polynômes de $[\mathsf{H}]$. En particulier, tous les polynômes de $[\mathsf{H}]$ appartiennent aux fonctions algébriques entières.} tandis que les indéterminées elles-mêmes jouent le rôle d'une base du module.

Que $\mathfrak G$ consiste en \emph{tous les éléments}
\begin{equation}
 z=\alpha+g_1(\eta)\cdot\eta_1+g_2(\eta)\cdot\eta_2+\cdots+g_t(\eta)\cdot\eta_t,
\tag{1}
\end{equation}
\emph{où $\alpha$ parcourt tous les entiers algébriques, $\eta_1\cdots\eta_t$ tous les éléments de base, et où, pour chaque combinaison fixe $\eta_1\cdots\eta_t$, les fonctions $g_1(\eta)\cdots g_t(\eta)$ parcourent individuellement toutes les fonctions algébriques entières des $\eta$.}\footnote{La propriété de module appartient aux éléments $z-\alpha$.}

Il est facile de vérifier que les propriétés I--IV valent pour ce domaine de module. D'abord I vaut, puisque la somme, la différence et le produit de deux éléments $z$ prennent de nouveau la forme (1). Le domaine $\mathfrak G$ contient aussi tous les éléments de base $\eta$, et en outre tous les éléments $\eta\cdot g(\eta)$, où $g(\eta)$ parcourt toutes les fonctions algébriques entières des $\eta$. Par conséquent toutes les fonctions algébriques entières des $\eta$, et donc toutes les fonctions algébriques des $\eta$ en général, peuvent être représentées comme quotients de deux éléments de $\mathfrak G$, ce qui prouve II. Dans la représentation (1), tous les entiers algébriques sont en particulier inclus --- pour $g_1=0\cdots g_t=0$ ---; mais $z$ ne peut pas être égal à un nombre algébrique fractionnaire. Car si $z$ représente un nombre algébrique $\beta$, alors de
\[
 \beta=\alpha+g_1(\eta)\eta_1+\cdots+g_t(\eta)\eta_t
\]
identiquement en les $\eta$, il résulte nécessairement $\beta=\alpha$. Cela se voit par la spécialisation
\[
 \eta_1=0\cdots \eta_t=0,
\]
sous laquelle les multiplicateurs $g_1(\eta)\cdots g_t(\eta)$ restent finis comme fonctions algébriques entières, puisqu'ils passent en des fonctions ou nombres algébriques entiers.\footnote{Cette preuve plus détaillée de IV est donnée à cause de l'exemple élargi à la fin du paragraphe. Ici suffirait la remarque que, par construction, $z$ est une fonction algébrique entière.} Ainsi les conditions I--IV sont satisfaites pour $\mathfrak G$.

D'autre part, sous aucun bon ordre on ne peut choisir dans $\mathfrak G$ une base telle que toutes les fonctions algébriques entières des éléments de base appartiennent à $\mathfrak G$. Soit $z(\eta)$ une fonction quelconque du domaine qui doit être choisie comme élément de base --- et qui doit donc effectivement contenir les indéterminées $\eta$. Nous montrerons qu'à partir d'une certaine valeur finie de $\nu$, dépendant de $z$, les fonctions
\[
 \xi=\sqrt[\nu]{z-\alpha}
\]
ne peuvent plus appartenir à $\mathfrak G$.

Supposons en effet que $\xi$ appartienne à $\mathfrak G$. Alors, par (1), il est de la forme
\[
 \zeta=\gamma+h_1(\eta)\eta_{i_1}+h_2(\eta)\eta_{i_2}
      +\cdots+h_\tau(\eta)\eta_{i_\tau},
\]
où $h_1(\eta)\cdots h_\tau(\eta)$ sont encore des fonctions algébriques entières des $\eta$, et où $\eta_{i_1}\cdots\eta_{i_\tau}$ sont des éléments de base quelconques, qui peuvent en particulier coïncider aussi avec $\eta_1\cdots\eta_t$. Il faut d'abord montrer que l'entier algébrique $\gamma$ est nul. Comme $h_1(\eta)\cdots h_\tau(\eta)$ sont des fonctions algébriques entières, $\xi$ devient égal à $\gamma$ pour $\eta_{i_1}=0\cdots\eta_{i_\tau}=0$, avec des valeurs arbitraires des autres $\eta$; en particulier, $\xi$ est alors égal à $\gamma$ pour
\[
 \eta_{i_1}=0\cdots \eta_{i_\tau}=0; \qquad
 \eta_1=0\cdots \eta_t=0.
\]
Mais sous cette spécialisation la fonction $(z-\alpha)$ s'annule par (1), et donc $\xi$ aussi; par conséquent $\gamma=0$, et l'on a
\[
 \zeta=h_1(\eta)\eta_{i_1}+h_2(\eta)\eta_{i_2}
      +\cdots+h_\tau(\eta)\eta_{i_\tau}.
\]
En élevant à la puissance $\nu$-ième, on obtient
\[
 \zeta^\nu=z-\alpha=F_\nu(\eta),
\]
où $F_\nu(\eta)$ désigne une forme homogène, non identiquement nulle, de dimension $\nu$-ième en $\eta_{i_1}\cdots\eta_{i_\tau}$, dont les coefficients sont des fonctions algébriques entières des $\eta$. Or $z-\alpha$, qui par (1) est une fonction algébrique entière des $\eta$, satisfait à une équation irréductible relativement à $[\mathsf{H}]$ :
\[
 (z-\alpha)^\varkappa+B(\eta)(z-\alpha)^{\varkappa-1}+\cdots+C(\eta)=0,
\]
où le dernier coefficient $C(\eta)$ --- produit de $(z-\alpha)$ avec ses conjugués --- est un polynôme; soit $\lambda$ son degré. D'autre part, de $z-\alpha=F_\nu(\eta)$, en multipliant $F_\nu(\eta)$ par les conjugués correspondants, on obtient
\[
 C(\eta)=G_{\nu\varkappa}(\eta),
\]
où $G_{\nu\varkappa}(\eta)$ désigne une forme homogène de dimension $\nu\varkappa$-ième en $\eta_{i_1}\cdots\eta_{i_\tau}$, dont les coefficients sont des polynômes en les $\eta$. Ainsi $G_{\nu\varkappa}$ est d'un degré au moins $\nu\varkappa$ en les $\eta$, ou $\lambda\geqq\nu\varkappa$; c'est-à-dire que les fonctions $\sqrt[\nu]{z-\alpha}$ pour lesquelles $\nu>\frac{\lambda}{\varkappa}$ n'appartiennent pas à $\mathfrak G$. La propriété de base 3) est donc exclue pour la totalité des domaines $\mathfrak G$.

\textbf{Remarque.} Une légère généralisation du domaine qui vient d'être construit conduit à un domaine $\mathfrak G$ qui, sous tout bon ordre, contient aussi des fonctionnelles fractionnaires. Que $\mathfrak G$ consiste en tous les éléments
\[
 z=\alpha+\gamma_1(\eta)\eta_1+\gamma_2(\eta)\eta_2
      +\cdots+\gamma_t(\eta)\eta_t,
\]
où maintenant $\gamma(\eta)$ parcourt toutes les fonctions algébriquement entières, mais non nécessairement entières, des $\eta$. La preuve des propriétés I--IV reste la même que ci-dessus. Mais puisque $\mathfrak G$, outre toute fonction $z$, contient aussi $a\cdot(z-\alpha)$, où $a$ désigne un nombre rationnel ou algébriquement fractionnaire quelconque, des fonctionnelles fractionnaires apparaissent également sous tout bon ordre.

\subsection*{\S{} 4. Les domaines les plus généraux d'entiers transcendants algébriquement-entiers.}

Pour formuler plus commodément ce qui suit, nous dirons qu'un élément $z$ est \emph{entier sur $\mathfrak T$} lorsque $z$ satisfait à une équation dont le coefficient dominant est l'unité et dont les autres coefficients sont des polynômes entiers en les éléments de $\mathfrak T$, où $\mathfrak T$ désigne un domaine quelconque prescrit.\footnote{Pour un domaine entier $\mathfrak T$, le coefficient dominant est l'unité et les autres coefficients sont des éléments de $\mathfrak T$.}

Un domaine $\mathfrak T$ s'appelle \emph{intégralement clos} s'il satisfait à la condition, comparer l'introduction :
\begin{enumerate}
\item[V.] Tout élément entier sur $\mathfrak T$ appartient à $\mathfrak T$.
\end{enumerate}

Nous montrons d'abord que la propriété abstraite V appartient au domaine de Zermelo $\mathfrak G_\eta$. Soit un élément $z$ entier sur $\mathfrak G_\eta$ en vertu d'une équation $f(z)=0$. Alors la multiplication de $f(z)$ par ses conjuguées relativement à $[\mathsf{H}]$ montre que $z$ est aussi entier sur $[\mathsf{H}]$, et appartient donc à $\mathfrak G_\eta$. La clôture intégrale du domaine fonctionnel considéré au \S{} 2 se démontre de la même manière.

Que la propriété V n'appartienne pas à tout domaine $\mathfrak G$ est montré par le domaine de module considéré au \S{} 3, qui ne possède précisément pas la propriété de base 3). Plus exactement, le \S{} 3 a montré que le domaine de module n'est intégralement clos relativement à aucun élément transcendant du domaine, tandis que, par III et IV, c'est bien entendu le cas relativement à tout élément algébrique du domaine.

Cela nous conduit d'abord à choisir, dans la totalité des domaines $\mathfrak G$, ceux qui possèdent aussi la propriété abstraite V; ceux-ci seront appelés \og{} domaines $\mathfrak H$ d'entiers transcendants algébriquement-entiers \fg{}.

Soit $\mathfrak H$ un tel domaine, et soit $\mathsf{H}$ une base algébrique de tous les nombres de $\mathfrak H$; d'après II, $\mathsf{H}$ est en même temps une base algébrique de tous les nombres. D'après III, I et V, $\mathfrak H$ contient, outre $\mathsf{H}$, toutes les fonctions algébriques entières des $\eta$, c'est-à-dire le domaine de Zermelo $\mathfrak G_\eta$. Ainsi, en analogie avec le \S{} 1 : tout domaine $\mathfrak H$ d'entiers transcendants algébriquement-entiers peut être construit au moyen d'une base algébrique de tous les nombres de telle manière que $\mathfrak H$ contienne le domaine de Zermelo $\mathfrak G_\eta$ comme diviseur.

Soit maintenant $\mathsf{H}$ une base algébrique quelconque de tous les nombres, et que $\mathfrak M(\mathfrak H)$ désigne la totalité de tous les domaines $\mathfrak H$ contenant $\mathsf{H}$. Tout domaine $\mathfrak H$ de $\mathfrak M(\mathfrak H)$ contient, comme on vient de le prouver, le domaine de Zermelo $\mathfrak G_\eta$ comme diviseur; et puisque $\mathfrak G_\eta$ lui-même est un domaine $\mathfrak H$, $\mathfrak G_\eta$ est le plus grand commun diviseur, ou l'intersection, de tous les domaines $\mathfrak H$ de $\mathfrak M(\mathfrak H)$. Ainsi $\mathfrak G_\eta$ est défini abstraitement. Il en résulte aussi directement que $\mathfrak M(\mathfrak H)$ est fermé pour l'intersection : c'est-à-dire que l'intersection de tous les domaines $\mathfrak H$ d'un sous-système quelconque de $\mathfrak M(\mathfrak H)$ appartient encore à $\mathfrak M(\mathfrak H)$. En effet, comme intersection, I et V sont satisfaites; II et III résultent de ce que $\mathfrak G_\eta$ est diviseur; IV résulte de ce que l'intersection est diviseur de $\mathfrak H$.

Comme le montre l'exemple du domaine fonctionnel au \S{} 2, les domaines $\mathfrak H$ contiennent en général d'autres fonctions encore que celles de $\mathfrak G_\eta$. Soit $\psi(\eta)$ une telle fonction rationnelle ou algébrique des $\eta$. Alors, par I, $\mathfrak H$ contient aussi toutes les combinaisons rationnelles entières de $\psi$ et, chaque fois, d'un nombre fini de fonctions de $\mathfrak G_\eta$. Le domaine entier ainsi obtenu --- constitué de tous les polynômes en $\psi$ à coefficients dans $\mathfrak G_\eta$ --- sera noté $[\mathfrak G_\eta,\psi]$, et le procédé qui y conduit sera appelé \og{} adjonction polynomiale de $\psi$ à $\mathfrak G_\eta$ \fg{}. Par V, $\mathfrak H$ contient en outre toutes les fonctions entières sur $[\mathfrak G_\eta,\psi]$; le domaine ainsi obtenu sera noté $\{\mathfrak G_\eta,\psi\}$, et le procédé qui y conduit \og{} adjonction intégrale de $\psi$ à $\mathfrak G_\eta$ \fg{}. Le domaine $\{\mathfrak G_\eta,\psi\}$ se compose de toutes les fonctions algébriques entières de $\psi$ à coefficients dans $\mathfrak G_\eta$ et est donc, par les arguments connus, un domaine entier.\footnote{Comparer par exemple Weber, \S{} 97, 6 et 7.}

Nous montrons que $\{\mathfrak G_\eta,\psi\}$ est aussi intégralement clos, c'est-à-dire satisfait à la condition V. En effet, que $z$ soit entier sur $\{\mathfrak G_\eta,\psi\}$ en vertu d'une équation
\[
 f(z;a\cdots c)=z^\varkappa+a z^{\varkappa-1}+\cdots+c=0,
\]
où $a\cdots c$, comme éléments de $\{\mathfrak G_\eta,\psi\}$, satisfont à des équations à coefficients dans $[\mathfrak G_\eta,\psi]$ :
\[
\begin{gathered}
 a^\lambda+A_1a^{\lambda-1}+\cdots+A_\lambda=0,\\
 \vdots\\
 c^\mu+C_1c^{\mu-1}+\cdots+C_\mu=0.
\end{gathered}
\]
Désignons leurs racines par $a,a'\cdots a^{(\lambda-1)}\cdots c,c'\cdots c^{(\mu-1)}$. Si l'on forme maintenant le produit sur toutes les combinaisons des racines,
\[
 g(z)=f(z;a\cdots c)f(z;a'\cdots c')\cdots
      f(z;a^{(\lambda-1)}\cdots c^{(\mu-1)}),
\]
alors $g(z)$ a pour coefficient dominant l'unité et pour autres coefficients des éléments de $[\mathfrak G_\eta,\psi]$. Par conséquent $z$ appartient à $\{\mathfrak G_\eta,\psi\}$.\footnote{Ainsi, en définissant l'intégralité algébrique au moyen d'une équation \emph{quelconque}, la notion d'irréductibilité relativement à un corps de base devient inutile pour la preuve de I et V. Comparer aussi le \S{} 2, où, seulement pour la preuve de IV, on a besoin de la proposition auxiliaire qui fait passer de la définition de l'intégralité algébrique par une équation arbitraire à l'équation irréductible. Comme cette proposition auxiliaire ne vaut plus en général dans le cas présent, IV doit être imposée à la fonction $\psi$ comme condition spéciale.}

Ainsi le domaine $\{\mathfrak G_\eta,\psi\}$ possède les propriétés I et V; et, puisque $\mathfrak G_\eta$ est diviseur du domaine, il possède aussi II et III. Que IV soit satisfaite dépend du choix de $\psi$; par exemple IV ne l'est pas pour $\psi=\frac{1}{n\cdot\eta}$, car alors $\frac1n$ appartient au domaine. Elle est toutefois satisfaite, d'après le \S{} 2, pour $\psi=\frac1\eta$, ou encore pour $\psi=\frac{\eta}{n}$. Dans ce dernier cas, en effet, si une fonction de $\{\mathfrak G_\eta,\psi\}$ représente un nombre algébrique, sa valeur reste inchangée pour $\eta=0$; elle passe ainsi en une fonction de $\mathfrak G_\eta$, et par conséquent en un entier algébrique.

Ainsi $\{\mathfrak G_\eta,\psi\}$ devient un domaine $\mathfrak H$ dès que l'on impose à $\psi$ la condition qu'aucun nombre algébriquement fractionnaire n'entre dans le domaine par l'adjonction algébriquement entière.

De manière tout à fait analogue, on définit l'adjonction polynomiale, respectivement intégrale, d'un système arbitraire $\mathfrak S$ de fonctions rationnelles ou algébriques des $\eta$ à $\mathfrak G_\eta$. L'adjonction polynomiale conduit au domaine entier $[\mathfrak G_\eta,\mathfrak S]$, constitué de tous les polynômes entiers en un nombre fini d'éléments de $\mathfrak G_\eta$ et de $\mathfrak S$ à la fois. L'adjonction intégrale donne le domaine $\{\mathfrak G_\eta,\mathfrak S\}$ constitué de tous les éléments entiers sur $[\mathfrak G_\eta,\mathfrak S]$. Exactement comme ci-dessus, on montre que le domaine entier $\{\mathfrak G_\eta,\mathfrak S\}$ est intégralement clos et satisfait aussi II et III. Ainsi, comme ci-dessus :

$\{\mathfrak G_\eta,\mathfrak S\}$ devient un domaine $\mathfrak H$ dès que l'on impose à $\mathfrak S$ la condition qu'aucun nombre algébriquement fractionnaire n'entre dans le domaine par l'adjonction intégrale --- c'est-à-dire dès que $\mathfrak S$ devient un système admissible.\footnote{Plus généralement, on montre de même : si $\mathfrak T$ est un domaine quelconque et si $\{\mathfrak T\}$ désigne la totalité des éléments entiers sur $\mathfrak T$, alors $\{\mathfrak T\}$ est intégralement clos.}

Il est maintenant important que la réciproque soit également vraie, de sorte que les domaines $\{\mathfrak G_\eta,\mathfrak S\}$ épuisent effectivement tous les domaines $\mathfrak H$ de $\mathfrak M(\mathfrak H)$ lorsque $\mathfrak S$ parcourt tous les systèmes admissibles. En effet, soit $\mathfrak H$ un domaine quelconque de $\mathfrak M(\mathfrak H)$, et désignons par $\mathfrak L=\mathfrak H-\mathfrak G_\eta$ le système constitué de toutes les fonctions de $\mathfrak H$ n'appartenant pas à $\mathfrak G_\eta$. Par V, $\mathfrak H$ contient le domaine $\{\mathfrak G_\eta,\mathfrak L\}$ comme diviseur; mais $\{\mathfrak G_\eta,\mathfrak L\}$ est aussi divisible par $\mathfrak H$, car il contient $\mathfrak G_\eta$ et $\mathfrak L$, et, puisque ceux-ci n'ont aucun élément commun, il contient aussi $\mathfrak H$. Donc $\mathfrak H=\{\mathfrak G_\eta,\mathfrak L\}$; et puisque IV est satisfaite pour $\mathfrak H$, $\mathfrak L$ est naturellement un système admissible.\footnote{L'adjonction d'un sous-système de $\mathfrak L$ peut déjà suffire; par exemple le domaine fonctionnel du \S{} 2 provient de l'adjonction algébriquement entière de toutes les fonctionnelles rationnelles entières --- à l'exception des polynômes --- à $\mathfrak G_\eta$.}

Lorsque $\mathsf{H}$ parcourt toutes les bases algébriques possibles, $\mathfrak M(\mathfrak H)$ parcourt, comme on l'a prouvé au début, la totalité de tous les domaines $\mathfrak H$. Par conséquent la question des domaines les plus généraux $\mathfrak H$ d'entiers transcendants algébriquement-entiers est complètement résolue par les résultats de ce paragraphe :
\begin{enumerate}
\item[a)] L'intersection de tous les domaines $\mathfrak H$ de $\mathfrak M(\mathfrak H)$ est donnée par le domaine de Zermelo $\mathfrak G_\eta$, qui est lui-même un domaine $\mathfrak H$; $\mathfrak M(\mathfrak H)$ est fermé relativement à la formation des intersections.
\item[b)] Tout domaine $\mathfrak H$ de $\mathfrak M(\mathfrak H)$ provient de l'adjonction intégrale d'un système admissible $\mathfrak S$ à $\mathfrak G_\eta$. Lorsque $\mathsf{H}$ parcourt toutes les bases algébriques possibles, $\mathfrak M(\mathfrak H)$ parcourt la totalité de tous les domaines $\mathfrak H$.\footnote{Les systèmes admissibles $\mathfrak S$ peuvent être trouvés comme suit. On place toutes les fonctions algébriquement fractionnaires des $\eta$ sous un bon ordre $\Omega$, et l'on prend dans $\mathfrak S$ toutes les fonctions sauf celles dont l'adjonction intégrale à $\mathfrak G_\eta$ et aux fonctions admises auparavant engendre un nombre algébriquement fractionnaire. Ce procédé peut être interrompu en un lieu arbitraire. Lorsque $\Omega$ parcourt tous les bons ordres, et que le procédé est interrompu en tout lieu arbitraire pour chaque $\Omega$ fixé, tous les systèmes admissibles $\mathfrak S$ sont ainsi épuisés.}
\end{enumerate}
\clearpage
\section*{\S{} 5. Les domaines les plus généraux formés d'entiers transcendants, sur la base d'une base algébrique.}

Soit $\mathsf{H}$ une base algébrique quelconque de tous les nombres, et soit $\mathfrak M(\mathfrak G)$ l'ensemble de tous les domaines $\mathfrak G$ qui contiennent $\mathsf{H}$ et donc $[\mathsf{H}]$. Lorsque $\mathsf{H}$ parcourt toutes les bases possibles, $\mathfrak M(\mathfrak G)$ parcourt l'ensemble de tous les domaines $\mathfrak G$; par conséquent, comme au \S{} 4, nous pouvons nous limiter à considérer les domaines $\mathfrak G$ appartenant à $\mathfrak M(\mathfrak G)$.

Les domaines $\mathfrak G$ de $\mathfrak M(\mathfrak G)$ contiennent tous $[\mathsf{H}]$ comme diviseur; mais, puisque $[\mathsf{H}]$ n'est pas lui-même un domaine $\mathfrak G$, on ne peut pas conclure comme au \S{} 4 que $[\mathsf{H}]$ soit le plus grand diviseur commun. Que cela soit néanmoins le cas est montré par une famille dénombrable de domaines $\mathfrak G$ de $\mathfrak M(\mathfrak G)$, obtenue par une légère généralisation du domaine de modules du \S{} 3, et n'ayant en effet aucun élément commun en dehors de $[\mathsf{H}]$.

\emph{Pour chaque $\nu$ fixé, soient les domaines $\mathfrak G_\nu$ $(\nu=1,2,\ldots\text{ à l'infini})$ composés de toutes les quantités}
\begin{equation}
z=f(\eta)+g_1(\eta)\eta_1^\nu+g_2(\eta)\eta_2^\nu+\cdots+g_t(\eta)\eta_t^\nu,
\tag{1}
\end{equation}
\emph{où $f(\eta)$ parcourt tous les polynômes entiers en les $\eta$; où $\eta_1\cdots\eta_t$ parcourent les quantités de base, et où, pour chaque combinaison fixée $\eta_1^\nu\cdots\eta_t^\nu$, les fonctions $g_1(\eta)\cdots g_t(\eta)$ parcourent indépendamment toutes les fonctions algébriques entières des $\eta$.}

Comme au \S{} 3, on voit que les domaines $\mathfrak G_\nu$ possèdent les propriétés abstraites I--IV; $\mathfrak G_\nu$ ne contient que des fonctions algébriquement entières et, pour $\nu=1$, il est identique au domaine du \S{} 3. Nous montrons maintenant que les domaines $\mathfrak G_\nu$ n'ont pas d'autre élément commun que les polynômes $f(\eta)$ de $[\mathsf{H}]$. Soit en effet $z$ une fonction algébriquement entière quelconque, mais non rationnelle, des $\eta$, satisfaisant à l'équation irréductible relativement à $[\mathsf{H}]$
\begin{equation}
z^\varkappa+a_1(\eta)z^{\varkappa-1}+a_2(\eta)z^{\varkappa-2}+\cdots+a_\varkappa(\eta)=0,
\tag{2}
\end{equation}
où $\varkappa>1$; et soit $\lambda$ le plus grand des degrés des polynômes $a_1(\eta)\cdots a_\varkappa(\eta)$. Si $z$ appartient à $\mathfrak G_\nu$, alors la quantité $\zeta=z-f(\eta)$ satisfait à l'équation, également irréductible relativement à $[\mathsf{H}]$,
\begin{equation}
\zeta^\varkappa+b_1(\eta)\zeta^{\varkappa-1}+b_2(\eta)\zeta^{\varkappa-2}+\cdots+b_\varkappa(\eta)=0,
\tag{3}
\end{equation}
où, d'après (1), $b_1(\eta),b_2(\eta)\cdots b_\varkappa(\eta)$, comme fonctions symétriques des $\zeta$, contiennent comme termes de plus basse dimension des termes d'au moins les dimensions $\nu,2\nu,\cdots,\varkappa\nu$. La comparaison de (2) et (3) donne
\[
 b_1(\eta)=a_1(\eta)+\varkappa\cdot f(\eta).
\]
Si donc la quantité $\xi=z+\frac{a_1(\eta)}{\varkappa}$ satisfait à
\begin{equation}
\xi^\varkappa+c_2(\eta)\xi^{\varkappa-2}+\cdots+c_\varkappa(\eta)=0,
\tag{4}
\end{equation}
alors
\begin{equation}
 \xi-\zeta=\dfrac{b_1(\eta)}{\varkappa};\qquad;
 \qquad
 c_i(\eta)\equiv b_i(\eta)\pmod{\dfrac{b_1(\eta)}{\varkappa}},
\tag{5}
\end{equation}
où $b_1(\eta)$ commence par des termes d'au moins la dimension $\nu$. Or les $c_i(\eta)$, comme le montre la comparaison de (2) avec (4), sont de degré $\chi$ dans les coefficients $a_i(\eta)$ de (2), donc de degré au plus $\varkappa\cdot\lambda$ dans les $\eta$; d'autre part tous les $b_i(\eta)$ commencent par des termes d'au moins la dimension $\nu$. Ainsi, si l'on choisit $\nu>\varkappa\cdot\lambda$, (5) donne
\[
 c_i(\eta)=0;\qquad
 \xi^\varkappa=0\quad\text{ou}\quad
 \left(z+\frac{a_1(\eta)}{\varkappa}\right)^\varkappa=0.
\]
Ainsi $z$, contrairement à l'hypothèse, devient rationnelle en les $\eta$. Donc :

\emph{Si $z$ est une fonction algébriquement entière, non rationnelle, des $\eta$, alors $z$ ne peut appartenir à aucun domaine $\mathfrak G_\nu$ pour lequel $\nu>\varkappa\lambda$, ou encore :}

\emph{L'intersection de tous les domaines $\mathfrak G$ de $\mathfrak M(\mathfrak G)$ est donnée par le domaine entier $[\mathsf{H}]$, qui lui-même n'est pas un domaine $\mathfrak G$; par conséquent $\mathfrak M(\mathfrak G)$ n'est pas fermé pour la formation des intersections.}

La construction des domaines $\mathfrak G$ les plus généraux au moyen d'une base algébrique est donc elle aussi moins facile à embrasser que celle des domaines $\mathfrak H$. Le domaine $[\mathsf{H}]$ possède les propriétés I, III et IV; si l'on adjoint polynomialement un système quelconque $\mathfrak S$, alors, pour le domaine entier résultant $[\mathsf{H},\mathfrak S]$, les propriétés I et III restent valables, tandis que IV constitue une condition spéciale imposée à $\mathfrak S$. Pour qu'un domaine $\mathfrak G$ naisse, il faut encore imposer II à $\mathfrak S$ comme condition; ou, exprimé autrement : pour tout corps $\mathfrak K$ fini sur $(\mathsf{H})$\footnote{$(\mathsf{H})$ désigne le corps formé de toutes les fonctions rationnelles des $\eta$ à coefficients algébriques.}, il doit exister un corps $\mathfrak L$ au-dessus de $\mathfrak K$, fini sur $\mathfrak K$, tel que $[\mathsf{H},\mathfrak S]$ contienne un élément primitif de $\mathfrak L$.\footnote{Cf. \S{} 7.} Réciproquement, tout domaine $\mathfrak G$ de $\mathfrak M(\mathfrak G)$ peut aussi être représenté sous la forme $[\mathsf{H},\mathfrak S]$; on obtient donc tous les domaines $\mathfrak G$ de $\mathfrak M(\mathfrak G)$ lorsque $\mathfrak S$ parcourt tous les systèmes ainsi restreints. La structure des systèmes $\mathfrak S$ apparaîtra plus simplement dans le paragraphe suivant au moyen de la \og{} base rationnelle \fg{}; en même temps on obtient une analogie complète avec les théorèmes du \S{} 4.

\section*{\S{} 6. Les domaines les plus généraux formés d'entiers transcendants, sur la base d'une base rationnelle.}

Par analogie avec la base algébrique, la base rationnelle doit être définie comme suit : \og{} Un système $\Theta$ de nombres $\vartheta$ est appelé base rationnelle de tous les nombres si $\Theta$ permet d'exprimer rationnellement tout nombre --- avec coefficients dans $\mathsf{K}$\footnote{Par fonctions rationnelles il faut toujours entendre de telles fonctions à coefficients algébriques; ces coefficients algébriques sont inclus dans la définition à cause de III et IV. Il serait plus adéquat d'exiger des coefficients rationnels aussi bien dans III et IV que dans la définition de la base rationnelle. Les raisonnements des \S{}\S{} 6 et 7 restent littéralement inchangés si l'on remplace seulement le corps $\mathsf{K}$ de tous les nombres algébriques par le corps de tous les nombres rationnels.} --- tandis que, sous au moins un bon ordre, aucun nombre de base ne peut être exprimé rationnellement --- avec coefficients dans $\mathsf{K}$ --- par un nombre fini de nombres de base précédents. \fg{}\footnote{À la différence des bases linéaire et algébrique, l'ordre des éléments joue ici un rôle. Par exemple, dans l'ordre $\eta,\sqrt{\eta}$ les deux éléments doivent être admis, tandis que dans l'ordre inverse seul $\sqrt{\eta}$ doit être admis.}

Soumettons à un bon ordre $\Omega$ un domaine $\mathfrak G$ formé d'entiers transcendants. Il peut alors arriver qu'une quantité $z$ de $\mathfrak G$ soit exprimable rationnellement par un nombre fini de quantités précédentes de $\mathfrak G$; c'est-à-dire qu'elle satisfasse à une équation $z=\varphi(\zeta)$, où $\varphi(\zeta)$ est une fonction rationnelle, à coefficients dans $\mathsf{K}$, de $\zeta_1,\ldots,\zeta_t$. En particulier, tout nombre algébrique $\alpha$ du domaine satisfait à l'équation $z=\alpha$. Les quantités restantes $\vartheta$ de $\mathfrak G$ forment un système $\Theta$ dont il faut montrer qu'il est une base rationnelle de toutes les quantités de $\mathfrak G$. En effet, par hypothèse, chaque quantité $\vartheta$ est rationnellement indépendante des quantités précédentes de $\Theta$. L'argument d'induction\footnote{Cf. Zermelo, loc. cit., \S{} 1.} montre en outre que toute relation $z=\varphi(\zeta)$ entraîne une relation $z=\psi(\vartheta)$; sinon il devrait y avoir un premier élément $z_0=\varphi_0(\zeta_1\cdots\zeta_t)$ qui ne puisse pas être exprimé rationnellement par les $\vartheta$, tandis que les $\xi_1,\ldots,\xi_t$ précédents seraient des fonctions rationnelles des $\vartheta$; cela donne une contradiction. Par II, $\Theta$ est en même temps une base rationnelle de tous les nombres; par III et I, outre $\Theta$, $\mathfrak G$ contient comme diviseur le domaine entier $[\Theta]$ constitué de tous les polynômes entiers en les $\vartheta$, tandis que, par IV, $[\Theta]$ ne contient aucun nombre algébriquement fractionnaire. Ainsi la base rationnelle $\Theta$ de tous les nombres satisfait à la condition accessoire que le domaine entier $[\Theta]$ qui en provient ne contienne aucun nombre algébriquement fractionnaire;\footnote{Que ce soit réellement une condition est montré par exemple par le fait que les deux quantités $\eta,\frac1{2\sqrt{\eta}}$ ne sont pas admissibles comme éléments de base, tandis que $\eta,\frac1{\sqrt{\eta}}$ le sont.} $\Theta$ est une \emph{base rationnelle avec condition accessoire} pour tous les nombres. Donc, par analogie avec le \S{} 2 : \emph{Tout domaine $\mathfrak G$ peut être construit au moyen d'une base rationnelle avec condition accessoire, de manière que $\mathfrak G$ contienne le domaine entier $[\Theta]$ comme diviseur.}

Soit $\Theta$ une base rationnelle quelconque avec condition accessoire pour tous les nombres --- dont l'existence est donc garantie par l'existence des domaines $\mathfrak G$. Soit $\mathfrak N(\mathfrak G)$ l'ensemble de tous les domaines $\mathfrak G$ qui contiennent $\Theta$ et donc $[\Theta]$. Lorsque $\Theta$ parcourt toutes les bases possibles avec condition accessoire, $\mathfrak N(\mathfrak G)$ parcourt, d'après ce qui précède, l'ensemble de tous les domaines $\mathfrak G$; par conséquent nous pouvons à nouveau nous limiter aux domaines $\mathfrak G$ de $\mathfrak N(\mathfrak G)$.

Or $[\Theta]$ lui-même est un domaine $\mathfrak G$; car tout nombre peut être exprimé rationnellement par $\Theta$, donc comme quotient de deux polynômes de $[\Theta]$; et, par sa construction, $[\Theta]$ satisfait aussi aux autres conditions I, III, IV. \emph{Ainsi $[\Theta]$ est le plus grand diviseur commun, ou l'intersection, de tous les domaines $\mathfrak G$ de $\mathfrak N(\mathfrak G)$; en outre $\mathfrak N(\mathfrak G)$ est fermé par rapport à l'intersection, c'est-à-dire que l'intersection de tous les domaines $\mathfrak G$ d'un sous-système quelconque de $\mathfrak N(\mathfrak G)$ appartient encore à $\mathfrak N(\mathfrak G)$.} En effet, I vaut pour l'intersection; II et III résultent du fait que $[\Theta]$ est un diviseur; et IV du fait que le domaine d'intersection est contenu comme diviseur dans un domaine $\mathfrak G$.

L'adjonction rationnellement entière d'un système quelconque $\mathfrak S$ de fonctions rationnelles des $\vartheta$ à $[\Theta]$ --- toute fonction algébrique des $\vartheta$ pouvant, par la propriété de la base rationnelle, être exprimée rationnellement au moyen de certains autres $\vartheta$ --- donne un domaine entier $[\Theta,\mathfrak S]$ possédant les propriétés I, II, III, et devient donc un domaine $\mathfrak G$ dès que $\mathfrak S$ est un \og{} système admissible \fg{}, c'est-à-dire dès qu'il satisfait à la condition qu'aucun nombre algébriquement fractionnaire n'entre dans le domaine par l'adjonction rationnellement entière. Réciproquement, tout domaine $\mathfrak G$ de $\mathfrak N(\mathfrak G)$ permet la représentation $[\Theta,\mathfrak S]$, pourvu seulement que $\mathfrak S$ désigne le système résiduel $\mathfrak G-[\Theta]$. Ainsi, pour les domaines $\mathfrak G$ les plus généraux, en complète analogie avec le \S{} 4, on a :
\begin{enumerate}
\item[a)] \emph{L'intersection de tous les domaines $\mathfrak G$ de $\mathfrak N(\mathfrak G)$ est donnée par le domaine entier $[\Theta]$, qui est lui-même un domaine $\mathfrak G$; $\mathfrak N(\mathfrak G)$ est fermé par rapport à la formation des intersections.}
\item[b)] \emph{Tout domaine $\mathfrak G$ de $\mathfrak N(\mathfrak G)$ naît par adjonction rationnellement entière d'un système admissible $\mathfrak S$ à $[\Theta]$. Lorsque $\Theta$ parcourt toutes les bases rationnelles possibles avec condition accessoire, $\mathfrak N(\mathfrak G)$ parcourt l'ensemble de tous les domaines $\mathfrak G$.}\footnote{Pour les systèmes admissibles $\mathfrak S$ les plus généraux, ce qui a été dit au \S{} 4 vaut en remplaçant \og{} algébrique \fg{} par \og{} rationnel \fg{}.}
\end{enumerate}

\emph{Remarque.} Si l'on extrait de $\Theta$ une base algébrique $\mathsf{H}$ de toutes les quantités de $\Theta$, alors $\mathsf{H}$ est en même temps une base algébrique de tous les nombres; réciproquement toute base algébrique peut être prolongée en une base rationnelle en plaçant $\mathsf{H}$ d'abord dans le bon ordre. En conséquence, la construction du \S{} 5 des domaines $\mathfrak G$ au moyen d'une base algébrique peut s'interpréter comme suit. Décomposons le système $\mathfrak S$ à adjoindre à $[\mathsf{H}]$ en deux sous-systèmes $\mathfrak S_1$ et $\mathfrak S_2$; l'adjonction de $\mathfrak S_1$ revient à prolonger $\mathsf{H}$ en une base rationnelle avec condition accessoire, à laquelle on adjoint ensuite encore polynomialement un système admissible $\mathfrak S_2$.

\section*{\S{} 7. Construction de la base rationnelle la plus générale avec condition accessoire.}

Les résultats du paragraphe précédent ont encore besoin d'un complément, car on y a prouvé l'\emph{existence} d'une base rationnelle avec condition accessoire pour tous les nombres, mais non un procédé pour la construire à partir d'un bon ordre prescrit. À cause de la condition accessoire, cette construction prend une forme plus compliquée lorsque, au lieu d'un domaine $\mathfrak G$, on a devant soi le corps de tous les nombres complexes.\footnote{Au \S{} 6 b), il suffirait de ne faire parcourir à $\Theta$ que les bases rationnelles spéciales qui se composent d'une base algébrique et de fonctions algébriquement entières des quantités de base; la condition accessoire est alors automatiquement satisfaite. Pour le voir, il suffit, d'après \S{} 6, de prolonger en une base rationnelle une base algébrique extraite de $\mathfrak G$, puis, dans la base ainsi obtenue, de multiplier chaque fonction algébriquement fractionnaire des $\eta$ par un polynôme convenablement choisi en les $\eta$, ce qui conserve la propriété de base. Cette construction se transpose sans changement au corps de tous les nombres complexes. La question, issue du \S{} 6, de la base rationnelle la plus générale avec condition accessoire possède toutefois aussi un intérêt propre.}

Soumettons le corps de tous les nombres complexes à un bon ordre $\Omega$. Alors des quantités de trois espèces peuvent se présenter :
\begin{enumerate}
\item[1.] Des quantités $y$ dont l'adjonction rationnellement entière aux quantités de base précédentes --- définies récursivement en 3.\footnote{Comme le montre l'argument d'induction, une telle définition récursive est possible.} --- force l'entrée d'un nombre algébriquement fractionnaire dans le domaine entier ainsi obtenu; si aucune quantité de base ne les précède, ce domaine consiste en tous les polynômes entiers en $y$. En particulier, tous les nombres algébriquement fractionnaires sont des quantités $y$, puisque le domaine entier obtenu contient les quantités $y=\beta$.
\item[2.] Des quantités $z$ qui ne possèdent pas la propriété 1, mais qui peuvent être exprimées rationnellement par un nombre fini de nombres précédents différents des $y$; elles satisfont donc à une relation $z=\varphi(\zeta_1\cdots\zeta_t)$, où $\varphi$ est une fonction rationnelle à coefficients dans $\mathsf{K}$ et où aucun des $\xi$ n'est un $y$. En particulier, tous les nombres algébriques entiers $\alpha$ appartiennent ici, puisque la propriété 1 ne vaut pas pour eux et qu'ils satisfont à la relation $z=\alpha$.
\item[3.] Des quantités $\vartheta$ pour lesquelles ni 1 ni 2 ne vaut et qui doivent être appelées quantités de base. En particulier, le premier nombre transcendant $\vartheta_0$ dans le bon ordre est une telle quantité de base; car les polynômes entiers en $\vartheta_0$ ne peuvent représenter aucun nombre algébriquement fractionnaire, et $\vartheta_0$ ne peut pas non plus dépendre rationnellement des nombres algébriques qui le précèdent.
\end{enumerate}

\emph{Il faut maintenant montrer que le système $\Theta$ de toutes les quantités de base $\vartheta$ est une base rationnelle avec condition accessoire de tous les nombres.}

Comme aucune quantité $y$ ne se trouve parmi les $\vartheta$, le domaine entier $[\Theta]$ ne peut contenir aucun nombre algébriquement fractionnaire; la condition accessoire est donc satisfaite. Comme, en outre, aucune quantité $z$ ne se trouve parmi les $\vartheta$, chaque $\vartheta$ est rationnellement indépendante des quantités de base précédentes. L'argument d'induction montre, exactement comme au \S{} 6, que toute relation $z=\varphi(\zeta)$ --- où aucun $y$ ne se trouve parmi les $\xi$ --- entraîne une relation $z=\psi(\vartheta)$, de sorte que toutes les quantités $z$ sont exprimables rationnellement par les $\vartheta$.

Il reste seulement à prouver l'assertion plus difficile que les $y$ sont eux aussi exprimables rationnellement par les $\vartheta$. Nous montrons d'abord que les $y$ dépendent algébriquement des $\vartheta$. Par hypothèse, il existe pour chaque $y$ au moins un nombre algébriquement fractionnaire $\beta$ admettant la représentation
\begin{equation}
\beta=f_0(\vartheta)+f_1(\vartheta)y+\cdots+f_\varkappa(\vartheta)y^\varkappa,
\tag{1}
\end{equation}
où $f_0(\vartheta)\cdots f_\varkappa(\vartheta)$ sont des polynômes de $[\Theta]$ et ne peuvent donc représenter aucun nombre algébriquement fractionnaire. Si $y$ était algébriquement indépendante des $\vartheta$, (1) vaudrait identiquement en $y$ et l'on aurait $\beta=f_0(\vartheta)$, en contradiction avec les propriétés de $[\Theta]$.

Pour conclure de la dépendance algébrique à la dépendance rationnelle, nous extrayons de $\Theta$ une base algébrique $\mathsf{H}$ de toutes les quantités de $\Theta$. Alors chaque $y$, comme fonction algébrique des $\vartheta$, dépend aussi algébriquement des $\eta$; par $y=y(\eta)$ est donc déterminé un corps algébrique $(\mathsf{H},y(\eta))$ au-dessus de $(\mathsf{H})$. Comme tous les $\eta$ sont en même temps des quantités $\vartheta$, la preuve de la dépendance rationnelle revient à montrer que pour chacun de ces corps il existe des éléments primitifs qui sont des fonctions rationnelles des $\vartheta$; plus précisément, on montrera que $y(\eta)$ perd la propriété 1 après multiplication par un polynôme en les $\eta$, et passe par conséquent en un tel élément.

On sait qu'entre $(\mathsf{H})$ et $(\mathsf{H},y)$ il n'existe qu'un nombre fini de corps intermédiaires. Soient $\Omega_0=(H),\Omega_1,\ldots,\Omega_\sigma$ ceux d'entre eux qui possèdent un élément primitif exprimable rationnellement par les $\vartheta$, de sorte que tout élément du corps devient une fonction rationnelle des $\vartheta$; cette condition est au moins toujours satisfaite pour $\mathfrak L_0=(\mathsf{H})$. Soit donnée la fonction primitive appartenant à $\mathfrak L_i$ sous la forme $\frac{g_i(\vartheta)}{h_i(\vartheta)}$, où $g_i,h_i$ sont des polynômes de $[\Theta]$. Alors --- $\alpha_i$ désignant un entier convenablement choisi --- le polynôme de $[\Theta]$
\[
 \zeta_i=g_i(\vartheta)+\alpha_i h_i(\vartheta)
\]
est une fonction primitive de $\mathfrak L_i$, ou d'un corps algébrique au-dessus de $\mathfrak L_i$,\footnote{Cf. la fin du \S{} 5.} et toute fonction $\psi(\eta)$ de $\mathfrak L_i$ admet la représentation
\begin{equation}
\psi(\eta)=
\frac{a_1(\eta)\zeta_i^{\varkappa-1}+a_2(\eta)\zeta_i^{\varkappa-2}+\cdots+a_\varkappa(\eta)}
     {a(\eta)}
=\frac{A(\eta,\zeta_i)}{a(\eta)},
\tag{2}
\end{equation}
où $a(\eta),a_1(\eta),\ldots,a_\varkappa(\eta)$ sont des polynômes entiers en les $\eta$, et donc $A(\eta,\zeta_i)$ et $a(\eta)$ sont des quantités de $[\Theta]$.

L'équation irréductible satisfaite par $y$ relativement à $\mathfrak L_i$ est, d'après (2), de la forme
\[
 f^{(i)}(\eta)y^{\lambda_i}+f_1^{(i)}(\eta,\zeta_i)y^{\lambda_i-1}
+\cdots+f_{\lambda_i}^{(i)}(\eta,\zeta_i)=0,
\]
et tous les coefficients de cette équation appartiennent à $[\Theta]$. Alors la fonction $y_i=f^{(i)}(\eta)\cdot y$ satisfait à l'équation également irréductible, relativement à $\mathfrak L_i$,
\[
 y_i^{\lambda_i}+f_1^{(i)}(\eta,\zeta_i)y_i^{\lambda_i-1}
+\cdots+\bigl(f^{(i)}(\eta)\bigr)^{\lambda_i-1}
 f_{\lambda_i}(\eta,\zeta_i)=0,
\]
en vertu de laquelle elle est entière sur $[\mathsf{H},\zeta_i]$; il en est de même du produit de $y_i$ par n'importe quel polynôme en les $\eta$. En particulier, la fonction
\begin{equation}
Y=f^{(0)}(\eta)f^{(1)}(\eta)\cdots f^{(\sigma)}(\eta)\cdot y
\tag{3}
\end{equation}
est entière sur chacun des domaines $[\mathsf{H},\zeta_i]$ en vertu de l'équation, irréductible dans chaque cas relativement à $\mathfrak L_i$,
\begin{equation}
Y^{\lambda_i}+F_1^{(i)}(\eta,\zeta_i)Y^{\lambda_i-1}
+\cdots+F_{\lambda_i}^{(i)}(\eta,\zeta_i)=0
\qquad (i=0,1,\ldots,\sigma).
\tag{4}
\end{equation}
Nous montrons maintenant que $Y$ est l'élément primitif recherché, exprimable rationnellement par les $\vartheta$. Supposons par exemple que le nombre algébrique $\gamma$ soit représenté par $Y$ et $[\Theta]$ :
\begin{equation}
\gamma=k_0(\vartheta)+k_1(\vartheta)Y+\cdots+k_\nu(\vartheta)Y^\nu,
\tag{5}
\end{equation}
où $k_0(\vartheta)\cdots k_\nu(\vartheta)$ sont des polynômes de $[\Theta]$. Déterminons maintenant l'équation irréductible satisfaite par $Y$ relativement au corps $\mathfrak K=(\mathsf{H};k_0(\vartheta)\cdots k_\nu(\vartheta))$ provenant du domaine de coefficients de (5). Elle est, comme on sait, donnée par l'équation irréductible de $Y$ relativement au corps d'intersection de $\mathfrak K$ et $(\mathsf{H},Y)$, situé entre $(\mathsf{H})$ et $(\mathsf{H},Y)$. Or $(\mathsf{H},Y)$ est identique à $(\mathsf{H},y)$ d'après (3); et comme tous les éléments de $\mathfrak K$, donc aussi tous les éléments du corps d'intersection, peuvent être exprimés rationnellement par un nombre fini de $\vartheta$, ce dernier corps est nécessairement l'un des corps $\mathfrak L_0,\mathfrak L_1,\cdots,\mathfrak L_\sigma$, disons $\mathfrak L_\tau$. L'équation irréductible cherchée est donc, d'après (4), de degré $\lambda_\tau$ et s'écrit
\begin{equation}
Y^{\lambda_\tau}+F_1^{(\tau)}(\eta,\zeta_\tau)Y^{\lambda_\tau-1}
+\cdots+F_{\lambda_\tau}^{(\tau)}(\eta,\zeta_\tau)=0
\tag{6}
\end{equation}
avec des polynômes de $[\Theta]$ comme coefficients. Au moyen de (6), (5) peut être réécrite sous la forme
\begin{equation}
\gamma=K_0(\vartheta)+K_1(\vartheta)Y+\cdots+
K_{\lambda_\tau-1}(\vartheta)Y^{\lambda_\tau-1},
\tag{7}
\end{equation}
où $K_0(\vartheta),K_1(\vartheta),\cdots$ appartiennent à $\mathfrak K$ et, d'autre part, sont des polynômes de $[\Theta]$. Mais puisque $\gamma$ est un nombre algébrique, la relation (7) n'atteint pas le degré $\lambda_\tau$ en $Y$ et est donc satisfaite \emph{identiquement} en $Y$. Ainsi
\[
 \gamma=K_0(\vartheta),
\]
c'est-à-dire que $\gamma$ appartient à $[\Theta]$ et est par conséquent un entier algébrique.

Lorsque $\gamma$ parcourt tous les nombres algébriques représentables par $Y$ et $[\Theta]$, le corps d'intersection de $(\mathsf{H},Y)$ et du corps correspondant $\mathfrak K$ ne parcourt toujours que les corps $\mathfrak L_0,\mathfrak L_1,\cdots,\mathfrak L_\sigma$, et tous les arguments précédents restent valables. Autrement dit, \emph{par adjonction rationnellement entière de $Y$ à $[\Theta]$, aucun nombre algébriquement fractionnaire ne peut être représenté}. Ainsi $Y$ ne possède plus la propriété 1; selon 2) ou 3) c'est une quantité $z$ ou $\vartheta$, et donc elle est rationnellement représentable par un nombre fini de $\vartheta$; d'après (3), il en est de même de $y$ : \emph{$\Theta$ est bien une base rationnelle de tous les nombres.} Réciproquement, puisque toute base rationnelle avec condition accessoire peut être construite suivant 1), 2), 3) en plaçant $\Theta$ d'abord dans le bon ordre, nous avons prouvé : \emph{La construction donnée par 1), 2), 3) conduit à la base rationnelle la plus générale avec condition accessoire de tous les nombres.}

\section*{\S{} 8. Répartition des domaines $\mathfrak G$ et $\mathfrak H$ en classes.\srcfn{*)}{Voir E. Steinitz, \emph{Algebraische Theorie der Körper}, Journal de Crelle 137 (1910), en particulier les \S\S{} 23 et 24. Le \S{} 23 établit déjà l'existence de la base algébrique.}}

À côté de la question des domaines $\mathfrak G$ et $\mathfrak H$ les plus généraux se pose la question supplémentaire des domaines \emph{essentiellement différents}, qui --- dans un sens à préciser ci-dessous --- ne peuvent pas être engendrés par le même principe de construction. Cela conduit à une répartition de ces domaines en classes.

\emph{Deux domaines seront dits appartenir à la même classe, ou être isomorphes, si leurs éléments peuvent être mis en correspondance biunivoque de telle sorte que la somme et le produit de deux éléments quelconques de l'un soient toujours associés à la somme et au produit des éléments correspondants de l'autre.} Il résulte de cette définition que, dans toute relation isomorphe, zéro et l'unité correspondent à eux-mêmes, et que toutes les relations algébriques entre un nombre fini d'éléments sont conservées pour les éléments correspondants, et réciproquement. En particulier, outre l'unité, tous les nombres rationnels correspondent à eux-mêmes; tandis que la totalité des nombres algébriques --- et de même la totalité des entiers algébriques --- se transforme en elle-même, bien qu'un nombre algébrique particulier puisse très bien correspondre à un conjugué.

Remarquons d'abord que tout domaine de l'ensemble $\mathfrak M(\mathfrak G)$ de tous les domaines $\mathfrak G$ contenant une base algébrique fixe $\mathsf{H}$ est isomorphe à un domaine de l'ensemble $\mathfrak M^*(\mathfrak G)$ appartenant à une autre base algébrique $\mathsf{H}^*$. En effet, le corps de tous les nombres complexes se déduit de chaque base algébrique par extension algébrique et a donc la même cardinalité que la base;\footnote{Steinitz, \S{}\S{} 23 et 24.} ainsi $\mathsf{H}$ et $\mathsf{H}^*$ ont la même cardinalité, et l'application isomorphe cherchée s'obtient en faisant correspondre les quantités de base $\eta_i$ aux $\eta_i^*$.\footnote{Cette conclusion ne vaut naturellement pas pour les ensembles $\mathfrak N(\mathfrak G)$ et $\mathfrak N^*(\mathfrak G)$ appartenant à deux bases rationnelles $\Theta$ et $\Theta^*$; l'isomorphisme de $\mathfrak N(\mathfrak G)$ et de $\mathfrak N^*(\mathfrak G)$ ne peut pas non plus être déduit de l'exprimabilité rationnelle mutuelle de $\Theta$ et $\Theta^*$. Ce que l'on obtient en revanche, c'est une relation isomorphe entre les corps $(\Theta)$ et $(\Theta^*)$.} D'autre part, $\mathfrak M(\mathfrak G)$ contient des domaines $\mathfrak G$ non isomorphes, ou essentiellement différents, comme il résulte des exemples des \S{}\S{} 2, 3 et 5. Car, puisque toutes les relations algébriques sont conservées dans l'application isomorphe, une base algébrique doit correspondre à nouveau à une base algébrique; or les \S{}\S{} 2 et 3 ont montré que les domaines de fonctionnelles et de modules ne possèdent aucune base algébrique pour laquelle toutes les propriétés de base du domaine de Zermelo seraient satisfaites. Comme le domaine de fonctionnelles du \S{} 2 est intégralement clos, le sous-ensemble $\mathfrak M(\mathfrak H)$ contient lui aussi des domaines $\mathfrak H$ essentiellement différents.

Nous montrons maintenant, de façon analogue à la construction de bases, comment on peut extraire de $\mathfrak M(\mathfrak G)$ un sous-ensemble $\mathfrak L(\mathfrak G)$ tel que tous les domaines $\mathfrak G$ de $\mathfrak L(\mathfrak G)$ soient essentiellement différents, tandis que tout domaine de $\mathfrak M(\mathfrak G)$ --- et donc tout domaine $\mathfrak G$ en général --- soit isomorphe à un domaine de $\mathfrak L(\mathfrak G)$. Soumettons $\mathfrak M(\mathfrak G)$ à un bon ordre $\Omega$; alors un domaine $\mathfrak G$ de $\mathfrak M(\mathfrak G)$ peut être isomorphe ou non à un domaine précédent. Les domaines $\mathfrak G$ isomorphes à aucun précédent forment un ensemble $\mathfrak L(\mathfrak G)$, qui, par construction, ne contient que des domaines $\mathfrak G$ essentiellement différents. L'argument d'induction montre aussi que tout domaine de $\mathfrak M(\mathfrak G)$ est isomorphe à un domaine de $\mathfrak L(\mathfrak G)$. Car si $\mathfrak G_1$ était le premier domaine pour lequel cela échouerait, alors $\mathfrak G_1$ appartiendrait lui-même à $\mathfrak L(\mathfrak G)$, ou bien serait isomorphe à un domaine précédent $\mathfrak G_0$, lequel, par hypothèse, est isomorphe à un domaine de $\mathfrak L(\mathfrak G)$; il en irait donc de même de $\mathfrak G_1$. Puisque tout domaine $\mathfrak G$ appartient à un certain ensemble $\mathfrak M^*(\mathfrak G)$ et est donc isomorphe à un domaine de $\mathfrak M(\mathfrak G)$, \emph{la totalité des classes de domaines $\mathfrak G$ est représentée par $\mathfrak L(\mathfrak G)$.}

Si, dans un domaine $\mathfrak G$ de $\mathfrak L(\mathfrak G)$, on remplace les quantités de base $\eta$ par celles $\eta^*$ d'une autre base, alors, comme on l'a montré plus haut, il en résulte un domaine isomorphe à $\mathfrak G$. Réciproquement, soit $\mathfrak G^*$ un domaine quelconque et supposons-le isomorphe à un domaine $\mathfrak G$ de $\mathfrak L(\mathfrak G)$. Si les quantités de base $\eta$ correspondent aux quantités $\eta^*$ de $\mathfrak G^*$, celles-ci forment à nouveau une base algébrique, et $\mathfrak G^*$ contient les mêmes fonctions algébriques de $\eta^*$ que $\mathfrak G$ contient de $\eta$ --- ou leurs conjuguées. \emph{Ainsi tout domaine isomorphe à $\mathfrak G$ s'obtient en faisant parcourir à la base algébrique $H$ toutes les bases algébriques possibles dans $\mathfrak G$ et dans ses conjugués relativement à $(\mathsf{H})$}\footnote{Les éléments de ces domaines conjugués à $\mathfrak G$ peuvent certes être exprimés rationnellement par les éléments de $\mathfrak G$ d'après II, mais pas nécessairement intégralement et rationnellement; les conjugués peuvent donc différer de $\mathfrak G$.} \emph{--- si ceux-ci diffèrent de $\mathfrak G$.} À partir de chaque domaine fixé $\mathfrak G_i$ dans $\mathfrak L(\mathfrak G)$, on obtient ainsi la classe $\mathfrak K_i$, qui contient tous et seulement les domaines $\mathfrak G$ isomorphes à $\mathfrak G_i$, bien que chaque domaine puisse y apparaître plusieurs fois, voire une infinité de fois. De $\mathfrak K_i$ on peut de nouveau extraire, par le procédé précédent, un sous-ensemble $\mathfrak K_i^*$ qui contient chaque domaine isomorphe à $\mathfrak G_i$ une et une seule fois. Lorsque $\mathfrak K_i^*$ parcourt toutes les classes, on obtient la totalité de tous les domaines $\mathfrak G$, chacun une et une seule fois. Pour caractériser la classe $\mathfrak K_i$, on peut naturellement employer, au lieu de $\mathfrak G_i$, tout autre domaine de la classe dont tous les domaines de $\mathfrak K_i$ peuvent être déduits comme ci-dessus. En résumé :

\emph{De $\mathfrak M(\mathfrak G)$ on peut extraire un sous-ensemble $\mathfrak L(\mathfrak G)$ tel que les domaines $\mathfrak G$ de $\mathfrak L(\mathfrak G)$ soient en correspondance biunivoque avec la totalité des classes (sans répétition). À partir de $\mathfrak G_i$, la classe correspondante $\mathfrak K_i$ naît en faisant parcourir à $H$ toutes les bases algébriques possibles dans $\mathfrak G_i$ et ses conjugués relativement à $(\mathsf{H})$. Si $\mathfrak K_i^*$ désigne tous les domaines mutuellement distincts de $\mathfrak K_i$, et si $\mathfrak K_i^*$ parcourt toutes les classes, alors la totalité de tous les domaines $\mathfrak G$ apparaît, chaque domaine une et une seule fois.}\footnote{Ici $\mathfrak M(\mathfrak G)$ doit être construit conformément aux \S{}\S{} 6 et 7 : d'après \S{} 7, on prolonge une base algébrique $\mathsf{H}$ en chaque base rationnelle $\Theta$ avec condition accessoire qui en provient, et, pour chaque $\Theta$, on forme d'après \S{} 6 b) la totalité $\mathfrak N(\mathfrak G)$ de tous les domaines $\mathfrak G$ contenant $\Theta$.}

L'énoncé analogue vaut pour les domaines $\mathfrak H$ formés d'entiers transcendants algébriquement entiers.

\section*{\S{} 9. Un exemple d'une infinité dénombrable de classes de domaines $\mathfrak G$.}

La cardinalité des constructions données au \S{} 4 b) et au \S{} 6 b) se lit, d'après les remarques qui les accompagnent, comme égale à la cardinalité de tous les bons ordres, donc comme $2^c$, si $c$ désigne la cardinalité du continu. Mais on ne peut rien en conclure sur la cardinalité de la totalité de tous les domaines $\mathfrak G$, chacun compté une seule fois, et encore moins sur la cardinalité de toutes les classes. Que le nombre de classes ne soit certainement pas fini sera maintenant montré par un exemple d'une infinité dénombrable de domaines $\mathfrak G$ --- analogues à ceux du \S{} 5 --- qui se révéleront tous essentiellement différents et fourniront ainsi une infinité dénombrable de classes.\footnote{Les domaines $\mathfrak G$ du \S{} 6 donnent eux aussi une infinité dénombrable de classes, mais la preuve est un peu plus compliquée.}

Par analogie avec \S{} 5, (1), soit le domaine $\mathfrak G_\sigma$ formé de la totalité des quantités
\begin{equation}
 z\equiv a_0+a_1(\eta)+\cdots+a_{\sigma-1}(\eta)
 \;\mathrm{mod.}\;\mathfrak M_\sigma(\eta),
\tag{1}
\end{equation}
\emph{où $a_i(\eta)$ désigne des formes homogènes de $i$-ième dimension en les $\eta$, et où le module $\mathfrak M_\sigma$ contient pour base tous les produits de puissances de dimension $\sigma$ en les $\eta$}\footnote{Naturellement, dans chaque $z$ individuel, seuls un nombre fini de produits de puissances des $\eta$ apparaissent dans le module.} \emph{et pour multiplicateurs toutes les fonctions algébriquement entières des $\eta$.} Nous montrons que, sous une relation isomorphe entre deux domaines $\mathfrak G_\sigma$ et $\mathfrak G_\tau$ $(\sigma\ne\tau)$, les modules $\mathfrak M_\sigma$ et $\mathfrak M_\tau$ doivent nécessairement se correspondre isomorphiquement; l'impossibilité s'ensuit alors facilement.

Désignons par $\xi$ les indéterminées de $\mathfrak G_\tau$; sous la relation isomorphe, supposons que les quantités $f(\eta)$ de $\mathfrak G_\sigma$ correspondent à ces $\xi$. Comme la relation isomorphe se conserve lors de la formation des quotients, le domaine $\mathfrak G_\zeta$ de toutes les fonctions algébriquement entières des $\xi$ correspond à un domaine intégralement clos $\mathfrak T$ formé de toutes les fonctions entières sur les $f(\eta)$, et donc aussi entières sur le domaine engendré par les $\eta$.

Supposons que (1) corresponde à une quantité de $\mathfrak M_\tau$. Alors, par la propriété de module de $\mathfrak M_\tau$, le produit de (1) par une quantité arbitraire $\psi(\eta)$ de $\mathfrak T$ doit lui aussi appartenir à $\mathfrak G_\sigma$; en formules :
\begin{equation}
(a_0+a_1(\eta)+\cdots+a_{\sigma-1}(\eta))\cdot\psi(\eta)
\equiv b_0+b_1(\eta)+\cdots+b_{\sigma-1}(\eta)
\;\mathrm{mod.}\;\mathfrak M_\sigma.
\tag{2}
\end{equation}
En annulant tous les $\eta$, on obtient $a_0\cdot\psi(0)=b_0$, donc (2) devient
\begin{equation}
(a_0+a_1(\eta)+\cdots+a_{\sigma-1}(\eta))(\psi(\eta)-\psi(0))
\equiv c_1(\eta)+\cdots+c_{\sigma-1}(\eta)
\;\mathrm{mod.}\;\mathfrak M_\sigma.
\tag{3}
\end{equation}
Or, outre $\psi(\eta)$, les fonctions
\[
 \chi_\nu(\eta)=\sqrt[\nu]{\psi(\eta)-\psi(0)}
\]
appartiennent aussi à $\mathfrak T$ pour tout $\nu$; et, puisque $\chi_\nu(0)=0$, les formules analogues à (3) sont :
\begin{equation}
(a_0+a_1(\eta)+\cdots+a_{\sigma-1}(\eta))\chi_\nu(\eta)
\equiv c_1^{(\nu)}(\eta)+\cdots+c_{\sigma-1}^{(\nu)}(\eta)
\pmod{\mathfrak M_\sigma}
\quad(\nu=1,2,\cdots).
\tag{4}
\end{equation}
Soit $A(\eta)$, de degré $\lambda$, le produit de $\chi_1=\psi(\eta)-\psi(0)$ par ses $\varkappa-1$ conjugués. Comme $\chi_1(0)=0$, $A(\eta)$ doit commencer par des termes d'au moins la première dimension. De (4) on tire
\begin{equation}
a_0\cdot\chi_\nu(\eta)\equiv0\pmod{\mathfrak M_1};\qquad
a_0^\nu\cdot\chi_1(\eta)\equiv0\pmod{\mathfrak M_\nu};\qquad
a_0^{\nu\varkappa}\cdot A(\eta)\equiv0\pmod{\mathfrak M_{\nu\varkappa}},
\tag{5}
\end{equation}
et donc, pour $a_0\ne0$, comme au \S{} 3, $\lambda>\nu\varkappa$. Mais comme $\nu<\frac{\lambda}{\varkappa}$ contredit la clôture algébrique-entière de $\mathfrak T$, il faut nécessairement $a_0=0$. De (4) on obtient alors
\begin{equation}
a_1(\eta)\cdot\chi_\nu(\eta)\equiv c_1^{(\nu)}(\eta)\pmod{\mathfrak M_2};\qquad
[a_1(\eta)]^{\nu\varkappa}\cdot A(\eta)
\equiv[c_1^{(\nu)}(\eta)]^{\nu\varkappa}\pmod{\mathfrak M_{\nu\varkappa+1}},
\tag{6}
\end{equation}
et de là, puisque $A(\eta)$ commence par des termes d'au moins la première dimension, $c_1^{(\nu)}(\eta)=0$. Ainsi, en complète analogie avec (5), pour tout $\nu$ :
\[
 a_1(\eta)\cdot\chi_\nu(\eta)\equiv0\pmod{\mathfrak M_2};\qquad
 [a_1(\eta)]^{\nu\varkappa}\cdot A(\eta)\equiv0\pmod{\mathfrak M_{2\nu\varkappa}},
\]
et donc $a_1(\eta)=0$. En continuant de cette manière, l'application alternée de (5) et (6) donne
\[
 a_0=0,\qquad a_1(\eta)=0,\quad \ldots,\quad a_{\sigma-1}(\eta)=0.
\]
Comme la considération analogue vaut aussi relativement à $\mathfrak G_\tau$ : \emph{Sous toute relation isomorphe entre $\mathfrak G_\sigma$ et $\mathfrak G_\tau$ $(\sigma\ne\tau)$, les modules $\mathfrak M_\sigma$ et $\mathfrak M_\tau$ se correspondent isomorphiquement.}

Supposons par exemple que $\sigma>\tau$. Aux indéterminées $\xi$ correspondaient les quantités $f(\eta)$ de $\mathfrak G_\sigma$; et, puisque toutes les quantités de $\mathfrak G_\tau$ peuvent être représentées par des polynômes en les $\xi$ mod. $\mathfrak M_\tau$, la correspondance de $\mathfrak M_\tau$ et de $\mathfrak M_\sigma$ implique que toutes les quantités de $\mathfrak G_\sigma$ s'obtiennent comme polynômes en les $f(\eta)$ mod. $\mathfrak M_\sigma$. Comme $\sigma>\tau\geqq1$, les indéterminées $\eta$ n'appartiennent certainement pas à $\mathfrak M_\sigma$, et donc, puisqu'elles sont exprimables intégralement et rationnellement par les $f(\eta)$ mod. $\mathfrak M_\sigma$, il doit exister un $f(\eta)$ ayant des termes linéaires non nuls. Soit donc
\[
 \zeta:f(\eta)\equiv a_0+a_1(\eta)\pmod{\mathfrak M_2}\quad[a_1(\eta)\ne0],
\]
\[
 \zeta^\tau:[f(\eta)]^\tau\equiv a_0^\tau\;\mathrm{mod.}\;\mathfrak M_1\quad\text{pour }a_0\ne0,
\]
\[
 \hphantom{\zeta^\tau:[f(\eta)]^\tau}\equiv [a_1(\eta)]^\tau
 \;\mathrm{mod.}\;\mathfrak M_{\tau+1}\quad\text{pour }a_0=0.
\]
Or $\zeta^\tau$ appartient à $\mathfrak M_\tau$, tandis que $[f(\eta)]^\tau$ commence par des termes non nuls de dimension zéro ou de dimension $\tau$ et, puisque $\tau<\sigma$, ne peut appartenir à $\mathfrak M_\sigma$, contrairement à la correspondance de $\mathfrak M_\sigma$ et $\mathfrak M_\tau$ obtenue ci-dessus. Le cas $\tau>\sigma$ est traité en même temps en échangeant $\eta$ et $\xi$. Il est donc montré que les domaines $\mathfrak G_\sigma$ et $\mathfrak G_\tau$ sont \emph{non isomorphes, ou essentiellement différents}. Ainsi \emph{les domaines $\mathfrak G_i$ $(i=1,2,\cdots\text{ à l'infini})$ donnent un exemple d'une infinité dénombrable de classes de domaines $\mathfrak G$.}

\emph{Remarque.} Bien que deux domaines quelconques $\mathfrak G_i$ soient non isomorphes, il peut très bien arriver, pour deux domaines, que chacun soit isomorphe à une partie de l'autre.\footnote{Ce comportement peut aussi se produire pour les corps; cf. Steinitz, loc. cit., conclusion.} Soient par exemple $\tau=1$ et $\sigma$ arbitraire :
\[
\mathfrak G_1:a_0+\mathfrak M_1(\zeta);\qquad
G_\sigma:a_0+a_1(\eta)+\cdots+a_{\sigma-1}(\eta)+\mathfrak M_\sigma(\eta).
\]
L'application $\zeta=\eta$ fait de $\mathfrak G_\sigma$ un diviseur propre de $\mathfrak G_1$; inversement, la correspondance biunivoque dans $\mathfrak G_1$ et $\mathfrak G_\sigma$ donnée par $\zeta=\eta^\sigma$, $\eta=\sqrt[\sigma]{\zeta}$ fait de $\mathfrak G_1$ un diviseur propre de $\mathfrak G_\sigma$. Ainsi chaque $\mathfrak G_\sigma$ est aussi isomorphe à un diviseur propre de lui-même. À la différence de la notion d'intersection, la notion de classe d'intersection --- c'est-à-dire d'une classe telle que chacun de ses domaines serait isomorphe à un diviseur propre de chaque domaine de chaque autre classe --- n'existe pas, du moins pas comme notion déterminée de manière unique.

\section*{\S{} 10. Les quantités entières d'un corps arbitraire.}

Les développements des paragraphes précédents ont reposé, pour les domaines $\mathfrak G$ --- pourvu que les nombres algébriques ne soient pas inclus dans les propriétés de définition III et IV des domaines $\mathfrak G$ ---, seulement sur le fait que la totalité de tous les nombres complexes forme un corps;\footnote{Ce n'est qu'au \S{} 7 qu'a été utilisée l'hypothèse supplémentaire que le \og{} corps premier \fg{} des nombres rationnels contient une infinité d'éléments; dans le cas d'un corps premier ne contenant qu'un nombre fini d'éléments, ce paragraphe disparaît simplement, comme nous allons le voir.} pour les domaines $\mathfrak H$, ils ont reposé en outre sur le fait que ce corps est algébriquement clos. Les résultats obtenus admettent donc directement la \emph{généralisation au corps le plus général}, qui, d'après Weber et E. Steinitz,\footnote{loc. cit., introduction.} est défini abstraitement comme \og{} un système d'éléments muni de deux opérations, addition et multiplication, soumises aux lois associative et commutative, liées par la loi distributive, et admettant des opérations inverses sans restriction et uniques \fg{}.\footnote{Seule la division par zéro est à exclure.}

Tout corps de ce genre contient un élément unité, et par conséquent le \og{} corps premier \fg{} qui en dérive par les opérations d'addition et de multiplication et leurs inverses. Ce corps premier est soit du type des nombres rationnels (caractéristique $0$), soit du type du système des classes modulo un nombre premier $p$ (caractéristique $p$).\footnote{Steinitz, \S{} 4.} Les \emph{quantités entières du corps premier} sont alors bien définies comme les quantités issues de l'élément unité par addition et soustraction; pour les corps premiers de caractéristique $p$, les quantités entières épuisent déjà le corps premier et il n'existe pas de quantités fractionnaires du corps premier. Tout corps algébriquement clos contient un \og{} corps absolument algébrique \fg{}, obtenu à partir du corps premier par adjonction de tous les éléments algébriques sur le corps premier; pour les corps de caractéristique $0$, il est donc du type du corps de tous les nombres algébriques. Les \emph{quantités algébriquement entières du corps absolument algébrique} sont alors bien définies comme toutes les quantités algébriquement entières sur les quantités entières du corps premier (ou, ce qui revient au même, sur l'unité); pour les corps de caractéristique $p$, il n'existe donc pas non plus de quantités algébriquement fractionnaires du corps absolument algébrique. Après ces remarques préliminaires, la définition des quantités \og{} entières \fg{} et \og{} algébriquement entières \fg{} se transpose aux corps arbitraires, respectivement aux corps algébriquement clos : \emph{Les quantités entières d'un corps $\Omega$ sont définies comme un domaine entier $\mathfrak G$ de quantités de $\Omega$ dont le corps des fractions\footnote{C'est-à-dire le corps formé de tous les quotients de deux quantités de $\mathfrak G$.} épuise $\mathfrak K$, et qui contient l'élément unité mais aucune quantité fractionnaire du corps premier.}

\emph{Les quantités entières au sens algébrique d'un corps algébriquement clos $\mathfrak L$ sont définies comme un domaine entier $\mathfrak H$ intégralement clos de quantités de $\Omega$ dont le corps des fractions épuise $\mathfrak L$, et qui contient l'élément unité mais aucune quantité algébriquement fractionnaire du corps absolument algébrique.}

Pour les corps de caractéristique zéro, les résultats obtenus dans les paragraphes précédents s'appliquent littéralement dès que l'on remplace seulement les \og{} nombres algébriques \fg{} par les quantités du corps premier, respectivement du corps absolument algébrique. Pour les corps de caractéristique $p$, les résultats se simplifient encore, parce que le corps premier ne contient aucune quantité fractionnaire, de sorte que la condition accessoire pour la base rationnelle et pour les systèmes à adjoindre disparaît simplement.\footnote{Cf. la première note de ce paragraphe.} Ainsi le résultat est le suivant : \emph{La définition abstraite permet comme quantités entières d'un corps de caractéristique $p$ : tout domaine entier dérivé d'une base rationnelle quelconque, le corps lui-même, et tout domaine entier situé entre ces domaines et le corps. De façon correspondante, comme quantités entières au sens algébrique d'un corps algébriquement clos de caractéristique première, on a : tout domaine entier intégralement clos situé entre le domaine dérivé d'une base algébrique et le corps lui-même --- y compris les deux domaines limites.} Pour les corps les plus généraux de caractéristique première, comme pour les corps premiers correspondants, la distinction entre entier et fractionnaire s'efface donc.

\medskip
\noindent Erlangen, 30 mars 1915.
\clearpage
\setcounter{footnote}{0}
\begin{center}
{\Large\bfseries \NoetherNavAnchor{P10}{10. Les équations fonctionnelles de l’application isomorphe}10. Les équations fonctionnelles de l'application\\ isomorphe}\par
\vspace{0.75em}
Par Emmy Noether, à Göttingen.\par
\vspace{1em}
\emph{Math. Ann.} 77 (1916), pp. 536--545
\end{center}

\vspace{2em}

D'après Dedekind,\footnote{Cf. par exemple Dirichlet-Dedekind, \emph{Leçons sur la théorie des nombres} (4e éd.), Supplément XI, \S\ 161. La désignation \og{}application isomorphe\fg{} ou \og{}isomorphisme\fg{}, calquée sur la théorie des groupes, ne se rencontre que dans la littérature postérieure; Dedekind parle des \og{}permutations d'un corps\fg{}.} on entend par \emph{application isomorphe du corps} $\Afield$ \emph{sur un système} $\Bfield$ --- qui se révèle après coup être lui aussi un corps --- \emph{une correspondance univoque telle qu'à chaque élément de} $\Afield$ \emph{corresponde un et un seul élément de} $\Bfield$; \emph{et qu'à la somme, à la différence, au produit et au quotient de deux éléments quelconques de} $\Afield$ \emph{soient de nouveau associés la somme, la différence, le produit et le quotient des éléments de} $\Bfield$ \emph{qui leur correspondent univoquement}.

Une telle application est réalisée par une fonction $f(z)$, univoque d'après ce qui précède. Si alors $z$ parcourt tous les éléments $x,y,\cdots$ du corps $\Afield$, $f(z)$ est donc caractérisée par les équations fonctionnelles
\[
\begin{array}{ll}
(1)\quad f(x+y)=f(x)+f(y); & (2)\quad f(x-y)=f(x)-f(y);\\[0.45em]
(3)\quad f(x\cdot y)=f(x)\cdot f(y); & (4)\quad f\!\left(\dfrac{x}{y}\right)=\dfrac{f(x)}{f(y)}.
\end{array}
\]

\emph{Dans ce qui suit seront indiquées les solutions univoques les plus générales} $f(z)$ \emph{de ces équations fonctionnelles, lorsque} $x,y,\cdots$ \emph{parcourent tous les nombres réels et complexes; c'est-à-dire la fonction la plus générale} $f(z)$ \emph{qui réalise une application isomorphe du corps de tous les nombres complexes}.\footnote{Cette question m'a été posée occasionnellement par M. Landau. Comme je l'ai remarqué après coup, une partie des solutions se trouve déjà chez H. Lebesgue, \emph{Sur les transformations ponctuelles, transformant les plans en plans}, Atti Acad. Torino, 1906/07; et aussi implicitement chez A. Ostrowski, \emph{Über einige Fragen der allgemeinen Körpertheorie}, Crelle's Journal 143 (1913), \S\ 2.} La solution pour des corps arbitraires s'obtient de façon tout à fait analogue.

G. Hamel\footnote{G. Hamel, \emph{Eine Basis aller Zahlen und die unstetigen Lösungen der Funktionalgleichung:} $f(x+y)=f(x)+f(y)$, Math. Ann. 60, p. 459 (1905).} a construit les solutions les plus générales de l'équation fonctionnelle (1) au moyen d'une base \emph{linéaire} de tous les nombres. La construction donnée ici pour les solutions de (1) à (4) --- donc pour la fonction d'application relativement à laquelle les opérations rationnelles et algébriques sont invariantes --- utilise la base \emph{rationnelle} et \emph{algébrique} de tous les nombres. Après que l'application isomorphe aura été discutée de plus près en 1, des conditions nécessaires pour $f(z)$ seront données en 2; en 3 elles seront reconnues comme suffisantes et conduiront à la construction effective.\footnote{Cf. les considérations parallèles dans mon travail: \emph{Die allgemeinsten Bereiche aus ganzen transzendenten Zahlen}. Math. Ann. 77, p. 103 (1915), en particulier \S\S\ 8 et 10.} Les solutions de (1) à (4), à l'exception connue de $f(z)=z$ ou $\overline z$, sont discontinues; en 4 on montre en outre qu'elles sont \og{}extrêmement discontinues\fg{}, fait que G. Hamel a démontré pour les solutions réelles de (1), mais qui n'a pas nécessairement lieu dans le domaine complexe.\footnote{L'expression de Hamel \og{}totalement discontinu\fg{} est employée dans le reste de la littérature en un autre sens.}

\medskip
\noindent\textbf{1.}\quad Nous tirons d'abord quelques conséquences des équations fonctionnelles (1) à (4), directement pour des corps généraux $\Afield$. D'après (1), l'univocité entraîne $f(0)=0$. Si $f(y)$ ne s'annule pas, alors le $y$ initial ne peut pas non plus s'annuler, et les équations fonctionnelles (1) à (4) montrent ainsi que \emph{le système d'application} $\Bfield$ \emph{de toutes les valeurs} $f(z)$ \emph{forme de nouveau un corps}. Inversement, si l'on prescrit la condition initiale $f(0)=0$, alors (2) entraîne l'univocité de la fonction $f(z)$: \emph{l'exigence d'univocité et la condition initiale} $f(0)=0$ \emph{sont donc équivalentes}. De (4) il résulte que $f(z)$ ne peut pas s'annuler identiquement; de (3), en outre, que $f(z)$ ne s'annule que pour $z=0$. Car de $f(y)=0$ et $y\ne0$ il résulterait, puisque $\frac{z}{y}$ appartient aussi à $\Afield$, que
\[
f(z)=f\!\left(\frac{z}{y}\cdot y\right)=f\!\left(\frac{z}{y}\right)\cdot f(y)=0,
\]
ce qui contredirait (4) en donnant l'annulation identique de $f(z)$. L'égalité $f(x)=f(y)$ entraîne donc $x=y$; c'est-à-dire qu'à chaque valeur de $f(z)$ correspond aussi une et une seule valeur de $z$: \emph{$f(z)$ est une fonction réciproquement univoque de} $z$. Comme le montrent les équations fonctionnelles, l'application définie par la fonction inverse est à son tour isomorphe. Puisque toute opération rationnelle est composée d'un nombre fini d'additions, de multiplications et de leurs inverses, on peut donc dire aussi: \emph{par l'application isomorphe est établie une relation biunivoque entre les éléments de} $\Afield$ \emph{et ceux de} $\Bfield$, \emph{de telle sorte que toutes les relations rationnelles existant entre un nombre fini quelconque d'éléments de} $\Afield$,
\[
F(x,y,\cdots,t)=0,
\]
\emph{sont conservées dans} $\Bfield$, \emph{et réciproquement}.

Il faut encore remarquer que, pour un corps, $f(z)$ est déjà caractérisée par les équations fonctionnelles (1) et (3), par l'exigence que $f(z)$ ne s'annule pas identiquement, et par la condition initiale $f(0)=0$. En effet, (2) résulte de (1) en remplaçant $x$ par l'élément $(x-y)$ appartenant à $\Afield$; puis (2) et $f(0)=0$ donnent l'univocité de $f(z)$. Comme on l'a montré plus haut, il résulte en outre de (3), en remplaçant $x$ par l'élément $\frac{x}{y}$ appartenant à $\Afield$, que $f(y)$ ne peut pas s'annuler pour $y\ne0$, donc la validité de (4) et en même temps la biunivocité. Si, au lieu d'un corps, on prend pour base un domaine d'intégrité $\Sdomain$, $\frac{x}{y}$ n'appartient pas nécessairement à $\Sdomain$, et l'on ne peut pas conclure de (3) à la biunivocité. Ici suffit cependant la formulation la plus connue: une application isomorphe est une correspondance \emph{biunivoque} entre les éléments de $\Sdomain$ et de $\Sdomain'$ telle qu'à la somme et au produit correspondent toujours la somme et le produit.\footnote{Cf. pour le n\textsuperscript{o} 1: Dedekind, loc. cit., \S\ 161.}

\medskip
\noindent\textbf{2.}\quad Désormais, pour simplifier, on prendra pour base, à la place de $\Afield$, le corps $\Cfield$ de tous les nombres réels et complexes. Il s'agit d'abord de discuter les systèmes de valeurs que prend $f(z)$. De (3) ou de (4), on déduit que $f(1)=1$, et dès lors, en vertu des équations fonctionnelles, pour tout \emph{nombre rationnel} $\alpha$, $f(\alpha)=\alpha$; \emph{ainsi tout nombre rationnel correspond à lui-même dans l'application}. Comme un \emph{nombre algébrique} $\beta$ satisfait à une équation irréductible à coefficients numériques rationnels $F(\beta,\alpha)=0$, le n\textsuperscript{o} 1 et la remarque précédente donnent, pour $f(\beta)$, l'équation $F(f(\beta),\alpha)=0$; \emph{tout nombre algébrique passe donc en lui-même ou en l'un de ses conjugués}, et, par biunivocité, à la totalité des nombres algébriques correspond de nouveau le corps de tous les nombres algébriques. Afin de reconnaître le comportement des nombres \emph{transcendants}, et en même temps de séparer les valeurs conjuguées, nous prenons pour base une \emph{base rationnelle de tous les nombres},\footnote{Cf. mon travail cité sur les \og{}nombres transcendants entiers\fg{}, \S\ 6.} c'est-à-dire un système $\Theta$ de nombres $\vartheta$ qui permet d'exprimer tous les nombres rationnellement --- avec des coefficients numériques rationnels\footnote{Par \og{}fonction rationnelle\fg{}, j'entendrai toujours une fonction dont les coefficients sont eux aussi des nombres rationnels.} --- tandis que, dans au moins un bon ordre, aucun nombre de base n'est exprimable rationnellement au moyen d'un nombre fini de précédents. L'existence de cette base résulte du théorème du bon ordre exactement comme celle des bases linéaires et algébriques.

Soumettons $\Cfield$ à un bon ordre $\Omega$. Alors tout nombre est ou bien exprimable rationnellement au moyen d'un nombre fini de nombres précédents, ou bien rationnellement indépendant des précédents. Ces derniers nombres forment la base $\Theta$, et le raisonnement par induction montre qu'effectivement tout nombre est exprimable rationnellement au moyen d'un nombre fini de $\vartheta$.

À chaque nombre de base $\vartheta$ est défini le \og{}corps de segment $\Kfield(\vartheta)$\fg{}: il est constitué de toutes les combinaisons rationnelles des nombres de base qui précèdent $\vartheta$ dans le bon ordre. Comme corps de segment du premier nombre de base, il faut considérer le corps $\Rfield$ de tous les nombres rationnels. Le nombre $\vartheta$ peut présenter deux sortes de comportement par rapport à $\Kfield(\vartheta)$: il peut être algébriquement dépendant de $\Kfield(\vartheta)$, ou algébriquement indépendant de lui. Puisque, d'après le n\textsuperscript{o} 1, toutes les relations valables dans $\Cfield$ sont aussi valables dans le corps image, et réciproquement, la totalité des valeurs $f(\vartheta)$ correspondant aux $\vartheta$ forme une base du domaine image, bien ordonnée au moyen du bon ordre de $\Theta$ lui-même. En particulier, au corps de segment $\Kfield(\vartheta)$ correspond le corps de segment $\Kfield(f(\vartheta))$; et selon que $\vartheta$ est algébriquement dépendant ou indépendant de $\Kfield(\vartheta)$, il en va de même de $f(\vartheta)$ relativement à $\Kfield(f(\vartheta))$. Dans le cas de dépendance, les équations irréductibles pour $\vartheta$ et pour $f(\vartheta)$ se correspondent aussi. La totalité des valeurs $\vartheta$ qui sont à chaque fois algébriquement indépendantes de $\Kfield(\vartheta)$ forment --- comme le montre encore l'induction --- une \emph{base algébrique de tous les nombres},\footnote{Cf. E. Zermelo, \emph{Über ganze transzendente Zahlen}, \S\ 1, Math. Ann. 75 (1914).} c'est-à-dire un système $\mathsf{H}$ de nombres $\eta$ qui permet d'exprimer algébriquement tous les nombres, tandis qu'aucune relation algébrique n'existe entre les $\eta$ eux-mêmes. À cette base algébrique $\mathsf{H}$ correspond dans le domaine image de nouveau une base algébrique $\mathsf{Z}$, qui, par biunivocité, a la même puissance que $\mathsf{H}$, à savoir la puissance du continu, puisque l'on sait qu'une extension algébrique n'augmente pas la puissance.\footnote{Cf. E. Steinitz, \emph{Algebraische Theorie der Körper}, Crelle's Journal 137 (1910), \S\ 24.} Mais de la connaissance des valeurs $f(\vartheta)$ on obtient immédiatement $f(z)$ pour tous les systèmes de valeurs, puisque l'égalité $z=\psi(\vartheta)$ entraîne toujours $f(z)=\psi(f(\vartheta))$.

\medskip
\noindent\textbf{3.}\quad Nous montrons maintenant que les conditions nécessaires trouvées ci-dessus fournissent effectivement aussi la construction de $f(z)$. La base $Z$ qui correspond biunivoquement à la base algébrique $H$ est construite et associée à $\mathsf{H}$ de la manière suivante. Posons un second bon ordre $\Omega'$, qui fournit une base algébrique $\mathsf{Z}'$ de tous les nombres. Que $\mathsf{Z}$, ayant pour éléments les $\xi$, soit un sous-ensemble de $\mathsf{Z}'$ de même puissance que $\mathsf{Z}'$, et donc que $\mathsf{H}$. À chaque élément $\eta_i$ de $\mathsf{H}$, associons un et un seul élément $\xi_i$, portant la même numérotation, par la règle suivante: $\xi_i$ est le \emph{premier}, dans le bon ordre $\Omega'$, de tous les éléments de $\mathsf{Z}$ qui n'ont été associés à aucun élément de $H$ précédant $\eta_i$.

Définissons maintenant $f(z)$ ainsi:
\begin{enumerate}
\item[(a)] \emph{pour tous les nombres rationnels} $\alpha$, \emph{par} $f(0)=0$, $f(\alpha)=\alpha$;
\item[(b)] \emph{pour toutes les grandeurs de la base algébrique} $\mathsf{H}$, \emph{par} $f(\eta_i)=\xi_i$;
\item[(c)] \emph{pour toutes les autres grandeurs} $\vartheta$ \emph{de la base rationnelle} $\Theta$ --- \emph{celles qui sont algébriquement dépendantes du corps de segment} $\Kfield(\vartheta)$, \emph{et satisfont donc à une équation} $F(\vartheta;\vartheta_{i_1},\cdots,\vartheta_{i_\tau})=0$ \emph{irréductible dans} $\Kfield(\vartheta)$ --- \emph{comme racine de l'équation} $F(f(\vartheta);f(\vartheta_{i_1}),\cdots,f(\vartheta_{i_\tau}))=0$;
\item[(d)] \emph{pour toutes les autres valeurs} $z$, \emph{qui, par la définition de} $\Theta$, \emph{sont représentables sous la forme} $z=\psi(\vartheta_{k_1}\cdots\vartheta_{k_\sigma})$, \emph{par} $f(z)=\psi(f(\vartheta_{k_1})\cdots f(\vartheta_{k_\sigma}))$.\footnote{Dans (d), (a) est contenu comme cas particulier.}
\end{enumerate}

Pour montrer que le $f(z)$ ainsi défini satisfait effectivement aux équations fonctionnelles (1) à (4), il suffit, d'après le n\textsuperscript{o} 1, de le prouver pour (1) et (3), puisque $\Cfield$ est un corps et que $f(z)$, à cause de (a) et (b), ne peut pas s'annuler identiquement et satisfait en outre à la condition initiale $f(0)=0$. Au moyen de la base rationnelle $\Theta$, exprimons $x,y,x+y$ sous la forme
\[
x=\psi_1(\vartheta_1\cdots\vartheta_t),\quad
 y=\psi_2(\vartheta_1\cdots\vartheta_t),\quad
 (x+y)=\psi_3(\vartheta_1\cdots\vartheta_t).
\]
Ainsi la preuve que $f(z)$ satisfait à l'équation fonctionnelle (1),
\[
f(x)+f(y)-f(x+y)=0,
\]
revient, d'après (d), à montrer que l'égalité
\[
\psi_1(\vartheta)+\psi_2(\vartheta)-\psi_3(\vartheta)
 =\frac{H_1(\vartheta)}{H_2(\vartheta)}=0
\]
entraîne toujours
\[
\psi_1(f(\vartheta))+\psi_2(f(\vartheta))-\psi_3(f(\vartheta))
 =\frac{H_1(f(\vartheta))}{H_2(f(\vartheta))}=0,
\]
où $H_1,H_2$ désignent des fonctions rationnelles entières de leurs arguments. L'équation fonctionnelle (3) conduit exactement à la même forme de condition; ainsi la satisfaction de toutes les équations fonctionnelles est identique à la condition suivante --- nécessaire aussi d'après les n\textsuperscript{os} 1 et 2: pour \emph{toute} fonction rationnelle entière $H(\vartheta)$, l'égalité $H(\vartheta)=0$ entraîne toujours $H(f(\vartheta))=0$, et $H(\vartheta)\ne0$ entraîne $H(f(\vartheta))\ne0$, ce dernier point à cause de l'apparition du dénominateur.

Que tel soit effectivement le cas se montre par induction. Supposons que, dans $H(\vartheta_1\cdots\vartheta_t)$, la grandeur de base $\vartheta_t$ soit, dans le bon ordre, la dernière parmi $\vartheta_1\cdots\vartheta_t$.\footnote{Les expressions $H(\vartheta)$ qui sont de degré zéro en les $\vartheta$ passent en elles-mêmes sous l'application et n'ont donc pas besoin d'être examinées davantage.} Alors $H(\vartheta_1\cdots\vartheta_t)=0$ peut être considéré comme une relation à laquelle satisfait $\vartheta_t$ relativement au corps de segment $\Kfield(\vartheta_t)$, et de même $H(f(\vartheta_1)\cdots f(\vartheta_t))=0$ comme une relation à laquelle satisfait $f(\vartheta_t)$ relativement à $\Kfield(f(\vartheta_t))$. Si toutes les relations $H(\vartheta)=0$ n'étaient pas également remplies par $f(\vartheta)$, et réciproquement, il devrait exister une première grandeur de base, disons $\vartheta_0$, pour laquelle au moins une relation valable relativement à $\Kfield(\vartheta_0)$ ne serait pas conservée dans le corps image, ou inversement, alors que les corps de segment précédents $\Kfield(\vartheta_0)$ et $\Kfield(f(\vartheta_0))$ seraient isomorphes. En particulier, d'après (a), le corps de segment de la première grandeur de base de $\Theta$, le corps $\Rfield$ de tous les nombres rationnels, est isomorphe à son corps image. Supposons maintenant que $H(\vartheta_0)$ ait la forme
\[
H(\vartheta_0)=a_0(\vartheta)\cdot\vartheta_0^{\varkappa}
+a_1(\vartheta)\cdot\vartheta_0^{\varkappa-1}+\cdots+a_{\varkappa}(\vartheta),
\]
où les $a_i(\vartheta)$ sont des grandeurs de $\Kfield(\vartheta_0)$. Il y a alors deux possibilités.

1) $\vartheta_0$ est algébriquement indépendant de $\Kfield(\vartheta_0)$. Alors $H(\vartheta_0)=0$ est satisfait identiquement en $\vartheta_0$,\footnote{Comme $x,y,\cdots$ peuvent être exprimés par les $\vartheta$, cette forme de relation peut très bien se produire.} de sorte que $a_i(\vartheta)=0$ pour tous les coefficients, et donc, par l'isomorphisme de $\Kfield(\vartheta_0)$ et $\Kfield(f(\vartheta_0))$, aussi $a_i(f(\vartheta))=0$. L'égalité $H(\vartheta_0)=0$ entraîne donc toujours $H(f(\vartheta_0))=0$. Réciproquement, supposons $H(f(\vartheta_0))=0$. Puisque $\vartheta_0$, étant algébriquement indépendant de $\Kfield(\vartheta_0)$, appartient à la base algébrique $\mathsf{H}$, $f(\vartheta_0)$ appartient par (b) à la base algébrique $\mathsf{Z}$ et est algébriquement indépendant de $\Kfield(f(\vartheta_0))$. Ainsi tous les $a_i(f(\vartheta))$ s'annulent, et par l'isomorphisme tous les $a_i(\vartheta)$ s'annulent. L'égalité $H(f(\vartheta_0))=0$ entraîne $H(\vartheta_0)=0$; par contraposition, $H(\vartheta_0)\ne0$ entraîne aussi $H(f(\vartheta_0))\ne0$.

2) $\vartheta_0$ est algébriquement dépendant de $\Kfield(\vartheta_0)$. Alors $H(\vartheta_0)$ est divisible par l'équation $F(\vartheta_0)$ irréductible relativement à $\Kfield(\vartheta_0)$; c'est-à-dire que, identiquement en $t$,
\[
H(t;a_i(\vartheta))=F(t,\vartheta)\cdot G(t,\vartheta_i).
\]
Comme tous les coefficients qui interviennent appartiennent à $\Kfield(\vartheta_0)$, la relation correspondante est remplie dans le corps isomorphe $\Kfield(f(\vartheta_0))$:
\[
H(t;a_i(f(\vartheta)))=F(t;f(\vartheta))\cdot G(t;f(\vartheta)).
\]
Mais, d'après (c), $f(\vartheta_0)$ a été choisi comme racine de $F(t;f(\vartheta))$. L'égalité $H(\vartheta_0)=0$ entraîne donc $H(f(\vartheta_0))=0$.

Réciproquement, si $H(f(\vartheta_0))=0$, alors la divisibilité de $H(t;a_i(f(\vartheta)))$ par la fonction $F(t;f(\vartheta))$, irréductible relativement à $\Kfield(f(\vartheta_0))$, entraîne, pour le corps isomorphe $\Kfield(\vartheta_0)$, la divisibilité de $H(t;a_i(\vartheta))$ par $F(t;\vartheta)$. Et puisque $F(t;f(\vartheta))$ provenait de l'équation irréductible de $\vartheta_0$ par la relation isomorphe, et que cette relation est biunivoque, $F(t;\vartheta)$ représente nécessairement la fonction irréductible dont $\vartheta_0$ est racine. Ainsi, $H(f(\vartheta_0))=0$ implique $H(\vartheta_0)=0$; de même, $H(\vartheta_0)\ne0$ implique $H(f(\vartheta_0))\ne0$.

De corps de segment isomorphes $\Kfield(\vartheta_0)$ et $\Kfield(f(\vartheta_0))$ naissent donc aussi, par adjonction de $\vartheta_0$ et de $f(\vartheta_0)$ respectivement, des corps isomorphes; l'hypothèse que cela échoue pour un certain $\vartheta_0$ conduit à une contradiction. Il est ainsi démontré que \emph{la fonction} $f(z)$ \emph{définie par (a) à (d) représente effectivement une solution des équations fonctionnelles (1) à (4), et, d'après 2, la plus générale}.

Toutes les considérations utilisées pour construire $f(z)$ n'ont employé que le fait que la totalité $\Cfield$ de tous les nombres forme un corps; elles valent donc pour tout corps $\Afield$ défini abstraitement. Il faut noter ici que le corps $\Rfield$ des nombres rationnels est remplacé par le \og{}corps premier\fg{} de $\Afield$, c'est-à-dire par le corps dérivé de l'élément unité de $\Afield$ par les opérations d'addition et de multiplication et leurs inverses. Si donc, dans (a) à (d), on remplace les nombres rationnels par les éléments du corps premier, alors \emph{la fonction} $f(z)$ \emph{ainsi obtenue réalise l'application isomorphe la plus générale d'un corps arbitraire défini abstraitement}.

\medskip
\noindent\textbf{4.}\quad Il reste à montrer que les fonctions $f(z)$ définies pour le corps $\Cfield$ de tous les nombres réels et complexes --- qui, à l'exception de $f(z)=z$ ou $\overline z$, sont connues pour être discontinues --- sont \emph{extrêmement discontinues}; c'est-à-dire que, dans tout voisinage de n'importe quel nombre réel ou complexe $z_0$, il existe des valeurs $z$ pour lesquelles $f(z)$ s'approche arbitrairement de n'importe quelle valeur prescrite $Z_0$. Cette extrême discontinuité appartient, comme G. Hamel l'a démontré loc. cit., aux solutions \emph{réelles} $\varphi(x)$, définies pour tous les nombres \emph{réels}, de l'équation fonctionnelle (1). Mais, comme le montrent les exemples ci-dessous, elle n'appartient pas nécessairement aux solutions définies sur les nombres complexes. Hamel raisonne ainsi. Si $\varphi(x)$ diffère de $C\cdot x$, il existe au moins deux valeurs de la base linéaire\footnote{La base linéaire de tous les nombres (réels) est donnée par un système de nombres (réels) qui permet d'exprimer tous les nombres (réels) linéairement, à coefficients rationnels, tandis qu'entre un nombre fini quelconque de nombres de base il n'existe aucune relation linéaire à coefficients rationnels.}, disons $x_1$ et $x_2$, telles que $\varphi(x_1):x_1\ne\varphi(x_2):x_2$. Alors les deux équations
\[
\alpha x_1+\beta x_2=a;\qquad \alpha\varphi(x_1)+\beta\varphi(x_2)=b
\]
ont une solution en nombres réels $\alpha,\beta$; par conséquent $a$ et $b$ peuvent être approchés arbitrairement par des nombres rationnels $\alpha,\beta$. Comme $\alpha,\beta$ sont rationnels, $\alpha\varphi(x_1)+\beta\varphi(x_2)$ représente alors effectivement $\varphi(\alpha x_1+\beta x_2)$. Ce raisonnement échoue pour les systèmes de valeurs complexes. En effet, on peut donner, aussi dans le domaine complexe, des solutions de (1) qui sont discontinues mais non extrêmement discontinues; par exemple
\[
\varphi(z)=\varphi(x+iy)=\varphi(x)+i(y+\varphi(x)),
\]
où $\varphi(x)$ désigne une solution réelle extrêmement discontinue dans le réel; ou, plus simplement encore, $\varphi(z)=\varphi(x)$. Dans les deux cas la partie réelle de $\varphi(z)$ est extrêmement discontinue, tandis que la partie imaginaire est déterminée par la partie réelle.

Ce comportement, et en même temps le comportement de la fonction $f(z)$, devient tout à fait clair si l'on introduit la notion de rang. Posons $z=x+iy$, $\varphi(z)=X+iY$, et appelons $\varphi(z)$ de rang quatre, trois, deux ou un, respectivement, selon qu'il n'existe aucune, une, deux ou trois relations linéairement indépendantes
\[
c_1x+c_2y+c_3X+c_4Y=0,
\]
où les $c$ sont naturellement des nombres réels.

1) Que $\varphi(z)$ soit de \emph{rang quatre}. Il existe alors au moins quatre valeurs de la base linéaire, disons $z_1,z_2,z_3,z_4$, telles que le déterminant
\[
\begin{vmatrix}
x_1&x_2&x_3&x_4\\
y_1&y_2&y_3&y_4\\
X_1&X_2&X_3&X_4\\
Y_1&Y_2&Y_3&Y_4
\end{vmatrix}
\]
ne s'annule pas. Les équations
\[
\begin{aligned}
\alpha x_1+\beta x_2+\gamma x_3+\delta x_4&=a,\\
\alpha y_1+{}\cdot{}\ \cdot{}\ \cdot{}+\delta y_4&=b,\\
\alpha X_1+{}\cdot{}\ \cdot{}\ \cdot{}+\delta X_4&=A,\\
\alpha Y_1+{}\cdot{}\ \cdot{}\ \cdot{}+\delta Y_4&=B
\end{aligned}
\]
ont donc une solution en nombres réels $\alpha,\cdots,\delta$, et $a,b,A,B$ peuvent être approchés arbitrairement par des nombres rationnels $\alpha,\beta,\gamma,\delta$. De plus on a
\[
\alpha(X_1+iY_1)+\beta(X_2+iY_2)+\gamma(X_3+iY_3)+\delta(X_4+iY_4)
 =\varphi(\alpha z_1+\beta z_2+\gamma z_3+\delta z_4).
\]
\og{}En général\fg{}, par conséquent, les solutions $\varphi(z)$ sont extrêmement discontinues aussi dans le domaine complexe; les \emph{valeurs de discontinuité} de $\varphi(z)$ au voisinage de $z_0=a+ib$ \emph{forment une variété linéaire bidimensionnelle}.

2) Que $\varphi(z)$ soit de \emph{rang trois}. Alors il existe au moins trois systèmes de valeurs de la base linéaire, disons $z_1,z_2,z_3$, tels que le déterminant ci-dessus ait rang trois. Les valeurs $a,b,A,B$ satisfaisant à la relation
\[
c_1a+c_2b+c_3A+c_4B=0,
\]
et celles-là seulement peuvent être approchées arbitrairement; les \emph{valeurs de discontinuité} forment une \emph{variété linéaire unidimensionnelle} déterminée par $z_0=a+ib$. Comme le montrent les exemples indiqués, ce cas peut effectivement se produire.

3) Que $\varphi(z)$ soit de \emph{rang deux}. Puisqu'on peut toujours choisir deux nombres complexes $z_1,z_2$ tels que
\[
\begin{vmatrix}x_1&x_2\\ y_1&y_2\end{vmatrix}\ne0,
\]
les relations ont la forme
\[
X=\lambda_1x+\lambda_2y;\qquad Y=\mu_1x+\mu_2y.
\]
L'expression qui en résulte
\[
\varphi(z)=(\lambda_1x+\lambda_2y)+i\cdot(\mu_1x+\mu_2y)
\]
est précisément la solution \emph{continue} la plus générale de l'équation fonctionnelle (1) dans le domaine des nombres complexes.

4) Que $\varphi(z)$ soit de \emph{rang un}. Ce cas ne peut pas se produire dans le domaine des nombres complexes; dans le domaine des nombres réels il signifie la solution \emph{continue} $\varphi(x)=C\cdot x$.

Nous montrons maintenant que les \emph{solutions} $f(z)$ \emph{des équations fonctionnelles (1) à (4), définies pour tous les nombres réels et complexes, ne peuvent avoir que le rang deux ou le rang quatre}. Ici le rang deux ne donne que les deux fonctions \emph{continues} $f(z)=z$ ou $\overline z$, comme le montre la spécialisation $f(1)=1$ et $f(i)=\pm i$, combinée avec 3). Le rang un étant déjà exclu, il reste à exclure le rang trois. Nous supposons donc qu'il existe au moins une relation
\[
c_1x+c_2y+c_3X+c_4Y=0,
\]
de sorte que $f(z)$ ait rang trois, ou peut-être un rang plus bas, et nous montrons que, sous cette hypothèse, c'est toujours le dernier cas qui se produit. Comme le montre de nouveau la spécialisation $f(1)=1$, $f(i)=\pm i$, la relation devient
\[
c_3(X-x)+c_4(Y\mp y)=0,
\]
où $c_3$ et $c_4$ ne s'annulent pas simultanément. Supposons d'abord, par exemple, $c_4=0$. La relation $(X-x)=0$ signifie alors qu'à un $z$ purement imaginaire correspond aussi un $f(z)$ purement imaginaire. Ainsi, pour tout $\mathfrak{y}$ réel, on a $f(i\mathfrak{y})=i\cdot\Yreal$, où $\Yreal$ est de nouveau réel. De là, puisque
\[
f(iy)=\pm i f(y),
\]
on obtient aussi $f(y)=\pm\Yreal$; par conséquent, à tout $z$ réel correspond aussi un $f(z)$ réel, et, à cause de $(X-x)$, toutes les valeurs réelles passent en elles-mêmes. On n'a donc que $f(z)=z$ ou $\overline z$, c'est-à-dire effectivement le rang deux. On conclut de façon tout à fait analogue lorsque $c_3=0$, $c_4\ne0$.

Supposons maintenant $c_3$ et $c_4$ différents de zéro, de sorte que la relation puisse s'écrire sous la forme
\[
(Y\mp y)=c\cdot(X-x)
\]
avec $c$ réel, non nul et fini. Cela signifie que pour $y=\pm cx$ on a aussi $Y=c\cdot X$. Ainsi, en tenant compte des équations fonctionnelles, pour tout $\mathfrak{x}$ réel on a
\[
f(\mathfrak{x}\cdot(1\pm ic))=\Xreal(1+ic)=f(\mathfrak{x})\cdot(1+i f(c)),
\]
où $\Xreal$ est de nouveau réel. Mais ici, d'après le n\textsuperscript{o} 1, le facteur de $f(\mathfrak{x})$ ne peut pas s'annuler identiquement, et la division donne
\[
f(\mathfrak{x})=\Xreal(\sigma+i\tau)
\]
avec $\sigma$ et $\tau$ réels, indépendants de $x$ et de $\Xreal$. En particulier, à $\mathfrak{x}=1$ appartient aussi un $\mathfrak{X}_0$ réel, et la condition $1=\mathfrak{X}_0(\sigma+i\tau)$ montre que nécessairement $\tau=0$, de sorte que $f(\mathfrak{x})=\sigma\cdot\Xreal$. Donc à tout $z$ réel correspond aussi un $f(z)$ réel; de façon équivalente, $y=0$ implique nécessairement $Y=0$. Par là la relation linéaire pour $z$ réel devient $X-x=0$;\footnote{On utilise ici le fait que $c\ne0$, de sorte que ni $c_3$ ni $c_4$ ne s'annule, tandis qu'auparavant seul $c_4\ne0$ avait été utilisé; il fallait donc traiter d'avance le cas où l'un de ces deux s'annule.} chaque valeur réelle correspond à elle-même. On a de nouveau seulement $f(z)=z$ ou $\overline z$, donc effectivement le rang deux. Le rang trois est ainsi exclu dans tous les cas. Mais, puisque --- comme le montrent les recherches des premiers numéros --- des solutions discontinues existent effectivement et doivent donc nécessairement avoir le rang quatre, on obtient d'après 1) le théorème: \emph{\og{}Toutes les solutions} $f(z)$ \emph{des équations fonctionnelles (1) à (4) --- à l'exception des solutions continues} $f(z)=z$ \emph{ou} $\overline z$ --- \emph{sont extrêmement discontinues dans le réel et dans le complexe}.\fg{}\footnote{Lebesgue montre loc. cit. que la partie réelle et la partie imaginaire de $f(z)$ sont toutes deux des fonctions \og{}non mesurables\fg{} de $x$ et de $y$; comme le montre l'exemple du début du n\textsuperscript{o} 4, $\varphi(z)=\varphi(x)+i(y+\varphi(x))$, on ne peut pas encore en conclure l'extrême discontinuité dans le domaine complexe.}

\medskip
\noindent Erlangen, 30 octobre 1915.
\clearpage
\setcounter{footnote}{0}
\begin{center}
{\Large\bfseries \NoetherNavAnchor{P11}{11. Équations à groupe prescrit}11. Équations à groupe prescrit}\par
\vspace{0.75em}
Par Emmy Noether, à Göttingen.\par
\vspace{1.5em}
\emph{Math. Ann.} 78 (1918), pp. 221--229
\end{center}

\vspace{1.5em}

Le problème de la construction d'équations à groupe prescrit peut être abordé dans deux directions, que l'on peut désigner brièvement comme la direction \og{}irrationnelle\fg{}, ou caractérisant les racines, et la direction \og{}rationnelle\fg{}, ou caractérisant les coefficients.

Dans la direction \og{}irrationnelle\fg{}, qui procède du point de vue de la théorie des fonctions et de l'arithmétique, se place le théorème de Kronecker suivant lequel tous les corps abéliens dans le domaine des nombres rationnels sont des corps cyclotomiques, ainsi que les théorèmes correspondants pour les corps relativement abéliens par rapport à des corps quadratiques. Ces théorèmes donnent, d'une part, la réalisabilité de tous les groupes abéliens par rapport à ces domaines spéciaux de rationalité; d'autre part, ils fournissent la totalité des équations et donnent en même temps un aperçu plus profond de la structure des corps de nombres ainsi définis. Mais chaque théorème particulier est restreint à un domaine spécial de rationalité, et chaque nouveau domaine de rationalité exige un traitement entièrement nouveau.

La direction \og{}rationnelle\fg{}, qui procède algébriquement, s'appuie sur le théorème d'irréductibilité de Hilbert, d'après lequel on peut, dans la représentation des coefficients, se restreindre à une représentation paramétrique. En relève par exemple l'existence, relativement à n'importe quel domaine de rationalité, d'une infinité d'équations ayant pour groupe le groupe alterné. Ici manque naturellement l'aperçu, décrit plus haut, sur la structure des corps définis par les équations; mais l'avantage de cette direction \og{}rationnelle\fg{} est que la réalisabilité des groupes est démontrée simultanément pour tous les domaines de rationalité, bien que les représentations paramétriques connues jusqu'ici ne donnent pas en général la totalité des équations.\footnote{Pour les équations résolubles, la représentation paramétrique des coefficients peut être remplacée par celle des racines. Récemment F. Mertens a donné les représentations les plus générales des racines pour certains groupes de degré 8: \og{}Gleichungen mit Quaternionengruppen\fg{}, Sitzb. d. Ak. d. Wiss., Wien, Abt. IIa, vol. 125, p. 735. Comme je l'apprends par cette note, G. Bucht avait déjà donné auparavant les expressions les plus générales des racines pour les formes normales toujours constructibles des équations de degrés 3 et 4 et des équations de degré 8 avec les groupes mentionnés: \og{}Über einige algebraische Körper achten Grades\fg{}, Arkiv for Math., Astron. och Physik, vol. 6, no. 30. [Addition à la correction.]}

Ce qui suit est une contribution à la direction \og{}rationnelle\fg{}: on donne des conditions suffisantes pour l'existence de représentations paramétriques qui, pour un domaine arbitraire de rationalité, fournissent la totalité des équations à groupe prescrit. Ces conditions consistent en ce que le corps des invariants appartenant au groupe --- le \emph{Gattungsbereich} de Lagrange --- soit représentable rationnellement par des fonctions indépendantes du corps, c'est-à-dire possède une \og{}base minimale\fg{}; ou, autrement dit, en ce qu'il soit isomorphe au corps de toutes les fonctions rationnelles en $n$ indéterminées.

Comme on le montrera encore au \S 2, ceci est vérifié par exemple pour tous les groupes qui se présentent dans les équations des degrés 3 et 4. La construction effective et la discussion plus précise de ces équations de degrés 3 et 4 se trouvent dans la dissertation de F. Seidelmann,\footnote{\emph{Die Gesamtheit der kubischen und biquadratischen Gleichungen mit Affekt bei beliebigem Rationalitätsbereich}. Erlangen, 1916.} dont un extrait suit immédiatement la présente communication.

La question de la base minimale des domaines de genre de Lagrange a été posée --- comme cela a déjà été mentionné dans le travail \og{}Körper und Systeme rationaler Funktionen\fg{} (Math. Ann. 76) --- indépendamment du lien exposé ici, par E. Fischer dans des conversations avec moi, et a donné l'impulsion aux présentes recherches.\footnote{Comparer la partie 2 de la communication préliminaire \og{}Rationale Funktionenkörper\fg{}, Jahresb. d. d. Math. Vereinig., vol. 22, 1913.}

\subsection*{\S 1. Conditions suffisantes pour la représentation paramétrique.}

\textbf{1.} La possibilité de ramener la construction d'équations à groupe prescrit à une représentation paramétrique découle du théorème d'irréductibilité de Hilbert, qui dit ceci. Soit l'équation
\begin{equation}\tag{1}
f(x)=x^n+F_1(\lambda_1,\ldots,\lambda_r)x^{n-1}+\cdots+F_n(\lambda_1,\ldots,\lambda_r)=0
\end{equation}
où $F_1(\lambda_1,
\ldots,\lambda_r),\ldots,F_n(\lambda_1,
\ldots,\lambda_r)$ sont des fonctions rationnelles des paramètres indépendants $\lambda_1,
\ldots,\lambda_r$, à coefficients dans un corps de nombres\footnote{Un domaine général de rationalité pourrait aussi prendre la place du corps de nombres.} $\Omega$, et où cette équation a le groupe $\Gamma$ par rapport au corps $\Omega(\lambda_1,
\ldots,\lambda_r)$ de toutes les fonctions rationnelles de $\lambda_1,
\ldots,\lambda_r$ à coefficients dans $\Omega$. Alors on peut remplacer les paramètres $\lambda$ par une infinité de systèmes de nombres rationnels de telle sorte que l'équation numérique ainsi obtenue
\begin{equation}\tag{2}
g(x)=x^n+a_1x^{n-1}+\cdots+a_n=0
\end{equation}
possède précisément le groupe $\Gamma$ par rapport à $\Omega$. Une telle équation $g(x)$ sera désignée plus précisément par $g_\Gamma$.

Il s'agit maintenant de savoir si l'on peut donner une telle représentation paramétrique (1) de sorte que l'équation numérique (2) qui en résulte parcoure la totalité de tous les $g_\Gamma$ lorsque $\lambda_1,
\ldots,\lambda_r$ parcourent toutes les quantités de $\Omega$, à l'exclusion de celles qui provoquent une réduction du groupe.\footnote{Il faut aussi compter ici l'apparition d'une racine double de $g(x)=0$.} Autrement dit, le système non linéaire d'équations
\begin{equation}\tag{3}
F_1(\lambda_1,
\ldots,\lambda_r)=a_1;\quad \ldots;\quad F_n(\lambda_1,
\ldots,\lambda_r)=a_n
\end{equation}
doit, pour chaque système de valeurs $a_1,
\ldots,a_n$ correspondant à un $g_\Gamma$, posséder une solution en quantités $\lambda_1,
\ldots,\lambda_r$ de $\Omega$. Comme le système (3) est non linéaire, l'exigence que la totalité des $g_\Gamma$ soit obtenue par une représentation paramétrique doit d'emblée être restreinte. En effet, pour de tels systèmes non linéaires, il peut toujours se produire des systèmes singuliers de valeurs $a_1,
\ldots,a_n$, satisfaisant à une relation algébrique $H(a_1,
\ldots,a_n)=0$, pour lesquels aucune solution n'existe;\footnote{Cela doit être développé prochainement en général dans un travail sur l'\og{}alternative pour les systèmes non linéaires d'équations\fg{}.} la représentation paramétrique (1) échoue alors et une représentation supplémentaire devient nécessaire. Dans le cas présent, toutefois, ces valeurs singulières de $a_1,
\ldots,a_n$ peuvent être caractérisées simplement.

\textbf{2.} Pour obtenir une telle représentation paramétrique, nous exigeons en outre que les $\lambda_i$ soient identiquement, comme fonctions des racines considérées comme indéterminées, des irrationalités naturelles appartenant au groupe. Nous entendons par là ce qui suit. Si, dans (1), les paramètres $\lambda_1,
\ldots,\lambda_r$ sont remplacés par des fonctions rationnelles convenablement choisies $\varphi_1(x),\ldots,\varphi_r(x)$ des indéterminées $x_1,
\ldots,x_n$, à coefficients dans $\Omega$, alors (1) doit se transformer en
\begin{equation}\tag{4}
h(x)=x^n-\sigma_1(x)x^{n-1}+\sigma_2(x)x^{n-2}\cdots \pm \sigma_n(x),
\end{equation}
où $\sigma_1(x),\ldots,\sigma_n(x)$ sont les fonctions symétriques élémentaires de $x_1,
\ldots,x_n$; et $h(x)$ doit avoir le groupe $\Gamma$ par rapport au corps de rationalité $\Omega(\varphi_1(x),\ldots,\varphi_r(x))$ qui naît ainsi de $\Omega(\lambda_1,
\ldots,\lambda_r)$. Les fonctions $\varphi_1(x),\ldots,\varphi_r(x)$ doivent aussi être algébriquement indépendantes, à cause de l'indépendance des paramètres.

Pour préciser cette exigence, il faut clarifier la relation de $\Omega(\varphi_1(x),\ldots,\varphi_r(x))$ avec le groupe $\Gamma$. Déterminons d'abord le corps de rationalité le plus général $K$ par rapport auquel $h(x)$ devient $h_\Gamma$. Soit $\mathfrak L_\Gamma$ le corps de toutes les fonctions rationnelles, à coefficients entiers rationnels, en $x_1,
\ldots,x_n$ qui admettent $\Gamma$ --- aussi appelé le \emph{Gattungsbereich} de Lagrange appartenant à $\Gamma$, ou le corps des invariants de $\Gamma$; soit $\Omega(\Omega_\Gamma)$ le corps correspondant obtenu en adjoignant $\Omega_\Gamma$ à $\Omega$.

Ainsi $\mathfrak L_\Gamma$ est un diviseur du corps $\mathfrak R(x_1,
\ldots,x_n)$ de toutes les fonctions rationnelles à coefficients entiers rationnels en $x_1,
\ldots,x_n$. Par définition du groupe, le corps le plus général $K$ est alors donné par tout corps qui contient $\Omega_\Gamma$ comme diviseur et dont l'intersection avec $\mathfrak R(x_1,
\ldots,x_n)$ est précisément $\Omega_\Gamma$. De même, pour tout corps $K$ contenant $\Omega$, l'intersection avec $\Omega(x_1,
\ldots,x_n)$ est donnée par $\Omega(\Omega_\Gamma)$. En particulier, puisque, d'après notre exigence, $\Omega(\varphi_1(x),\ldots,\varphi_r(x))$ devient un corps $K$, l'intersection de $\Omega(\varphi_1(x),\ldots,\varphi_r(x))$ avec $\Omega(x_1,
\ldots,x_n)$ est donnée par $\Omega(\Omega_\Gamma)$; d'autre part, puisque $\Omega(\varphi_1(x),\ldots,\varphi_r(x))$ est un diviseur de $\Omega(x_1,
\ldots,x_n)$, il lui est exactement égal. L'exigence signifie donc que le \emph{Gattungsbereich} de Lagrange $\Omega(\Omega_\Gamma)$ doit naître de $\Omega$ par adjonction des fonctions algébriquement indépendantes $\varphi_1(x),\ldots,\varphi_r(x)$. Il faut alors que $r=n$, puisque $\Omega_\Gamma$ contient les $n$ fonctions élémentaires algébriquement indépendantes, tandis que plus de $n$ fonctions ne peuvent pas être algébriquement indépendantes. Nous dirons que $\varphi_1(x),\ldots,\varphi_n(x)$ forment une \emph{base minimale} de $\Omega(\Omega_\Gamma)$; ou encore que $\Omega_\Gamma$ possède une base minimale à coefficients dans $\Omega$.

\textbf{3.} Il est maintenant essentiel que toute cette considération puisse être renversée, et qu'elle conduise donc effectivement à des conditions suffisantes\footnote{L'exigence formulée en 2 a été élargie par rapport à celle de 1; les conditions n'ont donc pas besoin d'être nécessaires.} pour la représentation paramétrique désirée. Supposons donc que $\varphi_1(x),\ldots,\varphi_n(x)$ soit une base minimale de $\Omega(\mathfrak L_\Gamma)$, et soit $\Omega$ le plus petit corps contenant les coefficients de $\varphi_1(x)\ldots\varphi_n(x)$. Alors $\Omega$ est nécessairement un corps de nombres algébriques; en particulier il peut être le corps $\mathfrak R$ de tous les nombres rationnels. Comme $\sigma_1(x),\ldots,\sigma_n(x)$ appartiennent à $\Omega(\mathfrak L_\Gamma)$, il existe des identités en $x_1,
\ldots,x_n$ --- où les $F_i$ désignent des fonctions rationnelles de leurs arguments ---
\begin{equation}\tag{5}
\sigma_1(x)=F_1(\varphi_1(x),\ldots,\varphi_n(x));\quad\ldots;\quad
\sigma_n(x)=F_n(\varphi_1(x),\ldots,\varphi_n(x)).
\end{equation}
Réciproquement, $\varphi_1(x),\ldots,\varphi_n(x)$, et donc aussi $\varphi(x)=\mu_1\varphi_1(x)+\cdots+\mu_n\varphi_n(x)$, dépendent algébriquement de $\sigma_1(x),\ldots,\sigma_n(x)$ au moyen de l'équation irréductible par rapport à $\Omega(\sigma_1(x),\ldots,\sigma_n(x))$,
\begin{equation}\tag{6}
G(\varphi(x);\sigma_1(x),\ldots,\sigma_n(x))=0,
\end{equation}
qui est encore une identité en les $x$.

Par (5), cette identité (6) se transforme en
\begin{equation}\tag{7}
G\bigl(\varphi(x);F_1(\varphi_1,
\ldots,\varphi_n),\ldots,F_n(\varphi_1,
\ldots,\varphi_n)\bigr)
=H(\varphi_1(x),\ldots,\varphi_n(x))=0.
\end{equation}
En raison de l'indépendance algébrique de $\varphi_1(x),\ldots,\varphi_n(x)$, (7) est satisfaite non seulement identiquement en $x$, mais aussi en les $\varphi_i$; c'est-à-dire identiquement en des paramètres indépendants $\lambda$:
\begin{equation}\tag{8}
H(\lambda_1,
\ldots,\lambda_n)=G\bigl(\mu_1\lambda_1+\cdots+\mu_n\lambda_n;
F_1(\lambda_1,
\ldots,\lambda_n),\ldots,F_n(\lambda_1,
\ldots,\lambda_n)\bigr)=0.
\end{equation}
De l'identité (8) il résulte que l'équation
\begin{equation}\tag{9}
f(x)=x^n-F_1(\lambda_1,
\ldots,\lambda_n)x^{n-1}+F_2(\lambda_1,
\ldots,\lambda_n)x^{n-2}\ldots \pm F_n(\lambda_1,
\ldots,\lambda_n)=0
\end{equation}
a le groupe $\Gamma$ par rapport à $\Omega(\lambda_1,
\ldots,\lambda_n)$. En effet, cette identité montre que la quantité rationnellement connue $\lambda=\mu_1\lambda_1+\cdots+\mu_n\lambda_n$ est une fonction des racines appartenant à $\Gamma$; mais aucune autre fonction appartenant à un sous-groupe de $\Gamma$ ne peut être rationnellement connue, comme le montre la substitution $\lambda_i=\varphi_i(x)$, tandis qu'après la même substitution $\lambda$ est aussi différente de toutes ses conjuguées. En outre, si $\Omega^*$ est un corps de nombres quelconque contenant $\Omega$, alors $f(x)$ a le groupe $\Gamma$ par rapport à $\Omega^*(\lambda_1,
\ldots,\lambda_n)$; car, d'après 2, sous la substitution $\lambda_i=\varphi_i(x)$ le groupe n'est pas réduit même après adjonction de $\Omega^*$.

Il reste à démontrer que, au sens modifié en 1, $f(x)$ parcourt effectivement la totalité de tous les $g_\Gamma$. À cette fin, remarquons que, d'après (5), les $F_i(\lambda)$ sont des fonctions algébriquement indépendantes de leurs arguments $\lambda$. Le système d'équations
\begin{equation}\tag{10}
F_1(\lambda_1,
\ldots,\lambda_n)=-a_1;\quad F_2(\lambda_1,
\ldots,\lambda_n)=a_2;\quad\ldots;\quad F_n(\lambda_1,
\ldots,\lambda_n)=\pm a_n
\end{equation}
possède donc des solutions pour tout système arbitraire de valeurs $a_1,
\ldots,a_n$, à l'exception des valeurs singulières qu'il reste à préciser. Pour tout système de valeurs $a_1,
\ldots,a_n$ correspondant à un $g_\Gamma$ par rapport à $\Omega^*$, les $\lambda_i$ correspondants, comme irrationalités naturelles appartenant au groupe,\footnote{Comparer aussi plus bas.} deviennent des quantités de $\Omega^*$. Si donc $\lambda_1,
\ldots,\lambda_n$ parcourent tous les systèmes de valeurs, $g(x)$ parcourt la totalité des équations, au sens modifié; est contenue ici la totalité des $g_\Gamma$ auxquels correspondent des valeurs de $\Omega^*$.

Réciproquement, comme, d'après le théorème d'irréductibilité de Hilbert, des valeurs $\lambda$ de $\Omega^*$ correspondent aussi \og{}en général\fg{} à des équations $g_\Gamma$, on obtient, au sens modifié en 1, que la représentation paramétrique (9) donne effectivement la totalité de tous les $g_\Gamma$ par rapport à $\Omega^*$.

Il reste maintenant à déterminer les systèmes singuliers de valeurs pour lesquels la représentation paramétrique échoue. Soient les fonctions de la base minimale, séparées en numérateur et dénominateur, données par
\[
\varphi_i(x)=\frac{\psi_i(x)}{\chi_i(x)}.
\]
En les substituant dans les identités (5), supposons que celles-ci se transforment, sans réduction, en
\begin{equation}\tag{11}
\sigma_k(x)=\frac{G_k(x)}{H_k(x)}\quad\text{ou}\quad H_k(x)\cdot\sigma_k(x)=G_k(x),
\end{equation}
de sorte que le produit de tous les $H_k(x)$ soit divisible par le produit de tous les dénominateurs $\chi_i(x)$. Pour tous les systèmes de racines $\alpha_1,
\ldots,\alpha_n$ pour lesquels
\[
H_1(\alpha)\cdot H_2(\alpha)\cdots H_n(\alpha)\ne0
\]
et pour lesquels, par conséquent, aucun dénominateur $\chi_i(\alpha)$ ne s'annule, on peut diviser (11) par $H_k(\alpha)$ et l'on obtient
\begin{equation}\tag{12}
\sigma_k(\alpha)=\frac{G_k(\alpha)}{H_k(\alpha)}=F_k(\varphi_1(\alpha),\ldots,\varphi_n(\alpha)),
\end{equation}
où, les dénominateurs ne s'annulant pas, les $\varphi_i(\alpha)$ prennent des valeurs finies, et deviennent pour les équations $g_\Gamma$ des quantités de $\Omega^*$. (12) représente le système de solutions de (10) dès que les $\alpha_1,
\ldots,\alpha_n$ sont déterminés de sorte que $a_k=(-1)^k\sigma_k(\alpha)$.\footnote{La résolution de l'équation est ainsi ramenée à la résolution du système constitué de fonctions indépendantes:
$\varphi_1(\alpha_1,
\ldots,\alpha_n)=\lambda_1;\ldots;\varphi_n(\alpha_1,
\ldots,\alpha_n)=\lambda_n$.}

Ainsi ne peuvent devenir singuliers que les systèmes de coefficients $a_k=(-1)^k\sigma_k(\alpha)$ pour lesquels
\[
H_1(\alpha)H_2(\alpha)\cdots H_n(\alpha)=0;
\]
c'est-à-dire ceux pour lesquels (11) se transforme en la représentation $0=0$. Une étude spéciale est alors nécessaire pour décider si, parmi les $a_1,
\ldots,a_n$ singuliers ainsi définis, il s'en trouve auxquels correspondent des équations $g_\Gamma$, et pour quels corps de nombres.\footnote{Par exemple, comme F. Seidelmann l'a montré loc. cit., de telles équations --- correspondant à des systèmes singuliers de valeurs --- ne se présentent que pour le groupe alterné des équations de degrés 3 et 4, et seulement pour les corps de nombres $\Omega^*$ qui contiennent le corps des racines cubiques de l'unité. Dans chaque cas on peut donner une représentation supplémentaire simple.} En résumé on obtient le

\smallskip
\noindent\textbf{Théorème.} \emph{Si le \emph{Gattungsbereich} de Lagrange appartenant au groupe $\Gamma$ possède une base minimale\footnote{L'existence d'une base minimale entraîne aussitôt celle d'une infinité d'autres, qui peuvent être déduites de la première par des transformations rationnelles inversibles.} à coefficients dans $\Omega$, alors, pour tout corps de nombres $\Omega^*$ contenant $\Omega$, la totalité des équations à groupe prescrit $\Gamma$ par rapport à $\Omega^*$ peut être construite rationnellement par la représentation paramétrique (2), à l'exception éventuelle de celles qui sont singulières par rapport à cette représentation paramétrique et qui peuvent être données séparément. Les paramètres doivent parcourir toutes les valeurs de $\Omega^*$ qui ne provoquent pas de réduction du groupe. La construction est possible pour tout corps de nombres si le corps des coefficients $\Omega$ de la base minimale est égal au corps $\mathfrak R$ de tous les nombres rationnels.}

En outre, d'après le théorème d'irréductibilité de Hilbert, la variété des paramètres n'est pas abaissée par l'exclusion des valeurs de $\Omega^*$ qui provoquent une réduction du groupe. Donc aussi: \emph{de l'existence d'une base minimale il résulte que la variété des équations ayant le groupe $\Gamma$ par rapport à $\Omega^*$ est égale à la variété des équations sans affect; la variété n'est pas abaissée par l'affect.}

\subsection*{\S 2. Application à des équations spéciales.}

Nous démontrons maintenant l'existence de la base minimale pour tous les groupes qui se présentent dans les équations des degrés 3 et 4; d'abord, tout à fait généralement, nous montrons que, pour les équations de degré $n$, le problème peut être ramené à un problème portant sur $n-2$ indéterminées. La première réduction correspond à la transformation linéaire de Tschirnhaus qui supprime le second terme. Nous prenons comme première fonction de base $\varphi_1(x)=\sigma_1(x)$, et nous remarquons encore que $\mathfrak L_\Gamma$ peut être dérivé rationnellement de toutes les formes homogènes en $x_1,
\ldots,x_n$ qui admettent $\Gamma$. Soit $p(x_1,
\ldots,x_n)$ une telle forme; alors
\[
q(x)=p\left(x_1-\frac{\sigma_1(x)}{n},\ldots,x_n-\frac{\sigma_1(x)}{n}\right)=p\left(x_i-\frac{\sigma_1(x)}{n}\right)
\]
appartient aussi à $\mathfrak L_\Gamma$, et de plus
\[
p(x)\equiv p\left(x_i-\frac{\sigma_1(x)}{n}\right)\pmod{\sigma_1(x)},
\quad\text{ou}\quad
p(x)=q(x)+\sigma_1(x)\cdot p_1(x),
\]
avec $p_1(x)$ appartenant à $\mathfrak L_\Gamma$; cette forme est aussi de degré plus bas que $p(x)$. En poursuivant le procédé, on voit que toutes les formes homogènes de $\mathfrak L_\Gamma$ --- et par suite toutes les fonctions rationnelles de $\mathfrak L_\Gamma$ tout court --- peuvent être exprimées rationnellement au moyen de $\varphi_1(x)=\sigma_1(x)$ et de formes homogènes
\[
q(x)=p\left(x_i-\frac{\sigma_1(x)}{n}\right)=p(y_i).
\]
Entre ces combinaisons $y_i$ existe toutefois la relation linéaire $\sum y_i=0$; les $q(x)$ ne dépendent donc que de $(n-1)$ arguments.

La seconde réduction repose sur le fait que les $q(x)$ sont des formes homogènes de leurs arguments. Nous prenons comme seconde fonction de base
\[
\varphi_2(x)=\frac{\tau_3(x)}{\tau_2(x)},
\]
où
\[
\tau_3(x)=\sigma_3\left(x_i-\frac{\sigma_1(x)}{n}\right)=\sigma_3(y_i),
\qquad
\tau_2(x)=\sigma_2\left(x_i-\frac{\sigma_1(x)}{n}\right)=\sigma_2(y_i).
\]
Ainsi $\varphi_2(x)$ est une fonction fractionnaire homogène de degré un des $(n-1)$ arguments, à savoir les $y_i$ avec $\sum y_i=0$. Au moyen de $\varphi_2(x)$, toutes les $q(x)$ peuvent s'exprimer rationnellement par les fonctions, appartenant elles aussi à $\mathfrak L_\Gamma$,
\[
r(x)=\frac{q(x)}{\varphi_2(x)^\lambda}
=\frac{q(x)\cdot\tau_2(x)^\lambda}{\tau_3(x)^\lambda}
=\frac{p(y_i)\cdot\sigma_2(y_i)^\lambda}{\sigma_3(y_i)^\lambda},
\]
où $\lambda$ désigne le degré de $q(x)$. Ainsi $r(x)$ devient homogène de degré zéro, et dépend donc seulement de $(n-2)$ arguments, les rapports $y_1:y_2:\cdots:y_n$ avec $\sum y_i=0$. Par conséquent $\mathfrak L_\Gamma$ peut être exprimé rationnellement par
\[
\varphi_1(x)=\sigma_1(x),
\qquad
\varphi_2(x)=\frac{\tau_3(x)}{\tau_2(x)},
\]
et par le corps de toutes les fonctions $r(x)$ de $\mathfrak L_\Gamma$ dépendant de ces $(n-2)$ arguments,\footnote{Pour l'équation de degré $n$ sans second terme,
\[
f(x)=x^n+a_2x^{n-2}+a_3x^{n-3}+\cdots+a_n=0,
\]
cette seconde réduction correspond à la substitution, toujours possible pour $a_2\ne0$, $a_3\ne0$,
\[
y=\frac{x\cdot a_2}{a_3},
\]
par laquelle $f(x)$, à un facteur près, se transforme en
\[
g(y)=y^n+\frac{a_2^3}{a_3^2}y^{n-2}+\frac{a_2^3}{a_3^2}y^{n-3}
+\frac{a_4\cdot a_2^4}{a_3^4}y^{n-4}+\cdots+a_n\cdot\frac{a_2^n}{a_3^n}=0,
\]
c'est-à-dire en une équation ne contenant que $n-2$ paramètres:
\[
g(y)=y^n+c\,y^{n-2}+c\,y^{n-3}+c_4y^{n-4}+\cdots+c_n=0.
\]} ce qui accomplit la réduction voulue.

Pour les équations des degrés 3 et 4, cela donne en particulier une réduction à des corps de fonctions d'un ou de deux arguments respectivement. Pour ceux-ci il existe toujours une base minimale, d'après les théorèmes de Lüroth et de Castelnuovo.\footnote{J. Lüroth: \og{}Beweis eines Satzes über rationale Kurven\fg{}, Math. Ann. 9 (1875); G. Castelnuovo: \og{}Sulla razionalità delle involuzioni piane\fg{}, Math. Ann. 44 (1893). Que, pour les corps de plus de deux indéterminées, une base minimale n'ait pas nécessairement besoin d'exister a été montré par F. Enriques par un contre-exemple, Rend. Acc. Linc., vol. 21, 21 Jan. 1912.} Dans le cas d'une indéterminée, la base minimale peut être déduite par des opérations rationnelles et contient donc, en particulier pour $\Omega_\Gamma$, seulement des coefficients numériques rationnels. Avec deux indéterminées, selon Castelnuovo, la base pourrait encore avoir ses coefficients dans un corps de nombres algébrique $\Omega$; l'étude effective (cf. F. Seidelmann loc. cit.) montre qu'ici encore on obtient pour chaque $\Omega_\Gamma$ une base à coefficients rationnels. On le voit presque sans calcul déjà du fait que, pour le groupe de Klein, on peut choisir une telle base à coefficients rationnels, à savoir
\[
\varphi_1(x)=\sigma_1(x);\quad
\varphi_2(x)=\frac{vw}{u};\quad
\varphi_3(x)=\frac{uw}{v};\quad
\varphi_4(x)=\frac{u\cdot v}{w},
\]
où
\[
u=x_1+x_2-x_3-x_4;\quad v=x_1-x_2+x_3-x_4;\quad w=x_1-x_2-x_3+x_4,
\]
et du fait que le groupe alterné se transforme en le groupe cyclique des trois fonctions de base $\varphi_2,\varphi_3,\varphi_4$, tandis que chaque groupe d'ordre huit se transforme en le groupe symétrique de deux de ces fonctions, ce qui produit une réduction aux groupes des équations de degrés 3 et 2. Pour le groupe cyclique, en revanche, une base à coefficients rationnels était déjà connue par Abel, bien qu'elle ne fût pas formulée comme telle.\footnote{Voir Weber, \emph{Algebra} (petite édition), \S 93, formule (12).} Ainsi est démontré: \emph{la totalité des équations de degrés 3 et 4 à groupe prescrit --- à l'exception éventuelle de celles qui sont singulières par rapport à la représentation paramétrique --- peut être construite rationnellement par représentation paramétrique pour tout domaine arbitraire de rationalité; leur variété est égale à la variété des équations sans affect.}

L'étude effective montre (cf. F. Seidelmann loc. cit.) qu'une représentation supplémentaire n'est nécessaire que pour le groupe alterné sur les domaines de rationalité contenant le corps des racines cubiques de l'unité, tandis que dans tous les autres cas la représentation paramétrique fournit directement la totalité.

Il faut encore remarquer que la base minimale est connue pour tous les groupes abéliens.\footnote{E. Fischer: \og{}Die Isomorphie der Invariantenkörper der endlichen Abelschen Gruppen linearer Transformationen\fg{}, Gött. Nachr. 1915, et \og{}Zur Theorie der endlichen Abelschen Gruppen\fg{}, Math. Ann. 77 (1915).} Ici toutefois il s'agit d'une base à coefficients dans le corps cyclotomique déduit des caractères du groupe. Ainsi, dans ce cas, on ne peut obtenir aucun résultat nouveau pour la construction d'équations à groupe prescrit.

\medskip
\noindent Göttingen, juillet 1916.
\clearpage
\setcounter{footnote}{0}
\providecommand{\bdelta}{\bar{\delta}}
\begin{center}
{\Large\bfseries \NoetherNavAnchor{P12}{12. Invariants d’expressions différentielles arbitraires}12. Invariants d'expressions différentielles arbitraires}\par
\vspace{0.8em}
Par Emmy Noether, à Göttingen.\par
\vspace{1.5em}
\emph{Nachrichten von der Gesellschaft der Wissenschaften zu Göttingen} 1918, pp. 37--44
\end{center}

\vspace{3em}
\begin{center}
Présenté dans la séance du 25 janvier 1918 par F. Klein.
\end{center}

Par expression différentielle j'entends une fonction
\[
f(x,dx)=f(x_1,\ldots,x_n;\,dx_1,\ldots,dx_n)
\]
des \(n\) variables \(x_1\ldots x_n\) et de leurs différentielles
\(dx_1\ldots dx_n\), que l'on suppose analytique dans les deux systèmes
de \(n\) arguments chacun, mais non nécessairement réelle. Je prends
maintenant pour base le groupe de toutes les transformations analytiques
des variables, auquel correspond le groupe de toutes les transformations
linéaires des différentielles:
\begin{equation}
x_i=x_i(y_1\ldots y_n);\qquad
dx_i=\sum\frac{\partial x_i}{\partial y_k}\,dy_k;\qquad
\delta x_i=\sum\frac{\partial x_i}{\partial y_k}\,\delta y_k,
\tag{1}
\end{equation}
par lesquelles \(f(x,dx)\) doit passer en \(g(y,dy)\). Par invariant de
\(f(x,dx)\), j'entends un invariant relativement à ce groupe, et relativement
au groupe ainsi induit pour les différentielles supérieures
\(d^2x,d\delta x\ldots\) et pour les dérivées de \(f(x,dx)\); donc une fonction
\(J\) (analytique) telle que, pour \(f\) arbitraire, et en vertu de (1),
on ait identiquement, dans les variables et dans toutes les différentielles
qui interviennent:
\[
\begin{aligned}
&J\!\left(f,\frac{\partial f}{\partial dx}\cdots
\frac{\partial^{\varrho+\sigma}f}{\partial x^\varrho\partial dx^\sigma}
\cdots dx,\delta x,d^2x,\ldots\right)\\
&\qquad =
J\!\left(g,\frac{\partial g}{\partial dy}\cdots
\frac{\partial^{\varrho+\sigma}g}{\partial y^\varrho\partial dy^\sigma}
\cdots dy,\delta y,d^2y\ldots\right)\footnotemark
\end{aligned}
\]
\footnotetext{\(\displaystyle
\frac{\partial^{\rho+\sigma}f}{\partial x^\rho\partial dx^\sigma}\)
sera partout employé comme abréviation de
\[
\frac{\partial^{\varrho+\sigma}f}
{\partial x_{i_1}\ldots\partial x_{t_\varrho}
 \partial dx_{k_1}\ldots\partial dx_{k_\sigma}},
\]
où les indices désignent quelconques des nombres \(1\) à \(n\).}

On définit exactement de même l'invariant d'un système simultané
d'expressions différentielles. Si, en particulier, un tel invariant ne contient
que les premières différentielles \(dx,\delta x\), et seulement des dérivées
par rapport à ces différentielles, non par rapport aux variables, alors
l'apparition des variables dans les fonctions et la transformation linéaire
des différentielles n'entrent pas en considération; ces invariants spéciaux
sont des invariants relativement au groupe de toutes les transformations
linéaires des différentielles \(dx,\delta x\), à coefficients indéterminés,
et seront appelés invariants projectifs du système simultané.

Pour les formes différentielles homogènes quadratiques, Christoffel
(Crelle 70) et Ricci (\emph{Math. Annalen} 54) ont établi un système de
formations invariantes tel que tous les invariants deviennent des invariants
projectifs de ce système simultané; par là, les questions relatives à la
totalité des invariants et à l'équivalence se ramènent à des questions de
théorie linéaire des invariants. C'est ce que je voudrais appeler le
\og{} théorème de réduction \fg{}. Le système complet se compose d'une infinité
dénombrable de formes; mais, pour les invariants d'un ordre fini quelconque
\(\varrho\), c'est-à-dire ceux qui ne contiennent pas de dérivées par rapport
aux variables d'ordre supérieur au \(\varrho\)-ième, il ne se compose que d'un
nombre fini de formes.

Dans ce qui suit, je montre la validité du \og{} théorème de réduction \fg{} pour
des expressions différentielles arbitraires; ici encore il s'agit d'un système
complet d'une infinité dénombrable de formations invariantes, mais, pour
chaque ordre fini, seulement d'un nombre fini d'entre elles. Toutefois, alors
que Christoffel et Ricci fondent la formation des invariants et la preuve du
théorème de réduction sur des éliminations difficiles à embrasser d'un coup
d'\oe{}il, je forme les invariants directement par des procédés de différentiation
invariante et de variation, ces derniers pouvant être regardés simplement
comme des procédés de différentiation par rapport à d'autres paramètres.
L'algorithme de ces procédés de différentiation invariante fut développé,
à la suite de Lagrange (\emph{Mécanique analytique}, deuxième partie,
section IV), par Riemann (\OE{}uvres XXII) et Lipschitz (Crelle 70, 72) pour
les formes différentielles quadratiques, sans toutefois aller jusqu'au
théorème de réduction. Les moyens auxiliaires pour prouver le théorème de
réduction sont fournis par les \og{} coordonnées normales \fg{} introduites par
Riemann pour les formes quadratiques (\OE{}uvres XIII); elles transforment
en droites les extrémales, issues d'un point, du problème variationnel
associé, mais elles n'existent que pour les expressions différentielles
homogènes
\(f(x,\varkappa\cdot dx)=\Phi(\varkappa)\cdot f(x,dx)\footnote{Il se voit
facilement que \(\varPhi(\varkappa)=\varkappa^c\), où \(c\) désigne une
grandeur réelle ou complexe quelconque.}\). Le cas inhomogène se ramène
cependant à celui-ci.

\begin{center}
\textbf{I. Formation des invariants.}
\end{center}

Une suite d'invariants projectifs de \(f(x,dx)\), qui en général ne s'arrête
pas, est fournie par les polaires. À cause de la linéarité de la transformation
des différentielles, de \(f(x,dx)=g(y,dy)\), on déduit aussi que
\(f(x,dx+\lambda\delta x)=g(y,dy+\lambda\delta y)\); et de là, par
développement suivant \(\lambda\), naissent les relations caractéristiques
de l'invariance:
\begin{equation}
\begin{gathered}
f(dx)=g(dy),\\
f_\delta=\sum\frac{\partial f(dx)}{\partial dx}\,\delta x
       =\sum\frac{\partial g(dy)}{\partial dy}\,\delta y,\\
\vdots\\
f_{\delta^\sigma}=
\sum\frac{\partial^\sigma f(dx)}{\partial dx^\sigma}\,\delta x^\sigma
=
\sum\frac{\partial^\sigma g(dy)}{\partial dy^\sigma}\,\delta y^\sigma,\\
\vdots
\end{gathered}
\tag{2}
\end{equation}

Chaque polaire, par application d'un procédé de variation \(\bdelta\),
donne naissance à une suite non terminante d'invariants généraux:
\begin{equation}
\begin{gathered}
\bdelta^\varrho f(x,dx)=\bdelta^\varrho g(y,dy),\\
\vdots\\
\bdelta^\varrho f_{\delta^\sigma}
=\bdelta^\varrho\sum\frac{\partial^\sigma f(dx)}{\partial dx^\sigma}\,
\delta x^\sigma
=\bdelta^\varrho\sum\frac{\partial^\sigma g(dy)}{\partial dy^\sigma}\,
\delta y^\sigma,\\
\vdots\\[0.2em]
(\varrho=0,1,2\ldots\ldots).
\end{gathered}
\tag{3}
\end{equation}
Au moyen de (1), les relations (2) et (3) permettent de lire les lois de
transformation des dérivées de \(f(x,dx)\); cela ne sera cependant pas
nécessaire dans la suite. À partir des invariants (3), on peut maintenant
obtenir par combinaison linéaire la \og{} forme normale de la \(\varrho\)-ième
variation \fg{}, que l'on trouve pour \(\varrho=1\) chez Lagrange et pour
\(\varrho=2\) chez Riemann; elle est caractérisée par le fait que les
différentielles mixtes d'ordre \(\varrho+1\), à savoir
\(d^\varrho\delta,d^{\varrho-1}\delta^2\ldots\), sont éliminées. On obtient:
\begin{equation}
\begin{gathered}
\varOmega_1=\delta f(dx)-d f_\delta(dx),\\
\varOmega_2=\delta^2 f-\delta d f_\delta+\frac12 d^2 f_{\delta^2},\\
\vdots\\
\varOmega_\varrho=\delta^\varrho f-\delta^{\varrho-1}d f_\delta
+\frac12\delta^{\varrho-2}d^2 f_{\delta^2}
+\cdots+\frac{(-1)^\varrho}{\varrho!}\,d^\varrho f_{\delta^\varrho},\\
\vdots
\end{gathered}
\tag{4}
\end{equation}

À partir des invariants (4), en généralisant un ansatz de Riemann, on peut
maintenant parvenir à des invariants qui ne contiennent que des premières
différentielles; de tels invariants seront appelés \og{} fonctions fondamentales \fg{}.
Si, en effet, dans \(\varOmega_1\), les différentielles sont remplacées par les
quotients différentiels, alors \(\varOmega_1=0\) (identiquement en \(\delta x\))
donne exactement les équations de Lagrange du problème variationnel
appartenant à
\[
f\!\left(\frac{dx}{dt}\right).
\]
J'impose maintenant la condition invariante que la direction
\((d+\lambda\delta)\) satisfasse à ces équations de Lagrange:
\begin{equation}
\varOmega_1(d+\lambda\delta)=
\bdelta f(dx+\lambda\delta x)
-(d+\lambda\delta)f_{\bdelta}(dx+\lambda\delta x)=0
\quad(\text{identiquement en }\bdelta x),
\tag{5}
\end{equation}
d'où, par la substitution \(\bdelta=d+\lambda\delta\), il résulte encore,
pour \(f(dx)\) homogène:
\begin{equation}
(d+\lambda\delta)f(dx+\lambda\delta x)=0,
\tag{6}
\end{equation}
sauf dans le cas de l'homogénéité du premier ordre. Dans ce dernier cas,
(6) doit être ajoutée comme nouvelle condition, car alors le système
d'équations (5) ne représente que \((n-1)\) conditions indépendantes. Si,
dans (5), respectivement dans (5) et (6), on annule les coefficients des
puissances zéro, première et seconde de \(\lambda\), on obtient ainsi
trois systèmes invariants d'équations
\(\varphi=0\), \(\psi=0\), \(\chi=0\) (identiquement en \(\bdelta x\)),
au moyen desquels \(d^2x,d\delta x,\delta^2x\) peuvent être exprimés par les
premières différentielles. Les autres systèmes invariants d'équations
\begin{equation}
\begin{gathered}
d^\sigma\varphi=0,\qquad d^\sigma\psi=0,\qquad
d^\sigma\chi=0,\qquad d^{\sigma-1}\delta\chi=0,\\
d^{\sigma-2}\delta^2\chi=0\ldots\delta^\sigma\chi=0
\qquad(\sigma=0,1,2\ldots)
\end{gathered}
\tag{7}
\end{equation}
suffisent alors précisément pour exprimer toutes les différentielles
supérieures par les premières.\footnote{Seule la forme linéaire
\(\sum A_i\,dx_i\) demande un traitement particulier, car ici le système
d'équations (5) ne contient pas de secondes différentielles.} Les fonctions
(4), liées au système invariant d'équations (7), conduisent donc à des
fonctions fondamentales \([\varOmega_\varrho]\) ne contenant que des premières
différentielles; de plus \([\varOmega_1]\) s'annule identiquement.

De même, à partir de la différentielle \(\delta h(dx,\delta x)\) d'une
fonction fondamentale \(h\), on peut, au moyen de (7), obtenir de nouveau
une fonction fondamentale, qui sera appelée la \og{} dérivée covariante \fg{}
\([h^{(1)}]\) de \(h\); en général, \([h^{(\sigma)}]\) désignera la dérivée
covariante de \([h^{(\sigma-1)}]\). Ces deux procédés conduisent ainsi à
une suite doublement infinie de fonctions fondamentales:
\begin{equation}
\begin{array}{cccccc}
[\varOmega_2], & [\varOmega_2^{(1)}], & [\varOmega_2^{(2)}] & \ldots,
& [\varOmega_2^{(\nu)}] & \ldots\\
\vdots &&&&&\\
{}[\varOmega_\varrho], & [\varOmega_\varrho^{(1)}], & \ldots & \ldots,
& [\varOmega_\varrho^{(\nu)}] & \ldots\\
\vdots & \vdots &&& \vdots&
\end{array}
\tag{8}
\end{equation}

Il apparaît en outre que les membres de gauche des équations qui expriment
les différentielles supérieures par les premières sont cogrédients aux
premières différentielles, c'est-à-dire qu'ils subissent les mêmes
transformations linéaires. Si donc, à la place des différentielles supérieures,
on introduit dans l'invariant ces membres de gauche \(p,q,r\ldots\)
comme arguments, alors, abstraction faite des dérivées, il ne dépend que
de grandeurs cogrédientes entre elles. La formation des fonctions (8)
indiquée ci-dessus revient simplement à l'introduction de ces nouvelles
grandeurs; chaque invariant passe ainsi en une somme d'invariants, et de
cette somme on choisit celui qui contient précisément \(dx,\delta x\), mais
non \(p,q,r\ldots\).

Dans la suite on montre que, sous l'hypothèse d'homogénéité
\[
f(\varkappa\cdot dx)=\varPhi(\varkappa)\cdot f(dx),
\]
les fonctions fondamentales (8) épuisent le système complet cherché.

\begin{center}
\textbf{II. Coordonnées normales, équivalence et théorème de réduction.}
\end{center}

J'arrive aux coordonnées normales en intégrant le problème variationnel
appartenant à \(f\!\left(\frac{dx}{dt}\right)\):
\begin{equation}
\begin{gathered}
\varOmega_1=\delta f\!\left(\frac{dx}{dt}\right)
-\frac{d}{dt}f_\delta\!\left(\frac{dx}{dt}\right)=0
\quad(\text{identiquement en }\delta x),\\
\frac{d}{dt}f\!\left(\frac{dx}{dt}\right)=0
\quad(\text{comme conséquence ou condition additionnelle}).
\end{gathered}
\tag{9}
\end{equation}
En raison de l'homogénéité supposée de \(f(dx)\), ce système d'équations
se transforme en lui-même par la substitution de paramètre
\(t=\varkappa\cdot\tau\). Par conséquent, si, au moyen de (9), on exprime
\(\frac{d^\varrho x_i}{dt^\varrho}\) comme fonction
\(\varphi_\varrho^{(i)}\!\left(\frac{dx}{dt}\right)\) des premières dérivées,
alors \(\varphi_\varrho^{(i)}\) est homogène d'ordre \(\varrho\):
\begin{equation}
\varphi_\varrho^{(i)}\!\left(\varkappa\cdot\frac{dx}{dt}\right)
=\varkappa\varphi_\varrho^{(i)}\!\left(\frac{dx}{dt}\right).
\tag{10}
\end{equation}
Si, dans l'intégration de (9), on prescrit maintenant, pour \(t=t_0\),
les valeurs initiales
\[
x=(x)_0,\qquad \frac{dx}{dt}=\left(\frac{dx}{dt}\right)_0
\]
et si l'on pose
\begin{equation}
u_i=(t-t_0)\left(\frac{dx_i}{dt}\right)_0,
\tag{11}
\end{equation}
où, en vertu de \(f\!\left(\frac{dx}{dt}\right)=\mathrm{const.}\), le
paramètre \((t-t_0)\) devient proportionnel à la \og{} longueur de l'extrémale \fg{},
alors le développement en puissances de \((t-t_0)\), en vertu de (10),
donne
\begin{equation}
x_i-(x_i)_0=u_i+\frac12\varphi_2^{(i)}(u)+\cdots+
\frac{1}{\varrho!}\varphi_\varrho^{(i)}(u)+\cdots.
\tag{12}
\end{equation}
Ce développement peut être regardé comme une transformation de variables
\(x_i=x_i(u)\), où les grandeurs \(u\) sont précisément les coordonnées
normales de Riemann. Ainsi, par (11) et (12), ces coordonnées transforment
les extrémales passant par \((x_0)\) en droites. De plus, à une transformation
arbitraire de variables \(x=x(y)\) correspond, par (11), une transformation
linéaire \(u=u(v)\) des coordonnées normales associées, dont les coefficients
ne dépendent que du système arbitraire de valeurs \((x)_0\); les différentielles
\(du,\delta u\) sont donc soumises à la même transformation que les \(u\)
eux-mêmes.

Au moyen de (12), que \(f(x,dx)\) passe en \(F(u,du)\), ou, développé en
puissances des \(u\),
\begin{equation}
F(u,du)=F_0(u,du)+F_1(u,du)+\cdots F_\varrho(u,du)+\cdots,
\tag{13}
\end{equation}
où \(F_\varrho(u,du)\) désigne une forme homogène de dimension \(\varrho\) en
\(u\). À cause de la linéarité de la transformation \(u=u(v)\), l'identité
\(F(u,du)=G(v,dv)\) transporte également les composantes homogènes
particulières dans les composantes correspondantes:
\begin{equation}
F_\varrho(u,du)=G_\varrho(v,dv).
\tag{14}
\end{equation}
Si donc on prend pour \((x)_0\) un système fixe quelconque de valeurs
constantes, (13) et (14) montrent que l'équivalence de \(f(dx)\) avec
\(g(dy)\) équivaut à l'équivalence des systèmes de fonctions associés
\(F_\varrho(u,du)\) relativement à une transformation linéaire. En outre,
les \(F_\varrho(u,du)\), ou les \(F_\varrho(\delta u,du)\) correspondantes,
représentent des invariants de \(f(dx)\), pris au point \(x=(x)_0\);
elles forment même un système complet de fonctions fondamentales pour
ce point \(x=(x)_0\). En effet,
\[
F_\varrho(\delta u,du)=\frac1{\varrho!}\sum
\frac{\partial F_\varrho}{\partial \delta u^\varrho}\,\delta u^\varrho;
\qquad
\frac{\partial F_\varrho}{\partial \delta u^\varrho}
=\left(\frac{\partial F(du)}{\partial u^\varrho}\right)_0,
\]
et cette dernière dérivée peut, au moyen de (12), s'exprimer de la même
manière par les dérivées de \(f(dx)\), prises pour \(x=(x)_0\), que la
dérivée correspondante de \(G(dv)\) peut s'exprimer par les dérivées de
\(g(dy)\). Comme
\[
\left(\frac{\partial^{\varrho+\sigma}F(du)}
{\partial u^\varrho\partial du^\sigma}\right)_0
=
\frac{\partial^{\varrho+\sigma}F_\varrho(\delta u,du)}
{\partial\delta u^\varrho\partial du^\sigma},
\]
chaque invariant de \(f(dx)=F(du)\), pris au point \(x=(x)_0\), devient
aussi un invariant projectif des \(F(\delta u,du)\), dès que les
différentielles supérieures ont seulement été remplacées par les
\(p,q,r\ldots\) cogrédients à \(dx,\delta x\).

Mais maintenant \((x)_0\) peut être regardé comme un paramètre; à chaque
invariant au point \(x=(x)_0\) correspond un invariant de \(f(dx)\), pourvu
seulement que \((x)_0\) soit de nouveau remplacé par \(x\). Ainsi, si les
\(F_\varrho(\delta u,du)\) deviennent des invariants
\(\varPsi_\varrho(\delta x,dx)\), pris pour \(x=(x)_0\) --- car, pour \(x=(x)_0\),
\(du,\delta u\) passent en \(dx,\delta x\) ---, alors les \(\varPsi_\varrho\)
elles-mêmes forment les fonctions fondamentales d'un système complet.
Il reste seulement à exprimer ce système de fonctions fondamentales par
des invariants explicites de \(f(dx)\), à savoir par les fonctions (8).
Que cela soit possible résulte de leur formation au moyen de procédés de
différentiation. Pour les termes d'ordre le plus élevé, ces procédés se
transforment en simples procédés de polarisation; et une transformation
analogue au développement en série de Clebsch--Gordan, combinée avec une
induction à cause des termes négligés, démontre l'assertion. Si l'on a
affaire aux invariants d'un système simultané, il faut encore ajouter aux
fonctions fondamentales (8) les dérivées covariantes des autres expressions
différentielles, comme le montre l'introduction, dans ces expressions, des
coordonnées normales appartenant à \(f(dx)\).

L'équivalence de \(f(dx)\) avec \(g(dy)\), qui avait été réduite à (14),
est donc aussi équivalente à l'équivalence des \(\varPsi_\varrho\), ou tout aussi
bien des fonctions (8), relativement à la transformation linéaire des
différentielles, sans qu'aucune condition spéciale d'intégrabilité soit
requise; cela résulte de la possibilité de la réduction à (14). Ainsi le
théorème de réduction est démontré dans toutes ses parties. En particulier,
l'annulation identique de la suite de fonctions
\([\varOmega_2],[\varOmega_3]\ldots[\varOmega_\varrho]\ldots\) est nécessaire et
suffisante pour que \(f(dx)\) puisse être transformée en une expression à
coefficients constants.

Enfin, si l'on prend pour base, au lieu du groupe de toutes les transformations
analytiques, un sous-groupe, alors le groupe des transformations linéaires
correspondantes des \(u\) passe lui aussi en un sous-groupe du groupe
projectif; les invariants de \(f(dx)\) redeviennent des invariants des
fonctions (8) relativement à une transformation linéaire, mais maintenant
relativement à ce sous-groupe. Ainsi, par homogénéisation, le cas des
fonctions non homogènes \(f(dx)\) peut être ramené au groupe affine; le
système complet peut alors être déduit des fonctions (8) formées pour
une variable de plus.

Du théorème de réduction on obtient le cas particulier des formes quadratiques,
et plus généralement des formes de dimension \(p\), du fait qu'ici le système
de fonctions s'arrête à \([\varOmega_2]\), plus généralement à \([\varOmega_p]\)
et à ses dérivées covariantes, puisque toutes les fonctions fondamentales
ultérieures deviennent des invariants projectifs de celles-ci. Cela explique
la position distinguée de la forme de courbure riemannienne \([\varOmega_2]\)
pour les formes quadratiques. La question de l'équivalence des formes
quadratiques, respectivement des formes de dimension \(p\), se ramène
précisément à l'équivalence de \([\varOmega_2]\), respectivement de
\([\varOmega_2]\ldots[\varOmega_p]\), et de leurs dérivées covariantes, par
transformation linéaire; et l'annulation identique de \(]\varOmega_2]\),
respectivement de \([\varOmega_2]\ldots[\varOmega_p]\), est nécessaire et
suffisante pour la transformation en formes à coefficients constants, ce
qui, pour les formes quadratiques, énonce l'un des résultats les plus connus.

Une présentation plus détaillée doit paraître prochainement dans les
\emph{Mathematische Annalen}.
\clearpage
\setcounter{footnote}{0}
\begin{center}
{\Large\bfseries \NoetherNavAnchor{P13}{13. Problèmes variationnels invariants}13. Problèmes variationnels invariants}\par
\vspace{0.6em}
(À F. Klein pour le cinquantenaire de son doctorat.)\par
\vspace{0.6em}
Par Emmy Noether, à Göttingen.\par
\vspace{1.5em}
\emph{Nachrichten von der Gesellschaft der Wissenschaften zu Göttingen} 1918, pp. 235--257
\end{center}

\vspace{3em}
\begin{center}
Présenté par F. Klein à la séance du 26 juillet 1918.\footnote{La rédaction définitive du manuscrit ne fut remise qu'à la fin de septembre.}
\end{center}

Il s'agit de problèmes variationnels qui admettent un groupe continu (au sens de Lie); les conséquences qui en résultent pour les équations différentielles correspondantes trouvent leur expression la plus générale dans les théorèmes formulés au \S{}1 et démontrés dans les sections suivantes. Pour ces équations différentielles issues de problèmes variationnels, on peut énoncer des assertions beaucoup plus précises que pour des équations différentielles quelconques admettant un groupe, lesquelles constituent l'objet des recherches de Lie. Ce qui suit repose donc sur une combinaison des méthodes du calcul formel des variations avec celles de la théorie des groupes de Lie. Pour des groupes et des problèmes variationnels particuliers, cette combinaison de méthodes n'est pas nouvelle; je mentionne Hamel et Herglotz pour des groupes finis particuliers, Lorentz et ses élèves (par exemple Fokker), Weyl et Klein pour des groupes infinis particuliers.\footnote{Hamel: \emph{Math. Ann.} vol. 59 et \emph{Zeitschrift f. Math. u. Phys.} vol. 50. Herglotz: \emph{Ann. d. Phys.} (4) vol. 36, en particulier \S{}9, p. 511. Fokker, Verslag d. Amsterdamer Akad., 27 janvier 1917. Pour la littérature ultérieure, comparer la seconde note de Klein: \emph{Göttinger Nachrichten}, 19 juillet 1918. Dans un travail de Kneser tout juste paru (\emph{Math. Zeitschrift} vol. 2), il s'agit de la construction d'invariants par une méthode analogue.} En particulier, la seconde note de Klein et les développements présents se sont influencés mutuellement; je puis renvoyer à ce sujet aux remarques finales de la note de Klein.

\section*{\S{}1. Préliminaires et formulation des théorèmes}

Toutes les fonctions qui interviennent ci-dessous sont supposées analytiques, ou au moins continues et continûment dérivables un nombre fini de fois, et uniformes dans le domaine considéré.

Par \og{}groupe de transformations\fg{}, on entend comme d'habitude un système de transformations tel qu'à chaque transformation corresponde une transformation inverse contenue dans le système, et tel que la composition de deux transformations quelconques du système appartienne encore au système. Le groupe s'appelle un groupe continu fini \(\G_\varrho\) lorsque ses transformations sont contenues dans une transformation la plus générale dépendant analytiquement de \(\varrho\) paramètres essentiels \(\eps\) (c'est-à-dire lorsque les \(\varrho\) paramètres ne peuvent pas être représentés comme \(\varrho\) fonctions d'un nombre moindre de paramètres). De façon correspondante, par groupe continu infini \(\G_{\infty\varrho}\), on entend un groupe dont les transformations les plus générales dépendent de \(\varrho\) fonctions arbitraires essentielles \(p(x)\) et de leurs dérivées, analytiquement ou au moins continûment et avec un nombre fini de dérivées continues. Un cas intermédiaire entre les deux est le groupe dépendant d'une infinité de paramètres, mais non de fonctions arbitraires. Enfin, on appelle groupe mixte un groupe qui dépend à la fois de fonctions arbitraires et de paramètres.\footnote{Lie définit, dans \emph{Grundlagen für die Theorie der unendlichen kontinuierlichen Transformationsgruppen} (rapports de la Société royale saxonne des sciences, 1891), le groupe continu infini comme un groupe de transformations dont les transformations sont données par les solutions les plus générales d'un système d'équations aux dérivées partielles, dès lors que ces solutions ne dépendent pas seulement d'un nombre fini de paramètres. On obtient ainsi l'un des types indiqués ci-dessus, différent du groupe fini; inversement cependant, le cas limite d'une infinité de paramètres n'a pas nécessairement à satisfaire un système d'équations différentielles.}

Soient \(x_1,\ldots,x_n\) des variables indépendantes, et soient \(u_1(x),\ldots,u_\mu(x)\) des fonctions qui en dépendent. Si l'on soumet les \(x\) et les \(u\) aux transformations d'un groupe, alors, puisque les transformations sont supposées inversibles, il doit de nouveau y avoir exactement \(n\) quantités indépendantes parmi les quantités transformées; on les note \(y_1,\ldots,y_n\). Les autres quantités, dépendant de celles-ci, sont notées \(v_1(y),\ldots,v_\mu(y)\). Les transformations peuvent aussi contenir les dérivées des \(u\) par rapport aux \(x\), donc \(\frac{\p u}{\p x},\ \frac{\p^2u}{\p x^2},\ \cdots\).\footnote{Je supprime les indices autant que possible, y compris dans les sommations; ainsi, par exemple, \(\frac{\p^2u}{\p x^2}\) désigne \(\frac{\p^2u_\alpha}{\p x_\beta\p x_\gamma}\), et ainsi de suite.} Une fonction s'appelle un invariant du groupe lorsqu'il existe une relation:
\[
P\!\left(x,u,\frac{\p u}{\p x},\frac{\p^2u}{\p x^2},\cdots\right)
 =
P\!\left(y,v,\frac{\p v}{\p y},\frac{\p^2v}{\p y^2},\cdots\right).
\]
En particulier, une intégrale \(I\) est donc un invariant du groupe lorsqu'il existe une relation:
\begin{equation}
\begin{aligned}
I&=\int\!\cdots\!\int f\!\left(x,u,\frac{\p u}{\p x},\frac{\p^2u}{\p x^2},\cdots\right)d x\\
 &=\int\!\cdots\!\int f\!\left(y,v,\frac{\p v}{\p y},\frac{\p^2v}{\p y^2},\cdots\right)d y\footnotemark
\end{aligned}
\tag{1}
\end{equation}
\footnotetext{J'écris brièvement \(d x,d y\) pour \(d x_1\ldots d x_n,d y_1\ldots d y_n\).}
où l'intégration s'étend à un domaine réel arbitraire en \(x\) et au domaine correspondant en \(y\).\footnote{Tous les arguments \(x,u,\eps,p(x)\) intervenant dans les transformations sont à prendre réels, tandis que les coefficients peuvent être complexes. Mais comme les résultats finaux sont des identités en les \(x,u\), les paramètres et les fonctions arbitraires, ils valent aussi pour des valeurs complexes dès que toutes les fonctions intervenantes sont supposées analytiques. En outre, une grande partie des résultats peut être fondée sans intégrales, de sorte que la restriction au réel n'est pas nécessaire même pour la justification. En revanche, les considérations à la fin du \S{}2 et au début du \S{}5 ne semblent pas pouvoir être conduites sans intégrales.}

D'autre part, pour une intégrale arbitraire \(I\), pas nécessairement invariante, je forme la première variation \(\delta I\), et je la transforme par intégration par parties suivant les règles du calcul des variations. Comme on sait, si \(\delta u\), ainsi que toutes ses dérivées qui interviennent au bord, est prise nulle sur le bord, mais reste sinon arbitraire, on obtient:
\begin{equation}
\delta I=\int\!\cdots\!\int \delta f\,d x
=\int\!\cdots\!\int\left(\sum\psi_i\!\left(x,u,\frac{\p u}{\p x},\cdots\right)\delta u_i\right)d x,
\tag{2}
\end{equation}
où les \(\psi\) désignent les expressions de Lagrange, c'est-à-dire les membres de gauche des équations de Lagrange du problème variationnel correspondant \(\delta I=0\). Cette relation intégrale correspond à une identité sans intégrales en \(\delta u\) et ses dérivées, qui apparaît lorsqu'on écrit aussi les termes de bord. Comme le montre l'intégration par parties, ces termes de bord sont des intégrales sur des divergences, c'est-à-dire sur des expressions
\[
\Div A=\frac{\p A_1}{\p x_1}+\cdots+\frac{\p A_n}{\p x_n},
\]
où \(A\) est linéaire en \(\delta u\) et ses dérivées. On obtient donc:
\begin{equation}
\sum\psi_i\,\delta u_i=\delta f+\Div A.
\tag{3}
\end{equation}
En particulier, si \(f\) ne contient que les premières dérivées des \(u\), alors, dans le cas de l'intégrale simple, l'identité (3) est identique à ce que von Heun appelle l'\og{}équation centrale lagrangienne\fg{}:
\begin{equation}
\sum \psi_i\,\delta u_i
=\delta f-\frac{d}{d x}\left(\sum\frac{\p f}{\p u_i'}\,\delta u_i\right),
\qquad \left(u_i'=\frac{d u_i}{d x}\right),
\tag{4}
\end{equation}
tandis que, pour l'intégrale \(n\)-uple, (3) devient:
\begin{equation}
\sum\psi_i\,\delta u_i
=
\delta f
-\frac{\p}{\p x_1}\left(\sum\frac{\p f}{\p \left(\frac{\p u_i}{\p x_1}\right)}\,\delta u_i\right)
-\cdots
-\frac{\p}{\p x_n}\left(\sum\frac{\p f}{\p \left(\frac{\p u_i}{\p x_n}\right)}\,\delta u_i\right).
\tag{5}
\end{equation}
Pour l'intégrale simple et \(\varkappa\) dérivées des \(u\), (3) est donné par:
\begin{equation}
\begin{aligned}
\sum\psi_i\,\delta u_i
={}&\delta f
-\frac{d}{d x}\Bigg\{\sum\left(
\binom{1}{1}\frac{\p f}{\p u_i^{(1)}}\delta u_i
+\binom{2}{1}\frac{\p f}{\p u_i^{(2)}}\delta u_i^{(1)}
+\cdots+
\binom{\varkappa}{1}\frac{\p f}{\p u_i^{(\varkappa)}}\delta u_i^{(\varkappa-1)}
\right]\Bigg\}\\
&+\frac{d^2}{d x^2}\Bigg\{\sum\left(
\binom{2}{2}\frac{\p f}{\p u_i^{(2)}}\delta u_i
+\binom{3}{2}\frac{\p f}{\p u_i^{(3)}}\delta u_i^{(1)}
+\cdots+
\binom{\varkappa}{2}\frac{\p f}{\p u_i^{(\varkappa)}}\delta u_i^{(\varkappa-2)}
\right]\Bigg\}\\
&+\cdots
+(-1)^\varkappa\frac{d^\varkappa}{d x^\varkappa}
\left\{\sum\binom{\varkappa}{\varkappa}
\frac{\p f}{\p u_i^{(\varkappa)}}\delta u_i\right\}.
\end{aligned}
\tag{6}
\end{equation}
Une identité correspondante vaut pour l'intégrale \(n\)-uple; en particulier \(A\) contient \(\delta u\) jusqu'à la dérivée d'ordre \((\varkappa-1)\). Que (4), (5) et (6) définissent effectivement les expressions de Lagrange \(\psi_i\) résulte du fait que, dans les combinaisons des membres de droite, toutes les dérivées supérieures de \(\delta u\) sont éliminées, tandis que la relation (2), à laquelle conduit l'intégration par parties, est vérifiée et qu'elle est obtenue de façon unique.

Dans la suite il s'agit des deux théorèmes suivants:

I. Si l'intégrale \(I\) est invariante par rapport à un \(\G_\varrho\), alors \(\varrho\) combinaisons linéairement indépendantes des expressions de Lagrange deviennent des divergences; réciproquement, il s'ensuit l'invariance de \(I\) par rapport à un \(\G_\varrho\). Le théorème vaut aussi dans le cas limite d'une infinité de paramètres.

II. Si l'intégrale \(I\) est invariante par rapport à un \(\G_{\infty\varrho}\) dans lequel les fonctions arbitraires interviennent jusqu'à la dérivée d'ordre \(\sigma\), alors il existe \(\varrho\) relations identiques entre les expressions de Lagrange et leurs dérivées jusqu'à l'ordre \(\sigma\); ici aussi la réciproque est valable.\footnote{Pour certains cas exceptionnels triviaux, comparer le \S{}2, seconde note.}

Pour les groupes mixtes, les assertions des deux théorèmes s'appliquent; il apparaît donc à la fois des dépendances et des relations de divergence indépendantes d'elles.

Si l'on passe de ces identités au problème variationnel correspondant, et si l'on pose donc \(\psi=0\),\footnote{Un peu plus généralement, on peut aussi poser \(\psi_i=T_i\); comparer le \S{}3, première note.} le théorème I dit, dans le cas unidimensionnel --- où la divergence devient une différentielle totale --- qu'il existe \(\varrho\) premières intégrales, bien que des dépendances non linéaires puissent se produire entre elles;\footnote{Comparer la fin du \S{}3.} dans le cas multidimensionnel on obtient les équations de divergence que l'on a récemment souvent appelées \og{}lois de conservation\fg{}. Le théorème II dit que \(\varrho\) des équations de Lagrange sont des conséquences des autres.

L'exemple le plus simple du théorème II --- sans la réciproque --- est donné par la représentation paramétrique de Weierstrass; ici, sous l'homogénéité du premier ordre, on sait que l'intégrale est invariante lorsqu'on remplace la variable indépendante \(x\) par une fonction arbitraire de \(x\), tout en laissant \(u\) inchangé \((y=p(x);\ v_i(y)=u_i(x))\). Ainsi une fonction arbitraire intervient, mais sans dérivées; et à cela correspond la relation linéaire bien connue entre les expressions de Lagrange elles-mêmes:
\[
\sum_i \psi_i\,\frac{\dd u_i}{\dd x}=0.
\]
Un autre exemple est fourni par la \og{}théorie générale de la relativité\fg{} des physiciens. On y considère le groupe de toutes les transformations des \(x\), \(y_i=p_i(x)\), tandis que les \(u\) (notés \(g_{\mu\nu}\) et \(q\)) sont soumis aux transformations induites pour les coefficients d'une forme différentielle quadratique et d'une forme différentielle linéaire, transformations qui contiennent les premières dérivées des fonctions arbitraires \(p(x)\). À cela correspondent les \(n\) dépendances connues entre les expressions de Lagrange et leurs premières dérivées.\footnote{Comparer, par exemple, l'exposé de Klein.}

Si, en particulier, on spécialise le groupe en n'admettant aucune dérivée de \(u(x)\) dans les transformations, et en faisant en outre dépendre les quantités indépendantes transformées seulement des \(x\) et non des \(u\), alors, comme on le montrera au \S{}5, l'invariance de \(I\) entraîne l'invariance relative de \(\sum\psi_i\delta u_i\),\footnote{C'est-à-dire que \(\sum\psi_i\delta u_i\) acquiert un facteur sous les transformations.} ainsi que celle des divergences intervenant dans le théorème I, dès que les paramètres sont soumis à des transformations appropriées. Il s'ensuit en outre que les premières intégrales mentionnées ci-dessus admettent elles aussi le groupe. Pour le théorème II on obtient de même l'invariance relative des membres de gauche des dépendances, combinés au moyen des fonctions arbitraires; et, comme conséquence, encore une fonction dont la divergence s'annule identiquement et qui admet le groupe --- la fonction qui, dans la théorie des physiciens de la relativité, assure le lien entre dépendances et théorème de l'énergie.\footnote{Comparer la seconde note de Klein.} Enfin, sous forme de théorie des groupes, le théorème II donne la preuve d'une assertion apparentée de Hilbert concernant la défaillance de théorèmes propres de l'énergie en \og{}relativité générale\fg{}. Avec ces remarques complémentaires, le théorème I contient tous les théorèmes sur les premières intégrales connus en mécanique et ainsi de suite, tandis que le théorème II peut être décrit comme la généralisation, en théorie des groupes, la plus grande possible de la \og{}théorie générale de la relativité\fg{}.

\section*{\S{}2. Relations de divergence et dépendances}

Soit \(\G\) un groupe continu, fini ou infini. On peut toujours faire en sorte que la transformation identique corresponde à la valeur zéro des paramètres \(\eps\), respectivement des fonctions arbitraires \(p(x)\).\footnote{Comparer par exemple Lie, \emph{Grundlagen}, p. 331. Si des fonctions arbitraires interviennent, les valeurs particulières \(a^\sigma\) des paramètres doivent être remplacées par des fonctions fixes \(p^\sigma,\ \frac{\p p^\sigma}{\p x},\ \cdots\); et de façon correspondante les valeurs \(a^\sigma+\eps\) par \(p^\sigma+p(x),\ \frac{\p p^\sigma}{\p x}+\frac{\p p}{\p x}\), et ainsi de suite.} La transformation la plus générale aura donc la forme
\[
y_i=A_i\!\left(x,u,\frac{\p u}{\p x},\cdots\right)=x_i+\varDelta x_i+\cdots,\qquad
v_i(y)=B_i\!\left(x,u,\frac{\p u}{\p x},\cdots\right)=u_i+\varDelta u_i+\cdots,
\]
où \(\varDelta x_i,\varDelta u_i\) désignent les termes de plus basse dimension en \(\eps\), respectivement en \(p(x)\) et ses dérivées; ils seront supposés linéaires en ces quantités. On montrera plus loin que cela ne restreint pas la généralité.

Supposons maintenant que l'intégrale \(I\) soit invariante par rapport à \(\G\), de sorte que la relation (1) soit vérifiée. En particulier, \(I\) est alors aussi invariante par rapport à la transformation infinitésimale contenue dans \(\G\):
\[
y_i=x_i+\varDelta x_i;\qquad v_i(y)=u_i+\varDelta u_i.
\]
Pour celle-ci, la relation (1) devient:
\begin{equation}
\begin{aligned}
0=\varDelta I
&=\int\!\cdots\!\int f\!\left(y,v(y),\frac{\p v}{\p y},\cdots\right)d y\\
&\qquad-\int\!\cdots\!\int f\!\left(x,u(x),\frac{\p u}{\p x},\cdots\right)d x.
\end{aligned}
\tag{7}
\end{equation}
La première intégrale doit s'étendre au domaine \(x+\varDelta x\) correspondant au domaine \(x\). Mais cette intégration peut aussi être transformée en une intégration sur le domaine \(x\) au moyen de la transformation valable pour des \(\varDelta x\) infinitésimaux:
\begin{equation}
\begin{aligned}
\int\!\cdots\!\int f\!\left(y,v(y),\frac{\p v}{\p y},\cdots\right)d y
&=\int\!\cdots\!\int f\!\left(x,v(x),\frac{\p v}{\p x},\cdots\right)d x\\
&\quad+\int\!\cdots\!\int \Div(f\cdot \varDelta x)\cdot d x.
\end{aligned}
\tag{8}
\end{equation}
Si l'on introduit donc, à la place de la transformation infinitésimale \(\varDelta u\), la variation
\begin{equation}
\bd u_i=v_i(x)-u_i(x)=\varDelta u_i-\sum\frac{\p u_i}{\p x_\lambda}\varDelta x_\lambda,
\tag{9}
\end{equation}
alors (7) et (8) deviennent:
\begin{equation}
0=\int\!\cdots\!\int\{\bd f+\Div(f\cdot\varDelta x)\}d x.
\tag{10}
\end{equation}
Comme la relation (10) est remplie lorsqu'on intègre sur tout domaine arbitraire, l'intégrande doit s'annuler identiquement. Les équations différentielles de Lie pour l'invariance de \(I\) passent ainsi à la relation
\begin{equation}
\bd f+\Div(f\cdot\varDelta x)=0.
\tag{11}
\end{equation}
Si l'on exprime ici \(\bd f\), d'après (3), au moyen des expressions de Lagrange, on obtient
\begin{equation}
\sum\psi_i\bd u_i=\Div B\qquad (B=A-f\cdot\varDelta x),
\tag{12}
\end{equation}
et cette relation représente donc, pour toute intégrale invariante \(I\), une identité en tous les arguments qui y interviennent; c'est la forme cherchée des équations différentielles de Lie pour \(I\).\footnote{(12) devient \(0=0\) dans le cas trivial --- qui ne peut se produire que si \(\varDelta x,\varDelta u\) dépendent aussi de dérivées de \(u\) --- où \(\Div(f\cdot\varDelta x)=0,\ \bd u=0\). Ces transformations infinitésimales doivent donc toujours être séparées des groupes, et dans la formulation des théorèmes seul le nombre des paramètres ou fonctions arbitraires restants doit être compté. Il faut laisser ouverte la question de savoir si les transformations infinitésimales restantes forment encore toujours un groupe.}

Supposons d'abord que \(\G\) soit un groupe continu fini \(\G_\varrho\). Puisque, par hypothèse, \(\varDelta u\) et \(\varDelta x\) sont linéaires en les paramètres \(\eps_1,\ldots,\eps_\varrho\), il en va de même, d'après (9), pour \(\bd u\) et ses dérivées; par conséquent \(A\) et \(B\) sont linéaires en les \(\eps\). Je pose donc
\[
B=B^{(1)}\eps_1+\cdots+B^{(\varrho)}\eps_\varrho;\qquad
\bd u=\bd u^{(1)}\eps_1+\cdots+\bd u^{(\varrho)}\eps_\varrho,
\]
où \(\bd u^{(1)},\ldots\) sont des fonctions de \(x,u,\frac{\p u}{\p x},\cdots\). De (12) résultent les relations de divergence cherchées:
\begin{equation}
\sum\psi_i\bd u_i^{(1)}=\Div B^{(1)};\ldots\qquad
\sum\psi_i\bd u_i^{(\varrho)}=\Div B^{(\varrho)}.
\tag{13}
\end{equation}
Ainsi \(\varrho\) combinaisons linéairement indépendantes des expressions de Lagrange deviennent des divergences. L'indépendance linéaire résulte du fait que, par (9), \(\bd u=0,\ \varDelta x=0\) impliquerait \(\varDelta u=0,\ \varDelta x=0\), et donc une dépendance entre les transformations infinitésimales. Or, par hypothèse, une telle dépendance n'existe pour aucune valeur des paramètres; sinon le \(\G_\varrho\) obtenu de nouveau par intégration des transformations infinitésimales dépendrait de moins de \(\varrho\) paramètres essentiels. La possibilité supplémentaire \(\bd u=0,\ \Div(f\varDelta x)=0\) a été exclue. Ces conclusions restent valables dans le cas limite d'une infinité de paramètres.

Soit maintenant \(\G\) un groupe continu infini \(\G_{\infty\varrho}\). Alors \(\bd u\) et ses dérivées, donc aussi \(B\), sont de nouveau linéaires en les fonctions arbitraires \(p(x)\) et leurs dérivées.\footnote{Que ce ne soit pas une restriction de supposer les \(p\) indépendantes de \(u,\frac{\p u}{\p x},\cdots\), cela est montré par la réciproque.} La substitution des valeurs de \(\bd u\), encore indépendamment de (12), donne:
\[
\sum\psi_i\bd u_i
=\sum_{\lambda,i}\psi_i\left\{
a_i^{(\lambda)}(x,u,\ldots)p^{(\lambda)}(x)
+b_i^{(\lambda)}(x,u,\ldots)\frac{\p p^{(\lambda)}}{\p x}
+\cdots+c_i^{(\lambda)}(x,u,\ldots)
\frac{\p^\sigma p^{(\lambda)}}{\p x^\sigma}\right\}.
\]
Or, de façon analogue à la formule d'intégration par parties, au moyen de l'identité
\[
\varphi(x,u,\ldots)\frac{\p^\tau p(x)}{\p x^\tau}
=(-1)^\tau\cdot\frac{\p^\tau\varphi}{\p x^\tau}\cdot p(x)
\quad \bmod \text{ divergences},
\]
les dérivées de \(p\) peuvent être remplacées par \(p\) lui-même et par des divergences linéaires en \(p\) et ses dérivées. Ainsi:
\begin{equation}
\begin{aligned}
\sum\psi_i\delta u_i
={}&\sum_\lambda\left\{
(a_i^{(\lambda)}\psi_i)-\frac{\p}{\p x}\bigl(b_i^{(\lambda)}\psi_i\bigr)
+\cdots+(-1^\sigma)
\frac{\p^\sigma}{\p x^\sigma}\bigl(c_i^{(\lambda)}\psi_i\bigr)
\right\}p^{(\lambda)}+\Div\varGamma,
\end{aligned}
\tag{14}
\end{equation}
avec sommation sur \(i\) dans l'accolade. En combinaison avec (12), cela donne:
\begin{equation}
\sum\left\{
(a_i^{(\lambda)}\psi_i)-\frac{\p}{\p x}\bigl(b_i^{(\lambda)}\psi_i\bigr)
+\cdots+(-1)^\sigma
\frac{\p^\sigma}{\p x^\sigma}\bigl(c_i^{(\lambda)}\psi_i\bigr)
\right\}p^{(\lambda)}
=\Div(B-\varGamma).
\tag{15}
\end{equation}
Je forme maintenant l'intégrale \(n\)-uple de (15), étendue à un domaine arbitraire, et je choisis les \(p(x)\) de telle façon qu'elles s'annulent au bord avec toutes les dérivées intervenant dans \((B-\varGamma)\). Comme l'intégrale d'une divergence se réduit à une intégrale de bord, l'intégrale du membre de gauche de (15) s'annule elle aussi pour des \(p(x)\) arbitraires s'annulant au bord avec un nombre suffisant de dérivées. Des raisonnements usuels il résulte que l'intégrande s'annule pour tout \(p(x)\), ce qui donne les \(\varrho\) relations
\begin{equation}
\sum\left\{
(a_i^{(\lambda)}\psi_i)-\frac{\p}{\p x}\bigl(b_i^{(\lambda)}\psi_i\bigr)
+\cdots+(-1)^\sigma
\frac{\p^\sigma}{\p x^\sigma}\bigl(c_i^{(\lambda)}\psi_i\bigr)
\right\}=0\qquad(\lambda=1,2\ldots\varrho).
\tag{16}
\end{equation}
Ce sont les dépendances cherchées entre les expressions de Lagrange et leurs dérivées lorsque \(I\) est invariante par rapport à \(\G_{\infty\varrho}\). L'indépendance linéaire se démontre comme ci-dessus, puisque la réciproque reconduit à (12); et puisque l'on peut de nouveau conclure des transformations infinitésimales aux transformations finies, comme cela est exécuté plus complètement au \S{}4. Par conséquent, pour un \(\G_{\infty\varrho}\), \(\varrho\) transformations arbitraires apparaissent déjà dans les transformations infinitésimales. De (15) et (16) il résulte en outre que \(\Div(B-\varGamma)=0\).

Si, conformément à un \og{}groupe mixte\fg{}, on prend \(\varDelta x\) et \(\varDelta u\) linéaires en les \(\eps\) et en les \(p(x)\), alors, en posant une fois les \(p(x)\) et une autre fois les \(\eps\) égaux à zéro, on voit qu'il existe à la fois les relations de divergence (13) et les dépendances (16).

\section*{\S{}3. La réciproque dans le cas du groupe fini}

Pour démontrer la réciproque, il faut d'abord parcourir essentiellement les considérations précédentes en sens inverse. De la validité de (13), en multipliant par les \(\eps\) et en additionnant, on obtient la validité de (12); et, au moyen de l'identité (3), on obtient une relation
\[
\bd f+\Div(A-B)=0.
\]
Ainsi, si l'on pose \(\varDelta x=\frac{1}{f}\cdot(A-B)\), on arrive par là à (11); par intégration, cela donne finalement (7), \(\varDelta I=0\), donc l'invariance de \(I\) par rapport à la transformation infinitésimale déterminée par \(\varDelta x,\varDelta u\), où les \(\varDelta u\) sont déterminés par (9) à partir de \(\varDelta x\) et de \(\bd u\), et où \(\varDelta x,\varDelta u\) deviennent linéaires en les paramètres. Mais \(\varDelta I=0\) implique, de la manière connue, l'invariance de \(I\) par rapport aux transformations finies issues de l'intégration du système simultané:
\begin{equation}
\frac{d x}{d t}=\varDelta x_i;\qquad
\frac{d u_i}{d t}=\varDelta u_i;\qquad
\left(\begin{array}{l}
x_i=y,\\
u_i=v_i
\end{array}\right\}
\quad\text{pour }t=0.
\tag{17}
\end{equation}
Ces transformations finies contiennent \(\varrho\) paramètres \(a_1\ldots a_\partial\), à savoir les combinaisons \(t\eps_1,\ldots,t\eps_\varrho\). De l'hypothèse qu'il doit y avoir \(\varrho\), et seulement \(\varrho\), relations de divergence linéairement indépendantes (13), il résulte en outre que les transformations finies, dès qu'elles ne contiennent pas les dérivées \(\frac{\p u}{\p x}\), forment toujours un groupe. Dans le cas contraire, au moins une transformation infinitésimale issue du processus de crochets de Lie ne serait pas une combinaison linéaire des \(\varrho\) restantes; et comme \(I\) admet aussi cette transformation, il y aurait plus de \(\varrho\) relations de divergence linéairement indépendantes. Ou bien cette transformation infinitésimale serait de la forme spéciale \(\bd u=0,\ \Div(f\cdot\varDelta x)=0\). Mais alors \(\varDelta x\) ou \(\varDelta u\) dépendrait de dérivées, contrairement à l'hypothèse. Il faut laisser ouverte la question de savoir si ce cas peut se produire lorsque des dérivées apparaissent dans \(\varDelta x\) ou \(\varDelta u\); il faut alors encore ajouter à la \(\varDelta x\) déterminée ci-dessus toutes les fonctions \(\varDelta x\) pour lesquelles \(\Div(f\cdot\varDelta x)=0\), afin de rétablir la propriété de groupe, mais, par convention, les paramètres ainsi ajoutés ne doivent pas être comptés. La réciproque est ainsi démontrée.

Il résulte aussi de cette réciproque que \(\varDelta x\) et \(\varDelta u\) peuvent effectivement être supposés linéaires en les paramètres. Si \(\varDelta u\) et \(\varDelta x\) étaient des formes de degré supérieur en \(\eps\), alors, en raison de l'indépendance linéaire des produits de puissances des \(\eps\), il en résulterait des relations entièrement correspondantes (13), mais en nombre plus grand; par la réciproque, celles-ci impliquent l'invariance de \(I\) par rapport à un groupe dont les transformations infinitésimales contiennent les paramètres linéairement. Si ce groupe doit contenir exactement \(\varrho\) paramètres, il doit donc exister des dépendances linéaires entre les relations de divergence obtenues à l'origine à partir des termes de degré supérieur en \(\eps\).

Il faut encore remarquer que, dans le cas où \(\varDelta x\) et \(\varDelta u\) contiennent aussi des dérivées des \(u\), les transformations finies peuvent dépendre d'une infinité de dérivées des \(u\); car dans ce cas l'intégration de (17), lorsqu'on détermine
\[
\frac{\dd^2x_i}{\dd t^2},\quad \frac{\dd^2u_i}{\dd t^2},
\]
conduit à
\[
\D\!\left(\frac{\p u}{\p x_\varkappa}\right)
=\frac{\p\,\D u}{\p x_\varkappa}
-\sum_\lambda\frac{\p u}{\p x_\lambda}
\frac{\p\,\D x_\lambda}{\p x_\varkappa},
\]
de sorte qu'en général le nombre de dérivées de \(u\) augmente à chaque pas. Un exemple est:
\[
f=\frac12 u'^2,\qquad
\psi=-u'',\qquad
\psi\cdot x=\frac{d}{d x}(u-u'x);\qquad
\bd u=x\cdot\eps,
\]
\[
\varDelta x=\frac{-2u}{u'^2}\eps;\qquad
\varDelta u=\left(x-\frac{2u}{u'}\right)\cdot\eps.
\]
Comme les expressions de Lagrange d'une divergence s'annulent identiquement, la réciproque montre finalement ceci: si \(I\) admet un \(\G_\varrho\), alors toute intégrale qui diffère de \(I\) seulement par une intégrale de bord, c'est-à-dire par une intégrale sur une divergence, admet elle aussi un \(\G_\varrho\) ayant les mêmes \(\bd u\), tandis que ses transformations infinitésimales contiendront en général des dérivées des \(u\). Ainsi, conformément à l'exemple précédent,
\[
f^*=\frac12\left\{u'^2-\frac{\dd}{\dd x}\left(\frac{u^2}{x}\right)\right\}
\]
admet la transformation infinitésimale \(\varDelta u=x\eps,\ \varDelta x=0\); tandis que des dérivées de \(u\) interviennent dans les transformations infinitésimales correspondant à \(f\).

Si l'on passe au problème variationnel, c'est-à-dire si l'on pose \(\psi_i=0\),\footnote{\(\psi_i=0\), ou un peu plus généralement \(\psi_i=T_i\), où les \(T_i\) sont des fonctions nouvellement adjointes, s'appellent en physique des \og{}équations de champ\fg{}. Dans le cas \(\psi_i=T_i\), les identités (13) passent aux équations \(\Div B^{(\lambda)}=\sum T_i\delta u_i^{(\lambda)}\), qui en physique sont encore appelées lois de conservation.} alors (13) devient les équations
\[
\Div B^{(1)}=0,\ldots,\qquad \Div B^{(\rho)}=0,
\]
qui sont souvent appelées \og{}lois de conservation\fg{}. Dans le cas unidimensionnel cela donne:
\[
B^{(1)}=\mathrm{const.},\ldots,\qquad B^{(\rho)}=\mathrm{const.}.
\]
Ici les \(B\) contiennent au plus les dérivées d'ordre \((2\varkappa-1)\) de \(u\) (d'après (6)), dès que \(\varDelta u\) et \(\varDelta x\) ne contiennent pas de dérivées d'ordre supérieur aux dérivées d'ordre \(\varkappa\) intervenant dans \(f\). Comme \(\psi\) contient en général des dérivées d'ordre \(2\varkappa\),\footnote{Dès que \(f\) n'est pas linéaire dans les dérivées d'ordre \(\varkappa\).} on a ainsi l'existence de \(\varrho\) premières intégrales. Que des dépendances non linéaires puissent exister entre elles est de nouveau montré par le \(f\) ci-dessus. Aux transformations linéairement indépendantes \(\varDelta u=\eps_1,\ \varDelta x=\eps_2\) correspondent les relations linéairement indépendantes
\[
u''=\frac{\dd}{\dd x}u',\qquad
u''\cdot u'=\frac12\frac{\dd}{\dd x}(u')^2,
\]
tandis qu'entre les premières intégrales \(u'=\mathrm{const.},\ u'^2=\mathrm{const.}\) il existe une dépendance non linéaire. C'est le cas élémentaire où \(\varDelta u,\varDelta x\) ne contiennent pas de dérivées de \(u\).\footnote{Sinon on a en outre \(u'^\lambda=\mathrm{const.}\) pour tout \(\lambda\), correspondant à
\[
u''\cdot (u')^{\lambda-1}=\frac1\lambda\frac{d}{d x}(u')^\lambda.
\]}
\section*{\S{}4. Réciproque dans le cas du groupe infini}

Il faut d'abord montrer que l'hypothèse de linéarité de \(\D x\) et de \(\D u\) ne constitue pas une restriction. Ici cela résulte, même sans la réciproque, du fait que \(\G_{\infty\rho}\) dépend formellement de \(\rho\), et seulement de \(\rho\), fonctions arbitraires. En effet, dans le cas non linéaire, lors de la composition des transformations, les termes de plus bas ordre s'additionnent et le nombre des fonctions arbitraires augmenterait. Supposons par exemple que
\[
\begin{aligned}
y&=A\!\left(x,u,\frac{\p u}{\p x},\cdots;p\right)\\
&=x+\sum a(x,u,\ldots)p^\nu
+b(x,u,\ldots)p^{\nu-1}\frac{\p p}{\p x}
+c\,p^{\nu-2}\left(\frac{\p p}{\p x}\right)^2+\cdots
+d\left(\frac{\p p}{\p x}\right)^\nu+\cdots,
\end{aligned}
\]
où
\[
p^\nu=(p^{(1)})^{\nu_1}\cdots(p^{(\rho)})^{\nu_\rho},
\]
et, de manière correspondante,
\[
v=B\!\left(x,u,\frac{\p u}{\p x},\ldots;p\right).
\]
Alors, par composition avec
\[
z=A\!\left(y,v,\frac{\p v}{\p y},\ldots;q\right),
\]
les termes de plus bas ordre donnent
\[
z=x+\sum a(p^\nu+q^\nu)+
b\left\{p^{\nu-1}\frac{\p p}{\p x}+q^{\nu-1}\frac{\p q}{\p x}\right\}
+c\left\{p^{\nu-2}\left(\frac{\p p}{\p x}\right)^2
+q^{\nu-2}\left(\frac{\p q}{\p x}\right)^2\right\}+\cdots.
\]
Si, ici, un coefficient différent de \(a\) et de \(b\) est non nul, de sorte qu'un terme
\[
p^{\nu-\sigma}\left(\frac{\p p}{\p x}\right)^\sigma
+q^{\nu-\sigma}\left(\frac{\p q}{\p x}\right)^\sigma
\]
apparaisse effectivement pour \(\sigma>1\), celui-ci ne peut pas s'écrire comme quotient différentiel d'une seule fonction, ni comme produit de puissances d'un tel quotient. Le nombre des fonctions arbitraires a donc augmenté, contrairement à l'hypothèse. Si tous les coefficients autres que \(a\) et \(b\) s'annulent, alors, suivant les valeurs des exposants \(\nu_1\ldots\nu_\varrho\), le second terme est le quotient différentiel du premier (comme par exemple toujours pour une \(\G_{\infty1}\)), de sorte que la linéarité se présente effectivement; ou bien le nombre des fonctions arbitraires augmente également dans ce cas. Les transformations infinitésimales satisfont donc, en raison de la linéarité en \(p(x)\), à un système d'équations différentielles partielles linéaires; et, puisque la propriété de groupe est vérifiée, elles forment un \og groupe infini de transformations infinitésimales \fg{} au sens de Lie (\emph{Grundlagen}, \S{}10).

La réciproque s'obtient maintenant de manière analogue au cas du groupe fini. L'existence des dépendances (16), après multiplication par \(p^{(\lambda)}(x)\) et addition, conduit, au moyen de la transformation identique (14), à
\[
\sum_i\psi_i\bd u_i=\Div\varGamma.
\]
De là on obtient, comme au \S{}3, la détermination de \(\varDelta x\) et de \(\varDelta u\), ainsi que l'invariance de \(I\) relativement à ces transformations infinitésimales, qui dépendent effectivement linéairement de \(\varrho\) fonctions arbitraires et de leurs dérivées jusqu'à l'ordre \(\sigma\). Que ces transformations infinitésimales, lorsqu'elles ne contiennent pas les dérivées \(\frac{\p u}{\p x},\cdots\), forment sûrement un groupe, résulte comme au \S{}3: sinon, par composition, apparaîtraient plus de fonctions arbitraires, alors que, par hypothèse, il ne doit y avoir que \(\rho\) dépendances (16). Elles forment donc un \og groupe infini de transformations infinitésimales \fg{}. Or un tel groupe consiste (\emph{Grundlagen}, théorème VII, p. 391) en les transformations infinitésimales les plus générales d'un certain \og groupe infini \(\G\) de transformations finies \fg{} ainsi défini au sens de Lie. Chaque transformation finie est alors engendrée à partir de transformations infinitésimales (\emph{Grundlagen}, \S{}7),\footnote{Il en résulte en particulier que le groupe \(\G\) engendré à partir des transformations infinitésimales \(\varDelta x,\varDelta u\) d'une \(\G_{\infty\varrho}\) ramène de nouveau à \(\G_{\infty\varrho}\). Car \(\G_{\infty\varrho}\) ne contient pas de transformations infinitésimales dépendant de fonctions arbitraires et différentes de \(\varDelta x,\varDelta u\), et ne peut pas non plus contenir de transformations indépendantes dépendant de paramètres, puisqu'autrement elle serait un groupe mixte. Or, d'après ce qui précède, les transformations finies sont déterminées par les infinitésimales.} et provient donc de l'intégration du système simultané:
\[
\frac{d x_i}{d t}=\varDelta x_i;\qquad
\frac{d u_i}{d t}=\varDelta u_i\quad
\left(\begin{array}{l}
x_i=y_i,\\
u_i=v
\end{array}\right\}
\quad\text{pour }t=0.
\]
Il peut toutefois être nécessaire de laisser les \(p(x)\) arbitraires dépendre encore de \(t\). Ainsi \(\G\) dépend effectivement de \(\varrho\) fonctions arbitraires. En particulier, s'il suffit de prendre \(p(x)\) indépendante de \(t\), cette dépendance est analytique dans les fonctions arbitraires \(q(x)=t\cdot p(x)\).\footnote{La question de savoir si ce dernier cas se produit peut-être toujours a été posée par Lie dans une autre formulation (\emph{Grundlagen}, \S{}7 et \S{}13, fin).} Si des dérivées \(\frac{\p u}{\p x},\cdots\) interviennent, il peut être nécessaire d'adjoindre encore des transformations infinitésimales \(\bd u=0,\ \Div(f\cdot\varDelta x)=0\) avant de tirer les mêmes conclusions.

À la suite d'un exemple de Lie (\emph{Grundlagen}, \S{}7), indiquons encore un cas assez général où l'on peut aller jusqu'à des formules explicites, qui montrent en même temps que les dérivées des fonctions arbitraires apparaissent jusqu'à l'ordre \(\sigma\); ici la réciproque est donc complète. Il s'agit des groupes de transformations infinitésimales auxquels correspond le groupe de toutes les transformations des \(x\), et des transformations des \(u\) qui en sont \og induites \fg{}; c'est-à-dire des transformations des \(u\) dans lesquelles \(\varDelta u\), et par conséquent \(u\), ne dépend que des fonctions arbitraires qui apparaissent dans \(\varDelta x\). On suppose encore que les dérivées \(\frac{\p u}{\p x},\cdots\) n'apparaissent pas dans \(\varDelta u\). On a donc:
\[
\varDelta x_i=p^{(i)}(x);\qquad
\varDelta u_i=
\sum_{\lambda=1}^{n}\left\{
a^{(\lambda)}(x,u)p^{(\lambda)}
+b^{(\lambda)}\frac{\p p^{(\lambda)}}{\p x}
+\cdots+
c^{(\lambda)}\frac{\p^\sigma p^{(\lambda)}}{\p x^\sigma}
\right\}.
\]
Comme la transformation infinitésimale \(\varDelta x=p(x)\) engendre toute transformation \(x=y+g(y)\) avec \(g(y)\) arbitraire, on peut en particulier déterminer \(p(x)\) en fonction de \(t\) de manière que soit engendré le groupe à un paramètre:
\begin{equation}
x_i=y_i+t\cdot g_i(y),
\tag{18}
\end{equation}
qui, pour \(t=0\), devient l'identité et, pour \(t=1\), devient la transformation cherchée \(x=y+g(y)\). En différentiant (18), on obtient en effet:
\begin{equation}
\frac{d x_i}{d t}=g_i(y)=p^{(i)}(x,t),
\tag{19}
\end{equation}
où \(p(x,t)\) est déterminé à partir de \(g(y)\) par inversion de (18); inversement, (18) résulte de (19) au moyen de la condition auxiliaire \(x_i=y_i\) pour \(t=0\), qui fixe l'intégrale de manière univoque. Au moyen de (18), les \(x\) qui figurent dans \(\varDelta u\) peuvent être remplacés par les constantes d'intégration \(y\) et par \(t\); les \(g(y)\) apparaissent alors exactement jusqu'à leur \(\sigma\)-ième dérivée, car dans
\[
\frac{\p p}{\p x}=\sum\frac{\p g}{\p y_\varkappa}\frac{\p y_\varkappa}{\p x}
\]
les \(\frac{\p y}{\p x}\) sont exprimées par \(\frac{\p x}{\p y}\), et, en général, \(\frac{\p^\sigma p}{\p x^\sigma}\) est remplacé par sa valeur dans
\[
\frac{\p g}{\p y},\ldots,\frac{\p x}{\p y},\ldots,\frac{\p^\sigma x}{\p y^\sigma}.
\]
Pour déterminer les \(u\), on obtient donc le système d'équations
\[
\frac{d u_i}{d t}
=
F_i\!\left(g(y),\frac{\p g}{\p y},\cdots\frac{\p^\sigma g}{\p y^\sigma},u,t\right),
\qquad (u_i=v_i\ \text{pour }t=0),
\]
où seules \(t\) et \(u\) sont variables, tandis que \(g(y),\ldots\) appartiennent au domaine des coefficients; l'intégration donne donc
\[
u_i=v_i+
B_i\!\left(v,g(y),\frac{\p g}{\p y},\cdots\frac{\p^\sigma g}{\p y^\sigma},t\right)_{t=1}.
\]
On obtient ainsi des transformations qui dépendent exactement de \(\sigma\) dérivées des fonctions arbitraires. L'identité y est contenue pour \(g(y)=0\), d'après (18); et la propriété de groupe résulte du fait que le procédé indiqué fournit chaque transformation \(x=y+g(y)\), ce par quoi la transformation induite des \(u\) est fixée univoquement, de sorte que le groupe \(\G\) est épuisé.

Il résulte encore incidemment de la réciproque que ce n'est pas une restriction de supposer les fonctions arbitraires dépendantes seulement des \(x\), et non de \(u,\frac{\p u}{\p x},\cdots\). Dans ce dernier cas, dans la transformation identique (14), donc aussi dans (15), apparaîtraient, en plus des \(p^{(\lambda)}\), les dérivées \(\frac{\p p^{(\lambda)}}{\p u},\ \frac{\p p^{(\lambda)}}{\p\frac{\p u}{\p x}},\cdots\). Si l'on prend alors successivement les \(p^{(\lambda)}\) comme étant de degré zéro, premier, etc. en \(u,\frac{\p u}{\p x},\cdots\), avec des fonctions arbitraires de \(x\) comme coefficients, des dépendances (16) apparaissent de nouveau, mais en plus grand nombre; par la réciproque précédente, elles ramènent toutefois au cas antérieur par regroupement avec des fonctions arbitraires qui ne dépendent que de \(x\). De même, on montre qu'au surgissement simultané de dépendances et de relations de divergence indépendantes de celles-ci correspondent des groupes mixtes.\footnote{Comme au \S{}3, il résulte encore ici de la réciproque qu'à côté de \(I\), toute intégrale \(I^*\) qui diffère de lui par une intégrale portant sur une divergence admet également un groupe infini avec les mêmes \(\bd u\); ici toutefois \(\varDelta x\) et \(\varDelta u\) contiendront en général des dérivées des \(u\). Einstein a introduit une telle intégrale \(I^*\) dans la relativité générale afin d'obtenir une formulation plus simple des théorèmes d'énergie; j'indique les transformations infinitésimales que cet \(I^*\) admet, en me conformant exactement à la notation de la seconde note de Klein. L'intégrale \(I=\int\!\ldots\!\int K\,d\omega=\int\!\ldots\!\int \mathfrak{R}\,d S\) admet le groupe de toutes les transformations des \(w\) et le groupe induit par là pour les \(g_{\mu\nu}\). À cela correspondent les dépendances ((30) chez Klein):
\[
\sum \mathfrak{R}_{\mu\nu}g_\tau^{\mu\nu}
+2\sum\frac{\p g^{\mu\nu}\mathfrak{R}_{\mu\tau}}{\p w^\sigma}=0.
\]
Or \(I^*=\int\!\ldots\!\int \mathfrak{R}^*\,d S\), où \(\mathfrak{R}^*=\mathfrak{R}+\Div\), et par conséquent \(\mathfrak{R}_{\mu\nu}^*=\mathfrak{R}_{\mu\nu}\), les \(\mathfrak{R}_{\mu\nu}^*,\mathfrak{R}_{\mu\nu}\) désignant respectivement les expressions de Lagrange. Les dépendances indiquées sont donc aussi des dépendances pour \(\mathfrak{R}_{\mu\nu}^*\); et, après multiplication par \(p^\tau\) et addition, en renversant les transformations de la différentiation d'un produit, on obtient
\[
\sum \mathfrak{R}_{\mu\nu}p^{\mu\nu}
+2\Div\!\left(\sum g^{\mu\sigma}\mathfrak{R}_{\mu\tau}p^\tau\right)=0,
\]
\[
\delta\mathfrak{R}^*+\Div\!\left(\sum 2g^{\mu\sigma}\mathfrak{R}_{\mu\tau}p^\tau
-\frac{\p\mathfrak{R}^*}{\p g^{\mu\nu}_{\sigma}}p^{\mu\nu}\right)=0.
\]
La comparaison avec l'équation différentielle de Lie
\(\delta\mathfrak{R}^*+\Div(\mathfrak{R}^*\varDelta w)=0\) donne
\[
\varDelta w^\sigma=\frac1{\mathfrak{R}^*}\cdot\left(\sum (2g^{\mu\sigma}\mathfrak{R}_{\mu\tau}p^\tau
-\frac{\p\mathfrak{R}^*}{\p g^{\mu\nu}_{\sigma}}p^{\mu\nu}\right),
\qquad
\varDelta g^{\mu\nu}=p^{\mu\nu}+\sum g^{\mu\nu}_{\sigma}\varDelta w^\sigma,
\]
comme transformations infinitésimales admises par \(I^*\). Ces transformations infinitésimales dépendent donc des premières et secondes dérivées des \(g^{\mu\nu}\), et contiennent les \(p\) arbitraires jusqu'à la première dérivée.}

\section*{\S{}5. Invariance des composantes individuelles des relations}

Si l'on spécialise le groupe \(\G\) au cas le plus simple, ordinairement considéré, en n'admettant dans les transformations aucune dérivée des \(u\), et en faisant dépendre les variables indépendantes transformées seulement des \(x\), non des \(u\), alors on peut conclure à l'invariance des composantes individuelles dans les formules. D'abord, par des raisonnements connus, on obtient l'invariance de
\[
\int\!\cdots\!\int\left(\sum\psi_i\delta u_i\right)\dd x,
\]
donc l'invariance relative de \(\sum\psi_i\delta u_i\),\footnote{C'est-à-dire que \(\sum\psi_i\delta u_i\) acquiert un facteur sous les transformations, ce que la théorie algébrique des invariants a toujours appelé invariance relative.} où \(\delta\) désigne une variation quelconque. En effet, d'une part,
\[
\delta I
=\int\!\cdots\!\int \delta f\!\left(x,u,\frac{\p u}{\p x},\cdots\right)d x
=\int\!\cdots\!\int \delta f\!\left(y,v,\frac{\p v}{\p y},\cdots\right)d y.
\]
D'autre part, pour \(\delta u,\delta\frac{\p u}{\p x},\cdots\) s'annulant au bord, auxquels, en raison de la transformation linéaire homogène des \(\delta u,\delta\frac{\p u}{\p x},\cdots\), correspondent des \(\delta v,\delta\frac{\p v}{\p y},\cdots\) s'annulant également au bord, on a
\[
\int\!\cdots\!\int \delta f\!\left(x,u,\frac{\p u}{\p x},\cdots\right)d x
=
\int\!\cdots\!\int\left(\sum\psi_i(u,\ldots)\delta u_i\right)d x,
\]
\[
\int\!\cdots\!\int \delta f\!\left(y,v,\frac{\p v}{\p x},\cdots\right)d y
=
\int\!\cdots\!\int\left(\sum\psi_i(v,\ldots)\delta v_i\right)d y.
\]
Par suite, pour \(\delta u,\delta\frac{\p u}{\p x},\cdots\) s'annulant au bord,
\[
\int\!\cdots\!\int\left(\sum\psi_i(u,\ldots)\delta u_i\right)d x
=
\int\!\cdots\!\int\left(\sum\psi_i(v,\ldots)\delta v_i\right)d y
\]
\[
=
\int\!\cdots\!\int\left(\sum\psi_i(v,\ldots)\delta v_i\right)
\left|\frac{\p y_i}{\p x_\varkappa}\right|d x.
\]
Si l'on exprime dans la troisième intégrale \(y,v,\delta v\) par \(x,u,\delta u\), et qu'on l'égale à la première, on obtient donc une relation
\[
\int\!\cdots\!\int\left(\sum\chi_i(u,\ldots)\delta u_i\right)d x=0
\]
pour \(\delta u\) s'annulant au bord, mais par ailleurs arbitraire. Il en résulte, comme on sait, l'annulation de l'intégrande pour \(\delta u\) arbitraire; la relation
\[
\sum\psi_i(u,\ldots)\delta u_i
=
\left|\frac{\p y_i}{\p x_\varkappa}\right|
\left(\sum\psi_i(v,\ldots)\delta v_i\right)
\]
vaut donc identiquement en \(\delta u\). Elle exprime l'invariance relative de \(\sum\psi_i\delta u_i\), et par conséquent l'invariance de
\[
\int\!\cdots\!\int\left(\sum\psi_i\delta u_i\right)\dd x.
\]
\footnote{Ces raisonnements échouent lorsque \(y\) dépend aussi des \(u\), car alors \(\delta f(y,v,\p v/\p y,\ldots)\) contient aussi des termes \(\sum(\p f/\p y)\delta y\), de sorte que la transformation de divergence ne conduit pas aux expressions de Lagrange. Ils échouent de même si l'on admet des dérivées des \(u\); alors les \(\delta v\) deviennent des combinaisons linéaires de \(\delta u,\delta(\p u/\p x),\ldots\), et ne conduisent donc qu'après une nouvelle transformation de divergence à une identité \(\int\!\cdots\!\int(\sum\chi_i(u,\ldots)\delta u_i)\dd x=0\), si bien qu'à droite les expressions de Lagrange n'apparaissent de nouveau pas. La question de savoir si l'invariance de \(\int\!\cdots\!\int(\sum\psi_i\delta u_i)\dd x\) permet déjà de conclure à l'existence de relations de divergence est, d'après la réciproque, équivalente à la question de savoir si elle entraîne l'invariance de \(I\) vis-à-vis d'un groupe qui ne conduit pas nécessairement aux mêmes \(\D u,\D x\), mais bien aux mêmes \(\bd u\). Dans le cas spécial de l'intégrale simple et de seules dérivées premières dans \(f\), on peut, pour le groupe fini, conclure de l'invariance des expressions de Lagrange à l'existence de premières intégrales (voir par exemple Engel, Gött. Nachr. 1916, p. 270).}

Pour appliquer ceci aux relations de divergence et aux dépendances dérivées, il faut d'abord montrer que le \(\bd u\) déduit de \(\varDelta u,\varDelta x\) satisfait effectivement aux lois de transformation de la variation \(\delta u\), pourvu seulement que les paramètres, respectivement les fonctions arbitraires, dans \(\bd v\) soient déterminés comme ils correspondent au groupe semblable des transformations infinitésimales en \(y,v\). Soit \(\mathfrak T_q\) la transformation qui fait passer \(x,u\) en \(y,v\), et soit \(\mathfrak T_p\) une transformation infinitésimale en \(x,u\). Alors la transformation semblable en \(y,v\) est donnée par \(\mathfrak T_r=\mathfrak T_q\mathfrak T_p\mathfrak T_q^{-1}\), les paramètres, respectivement les fonctions arbitraires, \(r\) étant donc déterminés par \(p\) et \(q\). En formules, cela s'exprime ainsi:
\[
\mathfrak T_p:\ \xi=x+\varDelta x(x,p);\qquad u^*=u+\varDelta u(x,u,p);
\]
\[
\mathfrak T_q:\ y=A(x,q);\qquad v=B(x,u,q);
\]
\[
\mathfrak T_q\mathfrak T_p:\ \eta=A(x+\varDelta x(x,p),q),\quad
v^*=B(x+\varDelta x(p),\,u+\varDelta u(p),q).
\]
Il en résulte \(\mathfrak T_r=\mathfrak T_q\mathfrak T_p\mathfrak T_q^{-1}\), donc
\[
\eta=y+\varDelta y(r);\qquad v^*=v+\varDelta v(r),
\]
où, au moyen de l'inverse de \(\mathfrak T_q\), on considère les \(x\) comme fonctions des \(y\) et l'on ne retient que les termes infinitésimaux. On a ainsi l'identité
\begin{equation}
\begin{gathered}
\eta=y+\varDelta y(r)=y+\sum\frac{\p A(x,q)}{\p x}\varDelta x(p),\\[0.55em]
v^*=v+\varDelta v(r)=v+\sum\frac{\p B(x,u,q)}{\p x}\varDelta x(p)
+\sum\frac{\p B(x,u,q)}{\p u}\varDelta u(p).
\end{gathered}
\tag{20}
\end{equation}
Si l'on y remplace \(\xi=x+\varDelta x\) par \(\xi-\varDelta \xi\), de sorte que \(\xi\) revienne à \(x\), et que donc \(\varDelta x\) s'annule, alors, d'après la première formule (20), \(\eta\) revient aussi à \(y=\eta-\varDelta\eta\). Si, sous cette substitution, \(\varDelta u(p)\) passe en \(\bd u(p)\), alors \(\varDelta v(r)\) passe aussi en \(\bd v(r)\), et la seconde formule (20) donne
\[
v+\bd v(y,v,\ldots r)=v+\sum\frac{\p B(x,u,q)}{\p u}\,\bd u(p),
\]
\[
\bd v(y,v,\ldots r)=\sum\frac{\p B}{\p u_\varkappa}\,\bd u_\varkappa(x,u,p).
\]
Les formules de transformation pour les variations sont donc effectivement satisfaites, pourvu seulement que \(\bd v\) soit pris comme dépendant des paramètres, respectivement des fonctions arbitraires, \(r\).\footnote{On voit de nouveau qu'il faut supposer \(y\) indépendant de \(u\), etc., pour que les conclusions soient valables. Comme exemple, on peut citer les \(\delta g^{\mu\nu}\) et \(\delta q_\varrho\) indiqués par Klein, qui satisfont aux transformations pour variations dès que les \(p\) sont soumis à une transformation vectorielle.}

Il s'ensuit en particulier l'invariance relative de \(\sum\psi_i\bd u_i\); donc aussi, d'après (12), puisque les relations de divergence sont également satisfaites en \(y,v\), l'invariance relative de \(\Div B\); et encore, d'après (14) et (13), l'invariance relative de \(\Div\varGamma\) et des membres de gauche des dépendances regroupés avec les \(p^{(\lambda)}\), où, dans les formules transformées, les \(p(x)\) arbitraires (respectivement les paramètres) doivent partout être remplacés par les \(r\). Il en résulte encore l'invariance relative de \(\Div(B-\varGamma)\), c'est-à-dire d'une divergence d'un système de fonctions \(B-\varGamma\) non identiquement nul, mais dont la divergence s'annule identiquement.

De l'invariance relative de \(\Div B\), dans le cas unidimensionnel et pour le groupe fini, on peut encore tirer une conclusion sur l'invariance des premières intégrales. La transformation des paramètres correspondant à la transformation infinitésimale devient, d'après (20), linéaire et homogène; et, parce que toutes les transformations sont inversibles, les \(\eps\) sont aussi linéaires et homogènes dans les paramètres transformés \(\eps^*\). Cette inversibilité demeure certainement conservée lorsqu'on pose \(\psi=0\), puisque aucune dérivée de \(u\) n'intervient dans (20). En égalant les coefficients des \(\eps^*\) dans
\[
\Div B(x,u,\ldots\eps)=\frac{d y}{d x}\cdot\Div B(y,v,\ldots\eps^*)
\]
les quantités \(\frac{d}{d y}B^{\lambda)}(y,v,\ldots)\) deviennent aussi des fonctions linéaires homogènes des \(\frac{d}{d x}B^{(\lambda)}(x,u,\ldots)\). Ainsi, de
\[
\frac{\dd}{\dd x}B^{(\lambda)}(x,u,\ldots)=0
\quad\text{ou}\quad B^{(\lambda)}(x,u)=\mathrm{const.}
\]
il résulte
\[
\frac{\dd}{\dd y}B^{(\lambda)}(y,v,\ldots)=0
\quad\text{ou}\quad B^{(\lambda)}(y,v)=\mathrm{const.}
\]
Les \(\rho\) premières intégrales qui correspondent à une \(\G_\varrho\) admettent donc chacune le groupe, de sorte que l'intégration ultérieure s'en trouve également simplifiée. L'exemple le plus simple en est le cas où \(f\) ne dépend pas de \(x\), ou d'un \(u\), ce qui correspond aux transformations infinitésimales \(\varDelta x=\eps,\varDelta u=0\), respectivement \(\varDelta x=0,\varDelta u=\eps\). Alors \(\bd u=-\eps\frac{d u}{d x}\), respectivement \(\eps\); et, comme \(B\) se déduit de \(f\) et de \(\bd u\) par différentiation et combinaisons rationnelles, il ne dépend donc pas non plus de \(x\), respectivement de \(u\), et admet les groupes correspondants.\footnote{Dans les cas où l'existence de premières intégrales résulte déjà de l'invariance de \(\int(\sum\psi_i\delta u_i)d x\), celles-ci n'admettent pas le groupe complet \(\G_\varrho\). Par exemple, \(\int(u''\delta u)d x\) admet la transformation infinitésimale \(\varDelta x=\eps_2;\ \varDelta u=\eps_1+x\eps_3\), tandis que la première intégrale \(u-u'x=\mathrm{const.}\), correspondant à \(\varDelta x=0,\varDelta u=x\eps_3\), n'admet pas les deux autres transformations infinitésimales, puisqu'elle contient explicitement à la fois \(u\) et \(x\). À cette première intégrale correspondent précisément des transformations infinitésimales pour \(f\) qui contiennent des dérivées. On voit donc que l'invariance de \(\int\!\ldots\!\int(\sum\psi_i\,d u_i)d x\) réalise en tout cas moins que l'invariance de \(I\), remarque à faire à propos de la question soulevée dans une note précédente.}

\section*{\S{}6. Une assertion de Hilbert}

De ce qui précède résulte enfin la preuve d'une assertion de Hilbert concernant le lien entre l'échec des théorèmes d'énergie propres et la \og relativité générale \fg{} (première note de Klein, Göttinger Nachr. 1917, Réponse, premier paragraphe), sous une forme généralisée du point de vue de la théorie des groupes.

Supposons que l'intégrale \(I\) admette une \(\G_{\infty\varrho}\), et que \(\G_\sigma\) soit un groupe fini quelconque provenant d'une spécialisation des fonctions arbitraires, donc un sous-groupe de \(\G_{\infty\varrho}\). Au groupe infini \(\G_{\infty\varrho}\) correspondent alors des dépendances (16), au groupe fini \(\G_\sigma\) des relations de divergence (13); et, réciproquement, de l'existence de relations de divergence quelconques résulte l'invariance de \(I\) par rapport à un groupe fini, qui est identique à \(\G_\sigma\) si et seulement si les \(\bd u\) sont des combinaisons linéaires de celles qui proviennent de \(\G_\sigma\). L'invariance par rapport à \(\G_\sigma\) ne peut donc conduire à aucune relation de divergence différente de (13). Mais comme de l'existence de (16) résulte l'invariance de \(I\) par rapport aux transformations infinitésimales \(\varDelta u,\varDelta x\) de \(\G_{\infty\varrho}\) pour \(p(x)\) arbitraire, il en résulte en particulier déjà l'invariance par rapport aux transformations infinitésimales de \(\G_\sigma\) qui naissent de celles-ci par spécialisation, et donc par rapport à \(\G_\sigma\). Les relations de divergence
\[
\sum\psi_i\bd u_i^{(\lambda)}=\Div B^{(\lambda)}
\]
doivent donc être des conséquences des dépendances (16), que l'on peut aussi écrire
\[
\sum\psi_i a_i^{(\lambda)}=\Div\chi^{(\lambda)},
\]
où les \(\chi^{(\lambda)}\) sont des combinaisons linéaires des expressions de Lagrange et de leurs dérivées. Comme les \(\psi\) interviennent linéairement aussi bien dans (13) que dans (16), les relations de divergence doivent donc, en particulier, être des combinaisons linéaires des dépendances (16). Ainsi
\[
\Div B^{(\lambda)}=\Div\left(\sum \alpha\cdot\chi^{(\varkappa)}\right),
\]
et les \(B^{(\lambda)}\) eux-mêmes se composent donc linéairement des \(\chi\), c'est-à-dire des expressions de Lagrange et de leurs dérivées, et de fonctions dont la divergence s'annule identiquement, comme par exemple les \(B-\varGamma\) qui apparaissent à la fin du \S{}2, pour lesquels \(\Div(B-\varGamma)=0\), et où la divergence possède en même temps une propriété d'invariant. J'appellerai \og impropres \fg{} les relations de divergence dans lesquelles les \(B^{(\;)}\) peuvent se composer de la manière indiquée à partir des expressions de Lagrange et de leurs dérivées; toutes les autres seront appelées \og propres \fg{}.

Inversement, si les relations de divergence sont des combinaisons linéaires des dépendances (16), donc \og impropres \fg{}, alors l'invariance par rapport à \(\G_\sigma\) résulte de celle par rapport à \(\G_{\infty\varrho}\); \(\G\) devient un sous-groupe de \(\G_{\infty\varrho}\). Les relations de divergence correspondant à un groupe fini \(\G_\sigma\) deviennent donc impropres si et seulement si \(\G_\sigma\) est un sous-groupe d'un groupe infini par rapport auquel \(I\) est invariant.

Par spécialisation des groupes, l'assertion originale de Hilbert s'en déduit. Par \og groupe des translations \fg{}, on entend le groupe fini
\[
y_i=x_i+\eps_i;\qquad v_i(y)=u_i(x);
\]
donc
\[
\varDelta x_i=\eps_i,\qquad \varDelta u_i=0,\qquad
\bd u_i=-\sum_\lambda\frac{\p u_i}{\p x_\lambda}\eps_\lambda.
\]
L'invariance par rapport au groupe des translations signifie, comme on sait, que dans
\[
I=\int\!\cdots\!\int f\!\left(x,u,\frac{\p u}{\p x},\ldots\right)\dd x
\]
les \(x\) n'apparaissent pas explicitement dans \(f\). Les \(n\) relations de divergence correspondantes
\[
\sum\psi_i\frac{\p u_i}{\p x_\lambda}=\Div B^{(\lambda)}
\qquad(\lambda=1,2,\ldots n)
\]
seront appelées \og relations d'énergie \fg{}, parce que les \og lois de conservation \fg{} \(\Div B^{(\lambda)}=0\) qui correspondent au problème variationnel répondent aux \og théorèmes d'énergie \fg{}, et les \(B^{(\lambda)}\) aux \og composantes d'énergie \fg{}. On obtient donc: si \(I\) admet le groupe des translations, les relations d'énergie deviennent impropres si et seulement si \(I\) est invariant par rapport à un groupe infini contenant le groupe des translations comme sous-groupe.\footnote{Les théorèmes d'énergie de la mécanique classique, de même que ceux de l'ancienne \og théorie de la relativité \fg{} (où \(\sum d x^2\) se transforme en lui-même), sont \og propres \fg{}, car aucun groupe infini n'intervient ici.}

Un exemple de tels groupes infinis est donné par le groupe de toutes les transformations des \(x\) et des transformations induites des \(u(x)\), dans lesquelles n'apparaissent que des dérivées des fonctions arbitraires \(p(x)\). Le groupe des translations naît par la spécialisation \(p^{(i)}(x)=\eps_i\); il faut toutefois laisser indécise la question de savoir si, par ce moyen --- et par les groupes qui naissent lorsqu'on modifie \(I\) par une intégrale de bord --- on a déjà donné les groupes les plus généraux de ce type. Des transformations induites du genre indiqué apparaissent par exemple lorsqu'on soumet les \(u\) aux transformations de coefficients d'une \og forme différentielle totale \fg{}, c'est-à-dire d'une forme
\[
\sum a\,\dd^\lambda x_i+\sum b\,\dd^{\lambda-1}x_i\,\dd x_\varkappa+\cdots,
\]
qui contient, en plus des \(d x\), des différentielles d'ordre supérieur. Des transformations induites plus spéciales, dans lesquelles les \(p(x)\) n'apparaissent qu'en première dérivée, sont données par les transformations de coefficients des formes différentielles ordinaires \(\sum c\,d x_{i_1}\ldots d x_{i_\lambda}\); ce sont celles qu'on a ordinairement considérées.

Un autre groupe du type indiqué --- qui, à cause de la présence du terme logarithmique, ne peut pas être une transformation de coefficients --- est par exemple le suivant:
\[
y=x+p(x),\qquad
v_i=u_i+\lg(1+p'(x))=u_i+\lg\frac{d y}{d x};
\]
\[
\varDelta x=p(x);\qquad
\varDelta u_i=p'(x)\srcfnmark{1)};\qquad
\bd u_i=p'(x)-u_i'p(x).
\]
\srcfntext{1)}{À partir de ces transformations infinitésimales, on calcule les transformations finies en remontant selon la méthode indiquée à la fin du \S{}4.}
Les dépendances (16) deviennent ici:
\[
\sum_i\left(\psi_i u_i'+\frac{d\psi_i}{d x}\right)=0,
\]
et les relations d'énergie impropres deviennent:
\[
\sum\left(\psi_i u_i'+
\frac{d(\psi_i+\mathrm{const.})}{d x}\right)=0.
\]
Une intégrale invariante la plus simple du groupe est:
\[
I=\int \frac{e^{-2u_1}}{u_1'-u_2'}\,d x.
\]
L'intégrale \(I\) la plus générale se détermine par intégration de l'équation différentielle de Lie (11)
\[
\bd f+\frac{d}{d x}(f\cdot \varDelta x)=0,
\]
qui, après substitution des valeurs de \(\varDelta x\) et de \(\bd u\), devient, dès que l'on suppose \(f\) dépendant seulement des dérivées premières des \(u\),
\[
\frac{\p f}{\p x}p(x)
+\left\{\sum\frac{\p f}{\p u_i}
-\frac{\p f}{\p u_i'}u_i'+f\right\}p'(x)
+\left\{\sum\frac{\p f}{\p u_i''}\right\}p''(x)=0
\]
(identiquement en \(p(x),p'(x),p''(x)\)). Déjà pour deux fonctions \(u(x)\), ce système d'équations possède des solutions qui contiennent effectivement les dérivées, à savoir
\[
f=(n_1'-u_2')\,
\varPhi\!\left(u_1-u_2,\frac{e^{-u_1}}{u_1'-u_2'}\right),
\]
où \(\varPhi\) désigne une fonction arbitraire des arguments indiqués.

Hilbert énonce son assertion en disant que l'échec des théorèmes d'énergie propres est un caractère distinctif de la \og théorie générale de la relativité \fg{}. Pour que cette assertion vaille littéralement, il faut donc comprendre la désignation \og relativité générale \fg{} plus largement que d'ordinaire, et l'étendre aussi aux groupes précédents dépendant de \(n\) fonctions arbitraires.\footnote{Ainsi se confirme de nouveau la justesse d'une remarque de Klein, selon laquelle la désignation \og relativité \fg{} usuelle en physique devrait être remplacée par \og invariance relativement à un groupe \fg{}. (\emph{Über die geometrischen Grundlagen der Lorentzgruppe}. \emph{Jhrber. d. d. Math. Vereinig.} vol. 19, p. 287, 1910; réimprimé dans la \emph{Physikalische Zeitschrift}.)}
\clearpage
\setcounter{footnote}{0}
\section*{\NoetherNavAnchor{P14}{14. La théorie arithmétique des fonctions algébriques d’une variable dans son rapport avec les autres théories et avec la théorie des corps de nombres}14. La théorie arithmétique des fonctions algébriques d'une variable dans son rapport avec les autres théories et avec la théorie des corps de nombres}
\begin{center}
Rapport d'Emmy Noether, à Göttingen.\\[0.5em]
\emph{Jahresbericht der Deutschen Mathematiker-Vereinigung} 28 (1919), pp. 182--203
\end{center}

Le rapport suivant est né du désir de disposer d'un lien entre les diverses théories des fonctions algébriques; en même temps, il peut être considéré comme un complément au rapport de Brill--Noether,\footnote{Le développement de la théorie des fonctions algébriques dans les temps anciens et récents. \emph{Jahresbericht der Deutschen Mathematiker-Vereinigung} vol. 3 (1894) [cité comme Brill--Noether].} où la théorie arithmétique est essentiellement absente.\footnote{À l'exception d'une discussion de la recherche de Kronecker sur le discriminant. (\emph{J. f. M.} vol. 91.)}

Pour les exposés des théories particulières, outre Brill--Noether, il faut prendre en considération les articles encyclopédiques suivants, avec la littérature qui y est indiquée:

Pour la théorie de Riemann: Wirtinger (II B 2), nos. 1--19;

pour la théorie de Weierstrass: Wirtinger (II B 2), nos. 23, 32, 33, 34;

pour la théorie algébro-géométrique de Brill--Noether: Wirtinger (II B 2), nos. 21--31; Berzolari (III C 4), nos. 23--38;

pour la théorie arithmétique de Dedekind--Weber et de Hensel--Landsberg: Hensel (II C 5), [cité comme Hensel];\footnote{Les parties concernant Dedekind--Weber sont dues à A. Ostrowski.}

plus généralement, pour la théorie arithmétique des grandeurs algébriques: Landsberg (I C 5), et pour la théorie arithmétique des fonctions algébriques de plusieurs variables: Jung (II C 6).

Pour la théorie algébrique des nombres, comparer le \og{}Zahlbericht\fg{} de Hilbert\footnote{La théorie des corps de nombres algébriques. \emph{Jahresbericht der Deutschen Mathematiker-Vereinigung} vol. 4 (1894/95).} et les leçons sur la théorie des nombres de Dirichlet--Dedekind [citées comme Dedekind--théorie des nombres].\footnote{Braunschweig, Vieweg, 4e éd. 1894.}

\subsection*{Table des matières}
\begin{enumerate}
\item Fondements des diverses théories.
\item Parallélisme entre nombres algébriques et fonctions algébriques (éléments conjugués, théorèmes de base, concept de module).
\item Parallélisme et différences dans la théorie des idéaux des nombres et des fonctions algébriques.
\item Rapport entre le fondement arithmétique du développement en séries et la théorie des idéaux.
\item Comparaison de la théorie arithmétique des fonctions algébriques avec la théorie géométrique. Concept d'invariance différent.
\item Suite de la comparaison. Points singuliers.
\item Comparaison entre le théorème des résidus et des théorèmes de la théorie des idéaux.
\item Questions transcendantes pour les fonctions algébriques et les nombres algébriques.
\end{enumerate}

\subsection*{1. Fondements des diverses théories}

La théorie de Riemann est la seule qui soit purement transcendante; elle repose sur des théorèmes d'existence (principe de Dirichlet, méthodes de Schwarz--Neumann). Les intégrales de fonctions algébriques y sont premières, les fonctions algébriques elles-mêmes secondaires. Les fonctions sur la surface de Riemann y sont caractérisées par leurs zéros et leurs pôles, correspondant aux polygones (respectivement aux diviseurs) de la théorie arithmétique; tandis que Weierstrass parle aussi des fonctions elles-mêmes avec leurs zéros et leurs pôles. Le problème principal, abstraction faite des théorèmes d'existence, est l'établissement de toutes les fonctions ayant des pôles donnés, ce qui est réalisé par les intégrales de première et de seconde espèce; s'y ajoute la question du nombre des constantes arbitraires qui interviennent encore lorsque les points à l'infini sont prescrits. Cette dernière question est résolue par le théorème de Riemann--Roch; dans la théorie arithmétique lui correspond la question de la dimension d'une classe de polygones (respectivement d'une classe de diviseurs), et dans la théorie algébro-géométrique la question de la dimension d'un système linéaire complet, à laquelle le théorème de Riemann--Roch donne également la réponse. La construction effective des fonctions est effectuée, dans la théorie arithmétique, au moyen des théorèmes de base; dans la théorie algébro-géométrique, par le théorème des résidus. En particulier, les zéros des intégrandes de première espèce correspondent à la classe différentielle dans la théorie arithmétique, et à la doctrine des groupes spéciaux (c'est-à-dire au faisceau des groupes de points découpés par les courbes adjointes d'ordre \((n-3)\), les courbes \(\varphi\)) dans la théorie algébro-géométrique.

De même que la théorie de Riemann a historiquement précédé les autres théories, bien que la preuve rigoureuse des théorèmes d'existence n'ait été obtenue que plus tard, de même des questions algébriques ont été résolues ensuite, avec ses auxiliaires transcendants, qui jusqu'à aujourd'hui ne sont pas encore accessibles à un traitement purement algébrique ou arithmétique; il faut mentionner ici avant tout la théorie de Hurwitz des correspondances singulières.\footnote{A. Hurwitz, Über algebraische Korrespondenzen und das erweiterte Korrespondenzprinzip. \emph{Math. Ann.} vol. 28. (Comparer Brill--Noether, section 10.)}

Les analogies avec les questions transcendantes dans la théorie des corps de nombres algébriques seront discutées dans la dernière section.

Pour les autres théories, les fonctions algébriques sont premières, les intégrales secondaires. À cet égard, la théorie de Weierstrass repose sur les fondements transcendants du développement en séries de puissances et du théorème des résidus; mais elle procède ensuite algébriquement, par l'indication explicite de toutes les fonctions considérées.

La théorie arithmétique sous la forme de Hensel--Landsberg utilise également les fondements transcendants du développement en séries. Mais en assignant, sur cette base, des diviseurs aux points individuels, elle procède ensuite de façon purement arithmétique et peut ainsi être regardée comme une fusion de Weierstrass et de Dedekind--Weber. Par contraste avec cette dernière, elle fournit en même temps des méthodes pour la construction effective de toutes les fonctions de base en question. Dans le détail, elle présente un certain nombre de simplifications par rapport à Dedekind--Weber, essentiellement par l'introduction de diviseurs \og{}fractionnaires\fg{}. Récemment Hensel\footnote{\emph{Mathematische Zeitschrift} vol. 4; ce fondement est aussi utilisé comme base de son rapport.} a donné un fondement arithmétique à la théorie des développements en séries (abstraction faite d'une méthode de majorantes pour la preuve de convergence), par quoi la théorie est devenue presque purement arithmétique.

La théorie algébro-géométrique de Brill--Noether ne présuppose que le concept de \og{}points successifs\fg{} d'une branche, qui est ramené au développement en série en un point ordinaire\footnote{Pour les points successifs \og{}singuliers\fg{}, comparer Brill--Noether, section 6, en particulier nos. 10 et 11. Des points successifs ordinaires sont par exemple les points d'intersection de la tangente avec la courbe ou des points de contact d'ordre supérieur (points de branchement).} et correspond au concept d'élément chez Weierstrass; elle continue ensuite par des méthodes algébriques rationnelles, essentiellement celles de la théorie de l'élimination ternaire; elle parvient elle aussi à des représentations explicites.

La théorie originelle, purement arithmétique, de Dedekind--Weber\footnote{\emph{J. f. M.} 92 (cité comme D.--W.).} est apparentée à celle de Riemann en ce que ses théorèmes ont de nouveau le caractère de preuves d'existence; elle conserve une pureté complète de méthode.\footnote{Pour cette raison, elle se prête mieux aussi à la comparaison avec les autres théories que, par exemple, la théorie de Hensel--Landsberg, qui n'interviendra qu'occasionnellement dans les comparaisons qui suivent. --- La théorie des intégrales abéliennes et de leur inversion, dans la mesure où elle présuppose la continuité, n'est cependant pas traitée dans D.--W., ni dans le développement ultérieur de la théorie algébrique (M. Noether, \emph{Ann.} 37).} Le fondement arithmétique de la théorie, où les variables ne jouent que le rôle d'indéterminées, est parallèle à la théorie des nombres algébriques, comme il sera expliqué plus précisément en 2.\footnote{Une comparaison de la théorie de Hensel--Landsberg avec la théorie des nombres, en particulier la théorie des nombres \(p\)-adiques, se trouve chez Hensel.} Ce fondement donne le moyen d'introduire arithmétiquement le concept de point (comparer Hensel no. 10a), et donc d'établir le concept de surface de Riemann absolue sans aucune considération de continuité. La théorie acquiert par là un caractère formel; par exemple, le concept de différentielle est également introduit de façon purement formelle. Sur la base obtenue par le concept de point, on peut démontrer la parenté étroite avec les méthodes et concepts de la théorie algébro-géométrique, puisque celle-ci aussi, abstraction faite du concept d'élément, est libre d'auxiliaires transcendants; cela sera fait dans les sections suivantes.

\subsection*{2. Parallélisme entre nombres algébriques et fonctions algébriques. (Éléments conjugués, théorèmes de base, concept de module)}

Un corps de nombres algébriques est défini comme un corps de degré \(n\) sur le corps \(R\) des nombres rationnels; un corps de fonctions algébriques d'une variable, comme un corps de degré \(n\) sur le corps \(K(z)\) de toutes les fonctions rationnelles d'une indéterminée \(z\) à coefficients complexes arbitraires. Ainsi, pour les deux corps valent en commun tous les théorèmes qui font complètement abstraction de la nature particulière du corps de base; ce sont les théorèmes sur les bases, les normes, les traces et les discriminants.\footnote{D.--W., \S{} 1 et 2; Hensel no. 4 (avec note 5).} Tous ces concepts sont introduits par Dedekind--Weber sans l'emploi des éléments conjugués. Il faut toutefois remarquer que le concept de corps conjugué peut être défini d'une façon entièrement formelle, comme Steinitz l'a montré à la suite de Kronecker.\footnote{Steinitz, Algebraische Theorie der Körper. (\emph{J. f. M.} 137, \S{} 6 et 8.)} En effet, si \(K\) est un corps arbitraire et \(\varphi(x)\) un polynôme en \(x\), irréductible sur \(K\), alors le système des classes modulo \(\varphi(x)\) forme un corps. Si l'on assigne à la classe de \(x\) modulo \(\varphi(x)\) un élément \(\xi\), alors, sous l'assignation isomorphe, l'élément \(f(\xi)\) correspond à la classe modulo \(\varphi(x)\) représentée par le polynôme \(f(x)\). Comme tous les polynômes divisibles par \(\varphi(x)\) représentent la classe nulle, \(\varphi(\xi)\) est assigné à la classe nulle; l'élément \(\xi\) peut donc être regardé symboliquement comme un zéro de \(\varphi(x)=0\) et engendre le corps algébrique \(K(\xi)\). La répétition du procédé conduit aux \(n\) corps conjugués \(K(j), K(j'),\ldots,K(j^{(n-1)})\). On peut ainsi, en particulier, parler d'éléments conjugués aussi pour les fonctions algébriques, sans quitter la théorie purement arithmétique.

Les deux corps de base \(R\) et \(K(z)\) ont en commun la propriété que le système de toutes les grandeurs entières y est défini, respectivement comme entiers rationnels et comme fonctions rationnelles entières d'une indéterminée. Ces grandeurs entières ont la propriété caractéristique que deux quelconques \(a\) et \(b\) possèdent un plus grand commun diviseur \(d\), représentable sous la forme \(ma+nb\), où \(d,m,n\) sont encore des grandeurs entières de \(R\) ou de \(K(z)\); et en fait \(d\) est le nombre absolument le plus petit, respectivement le polynôme de plus bas degré, représentable sous cette forme. C'est sur cette seule propriété, qui entraîne la décomposition unique en facteurs premiers pour les entiers rationnels et pour les fonctions d'une indéterminée, que reposent les théorèmes concernant les grandeurs entières du corps algébrique ou corps supérieur. L'argument qui prouve l'existence du plus grand commun diviseur donne aussi l'existence d'un système fondamental ou d'une base des grandeurs entières du corps algébrique. Si l'on ne considère que des bases du corps composées de grandeurs entières, alors leur discriminant devient un entier rationnel, respectivement une fonction rationnelle d'une indéterminée. Toute base qui correspond au nombre absolument le plus petit, respectivement à la fonction de plus bas degré, forme un système fondamental.

Une vue plus précise des questions de base, également pour des corps plus généraux, est donnée par l'introduction du concept de module, qui doit être formulé de manière à comprendre les définitions qui apparaissent diversement dans la littérature selon la nature des éléments. Un module --- relativement à un domaine intègre de base \(\mathfrak J\) --- est un système d'éléments tel que la différence de deux éléments appartienne encore au système, et de même le produit d'un élément par une grandeur arbitraire de \(\mathfrak J\). (Si \(\mathfrak J\) consiste en les entiers rationnels, la seconde exigence est conséquence de la première.)\footnote{Le concept de module est dû à Dedekind, qui introduisit d'abord le module de nombres (\(\mathfrak J=\) entiers rationnels) dans la deuxième édition de la théorie des nombres; le module des formes linéaires entières y est aussi contenu implicitement. Dans Dedekind--Weber apparaissent les \og{}modules de fonctions\fg{}, dont les éléments consistent en fonctions algébriques, tandis que \(\mathfrak J\) consiste en toutes les fonctions rationnelles entières d'une indéterminée. --- Le module des formes homogènes en \(n\) indéterminées, plus généralement des polynômes, apparaît d'abord chez Hilbert (\emph{Annalen} 36); \(\mathfrak J\) consiste en tous les polynômes, et puisque les éléments sont aussi des polynômes, ce module tombe sous notre définition d'idéal. Les \og{}systèmes de modules\fg{} de polynômes de Kronecker partent de la base, non de la définition abstraite.}

Tout module qui possède une base d'un nombre fini d'éléments, par lesquels tout élément peut être représenté linéairement à coefficients dans \(\mathfrak J\), est appelé fini (engendré finiment); en particulier, tout module à un terme est donc de la forme \(c\cdot\alpha\), où \(\alpha\) désigne l'élément de base et où \(c\) parcourt toutes les grandeurs de \(\mathfrak J\). \emph{Un module dont les éléments appartiennent eux-mêmes à \(\mathfrak J\) est appelé un idéal dans \(\mathfrak J\);} d'où résulte le concept d'idéal fini, et en particulier d'idéal à un terme (idéal principal).

Tous les théorèmes de base pour les grandeurs entières et les idéaux reposent maintenant sur le théorème suivant:

\emph{Si \(M\) est un module relativement à \(\mathfrak J\), dont les éléments sont des formes linéaires en \(n\) indéterminées à coefficients dans \(\mathfrak J\), et si tout idéal de \(\mathfrak J\) est fini, alors \(M\) est aussi un module fini. Si en particulier tout idéal de \(\mathfrak J\) est un idéal principal, alors \(M\) possède une base de formes linéairement indépendantes.}\footnote{Le théorème apparaît dans la littérature séparément pour les divers domaines de coefficients \(\mathfrak J\). Pour les entiers rationnels, comparer par exemple Hilbert, Zahlbericht \S{} 3 et 4; pour les polynômes en plusieurs indéterminées: Hilbert, Über die vollen Invariantensysteme: fin du \S{} 2 (\emph{Math. Ann.} 42).}

En effet, si les éléments de \(M\) sont donnés par \(l(x)=a_1x_1+\cdots+a_nx_n\), alors la totalité des coefficients \(a_n\) parcourt un idéal de \(\mathfrak J\), qui, par hypothèse, est de la forme \(b_1a_n^{(1)}+\cdots+b_\varrho a_n^{(\varrho)}\). La forme linéaire appartenant à \(M\),
\[
m(x)=l(x)-b_1l^{(1)}(x)-\cdots-b_\varrho l^{(\varrho)}(x),
\]
ne dépend donc plus que de \((n-1)\) indéterminées; l'assertion en résulte après un nombre fini de répétitions. Si dans chaque cas \(\varrho=1\), alors la base consiste en \(n\) formes en respectivement \(n,\ n-1\ldots1\) indéterminées, ce qui donne l'indépendance linéaire des formes non nulles.

Le théorème sur le système fondamental en résulte dans les divers cas de la manière suivante:

Si le corps algébrique est obtenu par adjonction de l'élément algébrique entier \(\delta\) au corps de base, et si \(\mathfrak J\) désigne la totalité de toutes les grandeurs entières du corps de base, alors toute grandeur entière du corps supérieur admet la représentation (comparer par exemple Zahlbericht \S{} 3)
\[
\omega=\frac{a_1\delta^{\,n-1}+a_2\delta^{\,n-2}\cdots+a_n}{\Delta(\delta)},
\]
où les \(a_i\) sont des grandeurs de \(\mathfrak J\) et où \(\Delta(\delta)\) désigne le discriminant de \(\delta\). Par la substitution
\[
x_i=\frac{\delta^{\,n-i}}{\Delta(\delta)}
\]
le module de toutes les grandeurs entières est transformé en un module de formes linéaires; d'où résulte l'existence du système fondamental dans tous les cas où tout idéal de \(\mathfrak J\) est fini. Le même argument montre plus généralement que tout idéal dans le domaine des grandeurs entières du corps supérieur est fini. Or, en particulier, pour les entiers rationnels et pour les fonctions d'une indéterminée, tout idéal est un idéal principal; il s'ensuit donc que les systèmes fondamentaux dans les corps de nombres algébriques et dans les corps de fonctions algébriques d'une variable consistent en exactement \(n\) grandeurs, lorsque \(n\) désigne le degré. Et puisque, dans ces corps, tout idéal est fini d'après ce qui précède, on obtient en même temps le système fondamental pour les corps relatifs, qui consistera en général en plus de grandeurs que le degré relatif ne l'indique.\footnote{D'après les recherches de Steinitz sur les modules de formes linéaires, où \(\mathfrak J\) consiste en les entiers d'un corps de nombres algébriques (fini), \(\nu+1\) éléments de base suffisent lorsque \(\nu\) désigne le degré relatif (\emph{Math. Ann.} 71, \S{} 4).} En outre, si \(\mathfrak J\) désigne le domaine de tous les polynômes en plusieurs indéterminées, tout idéal de \(\mathfrak J\) est aussi fini;\footnote{Par le théorème de la base de Hilbert (\emph{Math. Ann.} 36). Les modules de polynômes ici appelés idéaux par souci de cohérence sont ordinairement appelés \og{}modules de formes\fg{} ou simplement \og{}modules\fg{}. Le théorème de Hilbert selon lequel un tel module \(N\) est engendré par un système fini peut aussi être ramené au théorème sur les formes linéaires. Soit en effet \(f\) un polynôme de degré \(n\) contenu dans \(N\), en les \(r+1\) indéterminées \(x_1,\ldots,x_r,x_{r+1}\), contenant le terme \(x_{r+1}^{\,n}\) (ce qui peut toujours être obtenu par une transformation linéaire); alors tous les polynômes de \(N\) sont congrus mod. \(f\) à ceux qui sont représentables sous la forme \(a_1(x)x_{r+1}^{\,n-1}+\ldots a_n(x)\); par conséquent, en posant \(x_{r+1}^{\,n-i}=X_i\), ils forment un module de formes linéaires pour lequel on peut supposer fini tout idéal dans le domaine des \(r\) indéterminées, puisque c'est le cas pour une indéterminée.} par suite, l'existence du système fondamental est aussi obtenue pour les corps de fonctions algébriques de plusieurs indéterminées, et il consistera de nouveau en plus de grandeurs que le degré ne l'indique.

\subsection*{3. Parallélisme et différences dans la théorie des idéaux des nombres et des fonctions algébriques}

La preuve du théorème fondamental de la théorie des idéaux, selon lequel tout idéal peut être représenté de façon unique comme un produit de puissances d'idéaux premiers, peut être menée conjointement pour les nombres et les fonctions algébriques; uniquement sur la base de la propriété que, pour les entiers rationnels et pour les fonctions d'une indéterminée, tout idéal est principal.\footnote{Comparer une remarque dans l'introduction de Hurwitz: Über die Theorie der Ideale (\emph{Gött. Nachr.} 1894). Dans le manuel d'algèbre de Weber, la preuve, suivant Kronecker, est menée parallèlement pour les deux corps au moyen de l'introduction de fonctionnelles. [Le lien entre la théorie de Kronecker et la théorie des idéaux est assuré par le \og{}théorème de Gauss étendu\fg{}; Hurwitz, loc. cit.]} Le point principal de la preuve consiste à montrer que, de la divisibilité d'un idéal \(\frc\) par un idéal \(\fra\) (c'est-à-dire du fait que tout élément de \(\frc\) est contenu dans \(\fra\)), résulte la représentation en produit \(\frc=\fra\cdot\frb\).\footnote{Souligné en particulier par Dedekind (théorie des nombres); Supplément XI, théorème VII, \S{} 177; comparer la note de la p. 554. Le théorème utilise le fait qu'il s'agit de corps algébriques; il ne vaut plus pour les idéaux (modules) de polynômes, où le plus petit commun multiple prend la place du produit.} Chez Hurwitz, cette preuve repose sur le \og{}théorème de Gauss étendu\fg{}, qui dit qu'une grandeur entière \(\omega\) qui divise chaque coefficient du produit de deux polynômes \(\varphi\) et \(\psi\) divise aussi le produit de tous les coefficients de \(\varphi\) et de \(\psi\); chez Dedekind, elle repose sur un théorème de modules équivalent à ce théorème.\footnote{Théorie des nombres, Suppl. XI; théorème VI, \S{} 173. Dedekind signale l'équivalence des deux théorèmes (Über die Begründung der Idealtheorie, \emph{Göttinger Nachr.} 1895).}

En poursuivant, on peut définir conjointement le concept de module complémentaire (Hensel no. 12), et à partir de lui l'idéal de ramification (idéal fondamental, différente) \(\mathfrak z\) (sa norme devient le discriminant du corps \(D\)); et le théorème fondamental vaut que l'idéal de ramification \(\mathfrak z\) est le \emph{plus grand commun diviseur de tous les idéaux principaux \(F'(\delta)\)}; ici \(\delta\) parcourt toutes les grandeurs entières du corps et \(F(\delta)=0\) désigne à chaque fois l'équation irréductible correspondante. Il s'ensuit que le \emph{discriminant du corps contient tous, et seulement, les nombres premiers rationnels, respectivement les facteurs linéaires en \(z\), qui sont divisibles par le carré d'un idéal premier}. Pour \(F'(\delta)\) lui-même on obtient la représentation
\[
F'(\delta)=\frf\cdot\frD,
\]
et, en prenant les normes,
\[
\Delta(\delta)=R^2\cdot D;
\]
ici \(\frf\) désigne le conducteur de l'anneau dérivé de \(\delta\) (ordre, domaine intègre), c'est-à-dire que \(\frf\) consiste en la totalité des grandeurs (entières) qui, après multiplication par une grandeur entière arbitraire, sont divisibles par le module \([1,\delta,\ldots\delta^{n-1}]\). \(R^2=\operatorname{Norm}\frf\) représente le carré du déterminant de substitution qui transforme la base \(1,\delta\ldots\delta^{n-1}\) en une base des grandeurs entières.\footnote{Dedekind, Über die Diskriminante endlicher Körper (\emph{Abhandl. der Gött. Ges. d. Wissensch.} vol. 29, 1882). Les définitions qui y sont données pour les nombres algébriques peuvent être transférées directement aux fonctions algébriques; les preuves aussi peuvent être transférées presque exactement. L'ordre de preuve dans D.--W. est différent: voir plus bas. Pour les théorèmes cités, comparer aussi Zahlbericht, \S{} 31 et 32. (\(F'(\delta)\) y est appelé la \og{}différente\fg{} de l'entier \(\delta\).)}

Les \emph{différences} dans le développement ultérieur sont conditionnées par des propriétés particulières des corps de base. Ainsi, dans le cas des nombres algébriques, le fait que les \emph{éléments du corps de base sont en même temps des exposants} conduit à la possibilité que l'idéal de ramification soit divisible par une puissance plus élevée que la \((e-1)\)-ième puissance d'un idéal premier \(\frp\) qui apparaît dans \(p\) à la \(e\)-ième puissance (à savoir lorsque \(e\) est divisible par \(p\)); cela conduit aux concepts supplémentaires de corps d'inertie et de corps de ramification,\footnote{Comparer Zahlbericht, \S{} 39--45.} qui n'ont pas d'analogue dans la théorie des fonctions algébriques. Mais surtout, en raison de l'\emph{existence d'une infinité de constantes} dans le corps de base, la finitude du nombre de classes d'idéaux disparaît pour les fonctions algébriques, et avec elle toutes les questions liées à la détermination du nombre de classes. Néanmoins, en un certain sens, le domaine de toutes les formes premières (transcendantes)\footnote{Klein, Zur Theorie der Abelschen Funktionen (\emph{Ann.} 36, 1890).} peut être considéré comme un analogue transcendant du corps de classes; car ici chaque point est engendré par une seule forme, un idéal principal.

D'autre part, l'existence d'une infinité de constantes dans le corps de base \(K(z)\) conduit, pour les fonctions algébriques, à des théorèmes et à des questions qui n'ont pas d'analogue pour les nombres algébriques. Il en résulte d'abord une certaine simplification dans l'agencement de la preuve; par exemple dans la preuve donnée par D.--W. de la décomposition unique des idéaux en idéaux premiers, où le \og{}théorème de Gauss étendu\fg{} ou un théorème qui lui est équivalent peut être évité en utilisant le fait qu'il existe une infinité de classes modulo toute forme linéaire \((z-c)\), à savoir toutes les constantes (tandis que le nombre de classes modulo un nombre premier est fini).\footnote{D.--W. \S{} 9. Comparer les notes des pp. 212 et 213.}

Le fait supplémentaire que tout polynôme à coefficients fonctions rationnelles entières de \(z\) se décompose en facteurs linéaires mod. \((z-\alpha)\) (ce qui n'est pas le cas pour un polynôme entier mod. \(p\)) conduit à des résultats véritablement nouveaux; ce fait subsiste encore si, au lieu de toutes les constantes, on ne prend que tous les nombres algébriques.\footnote{D.--W. notent dans l'introduction que dans ce cas toute la théorie demeure intacte.} Il s'ensuit notamment que tout idéal premier est un idéal du premier degré;\footnote{D.--W., \S{} 9, no. 7.} et en outre que le discriminant du corps \(D\) devient l'intersection de tous les discriminants \(\Delta(\delta)\), lorsque \(\delta\) parcourt toutes les grandeurs entières;\footnote{D.--W., p. 224.} ces deux faits sont en contraste avec les nombres algébriques. Il en résulte ensuite de nouveau que l'idéal de ramification est l'intersection de tous les idéaux principaux \(F'(\delta)\).

D'autres différences sont conditionnées par le fait que, pour les nombres algébriques, le corps de base \(R\) des nombres rationnels est quelque chose d'absolument distingué, à savoir le corps premier contenu dans le corps (c'est-à-dire le corps dérivé de l'unité); tandis que le corps de base \(K(z)\) est quelque chose de relatif. Toute fonction non constante \(\xi\) peut être choisie comme variable indépendante; le corps devient alors un corps algébrique sur \(K(\xi)\). Le théorème selon lequel tout idéal premier \(\frp\) est un idéal du premier degré, donc toute fonction est congrue modulo \(\frp\) à une constante, combiné avec la possibilité de prendre une fonction arbitraire comme variable indépendante, conduit au concept le plus important qui dépasse les nombres algébriques, le concept de point sous forme purement arithmétique. (Comparer Hensel 10a.) \og{}Chaque point est ici représenté par un corps de constantes isomorphe au corps de fonctions.\fg{}\footnote{Deux corps sont, comme on sait, dits \og{}isomorphes\fg{} s'ils peuvent être mis en correspondance biunivoque de façon que somme, différence, produit et quotient se correspondent aussi.} La totalité des points ainsi définis forme la \og{}surface de Riemann absolue\fg{}. Cette définition du point ne dépend que du corps, non d'une variable spéciale. La preuve d'existence, toutefois, est effectuée en distinguant une certaine fonction comme variable indépendante \(z\); un point \(P\) est engendré par un idéal premier \(\frp\) en assignant à chaque fonction son reste constant modulo \(\frp\); ainsi, en particulier, les fonctions divisibles par \(\frp\) s'annulent en \(P\), et réciproquement. Lorsque \(\frp\) parcourt tous les idéaux premiers en \(z\), on obtient ainsi tous les points où \(z\) est fini; les points restants apparaissent par l'introduction de la variable indépendante \(\zeta=\frac{1}{z}\) et correspondent aux idéaux premiers qui divisent l'idéal principal \(\xi\). Naturellement, les concepts attachés au concept de point ne peuvent pas non plus avoir d'analogue pour les nombres algébriques; tels le concept de polygone, de classe de polygones, de dimension d'une classe de polygones, etc.

Il faut encore remarquer que le libre choix de la variable indépendante conduit à une autre interprétation de la décomposition \(F'(\vartheta)=\frf\cdot\mathfrak z\). Si l'on considère en effet \(\vartheta\) comme variable indépendante, on obtient en conséquence
\[
\left(\frac{\p F}{\p z}\right)=\frf\cdot\fra,
\]
puisque l'anneau dérivé de \(\vartheta\) avait été défini symétriquement par rapport à \(z\) et \(\vartheta\); et \(\frf\) devient le plus grand commun diviseur des deux idéaux principaux \(\frac{\p F}{\p z}\) et \(\frac{\p F}{\p\vartheta}\); le conducteur \(\frf\) devient en même temps l'\og{}idéal du point double\fg{} relativement à \(\vartheta,z\).\footnote{D.--W. p. 263.}

\subsection*{4. Rapport entre le fondement arithmétique du développement en séries et la théorie des idéaux}

Chez Hensel--Landsberg, l'auxiliaire transcendant du développement en séries et le concept de point emprunté à la théorie des fonctions prennent la place de la théorie des idéaux. Le fondement arithmétique récemment donné par Hensel pour le développement en séries des nombres et des fonctions algébriques\footnote{Eine neue Theorie der algebraischen Zahlen (\emph{Math. Zeitschr.} vol. 2, 1918) et Neue Begründung der arithmetischen Theorie der algebraischen Funktionen einer Variablen (\emph{Math. Zeitschr.} vol. 4, 1919). Pour un exposé plus détaillé, comparer Hensel no. 44 sq.} peut donc être regardé comme un équivalent de la théorie des idéaux. En fait, ce fondement arithmétique revient à passer, pour chaque idéal premier, à un tel corps d'extension (transcendant) dans lequel cet idéal premier devient un idéal principal, de sorte que les théorèmes ordinaires de décomposition valent; et il suffit d'effectuer cette considération pour tous les idéaux premiers apparaissant dans un nombre premier \(p\), respectivement dans une fonction linéaire \(z-a\).

On introduit d'abord un corps d'extension transcendant formé par tous les développements en séries selon les puissances entières de \(z-a\), respectivement de \(p\) (nombres \(p\)-adiques), avec dans chaque cas au plus un nombre fini de puissances négatives. Dans ce corps, l'équation irréductible qui définit le corps de nombres algébriques ou le corps de fonctions se décompose en le produit de \(\varrho\) facteurs irréductibles, correspondant à la décomposition de \(p\), respectivement de \(z-a\), en le produit des \(\varrho\) idéaux \og{}primaires\fg{} \(\mathfrak q_1,\ldots\mathfrak q_\varrho\). La représentation ultérieure de chaque idéal primaire comme puissance d'un idéal premier,
\[
\mathfrak q_i=\frp_i^{\,l_i},
\]
correspond chez Hensel à l'introduction d'un autre corps, algébrique relativement au corps d'extension transcendant. Et, dans le cas des fonctions algébriques, ce corps peut être défini par une équation pure; dans le cas des nombres algébriques, plus généralement, par une équation algébriquement résoluble. La simplification pour les fonctions algébriques tient au fait qu'ici n'interviennent pas de corps d'inertie et de ramification (comparer no. 3).

La \og{}série terminée\fg{} en un élément premier \(\pi\)\footnote{Neue Begründung \(\ldots\), formule (11).} qui intervient dans ces recherches (et qui n'exige pas encore de preuve transcendante de convergence)
\[
v=c_0+c_1\pi+\ldots+c_{\sigma-1}\pi^{\sigma-1}-v_\sigma\pi^\sigma,
\]
où \(c_0\ldots c_{\sigma-1}\) sont des constantes et où \(v_\sigma\) désigne un élément entier du corps, se trouve sous une forme exactement correspondante dans D.--W.\footnote{D.--W., p. 231 et \S{} 15.} Elle y est déduite de la théorie des idéaux et donne la base de la définition de l'ordre d'annulation (respectivement de devenir infini) d'une fonction en un point.
\clearpage
\subsection*{5. Comparaison de la théorie arithmétique des fonctions algébriques avec la théorie géométrique. Concept d'invariance différent}

Est \og{}invariant\fg{} tout concept caractérisé uniquement par le corps. Dans la théorie arithmétique, cette invariance est en général déjà donnée dans la définition, laquelle est énoncée par rapport à tous les éléments du corps, et non par rapport à une base quelconque ou à une variable distinguée; le concept de point indiqué plus haut en donne un exemple. En revanche, les autres théories partent d'une variable ou d'une base particulière et montrent après coup l'indépendance du concept à l'égard de ce choix particulier; l'invariance y devient une invariance vis-à-vis des transformations birationnelles.

Supposons en effet que le corps soit obtenu, d'une part, par adjonction de l'élément algébrique \(s\) à \(K(z)\), et, d'autre part, par adjonction de \(\sigma\) à \(K(\zeta)\), où \(f(z,s)=0\), respectivement \(F(\zeta,\sigma)=0\), désignent les équations irréductibles de \(s\), respectivement de \(\sigma\); ou encore plus généralement par adjonction d'un nombre fini d'éléments à \(K(\eta)\), avec les équations irréductibles correspondantes
\[
g_1(\eta,\tau_1)=0,\qquad g_2(\eta,\tau_1,\tau_2)=0,\ldots .
\]
Alors, par définition, tous les éléments peuvent s'exprimer rationnellement au moyen de \(z\) et \(s\), aussi bien qu'au moyen de \(\xi\) et \(\sigma\), et aussi au moyen de \(\eta,\tau_1,\tau_2\ldots\). En particulier, \(z\) et \(s\) sont donc exprimables rationnellement par \(\xi\) et \(\sigma\), et réciproquement; de même \(z\) et \(s\) le sont rationnellement par \(\eta,\tau_1,\tau_2\ldots\), et réciproquement. La courbe plane \(f(z,s)=0\) passe donc, par \og{}transformation birationnelle\fg{}, dans la courbe plane \(F(\zeta,\sigma)=0\), ou encore dans la courbe de l'espace des variables \(\eta,\tau_1,\tau_2\ldots\) définie par \(g_1(\eta,\tau_1)=0; g_2(\eta,\tau_1,\tau_2)=0,\ldots\). L'invariance se manifeste donc ainsi: la définition énoncée pour une courbe spéciale demeure valable pour toutes les courbes qui en résultent par transformation birationnelle (dans un espace d'un nombre arbitraire de dimensions); pour toutes les courbes de la \og{}classe\fg{} au sens de Riemann. Ici, un \og{}point\fg{} est défini comme un système de valeurs associé \(z=a,s=b\), de sorte que \(f(a,b)=0\); par transformation birationnelle, ce point passe en général de nouveau dans un système de valeurs associé \(\alpha,\beta\) de \(F(\zeta,\sigma)=0\), mais il peut, dans des cas particuliers, lorsque \(a,b\) était un point \og{}singulier\fg{}, correspondre à plusieurs \og{}points\fg{} de \(F(\zeta,\sigma)=0\); et l'on peut toujours indiquer des courbes sur lesquelles un point \og{}singulier\fg{} de \(f(z,s)=0\) correspond à un nombre fini de points simples, séparés ou voisins (résolution des singularités).\footnote{Comparer par exemple Brill--Noether, 6e section. Des précisions sur les points \og{}singuliers\fg{} sont données dans le numéro suivant.} De même, un groupe de points passe de nouveau dans un groupe de points; et le nombre de ses points est invariant dès que l'on compte un point singulier pour autant de points que de points simples lui correspondent sous une transformation appropriée. Puisque toute courbe peut être transformée birationnellement en une courbe n'ayant que des points \og{}multiples ordinaires\fg{},\footnotemark[\value{footnote}] la théorie géométrique développe ses concepts seulement pour de telles courbes spéciales et montre ensuite l'invariance vis-à-vis de toute transformation birationnelle, même de celle qui conduit aux points singuliers les plus généraux.

En particulier, comme courbe distinguée de la \og{}classe\fg{} --- le cas hyperelliptique excepté --- on peut introduire la \og{}courbe normale des \(\varphi\)\fg{} dans l'espace de \((p-1)\) dimensions, qui ne possède pas de points singuliers; ici les \(\varphi\) sont les \(p\) \og{}courbes adjointes d'ordre \((n-3)\)\fg{} pour \(f(z,s)=0\) rendu homogène. À une transformation birationnelle de \(f(z,s)=0\) en \(F(\zeta,\sigma)=0\) correspond une transformation linéaire des \(\varphi\), de sorte qu'il s'agit ici d'invariance au sens projectif. La représentation de toutes les fonctions du corps comme fonctions rationnelles des \(\varphi\) est appelée \og{}représentation invariante\fg{} dans la théorie géométrique.\footnote{Comparer par exemple Brill--Noether, 8e section. Pour les fonctions algébriques de plusieurs variables, la question de savoir si l'invariance peut être ramenée à une invariance vis-à-vis des transformations linéaires n'a pas été traitée.}

Le faisceau linéaire des \(\varphi\), ou encore des \og{}groupes spéciaux\fg{} qu'elles découpent, passe donc dans lui-même sous toute transformation birationnelle; il est par conséquent caractérisé uniquement par le corps. Il correspond à la classe différentielle dans la théorie arithmétique, où l'invariance se montre de nouveau directement dans la définition, laquelle est donnée par des propriétés relatives à chaque fonction du corps comme variable indépendante.

\subsection*{6. Suite de la comparaison. Points singuliers}

C'est à la différence de conception du concept de point que se rattache le fait que les \og{}points singuliers\fg{} des courbes, qui causent des difficultés particulières dans la théorie algébro-géométrique, n'apparaissent pas du tout dans la fondation de la théorie arithmétique; et, par suite, que les \og{}fonctions adjointes\fg{} ne sont pas requises pour construire la théorie arithmétique, comme c'est le cas dans la théorie géométrique, mais appartiennent à des parties spéciales et plus élevées de la théorie. D'autre part, les points de ramification n'apparaissent pas dans la fondation de la théorie géométrique; en fait, ils disparaissent aussi dans la représentation des intégrales sous la forme normale d'Aronhold (Hensel no 27, formule 80).

La définition du point singulier donnée en 5 peut être formulée ainsi: \(f(z,s)=0\) (respectivement une courbe dans un espace de dimension supérieure) possède en \(z=a,s=b\) (respectivement \(z=a; s_i=b_i\)) un point singulier si ce système de valeurs est pris par plus d'un point de la surface de Riemann absolue (point au sens arithmétique), ou s'il est pris en un ou plusieurs points avec un ordre plus élevé. Dans ce dernier cas, on a affaire à un point du polygone de ramification de \(z\) aussi bien que de \(s\); cela correspond aux \og{}points consécutifs\fg{} de la théorie géométrique; le point devient un point de rebroussement (supérieur) de \(f(z,s)=0\) (respectivement de la courbe dans l'espace supérieur).

Comme, par hypothèse, toute fonction arbitraire \(\eta\) du corps peut être représentée rationnellement au moyen de \(z\) et \(s\), tout point de la surface de Riemann absolue où \(z=a,s=b\) est obtenu, en vertu de l'isomorphie, en assignant à chaque fonction \(\eta\), représentée par \(z\) et \(s\), sa valeur pour \(z=a,s=b\). Pour que le point soit singulier, il doit donc exister au moins une fonction \(\eta\) qui, pour \(z=a,s=b\), prend plusieurs valeurs ou prend plusieurs fois un système de valeurs. Et cela, à cause de l'exprimabilité rationnelle au moyen de \(z\) et \(s\), ne peut avoir lieu que si, dans cette représentation, le numérateur et le dénominateur s'annulent simultanément pour \(z=a,s=b\). Il s'ensuit que la définition ci-dessus est identique à la définition usuelle, selon laquelle
\[
f,\quad \frac{\p f}{\p z},\quad \frac{\p f}{\p s}
\]
s'annulent pour \(z=a,s=b\). Car dans ce cas
\[
\eta=\frac{\frac{\p f}{\p z}}{\frac{\p f}{\p s}}
\]
est une fonction du genre indiqué qui est multivaluée pour \(z=a,s=b\), tandis que, dans le cas contraire, toute fonction, même si elle devient \(0/0\), ne prend qu'une seule valeur.

Avec le concept arithmétique de point, le point singulier ne peut pas apparaître, puisque deux points sont considérés comme distincts dès qu'une seule fonction se voit assigner des systèmes de valeurs différents. Mais même dans la théorie des idéaux qui sert à fonder le concept arithmétique de point, le concept de point singulier n'intervient pas, précisément parce que l'on considère toujours simultanément la totalité de toutes les grandeurs entières du corps. Tous les points singuliers situés dans le fini de \(F(z,\delta)=0\), où \(\delta\) est algébriquement entier relativement à \(z\), sont en effet engendrés (d'après la deuxième définition et la fin du no 3) par l'\og{}idéal du point double\fg{} \(\frf\). Or, puisque, d'après le théorème fondamental, l'idéal de ramification \(\mathfrak z\) est le plus grand commun diviseur de tous les idéaux principaux \(F'(\delta)\), on peut, pour chaque idéal premier \(\frp\), indiquer un \(\delta\) tel que le \(\frf\) correspondant ne soit pas divisible par \(\frp\); par suite --- puisque \(\frf\) est le plus grand commun diviseur de \(F'(\delta)\) et de \(F'(z)\) --- au moins l'un des deux idéaux principaux \(F'(\delta)\) et \(F'(z)\) n'est pas divisible par \(\frp\). Il n'existe donc \emph{aucun système de valeurs pour lequel tous les \(F'(\delta)\) et \(F'(z)\) s'annulent simultanément}, et par conséquent le système de valeurs \(z=a,\delta_i=b_i\) n'est pris qu'en un seul point de la surface de Riemann absolue; la \emph{totalité des fonctions algébriques entières du corps} (que l'on peut interpréter comme une courbe dans l'espace d'une infinité de dimensions) \emph{ne possède aucun point singulier dans le fini}; or seuls les valeurs finies entrent en considération dans la théorie des idéaux.

Mais les fonctions de base \(\omega_1\ldots\omega_n\) de toutes les grandeurs entières ne possèdent pas non plus de points singuliers dans le fini. D'après ce qui précède, tout point au sens arithmétique est déjà déterminé de façon univoque par un système de constantes assigné isomorphiquement à toutes les fonctions entières. Si donc une \og{}courbe spatiale\fg{} définie par certaines fonctions entières \(s_i\) possède, pour \(s_i=\alpha_i\), un point singulier dans le fini, il doit exister au moins une fonction entière \(\eta\) qui, pour \(s_i=\alpha_i\), prend plusieurs systèmes de valeurs. Si maintenant \(\omega_1,
\ldots,\omega_n\) désigne une base de toutes les grandeurs entières, alors, puisque \(\delta=a_1\omega_1\ldots+a_n\omega_n\) avec des \(a_i\) entiers rationnels, à chaque système de valeurs des \(\omega\) ne correspond qu'un seul système de valeurs de \(\delta\); par conséquent la courbe spatiale définie par \(\omega_1,
\ldots,\omega_n\) ne peut posséder aucun point singulier dans le fini; et \emph{chaque point de la surface de Riemann absolue où \(z\) est fini est déjà défini de façon univoque par un système de constantes assigné isomorphiquement aux \(\omega\). Les fonctions de base sont précisément les fonctions \(\eta\) qui deviennent multivaluées pour \(s_i=\alpha_i\).} Par l'introduction de la base des grandeurs entières, par exemple à la place d'une base \(1,\delta\ldots\delta^{n-1}\), la difficulté des points singuliers est donc contournée.

Dans la théorie géométrique, la tâche d'engendrer univoquement tous les points de la surface de Riemann absolue --- plus généralement, de former toutes les fonctions qui deviennent infinies en des points donnés --- est accomplie au moyen des formes adjointes et de leurs quotients. Une forme \(\psi\) est dite adjointe si \(\psi=0\) possède au moins un point \((i-1)\)-uple en chaque point \(i\)-uple de la courbe fondamentale \(f(z,s)=0\); un point singulier supérieur doit d'abord être résolu en points multiples ordinaires. Il a été montré plus haut que, pour la fonction la plus générale \(\eta=\frac{\psi}{\chi}\), le numérateur et le dénominateur doivent de toute façon s'annuler en un point singulier de \(f(z,s)=0\); qu'ils doivent être adjoints, cela ressort de l'examen des différentielles partout finies. Inversement, que les conditions d'adjonction soient également suffisantes pour former la fonction la plus générale, c'est ce que montre le \og{}théorème des résidus\fg{}, sur lequel on reviendra plus précisément. Les fonctions adjointes prennent donc la place des fonctions de base de la théorie arithmétique.

Arithmétiquement, à une forme adjointe correspond le numérateur d'une fonction du corps divisible par l'\og{}idéal du point double\fg{} \(\frf\), dès que l'on suppose que les points singuliers --- ce qui peut toujours être obtenu par transformation --- sont tous situés dans le fini. De là se montre formellement la représentation des grandeurs de base \(\omega_1,
\ldots,\omega_n\) comme quotients de formes adjointes, et par suite le rapport entre les différentes façons de les représenter, comme suit:

Soit \(R(z)\) le déterminant de substitution qui transforme une base \(1,\delta\ldots\delta^{n-1}\) en une base \(\omega_1,
\ldots,\omega_n\). Réciproquement, chaque \(\omega\) s'exprime alors comme quotient d'une fonction entière de \(z\) et \(\delta\) par \(R(z)\):
\[
\omega_i=\frac{G_i(z,\delta)}{R(z)}.
\]
Comme \(R(z)\omega_i\) appartient à l'anneau dérivé de \(\delta\), il résulte de la définition du conducteur que \(R(z)\) est divisible par \(\frf\) [de \((R(z))^2=\operatorname{Norm}\frf\), on ne peut conclure cette divisibilité que pour \((R(z))^2\)]. Puisque maintenant \(\omega_i\), en tant que fonction entière, reste finie pour tout \(z\) fini, la fonction numérateur \(G_i(z,\delta)\) doit aussi être divisible par \(\frf\); les numérateurs et les dénominateurs de tous les \(\omega_i\) sont donc adjoints. Cette représentation des \(\omega\) donne alors celle de toutes les fonctions du corps par quotients d'adjointes.

D'après ce qui précède, la théorie géométrique peut être caractérisée comme la \emph{théorie de l'anneau dérivé d'une grandeur \(\delta\)} (ordre), par opposition à la théorie arithmétique, qui considère \emph{toutes les grandeurs entières} simultanément. Il est donc clair que le conducteur de l'anneau --- les points singuliers --- joue un rôle décisif. On peut encore indiquer d'autres analogies entre la théorie arithmétique des anneaux (ordres) et la théorie géométrique. Ainsi, la théorie géométrique ne compte nulle part les points d'intersection qui tombent dans les points singuliers lors de l'intersection avec des courbes adjointes, mais ne considère que les groupes de points qui en sont distincts. À cela correspond le fait que Dedekind\footnote{Über die Anzahl der Idealklassen in den verschiedenen Ordnungen eines endlichen Körpers. (Festschrift Braunschweig 1877); \S{} 3--5. (Les théorèmes énoncés dans ces paragraphes pour les nombres algébriques valent littéralement pour les fonctions algébriques d'une variable.)} ne considère que les idéaux d'anneau qui sont premiers au conducteur \(\frf\), et qu'il établit, pour ce domaine d'idéaux, toutes les lois de divisibilité, y compris la décomposition unique en idéaux premiers.

Le \og{}théorème fondamental des fonctions algébriques\fg{}\footnote{M. Noether, Über einen Satz aus der Theorie der algebraischen Funktionen. \emph{Math. Ann.} vol. 6 (1873).} qui est à la base de la théorie géométrique, ou du moins une conséquence immédiate de celui-ci, trouve aussi, en un certain sens, un analogue dans la théorie arithmétique des anneaux. Ce théorème donne les conditions d'une représentabilité
\[
\psi(z,s)=A(z,s)f(z,s)+B(z,s)\varphi(z,s),
\]
où toutes les grandeurs qui interviennent sont entières rationnelles en \(s,z\) (plus généralement entières et homogènes en \(x_1x_2x_3\)). De ces conditions, il s'ensuit,\footnote{Brill--Noether, no 55.} en vertu de \(f(z,s)=0\), que le quotient \(\frac{\psi}{\varphi}\) est toujours égal à une fonction entière rationnelle \(B(z,s)\) lorsqu'il est adjoint à \(f\). À cela correspond le théorème du no 3, qui est à la base de toutes ces comparaisons: le conducteur \(\frf\) d'un anneau dérivé de \(\delta\) est égal à l'idéal du point double de \(f(z,\delta)=0\). Car, par la définition du conducteur, toute grandeur algébrique entière divisible par \(\frf\) appartient à l'anneau, et elle est donc entière et représentable rationnellement par \(z\) et \(\delta\); tandis que le théorème cité dit qu'une telle grandeur est une fonction adjointe.\footnote{On suppose ici constamment, à la suite de Dedekind, que \(\delta\) est une grandeur algébrique entière (ce que l'on peut toujours obtenir par transformation). Le cas correspondant plus exactement à la géométrie, \(f(z,s)=0\), où \(s\) peut aussi dépendre algébriquement de façon fractionnaire de \(z\), est traité par Hensel--Landsberg; comparer Hensel nos 26 et 29.}
\clearpage
\subsection*{7. Comparaison entre le théorème des résidus et les théorèmes de la théorie des idéaux}

La conséquence qui vient d'être mentionnée du \og{}théorème fondamental des fonctions algébriques\fg{} conduit directement, dans la théorie géométrique, au \og{}théorème des résidus\fg{}, et ainsi au fondement sur lequel toute la théorie géométrique est construite. Toutefois, à un certain égard, le théorème des résidus a de nouveau un caractère plus élémentaire que les autres fondements, dans la mesure où il peut être formulé arithmétiquement comme conséquence directe de la décomposition unique des idéaux en idéaux premiers, ou plutôt des faits qui sont à la base de ce théorème de décomposition, et ne repose donc pas sur la théorie supérieure des anneaux. On arrive à cette formulation arithmétique comme suit.

Supposons que la courbe de base \(f(z,\delta)=0\), comme on peut toujours l'obtenir par une transformation linéaire fractionnaire des variables, soit telle que les points singuliers et les groupes de points finis que l'on considère soient situés dans le fini; et supposons en outre que \(\delta\) dépende algébriquement et intégralement de \(z\), ce que l'on peut toujours obtenir par une transformation linéaire ultérieure entière dans les variables. Alors à un idéal premier \(\frp\) engendrant un point \(P\) de la surface de Riemann absolue, en lequel \(z=a,\delta=b\), correspond le point \(z=a,\delta=b\) de \(f(z,\delta)=0\); plus généralement, un idéal \(\fra=\frp_1^{\,e_1}\cdots\frp_\varrho^{\,e_\varrho}\) engendre le groupe de points constitué des points \(z=a_i,\delta=b_i\) avec la multiplicité correspondante; et tous les points de \(f(z,\delta)=0\) qui sont situés dans le fini sont ainsi engendrés. Comme, en un point \(P\), toutes les fonctions \(g_1,g_2\ldots\) divisibles par \(\frp\) s'annulent, et comme, réciproquement, cette totalité de fonctions \(g\) engendre l'idéal \(\frp\), le point correspondant sur \(f\) est l'intersection de
\[
g_1(z,\delta)=0,\quad g_2(z,\delta)=0,\ldots
\]
avec \(f=0\); il en va de même pour le groupe de points correspondant à un idéal \(\fra\). En fait, d'après le théorème selon lequel tout idéal est le plus grand commun diviseur de deux idéaux principaux, deux telles courbes suffisent toujours à découper un groupe de points. En particulier, le groupe de points correspondant à un idéal principal est découpé par une courbe \(g(z,\delta)=0\), et réciproquement; l'idéal principal correspond à l'intersection complète.

Deux groupes de points \(A\) et \(B\) s'appellent \emph{corésiduels} dans la théorie géométrique s'ils peuvent être complétés par le même groupe de points \(R\) en une intersection complète. Si ces groupes de points sont engendrés respectivement par les idéaux \(\fra,\frb,\frr\), on a entre ces idéaux la relation
\[
\fra\cdot\frr=(\alpha),\qquad \frb\cdot\frr=(\beta),
\]
qui montre que les idéaux \(\fra\) et \(\frb\) sont équivalents. Réciproquement, l'équivalence de deux idéaux, \(\fra=\eta\frb\), implique toujours cette relation en vertu du théorème selon lequel tout idéal peut être transformé en un idéal principal par multiplication par un autre idéal;\footnote{Comparer Dedekind, théorie des nombres \S{} 181; si l'on choisit \(\frr\) de sorte que \(\fra\frr\) devienne un idéal principal \((\alpha)\), alors \(\frb\frr=\eta\cdot(\alpha)\); et, comme \(\frb\frr\) ne contient que des grandeurs entières, il en va de même de \(\eta\cdot(\alpha)\), qui devient donc lui aussi un idéal principal \((\beta)\).} \emph{les idéaux équivalents et les groupes de points corésiduels sont donc des concepts identiques.}

Le théorème des résidus dit maintenant: \emph{si \(A\) et \(B\) sont corésiduels, alors ils peuvent être découpés par des courbes adjointes ayant le même résidu, en dehors des points singuliers.} Comme les zéros et les pôles d'une fonction sont toujours corésiduels, cela démontre en même temps que les fonctions les plus générales du corps sont représentables comme quotients de formes adjointes. Dans le langage de la théorie des idéaux, le théorème des résidus revient simplement à dire que, outre \(\fra\) et \(\frb\), les idéaux \(\fra\frf\) et \(\frb\frf\) sont eux aussi équivalents, où \(\frf\) désigne l'idéal des points doubles. Alors, d'après ce qui précède, il existe toujours un idéal \(\frm\) tel que
\[
\fra\frf\frm=(\alpha),\qquad \frb\frf\frm=(\beta)
\]
deviennent des idéaux principaux; \(\alpha=0\) et \(\beta=0\) sont les courbes adjointes qui découpent les groupes de points \(A\) et \(B\), respectivement leurs sous-groupes \(A'\) et \(B'\) distincts des points singuliers.\footnote{Si \(\fra=\frf'\fra_1,\ \frb=\frf'\frb_1,\ \frf=\frf'\frf_1\), où \(\fra_1,\frb_1,\frf_1\) sont premiers entre eux, alors déjà \(\fra_1\frf\mathfrak n,\ \frb_1\frf\mathfrak n\) deviennent des idéaux principaux, où \(\mathfrak n\) est premier à \(\frf\). Car tout idéal peut être transformé en un idéal principal par multiplication par un idéal premier à un idéal donné (D.--W. p. 214; ou théorie des nombres, \S{} 178, théorème XI).} La preuve du théorème des résidus est donc contenue arithmétiquement dans la possibilité d'assigner des points aux idéaux premiers, et dans la preuve que des idéaux équivalents peuvent être transformés en idéaux principaux par multiplication par un seul et même idéal.

Des groupes de points corésiduels n'ont pas besoin de contenir le même nombre de points, de même que des idéaux équivalents n'ont pas besoin d'avoir le même degré. Ainsi, par exemple, tous les idéaux principaux sont équivalents, et tous les groupes de points engendrés par des intersections complètes sont corésiduels. Cette différence par rapport aux polygones équivalents, c'est-à-dire aux diviseurs,\footnote{Un polygone \(A\) se compose d'un nombre fini de points de la surface de Riemann absolue; deux polygones \(A\) et \(B\) sont dits équivalents s'ils sont les zéros et les pôles d'une fonction \(\eta\); \(\eta\) admet alors la représentation par quotients de polygones \(\eta=\frac{A}{B}\). (D.--W. \S{} 17 et 18.)} s'explique par le fait que les points où \(z\) devient infini ne sont pas engendrés par des idéaux premiers en \(z\). Dans la représentation d'une fonction comme quotient d'idéaux \(\eta=\frac{\fra}{\frb}\), qui exprime l'équivalence de \(\fra\) et de \(\frb\), les zéros ou les pôles de \(\eta\) où \(z\) devient infini n'apparaissent donc pas. Par exemple, l'idéal principal \((\alpha)\) devient le produit des idéaux premiers qui correspondent aux zéros de la fonction entière \(\alpha\); les pôles de \(\alpha\) n'interviennent pas, puisque \(\alpha\) ne devient infini qu'aux points où \(z\) devient infini. En ce sens, la théorie des idéaux est l'analogue exact de la théorie des formes dans la théorie géométrique, où il n'y a de même que des zéros et pas de pôles.

Le passage des formes aux fonctions s'effectue géométriquement par formation de quotients de formes de même ordre; arithmétiquement, cela correspond au passage des classes d'idéaux aux classes de polygones, ou de diviseurs, définies de manière invariante. Car, comme aucune variable n'y est plus distinguée, la représentation d'une fonction par quotients de polygones donne tous les zéros et tous les pôles. Le rapport de la classe d'idéaux, classe de groupes de points corésiduels, à la classe de polygones est le suivant: si \(\calA\) et \(\calB\) sont les polygones engendrés par les idéaux \(\fra\) et \(\frb\), et si \(\calN\) est le polygone des points où \(z\) devient infini, c'est-à-dire le polygone dénominateur de \(z\) et des fonctions entières de \(z\), alors de l'équivalence de \(\fra\) et de \(\frb\) résulte l'équivalence de \(\calA\) et de \(\calB\) si et seulement si \(\fra\) et \(\frb\) ont le même degré; en général il n'en résulte que l'équivalence de \(\calA\calN^\nu\) avec \(\calB\calN^{\nu'}\), \((\nu\ne\nu')\). La division des idéaux, respectivement des groupes de points corésiduels, en classes est donc un regroupement d'une infinité de classes de polygones en une seule classe; elle peut aussi être conçue comme une division des polygones en classes modulo une puissance de \(\calN\).\footnote{Hensel--Landsberg, 25e leçon, p. 438. --- Si l'on représente les deux variables \(z\) et \(s\) par quotients de polygones, cela peut être compris comme une homogénéisation binaire de chacune de ces variables; si l'on choisit en particulier \(z\) et \(s\) de sorte qu'ils aient le même polygone dénominateur, on obtient l'homogénéisation ternaire usuelle en géométrie.}

Le théorème des résidus, pour le cas de groupes de points corésiduels composés de nombres différents de points, n'intervient que dans des questions purement géométriques. Dans l'application aux fonctions algébriques, c'est en fait seulement le cas correspondant aux classes de polygones qui intervient, celui des groupes de points corésiduels de même nombre. Le \og{}système linéaire complet\fg{} engendré par un groupe de points, c'est-à-dire la totalité des groupes de points corésiduels à ce groupe de points et ayant le même nombre de points, correspond exactement à la classe de polygones engendrée par le polygone correspondant; les concepts de \og{}somme de deux systèmes complets\fg{} et de \og{}produit de deux classes de polygones\fg{} coïncident eux aussi. Qu'un tel système complet ne contienne qu'un nombre fini de groupes de points linéairement indépendants est une conséquence directe du théorème des résidus, qui ramène la dimension du système, c'est-à-dire le nombre de groupes de points linéairement indépendants, à la dimension de toutes les courbes adjointes, d'ordre arbitraire mais fixe, passant par le groupe résiduel. Pour parvenir à la détermination effective de cette dimension, au théorème de Riemann--Roch, la théorie géométrique déduit d'abord du théorème des résidus le \og{}théorème de réduction\fg{}, ou théorème sur les points fixes;\footnote{Brill--Noether, no. 58.} puis, de celui-ci, le théorème de Riemann--Roch.

De façon tout à fait correspondante, la théorie arithmétique montre que chaque classe de polygones est de dimension finie, et détermine sa dimension par une base spéciale des grandeurs entières, une base normale définie au moyen du module complémentaire; elle obtient ainsi le théorème de Riemann--Roch comme conséquence du lien entre le module complémentaire et le polygone de ramification. Dans l'exposé de D.--W. intervient aussi le théorème de réduction,\footnote{D.--W., \S{} 30, no. 2.} mais non pas comme fondement, ainsi que dans la théorie géométrique. En conséquence de la correspondance mentionnée au no. 6 entre base et fonctions adjointes, on peut donc constater ici aussi une certaine correspondance des moyens auxiliaires dans les deux théories. En résumé, on peut dire que les formations de concepts dans les deux théories se recouvrent pour l'essentiel, même si, abstraction faite de la formulation, elles diffèrent par leur disposition. Les deux théories sont nées indépendamment l'une de l'autre, la théorie géométrique une décennie plus tôt. Il faut encore remarquer que le théorème de Riemann--Roch apparaît à l'origine, chez Riemann et Roch et aussi chez Weierstrass, comme théorème de rang; il s'agit du rang de la matrice \((\varphi_i(x_\varkappa))\), où les \(\varphi\) parcourent les \(p\) courbes adjointes d'ordre \(n-3\), et les \(x\) les points du groupe résiduel.

\subsection*{8. Questions transcendantes pour les fonctions algébriques et les nombres algébriques}

Hilbert a particulièrement signalé des analogies entre les questions transcendantes relatives aux fonctions algébriques et celles relatives aux nombres algébriques;\footnote{Mathematische Probleme, problème 12. (\emph{Göttinger Nachrichten} 1900).} il s'agit essentiellement de questions d'existence. Ainsi la question de Riemann concernant l'existence d'une fonction algébrique pour une surface de Riemann donnée, donc aussi pour des points de branchement donnés, correspond à la question de l'existence de corps de nombres algébriques de groupe et de discriminant prescrits.\footnote{Le traitement des fonctions algébriques ainsi suggéré, à partir du groupe, a été effectué par R. König comme cas spécial de questions plus générales; comparer surtout: Riemannsche Funktionen- und Differentialsysteme in der Ebene (\emph{J. f. M.} vol. 148) et \emph{Math. Ann.} 78, 79. À la place du corps algébrique intervient chez König un module fini spécial relativement à \(K(z)\): il n'y a qu'une \og{}propriété de classe\fg{}.} Alors que, dans le cas des fonctions algébriques, la preuve d'existence peut être menée en général, elle n'a réussi dans le domaine des corps de nombres algébriques que pour des corps de base spéciaux, les nombres rationnels et les corps quadratiques, et seulement pour des groupes abéliens.\footnote{Seul le cas simple des groupes qui apparaissent dans les équations du troisième et du quatrième degré, ainsi que quelques groupes tout à fait spéciaux pour des équations du degré 8, peut être traité pour des corps de base arbitraires: G. Bucht, Über einige algebraische Körper 8. Grades (\emph{Arkiv f. Math., Astron. och Physik}, vol. 6, no. 30). Comparer aussi E. Noether, Gleichungen mit vorgeschriebener Gruppe et la dissertation de Seidelmann qui y est citée (\emph{Math. Ann.} 78).} Il s'agit du théorème de Kronecker selon lequel tout corps de nombres abélien sur les nombres rationnels peut être composé à partir de corps de racines de l'unité, et des théorèmes correspondants de Fueter\footnote{\emph{Math. Ann.} 75.} et de Hecke\footnote{\emph{Math. Ann.} 76.} pour les corps de nombres relativement abéliens. Les corps de nombres cherchés sont donc effectivement fournis par un moyen transcendant, fonctions exponentielles et fonctions modulaires, par opposition à la preuve purement existentielle dans le cas des fonctions algébriques.

La preuve d'existence des fonctions algébriques est précédée par l'existence des intégrales partout finies. Un analogue en est l'existence, dans un corps de nombres algébriques, des \og{}nombres primaires singuliers\fg{} définis relativement à un nombre premier \(\ell\), dont les racines \(\ell\)-ièmes sont relativement non ramifiées et fournissent les premières briques du corps de classes.\footnote{Furtwängler, Reziprozitätsgesetze für Potenzreste mit Primzahlexponenten in algebraischen Zahlkörpern. \emph{Math. Ann.} 67 et 72.} Cette preuve d'existence est menée avec l'aide essentielle du théorème selon lequel il existe toujours, dans le corps de nombres, des idéaux premiers ayant des caractères de reste prescrits; ce théorème peut donc être considéré comme un analogue du problème aux limites sur lequel repose l'existence des intégrales partout finies.

Les intégrales partout finies fournissent un critère transcendant pour déterminer si deux polygones appartiennent à la même classe, c'est-à-dire si deux groupes de points de même nombre sont corésiduels, à savoir par le théorème d'Abel, qui dit que les intégrales de première espèce, étendues sur des \og{}chemins corésiduels\fg{}, s'annulent; ou par le théorème différentiel correspondant avec sa réciproque.\footnote{Comparer, par exemple, la formulation chez M. Noether (\emph{Ann.} 37). Comme le théorème des résidus est historiquement issu du théorème d'Abel, la théorie géométrique parle de la \og{}somme\fg{} de groupes de points et de systèmes linéaires, en suivant les sommes intégrales; par opposition au \og{}produit\fg{} des classes dans la théorie arithmétique.}

Comme analogue du théorème d'Abel pour les corps de nombres algébriques, il faut prendre la \og{}loi de réciprocité pour les résidus de puissance \(\ell\)-ième\fg{}, qui donne, pour le corps lui-même ou au moins pour un surcorps approprié, un critère disant que deux idéaux appartiennent à la même classe. Son contenu peut aussi se formuler en disant que, dans ce surcorps, les nombres primaires singuliers ont le même caractère de reste relativement à tous les idéaux de la même classe. Ce dernier fait vaut aussi dans le corps lui-même, mais il ne s'y confond pas avec la loi de réciprocité.\footnote{Furtwängler, loc. cit.}

Enfin, la théorie hensélienne des nombres \(p\)-adiques fournit l'analogue du développement en séries de puissances des fonctions algébriques; sur ce point, on renverra au rapport de Hensel.
\clearpage
\setcounter{footnote}{0}
\begin{center}
{\Large\bfseries \NoetherNavAnchor{P15}{15. La finitude du système des invariants entiers des formes binaires}15. La finitude du système des invariants entiers des formes binaires}\par
\vspace{1em}
\emph{Nachrichten von der Gesellschaft der Wissenschaften zu Göttingen} 1919, pp. 138--156
\end{center}

\vspace{2em}
\begin{center}
Par Emmy Noether, à Göttingen.\\[1em]
Présenté par F. Klein dans la séance du 27 mars 1919.
\end{center}

Dans ce qui suit, comme extension du théorème de finitude connu, on démontre le théorème suivant: \emph{tous les invariants polynomiaux entiers d'un système binaire de formes fondamentales sont des fonctions polynomiales rationnelles, à coefficients entiers, d'un nombre fini d'entre eux.} Ici, comme d'habitude, par \og invariant \fg{} on entend toujours un invariant polynomiale rationnel d'une forme fondamentale ou d'un système de formes fondamentales, c'est-à-dire un invariant relativement à la transformation des coefficients induite par le groupe de toutes les transformations linéaires.

La démonstration repose sur un raffinement d'intégralité de la démonstration de finitude de Mertens--Hilbert.\footnote{F. Mertens, \emph{Wiener Sitzungsberichte}, janv. 1889; D. Hilbert, \emph{Math. Ann.} vol. 33 (1889) [daté de mars 1888]. Les deux démonstrations, apparues indépendamment l'une de l'autre à la suite de Mertens, Crelle vol. 100, concordent quant au contenu.} Cette démonstration de finitude peut être caractérisée comme suit. On décompose les formes fondamentales en leurs facteurs linéaires -- plus précisément, on associe à chaque forme fondamentale un produit de formes linéaires -- ce qui transforme chaque invariant $I$ en un invariant de ces formes linéaires, donc en une fonction des racines homogénéisées. Cette fonction doit satisfaire à trois conditions: 1) la condition d'invariance, 2) des conditions de poids, et 3) des conditions de symétrie. Les conditions 2) et 3) sont nécessaires pour que $I$ devienne polynomiale et rationnelle dans les coefficients des formes fondamentales. La condition 1) est satisfaite en représentant $I$ comme une fonction polynomiale rationnelle des déterminants des formes linéaires; la condition 2), en prenant, au lieu des déterminants isolés, certains produits de leurs puissances comme nouveaux arguments. La condition 3) dit alors simplement que $I$ devient une fonction symétrique polynomiale rationnelle de rangées de quantités, et que réciproquement toute fonction symétrique de ce genre des rangées de quantités devient un invariant des formes fondamentales. Les fonctions symétriques élémentaires de ces rangées de quantités donnent donc le système complet.

Si l'on ne considère que les invariants entiers, la condition 1) peut être satisfaite exactement comme ci-dessus dès que, pour les invariants de formes linéaires, on a démontré le théorème plus précis selon lequel \emph{tous les invariants entiers des formes linéaires binaires sont des fonctions polynomiales entières des déterminants de ces formes linéaires}. Je montre au \S\ 3 que tel est effectivement le cas, sous la forme que l'ensemble de toutes les relations modulo un nombre premier quelconque entre ces déterminants coïncide avec l'ensemble de toutes les relations identiques; autrement dit, le module correspondant à ces relations reste un module premier modulo tout nombre premier. Cette vérification s'effectue en normalisant les classes résiduelles relativement au module des formes quadratiques correspondant aux relations quadratiques entre les déterminants; ces formes quadratiques se révèlent alors former une base entière du module, modulo tout nombre premier, tandis qu'auparavant elles n'étaient connues comme base du module que si l'on faisait abstraction de l'intégralité.

Pour satisfaire à la condition 2), aucune modification n'est nécessaire, puisque l'introduction de produits de puissances ne touche pas aux coefficients entiers. La condition 3) dit maintenant que $n_1!\ldots n_\mu! I$ (où $n_1,\ldots n_\mu$ sont les degrés des différentes formes fondamentales) devient une fonction symétrique entière de rangées de quantités, et que réciproquement toute fonction symétrique entière de ces rangées devient un invariant entier. Mais, comme on le voit facilement, ces fonctions symétriques entières de rangées possèdent une base entière, pour laquelle les seules fonctions élémentaires ne suffisent plus. L'application du théorème de la base de Hilbert à coefficients entiers fournit alors aussi une base entière pour les $I$ eux-mêmes.

Le \S\ 1 réunit les théorèmes particuliers de finitude correspondant aux conditions 1), 2) et 3), dont la démonstration de finitude découle au \S\ 2. Les théorèmes particuliers de finitude correspondant à la condition 3) donnent aussi la finitude des invariants entiers des groupes finis de substitutions entières ou algébriquement entières, comme il est brièvement exposé au \S\ 4.

\section*{\S\ 1. Théorèmes particuliers de finitude}

J'énonce d'abord les théorèmes particuliers de finitude suivants; ils serviront à satisfaire les conditions d'invariance, de poids et de symétrie mentionnées dans l'introduction.

\textbf{I.} Tout invariant entier d'un système de formes linéaires binaires est une fonction polynomiale rationnelle, à coefficients entiers, des déterminants de ces formes linéaires. La démonstration sera donnée au \S\ 3.

\textbf{II.} Tout système $\frS$ de produits de puissances de $n$ variables à exposants non négatifs, qui, avec deux produits de puissances quelconques $f$ et $g$, contient aussi leur quotient dès que ce quotient est polynomiale dans les variables, possède une base multiplicative finie $f_1,\ldots,f_k$, composée de fonctions de $\frS$, telle que pour tout $f$ de $\frS$ on ait une représentation
\[
        f=f_1^{\lambda_1}\cdots f_k^{\lambda_k}
\]
à exposants non négatifs $\lambda_1,\ldots,\lambda_k$.

En effet, d'après le théorème de la base de Hilbert, il existe un nombre fini de fonctions de $\frS$, disons $f_1,\ldots,f_k$, dont chacune contient effectivement les $x$, telles que tout $f$ non constant de $\frS$ appartienne au module qu'elles engendrent:
\[
        f\equiv0\pmod{(f_1,\ldots,f_k)}.
\]
Il s'ensuit que, pour tout $f=x_1^{i_1}\cdots x_n^{i_n}$, il existe au moins un $f_\nu=x_1^{\nu_1}\cdots x_n^{\nu_n}$ tel que $\nu_1\leqq i_1;\ldots\nu_n\leqq i_n$. Dans la représentation
\[
        f=x_1^{i_1-\nu_1}\cdots x_n^{i_n-\nu_n} f_\nu = A f_\nu
\]
le facteur $A$ appartient à $\frS$ par hypothèse, mais il est de degré plus petit dans les variables que $f$. En répétant donc l'argument un nombre fini de fois, on en déduit l'assertion. Pour que la représentation vaille aussi pour $f=1$, il suffit de poser tous les exposants $\lambda$ égaux à zéro.\footnote{Cette représentation peut aussi être démontrée directement, et réciproquement le théorème de Hilbert en découle: comparer P. Gordan, \emph{Neuer Beweis des Hilbertschen Satzes über homogene Funktionen}, Gött. Nachr. 1899. Une autre démonstration directe de cette représentation, dont le théorème II est également déduit, se trouve chez A. Ostrowski, \emph{Über die Existenz einer endlichen Basis bei gewissen Funktionensystemen}, \emph{Math. Ann.} 78 (1916), \S\ 2,1 et \S\ 3,2. Le cas particulier du théorème II employé au \S\ 2, où les exposants parcourent toutes les solutions non négatives d'un système diophantien d'équations, se trouve déjà dans Gordan--Kerschensteiner, \emph{Vorlesungen über Invariantentheorie}, vol. I, p. 199.}

\textbf{III.} Toutes les fonctions symétriques polynomiales entières de $n$ rangées de quantités sont des fonctions polynomiales entières d'un nombre fini d'entre elles.

Soient $\tau_1(x),\ldots,\tau_n(x)$ les fonctions symétriques élémentaires des indéterminées $x_1,\ldots,x_n$. Alors, pour tout exposant $k$, il existe une identité
\[
        x_i^k = G_0^{(k)}(\tau)+x_iG_1^{(k)}(\tau)+\cdots+x_i^{n-1}G_{n-1}^{(k)}(\tau),
        \qquad (i=1,2,\ldots,n),
\]
où $G_0^{(k)},\ldots,G_{n-1}^{(k)}$ sont des fonctions polynomiales entières des $\tau(x)$. Si $x_i,y_i,\ldots,z_i$ sont les éléments de la $i$-ième rangée, alors, par ces identités et par les identités correspondantes pour $y,\ldots,z$, toutes les fonctions symétriques spéciales à une seule rangée
\[
        \sum_{i=1}^n x_i^{\alpha}y_i^{\beta}\cdots z_i^{\gamma}
\]
deviennent des fonctions polynomiales entières du nombre fini des fonctions
\[
        \sum_{i=1}^n x_i^{a}y_i^{b}\cdots z_i^{c}
        \qquad
        (0\le a<n;\;0\le b<n;\;\ldots\;0\le c<n)
\]
et des fonctions symétriques élémentaires $\tau(x),\tau(y),\ldots,\tau(z)$.\footnote{Hilbert, loc. cit., prend dans la base les sommes de puissances au lieu de $\tau(x),\ldots,\tau(z)$; la base se compose ainsi seulement de fonctions à une rangée, mais l'intégralité de la représentation est alors perdue.} Si l'on a les fonctions symétriques les plus générales des $n$ rangées -- cas qui n'intervient pas ci-dessous -- on montre de la même manière que les fonctions
\[
        \sum x_{i_1}^{a_1}y_{i_1}^{b_1}\cdots z_{i_1}^{c_1}
              x_{i_2}^{a_2}y_{i_2}^{b_2}\cdots z_{i_2}^{c_2}
              \cdots
              x_{i_n}^{a_n}y_{i_n}^{b_n}\cdots z_{i_n}^{c_n}
\]
\[
        (0\le a_\lambda<n,\;0\le b_\lambda<n,\;\ldots,\;0\le c_\lambda<n),
\]
où $i_1,i_2,\ldots,i_n$ parcourt toutes les permutations, forment, avec $\tau(x),\tau(y),\ldots,\tau(z)$, une base entière.

\textbf{IV.} Soient $F_1(x),\ldots,F_k(x)$ des fonctions polynomiales entières de $x_1,\ldots,x_n$, et soit $a$ un entier fixé. Alors toutes les fonctions
\[
        \frac{G(F)}{a},
\]
où $G$ désigne une fonction polynomiale entière, qui deviennent entières dans les $x$, sont des fonctions polynomiales entières d'un nombre fini d'entre elles.

En effet, d'après le théorème de la base de Hilbert à coefficients entiers, on a, pour chaque tel $G(F)/a$, une représentation
\[
 \frac{G(F)}{a}
 =
 A_1(F_1\ldots F_k)\frac{G_1(F)}{a}
 +\cdots+
 A_\varrho(F_1\ldots F_k)\frac{G_\varrho(F)}{a},
\]
où $A_1,\ldots,A_\varrho$ sont des fonctions polynomiales entières des $F$, et
\[
        \frac{G_1(F)}{a},\ldots,\frac{G_\varrho(F)}{a}
\]
appartiennent au système. Les fonctions
\[
        F_1=\frac{aF_1}{a},\ldots,F_k=\frac{aF_k}{a};
        \quad
        \frac{G_1(F)}{a},\ldots,\frac{G_\varrho(F)}{a}
\]
forment donc la base cherchée.

\setcounter{equation}{0}
\section*{\S\ 2. Finitude des invariants binaires entiers}

Soient $x,y$ des variables binaires soumises à une transformation à coefficients indéterminés
\begin{equation}
        x=s_{11}x'+s_{12}y';\qquad y=s_{21}x'+s_{22}y',
\end{equation}
et soit $f$ une forme binaire fondamentale dont les coefficients sont des indéterminées. Alors, par (1),
\begin{equation}
        f=a_0x^n+a_1x^{n-1}y+\cdots+a_ny^n
        =a'_0x'^n_n+a'_1x'^{\,n-1}y'+\cdots+a'_ny'^n.
\end{equation}
Par invariant entier de $f$, on entend, comme on sait, une fonction polynomiale entière $I(a)$ telle que
\begin{equation}
        I(a')=|s_{ik}|^\varrho I(a),\qquad \text{identiquement en $a,s_{ik}$}.
\end{equation}
Les invariants entiers d'un système de formes fondamentales sont définis de manière correspondante; on peut inclure dans la définition l'homogénéité par rapport aux coefficients des formes fondamentales prises séparément, puisque les autres invariants se composent de tels invariants individuellement homogènes de manière linéaire entière.

Outre (2), considérons une forme binaire spéciale $g$, à savoir le produit de $n$ formes linéaires dont les coefficients sont de nouveau des indéterminées:\footnote{Si, comme on le fait souvent en théorie des invariants, on écrit $f$ avec des coefficients binomiaux, $f=b_0x^n+\binom n1b_1x^{n-1}y+\cdots+b_ny^n$, alors le domaine de tous les invariants entiers $I(b)$ est agrandi par rapport à celui des $I(a)$, et les arguments de finitude employés ici échouent; $I(b)$ n'est plus une fonction entière des racines.}
\begin{align}
        g&=(\alpha_1x+\beta_1y)(\alpha_2x+\beta_2y)\cdots(\alpha_nx+\beta_ny)\notag\\
         &=\sigma_0(\alpha,\beta)x^n+\sigma_1(\alpha,\beta)x^{n-1}y+\cdots\sigma_n(\alpha,\beta)y^n\notag\\
         &=(\alpha_1'x'+\beta_1'y')(\alpha_2'x'+\beta_2'y')\cdots(\alpha_n'x'+\beta_n'y')\notag\\
         &=\sigma_0'x'^n+\sigma_1'x'^{\,n-1}y'+\cdots\sigma_n'y'^n.
\end{align}
Ainsi $\sigma_0(\alpha,\beta),\sigma_1(\alpha,\beta),\ldots,\sigma_n(\alpha,\beta)$ désignent les fonctions symétriques élémentaires homogénéisées, et les coefficients transformés $\sigma'_i$ deviennent égaux à $\sigma_i(\alpha',\beta')$.

En raison de l'indépendance algébrique de ces fonctions élémentaires, à toute relation
\begin{equation}
        G(\sigma_0,\sigma_1,\ldots,\sigma_n)=0
        \quad \text{identiquement en $\alpha,\beta$}
\end{equation}
correspond une autre relation
\begin{equation}
        G(\sigma_0,\sigma_1,\ldots,\sigma_n)=0
        \quad \text{identiquement en $\sigma_0,\ldots,\sigma_n$},
\end{equation}
et par suite, pour des indéterminées $a_0,\ldots,a_n$,
\begin{equation}
        G(a_0,a_1,\ldots,a_n)=0
        \quad \text{identiquement en $a_0,\ldots,a_n$}.
\end{equation}

Maintenant tout invariant entier $I(a)$ de $f$, par la substitution $a_i=\sigma_i(\alpha,\beta)$, donne un invariant entier $I(\sigma)=H(\alpha,\beta)$ de $g$, où $H$ devient une fonction entière de $\alpha,\beta$. Pour cet invariant, la relation (3) devient
\begin{equation}
        I(\sigma')=|s_{ik}|^\varrho I(\sigma)
        \quad \text{identiquement en $\sigma,s_{ik}$},
\end{equation}
et comme $\sigma'=\sigma(\alpha',\beta')$, (8) donne
\begin{equation}
        H(\alpha',\beta')=|s_{ik}|^\varrho H(\alpha,\beta)
        \quad \text{identiquement en $\alpha,\beta;s_{ik}$}.
\end{equation}
Ainsi $I(\sigma)=H(\alpha,\beta)$ devient un invariant entier des $n$ formes linéaires $(\alpha_i x+\beta_i y)$. Réciproquement, supposons que $H(\alpha,\beta)$ soit un invariant entier de ces $n$ formes linéaires, donc satisfasse (9), et soit en même temps égal à une fonction entière $I(\sigma)$ des $\sigma$. Puisque alors
\[
        H(\alpha',\beta)=I(\sigma(\alpha',\beta'))=I(\sigma'),
\]
l'équation (9) peut aussi s'écrire
\begin{equation*}
        I(\sigma')=|s_{ik}|^\varrho I(\sigma)
        \quad \text{identiquement en $\alpha,\beta;s_{ik}$}.\tag{9a}
\end{equation*}
Il résulte de (5) et (6) que (8) vaut, et de (7) que (3) vaut. Par conséquent, à tout invariant entier $I(a)$ correspond un invariant entier $H(\alpha,\beta)=I(\sigma)$ des $n$ formes linéaires, et réciproquement. De plus, si $I_1(a),\ldots,I_k(a)$ est une base entière de tous les $I(\alpha)$, alors $H_1(\alpha,\beta),\ldots,H_k(\alpha,\beta)$ est une base entière de tous les $H(\alpha,\beta)$ qui sont en même temps des fonctions entières $I(\sigma)$; ici aussi la réciproque vaut par (5), (6) et (7). Ainsi, pour démontrer l'existence d'une base entière de tous les $I(a)$, il suffit de la démontrer pour tous les invariants entiers $H(\alpha,\beta)$ qui sont en même temps des fonctions entières $I(\sigma)$; et la base
\[
        H_1(\alpha,\beta)=I_1(\sigma),\ldots,H_k(\alpha,\beta)=I_k(\sigma)
\]
fournit alors la base $I_1(a),\ldots,I_k(a)$ pour les $I(a)$.

Faisons parcourir à $H(\alpha,\beta)=I(\sigma)$ la totalité de tous les invariants entiers de (4). Comme, par (3), tout invariant est homogène dans les $\sigma$, et comme, par (4), chaque $\sigma$ est homogène et linéaire dans les coefficients de chaque forme linéaire, $H(\alpha,\beta)$ est homogène dans chaque rangée $\alpha_i\beta_i$ et de même dimension dans chaque rangée; en outre, c'est naturellement une fonction symétrique des $n$ rangées. Réciproquement, par le théorème fondamental sur les fonctions symétriques, toute fonction symétrique entière $H(\alpha,\beta)$ des $n$ rangées, homogène séparément et de même dimension dans chaque rangée, est une fonction polynomiale entière des $\sigma$. Ainsi, d'après l'argument précédent, elle est un invariant $I(\sigma)$ dès que $H(\alpha,\beta)$ est invariant. La totalité des invariants entiers de (4) est donc caractérisée comme la totalité des fonctions entières $H(\alpha,\beta)$ ayant les propriétés suivantes:
\[
\begin{array}{ll}
1) & \text{invariant entier des $n$ formes linéaires $\alpha_i x+\beta_i y$;}\\
2) & \text{homogène séparément et de même dimension dans chacune des $n$ rangées $\alpha_i,\beta_i$;}\\
3) & \text{symétrique dans les $n$ rangées $\alpha_i,\beta_i$.}
\end{array}
\]
Par le théorème particulier de finitude I, $H(\alpha,\beta)$, en raison de 1), devient une fonction polynomiale entière des déterminants
\[
        (ik)=\alpha_i\beta_k-\beta_i\alpha_k;
\]
donc
\[
        H(\alpha,\beta)=G((ik)).
\]
Dans cette représentation, la symétrie peut aussi être exprimée formellement en appliquant toutes les permutations des $n$ rangées à $G((ik))$:
\begin{equation}
        n!\,H(\alpha,\beta)
        =
        G^{(1)}((ik))+G^{(2)}((ik))+\cdots+G^{(n!)}((ik))
        =G^*((ik)).
\end{equation}
Avec chaque produit de puissances $P$ des $(ik)$, $G^*$ contient donc aussi la somme $Q=P^{(1)}+\cdots+P^{(n!)}$ sur toutes les permutations de $P$. En fait $G^*$ devient une combinaison linéaire entière d'un nombre fini de $Q$, $G^*=\sum cQ$. Comme les $Q$ satisfont aux conditions 1), 2) et 3), et deviennent donc eux-mêmes des invariants, il suffit de considérer les $Q$ à la place de $G^*((ik))$.

Faisons maintenant parcourir à $P$ le système $\frS$ de tous les produits de puissances des $(ik)$ qui interviennent dans les $Q$, c'est-à-dire qui satisfont à la condition 2). Si le quotient de deux $P$ est polynomiale dans les $(ik)$, il satisfait lui aussi à la condition 2), et appartient donc à $\frS$. Par le théorème II, il existe donc un nombre fini de tels $P$, disons $P_1,\ldots,P_\varrho$, tels que tout $P$ soit représentable sous la forme
\[
        P=P_1^{\lambda_1}\cdots P_\varrho^{\lambda_\varrho}
\]
avec des exposants non négatifs $\lambda$. En appliquant toutes les permutations, on obtient
\[
        Q=
        P_1^{(1)\lambda_1}\cdots P_\varrho^{(1)\lambda_\varrho}
        +\cdots+
        P_1^{(n!)\lambda_1}\cdots P_\varrho^{(n!)\lambda_\varrho}.
\]
Ainsi $Q$ devient une fonction symétrique entière des $n!$ rangées de quantités, chacune à $\varrho$ éléments:
\[
        P_1^{(1)}\ldots P_\varrho^{(1)};
        \quad
        P_1^{(2)}\ldots P_\varrho^{(2)};
        \quad \ldots;\quad
        P_1^{(n!)}\ldots P_\varrho^{(n!)};
\]
donc, par le théorème III, c'est une fonction polynomiale entière d'un nombre fini de fonctions symétriques de ces rangées. Ces fonctions de base sont elles-mêmes des quantités $Q$ ou des fonctions symétriques élémentaires des $n!$ éléments $P_i^{(1)},P_i^{(2)},\ldots,P_i^{(n!)}$, et satisfont donc aux conditions 1), 2) et 3) et deviennent des invariants entiers
\begin{equation}
        H_1(\alpha,\beta)=I_1(\sigma);\ldots;\quad
        H_l(\alpha,\beta)=I_l(\sigma).
\end{equation}
Par (10), puisque $G^*=\sum cQ$, tout $H(\alpha,\beta)=I(\sigma)$ est donc de la forme
\[
        H(\alpha,\beta)=I(\sigma)=\frac1{n!}\Phi(I_1,\ldots,I_l),
\]
où $\Phi$ désigne une fonction polynomiale entière; réciproquement, $\frac1{n!}\Phi$ est un invariant entier dès qu'il est polynomiale en $\alpha,\beta$. Par le théorème IV, on peut donc adjoindre aux invariants (11) un nombre fini d'autres invariants $I_{l+1}(\sigma),\ldots,I_k(\sigma)$ de telle manière que
\[
        H_1(\alpha,\beta)=I_1(\sigma);\ldots;H_l(\alpha,\beta)=I_l(\sigma);
\]
\[
        H_{l+1}(\alpha,\beta)=I_{l+1}(\sigma);\ldots;H_k(\alpha,\beta)=I_k(\sigma)
\]
forment une base entière de tous les invariants $H(\alpha,\beta)=I(\sigma)$. Ceci démontre le théorème dans le cas d'une seule forme fondamentale.

Si l'on a affaire à un système de formes fondamentales, alors, comme on l'a remarqué plus haut, on peut supposer l'homogénéité des invariants dans les coefficients de chaque forme fondamentale. On remplace de nouveau les formes fondamentales générales par des formes qui sont respectivement des produits de formes linéaires. Leurs invariants deviennent alors: 1) des invariants de ces formes linéaires; 2) homogènes séparément dans les rangées formées par les coefficients des formes linéaires, et de dimension égale dans les rangées appartenant à chaque forme fondamentale séparée; 3) symétriques dans les rangées appartenant à chaque forme fondamentale séparée. Réciproquement, chaque invariant entier de toutes les formes linéaires satisfaisant à ces conditions correspond de nouveau à un invariant entier du système de formes fondamentales. La démonstration est absolument parallèle à celle qui vient d'être donnée. À chaque produit de puissances $P=P_1^{\lambda_1}\cdots P_\varrho^{\lambda_\varrho}$ des déterminants individuels satisfaisant aux conditions de poids (théorèmes I et II), il faut appliquer toutes les
\[
        N=n_1!n_2!\cdots n_\mu!
\]
permutations, où $n_1,\ldots n_\mu$ désignent les degrés des formes fondamentales. Ainsi $N\cdot I$ est transformé en une fonction symétrique entière de $N$ rangées, chacune à $\varrho$ éléments. On en obtient la base entière de tous les $N\cdot I$, et par conséquent, par le théorème IV, la base entière de tous les $I$.
\setcounter{equation}{0}
\section*{\S\ 3. La base des invariants entiers des formes linéaires binaires}

Il reste maintenant à compléter la démonstration du théorème particulier de finitude I. D'après un théorème connu, démontré d'abord par Clebsch, tous les invariants rationnels polynomiaux des $n$ formes linéaires $\alpha_i x+\beta_i y$ sont des fonctions rationnelles polynomiales des déterminants $(ik)$. En particulier, tous les invariants entiers de ces formes linéaires deviennent donc des fonctions rationnelles polynomiales, quoique non nécessairement entières, de ces déterminants.

On montrera ci-dessous que, grâce aux relations identiques qui existent entre les $(ik)$, cette représentation est toujours possible intégralement.

En formules, et dans le langage de la théorie des modules, cela s'exprime ainsi. Soient $\xi_{ik}$, $i<k$, des indéterminées. Alors l'ensemble de tous les polynômes entiers $F(\xi_{ik})$ possédant la propriété que $F((ik))$ s'annule identiquement en $\alpha,\beta$ forme un module $\frM$; la somme de deux tels $F$, et aussi le produit d'un tel $F$ par un polynôme entier quelconque en les $\xi_{ik}$, appartient encore à cet ensemble. En formules:
\[
        F((ik))=0\quad \text{identiquement en $\alpha,\beta$}
\]
implique
\[
        F(\xi_{ik})\equiv0 \bmod \frM\quad [\text{identiquement en $\xi_{ik}$}],
\]
et réciproquement.

Je montre maintenant que l'énoncé tout à fait analogue suivant est vrai. De
\[
        F((ik))\equiv0\pmod p\quad \text{identiquement en $\alpha,\beta$}
\]
\begin{equation}
        \text{il s'ensuit que}\qquad
        F(\xi_{ik})\equiv0\pmod{(\frM,p)}
        \quad \text{identiquement en $\xi_{ik}$},
\end{equation}
et réciproquement, pour tout nombre premier $p$. Ainsi $\frM$ est aussi identique au module $\frM_p$ correspondant à l'ensemble de toutes les relations modulo $p$ entre les $(ik)$; ce module est constitué de tous les polynômes entiers $F(\xi_{ik})$ pour lesquels $F((ik))$ s'annule identiquement modulo $p$ en $\alpha,\beta$.

Que (1) exprime la représentabilité entière se voit comme suit. On peut aussi l'écrire sous la forme: de $F((ik))=p\cdot H(\alpha,\beta)$, identiquement en $\alpha,\beta$, on déduit que
\[
        F(\xi_{ik})-p\cdot G(\xi_{ik})\equiv0\pmod{\frM}
        \quad \text{identiquement en $\xi_{ik}$},
\]
et donc, par définition de $\frM$,
\[
        F((ik))=p\cdot G((ik))\quad [\text{identiquement en $\alpha,\beta$}]\footnote{Cette dernière conclusion démontre en même temps la réciproque de (1).}
\]
Ainsi $p\cdot H(\alpha,\beta)=p\cdot G((ik))$, ou
\[
        H(\alpha,\beta)=G((ik)).
\]
Si donc on a une représentation
\[
        H(\alpha,\beta)
        =
        \frac1{p_1^{\lambda_1}\cdots p_\varrho^{\lambda_\varrho}}
        F((ik)),
\]
où $p_1,\ldots p_\varrho$ désignent des nombres premiers, il en résulte aussi que
\[
        H(\alpha,\beta)
        =
        \frac1{p_1^{\lambda_1-1}\cdots p_\varrho^{\lambda_\varrho}}
        G((ik)),
\]
et, en répétant ce procédé un nombre fini de fois, on obtient l'intégralité de la représentation.

Pour la démonstration de (1), je désigne par
\[
        f_{ik,rs}=\xi_{ik}\xi_{rs}-\xi_{ir}\xi_{ks}+\xi_{is}\xi_{kr},
        \qquad i<k<r<s,
\]
les formes des indéterminées $\xi_{ik}$ qui correspondent aux relations quadratiques entre les $(ik)$,
\[
        (ik)(rs)-(ir)(ks)+(is)(kr)=0,
\]
et par $\frN=(f_{ik,rs})$ le module des polynômes entiers engendré par elles. De
\[
        F(\xi_{ik})\equiv0\pmod{\frN}
\]
il s'ensuit que
\[
        F((ik))=0\quad [\text{identiquement en $\alpha,\beta$}],
\]
et
\begin{equation}
\begin{array}{rcl}
        \text{de }F(\xi_{ik})&\equiv&0\pmod{(\frN,p)}\\
        \text{on déduit que }F((ik))&\equiv&0\pmod p\quad[\text{identiquement en $\alpha,\beta$}].
\end{array}
\end{equation}
Je vais maintenant montrer, en normalisant les classes résiduelles par rapport à $\frN$, que la réciproque de (2) vaut aussi. Alors $\frN=\frM=\frM_p$, et (1) sera démontré. L'égalité $\frN=\frM$ s'obtient comme cas particulier $\frN=\frM_p$ pour $p=0$; le cas $p=0$ est toujours compris dans les considérations suivantes. Autrement dit, les $f_{ik,rs}$ forment une base entière du module $\frM_p$, et en particulier aussi de $\frM_0=\frM$.

\subsection*{a) Cas de deux ou trois rangées $\alpha_i\beta_i$}

Comme il n'apparaît pas encore quatre indices distincts, $\frN$ est le module nul. Il faut donc montrer que de
\begin{equation}
        F((ik))\equiv0\pmod p\quad \text{identiquement en $\alpha,\beta$}
\end{equation}
il s'ensuit que
\[
        F(\xi_{ik})\equiv0\pmod p\quad \text{identiquement en $\xi_{ik}$}.
\]
Cela démontre (1) avec $\frM=\frM_p=0$.

Comme la condition $F(ik)=p\cdot H(\alpha,\beta)$ doit être satisfaite séparément par les composantes homogènes dans les rangées individuelles, il suffit partout de supposer $F((ik))$ homogène en chaque rangée, et donc $F(\xi_{ik})$ homogène en chaque indice. Dans le cas de deux rangées il n'y a que le déterminant $(12)$ et l'indéterminée correspondante $\xi_{12}$. Par l'homogénéité supposée,
\[
        F=c\cdot\,\xi_{12}^x.
\]
De
\[
        c\cdot\,(12)^\varkappa\equiv0\pmod p\quad \text{identiquement en $\alpha,\beta$}
\]
il s'ensuit, puisque $(12)\not\equiv0\pmod p$ identiquement en $\alpha,\beta$, que $c\equiv0\pmod p$; donc
\[
        F(\xi_{ik})\equiv0\pmod p\quad [\text{identiquement en $\xi_{ik}$}].
\]
Cela démontre (3), et par conséquent (1).

Dans le cas de trois rangées il y a les déterminants $(12),(13),(23)$, auxquels correspondent les indéterminées $\xi_{12},\xi_{13},\xi_{23}$. On a donc
\[
        F=\sum c\,\xi_{12}^{\tau_{12}}\xi_{13}^{\tau_{13}}\xi_{23}^{\tau_{23}}.
\]
Soient $\sigma_1,\sigma_2,\sigma_3$ les degrés de $F$ dans les indices $1,2,3$, respectivement. Alors
\[
        \sigma_1=\tau_{12}+\tau_{13},\qquad
        \sigma_2=\tau_{12}+\tau_{23},\qquad
        \sigma_3=\tau_{13}+\tau_{23},
\]
avec l'inversion unique
\[
        \tau_{12}=\frac12(\sigma_1+\sigma_2-\sigma_3),\quad
        \tau_{13}=\frac12(\sigma_1-\sigma_2+\sigma_3),\quad
        \tau_{23}=\frac12(-\sigma_1+\sigma_2+\sigma_3).
\]
Ainsi chaque $F$ séparément homogène contient au plus un terme,
\[
        F=c\,\xi_{12}^{\tau_{12}}\xi_{13}^{\tau_{13}}\xi_{23}^{\tau_{23}},
\]
d'où, comme ci-dessus, (3), et par conséquent (1), s'ensuit.

\subsection*{b) Cas de quatre rangées $\alpha_i,\beta_i$}

Les six indéterminées $\xi_{12},\xi_{13},\xi_{14},\xi_{23},\xi_{24},\xi_{34}$, et les déterminants correspondants, apparaissent; et $\frN=(f_{12,34})$. Comme
\begin{equation}
        \xi_{14}\xi_{23}\equiv \xi_{13}\xi_{24}-\xi_{12}\xi_{34}\pmod{\frN},
\end{equation}
tout produit de puissances des $\xi_{ik}$ est congru modulo $\frN$ à une combinaison linéaire entière de tels produits de puissances où le produit $\xi_{14}\xi_{23}$ n'apparaît pas. Ainsi, pour tout polynôme $F(\xi_{ik})$,
\begin{equation}
        F(\xi_{ik})\equiv F_0(\xi_{ik})\pmod{\frN},
\end{equation}
où $F_0$ ne contient pas le produit $\xi_{14}\xi_{23}$. Le polynôme $F_0$ sera appelé le polynôme normalisé de la classe résiduelle de $F$. Comme (4) est séparément homogène, à un $F$ séparément homogène correspond de nouveau un $F_0$ séparément homogène.

Je pose maintenant
\begin{equation}
        F_0(\xi_{ik})=G(\xi_{ik})+\xi_{34}\cdot F_1(\xi_{ik}),
\end{equation}
où $G$ ne contient aucun produit de puissances divisible par $\xi_{34}$, et ensuite
\begin{equation}
        [G]_{\xi_{i4}=\xi_{i3}}=K(\xi_{12},\xi_{13},\xi_{23}).
\end{equation}
De
\[
        F((ik))\equiv0\pmod p\quad \text{identiquement en $\alpha,\beta$}
\]
il s'ensuit maintenant, par (5) et en tenant compte de (2), que
\begin{equation}
        F_0((ik))\equiv0\pmod p\quad \text{identiquement en $\alpha,\beta$}.
\end{equation}
Par la spécialisation $\alpha_4=\alpha_3,\;\beta_4=\beta_3$, les formules (6) et (7) donnent
\[
        K(ik)\equiv0\pmod p\quad [\text{identiquement en $\alpha,\beta$}].
\]
D'après ce qui a été démontré sous a), on a donc
\[
        K(\xi_{ik})\equiv0\pmod p\quad \text{identiquement en $\xi_{ik}$}.
\]
Il en résulte ce qui reste à démontrer ci-dessous:
\[
        G(\xi_{ik})\equiv0\pmod p\quad \text{identiquement en $\xi_{ik}$}.
\]
Ainsi (6) donne
\[
        F_0(\xi_{ik})\equiv \xi_{34}\cdot F_1(\xi_{ik})\pmod p
        \quad \text{identiquement en $\xi_{ik}$}.
\]
Comme $(34)\not\equiv0\pmod p$ identiquement en $\alpha,\beta$, il résulte de (8) que
\[
        F_1((ik))\equiv0\pmod p\quad [\text{identiquement en $\alpha,\beta$}],
\]
où $F_1(\xi_{ik})$ est de nouveau normalisé et de degré plus bas en $\xi_{34}$. En répétant le procédé un nombre fini de fois, on obtient
\[
        F_0(\xi_{ik})\equiv\xi_{34}^{\lambda}F_\lambda(\xi_{ik})\pmod p,
\]
où $F_\lambda$ est de degré nul en $\xi_{34}$, et donc s'annule modulo $p$ d'après ce qui vient d'être démontré. Il en résulte donc
\begin{equation}
\begin{array}{rcl}
        F_0(\xi_{ik})&\equiv&0\pmod p\quad [\text{identiquement en $\xi_{ik}$}],\quad\text{ou}\\
        F(\xi_{ik})&\equiv&0\pmod{(\frN,p)}\quad [\text{identiquement en $\xi_{ik}$}].
\end{array}
\end{equation}
Ceci démontre la réciproque de (2), et donc (1), en même temps sous la forme plus précise pour $F_0(\xi_{ik})$.\footnote{Cette assertion plus précise signifie que chaque classe résiduelle contient un seul polynôme normalisé.}

Il reste seulement à ajouter que l'on déduit $G(\xi_{ik})\equiv0\pmod p$ de $K(\xi_{ik})\equiv0\pmod p$; ou, de façon équivalente, que de $G(\xi_{ik})\not\equiv0\pmod p$, on déduit $K(\xi_{ik})\not\equiv0\pmod p$. La normalisation n'est nécessaire que pour cette partie de l'induction. Sous la substitution $\xi_{i4}=\xi_{i3}$, les produits de puissances individuels de $G$ ne s'annulent pas. Une annulation de $K$ ne pourrait donc se produire que parce que des produits de puissances distincts de $G$ correspondent à des produits égaux dans $K$; et cela ne peut arriver que si ces produits de puissances ne diffèrent que par une permutation des indices 3 et 4. Supposons, par exemple, qu'un tel produit de puissances normalisé soit donné par
\[
        \xi_{12}^{\tau_{12}}\xi_{13}^{\tau_{13}}\xi_{14}^{\tau_{14}}\xi_{24}^{\tau_{24}}.
\]
En raison de l'homogénéité séparée supposée en chaque indice, le remplacement de 3 par 4 dans un $\xi_{ik}$ doit s'accompagner du remplacement de 4 par 3 dans un autre; un terme correspondant contiendrait donc, au lieu de $\xi_{13}\xi_{24}$, le produit $\xi_{14}\xi_{23}$, ce qui est exclu par la normalisation. De la même manière, contre
\[
        \xi_{12}^{\tau_{12}}\xi_{13}^{\tau_{13}}\xi_{14}^{\tau_{14}}\xi_{24}^{\tau_{24}}
\]
seul un terme contenant le facteur $\xi_{14}\xi_{23}$ pourrait se compenser, ce qui est encore exclu par la normalisation. Ainsi $K(\xi_{ik})\equiv0\pmod p$ implique $G(\xi_{ik})\equiv0\pmod p$.

\subsection*{c) Le cas de $n$ rangées $\alpha_i,\beta_i$}

Le cas de $n$ rangées se règle maintenant par induction complète, comme généralisation exacte du cas de quatre rangées qui vient d'être traité.\footnote{La démonstration vaut encore pour $n=3$.} Je montre d'abord de nouveau que, pour chaque polynôme $F$, il existe une congruence
\begin{equation}
        F(\xi_{ik})\equiv F_0(\xi_{ik})\pmod{\frN}
        \quad \text{identiquement en $\xi_{ik}$},
\end{equation}
où $\frN=(f_{ik,rs})$ et où $F_0(\xi_{ik})$ est normalisé. La normalisation doit consister en ceci: si $i,k,r,s$ sont quatre indices distincts et $i<k<r<s$, alors $F_0$ ne doit pas contenir le produit $\xi_{is}\xi_{kr}$; c'est-à-dire que les indices d'un $\xi$ ne doivent pas se trouver tous deux entre ceux de l'autre. Cette normalisation n'impose encore aucune restriction dans le cas de deux et trois rangées, où tout polynôme est donc normalisé; dans le cas de quatre rangées, comme on l'a montré, il existe aussi un polynôme normalisé pour chaque classe résiduelle. L'existence d'un polynôme normalisé pour chaque classe résiduelle peut donc être admise démontrée pour $(n-1)$ rangées.

Qu'un produit de puissances arbitraire dans $F$ soit écrit sous la forme
\[
        \xi_{i_1n}\cdots\xi_{i_\sigma n}\,P,
\]
où $P$ ne contient plus l'indice $n$. Par hypothèse d'induction, $P$ est congru modulo $\frN$ à un polynôme entier normalisé $P_0$:
\[
        \xi_{i_1n}\cdots\xi_{i_\sigma n}P
        \equiv
        \xi_{i_1n}\cdots\xi_{i_\sigma n}P_0
        \pmod{\frN}.
\]
Si $P_0$ ne contient aucun $\xi_{rs}$ tel que, pour au moins un indice $i$, les deux indices $r,s$ soient compris entre $i_\lambda$ et $n$, c'est-à-dire $i_\lambda<r<s<n$, alors la normalisation est déjà atteinte, puisque les hypothèses sont remplies pour le produit de deux $\xi$ quelconques provenant de $P_0$. Sinon, soit $i_\lambda$ un tel indice, $i_\lambda<r<s<n$, et soit $P_0^*$ le produit de puissances provenant de $P_0$ qui contient $\xi_{rs}$. Comme
\[
        \xi_{i_\lambda n}\xi_{rs}\equiv
        \xi_{i_\lambda s}\xi_{rn}-\xi_{i_\lambda r}\xi_{sn}
        \pmod{\frN},
\]
il vient
\begin{align*}
 &\xi_{i_1n}\cdots\xi_{i_\lambda n}\cdots\xi_{i_\sigma n}P_0^*\\
 &\equiv
 \xi_{i_1n}\cdots\xi_{rn}\cdots\xi_{i_\sigma n}Q
 +
 \xi_{i_1n}\cdots\xi_{sn}\cdots\xi_{i_\sigma n}R
 \pmod{\frN},\\
 &\equiv
 \xi_{i_1n}\cdots\xi_{rn}\cdots\xi_{i_\sigma n}Q_0
 +
 \xi_{i_1n}\cdots\xi_{sn}\cdots\xi_{i_\sigma n}R_0
 \pmod{\frN},
\end{align*}
où $Q_0$ et $R_0$ désignent de nouveau les polynômes normalisés des classes résiduelles de $Q$ et de $R$. Ils existent par hypothèse d'induction, puisque $Q$ et $R$ ne contiennent pas l'indice $n$, et ce sont assurément des polynômes entiers, puisque $Q$ et $R$ sont simplement des produits de puissances. Comme $i_\lambda<r<s<n$, on a aussi
\[
 ((n-i_1)+\cdots(n-r)+\cdots(n-i_\sigma))
 <
 ((n-i_1)+\cdots(n-i_\lambda)+\cdots(n-i_\sigma)),
\]
\[
 ((n-i_1)+\cdots(n-s)+\cdots(n-i_\sigma))
 <
 ((n-i_1)+\cdots(n-i_\lambda)+\cdots(n-i_\sigma)).
\]
Ainsi, en passant de $P_0^*$ à $Q_0,R_0$, cette somme de différences a diminué; comme la somme est nécessairement $\ge \sigma$, le procédé doit s'arrêter après un nombre fini de pas pour chaque produit de puissances $P_0^*$. Par conséquent
\[
        F(\xi_{ik})
        =
        \sum \xi_{i_1n}\cdots\xi_{i_\sigma n}\cdot P^{(\nu)}
        \equiv
        \sum \xi_{j_1n}\cdots\xi_{j_\sigma n}\cdot S_0^{(\nu)}
        =
        F_0(\xi_{ik})\pmod{\frN},
\]
où le polynôme entier $F_0(\xi_{ik})$ est nécessairement normalisé, car sinon la somme des différences pourrait encore être abaissée. Ceci démontre (10). Comme, par hypothèse d'induction pour $(n-1)$ indices, à chaque $F$ séparément homogène correspond un $F_0$ normalisé séparément homogène, le procédé donné montre qu'il en est de même pour $n$ indices. Il ne s'agira donc ci-dessous que de polynômes séparément homogènes.

Comme en b), posons
\begin{equation}
        F_0(\xi_{ik})=G(\xi_{ik})+\xi_{n-1,n}F_1(\xi_{ik}),\footnote{L'identité correspondant à (10) et (11),
        \[
          F((ik))=F_0((ik))=G((ik))+(n-1,n)\cdot F_1((ik)),
        \]
        représente un analogue du développement en série de Clebsch--Gordan dès que les deux rangées $(n-1)$ et $(n)$ sont remplacées par des rangées de variables $x,y$, de sorte que $(n-1,n)$ soit remplacé par $(xy)$. Mais tandis que Clebsch--Gordan procède par un développement suivant les polaires, la normalisation de $F_0$ correspond à des \og termes initiaux \fg{} de polaires, et force par là l'intégralité qui ne vaut pas chez Clebsch--Gordan. La démonstration connue du fait que, si l'on fait abstraction de l'intégralité, les $f_{ik,rs}$ forment une base de $\frM$, se fait précisément au moyen du développement en série de Clebsch--Gordan (Gordan--Kerschensteiner, \emph{Vorlesungen über Invariantentheorie}, p. 132).}
\end{equation}
où $G$ ne contient aucun produit de puissances divisible par $\xi_{n-1,n}$, et posons en outre
\begin{equation}
        [G]_{(\xi_{in}=\xi_{i,n-1})}=K(\xi_{ik}).
\end{equation}
Comme le montre la définition de la normalisation, $K$, aussi bien que $G$, est normalisé.

Il existe de nouveau une correspondance biunivoque entre $G$ et $K$. Des produits de puissances distincts de $G$ ne peuvent correspondre au même produit dans $K$ que s'ils ne diffèrent que par une permutation des indices $n$ et $(n-1)$. Supposons, par exemple, que $G$ soit de degré $\varrho$ en l'indice $(n-1)$ et de degré $\sigma$ en l'indice $n$, et que
\[
        \varkappa_1\leqq\varkappa_2\cdots\leqq\varkappa_\varrho\leqq\varkappa_{\varrho+1}\cdots\leqq\varkappa_{\varrho+\sigma}
\]
soient les indices liés à $(n-1)$ et à $n$ dans un produit de puissances de $G$. Leur nombre doit être $\varrho+\sigma$, puisque, d'après (11), l'indéterminée $\xi_{n-1,n}$ n'apparaît pas dans $G$. Par la normalisation, le produit de puissances est donc de la forme
\[
        \xi_{\varkappa_1,n-1}\cdots\xi_{\varkappa_\varrho,n-1}
        \xi_{\varkappa_{\varrho+1},n}\cdots\xi_{\varkappa_{\varrho+\sigma},n}\cdot P(\xi_{ik}),
\]
où $P(\xi_{ik})$ ne contient plus les indices $n$ et $(n-1)$. Chaque permutation de $n$ avec $(n-1)$ donnerait, pour $i\ne j$, le facteur $\xi_{in}\xi_{j,n-1}$ avec $i<j$; alors $j,n-1$ est situé entre $i$ et $n$, et le terme ne serait pas normalisé. Ainsi, pour $\varkappa_1,\varkappa_\varrho,\ldots,\varkappa_{r+\sigma}$ fixes et $P$ fixe, $G$ ne contient qu'un produit de puissances. Il est donc montré que des produits de puissances distincts de $G$ correspondent à des produits distincts de $K$. Par conséquent
\begin{equation}
        K(\xi_{ik})\equiv0\pmod p\quad\Longrightarrow\quad
        G(\xi_{ik})\equiv0\pmod p.
\end{equation}

De
\[
        F((ik))\equiv0\pmod p\quad \text{identiquement en $\alpha,\beta$}
\]
il s'ensuit maintenant, par (10) et en tenant compte de (2), que
\begin{equation}
        F_0((ik))\equiv0\pmod p\quad \text{identiquement en $\alpha,\beta$}.
\end{equation}
En spécialisant $\alpha_n=\alpha_{n-1},\;\beta_n=\beta_{n-1}$, (11) et (12) donnent
\begin{equation}
        K((ik))\equiv0\pmod p\quad [\text{identiquement en $\alpha,\beta$}].
\end{equation}
Comme $K(\xi_{ik})$ est normalisé et ne contient que $(n-1)$ indices, il résulte de ce qui a été démontré sous a) et sous b), formule (9), que
\[
        K(\xi_{ik})\equiv0\pmod p\quad \text{identiquement en $\xi_{ik}$}.
\]
Alors, par (13),
\[
        G(\xi_{ik})\equiv0\pmod p\quad \text{identiquement en $\xi_{ik}$}.
\]
Donc (11) donne
\[
        F_0(\xi_{ik})\equiv \xi_{n-1,n}F_1(\xi_{ik})\pmod p
        \quad \text{identiquement en $\xi_{ik}$}.
\]
Puisque $(n,n-1)\not\equiv0\pmod p$ identiquement en $\alpha,\beta$, il résulte de (14) que
\[
        F_1((ik))\equiv0\pmod p\quad [\text{identiquement en $\alpha,\beta$}],
\]
où $F_1(\xi_{ik})$ est encore normalisé et de degré plus bas en $\xi_{n-1,n}$ que $F_0$. La répétition un nombre fini de fois du procédé donne
\[
        F_0(\xi_{ik})\equiv\xi_{n-1,n}^{\lambda}F_\lambda(\xi_{ik})\pmod p,
\]
où $F_\lambda$ est de degré nul en $\xi_{n-1,n}$, et donc s'annule modulo $p$ par ce qui vient d'être démontré. On obtient ainsi
\begin{equation}
\begin{array}{rcl}
        F_0(\xi_{ik})&\equiv&0\pmod p\quad \text{identiquement en $\xi_{ik}$},\quad\text{ou}\\
        F(\xi_{ik})&\equiv&0\pmod{(\frN,p)}\quad \text{identiquement en $\xi_{ik}$}.
\end{array}
\end{equation}
La réciproque de (2), et donc (1), est ainsi démontrée. En même temps est démontrée, dans le cas de $n$ indices, et par conséquent en général, l'assertion plus précise utilisée dans le pas d'induction pour $F_0(\xi_{ik})$.

\setcounter{equation}{0}
\section*{\S\ 4. La finitude des invariants entiers des groupes finis}

La preuve élémentaire de la finitude que j'ai donnée dans ma note \og Le théorème de finitude pour les invariants des groupes finis \fg{}\footnote{\emph{Math. Ann.} 77, p. 89 (1915).} admet, au moyen des théorèmes particuliers de finitude III et IV du \S\ 1, le renforcement à l'intégralité, tandis que la preuve ordinaire fondée sur le théorème de la base de Hilbert ne fournit pas cette intégralité. Même l'emploi d'une base finie à coefficients entiers ne donne qu'une représentation où apparaît au dénominateur une puissance arbitrairement élevée de l'entier $h$, l'ordre du groupe.

Je me rattache partout aux notations de la note mentionnée. Soit le groupe fini $\frH$ formé des $h$ transformations linéaires $A_1,\ldots,A_h$ de déterminant non nul,\footnote{Il suffirait de supposer que $\frH$ est simplement isomorphe à un groupe fini abstrait; c'est-à-dire que, si un déterminant de transformation des $A$ s'annule, $\frH$ soit semblable à un groupe $\begin{pmatrix}\frS&0\\0&0\end{pmatrix}$, où les déterminants des substitutions de $\frS$ ne s'annulent pas. Dans ce cas les invariants de $\frH$ coïncident avec les invariants de $\frS$, de sorte que la condition de non-annulation des déterminants ne constitue en fait aucune restriction de généralité.} où $A_k$ désigne la transformation linéaire
\[
        x_1^{(k)}=\sum_{\nu=1}^{n}a_{1\nu}^{(k)}x_\nu,
        \quad\ldots,\quad
        x_n^{(k)}=\sum_{\nu=1}^{n}a_{n\nu}^{(k)}x_\nu,
\]
ou, brièvement, $(x^{(k)})=A_k(x)$. Le groupe $\frH$ envoie donc la rangée $(x)$, d'éléments $x_1,\ldots,x_n$, sur les rangées $(x^{(k)})$, d'éléments $x_1^{(k)},\ldots,x_n^{(k)}$. On suppose que les coefficients $a_{i\nu}^{(k)}$ sont des entiers algébriques; cela ne restreint pas la généralité, car, d'après un théorème de J. Schur,\footnote{\emph{Über Gruppen linearer Substitutionen mit Koeffizienten aus einem algebraischen Zahlkörper}. Satz VII. \emph{Math. Ann.} 71, p. 555 (1912).} tout groupe fini de transformations linéaires est semblable à un groupe pour lequel les $a_{i\nu}^{(k)}$ deviennent des entiers algébriques.

Par invariant absolu entier du groupe, on entend une fonction rationnelle polynomiale de $x_1,\ldots,x_n$, à coefficients entiers algébriques du corps de nombres $\frK$ engendré par les $a_{i\nu}^{(k)}$, qui reste identiquement inchangée sous l'application de $A_1,\ldots,A_h$. Pour un tel invariant $f(x)$ on a
\begin{equation}
        f(x)=f(x^{(1)})=\cdots=f(x^{(h)})
        =\frac1h\cdot\sum_{k=1}^{h}f(x^{(k)}).
\end{equation}
Réciproquement, toute fonction symétrique entière des $h$ rangées de grandeurs $(x^{(k)})$ devient aussi un invariant entier du groupe.

Soit maintenant $\omega_1,\omega_2,\ldots,\omega_\varrho$ une base de tous les entiers de $\frK$. Alors chaque invariant peut être représenté sous la forme
\[
        f(x)=\omega_1 f_1(x)+\cdots+\omega_\varrho f_\varrho(x),
\]
où $f_1(x),\ldots,f_\varrho(x)$ ont des coefficients entiers rationnels. Par (1) il vient donc
\[
        h.f(x)
        =
        \omega_1\sum_{k=1}^{h}f_1(x^{(k)})
        +\cdots+
        \omega_\varrho\sum_{k=1}^{h}f_\varrho(x^{(k)}).
\]
Ici les sommes individuelles sont des fonctions symétriques des $h$ rangées de grandeurs $(x^{(k)})$, et même à coefficients entiers rationnels. Par le théorème III, elles peuvent donc être exprimées polynomialement et rationnellement, avec des coefficients entiers rationnels, au moyen d'un nombre fini de fonctions symétriques entières-rationnelles de ces rangées. D'après ce qui précède, ces fonctions de base $I_1(x),\ldots,I_\sigma(x)$ deviennent des invariants entiers du groupe, en général à coefficients entiers algébriques, puisque $\frH$ ne contient pas nécessairement avec $A$ toutes ses transformations algébriquement conjuguées.

Il existe donc une représentation
\begin{equation}
        h\cdot f(x)=\omega_1G_1(I)+\cdots+\omega_\varrho G_\varrho(I),
\end{equation}
où $G_1,\ldots,G_\varrho$ ont des coefficients entiers rationnels.

Soient maintenant $w_1,\ldots,w_\varrho$ des indéterminées, et que
\[
        L=w_1G_1(I)+\cdots+w_\varrho G_\varrho(I)
\]
parcoure toutes les formes linéaires en les $w$ qui, par la substitution $w_i=\omega_i/h$, correspondent à des invariants $f(x)$ dans la représentation (2). Le théorème de la base de Hilbert à coefficients entiers donne
\begin{equation}
        L=A_1(I)L_1+\cdots+A_\tau(I)L_\tau,
\end{equation}
où les $A_i$ ont des coefficients entiers rationnels, et où, puisque tous les $L$ sont linéaires en les $w$, les $A_i$ ne contiennent plus les $w$. Si, par la substitution $\omega_i=\frac{w_i}{h}$, $L_1,\ldots,L_\tau$ deviennent les invariants $\mathfrak I_1(x),\ldots,\mathfrak I_\tau(x)$, alors (3) devient
\begin{equation}
        f(x)=A_1(I)\cdot\mathfrak I_1(x)+\cdots+A_\tau(I)\mathfrak I_\tau(x).
\end{equation}
Cette représentation montre que
\[
        I_1(x),\ldots,I_\sigma(x);\mathfrak I_1(x),\ldots,\mathfrak I_\tau(x)
\]
forment une base de tous les invariants entiers du groupe $\frH$, telle que chaque invariant soit représentable comme une fonction rationnelle polynomiale, à coefficients entiers rationnels, de ces invariants en nombre fini.

Si, en particulier, le groupe $\frH$ a la propriété de contenir avec chaque transformation $A$ aussi toutes les transformations $A',A'',\ldots$ à coefficients algébriquement conjugués, alors les fonctions symétriques entières-rationnelles des $h$ rangées $(x^{(k)})$ deviennent des fonctions de $x$ à coefficients entiers rationnels. Cela vaut en particulier aussi pour les fonctions de base des fonctions symétriques. Si l'on se restreint donc aux invariants $f(x)$ à coefficients entiers rationnels, alors les fonctions de base
\[
        I_1(x),\ldots,I_\sigma(x);\mathfrak I_1(x),\ldots,\mathfrak I_\tau(x)
\]
ont elles aussi des coefficients entiers rationnels. Sous cette hypothèse spéciale, qui en général n'est pas satisfaite par les groupes $\frH$, la finitude entière vaut donc aussi pour le sous-domaine des invariants à coefficients entiers rationnels.
\clearpage
\setcounter{footnote}{0}
\section*{\NoetherNavAnchor{P16}{16. Sur le développement en série dans la théorie des formes}16. Sur le développement en série dans la théorie des formes}
\begin{center}
\textbf{(Addendum et rectification au n\textsuperscript{o} III de l'article
\og{}Sur la représentation entière rationnelle des invariants d'un système d'un nombre arbitraire de formes fondamentales\fg{},
\emph{Math. Ann.} 77, p. 93.)}\\[0.65em]
Par Emmy Noether, à Göttingen.\\[0.65em]
\emph{Math. Ann.} 81 (1920), p. 25--30
\end{center}

Par \og{} développement en série de la théorie des formes \fg{}, on entend, comme on sait, d'une part la généralisation du développement en série de Clebsch--Gordan --- c'est-à-dire du développement d'une forme contenant deux séries binaires suivant les puissances du déterminant de ces séries, les coefficients étant des polaires de formes ne dépendant que d'une seule série de variables, ou, ce qui revient au même, s'annulant par application du procédé $\Omega$ --- à des formes contenant un nombre arbitraire de séries de $n$ variables chacune; et, d'autre part, comme généralisation des formes normales qui se présentent pour les \og{} coordonnées complexes \fg{} $t_{ik}$, un développement suivant des formes normales pour des formes qui contiennent simultanément des grandeurs de différents degrés $t_i, t_{ik}, t_{ikl}\ldots$. Cette seconde sorte de développement peut toutefois être obtenue à partir de la première par des procédés de différentiation invariante, comme l'ont exécuté surtout Mertens et Deruyts.\footnote{F. Mertens, \emph{Über invariante Gebilde quaternärer Formen} (Sitzber. d. Ak. d. Wiss. Wien 98, Abt. IIa (1889)) pour le cas quaternaire. Les formes normales générales se trouvent chez J. Deruyts: \emph{Sur la réduction des fonctions invariantes dans le système des variables géométriques}, Bulletins de l'Ac. R. de Belgique 24 (1892). Sans connaître ce travail, et en suivant de près Mertens, j'ai donné le développement général suivant les formes normales au \S~7 du travail: \emph{Zur Invariantentheorie der Formen von $n$ Variabeln}, J. f. M. 139 (1910).}

Les divers développements de la première sorte résultent tous --- comme cela est brièvement esquissé, par exemple, dans \emph{Math. Ann.} 77, loc. cit. III --- du développement d'une forme de $n$ séries de $n$ variables chacune suivant les puissances du déterminant de ces séries, les coefficients étant des polaires de formes de seulement $(n-1)$ séries, elles-mêmes déduites de la forme de départ par formation de polaires et par le procédé $\Omega$.\footnote{Comparer les indications bibliographiques loc. cit. A paru depuis, et repose également sur la méthode formelle: Schouten, \emph{Over reeksontwikkelingen van algebraische vormen met verschillende rijen van variabelen van verschillenden graad}, Verslag Ak. Amsterdam 27 (1919), p. 1481.} Le développement originel de Clebsch--Gordan $(n=2)$ était allé encore un pas plus loin: les procédés polaires spéciaux qui y apparaissent y sont indiqués explicitement, même avec les coefficients numériques.

Ce que je voudrais maintenant ajouter à la note indiquée --- où le développement en série est conçu de manière plus conceptuelle, comme un problème de faisceaux linéaires de formes --- est une représentation explicite à coefficients numériques près, qui s'obtient par spécialisation des recherches d'E. Fischer relatives aux systèmes généraux de modules.\footnote{E. Fischer, \emph{Über die Differentiationsprozesse der Algebra}, J. f. M. 148 (1917).}

On est conduit à cette insertion par l'interprétation de la forme normale dans les $t_{ik}$. En effet, si l'on remplace les \og{} coordonnées complexes \fg{} $t_{ik}$ par les \og{} coordonnées de droites \fg{}
\[
        p_{ik}=(x_i y_k-x_k y_i),
\]
alors, en vertu des relations identiques qui existent entre les $p_{ik}$, la représentation d'une forme $F(p_{ik})$ n'est plus univoque. On peut obtenir l'univocité si l'on exige qu'entre les coefficients de $F$ subsistent les mêmes relations linéaires qu'entre les produits de puissances des $p_{ik}$ qui apparaissent dans $F$.\footnote{Comparer Clebsch, \emph{Über eine Fundamentalaufgabe der Invariantentheorie}, \S~4, Abh. d. K. G. d. Wissensch. zu Göttingen 17 (1872).} Pour cette forme normale $\Phi$, on a donc:
\begin{equation}
\begin{aligned}
        F(p_{ik})&=\Phi(p_{ik}) &&[\text{identiquement en }x,y],\quad\text{ou}\notag\\
        F(t_{ik})&\equiv\Phi(t_{ik}) &&\modu{M},
\end{aligned}
\tag{1}
\end{equation}
où $M$ désigne le module de toutes les relations identiques entre les $p_{ik}$, c'est-à-dire qu'il consiste en la totalité des formes $G(t_{ik})$ pour lesquelles $G(p_{ik})$ s'annule identiquement en $x,y$. $\Phi(t_{ik})$ est la \og{} forme normale \fg{} en coordonnées complexes, tandis que les fonctions de base de $M$ deviennent des formes quadratiques dans les $t_{ik}$. Par itération du procédé relativement aux facteurs de ces fonctions de base, on obtient le développement en série complet. Des considérations correspondantes valent lorsque $t_i,t_{ik},t_{ikl},\ldots$ apparaissent simultanément.

Par une marche de pensée analogue, on peut ramener le développement suivant les polaires. Si, par exemple, dans une forme $F(x,y,z)$, on remplace les trois séries ternaires $x=(x_1,x_2,x_3),y,z$ par des séries $\xi,\eta,\zeta$ qui sont composées linéairement de deux seulement d'entre elles, par exemple $\xi=x,\eta=y;\zeta=\lambda_1x+\lambda_2y$, alors la représentation de $F(\xi,\eta,\zeta)$ n'est de nouveau plus univoque. On peut de nouveau prescrire pour les coefficients les mêmes relations linéaires que pour les produits de puissances de $\xi,\eta,\zeta$ qui apparaissent dans $F$. Puisqu'ici le module $M$ de toutes les relations identiques possède pour fonction de base le déterminant $\Delta=(xyz)$ des trois séries (loc. cit. III, lemme auxiliaire a)), [1] est remplacée par:
\begin{equation}
\begin{aligned}
        F(\xi,\eta,\zeta)&=\Phi(\xi,\eta,\zeta)&&[\text{identiquement en }x,y],\quad\text{ou}\notag\\
        F(x,y,z)&\equiv\Phi(x,y,z)&&\modu{\Delta}.
\end{aligned}
\tag{2}
\end{equation}
Il apparaît que $\Phi$ provient de la formation de polaires (comparer la formule (2) loc. cit.); le développement en série complet s'obtient de nouveau par itération.

La formation des formes normales $\Phi$ se subordonne donc, dans les deux cas, au problème suivant: si, dans $F(z)$, on remplace les indéterminées $z_i$ par des polynômes $f_i(\xi)$ en des paramètres $\xi$, de sorte que la représentation $F(f_i(\xi))$ ne soit donc plus univoque, l'univocité doit être fixée en exigeant qu'entre les coefficients subsistent les mêmes relations linéaires qu'entre les produits de puissances des $f_i(\xi)$ qui apparaissent dans $F$. On a donc
\begin{equation}
\begin{aligned}
        F(f_i(\xi))&=\Phi(f_i(\xi))&&[\text{identiquement en }\xi],\quad\text{ou}\notag\\
        F(z)&\equiv\Phi(z)&&\modu{M},
\end{aligned}
\tag{3}
\end{equation}
où $M$ désigne le module de toutes les relations identiques entre les $f_i(\xi)$.

Pour cette forme normale $\Phi(z)$, E. Fischer, au \S~11, I du travail cité, a donné un développement explicite à coefficients numériques près, à savoir au moyen d'opérations linéaires. En même temps, il a cependant aussi donné, dans l'introduction III et au \S~6 II, pour la différence $F-\Phi$ appartenant à $M$ --- en supposant connue une base du module $M$ --- une expression explicite dans le même sens; et, en combinant ces développements au \S~11 II, il a ainsi remplacé la formule [3] par une représentation explicite. La question des coefficients numériques, encore ouverte ici, est, d'après le \S~12 du même travail, identique à la détermination des valeurs propres et des fonctions propres des opérations linéaires (formule 82).

Je spécialise maintenant les formules de Fischer, en 1., au cas du développement en série suivant les polaires, et j'en obtiens, par itération, le développement en série. L'opération destinée à déterminer $\Phi$ se transforme alors en procédés polaires déterminés, ce qui est considérablement plus précis que la formulation antérieure rappelée plus haut. Dans l'opération destinée à déterminer $F-\Phi$, la multiplication par le déterminant $\Delta$ et le procédé $\Omega$ apparaissent combinés. Pour parvenir à la forme usuelle, non explicite, il faut encore remarquer que $\Omega\Delta$ peut s'exprimer par des procédés polaires.

Il faut aussi ajouter ici une correction. Dans le travail cité plus haut (p. 101, ligne 6 du haut), j'ai supposé par inadvertance comme généralement valable la transformation identique
\[
        \Omega(\Delta\psi)=c_1\psi+c_2\Delta\cdot\Omega(\psi),
\]
qui ne vaut que dans le domaine binaire, et j'ai ainsi effectué le passage de la formule fondamentale (2), qui donne l'expression de la forme normale $\Phi$ dans [2], au développement en série (3). Les considérations de cet endroit ne donnent donc que la formule (2), et, comme cas spécial de celle-ci, le théorème de réduction (1); mais, puisqu'il manque une expression explicite pour la différence $F-\Phi$, elles ne suffisent pas à fonder le développement en série.

En 2., j'obtiens de même, par spécialisation, la forme normale $\Phi(t_{ik})$ de [1]; plus généralement, la forme normale $\Phi(t_i,t_{ik},t_{ikl},\ldots)$, ce qui conduit à des formules explicites connues. Mais comme l'établissement du développement en série complet exige la connaissance d'une base du module $M$, la voie de la théorie des invariants (comparer la littérature citée) est ici préférable; elle n'exige pas cette connaissance, mais fournit au contraire une base du module au moyen du développement en série.

1. D'après Fischer, \S~11 (formule (19) et les lignes suivantes), la forme normale $\Phi(z)$ de [3] --- notée là $N(\zeta)$ --- est donnée par:
\begin{equation}
        \Phi(z)=(a_1S+a_2S^2+\cdots+a_rS^r)F(z),
\tag{4}
\end{equation}
où les constantes $a_1\ldots a_r$ appartiennent au domaine de coefficients des $f(\xi)$ et ne dépendent que du degré de $F(z)$; l'opérateur $S$ est défini par
\begin{equation}
\begin{gathered}
        S=AB;\qquad BF(z)=F(f(\xi));\\
        AF(f(\xi))=\left\{\sum z_i f_i\!\left(\frac{\p}{\p\xi}\right)\right\}^{p}F(f(\xi))\footnotemark.
\end{gathered}
\tag{5}
\end{equation}
\footnotetext{Comparer la formule (75). Dans le cas de [3], la somme $\alpha$ parcourt tous les produits de puissances de dimension $p$, et (75) se spécialise donc en [5].}
Si, dans [2], on considère $F$ comme une forme de $p$-ième dimension en $z$ seulement, alors $f_i(\xi)$ se spécialise en $\lambda_1x_i+\lambda_2y_i$; on obtient par conséquent:
\[
\begin{gathered}
        BF=F(x,y,\lambda_1x+\lambda_2y),\\
        SF=\left\{\sum z_i\left(\frac{\p}{\p\lambda_1}\frac{\p}{\p x_i}+
        \frac{\p}{\p\lambda_2}\frac{\p}{\p y_i}\right)\right\}^{p}
        F(x,y,\lambda_1x+\lambda_2y),
\end{gathered}
\]
ou, écrit en procédés polaires (comparer loc. cit. II, 2 et 3, aussi pour la notation abrégée):
\begin{equation}
        SF(x,y,z)=\frac1{p!}\left[z\left(\frac{\p}{\p\lambda_1}\frac{\p}{\p x}+
        \frac{\p}{\p\lambda_2}\frac{\p}{\p y}\right)\right]^p
        \left[(\lambda_1x+\lambda_3y)\frac{\p}{\p z}\right]^p F(x,y,z).
\tag{6}
\end{equation}
\emph{Les équations} [4] \emph{et} [6] \emph{donnent l'expression explicite cherchée pour} $\Phi(x,y,z)$; en même temps $\Phi(x,y,z)$ rentre dans la forme générale mentionnée au début ($\sum PH$ dans (2), loc. cit.) --- polaires d'expressions qui ne dépendent que de deux séries et qui sont obtenues à partir de $F$ par des procédés polaires. Les coefficients encore indéterminés $a_i$ sont des nombres rationnels, puisque les $f_i(\xi)$ ont ici des coefficients numériques rationnels entiers. Si [2] est écrit sous la forme
\begin{equation}
        F=\Phi+\Delta F_1,
\tag{2a}
\end{equation}
alors la spécialisation de la formule donnée par Fischer dans l'introduction III et au \S~6 II (p. 41) donne
\begin{equation}
        F_1=(\alpha_1\Omega+\alpha_2\Omega\Delta\Omega+\cdots+\alpha_\lambda\Omega\Delta\Omega\cdots\Delta\Omega)F\footnotemark,
        \qquad \left[\Omega=\left(\frac{\p}{\p x}\frac{\p}{\p y}\frac{\p}{\p z}\right)\right],
\tag{7}
\end{equation}
\footnotetext{Kerim utilise cette formule spécialisée dans sa dissertation d'Erlangen, à paraître prochainement, dans l'application aux \og{} formes d'inertie \fg{}.}
où les $\alpha_i$ sont de nouveau des nombres rationnels ne dépendant que du degré de $F$.

L'itération conduit à:
\[
        F_1=\Phi_1+\Delta F_2;\quad F_2=\Phi_2+\Delta F_3;\quad\ldots;\quad
        F_i=\Phi_i+\Delta F_{i+1};\quad\ldots,
\]
où, chaque fois, les $\Phi_i$ sont les formes normales correspondant aux $F_i$. Comme les $F_i$ sont de degré $(p-i)$ en $z$, [4] et [6] donnent:
\begin{equation}
\begin{gathered}
        \Phi_i=(a_{i1}S_i+a_{i2}S_i^2+\cdots)F_i,\\
        S_iF_i=\frac1{(p-i)!}
        \left[z\left(\frac{\p}{\p\lambda_1}\frac{\p}{\p x}+
        \frac{\p}{\p\lambda_2}\frac{\p}{\p y}\right)\right]^{p-i}
        \left[(\lambda_1x+\lambda_2y)\frac{\p}{\p z}\right]^{p-i}F_i.
\end{gathered}
\tag{8}
\end{equation}
À la formule [7] correspond
\[
        F_i=(\alpha_{i1}\Omega+\alpha_{i2}\Omega\Delta\Omega+\cdots)F_{i-1},
        \quad\text{et, par itération,}
\]
\begin{equation}
        F_i=(\beta_1\Omega^i+\beta_2\Omega^i\Delta\Omega+\cdots
        \beta_\rho\Omega^{k_1}\Delta\Omega^{k_2}\cdots\Delta\Omega^{k_\lambda}+\cdots)F
        \quad(k_1+k_2+\cdots+k_\lambda-\lambda+1=i).
\tag{9}
\end{equation}
\emph{Au moyen de} [8] \emph{et} [9], \emph{le développement}
\[
        F=\Phi+\Delta\Phi_1+\Delta^2\Phi_2+\cdots+\Delta^i\Phi_i+\cdots\Delta^p\Phi_p
\]
\emph{a atteint une représentation explicite dans le sens indiqué.}

Si l'on applique à [7] la transformation identique $\Omega\Delta F=PF$, où $P$ désigne des procédés polaires,\footnote{Comparer par exemple F. Mertens, \emph{Über ganze Funktionen von $m$ Systemen von je $n$ Unbestimmten}, Monatsh. f. Math. u. Phys. 4, formule (5).} alors, comme ces procédés commutent avec le procédé $\Omega$, on obtient la formule connue
\[
        F_1=P\Omega F;\quad\ldots;\quad F_i=P\Omega^iF;
\]
et de là, conjointement avec [8] ou avec les considérations loc. cit., résulte la formulation connue du développement en série, par exemple celle de (3) loc. cit. Si l'on a affaire à $n$ séries de $n$ variables, $x^{(1)},x^{(2)},\ldots,x^{(n-1)},z$, alors correspondamment
\[
        f_i(\xi)=\lambda_1x_i^{(1)}+\lambda_2x_i^{(2)}+\cdots+\lambda_{n-1}x_i^{(n-1)};
\]
et de même $\Delta$ et $\Omega$ doivent être définis pour $n$ séries.

2. Pour la forme normale de $F(t_i,t_{ik},t_{ikl},\ldots)$, les fonctions $f_i(\xi),f_{ik}(\xi),\ldots$ sont données par les déterminants à une série, à deux séries, jusqu'à $(n-1)$ séries:
\[
        \xi_i^{(1)},\quad (\xi^{(1)}\xi^{(2)})_{ik},\quad\ldots,\quad
        (\xi^{(1)}\xi^{(2)}\cdots\xi^{(n-1)})_{i_1i_2\cdots i_{n-1}}.
\]
Ainsi, comme spécialisation de [5], on obtient:
\[
\begin{gathered}
        S=AB;\qquad
        BF=F\bigl(\xi_i^{(1)},(\xi^{(1)}\xi^{(2)})_{ik},\ldots\bigr),\\
        ABF=\left(\sum t_i\frac{\p}{\p\xi_i^{(1)}}\right)^{\lambda_1}
        \left(\sum t_{ik}\left(\frac{\p}{\p\xi^{(1)}}\frac{\p}{\p\xi^{(2)}}\right)_{ik}\right)^{\lambda_2}
        \cdots F\bigl(\xi_i^{(1)},(\xi^{(1)}\xi^{(2)})_{ik},\ldots\bigr),
\end{gathered}
\]
où $\lambda_1,\lambda_2\ldots$ désignent les degrés en $t_i,t_{ik},\ldots$.

D'après des considérations de théorie des invariants (à savoir que tous les invariants de $F$ qui sont linéaires dans les coefficients peuvent être composés linéairement à partir de $F$ et de formes de $M$, donc en particulier aussi de $SF$, qui n'appartient pas à $M$)\footnote{Cela résulte, par exemple, des \S~6 et \S~7, 2. de mon travail cité \emph{Zur Invariantentheorie der Formen von $n$ Variabeln}.} on obtient ici la congruence $\Phi\equiv a_1SF\modu{M}$; et, à cause de la propriété de forme normale, on en tire, à la place de [4], la relation simple
\begin{equation}
        \Phi=a_1SF=a_1S\Phi.
\tag{10}
\end{equation}
La seconde égalité vaut parce que l'application de $S$ fait disparaître toute forme de $M$. La formule [10] montre que la forme normale $\Phi$ est elle-même en même temps une fonction propre; et, d'après la considération de théorie des invariants ci-dessus, il n'y a précisément pas d'autres fonctions propres qui soient en même temps des invariants de $F$ (comparer, au sujet de [10], Deruyts, formules (10) et (10')).

\begin{center}
\textbf{Rectifications.}
\end{center}

1. Dans le travail \og{} Die allgemeinsten Bereiche aus ganzen transzendenten Zahlen \fg{} (\emph{Math. Ann.} 77), p. 121, ligne 8 du haut, ligne 9 du haut et ligne 11 du bas, le corps $(H,Y)$ doit être remplacé par le \og{} corps $(H,Y,Y',\ldots,Y^{(\lambda_0-1)})$, où $Y',\ldots,Y^{(\lambda_0-1)}$ désignent les conjugués de $Y$ relativement à $H$. \fg{} De même, p. 120, ligne 5 du haut, et p. 121, ligne 9 du haut, $(H,y)$ doit être remplacé par $(H,y,y',\ldots,y^{(\lambda_0-1)})$.

2. Dans le travail \og{} Gleichungen mit vorgeschriebener Gruppe \fg{} (\emph{Math. Ann.} 78), dans la formulation du théorème d'irréductibilité de Hilbert --- p. 222, ligne 3 du bas --- au lieu de \og{} corps de nombres $\Omega$ \fg{}, il faut dire plus précisément: \og{} corps de nombres algébrique fini $\Omega$ \fg{}, comme M. Ostrowski me l'a fait remarquer. En effet, par exemple, relativement au corps de tous les nombres algébriques, le groupe de toute équation à coefficients numériques algébriques est le groupe identique. --- De même, dans l'introduction, par \og{} domaine de rationalité arbitraire \fg{}, il faut entendre un corps de nombres algébrique fini auquel on peut encore adjoindre un nombre fini de paramètres, ou l'un des corps intermédiaires ainsi obtenus.

\begin{center}
(Accepté en septembre 1919.)
\end{center}
\clearpage
\setcounter{footnote}{0}
\section*{\NoetherNavAnchor{P17}{17. Modules dans des domaines non commutatifs, provenant notamment d’expressions différentielles et d’expressions aux différences}17. Modules dans des domaines non commutatifs, provenant notamment d'expressions différentielles et d'expressions aux différences}
\begin{center}
Emmy Noether et Werner Schmeidler, à Göttingen.\\[0.65em]
\emph{Mathematische Zeitschrift} 8 (1920), p. 1--35
\end{center}

\begin{center}
\textbf{Table des matières.}
\end{center}
\begin{center}
\begin{minipage}{0.84\linewidth}
Introduction.\\
\S~1. Généralités formelles sur les expressions différentielles et aux différences.\\
\S~2. L'anneau de polynômes général.\\
\S~3. Modules d'un seul côté et leurs groupes quotients.\\
\S~4. Réductibilité des groupes quotients en deux sous-groupes et leurs relations avec le module.\\
\S~5. Calcul avec les unités et réductibilité en $k$ sous-groupes.\\
\S~6. Théorèmes de finitude.\\
\S~7. Un cas de décomposition univoque.\\
\S~8. Décomposition additive de groupes complètement réductibles en groupes premiers. Le théorème de l'isomorphie.\\
\S~9. Rapport entre les notions \og{} isomorphe \fg{} et \og{} de même espèce \fg{}.\\
\S~10. Existence d'une infinité de décompositions d'un groupe.\\
\S~11. Relations avec les systèmes d'équations différentielles et leurs intégrales.\\
\S~12. Exemples.
\end{minipage}
\end{center}

\begin{center}\textbf{Introduction.}\end{center}

La théorie formelle des expressions différentielles d'une variable conduit jusqu'à une décomposition en facteurs irréductibles qui est univoque en un certain sens.\footnote{Cf. E. Landau, \emph{Ein Satz über die Zerlegung homogener linearer Differentialausdrücke in irreduzible Faktoren}, \emph{Journal für die reine und angewandte Mathematik} 124 (1902), p. 115--120, et, à sa suite, A. Loewy, \emph{Über reduzible lineare homogene Differentialgleichungen}, \emph{Mathematische Annalen} 56 (1903), p. 549--584.} Mais tandis que le théorème correspondant de l'algèbre se transporte aux polynômes de plusieurs variables, cela n'est plus possible pour les expressions différentielles partielles.\footnote{Cf. un contre-exemple de Landau publié dans H. Blumberg, \emph{Über algebraische Eigenschaften von linearen homogenen Differentialausdrücken}, dissertation inaugurale, Göttingen 1912, 52 pages, p. 51--52.} Or, pour les expressions différentielles d'une variable, à côté de la représentation en produit apparaît encore une représentation comme plus petit multiple commun, où des lois analogues valent: dans deux décompositions différentes, les constituants particuliers peuvent en effet être associés deux à deux; ils sont \og{} de même espèce \fg{}.\footnote{Cf., outre les travaux cités, A. Loewy, \emph{Über vollständig reduzible lineare homogene Differentialgleichungen}, \emph{Mathematische Annalen} 62 (1906), p. 89--117; cité comme Loewy.}

Cette notion de plus petit multiple commun peut maintenant être comprise comme une notion de module et être ainsi transportée à $n$ variables; on obtient ainsi, dans le présent travail, la généralisation naturelle des théorèmes mentionnés. Au-delà de cela, nous remarquons encore que, des expressions différentielles, on n'utilise que le fait suivant: on peut les considérer comme des polynômes à multiplication non commutative, pour lesquels le degré du produit est égal à la somme des degrés des facteurs. Ces conditions sont satisfaites, par exemple, aussi par les expressions aux différences; les théorèmes valent donc aussi pour celles-ci, où ils n'étaient également connus jusqu'ici que pour une variable.\footnote{Cf. G. Wallenberg--A. Guldberg, \emph{Theorie der linearen Differenzengleichungen}, Leipzig et Berlin 1911.}

Nous développons donc en général une théorie des modules dont les éléments sont des polynômes à multiplication non commutative, et nous ramenons pour cela l'étude des modules à celle des classes modulo ces modules.\footnote{Pour le cas des domaines commutatifs, cf. W. Schmeidler, \emph{Über Moduln und Gruppen hyperkomplexer Größen}, \emph{Mathematische Zeitschrift} 3 (1919), p. 29--42.} À la différence du cas spécial commutatif, on peut certes parler ici aussi de la somme de classes modulo un module, mais non de leur produit; on peut seulement parler du produit d'une telle classe par un polynôme. La notion de \og{} groupe quotient \fg{} (historiquement \emph{Restgruppe}: groupe des classes modulo le module) doit donc être modifiée par rapport au cas commutatif. Dans le cas spécial des polynômes d'une variable, le groupe quotient est en outre \og{} fini \fg{}, c'est-à-dire qu'il ne possède qu'un nombre fini de classes modulo le module, linéairement indépendantes. En général, un tel \og{} ordre \fg{} fini n'existe pas; les groupes sont \og{} infinis \fg{}. Pour nos méthodes, toutefois, cette différence n'entre pas en considération.

Comme dans le cas commutatif, à la représentation d'un module comme plus petit multiple commun de modules premiers entre eux correspond le procédé plus simple de la décomposition additive du groupe quotient. Ici l'isomorphie des groupes quotients se révèle essentielle; quand on spécialise aux expressions différentielles d'une variable, cette isomorphie devient la notion de même espèce pour les modules associés, tandis que, quand on spécialise aux modules commutatifs, elle devient l'identité. D'ailleurs, la notion de même espèce se transporte aussi à $n$ variables, et son accord avec la notion d'isomorphisme peut être démontré. Que, dans la décomposition additive du groupe quotient, il n'apparaisse toujours qu'un nombre fini de termes, cela résulte du transfert du théorème de la base de Hilbert, qui repose sur la méthode de démonstration de Gordan; la méthode de démonstration originelle de Hilbert exigerait des restrictions sur les lois de multiplication (dans la formule (6), à droite, $a$ au lieu de $a_i$, ce qui est réalisé pour les expressions différentielles, mais non pour les expressions aux différences).

Pour parvenir maintenant à un théorème d'unicité, il faut spécialiser les modules --- comme le fait aussi Loewy --- en modules \og{} complètement réductibles \fg{}, c'est-à-dire en ceux qui sont représentables comme plus petits multiples communs de modules premiers. Pour deux décompositions différentes, les constituants sont alors ou bien identiques deux à deux, ou bien du moins à chaque constituant de l'une des décompositions correspond un constituant isomorphe de l'autre, et réciproquement; en même temps chaque décomposition contient alors au moins deux constituants isomorphes. En outre, de l'apparition de deux constituants isomorphes dans une seule et même décomposition résulte l'existence d'une infinité de décompositions, et cela vaut même sans la restriction aux modules premiers, ce qui n'était jusqu'ici pas connu même dans le cas spécial des expressions différentielles d'une variable.

Enfin, pour les modules provenant d'expressions différentielles, nous abordons le rapport avec les intégrales et montrons que, pour les groupes quotients finis, l'ordre est égal au nombre maximal d'intégrales linéairement indépendantes; ainsi des fonctions arbitraires apparaissent dans l'intégrale générale exactement lorsque le groupe quotient n'est pas fini. La notion selon laquelle deux modules sont de même espèce signifie ici, comme dans le cas d'une variable, une parenté \og{} birationnelle \fg{} entre les intégrales. Quelques exemples terminent le travail.

La première impulsion de ce travail fut donnée, il y a plusieurs années, par une question de E. Landau à E. Noether au sujet de la généralisation du théorème du produit. E. Noether remarqua alors seulement que la question relevait de la théorie des modules dans l'anneau de polynômes non commutatif, mais elle ne pouvait pas encore être abordée. Les théorèmes de Lasker sur les modules, connus alors, n'étaient pas directement transférables à cet anneau, parce que la méthode de démonstration de Lasker repose sur le théorème selon lequel un polynôme se décompose univoquement en facteurs irréductibles.\footnote{E. Noether a remarqué récemment que les théorèmes de Lasker peuvent eux aussi être établis, sans utiliser ce théorème, par des méthodes apparentées à celles employées ici.} Seules les méthodes de Schmeidler offraient la possibilité d'un transfert, que nous avons alors réalisé en commun. Dans le détail, mentionnons que la généralisation du groupe quotient, le transfert et l'application du théorème de finitude de Hilbert, ainsi que l'existence d'une infinité de décompositions, sont dus à E. Noether; la généralisation de la notion loewyenne de réductibilité complète, la notion d'isomorphisme et son rapport avec la notion de même espèce sont dus à W. Schmeidler.

\subsection*{\S~1. Généralités formelles sur les expressions différentielles et aux différences.}

1. Pour les expressions différentielles d'une variable, on a introduit une formation symbolique du produit. Soient
\[
 a_0(x)\frac{\dd^n y}{\dd x^n}+a_1(x)\frac{\dd^{n-1}y}{\dd x^{n-1}}+\cdots+a_n(x)y=A(y),
\]
\[
 b_0(x)\frac{\dd^m y}{\dd x^m}+b_1(x)\frac{\dd^{m-1}y}{\dd x^{m-1}}+\cdots+b_m(x)y=B(y).
\]
On définit alors le produit symbolique $AB(y)$ comme l'expression différentielle
\[
 a_0(x)\frac{\dd^n B(y)}{\dd x^n}+a_1(x)\frac{\dd^{n-1}B(y)}{\dd x^{n-1}}+\cdots+a_n(x)B(y).
\]
Cette formation de produit est en réalité une multiplication de deux opérateurs; on l'écrit le plus simplement en remplaçant, comme d'ordinaire, $\frac{\dd^n}{\dd x^n}$ par $\left(\frac{\dd}{\dd x}\right)^n$. Alors $AB$ devient
\[
 \left(a_0\left(\frac{\dd}{\dd x}\right)^n+\cdots+a_n\right)
 \left(b_0\left(\frac{\dd}{\dd x}\right)^m+\cdots+b_m\right)y,
\]
où, dans l'exécution du produit, valent les règles ordinaires de différentiation des sommes et des produits. Tout tel opérateur peut maintenant être considéré comme un polynôme en la seule variable $\frac{\dd}{\dd x}=\xi$, à condition de ne pas écrire explicitement la fonction à laquelle l'opérateur est appliqué et de fixer, pour la variable $\xi$, la règle de multiplication correspondant à la différentiation d'un produit:
\[
 \frac{\dd}{\dd x}\bigl(a(x)\cdot y\bigr)=a(x)\frac{\dd}{\dd x}(y)+a'(x)y,
\]
ou, dans notre notation $(n=1,m=0)$,
\begin{equation}
 \xi\cdot a=a\cdot\xi+a'.
\tag{1}
\end{equation}

On procède tout à fait de même pour les expressions différentielles partielles linéaires
\[
 \left(\sum a_{\alpha_1\ldots\alpha_n}(x_1,\ldots,x_n)
 \frac{\p^{\alpha_1+\cdots+\alpha_n}}
 {\p x_1^{\alpha_1}\cdots\p x_n^{\alpha_n}}\right)y
\]
d'une fonction $y$. Pour les $n$ opérateurs $\frac{\p}{\p x_1},\ldots,\frac{\p}{\p x_n}$, nous introduisons des variables $\xi_1,\ldots,\xi_n$ et obtenons ainsi le polynôme
\[
 \sum a_{\alpha_1\ldots\alpha_n}(x_1,\ldots,x_n)
 \xi_1^{\alpha_1}\cdots\xi_n^{\alpha_n}.
\]
Le produit de deux polynômes se calcule alors suivant les lois
\begin{equation}
 \xi_i a=a\xi_i+a_i \qquad (i=1,\ldots,n),
\tag{2}
\end{equation}
où $a_i$ désigne la dérivée partielle de $a$ par rapport à $x_i$; les produits des $\xi$ entre eux sont commutatifs. À l'exception de la loi commutative de la multiplication, toutes les lois de l'addition et de la multiplication valent pour les polynômes ainsi définis comme dans l'algèbre ordinaire. En particulier, pour le produit de deux tels polynômes, le degré total, ainsi que le degré en chaque variable particulière, est égal à la somme des degrés correspondants des deux facteurs.

2. Pour les expressions aux différences (cf. Wallenberg, loc. cit.) de la forme
\[
 a_x^{(0)}y_{x+n}+a_x^{(1)}y_{x+n-1}+\cdots+a_x^{(n)}y_x=A_x,
\]
\[
 b_x^{(0)}y_{x+m}+b_x^{(1)}y_{x+m-1}+\cdots+b_x^{(m)}y_x=B_x,
\]
on a une définition correspondante du produit
\[
 (AB)_x=a_x^{(0)}B_{x+n}+a_x^{(1)}B_{x+n-1}+\cdots+a_x^{(n)}B_x,
\]
qui peut être effectuée à l'aide de la règle de multiplication
\[
 (ay)_{x+1}=a_{x+1}y_{x+1}.
\]
Si nous introduisons l'opérateur $\xi$ pour le passage de $x$ à $x+1$, alors nous obtenons $\xi\{ay\}=a^*\xi\{y\}$, où $a_x^*=a_{x+1}$. De nouveau on peut supprimer la fonction $y$ à laquelle l'opérateur est appliqué et l'on obtient la loi
\begin{equation}
 \xi a=a^*\xi,
\tag{3}
\end{equation}
où, avec $a$, $a^*$ ne s'annule pas non plus identiquement. De façon tout à fait analogue, pour $n$ variables,
\begin{equation}
 \xi_i a=a_i\xi_i\qquad (i=1,\ldots,n),
\tag{4}
\end{equation}
où $\xi_i$ désigne le passage de $x_i$ à $x_i+1$, et où $a_i$ désigne la fonction obtenue à partir de $a$ en remplaçant $x_i$ par $x_i+1$; elle ne s'annule donc pas identiquement si $a$ ne le fait pas. Avec ces conventions, toutes les règles d'addition et de multiplication, y compris la règle concernant le degré d'un produit, valent à nouveau, à l'exception de la loi commutative.

Les deux cas seront subordonnés, dans le paragraphe suivant, à un anneau de polynômes défini de façon tout à fait générale avec des coefficients arbitraires.

\subsection*{\S~2. L'anneau de polynômes général.}

Nous partons d'un corps $\mathsf{P}$ défini abstraitement de manière arbitraire (dont nous désignons les éléments par des minuscules latines), pour lequel valent toutes les règles de l'algèbre, en particulier aussi la loi commutative de la multiplication avec inversion univoque. À ce corps nous adjoignons $n$ indéterminées $\xi_1,\ldots,\xi_n$; c'est-à-dire que nous considérons le domaine d'intégrité $J$ constitué par les produits $a\xi_{\nu_1}b\xi_{\nu_2}\ldots c\xi_{\nu_\rho}$ et par leurs sommes finies, toutes les lois d'addition et de multiplication devant valoir à l'exception de la loi commutative de la multiplication. En particulier, nous appelons \og{} polynômes \fg{} les sommes finies de la forme
\[
 F(\xi_1,\ldots,\xi_n)=\sum a_{\nu_1\ldots\nu_\rho}\xi_{\nu_1}\cdots\xi_{\nu_\rho}
\]
(elles seront toujours désignées par des majuscules latines). À la place de la loi commutative, nous soumettons maintenant les éléments de $J$ à des lois de produit telles que \emph{le produit de deux polynômes soit de nouveau un polynôme}; le domaine de tous les polynômes de $J$ épuise donc déjà $J$. Pour cela il est manifestement nécessaire et suffisant que, pour les éléments $\xi_1,\ldots,\xi_n$ et un élément arbitraire $a$ de $\mathsf{P}$, valent des règles de produit de la forme
\begin{equation}
 \xi_i a=F_i(\xi_1,\ldots,\xi_n)\qquad (i=1,\ldots,n),
\tag{5}
\end{equation}
où $F_i(\xi_1,\ldots,\xi_n)$ désigne un polynôme quelconque dépendant naturellement de $\xi_i$ et de $a$. (En particulier, les règles (2) et (4) pour les expressions différentielles et pour les expressions aux différences sont respectivement de cette forme.) Car, puisque $\xi_i a$ n'est pas en soi un polynôme, mais appartient par hypothèse au domaine, l'équation est nécessaire. Réciproquement, par application finiment répétée de la formule (5), chaque élément du domaine $J$ peut être transformé en un polynôme.

Pour l'unité du corps, que nous désignons par 1, nous exigerons en outre $\xi_i\cdot1=1\cdot\xi_i=\xi_i$. Nous stipulerons aussi que la multiplication de $\xi_1,\ldots,\xi_n$ entre eux est commutative.\footnote{Cela n'est nécessaire qu'à partir de \S~6 pour le théorème de finitude et ses conséquences.} Chaque polynôme de $J$ s'écrit alors exactement comme un polynôme ordinaire en $\xi_1,\ldots,\xi_n$ à coefficients dans $\mathsf{P}$, et deux polynômes ne coïncident que lorsque tous leurs coefficients coïncident.

Plus précisément, dans ce qui suit,\footnote{Ceci aussi n'est utilisé qu'à partir de \S~6.} nous exigeons qu'à la place de (5) vaille la règle
\begin{equation}
 \xi_i a=a_i\xi_i+b_i\qquad (a_i\ne0)\quad(i=1,\ldots,n).
\tag{6}
\end{equation}
Ici encore les lois (2) et (4) sont en particulier incluses. Alors le degré d'un produit est égal à la somme des degrés des facteurs; ceci vaut tant pour le degré total que pour le degré en chaque variable particulière, de sorte que le produit de deux polynômes qui ne s'annulent pas identiquement est également différent de 0.

\subsection*{\S~3. Modules d'un seul côté et leurs groupes quotients.}

Pour notre domaine $J$, puisque la multiplication n'est pas commutative, les espèces suivantes de modules\footnote{Nous désignons les modules par des majuscules allemandes.} peuvent être distinguées:\footnote{Les idéaux d'un seul côté ont déjà été considérés par A. Hurwitz; cf. \emph{Vorlesungen über die Zahlentheorie der Quaternionen} (Berlin, Springer 1919), sixième leçon. Toutefois, il s'agit là essentiellement d'idéaux principaux; il en est de même des travaux de Du Pasquier cités dans la note 1 de ce livre.}

1. Le \emph{module du côté droit}, qui, avec $M$, contient aussi $FM$ pour un polynôme arbitraire $F$, et qui, avec $M_1$ et $M_2$, contient aussi $M_1+M_2$.

2. Le \emph{module du côté gauche}, qui, avec $M$, contient aussi $MF$, et qui, avec $M_1$ et $M_2$, contient aussi $M_1+M_2$.

Nous désignons ces deux espèces par le nom commun de \emph{modules d'un seul côté}; les modules qui sont à la fois du côté droit et du côté gauche seront appelés \emph{bilatères}. Dans ce qui suit nous nous limitons aux modules du côté droit et remarquons que la théorie de ceux du côté gauche peut être édifiée exactement de la même façon.

Nous disons que deux polynômes $F$ et $G$ sont congrus modulo $\mM$, et écrivons $F\equiv G(\mM)$, si leur différence est divisible par $\mM$, c'est-à-dire si elle appartient à $\mM$. L'ensemble de tous les polynômes congrus à un polynôme donné forme une \emph{classe modulo $\mM$}. (Les classes modulo $\mM$ seront désignées par de petites lettres allemandes.) Puisque de
\[
 F_1\equiv F_2(\mM)\qquad\text{et}\qquad G_1\equiv G_2(\mM)
\]
on déduit que $F_1+G_1\equiv F_2+G_2(\mM)$, la somme de deux classes modulo $\mM$, $\mx$ et $\my$, est définie et est de nouveau une classe modulo $\mM$, $\mx+\my$. En revanche, parce que le module est d'un seul côté, on ne peut pas parler du produit de deux telles classes. En effet, de
\[
 F_1=F_2+M_1,\qquad G_1=G_2+M_2
\]
on obtient seulement
\[
 F_1G_1=F_2G_2+F_2M_2+M_1G_2+M_1M_2;
\]
mais comme $M_1G_2$ n'a pas à appartenir à $\mM$ lorsque $M_1$ appartient à $\mM$, $F_1G_1$ et $F_2G_2$ ne sont pas nécessairement congrus en général. En revanche, si $M_1=0$, donc $F_1=F_2=F$, on voit que $G_1\equiv G_2$ implique $FG_1\equiv FG_2$. Ainsi une classe arbitraire $\my$ modulo $\mM$ peut être multipliée par devant par un polynôme $F$, et cela donne une classe modulo $\mM$ bien définie, $F\my=\mz$.

Pour le calcul avec les classes modulo $\mM$ valent les lois commutative et associative de l'addition, ainsi que la loi de réversibilité univoque de l'addition; en outre, pour la multiplication de ces classes par des polynômes et l'addition, vaut la loi distributive $F(\my_1+\my_2)=F\my_1+F\my_2$, comme il résulte directement du calcul avec les polynômes modulo $\mM$.

La classe modulo $\mM$ à laquelle appartient l'unité sera appelée classe unité, ou unité, et notée $\me$. Alors $F\me=\mx$ est la classe modulo $\mM$ à laquelle appartient le polynôme $F$. Contrairement aux classes modulo $\mM$ en général, le produit $\mx\me$ est cependant ici aussi bien défini et égal à $\mx$; car de
\[
 F_1=F_2+M_1,\qquad G_1=1+M_2
\]
et de $M_1\cdot1=M_1$ on déduit que
\[
 F_1G_1=F_2+M.
\]
L'ensemble des classes modulo $\mM$ s'appelle le \emph{groupe quotient de} $\mM$. Un groupe quotient possède donc les propriétés suivantes: 1. avec chaque classe $\mx$ et un polynôme arbitraire $F$, il contient la classe $F\mx$; et avec les classes $\mx_1$ et $\mx_2$, il contient la classe $\mx_1+\mx_2$.\footnote{Cette propriété pourrait, en un sens plus général, être appelée la \og{} propriété de module \fg{} du groupe quotient.} 2. Il possède une unité $\me$, de sorte que $F\me$ représente, pour chaque polynôme $F$, la classe $\mx$ de $F$ modulo $\mM$. \emph{Ainsi, lorsque} $F$ \emph{parcourt tous les polynômes,} $F\me$ \emph{parcourt tout le groupe quotient.}

\subsection*{\S~4. Réductibilité du groupe quotient en deux sous-groupes et sa relation avec le module.}

Pour les modules d'un seul côté, les notions de \emph{divisibilité}, de \emph{plus grand commun diviseur} et de \emph{plus petit multiple commun} demeurent les mêmes que dans le cas bilatère, comme le montrent les définitions suivantes.

Un module $\mM$ est dit divisible par un autre module $\mN$ si $M\equiv0(\mM)$ entraîne aussi $M\equiv0(\mN)$.

Pour deux modules $\mM$ et $\mN$, le \emph{plus grand commun diviseur} $(\mM,\mN)$ est défini comme l'ensemble de tous les polynômes représentables sous la forme $M+N$, où $M\equiv0(\mM)$ et $N\equiv0(\mN)$; c'est de nouveau un module. Il en va de même pour plusieurs modules.

L'ensemble de tous les polynômes divisibles à la fois par $\mM$ et par $\mN$ forme un module $[\mM,\mN]$, le \emph{plus petit multiple commun} de $\mM$ et de $\mN$. Cette définition s'applique aussi de façon correspondante à plusieurs modules, et même à un ensemble de modules de cardinalité arbitraire.

Deux modules sont dits premiers entre eux lorsque leur plus grand commun diviseur est le module unité, c'est-à-dire l'ensemble de tous les polynômes.

Nous partons maintenant d'un module $\mM$ qui est représentable comme plus petit multiple commun de deux modules premiers entre eux $\mN$ et $\mL$, supposés tous deux être des diviseurs propres:
\begin{equation}
 \mM=[\mN,\mL].
\tag{7}
\end{equation}
Nous examinons le groupe quotient $\mG$ de $\mM$. Par hypothèse,
\begin{equation}
 1\equiv N+L(\mM)
\tag{8}
\end{equation}
pour deux polynômes $N\equiv0(\mN)$ et $L\equiv0(\mL)$. De là nous montrons
\begin{equation}
 \mL=(\mM,L),\qquad \mN=(\mM,N),
\tag{9}
\end{equation}
où, par exemple, $(\mM,L)$ désigne le module formé de tous les polynômes de la forme $M+FL$. En effet, d'abord $(\mM,L)\equiv0(\mL)$. Réciproquement, il résulte de (8), pour un polynôme arbitraire $F$, que
\begin{equation}
 F\equiv FN(\mL).
\tag{10}
\end{equation}
Si maintenant $F\equiv0(\mL)$, alors $FN\equiv0(\mL)$, donc $\equiv0(\mM)$, et par conséquent $F\equiv FL(\mM)$. On démontre de la même manière $\mN=(\mM,N)$. Il en résulte incidemment, puisque $\mN$ et $\mL$ sont des diviseurs propres de $\mM$, que ni $L\equiv0$ ni $N\equiv0(\mM)$ ne peut avoir lieu.

Soit $\ma$ la classe de $N$ modulo $\mM$, et $\mb$ la classe de $L$ modulo $\mM$. Alors, d'après (8),
\[
 \me=\ma+\mb\qquad(\ma\ne0,\ \mb\ne0).
\]
Comme chaque classe modulo $\mM$ est de la forme $F\me$, la loi distributive donne, pour chaque classe modulo $\mM$,
\begin{equation}
 F\me=F\ma+F\mb\qquad(\ma\ne0,\ \mb\ne0).
\tag{11}
\end{equation}
Cette représentation d'une classe arbitraire modulo $\mM$ comme somme de deux classes de la forme $G\ma$ et $H\mb$ est cependant \emph{unique}. Car, à partir d'une relation $0=G\ma+H\mb$, si $K$ est un polynôme représentant la classe $G\ma=-H\mb$, alors manifestement $K\equiv GN(\mM)$, $K\equiv-HL(\mM)$, donc $K\equiv0(\mN)$, $K\equiv0(\mL)$, et par suite $K\equiv0(\mM)$; c'est-à-dire $G\ma=0$, $H\mb=0$.

Réciproquement, supposons maintenant que la relation (11) soit satisfaite pour chaque polynôme $F$, et cela de manière unique au sens précédent. Soient alors $N$ un polynôme de $\ma$ et $L$ un polynôme de $\mb$. Nous formons les deux modules (9) et affirmons (7) et (8). Cette dernière relation résulte de (11) pour $F=1$. En outre, par la définition de $\mN$ et de $\mL$, $\mM$ est un multiple commun de ces deux modules, et même le plus petit. Car $F\equiv0(\mL)$ et $F\equiv0(\mN)$ entraînent $F\equiv HL\equiv GN(\mM)$; donc $G\ma=H\mb$, $G\ma-H\mb=0$, et par suite, par l'unicité de la représentation, $G\ma=0$, $H\mb=0$, $F\equiv0(\mM)$. Il est donc montré que la \emph{représentation d'un module} $\mM$ \emph{comme plus petit multiple commun de deux modules premiers entre eux est équivalente à la décomposition additive unique du groupe quotient} $\mG$ \emph{de} $\mM$, \emph{telle qu'elle est donnée par la formule} (11).

Nous examinons maintenant le système $\mA$ des classes modulo $\mM$ de la forme $F\ma$, que nous voulons mettre en relation avec le groupe quotient $\mbarA$ de $\mL$. D'après (11), chaque polynôme $F$ de $\mL$ appartient à la classe $F\ma$ modulo $\mM$, tandis que $\bara$ est la classe de $N$ modulo $\mL$, et donc la classe unité de $\mbarA$. Si nous associons maintenant aux classes $F\ma$ de $\mA$ les classes $F\bara$ de $\mbarA$, alors à chaque classe de $\mA$ est assignée une classe de $\mbarA$, de telle sorte que la somme de deux classes de $\mA$ corresponde à la somme des classes assignées de $\mbarA$, et que le produit d'une classe de $\mA$ par un polynôme corresponde au produit de la classe assignée de $\mbarA$ par le même polynôme. Cette assignation est en outre biunivoque: $F\ma=0$ entraîne $F\equiv0(\mL)$ d'après (11), donc $F\bara=0$; réciproquement, $F\bara=0$, donc $F\equiv0(\mathfrak L)$, signifie aussi, d'après (10), que $FN\equiv0(\mL)$, donc $FN\equiv0(\mM)$, c'est-à-dire $F\ma=0$.

Cela nous conduit à une notion fondamentale, que nous définissons comme suit:

\emph{Une correspondance biunivoque entre deux systèmes $\mA$ et $\mbarA$ de classes modulo leurs modules respectifs s'appelle isomorphe si, à la somme de deux classes de $\mA$, est assignée la somme des classes correspondantes de $\mbarA$, et si, au produit d'une classe arbitraire de $\mA$ par un polynôme, est assigné le produit de la classe correspondante de $\mbarA$ par le même polynôme.}

Un sous-système $\mA$ de classes modulo $\mM$ dans $\mG$ qui est isomorphe au groupe quotient d'un module s'appelle un \emph{sous-groupe} de $\mG$. Comme le montre la preuve précédente, l'assignation est ici plus précisément telle que chaque classe de $\bar{\mathfrak A}$ provient de la classe correspondante de $\mbarA$ par adjonction de certains polynômes, à savoir de tous les polynômes appartenant à $\mL$ mais non à $\mM$.

De même, on montre que l'ensemble de toutes les classes modulo $\mM$ de la forme $F\mb$ constitue un sous-groupe $\mB$ de $\mG$ qui est isomorphe au groupe quotient $\mbarB$ de $\mN$.

Un groupe quotient $\mG$ dont les classes admettent une décomposition additive unique (11) sera dit \emph{réductible} en les deux sous-groupes $\mA$ et $\mB$, ce que l'on écrit $\mG=\mA+\mB$; un groupe qui n'est pas réductible s'appelle \emph{irréductible}. Nous appelons aussi parfois les modules associés réductibles ou irréductibles. Il s'ensuit donc ce qui suit.

\textbf{Théorème 1.}\footnote{Cf. Schmeidler, loc. cit., p. 39.} \emph{Un module $\mM$ est représentable comme plus petit multiple commun de deux diviseurs propres premiers entre eux $\mL$ et $\mN$ si et seulement si son groupe quotient $\mG$ est réductible en deux sous-groupes $\mA$ et $\mB$. Dans ce cas, $\mA$ et $\mB$ sont respectivement isomorphes aux groupes quotients de $\mL$ et de $\mN$.}
\subsection*{\S~5. Calcul avec les unités et réductibilité en $k$ sous-groupes.}

Les classes résiduelles $\ma$ et $\mb$ de $\mG$ ont en commun avec l'unité $\me$ la propriété que le produit de $\ma$ et de $\mb$ avec une classe résiduelle arbitraire $\mx$ est défini; nous appelons donc également $\ma$ et $\mb$ des unités de $\mG$. En effet, de (11), pour $F=M$, il résulte
\[
 0=M\me=M\ma+M\mb,
\]
et par suite, d'après l'unicité, $M\ma=0$, $M\mb=0$, et donc
\[
 (F+M)\ma=F\ma=\mx\ma,
\]
où $\mx$ désigne la classe résiduelle de $F$. À la place de la relation (11), nous pouvons donc écrire la relation de classes équivalente
\[
 \mx\me=\mx\ma+\mx\mb.
\]
En particulier, en posant $\mx=\ma$, on obtient
\[
 \ma=\ma^2+\ma\mb,
\]
et de là, par unicité, $\ma^2=\ma$, $\ma\mb=0$. De même, on obtient $\mb^2=\mb$, $\mb\ma=0$.

Correspondant à la représentation d'un groupe quotient comme somme de deux constituants, nous pouvons maintenant définir aussi une telle représentation comme somme de $k$ constituants. Supposons donc donnée une représentation univoque
\[
 F\me=F\ma_1+F\ma_2+\cdots+F\ma_k\qquad(\ma_1\ne0,\ldots,\ma_k\ne0),
\]
où donc de $0=G_1\ma_1+\cdots+G_k\ma_k$ il résulte aussi $G_i\ma_i=0$. Alors, en particulier pour $F=M$, on obtient la relation $M\ma_i=0$, de sorte que la représentation peut aussi s'écrire sous la forme\footnote{Pour la décomposition additive des groupes quotients, les lois commutative et associative sont valables; cependant, contrairement au cas commutatif, $\mathfrak A_1+\mathfrak A_2=\mathfrak A_1+\mathfrak A_3$ n'entraîne pas $\mathfrak A_2=\mathfrak A_3$; cf. l'exemple 1 du \S~12.}
\begin{equation}
 \mx\me=\mx\ma_1+\cdots+\mx\ma_k\qquad(\ma_1\ne0,\ldots,\ma_k\ne0).
\tag{12}
\end{equation}
Cette représentation peut également être considérée comme issue d'une décomposition successive en deux constituants. Si l'on peut la poursuivre indéfiniment, avec tous les $\ma_i\ne0$, nous parlons d'une somme d'une infinité de constituants.

En posant $\mx=\ma_i$ dans (12), on obtient
\begin{equation}
 \ma_i^2=\ma_i,
 \qquad \ma_i\ma_j=0\quad(i\ne j),
 \qquad i,j=1,\ldots,k,
\tag{13}
\end{equation}
avec $\ma_i\ne0$. Réciproquement, supposons donné dans $\mG$ un système de classes résiduelles $\ma_1\ldots\ma_k$ pour lequel les conditions d'orthogonalité (13) sont satisfaites et pour lequel
\begin{equation}
 \me=\ma_1+\cdots+\ma_k
\tag{14}
\end{equation}
est valable. Nous allons montrer que cela donne une représentation du groupe comme somme de $k$ constituants. D'abord, il résulte de l'hypothèse, pour tout polynôme $A_i$ de $\ma_i$, puisque $A_i\ma_i=(A_i+M)\ma_i$, que $M\ma_i=0$; donc $\ma_i$ peut être multipliée par toute classe. Ainsi (14) entraîne la représentation (12), et cette représentation est univoque. Car si $0=\mx_1\ma_1+\cdots+\mx_k\ma_k$, alors la multiplication à droite par $\ma_i$ donne, d'après (13), $0=\mx_i\ma_i$, ce qui prouve l'unicité et donc l'assertion. Les classes résiduelles $\ma_1\ldots\ma_k$ qui interviennent dans la décomposition de la classe unité seront appelées les unités appartenant à la décomposition; de façon équivalente, elles sont caractérisées par les relations d'orthogonalité (13) et la condition (14). Les développements ultérieurs consisteront essentiellement à calculer avec les unités.

Il reste à établir le lien entre la décomposition additive du groupe en $k$ constituants et la décomposition correspondante du module. Comme dans le cas $k=2$, il en résultera par induction une représentation du module comme plus petit commun multiple de $k$ modules premiers entre eux.

Soit donnée la décomposition (12), que nous écrivons sous la forme
\begin{align}
 \mx\me&=\mx\mb+\mx\ma_k,\tag{15}\\
 \mx\mb&=\mx\ma_1+\cdots+\mx\ma_{k-1}.\tag{16}
\end{align}
D'après le théorème I, on déduit de (15) que les classes résiduelles $\mx\mb$ sont isomorphes au groupe quotient du module
\begin{equation}
 \mN=(\mM,A_k).
\tag{17}
\end{equation}
Pour leurs classes résiduelles $\mx\overline{\mb}$, la représentation correspondant à (16) est donc également valable:
\begin{equation}
 \mx\overline{\mb}=\mx\overline{\ma}_1+\cdots+\mx\overline{\ma}_{n-1}.
\tag{18}
\end{equation}
Supposons maintenant déjà démontré que
\begin{equation}
 \mN=[\mM_1,\ldots,\mM_{k-1}],
\tag{19}
\end{equation}
où
\begin{equation}
\left\{
\begin{aligned}
 \mM_1&=(\mN,A_2+\cdots+A_{k-1}),\\
 \mM_2&=(\mN,A_1+A_3+\cdots+A_{k-1}),\\
 &\hspace{1.8cm}\vdots\\
 \mM_{k-1}&=(\mN,A_1+\cdots+A_{k-2}).
\end{aligned}\right.
\tag{20}
\end{equation}
Pour deux termes de somme, cette assertion coïncide avec le résultat antérieur. En outre, les polynômes $A_i$ peuvent ici être choisis spécialement comme polynômes de $\overline{\ma}_i$ qui sont en même temps contenus dans $\ma_i$, puisque, d'après le \S~4, les polynômes de cette dernière classe résiduelle sont contenus dans la première. Alors, d'après (15),
\begin{equation}
 \mM=[\mN,\mM_k],
\tag{21}
\end{equation}
où
\[
 \mM_k=(\mM,A_1+\cdots+A_{k-1}).
\]
De (19) et (21) on obtient
\[
 \mM=[[\mM_1,
 \ldots,\mM_{k-1}],\mM_k]=[\mM_1,\ldots,\mM_{k-1},\mM_k]
\]
par la définition du plus petit commun multiple; et de (17) et (20) on obtient
\[
 \mM_1=(\mM,A_k,A_2+\cdots+A_{k-1}).
\]
Il en résulte que
\[
 \mM_1=(\mM,A_2+\cdots+A_k),
\]
car ce dernier module contient, d'après (13), aussi le polynôme
\[
 A_k\equiv A_k(A_2+\cdots+A_k)(\mM).
\]
On obtient des expressions correspondantes pour $\mM_2,\ldots,\mM_{k-1}$. En outre, les modules $\mM_1,\ldots,\mM_k$ sont premiers entre eux, puisque d'après (14)
\[
 \frac{1}{k-1}\Bigl((A_2+\cdots+A_k)+(A_1+A_3+\cdots+A_k)+\cdots+(A_1+\cdots+A_{k-1})\Bigr)\equiv1\pmod{\mM}.
\]
Ils forment de plus un système totalement premier entre eux, c'est-à-dire que le plus petit commun multiple d'un choix quelconque d'entre eux est premier avec le plus petit commun multiple des autres.\footnote{C'est une généralisation d'un concept introduit par Loewy.} Cela résulte directement en regroupant dans la formule (12) les termes de somme en deux groupes convenables.

Il est donc montré qu'à une décomposition du groupe quotient en $k$ termes de somme correspond une représentation du module comme plus petit commun multiple de $k$ modules qui forment un système totalement premier entre eux. Nous prouvons maintenant la réciproque, également par induction.

Supposons donc que
\[
 \mM=[\mM_1,\ldots,\mM_k]=[[\mM_1,\ldots,\mM_{k-1}],\mM_k]=[\mN,\mM_k],
\]
où $\mM_1,\ldots,\mM_k$ forment un système totalement premier entre eux. En combinaison avec le théorème I, il en résulte pour le groupe quotient la représentation (15). Les modules $\mM_1,\ldots,\mM_{k-1}$ forment alors eux aussi un système totalement premier entre eux. Car si, dans une certaine division des modules de ce système en deux groupes, le plus petit commun multiple de l'un des groupes et le plus petit commun multiple de l'autre avaient un diviseur commun, celui-ci resterait encore si l'on ajoutait le module $\mM_k$ à l'un des deux groupes et que l'on formait ensuite le plus petit commun multiple.\footnote{Si l'on poursuit ce raisonnement, on voit que parmi $\mM_1,\ldots,\mM_k$ deux quelconques sont premiers entre eux; la réciproque, à savoir que cela impliquerait que $\mM_1,\ldots,\mM_k$ forment un système totalement premier entre eux, n'est en général valable que dans le cas commutatif; cf. le contre-exemple 1 du \S~12.} Pour $k=2$, \og{}totalement premier entre eux\fg{} signifie la même chose que \og{}premier entre eux\fg{}; notre assertion peut donc être supposée démontrée jusqu'à $k-1$. Ainsi, pour la classe résiduelle $\mx\overline{\mb}$, la représentation univoque (18) est valable, donc aussi (16) par l'isomorphisme, et l'assertion s'ensuit.

\textbf{Théorème II.} \emph{Un module $\mM$ est représentable comme plus petit commun multiple de $k$ modules $\mM_1,\ldots,\mM_k$ qui forment un système totalement premier entre eux si et seulement si son groupe quotient $\mG$ est réductible en $k$ sous-groupes $\mU_1,\ldots,\mU_k$. Avec une notation convenable, $\mU_i$ est alors isomorphe au groupe quotient de $\mM_i$.}

\subsection*{\S~6. Théorèmes de finitude.}

Dans ce paragraphe, nous montrerons que tout groupe quotient peut être représenté comme somme d'un nombre fini de constituants irréductibles. Pour cela, nous transposons d'abord le théorème de la base de Hilbert.

Nous utiliserons désormais les hypothèses plus spéciales formulées au \S~2, à savoir que la multiplication des $\xi_1,\ldots,\xi_n$ entre eux est commutative et que la loi de multiplication des polynômes a la forme
\begin{equation}
 \xi_i a=a_i\xi_i+b_i,
 \qquad (a_i\ne0)\quad(i=1,\ldots,n).
\tag{6}
\end{equation}
Il suffit de mener la preuve seulement pour les modules; en effet, si $\mS$ est un système quelconque de polynômes, nous formons le module $\mM$ qui en dérive en adjoignant à $\mS$ tous les polynômes qui peuvent être composés à partir d'un nombre fini de polynômes $G_1,\ldots,G_k$ de $\mS$ sous la forme $A_1G_1+\cdots+A_kG_k$. Si $F_1,\ldots,F_\varrho$ est alors une base du module $\mM$, et si
\[
 F_i=\sum_j A_{ij}G_j\qquad(i=1,\ldots,l),
\]
alors les polynômes $G_j$ en nombre fini qui apparaissent dans ces représentations forment une base de $\mS$.

Nous prouvons le théorème de la base de Hilbert en suivant Gordan (\emph{Göttinger Nachrichten} 1899). La preuve se divise en deux parties. La première concerne un système $\mP$ d'une infinité de produits de puissances et reste essentiellement inchangée grâce à la commutativité de $\xi_1,\ldots,\xi_n$. Les produits de puissances
\[
 P=\xi_1^{a_1}\cdots\xi_n^{a_n},
\]
chacun d'eux n'étant inscrit dans le système qu'une seule fois, doivent être ordonnés par degré croissant et, à degré égal, d'une manière fixe quelconque. Nous considérons maintenant le sous-système $\mathfrak{Q}$ de $\mP$ constitué de tous les produits de puissances $P=Q$ qui ne sont divisibles par aucun $P$ précédent. Alors aucun $Q$ n'est divisible non plus par un $Q$ ultérieur, car les termes ultérieurs sont soit de degré plus élevé, soit, à degré égal, différents des précédents. D'autre part, tout $P$ de $\mP$ est divisible par au moins un $Q$. En effet, si au contraire $P_m$ était le premier $P$ qui n'est divisible par aucun $Q$, alors $P_m$ ne pourrait pas être lui-même un $Q$ et serait donc divisible par au moins un $P$ précédent; puisque, par hypothèse, ce $P$ précédent est divisible par un $Q$, $P_m$ lui-même serait divisible par un $Q$, contradiction.

Nous affirmons maintenant que le nombre des $Q$ est fini. En effet, soit
\[
 Q_0=\xi_1^{h_1}\cdots\xi_n^{h_n}
\]
le premier $Q$. Comme un $Q_\rho=\xi_1^{h_{\rho1}}\cdots\xi_n^{h_{\rho n}}$ arbitraire n'est pas divisible par $Q_0$, pour chaque $Q$ au moins une des inégalités $h_{\rho\lambda}<h_\lambda$ doit être satisfaite; que $\lambda$ désigne le plus petit indice pour lequel c'est le cas. Nous réunissons tous les $Q$ pour lesquels $h_{\rho\lambda}$ a la même valeur fixe $c_\lambda$ en une classe $K(c_\lambda)$. Comme $c_\lambda<h_\lambda$, le nombre de ces classes est fini. Mais chaque classe ne contient qu'un nombre fini d'éléments. Car si l'on écrit un $Q$ de la classe $K(c_\lambda)$ sous la forme $\xi_\lambda^{c_\lambda}Q'$, les $Q'$ forment encore un système de produits de puissances dont aucun n'est divisible par un autre, et ils ne contiennent plus que $n-1$ variables. Pour $n=2$, les $Q'$ ne contiennent qu'une variable; leur nombre est alors $=1$, donc fini, de sorte que l'on peut supposer le nombre fini aussi pour $n-1$ variables. L'assertion est ainsi prouvée.

Deuxièmement, soit $\mM$ un module de polynômes dont les termes sont ordonnés lexicographiquement. Soit $\mP$ le système des produits de puissances qui apparaissent dans le terme le plus élevé, chacun étant écrit une seule fois. Alors au système $\mathfrak{Q}=(Q_1,\ldots,Q_\varrho)$ correspond un système d'un nombre fini de polynômes $F_1,\ldots,F_\varrho$, en assignant à chaque $Q_i$ un polynôme quelconque de $\mM$ dont le terme le plus élevé est égal à $b_iQ_i$, de sorte que $F_i=b_iQ_i+$ avec des termes inférieurs.

Soit maintenant $F=bP+\cdots$ un polynôme quelconque de $\mM$, où $P=\xi_1^{\beta_1}\cdots\xi_n^{\beta_n}Q_i$. Formons le produit
\[
 \xi_1^{\beta_1}\cdots\xi_n^{\beta_n}F_i
 =\xi_1^{\beta_1}\cdots\xi_n^{\beta_n}(b_iQ_i+\cdots)
 =c\,\xi_1^{\beta_1}\cdots\xi_n^{\beta_n}Q_i+\cdots
 =cP+\cdots.
\]
Ici $c\ne0$ d'après (6), et tous les termes suivants sont des termes inférieurs dans l'ordre lexicographique, puisque la multiplication, en ce qui concerne les exposants, est commutative à des termes inférieurs près.\footnote{On voit par là qu'à la place de la loi (6) on pourrait aussi utiliser la loi un peu moins restrictive
\[
 \xi_i a=a_i\xi_i+a_{i,i+1}\xi_{i+1}+\cdots+a_{in}\xi_n+b_i\qquad(a_i\ne0)
\]
(pour une numérotation fixe quelconque des $\xi$).} La différence
\[
 F-\frac bc\,\xi_1^{\beta_1}\cdots\xi_n^{\beta_n}F_i
\]
ne contient alors que des termes inférieurs et appartient de nouveau à $\mM$. Le procédé s'arrête donc après un nombre fini d'étapes, de sorte que $F_1,\ldots,F_\varrho$ sont reconnus comme base du module $\mM$.

\textbf{Théorème III.} \emph{Tout système $\mS$ d'une infinité de polynômes possède une base finie comme module, $G_1,\ldots,G_k$, telle que tout polynôme $G$ de $\mS$ admette une représentation $G=A_1G_1+\cdots+A_kG_k$.}

Il en résulte immédiatement le

\textbf{Corollaire.} \emph{Tout système $\mS$ d'une infinité de classes résiduelles modulo $\mM$ possède une base $\mx_1,\ldots,\mx_k$, telle que tout $\mx$ de $\mS$ soit représentable sous la forme $\mx=A_1\mx_1+\cdots+A_k\mx_k$.}

Il suffit en effet d'associer à chaque classe résiduelle un représentant et de considérer le système formé par ces représentants et par tous les polynômes de $\mM$.

Nous passons maintenant à la preuve du théorème suivant.

\textbf{Théorème IV.} \emph{Tout groupe quotient peut être représenté comme somme d'un nombre fini de constituants irréductibles; autrement dit, d'après le théorème II, tout module peut être représenté comme plus petit commun multiple d'un nombre fini de modules irréductibles qui forment un système totalement premier entre eux.}

Supposons, par contradiction, que $\mG$ soit une somme d'une infinité de constituants $\mU_1,\mU_2,\ldots$ (cf. \S~5), et soient $\ma_1,\ma_2,\ldots$ les unités correspondantes, toutes non nulles par définition. Nous formons le système $\mS$ de toutes ces unités; d'après le corollaire au théorème III, il possède une base finie $\ma_{i_1},\ldots,\ma_{i_\nu}$ avec $i_1<\cdots<i_\nu$. Soit $\mu>i_\nu$. Alors
\[
 \ma_\mu=\mx_1\ma_{i_1}+\cdots+\mx_\nu\ma_{i_\nu}.
\]
C'est une relation linéaire entre unités, exclue par la définition. Ainsi, dans toute décomposition additive d'un groupe quotient, un constituant irréductible apparaît après un nombre fini d'étapes; grâce à la commutativité de la décomposition, nous pouvons le placer au début. Nous pouvons donc, sans restreindre la généralité, supposer que les $\mU_i$ eux-mêmes sont irréductibles. Mais d'après l'argument précédent, seuls un nombre fini de constituants $\mU_i$ peuvent apparaître. L'assertion est démontrée.

\subsection*{\S~7. Un cas de décomposition unique.}

Soit
\[
 \mG=\mU_1+\cdots+\mU_k=\mB_1+\cdots+\mB_l
\]
un groupe quotient dont les constituants $\mU_1,\ldots,\mU_k$, $\mB_1,\ldots,\mB_l$ soient irréductibles. Nous avons alors
\begin{equation}
 \mx\me=\mx\ma_1+\cdots+\mx\ma_k
 =\mx\mb_1+\cdots+\mx\mb_l,
\tag{22}
\end{equation}
où $\ma_\varkappa$ est la classe résiduelle de $\mG$ associée isomorphiquement à la classe unité de $\mU_\varkappa$, et $\mb_\lambda$ celle qui est associée à la classe unité de $\mB_\lambda$. En remplaçant maintenant $\mx$ par $\mx\mb_\lambda$, où $\lambda$ est un indice fixe de la suite $1,\ldots,l$, on obtient
\begin{equation}
 \mx\mb_\lambda=\mx\mb_\lambda\ma_1+\cdots+\mx\mb_\lambda\ma_k.
\tag{23}
\end{equation}
Cette représentation du groupe $\mB_\lambda$ est univoque, mais elle n'est pas une décomposition additive de $\mB_\lambda$, car les classes résiduelles individuelles du membre de droite n'ont pas nécessairement à appartenir à $\mB_\lambda$. Comme $\mb_\lambda$ ne s'annule pas, tous les produits $\mb_\lambda\ma_\varkappa$ ne peuvent pas être nuls. Supposons donc, sans restreindre la généralité,
\[
 \mb_\lambda\ma_1\ne0,
 \ldots,
 \mb_\lambda\ma_\rho\ne0,
 \qquad
 \mb_\lambda\ma_{\rho+1}=0,
 \ldots,
 \mb_\lambda\ma_k=0.
\]
Alors
\[
 \mx\mb_\lambda=\mx\mb_\lambda\ma_1+\cdots+\mx\mb_\lambda\ma_\rho.
\]

Pour un $\varkappa$ fixé dans la suite $1,\ldots,\rho$, considérons l'ensemble de ces classes résiduelles $\mx\mb_\lambda$ qui ne sont pas nulles et pour lesquelles aussi $\mx\mb_\lambda\ma_\varkappa\ne0$. Les deux systèmes $\mx\mb_\lambda$ et $\mx\mb_\lambda\ma_\varkappa$ sont isomorphes. D'abord, l'association est biunivoque: $\mx\mb_\lambda\ma_\varkappa=0$ entraîne, par hypothèse, $\mx\mb_\lambda=0$, et cette dernière relation entraîne, par l'unicité de la représentation (23), $\mx\mb_\lambda\ma_\varkappa=0$. En outre, aux sommes correspondent les sommes, et aux produits avec des polynômes arbitraires correspondent les produits correspondants. En maintenant $\varkappa$ fixe, on peut faire la même considération pour chaque $\lambda$ tel que $\mb_\lambda\ma_\varkappa\ne0$. Pour chaque $\ma_\varkappa$, il doit exister au moins un tel $\mb_\lambda$, puisque $\ma_\varkappa=(\mb_1+\cdots+\mb_l)\ma_\varkappa\ne0$.

Nous supposons d'abord que pour chaque $\varkappa$ il existe seulement un $\lambda$, et pour chaque $\lambda$ seulement un $\varkappa$, tels que $\mb_\lambda\ma_\varkappa\ne0$.

De cette correspondance il résulte $k=l$; en outre, la désignation peut être choisie de telle sorte que $\mb_\varkappa\ma_\varkappa\ne0$ et $\mb_\lambda\ma_\varkappa=0$ pour $\lambda\ne\varkappa$. Nous disons alors que les produits $\mb_\lambda\ma_\varkappa$ forment un schéma diagonal. Dans ce cas, d'après (23), l'unité est $\mb_\varkappa=\mb_\varkappa\ma_\varkappa$, et, à cause de
\begin{equation}
 \ma_\varkappa=\mb_1\ma_\varkappa+\cdots+\mb_l\ma_\varkappa=\mb_\varkappa\ma_\varkappa=\mb_\varkappa,
\tag{24}
\end{equation}
on a en général $\mU_\varkappa=\mB_\varkappa$. Il existe donc unicité de la décomposition, et les produits $\ma_\varkappa\mb_\lambda$ forment un schéma diagonal.

Si, par exemple, la loi commutative vaut spécialement pour les classes résiduelles $\ma_\varkappa\mb_\lambda$, l'hypothèse précédente est toujours remplie. Car dans (23), chaque constituant $\mx\mb_\lambda\ma_\varkappa=\mx\ma_\varkappa\mb_\lambda$ est alors contenu dans $\mB_\lambda$, de sorte que (23) donne une décomposition additive du groupe $\mB_\lambda$. En raison de l'irréductibilité de $\mB_\lambda$, tous les $\mb_\lambda\ma_\varkappa$ sauf un sont donc nuls. D'après (24), pour $\varkappa$ fixé il n'y a en outre, de façon correspondante, qu'un seul $\mb_\lambda\ma_\varkappa$ différent de zéro; sinon (24), par l'interchangeabilité de $\ma_\varkappa$ et de $\mb_\lambda$, fournirait une décomposition additive de $\mU_\varkappa$, qui devrait pourtant être irréductible.\footnote{Il en va de même dans le cas bilatère non commutatif traité par Schmeidler, a. a. O., où le produit de toute classe résiduelle de l'un des sous-groupes par toute classe résiduelle de l'autre sous-groupe disparaît. Alors, en effet, comme $\mb_\lambda\ma_\varkappa=\mb_\lambda\ma_\varkappa\mb_1+\cdots+\mb_\lambda\ma_\varkappa\mb_l$ et comme $\mb_\lambda(\ma_\varkappa\mb_\mu)=0$ pour $\lambda\ne\mu$, on a aussi $\mb_\lambda\ma_\varkappa=\mb_\lambda\ma_\varkappa\mb_\lambda$, donc cette classe est contenue dans $\mB_\lambda$; par l'irréductibilité de $\mB_\lambda$, il s'ensuit que pour un seul, et donc exactement un, $\varkappa$, le produit $\mb_\lambda\ma_\varkappa\ne0$. De la même manière, on montre qu'à chaque $\ma_\varkappa$ il existe un et un seul $\mb_\lambda$ tel que $\ma_\varkappa\mb_\lambda\ne0$, d'où $k=l$. Ainsi les $\mb_\lambda\ma_\varkappa$ forment un schéma diagonal, ce qui prouve l'unicité démontrée là.}

Les mêmes conclusions valent dans le cas où les conditions $\mb_\lambda\ma_\lambda\ne0$, $\mb_\lambda\ma_\varkappa=0$ pour $\lambda\ne\varkappa$, $\varkappa=1,\ldots,k$, et de même, pour chaque $\mu=1,\ldots,l$, $\mb_\mu\ma_\varkappa=0$, $\mb_\varkappa\ma_\varkappa\ne0$, $\varkappa=1,\ldots,\sigma$, sont satisfaites non pas pour tous les $\mb_\lambda$, mais seulement pour une partie d'entre eux ($\lambda=1,\ldots,\sigma$). On trouve $\mU_\varkappa=\mB_\varkappa$ pour $\varkappa=1,\ldots,\sigma$; et si l'on pose
\[
 \mU_{\sigma+1}+\cdots+\mU_k=\mU,
 \qquad
 \mB_{\sigma+1}+\cdots+\mB_l=\mB,
\]
alors, aussi pour les unités $\ma$ et $\mb$ de $\mU$ et $\mB$, les conditions $\mb\ma_\varkappa=0$, $\mb_\varkappa\ma=0$ ($\varkappa=1,\ldots,\sigma$), et $\mb\ma\ne0$ sont satisfaites, d'où $\mU=\mB$.
\subsection*{\S~8. Décomposition additive des groupes complètement réductibles en groupes premiers. Le théorème d'isomorphie.}

Nous considérons maintenant un autre cas, où il n'y a certes pas unicité, mais où il y a isomorphie des différentes décompositions. Il s'agit ici de la généralisation du cas considéré par Loewy.

\textbf{Définition.} \emph{Un module est appelé module premier s'il ne possède aucun diviseur autre que lui-même et le module unité contenant tous les polynômes. Le groupe quotient d'un module premier est appelé groupe premier.}

Tout groupe premier est donc \emph{a fortiori} irréductible. Il existe aussi des groupes premiers infinis (voir la définition du \S~11 et l'exemple \S~12, 2).

Supposons maintenant que le groupe quotient donné $\mG$ soit représentable de deux façons comme somme d'un nombre fini de groupes premiers $\mU_1,\ldots,\mU_k$ et $\mB_1,\ldots,\mB_l$. Un groupe quotient qui admet au moins une telle représentation comme somme de groupes premiers, ainsi que le module correspondant $\mM$, sera appelé complètement réductible, conformément à la terminologie de Loewy.

Nous considérons les produits $\ma_\varkappa\mb_\lambda$ ($\varkappa=1,\ldots,k$, $\lambda=1,\ldots,l$), et nous voulons montrer que ou bien $\ma_\varkappa\mb_\lambda=0$, ou bien $\mx\ma_\varkappa\mb_\lambda$ parcourt tout le groupe $\mB_\lambda$ lorsque $\mx$ parcourt toutes les classes résiduelles de $\mG$. En effet, le module appartenant à $\mB_\lambda$ est
\[
 (\mM,\mb_1+\cdots+\mb_{\lambda-1}+\mb_{\lambda+1}+\cdots+\mb_l)=\mM_{\mb_\lambda};
\]
il provient de $\mM$ en adjoignant tous les polynômes de la classe résiduelle $\mb_1+\cdots+\mb_{\lambda-1}+\mb_{\lambda+1}+\cdots+\mb_l$. Ce module possède le diviseur
\[
 (\mM,\mb_1+\cdots+\mb_{\lambda-1}+\mb_{\lambda+1}+\cdots+\mb_l,
 \ma_\varkappa\mb_\lambda).
\]
Mais puisque $\mM_{\mb_\lambda}$ est, par hypothèse, un module premier, ce diviseur est ou bien égal à $\mM_{\mb_\lambda}$, ou bien au module unité. Dans le premier cas $\ma_\varkappa\mb_\lambda$ appartient à $\mM_{\mb_\lambda}$; il existe donc une relation
\[
 \ma_\varkappa\mb_\lambda=\mx(\mb_1+\cdots+\mb_{\lambda-1}+\mb_{\lambda+1}+\cdots+\mb_l),
\]
d'où $\ma_\varkappa\mb_\lambda=0$. Dans l'autre cas il existe deux classes $\mx_0$ et $\mx_1$ telles que
\[
 \mx_0\ma_\varkappa\mb_\lambda
 +\mx_1(\mb_1+\cdots+\mb_{\lambda-1}+\mb_{\lambda+1}+\cdots+\mb_l)
 =\me=\mb_1+\cdots+\mb_{\lambda-1}+\mb_{\lambda+1}+\cdots+\mb_l+\mb_\lambda,
\]
et donc, par unicité, $\mx_0\ma_\varkappa\mb_\lambda=\mb_\lambda$. Ainsi $F\mx_0\ma_\varkappa\mb_\lambda$, et donc aussi $\mx\ma_\varkappa\mb_\lambda$, parcourt tout le groupe $\mB_\lambda$ lorsque $F$, respectivement $\mx$, parcourt tous les polynômes, respectivement toutes les classes résiduelles de $\mG$. L'assertion est ainsi démontrée.

Si $\ma_\varkappa\mb_\lambda\ne0$, alors $\mU_\varkappa$ est en outre isomorphe à $\mB_\lambda$. Car, d'après le \S~7, l'ensemble de toutes les classes $\mx\ma_\varkappa\ne0$ pour lesquelles $\mx\ma_\varkappa\mb_\lambda\ne0$ est isomorphe au système des classes $\mx\ma_\varkappa\mb_\lambda$; or celles-ci épuisent le sous-groupe $\mB_\lambda$. Il reste seulement à montrer que les classes indiquées $\mx\ma_\varkappa$ épuisent $\mU_\varkappa$. Pour cela, considérons toutes les classes résiduelles $\mx\ma_\varkappa$ pour lesquelles $\mx\ma_\varkappa\mb_\lambda=0$. Ce système possède la propriété que la somme de deux de ses classes et le produit avec un polynôme arbitraire appartiennent de nouveau au système. Si l'on adjoint toutes ces classes résiduelles au module $\mM_{\ma_\varkappa}$, on obtient de nouveau un module qui est diviseur de $\mM_{\ma_\varkappa}$; il est donc soit $\mM_{\ma_\varkappa}$, soit le module unité. Dans ce dernier cas toutefois $\ma_\varkappa$ lui-même y serait contenu, de sorte que $\ma_\varkappa\mb_\lambda=0$, contrairement à l'hypothèse. Il n'existe donc aucune classe $\mx\ma_\varkappa\ne0$ pour laquelle $\mx\ma_\varkappa\mb_\lambda=0$, et l'isomorphie de $\mU_\varkappa$ avec $\mB_\lambda$ est prouvée.

Supposons maintenant, pour un $\varkappa$ fixé, que $\ma_\varkappa\mb_{\lambda_1}\ne0,\ldots,\ma_\varkappa\mb_{\lambda_i}\ne0$ avec $i>1$. Alors $\mU_\varkappa$ est isomorphe aux groupes $\mB_{\lambda_1},\ldots,\mB_{\lambda_i}$; ces groupes sont donc isomorphes entre eux.

Nous montrons maintenant qu'au moins deux des $\mU$ doivent alors aussi être isomorphes. En effet, si dans les produits $\mb_\lambda\ma_\varkappa$ il ne correspondait à chaque $\mb_\lambda$ qu'un seul $\ma_\varkappa$ avec $\mb_\lambda\ma_\varkappa\ne0$, de sorte que $\mb_\lambda=\mb_\lambda\ma_\varkappa$, alors $\mB_\lambda=\mU_\varkappa$, puisque $\mU_\varkappa$, en tant que groupe premier, ne peut contenir de sous-groupe propre. Ainsi $k=l$ et, après une numérotation correcte, $\ma_\varkappa=\mb_\varkappa$; mais alors le schéma $\ma_\varkappa\mb_\lambda$ serait un schéma diagonal, ce qui n'est pas le cas. Le théorème suivant est donc démontré.

\textbf{Théorème V.} \emph{Si un groupe complètement réductible $\mG$ est représenté de deux façons comme somme de groupes premiers,
\[
 \mG=\mU_1+\cdots+\mU_k=\mB_1+\cdots+\mB_l,
\]
alors tout groupe $\mU$ est isomorphe à au moins un groupe $\mB$, et réciproquement. Si chaque $\mU$ est isomorphe à exactement un $\mB$, les décompositions sont identiques. Dans l'autre cas, au moins deux groupes $\mU$ sont isomorphes entre eux, et de même au moins deux groupes $\mB$ sont isomorphes entre eux.}

\subsection*{\S~9. Rapport entre les notions \og{} isomorphe \fg{} et \og{} de même espèce \fg{}.}

Dans ce paragraphe nous montrons que la notion \og{} de même espèce \fg{}, connue pour les expressions différentielles d'une variable, conduit, lorsqu'on la généralise conformément au sens, à l'isomorphie des groupes quotients correspondants.

\textbf{Définition.}\footnote{Pour les expressions différentielles d'une variable, voir la définition formelle de Blumberg, loc. cit., p. 25.} \emph{Deux modules $\mM$ et $\mN$ sont dits de même espèce s'il existe deux polynômes $P$ et $Q$, avec $P$ premier à $\mM$ et $Q$ premier à $\mN$, tels que, pour tout $M\equiv0(\mM)$ et tout $N\equiv0(\mN)$,}
\begin{align}
 MQ&\equiv0(\mN),\tag{25}\\
 NP&\equiv0(\mM)\tag{26}
\end{align}
\emph{soient satisfaites. La primalité relative signifie ici l'existence de deux polynômes $X$ et $Y$ tels que}
\begin{align}
 YQ&\equiv1(\mN),\tag{27}\\
 XP&\equiv1(\mM)\tag{28}
\end{align}
\emph{soient satisfaites.}

Nous avons alors le théorème suivant.

\textbf{Théorème VI.} \emph{Deux modules sont de même espèce si et seulement si leurs groupes quotients sont isomorphes.}

1. Soit $\mA$ le groupe quotient de $\mM$, $\mB$ le groupe quotient de $\mN$, et soit $\mA$ isomorphe à $\mB$; soient $\ma$ et $\mb$ les classes unités de $\mA$ et de $\mB$. Soit en outre $Q\mb$ la classe de $\mB$ qui correspond à l'unité $\ma$ de $\mA$, et de même $P\ma$ la classe de $\mA$ qui correspond à l'unité $\mb$ de $\mB$. Nous écrivons cette correspondance sous la forme
\begin{align}
 \ma&\sim Q\mb,\tag{29}\\
 P\ma&\sim\mb.\tag{30}
\end{align}
Il en résulte
\[
 0=M\ma\sim MQ\mb,
 \qquad\text{c'est-à-dire}\qquad MQ\equiv0(\mN),
\]
ce qui démontre (25). De même on a
\[
 NP\ma\sim N\mb=0,
 \qquad\text{c'est-à-dire}\qquad NP\equiv0(\mM),
\]
donc la relation (26). En outre, d'après (29), $P\ma\sim PQ\mb$, et comme d'après (30) on avait supposé $P\ma\sim\mb$, l'unicité de la correspondance donne $PQ\mb=\mb$, donc $PQ\equiv1(\mN)$, et de même $QP\equiv1(\mM)$. Ainsi (27) et (28) sont démontrées avec $X=Q$, $Y=P$.

2. Réciproquement, supposons que $\mM$ et $\mN$ soient de même espèce, de sorte que les relations (25) à (28) soient satisfaites. Nous associons à une classe arbitraire $F\ma$ de $\mA$ la classe $FQ\mb$ de $\mB$. De cette manière, à chaque classe de $\mA$ ne correspond qu'une seule classe de $\mB$; car de $F\ma=0$, donc de $F=M$, il résulte par (25) que $FQ\mb=0$. En outre, tout le groupe $\mB$ est épuisé par cette correspondance, puisque la classe spéciale $Y\ma$ correspond à $YQ\mb$, qui est égale à $\mb$ par (27). Nous avons donc ainsi une application univoque $\Phi$ de $\mA$ sur $\mB$, dont il n'est pas encore montré qu'elle soit réciproquement univoque. De même, les équations (26) et (28) donnent une application univoque $\Psi$ de $\mB$ sur $\mA$, qui n'est pas non plus encore montrée réciproquement univoque.

Si l'une des deux applications est réciproquement univoque, l'isomorphie est démontrée, car les exigences relatives à la somme et au produit sont satisfaites. Mais si les deux applications $\Phi$ et $\Psi$ n'étaient pas réciproquement univoques, il y aurait alors dans $\mA$ des classes $G\ma$ auxquelles correspond dans $\mB$ la classe $GQ\mb=0$. Les classes $FQ\mb$ correspondant à toutes les autres classes résiduelles $F\ma$ épuiseraient alors déjà le groupe $\mB$, et par conséquent l'application $\Psi$ rendrait tout le groupe $\mA$ dans les classes $FQP\ma$. Comme nous le savons, il y a aussi dans $\mB$ des classes non nulles auxquelles correspond, sous $\Psi$, la classe nulle de $\mA$; désignons par $G_1\ma$ toutes les classes pour lesquelles $G_1\ma\ne0$ mais $G_1QP\ma=0$. Tous les autres polynômes $F$ ont la propriété que $FQP\ma$ épuise le groupe $\mA$. Les polynômes $F$ pour lesquels $FQP\ma=G_1\ma$ seront désignés par $G_2$; de tels polynômes existent dès qu'il existe des polynômes $G_1$. Pour chaque $G_2$, on a $G_2QP\ma\ne0$, mais $G_2(QP)^2\ma=0$. De même, pour tout entier $\nu$, il existe des classes $G_\nu\ma$ telles que $G_\nu(QP)^{\nu-1}\ma\ne0$ mais $G_\nu(QP)^\nu\ma=0$.

Nous considérons maintenant le système $\mS$ formé de tous les polynômes $G_1,G_2,\ldots$. Il possède une base finie comme module, $G_{\lambda_1},\ldots,G_{\lambda_\varrho}$. Choisissons $\nu$ plus grand que le plus grand des indices $\lambda_1,\ldots,\lambda_\varrho$, que nous notons $\lambda$. Alors
\[
 G_\nu=A_1G_{\lambda_1}+\cdots+A_\varrho G_{\lambda_\varrho}.
\]
En multipliant par $(QP)^\lambda$, nous obtenons
\[
 G_\nu(QP)^\lambda\equiv0(\mM),
\]
ce qui contredit notre hypothèse. Il s'ensuit l'isomorphie, et donc, d'après 1, aussi la validité des relations plus précises (27), (28) avec $X=Q$, $Y=P$.

En vertu du théorème VI, le théorème V peut aussi être formulé comme suit.

\textbf{Théorème VII.}\footnote{Voir Loewy, p. 100--101.} \emph{Si un module complètement réductible $\mM$ est représentable de deux façons comme plus petit commun multiple de modules premiers $\mP_1,\ldots,\mP_k$ et $\mathfrak{Q}_1,\ldots,\mathfrak{Q}_l$, formant chacun un système totalement premier entre eux, alors tout $\mP$ est de même espèce qu'au moins un $\mathfrak{Q}$, et réciproquement. Si chaque $\mP$ est de même espèce qu'exactement un $\mathfrak{Q}$, les décompositions sont identiques. Dans l'autre cas, au moins deux des modules $\mP$ sont de même espèce l'un que l'autre, et de même au moins deux des modules $\mathfrak{Q}$ sont de même espèce l'un que l'autre.}
\subsection*{\S~10. Existence d'une infinité de décompositions d'un groupe.}

Si la représentation comme somme d'un groupe complètement réductible n'est pas unique, alors, d'après le théorème V, il existe dans chaque représentation au moins deux sous-groupes isomorphes entre eux. Nous considérons maintenant le sous-groupe formé à partir d'un tel système de sous-groupes isomorphes,
\[
  \mathfrak H=\mA_1+\cdots+\mA_\alpha,
\]
où donc $\mA_1,\mA_2,\ldots,\mA_\alpha$ sont tous isomorphes, et nous affirmons que $\mathfrak H$ admet alors une infinité de décompositions distinctes de ce genre, pourvu seulement que le domaine de rationalité $\mathsf{P}$ contienne une infinité de \og{} constantes \fg{}, c'est-à-dire de grandeurs $a$ pour lesquelles $\xi_i a=a\xi_i$. Ce théorème vaut même lorsque l'on ne suppose rien de plus sur les groupes $\mA_i$, pas même l'irréductibilité. On a donc:

\textbf{Théorème VIII.} \emph{Si $\mG$ est représentable comme somme d'un nombre fini de sous-groupes isomorphes entre eux, alors cette représentation est possible d'une infinité de façons distinctes, pourvu seulement que le domaine de rationalité $\mathsf{P}$ contienne une infinité de constantes. Par constantes, on entend ici les grandeurs $a$ du domaine pour lesquelles $\xi_i a=a\xi_i$.}

\textbf{Démonstration.} Soit donc
\[
  \mG=\mA_1+\cdots+\mA_\alpha,
\]
et soit $\mA_i$ isomorphe à $\mA_\varkappa$ pour $i,\varkappa=1,\ldots,\alpha$ ($\alpha\ge2$). Pour indiquer une correspondance isomorphe déterminée, nous distinguons d'abord un sous-groupe, par exemple $\mA_1$. Que l'isomorphie entre $\mA_1$ et $\mA_i$ soit fixée par
\begin{equation}
 a_1\sim t_{1\varkappa}a_\varkappa,
 \qquad
 a_i\sim t_{i1}a_1,
 \qquad
 t_{11}a_1=a_1
 \quad(i,\varkappa=1,\ldots,\alpha),
\tag{31}
\end{equation}
c'est-à-dire que, dans le groupe $\mA_1$, chaque classe résiduelle se corresponde à elle-même. De $M\mathfrak{a}_1=0$ et $M\mathfrak{a}_i=0$ il résulte alors $M\mathfrak{t}_{1\varkappa}\mathfrak{a}_\varkappa=0$ et $M\mathfrak{t}_{i1}\mathfrak{a}_1=0$; les classes $\mathfrak{t}_{1\varkappa}\mathfrak{a}_\varkappa$ et $\mathfrak{t}_{i1}\mathfrak{a}_1$ admettent donc, comme les unités, la multiplication par des classes, et non seulement par des polynômes. Il résulte alors de la définition de l'isomorphie que, pour deux classes correspondantes $\my$ et $\bar{\my}$ qui admettent cette multiplication par classes, $\mathfrak x\mathfrak y$ et $\mathfrak x\bar{\mathfrak y}$ sont de nouveau des classes correspondantes. Par conséquent les relations $\mathfrak{a}_r\mathfrak{a}_s=0$ $(r\ne s)$ ($r=1,\ldots,\alpha$; $s=1,\ldots,\alpha$) donnent, pour $s=1$,
\begin{equation}
\left\{
\begin{aligned}
 \mathfrak{a}_r \mathfrak{t}_{1\varkappa}\mathfrak{a}_\varkappa&=0 &&(\varkappa=1,\ldots,\alpha,\ r\ne1),\\
 \text{et pour }s>1\qquad \mathfrak{a}_r \mathfrak{t}_{s1}\mathfrak{a}_1&=0 &&(r\ne s),\\
 \text{donc aussi}\qquad \mathfrak{t}_{1\varkappa}\mathfrak{a}_\varkappa&=\mathfrak{a}_1\mathfrak{t}_{1\varkappa}\mathfrak{a}_\varkappa,
 & \mathfrak{t}_{i1}\mathfrak{a}_1&=\mathfrak{a}_i\mathfrak{t}_{i1}\mathfrak{a}_1.
\end{aligned}\right.
\tag{32}
\end{equation}
En outre, de (31), en généralisant (27) et (28), l'unicité de la correspondance donne
\begin{equation}
 \mathfrak{t}_{1i}\mathfrak{t}_{i1}\mathfrak{a}_1=\mathfrak{a}_1,
 \qquad
 \mathfrak{t}_{\varkappa1}\mathfrak{t}_{1\varkappa}\mathfrak{a}_\varkappa=\mathfrak{a}_\varkappa.
\tag{33}
\end{equation}

Comme isomorphie entre $\mA_i$ et $\mA_\varkappa$, soit maintenant fixée la relation qui résulte de la composition de (31); une telle fixation est nécessaire, puisque les groupes individuels peuvent admettre des isomorphies en eux-mêmes. De (31), (32) et (33), on obtient donc
\begin{equation}
 \mathfrak{a}_i\sim \mathfrak{t}_{i1}\mathfrak{t}_{1\varkappa}\mathfrak{a}_\varkappa=\mathfrak{t}_{i\varkappa}\mathfrak{a}_\varkappa,
 \qquad
 \mathfrak{t}_{ii}\mathfrak{a}_i=\mathfrak{a}_i,
 \qquad
 \mathfrak{a}_r\mathfrak{t}_{s\varkappa}\mathfrak{a}_\varkappa=0\quad (r\ne s),
\tag{34}
\end{equation}
où l'on a posé $\mathfrak{t}_{i1}\mathfrak{t}_{1\varkappa}=\mathfrak{t}_{i\varkappa}$. De (33) et (34), il résulte encore
\[
 \mathfrak{t}_{i\varkappa}\mathfrak{t}_{\varkappa l}\mathfrak{a}_l=\mathfrak{t}_{i1}\mathfrak{t}_{1\varkappa}\mathfrak{t}_{\varkappa1}\mathfrak{t}_{1l}\mathfrak{a}_l
 =\mathfrak{t}_{i1}\mathfrak{t}_{1l}\mathfrak{a}_l=\mathfrak{t}_{il}\mathfrak{a}_l;
\]
par là, le privilège accordé au groupe $\mA_1$ est de nouveau supprimé. La relation isomorphe entre $\mA_i$ et $\mA_\varkappa$ reste la même, quel que soit le groupe $\mA_\varkappa$ utilisé à la place de $\mA_1$ pour la définition. Les $\mathfrak{t}_{i\varkappa}$ sont donc des classes qui, en résumé, satisfont aux relations
\begin{equation}
 M\mathfrak{t}_{i\varkappa}\mathfrak{a}_\varkappa=0;
 \quad \mathfrak{a}_r\mathfrak{t}_{i\varkappa}\mathfrak{a}_\varkappa=0\ (r\ne i);
 \quad \mathfrak{t}_{i\varkappa}\mathfrak{t}_{\varkappa l}\mathfrak{a}_l=\mathfrak{t}_{il}\mathfrak{a}_l;
 \quad \mathfrak{t}_{\varkappa\varkappa}\mathfrak{a}_\varkappa=\mathfrak{t}_{\varkappa i}\mathfrak{t}_{i\varkappa}\mathfrak{a}_\varkappa=\mathfrak{a}_\varkappa.
\tag{35}
\end{equation}
Ces relations disent exactement que les modules $\mM_i$ et $\mM_\varkappa$ appartenant à $\mA_i$ et $\mA_\varkappa$ sont de même espèce (voir la définition du \S~9); à la place des $P$ et $Q$ de là-bas, on a ici les classes $\mathfrak{t}_{i\varkappa}$ et $\mathfrak{t}_{\varkappa i}$, et les deux premières relations (35), prises ensemble, correspondent à la formule (25), respectivement (26). Réciproquement, ces relations impliquent l'isomorphie d'après le théorème VI.\footnote{Et cela sans application du théorème de finitude, puisque les formules représentent déjà respectivement une application $\varphi_{i\varkappa}$ et son inverse $\varphi_{\varkappa i}$. En effet, de $\mathfrak{x}\mathfrak{a}_i=\mathfrak{x}\mathfrak{a}_i \mathfrak{t}_{i\varkappa}\mathfrak{a}_\varkappa \mathfrak{t}_{\varkappa i}\mathfrak{a}_i$ il résulte que $\mathfrak{x}\mathfrak{a}_i\ne0$ implique aussi $\mathfrak{x}\mathfrak{a}_i \mathfrak{t}_{i\varkappa}\mathfrak{a}_\varkappa\ne0$.}

À partir des classes $\mathfrak{t}_{i\varkappa}\mathfrak{a}_\varkappa$, on peut maintenant composer des classes $\mathfrak{b}_\rho$ qui se révéleront de nouveau être les unités de groupes $\mB_\rho$, d'où résultera une représentation
\[
 \mG=\mB_1+\cdots+\mB_\alpha.
\]
Soient en effet $\lambda_{i\varkappa}$ des constantes dont le déterminant $|\lambda_{i\varkappa}|\ne0$, et soient $\lambda^{i\varkappa}$ les cofacteurs correspondants, divisés par $|\lambda_{i\varkappa}|$, de sorte que les relations suivantes soient satisfaites:
\begin{equation}
\left\{
\begin{aligned}
 \sum_{r=1}^{\alpha}\lambda_{ri}\lambda^{r\varkappa}&=\delta_{i\varkappa}=\begin{cases}0,&i\ne\varkappa,\\1,&i=\varkappa,
 \end{cases}\\
 \sum_{r=1}^{\alpha}\lambda_{ir}\lambda^{\varkappa r}&=\delta_{i\varkappa}.
\end{aligned}\right.
\tag{36}
\end{equation}
Nous définissons alors
\begin{equation}
 \mathfrak{b}_\rho=\sum_{i,\varkappa}\lambda_{\rho i}\lambda^{\rho\varkappa}\mathfrak{t}_{i\varkappa}\mathfrak{a}_\varkappa
 \qquad(\rho=1,\ldots,\alpha).
\tag{37}
\end{equation}
De (35), (36) et (37), on déduit $M\mathfrak{b}_\rho=0$ et
\begin{equation}
 \sum_{\rho=1}^{\alpha}\mathfrak{b}_\rho
 =\sum_{i,\varkappa,\rho}\lambda_{\rho i}\lambda^{\rho\varkappa}\mathfrak{t}_{i\varkappa}\mathfrak{a}_\varkappa
 =\sum_{i,\varkappa}\delta_{i\varkappa}\mathfrak{t}_{i\varkappa}\mathfrak{a}_\varkappa
 =\sum_i \mathfrak{t}_{ii}\mathfrak{a}_i=\sum_i \mathfrak{a}_i=\mathfrak{e},
\tag{38}
\end{equation}
et en outre
\[
\begin{aligned}
 \mathfrak{b}_\rho \mathfrak{b}_\sigma
 &=\sum_{\substack{i,\varkappa\\ \mu,\nu}}
  \lambda_{\rho i}\lambda^{\rho\varkappa}\lambda_{\sigma\mu}\lambda^{\sigma\nu}
  \mathfrak{t}_{i\varkappa}\mathfrak{a}_\varkappa \mathfrak{t}_{\mu\nu}\mathfrak{a}_\nu \\
 &=\sum_{i,\mu,\nu}\lambda_{\rho i}\lambda^{\rho\mu}\lambda_{\sigma\mu}\lambda^{\sigma\nu}\mathfrak{t}_{i\nu}\mathfrak{a}_\nu
 =\delta_{\sigma\rho}\sum_{i,\nu}\lambda_{\rho i}\lambda^{\sigma\nu}\mathfrak{t}_{i\nu}\mathfrak{a}_\nu
 =\begin{cases}0&(\rho\ne\sigma),\\ \mathfrak{b}_\rho&(\rho=\sigma),\end{cases}
\end{aligned}
\]
$(\rho,\sigma=1,\ldots,\alpha)$. La représentation $\mathfrak r e=\mathfrak r b_1+\cdots+\mathfrak r b_\alpha$ est donc une décomposition additive du groupe $\mG$. Les groupes $\mB_\rho$ sont distincts des groupes $\mA_\sigma$ dès que les $\lambda_{i\varkappa}$ ne forment pas seulement une matrice diagonale. Car, si nous formons $a_\sigma b_\rho$, il vient d'après (35):
\[
 a_\sigma b_\rho
 =\sum_{i,\varkappa}\lambda_{\rho i}\lambda^{\rho\varkappa}a_\sigma t_{i\varkappa}a_\varkappa
 =\lambda_{\rho\sigma}
 \sum_{\varkappa}\lambda^{\rho\varkappa}\mathfrak{t}_{\sigma\varkappa}\mathfrak{a}_\varkappa\ne0
 \quad\text{pour }\lambda_{\rho\sigma}\ne0,
\]
puisque, sinon, dans la dernière somme, chaque terme devrait disparaître séparément, ce qui n'est pas le cas. S'il existe donc, pour un $\sigma$ fixé, au moins deux $\rho$ tels que $\lambda_{\rho\sigma}\ne0$, alors $\mA_\sigma$ ne peut être identique à aucun $\mB_\rho$.

Mais les groupes $\mB_\rho$ sont aussi isomorphes entre eux et aux groupes $\mA_\sigma$. Pour cela, nous montrons d'abord l'existence de classes $\mathfrak{r}_{\rho\sigma}\mathfrak{a}_\sigma$ et $\mathfrak{r}^{\sigma\rho}\mathfrak{b}_\rho$ qui réalisent l'isomorphie de $\mA_\sigma$ et $\mB_\rho$ ($\rho,\sigma=1,\ldots,\alpha$). Nous vérifions en effet que, pour les classes
\[
 \mathfrak{r}_{\rho\sigma}\mathfrak{a}_\sigma=\sum_i\lambda_{\rho i}\mathfrak{t}_{i\sigma}\mathfrak{a}_\sigma,
 \qquad
 \mathfrak{r}^{\sigma\rho}=\sum_\varkappa\lambda^{\rho\varkappa}\mathfrak{t}_{\sigma\varkappa}\mathfrak{a}_\varkappa
 =\mathfrak{r}^{\sigma\rho}\mathfrak{b}_\rho\footnotemark
\]
\footnotetext{À cause de (37), on a
\[
 r^{\sigma\rho}b_\rho=\sum_{\varkappa,i,\mu}\lambda^{\rho\varkappa}t_{\sigma\varkappa}a_\varkappa\lambda_{\rho i}\lambda^{\rho\mu}t_{i\mu}a_\mu
 =\sum_{\varkappa,\mu}\lambda^{\rho\varkappa}\lambda_{\rho\varkappa}\lambda^{\rho\mu}t_{\sigma\mu}a_\mu=\sum_\mu\lambda^{\rho\mu}t_{\sigma\mu}a_\mu=r^{\sigma\rho}.
\]}
les relations \og{} de même espèce \fg{} entre les modules
\[
 (\mathfrak M,\mathfrak{b}_1+\cdots+\mathfrak{b}_{\rho-1}+\mathfrak{b}_{\rho+1}+\cdots+\mathfrak{b}_\alpha)
\]
et
\[
 (\mathfrak M,\mathfrak{a}_1+\cdots+\mathfrak{a}_{\sigma-1}+\mathfrak{a}_{\sigma+1}+\cdots+\alpha_\sigma)
\]
sont satisfaites:
\begin{equation}
\left\{
\begin{aligned}
 \mathfrak{b}_\varkappa \mathfrak{r}_{\rho\sigma}\mathfrak{a}_\sigma&=0 &&(\varkappa\ne\rho);&
 \mathfrak{r}^{\sigma\rho}\mathfrak{r}_{\rho\sigma}\mathfrak{a}_\sigma&=\mathfrak{a}_\sigma\footnotemark,\\
 \mathfrak{a}_\varkappa \mathfrak{r}^{\sigma\rho}\mathfrak{b}_\rho&=0 &&(\varkappa\ne\sigma);&
 \mathfrak{r}_{\rho\sigma}\mathfrak{r}^{\sigma\rho}\mathfrak{b}_\rho&=\mathfrak{b}_\rho,
\end{aligned}\right.
\tag{39}
\end{equation}
\footnotetext{Le nombre des formules est doublé par rapport à (35), parce que l'inverse ne s'obtient pas ici, comme là-bas, par permutation des indices.}
d'où, comme dans (35), on obtient l'isomorphie au moyen de la correspondance
\begin{equation}
 \mathfrak{b}_\rho\sim \mathfrak{r}_{\rho\sigma}\mathfrak{a}_\sigma;\qquad,
 \qquad
 \mathfrak{a}_\sigma\sim \mathfrak{r}^{\sigma\rho}\mathfrak{b}_\rho.
\tag{40}
\end{equation}
En effet, il vient:
\[
\begin{aligned}
 \mathfrak{b}_\varkappa \mathfrak{r}_{\rho\sigma}\mathfrak{a}_\sigma
 &=\sum_{\mu,\nu,i}\lambda_{\varkappa\mu}\lambda^{\varkappa\nu}\mathfrak{t}_{\mu\nu}\mathfrak{a}_\nu\cdot\lambda_{\rho i}\mathfrak{t}_{i\sigma}\mathfrak{a}_\sigma\\
 &=\sum_{\mu,i}\lambda_{\varkappa\mu}\lambda^{\varkappa i}\lambda_{\rho i}\mathfrak{t}_{\mu\sigma}\mathfrak{a}_\sigma
 =\delta_{\rho\varkappa}\sum_\mu\lambda_{\varkappa\mu}\mathfrak{t}_{\mu\sigma}\mathfrak{a}_\sigma
 =\delta_{\rho\varkappa}\mathfrak{r}_{\varkappa\sigma}\mathfrak{a}_\sigma=0\quad(\varkappa\ne\rho),
\end{aligned}
\]
puis
\[
 \mathfrak{r}^{\sigma\rho}\mathfrak{r}_{\rho\sigma}\mathfrak{a}_\sigma
 =\sum_{\varkappa,i}\lambda^{\rho\varkappa}\mathfrak{t}_{\sigma\varkappa}\mathfrak{a}_\varkappa\cdot\lambda_{\rho i}\mathfrak{t}_{i\sigma}\mathfrak{a}_\sigma
 =\sum_i\lambda^{\rho i}\lambda_{\rho i}\mathfrak{t}_{\sigma\sigma}\mathfrak{a}_\sigma=\mathfrak{a}_\sigma.
\]
De même, inversement, on trouve
\[
 \mathfrak{a}_\varkappa \mathfrak{r}^{\sigma\rho}\mathfrak{b}_\rho
 =\mathfrak{a}_\varkappa\sum_\mu\lambda^{\rho\mu}\mathfrak{t}_{\sigma\mu}\mathfrak{a}_\mu=0
 \quad(\varkappa\ne\sigma),
\]
et
\[
 \mathfrak{r}_{\rho\sigma}\mathfrak{r}^{\sigma\rho}\mathfrak{b}_\rho
 =\sum_{i,\mu}\lambda_{\rho i}\mathfrak{t}_{i\sigma}\mathfrak{a}_\sigma\lambda^{\rho\mu}\mathfrak{t}_{\sigma\mu}\mathfrak{a}_\mu
 =\sum_{i,\mu}\lambda_{\rho i}\lambda^{\rho\mu}\mathfrak{t}_{i\mu}\mathfrak{a}_\mu=\mathfrak{b}_\rho.
\]
Des relations de même espèce ainsi démontrées résulte donc l'isomorphie de $\mA_\sigma$ et de $\mB_\rho$ (à nouveau sans le théorème de finitude). L'isomorphie entre $\mB_\rho$ et $\mB_\tau$ s'obtient de là par composition; nous écrivons encore les formules. Posons
\[
 \mathfrak{s}_{\rho\tau}\mathfrak{b}_\tau=\mathfrak{r}_{\rho\varkappa}\mathfrak{r}^{\varkappa\tau}\mathfrak{b}_\tau
 =\sum_{i,j}\lambda_{\rho i}\mathfrak{t}_{i\varkappa}\mathfrak{a}_\varkappa\cdot\lambda^{\tau j}\mathfrak{t}_{\varkappa j}\mathfrak{a}_j
 =\sum_{i,j}\lambda_{\rho i}\lambda^{\tau j}\mathfrak{t}_{ij}\mathfrak{a}_j;
\]
alors la correspondance est $\mathfrak{b}_\rho\sim\mathfrak{s}_{\rho\tau}\mathfrak{b}_\tau$. Ici encore les relations de même espèce se vérifient de nouveau formellement; si l'on met les $\mathfrak{s}_{\rho\tau}$ à la place des $\mathfrak{t}_{i\varkappa}$, elles satisfont à (35). En effet,
\[
 \mathfrak{b}_\sigma\mathfrak{s}_{\rho\tau}\mathfrak{b}_\tau=\mathfrak{b}_\sigma \mathfrak{r}_{\rho\varkappa}\mathfrak{r}^{\varkappa\tau}\mathfrak{b}_\tau=0\quad(\sigma\ne\rho),
\]
\[
 \mathfrak{s}_{\rho\sigma}\mathfrak{s}_{\sigma\tau}\mathfrak{b}_\tau=\mathfrak{r}_{\rho\varkappa}\mathfrak{r}^{\varkappa\sigma}\mathfrak{r}_{\sigma\varkappa}\mathfrak{r}^{\varkappa\tau}\mathfrak{b}_\tau
 =\mathfrak{r}_{\rho\varkappa}\mathfrak{a}_\varkappa \mathfrak{r}^{\varkappa\tau}\mathfrak{b}_\tau=\mathfrak{s}_{\rho\tau}\mathfrak{b}_\tau.
\]

Ainsi le théorème VIII est démontré, car il suffit de choisir les $\lambda_{i\varkappa}$ d'une infinité de manières de sorte qu'ils ne forment pas une matrice diagonale et qu'ils ne puissent pas non plus être transformés les uns dans les autres par des matrices diagonales.

Une conséquence du théorème VIII, en liaison avec le théorème V ou VII, est la suivante:

\textbf{Théorème IX.}\footnote{Voir Loewy, p. 106--107.} \emph{Soit le module $\mM$ complètement réductible et plus petit commun multiple de modules premiers $\mP_1,\ldots,\mP_k$ qui forment un système totalement premier entre eux, et supposons que le domaine de rationalité $P$ contienne une infinité de constantes. Alors la condition nécessaire et suffisante pour que $\mM$ soit représentable d'une infinité de façons différentes comme plus petit commun multiple de modules premiers est qu'au moins deux des modules $\mP_1,\ldots,\mP_k$ soient de même espèce.}

Nous voulons maintenant déduire des critères généraux pour qu'un module $\mM$ soit divisible par une infinité de modules premiers, et nous démontrons d'abord le suivant:

\textbf{Lemme.}\footnote{Voir Loewy, p. 96.} \emph{Si un module complètement réductible $\mM=[\mP_1,\ldots,\mP_k]$ est divisible par un module premier $\mP'$ différent de tous les $\mP_i$, alors au moins deux des $\mP_i$ sont de même espèce.}

Nous extrayons d'abord de $\mP_1,\ldots,\mP_k$ un système totalement premier entre eux en biffant successivement chacun de ces modules premiers qui est contenu dans le plus petit commun multiple de ceux qui ne sont pas encore biffés. On peut donc supposer maintenant que $\mP_1,\ldots,\mP_k$ eux-mêmes forment un système totalement premier entre eux. Soit $a'$ une classe de $\mM$ correspondant à la classe unité modulo $\mP'$. Alors
\[
 F\mathfrak{a}'=F\mathfrak{a}'\mathfrak{a}_1+\cdots+F\mathfrak{a}'\mathfrak{a}_k
\]
pour tout polynôme $F$. Parmi les produits $\mathfrak{a}'\mathfrak{a}_i$, au moins deux sont non nuls, disons $\mathfrak{a}'\mathfrak{a}_1$ et $\mathfrak{a}'\mathfrak{a}_2$; car si, par exemple, seul $\mathfrak{a}'\mathfrak{a}_1\ne0$, alors $F\mathfrak{a}'=F\mathfrak{a}'\mathfrak{a}_1$, et les classes résiduelles $F\mathfrak{a}'$ formeraient un sous-groupe de $\mA_1$, donc, puisque $\mA_1$ est un groupe premier, seraient identiques à $\mA_1$, de sorte que les modules $\mP'$ et $\mP_1$ coïncideraient aussi. Le raisonnement du \S~8 montre alors que $\mA'$ est isomorphe à $\mA_1$ et à $\mA_2$.

Soit maintenant $\mM$ un module arbitraire. Nous considérons le plus petit commun multiple $\mN$ de tous les modules premiers par lesquels $\mM$ est divisible.\footnote{S'ils sont en nombre infini, alors, d'après le théorème IV, aucune décomposition additive du groupe quotient ne peut correspondre à cette représentation. Voir d'ailleurs l'exemple \S~12, 4.} Deux cas sont alors possibles. Ou bien $\mN$ n'est pas représentable comme plus petit commun multiple d'un nombre fini de modules premiers; dans ce cas $\mN$, et donc aussi $\mM$, possède une infinité de modules premiers comme diviseurs. Ou bien $\mN$ est représentable comme plus petit commun multiple d'un nombre fini de modules premiers; alors $\mN$ s'appelle le plus grand module complètement réductible de $\mM$ (à la suite d'une construction correspondante de Loewy). Si, parmi ces modules premiers en nombre fini, deux sont de même espèce, alors, par le théorème IX, le module $\mN$, et par conséquent $\mM$, est divisible par une infinité de modules premiers. Réciproquement, si ce dernier cas se produit, il existe au moins un diviseur premier de $\mN$ distinct de tous les modules en nombre fini dont le plus petit commun multiple représente $\mN$. Par le lemme, $\mN$, et donc $\mM$, possède deux diviseurs premiers de même espèce. Nous avons ainsi démontré le théorème suivant:

\textbf{Théorème X.} \emph{Si le domaine de rationalité $\mathsf{P}$ contient une infinité de constantes, alors la condition nécessaire et suffisante pour qu'un module soit divisible par une infinité de modules premiers distincts est que, ou bien il n'existe pas de plus grand module complètement réductible, ou bien, s'il en existe un, celui-ci soit divisible par au moins deux modules premiers distincts de même espèce.}

\subsection*{\S~11. Relations avec les systèmes d'équations différentielles et leurs intégrales.}

En vertu du théorème de finitude, on peut associer à tout module d'expressions différentielles un système composé d'un nombre fini d'expressions différentielles comme base, de telle sorte que l'ensemble des intégrales simultanées de toutes les équations différentielles obtenues en annulant une expression différentielle du module coïncide avec l'ensemble des intégrales du système fini d'équations différentielles que l'on obtient en annulant les expressions différentielles de base. C'est ainsi que se donne le rattachement à la théorie des systèmes d'équations différentielles partielles linéaires; nous poursuivons encore un peu ce rattachement dans le cas où le groupe quotient du module est \og{} fini \fg{}.

\textbf{Définition.} \emph{Un groupe quotient s'appelle fini s'il existe un nombre $\mu$, l'ordre du groupe, tel qu'entre $\mu+1$ classes résiduelles quelconques il existe une relation linéaire à coefficients dans $\mathsf{P}$, tandis qu'il existe au moins un système de $\mu$ classes résiduelles entre lesquelles aucune relation n'existe, système que nous appellerons une base linéaire du groupe quotient. Dans l'autre cas, le groupe quotient s'appelle infini.}

La somme de deux groupes quotients finis a un ordre fini, à savoir la somme de leurs ordres. La somme de deux groupes infinis, ou d'un groupe fini et d'un groupe infini, est un groupe infini. Des exemples de groupes finis sont fournis par tous les groupes quotients de modules à une variable; mais, même pour plusieurs variables, l'ordre peut devenir fini, par exemple pour le module de base $\xi_1^2$, $\xi_2^2$, où toutes les classes résiduelles peuvent être composées linéairement à partir de celles de $1,\xi_1,\xi_2,\xi_1\xi_2$. Pour des exemples de groupes infinis, voir \S~12.

On a maintenant le théorème suivant.

\textbf{Théorème XI.} \emph{Si un module possède un groupe quotient fini, alors tout système d'équations différentielles qui naît de l'annulation d'une base du module admet exactement autant d'intégrales linéairement indépendantes que l'indique l'ordre de son groupe quotient. Réciproquement, pour un système donné d'un nombre fini d'équations différentielles partielles linéaires qui ne possèdent qu'un nombre fini d'intégrales linéairement indépendantes, ce nombre est égal à l'ordre du groupe quotient du module engendré par les expressions différentielles.}

Le domaine de rationalité $\mathsf{P}$ est ici constitué de fonctions analytiques de $n$ variables $x_1,\ldots,x_n$, dont les singularités éventuelles, à l'intérieur d'un certain domaine $n$-dimensionnel, forment tout au plus des variétés de dimensions inférieures, de sorte qu'il existe un nombre arbitraire de points ayant des voisinages réguliers $n$-dimensionnels.

Soit maintenant $\mM$ un module à groupe quotient fini; nous montrons alors qu'il existe une base linéaire formée de produits de puissances qui possède la propriété suivante: \emph{dans l'expression d'une classe résiduelle quelconque au moyen de cette base, seules des puissances d'un nombre fini de grandeurs différentes de $P$ peuvent apparaître au dénominateur}.

À cette fin, nous ordonnons lexicographiquement tous les produits de puissances $\xi_1^{a_1}\cdots\xi_n^{a_n}$ et prenons comme représentants $Q_1,\ldots,Q_m$ des classes résiduelles tous ceux qui sont linéairement indépendants modulo $\mM$ de ceux qui les précèdent. Par hypothèse, il n'y en a qu'un nombre fini. Tous les autres peuvent s'exprimer linéairement par ceux-ci.

Soit $\mathfrak S$ le système de tous les produits de puissances non admis dans la base linéaire. Alors, d'après la première partie de la démonstration du théorème III, $\mathfrak S$ possède un sous-système fini $\mathfrak T$ de tels produits de puissances, tel que tout produit de puissances de $\mathfrak S$ soit divisible par au moins un produit de $\mathfrak T$. Soient $P_1,\ldots,P_\rho$ les produits de puissances de $\mathfrak T$, et soient $M_1\equiv0,\ldots,M_\rho\equiv0\;(\mM)$ les relations au moyen desquelles $P_1,\ldots,P_\rho$ sont chaque fois exprimés, modulo $\mM$, par une combinaison linéaire de produits de puissances inférieurs; par exemple
\[
 M_i=b_iP_i+\text{termes inférieurs}.
\]
Si maintenant $P$ est un produit de puissances quelconque de $\mathfrak S$, donc de la forme
\[
 \xi_1^{h_1}\cdots\xi_n^{h_n}P_i,
\]
il vient
\[
\begin{aligned}
 \xi_1^{h_1}\cdots\xi_n^{h_n}M_i
 &=b_i\xi_1^{h_1}\cdots\xi_n^{h_n}P_i+\text{termes inférieurs}\\
 &=b_iP+\text{termes inférieurs}.
\end{aligned}
\]
Admettons maintenant comme démontré que, pour tous les produits de puissances inférieurs à $P$, seules des puissances de $b_1,\ldots,b_\rho$ apparaissent au dénominateur, tandis que les numérateurs sont formés à partir des coefficients des $M_i$ par addition, multiplication et différentiation; alors cela vaut aussi pour $P$ lui-même. Or cette hypothèse est permise, puisqu'elle est satisfaite pour $P_1$, le premier élément de $\mathfrak S$ et de $\mathfrak T$.

D'après le théorème III, $M_1,\ldots,M_\rho$ est en outre une base du module $\mM$. Considérons maintenant un système de valeurs $\beta_1,\ldots,\beta_n$ de $x_1,\ldots,x_n$ pour lequel $b_1\ne0,\ldots,b_\rho\ne0$ et les autres coefficients des $M_i$ sont réguliers, c'est-à-dire un point régulier du système d'équations différentielles $M_1=0,\ldots,M_\rho=0$. Aux produits de puissances $Q_i=\xi_1^{h_{i1}}\cdots\xi_n^{h_{in}}$ $(i=1,\ldots,m)$ correspondent alors les dérivées partielles
\[
 \frac{\partial^{h_{i1}+\cdots+h_{in}}}{\partial x_1^{h_{i1}}\cdots\partial x_n^{h_{in}}}y,
\]
pour lesquelles nous prescrivons des valeurs constantes arbitraires $(\gamma_1,\ldots,\gamma_m)$ au point $(\beta_1,\ldots,\beta_n)$. Alors les équations différentielles $M_i=0$ déterminent de façon unique, comme valeurs finies en ce point, les valeurs de toutes les dérivées supérieures de $y$, puisque seuls les $b_i$ apparaissent au dénominateur; et l'on sait que cela implique l'existence de $m$ intégrales analytiques linéairement indépendantes.

D'autre part, sous nos hypothèses, il ne peut exister plus de $m$ intégrales linéairement indépendantes; sinon, en un point régulier, on pourrait prescrire arbitrairement les valeurs de $m+1$ dérivées convenablement choisies, contrairement au fait qu'entre $(m+1)$ produits de puissances quelconques doivent exister des dépendances linéaires modulo $\mM$.

Réciproquement, si un module possède un groupe quotient infini, on montre par la même voie qu'il existe une infinité d'intégrales linéairement indépendantes; dans le cas où il n'y a que $m$ intégrales linéairement indépendantes, le groupe quotient est donc fini, et son ordre est, d'après ce qui précède, égal à $m$.

Pour terminer ce paragraphe, nous montrons que, pour deux modules d'expressions différentielles, le concept d'isomorphie, ou ce qui revient au même celui de même espèce, a pour les intégrales des systèmes d'équations différentielles associés la même signification que le concept d'espèce utilisé par Poincaré pour une variable (cf. Blumberg, loc. cit., appendice I).

\textbf{Théorème XII.} \emph{Si deux modules $\mM$ et $\mN$ d'expressions différentielles sont de même espèce, alors il existe deux expressions différentielles $P$ et $Q$ telles que $P(y)$ et $Q(z)$ parcourent respectivement toutes les intégrales de $\mN$ et de $\mM$, lorsque $y$ parcourt toutes les intégrales de $\mM$ et que $z$ parcourt toutes les intégrales de $\mN$.}

Par hypothèse, en effet (d'après \S~9), les relations suivantes existent pour deux polynômes $P$ et $Q$ convenablement choisis:
\[
 MQ\equiv0(\mathfrak N)\qquad PQ\equiv1(\mathfrak N),
\]
\[
 NP\equiv0(\mathfrak M)\qquad QP\equiv1(\mathfrak M).
\]
Ainsi, puisque $NP(y)=0$, la fonction $P(y)$ est une intégrale de $\mN$. De même, $Q(z)$ est une intégrale de $\mM$. Puisque $PQ(z)=z$ et $QP(y)=y$, $P(y)$ parcourt toutes les intégrales de $\mN$, et $Q(z)$ toutes les intégrales de $\mM$. De cette manière, à tout $y$ non nul correspond aussi un $z$ non nul, et réciproquement.

\subsection*{\S~12. Exemples.}

1. Comme premier exemple, considérons une équation différentielle linéaire ordinaire du second ordre dont les coefficients soient représentés comme fonctions symétriques des intégrales. Le domaine de rationalité $\mathsf{P}$ consiste ici en toutes les fonctions rationnelles des intégrales et de leurs dérivées; nous nous restreignons à un voisinage suffisamment petit d'un point régulier de l'équation différentielle. Nous définissons le module $\mM$ comme l'ensemble de toutes les expressions différentielles qui ont $y_1$ et $y_2$ pour intégrales, et pour lesquelles nous écrivons de nouveau des polynômes en une indéterminée $\xi$. La base du module $\mM$ est alors à un seul élément, et elle est donnée par le polynôme
\[
 M=
 \begin{vmatrix} y_1&y_2\\ y_1'&y_2'\end{vmatrix}\xi^2
 -\begin{vmatrix} y_1&y_2\\ y_1''&y_2''\end{vmatrix}\xi
 +\begin{vmatrix} y_1'&y_2'\\ y_1''&y_2''\end{vmatrix},
\]
qui est déterminé de façon unique à un facteur de $\mathsf{P}$ près. On a alors
\begin{equation}
 \mM=[\mM_1,\mM_2],
\tag{41}
\end{equation}
où $\mM_i$ est constitué de toutes les expressions différentielles possédant l'intégrale $y_i$, et dont la base est déterminée à un facteur de $\mathsf{P}$ près par
\[
 M_i=y_i\xi-y_i'\qquad(i=1,2).
\]
En effet, $\mM$ est divisible par $\mM_1$ et $\mM_2$, puisqu'il s'annule pour les intégrales $y_1$ et $y_2$, donc il est divisible par $[\mM_1,\mM_2]$; et, comme l'ordre du polynôme générateur de $[\mM_1,\mM_2]$ doit être au moins $2$, on a $\mM=[\mM_1,\mM_2]$. De plus, $\mM_1$ et $\mM_2$ sont premiers entre eux, car
\begin{equation}
 \frac{y_2}{D}(y_1\xi-y_1')-\frac{y_1}{D}(y_2\xi-y_2')=1,
 \qquad
 D=D(y)=\begin{vmatrix}y_1&y_2\\ y_1'&y_2'\end{vmatrix}.
\tag{42}
\end{equation}
Ils sont en outre des modules premiers, puisque leur groupe quotient a l'ordre 1.

Mais, outre $y_1$ et $y_2$, $\mM$ possède aussi toutes les intégrales
\[
 z_1=\alpha_{11}y_1+\alpha_{12}y_2,
 \qquad
 z_2=\alpha_{21}y_1+\alpha_{22}y_2,
 \qquad
 A=|\alpha_{ik}|\ne0,
\]
et il est donc identique au plus petit commun multiple des deux modules premiers entre eux $\mN_1$ et $\mN_2$, qui possèdent respectivement les intégrales $z_1$ et $z_2$, et dont les polynômes générateurs sont donnés par
\[
 N_i=z_i\xi-z_i'=(\alpha_{i1}y_1+\alpha_{i2}y_2)\xi-(\alpha_{i1}y_1'+\alpha_{i2}y_2').
\]
Ainsi $\mM$ admet une infinité de représentations distinctes comme plus petit commun multiple de modules premiers entre eux, et, d'après le théorème VII, les deux modules $\mN_1$ et $\mN_2$ sont donc de même espèce entre eux et avec tous les $\mM_1$ et $\mM_2$. En fait, les groupes quotients de tous ces modules ont l'ordre 1, donc ils sont isomorphes.

Nous voulons maintenant établir la décomposition additive du groupe quotient $\mG=\mA_1+\mA_2$ correspondant à la représentation (41). À cette fin, nous partons, comme au \S~4, de la relation (42), qui exprime que $\mM_1$ et $\mM_2$ sont premiers entre eux. Pour les unités $\mathfrak{a}_1$ et $\mathfrak{a}_2$, nous obtenons les classes résiduelles engendrées par les polynômes
\[
 -\frac{y_1}{D}(y_2\xi-y_2')
 \quad\text{et}\quad
 \frac{y_2}{D}(y_1\xi-y_1').
\]
Ainsi
\[
 \mM_1=\left(\mM,\frac{y_2}{D}(y_1\xi-y_1')\right),
 \qquad
 \mM_2=\left(\mM,-\frac{y_1}{D}(y_2\xi-y_2')\right).
\]
Si nous formons maintenant de même les unités $\mathfrak{b}_1$ et $\mathfrak{b}_2$ qui correspondent à la décomposition $\mM=[\mN_1,\mN_2]$, en remplaçant $y_i$ par $z_i$, il vient
\[
 -\frac{z_1}{D(z)}(z_2\xi-z_2')
 =-\frac{\alpha_{11}y_1+\alpha_{12}y_2}{AD(y)}
 ((\alpha_{21}y_1+\alpha_{22}y_2)\xi-(\alpha_{21}y_1'+\alpha_{22}y_2')),
\]
\[
 \frac{z_2}{D(z)}(z_1\xi-z_1')
 =\frac{\alpha_{21}y_1+\alpha_{22}y_2}{AD(z)}
 ((\alpha_{11}y_1+\alpha_{12}y_2)\xi-(\alpha_{11}y_1'+\alpha_{12}y_2')),
\]
donc
\begin{equation}
\left\{
\begin{aligned}
 \mathfrak{b}_1&=\frac{\alpha_{11}y_1+\alpha_{12}y_2}{y_1}
 \left(\frac{\alpha_{22}}{A}\right)\mathfrak{a}_1
 +\frac{\alpha_{11}y_1+\alpha_{12}y_2}{y_2}
 \left(-\frac{\alpha_{21}}{A}\right)\mathfrak{a}_2
 =\frac{z_1}{y_1}\alpha^{11}\cdot \mathfrak{a}_1+\frac{z_1}{y_2}\alpha^{12}\cdot \mathfrak{a}_2,\\
 \mathfrak{b}_2&=\frac{\alpha_{21}y_1+\alpha_{22}y_2}{y_1}
 \left(-\frac{\alpha_{12}}{A}\right)\mathfrak{a}_1
 +\frac{\alpha_{21}y_1+\alpha_{22}y_2}{y_2}
 \left(\frac{\alpha_{11}}{A}\right)\mathfrak{a}_2
 =\frac{z_2}{y_1}\alpha^{21}\mathfrak{a}_1+\frac{z_2}{y_2}\alpha^{22}\mathfrak{a}_2.
\end{aligned}\right.
\tag{43}
\end{equation}
Ici $(\alpha^{ik})$ est l'inverse de la matrice $(\alpha_{ik})$.

Inversement, en échangeant $z_i$ avec $y_i$, on obtient
\begin{equation}
\left\{
\begin{aligned}
 \mathfrak{a}_1&=\frac{y_1}{z_1}\alpha_{11}\mathfrak{b}_1+\frac{y_1}{z_2}\alpha_{12}\mathfrak{b}_2,\\
 \mathfrak{a}_2&=\frac{y_2}{z_1}\alpha_{21}\mathfrak{b}_1+\frac{y_2}{z_2}\alpha_{22}\mathfrak{b}_2.
\end{aligned}\right.
\tag{44}
\end{equation}
Ces formules font voir l'isomorphie des groupes $\mA$ et $\mB$. D'après nos considérations générales, en effet, si $\mathfrak{a}_1\mathfrak{b}_\sigma\ne0$, la correspondance en question est $\mathfrak{a}_1\sim \mathfrak{a}_1\mathfrak{b}_\sigma$; donc, d'après (44),
\[
 \mathfrak{a}_1\sim\frac{y_1}{z_\sigma}\alpha_{1\sigma}\mathfrak{b}_\sigma,
 \qquad\text{si }\alpha_{1\sigma}\ne0,
\]
et, d'après (43), de même
\[
 \mathfrak{b}_\sigma\sim\frac{z_\sigma}{y_2}\alpha^{\sigma2}\mathfrak{a}_2,
 \qquad\text{si }\alpha^{\sigma2}\ne0,
\]
d'où
\[
 \frac{y_1}{z_\sigma}\alpha_{1\sigma}\mathfrak{b}_\sigma
 \sim
 \frac{y_1}{z_\sigma}\alpha_{1\sigma}\frac{z_\sigma}{y_2}\alpha^{\sigma2}\mathfrak{a}_2
 =\frac{y_1}{y_2}\alpha_{1\sigma}\alpha^{\sigma2}\mathfrak{a}_2.
\]
Ainsi
\[
 \mathfrak{a}_1\sim\alpha_{1\sigma}\alpha^{\sigma2}\frac{y_1}{y_2}\mathfrak{a}_2
 \quad\text{pour }\alpha_{1\sigma}\alpha^{\sigma2}\ne0,
 \qquad
 \frac{1}{\alpha_{1\sigma}\alpha^{\sigma2}}\frac{y_2}{y_1}\mathfrak{a}_1\sim \mathfrak{a}_2,
\]
ce qui représente l'inverse. Si maintenant $(\alpha_{i\varkappa})$ n'est pas une matrice diagonale, cette condition $\alpha_{1\sigma}\alpha^{\sigma2}\ne0$ est toujours satisfaite pour un $\sigma$. Le cas de la matrice diagonale doit être exclu, car alors les modules sont identiques. Si, par exemple, seul $\alpha_{12}=0$, alors $\mA_1=\mB_1$, mais non $\mA_2=\mB_2$. De $\mA_1+\mA_2=\mB_1+\mB_2$ et $\mA_1=\mB_1$ on ne peut donc pas conclure que $\mA_2=\mB_2$ (cf. la remarque 12). Les trois modules $\mM_1$, $\mM_2$ et $\mN_2$ sont alors deux à deux premiers entre eux, mais ne forment pas un système totalement premier entre eux, puisque $[\mM_1,\mM_2]$ est divisible par $\mN_2$ (cf. la remarque 15).

En posant encore
\[
 \mathfrak{t}_{12}=\alpha_{1\sigma}\alpha^{\sigma2}\frac{y_1}{y_2}\mathfrak{a}_2,
 \qquad
 \mathfrak{t}_{21}=\frac{1}{\alpha_{1\sigma}\alpha^{\sigma2}}\frac{y_2}{y_1}\mathfrak{a}_1,
\]
les relations (35), comme on le confirme par le calcul, sont satisfaites.

2. Un exemple de groupe infini est fourni par tout module à plus d'une variable dont la base ne se compose que d'un seul polynôme. L'exemple suivant, un peu plus compliqué, doit montrer que des groupes infinis peuvent aussi apparaître pour des modules à plusieurs éléments, c'est-à-dire pour ceux dont la base du module comprend plus d'un polynôme.\footnote{Comparer l'exemple de Landau communiqué par Blumberg (loc. cit., p. 51--52). Ici et dans la suite nous désignons les variables indépendantes par $x$ et $y$.}
\[
 \mM=[(\xi+x\eta),(\xi+1)]=[(P),(Q)]
 \qquad\left(\xi=\frac{\partial}{\partial x},\ \eta=\frac{\partial}{\partial y}\right).
\]
La règle de multiplication est celle des expressions différentielles. Le groupe quotient est ici infini parce que les groupes quotients de $(\xi+x\eta)$ et de $(\xi+1)$ sont tous deux infinis, tandis que les modules $(\xi+x\eta)$ et $(\xi+1)$ sont premiers entre eux, car
\[
 1=(-x\xi-x^2\eta+1)Q+(x\xi+x-1)P.
\]
Il est maintenant très facile de montrer que $\mM$ ne possède aucun polynôme du second degré, en posant
\[
 (u_1\xi+u_2\eta+u_3)(\xi+x\eta)=(v_1\xi+v_2\eta+v_3)(\xi+1).
\]
On en tire $u_1=u_2=u_3=v_1=v_2=v_3=0$. Si l'on cherche ensuite des polynômes du troisième degré dans $\mM$, on en trouve deux linéairement indépendants, par exemple
\[
 (\xi^2+x\xi\eta+\xi+(2+x)\eta)Q=(\xi^2+2\xi+1)P
\]
et
\[
 (\xi^2+2x\xi\eta+x^2\eta^2+\eta)Q
 =(\xi^2+x\xi\eta+\xi+(x-1)\eta)P.
\]
Si $\mM$ était un module à un seul élément, les deux polynômes devraient être divisibles par le polynôme générateur et, puisque $\mM$ ne contient aucun polynôme du second degré, par un polynôme du troisième degré; ils devraient donc être proportionnels, ce qui n'est pas le cas. Le fait que $\mM$ ait plusieurs éléments est la raison interne de l'échec des théorèmes de produit de Landau, valables dans le cas d'une seule variable.

3. Un autre exemple montre que des groupes infinis peuvent aussi apparaître pour des modules premiers.\footnote{Notre définition de \og{} module premier \fg{} est différente de ce que l'on appelle module premier dans le cas commutatif, car il existe alors toujours des diviseurs de multiplicité inférieure, sauf lorsque la multiplicité est $0$, c'est-à-dire lorsque le groupe quotient est fini.} Soit $P$ le domaine de toutes les fonctions rationnelles de deux variables $x,y$. Le module $\mM$ de base
\[
 \xi-\frac{y}{x}=\xi-a
\]
possède un groupe infini, puisque celui-ci contient la classe résiduelle de chaque puissance de $y$; d'autre part, $\mM$ est un module premier. En effet, nous montrons que tout diviseur
\[
 \mM_1=(\xi-a,F(\xi,\eta))
\]
est égal au module unité. Tout polynôme $F(\xi,\eta)$ peut en effet être représenté modulo $(\xi-a)$ par un polynôme $G(\eta)$:
\[
 F(\xi,\eta)\equiv c_0\eta^\nu+c_1\eta^{\nu-1}+\cdots+c_\nu\quad(\mM).
\]
Sans restreindre la généralité, soit $c_0=1$. Or, d'après la loi de multiplication, $\xi\eta^\nu-\eta^\nu\xi=0$; donc, puisque
\[
 \xi\equiv a(\mM_1),
 \qquad
 \eta^\nu\equiv F_1(\mM_1),
\]
où $F_1=-(c_1\eta^{\nu-1}+\cdots+c_\nu)$, on a
\[
 \xi F_1-\eta^\nu a\equiv0(\mM_1).
\]
En développant
\[
 \eta^\nu a=a\eta^\nu+\nu\frac{\partial a}{\partial y}\eta^{\nu-1}+\cdots,
\]
et en remplaçant de nouveau ici $\eta^\nu$ par $F_1$, on obtient
\[
 \xi(c_1\eta^{\nu-1}+\cdots+c_\nu)-a(c_1\eta^{\nu-1}+\cdots+c_\nu)
 +\nu\frac{\partial a}{\partial y}\eta^{\nu-1}+\cdots\equiv0(\mM_1),
\]
ou
\[
 c_1\eta^{\nu-1}a+\frac{\partial c_1}{\partial x}\eta^{\nu-1}+\cdots
 -ac_1\eta^{\nu-1}+\cdots+
 \nu\frac{\partial a}{\partial y}\eta^{\nu-1}+\cdots
 \equiv0(\mM_1),
\]
ou enfin
\[
 \left(\frac{\partial c_1}{\partial x}+\nu\frac{\partial a}{\partial y}\right)\eta^{\nu-1}+\cdots
 \equiv0(\mM_1).
\]
À gauche se trouve maintenant un polynôme en $\eta$ de degré exact $\nu-1$, qui ne s'annule pas identiquement; car son coefficient dominant
\[
 \frac{\partial c_1}{\partial x}+\nu\frac{\partial a}{\partial y}
 =\frac{\partial c_1}{\partial x}+\nu\frac{1}{x}
 \quad\text{est }\ne0,
\]
puisque sinon $c_1=-\nu\log x+\varphi(y)$ ne serait pas rationnel. Le procédé peut donc être poursuivi et conduit à un polynôme non identiquement nul de degré zéro, ce qui montre que l'unité est aussi contenue dans $\mM_1$.\footnote{Ce procédé revient à démontrer que les conditions d'intégrabilité nécessaires pour un système d'équations différentielles ne sont pas satisfaites.}

4. Si, en revanche, nous adjoignons encore au domaine de rationalité la fonction $\log x$, alors $\mM$ possède les diviseurs
\[
 \left(\xi-\frac{y}{x},\ \eta-\log x+\varphi(y)\right),
\]
où $\varphi(y)$ parcourt toutes les fonctions rationnelles de $y$. Ces diviseurs sont tous distincts les uns des autres; pour chacun la condition d'intégrabilité est satisfaite, et l'intégrale correspondante est
\[
 y\log x-\int\varphi(y)\,dy+\alpha.
\]
Par conséquent, l'unité ne peut pas être contenue dans le module, car sinon la fonction nulle serait la seule intégrale possible. D'autre part, $\xi$ et $\eta$ sont, relativement au module, congrues à une grandeur du domaine; le groupe quotient est donc fini et d'ordre 1. De plus, deux quelconques de ces diviseurs en nombre infini sont premiers entre eux, et le module donné $\mM$ est égal au plus petit commun multiple $\mN$ de tous ces diviseurs. En effet, il est divisible par tous les diviseurs, donc aussi par $\mN$. Si maintenant $\mN$ était un diviseur propre de $\mM$, alors $\mN$ devrait encore contenir au moins un polynôme $F$, disons de degré $\rho$, qui ne dépende que de $\eta$. Faisons parcourir à $\varphi(y)$ les fonctions $y,\ldots,y^{\rho+1}$; il n'existe alors aucune relation linéaire, à coefficients indépendants de $y$, entre les intégrales correspondantes
\[
 y\log x-\frac{y^{i+1}}{i+1}+c_i\qquad(i=1,\ldots,\rho+1).
\]
Mais l'équation différentielle d'ordre $\rho$, $F=0$, ne peut pas posséder $\rho+1$ intégrales linéairement indépendantes.

Le module $\mM=\mN$ n'est cependant pas complètement réductible.\footnote{Nous avons ainsi un exemple du premier cas du théorème X.} Sinon, il serait le plus petit commun multiple d'un nombre fini de ces modules premiers, et son groupe quotient serait donc d'ordre fini, tandis que le groupe quotient de $\xi-y/x$ contient bien les classes résiduelles de toutes les puissances de $\eta$, et est donc infini.

Göttingen, le 1er août 1919.

\begin{center}
(Reçu le 4 août 1919.)
\end{center}
\clearpage
\setcounter{footnote}{0}
\begin{center}
{\Large\bfseries \NoetherNavAnchor{P18}{18. Sur un travail de K. Hentzelt, tombé à la guerre, sur la théorie de l’élimination}18. Sur un travail de K. Hentzelt, tombé à la guerre,\par sur la théorie de l'élimination}\par
\vspace{1.2em}
J. Ber. d. DMV 30 (1921), p. 101
\end{center}

\begin{center}
\textbf{11\textsuperscript{e} séance, vendredi 23 septembre 1921, à 11 heures du matin.}\\
(Présidence Loewy.)
\end{center}

\textbf{1. E. Noether: Sur un travail de K. Hentzelt, tombé à la guerre, sur la théorie de l'élimination.}

Kurt Hentzelt est porté disparu depuis octobre 1914 devant Dixmuiden et doit être compté parmi les morts. Les parties essentielles de sa dissertation d'Erlangen --- qu'il a laissée construite sans lacune, mais présentée de façon presque incompréhensible --- je les publierai prochainement, dans une libre adaptation, aux \emph{Mathematische Annalen}.

Il s'agit du problème général de la \textbf{théorie de l'élimination}: déterminer les zéros communs de tous les polynômes d'un idéal $\mathfrak M$ formé de polynômes. Si l'on soumet d'abord les variables à une transformation linéaire à coefficients indéterminés, le résultat principal peut s'énoncer ainsi:

À l'idéal $\mathfrak M$ on peut associer de manière univoque une forme résultante:
\[
 R_{\mathfrak M}=R^{(1)}(x_1,\ldots,x_n)\cdots R^{(i)}(x_i,\ldots,x_n)\cdots R^{(n)}(x_n)\equiv0\pmod{\mathfrak M},
\]
 de telle sorte que $R_{\mathfrak M}$ s'annule pour tous les zéros de $\mathfrak M$ et seulement pour ceux-ci. Si $\mathfrak N$ est un diviseur de $\mathfrak M$, et si $R_{\mathfrak N}$ et $R_{\mathfrak M}$ coïncident, alors $\mathfrak N$ et $\mathfrak M$ eux-mêmes coïncident également.

La seconde partie du théorème montre que la forme résultante fournit les zéros avec leur multiplicité caractéristique. Elle repose sur le fait que $\mathfrak M$ peut être considéré comme un module de formes linéaires dans les produits de puissances de $x_1,\ldots x_{i-1}$; les différents facteurs $R^{(i)}(x_i\ldots x_n)$ de $R_{\mathfrak M}$ deviennent alors les normes du \og{}module fondamental\fg{} correspondant par rapport à ce module, lorsque $x_{i+1},\ldots x_n$ sont adjoints au domaine des coefficients. Si l'on fait appel à la théorie abstraite des idéaux, on obtient pour cela l'interprétation suivante: si
\[
 \mathfrak M=[\mathfrak Q,\mathfrak Q_1,\ldots,\mathfrak Q_r]=[\mathfrak Q,\mathfrak L]
\]
est une décomposition de $\mathfrak M$ en idéaux primaires $\mathfrak Q_\lambda$, et si $\mathfrak L$ est le complément de $\mathfrak Q$, alors à l'idéal $\mathfrak Q$ correspond un facteur primaire $Q^{(i)}(x_i\ldots x_n)$ de $R^{(i)}(x_i\ldots x_n)$, qui représente, au sens indiqué ci-dessus, la norme de $\mathfrak L$ relativement à $\mathfrak Q$.
\clearpage
\setcounter{footnote}{0}
\section*{\NoetherNavAnchor{P19}{19. Théorie des idéaux dans les anneaux}19. Théorie des idéaux dans les anneaux}
\begin{center}
Par Emmy Noether, à Göttingen.\\[0.9em]
Math. Ann. 83 (1921), pp. 24--66
\end{center}

\begin{center}\textbf{Table des matières.}\end{center}
\noindent Introduction.\\
\S~1. Anneau, idéal, condition de finitude.\\
\S~2. Représentation d'un idéal comme plus petit commun multiple d'un nombre fini d'idéaux irréductibles.\\
\S~3. Égalité du nombre des composantes dans deux décompositions différentes en idéaux irréductibles.\\
\S~4. Idéaux primaires. Unicité des idéaux premiers associés dans deux décompositions différentes en idéaux irréductibles.\\
\S~5. Représentation d'un idéal comme plus petit commun multiple des plus grands idéaux primaires. Unicité des idéaux premiers associés.\\
\S~6. Représentation unique d'un idéal comme plus petit commun multiple d'idéaux relativement premiers irréductibles.\\
\S~7. Unicité des idéaux isolés.\\
\S~8. Représentation unique d'un idéal comme produit d'idéaux étrangers-irréductibles.\\
\S~9. Extension de l'étude aux modules. Égalité du nombre des composantes dans les décompositions en modules irréductibles.\\
\S~10. Cas particulier de l'anneau de polynômes.\\
\S~11. Exemples tirés de la théorie des nombres et de la théorie des expressions différentielles.\\
\S~12. Exemple tiré de la théorie des diviseurs élémentaires.

\subsection*{Introduction.}

Le contenu du présent article consiste dans le \emph{transfert des théorèmes de décomposition des entiers rationnels, respectivement des idéaux dans les corps de nombres algébriques, aux idéaux dans des anneaux intègres quelconques, et plus généralement dans des anneaux quelconques}. Pour comprendre ce transfert, indiquons d'abord les théorèmes de décomposition des entiers rationnels sous une forme quelque peu différente de la formulation usuelle.

Si, dans
\[
 a=p_1^{\varrho_1}p_2^{\varrho_2}\cdots p_\sigma^{\varrho_\sigma}=q_1q_2\cdots q_\sigma,
\]
on considère les facteurs puissances de nombres premiers $q_i$ comme les composantes de la décomposition, alors ces composantes possèdent les propriétés caractéristiques suivantes.

1. Elles sont \emph{premières entre elles deux à deux}; mais aucun $q$ ne peut être représenté comme produit de nombres premiers entre eux deux à deux; il y a donc, en ce sens, irréductibilité. De la primalité deux à deux il résulte encore que le produit $q_1\cdots q_\sigma$ est égal au plus petit commun multiple $[q_1\cdots q_\sigma]$.

2. Deux quelconques des composantes, $q_i$ et $q_k$, sont \emph{relativement premières}; c'est-à-dire que, si $b\cdot q_i$ est divisible par $q_k$, alors $b$ est divisible par $q_k$. En ce sens aussi il y a irréductibilité.

3. Chaque $q$ est \emph{primaire}; c'est-à-dire que, si un produit $b\cdot c$ est divisible par $q$, mais que $b$ ne l'est pas, alors une puissance\footnote{Si cette puissance est toujours la première, il s'agit, comme on sait, de nombres premiers.} de $c$ est divisible. La représentation est en outre une représentation par \emph{plus grandes composantes primaires}, puisque le produit de deux $q$ distincts n'est plus primaire. Les $q$ sont également irréductibles relativement à la décomposition en plus grandes composantes primaires.

4. Chaque $q$ est \emph{irréductible} au sens où il ne peut pas être représenté comme plus petit commun multiple de deux diviseurs propres.

Le lien entre ces nombres primaires $q$ et les nombres premiers $p$ consiste en ceci: à chaque $q$ correspond un et, au signe près, un seul $p$ qui est diviseur de $q$ et dont une puissance est divisible par $q$: le nombre premier associé. Si $p^\varrho$ est la plus petite puissance de cette espèce --- $\varrho$ étant l'exposant de $q$ ---, alors ici, en particulier, $p^\varrho$ est égal à $q$. Le théorème d'unicité peut maintenant s'énoncer ainsi:

\emph{Dans deux décompositions différentes d'un entier rationnel en les composantes irréductibles, plus grandes primaires $q$, le nombre des composantes, les nombres premiers associés (au signe près) et les exposants coïncident. Comme $p^\varrho=q$, il en résulte aussi que les $q$ eux-mêmes coïncident (au signe près).}

L'indétermination donnée par le signe est, comme on sait, levée lorsqu'au lieu des nombres on considère les idéaux qui en dérivent (tous les nombres divisibles par $a$); alors la formulation vaut exactement de même pour la décomposition unique des idéaux des corps de nombres algébriques (finis) en puissances d'idéaux premiers.

Dans la suite on prend donc pour base (\S{}~1) un anneau général, qui ne doit satisfaire qu'à la condition de finitude suivant laquelle tout idéal du domaine possède une base d'idéal finie. Sans une telle condition de finitude, il ne serait pas nécessaire qu'il existe des idéaux irréductibles et premiers, comme le montre le domaine de tous les entiers algébriques, dans lequel il n'y a pas de décomposition en idéaux premiers.

Il apparaît que --- conformément aux quatre propriétés caractéristiques des composantes $q$ --- il existe en général quatre décompositions séparées, qui procèdent l'une de l'autre chaque fois par subdivision. Il s'agit, pour la décomposition en idéaux étrangers-irréductibles, d'une représentation comme produit; pour les trois autres décompositions, d'une représentation réduite (\S{}~2) comme plus petit commun multiple. Le lien entre idéal primaire --- les idéaux irréductibles eux aussi sont primaires --- et idéal premier associé est également conservé: à chaque idéal primaire $\ideal{Q}$ est déterminé de façon unique un idéal premier associé $\ideal{P}$ qui est diviseur de $\ideal{Q}$ et dont une puissance devient divisible par $\ideal{Q}$. Si $\ideal{P}^{\varrho}$ est la plus petite puissance de cette espèce --- $\varrho$ étant l'exposant de $\ideal{Q}$ ---, alors ici $\ideal{P}^{\varrho}$ n'a pas besoin de coïncider avec $\ideal{Q}$. Le théorème d'unicité s'exprime alors ainsi.

Les décompositions 1 et 2 sont uniques; dans deux décompositions différentes 3 ou 4, le nombre des composantes et les idéaux premiers associés coïncident;\footnote{Vraisemblablement, au-delà de cela, les exposants coïncident aussi, et plus généralement encore les composantes correspondantes sont isomorphes.} les idéaux isolés qui apparaissent parmi les composantes (\S{}~7) sont déterminés de façon unique.

Pour la preuve des théorèmes de décomposition, la condition de finitude fournit d'abord le \og{} théorème de la chaîne finie \fg{}, énoncé pour la première fois par Dedekind pour les modules de nombres finis; on en déduit la représentation 4 de chaque idéal comme plus petit commun multiple d'un nombre fini d'idéaux irréductibles. En transformant la notion de réductibilité d'une composante, on obtient à partir de là le théorème fondamental d'unicité pour la décomposition 4 en idéaux irréductibles. En groupant à chaque fois un nombre fini de composantes, on passe aux autres décompositions, dont les théorèmes d'unicité résultent alors du théorème d'unicité 4.

Enfin il est montré (\S{}~9) que la représentation par un nombre fini de composantes irréductibles a encore lieu sous des hypothèses plus faibles; la commutativité de l'anneau n'est pas nécessaire, et il suffit de considérer, à la place d'un idéal, un module relativement au domaine. Dans ce cas plus général subsiste encore l'égalité du nombre des composantes dans deux décompositions différentes, tandis que les notions de premier et de primaire sont liées à la commutativité et à la notion d'idéal; en revanche, la notion d'idéaux étrangers reste valable dans les domaines non commutatifs.

L'anneau le plus simple pour lequel les quatre décompositions séparées se présentent effectivement est le domaine de tous les polynômes en $n$ variables à coefficients complexes arbitraires. Les diverses décompositions s'interprètent ici de façon irrationnelle par le comportement des configurations algébriques, et le théorème d'unicité pour les idéaux premiers associés correspond au théorème fondamental de la théorie de l'élimination sur la décomposabilité unique des configurations algébriques en irréductibles. D'autres exemples sont donnés par tous les domaines d'intégrité finis formés de polynômes (\S{}~10). Mais même le domaine simple de tous les nombres pairs, plus généralement de tous les nombres divisibles par un nombre fixe, fournit déjà un exemple de décompositions partiellement séparées (\S{}~11). Un exemple de théorie des idéaux dans des domaines non commutatifs est fourni par la théorie des diviseurs élémentaires (\S{}~12), où il existe une décomposition unique en idéaux, respectivement en classes, irréductibles. Ces classes irréductibles caractérisent complètement les constituants irréductibles des diviseurs élémentaires, et peuvent peut-être être regardées comme leur équivalent pour les domaines où la théorie ordinaire des diviseurs élémentaires échoue.

Voici ce qu'il faut remarquer au sujet de la littérature existante. La décomposition en plus grands idéaux primaires a été donnée pour l'anneau de polynômes à coefficients complexes arbitraires, respectivement entiers, par Lasker, et poursuivie par Macaulay sur certains points.\footnote{E. Lasker, \emph{Zur Theorie der Moduln und Ideale}. Math. Ann. 60 (1905), p. 20, théorèmes VII et XIII. --- F. S. Macaulay, \emph{On the Resolution of a given Modular System into Primary Systems including some Properties of Hilbert Numbers}. Math. Ann. 74 (1913), p. 66.} Tous deux s'appuient sur la théorie de l'élimination; ils utilisent donc le fait qu'un polynôme se représente de façon unique comme produit de polynômes irréductibles. En réalité les théorèmes de décomposition des idéaux sont indépendants de cette hypothèse, comme le laisse supposer la théorie des idéaux dans les corps de nombres algébriques et comme le montre le présent travail. Chez Lasker et Macaulay aussi, l'idéal primaire est défini à partir de notions de la théorie de l'élimination.

La décomposition en idéaux irréductibles et la décomposition en idéaux irréductibles relativement premiers ne semblent pas avoir été remarquées dans la littérature, même pour l'anneau de polynômes; chez Macaulay seulement on trouve une remarque sur l'unicité des idéaux primaires isolés.

La décomposition en idéaux étrangers-irréductibles a été donnée pour l'anneau de polynômes par Schmeidler,\footnote{W. Schmeidler, \emph{Über Moduln und Gruppen hyperkomplexer Größen}. Math. Zeitschr. 3 (1919), p. 29.} en utilisant la théorie de l'élimination pour la preuve de finitude. Mais le théorème d'unicité n'y est énoncé que pour des classes d'idéaux, non pour les idéaux eux-mêmes. Ce dernier théorème d'unicité se trouve dans un travail commun,\footnote{E. Noether--W. Schmeidler, \emph{Moduln in nichtkommutativen Bereichen, insbesondere aus Differential- und Differenzausdrücken}. Math. Zeitschr. 8 (1920), p. 1.} où il s'agit d'idéaux dans des anneaux de polynômes non commutatifs. Là on n'utilise que la base d'idéal finie; les théorèmes et les méthodes subsistent donc pour des anneaux généraux, et le présent travail permet de les renforcer en ce qui concerne l'égalité du nombre (\S{}~11).

Les présentes recherches constituent une forte généralisation et un développement ultérieur des formations conceptuelles qui sont à la base de ces deux travaux. L'essentiel des deux travaux est le passage de la représentation comme plus petit commun multiple à une décomposition additive du système des classes modulo l'idéal considéré. Ici, pour la simplicité de l'exposé, on reste de nouveau au plus petit commun multiple; à la décomposition additive correspond alors la transformation de la notion de réductibilité en une propriété du complément (\S{}~3). Toutefois, selon des considérations qui correspondent essentiellement à celles du travail commun, tous les théorèmes indiqués peuvent aussi être conçus comme des théorèmes de décomposition additive pour ce système de classes et pour certains de ses sous-systèmes. L'anneau quotient ainsi formé est d'une généralité égale à celle de l'anneau pris au départ; en effet, chaque anneau peut être conçu comme l'anneau quotient par l'idéal qui correspond à l'ensemble des relations identiques entre ses éléments; ou encore par un sous-système de ces relations, lorsqu'on suppose les autres relations déjà satisfaites dans le domaine.

Cette remarque donne aussi la place des travaux de Fraenkel.\footnote{A. Fraenkel, \emph{Über die Teiler der Null und die Zerlegung von Ringen}. J. f. M. 145 (1914), p. 139. \emph{Über gewisse Teilbereiche und Erweiterungen von Ringen}. Habilitationsschrift, Leipzig, Teubner, 1916. \emph{Über einfache Erweiterungen zerlegbarer Ringe}. J. f. M. 151 (1920), p. 121.} Fraenkel considère des décompositions additives d'anneaux qui sont soumises à des conditions restrictives telles (existence d'éléments réguliers, division par ceux-ci, condition de décomposabilité) que, pour l'idéal correspondant, les quatre décompositions coïncident. À cause de cette coïncidence, sa condition de finitude, selon laquelle l'idéal ne doit posséder qu'un nombre fini de diviseurs propres --- condition dont il s'écarte d'ailleurs partiellement ---, ne signifie pas une restriction plus forte que la nôtre. Le point de départ de Fraenkel est conditionné par les buts différents, essentiellement algébriques, de ses travaux; par extension algébrique il parvient ensuite à des anneaux plus généraux, avec des conditions moins restrictives.

\subsection*{\S~1. Anneau, idéal, condition de finitude.}

\textbf{1.} Le domaine $\Sigma$ pris pour base est un anneau (commutatif) dans la définition abstraite;\footnote{La définition est empruntée à la thèse d'habilitation de Fraenkel, en omettant ses conditions restrictives 6, I et II; pour cela il a fallu inclure la loi commutative de l'addition. Il s'agit donc des lois définissant un corps, en omettant l'inversibilité de la multiplication.} c'est-à-dire que $\Sigma$ consiste en un système d'éléments $a,b,c,\ldots,f,g,h,\ldots$ dans lequel une relation satisfaisant aux conditions usuelles est définie comme égalité; et dans lequel, par deux opérations (modes de composition), l'addition et la multiplication, deux éléments quelconques de l'anneau $a$ et $b$ donnent toujours de façon unique un troisième élément, respectivement comme somme $a+b$ et comme produit $a\cdot b$. L'anneau et les opérations, par ailleurs tout à fait arbitraires, doivent satisfaire aux lois suivantes:
\begin{enumerate}
\item la loi associative de l'addition: $(a+b)+c=a+(b+c)$;
\item la loi commutative de l'addition: $a+b=b+a$;
\item la loi associative de la multiplication: $(a\cdot b)\cdot c=a\cdot(b\cdot c)$;
\item la loi commutative de la multiplication: $a\cdot b=b\cdot a$;
\item la loi distributive: $a\cdot(b+c)=a\cdot b+a\cdot c$;
\item la loi de la soustraction sans restriction et unique. Il existe dans $\Sigma$ un seul élément $x$ qui satisfait l'équation $a+x=b$. (On note $x=b-a$.)
\end{enumerate}

De ces propriétés résulte l'existence du zéro; mais un anneau n'a pas besoin de posséder une unité; et le produit de deux éléments peut s'annuler sans que l'un des facteurs s'annule. Les anneaux pour lesquels l'annulation d'un produit entraîne toujours l'annulation d'un facteur, et qui possèdent en outre une unité, seront appelés domaines d'intégrité proprement dits. Pour la somme finie $a+a\ldots+a$ nous introduisons l'abréviation usuelle $na$, où les entiers $n$ doivent être regardés seulement comme des signes abréviatifs, non comme des éléments de l'anneau, et sont définis récursivement par $a=1\cdot a$, $na+a=(n+1)a$.

\textbf{2.} Par idéal $\ideal{M}$\footnote{Les idéaux sont désignés par de grandes lettres allemandes. $\ideal{M}$ doit rappeler l'exemple de l'idéal dans les polynômes ordinairement appelé \og{} module \fg{} ou module de formes. D'ailleurs les \S{}\S{}~1--3 n'utilisent que la propriété de module, et non la propriété d'idéal; comparer \S{}~9.} dans $\Sigma$, on entend un système d'éléments de $\Sigma$ qui satisfait aux deux conditions:
\begin{enumerate}
\item $\ideal{M}$ contient, avec $f$, aussi $a\cdot f$, où $a$ est un élément quelconque de $\Sigma$;
\item $\ideal{M}$ contient, avec $f$ et $g$, aussi la différence $f-g$; donc, avec $f$, aussi $nf$ pour tout entier $n$.
\end{enumerate}

Si $f$ est un élément de $\ideal{M}$, nous l'exprimons comme d'ordinaire par $f\equiv0(\ideal{M})$ et nous disons que $f$ est divisible par $\ideal{M}$. Si tout élément de $\ideal{N}$ est en même temps un élément de $\ideal{M}$, et donc divisible par $\ideal{M}$, nous disons: $\ideal{N}$ est divisible par $\ideal{M}$; en signes, $\ideal{N}\equiv0(\ideal{M})$. $\ideal{M}$ s'appelle diviseur propre de $\ideal{N}$ lorsqu'il contient des éléments différents de ceux de $\ideal{N}$, et n'est donc pas inversement divisible par $\ideal{N}$. De $\ideal{N}\equiv0(\ideal{M})$ et $\ideal{M}\equiv0(\ideal{N})$ résulte $\ideal{N}=\ideal{M}$.

Les autres notions connues sont conservées mot pour mot. Par plus grand commun diviseur de deux idéaux $\ideal{A}$ et $\ideal{B}$ --- $\ideal{D}=(\ideal{A},\ideal{B})$ ---, nous entendons l'ensemble des éléments qui peuvent être représentés sous la forme $a+b$, où $a$ parcourt tous les éléments de $\ideal{A}$ et $b$ tous ceux de $\ideal{B}$; $\ideal{D}$ est de nouveau un idéal. De même, le plus grand commun diviseur d'une infinité d'idéaux --- $\ideal{D}=(\ideal{A}_1,\ideal{A}_2,\ldots,\ideal{A}_\nu,\ldots)$ --- est défini comme l'ensemble des éléments $d$ représentables comme sommes d'éléments provenant à chaque fois d'un nombre fini d'idéaux:
$d=a_{i_1}+a_{i_2}+\cdots+a_{i_n}$; ici aussi $\ideal{D}$ est de nouveau un idéal.

Si, en particulier, l'idéal $\ideal{M}$ contient un nombre fini d'éléments $f_1,f_2,\ldots,f_\varrho$ tels que
\[
 \ideal{M}=(f_1,\ldots,f_\varrho),\qquad
 f=a_1f_1+\cdots+a_\varrho f_\varrho+n_1f_1+\cdots+n_\varrho f_\varrho
\]
pour tout $f\equiv0(\ideal{M})$, où les $a_i$ sont des éléments de l'anneau et les $n_i$ des entiers, alors $\ideal{M}$ s'appelle un idéal fini; $f_1,\ldots,f_\varrho$ une base d'idéal.

Dans la suite nous ne prendrons pour base que des anneaux $\Sigma$ qui satisfont à la condition de finitude: \emph{tout idéal de $\Sigma$ est fini, donc possède une base d'idéal.}

\textbf{3.} De la condition de finitude résulte directement le résultat sur lequel reposent toutes les considérations suivantes.

\textbf{Théorème I (théorème de la chaîne finie).}\footnote{Énoncé d'abord pour les modules de nombres par Dedekind: \emph{Zahlentheorie}, Suppl. XI, \S~172, théorème VIII (4e éd.); notre preuve et la désignation \og{} chaîne \fg{} sont reprises de là. Pour les idéaux de polynômes, voir Lasker, loc. cit., p. 56 (lemme). Mais, dans les deux cas, le théorème ne trouve qu'une application isolée. Nos applications reposent partout sur le postulat du choix.} \emph{Soit $\ideal{M},\ideal{M}_1,\ideal{M}_2,\ldots,\ideal{M}_\nu,\ldots$ un système dénombrable infini d'idéaux de $\Sigma$, dont chacun est divisible par le suivant. Alors, à partir d'un indice fini $n$, tous les idéaux sont identiques, $\ideal{M}_n=\ideal{M}_{n+1}=\cdots$. Autrement dit: si $\ideal M,\ideal M_1,\ideal M_2,\ldots,\ideal M_\nu,\ldots$ forment une chaîne simplement ordonnée d'idéaux telle que chaque idéal soit un diviseur propre de celui qui le précède immédiatement, alors la chaîne s'interrompt au bout d'un nombre fini de pas.}

En effet, soit $\ideal{D}=(\ideal{M}_1,\ideal{M}_2,\ldots,\ideal{M}_\nu,\ldots)$ le plus grand commun diviseur du système, et soit $f_1\ldots f_k$ une base de $\ideal{D}$, qui existe toujours en vertu de la condition de finitude. De l'hypothèse de divisibilité il résulte que chaque élément de $\ideal{D}$ est en même temps élément d'un idéal de la chaîne; car de
\[
 f=g+h,\qquad g\equiv0(\ideal{M}_r),\qquad h\equiv0(\ideal{M}_s),\qquad (r\le s)
\]
on obtient $g\equiv0(\ideal{M}_s)$ et donc $f\equiv0(\ideal{M}_s)$. Il en va de même lorsque $f$ est une somme de plusieurs constituants. Il existe donc un indice fini $n$ tel que
\[
 f_1\equiv0(\ideal{M}_n);\quad\ldots;\quad f_k\equiv0(\ideal{M}_n);\quad
 \ideal{D}=(f_1\ldots f_k)\equiv0(\ideal{M}_n).
\]
Comme, inversement, $\ideal{M}_n\equiv0(\ideal{D})$, on a $\ideal{M}_n=\ideal{D}$; et comme en outre
\[
 \ideal{M}_{n+i}\equiv0(\ideal{D});\qquad
 \ideal{D}=\ideal{M}_n\equiv0(\ideal{M}_{n+i}),
\]
on a aussi $\ideal{M}_{n+i}=\ideal{D}=\ideal{M}_n$ pour tout $i$, ce qui démontre le théorème.

Remarquons que, réciproquement, l'existence d'une base d'idéal résulte de nouveau de ce théorème, de sorte que la condition de finitude aurait aussi pu être énoncée sous cette forme indépendante d'une base.
\subsection*{\S~2. Représentation d'un idéal comme plus petit commun multiple d'un nombre fini d'idéaux irréductibles.}

Le plus petit commun multiple $[\ideal B_1,\ideal B_2,\ldots,\ideal B_k]$ des idéaux $\ideal B_1,\ideal B_2,\ldots,\ideal B_k$ est défini comme d'ordinaire comme l'ensemble des éléments qui sont divisibles aussi bien par $\ideal B_1$ que par $\ideal B_2,\ldots,\ideal B_k$; en symboles:
\[
 \text{de } f\equiv0(\ideal B_i),\quad (i=1,2,\ldots,k),
 \quad\text{il résulte}\quad
 f\equiv0([\ideal B_1,\ideal B_2,\ldots,\ideal B_k]),
\]
et réciproquement. Le plus petit commun multiple est de nouveau un idéal; nous appelons aussi les idéaux $\ideal B_i$ les composantes de la décomposition.

\textbf{Définition I.} Une représentation $\ideal M=[\ideal B_1\cdots\ideal B_k]$ s'appelle une représentation réduite si aucun $\ideal B_i$ ne divise le plus petit commun multiple $\ideal A_i$ des autres idéaux, et si aucun $\ideal B_i$ ne peut être remplacé par un diviseur propre.\footnote{Un exemple de représentation non réduite est
\[
 (x^2,xy)=[(x),(x^2,xy,y^\lambda)]
\]
pour tout exposant $\lambda\ge2$; la représentation $[(x),(x^2,y)]$ correspondant à $\lambda=1$ est une représentation réduite associée. Cette représentation m'a été indiquée par K. Hentzelt, tombé pendant la guerre, comme l'exemple le plus simple d'une décomposition non unique en idéaux primaires.} Si les conditions ne sont satisfaites que pour l'idéal $\ideal B_i$, la représentation s'appelle réduite par rapport à $\ideal B_i$. Le plus petit commun multiple $\ideal A_i=[\ideal B_1,\ldots,\ideal B_{i-1},\ideal B_{i+1},\ldots,\ideal B_k]$ est appelé le complément de $\ideal B_i$. Les représentations dans lesquelles seule la première condition est satisfaite s'appellent les représentations les plus courtes.

Pour représenter un idéal comme plus petit commun multiple, il suffit donc de se limiter aux représentations réduites, d'après le résultat suivant.

\textbf{Lemme I.} \emph{Toute représentation d'un idéal comme plus petit commun multiple d'un nombre fini d'idéaux peut être remplacée, d'au moins une manière, par une représentation réduite; une telle représentation peut en particulier être atteinte par décomposition successive.}

En effet, soit $\ideal M=[\ideal B_1^*\cdots\ideal B_l^*]$ une représentation quelconque de $\ideal M$. Nous omettons successivement les $\ideal B_i^*$ qui divisent le plus petit commun multiple des idéaux laissés en place. Puisque les idéaux restants donnent encore $\ideal M$, dans la représentation ainsi obtenue
\[
 \ideal M=[\ideal B_1\cdots\ideal B_k]=[\ideal A_i,\ideal B_i]
\]
la première condition est satisfaite; elle est donc une représentation la plus courte, et cette condition reste satisfaite si l'on remplace un $\ideal B_i$ quelconque par un diviseur propre. La seconde condition, en revanche, est toujours accessible par le théorème de la chaîne finie (théorème I). Supposons en effet que
\[
 \ideal M=[\ideal A_i,\ideal B_i]=[\ideal A_i,\ideal B_i^{(1)}]
 =\cdots=[\ideal A_i,\ideal B_i^{(\nu)}],\ldots,
\]
où chaque $\ideal B_i^{(\nu)}$ est un diviseur propre de celui qui le précède immédiatement. Alors, d'après ce théorème, la chaîne $\ideal B_i,\ideal B_i^{(1)},\ldots,\ideal B_i^{(\nu)},\ldots$ doit s'interrompre après un nombre fini d'étapes; dans la représentation $\ideal M=[\ideal A_i,\ideal B_i^{(n)}]$, l'idéal $\ideal B_i^{(n)}$ ne peut donc plus être remplacé par un diviseur propre; et cela vaut a fortiori lorsque $\ideal A_i$ est remplacé par un diviseur propre. En appliquant le procédé successivement à chaque $\ideal B_i$, en formant chaque fois le complément avec les $\ideal B$ déjà réduits, on obtient une représentation réduite.\footnote{Que cette représentation ne soit pas déterminée de manière unique par la représentation donnée, c'est ce que montre l'exemple précédent. Pour $(x^2,xy)=[(x),(x^2,xy,y^\lambda)]$, où $\lambda\ge2$, il y a, outre $[(x),(x^2,y)]$, aussi $[(x),(x^2,\mu x+y)]$, pour un $\mu$ arbitraire, comme représentation réduite.}

Pour obtenir une telle représentation successivement, il reste à montrer que des représentations réduites individuelles
\[
 \ideal M=[\ideal B_1,\ideal C_1],\qquad
 \ideal C_1=[\ideal B_2,\ideal C_2],\quad\ldots,\quad
 \ideal C_{k-1}=[\ideal B_k,\ideal C_k]
\]
il résulte que la représentation obtenue $\ideal M=[\ideal B_1\cdots\ideal B_k,\ideal C_k]$ est réduite. Pour cela, il suffit de montrer que, des représentations réduites
\[
 \ideal M=[\ideal B_1\cdots\ideal B_l,\ideal C],\qquad \ideal C=[\ideal C_1,\ideal C_2],
\]
on obtient une représentation réduite $\ideal M=[\ideal B_1\cdots\ideal B_\varrho,\ideal C_1,\ideal C_2]$. En effet, par hypothèse aucun des $\ideal B_i$ ne divise son complément; si cela arrivait pour un $\ideal C_i$, alors, contrairement à l'hypothèse dans la première représentation, $\ideal C$ pourrait être remplacé par un diviseur propre, puisque, par la seconde représentation réduite, $\ideal C_1$ et $\ideal C_2$ sont des diviseurs propres de $\ideal C$; la représentation est donc une représentation la plus courte. De plus, par hypothèse, aucun $\ideal B$ ne peut être remplacé par un diviseur propre; si cela était possible pour un $\ideal C_i$, cela correspondrait, contrairement à l'hypothèse, au remplacement de $\ideal C$ par un diviseur propre, puisque la représentation de $\ideal C$ est réduite. Le lemme est ainsi démontré.

\textbf{Définition II.} Un idéal $\ideal M$ s'appelle réductible s'il peut être représenté comme plus petit commun multiple de deux diviseurs propres; dans le cas contraire, $\ideal M$ s'appelle irréductible.

Nous démontrons maintenant, au moyen du théorème I de la chaîne finie et en utilisant les représentations réduites:

\textbf{Théorème II.} \emph{Tout idéal est représentable comme plus petit commun multiple d'un nombre fini d'idéaux irréductibles.}\footnote{Que, en général, une telle représentation ne soit pas unique, c'est ce que montre l'exemple précédent: $(x^2,xy)=[(x),(x^2,\mu x+y)]$ pour un $\mu$ arbitraire. Les deux composantes sont irréductibles pour tout $\mu$. En effet, tous les diviseurs de $(x)$ sont de la forme $(x,g(y))$, où $g(y)$ désigne un polynôme en $y$; le plus petit commun multiple de deux quelconques d'entre eux a donc aussi cette forme, et est par suite un diviseur propre de $(x)$. L'idéal $(x^2,\mu x+y)$ ne possède que l'unique diviseur $(x,y)$, et il est donc nécessairement lui aussi irréductible.}

Un idéal arbitraire $\ideal M$ est en effet soit irréductible; alors $\ideal M=[\ideal M]$ est une représentation du type exigé par le théorème II; soit $\ideal M=[\ideal B_1,\ideal C_1]$, où $\ideal B_1$ et $\ideal C_1$ sont des diviseurs propres de $\ideal M$, et où la représentation peut, d'après le lemme I, être supposée réduite. Pour $\ideal C_1$ vaut la même alternative: ou bien il est irréductible, ou bien il existe une représentation réduite $\ideal C_1=[\ideal B_2,\ideal C_2]$.

En poursuivant ainsi, on obtient la suite de représentations réduites
\begin{equation}
\begin{aligned}
 \ideal M&=[\ideal B_1,\ideal C_1];\qquad
 \ideal C_1=[\ideal B_2,\ideal C_2];\quad\ldots;\\
 \ideal C_{\nu-1}&=[\ideal B_\nu,\ideal C_\nu];\quad\ldots
\end{aligned}
\tag{1}
\end{equation}
Dans la chaîne $\ideal C_1,\ideal C_2,\ldots,\ideal C_\nu,\ldots$, chaque $\ideal C_i$ est un diviseur propre de celui qui le précède immédiatement; la chaîne s'arrête donc après un nombre fini d'étapes. Il existe un indice $n$ tel que $\ideal C_n$ soit irréductible. D'après le lemme I, en outre, la représentation $\ideal M=[\ideal B_1\cdots\ideal B_n\ideal C_n]$ est réduite; $\ideal C_n$ ne peut donc pas diviser son complément $\ideal A_n$, et, dans la représentation $\ideal M=[\ideal A_n,\ideal C_n]$, l'idéal $\ideal C_n$ ne peut pas être remplacé par un diviseur propre. En remplaçant, si nécessaire, $\ideal A_n$ par un diviseur propre,\footnote{En fait, la représentation est aussi réduite par rapport à $\ideal A_n$, comme il sera montré au \S~3 (lemme IV), en réciproque du lemme I.} de façon que la représentation devienne réduite, on a montré que tout idéal réductible admet une représentation réduite comme plus petit commun multiple d'un idéal irréductible et d'un idéal complémentaire. Dans la suite (1), on peut donc supposer sans restreindre la généralité que tous les $\ideal B_i$ sont irréductibles; la répétition du raisonnement précédent donne l'existence d'un $\ideal C_n$ irréductible, ce qui démontre le théorème II.
\subsection*{\S~3. Égalité du nombre des composantes dans deux décompositions différentes en idéaux irréductibles.}

Pour démontrer l'égalité du nombre des composantes, il faut d'abord exprimer la réductibilité, ou l'irréductibilité, d'un idéal par des propriétés de son complément, selon le résultat suivant.

\textbf{Théorème III.}\footnote{Le théorème III correspond au passage des modules aux groupes quotients dans les travaux de Schmeidler et de Noether--Schmeidler (voir l'introduction). À $\ideal A$ correspond le groupe quotient; $\ideal N_1$ et $\ideal N_2$ correspondent aux sous-groupes dans lesquels le groupe quotient se décompose.} \emph{Soit la représentation la plus courte $\ideal M=[\ideal A,\ideal C]$, réduite par rapport à $\ideal C$. Alors la condition nécessaire et suffisante pour que $\ideal C$ soit réductible est qu'il existe deux idéaux $\ideal N_1$ et $\ideal N_2$, diviseurs propres de $\ideal M$, tels que}
\begin{equation}
 \ideal N_1\equiv0(\ideal A),\qquad
 \ideal N_2\equiv0(\ideal A),\qquad
 [\ideal N_1,\ideal N_2]=\ideal M.
\tag{2}
\end{equation}
\emph{Il en résulte encore ceci: si les conditions (2) sont satisfaites et si $\ideal C$ est irréductible, alors au moins l'un des $\ideal N_i$ n'est pas un diviseur propre de $\ideal M$; on a $\ideal N_i=\ideal M$.}

Soit $\ideal C=[\ideal C_1,\ideal C_2]$, où $\ideal C_1$ et $\ideal C_2$ sont des diviseurs propres de $\ideal C$. Alors
\[
 \ideal M=[\ideal A,\ideal C]=[\ideal A,\ideal C_1,\ideal C_2]
      =[[\ideal A,\ideal C_1],[\ideal A,\ideal C_2]].
\]
Ici les idéaux $[\ideal A,\ideal C_i]$ sont des diviseurs propres de $\ideal M$; sinon $[\ideal A,\ideal C]$ ne serait pas réduite par rapport à $\ideal C$. Comme la divisibilité par $\ideal A$ est elle aussi satisfaite, la nécessité de la condition (2) est démontrée. La représentation (2) n'est pas réduite, puisqu'un $[\ideal A,\ideal C_i]$ peut être remplacé par $\ideal C_i$.

Supposons réciproquement que (2) soit satisfaite. Nous formons les idéaux
\[
 \ideal C_1=(\ideal C,\ideal N_1),\qquad
 \ideal C_2=(\ideal C,\ideal N_2),\qquad
 \ideal C^*=[\ideal C_1,\ideal C_2].
\]
Alors $\ideal C$ est divisible aussi bien par $\ideal C_1$ que par $\ideal C_2$, donc aussi par leur plus petit commun multiple $\ideal C^*$. Pour montrer que $\ideal C^*$ est divisible par $\ideal C$, soit
\[
 f\equiv0(\ideal C^*);\quad f\equiv0(\ideal C_1);\quad f\equiv0(\ideal C_2),
 \quad\text{ou encore}\quad f=c+n_1;\quad f=\bar c+n_2,
\]
où $c,\bar c,n_1,n_2$ sont respectivement des grandeurs de $\ideal C,\ideal N_1,\ideal N_2$; en particulier $n_i$ est divisible par $\ideal N_i$. Donc la différence
\[
 g=c-\bar c=n_2-n_1
\]
est divisible aussi bien par $\ideal C$ que par $\ideal A$, donc par $\ideal M$. Puisque $n_1=n_2+m$, $n_1$ (et de même $n_2$) est en outre divisible aussi bien par $\ideal N_1$ que par $\ideal N_2$, donc par $\ideal M$. Ainsi
\[
 f=c+m;\qquad f\equiv0(\ideal C);\qquad \ideal C^*=\ideal C.
\]
Les idéaux $\ideal C_1$ et $\ideal C_2$ sont alors des diviseurs propres de $\ideal C$; car si $\ideal C_i=(\ideal C,\ideal N_i)=\ideal C$, alors $\ideal N_i$ serait divisible par $\ideal C$, et donc, en raison de sa divisibilité par $\ideal A$, égal à $\ideal M$, contrairement à l'hypothèse. Ainsi $\ideal C=[\ideal C_1,\ideal C_2]$ est reconnu réductible; le théorème III est démontré.

Remarquons que presque le même raisonnement montre encore ce qui suit.

\textbf{Lemme II.} \emph{Si, dans une représentation la plus courte $\ideal M=[\ideal A,\ideal C]$, l'idéal $\ideal C$ peut être remplacé par un diviseur propre, alors $\ideal C$ est réductible.}

Soit donc
\[
 \ideal M=[\ideal A,\ideal C]=[\ideal A,\ideal C_1],
\]
et posons
\[
 \ideal C^*=[\ideal C_1,(\ideal A,\ideal C)].
\]
Alors $\ideal C$ est de nouveau divisible par $\ideal C^*$. De $f\equiv0(\ideal C^*)$ il résulte en outre
\[
 f=c_1=a+c.
\]
La différence $a=c_1-c$ est donc divisible aussi bien par $\ideal A$ que par $\ideal C_1$, par conséquent par $\ideal M$; ainsi $c_1=c+m$, $f\equiv0(\ideal C)$, $\ideal C^*=\ideal C$. Comme, par hypothèse, $\ideal C_1$ et $(\ideal A,\ideal C)$ sont tous deux des diviseurs propres de $\ideal C$, $\ideal C=\ideal C^*$ est reconnu réductible.

Un $\ideal C$ irréductible ne peut donc pas être remplacé par un diviseur propre.

Supposons maintenant données deux représentations les plus courtes différentes de $\ideal M$ comme plus petit commun multiple d'un nombre fini d'idéaux irréductibles:
\[
 \ideal M=[\ideal B_1\cdots\ideal B_k]=[\ideal D_1\cdots\ideal D_l].
\]
D'après la remarque qui suit le lemme II, ces représentations sont en même temps réduites. Nous démontrons d'abord le lemme suivant.

\textbf{Lemme III.} \emph{À tout complément $\ideal A_i=[\ideal B_1\cdots\ideal B_{i-1}\ideal B_{i+1}\cdots\ideal B_k]$ correspond un idéal $\ideal D_j$ tel que $\ideal M=[\ideal A_i,\ideal D_j]$.}

En effet, si l'on pose $\ideal M=[\ideal D_1,\ideal C_1]$, $\ideal C_1=[\ideal D_2,\ideal C_{12}]$, etc., on obtient
\[
 \ideal M=[\ideal A_i,\ideal M]=[\ideal A_i,\ideal D_1,\ideal C_1]
       =[[\ideal A_i\ideal D_1],[\ideal A_i\ideal C_1]].
\]
Ici, pour $\ideal N_1=[\ideal A_i,\ideal D_1]$ et $\ideal N_2=[\ideal A_i,\ideal C_1]$, les conditions (2) du théorème III sont satisfaites, puisque $\ideal M=[\ideal A_i,\ideal B_i]$ est réduite par rapport à $\ideal B_i$ et que la représentation est la plus courte. Comme $\ideal B_i$ était supposé irréductible, il faut nécessairement que l'un des $\ideal N_i$ devienne égal à $\ideal M$.

Si $\ideal N_1=\ideal M$, le lemme est démontré. Si $\ideal N_2=\ideal M$, on obtient de même $\ideal M=[[\ideal A_i\ideal D_2],[\ideal A_i\ideal C_{12}]]$, où, par le même raisonnement, une composante doit de nouveau devenir égale à $\ideal M$. En poursuivant ainsi, ou bien on obtient $\ideal M=[\ideal A_i,\ideal D_j]$, avec $j<l$, ou bien $\ideal M=[\ideal A_i,\ideal C_{1\ldots l-1}]$; puisque $\ideal C_{1\ldots l-1}=\ideal D_l$, le lemme est démontré.

Il en résulte maintenant le théorème suivant.

\textbf{Théorème IV.} \emph{Dans deux représentations les plus courtes différentes d'un idéal comme plus petit commun multiple d'idéaux irréductibles, le nombre des composantes est le même.}

Du lemme il résulte en effet, pour $i=1$,
\[
 \ideal M=[\ideal A_1,\ideal B_1]=[\ideal A_1,\ideal D_{j_1}]
          =[\ideal D_{j_1},\ideal B_2,\ldots,\ideal B_k].
\]
Considérons maintenant les deux décompositions
\[
 \ideal M=[\ideal D_{j_1},\ideal B_2,\ldots,\ideal B_k]
          =[\ideal D_1,\ideal D_2,\ldots,\ideal D_l],
\]
et répétons le raisonnement précédent relativement au complément $\bar{\ideal A}_2=[\ideal D_{j_1},\ideal B_3,\ldots,\ideal B_k]$ de $\ideal B_2$; on obtient
\[
 \ideal M=[\bar{\ideal A}_2,\ideal B_2]=[\bar{\ideal A}_2,\ideal D_{j_2}]
          =[\ideal D_{j_1},\ideal D_{j_2}\ideal B_3,\ldots,\ideal B_k],
\]
et, en poursuivant le procédé,
\[
 \ideal M=[\ideal D_{j_1},\ideal D_{j_2},\ldots,\ideal D_{j_k}].
\]
Puisque, par hypothèse, la représentation $\ideal M=[\ideal D_1\cdots\ideal D_l]$ est la plus courte, aucun $\ideal D$ ne peut être omis; les idéaux distincts parmi les $\ideal D_{j_i}$ doivent donc épuiser tous les $\ideal D$, et ainsi $k\ge l$. Si, dans le lemme et dans les raisonnements qui suivent, on échange partout les $\ideal B$ et les $\ideal D$, on obtient de même $l\ge k$, et donc $k=l$, ce qui démontre l'égalité du nombre. Il en résulte encore que les idéaux $\ideal D_{j_i}$ sont tous distincts entre eux; sinon, dans une représentation la plus courte par les $\ideal D_{j_i}$, il apparaîtrait moins de $k$ composantes. On peut donc choisir la notation de sorte que $\ideal D_{j_i}=\ideal D_i$. Pour la même raison, toutes les représentations intermédiaires $\ideal M=[\ideal D_{j_1}\cdots\ideal D_{j_i},\ideal B_{i+1}\cdots\ideal B_k]$ sont les plus courtes et donc, d'après la remarque qui suit le lemme II, réduites.

L'égalité du nombre conduit à une réciproque du lemme I.

\textbf{Lemme IV.} \emph{Si, dans une représentation réduite, on réunit les composantes en groupes et que l'on forme leurs plus petits communs multiples, la représentation obtenue est réduite. Autrement dit, d'une représentation réduite}
\[
 \ideal M=[\ideal C_{11}\cdots\ideal C_{1\mu_1};\;\ldots;\;,
           \ideal C_{\sigma1}\cdots\ideal C_{\sigma\mu_\sigma}]
\]
\emph{il résulte que $\ideal M=[\ideal N_1\cdots\ideal N_\sigma]=[\ideal N_i,\ideal L_i]$ est aussi réduite, si l'on pose $\ideal N_i=[\ideal C_{i1}\cdots\ideal C_{i\mu_i}]$.}

Il faut d'abord remarquer que $\ideal N_i$ ne peut pas diviser son complément $\ideal L_i$, puisque aucun de ses diviseurs $\ideal C_{ij}$ ne divise $\ideal L_i$; la représentation est donc la plus courte. Pour montrer que $\ideal N_i$ ne peut être remplacé par aucun diviseur propre, nous décomposons les $\ideal C$ en leurs idéaux irréductibles $\ideal B$\footnote{Il faut toujours entendre par là une représentation par les $\ideal B$ qui soit la plus courte, donc réduite.}, de sorte que les représentations réduites (d'après le lemme I) suivantes apparaissent:
\[
 \ideal M=[\ideal B_{11}\cdots\ideal B_{1\lambda_1};\;\ldots;\;,
           \ideal B_{\sigma1}\cdots\ideal B_{\sigma\lambda_\sigma}],
 \qquad
 \ideal N_i=[\ideal B_{i1}\cdots\ideal B_{i\lambda_i}].
\]

Soit maintenant $\ideal M=[\ideal N_i^*,\ideal L_i]$ réduite par rapport à $\ideal N_i^*$, où $\ideal N_i^*$ est un diviseur propre de $\ideal N_i$. D'après le lemme II,
\[
 \ideal N_i=[\ideal N_i^*, (\ideal N_i,\ideal L_i)];
\]
et cette représentation est nécessairement réduite par rapport à $\ideal N_i^*$, car sinon $\ideal N_i^*$ pourrait aussi être remplacé dans $\ideal M$ par un diviseur propre. Remplaçons au besoin aussi $(\ideal N_i,\ideal L_i)$ par un diviseur propre, de sorte qu'une représentation réduite de $\ideal N_i$ soit obtenue. Si l'on décompose maintenant les deux composantes de $\ideal N_i$ en idéaux irréductibles, le nombre $\lambda_i$ des idéaux irréductibles de $\ideal N_i$ se compose additivement des nombres correspondant aux composantes; le nombre d'idéaux irréductibles correspondant au diviseur propre $\ideal N_i^*$ est donc nécessairement plus petit que $\lambda_i$. Mais alors la décomposition de $\ideal M=[\ideal N_i^*,\ideal L]$ en idéaux irréductibles conduirait aussi à moins de $\sum_i\lambda_i$ idéaux, en contradiction avec l'égalité du nombre. Comme cas particulier $\sigma=2$, il en résulte encore que la représentation $\ideal M=[\ideal N_i,\ideal L_i]$ est aussi réduite par rapport au complément $\ideal L_i$.
\subsection*{\S~4. Idéaux primaires. Unicité des idéaux premiers associés dans deux décompositions différentes en idéaux irréductibles.}

Il s'agit, dans ce qui suit, du rapport entre idéaux primaires et idéaux irréductibles.

\textbf{Définition III.} Un idéal $\ideal Q$ est dit primaire si de $a\cdot b\equiv0(\ideal Q)$ et $a\not\equiv0(\ideal Q)$, il s'ensuit nécessairement que $b^\varkappa\equiv0(\ideal Q)$, où l'exposant $\varkappa$ est un nombre fini.

La définition peut aussi s'énoncer ainsi: si un produit $a\cdot b$ est divisible par $\ideal Q$, alors ou bien un facteur est divisible, ou bien une puissance de chacun des facteurs l'est. Si en particulier $\varkappa$ est toujours égal à $1$, l'idéal s'appelle un idéal premier.

De la définition de l'idéal primaire (respectivement premier), il résulte, grâce à l'existence d'une base, la définition qui ne porte que sur les produits d'idéaux\footnote{Par le produit $\ideal A\cdot\ideal B$ de deux idéaux, on entend, comme d'usage, l'idéal constitué par l'ensemble des grandeurs $a\cdot b$ et de leurs sommes finies.}:

\textbf{Définition IIIa.} Un idéal $\ideal Q$ est dit primaire si de $\ideal A\cdot\ideal B\equiv0(\ideal Q)$ et $\ideal A\not\equiv0(\ideal Q)$, il s'ensuit nécessairement que $\ideal B^\lambda\equiv0(\ideal Q)$. Si $\lambda$ est toujours égal à $1$, l'idéal s'appelle un idéal premier. Pour un idéal premier $\ideal P$, il résulte donc toujours de $\ideal A\cdot\ideal B\equiv0(\ideal P)$ et de $\ideal A\not\equiv0(\ideal P)$ que $\ideal B\equiv0(\ideal P)$.

En effet, puisque IIIa contient la définition III comme cas particulier pour $\ideal A=(a)$, $\ideal B=(b)$, tout idéal primaire d'après IIIa est aussi primaire d'après III. Réciproquement, soit $\ideal Q$ primaire d'après III, et supposons remplie l'hypothèse de IIIa: $\ideal A\cdot\ideal B\equiv0(\ideal Q)$. Alors ou bien $\ideal A\equiv0(\ideal Q)$ s'ensuit, ou bien il existe au moins une grandeur $a\equiv0(\ideal A)$ telle que $a\cdot\ideal B\equiv0(\ideal Q)$ et $a\not\equiv0(\ideal Q)$. Si maintenant $b_1\ldots b_r$ est une base idéale de $\ideal B$, alors, d'après la définition III, puisque $a\cdot b_i\equiv0(\ideal Q)$, on obtient
\[
 b_1^{\varkappa_1}\equiv0(\ideal Q),\quad \ldots,\quad b_r^{\varkappa_r}\equiv0(\ideal Q).
\]
À cause de
\[
 b=f_1b_1+\cdots+f_rb_r+n_1b_1+\cdots+n_rb_r,
\]
il s'ensuit donc que, pour $\lambda=\varkappa_1+\cdots+\varkappa_r$, le produit de $\lambda$ grandeurs $b$ quelconques est divisible par $\ideal Q$; ainsi, pour les idéaux primaires au sens de III, l'accomplissement de la définition IIIa est démontré. Dans le cas spécial des idéaux premiers $\ideal P$, il résulte de $a\cdot\ideal B\equiv0(\ideal P)$, donc aussi de $a\cdot b\equiv0(\ideal P)$ pour chaque $b\equiv0(\ideal B)$, avec $a\not\equiv0(\ideal P)$, que $b\equiv0(\ideal P)$, et par suite $\ideal B\equiv0(\ideal P)$. L'équivalence des deux définitions est ainsi établie.

Le rapport entre idéaux primaires et idéaux premiers est établi par la remarque suivante: l'ensemble $\ideal P$ de tous les éléments $p$ qui ont la propriété qu'une puissance de $p$ est divisible par $\ideal Q$ forme un idéal premier. Il est d'abord clair que $\ideal P$ est un idéal, car avec $p_1$ et $p_2$, les éléments $ap_1$ et $(p_1-p_2)$ possèdent aussi la propriété indiquée. Par le raisonnement de base employé dans la définition IIIa, il résulte en outre qu'il existe un nombre $\lambda$ tel que $\ideal P^\lambda\equiv0(\ideal Q)$.

Soit maintenant
\[
 a\cdot b\equiv0(\ideal P);\qquad a\not\equiv0(\ideal P).
\]
Par la définition de $\ideal P$, on obtient alors
\[
 a^\lambda\cdot b^\lambda\equiv0(\ideal Q);\qquad a^\lambda\not\equiv0(\ideal Q);
\]
donc, par la définition de $\ideal Q$,
\[
 b^{\lambda\varkappa}\equiv0(\ideal Q)\qquad\text{et par conséquent}\qquad b\equiv0(\ideal P),
\]
ce qui démontre que $\ideal P$ est un idéal premier. L'idéal $\ideal P$ est aussi défini comme le plus grand commun diviseur de tous les idéaux $\ideal B$ ayant la propriété qu'une puissance de $\ideal B$ est divisible par $\ideal Q$. En effet, chacun de ces $\ideal B$ est divisible par $\ideal P$, par définition; donc aussi leur plus grand commun diviseur $\ideal D$. Inversement, $\ideal P$ lui-même est un idéal $\ideal B$, donc divisible par $\ideal D$, ce qui prouve $\ideal P=\ideal D$.

Ainsi $\ideal P$ est un idéal premier, diviseur de $\ideal Q$, et dont une puissance est divisible par $\ideal Q$; par cette propriété il est défini de façon univoque. Car de
\[
 \ideal Q\equiv0(\ideal P);\qquad \ideal P^\lambda\equiv0(\ideal Q)
 \ideal Q\equiv0(\bar{\ideal P});\qquad \bar{\ideal P}^{\mu}\equiv0(\ideal Q)
\]
il s'ensuit
\[
 \ideal P^\lambda\equiv0(\bar{\ideal P});\qquad
 \bar{\ideal P}^{\mu}\equiv0(\ideal P),
\]
donc, d'après la propriété des idéaux premiers,
\[
 \ideal P\equiv0(\bar{\ideal P});\qquad
 \bar{\ideal P}\equiv0(\ideal P);\qquad
 \ideal P=\bar{\ideal P}.
\]

En résumé, nous avons:

\textbf{Théorème V.} \emph{À chaque idéal primaire $\ideal Q$ correspond un et un seul idéal premier $\ideal P$, qui est diviseur de $\ideal Q$ et dont une puissance est divisible par $\ideal Q$; $\ideal P$ sera appelé \og{} idéal premier associé \fg{}\footnote{Que la réciproque ne vaille pas est montré par l'exemple $\ideal M=(x^2,xy)$. L'idéal premier $(x)$ satisfait à toutes les conditions, mais $\ideal M$ n'est pas primaire.}. $\ideal P$ est défini comme le plus grand commun diviseur de tous les idéaux $\ideal B$ ayant la propriété qu'une puissance de $\ideal B$ est divisible par $\ideal Q$. Si $\varrho$ est le plus petit nombre tel que $\ideal P^\varrho\equiv0(\ideal Q)$, alors $\varrho$ sera appelé l'exposant de $\ideal Q$\footnote{Le fait qu'en général, contrairement au domaine des entiers rationnels ou algébriques, on n'ait pas $\ideal P^\varrho=\ideal Q$ est montré par l'exemple
\[
 \ideal Q=(x^2,y),\qquad \ideal P=(x,y),\qquad
 \ideal P^2=(x^2,xy,y^2)\equiv0(\ideal Q),
 \quad \text{mais }\ideal Q\not\equiv0(\ideal P^2);
\]
donc $\ideal Q$ est différent de $\ideal P^2$.}.}

Nous démontrons maintenant, comme lien entre primaire et irréductible:

\textbf{Théorème VI.} \emph{Tout idéal non primaire est réductible; en d'autres termes, tout idéal irréductible est primaire\footnote{Que la réciproque ne vaille pas ici est montré par exemple par
\[
 \ideal Q=(x^2,xy,y^\lambda)=[(x^2,y),(x,y^\lambda)],
\]
où $\lambda\ge2$. Ici $\ideal Q$ est primaire, mais réductible. Le fait que $\ideal Q$ soit primaire résulte de ce qu'il contient tous les produits de puissances de dimension $\lambda$ en $x,y$; pour chaque polynôme sans terme constant, une puissance est donc divisible par $\ideal Q$. Mais si, dans $a\cdot b\equiv0(\ideal Q)$, le polynôme $b$ contient un terme constant --- donc si $b^\varkappa\not\equiv0(\ideal Q)$ pour tout $\varkappa$ ---, alors, à cause de l'homogénéité des polynômes générateurs de $\ideal Q$, chaque composante homogène de $a\cdot b$ étant divisible par $\ideal Q$, il faut que $a$ soit divisible par $\ideal Q$.}.}

Soit $\ideal K$ un idéal non primaire; d'après la définition III, il existe donc au moins un couple de grandeurs $a,b$ tel que
\begin{equation}
 a\cdot b\equiv0(\ideal K),\qquad
 a\not\equiv0(\ideal K),\qquad
 b^\varkappa\not\equiv0(\ideal K)\quad \text{pour tout }\varkappa.
\tag{3}
\end{equation}

Nous formons maintenant les deux idéaux
\[
 \ideal L_0=(\ideal K,a);\qquad \ideal N_0=(\ideal K,b),
\]
qui, d'après (3), sont des diviseurs propres de $\ideal K$ et pour lesquels, d'après (3), on a
\begin{equation}
 \ideal L_0\cdot\ideal N_0\equiv0(\ideal K).
\tag{4}
\end{equation}

Pour les éléments $f$ du plus petit commun multiple $\ideal K_0=[\ideal L_0,\ideal N_0]$, l'alternative suivante vaut. Ou bien de
\[
 f\equiv0(\ideal L_0),\qquad f\equiv0(\ideal N_0),\qquad \text{c'est-à-dire}\qquad f\equiv a_1\cdot b\pmod{\ideal K}
\]
il résulte toujours une représentation
\[
 f\equiv l_0\cdot b\pmod{\ideal K},\qquad l_0\equiv0(\ideal L_0);
\]
alors, d'après (4), $f\equiv0(\ideal K)$; ainsi $\ideal K_0\equiv0(\ideal K)$, et, puisque $\ideal K\equiv0(\ideal K_0)$, on a aussi $\ideal K=\ideal K_0$, ce qui démontre que $\ideal K$ est réductible.

Ou bien il existe au moins un $f\equiv0(\ideal K_0)$ pour lequel aucun tel $l_0$ n'existe. Nous formons alors, avec le $a_1$ appartenant à ce $f$,
\[
 \ideal L_1=(\ideal L_0,a_1)=(\ideal K,a,a_1),\qquad
 \ideal N_1=(\ideal K,b^2).
\]
Alors, puisque $a_1\cdot b\equiv0(\ideal L_0)$, il résulte aussi de (4) que
\begin{equation}
 \ideal L_1\cdot\ideal N_1\equiv0(\ideal K),
\tag{4'}
\end{equation}
et $\ideal L_1$ devient un diviseur propre de $\ideal L_0$.

Pour les éléments $f$ de $\ideal K_1=[\ideal L_1,\ideal N_1]$, la même alternative vaut. Ou bien de
\[
 f\equiv0(\ideal L_1),\qquad f\equiv0(\ideal N_1),
\]
c'est-à-dire de $f\equiv a_2\cdot b^2\pmod{\ideal K}$, il résulte toujours
\[
 f\equiv l_1\cdot b^2,
 \qquad l_1\equiv0(\ideal L_1),
\]
et alors, d'après (4'),
\[
 \ideal K=\ideal K_1.
\]

Ou bien, pour au moins un $f$, il n'existe aucun tel $l_1$; cela conduit à former $\ideal L_2=(\ideal L_1,a_2)$, $\ideal N_2=(\ideal K,b^{2^2})$, avec $\ideal L_2\cdot\ideal N_2\equiv0(\ideal K)$, où $\ideal L_2$ est un diviseur propre de $\ideal L_1$. En poursuivant ainsi, nous définissons en général
\[
 \ideal L_0=(\ideal K,a),\quad
 \ideal L_1=(\ideal L_0,a_1);\quad\ldots;\quad
 \ideal L_\nu=(\ideal L_{\nu-1},a_\nu);\quad\ldots;
\]
\[
 \ideal N_0=(\ideal K,b),\quad
 \ideal N_1=(\ideal K,b^2),\quad \ldots,\quad
 \ideal N_\nu=(\ideal K,b^{2^\nu}),\quad \ldots,
\]
où les $a_i$ sont définis par le fait qu'il existe un $f$ tel que
\[
 f\equiv0(\ideal L_{i-1}),\qquad f\equiv0(\ideal N_{i-1}),
\]
c'est-à-dire
\[
 f\equiv a_i\cdot b^{2^{i-1}}\pmod{\ideal K},\qquad
 a_i\not\equiv0(\ideal L_{i-1}).
\]
D'après cela, on a en général $\ideal L_i\cdot\ideal N_i\equiv0(\ideal K)$; $\ideal N_i$ est, par (3), un diviseur propre de $\ideal K$, et $\ideal L_i$ un diviseur propre de $\ideal L_{i-1}$. D'après le théorème I de la chaîne finie, la chaîne des $\ideal L$ doit donc s'arrêter au bout d'un nombre fini de pas, disons avec $\ideal L_n$. Pour chaque $f\equiv0(\ideal L_n)$, $f\equiv0(\ideal N_n)$, on a donc $f\equiv l_n\cdot b^{2^n}\;(\ideal K)$ avec $l_n\equiv0(\ideal L_n)$; par conséquent, d'après le raisonnement ci-dessus, on obtient $\ideal K=[\ideal L_n,\ideal N_n]$, ce qui démontre que $\ideal K$ est réductible.

De ce qui vient d'être démontré résulte l'unicité des idéaux premiers associés, comme suit. Soient
\[
 \ideal M=[\ideal B_1\cdots\ideal B_k]=[\ideal D_1\cdots\ideal D_k]
\]
deux représentations les plus courtes, donc réduites, de $\ideal M$ comme plus petit commun multiple d'idéaux irréductibles, dont le nombre concorde d'après le théorème IV. D'après ce théorème, les représentations intermédiaires qui y interviennent, où, comme il y est remarqué, on peut poser l'indice $j_i=i$, sont aussi les représentations les plus courtes:
\[
\begin{aligned}
 \ideal M&=[\ideal D_1\cdots\ideal D_{i-1}\ideal B_i\ideal B_{i+1}\cdots\ideal B_k]
       =[\ideal D_1\cdots\ideal D_{i-1}\ideal D_i\ideal B_{i+1}\cdots\ideal B_k]\\
       &=[\bar{\ideal A}_i,\ideal B_i]=[\bar{\ideal A}_i,\ideal D_i].
\end{aligned}
\]
Il vient donc
\[
 \overline{\ideal A}_i\cdot\ideal B_i\equiv0(\ideal D_i),\qquad
 \overline{\ideal A}_i\not\equiv0(\ideal D_i),
 \qquad
 \overline{\ideal A}_i\cdot\ideal D_i\equiv0(\ideal B_i),\qquad
 \overline{\ideal A}_i\not\equiv0(\ideal B_i).
\]

Or, d'après le théorème VI, les idéaux irréductibles $\ideal B_i$ et $\ideal D_i$ sont primaires; il en résulte l'existence de deux nombres $\lambda_i$ et $\mu_i$ tels que
\begin{equation}
 \ideal B_i^{\lambda_i}\equiv0(\ideal D_i),
 \qquad
 \ideal D_i^{\mu_i}\equiv0(\ideal B_i).
\tag{5}
\end{equation}
Désignons par $\ideal P_i$ et $\bar{\ideal P}_i$ les idéaux premiers associés de $\ideal B_i$ et de $\ideal D_i$; ainsi $\ideal P_i^{\varrho_i}\equiv0(\ideal B_i)$ et $\bar{\ideal P}_i^{\sigma_i}\equiv0(\ideal D_i)$. Alors, d'après (5),
\[
 \ideal P_i^{\lambda_i\varrho_i}\equiv0(\bar{\ideal P}_i),
 \qquad
 \bar{\ideal P}_i^{\mu_i\sigma_i}\equiv0(\ideal P_i),
\]
et, par la propriété des idéaux premiers,
\[
 \ideal P_i\equiv0(\bar{\ideal P}_i),
 \qquad
 \bar{\ideal P}_i\equiv0(\ideal P_i),
 \qquad
 \ideal P_i=\bar{\ideal P}_i.
\]
Cela démontre:

\textbf{Théorème VII.} \emph{Dans deux représentations les plus courtes différentes d'un idéal comme plus petit commun multiple d'idéaux irréductibles, les idéaux premiers associés concordent; parmi eux peuvent aussi apparaître des idéaux égaux\footnote{C'est ce que montre par exemple l'exemple de la note 19:
\[
 (x^2,xy,y^\lambda)=[(x^2,y),(x,y^\lambda)],
\]
où $\lambda\ge2$. Les deux idéaux premiers associés sont ici $(x,y)$.}, et ils apparaissent aussi le même nombre de fois dans chaque décomposition. Les idéaux eux-mêmes peuvent donc être associés deux à deux d'au moins une manière telle qu'à chaque fois une puissance de l'idéal $\ideal B_i$ soit divisible par l'idéal $\ideal D_i$ qui lui est associé, et réciproquement. Leur nombre concorde d'après le théorème IV\footnote{Pour l'unicité des idéaux \og{} isolés \fg{} contenus parmi les idéaux irréductibles, voir \S~7.}.}
\subsection*{\S~5. Représentation d'un idéal comme plus petit commun multiple de plus grands idéaux primaires. Unicité des idéaux premiers associés.}

\textbf{Définition IV.} \emph{Une représentation la plus courte $\ideal M=[\ideal Q_1\cdots\ideal Q_\alpha]$ sera dite représentation comme plus petit commun multiple de plus grands idéaux primaires lorsque tous les $\ideal Q$ sont primaires, mais que le plus petit commun multiple de deux des $\ideal Q$ n'est plus primaire.}

Qu'il existe toujours au moins une telle représentation résulte de la représentation de $\ideal M$ comme plus petit commun multiple d'idéaux irréductibles. Car ces idéaux sont primaires; ou bien on a déjà une représentation par de plus grands primaires, ou bien le plus petit commun multiple de deux quelconques de ces idéaux redevient primaire. Comme le nombre des idéaux a alors diminué d'une unité, la répétition du procédé conduit, après un nombre fini d'étapes, à la représentation cherchée.

Cette représentation est réduite d'après le lemme IV. Réciproquement, toute représentation réduite par de plus grands idéaux primaires naît de cette manière, comme le montre la résolution des $\ideal Q$ en idéaux irréductibles.

Pour déduire ici du théorème VII un théorème d'unicité correspondant, il faut étudier le lien avec les idéaux premiers associés, d'après le théorème suivant.

\textbf{Théorème VIII.} \emph{Si les idéaux primaires $\ideal N_1,\ideal N_2,\ldots,\ideal N_\lambda$ ont tous le même idéal premier associé $\ideal P$, alors leur plus petit commun multiple $\ideal Q=[\ideal N_1,\ideal N_2,\ldots,\ideal N_\lambda]$ est lui aussi primaire et a $\ideal P$ pour idéal premier associé. Réciproquement, si $\ideal Q=[\ideal N_1\cdots\ideal N_\lambda]$ est une représentation réduite de l'idéal primaire $\ideal Q$, alors tous les $\ideal N_i$ sont primaires et ont pour idéal premier associé l'idéal premier associé $\ideal P$ de $\ideal Q$.}

Pour démontrer la première partie de l'assertion, remarquons d'abord que de $\ideal P^{\varrho_i}\equiv0(\ideal N_i)$ pour chaque $i$, on déduit aussi que $\ideal P^\tau\equiv0(\ideal Q)$, où $\tau$ désigne le plus grand des indices $\varrho_i$. Comme $\ideal P$ est en outre un diviseur de $\ideal Q$, $\ideal P$ est nécessairement l'idéal premier associé si $\ideal Q$ est primaire. De
\[
 \ideal A\cdot\ideal B\equiv0(\ideal Q);\qquad
 \ideal B^k\not\equiv0(\ideal Q)\quad\text{(pour tout }k)
\]
il s'ensuit donc
\[
 \ideal B\not\equiv0(\ideal P);
 \quad\text{donc}\quad
 \ideal B^k\not\equiv(\ideal N_i)\quad\text{(pour tout }k),
\]
et par suite $\ideal A\equiv0(\ideal N_i)$; donc $\ideal A\equiv0(\ideal Q)$, ce qui prouve que $\ideal Q$ est primaire et que $\ideal P$ est l'idéal premier associé.

Réciproquement, soit d'abord
\[
 \ideal Q=[\ideal N_1\cdots\ideal N_\lambda]=[\ideal N_i,\ideal L_i]
\]
une représentation la plus courte de $\ideal Q$ par des idéaux primaires $\ideal N_i$, et soient respectivement $\ideal P_i$ les idéaux premiers associés. De
\[
 \ideal L_i\cdot\ideal N_i\equiv0(\ideal Q);\qquad
 \ideal L_i\not\equiv0(\ideal Q)
\]
(à cause de la représentation la plus courte), il vient alors
\[
 \ideal N_i^{\sigma_i}\equiv0(\ideal Q);
 \qquad\text{ou}\qquad
 \ideal P_i^{\varrho_i\sigma_i}\equiv0(\ideal Q).
\]
Puisque $\ideal P_i$ est en même temps un diviseur de $\ideal Q$, $\ideal P_i$ coïncide donc, pour chaque $i$, avec l'idéal premier associé $\ideal P$ de $\ideal Q$.

Il reste à montrer que, dans toute représentation réduite $\ideal Q=[\ideal N_1\cdots\ideal N_\lambda]$, les $\ideal N_i$ sont primaires.\footnote{Que la représentation réduite soit ici essentielle, c'est ce que montre l'exemple de la représentation non réduite:
\[
 \ideal Q=[x^2,xy,y^\lambda]=[(x^2,xy,y^2,yz),(x,y^\lambda)],
 \qquad \lambda\ge2.
\]
Ici $(x^2,xy,y^2,yz)=[(x^2,y),(x,y^2,z)]$ n'est pas primaire d'après ce qui vient d'être prouvé; car cette dernière représentation est la plus courte par idéaux primaires, mais les idéaux premiers associés $(x,y)$ et $(x,y,z)$ sont différents. ($\ideal Q$ est primaire d'après la note 19.)} Pour cela, on résout les $\ideal N_i$ en leurs idéaux irréductibles $\ideal B$; dans la représentation la plus courte qui en résulte,
\[
 \ideal Q=[\ideal B_1\cdots\ideal B_\mu],
\]
chaque idéal est alors primaire et possède, d'après ce qui vient d'être prouvé, $\ideal P$ pour idéal premier associé. Mais alors, d'après la première partie du théorème démontrée plus haut, il en va de même pour chaque $\ideal N_i$, ce qui achève la preuve du théorème VIII.

\textbf{Addition.} Il en résulte encore qu'un idéal premier est nécessairement irréductible. Car, d'après le théorème VIII, la représentation réduite $\ideal P=[\ideal N,\ideal L]$ donne $\ideal N\equiv0(\ideal P)$; donc, puisque $\ideal P\equiv0(\ideal N)$, aussi $\ideal P=\ideal N$, et de même $\ideal P=\ideal L$. L'irréductibilité de $\ideal P$ s'obtient aussi directement: de $\ideal P=[\ideal N,\ideal L]$, il résulte en effet que $\ideal N\cdot\ideal L\equiv0(\ideal P)$, $\ideal N\not\equiv(\ideal P)$, $\ideal L\not\equiv0(\ideal P)$, en contradiction avec la propriété définissant l'idéal premier.

Soient maintenant
\[
 \ideal M=[\ideal Q_1\cdots\ideal Q_\alpha]=[\bar{\ideal Q}_1\cdots\bar{\ideal Q}_\beta]
\]
deux représentations réduites de $\ideal M$ comme plus petit commun multiple de plus grands idéaux primaires. En résolvant les $\ideal Q$ en leurs idéaux irréductibles $\ideal B$, et les $\bar{\ideal Q}$ de même en idéaux irréductibles $\ideal D$, on obtient deux représentations réduites de $\ideal M$ comme plus petit commun multiple d'idéaux irréductibles, dans lesquelles, d'après le théorème VII, le nombre des composantes aussi bien que les idéaux premiers associés concordent. D'après le théorème VIII, tous les idéaux irréductibles $\ideal B$ qui apparaissent dans un $\ideal Q_i$ fixé possèdent le même idéal premier associé $\ideal P_i$; tandis que le $\ideal P_k$ correspondant à $\ideal Q_k$ en est nécessairement distinct, car autrement, d'après le théorème VIII, on n'aurait pas une représentation par de plus grands idéaux primaires. Le nombre $\alpha$ des $\ideal Q$ est donc égal au nombre des différents idéaux premiers associés $\ideal P$ des $\ideal B$; ces différents $\ideal P$ forment les idéaux premiers associés des $\ideal Q$. Il en va de même des $\bar{\ideal Q}$ par rapport à leur résolution en les $\ideal D$. Il résulte donc du théorème VII que le nombre des $\ideal Q$ égale celui des $\bar{\ideal Q}$ et que leurs idéaux premiers associés concordent. En même temps, il apparaît que le groupement, indiqué au début du paragraphe, des idéaux irréductibles en de plus grands primaires consiste à réunir tous ceux, et seulement ceux, qui ont le même idéal premier associé. Le théorème VIII montre encore la propriété d'irréductibilité des plus grands idéaux primaires: ils n'admettent aucune représentation réduite comme plus petit commun multiple de plus grands primaires.

En résumé, on a prouvé:

\textbf{Théorème IX.} \emph{Dans deux représentations réduites d'un idéal comme plus petit commun multiple de plus grands idéaux primaires, le nombre des composantes et les idéaux premiers associés, qui sont tous distincts les uns des autres, concordent. Autrement dit: à chaque $\ideal Q$ on peut faire correspondre de manière biunivoque un $\bar{\ideal Q}$ de telle sorte qu'une puissance de $\ideal Q$ soit divisible par $\ideal Q'$, et réciproquement.}\footnote{Un exemple de représentations différentes est celui donné dans la note 12 au théorème II:
\[
 (x^2,xy)=[(x),(x^2,\mu x+y)]
\]
pour $\mu$ arbitraire. Comme les idéaux premiers associés $\ideal P_1=(x)$, $\ideal P_2=(x,y)$ sont distincts, il s'agit de plus grands idéaux primaires. --- Pour l'unicité des plus grands idéaux primaires \og{} isolés \fg{}, voir \S~7.} \emph{--- Les $\ideal Q$ et les $\ideal Q'$ possèdent la propriété d'irréductibilité relativement à la décomposition en plus grands idéaux primaires.}

\textbf{Addition.} Remarquons que le théorème IX subsiste pour l'essentiel si, au lieu d'une représentation réduite, on suppose seulement une représentation la plus courte. Si alors, par exemple,
\[
 \ideal M=[\ideal Q_1\cdots\ideal Q_i^*\cdots\ideal Q_\alpha]
\]
est réduite relativement à $\ideal Q_i^*$, et si $\ideal Q_i^*$ est un diviseur propre de $\ideal Q_i$, $\ideal L_i$ étant le complément de $\ideal Q_i$, il vient, d'après le lemme IV,
\[
 \ideal Q_i=[\ideal Q_i^*,(\ideal L_i,\ideal Q_i)],
\]
et cette représentation est réduite relativement à $\ideal Q_i^*$. D'après le théorème VIII, dans l'application duquel il faut au besoin remplacer $(\ideal L_i,\ideal Q_i)$ par un diviseur propre, $\ideal Q_i^*$ est donc primaire et possède le même idéal premier associé $\ideal P_i$ que $\ideal Q_i$. La continuation du procédé montre qu'à toute telle représentation on peut associer une représentation réduite par de plus grands idéaux primaires, de telle sorte que le nombre des composantes et les idéaux premiers associés concordent. Il vaut donc aussi pour la représentation la plus courte que, dans deux représentations différentes, le nombre des composantes et les idéaux premiers associés concordent.

Les idéaux premiers associés ainsi définis de manière univoque et distincts les uns des autres seront brièvement appelés \og{} les idéaux premiers associés de $\ideal M$ \fg{}.
\subsection*{\S~6. Représentation unique d'un idéal comme plus petit commun multiple d'idéaux irréductibles relativement premiers.}

\textbf{Définition V.} \emph{Un idéal $\ideal R$ est dit relativement premier à $\ideal S$ si de
\[
 \ideal T\cdot\ideal R\equiv0(\ideal S)
\]
il s'ensuit nécessairement que $\ideal T\equiv0(\ideal S)$. Si $\ideal R$ est relativement premier à $\ideal S$, et $\ideal S$ relativement premier à $\ideal R$, alors $\ideal R$ et $\ideal S$ sont dits mutuellement relativement premiers.}\footnote{La condition d'être relativement premier n'est pas symétrique. Par exemple $\ideal R=(x^2,y)$ est relativement premier à $\ideal S=(x)$; mais $\ideal S$ n'est pas relativement premier à $\ideal R$, puisque $\ideal S^2\equiv0(\ideal R)$, tandis que $\ideal T=\ideal S\not\equiv0(\ideal R)$.}

\emph{Un idéal est dit irréductible relativement premier s'il ne peut pas être représenté comme plus petit commun multiple de diviseurs propres mutuellement relativement premiers.}

Si, en particulier, à la place de $\ideal T$, on prend le plus grand commun diviseur $\ideal T_0$ de tous les $\ideal T$ pour lesquels $\ideal T\cdot\ideal R\equiv0(\ideal S)$, alors on a aussi $\ideal T_0\cdot\ideal R\equiv0(\ideal S)$ et $\ideal S\equiv0(\ideal T_0)$. Ainsi $\ideal T_0=\ideal S$ lorsque $\ideal R$ est relativement premier à $\ideal S$; et $\ideal T_0$ est un diviseur propre de $\ideal S$ lorsque $\ideal R$ n'est pas relativement premier à $\ideal S$.\footnote{Le $\ideal T_0$ ainsi défini coïncide avec le \og{} module résiduel \fg{} de Lasker; et, dans son extension aux modules de nombres au lieu des idéaux, avec le \og{} quotient \fg{} de deux modules chez Dedekind. Lasker, loc. cit., p. 49; Dedekind (\emph{Zahlentheorie}), p. 504.}

La preuve d'unicité repose sur le théorème suivant.

\textbf{Théorème X.} \emph{1. Si $\ideal R$ est relativement premier aux idéaux $\ideal S_1\ldots\ideal S_\lambda$, alors $\ideal R$ est aussi relativement premier à leur plus petit commun multiple $\ideal S$.}

\emph{2. Si les idéaux $\ideal S_1,\ldots,\ideal S_\lambda$ sont relativement premiers à $\ideal R$, alors leur plus petit commun multiple $\ideal S$ est aussi relativement premier à $\ideal R$.}

\emph{3. Si $\ideal R$ est relativement premier à $\ideal S$ et si $\ideal S=[\ideal S_1\cdots\ideal S_\lambda]$ est une représentation réduite de $\ideal S$, alors $\ideal R$ est aussi relativement premier à chacun des $\ideal S_i$.}

\emph{4. Si $\ideal S$ est relativement premier à $\ideal R$, alors tout diviseur $\ideal S_i$ de $\ideal S$ est aussi relativement premier à $\ideal R$.}

1. Comme $\ideal T\cdot\ideal R\equiv0(\ideal S)$ entraîne nécessairement $\ideal T\cdot\ideal R\equiv0(\ideal S_i)$, l'hypothèse donne $\ideal T\equiv0(\ideal S_i)$, et donc $\ideal T\equiv0(\ideal S)$.

2. Posons
\[
 \ideal C_1=[\ideal S_2\cdots\ideal S_\lambda]
 \ideal C_{12}=[\ideal S_3\cdots\ideal S_\lambda];\ \ldots
 \ideal C_{12\ldots\lambda-1}=\ideal S_\lambda.
\]
De $\ideal T\cdot\ideal S\equiv0(\ideal R)$, on déduit alors $\ideal T\cdot\ideal C_1\cdot\ideal S_1\equiv0(\ideal R)$; donc, par l'hypothèse, $\ideal T\cdot\ideal C_1\equiv0(\ideal R)$. On obtient ensuite $\ideal T\cdot\ideal C_{12}\cdot\ideal S_2\equiv0(\ideal R)$, donc $\ideal T\cdot\ideal C_{12}\equiv0(\ideal R)$, et finalement $\ideal T\cdot\ideal C_{12\ldots\lambda-1}=\ideal T\cdot\ideal S_\lambda\equiv0(\ideal R)$; donc $\ideal T\equiv0(\ideal R)$.

3. Soit $\widehat{\ideal S}_i$ le complément de $\ideal S_i$; de $\ideal T_0\cdot\ideal R\equiv0(\ideal S_i)$, et aussi de $\ideal C_i\cdot\ideal R\equiv0(\ideal C_i)$, on déduit que $[\ideal T_0,\ideal C_i]\cdot\ideal R\equiv0(\ideal S)$. Si $\ideal T_0$ était un diviseur propre de $\ideal S_i$, alors, puisque la représentation $\ideal S=[\ideal S_i,\widehat{\ideal S}_i]$ est réduite, $[\ideal T_0,\widehat{\ideal S}_i]$ serait aussi un diviseur propre de $\ideal S$, contrairement à l'hypothèse.

4. De $\ideal T_0\ideal S_i\equiv0(\ideal R)$, où $\ideal T_0$ est un diviseur propre de $\ideal R$, on déduit aussi $\ideal T_0\cdot\ideal S\equiv0(\ideal R)$; ainsi $\ideal T_0$ serait un diviseur propre de $\ideal R$, contrairement à l'hypothèse.

La définition V montre en particulier que tout idéal est relativement premier à l'idéal unité $\ideal O$ formé de tous les éléments de $\Sigma$;\footnote{$\ideal O$ joue le rôle de l'unité seulement relativement à la divisibilité et au plus petit commun multiple, non relativement à la formation de produits. On a par exemple $\ideal O=(x)$ pour le domaine de tous les polynômes entiers en $x$ sans terme constant; $\ideal O=(2)$ pour le domaine de tous les nombres pairs.} cependant les théorèmes suivants ne valent que pour des idéaux différents de $\ideal O$.

Du théorème X que l'on vient de démontrer découle d'abord:

\textbf{Lemme V.} \emph{Toute représentation d'un idéal comme plus petit commun multiple d'idéaux mutuellement relativement premiers, distincts de $\ideal O$, est réduite.}

Soit
\[
 \ideal M=[\ideal R_1\ideal R_2\cdots\ideal R_\sigma]=[\ideal R_i,\ideal L_i]
\]
une telle représentation. D'après le théorème X, 2, $\ideal L_i$ est alors lui aussi relativement premier à $\ideal R_i$; $\ideal R_i$ ne peut donc pas diviser $\ideal L_i$, et la représentation devient la plus courte. En effet, de $\ideal L_i\equiv0(\ideal R_i)$, où $\ideal L_i$ est relativement premier à $\ideal R_i$, il suivrait, puisque $\ideal O\cdot\ideal L_i\equiv0(\ideal R_i)$, que $\ideal O\equiv0(\ideal R_i)$; donc $\ideal R_i=\ideal O$, ce qui est exclu par hypothèse.

Supposons maintenant que $\ideal R_i$ puisse être remplacé par le diviseur propre $\ideal R_i^*$. Alors, d'après le lemme II, $\ideal R_i$ est réductible et l'on obtient
\[
 \ideal R_i=[\ideal R_i^*,(\ideal R_i,\ideal L_i)].
\]
Si l'on remplace ici, au besoin, $(\ideal R_i,\ideal L_i)$ par un diviseur propre $(\ideal R_i,\ideal L_i)^*$, et de même $\ideal R_i^*$ par un diviseur propre, de sorte qu'il en résulte une représentation réduite de $\ideal R_i$, le théorème X, 3, montre que $\ideal L_i$ devient aussi relativement premier à $(\ideal R_i,\ideal L_i)^*$; et celui-ci est différent de $\ideal O$ à cause de la représentation la plus courte. Mais comme $(\ideal R_i,\ideal L_i)^*$ divise $(\ideal R_i,\ideal L_i)$, et $(\ideal R_i,\ideal L_i)$ divise $\ideal L_i$, on a ainsi obtenu une contradiction.

Du théorème X résulte en outre le théorème qui sert d'intermédiaire avec les idéaux premiers associés.

\textbf{Théorème XI.} \emph{Si $\ideal R$ est relativement premier à $\ideal S$ et si $\ideal S$ est différent de $\ideal O$, alors aucun idéal premier associé de $\ideal R$\footnote{Définition des idéaux premiers associés de $\ideal R$ à la fin de \S~5.} n'est divisible par un idéal premier associé de $\ideal S$. Réciproquement, s'il n'existe aucune telle divisibilité, alors $\ideal R$ est relativement premier à $\ideal S$, et bien entendu $\ideal S$ est différent de $\ideal O$.}

Soient
\[
 \ideal R=[\ideal Q_1\cdots\ideal Q_\alpha],\qquad
 \ideal S=[\ideal Q_1^*\cdots\ideal Q_\beta^*]
\]
des représentations réduites de $\ideal R$ et $\ideal S$ par de plus grands idéaux primaires; soient $\ideal P_1,\ldots,\ideal P_\alpha$ et $\ideal P_1^*,\ldots,\ideal P_\beta^*$ les idéaux premiers associés. Nous montrons l'assertion sous la forme suivante: si un $\ideal P$ est divisible par un $\ideal P^*$, alors $\ideal R$ ne peut pas être relativement premier à $\ideal S$, et réciproquement.

Soit donc $\ideal P_\mu\equiv0(\ideal P_\nu^*)$, et par suite aussi $\ideal Q_\mu\equiv0(\ideal P_\nu^*)$; de là, d'après la définition de $\ideal P_\nu^*$, on déduit que $\ideal Q_\mu^{\sigma_\nu}\equiv0(\ideal Q_\nu^*)$, donc aussi $\ideal R^{\sigma_\nu}\equiv0(\ideal Q_\nu^*)$. Soit maintenant $\ideal R^\tau$ la plus basse puissance de $\ideal R$ qui soit divisible par $\ideal Q_\nu^*$. Pour $\tau=1$, $\ideal R$ est divisible par $\ideal Q_\nu^*$; mais comme, par hypothèse, $\ideal S$ est différent de $\ideal O$, $\ideal R$ n'est donc pas relativement premier à $\ideal Q_\nu^*$. Pour $\tau\ge2$, on a $\ideal R^{\tau-1}\cdot\ideal R\equiv0(\ideal Q_\nu^*)$, $\ideal R^{\tau-1}\not\equiv0(\ideal Q_\nu^*)$, donc $\ideal R$ n'est pas relativement premier à $\ideal Q_\nu^*$;\footnote{$\ideal R^0$ n'est pas défini, puisque $\Sigma$ n'a pas nécessairement d'unité; c'est pourquoi le cas $\tau=1$ a dû être traité d'abord. Le cas $\tau=0$ est aussi exclu, lorsque $\Sigma$ possède une unité, par l'hypothèse faite sur $\ideal S$.} et par conséquent, dans les deux cas, d'après le théorème X, 3, $\ideal R$ n'est pas non plus relativement premier à $\ideal S$.

Réciproquement, si $\ideal R$ n'est pas relativement premier à $\ideal S$, alors, d'après le théorème X, 1, $\ideal R$ n'est pas relativement premier à au moins un $\ideal Q_\nu^*$. On a donc
\[
 \ideal T_0\cdot\ideal R\equiv0(\ideal Q_\nu^*)
 \ideal T_0\not\equiv0(\ideal Q_\nu^*),
\]
et, puisque $\ideal Q_\nu^*$ est primaire,
\[
 \ideal R^\tau\equiv0(\ideal Q_\nu^*),\qquad
 \text{donc aussi}\quad
 \ideal Q_1^{\tau}\cdots\ideal Q_\alpha^{\tau}\equiv0(\ideal Q_\nu^*).
\]
Mais il en résulte, pour les idéaux premiers associés,
\[
 \ideal P_1^{\tau\varrho_1}\cdots\ideal P_\alpha^{\tau\varrho_\alpha}\equiv0(\ideal P_\nu^*),
\]
et, d'après la propriété des idéaux premiers, $\ideal P_\nu^*$ divise donc au moins un $\ideal P$. Ainsi le théorème XI est démontré.

Les théorèmes X et XI donnent maintenant l'existence\footnote{L'existence de la décomposition se démontre aussi directement, en analogie exacte avec l'existence, établie en \S~2, de la décomposition en un nombre fini d'idéaux irréductibles.} et l'unicité de la décomposition en idéaux irréductibles relativement premiers, comme suit.

Soit $\ideal M=[\ideal Q_1\cdots\ideal Q_\alpha]$ une représentation réduite (ou du moins la plus courte) de $\ideal M$ par de plus grands idéaux primaires; soient $\ideal P_1,\ldots,\ideal P_\alpha$ les idéaux premiers associés. Nous regroupons les $\ideal P$ en groupes de telle sorte qu'aucun idéal d'un groupe ne soit divisible par un idéal d'un groupe différent, tandis que le groupe individuel ne puisse pas se décomposer en deux sous-groupes qui posséderaient tous deux cette propriété.

Pour construire une telle division en groupes, il faut remarquer que, par définition, chaque groupe $G$ doit contenir, avec chaque idéal $\ideal P$, tous ses diviseurs et multiples apparaissant parmi $\ideal P_1,\ldots,\ideal P_\alpha$ (c'est-à-dire ceux qui sont divisibles par $\ideal P$). Soient par exemple $\ideal P^{(i_1)}$ tous les multiples de $\ideal P$, $\ideal P_{j_1}^{(i_1)}$ tous les diviseurs de $\ideal P^{(i_1)}$, $\ideal P_{j_1}^{(i_1i_2)}$ tous les multiples de $\ideal P_{j_1}^{(i_1)}$, et ainsi de suite; en général, $\ideal P_{j_1\ldots j_{\lambda-1}}^{(i_1\ldots i_\lambda)}$ tous les multiples de $\ideal P_{j_1\ldots j_{\lambda-1}}^{(i_1\ldots i_{\lambda-1})}$, et $\ideal P_{j_1\ldots j_\lambda}^{(i_1\ldots i_\lambda)}$ tous les diviseurs de $\ideal P_{j_1\ldots j_{\lambda-1}}^{(i_1\ldots i_\lambda)}$. Comme il ne s'agit au total que d'un nombre fini d'idéaux $\ideal P$, le procédé doit s'arrêter après un nombre fini d'étapes, c'est-à-dire ne fournir aucun idéal différent de tous ceux qui précèdent. L'ensemble des idéaux $\ideal P$ ainsi obtenu forme effectivement un groupe $G$ ayant les propriétés voulues.

En effet, par définition, $G$ contient avec chaque $\ideal P_{j_1\ldots j_{\lambda-1}}^{(i_1\ldots i_{\lambda-1})}$ aussi tous les multiples $\ideal P_{j_1\ldots j_{\lambda-1}}^{(i_1\ldots i_\lambda)}$, et avec chaque $\ideal P_{j_1\ldots j_{\lambda-1}}^{(i_1\ldots i_\lambda)}$ aussi tous les diviseurs $\ideal P_{j_1\ldots j_\lambda}^{(i_1\ldots i_\lambda)}$. Parmi ceux-ci sont cependant aussi contenus tous les diviseurs des $\ideal P_{j_1\ldots j_{\lambda-1}}^{(i_1\ldots i_{\lambda-1})}$, tandis que les multiples des $\ideal P_{j_1\ldots j_{\lambda-1}}^{(i_1\ldots i_\lambda)}$ sont eux-mêmes de nouveau des $\ideal P_{j_1\ldots j_{\lambda-1}}^{(i_1\ldots i_\lambda)}$. Les idéaux qui ne sont pas contenus dans $G$ ne peuvent donc être ni diviseurs ni multiples de ceux qui sont contenus dans $G$.

Mais $G$ satisfait aussi à la condition d'irréductibilité. S'il existe en effet une division en deux sous-groupes $G^{(1)}$ et $G^{(2)}$, et si $G^{(1)}$ contient par exemple $\ideal P_{j_1\ldots j_\lambda}^{(i_1\ldots i_\lambda)}$ (respectivement $\ideal P_{j_1\ldots j_{\lambda-1}}^{(i_1\ldots i_\lambda)}$), alors il contient aussi tous ceux qui précèdent, car ceux-ci sont alternativement multiples et diviseurs (respectivement diviseurs et multiples); donc il contient aussi $\ideal P$ et ainsi tout le groupe $G$. En procédant de même avec les idéaux qui ne sont pas contenus dans $G$, on obtient ainsi une division en groupes $G_1,\ldots,G_\sigma$ de tous les $\ideal P$, ayant les propriétés voulues. Une telle division en groupes est unique; car si $G'_1,\ldots,G'_\tau$ est une seconde division et si $\ideal P_{j_1\ldots j_\lambda}^{(i_1\ldots i_\lambda)}$ (respectivement $\ideal P_{j_1\ldots j_{\lambda-1}}^{(i_1\ldots i_\lambda)}$) est un élément de $G'_i$, alors, d'après ce qui précède, $G'_i$ contient tout le groupe $G$ et, à cause de la propriété d'irréductibilité, il est donc identique à $G$.

Si maintenant $\ideal P_{i\mu}$ (où $\mu$ parcourt les valeurs de $1$ à $\lambda_i$) désignent les idéaux $\ideal P$ réunis dans un groupe $G_i$, alors, puisque les $\ideal P$ sont tous distincts, les idéaux primaires $\ideal Q_{i\mu}$ correspondant à une représentation la plus courte sont déterminés de façon unique. Posons
\[
 \ideal R_i=[\ideal Q_{i1}\cdots\ideal Q_{i\lambda_i}],
 \qquad
 \text{on obtient alors}\quad
 \ideal M=[\ideal R_1\cdots\ideal R_\sigma].
\]
Nous montrons que l'on a obtenu ainsi une décomposition de $\ideal M$ en idéaux irréductibles relativement premiers. Il faut d'abord remarquer que le théorème XI reste applicable même lorsque $\ideal R_i=[\ideal Q_{i_1}\cdots\ideal Q_{i\lambda_i}]$ est seulement une représentation la plus courte, car, d'après l'addition au théorème IX, les idéaux premiers associés sont déjà définis de manière univoque par là. Comme aucun idéal premier associé $\ideal P_{i\mu}$ de $\ideal R_i$ n'est divisible par un idéal premier associé $\ideal P_{j\nu}$ de $\ideal R_j$, ni inversement, les théorèmes XI montrent que $\ideal R_i$ et $\ideal R_j$ sont mutuellement relativement premiers; et chaque $\ideal R_i$ est, d'après le théorème XI, irréductible relativement premier grâce à la propriété d'irréductibilité du groupe $G_i$, car en cas de réductibilité seuls des idéaux différents de $\ideal O$ entrent en considération. D'après le lemme V il s'agit en outre d'une représentation réduite, même si l'on était parti à l'origine seulement d'une représentation la plus courte par les $\ideal Q$. Réciproquement, d'après le théorème XI, toute décomposition en idéaux irréductibles relativement premiers conduit à la division en groupes des $\ideal P_{i\mu}$ indiquée ci-dessus.

Soit maintenant
\[
 \ideal M=[\bar{\ideal R}_1\cdots\bar{\ideal R}_\tau]
\]
une seconde représentation de $\ideal M$ par idéaux irréductibles relativement premiers; elle est réduite d'après le lemme V. Comme alors, ainsi que le montre la résolution des $\bar{\ideal R}$ en plus grands idéaux primaires, les idéaux premiers associés concordent, les divisions en groupes de ces idéaux premiers, dont l'unicité vient d'être prouvée, concordent également. On a donc $\tau=\sigma$; et la notation peut être choisie de sorte que $\ideal R_i$ et $\bar{\ideal R}_i$ appartiennent au même groupe. Si alors
\[
 \ideal M=[\ideal R_i,\ideal L_i]=[\bar{\ideal R}_i,\bar{\ideal L}_i]
\]
est la représentation par l'idéal et son complément, alors, puisque $\bar{\ideal R}_i$ est associé au même groupe que $\ideal R_i$, d'après le théorème XI $\ideal L_i$ est aussi relativement premier à $\bar{\ideal R}_i$, et $\bar{\ideal L}_i$ relativement premier à $\ideal R_i$. À cause de
\[
 \ideal R_i\cdot\ideal L_i\equiv0(\bar{\ideal R}_i),
 \qquad
 \bar{\ideal R}_i\cdot\bar{\ideal L}_i\equiv0(\ideal R_i)
\]
il vient donc
\[
 \ideal R_i\equiv0(\bar{\ideal R}_i),
 \qquad
 \bar{\ideal R}_i\equiv0(\ideal R_i),
 \qquad
 \ideal R_i=\bar{\ideal R}_i.
\]
Ainsi est démontré:

\textbf{Théorème XII.} \emph{Tout idéal se représente de manière unique comme plus petit commun multiple d'un nombre fini d'idéaux mutuellement relativement premiers et irréductibles relativement premiers.}

\setcounter{footnote}{29}
\subsection*{\S~7. Unicité des idéaux isolés.}

\textbf{Définition VI.} \emph{Si la représentation la plus courte $\ideal M=[\ideal R,\ideal L]$ est réduite par rapport à $\ideal L$, on dit que $\ideal R$ est un idéal isolé lorsqu'aucun idéal premier associé de $\ideal R$ ne divise un idéal premier associé de $\ideal L$, autrement dit lorsque $\ideal L$ est relativement premier à $\ideal R$.}

D'après cela, la représentation $\ideal M=[\ideal R,\ideal L]$ satisfait aux conditions de la représentation $\ideal M=[\ideal R_i,\ideal L_i]$ dans le lemme V, et le lemme V démontre:

\textbf{Lemme VI.} \emph{Si $\ideal R$ est un idéal isolé de la représentation la plus courte $\ideal M=[\ideal R,\ideal L]$, réduite par rapport à $\ideal L$, alors cette représentation est aussi réduite par rapport à $\ideal R$.}

Comme il s'agit donc, pour les idéaux isolés, d'une représentation réduite, les idéaux premiers associés qui apparaissent dans la décomposition de $\ideal R$ et de $\ideal L$ en idéaux irréductibles se complètent pour donner les idéaux premiers associés, déterminés de manière unique, qui apparaissent dans la décomposition correspondante de $\ideal M$. Ainsi aucun idéal premier associé de $\ideal R$, appartenant à la décomposition en idéaux irréductibles, ne divise l'un des autres idéaux premiers associés appartenant à la décomposition correspondante de $\ideal M$. Réciproquement, si cette condition est remplie, et si $\ideal R$ apparaît dans au moins une représentation $\ideal M=[\ideal R,\ideal L]$ réduite par rapport à $\ideal R$, donc aussi dans une représentation réduite $\ideal M=[\ideal R,\ideal L^*]$, alors $\ideal R$ est isolé au sens de la définition VI. Il en résulte la définition suivante, indépendante du complément spécial $\ideal L$:

\textbf{Définition VIa.} \emph{$\ideal R$ est dit idéal isolé lorsque les idéaux premiers de $\ideal R$ associés à la décomposition en idéaux irréductibles ne divisent aucun des autres idéaux premiers associés à la décomposition correspondante de $\ideal M$, et lorsque $\ideal R$ apparaît dans au moins une représentation $\ideal M=[\ideal R,\ideal L]$ réduite par rapport à $\ideal R$.}\footnote{Si l'on introduisait les idéaux premiers associés ordinaires (fin de \S~5), il faudrait ajouter, comme condition spéciale, que ceux de $\ideal L$ soient tous distincts de ceux de $\ideal R$. D'après la définition VIa, il n'est donc plus nécessaire de supposer la représentation réduite par rapport au complément.}

Soient maintenant
\[
 \ideal M=[\ideal R,\ideal L]=[\bar{\ideal R},\bar{\ideal L}]
\]
deux représentations de $\ideal M$ par des idéaux isolés $\ideal R$ et $\bar{\ideal R}$ et par des compléments $\ideal L$ et $\bar{\ideal L}$, telles que les idéaux premiers associés de $\ideal R$ et de $\bar{\ideal R}$ concordent. Si l'on remplace $\ideal L$ et $\bar{\ideal L}$ par des diviseurs $\ideal L^*$ et $\bar{\ideal L}^{*}$ tels que les représentations deviennent réduites, alors les idéaux premiers associés de $\ideal L^*$ et de $\bar{\ideal L}^{*}$ concordent eux aussi; d'après le théorème XI, $\ideal L^*$ est donc relativement premier à $\bar{\ideal R}$, et $\bar{\ideal L}^{*}$ relativement premier à $\ideal R$. À cause de
\[
 \ideal R\cdot\ideal L^*\equiv0(\bar{\ideal R}),
 \qquad
 \bar{\ideal R}\bar{\ideal L}^{*}\equiv0(\ideal R),
\]
il vient donc
\[
 \ideal R\equiv0(\bar{\ideal R}),
 \qquad
 \bar{\ideal R}\equiv0(\ideal R),
 \qquad
 \ideal R=\bar{\ideal R};
\]
les idéaux isolés sont ainsi déterminés de manière unique par leurs idéaux premiers associés. Cela donne en particulier un renforcement des théorèmes VII et IX sur la décomposition en idéaux irréductibles, respectivement en plus grands idéaux primaires; d'après l'addition au théorème IX, il suffit ici de supposer une représentation la plus courte. En résumé, on obtient:

\textbf{Théorème XIII.} \emph{Dans toute représentation la plus courte d'un idéal comme plus petit commun multiple d'idéaux irréductibles, respectivement de plus grands idéaux primaires, les idéaux isolés irréductibles, respectivement les plus grands idéaux primaires isolés, sont déterminés de manière unique; l'ambiguïté ne concerne que les idéaux irréductibles non isolés, respectivement les plus grands idéaux primaires non isolés.}\footnote{Pour les idéaux dans les anneaux de polynômes, ce théorème a déjà été communiqué sans démonstration par Macaulay dans le cas de la décomposition en plus grands idéaux primaires; sa définition des idéaux primaires isolés et non isolés (imbedded) peut être considérée comme une formulation irrationnelle de celle donnée ci-dessous.} \emph{En général, les idéaux isolés sont déterminés de manière unique par leurs idéaux premiers associés.}

Si, dans une telle représentation la plus courte par idéaux irréductibles, respectivement par plus grands idéaux primaires, les idéaux $\ideal B_i$, respectivement $\ideal Q_j$, sont non isolés, alors, par définition, leurs compléments $\ideal A_i$, respectivement $\ideal L_j$, sont divisibles par $\ideal P_i$, respectivement $\ideal P_j$. Il vient donc
\[
 \ideal A_i^{\varrho_i}\equiv0(\ideal B_i)
 \qquad \text{respectivement} \qquad
 \ideal L_j^{\sigma_j}\equiv0(\ideal Q_j).
\]
Réciproquement, si ces relations sont satisfaites, alors $\ideal P_i$, respectivement $\ideal P_j$, divise au moins un idéal premier associé du complément, donc aussi, d'après l'addition au théorème IX, un idéal premier associé du diviseur $\ideal L_j^*$ de $\ideal L_j$ qui donne une représentation réduite par rapport à $\ideal L_j^*$. Ainsi $\ideal B_i$, respectivement $\ideal Q_j$, est non isolé. En particulier, les idéaux irréductibles $\ideal B_i$ dont l'idéal premier associé apparaît plus d'une fois dans la décomposition de $\ideal M$ sont toujours non isolés. \emph{Les idéaux primaires non isolés sont donc aussi caractérisés par le fait qu'une puissance de tout complément devient divisible par eux; les idéaux primaires isolés, par le fait que cela ne peut pas se produire.}

\setcounter{footnote}{31}
\subsection*{\S~8. Représentation unique d'un idéal comme produit d'idéaux étrangers-irréductibles.}

Si l'anneau $\Sigma$ considéré possède une unité, c'est-à-dire un élément $\varepsilon$ tel que $\varepsilon\cdot a=a$ pour tout élément de $\Sigma$,\footnote{On sait que, du fait de la commutativité de la multiplication, $\Sigma$ ne peut posséder plus d'une unité; car, pour deux unités quelconques $\varepsilon_1$ et $\varepsilon_2$, on a $\varepsilon_1\varepsilon_2=\varepsilon_2=\varepsilon_1$.} les idéaux étrangers se définissent ainsi:

\textbf{Définition VIII.} \emph{Deux idéaux $\ideal R$ et $\ideal S$ sont dits étrangers lorsque leur plus grand commun diviseur est égal à l'idéal unité $\ideal O=(\varepsilon)$ formé de tous les éléments de $\Sigma$. Un idéal est dit étranger-irréductible lorsqu'il ne peut pas être représenté comme plus petit commun multiple d'idéaux deux à deux étrangers.}

Remarquons que deux idéaux étrangers sont toujours mutuellement relativement premiers. En effet, par définition, il existe deux éléments $r\equiv0(\ideal R)$, $s\equiv0(\ideal S)$ tels que $\varepsilon=r+s$. Or de $\ideal T\cdot\ideal R\equiv0(\ideal S)$, on déduit que $\ideal T\cdot r\equiv0(\ideal S)$, donc que $\ideal T\cdot\varepsilon=\ideal T\equiv0(\ideal S)$; et, de même, de $\bar{\ideal T}\cdot\ideal S\equiv0(\ideal R)$, on déduit aussi que $\bar{\ideal T}\equiv0(\ideal R)$.\footnote{La réciproque, en revanche, n'est pas vraie; par exemple les idéaux $\ideal R=(x)$, $\ideal S=(y)$ sont mutuellement relativement premiers, mais ne sont pas étrangers.} D'après le lemme V, toute représentation par idéaux deux à deux étrangers est donc en même temps réduite.

La démonstration d'unicité repose sur le théorème correspondant au théorème X:

\textbf{Théorème XIV.} \emph{Si $\ideal R$ est étranger à chacun des idéaux $\ideal S_1,\ldots,\ideal S_\lambda$, alors $\ideal R$ est aussi étranger à $\ideal S=[\ideal S_1\cdots\ideal S_\lambda]$. Réciproquement, de ce que $\ideal R$ et $\ideal S$ sont étrangers, on déduit aussi que $\ideal R$ est étranger à chaque $\ideal S_j$. Si $\ideal R=[\ideal R_1\cdots\ideal R_\mu]$ et si chaque $\ideal R_i$ est étranger à chaque $\ideal S_j$, alors $\ideal R$ et $\ideal S$ sont étrangers; ici encore, la réciproque vaut.}

Si, en effet, $\ideal R$ est étranger à chaque $\ideal S_j$, il existe des éléments $s_j$ tels que
\[
 s_j\equiv0(\ideal S_j),
 \qquad
 s_j\equiv\varepsilon(\ideal R).
\]
Par conséquent,
\[
 s_1\cdot s_2\cdots s_\lambda\equiv0(\ideal S),
 \qquad
 s_1\cdot s_2\cdots s_\lambda\equiv\varepsilon(\ideal R),
 \qquad
 (\ideal R,\ideal S)=(\varepsilon).
\]
Mais comme $(\ideal R,\ideal S)$ est divisible par chacun des $(\ideal R,\ideal S_j)$, la réciproque vaut aussi. L'application répétée de ce raisonnement donne la seconde partie de l'assertion. En effet, si $\ideal R_i$ est, pour un $i$ fixé, étranger à $\ideal S_1,\ldots,\ideal S_\lambda$, alors $\ideal R_i$ est étranger à $\ideal S$. Si cela vaut pour tout $i$, alors, puisque la notion d'être étranger est réciproque, $\ideal S$ est étranger à $\ideal R$. Réciproquement, du fait que $\ideal R$ et $\ideal S$ sont étrangers, il s'ensuit que $\ideal S$ est étranger à $\ideal R_i$, et par conséquent que $\ideal R_i$ est étranger à $\ideal S$.

La démonstration de l'existence et de l'unicité\footnote{L'existence de la décomposition se démontre de nouveau directement par analogie avec \S~2; la démonstration d'unicité se conduit également de façon directe (voir ce qui est dit dans l'introduction à propos de Schmeidler et de Noether--Schmeidler). La démonstration donnée ici fournit en même temps un aperçu de la structure des idéaux étrangers-irréductibles.} de la décomposition en idéaux étrangers-irréductibles aboutit, comme la démonstration correspondante pour les idéaux relativement premiers, à une division unique en groupes. Mais, puisque les idéaux relativement premiers-irréductibles $\ideal R_1,\ldots,\ideal R_\sigma$ de $\ideal M$ sont définis de manière unique d'après le théorème XII, il est inutile ici de remonter aux idéaux premiers associés.

Nous réunissons les idéaux relativement premiers-irréductibles $\ideal R_1,\ldots,\ideal R_\sigma$ de $\ideal M$, définis de manière unique, en groupes de telle sorte que \emph{chaque idéal d'un groupe soit étranger à chaque idéal d'un autre groupe, tandis que le groupe pris isolément ne puisse pas se scinder en deux sous-groupes de telle sorte que chaque idéal d'un sous-groupe soit étranger à chaque idéal de l'autre sous-groupe}. Une telle division en groupes s'obtient comme suit. Par définition, le groupe individuel $G$ doit contenir, avec chaque idéal $\ideal R$, tous les idéaux qui ne sont pas étrangers à $\ideal R$. Soient ceux-ci donnés, par exemple, par $\ideal R_{i_1}$; soient ensuite $\ideal R_{i_1i_2}$ les idéaux qui ne sont pas étrangers à ces derniers; et, en général, soit $\ideal R_{i_1\ldots i_\lambda}$ non étranger à $\ideal R_{i_1\ldots i_{\lambda-1}}$. Comme il n'y a au total qu'un nombre fini d'idéaux, le procédé doit s'arrêter après un nombre fini d'étapes, c'est-à-dire ne plus fournir d'idéaux distincts de tous ceux qui précèdent; l'ensemble d'idéaux ainsi obtenu forme un groupe $G$ ayant les propriétés voulues. En effet, par la construction de $G$, tous les idéaux qui ne sont pas contenus dans $G$ sont étrangers à tous ceux qui y sont contenus. Supposons encore donnée une division en deux sous-groupes $G^{(1)}$ et $G^{(2)}$, et que, par exemple, $\ideal R_{i_1\ldots i_\lambda}$ soit un élément de $G^{(1)}$. Puisque la propriété d'être étranger est réciproque, $G^{(1)}$ doit alors aussi contenir $\ideal R_{i_1\ldots i_{\lambda-1}},\ldots,\ideal R_{i_1}$, donc aussi $\ideal R$, et par suite tout le groupe $G$; l'irréductibilité est ainsi démontrée. En procédant de même avec les groupes non contenus dans $G$, on obtient donc une division en groupes $G_1,\ldots,G_\tau$ de tous les $\ideal R$.

Cette division en groupes est unique. En effet, soit $G'_1,\ldots,G'_\tau$ une seconde division, et soit $\ideal R_{i_1\ldots i_\lambda}$ un élément de $G'_i$. Alors $G'_i$ contient aussi $\ideal R$, et par conséquent $G_1$, et il ne peut contenir, en raison de la condition d'irréductibilité, aucun élément distinct de ceux de $G_1$; ainsi $G'_i$ est égal à $G_1$.

Soit maintenant $\ideal T_i$ le plus petit commun multiple des idéaux $\ideal R$ réunis dans un groupe $G_i$. Nous montrons que
\[
 \ideal M=[\ideal T_1\cdots\ideal T_\tau]
\]
devient une représentation de $\ideal M$ par idéaux étrangers-irréductibles. D'abord, la résolution des $\ideal T$ dans les $\ideal R$ montre que $[\ideal T_1\cdots\ideal T_\tau]$ représente effectivement $\ideal M$. Ensuite, d'après le théorème XIV, les $\ideal T$ sont deux à deux étrangers, et chacun des $\ideal T$ est étranger-irréductible. L'existence d'une telle représentation est donc démontrée.

Pour démontrer l'unicité, soit $\ideal M=[\bar{\ideal T}_1\cdots\bar{\ideal T}_{\bar\tau}]$ une seconde représentation de ce type. Si l'on résout les $\bar{\ideal T}$ en leurs idéaux relativement premiers-irréductibles $\ideal R$, alors les $\ideal R$ qui apparaissent dans des $\bar{\ideal T}$ différents sont, d'après le théorème XIV, étrangers les uns aux autres, donc aussi mutuellement relativement premiers; ils coïncident donc avec les idéaux relativement premiers-irréductibles $\ideal R$ de $\ideal M$, définis de manière unique. Les $\bar{\ideal T}_i$ engendrent en outre, d'après le théorème XIV, une division en groupes $G'_i$ des $\ideal R$ avec les propriétés indiquées. Mais puisque cette division en groupes est unique et que chaque $\bar{\ideal T}_i$ est déterminé de manière unique par le groupe $G'_i=G_i$, on a $\bar{\ideal T}_i=\ideal T_i$; l'unicité est ainsi démontrée.

Pour des idéaux deux à deux étrangers, le plus petit commun multiple devient en outre égal au produit. En effet, d'après le théorème XIV, pour $\ideal M=[\ideal T_1\cdots\ideal T_\tau]$, le complément $\ideal L_i$ est aussi étranger à $\ideal T_i$; il existe donc des éléments
\[
 t_i\equiv0(\ideal T_i),
 \qquad
 l_i\equiv0(\ideal L_i),
 \qquad
 \varepsilon=t_i+l_i.
\]
Si $f\equiv0(\ideal T_i)$ et $f\equiv0(\ideal L_i)$, on a donc
\[
 f=f\varepsilon=ft_i+fl_i,
 \quad\text{et par conséquent}\quad
 f\equiv0(\ideal L_i\cdot\ideal T_i).
\]
Comme, réciproquement, $\ideal L_i\cdot\ideal T_i$ est divisible par $[\ideal L_i,\ideal T_i]$, on obtient $[\ideal L_i,\ideal T_i]=\ideal L_i\cdot\ideal T_i$; et, en poursuivant finalement le procédé sur $\ideal L_i$:
\[
 \ideal M=[\ideal T_1\cdots\ideal T_\tau]
       =\ideal T_1\cdot\ideal T_2\cdots\ideal T_\tau.
\]
Ceci démontre:

\textbf{Théorème XV.} \emph{Tout idéal se représente de manière unique comme produit d'un nombre fini d'idéaux deux à deux étrangers et étrangers-irréductibles.}

\setcounter{footnote}{34}
\subsection*{\S~9. Extension de l'étude aux modules. Égalité du nombre des composantes dans les décompositions en modules irréductibles.}

Nous montrons maintenant que le contenu des trois premiers paragraphes, qui se rapporte aux idéaux irréductibles et non aux idéaux primaires ni aux idéaux premiers, subsiste sous des hypothèses plus faibles. Ces paragraphes n'utilisent en effet pas la loi commutative de la multiplication et ne se rapportent qu'à la propriété qu'ont les idéaux d'être des modules; ils restent donc valables pour des modules relativement à des domaines non commutatifs, qu'il faut maintenant définir. La définition des modules doit avoir pour fondement un double domaine $(\Sigma,T)$ possédant les propriétés suivantes:

$\Sigma$ est un anneau non commutatif défini abstraitement, c'est-à-dire que $\Sigma$ est un système d'éléments $a,b,c,\ldots$ pour lesquels deux sortes de compositions sont définies, l'addition de l'anneau $(\#)$ et la multiplication de l'anneau $(\times)$, qui satisfont aux lois énoncées au \S~1, à l'exception de la loi commutative 4 de la multiplication de l'anneau.

$T$ est un système d'éléments $\alpha,\beta,\gamma,\ldots$ pour lequel, en liaison avec $\Sigma$, deux compositions sont également définies: l'addition, qui, de deux éléments $\alpha,\beta$, engendre de façon univoque un troisième élément $\alpha+\beta$; la multiplication d'un élément $\alpha$ de $T$ par un élément $c$ de $\Sigma$, qui engendre de façon univoque un élément $c\cdot\alpha$ de $T$.\footnote{Il s'agit donc ici d'une multiplication \og à droite \fg{}, d'un domaine $T$ \og à droite \fg{} et, par conséquent, de modules et d'idéaux \og à droite \fg{}. Si l'on prenait pour $T$ une multiplication à gauche $\alpha\cdot c$, on obtiendrait une théorie correspondante des modules et idéaux à gauche; $M$ contiendrait alors, avec $\alpha$, aussi $\alpha\cdot c$. La loi associative serait ici de la forme $(\gamma\cdot b)\cdot a=\gamma\cdot(b\times a)$.}

Pour ces compositions valent les lois suivantes:
\begin{enumerate}
\item La loi associative de l'addition: $(\alpha+\beta)+\gamma=\alpha+(\beta+\gamma)$.
\item La loi commutative de l'addition: $\alpha+\beta=\beta+\alpha$.
\item La loi de la soustraction illimitée et univoque: il existe dans $T$ un et un seul élément $\xi$ qui satisfait à l'équation $\alpha+\xi=\beta$ (on note $\xi=\beta-\alpha$).
\item La loi associative de la multiplication: $a\cdot(b\cdot\gamma)=(a\times b)\cdot\gamma$.
\item La loi distributive:
\[
 (a\# b)\cdot\gamma=a\cdot\gamma+b\cdot\gamma;\qquad c\cdot(\alpha+\beta)=c\cdot\alpha+c\cdot\beta.
\]
\end{enumerate}
De ces conditions résulte, comme on sait, l'existence de l'élément nul, ainsi que la validité de la loi distributive pour la combinaison de la soustraction et de la multiplication:
\[
 (a-b)\gamma=a\gamma-b\gamma,\qquad c(\alpha-\beta)=c\alpha-c\beta,
\]
où $(\dot{-})$ désigne la soustraction dans $\Sigma$. Si $\Sigma$ contient une unité $\varepsilon$, on doit avoir $\varepsilon\cdot\alpha=\alpha$ pour tout élément $\alpha$ de $T$.

Par module $M$ dans $(\Sigma,T)$, on entendra un système d'éléments de $T$ qui satisfait aux deux conditions:
\begin{enumerate}
\item $M$ contient avec $\alpha$ aussi $c\cdot\alpha$, où $c$ est un élément quelconque de $\Sigma$.
\item $M$ contient avec $\alpha$ et $\beta$ aussi la différence $\alpha-\beta$, donc avec $\alpha$ aussi $n\alpha$ pour tout entier $n$.\footnote{Les entiers doivent de nouveau être considérés comme des signes abréviatifs, et non comme des éléments de l'anneau.}
\end{enumerate}
D'après cette définition, $T$ lui-même forme un module dans $(\Sigma,T)$. Si, en particulier, le domaine $T$ et les compositions qui y sont fixées coïncident avec le domaine $\Sigma$ et les compositions qui y valent, alors le module $M$ devient un idéal à droite $M$ dans $\Sigma$. Si $\Sigma$ est en outre supposé commutatif, on obtient le concept ordinaire d'idéal, qui apparaît ainsi comme un cas particulier du concept de module.\footnote{L'exemple le plus simple d'un module est donné par le module des formes linéaires entières: $\Sigma$ se compose ici de tous les nombres rationnels entiers, $T$ de toutes les formes linéaires entières. Un module un peu plus général apparaît si l'on prend dans $\Sigma$ et $T$ des nombres algébriques entiers au lieu des nombres rationnels entiers, ou encore tous les nombres pairs. Si, à la place des formes linéaires, on considère chaque fois le complexe de tous les coefficients comme un seul élément, les compositions dans $\Sigma$ et dans $T$ sont effectivement différentes. Les idéaux dans les anneaux de polynômes non commutatifs font l'objet du travail commun de Noether--Schmeidler. Les idéaux dans d'autres domaines non commutatifs spéciaux sont traités dans les leçons de Hurwitz sur la théorie des nombres de quaternions (Berlin, Springer 1919) et dans les travaux de Du Pasquier qui y sont cités.}

Pour les modules, toutes les définitions du \S~1 demeurent valables. Ainsi $\alpha\eqzero{M}$, respectivement $N\eqzero{M}$, signifie que $\alpha$, respectivement chaque élément de $N$, est élément de $M$; autrement dit, $\alpha$, respectivement $N$, est divisible par $M$. $M$ est un diviseur propre de $N$ lorsque $M$ contient des éléments différents de ceux de $N$; de $N\eqzero{M}$ et $M\eqzero{N}$, on déduit que $M=N$. La définition du plus grand commun diviseur et celle du plus petit commun multiple restent également valables mot pour mot. Si, en particulier, le module $M$ contient un nombre fini d'éléments $\alpha_1\ldots\alpha_\varrho$ tels que
\[
 M=(\alpha_1,
 \ldots,\alpha_\varrho),
\]
c'est-à-dire si
\[
 \alpha=c_1\alpha_1+\cdots+c_\varrho\alpha_\varrho
       +n_1\alpha_1+\cdots+n_\varrho\alpha_\varrho
\]
pour tout $\alpha\eqzero{M}$, les $c_i$ étant des grandeurs de $\Sigma$ et les $n_i$ des entiers, alors $M$ s'appelle un module fini, et $\alpha_1\ldots\alpha_\varrho$ une base du module $M$.

Nous ne prendrons donc désormais pour fondement, de manière analogue au \S~1, que des domaines $(\Sigma,T)$ qui satisfont à la condition de finitude: tout module dans $(\Sigma,T)$ est fini, c'est-à-dire possède une base finie.

Alors, pour ce domaine $(\Sigma,T)$, le théorème I de la chaîne finie vaut aussi pour les modules, comme le montre la démonstration donnée là, et toutes les hypothèses des \S\S~2 et 3 sont ainsi remplies. La définition I et le lemme I sur les représentations les plus courtes et réduites se transposent directement; de même le théorème II sur la représentabilité de tout module comme plus petit commun multiple d'un nombre fini de modules irréductibles, le lemme II montrant que toute représentation la plus courte de ce type est en même temps réduite. Le théorème III subsiste encore: il exprime la réductibilité d'un module par les propriétés de son complément; il en résulte le lemme III, et finalement le théorème IV, qui affirme l'égalité du nombre des composantes dans deux représentations les plus courtes différentes d'un module comme plus petit commun multiple de modules irréductibles. Le théorème IV donne encore le lemme IV comme réciproque du lemme I sur la représentation réduite.

\emph{Addition.} Le même raisonnement montre que tous ces théorèmes et définitions subsistent lorsque, pour des domaines $T$ bilatères, c'est-à-dire à la fois à droite et à gauche, on entend partout par module un module bilatère, c'est-à-dire un module qui contient, avec $\alpha$, aussi $c\cdot\alpha$ et $\alpha\cdot c$, et avec $\alpha$ et $\beta$ aussi $\alpha-\beta$.

Tandis que tous ces théorèmes ne s'appuient que sur le concept de divisibilité et de plus petit commun multiple, les autres théorèmes d'unicité reposent essentiellement sur le concept de produit et n'admettent donc pas de transposition directe. Ainsi la définition de l'idéal primaire et de l'idéal premier ne peut pas être transposée aux modules, puisque le produit de deux grandeurs de $T$ n'est pas défini. La transposition formellement possible aux anneaux non commutatifs perd toutefois son sens, parce qu'ici on ne peut pas démontrer l'existence de l'idéal premier associé\footnote{En effet, les congruences $p_1^{\lambda_1}\eqzero{\ideal Q}$ et $p_2^{\lambda_2}\eqzero{\ideal Q}$ n'entraînent pas ici $(p_1-p_2)^{\lambda_1+\lambda_2}\eqzero{\ideal Q}$; de même, $(a\cdot b)^\lambda\eqzero{\ideal Q}$ n'entraîne pas $a^\lambda\cdot b^\lambda\eqzero{\ideal Q}$. Ainsi $\ideal P$ ne peut être établi ni comme idéal, ni avec la propriété des idéaux premiers. Pour les idéaux bilatères, $\ideal P$ peut certes être établi comme idéal, mais non comme idéal premier.} et parce que la démonstration qu'un idéal irréductible est primaire échoue également. Or ces deux faits forment la base des théorèmes d'unicité suivants. En revanche, lorsque l'anneau non commutatif possède une unité, on peut définir les idéaux étrangers et étrangers-irréductibles, et les raisonnements utilisés pour démontrer le théorème II montrent que tout idéal peut être représenté comme plus petit commun multiple d'un nombre fini d'idéaux deux à deux étrangers et étrangers-irréductibles.\footnote{Pour des idéaux spéciaux, \og complètement réductibles \fg{}, des théorèmes d'unicité peuvent être établis ici aussi; voir le travail commun de Noether--Schmeidler.}

Mentionnons enfin encore un critère suffisant pour que la condition de finitude soit remplie dans $(\Sigma,T)$: si $\Sigma$ contient une unité et satisfait à la condition de finitude, et si $T$ lui-même est un module fini dans $(\Sigma,T)$, alors tout module dans $(\Sigma,T)$ est fini.

De l'existence de l'unité dans $\Sigma$ il résulte en effet, pour
\[
 \ideal M=(f_1,
\ldots f_\varrho);\qquad
 f=\bar b_1f_1+\cdots+\bar b_\varrho f_\varrho+n_1f_1+\cdots+n_\varrho f_\varrho,
\]
où les $\bar b_i$ sont des éléments de $\Sigma$ et les $n_i$ des entiers, une représentation
\[
 f=b_1f_1+\cdots+b_\varrho f_\varrho,
\]
où les $b_i$ sont des éléments de $\Sigma$. Car, puisque $f_i=\varepsilon f_i$, on a $n_if_i=n_i\varepsilon f_i$, et, comme $n_i\varepsilon=(\varepsilon+\cdots+\varepsilon)$ appartient à $\Sigma$, $b_i=\bar b_i+n_i\varepsilon$ devient un élément de $\Sigma$. L'hypothèse sur $T$ signifie maintenant que tout élément $\alpha$ de $T$ admet une représentation
\[
 \alpha=\bar a_1\tau_1+\cdots+\bar a_k\tau_k+n_1\tau_1+\cdots+n_k\tau_k,
\]
et par conséquent
\[
 \alpha=a_1\tau_1+\cdots+a_k\tau_k,
\]
où la seconde représentation, à cause de $\varepsilon\tau_i=\tau_i$, s'obtient à partir de la première comme ci-dessus.

Lorsque les $\alpha$ parcourent maintenant les éléments d'un module $M$ de $(\Sigma,T)$, les coefficients $a_k$ de $\tau_k$ parcourent un idéal $\ideal M_k$ de $\Sigma$; d'après ce qui précède, pour tout $a_k\eqzero{\ideal M_k}$, on a donc aussi
\[
 a_k=b_1a_k^{(1)}+\cdots+b_\varrho a_k^{(\varrho)}.
\]
Si $\alpha^{(i)}$ désigne un élément de $M$ pour lequel le coefficient de $\tau_k$ est égal à $a_k^{(i)}$, alors $\alpha-b_1\alpha^{(1)}\ldots-b_\varrho\alpha^{(\varrho)}$ devient un élément appartenant à $M$ qui ne dépend plus que de $\tau_1,
\ldots,\tau_{k-1}$. On peut répéter le procédé sur la totalité de ces éléments qui ne contiennent plus $\tau_k$, et qui forment un module $M'$; ainsi, après un nombre fini de répétitions, l'assertion est démontrée.

\setcounter{footnote}{39}
\subsection*{\S~10. Cas particulier de l'anneau de polynômes.}

1. Supposons que l'anneau $\Sigma$ considéré se compose de tous les polynômes en $x_1,
\ldots,x_n$ à coefficients complexes arbitraires; pour ce domaine, d'après le théorème de la base de Hilbert (Ann. 36), la condition de finitude est remplie. Il s'agit du rapport de nos théorèmes avec les théorèmes connus de la théorie de l'élimination et de la théorie des modules.

Ce rapport est établi par le cas particulier suivant d'un théorème bien connu de Hilbert:\footnote{Sur les systèmes complets d'invariants. Math. Ann. 42 (1893), \S~3, p. 313.}

Si $f$ s'annule pour tout système fini de valeurs de $x_1,
\ldots,x_n$ qui est un zéro de tous les polynômes d'un idéal premier $\ideal P$ (un zéro de $\ideal P$), alors $f$ est divisible par $\ideal P$. En d'autres termes: un idéal premier $\ideal P$ se compose de la totalité des polynômes qui s'annulent en ses zéros.\footnote{Comme Lasker [Math. Ann. 60 (1905), p. 607] l'a montré, ce cas particulier peut aussi se démontrer directement dans le cas des formes homogènes; et, inversement, il en résulte de nouveau le théorème de Hilbert dans le cas homogène et inhomogène, que l'on peut énoncer ainsi: si un idéal $\ideal R$ s'annule en tous les zéros de $\ideal M$, alors une puissance de $\ideal R$ est divisible par $\ideal M$. Ce théorème, comme son cas particulier, ne vaut toutefois que lorsque le domaine des valeurs des $x$ est algébriquement clos; il ne peut donc pas résulter de nos théorèmes seuls, mais doit utiliser l'existence des racines, par exemple à l'endroit où un idéal dont les zéros ne consistent qu'en $x_1=0,
\ldots,x_n=0$ contient tous les produits de puissances des $x$ à partir d'une certaine dimension. Le reste de la démonstration se simplifie quelque peu, par rapport à Lasker, en utilisant nos théorèmes. Lasker doit en effet faire appel au théorème de Hilbert aussi pour établir la décomposition d'un idéal en plus grands idéaux primaires.}

Si donc un produit $fg$ s'annule pour tous les zéros de $\ideal P$, alors au moins l'un des facteurs s'annule; les zéros forment une variété algébrique irréductible. Si, inversement, on prend pour base cette définition de la variété irréductible, on voit que la totalité des polynômes qui s'annulent dans une variété irréductible forme un idéal premier; idéal premier et variété irréductible se correspondent donc de manière biunivoque. À cause de $\ideal Q\eqzero{\ideal P}$, $\ideal P^e\eqzero{\ideal Q}$, les zéros d'un idéal primaire coïncident en outre avec ceux de son idéal premier associé.\footnote{Cette propriété d'un idéal primaire, celle de posséder une variété irréductible, est prise par Macaulay (voir l'introduction) comme définition, tandis que Lasker ne fait entrer dans la définition que le concept de multiplicité d'une variété, et définit le reste abstraitement. Les idéaux primaires qui ne s'annulent que pour $x_1=0,
\ldots,x_n=0$ occupent chez Lasker une position spéciale.} La représentation d'un idéal comme plus petit commun multiple de plus grands idéaux primaires donne donc une résolution de tous les zéros de l'idéal en variétés irréductibles; et, comme Lasker l'a montré, la réciproque vaut également. La démonstration d'unicité des idéaux premiers associés correspond donc ici au théorème fondamental de la théorie de l'élimination sur la décomposabilité unique d'une variété algébrique en variétés irréductibles; pour des anneaux de polynômes particuliers, où il n'existe pas de représentation unique des polynômes comme produit de polynômes irréductibles de l'anneau, et par conséquent pas non plus de théorie de l'élimination, elle peut valoir comme équivalent de ce théorème de la théorie de l'élimination.

Les variétés irréductibles qui correspondent aux idéaux primaires isolés sont précisément celles qui apparaissent dans la \og résolvante minimale \fg{}\footnote{Voir par exemple J. König, \emph{Einleitung in die allgemeine Theorie der algebraischen Größen} (Leipzig, Teubner, 1903), p. 235.}; car les zéros de tout plus grand idéal primaire non isolé sont en même temps zéros d'au moins un idéal isolé, à savoir d'un idéal dont l'idéal premier associé est divisible par celui de l'idéal non isolé. L'unicité des idéaux primaires isolés donne donc, dans les exposants, de nouveaux nombres invariants de multiplicité. Les théorèmes d'unicité dans la résolution des idéaux primaires en idéaux irréductibles peuvent eux aussi être regardés comme un complément de la théorie de l'élimination au sens de la multiplicité.

Après ces remarques, les différentes décompositions peuvent s'interpréter dans leur comportement par rapport aux variétés algébriques. Aux idéaux deux à deux étrangers correspondent des variétés qui ne possèdent aucun zéro commun; pour les idéaux mutuellement relativement premiers, aucune variété irréductible de l'un des idéaux n'est en même temps un zéro commun; les plus grands idéaux primaires ne s'annulent que dans des variétés irréductibles toutes distinctes les unes des autres; dans la décomposition en idéaux irréductibles, les mêmes variétés irréductibles peuvent aussi apparaître plusieurs fois.

Remarquons encore qu'à la place du domaine général des polynômes on peut aussi prendre pour base le domaine de toutes les formes homogènes, car on se convainc facilement que, même pour les compositions qui y valent --- l'addition n'étant définie que pour des formes de même dimension ---, les théorèmes généraux subsistent.\footnote{Mais le fait qu'ici, en cas d'ambiguïté, il puisse exister à côté de décompositions homogènes aussi des décompositions inhomogènes est montré par l'exemple $ (x^3,xy,y^3)=[(x^3,y);(y^3,x)]=[(xy,x^3,y^3,x+y^2);(xy,x^3,y^3,y+x^2)]$.}

Un exemple le plus simple des quatre décompositions différentes --- dont les formules suivent ci-dessous --- est donné, d'après ce qui précède, par une droite et deux droites qui lui sont gauches et se coupent entre elles, dont l'une contient, en multiplicité supérieure, un point distinct du point d'intersection. À la décomposition en idéaux étrangers-irréductibles correspond la décomposition en la droite et la variété qui lui est gauche; cette variété se décompose, dans la décomposition en idéaux relativement premiers-irréductibles, en les deux droites; à la décomposition en plus grands idéaux primaires correspond un détachement du point de multiplicité supérieure, tandis que la décomposition en idéaux irréductibles impose une résolution de ce point.

Si l'on prend ce point pour origine, la droite qui le parcourt pour axe des $y$, la droite qui coupe celle-ci parallèle à l'axe des $x$, et la droite qui lui est gauche parallèle à l'axe des $z$, une telle configuration est représentée par exemple par les idéaux irréductibles suivants:\footnote{Parmi eux, les trois premiers sont irréductibles comme idéaux premiers; $\ideal B_4$, parce qu'il ne possède que les diviseurs $(x^2,y,z)$ et $(x,y,z)$; $\ideal B_5$, parce que tout diviseur contient le polynôme $xy$.}
\[
\begin{gathered}
 \ideal B_1=(x-1,y),\quad
 \ideal B_2=(y-1,z),\quad
 \ideal B_3=(x,z),\quad
 \ideal B_4=(x^3,y,z),\quad
 \ideal B_5=(x^2,y^2,z).
\end{gathered}
\]
Les idéaux premiers associés sont:
\[
 \ideal P_1=\ideal B_1,
 \quad \ideal P_2=\ideal B_2,
 \quad \ideal P_3=\ideal B_3,
 \quad \ideal P_4=\ideal P_5=(x,y,z).
\]
Les plus grands idéaux primaires sont:
\[
 \ideal Q_1=\ideal B_1,
 \quad \ideal Q_2=\ideal B_2,
 \quad \ideal Q_3=\ideal B_3,
 \quad \ideal Q_4=[\ideal B_4,\ideal B_5]=(x^3,y^2,x^2y,z).
\]
Les idéaux relativement premiers-irréductibles sont:
\[
 \ideal R_1=\ideal Q_1,
 \quad \ideal R_2=\ideal Q_2,
 \quad \ideal R_3=[\ideal Q_3,\ideal Q_4]=(x^3,x^2y,xy^2,z).
\]
Les idéaux étrangers-irréductibles sont:
\[
 \ideal S_1=\ideal R_1,
 \quad
 \ideal S_2=[\ideal R_2,\ideal R_3]=((y-1)x^3,(y-1)x^2y,(y-1)xy^2,z).
\]
Il en résulte l'idéal total:
\[
\begin{aligned}
 \ideal M&=[\ideal S_1,\ideal S_2]=\ideal S_1\ideal S_2 \\
 &=((x-1)(y-1)x^3,(y-1)x^2y,(y-1)xy^2,(x-1)z,y(y-1)x^3,yz),
\end{aligned}
\]
où
\[
 1=-(y-1)x^3+(y-1)(x^3-1)+y,
\]
\[
 (y-1)(x^3-1)+y\eqzero{\ideal S_1},
 \qquad -(y-1)x^3\eqzero{\ideal S_2}.
\]
Ici les idéaux $\ideal B_1$, $\ideal B_2$, $\ideal B_3$ sont isolés, et donc déterminés de manière unique aussi dans les décompositions en idéaux irréductibles et en plus grands idéaux primaires. Les idéaux $\ideal B_4$ et $\ideal B_5$, respectivement $\ideal Q_4$, sont non isolés; ils ne sont pas déterminés de manière unique, mais peuvent par exemple être remplacés par
\[
 \ideal D_4=(x^3,y+\lambda x^2,z),
 \qquad
 \ideal D_5=(x^2+\mu xy,y^2,z),
\]
et de même $\ideal Q_4$ peut être remplacé par
\[
 \overline{\ideal Q}_4=(x^3,x^2y,y^2+\lambda xy,z).
\]

2. Tout comme l'anneau général de polynômes (et l'anneau de polynômes à coefficients entiers), tout domaine d'intégrité fini formé de polynômes satisfait aussi à la condition de finitude, comme le montre le théorème de la base de Hilbert,\footnote{Réciproquement, si, pour un anneau de polynômes, la condition de finitude est remplie et si tout polynôme admet au moins une représentation où les multiplicateurs sont de degré inférieur en $x$, alors l'anneau est un domaine d'intégrité fini.} les coefficients pouvant être attribués à un corps arbitraire. Nous donnons encore un exemple pour l'anneau de tous les polynômes pairs, comme anneau le plus simple où, à cause de $x^2\cdot y^2=(xy)^2$, il n'existe pas de représentation unique des polynômes comme produit de polynômes irréductibles de l'anneau. Il s'agit de la même configuration que dans l'exemple ci-dessus, maintenant donnée par les idéaux irréductibles suivants:
\[
\begin{gathered}
 \ideal B_1=(x^2-1,xy,y^2,yz),\quad
 \ideal B_2=(y^2-1,xz,yz,z^2),\\
 \ideal B_3=(x^2,xy,xz,yz,z^2),\quad
 \ideal B_4=(x^4,xy,y^2,xz,yz,z^2),\\
 \ideal B_5=(x^2,y^2,xz,yz,z^2).
\end{gathered}
\]
Les idéaux premiers associés sont:
\[
 \ideal P_1=\ideal B_1,
 \quad \ideal P_2=\ideal B_2,
 \quad \ideal P_3=\ideal B_3,
 \quad \ideal P_4=\ideal P_5=(x^2,xy,y^2,xz,yz,z^2).
\]
Le fait que $\ideal P_1$ soit un idéal premier résulte de ce que tout polynôme du domaine devient de la forme
\[
 f\equiv\varphi(z^2)+xz\,\psi(z^2)\pmod{\ideal P_1}.
\]
Soient donc
\[
 f_1\equiv\varphi_1(z^2)+xz\,\psi_1(z^2),
 \qquad
 f_2\equiv\varphi_2(z^2)+xz\,\psi_2(z^2),
\]
alors de $f_1f_2\eqzero{\ideal P_1}$, on déduit que les équations
\[
 \varphi_1\varphi_2+z^2\psi_1\psi_2=0,
 \qquad
 \varphi_2\psi_1+\varphi_1\psi_2=0
\]
sont satisfaites, et de là $f_1\eqzero{\ideal P_1}$ ou $f_2\eqzero{\ideal P_1}$. Exactement de même, on voit que $\ideal P_2$ est un idéal premier; $\ideal P_3$ est un idéal premier, puisque tout polynôme du domaine est congru, modulo $\ideal P_3$, à un polynôme en $y^2$; $\ideal P_4$ se compose de tous les polynômes du domaine. Les idéaux $\ideal B_1$, $\ideal B_2$ et $\ideal B_3$ sont donc irréductibles eux aussi comme idéaux premiers; mais $\ideal B_4$ et $\ideal B_5$ ne possèdent chacun que l'unique diviseur propre $\ideal P_4$, et sont donc nécessairement irréductibles.

Des idéaux irréductibles on obtient les plus grands idéaux primaires:
\[
\begin{gathered}
 \ideal Q_1=\ideal B_1,
 \quad \ideal Q_2=\ideal B_2,
 \quad \ideal Q_3=\ideal B_3,
 \quad
 \ideal Q_4=[\ideal B_4,\ideal B_5]=(x^4,x^3y,y^2,xz,yz,z^2),
\end{gathered}
\]
les idéaux relativement premiers-irréductibles:
\[
 \ideal R_1=\ideal Q_1,
 \quad \ideal R_2=\ideal Q_2,
 \quad
 \ideal R_3=[\ideal Q_3,\ideal Q_4]=(x^4,x^3y,x^2y^2,xy^3,xz,yz,z^2),
\]
les idéaux étrangers-irréductibles:
\[
 \ideal S_1=\ideal R_1,
 \quad \ideal S_2=[\ideal R_2,\ideal R_3]
 =((y^2-1)x^4,(y^2-1)x^3y,(y^2-1)x^2y^2,
 (y^2-1)xy^3,xz,yz,z^2).
\]
Comme dans le premier exemple, on a:
\[
 [\ideal B_4,\ideal B_5]=[\ideal D_4,\ideal D_5],
 \qquad
 [\ideal Q_3,\ideal Q_4]=[\ideal Q_3,
\overline{\ideal Q}_4],
\]
où
\[
 \ideal D_4=(x^4,xy+\lambda x^2,\ldots),\qquad
 \ideal D_5=(x^2+\mu xy,\ldots),\qquad \lambda\mu\ne1,
\]
et
\[
 \overline{\ideal Q}_4=(x^4,x^3y,y^2+\lambda xy,\ldots).
\]
Les autres idéaux irréductibles, respectivement plus grands idéaux primaires, $\ideal B_1$, $\ideal B_2$, $\ideal B_3$, sont déterminés de manière unique comme idéaux isolés.

\setcounter{footnote}{46}
\subsection*{\S~11. Exemples tirés de la théorie des nombres et de la théorie des expressions différentielles.}

1. Supposons que le domaine $\Sigma$ considéré se compose de tous les nombres entiers rationnels pairs. On peut donc faire correspondre $\Sigma$ de manière biunivoque au domaine de tous les nombres entiers rationnels, en faisant correspondre à chaque nombre $2a$ de $\Sigma$ le nombre $a$. Il en résulte immédiatement que tout idéal dans $\Sigma$ est un idéal principal $(2a)$; dans la représentation par base $2c=n\cdot2a$, pour chaque élément $2c$ de l'idéal, les $n$ impairs ne signifient que des abréviations pour des sommes finies.

Les idéaux premiers du domaine sont donnés par $\ideal P_0=\ideal O=(2)$ et par $\ideal P=(2p)$, où $p$ désigne un nombre premier impair; tout idéal premier est donc divisible par $\ideal P_0$, mais par aucun autre idéal premier. Les idéaux primaires sont donnés par
\[
 \ideal Q_{\varrho_0}=(2\cdot2^{\varrho_0})
 \qquad\text{et}\qquad
 \ideal Q_\varrho=(2p^\varrho);
\]
ils sont en même temps des idéaux irréductibles et, d'après ce qui a été dit des idéaux premiers, deux idéaux appartenant à deux nombres premiers impairs distincts sont mutuellement relativement premiers, mais aucun $\ideal Q$ n'est relativement premier à un $\ideal Q_{\varrho_0}$ quelconque.\footnote{En effet, de $2b\cdot2p_1^{\varrho_1}\eqzero{2p_2^{\varrho_2}}$, pour $p_1\ne p_2$ impairs, on déduit toujours que $2b\eqzero{2p_2^{\varrho_2}}$; mais de $2b\cdot2p^\varrho\eqzero{2\cdot2^{\varrho_0}}$, on déduit seulement $2b\equiv2\cdot2^{\varrho_0-1}\pmod{2\cdot2^{\varrho_0}}$.} Les $\ideal Q_{\varrho_0}$ sont donc les seuls idéaux primaires non isolés. À la décomposition unique de $a$ en puissances de nombres premiers correspond la représentation unique de l'idéal $(2a)$ par de plus grands idéaux primaires, en même temps irréductibles:
\[
 (2a)=[(2\cdot2^{\varrho_0}),(2p_1)^{\varrho_1},
\ldots,(2p_a)^{\varrho_a}].
\]
Ainsi, contrairement aux exemples de l'anneau de polynômes, les plus grands idéaux primaires non isolés sont eux aussi déterminés de manière unique. Pour $\varrho_0=0$, il s'agit en même temps d'une représentation par idéaux mutuellement relativement premiers, tandis que pour $\varrho_0>0$ l'idéal est relativement premier-irréductible.

Alors que, dans le domaine de tous les entiers, les quatre décompositions différentes coïncident, tel n'est ici le cas que, d'une part, pour les deux décompositions en plus grands idéaux primaires et en idéaux irréductibles, et, pour $\varrho_0>0$, d'autre part, pour les décompositions en idéaux étrangers-irréductibles (tout idéal est étranger-irréductible, puisque le domaine ne possède pas d'unité) et en idéaux relativement premiers-irréductibles; pour $\varrho_0=0$, en revanche, les décompositions étrangères-irréductibles et relativement premières-irréductibles deviennent distinctes l'une de l'autre. En même temps se présente déjà ici un exemple du fait qu'un idéal premier peut être divisible par un autre sans lui être identique; plus généralement, du fait que de la divisibilité ne résulte pas la représentation comme produit. Ce dernier fait --- conséquence de ce que le domaine ne contient pas d'unité --- est aussi la raison pour laquelle il n'existe pas de représentation unique des nombres du domaine comme produit de nombres irréductibles du domaine, bien que tout idéal devienne un idéal principal; l'introduction du plus petit commun multiple se révèle donc ici nécessaire. Remarquons encore que les rapports restent exactement les mêmes si, au lieu de tous les nombres pairs, on prend pour base tous les nombres divisibles par un nombre premier fixé ou par une puissance d'un nombre premier.

En revanche, les idéaux irréductibles et primaires deviennent eux aussi différents lorsque $\Sigma$ se compose de tous les nombres divisibles par un nombre composé $g=\mathfrak p_1^{\sigma_1}\ldots\mathfrak p_\nu^{\sigma_\nu}$. Tout idéal redevient un idéal principal $(g\cdot a)$, et les idéaux premiers sont de nouveau donnés par $\ideal P_0=\ideal O=(g)$; $\ideal P=(g\cdot p)$, où $p$ désigne un nombre premier distinct de ceux qui entrent dans $g$. En revanche, les idéaux irréductibles sont donnés par $\ideal B_{\lambda_i}=(g\cdot p_i^{\lambda_i})$ et $\ideal B_e=(g\cdot p^e)$; les idéaux primaires par $\ideal Q_e=\ideal B_e$ et par ceux qui diffèrent des idéaux irréductibles,
\[
 \ideal Q_{\lambda_1\ldots\lambda_\nu}=(g\cdot p_1^{\lambda_1}\cdots p_\nu^{\lambda_\nu}),
\]
où les $\ideal B_{\lambda_i}$ et $\ideal Q_{\lambda_1\ldots\lambda_\nu}$ possèdent tous le même idéal premier associé $\ideal P_0=(g)$. Ici encore subsiste l'unicité de la décomposition en idéaux irréductibles, et par conséquent aussi de la décomposition en plus grands idéaux primaires; ainsi les idéaux non isolés sont de nouveau déterminés de manière unique.

2. Un exemple d'anneau non commutatif est fourni par la théorie des idéaux dans les anneaux de polynômes non commutatifs traitée dans le travail Noether--Schmeidler. Il s'agit en particulier des idéaux \og complètement réductibles \fg{}, c'est-à-dire de ceux pour lesquels les composantes de la décomposition sont deux à deux étrangères et ne possèdent aucun diviseur propre; les composantes sont donc a fortiori irréductibles. On obtient ainsi, en complément de l'isomorphie démontrée là d'après le \S~9, l'égalité du nombre des composantes dans deux décompositions différentes. Pour la décomposition des systèmes d'expressions différentielles linéaires partielles ou ordinaires, qui apparaît comme cas particulier de ce travail, on a ainsi obtenu un résultat qui, même dans le cas connu d'une expression différentielle linéaire ordinaire, ne semble pas avoir été remarqué.

En même temps, le système $T$ de toutes les classes modulo un idéal fixé $\ideal M$, avec l'anneau de polynômes non commutatifs $\Sigma$, fournit un double domaine $(\Sigma,T)$ où $T$ possède la propriété de module relativement à $\Sigma$. En effet, la différence de deux classes modulo $\ideal M$ est de nouveau une classe modulo $\ideal M$, de même que le produit d'une classe modulo $\ideal M$ par un polynôme quelconque, tandis que le produit de deux classes modulo $\ideal M$ n'existe pas (loc. cit. \S~3). Les systèmes de classes modulo $\ideal M$ qui y sont désignés comme \og sous-groupes \fg{} forment ainsi des exemples de modules dans des doubles domaines $(\Sigma,T)$ où l'anneau $\Sigma$ pris pour base est non commutatif.

\setcounter{footnote}{47}
\subsection*{\S~12. Exemple tiré de la théorie des diviseurs élémentaires.}

Il s'agit d'une conception de la théorie des diviseurs élémentaires imposée par les développements généraux, mais qui est elle-même supposée connue.

Soit $\Sigma$ le domaine de toutes les matrices entières à $n^2$ éléments, pour lesquelles l'addition et la multiplication sont définies au sens usuel pour les matrices. $\Sigma$ devient donc un anneau non commutatif; les idéaux sont ainsi en général unilatères, les idéaux bilatères n'y étant contenus que comme cas particulier.\footnote{La théorie des idéaux de ces domaines fait l'objet des travaux de Du Pasquier: \emph{Zahlentheorie der Tettarionen}, Dissertation Zürich, Vierteljahrsschr. d. Naturf. Ges. Zürich, 51 (1906); \emph{Zur Theorie der Tettarionenideale}, ibid., 52 (1907). Le contenu du second travail est la démonstration que tout idéal devient un idéal principal.}

Nous montrons d'abord que tout idéal devient un idéal principal. Pour cela, dans le cas des idéaux à droite, nous associons à chaque matrice $A=(a_{ik})$ un module
\[
 \mathsf{A}=(a_{11}\xi_1+\cdots+a_{1n}\xi_n,
 \ldots,
 a_{n1}\xi_1+\cdots+a_{nn}\xi_n)
\]
de formes linéaires entières. Réciproquement, à ce module correspond toute matrice qui fournit une base de $\mathsf{A}$, donc, avec $A$, aussi $UA$, où $U$ est unimodulaire. Plus généralement, au produit $PA$ correspond un module $\mathsf{B}$ qui devient multiple de $\mathsf{A}$. Une forme linéaire isolée de $\mathsf{A}$ est donnée par les $P$ qui ne contiennent qu'une seule ligne non nulle.

Soient maintenant $A_1,A_2,
\ldots,A_\nu,
\ldots$ tous les éléments d'un idéal $\ideal M$; soient $A_1,A_2,
\ldots,A_\nu,
\ldots$ les modules qui leur sont associés, $A$ leur plus grand commun diviseur, et $UA$ la matrice la plus générale associée à ce module $A$. À chaque forme linéaire isolée de $A$ correspond alors, d'après la définition du plus grand commun diviseur, une matrice $P_1A_{i_1}+\cdots+P_\sigma A_{i_\sigma}$, où, d'après ce qui précède, les $P$ ne possèdent qu'une seule ligne non nulle. Il en résulte que la matrice $A$ correspondant à une base de $A$ admet elle aussi une telle représentation, cette fois avec un $P$ général, et devient par conséquent un élément de $\ideal M$. Comme, en outre, chaque module $A_i$ est divisible par $A$, toute matrice $A_i$ devient divisible par $A$; $A$, et plus généralement $UA$, forme une base de $\ideal M$. S'il s'agit d'idéaux à gauche, il faut considérer de manière correspondante les colonnes de chaque matrice comme base d'un module; tout idéal devient un idéal principal, pour lequel, avec $A$, aussi $AV$ devient une base, $V$ désignant une matrice unimodulaire arbitraire.

Dans ce qui suit, pour le rapport avec la théorie des diviseurs élémentaires, nous prenons pour base des idéaux bilatères, pour lesquels donc, en particulier, $PAQ$ appartient à l'idéal avec $A$. Alors la base la plus générale d'un tel idéal est donnée, d'après ce qui précède, par $UAV$, où $U$ et $V$ sont unimodulaires. Les éléments de base épuisent donc une classe de matrices équivalentes,\footnote{Aux éléments de base des idéaux unilatères correspondent des classes à droite, respectivement à gauche.} et il existe une relation biunivoque entre idéal et classe, par conséquent aussi entre idéal et système de diviseurs élémentaires $(a_1\mid a_2\mid\cdots\mid a_n)$ de la classe, où les $a_i$ sont, comme on sait, des entiers non négatifs dont chacun divise le suivant. La matrice de la classe qui apparaît dans la forme normale déterminée par les diviseurs élémentaires peut ainsi être conçue comme une base spéciale de l'idéal; de la divisibilité des idéaux, respectivement des classes, résulte celle des diviseurs élémentaires, et réciproquement.

Mais Du Pasquier a montré, loc. cit. \S~11, que, pour tout idéal bilatère, le rang est $n$ et tous les diviseurs élémentaires coïncident. Pour pouvoir considérer le cas de diviseurs élémentaires généraux, nous devons donc partir non pas des idéaux, mais directement des classes bilatères (qui seront elles aussi désignées par de grandes lettres gothiques). Une classe $\mathfrak A=UAV$ est alors divisible par une autre classe $\mathfrak B$ lorsque $A=PBQ$.

En général, on a: \emph{le plus petit commun multiple (le plus grand commun diviseur) de deux classes s'obtient en formant le plus petit commun multiple (le plus grand commun diviseur) des systèmes de diviseurs élémentaires correspondants}.\footnote{Dans le travail \emph{Zur Theorie der Moduln}, Math. Ann. 52 (1899), p. 1, E. Steinitz définit le plus petit commun multiple (le plus grand commun diviseur) de classes par les plus petits communs multiples et les plus grands communs diviseurs des systèmes élémentaires. Indépendamment de cela, le plus petit commun multiple de classes se trouve comme \og composition de congruence \fg{} chez H. Brandt, \emph{Komposition der binären quadratischen Formen relativ einer Grundform}, J. f. M. 150 (1919), p. 1.} En effet, soient
\[
 (a_1\mid a_2\mid\cdots\mid a_n),\quad
 (b_1\mid b_2\mid\cdots\mid b_n),\quad
 (c_1\mid c_2\mid\cdots\mid c_n),
\]
où $c_i=[a_i,b_i]$, les systèmes de diviseurs élémentaires de $\mathfrak A$, $\mathfrak B$, $\mathfrak C^*$, et soit $\mathfrak C=[\mathfrak A,\mathfrak B]$. Alors $\mathfrak C^*$ est divisible par $\mathfrak A$ et $\mathfrak B$, donc par $\mathfrak C$; réciproquement, les diviseurs élémentaires de $\mathfrak C$ sont divisibles par ceux de $\mathfrak C^*$, donc $\mathfrak C$ est divisible par $\mathfrak C^*$, et par suite $\mathfrak C=\mathfrak C^*$. La démonstration pour le plus grand commun diviseur s'obtient de manière correspondante.

À la représentation unique des diviseurs élémentaires $a_i$ comme plus petit commun multiple de puissances de nombres premiers correspond ainsi une représentation de $\mathfrak A$ comme plus petit commun multiple de classes $\mathfrak Q$, dont les diviseurs élémentaires sont donnés par les puissances d'un nombre premier, en symboles
\[
 \mathfrak Q\sim(p^{r_1}\mid p^{r_2}\mid\cdots\mid p^{r_\varrho}\mid0\mid\cdots\mid0),
 \qquad r_1\leqq r_2\leqq\cdots\leqq r_\varrho.
\]
Si, en particulier, le rang $\varrho$ est égal à $n$, on a ici une décomposition en classes étrangères et étrangères-irréductibles, qui est unique malgré le caractère non commutatif du domaine.

Les classes $\mathfrak Q$ se décomposent encore en classes qui correspondent au système de diviseurs élémentaires:
\[
\begin{gathered}
 \mathfrak B_1\sim(p^{r_1}\mid\cdots\mid p^{r_1}),\quad
 \mathfrak B_2\sim(1\mid p^{r_2}\mid\cdots\mid p^{r_2});\quad \ldots,\\
 \mathfrak B_\varrho\sim(1\mid\cdots\mid1\mid p^{r_\varrho}\mid\cdots\mid p^{r_\varrho});\quad
 \mathfrak B_{\varrho+1}\sim(1\mid\cdots\mid1\mid0\mid\cdots\mid0),
\end{gathered}
\]
où, dans $\mathfrak B_\nu$, le nombre $1$ apparaît chaque fois $(\nu-1)$ fois. Si, par exemple, $r_1=r_2=\cdots=r_\mu=0$, alors $\mathfrak B_1,
\ldots,\mathfrak B_\mu$ deviennent égales à la classe unité et sont par conséquent à omettre dans la décomposition; il en va de même pour $\mathfrak B_{e+1}$ lorsque $\varrho=n$. Si, de plus, $r_\nu=r_{\nu+1}=\cdots=r_{\nu+\lambda}$, alors $\mathfrak B_{\nu+1},
\ldots,\mathfrak B_{\nu+\lambda}$ deviennent des diviseurs propres de $\mathfrak B_\nu$, et doivent donc également être éliminées. Désignons par $\mathfrak B_{i_1},
\ldots,\mathfrak B_{i_k}$ celles qui restent et qui donnent maintenant une représentation la plus courte; alors
\[
 \mathfrak Q=[\mathfrak B_{i_1},
\ldots,\mathfrak B_{i_k}]
\]
est la décomposition unique de $\mathfrak Q$ en classes irréductibles. En effet, soit
\[
 \mathfrak B_\nu\sim(1\mid\cdots\mid1\mid p^{r_\nu}\mid\cdots\mid p^{r_\nu})
\]
représentable comme plus petit commun multiple de
\[
 \mathfrak C\sim(1\mid\cdots\mid1\mid p^{s_1}\mid\cdots\mid p^{s_\lambda})
 \quad\text{et}\quad
 \mathfrak D\sim(1\mid\cdots\mid1\mid p^{t_1}\mid\cdots\mid p^{t_\mu}).
\]
Alors, dans le système de diviseurs élémentaires de $\mathfrak C$ et de $\mathfrak D$, le nombre $1$ doit se trouver aux premières $(\nu-1)$ places; à la $\nu$-ième place, l'exposant $s_1$ ou $t_1$, disons $s_1$, doit devenir égal à $r_\nu$. Mais comme
\[
 r_\nu=s_1\leqq s_2\leqq\cdots\leqq s_\lambda\leqq r_\nu
\]
on obtient $\mathfrak C=\mathfrak B_\nu$, donc $\mathfrak B_\nu$ est irréductible.\footnote{En revanche, les $\mathfrak B$ sont encore réductibles en classes unilatères; ici il n'existe plus de relation biunivoque entre diviseurs élémentaires et classe, et donc pas non plus de décomposition unique en classes unilatères irréductibles. C'est ce que montre, par exemple, l'exemple suivant, indiqué par H. Brandt, de décomposition en classes à droite (où les classes sont représentées par une matrice de base, respectivement par le module correspondant):
\[
\begin{gathered}
 \mathfrak B=[\mathfrak C_1,\mathfrak C_2]=[\mathfrak D_1,\mathfrak D_2];\quad
 \mathfrak B\sim\begin{pmatrix}p&0\\0&p\end{pmatrix};\quad
 \mathfrak C_1\sim\begin{pmatrix}1&0\\0&p\end{pmatrix};\quad
 \mathfrak C_2\sim\begin{pmatrix}p&0\\0&1\end{pmatrix};\\
 \mathfrak D_1\sim\begin{pmatrix}1&0\\0&p\end{pmatrix};
 \quad
 \mathfrak D_2\sim\begin{pmatrix}p&0\\(p-1)&1\end{pmatrix}.
\end{gathered}
\]
En effet, les modules $(\xi,p\eta)$; $(p\xi,\eta)$; $(p\xi,(p-1)\xi+\eta)$ ne possèdent chacun comme diviseur propre que le module $(\xi,\eta)$, et sont donc irréductibles; ils sont en outre distincts les uns des autres.} La même chose vaut pour $\mathfrak B_{\varrho+1}$, où $p$ doit être remplacé par $0$. Mais chacune de ces classes irréductibles donne, comme le montre la formation du plus petit commun multiple, un exposant déterminé et la place où cet exposant apparaît pour la première fois dans le système de diviseurs élémentaires de $\mathfrak Q$, respectivement de $\mathfrak A$, tandis que $\mathfrak B_{\varrho+1}$ indique le rang. Comme ces nombres sont fixés de manière unique par le système de diviseurs élémentaires de $\mathfrak Q$, respectivement de $\mathfrak A$, et comme la relation entre diviseurs élémentaires et classe est biunivoque, la décomposition de $\mathfrak Q$, et de même celle d'une classe quelconque, en classes irréductibles est donc unique.

En résumé, on peut dire: \emph{Toute classe bilatère $\mathfrak A$ de matrices entières ayant un nombre borné d'éléments se représente de manière unique comme plus petit commun multiple d'un nombre fini de classes bilatères irréductibles. Chaque classe irréductible représente alors un diviseur premier fixe du système de diviseurs élémentaires de $\mathfrak A$, un exposant associé et la place où cet exposant apparaît pour la première fois. La classe irréductible correspondant au diviseur $0$ indique le rang de $\mathfrak A$.}

\begin{flushleft}
Erlangen, octobre 1920.
\end{flushleft}
\begin{center}
(Reçu le 16. 10. 1920.)
\end{center}
\clearpage
\begingroup
\def\NoetherPaperTwentyBodyOnly{1}
\ifdefined\NoetherPaperTwentyBodyOnly
\else
\documentclass[11pt]{article}
\usepackage[a4paper,margin=1.45cm]{geometry}
\usepackage[T1]{fontenc}
\usepackage[utf8]{inputenc}
\usepackage{lmodern}
\usepackage[french]{babel}
\usepackage{amsmath,amssymb,mathtools}
\usepackage[protrusion=true,expansion=false]{microtype}
\setlength{\parindent}{1.25em}
\setlength{\parskip}{0.10em}
\allowdisplaybreaks
\setlength{\emergencystretch}{12em}
\sloppy\hbadness=10000\vbadness=10000\hfuzz=2pt\vfuzz=10pt\raggedbottom
\makeatletter\renewcommand\footnotesize{\@setfontsize\footnotesize{7.5}{8.7}\selectfont}\makeatother

\begin{document}
\fi


\setcounter{footnote}{0}
\section*{\NoetherNavAnchor{P20}{20. Un critère algébrique d’irréductibilité absolue}20. Un critère algébrique d'irréductibilité absolue}
\begin{center}
Par Emmy Noether, à Göttingen.\\[0.6em]
Math. Ann. 85 (1922), p.~26--33
\end{center}

Un polynôme -- à coefficients dans un corps arbitraire, défini abstraitement -- est dit \emph{absolument irréductible} s'il demeure irréductible dans le corps algébriquement clos auquel le domaine des coefficients peut être étendu.\footnote{Que tout corps puisse être étendu, et cela essentiellement de manière unique, en un corps algébriquement clos, c'est-à-dire en un corps dans lequel tout polynôme d'une variable se décompose en facteurs linéaires, a été montré par E. Steinitz dans son \emph{Algebraische Theorie der Körper} (J. f. M. 137 (1910), p.~167); il a ainsi donné l'équivalent rationnel du théorème fondamental de l'algèbre.} Je montre dans ce qui suit que l'irréductibilité absolue, par opposition à l'irréductibilité relativement à un corps donné, peut être saisie algébriquement; plus précisément, on a le théorème suivant:

\emph{À tout polynôme de $n\geqq2$ variables à coefficients indéterminés on peut associer une fonction entière rationnelle, à coefficients entiers, de ces coefficients et d'autres indéterminées, la forme de réductibilité; de telle sorte que, pour tout polynôme spécial de même degré, la condition nécessaire et suffisante de l'irréductibilité absolue soit donnée par le non-évanouissement de la forme de réductibilité. En d'autres termes: pour tout système spécial de valeurs des coefficients pour lequel la forme de réductibilité s'annule identiquement en les indéterminées et qui n'entraîne pas d'abaissement du degré, le polynôme spécial devient réductible, et réciproquement.}

Si, au lieu de polynômes, on considère des formes homogènes en une variable de plus, la condition de degré est remplacée par la condition plus simple que le système de coefficients ne s'annule pas identiquement.

Comme conséquence directe du théorème, on obtient encore le résultat énoncé d'abord par A. Ostrowski\footnote{\emph{Zur arithmetischen Theorie der algebraischen Größen}. Gött. Nachr. 1919, p.~279 (lemme, p.~296). Dans le cas des corps formés de nombres ordinaires, Ostrowski, autant que je sache, est aussi parvenu au théorème principal ci-dessus; il publiera prochainement les recherches qui s'y rapportent. Que, dans la disposition de ma démonstration, le théorème principal vaille pour des corps arbitraires, définis abstraitement, tient au fait que la forme de réductibilité construite pour des indéterminées possède des coefficients numériques indépendants du corps particulier pris pour base.}:

\emph{Tout polynôme absolument irréductible dont les coefficients sont des nombres algébriques reste irréductible modulo tout idéal premier d'un corps algébrique fini $\Omega$ contenant le domaine des coefficients, à l'exception d'au plus un nombre fini d'idéaux premiers de $\Omega$. En particulier, $\Omega$ peut aussi être le corps des nombres rationnels.}

Je compte revenir ailleurs sur d'autres applications à la théorie des fonctions relativement entières.\footnote{D. Hilbert, \emph{Mathematische Probleme}. Gött. Nachr. 1900, p.~253, problème 14.}

1. Montrons d'abord que \emph{tout polynôme de $n\geqq2$ variables à coefficients indéterminés est absolument irréductible}. Posons en effet:
\begin{equation}
\begin{aligned}
F(x)&=\sum A_{i_1\ldots i_n}x_1^{i_1}\ldots x_n^{i_n} &&(i_1+\ldots+i_n\leqq l),\\
G(x)&=\sum B_{j_1\ldots j_n}x_1^{j_1}\ldots x_n^{j_n} &&(j_1+\ldots+j_n\leqq m),\\
H(x)&=F(x)G(x)=\sum C_{k_1\ldots k_n}(A,B)x_1^{k_1}\ldots x_n^{k_n} &&(k_1+\ldots+k_n\leqq l+m),
\end{aligned}
\tag{1}
\end{equation}
où les $A$ et les $B$ désignent des indéterminées, et où les $C(A,B)$ deviennent des combinaisons bilinéaires entières des $A$ et des $B$. Les quotients $D=C(A,B)/C_{0\ldots0}(A,B)$
sont donc des fonctions rationnelles des quotients $A/A_{0\ldots0}$ et $B/B_{0\ldots0}$; mais le nombre de ces fonctions est plus grand que celui de leurs arguments; il est respectivement $\binom{l+m+n}{n}-1,\ \binom{l+n}{n}-1,\ \binom{m+n}{n}-1,$ et l'on a
\[
\binom{l+m+n}{n}+1>
\binom{l+n}{n}+\binom{m+n}{n}
\quad\text{pour } n\geqq2,\ l>0,\ m>0.
\]
Il existe donc au moins \emph{une} relation algébrique entre les $D(A,B)$, et par conséquent aussi entre les $C(A,B)$:
\begin{equation}
 \varPhi(Z)\ne0\quad[\text{identiquement en }Z_{k_1\ldots k_n}],\
 \varPhi(C(A,B))=0\quad[\text{identiquement en }A,B],
\tag{2}
\end{equation}
où $\varPhi$ désigne un polynôme entier.

Un polynôme ayant des indéterminées $Z$ pour coefficients,
\begin{equation}
 E(x)=\sum Z_{k_1\ldots k_n}x_1^{k_1}\ldots x_n^{k_n}
 \qquad(k_1+\ldots+k_n\leqq l+m),
\tag{3}
\end{equation}
est donc nécessairement absolument irréductible pour $n\geqq2$.\footnote{Car les indéterminées $Z$ ne sont pas, dans la terminologie de 2, des zéros de $\mathfrak P$.}

2. L'ensemble de ces polynômes $\varPhi(Z)$ qui, sous la substitution $Z=C(A,B)$, s'annulent identiquement en $A,B$, forme un \emph{idéal de polynômes},\footnote{Ces ensembles sont ordinairement appelés modules; en fait, toutefois, la propriété qui les caractérise -- contenir, avec $\varPhi_1$ et $\varPhi_2$, aussi $\varPhi_1\mp\varPhi_2$, et, avec $\varPhi$, aussi $\Psi\varPhi$, où $\Psi$ désigne un polynôme entier arbitraire en $Z$ -- est la propriété d'idéal.} plus précisément un idéal premier $\mathfrak P$; car si un produit $\varPhi_1\varPhi_2$ s'annule par la substitution, l'un au moins des facteurs s'annule. Comme (2) est satisfait identiquement en $A,B$, cela subsiste pour tout système spécial de valeurs $\mathsf{A},\mathsf{B}$, où $\mathsf{A},\mathsf{B}$ et par suite aussi $\mathsf{\Gamma}=C(\mathsf{A},\mathsf{B})$ désignent des grandeurs d'un corps quelconque. \emph{Les coefficients $\mathsf{\Gamma}$ de tout polynôme réductible dans un corps d'extension satisfont donc aux conditions $\varPhi(\mathsf{\Gamma})=0$; ils sont des zéros de $\mathfrak P$,} ou encore de ses polynômes générateurs, en nombre fini. Il reste maintenant à démontrer la réciproque.

À cet effet il faut remarquer que toutes les relations entre $A,B,C$ se conservent si, par l'introduction d'une nouvelle variable $y$, on transforme les polynômes (1) en formes homogènes $F(x,y)$, $G(x,y)$, $H(x,y)$ de dimensions $l,m,l+m$. Ainsi, des coefficients $\mathsf{\Gamma}$ deviennent aussi des zéros de $\mathfrak P$ lorsque, par exemple, $F(x,y)$\footnote{Les polynômes et les formes qui naissent en remplaçant les indéterminées dans $F,G,\ldots,f,g,\ldots$ par des systèmes spéciaux de valeurs seront partout désignés par une barre: $\bar F,\bar G,\ldots,\bar f,\bar g,\ldots$.} devient une puissance de $y$; à ceux-ci correspond donc, par le passage aux polynômes inhomogènes, un abaissement du degré de $\bar H(x)$ et non une décomposition. On montrera qu'à l'exception de cet abaissement du degré la réciproque vaut toujours; en particulier, que pour les \emph{formes homogènes la réciproque vaut sans exception dès que tous les $\mathsf{\Gamma}$ ne s'annulent pas; autrement dit, que tout zéro de $\mathfrak P$ est fourni par la représentation paramétrique $\mathsf{\Gamma}=C(\mathsf{A},\mathsf{B})$ avec des grandeurs $\mathsf{A},\mathsf{B}$ d'un corps d'extension}.\footnote{Que, \og en général \fg{}, les zéros de $\mathfrak P$ soient fournis par la représentation paramétrique résulte de la propriété de l'idéal premier $\mathfrak P$ de définir une variété algébrique irréductible. Mais, comme on sait, il n'en résulte pas -- c'est ce que j'avais d'abord négligé et ce sur quoi A. Ostrowski a attiré mon attention il y a longtemps -- que la représentation paramétrique ne puisse jamais faire défaut pour un système spécial de valeurs. La démonstration donnée ici repose sur des notions plus élémentaires.}

Je mène cette démonstration par la construction directe d'un nombre fini de fonctions $\varPhi(Z)$; en effet, au moyen de la substitution de Kronecker
\[
 x_\lambda=\xi^{d^{\,n-\lambda}}
\]
j'associe aux polynômes (1) des polynômes d'\emph{une} variable; je montre que des lacunes apparaissent alors dans les exposants, et, par la formation de la norme des coefficients correspondants, je parviens aux fonctions $\varPhi(Z)$ et au critère cherché.

3. Posons:\footnote{\emph{Festschrift}, p.~11.}
\begin{equation}
 d=l+m+1,\qquad
 i=i_1d^{n-1}+i_2d^{n-2}+\ldots+i_n,
\quad
 j=j_1d^{n-1}+\ldots+j_n,
\quad
 k=k_1d^{n-1}+\ldots+k_n.
\tag{4}
\end{equation}
Alors, par la substitution de Kronecker, correspondent aux polynômes $F,G,H$ de (1), respectivement, les polynômes
\begin{equation}
\begin{aligned}
 f(\xi)&=\sum A_{i_1\ldots i_n}\xi^i &&(i_1+\ldots+i_n\leqq l),\\
 g(\xi)&=\sum B_{j_1\ldots j_n}\xi^j &&(j_1+\ldots+j_n\leqq m),\\
 h(\xi)&=\sum C_{k_1\ldots k_n}(A,B)\xi^k &&(k_1+\ldots+k_n\leqq l+m).
\end{aligned}
\tag{5}
\end{equation}
Cette correspondance est, comme on sait, biunivoque, puisque pour tous les exposants induits qui apparaissent dans (5), $k=0$ implique aussi $k_1=0,\ldots,k_n=0$; et donc $i+j=i'+j'$ implique aussi $i_1+j_1=i'_1+j'_1;\ldots;i_n+j_n=i'_n+j'_n$. Ainsi des produits distincts $\xi^i\xi^j$ ne peuvent fournir le même exposant $\xi^k$ que si les produits correspondants $x_1^{i_1}\ldots x_n^{i_n}x_1^{j_1}\ldots x_n^{j_n}$ conduisent eux aussi au même exposant. Il vient donc:
\begin{equation}
 h(\xi)=f(\xi)g(\xi),
\tag{6}
\end{equation}
et de la validité d'une relation (6), de telle sorte que dans $f(\xi)$ et $g(\xi)$ n'apparaissent aucun exposant autre que les exposants induits, on obtient inversement
\begin{equation}
 H(x)=F(x)G(x).
\tag{7}
\end{equation}
Or \emph{des lacunes apparaissent effectivement dans les exposants induits de $f(\xi)$ et de $g(\xi)$}. En effet, le plus grand exposant de $f(\xi)$ est $l\cdot d^{n-1}$, le second plus grand est $(l-1)d^{n-1}+d^{n-2}$; leur différence est $d^{n-2}\cdot(d-1)>1$, puisque $d=l+m+1\geqq3$, $n\geqq2$; il en va de même pour $g(\xi)$.

Les relations (6) et (7), valables pour les indéterminées $A,B$, se conservent pour tout système spécial de valeurs $\mathsf{A},\mathsf{B}$. Pour que les degrés se conservent également, j'homogénéise (6) et (7) au moyen des variables $\eta$ et $y$. \emph{Une forme homogène $\bar{H}(x,y)$ se décompose donc dans un corps d'extension algébrique si et seulement si, dans ce corps, la forme binaire correspondante $\bar{h}(\xi,\eta)$ peut se décomposer en deux facteurs de telle manière qu'aucun exposant autre que les exposants induits n'apparaisse dans les facteurs.}

4. Afin d'effectuer la formation de la norme mentionnée à la fin de 2, soient
\[
 a_1,b_1;\ldots; a_p,b_p;\quad a_{p+1},b_{p+1};\ldots; a_{p+q},b_{p+q}
\]
de nouvelles indéterminées. Posons
\begin{equation}
\begin{aligned}
 \varphi(\xi,\eta)&=(a_1\xi+b_1\eta)\ldots(a_p\xi+b_p\eta)
     =r_p\xi^p+r_{p-1}\xi^{p-1}\eta+\ldots+r_0\eta^p,\\
 \psi(\xi,\eta)&=(a_{p+1}\xi+b_{p+1}\eta)\ldots(a_{p+q}\xi+b_{p+q}\eta)
     =s_q\xi^q+s_{q-1}\xi^{q-1}\eta+\ldots+s_0\eta^q,\\
 \chi(\xi,\eta)&=\varphi(\xi,\eta)\psi(\xi,\eta)
     =t_{p+q}\xi^{p+q}+t_{p+q-1}\xi^{p+q-1}\eta+\ldots+t_0\eta^{p+q},
\end{aligned}
\tag{8}
\end{equation}
où $r,s,t$ désignent donc les fonctions symétriques élémentaires rendues homogènes; c'est-à-dire ces fonctions symétriques des rangées $a,b$, linéaires dans chaque rangée prise isolément, à partir desquelles toutes les fonctions symétriques entières qui sont homogènes et de même dimension dans chaque rangée isolée peuvent être composées intégralement et à coefficients entiers -- et cela d'une seule manière.

Je forme maintenant, avec d'autres indéterminées $u_{\mu\nu}$, la forme linéaire
\begin{equation}
 \zeta(u,a,b)=\sum u_{\mu\nu}r_\mu s_\nu,
\tag{9}
\end{equation}
où la sommation s'étend sur des indices prescrits $\mu$ et $\nu$. Le produit de $\zeta(u)$ par tous ses conjugués, distincts de $\zeta(u)$ et distincts entre eux, qui naissent de la permutation de la rangée $a,b$ --- la norme $N(\zeta(u))$ de $\zeta(u)$ relativement au corps des $t(a,b)$, lorsque $\zeta(u)$ est considéré comme fonction du corps des $r_\mu(a,b)$ et des $s_\nu(a,b)$ --- devient alors une fonction symétrique entière de toutes les rangées $a,b$, homogène et de même dimension dans chaque rangée; donc, d'après ce qui précède, une fonction entière à coefficients entiers, déterminée de manière unique, des $t(a,b)$ et des $u_{\mu\nu}$:
\begin{equation}
 N(\zeta(u,a,b))=T(t(a,b),u)\qquad[\text{identiquement en }a,b,u].
\tag{10}
\end{equation}
Le même raisonnement montre que l'on a aussi
\[
N(z-\zeta(u))=z^\delta+T_{\delta-1}(t,u)z^{\delta-1}+\ldots+T(t,u)
\qquad[\text{identiquement en }a,b,u,z],
\]
où tous les $T$ désignent des fonctions entières à coefficients entiers des $t$ et des $u$. Les deux côtés de cette identité s'annulent pour $z=\zeta(u)$; et même, à cause de l'indépendance algébrique des fonctions symétriques élémentaires $r(a,b)$, $s(a,b)$, identiquement en $r,s$, si l'on pose $t(a,b)$, d'après (8), égal à une combinaison bilinéaire entière $w(r,s)$ des $r$ et des $s$, et que l'on considère $\zeta(u)$ comme une fonction de $r$ et de $s$. Il vient donc
\begin{equation}
\zeta(u)^\delta+T_{\delta-1}(w(r,s),u)\zeta(u)^{\delta-1}+\ldots+T(w(r,s),u)=0
\quad[\text{identiquement en }r,s,u].
\tag{11}
\end{equation}\footnote{(11), lorsque la sommation dans (9) s'étend à tous les indices, représente l'extension kroneckerienne du théorème de Gauss, et cela essentiellement dans la disposition de la démonstration de Kronecker (\emph{Zur Theorie der Formen höherer Stufen}, Berliner Ber. 1883, Werke, 2, p.~417).}

Les identités (10) et (11) se conservent pour tous les systèmes spéciaux de valeurs $\alpha,\beta$ de $a,b$; respectivement $\varrho,\sigma$ de $r,s$. Si $\varrho,\sigma$ sont choisis en particulier de telle sorte que tous les produits $\varrho_\mu\sigma_\nu$ s'annulent -- où $\mu,\nu$ sont les indices apparaissant dans (9) -- et donc que $\zeta(u)$ s'annule par la substitution, alors, en posant $w(\varrho,\sigma)=\tau$, on obtient d'après (11):
\begin{equation}
 T(\tau,u)=0\qquad[\text{identiquement en }u].
\tag{12}
\end{equation}

Réciproquement, soient $\tau_0,\ldots,\tau_{p+q}$ des systèmes de valeurs non tous nuls pour lesquels (12) est satisfait. Comme il existe, dans un corps d'extension algébrique, une décomposition
\begin{equation}
 \bar\chi(\xi,\eta)=\tau_{p+q}\xi^{p+q}+\ldots+\tau_0\eta^{p+q}
=(\alpha_1\xi+\beta_1\eta)\ldots(\alpha_{p+q}\xi+\beta_{p+q}\eta),
\tag{12}
\end{equation}
où, dans aucun facteur, $\alpha$ et $\beta$ ne s'annulent simultanément,\footnote{Ce \og théorème fondamental de l'algèbre \fg{} rationnel est le théorème d'existence sur lequel repose la validité sans exception de la réciproque énoncée à la fin de 2.} bien que certains $\alpha$ ou $\beta$ puissent s'annuler, on a $\tau=t(\alpha,\beta)$. En substituant ces valeurs dans (10), on voit, grâce à (12), que $\zeta(u,\alpha,\beta)$ ou l'un de ses facteurs conjugués s'annule identiquement en $u$. Ainsi, après une numérotation convenable de $\alpha,\beta$, si l'on pose maintenant $r(\alpha,\beta)=\varrho$, $s(\alpha,\beta)=\sigma$, cela signifie que $\bar\chi(\xi,\eta)$ admet au moins une décomposition
\begin{equation}
\bar\chi(\xi,\eta)=\bar\varphi(\xi,\eta)\bar\psi(\xi,\eta)
=\sum \varrho_\mu\xi^\mu\eta^{p-\mu}\cdot\sum \sigma_\nu\xi^\nu\eta^{q-\nu},
\tag{13}
\end{equation}
de telle sorte que tous les produits $\varrho_\mu\sigma_\nu$ apparaissant dans $\zeta(u)$ s'annulent, mais non pas tous les $\varrho$ ni tous les $\sigma$.

5. Les degrés $p,q$ dans (8) sont maintenant pris égaux à $ld^{n-1}$ et $md^{n-1}$, c'est-à-dire aux degrés de $f(\xi)$ et de $g(\xi)$ dans (5). Les indices $\mu,\nu$ sont décomposés en $\mu^{(1)},\nu^{(1)},\mu^{(2)},\nu^{(2)}$; ici $\mu^{(1)}$ parcourt tous les exposants absents de $f(\xi)$, $\nu^{(1)}$ tous les exposants de $\xi$ dans $\psi(\xi,\eta)$; de même, $\nu^{(2)}$ parcourt tous les exposants absents de $g(\xi)$, et $\mu^{(2)}$ tous les exposants de $\xi$ dans $\varphi(\xi,\eta)$. Désignons en outre par
\[
 e(\xi)=\sum Z_{k_1\ldots k_n}\xi^k\qquad(k_1+\ldots+k_n\leqq l+m)
\]
le polynôme d'une variable qui correspond au polynôme (3), et par
\[
 S(Z,u)
\]
le polynôme entier en $Z$ et $u$ qui résulte de $T(t,u)$ --- le membre de droite de (10), déterminé de manière unique, lorsque la sommation dans (9) s'étend aux indices indiqués $\mu^{(1)},\nu^{(1)},\mu^{(2)},\nu^{(2)}$ --- en remplaçant les $t$ par les coefficients $Z$ et les zéros correspondants de $e(\xi)$.

Si l'on remplace maintenant les $r,s$ par les coefficients $A,B$ et par les zéros correspondants de $f(\xi)$ et de $g(\xi)$, alors $\zeta(u)$ s'annule par cette substitution pour la sommation indiquée. De plus, $t=w(r,s)$ passe aux coefficients $C(A,B)$ et aux zéros correspondants de $h(\xi)$; donc $T(t,u)$ passe à $S(C(A,B),u)$. D'après l'identité (11), il vient ainsi
\begin{equation}
 S(C(A,B),u)=0\qquad[\text{identiquement en }A,B,u].
\tag{14}
\end{equation}
Comme (14) se conserve pour tout système spécial de valeurs, les coefficients $\mathsf{\Gamma}=C(\mathsf{A},\mathsf{B})$ de tels $\bar H(x)$ spéciaux, ou plus généralement de tels $\bar H(x,y)$, qui se décomposent en deux facteurs dans un corps d'extension, satisfont donc à la condition
\begin{equation}
 S(\mathsf{\Gamma},u)=0\qquad[\text{identiquement en }u].
\tag{15}
\end{equation}

Réciproquement, si, pour les coefficients $\mathsf{\Gamma}$, non tous nuls, d'un polynôme $\bar H(x)$, respectivement d'une forme $\bar H(x,y)$, la condition (15) est satisfaite, cela équivaut, d'après ce qui précède, à $T(\tau,u)=0$, où l'on entend par $\tau$ tous les coefficients de $\bar h(\xi,\eta)$. D'après (13), il existe donc, dans un corps d'extension des $\mathsf{\Gamma}$, une décomposition:
\[
 \bar h(\xi,\eta)=\bar f(\xi,\eta)\cdot\bar g(\xi,\eta),
\]
où, à cause de la sommation indiquée, aucun exposant autre qu'induit n'apparaît dans $\bar f$ et $\bar g$; par conséquent il existe une relation
\[
 \bar H(x,y)=\bar F(x,y)\cdot\bar G(x,y).
\]
Il s'ensuit, dès que pour $y=1$ aucun facteur ne devient de degré nul, donc en particulier si les $\mathsf{\Gamma}$ n'entraînent pas d'abaissement du degré de $\bar H(x)$, la relation supplémentaire
\[
 \bar H(x)=\bar F(x)\cdot\bar G(x).
\]

\emph{La condition (15) est donc nécessaire et suffisante pour que $\bar H(x,y)$, et, si l'on exclut l'abaissement du degré, aussi $\bar H(x)$, se décompose en deux facteurs dans un corps d'extension.}

Il en résulte encore que $S(Z,u)$ ne peut pas s'annuler identiquement en $Z$ et $u$, puisqu'en 1 on a montré que, pour $n\geqq2$, les polynômes à coefficients indéterminés, donc aussi les formes homogènes correspondantes en une variable de plus, sont absolument irréductibles. Il vient donc, en entendant par $U$ des produits de puissances des $u$,
\[
 S(Z,u)=\sum \varPhi_\lambda(Z)U_\lambda,
\]
où les $\varPhi_\lambda(Z)$ deviennent, d'après (14), des polynômes de $\mathfrak P$. Il est ainsi montré que $\mathfrak P$ ne possède pas d'autres zéros que ceux donnés par la représentation paramétrique $\mathsf{\Gamma}=C(\mathsf{A},\mathsf{B})$, où $\mathsf{A}$ et $\mathsf{B}$ appartiennent à un corps d'extension algébrique des $\mathsf{\Gamma}$. La réciproque énoncée à la fin de 2 est donc démontrée.

6. \emph{La condition de l'irréductibilité absolue} peut maintenant être formulée directement. Il suffit de décomposer, de toutes les manières possibles, le degré de $E(x)$ dans (3) en deux sommants positifs $l+m$, et de former pour chaque décomposition la forme $S(Z,u)$. Le produit de ces formes
\begin{equation}
 R(Z,u)=\prod S(Z,u)
\tag{16}
\end{equation}
constitue alors la forme de réductibilité de $E(x)$; car, d'après ce qui précède, à chaque décomposition de $\bar E(x)$, respectivement de $\bar E(x,y)$, correspond l'annulation d'un facteur de (16), et réciproquement. Ainsi le théorème principal énoncé dans l'introduction est démontré, et il est montré en même temps comment la forme de réductibilité peut effectivement être établie en un nombre fini d'étapes.

7. De l'existence de la forme de réductibilité (16) résulte immédiatement le théorème d'Ostrowski mentionné dans l'introduction. Soit
\[
 \bar E(x)=\sum \mathsf{\Gamma}_{k_1\ldots k_n}x_1^{k_1}\ldots x_n^{k_n}
 \qquad(k_1+\ldots+k_n\leqq v)
\]
un polynôme absolument irréductible à coefficients entiers algébriques $\mathsf{\Gamma}$; que $\frK$ désigne un corps de nombres algébriques fini contenant tous les $\mathsf{\Gamma}$, et que $\mathfrak p$ soit un idéal premier de $\frK$.

Alors la forme de réductibilité $R(Z,u)$ de degré exact, formée pour $E(x)$, reste différente de zéro sous la substitution $Z=\mathsf{\Gamma}$; on a
\[
 R(\mathsf{\Gamma},u)=\sum \mathsf{P}_\lambda U_\lambda\ne0\qquad[\text{identiquement en }u].
\]
L'idéal $\mathfrak r=(\mathsf{P}_1,\mathsf{P}_2,\ldots)$ n'est donc divisible que par un nombre fini d'idéaux premiers de $\frK$. Si $\mathfrak p$ est distinct de ces idéaux en nombre fini, alors $R(\mathsf{\Gamma},u)$ ne s'annule pas identiquement modulo $\mathfrak p$; et comme les classes modulo $\mathfrak p$ forment de nouveau un corps, il s'ensuit que $\bar E(x)$ modulo $\mathfrak p$ est irréductible, plus exactement absolument irréductible. Il en va de même de $\bar E(x,y)$, de sorte qu'aucun abaissement du degré modulo $\mathfrak p$ ne se produit.

Si les coefficients de $\bar E(x)$ sont des nombres algébriques fractionnaires, il faut aussi supposer $\mathfrak p$ distinct des idéaux premiers, en nombre fini, qui divisent le dénominateur commun $\varDelta$; modulo tous les autres idéaux premiers, $\bar E(x)$ peut être remplacé par $\varDelta\cdot\bar E(x)$, de sorte que les hypothèses ci-dessus sont satisfaites. Le théorème est ainsi démontré.

\begin{flushleft}
Erlangen, août 1921.
\end{flushleft}
\begin{center}
(Reçu le 28. 8. 1921.)
\end{center}

\clearpage

\ifdefined\NoetherPaperTwentyBodyOnly
\else
\end{document}
\fi
\endgroup
\clearpage
\begingroup
\def\NoetherPaperTwentyOneBodyOnly{1}
\ifdefined\NoetherPaperTwentyOneBodyOnly
\else
\documentclass[11pt]{article}
\usepackage[a4paper,margin=1.45cm]{geometry}
\usepackage[T1]{fontenc}
\usepackage[utf8]{inputenc}
\usepackage{lmodern}
\usepackage[french]{babel}
\usepackage{amsmath,amssymb,mathtools}
\usepackage[protrusion=true,expansion=false]{microtype}
\setlength{\parindent}{1.25em}
\setlength{\parskip}{0.10em}
\allowdisplaybreaks
\setlength{\emergencystretch}{12em}
\sloppy\hbadness=10000\vbadness=10000\hfuzz=2pt\vfuzz=10pt\raggedbottom
\makeatletter\renewcommand\footnotesize{\@setfontsize\footnotesize{7.5}{8.7}\selectfont}\makeatother

\begin{document}
\fi


\section*{\NoetherNavAnchor{P21}{21. Calcul des variations formel et invariants différentiels}21. Calcul des variations formel et invariants différentiels}
\begin{center}
Encyklopädie d. math. Wiss. III, 3 (1922), p.~68--71 (dans: R. Weitzenböck, \emph{Differentialinvarianten})
\end{center}

\setcounter{footnote}{148}
\subsection*{28. Calcul des variations formel et invariants différentiels.\footnote{Ce numéro est dû à E. Noether.}}

La production de tous les invariants différentiels et la déduction des théorèmes principaux peuvent être menées de la manière la plus uniforme par les méthodes formelles du calcul des variations. Ce principe remonte à Riemann\textsuperscript{78)} et a reçu son développement principal dans la théorie des formes par Lipschitz\textsuperscript{73)}.

Si l'on part du problème variationnel appartenant à une forme homogène $f(dx)$ -- où l'on suppose
\[
\left|\frac{\partial^2 f}{\partial dx_i\,\partial dx_k}\right|\ne0
\]
-- à savoir
\begin{equation}
 \delta\int_{t_0}^{t_1} f(x')\,dt=0,
 \qquad \left(x_i'=\frac{dx_i}{dt}\right),
\tag{140}
\end{equation}
on est conduit, par intégration partielle, au système invariant d'équations
\[
 2\psi_i=\frac{d}{dt}\frac{\partial f}{\partial x_i'}-\frac{\partial f}{\partial x_i}=0,
\]
où les expressions de Lagrange $\psi_i$ sont contragrédientes aux $dx$. Formellement, cela se voit le plus simplement si l'on définit les $\psi_i$ par l'identité formelle sans intégrale correspondant à l'intégration partielle (équation centrale de Lagrange\footnote{La désignation, pour un $f$ spécial, est due à Heun; voir son article IV, 1 II, p.~447.})
\[
 \delta f-df_\delta=-2\sum\psi_i(d,d)\,\delta x_i,
 \qquad
 \left(f_\delta=\sum\frac{\partial f}{\partial dx_i}\,\delta x_i\right),
\]
où $f$ aussi bien que $df_\delta$ sont des invariants, et par conséquent aussi le membre de droite. En particulier, pour une forme quadratique, il vient
\[
 \delta\sum g_{ik}dx_i dx_k-2d\sum g_{ik}dx_i\delta x_k
 =-2\sum\psi_\mu(d,d)\,\delta x_\mu.
\]
Si l'on remplace ici $d$ par le cogrédient $(d+\lambda\delta)$ et que l'on compare les puissances de $\lambda$, on obtient le système correspondant d'identités:
\begin{equation}
\left\{
\begin{aligned}
D\sum g_{ik}dx_i dx_k-2d\sum g_{ik}dx_iDx_k
   &=-2\sum\psi_\mu(d,d)\,Dx_\mu,\\
D\sum g_{ik}dx_i\delta x_k-\delta\sum g_{ik}dx_iDx_k
   -d\sum g_{ik}\delta x_iDx_k
   &=-2\sum\psi_\mu(d,\delta)\,Dx_\mu,\\
D\sum g_{ik}\delta x_i\delta x_k-2\delta\sum g_{ik}\delta x_iDx_k
   &=-2\sum\psi_\mu(\delta,\delta)\,Dx_\mu.
\end{aligned}
\right.
\tag{141}
\end{equation}
Il en résulte que $\psi_\mu(d,\delta)$, donc en particulier aussi $\psi_\mu(d,d)$ et $\psi_\mu(\delta,\delta)$, se présentent comme des vecteurs contragrédients. De (141) on calcule
\[
 \psi_\mu(d,\delta)=\sum g_{i\mu}\,d\delta x_i+
 \sum\left[\begin{smallmatrix}ik\\ \mu\end{smallmatrix}\right]dx_i\delta x_k,
\]
et il en résulte le vecteur cogrédient:
\begin{equation}
 p^\sigma(d,\delta)=\sum g^{\sigma\mu}\psi_\mu
 =d\delta x_\sigma+
 \sum\left\{\begin{smallmatrix}ik\\ \sigma\end{smallmatrix}\right\}dx_i\delta x_k.
\tag{142}
\end{equation}
Poser $p^\sigma(d,d)=0$ donne évidemment les équations des lignes géodésiques.

La connaissance du vecteur cogrédient $p^\sigma(d,\delta)$ conduit directement à la dérivée covariante (cf. n\textsuperscript{o}~19). Si, par exemple, $h(d,\delta)$ est une forme des deux rangées $dx$ et $\delta x$, alors l'expression
\begin{equation}
 h^{(1)}(dx,\delta x)=\delta h-
 \sum\frac{\partial h}{\partial dx_\sigma}p^\sigma(d,\delta)-
 \sum\frac{\partial h}{\partial \delta x_\sigma}p^\sigma(\delta,\delta)
\tag{143}
\end{equation}
est elle-même un invariant, comme différence d'invariants, et elle est formée de telle sorte que les différentielles secondes s'éliminent. Ainsi $h^{(1)}(dx,\delta x)$ représente une forme ne contenant que des différentielles premières, la dérivée covariante de $h$. Il en va de même lorsque $h$ contient davantage de rangées $d^{(1)}x,\ldots,d^{(\lambda)}x$.\footnote{Lipschitz, J. f. Math. 72 (1870), p.~17.}

La formation de la forme de courbure chez Riemann et Lipschitz repose sur le même principe. On passe de $\delta^2 f$ à la \og forme normale de la seconde variation \fg{}
\begin{equation}
 \Omega=\delta^2\sum g_{ik}dx_i dx_k
 -2d\delta\sum g_{ik}dx_i\delta x_k
 +d^2\sum g_{ik}\delta x_i\delta x_k,
\tag{144}
\end{equation}
qui, comme agrégat d'invariants, représente elle-même encore un invariant, et dans laquelle, maintenant par généralisation de l'équation centrale de Lagrange, les différentielles troisièmes sont éliminées. L'élimination des différentielles secondes se fait par analogie avec la dérivation covariante, en formant la différence
\begin{equation}
 K=\Omega-2\sum\{p^\sigma(d,\delta)\psi_\sigma(d,\delta)-p^\sigma(\delta,\delta)\psi_\sigma(d,d)\},
\tag{145}
\end{equation}
ce qui peut aussi s'écrire\footnote{Lipschitz, J. f. Math. 82 (1876), \S~1.}
\begin{equation}
\begin{aligned}
 \frac12 K={}&d\sum\psi_\sigma(d,\delta)\delta x_\sigma
 -\delta\sum\psi_\sigma(d,d)\delta x_\sigma\\
 &-\sum\{p^\sigma(d,\delta)\psi_\sigma(d,\delta)-p^\sigma(\delta,\delta)\psi_\sigma(d,d)\}.
\end{aligned}
\tag{146}
\end{equation}
La loi de formation de $K$ peut encore s'énoncer en disant que, dans $\Omega$, les différentielles secondes sont éliminées par l'annulation de $p(d,\delta)$ et de $p(\delta,\delta)$ -- respectivement des membres de droite de (141). C'est la définition de la forme de courbure par Riemann (cf. n\textsuperscript{o}~21), mais il faut remarquer expressément -- ce que (146) montre avec une clarté particulière -- que cette annulation ne peut avoir lieu qu'après la différentiation; le mode de formation formel de Lipschitz contourne cette difficulté. De même, la dérivée covariante $h^{(1)}$ (143) pourrait être définie en éliminant dans $\delta h$ les différentielles secondes par l'annulation de $p(d,\delta),\ldots$.

De la forme de courbure résulte, par formation successive des dérivées covariantes, une suite de formes $[\Phi_4,\Phi_5,\ldots$ chez Christoffel], dont les invariants projectifs, après adjonction de $f$, épuisent selon Christoffel et Ricci la totalité des invariants différentiels de la forme quadratique, et dont l'équivalence vis-à-vis des transformations linéaires, sans autre condition d'intégrabilité, exprime l'équivalence de deux formes différentielles quadratiques (cf. n\textsuperscript{o}~22). La raison interne de ce théorème de réduction repose sur l'existence des coordonnées normales de Riemann issues du problème variationnel (140) (cf. n\textsuperscript{o}~20), lesquelles transforment en droites les lignes géodésiques passant par un point, et se transforment par conséquent linéairement sous une transformation arbitraire des $x$. La formation de tous les invariants et la question de l'équivalence sont donc ramenées à la question correspondante pour les composantes homogènes individuelles dans le développement de $f$ suivant les coordonnées normales, et l'on montre que ces composantes se composent linéairement à partir de $f,\Phi_4,\Phi_5,\ldots$. Pour les formes quadratiques, cela a été exécuté en détail par Vermeil\textsuperscript{92)}, en s'appuyant sur une note de E. Noether\textsuperscript{93)}. D'après cette note, les théorèmes et les méthodes subsistent lorsque l'on prend pour base une expression différentielle arbitraire (cf. n\textsuperscript{os}~22 et 23); ainsi apparaît la portée des méthodes du calcul des variations formel par rapport à celles de la pure théorie de l'élimination.\footnote{Comparer, pour ces développements et pour d'autres interprétations géométriques, les conférences de séminaire de F. Klein indiquées sous \og Littérature \fg{}.}

L'annulation de $p(d,\delta)$ a été interprétée géométriquement par Levi-Civita\textsuperscript{134)} comme \og déplacement parallèle d'un vecteur \fg{} (cf. aussi Hessenberg\textsuperscript{133)}). Cette notion, que Weyl a saisie axiomatiquement, est essentiellement restreinte aux formes quadratiques, mais elle admet une généralisation d'un autre genre. Elle subsiste lorsqu'on élargit le groupe (multiplication des $g_{ik}$ par une fonction arbitraire), dès que l'on prend pour base, outre la forme quadratique, encore une forme linéaire. $p(d,\delta)$ est alors déterminé de manière unique, et l'on peut effectuer des formations d'invariants correspondant à celles qui précèdent et définir des coordonnées géodésiques; mais $p(d,d)=0$ ne provient alors plus d'un problème variationnel!\footnote{H. Weyl, \emph{Reine Infinitesimalgeometrie}, Math. Ztschr. 2 (1918), p.~384--411, et \emph{Raum, Zeit, Materie} (3\textsuperscript{e} et 4\textsuperscript{e} éd.).} Les questions concernant le théorème de réduction et l'équivalence ne sont pas encore abordées ici de manière générale; R. Weitzenböck n'a donné un théorème de réduction que pour les invariants d'ordre le plus bas.\footnote{Wien. Ber. 129 (1920), p.~683 et 697.}

Enfin, il faut signaler les \og problèmes variationnels invariants \fg{} généraux, c'est-à-dire ceux dans lesquels l'intégrale $J$ (simple ou multiple) est invariante vis-à-vis d'un groupe quelconque au sens de Lie. Des cas spéciaux, intéressants pour la physique, font l'objet de travaux de Hamel\footnote{Math. Ann. 59 (1904), p.~416--434; Ztschr. Math. Phys. 50 (1904), p.~1--54.}, Herglotz\footnote{Ann. d. Phys. (4) 36 (1911), p.~511, en particulier \S~9.}, Lorentz\footnote{Verslag der Amsterd. Akad., avril, sept. 1916, oct. 1916, mai 1917.}, Fokker\footnote{Ibid., janv. 1917.}, Weyl\footnote{Ann. d. Phys. (4) 54 (1917), p.~117, \S~2.} et Klein\footnote{Gött. Nachr. 1917, p.~469--482; ibid. 1918, p.~171--189.}. La formulation principielle chez E. Noether\footnote{Gött. Nachr. 1918, p.~235--257.} montre qu'à l'invariance de $J$ vis-à-vis d'un $G_\varrho$ (groupe fini à $\varrho$ paramètres essentiels) correspondent $\varrho$ divergences linéairement indépendantes; à l'invariance vis-à-vis d'un groupe infini qui contient $\varrho$ fonctions arbitraires jusqu'à la $\sigma$-ième dérivée correspondent $\varrho$ dépendances entre les expressions de Lagrange et leurs dérivées jusqu'à l'ordre $\sigma$. Dans les deux cas, la réciproque vaut. Comme les expressions de Lagrange deviennent des invariants relatifs du groupe, on possède en même temps un procédé générateur d'invariants.

\begin{center}
(Achevé en mars 1921.)
\end{center}

\clearpage

\ifdefined\NoetherPaperTwentyOneBodyOnly
\else
\end{document}
\fi
\endgroup
\clearpage
\setcounter{footnote}{0}
\begingroup
\providecommand{\Amod}{\mathfrak A}
\providecommand{\Bmod}{\mathfrak B}
\providecommand{\Gmod}{\mathfrak G}
\providecommand{\Mmod}{\mathfrak M}
\providecommand{\Nmod}{\mathfrak N}
\providecommand{\Rmod}{\mathfrak R}
\providecommand{\Smod}{\mathfrak S}
\providecommand{\mideal}{\mathfrak m}
\providecommand{\nideal}{\mathfrak n}
\providecommand{\gideal}{\mathfrak g}
\providecommand{\hideal}{\mathfrak h}
\providecommand{\jideal}{\mathfrak j}
\providecommand{\tideal}{\mathfrak t}
\providecommand{\aideal}{\mathfrak a}
\providecommand{\bideal}{\mathfrak b}
\providecommand{\bb}{\mathfrak b}
\providecommand{\cc}{\mathfrak c}
\providecommand{\zero}[1]{\equiv 0\pmod{#1}}
\providecommand{\Norm}{N}
\def\NoetherPaperTwentyTwoSOneBodyOnly{1}
\ifdefined\NoetherPaperTwentyTwoSOneBodyOnly
\else
\documentclass[11pt]{article}
\usepackage[a4paper,margin=1.45cm]{geometry}
\usepackage[T1]{fontenc}
\usepackage[utf8]{inputenc}
\usepackage{lmodern}
\usepackage[french]{babel}
\usepackage{amsmath,amssymb,mathtools}
\usepackage[protrusion=true,expansion=false]{microtype}
\setlength{\parindent}{1.25em}
\setlength{\parskip}{0.10em}
\allowdisplaybreaks
\setlength{\emergencystretch}{12em}
\sloppy\hbadness=10000\vbadness=10000\hfuzz=2pt\vfuzz=10pt\raggedbottom
\makeatletter\renewcommand\footnotesize{\@setfontsize\footnotesize{7.5}{8.7}\selectfont}\makeatother
\providecommand{\Amod}{\mathfrak A}
\providecommand{\Bmod}{\mathfrak B}
\providecommand{\Gmod}{\mathfrak G}
\providecommand{\Mmod}{\mathfrak M}
\providecommand{\Nmod}{\mathfrak N}
\providecommand{\Rmod}{\mathfrak R}
\providecommand{\Smod}{\mathfrak S}
\providecommand{\mideal}{\mathfrak m}
\providecommand{\nideal}{\mathfrak n}
\providecommand{\gideal}{\mathfrak g}
\providecommand{\hideal}{\mathfrak h}
\providecommand{\jideal}{\mathfrak j}
\providecommand{\tideal}{\mathfrak t}
\providecommand{\aideal}{\mathfrak a}
\providecommand{\bideal}{\mathfrak b}
\providecommand{\bb}{\mathfrak b}
\providecommand{\cc}{\mathfrak c}
\providecommand{\zero}[1]{\equiv 0\pmod{#1}}
\providecommand{\Norm}{N}

\begin{document}
\fi


\setcounter{footnote}{0}
\section*{\NoetherNavAnchor{P22}{22. Sur la théorie des idéaux polynomiaux et des résultantes}22. Sur la théorie des idéaux polynomiaux et des résultantes}
\begin{center}
Math. Ann. 88 (1923), p.~53--79\\[0.9em]
{\Large\bfseries Sur la théorie des idéaux polynomiaux et des résultantes.}\\[1.0em]
Par\\[0.45em]
Kurt Hentzelt \(\dagger\).\\[0.8em]
Édité par Emmy Noether à Göttingen\footnotemark.\\[0.6em]
\rule{4em}{0.4pt}
\end{center}

Il s'agit dans ce qui suit du \emph{problème général de la théorie de l'élimination, déterminer les zéros communs de tous les polynômes d'un idéal de polynômes}; ou, ce qui revient au même, les zéros d'un nombre fini quelconque de polynômes qui forment une base de l'idéal; en même temps apparaîtra une interprétation conceptuelle des multiplicités qui interviennent. Les coefficients des polynômes peuvent être attribués à un corps quelconque; les zéros appartiennent alors au corps algébriquement clos qui en est dérivé. En outre, l'idéal peut être supposé transformé (\S{}~3) sans restriction de généralité. Le résultat essentiel est alors le suivant:

\emph{À tout idéal $\mideal$ de polynômes on peut associer une forme résultante:}
\[
 R_{\mideal}=R^{(1)}(x_1\ldots x_n)\cdot\ldots\cdot
 R^{(i)}(x_i\ldots x_n)\cdot\ldots\cdot R^{(n)}(x_n)\zero{\mideal},
\]
\footnotetext{Kurt Hentzelt est porté disparu depuis octobre 1914 devant Dixmuiden et doit être compté parmi les morts. Le mémoire donné ici est une adaptation entièrement libre de la partie la plus essentielle de sa dissertation \og Zur Theorie der Formenmoduln und Resultanten \fg{}, avec laquelle il obtint son doctorat à l'été 1914 auprès de E. Fischer à Erlangen. Cette dissertation, composée tout entière sur le fondement de ses propres idées, est construite sans lacunes; mais lemme succède à lemme, tous les concepts sont décrits par des formules à quatre et cinq indices, le texte fait presque entièrement défaut, de sorte que les plus grandes difficultés s'opposent à la compréhension. Il n'a lui-même plus eu le temps d'accomplir le remaniement projeté. Je reproduis le travail sous une forme purement conceptuelle; il en résulte une grande simplification des démonstrations, qui remontent partout, dans leurs idées fondamentales, à Hentzelt, et, je l'espère, la beauté du travail devient manifeste. -- Les parties de la dissertation qui, une base étant donnée, concernent la question de la formation, en un nombre fini d'étapes, des fonctions qui apparaissent doivent être réservées à une publication séparée. (E. N.)}
\emph{de telle sorte que $R_{\mideal}$ s'annule pour tous les zéros de $\mideal$, et seulement pour ceux-ci. Si $\nideal$ est un diviseur de $\mideal$, et si les formes résultantes $R_{\nideal}$ et $R_{\mideal}$ coïncident, alors les idéaux $\nideal$ et $\mideal$ coïncident aussi.}

La seconde partie du théorème principal montre que la forme résultante fournit les zéros avec leur multiplicité caractéristique. Cette seconde partie est analogue au fait que deux idéaux d'un corps de nombres algébrique, dont l'un est un diviseur de l'autre, coïncident si et seulement si leurs \emph{normes} coïncident; la raison interne des deux théorèmes est également la même.

En effet, $\mideal$ peut être conçu comme un module de formes linéaires $\mathfrak M_{i-1}$, en considérant chaque polynôme comme une forme linéaire dans les produits de puissances de $x_1\ldots x_{i-1}$ et en adjoignant $x_{i+1}\ldots x_n$ au domaine des coefficients. Le facteur résultant $R^{(i)}$ devient alors égal à la \emph{norme} du module fondamental correspondant (\S{}~1) $\Gmod_{i-1}$ par rapport à $\mathfrak M_{i-1}$; on obtient ainsi aussitôt une \emph{interprétation de la multiplicité -- produit du degré et de l'exposant d'un facteur de $R^{(i)}$ -- par le nombre de classes modulo $\mathfrak M_{i-1}$ linéairement indépendantes}\footnote{Ces derniers résultats montrent leur véritable signification lorsqu'on fait intervenir la décomposition des idéaux en idéaux primaires; la multiplicité correspondant à un idéal primaire est alors définie comme le nombre de classes modulo cet idéal, linéairement indépendantes, du complément; cela sera traité ailleurs. (E. N.)}, fait qui n'était connu jusqu'à présent que dans le cas spécial où la résultante peut être définie pour des coefficients indéterminés\footnote{Cf. Macaulay, \emph{The algebraic theory of modular systems}, Cambridge Tracts 19 (Cambridge University Press, 1916), no~67; voir aussi Lasker, Zur Theorie der Moduln und Ideale, Math. Ann. 60 (1905), p.~20--116; p.~98. Que, dans ce cas, la résultante et la forme résultante coïncident est montré dans les théorèmes encore non publiés de Hentzelt mentionnés sous 1). (E. N.)}.

Pour déduire ces résultats, on donne aux \S{}~1 et \S{}~2 des théorèmes sur les modules de formes linéaires dont les coefficients sont des nombres rationnels entiers ou des polynômes d'une indéterminée. Le lien avec les idéaux de polynômes introduits au \S{}~3 est assuré par le théorème d'isomorphie du \S{}~4. Le \S{}~5 donne la seconde partie du théorème principal mentionnée ci-dessus, le \S{}~7 le théorème sur les zéros, que précède au \S{}~6 une théorie générale de l'élimination. Là, en même temps, lors de l'introduction de nouvelles indéterminées, est démontrée la décomposition de la résolvante en formes linéaires de ces indéterminées; et ainsi, pour l'élimination donnée ici, est fournie une démonstration irréprochable d'une affirmation de Kronecker\footnote{La tentative de démonstration chez König, \emph{Algebraische Größen} (Leipzig 1903, Teubner), V, \S{}~4 -- pour la théorie de l'élimination de Kronecker -- est notoirement restée sans succès. Que, même pour les multiplicités introduites par König dans la théorie de l'élimination de Kronecker, la seconde partie du théorème principal de Hentzelt ne vaille pas, est montré par Macaulay, loc. cit., p.~23; il y est en outre montré que la théorie de l'élimination de Kronecker ne correspond pas à la décomposition en idéaux primaires. (E. N.)}.

\section*{\S{} 1. Modules de formes linéaires.}

Par grandeurs entières $a,b,c,\ldots$ on entend, dans les deux premiers paragraphes, des nombres rationnels entiers ou des polynômes d'une indéterminée à coefficients dans un corps $\mathsf{P}$ quelconque, défini abstraitement; par formes linéaires $a(\xi),b(\xi),\ldots$, de telles formes en un nombre fini d'indéterminées $\xi_1\ldots\xi_k$ avec des grandeurs entières pour coefficients.

Un \emph{module $\Amod$ de formes linéaires} est défini par les conditions suivantes: avec $a(\xi)$ et $b(\xi)$, $\Amod$ contient aussi la différence $a(\xi)-b(\xi)$; avec $a(\xi)$, il contient aussi $c\cdot a(\xi)$, où $c$ désigne une grandeur entière arbitraire.

Par le \emph{rang de $\Amod$}, nous entendons, comme d'habitude, le nombre maximal de formes linéaires linéairement indépendantes de $\Amod$; par le $\lambda$-ième \emph{diviseur déterminantiel} $d_\lambda$ de $\Amod$, le plus grand diviseur commun de tous les déterminants de degré $\lambda$ de $\Amod$ (c'est-à-dire de tous les déterminants de degré $\lambda$ que l'on peut former à partir de la matrice des coefficients de $\lambda$ formes linéaires quelconques de $\Amod$). Si $\varrho$ est le rang de $\Amod$, de sorte que $d_1,\ldots,d_\varrho$ sont les diviseurs déterminantiels non nuls, les \emph{diviseurs élémentaires} de $\Amod$ sont définis par $e_1=d_1$, $e_2=d_2/d_1,\ldots,e_\varrho=d_\varrho/d_{\varrho-1}$, $e_{\varrho+i}=0$. On sait que $\Amod$ est un module fini qui possède une base du module formée d'exactement $\varrho$ formes linéaires $a_1(\xi),\ldots,a_\varrho(\xi)$, au moyen desquelles toute forme linéaire de $\Amod$ se représente linéairement, avec des grandeurs entières comme coefficients: $\Amod=(a_1(\xi),\ldots,a_\varrho(\xi))$. Si $A$ désigne la matrice des coefficients de ces formes linéaires, les diviseurs déterminantiels et élémentaires de $\Amod$ coïncident donc avec ceux de $A$; il en résulte d'abord que \emph{chaque diviseur élémentaire $e_\lambda$ divise le suivant $e_{\lambda+1}$}. L'équation matricielle existant d'après la théorie des diviseurs élémentaires,
\[
 P A Q=\begin{pmatrix}
 e_1&&\\
 &\ddots&\\
 &&e_\varrho
 %
 \end{pmatrix}
\]
montre en outre -- lorsque, par une transformation unimodulaire $\xi=Q(\eta)$, on introduit de nouvelles indéterminées $\eta_1\ldots\eta_k$ -- l'existence de l'équation de modules:
\begin{equation}
 \Amod=(e_1\eta_1,e_2\eta_2,\ldots,e_\varrho\eta_\varrho).
\tag{1}
\end{equation}

Le concept le plus essentiel pour ce qui suit est celui de module fondamental\footnote{Cf. Steinitz, Rechteckige Systeme und Moduln in algebraischen Zahlkörpern II. Math. Ann. 72 (1912), p.~297--345, no~36. Hentzelt semble en tout cas être parvenu, indépendamment de Steinitz -- avec qui l'on trouve par ailleurs aussi des points de contact --, pour son cas plus simple, à ce concept sous la forme de la formule (4). (E. N.)}, selon la

\medskip
\noindent\textbf{Définition I.} \emph{Un module $\Gmod$ s'appelle module fondamental s'il n'a aucun diviseur propre de même rang}\footnote{Diviseur entendu au sens des modules: $\Amod$ est divisible par $\Bmod$, $\Amod\zero{\Bmod}$, si tout élément de $\Amod$ est contenu dans $\Bmod$; $\Bmod$ s'appelle diviseur propre s'il contient des éléments différents de ceux de $\Amod$.}.

$\Gmod$ est donc un module fondamental si et seulement si, de
\begin{equation}
 c\cdot g(\xi)\zero{\Gmod};\quad c\ne0 \quad\text{on déduit toujours:}\quad g(\xi)\zero{\Gmod}.
\tag{2}
\end{equation}

Le lien entre un module arbitraire et un module fondamental est donné par le

\medskip
\noindent\textbf{Théorème I.} \emph{Tout module $\Amod$ est divisible par un et un seul module fondamental $\Gmod$ de même rang; $\Gmod$ est défini comme le plus petit diviseur de même rang de $\Amod$ et représenté par}
\begin{equation}
 \Gmod=(\eta_1,\ldots,\eta_\varrho);
\tag{3}
\end{equation}
\emph{$\Gmod$ sera appelé le module fondamental de $\Amod$.}

La condition que $\Gmod$ doive être le plus petit diviseur de même rang s'exprime ainsi: $\Gmod$ contient toutes et seulement les formes linéaires $g(\xi)$ qui satisfont à une relation
\begin{equation}
 b\cdot g(\xi)\zero{\Amod};\quad b\ne0.
\tag{4}
\end{equation}
Il en résulte d'abord que $\Gmod$ est un module, puisque, avec
\[
 b_1g_1(\xi)\zero{\Amod};\quad b_1\ne0;
 \qquad
 b_2g_2(\xi)\zero{\Amod};\quad b_2\ne0
\]
on a aussi
\[
 b_1b_2\bigl(g_1(\xi)-g_2(\xi)\bigr)\zero{\Amod};\quad b_1b_2\ne0
\]
et bien entendu aussi $b_1\cdot c g_1(\xi)\zero{\Amod}$. Mais $\Gmod$ est également un module fondamental; car de
\[
 c\cdot g(\xi)\zero{\Gmod};\quad c\ne0
\]
on déduit, d'après (4),
\[
 b c g(\xi)\zero{\Amod};\quad bc\ne0,
\]
donc, de nouveau d'après (4), aussi $g(\xi)\zero{\Gmod}$.

Il n'existe en outre aucun module fondamental $\Gmod_1$ différent du plus petit diviseur de même rang $\Gmod$ qui satisfasse aux conditions. Car $\Gmod_1$ devrait être divisible par $\Gmod$, mais, comme module fondamental, il n'a aucun diviseur propre de même rang, et il est donc identique à $\Gmod$. Enfin, le module représenté par le membre de droite de (3) est fondamental, puisque tout diviseur propre doit contenir une indéterminée supplémentaire $\eta$ et possède donc un rang plus élevé; ainsi (3) donne le module fondamental de $\Amod$ défini de manière unique, ce qui \emph{démontre le Théorème I dans toutes ses parties}.

Un autre concept essentiel pour la suite est le concept dedekindien du quotient de deux modules\footnote{Chez Dedekind, il s'agit de modules de nombres, pour lesquels une multiplication est aussi définie, de sorte que le quotient de Dedekind ne coïncide pas avec le quotient défini ici. (E. N.)}, selon la

\medskip
\noindent\textbf{Définition II.} \emph{Le quotient $\cc=\Amod/\Bmod$ de deux modules de formes linéaires est défini comme l'ensemble de toutes les grandeurs entières $c$ qui satisfont à la condition}
\[
 c\cdot\Bmod\zero{\Amod}
\]
\emph{$\cc$ représente un idéal de grandeurs entières, donc un idéal principal. De même, le quotient $\mathfrak C=\Amod/\bb$ d'un module $\Amod$ de formes linéaires par un idéal $\bb$ de grandeurs entières est défini comme l'ensemble des formes linéaires $c(\xi)$ qui satisfont à la condition}
\[
 c(\xi)\cdot\bb\zero{\Amod}
\]
\emph{$\mathfrak C$ représente de nouveau un module de formes linéaires, et même, dès que $\bb$ est différent de l'idéal nul, un module de même rang que $\Amod$.}

Il faut seulement remarquer qu'il est clair, d'après la définition, que la propriété de module appartient chaque fois au quotient; or les systèmes de grandeurs entières qui possèdent la propriété de module sont des idéaux.

Le concept de quotient conduit à un lien entre module fondamental et plus haut diviseur élémentaire, selon le

\medskip
\noindent\textbf{Théorème II.} \emph{Entre le module fondamental et le plus haut diviseur élémentaire de $\Amod$ subsiste la relation réciproque}
\begin{equation}
 (e_\varrho)=\Amod/\Gmod;\qquad \Gmod=\Amod/(e_\varrho),
\tag{5}
\end{equation}
\emph{où $(e_\varrho)$ désigne l'idéal principal dérivé de $e_\varrho$.}

En effet, puisque chaque diviseur élémentaire divise $e_\varrho$, on obtient, d'après (1) et (3),
\begin{equation}
 e_\varrho\Gmod\zero{\Amod}\quad \text{et par conséquent}\quad (e_\varrho)\cdot\Gmod\zero{\Amod};
\tag{6}
\end{equation}
ainsi $(e_\varrho)$ est divisible par le quotient $\Amod/\Gmod$. Mais la réciproque vaut également; car de
\[
 b\Gmod\zero{\Amod}
 \quad\text{on déduit}\quad b\eta_\varrho\zero{\Amod}
 \quad\text{et donc}\quad b\equiv 0(e_\varrho),
\]
ce qui démontre la première formule (5)\footnote{Si l'on pose $\Amod_0=\Amod,\ldots,\Amod_\lambda=(\Amod,\eta_\varrho,\ldots,\eta_{\varrho-\lambda+1})$, de sorte que chaque $\Amod_\lambda$ possède le plus haut diviseur élémentaire $e_{\varrho-\lambda}$, les mêmes raisonnements donnent:
\[
 (e_\varrho)=\Amod_0/\Amod_1,\ldots,(e_{\varrho-\lambda})=\Amod_\lambda/\Amod_{\lambda+1},\ldots,(e_1)=\Amod_{\varrho-1}/\Amod_\varrho.
\]
Plus généralement, soit posé $\zeta_\lambda=b_{\lambda1}\eta_1+\cdots+b_{\lambda\lambda}\eta_\lambda$, soit $(b_{\lambda\lambda},e_\lambda)=1$, et soit
\[
 \Bmod_0=\Amod,\ldots,\Bmod_\lambda=(\Amod,\zeta_\varrho,\ldots,\zeta_{\varrho-\lambda+1});
\]
alors $\Bmod_\lambda$ possède aussi le plus haut diviseur élémentaire $e_{\varrho-\lambda}$, et l'on obtient donc:
\[
 (e_{\varrho-\lambda})=\Bmod_\lambda/\Bmod_{\lambda+1}.
\]}. (6) montre en outre que $\Gmod$ est divisible par le quotient $\Amod/(e_\varrho)$; ce quotient est un diviseur de même rang de $\Amod$, donc, d'après la définition du module fondamental de $\Amod$, il est réciproquement divisible par $\Gmod$, ce qui démontre aussi la seconde formule (5).

Si $d$ désigne un déterminant quelconque de degré $\varrho$ de $\Amod$ qui ne s'annule pas identiquement, alors $d$ est divisible par le $\varrho$-ième diviseur déterminantiel $d_\varrho$, et par conséquent par $e_\varrho$; on a donc $d\Gmod\zero{\Amod}$ et, d'après le raisonnement précédent, il s'ensuit que $\Gmod=\Amod/(d)$. Mais il est essentiel pour la suite que cette dernière représentation de $\Gmod$ comme quotient puisse aussi se démontrer sans revenir aux diviseurs élémentaires, et qu'elle ait par conséquent une validité plus générale, selon le

\medskip
\noindent\textbf{Théorème III.} \emph{Si l'on entend par grandeurs entières des polynômes en plusieurs indéterminées}\footnote{En fait on utilise seulement que les grandeurs entières se reproduisent par addition, soustraction et multiplication, avec validité des règles ordinaires du calcul, c'est-à-dire qu'elles forment un anneau; et de plus que, dans cet anneau, un produit ne s'annule que lorsqu'un facteur s'annule, de sorte que l'anneau peut être étendu à un corps par formation de quotients (adjonction de couples d'éléments). En revanche, aucune représentation du module par une base n'est utilisée; (7) vaut donc par exemple aussi pour le domaine de tous les nombres algébriques entiers comme coefficients. (E. N.)}, \emph{alors la représentation comme quotient}
\begin{equation}
 \Gmod=\Amod/(d),
\tag{7}
\end{equation}
\emph{vaut encore, où $d$ désigne un déterminant quelconque de degré $\varrho$ de $\Amod$ qui ne s'annule pas identiquement.}

Il est d'abord clair que les concepts de rang, de module fondamental et de quotient de modules -- qui maintenant ne devient seulement plus un idéal principal -- se conservent dans cette extension, de même que le Théorème I, à l'exception de la représentation par une base.

Soient maintenant
\[
 a_1(\xi)=a_{11}\xi_1+\cdots+a_{1k}\xi_k;\ \ldots;\
 a_\varrho(\xi)=a_{\varrho1}\xi_1+\cdots+a_{\varrho k}\xi_k
\]
$\varrho$ formes linéaires linéairement indépendantes de $\Amod$, qui en général ne formeront pas une base du module; les indéterminées $\xi$ soient désignées de telle manière que
\[
 d=|a_{ij}|\ne0;\qquad i,j=1,2,\ldots,\varrho.
\]
Il en résulte que $\varrho$ formes linéaires de forme spéciale appartiennent à $\Amod$:
\begin{equation}
 d\xi_\lambda+b_{\lambda,\varrho+1}\xi_{\varrho+1}+\cdots+b_{\lambda,k}\xi_k\zero{\Amod}
 \quad(\lambda=1,2,\ldots,\varrho).
\tag{8}
\end{equation}
Mais $\Amod$ ne peut contenir aucune forme linéaire dépendant seulement de $\xi_{\varrho+1}\ldots\xi_k$, par exemple $c_{\varrho+1}\xi_{\varrho+1}+\cdots+c_k\xi_k$, car sinon au moins un déterminant à $(\varrho+1)$ rangées, à savoir $d\cdot c_\sigma$, serait différent de zéro.

Or le quotient $\Amod/(d)$, comme diviseur de même rang de $\Amod$, est divisible par $\Gmod$; mais la réciproque vaut aussi. En effet, pour tout $g(\xi)\zero{\Gmod}$, on a par définition:
\[
 c\cdot g(\xi)\zero{\Amod};\quad c\ne0
 \quad\text{et par conséquent}\quad d\cdot c\cdot g(\xi)\zero{\Amod}.
\]
D'après (8), cependant,
\[
 d\cdot g(\xi)\equiv g_{\varrho+1}\xi_{\varrho+1}+\cdots+g_k\xi_k\,(\Amod),
\]
d'où il vient
\[
 c g_{\varrho+1}\xi_{\varrho+1}+\cdots+c g_k\xi_k\zero{\Amod},
\]
donc, d'après ce qui vient d'être remarqué, $c\cdot g_{\varrho+1}=0,\ldots,c\cdot g_k=0$, et ainsi, puisque $c\ne0$, aussi $g_{\varrho+1}=0,\ldots,g_k=0$, donc $d\cdot g(\xi)\zero{\Amod}$, ce qui démontre (7).

Il faut remarquer que $d\cdot g(\xi)\zero{\Amod}$ résulte aussi directement du fait que les deux modules $(a_1(\xi),\ldots,a_\varrho(\xi))$ et $(g(\xi),a_1(\xi),\ldots,a_\varrho(\xi))$ ont le même rang $\varrho$, de sorte que -- en posant $g(\xi)=c_1\xi_1+\cdots+c_k\xi_k$ -- le déterminant à $(\varrho+1)$ rangées
\[
\begin{vmatrix}
 a_1(\xi)&a_{11}&\cdots&a_{1\varrho}\\
 \vdots&\vdots&&\vdots\\
 a_\varrho(\xi)&a_{\varrho1}&\cdots&a_{\varrho\varrho}\\
 g(\xi)&c_1&\cdots&c_\varrho
\end{vmatrix}
\]
s'annule. Mais la démonstration donnée ci-dessus revient au \S{}~4, à la suite de la formule (30).

\clearpage
\ifdefined\NoetherPaperTwentyTwoSOneBodyOnly
\else
\end{document}
\fi
\endgroup
\begingroup
\providecommand{\Amod}{\mathfrak A}
\providecommand{\Bmod}{\mathfrak B}
\providecommand{\Gmod}{\mathfrak G}
\providecommand{\Mmod}{\mathfrak M}
\providecommand{\Nmod}{\mathfrak N}
\providecommand{\Rmod}{\mathfrak R}
\providecommand{\Smod}{\mathfrak S}
\providecommand{\mideal}{\mathfrak m}
\providecommand{\nideal}{\mathfrak n}
\providecommand{\gideal}{\mathfrak g}
\providecommand{\hideal}{\mathfrak h}
\providecommand{\jideal}{\mathfrak j}
\providecommand{\tideal}{\mathfrak t}
\providecommand{\aideal}{\mathfrak a}
\providecommand{\bideal}{\mathfrak b}
\providecommand{\bb}{\mathfrak b}
\providecommand{\cc}{\mathfrak c}
\providecommand{\zero}[1]{\equiv 0\pmod{#1}}
\providecommand{\Norm}{N}
\def\NoetherPaperTwentyTwoSTwoBodyOnly{1}
\ifdefined\NoetherPaperTwentyTwoSTwoBodyOnly
\else
\documentclass[11pt]{article}
\usepackage[a4paper,margin=1.45cm]{geometry}
\usepackage[T1]{fontenc}
\usepackage[utf8]{inputenc}
\usepackage{lmodern}
\usepackage[french]{babel}
\usepackage{amsmath,amssymb,mathtools}
\usepackage[protrusion=true,expansion=false]{microtype}
\setlength{\parindent}{1.25em}
\setlength{\parskip}{0.10em}
\allowdisplaybreaks
\setlength{\emergencystretch}{12em}
\sloppy\hbadness=10000\vbadness=10000\hfuzz=2pt\vfuzz=10pt\raggedbottom
\makeatletter\renewcommand\footnotesize{\@setfontsize\footnotesize{7.5}{8.7}\selectfont}\makeatother
\providecommand{\Amod}{\mathfrak A}
\providecommand{\Bmod}{\mathfrak B}
\providecommand{\Gmod}{\mathfrak G}
\providecommand{\Mmod}{\mathfrak M}
\providecommand{\Nmod}{\mathfrak N}
\providecommand{\Rmod}{\mathfrak R}
\providecommand{\Smod}{\mathfrak S}
\providecommand{\mideal}{\mathfrak m}
\providecommand{\nideal}{\mathfrak n}
\providecommand{\gideal}{\mathfrak g}
\providecommand{\hideal}{\mathfrak h}
\providecommand{\jideal}{\mathfrak j}
\providecommand{\tideal}{\mathfrak t}
\providecommand{\aideal}{\mathfrak a}
\providecommand{\bideal}{\mathfrak b}
\providecommand{\bb}{\mathfrak b}
\providecommand{\cc}{\mathfrak c}
\providecommand{\zero}[1]{\equiv 0\pmod{#1}}
\providecommand{\Norm}{N}

\begin{document}
\fi


\setcounter{footnote}{7}
\section*{\S{} 2. Théorèmes de décomposition pour les normes et les modules.}

Il s'agit maintenant de nouveau de modules de formes linéaires relativement aux nombres rationnels entiers ou aux polynômes d'une indéterminée à coefficients dans $\mathsf{P}$.

\medskip
\noindent\textbf{Définition III.} \emph{Si, comme d'habitude, on réunit en une même classe modulo $\Amod$ toutes les formes linéaires en $\xi$ qui sont congruentes entre elles, le symbole $\Gmod\mid\Amod$ désigne alors le système des classes modulo $\Amod$ représentées par des éléments de $\Gmod$. La norme de $\Gmod$ par rapport à $\Amod$ -- en symboles $\Norm(\Gmod\mid\Amod)$ -- est définie comme le déterminant de la substitution de passage de $\Gmod$ à $\Amod$, c'est-à-dire comme le déterminant de la substitution qui exprime une base quelconque du module $\Amod$, linéairement indépendante, au moyen d'une telle base du module fondamental $\Gmod$ de $\Amod$.}

D'après (1) et (3), respectivement d'après la remarque au Théorème II, on obtient:
\begin{equation}
 \Norm(\Gmod\mid\Amod)=e_1e_2\cdots e_\varrho=d_\varrho;
 \qquad
 \Gmod=\Amod/\bigl(\Norm(\Gmod\mid\Amod)\bigr).
\tag{9}
\end{equation}

Des relations (1), (3) et (9), il résulte, de manière classique, que, dans le cas des nombres rationnels entiers, $\Norm(\Gmod\mid\Amod)$ représente le nombre des classes de $\Gmod$ modulo $\Amod$, tandis que, dans le cas des polynômes d'une indéterminée -- où le nombre de ces classes devient en général infini -- le degré de $\Norm(\Gmod\mid\Amod)$ devient égal au \emph{nombre des classes modulo $\Amod$ linéairement indépendantes sur $\mathsf{P}$}. Si $\Bmod$ est un diviseur de même rang de $\Amod$, alors $\Bmod$ contient, avec un représentant d'une telle classe modulo $\Amod$, tous les éléments de cette classe; il naît donc de $\Amod$ par adjonction d'un nombre fini de classes modulo $\Amod$, respectivement de leurs combinaisons linéaires. Il en résulte immédiatement:

\medskip
\noindent\textbf{Théorème IV.} \emph{Si $\Bmod$ est un diviseur de même rang de $\Amod$, de sorte que les modules fondamentaux coïncident, et si l'on a en outre $\Norm(\Gmod\mid\Bmod)=\Norm(\Gmod\mid\Amod)$, alors on a aussi $\Bmod=\Amod$.}

Sur le lien entre la décomposition des normes et celle des modules, on a le

\medskip
\noindent\textbf{Théorème V.} \emph{Soit}
\begin{equation}
 \Norm(\Gmod\mid\Amod)=r\cdot s,\qquad (r,s)=1;
\tag{10}
\end{equation}
\emph{alors il existe un et un seul module $\Rmod$, de même un et un seul module $\Smod$, tous deux diviseurs de même rang de $\Amod$, tels que}
\[
 r=\Norm(\Gmod\ \Rmod),\qquad s=\Norm(\Gmod\mid\Smod)
\]
\emph{devienne vrai. $\Rmod$ et $\Smod$ sont définis par}
\[
 \Rmod=\Amod/(s);\qquad \Smod=\Amod/(r),
\]
\emph{et $\Amod$ devient égal au plus petit commun multiple de $\Rmod$ et $\Smod$.}\footnote{Plus petit commun multiple $[\Rmod,\Smod]$ entendu au sens des modules: l'ensemble des formes linéaires qui sont divisibles aussi bien par $\Rmod$ que par $\Smod$. De manière correspondante, le plus grand commun diviseur $(\Rmod,\Smod)$ doit être entendu au sens des modules comme l'ensemble des formes linéaires qui peuvent se représenter comme somme d'un élément de $\Rmod$ et d'un élément de $\Smod$.}
\emph{Si l'on pose $r_\varrho=(r,e_\varrho)$, $s_\varrho=(s,e_\varrho)$, alors valent les réciprocités}
\begin{equation}
 (s_\varrho)=\Amod/(\Rmod);\quad \Rmod=\Amod/(s_\varrho);
 \qquad
 (r_\varrho)=\Amod/\Smod;\quad \Smod=\Amod/(r_\varrho).
\tag{11}
\end{equation}

En effet, si l'on pose en général $r_i=(r,e_i)$, $s_i=(s,e_i)$, il vient d'après (9) et (10):
\[
 (r_i,s_i)=1;\qquad e_i=r_is_i;\qquad r=r_1\cdots r_\varrho;\qquad s=s_1\cdots s_\varrho,
\]
et chaque $r_i$, respectivement $s_i$, divise le suivant $r_{i+1}$, respectivement $s_{i+1}$. Posons donc
\[
 \Rmod=(r_1\eta_1,\ldots,r_\varrho\eta_\varrho);\qquad
 \Smod=(s_1\eta_1,\ldots,s_\varrho\eta_\varrho).
\]
Alors $\Rmod$ et $\Smod$ deviennent des diviseurs de même rang de $\Amod$; à cause du fait que $r_i$ et $s_i$ sont premiers entre eux, $\Amod$ devient égal au plus petit commun multiple de $\Rmod$ et $\Smod$, et il vient en outre
\[
 \Norm(\Gmod\mid\Rmod)=r_1\cdots r_\varrho=r;\qquad
 \Norm(\Gmod\mid\Smod)=s_1\cdots s_\varrho=s.
\]
De plus on a
\[
 s_\varrho\Rmod\zero{\Amod};\qquad r_\varrho\Smod\zero{\Amod};
\]
de $c\Rmod\zero{\Amod}$ on déduit $cr_\varrho\eta_\varrho\zero{\Amod}$, donc $c\zero{(s_\varrho)}$, ce qui démontre $(s_\varrho)=\Amod/\Rmod$, et de même $(r_\varrho)=\Amod/\Smod$. De
\[
 b(\eta)=b_1\eta_1+\cdots+b_\varrho\eta_\varrho\zero{\Amod/(s_\varrho)}
\]
on déduit en outre, parce que tous les $r_i$ sont premiers avec $s_\varrho$, que $b_i\zero{(r_i)}$; donc $\Rmod=\Amod/(s_\varrho)$, et de même $\Smod=\Amod/(r_\varrho)$. Ainsi les réciprocités (11) sont démontrées; dans le cas spécial $r=1$, elles passent en (5).

Il reste encore à montrer la \emph{détermination unique} de $\Rmod$ et de $\Smod$. Soient $\overline{\Rmod}$ et $\overline{\Smod}$ des modules satisfaisant aux conditions du Théorème V, et soient $\bar r_i$ et $\bar s_i$ leurs diviseurs élémentaires. Puisque $\bar r_i$ et $\bar s_i$ sont des diviseurs de $e_i$, on obtient, à cause de
\[
 r=\bar r_1\cdots\bar r_\varrho,
 \qquad
 s=\bar s_1\cdots\bar s_\varrho,
 \qquad
 (\bar r_i,\bar s_i)=1,
\]
aussi
\[
 e_i=\bar r_i\bar s_i,
 \qquad
 \bar r_i=(r,e_i)=r_i,
 \qquad
 \bar s_i=(s,e_i)=s_i.
\]
Ainsi $\overline{\Rmod}$ et $\overline{\Smod}$ coïncident avec $\Rmod$ et $\Smod$ quant aux diviseurs élémentaires et au module fondamental. Si donc, au moyen d'une substitution unimodulaire $\eta=Q(\zeta)$, on introduit de nouvelles indéterminées $\xi$, on a
\[
 \overline{\Rmod}=(r_1\zeta_1,\ldots,r_\varrho\zeta_\varrho);
\]
\[
 \Amod=\bigl(r_1s_1(q_{11}\zeta_1+\cdots+q_{1\varrho}\zeta_\varrho),\ldots,
 r_\varrho s_\varrho(q_{\varrho1}\zeta_1+\cdots+q_{\varrho\varrho}\zeta_\varrho)\bigr).
\]
De $\Amod\zero{\overline{\Rmod}}$ il vient donc
\[
 r_i s_i(q_{i1}\zeta_1+\cdots+q_{i\varrho}\zeta_\varrho)\zero{\overline{\Rmod}};
\]
donc $r_is_iq_{ij}\zero{(r_j)}$. Comme $(r,s)=1$, cela donne $r_iq_{ij}\zero{(r_j)}$, ou encore $\Rmod\zero{\overline{\Rmod}}$. Mais puisque $\Norm(\Gmod\mid\Rmod)=\Norm(\Gmod\mid\overline{\Rmod})$, il vient $\Rmod=\overline{\Rmod}$ d'après le Théorème IV, et de même $\Smod=\overline{\Smod}$, ce qui démontre le Théorème V dans toutes ses parties.

\clearpage
\ifdefined\NoetherPaperTwentyTwoSTwoBodyOnly
\else
\end{document}
\fi
\endgroup
\begingroup
\providecommand{\Amod}{\mathfrak A}
\providecommand{\Bmod}{\mathfrak B}
\providecommand{\Gmod}{\mathfrak G}
\providecommand{\Mmod}{\mathfrak M}
\providecommand{\Nmod}{\mathfrak N}
\providecommand{\Rmod}{\mathfrak R}
\providecommand{\Smod}{\mathfrak S}
\providecommand{\mideal}{\mathfrak m}
\providecommand{\nideal}{\mathfrak n}
\providecommand{\gideal}{\mathfrak g}
\providecommand{\hideal}{\mathfrak h}
\providecommand{\jideal}{\mathfrak j}
\providecommand{\tideal}{\mathfrak t}
\providecommand{\aideal}{\mathfrak a}
\providecommand{\bideal}{\mathfrak b}
\providecommand{\bb}{\mathfrak b}
\providecommand{\cc}{\mathfrak c}
\providecommand{\zero}[1]{\equiv 0\pmod{#1}}
\providecommand{\Norm}{N}
\def\NoetherPaperTwentyTwoSThreeBodyOnly{1}
\ifdefined\NoetherPaperTwentyTwoSThreeBodyOnly
\else
\documentclass[11pt]{article}
\usepackage[a4paper,margin=1.45cm]{geometry}
\usepackage[T1]{fontenc}
\usepackage[utf8]{inputenc}
\usepackage{lmodern}
\usepackage[french]{babel}
\usepackage{amsmath,amssymb,mathtools}
\usepackage[protrusion=true,expansion=false]{microtype}
\setlength{\parindent}{1.25em}
\setlength{\parskip}{0.10em}
\allowdisplaybreaks
\setlength{\emergencystretch}{12em}
\sloppy\hbadness=10000\vbadness=10000\hfuzz=2pt\vfuzz=10pt\raggedbottom
\makeatletter\renewcommand\footnotesize{\@setfontsize\footnotesize{7.5}{8.7}\selectfont}\makeatother
\providecommand{\Amod}{\mathfrak A}
\providecommand{\Bmod}{\mathfrak B}
\providecommand{\Gmod}{\mathfrak G}
\providecommand{\Mmod}{\mathfrak M}
\providecommand{\Nmod}{\mathfrak N}
\providecommand{\Rmod}{\mathfrak R}
\providecommand{\Smod}{\mathfrak S}
\providecommand{\mideal}{\mathfrak m}
\providecommand{\nideal}{\mathfrak n}
\providecommand{\gideal}{\mathfrak g}
\providecommand{\hideal}{\mathfrak h}
\providecommand{\jideal}{\mathfrak j}
\providecommand{\tideal}{\mathfrak t}
\providecommand{\aideal}{\mathfrak a}
\providecommand{\bideal}{\mathfrak b}
\providecommand{\bb}{\mathfrak b}
\providecommand{\cc}{\mathfrak c}
\providecommand{\zero}[1]{\equiv 0\pmod{#1}}
\providecommand{\Norm}{N}

\begin{document}
\fi


\setcounter{footnote}{8}
\section*{\S{} 3. Idéaux de polynômes.}

Soit $\overline{\mideal}$ un \emph{idéal de polynômes} des $n$ indéterminées $y_1\ldots y_n$, à coefficients dans un corps $\mathsf{P}$ quelconque, défini abstraitement; c'est-à-dire que $\overline{\mideal}$ contient, avec $\Phi_1(y)$ et $\Phi_2(y)$, aussi la différence $\Phi_1(y)-\Phi_2(y)$, et, avec $\Phi(y)$, aussi $A(y)\cdot\Phi(y)$, où $A(y)$ désigne un polynôme quelconque à coefficients dans $\mathsf{P}$\footnote{La désignation \og idéal \fg{}, qui résulte conceptuellement, au lieu de l'ancien \og module \fg{} ou \og module de formes \fg{} jadis généralement usité, se révèle ici déjà nécessaire pour distinguer ces objets des modules de formes linéaires. (E. N.)}.

Que les $y$ soient soumis à une transformation dont les coefficients sont des indéterminées,
\begin{equation}
 y=U(x);\quad \text{par exemple}\quad
 \begin{cases}
 y_1=x_1,\\
 y_2=u_{21}x_1+x_2,\\
 \vdots\\
 y_n=u_{n1}x_1+\cdots+u_{n,n-1}x_{n-1}+x_n,
 \end{cases}
\tag{12}
\end{equation}
et que les indéterminées $u_{\mu\nu}$ soient adjointes au domaine des coefficients des polynômes, lequel passe ainsi à un corps $\mathsf{P}(u)$. Par cette adjonction, que $\overline{\mideal}$ passe à $\overline{\mideal}_{(u)}$; ainsi $\overline{\mideal}_{(u)}$ contient, outre les polynômes $\Phi_\lambda(y)$ de $\overline{\mideal}$, aussi toutes les combinaisons linéaires $\sum \alpha_\lambda(u)\Phi_\lambda(y)$ d'un nombre fini quelconque de $\Phi_\lambda$, où les $\alpha_\lambda(u)$ sont des fonctions rationnelles des $u_{\mu\nu}$ à coefficients dans $\mathsf{P}$. Par multiplication par un polynôme convenable en $u$, les $\alpha_\lambda(u)$ peuvent être supposés, sans restriction de généralité, être des produits de puissances en $u$. On a donc:
\begin{equation}
 \sum \alpha_\lambda(u)\Phi_\lambda(y)\zero{\mideal_{(u)}};
 \qquad
 \Phi_\lambda(y)\zero{\overline{\mideal}}\quad\text{et réciproquement}.
\tag{13}
\end{equation}
En vertu de (12), $\overline{\mideal}_{(u)}$ passe à un idéal $\mideal$ de polynômes en $x_1\ldots x_n$, à coefficients dans $\mathsf{P}(u)$:
\begin{equation}
 \overline{\mideal}_{(u)}=\mideal\quad \text{en vertu de}\quad y=U(x).
\tag{14}
\end{equation}
La combinaison des relations (14) et (13) exprime ce qui suit:
\begin{equation}
 \text{De } F(x)\zero{\mideal};\quad
 F(x)=\Phi(y)=\sum \alpha_\lambda(u)\Phi_\lambda(y)=\sum \alpha_\lambda(u)F_\lambda(x)
 \text{ on déduit } F_\lambda(x)\zero{\overline{\mideal}};
\tag{15}
\end{equation}
tandis que, réciproquement, de (14) et (15) on retrouve (13).

\medskip
\noindent\textbf{Définition IV.} \emph{Les idéaux $\mideal$ de polynômes en $x_1\ldots x_n$ à coefficients dans $\mathsf{P}(u)$ qui satisfont aux relations (15), donc qui naissent de $\overline{\mideal}_{(u)}$ en vertu de $y=U(x)$, seront appelés idéaux transformés.}

Les idéaux transformés se distinguent par l'existence de polynômes réguliers; \emph{un polynôme de degré $r$ s'appelle régulier par rapport à $x_i$ s'il contient le terme $x_i^r$ avec un coefficient non nul}; on voit que les polynômes $F_\lambda(x)$ sont réguliers par rapport à $x_1$. Dans ce qui suit il s'agira essentiellement de concevoir ces idéaux transformés comme des modules de formes linéaires en certains produits de puissances des $x$. Pour cela il faut fixer le concept d'idéal fondamental, qui passera plus tard au module fondamental. Nous faisons d'abord précéder une notation qui sera désormais maintenue partout:

\medskip
\noindent\textbf{Notation.} \emph{Par $F^{(i)},f^{(i)},a^{(i)},\ldots$ on entendra toujours des polynômes des $x$ qui sont libres de $x_1,
\ldots,x_{i-1}$.}

\medskip
\noindent\textbf{Définition V.} \emph{À chaque idéal $\mideal$ sont associés $n$ idéaux fondamentaux $\gideal_0,\gideal_1,\ldots,\gideal_{n-1}$, définis par la stipulation suivante: l'idéal fondamental de degré $(i-1)$, $\gideal_{i-1}$, contient tous et seulement les polynômes $G(x)$ pour lesquels il existe un polynôme $b^{(i)}$ -- variant en général avec $G(x)$ -- tel que}
\begin{equation}
 b^{(i)}G(x)\zero{\mideal};\qquad b^{(i)}\ne0.
\tag{16}
\end{equation}

Que les $\gideal$ soient effectivement des idéaux correspond exactement à la démonstration de la propriété de module pour $\Gmod$ à la suite de la formule (4), dont (16) est l'analogue. L'idéal $\gideal_i$ est divisible par $\gideal_{i-1}$, car tout polynôme $b^{(i+1)}$ est en même temps un $b^{(i)}$.

Pour les idéaux fondamentaux, on a:

\medskip
\noindent\textbf{Théorème VI.} \emph{Les idéaux fondamentaux d'idéaux transformés sont des idéaux transformés}\footnote{Le Théorème VI n'est utilisé qu'au \S{}~7.}.

La démonstration repose sur un théorème connu de Dedekind--Mertens\footnote{Dedekind: Über einen arithmetischen Satz von Gauß, Mitt. d. deutsch. math. Ges. zu Prag 1892. F. Mertens: Über einen algebraischen Satz, Ber. d. Ak. d. Wissensch. Wien 101 (1892), p.~1560--1566. -- L'extension kroneckerienne du théorème de Gauss aux polynômes à coefficients indéterminés est une conséquence immédiate de ce théorème.}: soient $a$ et $b$ des indéterminées, et soit posé
\[
 \sum a_{i_1\ldots i_s}z_1^{i_1}\cdots z_s^{i_s}\cdot
 \sum b_{j_1\ldots j_s}z_1^{j_1}\cdots z_s^{j_s}
 =
 \sum c(a,b)_{k_1\ldots k_s}z_1^{k_1}\ldots z_s^{k_s},
\]
\[
 \sum i\leqq \nu,
 \qquad
 \sum j\leqq \mu,
 \qquad
 \sum k\leqq \nu+\mu;
\]
si en outre $\Amod,\Bmod,\mathfrak C$ désignent les modules respectivement dérivés des $a,b,c$ au moyen de nombres rationnels entiers, il existe un exposant $q$ tel que l'identité
\begin{equation}
 \Amod^q\Bmod=\Amod^{q-1}\mathfrak C
\tag{17}
\end{equation}
vaille. Ici $\Amod^q$ se compose des produits de puissances de dimension $q$ des $a$ et de leurs combinaisons linéaires entières.

Pour démontrer le Théorème VI, posons maintenant:
\begin{equation}
\begin{cases}
 G(x)=\Gamma(y)=\sum \alpha_\lambda(u)\Gamma_\lambda(y)=\sum \alpha_\lambda(u)G_\lambda(x),\\
 b^{(i)}(x)=\varphi(y)=\sum \beta_\mu(u)\varphi_\mu(y),
\end{cases}
\tag{18}
\end{equation}
où les $\alpha_\lambda(u)$ et les $\beta_\mu(u)$ désignent donc des produits de puissances des $u$. Alors, d'après (16), (14) et (18), on a:
\begin{equation}
 \sum\beta_\mu(u)\varphi_\mu(y)\cdot
 \sum\alpha_\lambda(u)\Gamma_\lambda(y)\zero{\overline{\mideal}_{(u)}};
 \qquad
 \sum\beta_\mu(u)\varphi_\mu(y)\ne0.
\tag{19}
\end{equation}
À gauche se trouve le produit de deux polynômes en $u$, dont les coefficients sont les polynômes $\varphi_\mu(y)$ et $\Gamma_\lambda(y)$; les coefficients du polynôme-produit en $u$ sont, d'après le membre de droite de (19), des polynômes de $\overline{\mideal}$. Si donc, dans (17), on remplace les $a,b,c$ par ces polynômes en $y$, il en résulte
\[
 \left\{\sum\beta_\mu(u)\varphi_\lambda(y)\right\}^q\cdot \Gamma_\lambda(u)\zero{\overline{\mideal}_{(u)}},
\]
ce qui, d'après (18), équivaut à:
\begin{equation}
 b^{(i)}(x)^q\cdot G_\lambda(x)\zero{\mideal};
 \qquad
 G_\lambda(x)\zero{\gideal_{i-1}}.
\tag{20}
\end{equation}
Les relations (16), (18) et (20) montrent que les \emph{idéaux fondamentaux $\gideal_{i-1}$ sont effectivement des idéaux transformés.}

\clearpage
\ifdefined\NoetherPaperTwentyTwoSThreeBodyOnly
\else
\end{document}
\fi
\endgroup
\begingroup
\providecommand{\Amod}{\mathfrak A}
\providecommand{\Bmod}{\mathfrak B}
\providecommand{\Gmod}{\mathfrak G}
\providecommand{\Mmod}{\mathfrak M}
\providecommand{\Nmod}{\mathfrak N}
\providecommand{\Rmod}{\mathfrak R}
\providecommand{\Smod}{\mathfrak S}
\providecommand{\mideal}{\mathfrak m}
\providecommand{\nideal}{\mathfrak n}
\providecommand{\gideal}{\mathfrak g}
\providecommand{\hideal}{\mathfrak h}
\providecommand{\jideal}{\mathfrak j}
\providecommand{\tideal}{\mathfrak t}
\providecommand{\aideal}{\mathfrak a}
\providecommand{\bideal}{\mathfrak b}
\providecommand{\bb}{\mathfrak b}
\providecommand{\cc}{\mathfrak c}
\providecommand{\zero}[1]{\equiv 0\pmod{#1}}
\providecommand{\Norm}{N}
\def\NoetherPaperTwentyTwoSFourBodyOnly{1}
\ifdefined\NoetherPaperTwentyTwoSFourBodyOnly
\providecommand{\NoetherPaperTwentyTwoSFourEnd}{}
\else
\documentclass[11pt]{article}
\usepackage[a4paper,margin=1.45cm]{geometry}
\usepackage[T1]{fontenc}
\usepackage[utf8]{inputenc}
\usepackage{lmodern}
\usepackage[french]{babel}
\usepackage{amsmath,amssymb,mathtools}
\usepackage[protrusion=true,expansion=false]{microtype}
\setlength{\parindent}{1.25em}
\setlength{\parskip}{0.10em}
\allowdisplaybreaks
\setlength{\emergencystretch}{12em}
\sloppy\hbadness=10000\vbadness=10000\hfuzz=2pt\vfuzz=10pt\raggedbottom
\makeatletter\renewcommand\footnotesize{\@setfontsize\footnotesize{7.5}{8.7}\selectfont}\makeatother
\providecommand{\Amod}{\mathfrak A}
\providecommand{\Bmod}{\mathfrak B}
\providecommand{\Gmod}{\mathfrak G}
\providecommand{\Mmod}{\mathfrak M}
\providecommand{\Nmod}{\mathfrak N}
\providecommand{\Rmod}{\mathfrak R}
\providecommand{\Smod}{\mathfrak S}
\providecommand{\mideal}{\mathfrak m}
\providecommand{\nideal}{\mathfrak n}
\providecommand{\gideal}{\mathfrak g}
\providecommand{\hideal}{\mathfrak h}
\providecommand{\jideal}{\mathfrak j}
\providecommand{\tideal}{\mathfrak t}
\providecommand{\aideal}{\mathfrak a}
\providecommand{\bideal}{\mathfrak b}
\providecommand{\bb}{\mathfrak b}
\providecommand{\cc}{\mathfrak c}
\providecommand{\zero}[1]{\equiv 0\pmod{#1}}
\providecommand{\Norm}{N}

\begin{document}
\providecommand{\NoetherPaperTwentyTwoSFourEnd}{\end{document}}
\fi

\setcounter{footnote}{11}

\section*{\S 4. Rapport entre les idéaux de polynômes et les modules de formes linéaires.}

Pour concevoir un idéal $\mideal$ de polynômes, que l'on supposera toujours transformé, comme un module $\Mmod_{i-1}$ $(i=1\ldots n)$ de formes linéaires, il faut regarder les polynômes individuels $F(x)$ comme des formes linéaires dans les produits de puissances $\xi_\lambda$ de $x_1\ldots x_{i-1}$:
\[
  F(x)=\sum a^{(i)}_\lambda\xi_\lambda,
\]
où, suivant la notation du \S~3, les $a^{(i)}_\lambda$ désignent des polynômes en $x_i\ldots x_n$. Il résulte aussitôt de la définition de l'idéal que $\Mmod_{i-1}$ possède la propriété de module relativement à ces grandeurs entières $a^{(i)},b^{(i)},\ldots$; en outre, comme les infiniment nombreux produits de puissances $\xi$ de $x_1\ldots x_{i-1}$ n'interviennent que \emph{linéairement}, ils jouent, en raison de leur indépendance linéaire, le rôle d'indéterminées pour $\Mmod_{i-1}$. Chaque module $\Mmod_{i-1}$ contient infiniment beaucoup d'indéterminées $\xi$, mais dans chaque forme linéaire particulière il n'intervient qu'un nombre fini d'indéterminées. De même, chaque idéal fondamental $\gideal_{i-1}$ doit être conçu comme un module $\Gmod_{i-1}$ de formes linéaires, où les indéterminées sont de nouveau les produits de puissances de $x_1\ldots x_{i-1}$. Pour $i=1$, il ne s'agit que d'\emph{une} indéterminée $\xi_0$, qui correspond à l'unité.

On peut maintenant, et c'est là-dessus que repose tout ce qui suit, se limiter, malgré ces infiniment nombreuses indéterminées $\xi$, aux modules de formes linéaires en un nombre fini d'indéterminées, d'après le théorème suivant.

\medskip
\noindent\textbf{Théorème VII.} \emph{Le système des classes de $\Gmod_{i-1}$ modulo $\Mmod_{i-1}$ est isomorphe au système des classes de $\Gmod^*_{i-1}$ modulo $\Mmod^*_{i-1}$, où $\Mmod^*_{i-1}$ désigne un module en un nombre fini d'indéterminées et $\Gmod^*_{i-1}$ son module fondamental. Le module $\Mmod_{i-1}$ ne possède, après adjonction de $x_{i+1}\ldots x_n$ à $\mathsf{P}(u)$, qu'un nombre fini de diviseurs élémentaires différents de $1$, à savoir les diviseurs élémentaires de $\Mmod^*_{i-1}$ qui sont différents de $1$.}

Les diviseurs élémentaires de $\Mmod_{i-1}$ doivent ici être définis comme dans la note 8). L'isomorphie qui figure dans le Théorème VII repose sur la définition suivante.

\medskip
\noindent\textbf{Définition V.} \emph{Deux systèmes de classes modulo leurs modules respectifs s'appellent isomorphes lorsqu'ils peuvent être mis en correspondance biunivoque de telle sorte qu'à la différence de deux classes corresponde la différence des classes correspondantes, et qu'au produit d'une classe par une grandeur entière corresponde le produit de la classe correspondante par la même grandeur entière.}

La démonstration du Théorème VII s'obtient par induction complète. Pour $i=1$, le Théorème VII est manifestement exact; car $\Mmod_0$ est un module de formes linéaires en une indéterminée $\xi_0$, qui correspond au produit de puissances de dimension zéro des $x$, donc à l'unité; les grandeurs entières sont tous les polynômes en $x_1\ldots x_n$. L'idéal fondamental $\gideal_0$ devient égal à l'idéal-unité formé de tous les polynômes en $x_1\ldots x_n$, et $\Gmod_0$ est donc égal au module $(\xi_0)$, le module fondamental de $\Mmod_0$, de sorte qu'ici $\Gmod^*_0\mid\Mmod^*_0$ coïncide avec $\Gmod_0\mid\Mmod_0$. Enfin, comme $\Mmod_0$ est de rang $1$, il ne peut posséder au plus qu'un diviseur élémentaire différent de $1$.

Nous passons d'abord de $i=1$ à $i=2$, afin de conduire ensuite la conclusion inductive générale au moyen de l'intuition ainsi obtenue sur la structure de $\Gmod_1\mid\Mmod_1$. Pour cela, nous montrons d'abord que, pour $\Gmod_0\mid\Mmod_0$, on peut prendre un système de représentants borné en $x_1$; en raison de la divisibilité de $\gideal_1$ par $\gideal_0$, cela conduira alors aux indéterminées en nombre fini de $\Gmod^*_1$.

Comme idéal transformé, $\mideal$ possède d'après le \S~3 au moins un polynôme $C^{(1)}(x)$ régulier en $x_1$, disons de dimension $k$; $\Mmod_0$ contient donc la forme linéaire $C^{(1)}\xi_0$. En vertu de la régularité de $C^{(1)}(x)$, tout polynôme $H(x)$ admet une représentation:
\begin{equation}
 H(x)\equiv b^{(2)}_0+b^{(2)}_1x_1+\cdots+b^{(2)}_{k-1}x_1^{k-1}\quad (C^{(1)}(x));
\tag{21}
\end{equation}
ainsi, à cause de l'appartenance de $C^{(1)}\xi_0$ à $\Mmod_0$, on peut effectivement choisir pour $\Gmod_0\mid\Mmod_0$ un système de représentants qui n'atteint pas la $k$-ième dimension en $x_1$.

Si maintenant, en particulier, on pose pour les polynômes $G(x)$ de $\gideal_1$ et $F(x)$ de $\mideal$:
\begin{equation}
 \left\{\begin{aligned}
  G(x)&\equiv g^{(2)}_0+g^{(2)}_1x_1+\cdots+g^{(2)}_{k-1}x_1^{k-1} &&(C^{(1)}(x)),\\
  F(x)&\equiv a^{(2)}_0+a^{(2)}_1x_1+\cdots+a^{(2)}_{k-1}x_1^{k-1} &&(C^{(1)}(x)),
 \end{aligned}\right.
\tag{22}
\end{equation}
et si $\Gmod^*_1$ désigne le système de toutes les formes linéaires $g(\xi)=g^{(2)}_0\xi_0+\cdots+g^{(2)}_{k-1}\xi_{k-1}$, tandis que $\Mmod^*_1$ désigne de façon correspondante le système des formes linéaires $a(\xi)=a^{(2)}_0\xi_0+\cdots+a^{(2)}_{k-1}\xi_{k-1}$, il résulte de la propriété d'idéal de $\gideal_1$ et de $\mideal$ que $\Gmod^*_1$ et $\Mmod^*_1$ sont des modules de formes linéaires en $\xi_0\ldots\xi_{k-1}$ relativement aux grandeurs entières $b^{(2)}, c^{(2)}\ldots$. De même, l'idéal principal dérivé de $C^{(1)}(x)$ donne un module relativement à ces grandeurs entières:
\[
 \mathfrak C_1=(\zeta_0,\zeta_1,\ldots,\zeta_\nu,\ldots),
\]
où l'on a posé $\zeta_\nu=x_1^\nu C^{(1)}(x)$. Les $\zeta_\nu$ sont linéairement indépendants relativement à ces grandeurs entières $b^{(2)},c^{(2)},\ldots$, aussi bien entre eux que vis-à-vis des $\xi$ en nombre fini. De la divisibilité de $\mathfrak C_1$ par $\Mmod_1$, et par conséquent par $\Gmod_1$, il résulte que $\Gmod^*_1$ est divisible par $\Gmod_1$, respectivement $\Mmod^*_1$ par $\Mmod_1$. En tenant compte de (22), on obtient donc:
\begin{equation}
 \Gmod_1=(\Gmod^*_1,\mathfrak C_1),\qquad
 \Mmod_1=(\Mmod^*_1,\mathfrak C_1).
\tag{23}
\end{equation}

De (23) et de l'indépendance linéaire des $\xi$ et des $\zeta$, le Théorème VII pour $i=2$ résulte comme suit. À cause de la divisibilité de $\mathfrak C_1$ par $\Mmod_1$, on a d'abord $\Gmod_1\mid\Mmod_1=\Gmod^*_1\mid\Mmod_1$. Pour démontrer l'isomorphie avec $\Gmod^*_1\mid\Mmod^*_1$, supposons qu'un système de représentants de $\Gmod^*_1\mid\Mmod^*_1$ soit donné par certaines formes linéaires $g^*(\xi)$ de $\Gmod^*_1$. En raison de l'indépendance linéaire des $\xi$ et des $\zeta$, on déduit alors de $g^*(\xi)\equiv 0(\Mmod_1)$ que $g^*(\xi)\equiv 0(\Mmod^*_1)$; les $g^*(\xi)$ sont donc aussi incongrus modulo $\Mmod_1$, forment un système de représentants de $\Gmod^*_1\mid\Mmod_1$, et la correspondance de $\Gmod^*_1\mid\Mmod_1$ vers $\Gmod^*_1\mid\Mmod^*_1$ transmise par ce système de représentants est isomorphe d'après la Définition V.

Tout à fait de même, on obtient: de $b^{(2)}g(\xi)\equiv0(\Mmod_1)$, $b^{(2)}\ne0$, on déduit que $b^{(2)}g(\xi)\equiv0(\Mmod^*_1)$, $b^{(2)}\ne0$. Comme cela est satisfait pour tous les $g(\xi)$ de $\Gmod^*_1$ et seulement pour ceux-ci -- car $\Gmod^*_1$ consiste, d'après (23), en la totalité des polynômes de l'idéal fondamental $\gideal_1$ qui sont en même temps des formes linéaires en $\xi$ --, $\Gmod^*_1$ est ainsi caractérisé comme module fondamental de $\Mmod^*_1$.

Si l'on adjoint enfin $x_3\ldots x_n$ à $\mathsf{P}(u)$, on obtient:
\begin{equation}
\left\{\begin{aligned}
 \Gmod^*_1&=(\eta_1\ldots\eta_\rho);&
 \Mmod^*_1&=(e_1\eta_1,\ldots,e_\rho\eta_\rho);\\
 \Gmod_1&=(\eta_\rho\ldots\eta_1,\zeta_0\ldots\zeta_\nu\ldots);&
 \Mmod_1&=(e_\rho\eta_\rho,\ldots,e_1\eta_1,\zeta_0\ldots\zeta_\nu\ldots);
\end{aligned}\right.
\tag{24}
\end{equation}
et donc:
\[
 (e_\rho)=\Mmod_1/\Gmod_1=\Mmod_1/(\Mmod_1,\eta_\rho),\qquad
 (e_{\rho-1})=(\Mmod_1,\eta_\rho)/\Gmod_1\quad\hbox{etc.}
\]
Ainsi, en tenant compte de la note 8), $e_\rho,e_{\rho-1},\ldots,e_1,1,\ldots,1,\ldots$ sont reconnus comme les diviseurs élémentaires de $\Mmod_1$. Le Théorème VII est donc démontré en toutes ses parties pour $i=2$.

Il faut remarquer que (22) et (23) n'utilisent que le fait que $C^{(1)}(x)$ est régulier en $x_1$. Ces formules, et les considérations qui s'y rattachent et conduisent au Théorème VII, restent donc valables lorsque, à la place de (12), on prend pour base la transformation spéciale
\begin{equation}
 y_1=x_1;\quad y_2=u_{21}x_1+z_2;\quad\ldots;\quad y_n=u_{n1}x_1+z_n
\tag{25}
\end{equation}
qui, composée avec
\begin{equation}
\left\{\begin{array}{l}
 z=U_1(x)\quad\hbox{ou}\\[2mm]
 z_2=x_2;\quad z_3=u_{32}x_2+x_3;\quad\ldots;\quad z_n=u_{n2}x_2+\cdots+u_{n,n-1}x_{n-1}+x_n
\end{array}\right.
\tag{26}
\end{equation}
redonne (12). Si $\overline{\mideal}_{(u)}$ passe, en vertu de (25), en $\widetilde{\mideal}$, et si $\widetilde{\gideal}_1,\widetilde{\Gmod}_1\ldots$ ont pour $\widetilde{\mideal}$ la même signification que $\gideal_1,\Gmod_1\ldots$ ont pour $\mideal$, alors on a donc:
\begin{equation}
 \widetilde{\Gmod}_1=(\widetilde{\Gmod}^{*}_1,\widetilde{\mathfrak C}_1),\qquad
 \widetilde{\Mmod}_1=(\widetilde{\Mmod}^{*}_1,\widetilde{\mathfrak C}_1).
\tag{27}
\end{equation}

Ici (27) naît de (23) simplement par l'inversion de (26). Car cela est manifestement vérifié pour $\mideal$ et $C^{(1)}(x)$, donc aussi pour $\widetilde{\mathfrak C}_1$ et $\widetilde\Mmod_1$. Cela vaut aussi pour $\gideal_1$, puisque
\[
 b^{(2)}(x)G(x)\equiv0(\mideal),\quad b^{(2)}\ne0,
 \qquad
 \widetilde b^{(2)}(z)\widetilde G(x,z)\equiv0(\widetilde\mideal),\quad \widetilde b^{(2)}\ne0
\]
s'entraînent réciproquement en vertu de $z=U_1(x)$. Ainsi $\widetilde\gideal_1=\gideal_1$ en vertu de (26); donc aussi $\widetilde\Gmod_1=\Gmod_1$, et, par suite, d'après la définition de $\Gmod_1^*$ et de $\Mmod_1^*$ donnée par (22), il en va de même pour $\widetilde\Gmod_1^*$ et $\widetilde\Mmod_1^*$, et donc pour toute la représentation (27).

Pour la conclusion inductive générale, supposons satisfaites les hypothèses:
\begin{equation}
\left\{\begin{aligned}
 \Gmod_{i-1}&=(\Gmod_{i-1}^*,\mathfrak C_{i-1}),&
 \Mmod_{i-1}&=(\Mmod_{i-1}^*,\mathfrak C_{i-1}),\\
 \mathfrak C_{i-1}&=(\zeta_0,\zeta_1,\ldots,\zeta_\nu,\ldots),
\end{aligned}\right.
\tag{28}
\end{equation}
où $\Gmod^*_{i-1}$ et $\Mmod^*_{i-1}$ désignent des modules -- relativement aux grandeurs entières $b^{(i)},c^{(i)},\ldots$ -- de formes linéaires en un nombre fini d'indéterminées $\xi$, certains produits de puissances de $x_1,\ldots,x_{i-1}$; et où les $\zeta$ sont linéairement indépendants, aussi bien entre eux que vis-à-vis des $\xi$, relativement à ces grandeurs entières. Une représentation correspondant à (28), ainsi que l'indépendance linéaire des $\xi$ et des $\zeta$, doit également valoir lorsqu'à la place de (12) on prend pour base la transformation spéciale
\begin{equation}
\begin{aligned}
 y_1&=x_1;\quad \ldots;\quad y_{i-1}=u_{i-1,1}x_1+\cdots+x_{i-1};\\
 y_i&=u_{i,1}x_1+\cdots+u_{i,i-1}x_{i-1}+z_i;\quad\ldots;\quad
 y_n=u_{n,1}x_1+\cdots+u_{n,i-1}x_{i-1}+z_n
\end{aligned}
\tag{29}
\end{equation}
qui peut être composée pour donner (12) au moyen de
\[
 z=U_{i-1}(x)\quad\hbox{ou}\quad
 z_i=x_i;\quad z_{i+1}=u_{i+1,i}x_i+x_{i+1};\quad\ldots;\quad
 z_n=u_{n,i}x_i+\cdots+u_{n,n-1}x_{n-1}+x_n,
\]
et cette représentation spéciale doit sortir de (28) simplement par l'inversion de $z=U_{i-1}(x)$.

Des hypothèses (28) et de l'indépendance linéaire des $\xi$ et des $\zeta$ résulte l'isomorphie de $\Gmod_{i-1}\mid\Mmod_{i-1}$ avec $\Gmod^*_{i-1}\mid\Mmod^*_{i-1}$, par correspondance avec le même système de représentants, au moyen des mêmes considérations exactement que celles qui se rattachaient ci-dessus à (23); de même, il en résulte que $\Gmod^*_{i-1}$ devient module fondamental de $\Mmod^*_{i-1}$, et que $\Mmod_{i-1}$ ne possède qu'un nombre fini de diviseurs élémentaires différents de $1$, à savoir les diviseurs élémentaires de $\Mmod^*_{i-1}$ qui sont différents de $1$.

Pour effectuer la conclusion inductive, il faut d'abord montrer de nouveau que, pour $\Gmod^*_{i-1}\mid\Mmod^*_{i-1}$, donc aussi pour $\Gmod_{i-1}\mid\Mmod_{i-1}$, on peut prendre un système de représentants borné en $x_i$. Cela résulte de la représentation du quotient démontrée en (7), Théorème III:
\begin{equation}
 \Gmod^*_{i-1}=\Mmod^*_{i-1}/(C^{(i)}),
\tag{30}
\end{equation}
où $C^{(i)}$ désigne un déterminant quelconque, non identiquement nul, de degré $\rho$ de $\Mmod^*_{i-1}$, $\rho$ étant le rang de $\Mmod^*_{i-1}$. De la partie des hypothèses relative à la transformation spéciale (29), il résulte que $C^{(i)}$ peut toujours être supposé régulier relativement à $x_i$. En effet, le module $\widetilde\Mmod^*_{i-1}$ correspondant à $\Mmod^*_{i-1}$ a le même rang; si donc $\widetilde C^{(i)}$ désigne un déterminant quelconque, non identiquement nul, de degré $\rho$ de $\widetilde\Mmod^*_{i-1}$ qui, par $z=U_{i-1}(x)$, passe en un $C^{(i)}$ régulier, alors, par hypothèse, $C^{(i)}$ devient un déterminant de degré $\rho$ de $\Mmod^*_{i-1}$.

L'existence de $C^{(i)}$ implique, comme en (8), avec une numérotation convenable des $\xi$, l'appartenance à $\Mmod^*_{i-1}$ de $\rho$ formes linéaires de la forme spéciale
\[
 L_\lambda(\xi)=C^{(i)}\xi_\lambda+c^{(i)}_{\lambda,1}\xi_{\rho+1}+\cdots+c^{(i)}_{\lambda,s-\rho}\xi_s
 \qquad(\lambda=1\ldots\rho).
\]
Toute forme linéaire en $\xi$ peut alors être ramenée à la forme:
\begin{equation}
 H(\xi)\equiv a^{(i)}_1\xi_1+\cdots+a^{(i)}_\rho\xi_\rho+b^{(i)}_1\xi_{\rho+1}+\cdots+b^{(i)}_{s-\rho}\xi_s\quad (L_1(\xi),\ldots,L_\rho(\xi)),
\tag{31}
\end{equation}
où $a^{(i)}_1\ldots a^{(i)}_\rho$ sont bornés en $x_i$ et n'atteignent pas le degré $k$ de $C^{(i)}$. Cette représentation est unique; car de
\[
 a^{(i)}_1\xi_1+\cdots+a^{(i)}_\rho\xi_\rho+b^{(i)}_1\xi_{\rho+1}+\cdots+b^{(i)}_{s-\rho}\xi_s\equiv0\quad (L_1(\xi),\ldots,L_\rho(\xi))
\]
on déduit, par comparaison des coefficients de $\xi_1\ldots\xi_\rho$ et en tenant compte des degrés en $x_i$, l'annulation identique de tous les $a^{(i)}$ et $b^{(i)}$.

Soit maintenant en particulier
\[
 H(\xi)\equiv0(G^*_{i-1});\qquad\hbox{donc}\qquad C^{(i)}H(\xi)\equiv0(\Mmod^*_{i-1})\quad\hbox{d'après (30).}
\]
On obtient alors, comme dans la démonstration du Théorème III:
\[
 C^{(i)}H(\xi)\equiv\sum_{\tau=1}^{s-\rho}\{C^{(i)}b^{(i)}_\tau-\sum_\lambda a^{(i)}_\lambda c^{(i)}_{\lambda,\tau}\}\xi_{\rho+\tau}
 \equiv0(\Mmod^*_{i-1}),
\]
et donc, puisque, comme dans le Théorème III, les coefficients de $\xi_{\rho+1}\ldots\xi_s$ doivent s'annuler:
\[
 C^{(i)}b^{(i)}_\tau=\sum_\lambda a^{(i)}_\lambda c^{(i)}_{\lambda,\tau}\qquad(\tau=1\ldots s-\rho),
\]
d'où il résulte que les $b^{(i)}_\tau$ sont eux aussi bornés et n'atteignent pas le degré maximal des $c^{(i)}_{\lambda,\tau}$. L'existence d'un système de représentants borné en $x_i$ pour $\Gmod^*_{i-1}\mid\Mmod^*_{i-1}$, c'est-à-dire pour $\Gmod_{i-1}\mid\Mmod_{i-1}$, qui n'atteint pas un certain degré fixe $r$, est ainsi démontrée.

Désignons maintenant par $\vartheta_1\ldots\vartheta_t$ les produits de puissances en nombre fini
\[
 x_i^\alpha\xi_\lambda\quad(\alpha=0,1,\ldots,k-1;\ \lambda=1,2,\ldots,\rho),\qquad
 x_i^\beta\xi_{\rho+\tau}\quad(\beta=0,1,\ldots,r-1;\ \tau=1,2,\ldots,s-\rho)
\]
et rangeons les expressions dénombrablement infinies
\[
 x_i^\gamma L_\lambda\quad(\gamma=0,1\ldots\text{ in inf.}\ldots;\ \lambda=1,2,\ldots,\rho),
 \qquad
 x_i^\gamma\zeta_\nu\quad(\gamma,\nu=0,1\ldots\text{ in inf.}\ldots)
\]
en une suite simplement infinie $\omega_0,\omega_1,\ldots,\omega_\mu\ldots$; alors les $\vartheta$ et les $\omega$ sont linéairement indépendants, aussi bien séparément entre eux que les uns des autres, relativement aux grandeurs entières $b^{(i+1)},c^{(i+1)},\ldots$. Car de
\[
 \sum c^{(i+1)}_{\alpha\lambda} x_i^\alpha\xi_\lambda+
 \sum c^{(i+1)}_{\beta\tau} x_i^\beta\xi_{\rho+\tau}+
 \sum d^{(i+1)}_{\gamma\lambda} x_i^\gamma L_\lambda(\xi)+
 \sum e^{(i+1)}_{\gamma\nu}x_i^\gamma\zeta_\nu=0,
\]
chaque somme ne portant que sur un nombre fini de termes, il résulte, à cause de l'indépendance linéaire supposée des $\xi$ et des $\zeta$, ainsi que des $\xi$ entre eux relativement aux grandeurs entières $b^{(i)},\ldots$, que tous les $e^{(i+1)}_{\gamma\nu}$ s'annulent. De l'unicité de la représentation (31), on déduit ensuite l'annulation des $c^{(i+1)}_{\alpha\lambda}$ et des $c^{(i+1)}_{\beta\tau}$, puis enfin, en raison de l'indépendance linéaire des $L_\lambda(\xi)$, l'annulation des $d^{(i+1)}_{\gamma\lambda}$.

Si l'on conçoit maintenant le module $(\mathfrak C_{i-1},L_1(\xi),\ldots,L_\rho(\xi))$ comme un module $\mathfrak C_i$ relativement aux grandeurs entières $b^{(i+1)}$, alors $\mathfrak C_i$ devient divisible par $\Mmod_i$, à cause de la divisibilité de $\mathfrak C_{i-1},L_1(\xi)\ldots L_\rho(\xi)$ par $\Mmod_{i-1}$, et l'on obtient
\[
 \mathfrak C_i=(\omega_0,\omega_1,\ldots,\omega_\mu\ldots).
\]
Pour tout $H(x)\equiv0(\gideal_{i-1})$, on a donc, d'après (28), (31) et ce qui vient d'être démontré:
\[
 H(x)\equiv b^{(i+1)}\vartheta_1+\cdots+b^{(i+1)}_t\vartheta_t\quad(\mathfrak C_i)
\]
comme équation de modules relativement aux grandeurs entières $b^{(i+1)},c^{(i+1)}$. Or $\gideal_i$ est divisible par $\gideal_{i-1}$, et $\mideal$ divisible par $\gideal_i$; si donc, en particulier, on pose pour tout $G(x)$ de $\mathfrak C_i$ et tout $F(x)$ de $\Mmod_i$:
\[
 G(x)\equiv g^{(i+1)}_1\vartheta_1+\cdots+g^{(i+1)}_t\vartheta_t\quad(\mathfrak C_i),\qquad
 F(x)\equiv a^{(i+1)}\vartheta_1+\cdots+a^{(i+1)}\vartheta_t\quad(\mathfrak C_i),
\]
et si le module formé de tous les $g(\vartheta)$ est désigné par $\Gmod_i^*$, celui formé de tous les $a(\vartheta)$ par $\Mmod_i^*$, les mêmes considérations exactement que celles qui conduisaient à (23) donnent:
\begin{equation}
 \Gmod_i=(\Gmod_i^*,\mathfrak C_i),\qquad
 \Mmod_i=(\Mmod_i^*,\mathfrak C_i),\qquad
 \mathfrak C_i=(\omega_0,\omega_1,\ldots,\omega_\mu\ldots).
\tag{32}
\end{equation}
Dans la dérivation de (32), seule la régularité de $C^{(i)}$ en $x_i$ a de nouveau été utilisée; toute la considération reste donc valable si l'on prend pour base la transformation spéciale (29) avec un indice augmenté d'une unité. Alors les mêmes considérations que celles qui se rattachent à (27) montrent que cette représentation spéciale sort de (32) simplement par l'inversion de $z=U_i(x)$.

Ainsi toutes les hypothèses (28) etc. de la conclusion inductive générale sont démontrées pour l'indice suivant; et puisque, comme cela a été montré là, le Théorème VII découle immédiatement de ces hypothèses, le Théorème VII est donc démontré en général. En même temps, il est montré que l'arbitraire de $\Gmod^*_{i-1},\Mmod^*_{i-1}$ donné par l'arbitraire des $C^{(i)}$ est à nouveau supprimé par l'isomorphie du système des classes de $\Gmod^*_{i-1}$ modulo $\Mmod^*_{i-1}$.

En particulier, si le corps $\mathsf{P}$ pris au \S~3 comme domaine de coefficients est de caractéristique nulle -- c'est-à-dire si le corps premier contenu dans $\mathsf{P}$ et dérivé de l'unité est du type des nombres rationnels --, les indéterminées $u_{\mu\nu}$ peuvent être spécialisées en des grandeurs $\bar u_{\mu\nu}$ de $\mathsf{P}$ de telle sorte que, pour $i=1\ldots n$, chaque $C^{(i)}$ reste régulier en $x_i$. Les considérations ci-dessus montrent que le Théorème VII subsiste aussi sous cette spécialisation.\footnote{Hentzelt prend, au lieu d'un $\mathsf{P}$ quelconque, seulement le corps de tous les nombres complexes et traite principalement ce dernier cas, tandis qu'il ne prend des indéterminées pour base que dans la théorie de l'élimination proprement dite. Hentzelt montre ensuite, essentiellement par les conclusions données ici, qu'il suffit, pour former la forme résultante, d'aller dans chaque $\Mmod_{i-1}$ seulement jusqu'à un certain degré fini dépassant une borne fixe dans les $x$. Il manque cependant l'isomorphie, donnée dans le Théorème VII, de $\Gmod_{i-1}\mid\Mmod_{i-1}$ avec $\Gmod^*_{i-1}\mid\Mmod^*_{i-1}$, qui expose la raison de cette indépendance à l'égard des nombres de degrés. (E. N.)} Le module fondamental et les diviseurs élémentaires n'ont cependant pas nécessairement à provenir du cas général par la spécialisation $u_{\mu\nu}=\bar u_{\mu\nu}$.\footnote{Pour $\mideal=(x^2,ux+y)$, on a $\gideal_0=(1)$, $\gideal_1=(1)$. Si l'on choisit $C^{(1)}=x^2$, alors $\Gmod_1^*=(1,x)$, $\Mmod_1^*=(ux+y,xy)$; ce dernier $\Mmod_1^*$ possède les diviseurs élémentaires $y^2$ et $1$, et l'on peut choisir $C^{(2)}=y^2$. Pour $u=0$, $C^{(1)}$ et $C^{(2)}$ restent réguliers respectivement en $x$ et en $y$; mais $\widetilde{\Mmod}_1^*=(y,xy)$ possède les diviseurs élémentaires $y,y$. Donc $\widetilde{\Gmod}_1\mid\widetilde{\Mmod}_1$ devient encore isomorphe au système des classes modulo un module fini, mais n'est pas isomorphe à $\Gmod_1\mid\Mmod_1$. Si, en revanche, on avait choisi $C^{(1)}=ux+y$ et spécialisé de façon que $C^{(1)}$ fût resté régulier, alors l'isomorphie aurait elle aussi été conservée. (E. N.)}

\clearpage
\NoetherPaperTwentyTwoSFourEnd
\endgroup
\begingroup
\providecommand{\Amod}{\mathfrak A}
\providecommand{\Bmod}{\mathfrak B}
\providecommand{\Gmod}{\mathfrak G}
\providecommand{\Mmod}{\mathfrak M}
\providecommand{\Nmod}{\mathfrak N}
\providecommand{\Rmod}{\mathfrak R}
\providecommand{\Smod}{\mathfrak S}
\providecommand{\mideal}{\mathfrak m}
\providecommand{\nideal}{\mathfrak n}
\providecommand{\gideal}{\mathfrak g}
\providecommand{\hideal}{\mathfrak h}
\providecommand{\jideal}{\mathfrak j}
\providecommand{\fraks}{\mathfrak s}
\providecommand{\tideal}{\mathfrak t}
\providecommand{\aideal}{\mathfrak a}
\providecommand{\bideal}{\mathfrak b}
\providecommand{\bb}{\mathfrak b}
\providecommand{\cc}{\mathfrak c}
\providecommand{\zero}[1]{\equiv 0\pmod{#1}}
\providecommand{\Norm}{N}
\def\NoetherPaperTwentyTwoSFiveBodyOnly{1}
\ifdefined\NoetherPaperTwentyTwoSFiveBodyOnly
\providecommand{\NoetherPaperTwentyTwoSFiveEnd}{}
\else
\documentclass[11pt]{article}
\usepackage[a4paper,margin=1.45cm]{geometry}
\usepackage[T1]{fontenc}
\usepackage[utf8]{inputenc}
\usepackage{lmodern}
\usepackage[french]{babel}
\usepackage{amsmath,amssymb,mathtools}
\usepackage[protrusion=true,expansion=false]{microtype}
\setlength{\parindent}{1.25em}
\setlength{\parskip}{0.10em}
\allowdisplaybreaks
\setlength{\emergencystretch}{12em}
\sloppy\hbadness=10000\vbadness=10000\hfuzz=2pt\vfuzz=10pt\raggedbottom
\makeatletter\renewcommand\footnotesize{\@setfontsize\footnotesize{7.5}{8.7}\selectfont}\makeatother
\providecommand{\Amod}{\mathfrak A}
\providecommand{\Bmod}{\mathfrak B}
\providecommand{\Gmod}{\mathfrak G}
\providecommand{\Mmod}{\mathfrak M}
\providecommand{\Nmod}{\mathfrak N}
\providecommand{\Rmod}{\mathfrak R}
\providecommand{\Smod}{\mathfrak S}
\providecommand{\mideal}{\mathfrak m}
\providecommand{\nideal}{\mathfrak n}
\providecommand{\gideal}{\mathfrak g}
\providecommand{\hideal}{\mathfrak h}
\providecommand{\jideal}{\mathfrak j}
\providecommand{\fraks}{\mathfrak s}
\providecommand{\tideal}{\mathfrak t}
\providecommand{\aideal}{\mathfrak a}
\providecommand{\bideal}{\mathfrak b}
\providecommand{\bb}{\mathfrak b}
\providecommand{\cc}{\mathfrak c}
\providecommand{\zero}[1]{\equiv 0\pmod{#1}}
\providecommand{\Norm}{N}

\begin{document}
\providecommand{\NoetherPaperTwentyTwoSFiveEnd}{\end{document}}
\fi

\setcounter{footnote}{13}

\section*{\S 5. Forme résultante et forme des diviseurs élémentaires d'un idéal de polynômes.}

Soient $x_{i+1}\ldots x_n$ adjoints à $\mathsf{P}(u)$, de sorte que $\Gmod_{i-1}$ et $\Mmod_{i-1}$, et par conséquent aussi $\Gmod^*_{i-1}$ et $\Mmod^*_{i-1}$, passent en modules de formes linéaires où les grandeurs entières peuvent être regardées comme des polynômes d'\emph{une} indéterminée $x_i$. D'après le Théorème VII, les diviseurs élémentaires de $\Mmod_{i-1}$ qui sont différents de $1$, ainsi qu'au plus un nombre fini de diviseurs élémentaires $1$ de $\Mmod_{i-1}$, coïncident alors avec ceux de $\Mmod^*_{i-1}$. Le produit de ces diviseurs élémentaires, le $\rho$-ième diviseur déterminantiel de $\Mmod^*_{i-1}$, était égal au déterminant de la substitution de passage de $\Gmod^*_{i-1}$ à $\Mmod^*_{i-1}$, donc égal à $N(\Gmod^*_{i-1}\mid\Mmod^*_{i-1})$ d'après la Définition III. D'après (24), respectivement d'après la formule analogue qui résulte de (28) pour $i$ général, il apparaît que le produit de ces diviseurs élémentaires devient aussi égal au déterminant de la substitution de passage de $\Gmod_{i-1}$ à $\Mmod_{i-1}$; il doit donc être désigné par $N(\Gmod_{i-1}\cdot\Mmod_{i-1})$. On obtient ainsi
\[
 N(\Gmod_{i-1}\mid\Mmod_{i-1})=N(\Gmod^*_{i-1}\mid\Mmod^*_{i-1})=R^{(i)}(x),
\]
où le polynôme $R^{(i)}(x)$, c'est-à-dire le plus grand commun diviseur, au sens polynomial, de tous les déterminants à $\rho$ lignes de $\Mmod^*_{i-1}$, peut être supposé entier et primitif en $x_{i+1}\ldots x_n$ et, par conséquent, comme diviseur de $C^{(i)}$, devient régulier en $x_i$. De même, le diviseur élémentaire le plus élevé $E^{(i)}(x)$ de $\Mmod_{i-1}$, respectivement de $\Mmod^*_{i-1}$, peut être supposé entier et primitif en $x_{i+1}\ldots x_n$, et par conséquent régulier en $x_i$. Cela conduit à la définition suivante.

\medskip
\noindent\textbf{Définition VI.} \emph{Le polynôme $R^{(i)}(x)$ est appelé la résultante de $i$-ième niveau de $\mideal$, et $E^{(i)}(x)$ le diviseur élémentaire de $i$-ième niveau. La résultante de $i$-ième niveau est donc la norme de l'idéal fondamental $\gideal_{i-1}$ par rapport à $\mideal$, conçue comme norme du module $\Gmod_{i-1}$ par rapport à $\Mmod_{i-1}$, où $x_{i+1}\ldots x_n$ sont adjoints à $\mathsf{P}(u)$; de même, $(E^{(i)}(x))$, sous la même adjonction, devient égal au quotient $\Mmod_{i-1}/\Gmod_{i-1}$. Comme forme résultante, respectivement forme des diviseurs élémentaires de $\mideal$, on désigne les produits suivants:}
\[
 R_\mideal=R^{(1)}\cdot R^{(2)}\cdot\ldots\cdot R^{(n)};\qquad
 E_\mideal=E^{(1)}\cdot E^{(2)}\cdot\ldots\cdot E^{(n)}.
\]

Pour la forme résultante et la forme des diviseurs élémentaires, on a:

\medskip
\noindent\textbf{Théorème VIII.} \emph{La forme résultante est divisible par la forme des diviseurs élémentaires; inversement, une puissance de $E_\mideal$ est divisible par $R_\mideal$. $E_\mideal$ et $R_\mideal$ sont divisibles par $\mideal$; plus généralement, $E^{(i)}\cdots E^{(n)}\gideal_{i-1}\equiv0(\mideal)$, et par conséquent $R^{(i)}\cdots R^{(n)}\gideal_{i-1}\equiv0(\mideal)$.}

En effet, comme la norme est divisible par le plus haut diviseur élémentaire et qu'inversement une puissance de ce plus haut diviseur élémentaire est divisible par la norme, on obtient cette divisibilité pour $R^{(i)}(x)$ et $E^{(i)}(x)$ après adjonction de $x_{i+1}\ldots x_n$. Mais puisque $R^{(i)}(x)$ et $E^{(i)}(x)$ sont supposés primitifs comme polynômes relativement à $x_{i+1}\ldots x_n$, la divisibilité vaut relativement à toutes les indéterminées $x_i,x_{i+1},\ldots,x_n$; ainsi vaut, d'après la définition de ces expressions relativement à toutes les indéterminées $x$, la divisibilité de $R_\mideal$ par $E_\mideal$, et d'une puissance de $E_\mideal$ par $R_\mideal$.

Le plus haut diviseur élémentaire $E^{(i)}(x)$ est en outre, d'après le Théorème VII et le Théorème II, défini, après adjonction de $x_{i+1}\ldots x_n$, comme le quotient $\Mmod_{i-1}/\Gmod_{i-1}$. Par conséquent, si l'on revient à des polynômes en toutes les indéterminées $x$, il existe pour tout
\begin{equation}
 G(x)\equiv0(\gideal_{i-1})
 \quad\hbox{un}\quad b^{(i+1)}\ne0,
 \quad\hbox{tel que}\quad b^{(i+1)}E^{(i)}(x)G(x)\equiv0(\mideal)
\tag{33}
\end{equation}
soit vérifié. Or $\gideal_0$ devient en particulier égal à l'idéal-unité, donc:
\[
 b^{(2)}E^{(1)}(x)\equiv0(\mideal);\quad b^{(2)}\ne0,
 \qquad\hbox{et donc}\qquad E^{(1)}(x)\equiv0(\gideal_1).
\]
En général, d'après (33), appliqué à
\[
 E^{(1)}(x)\cdots E^{(i-1)}(x)\equiv0(\gideal_{i-1}),
\]
il existe un $b^{(i+1)}\ne0$ tel que $b^{(i+1)}E^{(1)}\cdots E^{(i)}\equiv0(\mideal)$; donc $E^{(1)}\cdots E^{(i)}\equiv0(\gideal_i)$. Comme $\gideal_n=\mideal$, il vient ainsi
\begin{equation}
 E_\mideal=E^{(1)}\cdots E^{(n)}\equiv0(\mideal)
 \quad\hbox{et par conséquent}\quad
 R_\mideal=R^{(1)}\cdots R^{(n)}\equiv0(\mideal).
\tag{34}
\end{equation}

Si l'on fait parcourir à $G(x)$, dans (33), les polynômes en nombre fini d'une base d'idéal de $\gideal_{i-1}$, et si $c^{(i+1)}(x)$ désigne le produit des $b^{(i+1)}(x)$ qui leur correspondent, on obtient
\[
 c^{(i+1)}E^{(i)}(x)\gideal_{i-1}\equiv0(\mideal);
 \qquad c^{(i+1)}\ne0
 \quad\hbox{et donc}\quad E^{(i)}(x)\gideal_{i-1}\equiv0(\gideal_i),
\]
et de là, comme ci-dessus par répétition un nombre fini de fois:
\[
 E^{(n)}(x)\cdots E^{(i)}(x)\gideal_{i-1}\equiv0(\mideal)
\]
et par conséquent aussi
\[
 R^{(n)}(x)\cdots R^{(i)}(x)\gideal_{i-1}\equiv0(\mideal),
\]
ce qui démontre le Théorème VIII dans toutes ses parties.

Le Théorème VIII repose essentiellement sur des propriétés de la forme des diviseurs élémentaires; mais que la forme résultante ait une signification caractéristique pour $\mideal$ est montré par le théorème suivant.

\medskip
\noindent\textbf{Théorème IX.} \emph{Si $\nideal$ est un diviseur de $\mideal$ -- diviseur entendu au sens des idéaux -- et si les formes résultantes $R_\nideal$ et $R_\mideal$ coïncident, alors les idéaux $\nideal$ et $\mideal$ coïncident.}

Soient en effet $\gideal_0\ldots\gideal_{n-1},\gideal_n=\mideal$ et $\hideal_0\ldots\hideal_{n-1},\hideal_n=\nideal$ les idéaux fondamentaux de $\mideal$ et de $\nideal$. D'abord $\gideal_0$ et $\hideal_0$ deviennent égaux à l'idéal-unité, donc $\gideal_0=\hideal_0$. Supposons maintenant $\gideal_{i-1}=\hideal_{i-1}$, de sorte que $\Mmod_{i-1}$ et $\Nmod_{i-1}$ possèdent le même module fondamental $\Gmod_{i-1}$. L'hypothèse $R_\nideal=R_\mideal$ admet, après adjonction de $x_{i+1}\ldots x_n$, la forme
\[
 N(\Gmod_{i-1}\mid\Nmod_{i-1})=N(\Gmod_{i-1}\mid\Mmod_{i-1})
 \quad\hbox{ou encore}\quad
 N(\Gmod^*_{i-1}\mid\Nmod^*_{i-1})=N(\Gmod^*_{i-1}\mid\Mmod^*_{i-1});
\]
ici on a posé, conformément à (32): $\Nmod_{i-1}=(\Nmod^*_{i-1},\mathfrak C_{i-1})$, de sorte que $\Nmod^*_{i-1}$ devient donc aussi un module de formes linéaires en $\xi$, et que $\Gmod^*_{i-1}$ en est le module fondamental. En raison de la divisibilité de $\Mmod_{i-1}$ par $\Nmod_{i-1}$, qui entraîne la divisibilité de $\Mmod^*_{i-1}$ par $\Nmod^*_{i-1}$, il résulte alors du Théorème IV que, sous cette adjonction, les deux modules $\Nmod^*_{i-1}$ et $\Mmod^*_{i-1}$ coïncident, et par conséquent aussi $\Nmod_{i-1}$ et $\Mmod_{i-1}$. Mais l'adjonction de $x_{i+1}\ldots x_n$ signifie que l'on ajoute à $\Mmod_{i-1}$, respectivement à $\mideal$, tous les polynômes $G(x)$ tels que
\[
 b^{(i+2)}(x)G(x)\equiv0(\mideal);
 \qquad b^{(i+1)}\ne0
\]
soit vérifié, c'est-à-dire que, par l'adjonction, $\Mmod_{i-1}$, respectivement $\mideal$, passe en $\gideal_i$, et de façon correspondante $\nideal$ en $\hideal_i$; ainsi $\gideal_i=\hideal_i$, et par conséquent $\gideal_n=\hideal_n$, donc $\mideal=\nideal$ est démontré.

Enfin, le même principe de transfert, appliqué au Théorème V, donne encore:

\medskip
\noindent\textbf{Théorème X.} \emph{Soit $R^{(i)}=S^{(i)}T^{(i)}$ une décomposition de $R^{(i)}$ en deux polynômes $S^{(i)}$ et $T^{(i)}$ premiers entre eux -- au sens polynomial. Alors il existe un et un seul idéal $\fraks$, et de même un et un seul idéal $\tideal$, tous deux diviseurs de $\mideal$ avec même idéal fondamental de niveau $(i-1)$, $\gideal_{i-1}$, tels que $S^{(i)}$ devienne égal à $N(\Gmod_{i-1}\mid\Smod_{i-1})$ et $T^{(i)}$ égal à $N(\Gmod_{i-1}\mid\mathfrak T_{i-1})$. Les idéaux $\fraks$ et $\tideal$ sont définis comme la totalité des polynômes $S(x)$, respectivement $T(x)$, tels que}
\[
 s^{(i+1)}(x)T^{(i)}(x)S(x)\equiv0(\mideal);\quad s^{(i+1)}\ne0,
\]
\[
 t^{(i+1)}(x)S^{(i)}(x)T(x)\equiv0(\mideal);\quad t^{(i+1)}\ne0
\]
\emph{soit vérifié, et $\gideal_i$ devient le plus petit commun multiple de $\fraks$ et $\tideal$,}
\[
 \gideal_i=[\fraks,\tideal].
\]

Il reste seulement à remarquer que $s^{(i+1)},t^{(i+1)}$ peuvent être choisis fixes pour $\fraks$ et $\tideal$, comme produits des $s^{(i+1)},t^{(i+1)}$ correspondant aux éléments de base de $\fraks$, respectivement de $\tideal$, et que, comme on l'a remarqué ci-dessus, l'adjonction de $x_{i+1}\ldots x_n$ transforme $\mideal$ en $\gideal_i$. En particulier, la décomposition de $R^{(i)}$ peut être poursuivie jusqu'aux puissances de polynômes irréductibles -- facteurs primaires --, ce qui entraîne une représentation correspondante de $\gideal_i$ comme plus petit commun multiple.

\clearpage
\NoetherPaperTwentyTwoSFiveEnd
\endgroup
\begingroup
\providecommand{\Amod}{\mathfrak A}
\providecommand{\Bmod}{\mathfrak B}
\providecommand{\Gmod}{\mathfrak G}
\providecommand{\Mmod}{\mathfrak M}
\providecommand{\Nmod}{\mathfrak N}
\providecommand{\Rmod}{\mathfrak R}
\providecommand{\Smod}{\mathfrak S}
\providecommand{\mideal}{\mathfrak m}
\providecommand{\nideal}{\mathfrak n}
\providecommand{\gideal}{\mathfrak g}
\providecommand{\hideal}{\mathfrak h}
\providecommand{\jideal}{\mathfrak j}
\providecommand{\fraks}{\mathfrak s}
\providecommand{\tideal}{\mathfrak t}
\providecommand{\aideal}{\mathfrak a}
\providecommand{\bideal}{\mathfrak b}
\providecommand{\bb}{\mathfrak b}
\providecommand{\cc}{\mathfrak c}
\providecommand{\zero}[1]{\equiv 0\pmod{#1}}
\providecommand{\Norm}{N}
\def\NoetherPaperTwentyTwoSSixSevenBodyOnly{1}
\ifdefined\NoetherPaperTwentyTwoSSixSevenBodyOnly
\providecommand{\NoetherPaperTwentyTwoSSixSevenEnd}{}
\else
\documentclass[11pt]{article}
\usepackage[a4paper,margin=1.45cm]{geometry}
\usepackage[T1]{fontenc}
\usepackage[utf8]{inputenc}
\usepackage{lmodern}
\usepackage[french]{babel}
\usepackage{amsmath,amssymb,mathtools}
\usepackage[protrusion=true,expansion=false]{microtype}
\setlength{\parindent}{1.25em}
\setlength{\parskip}{0.10em}
\allowdisplaybreaks
\setlength{\emergencystretch}{12em}
\sloppy\hbadness=10000\vbadness=10000\hfuzz=2pt\vfuzz=10pt\raggedbottom
\makeatletter\renewcommand\footnotesize{\@setfontsize\footnotesize{7.5}{8.7}\selectfont}\makeatother
\providecommand{\Amod}{\mathfrak A}
\providecommand{\Bmod}{\mathfrak B}
\providecommand{\Gmod}{\mathfrak G}
\providecommand{\Mmod}{\mathfrak M}
\providecommand{\Nmod}{\mathfrak N}
\providecommand{\Rmod}{\mathfrak R}
\providecommand{\Smod}{\mathfrak S}
\providecommand{\mideal}{\mathfrak m}
\providecommand{\nideal}{\mathfrak n}
\providecommand{\gideal}{\mathfrak g}
\providecommand{\hideal}{\mathfrak h}
\providecommand{\jideal}{\mathfrak j}
\providecommand{\fraks}{\mathfrak s}
\providecommand{\tideal}{\mathfrak t}
\providecommand{\aideal}{\mathfrak a}
\providecommand{\bideal}{\mathfrak b}
\providecommand{\bb}{\mathfrak b}
\providecommand{\cc}{\mathfrak c}
\providecommand{\zero}[1]{\equiv 0\pmod{#1}}
\providecommand{\Norm}{N}

\begin{document}
\providecommand{\NoetherPaperTwentyTwoSSixSevenEnd}{\end{document}}
\fi

\setcounter{footnote}{13}

\section*{\S 6. Théorie de l'élimination proprement dite.}

D'après le Théorème VIII, respectivement la formule (34), la forme résultante et la forme des diviseurs élémentaires sont divisibles par $\mideal$, et s'annulent donc pour tous les zéros de $\mideal$; par ``zéros de $\mideal$'', on entend ici les systèmes de valeurs de $x_1\ldots x_i$ qui appartiennent au corps algébriquement clos dérivé de $\mathsf{P}(u;x_{i+1}\ldots x_n)$\footnote{Steinitz a montré dans son \emph{Algebraische Theorie der Körper} (J. f. M. 137 (1910), p.~167--309) que tout corps peut être étendu, essentiellement de façon unique, en un corps algébriquement clos, donnant ainsi l'équivalent rationnel du théorème fondamental de l'algèbre. En fait, pour chaque $i=1,\ldots,n$, il s'agit d'un corps d'extension fini de $\mathsf{P}(u,x_{i+1},\ldots,x_n)$. (E. N.)} et tels que, pour ces systèmes de valeurs, tout polynôme de $\mideal$ s'annule, tandis que l'on pense $x_{i+1}\ldots x_n$ comme adjoints au domaine de coefficients $(i=1,\ldots,n)$. Ces $x_{i+1}\ldots x_n$ adjoints pourront, comme on le montrera, être aussi remplacés par des grandeurs de $\mathsf{P}(u)$. Le contenu de la théorie de l'élimination est inversement la déduction des zéros de $\mideal$ à partir de ceux de la forme résultante. Cette réciproque repose sur l'élimination successive à développer dans ce paragraphe; comme, dans le paragraphe suivant, on démontrera la divisibilité de la forme $D^{(i)}(x)$ qui y intervient par $E^{(i)}(x)$, la réciproque désirée suivra de la divisibilité d'une puissance de $E^{(i)}(x)$ par $R^{(i)}(x)$.

La théorie de l'élimination successive à donner ici repose sur le résultat suivant.

\medskip
\noindent\textbf{Théorème XI.} \emph{Soit $\bideal$ un idéal de polynômes en une indéterminée $t$ à coefficients dans $\mathsf{P}$, donc un idéal principal; soit $h(t)$ un polynôme de $\bideal$ de degré $k$, et soit $\Bmod$ le module de formes linéaires en $1,t,\ldots,t^{k-1}$ dans lequel $\bideal$ passe modulo $h(t)$. Alors $\Bmod$ est de rang $k-p$, où $p$ désigne le degré du polynôme générateur $f(t)$ de $\bideal$.}

Par hypothèse, $h(t)=h_1(t)\cdot f(t)$, où $h_1(t)$ est de degré $k-p$. Les classes de $\bideal$ modulo $h(t)$ représentées par les polynômes $f(t),tf(t),\ldots,t^{k-p-1}f(t)$ sont donc linéairement indépendantes; et chaque classe de $\bideal$ modulo $h(t)$ se représente linéairement, avec coefficients dans $\mathsf{P}$, au moyen de ces $k-p$ classes spéciales. Mais modulo $h(t)$, $\bideal$ passe en un module de formes linéaires en $1,t,\ldots,t^{k-1}$, dont le rang est par conséquent exactement $k-p$.

Soit maintenant $\aideal=\aideal_1$ un idéal transformé de polynômes, soit $A^{(1)}(x)$ un polynôme de $\aideal_1$ régulier en $x_1$ et de degré $k_1$, et soit $\mathfrak A_1$ le module de formes linéaires en $1,x_1,\ldots,x_1^{k_1-1}$, avec polynômes $a^{(2)}(x)$ comme coefficients, dans lequel $\aideal_1$ passe modulo $A^{(1)}(x)$. Que $\aideal_2$ désigne en outre l'idéal en $x_2\ldots x_n$ dérivé de tous les déterminants à $k_1$ lignes de $\mathfrak A_1$. D'après le Théorème XI, $\aideal_2$ est alors, et seulement alors, égal à l'idéal nul lorsque les polynômes de $\aideal_1$ possèdent un diviseur commun -- diviseur entendu au sens polynomial -- de degré $p_1>0$ en $x_1$. Si $\aideal_2$ est différent de l'idéal nul, il contient, puisque $\aideal$ était supposé transformé, un polynôme $A^{(2)}(x)$ régulier en $x_2$ et de degré $k_2$, comme on le voit directement en composant la transformation (12) à partir des transformations spéciales (25) et (26). Que $\mathfrak A_2$ désigne de nouveau le module de formes linéaires en $1,x_2,\ldots,x_2^{k_2-1}$ dans lequel $\aideal_2$ passe modulo $A^{(2)}(x)$, et $\aideal_3$ l'idéal en $x_3\ldots x_n$ dérivé de tous les déterminants à $k_2$ lignes de $\mathfrak A_2$; cet idéal doit contenir un polynôme $A^{(3)}(x)$ régulier en $x_3$.

On poursuit le procédé de cette manière jusqu'à ce que l'on arrive à un idéal nul ou jusqu'à ce que $\aideal_{n+1}$ devienne l'idéal-unité; car, comme $\aideal_{n+1}$ est libre des $x$, il ne peut être que l'idéal nul ou l'idéal-unité. Il apparaît que les idéaux $\aideal_1,\aideal_2,\ldots$ sont déterminés univoquement par $\aideal$, indépendamment des polynômes $A^{(\lambda)}(x)$ pris pour base. Cela est clair pour $\aideal_1=\aideal$ et peut donc être supposé pour $\aideal_1,\ldots,\aideal_\lambda$. Si maintenant $\aideal_\lambda$ passe modulo $A^{(\lambda)}(x)$ dans le module de formes linéaires $\mathfrak A_\lambda$, alors $\mathfrak A_\lambda$ peut aussi être caractérisé comme la totalité des polynômes de $\aideal_\lambda$ qui n'atteignent pas le degré $k_\lambda$ en $x_\lambda$; $\mathfrak A_\lambda$ ne dépend donc que du degré $k_\lambda$ de $A^{(\lambda)}(x)$. Soit $\bar A^{(\lambda)}(x)$ un polynôme de $\aideal_\lambda$, régulier en $x_\lambda$ et de degré $\bar k_\lambda>k_\lambda$, soit $\overline{\mathfrak A}_\lambda$ le module de formes linéaires correspondant, et soit $\bar\aideal_{\lambda+1}$ l'idéal dérivé des déterminants à $\bar k_\lambda$ lignes de $\overline{\mathfrak A}_\lambda$. Alors on a l'équation de modules
\[
 \overline{\mathfrak A}_\lambda=(\mathfrak A_\lambda,A^{(\lambda)},x_\lambda A^{(\lambda)},\ldots,x^{\bar k_\lambda-k_\lambda-1}A^{(\lambda)}).
\]
Comme, en raison de la régularité de $A^{(\lambda)}(x)$ en $x_\lambda$, les polynômes $A^{(\lambda)},x_\lambda A^{(\lambda)},\ldots$ peuvent être introduits comme nouvelles indéterminées des formes linéaires, il s'ensuit que les déterminants à $k_\lambda$ lignes de $\mathfrak A_\lambda$ coïncident avec ceux à $\bar k_\lambda$ lignes de $\overline{\mathfrak A}_\lambda$; ainsi $\bar\aideal_{\lambda+1}=\aideal_{\lambda+1}$.\footnote{C'est, en principe, la même conclusion que celle sur laquelle reposait l'isomorphie de $\Gmod^*_{i-1}\mid\Mmod^*_{i-1}$ avec $\Gmod_{i-1}\mid\Mmod_{i-1}$.}

En outre, $\aideal_{\lambda+1}$ est toujours divisible par $\aideal_\lambda$. C'est clair si $\aideal_{\lambda+1}$ est l'idéal nul; dans le cas contraire, $\mathfrak A_\lambda$ est de rang $k_\lambda$, et $\aideal_{\lambda+1}$ est donc l'idéal qui, par définition, se compose des déterminants à $k_\lambda$ lignes de $\mathfrak A_\lambda$, et qui est divisible par $\mathfrak A_\lambda$, donc par $\aideal_\lambda$. Chaque $\aideal_{\lambda+1}$ est ainsi divisible par $\aideal$; si donc $\aideal_{n+1}$ devient l'idéal-unité, alors $\aideal$ lui-même est l'idéal-unité.

Soit maintenant $\aideal$ différent de l'idéal-unité; $\aideal_1\ne0,\ldots,\aideal_i\ne0;\ \aideal_{i+1}=0$. Que $D^{(i)}(x)$ soit le plus grand commun diviseur -- diviseur entendu au sens polynomial -- de tous les polynômes de $\aideal_i$, dont l'existence est donnée par le Théorème XI; comme diviseur de $A^{(i)}(x)$, $D^{(i)}(x)$ devient lui-même régulier en $x_i$ et de degré $p_i>0$. On adjoint $x_{i+1}\ldots x_n$ à $\mathsf{P}(u)$; alors $D^{(i)}(x)$ admet, dans un corps d'extension algébrique de $\mathsf{P}(u;x_{i+1}\ldots x_n)$, une décomposition en $p_i$ facteurs linéaires. Soit $x_i-\bar x_i$ un tel facteur linéaire, de sorte que, pour $x_i=\bar x_i$, tous les polynômes de $\aideal_i$ s'annulent, et que par conséquent, d'après le Théorème XI, les polynômes de $\aideal_{i-1}$ reçoivent, sous cette spécialisation, un plus grand commun diviseur $D^{(i-1)}(x)$. Comme diviseur de $A^{(i-1)}(x)$, $D^{(i-1)}(x)$ devient lui-même de nouveau régulier en $x_{i-1}$ et de degré $p_{i-1}>0$; il ne peut donc pas, en particulier, s'annuler identiquement, et il admet dans un corps d'extension de $\mathsf{P}(u,\bar x_i,x_{i+1},\ldots,x_n)$ une décomposition en $p_{i-1}$ facteurs linéaires. Si $x_{i-1}-\bar x_{i-1}$ est un tel facteur linéaire, il lui correspond un plus grand commun diviseur provenant de $\aideal_{i-2}$. En continuant ainsi, on parvient à un nombre fini de zéros communs de $\aideal$, situés dans un corps d'extension algébrique de $\mathsf{P}(u;x_{i+1}\ldots x_n)$. Inversement, en raison de la divisibilité de $\aideal_\lambda$ par $\aideal$, tout zéro de $\aideal$ est aussi un zéro de $\aideal_\lambda$. Si donc, en particulier, $x_1=\bar x_1,\ldots,x_i=\bar x_i,x_{i+1}=x_{i+1},\ldots,x_n=x_n$ est un zéro de $\aideal$, on en déduit que $\aideal_{i+1}=0$; et les polynômes de $\aideal_i$ reçoivent un plus grand commun diviseur avec le facteur linéaire $x_i-\bar x_i$. En résumé:

\medskip
\noindent\textbf{Théorème XII.} \emph{Soit $\aideal$ différent de l'idéal-unité, $\aideal_1\ne0\ldots;\ \aideal_i\ne0;\ \aideal_{i+1}=0$. Alors, après adjonction de $x_{i+1}\ldots x_n$ à $\mathsf{P}(u)$, il existe un nombre fini de systèmes de valeurs associés $\bar x_1\ldots\bar x_i$, situés dans un corps d'extension algébrique de $\mathsf{P}(u,x_{i+1},\ldots,x_n)$, tels que tous les polynômes de $\aideal$ s'annulent pour $x_1=\bar x_1,\ldots,x_i=\bar x_i,x_{i+1}=x_{i+1},\ldots,x_n=x_n$. Inversement, si un tel zéro de $\aideal$ est présent, il en résulte $\aideal_{i+1}=0$, ainsi que l'apparition d'un diviseur commun $D^{(i)}(x)$ de tous les polynômes de $\aideal_i$, possédant le facteur linéaire $x_i-\bar x_i$.}

Le Théorème XII montre en particulier que tout idéal $\aideal$ qui ne possède aucun zéro devient l'idéal-unité.

Pour effectuer une décomposition qui montre aussi explicitement l'association des systèmes de valeurs, on introduit de nouvelles indéterminées $v$, que l'on adjoint à $\mathsf{P}(u,x_{i+1},\ldots,x_n)$. Posons:
\begin{equation}
 x_i=-v_1x_1-\cdots-v_{i-1}x_{i-1}+z_i,
\tag{35}
\end{equation}
de sorte que la composition de (12) avec (35) donne de nouveau une substitution du type (12), où seuls les $u_{i\kappa}$ sont maintenant remplacés par $w_{i\kappa}=u_{i\kappa}-v_\kappa$, et plus généralement $u_{i+\lambda,\kappa}$ par $w_{i+\lambda,\kappa}=u_{i+\lambda,\kappa}-u_{i+\lambda,i}v_\kappa$. Si donc la substitution (35) fait passer l'idéal $\aideal$, transformé par hypothèse, en $\bideal$, alors $\bideal$ naît simplement de $\aideal$ par remplacement des $u_{\mu\nu}$ par les $w_{\mu\nu}$ et par adjonction des $v$; et les idéaux $\bideal_1,\bideal_2,\ldots$ définis au moyen de $\bideal$ naissent par le même processus de $\aideal_1,\aideal_2,\ldots$. Inversement, $\aideal_\lambda$ provient de $\bideal_\lambda$ par $v=0$; les idéaux $\aideal_\lambda$ et $\bideal_\lambda$ sont donc simultanément des idéaux nuls, ou simultanément différents de l'idéal nul.

Soit maintenant de nouveau $\aideal_1\ne0;\ldots;\aideal_i\ne0;\aideal_{i+1}=0$; donc aussi $\bideal_1\ne0;\ldots;\bideal_i\ne0;\bideal_{i+1}=0$; et soit $\bar x_1\ldots\bar x_i,x_{i+1}\ldots x_n$ un système de zéros associés de $\aideal$. Si l'on définit $\bar z_i=\bar x_i+v_1\bar x_1+\cdots+v_{i-1}\bar x_{i-1}$, alors $\bar x_1\ldots\bar x_{i-1},\bar z_i,x_{i+1}\ldots x_n$ devient, en vertu de (35), un zéro de $\bideal$. Si donc $H^{(i)}(z,x)$ désigne le plus grand commun diviseur de tous les polynômes de $\bideal_i$, que l'on peut supposer entier et primitif en les $v$, alors $H^{(i)}(z,x)$ reçoit, d'après la dernière partie du Théorème XII, le facteur
\[
 z_i-\bar z_i=z_i-(\bar x_i+v_1\bar x_1+\cdots+v_{i-1}\bar x_{i-1}),
\]
et à chaque zéro de $\aideal$ qui apparaît après adjonction de $x_{i+1},\ldots,x_n$ correspond un tel facteur. Il faut montrer que la décomposition de $H^{(i)}(z,x)$ est ainsi épuisée, c'est-à-dire qu'on ne peut pas détacher de $H^{(i)}$ un facteur $H^{(i)}_1$ qui ne se décompose pas en facteurs linéaires en $z$ \emph{et} $v$. Supposons que ce soit le cas. Alors, en raison de la primitivité en $v$ supposée, $H^{(i)}$ contient au moins un facteur linéaire $z_i-\bar{\bar z}_i$, où $\bar{\bar z}_i$ appartient à un corps d'extension algébrique de $\mathsf{P}(u,v,x_{i+1}\ldots x_n)$. Si $\bar{\bar x}_1\ldots\bar{\bar x}_{i-1},\bar{\bar x}_i,x_{i+1}\ldots x_n$ est le zéro correspondant de $\aideal$, celui-ci doit coïncider avec l'un des zéros en nombre fini de $\aideal$ qui apparaissent après adjonction de $x_{i+1}\ldots x_n$, et donc être indépendant de $v$. Ainsi $\bar{\bar z}_i$ contient nécessairement les $v$ linéairement, non algébriquement ni rationnellement de façon non linéaire; contrairement à l'hypothèse, on peut donc détacher de $H^{(i)}_1$ un autre facteur linéaire $z_i-(\bar{\bar x}_i+v_1\bar x_1+\cdots+v_{i-1}\bar x_{i-1})$ en $z$ et $v$. On obtient donc la décomposition explicite:
\[
 H^{(i)}(z,x)=\prod \{z_i-(\bar x_{i\lambda}+v_1\bar x_{1\lambda}+\cdots+v_{i-1}\bar x_{i-1,\lambda})\}^{\alpha_\lambda}.
\]
Comme le degré de $H^{(i)}$ et les exposants individuels $\alpha_\lambda$ doivent coïncider avec ceux qui leur correspondent pour $D^{(i)}(x)$ -- car $H^{(i)}$ naît de $D^{(i)}$ par la spécialisation $u_{\mu\nu}=w_{\mu\nu}$, et $D^{(i)}$ de $H^{(i)}$ par la spécialisation $v=0$ --, cette décomposition montre aussi que, puisque les indéterminées $u_{\mu\nu}$ sont adjointes à $\mathsf{P}$, à chaque zéro $\bar x_i$ de $D^{(i)}$ apparaissant dans l'élimination à partir de $\aideal$ correspond un plus grand commun diviseur provenant de $\aideal_{i-1}\ldots\aideal_1$ qui est chaque fois linéaire, ou une puissance d'un tel diviseur; et donc, dans chaque cas, on n'obtient qu'un seul système de zéros associés $\bar x_1\ldots\bar x_i,x_{i+1}\ldots x_n$ de $\aideal$.\footnote{Il faut souligner que l'élimination successive donnée ici, contrairement à la méthode de Kronecker, ne dépend que de l'idéal donné $\aideal$ et est indépendante de toute base d'idéal; cela vaut en particulier aussi pour les exposants $\alpha_\lambda$. Une exécution complète de l'élimination par cette méthode exigerait, selon Hentzelt, la continuation du procédé sur le quotient $\aideal_i/D^{(i)}$, ce que l'on peut omettre en vue du paragraphe suivant. (E. N.)}

\section*{\S 7. Forme résultante et théorie de l'élimination.}

Pour établir le rapport entre la forme résultante et la théorie de l'élimination, on appliquera spécialement l'élimination successive donnée au \S~6 au quotient $\aideal=\mideal/\gideal_{i-1}$ de l'idéal par l'idéal fondamental de niveau $i-1$. Il faut d'abord montrer que ce quotient représente un idéal transformé. Or $\mideal$ est transformé et donc, d'après le Théorème VI du \S~3, $\gideal_{i-1}$ l'est également. De
\[
 H(x)\equiv0(\mideal/\gideal_{i-1});\qquad
 H(x)=\overline H(y)=\sum\alpha_\lambda(u)\overline H_\lambda(y)=\sum\alpha_\lambda(u)H_\lambda(x)
\]
on en déduit donc:
\[
 \overline H(y)\cdot\overline\gideal_{i-1(u)}\equiv0(\overline\mideal_{(u)});
 \quad\hbox{donc aussi}\quad
 \overline H(y)\overline\gideal_{i-1}\equiv0(\overline\mideal_{(u)}),
\]
ou
\[
 \overline H_\lambda(y)\overline\gideal_{i-1}\equiv0(\overline\mideal),
 \quad\hbox{donc}\quad H_\lambda(x)\gideal_{i-1}\equiv0(\mideal),
\]
ce qui démontre la divisibilité de $H_\lambda(x)$ par $\mideal/\gideal_{i-1}$, et par là que cet idéal est transformé; la méthode d'élimination du \S~6 devient donc applicable.

Mais (30), en tenant compte de (28) -- \S~4 --, montre que $C^{(i)}(x)$ est divisible par $\mideal/\gideal_{i-1}$, où $C^{(i)}(x)$ désigne un déterminant quelconque, non identiquement nul, de degré $\rho$ de $\Mmod^*_{i-1}$. Comme $C^{(i)}$, pour $i>1$, est de degré nul en $x_1$, le module $\mathfrak A_1$ correspondant à l'idéal $\aideal=\aideal_1$ contient donc les polynômes $C^{(i)},x_1C^{(i)},\ldots,x_1^{k_1-1}C^{(i)}$; par conséquent $(C^{(i)})^{k_1}$ est divisible par $\aideal_2$; de même $(C^{(i)})^{k_1k_2}$ devient divisible par $\aideal_3$, et finalement $(C^{(i)})^{k_1k_2\cdots k_{i-1}}$ divisible par $\aideal_i$. Ainsi $\aideal_1,\ldots,\aideal_i$ sont différents de l'idéal nul, ce qui vaut aussi pour $i=1$. Soit maintenant $D^{(i)}(x)$ le plus grand commun diviseur de tous les polynômes de $\aideal_i$, donc de degré $p_i>0$ seulement lorsque $\aideal_{i+1}$ devient l'idéal nul. D'après la définition de $D^{(i)}(x)$, en tenant compte de la divisibilité de $\aideal_i$ par $\aideal$, on obtient:
\[
 d^{(i+1)}D^{(i)}(x)\equiv0(\mideal/\gideal_{i-1});\qquad d^{(i+1)}\ne0;
\]
ainsi $d^{(i+1)}D^{(i)}$ devient un polynôme de $\mideal/\gideal_{i-1}$ libre de $x_1\ldots x_{i-1}$. Or $E^{(i)}(x)$ avait été défini, comme plus haut diviseur élémentaire de $\Mmod_{i-1}$, respectivement de $\Mmod^*_{i-1}$ (Théorème II et \S~4), comme le plus grand commun diviseur -- au sens polynomial -- de tous les polynômes de $\mideal/\gideal_{i-1}$ libres de $x_1\ldots x_{i-1}$. Ainsi $d^{(i+1)}D^{(i)}$, et, en raison de la régularité de $E^{(i)}(x)$ en $x_i$, aussi $D^{(i)}(x)$, devient divisible par $E^{(i)}(x)$.

Le même raisonnement s'applique lorsque, dès le départ, la transformation (12) avait été composée avec (35), c'est-à-dire lorsque les indéterminées $u_{\mu\nu}$ avaient été remplacées par $w_{\mu\nu}$; cela donne la divisibilité de $H^{(i)}(z,x)$ par le $\bar E^{(i)}(z,x)$ correspondant. Comme, de même que $D^{(i)}(x)$ et $H^{(i)}(z,x)$, les polynômes $E^{(i)}(x)$ et $\bar E^{(i)}(z,x)$, et de même $R^{(i)}(x)$ et $\bar R^{(i)}(z,x)$, proviennent réciproquement les uns des autres par spécialisation, les nombres de multiplicité dans la décomposition en facteurs linéaires doivent ici aussi coïncider réciproquement. En tenant encore compte du fait qu'une puissance de $E^{(i)}(x)$ est divisible par $R^{(i)}(x)$, on obtient en résumé:

\medskip
\noindent\textbf{Théorème XIII.} \emph{Si l'on décompose $R^{(i)}(x)$, dans un corps d'extension algébrique de $\mathsf{P}(u,x_{i+1}\ldots x_n)$, en facteurs linéaires $x_i-\bar x_i$, alors $\bar x_i$ se complète d'une et une seule manière en un système de valeurs associées $\bar x_1\ldots\bar x_i,x_{i+1}\ldots x_n$ qui devient zéro de $\mideal/\gideal_{i-1}$ et par conséquent de $\mideal$. Les zéros de $\mideal$ en nombre fini qui naissent ainsi des facteurs linéaires de $R^{(i)}$ -- en nombre fini après adjonction de $x_{i+1}\ldots x_n$ -- sont réunis par la décomposition explicite:}
\begin{equation}
 \bar R^{(i)}(z,x)=\prod\{z_i-(\bar x_{i\nu}+v_1\bar x_{1\nu}+\cdots+v_{i-1}\bar x_{i-1,\nu})\}^{\lambda_\nu}.
\tag{36}
\end{equation}
\emph{Cette décomposition subsiste, en raison de la régularité de $\bar R^{(i)}(z,x)$ en $z_i$, si les $x_{i+1}\ldots x_n$ adjoints à $\mathsf{P}(u)$ sont remplacés par des systèmes de valeurs spéciaux quelconques $\bar x_{i+1}\ldots\bar x_n$ de $\mathsf{P}(u)$ ou du corps algébriquement clos qui en est dérivé.}

On a ainsi effectué en même temps une séparation des zéros de $\mideal$ suivant les différents facteurs de la forme résultante; le nombre des indéterminées $x$ adjointes à $\mathsf{P}(u)$ est aussi appelé la dimension de la formation algébrique correspondante.

Si, dans la décomposition de $\bar R^{(i)}(z,x)$ ou dans celle correspondante de $R^{(i)}(x)$, on rassemble les facteurs conjugués relativement à $\mathsf{P}(u,x_{i+1}\ldots x_n)$, qui, pour un $\mathsf{P}$ général, peuvent être entièrement ou partiellement identiques, on obtient une décomposition de $R^{(i)}(x)$ en puissances de polynômes irréductibles relativement à $\mathsf{P}(u,x_{i+1}\ldots x_n)$:
\[
 R^{(i)}(x)=S^{(i)}_1(x)^{\beta_1}\cdots S^{(i)}_\varkappa(x)^{\beta_\varkappa}.
\]
À cette décomposition correspond, d'après le Théorème X, une représentation $\gideal_i=[\fraks_1\ldots\fraks_\varkappa]$, telle que $S^{(i)}_\mu(x)^{\beta_\mu}$ devient égal à la norme de $\gideal_{i-1}$ par rapport à l'idéal $\fraks_\mu$, respectivement à la norme du module $\Gmod_{i-1}$ par rapport au module correspondant à $\fraks_\mu$. Ainsi la signification des exposants $\beta_\mu$ est également éclaircie; en particulier -- si $\gamma_\mu$ désigne le degré de $S^{(i)}_\mu$ -- la multiplicité $\beta_\mu\gamma_\mu$ devient égale au nombre des classes de $\gideal_{i-1}$ modulo $\fraks_\mu$ qui sont linéairement indépendantes relativement à $\mathsf{P}(u,x_{i+1}\ldots x_n)$.

Considérons encore brièvement le cas spécial où $\mathsf{P}$ est de caractéristique nulle. Si l'on spécialise ici les $u_{\mu\nu}$ en des grandeurs $\bar u_{\mu\nu}$ de $\mathsf{P}$ de telle sorte que les $n$ déterminants $C^{(i)}(x)$ $(i=1,2,\ldots,n)$ restent chacun réguliers en $x_i$,\footnote{Hentzelt effectue d'emblée cette spécialisation des $u_{\mu\nu}$ et doit par conséquent recourir à des considérations beaucoup plus compliquées pour démontrer la décomposition de $H^{(i)}(z,x)$, respectivement de $\bar R^{(i)}(z.x)$, en facteurs linéaires en $z$ et $v$. Les exposants qui y apparaissent ne coïncident pas nécessairement avec ceux ci-dessus. (E. N.)} alors la régularité de $R^{(i)}$ est également conservée. Sous cette spécialisation, $R^{(i)}$ et de même $R^{(i)}(z,x)$ deviennent un diviseur commun de tous les déterminants à $\rho$ lignes de $\Mmod^*_{i-1}$, mais non nécessairement le plus grand commun diviseur. On peut cependant toujours effectuer la spécialisation de sorte que cette dernière propriété soit elle aussi conservée. Car $\bar R^{(i)}(z,x)$ peut -- comme le montre l'existence de la base du module -- être aussi caractérisé comme le plus grand commun diviseur d'un nombre fini de déterminants à $\rho$ lignes de $\Mmod^*_{i-1}$. La décomposition de $\bar R^{(i)}(z,x)$ ainsi spécialisé s'obtient alors simplement en remplaçant dans (36) les $u_{\mu\nu}$ par les $\bar u_{\mu\nu}$, de sorte que toute la théorie de l'élimination est conservée.

\begin{center}
(Reçu le 17.3.1922.)
\end{center}

\clearpage
\NoetherPaperTwentyTwoSSixSevenEnd
\endgroup
\clearpage
\setcounter{footnote}{0}
\providecommand{\mJ}{\mathfrak{J}}
\providecommand{\mS}{\mathfrak{S}}
\providecommand{\mo}{\mathfrak{o}}
\providecommand{\ma}{\mathfrak{a}}
\providecommand{\mb}{\mathfrak{b}}
\providecommand{\dd}{\mathrm d}
\providecommand{\D}{}
\renewcommand{\D}{\Delta}
\providecommand{\dx}{\dd x}
\providecommand{\dy}{\dd y}

\section*{\NoetherNavAnchor{P23}{23. Invariants algébriques et différentiels}23. Invariants algébriques et différentiels}
\begin{center}
J. Ber. d. DMV 32 (1923), p. 177--184
\end{center}

\begin{center}
{\Large\bfseries Invariants algébriques et différentiels.}\par
\vspace{0.45em}
Compte rendu d'après la conférence donnée à Leipzig le 18 septembre 1922.\par
\vspace{1.0em}
Par \textsc{Emmy Noether} à Göttingen.
\end{center}

Si je dois rendre compte du développement des invariants algébriques et différentiels, je voudrais me limiter, dans les deux domaines, à la période critique qui, selon une remarque de Hilbert, succède à la période naïve et formelle. Cette période critique est caractérisée, pour les invariants algébriques, par le nom de Hilbert lui-même, et, pour les invariants différentiels, par le nom de Riemann; ou, du point de vue du contenu, pour les invariants algébriques, par les méthodes arithmétiques de l'algèbre, qui se sont essentiellement développées au contact de ces questions et qui purent y montrer toute leur acuité, et, pour les invariants différentiels, par les méthodes du calcul formel des variations. Je voudrais donc parler des méthodes arithmétiques, avec leur portée qui dépasse l'objet particulier, tandis que, dans la seconde partie, je pourrai me borner à la subordination des invariants différentiels aux invariants algébriques, subordination qui est accomplie précisément dans le prolongement des méthodes du calcul des variations.

\textbf{1.} Les \emph{invariants algébriques} reposent, comme on sait, sur le \emph{groupe linéaire général} en $x_1,\ldots,x_n$:
\[
\tag{1}
x_i=\sum s_{ik}x'_k
\qquad\text{ou brièvement:}\qquad
x=S(x'),
\]
où les $s_{ik}$ désignent des indéterminées, de sorte que $\D=|s_{ik}|$ est non nul. Si maintenant $f(a,x)$, $g(b,x),\ldots$ sont des formes en $x$, de dimensions $\alpha,\beta,\ldots$, avec des indéterminées $a,b,\ldots$ comme coefficients, alors, au moyen de (1), naît la \emph{transformation induite} des $a,b,\ldots$; car de
\[
\begin{aligned}
f(a,x)&=\sum a_{i_1\ldots i_n}x_1^{i_1}\cdots x_n^{i_n}
      =\sum a'_{j_1\ldots i_n}{x'_1}^{j_1}\cdots {x'_n}^{j_n}
      =f(a',x'),\\
g(b,x)&=\sum b_{r_1\ldots r_n}x_1^{r_1}\cdots x_n^{r_n}
      =\sum b'_{s_1\ldots s_n}{x'_1}^{s_1}\cdots {x'_n}^{s_n}
      =g(b',x')
\end{aligned}
\]
il résulte que les $a',b',\ldots$ deviennent des fonctions linéaires homogènes des $a,b,\ldots$, dont les coefficients sont de degrés $\alpha,\beta,\ldots$ en les $s_{ik}$:
\[
\tag{2}
a'=A_\alpha(a);\qquad b'=B_\beta(b)\ldots .
\]
Par \emph{invariant}, on entend alors un invariant relativement aux transformations (1) et (2); c'est-à-dire une fonction rationnelle entière des indéterminées $a,b,\ldots,x,y,\ldots$ -- où $y=S(y')$ aussi -- qui, sous l'application des transformations, se reproduit à un facteur près, ce facteur étant une puissance du déterminant de substitution:
\[
\tag{3}
I(a',b',\ldots,x',y',\ldots)=\D^\rho I(a,b,\ldots,x,y,\ldots).
\]
Les coefficients de $I$ peuvent être supposés rationnels, ou même entiers rationnels, puisque tout invariant à coefficients dans un corps de nombres arbitraire $P$ peut s'exprimer linéairement, avec des coefficients de $P$, au moyen d'un nombre fini de tels invariants spéciaux. De même, ce n'est pas une restriction de supposer $I$ séparément homogène dans chaque série, car, là encore, tout invariant se compose comme somme d'un nombre fini de tels invariants spéciaux.

Il faut remarquer que, réciproquement, (1) et (2) sont de nouveau déterminées par l'exigence (3). Comme Ostrowski et Schur l'ont montré récemment,\footnote{A. Ostrowski et I. Schur, \emph{Über eine fundamentale Eigenschaft der Invarianten einer binären Form}, Math. Zeitschr. 15 (1922), et travaux connexes encore inédits d'Ostrowski.} les transformations induites peuvent aussi être caractérisées comme la totalité des transformations linéaires, séparées dans chaque série individuelle, qui -- à certaines exceptions indiquables près -- admettent un invariant arbitraire $I$ libre de $x,y,\ldots$.

Le produit de deux invariants est de nouveau un invariant, mais la somme ne l'est que si, et seulement si, le poids -- l'exposant $\rho$ dans (3) -- coïncide. Si l'on se restreint toutefois aux transformations de déterminant un, alors une somme arbitraire est de nouveau un invariant et épuise aussi tous les invariants relativement au groupe ainsi restreint. On obtient donc un domaine d’intégrité, et plus précisément un \emph{domaine d’intégrité homogène}, c'est-à-dire un domaine tel que l'appartenance d'un polynôme entraîne l'appartenance séparée de ses composantes homogènes; ce qui revient donc ici encore aux invariants séparément homogènes selon (3).

Ici naît le problème fondamental autour duquel se sont développées la théorie des invariants et, plus généralement, l'algèbre arithmétique: \emph{ce domaine d’intégrité est-il fini}, c'est-à-dire possède-t-il une \emph{base d'intégrité finie} $I_1,\ldots,I_k$ telle que tout invariant puisse se représenter rationnellement et intégralement par $I_1,\ldots,I_k$, $I=G(I_1,\ldots,I_k)$? La solution hilbertienne du problème -- la démonstration de Gordan pour le cas binaire ne permet pas d'extension à cause des calculs symboliques impossibles à dominer, et celle de Mertens--Hilbert ne le permet pas non plus à cause de son emploi de la décomposition d'une forme binaire en formes linéaires -- repose sur la réduction de la question de la base d'intégrité à celle de la base d'idéal. Je voudrais esquisser ici la marche des idées, en soulignant les idées arithmétiques fondamentales, puis aborder encore trois cercles de questions qui s'y rattachent: la formation effective de la base par un nombre fini d'étapes, les questions de finitude pour les sous-groupes, et les questions de finitude compte tenu de l'intégralité.

\textbf{2.} Je rappelle la \emph{définition de l'idéal} dans les corps de nombres algébriques finis. Un système $\ma$ de nombres du corps s'appelle idéal lorsque
\[
\begin{array}{ll}
1.& \ma\ \text{appartient au domaine }\mo\text{ de tous les entiers du corps},\\
2.& \text{si, avec }\alpha\text{ et }\beta,\text{ la différence }\alpha-\beta\text{ appartient aussi à }\ma,\\
3.& \text{si, avec }\alpha,\text{ appartient aussi }\lambda\alpha\text{ à }\ma,\text{ où }\lambda\text{ désigne un élément arbitraire de }\mo.
\end{array}
\]
Et l'on sait que vaut le théorème selon lequel chaque idéal possède une \emph{base d'idéal}: $\ma=(\alpha_1,\ldots,\alpha_k)$; c'est-à-dire qu'il existe un nombre fini d'éléments $\alpha_1,\ldots,\alpha_k$ de $\ma$ tels que $\ma$ soit dérivable de $\alpha_1,\ldots,\alpha_k$ au moyen de 2. et 3.; autrement dit, $\alpha=\lambda_1\alpha_1+\cdots+\lambda_k\alpha_k$ épuise tous les éléments de $\ma$.\footnote{Que $k$ puisse être abaissé à deux n'importe pas pour ce qui suit.}

Cette définition de l'idéal repose seulement sur le fait que $\mo$ forme un domaine d’intégrité; la notion d'idéal, et de même la notion de base d'idéal, demeure donc littéralement conservée si l'on remplace $\mo$ par un domaine d’intégrité arbitraire.

La démonstration hilbertienne de finitude se décompose alors en trois théorèmes:
\[
\begin{array}{ll}
1.& \text{Dans l'anneau de polynômes en }n\text{ indéterminées à coefficients dans un}\\
& \text{corps de rationalité, le théorème de la base d'idéal vaut pour tout idéal}\\
& \text{-- et par conséquent aussi pour tout système arbitraire }\mS.\\[2pt]
2.& \text{Un domaine d’intégrité homogène }\mJ\text{ de polynômes est fini si et seulement si}\\
& \text{le théorème de la base d'idéal vaut dans }\mJ.\\[2pt]
3.& \text{Pour le domaine des invariants }\mJ,\text{ le théorème de la base d'idéal vaut dans la}\\
& \text{forme renforcée selon laquelle toute base d'idéal formée relativement à l'anneau}\\
& \text{de tous les polynômes est en même temps une base d'idéal dans }\mJ,\text{ de sorte que, à côté de}
\end{array}
\]
\[
 I=A_1I_1+\cdots+A_kI_k
 \qquad\text{on ait toujours}\qquad
 I=i_1I_1+\cdots+i_kI_k,
\]
où les $A$ désignent des polynômes en $a,b,\ldots,x,y,\ldots$, et les $i$ des invariants.

La preuve de 1. repose, en principe, sur les mêmes raisonnements que dans le corps de nombres algébriques; dans les deux cas, l'existence de la base d'idéal résulte directement d'un théorème général sur les modules de formes linéaires.\footnote{Voir à ce sujet mon rapport comparatif sur la théorie arithmétique des fonctions algébriques, Jahresber. 28 (1920), p. 187 et note 2), p. 188.} De 1. résulte directement l'existence de la base d'idéal pour tout domaine d’intégrité fini de polynômes; il est facile de voir que, pour les domaines homogènes, la réciproque vaut aussi. Ce n'est qu'en 3. que des propriétés spéciales des invariants sont utilisées; et, chez Hilbert, 3. résulte d'un théorème connu sur le procédé $\Omega$, tandis que E. Fischer a montré récemment que 3. peut déjà se déduire -- en partant d'autres théorèmes généraux -- de la propriété du groupe linéaire général de contenir, avec chaque transformation, aussi la transformation contragrédiente correspondante, respectivement sa conjuguée complexe.\footnote{E. Fischer, \emph{Über die Endlichkeit der Invarianten}. Göttinger Nachrichten 1915, et conférence de Leipzig.}

\textbf{3.} La \emph{formation effective de la base} repose sur une autre transformation arithmétique de la notion de finitude, par recours à la théorie des \emph{corps de fonctions}. Il vaut:
\[
\begin{array}{ll}
4.& \text{Un domaine d’intégrité }\mJ\text{ de polynômes est fini si et seulement si}\\
& \text{il contient un sous-domaine fini }\mJ'\text{ tel que }\mJ\\
& \text{dépende algébriquement et intégralement de }\mJ';\text{ toute base d'intégrité de }\mJ'\\
& \text{est alors complétée, par un système fondamental de ces grandeurs}\\
& \text{algébriquement entières, en une base de }\mJ.
\end{array}
\]
\footnote{Hilbert montre (\emph{Über die vollen Invariantensysteme}, Math. Annalen 42 (1893), \S\S{} 1 et 2) que ces faits résultent de la finitude supposée; inversement, que la finitude résulte de la dépendance algébriquement entière se déduit de l'existence de la base rationnelle (E. Noether, \emph{Körper und Systeme rationaler Funktionen}, Math. Annalen 76 (1915), \S\ 12).}
S'il s'agit spécialement, comme dans le cas des invariants, de domaines homogènes $\mJ$ pour lesquels le théorème de la base d'idéal vaut dans la forme plus forte 3., un tel domaine $\mJ'$ peut être déduit de certains polynômes $I_1,\ldots,I_s$, toujours existants et en nombre fini, comme base d'intégrité, polynômes dont l'annulation entraîne l'annulation de tous les polynômes de $\mJ$. Ce fait se fonde de nouveau sur un théorème général de la théorie des idéaux:
\[
\begin{array}{ll}
5.& \text{À tout idéal }\ma\text{ de polynômes appartient un entier rationnel ordinaire }r\\
& \text{tel que, pour tout idéal }\mb\text{ qui s'annule en tous les zéros de }\ma,\\
& \mb^r\text{ soit divisible par }\ma.
\end{array}
\]

Dans le cas des invariants, on peut déterminer une \emph{borne supérieure pour les degrés} de $I_1,\ldots,I_s$, et donc aussi pour leur nombre $s$, qui dépend seulement des degrés des formes fondamentales prises pour base. Ici encore interviennent des propriétés spéciales des invariants; il s'agit du théorème selon lequel un système de formes fondamentales numériquement spécialisées possède un invariant non nul si et seulement si $\D=|s_{ik}$ dépend algébriquement et intégralement des coefficients transformés.

À partir de cette borne supérieure, K. Hentzelt\footnote{Je publierai prochainement le travail d'après les papiers laissés par lui.} a déterminé une \emph{borne supérieure pour les degrés des invariants de la base d'intégrité}, qui dépend de nouveau seulement des degrés des formes fondamentales. Il y parvient en donnant, pour l'exposant $r$ en 5., une borne supérieure qui dépend seulement du nombre et du plus haut degré des polynômes d'une base d'idéal de $\ma$, qui est indépendante des coefficients de ces polynômes, et qui se calcule à partir des nombres indiqués en un nombre fini d'étapes. En principe, le problème posé est ainsi entièrement réglé; toutefois la borne indiquée est très haute: ainsi, pour une forme quadratique ternaire, où le discriminant, c'est-à-dire un invariant de degré 3, forme déjà la base de tous les invariants libres des $x$, on monte, selon Hilbert et Hentzelt, jusque dans les millions.

\textbf{4.} Pour la question de finitude des invariants de \emph{sous-groupes} du groupe linéaire général, on n'est parvenu jusqu'ici qu'à des résultats partiels. La méthode repose presque exclusivement sur la réduction des invariants de sous-groupes à des invariants simultanés du groupe linéaire général, suivant le procédé adopté par Study pour le groupe orthogonal. Weitzenböck l'a réalisée pour les invariants de mouvement et les invariants affines, Deruyts pour les semi-invariants et les invariants de cisaillement;\footnote{Pour les indications bibliographiques, voir l'article d'encyclopédie de Weitzenböck sur la théorie des invariants; III E 1.} la question de la portée de cette méthode demeure indécise. La remarque de Fischer mentionnée à la fin de 2. résout également le problème de finitude pour une série de sous-groupes; en particulier, elle fournit ainsi la démonstration la plus simple pour le groupe orthogonal. Mais l'exemple des semi-invariants montre -- comme Fischer l'a récemment signalé\footnote{Conférence de Leipzig et Math. Zeitschrift.} -- que, pour des sous-groupes, malgré la finitude, la forme plus forte 3. pour la base d'idéal n'a pas besoin d'être satisfaite.

Saisie au sens de la théorie des corps, la question de finitude présente ici conduit à la question hilbertienne des \emph{fonctions relativement entières}, c'est-à-dire à la question de savoir si de tels domaines d’intégrité, constitués par la totalité des polynômes d'un corps de fonctions rationnelles, sont toujours finis. Dans cette direction se trouve le fait -- démontrable de façon élémentaire -- que, pour les invariants des groupes finis, une base d'intégrité distinguée peut être donnée directement: le système des coefficients de la résolvante de Galois.\footnote{E. Noether, \emph{Der Endlichkeitssatz der Invarianten endlicher Gruppen}. Math. Ann. 77 (1916).}

\textbf{5.} La question de savoir si le théorème de finitude des invariants peut être renforcé relativement à l'\emph{intégralité} n'est pas non plus réglée en général. Certes, le théorème de la base d'idéal vaut aussi, comme Hilbert l'a montré, pour l'anneau des polynômes à coefficients entiers; et le lien 2. entre base d'idéal et base d'intégrité demeure également valable;\footnote{De ce lien résulte en particulier que, pour tous les domaines d’intégrité finis, valent les théorèmes de décomposition de la théorie générale des idéaux, tels que je les ai développés dans Math. Ann. 83 (1921).} mais la preuve de la forme renforcée 3. -- qui n'est pas nécessaire -- échoue. Pour le cas des invariants binaires, la finitude entière peut se démontrer;\footnote{E. Noether, \emph{Die Endlichkeit des Systems der ganzzahligen Invarianten binärer Formen}. Göttinger Nachrichten 1919.} la preuve, qui représente un renforcement d'intégralité de la preuve de finitude binaire de Mertens--Hilbert, utilise toutefois, en plus du théorème de la base d'idéal entière, la décomposition des formes binaires en formes linéaires, et ne permet donc aucun transfert à un plus grand nombre d'indéterminées. En revanche, sur des fondements analogues, la finitude entière s'obtient pour les invariants des groupes finis, comme renforcement de la preuve mentionnée plus haut, bien qu'il n'en sorte alors de nouveau qu'une preuve d'existence et aucune indication effective de la base.\footnote{\S\ 4 du travail cité sous 11).} Ici aussi, seule une poursuite du développement de la théorie des corps de fonctions et de la théorie générale des idéaux conduira probablement au but.

\textbf{6.} Je passe maintenant aux \emph{invariants différentiels} et à la question posée dans l'introduction, celle de leur réduction à la théorie des invariants algébriques. Le groupe pris pour base est ici, comme on sait, défini au moyen de $x_i=x_i(y)$ par:
\[
\tag{4}
d x_i=\sum \frac{\partial x_i}{\partial y_k}d y_k;\qquad
d\delta x_i=\sum \frac{\partial x_i}{\partial y_k}d\delta y_k+
\sum\frac{\partial^2 x_i}{\partial y_k\partial y_l}d y_k\,\delta y_l;\ \ldots .
\]
Pour passer aux transformations induites par (4), on peut prendre pour point de départ une fonction arbitraire $f(x,d x)=g(y,d y)$; en exigeant l'invariance de $\delta f,\delta^2f,\ldots$ relativement à (4), on obtient les transformations des dérivées de $f$ par rapport à $x$ et à $d x$. Si l'on se restreint par exemple à des formes homogènes en les $d x$:
\[
 f(x,d x)=\sum a_{i_1\ldots i_n}(x)d x_1^{i_1}\cdots d x_n^{i_n}
 =\sum a'_{j_1\ldots j_n}(y)d y_1^{j_1}\cdots d y_n^{j_n}
 =g(y,d y),
\]
alors les $a'$ sont linéaires en les $a$, les $\frac{\partial a'}{\partial y}$ sont linéaires en les $a$ et les $\frac{\partial a}{\partial x}$, et ainsi de suite:
\[
\tag{5}
a'=A_0(a),\qquad
\frac{\partial a'}{\partial y}=A_1\!\left(a,\frac{\partial a}{\partial x}\right),\qquad
\frac{\partial 2a'}{\partial y^2}=A_2\!\left(a,\frac{\partial a}{\partial x},\frac{\partial 2a}{\partial x^2}\right),\ldots
\]
et il s'agit d'invariance relativement aux transformations (4) et (5). Toutes les grandeurs qui interviennent
\[
a,\ \frac{\partial a}{\partial x},\ldots,\ \dd x,\ \delta x,\ \dd\delta x,\ldots;\ 
\frac{\partial x}{\partial y},\ \frac{\partial^2 x}{\partial y^2},\ldots
\]
doivent alors être considérées comme des indéterminées, de sorte que le déterminant de chaque transformation individuelle est certainement non nul. Ce n'est qu'au moment d'appliquer la théorie formellement achevée à des fonctions spéciales que surgit la question de la restriction à de telles fonctions que les transformations restent réversibles de manière univoque dans un certain domaine, et qu'il existe un nombre suffisant de quotients différentiels, exactement comme, dans le cas algébrique, lors d'une spécialisation numérique, il faut supposer non nul le déterminant $|s_{ik}|$.

Dans cette considération de tous les arguments qui apparaissent comme indéterminées, on a donc affaire à un groupe linéaire spécial, non directement accessible aux méthodes algébriques ordinaires, dans lequel les transformations de $d x$ et de $a'=A_0(a)$ coïncident avec (1) et avec la transformation induite (2). La question se pose alors de savoir si, peut-être en général, par \emph{introduction de nouveaux arguments}, les transformations (4) et (5) peuvent se ramener au groupe linéaire général (1) et aux transformations (2) \emph{induites} par lui.

C'est effectivement le cas; les différentielles supérieures $d^2x,\ d\delta x\ldots$ doivent être remplacées par les dérivées de Lagrange qui naissent du problème variationnel appartenant à $f(x,d x)$, et par d'autres combinaisons qui en résultent par polarisation (connexions au sens de Weyl et de Schouten), toutes cogrédientes aux $d x$, et dont la transformation est donc donnée par le groupe linéaire général (1). Les grandeurs $\frac{\partial a}{\partial x},\frac{\partial^2a}{\partial x^2},\ldots$ (ou, plus généralement, les dérivées de $f$, supposé homogène, par rapport à $x$ et à $d x$) doivent en revanche être remplacées par les coefficients de certaines formes provenant des "formes normales des variations supérieures" par l'élimination indiquée plus haut des $d^2x,\ d\delta x\ldots$, ainsi que de leurs "dérivées covariantes"
\[
 [\Omega_i],\quad [\Omega_i^{(1)}],\quad [\Omega_i^{(2)}],\ldots\qquad (i=1,2,\ldots);
\]
et comme les $[\Omega_i^{(k)}]$ sont des invariants relativement à (4) et (5) qui ne dépendent que des $d x$ et des $\delta x$, ces coefficients satisfont aux transformations algébriquement induites (2).\footnote{E. Noether, \emph{Invarianten beliebiger Differentialausdrücke}. Göttinger Nachrichten 1918.} Par ce \emph{théorème de réduction}, la théorie des invariants différentiels est subordonnée à la théorie algébrique des invariants.

La série des $[\Omega_i]$ s'arrête avec $[\Omega_\alpha]$ lorsque $f(x,d x)$ est une forme homogène de dimension $\alpha$ en les $d x$. Pour les formes quadratiques, $[\Omega_2]$ devient identique à la forme de courbure de Riemann, et le procédé de formation indiqué ici est exactement celui que Riemann a donné dans son mémoire latin pour le prix, procédé que Lipschitz a développé lui aussi, indépendamment de cela, à la suite de la leçon d'habilitation de Riemann.\footnote{Voir aussi mon exposé sur Riemann et Lipschitz au n\textsuperscript{o} 28 de la seconde partie de l'article encyclopédique de Weitzenböck mentionné sous 7).} La preuve du théorème de réduction, que Christoffel et Ricci avaient menée par le calcul pour les formes différentielles quadratiques, repose sur l'introduction des "coordonnées normales de Riemann", qui transforment en droites les extrémales issues d'un point; puisque, par là, à une transformation arbitraire des $x$ correspond une transformation linéaire des coordonnées normales, celles-ci déterminent le passage au groupe linéaire général.

La question se pose de savoir si un tel théorème de réduction existe aussi lorsque, suivant Weyl et Schouten, on ne prend pas pour base un problème variationnel afin de définir le groupe, mais seulement une connexion; c'est-à-dire lorsque, conformément aux dérivées de Lagrange, on établit des expressions $\psi_i(d\delta x,d x,\delta x)$ par lesquelles l'élimination des différentielles supérieures est obtenue, en ajoutant, par analogie avec les formes différentielles quadratiques, la restriction en soi inutile selon laquelle les $\psi_i$ doivent être linéaires en les $d x$ (pour un $f(x,d x)$ homogène arbitraire, les dérivées de Lagrange polarisées ne deviennent homogènes du premier ordre qu'en les $d x$). Par l'exigence de cogrédience des $\psi$ relativement aux $d x$, la transformation de leurs coefficients est induite, et ainsi, par analogie avec (5), le groupe est fixé. Pour les invariants des ordres les plus bas, Weitzenböck a confirmé par le calcul le théorème de réduction dans le cas de Weyl; un point de départ général manque encore.\footnote{Voir à ce sujet: Weyl, \emph{Raum, Zeit, Materie}; Schouten, Math. Zeitschr. 13 (1922); Weitzenböck, Sitzungsber. d. Wiener Akademie, 1920 et 1921.}

\begin{center}
(Reçu le 26 octobre 1922.)
\end{center}
\clearpage
\setcounter{footnote}{0}
\providecommand{\frakS}{\mathfrak S}
\providecommand{\fraka}{\mathfrak a}
\providecommand{\frakb}{\mathfrak b}
\providecommand{\frakc}{\mathfrak c}
\providecommand{\frakg}{\mathfrak g}
\providecommand{\frakG}{\mathfrak G}
\providecommand{\frakh}{\mathfrak h}
\providecommand{\frakm}{\mathfrak m}
\providecommand{\frakM}{\mathfrak M}
\providecommand{\frakn}{\mathfrak n}
\providecommand{\frako}{\mathfrak o}
\providecommand{\frakp}{\mathfrak p}
\providecommand{\frakq}{\mathfrak q}
\providecommand{\frakr}{\mathfrak r}
\providecommand{\barfrakm}{\overline{\mathfrak m}}

\section*{\NoetherNavAnchor{P24}{24. Théorie de l’élimination et théorie générale des idéaux}24. Théorie de l'élimination et théorie générale des idéaux}
\begin{center}
Math. Ann. 90 (1923), p. 229--261\\[0.9em]
{\Large\bfseries Théorie de l'élimination et théorie générale des idéaux.}\\[1.0em]
Par\\[0.45em]
Emmy Noether à Göttingen.\\[0.6em]
\rule{4em}{0.4pt}
\end{center}

Il s'agit, dans ce qui suit, d'inscrire la théorie de l'élimination -- dans la forme arithmétique que j'ai donnée à l'exposé de Hentzelt\footnote{Hentzelt, Kurt: \emph{Zur Theorie der Polynomideale und Resultanten}. Édité par E. Noether. Math. Ann. 88 (1922), p. 53--79; cité H.-N.} et qui peut se désigner comme une théorie des idéaux dans l'anneau de polynômes -- dans la théorie générale des idéaux,\footnote{Noether, E.: \emph{Idealtheorie in Ringbereichen}, Math. Ann. 83 (1921), p. 24--66; cité \emph{Idealtheorie}.} et en même temps de donner une refondation des parties relatives aux zéros. Les notions fondamentales des deux théories sont brièvement réunies au \S{} 1, de même que celles de la théorie des corps,\footnote{Steinitz, E.: \emph{Algebraische Theorie der Körper}, J. f. M. 137 (1910), p. 167--309; cité Steinitz.} de sorte que je pourrai m'y référer ici.

L'insertion et la refondation se groupent autour de quatre questions: formulation arithmétique de la notion de dimension; théorie des zéros; rapport des théorèmes de décomposition pour les normes et les diviseurs élémentaires avec ceux de la théorie générale des idéaux; caractérisation des idéaux premiers et primaires par la norme et la forme des diviseurs élémentaires.

La formulation arithmétique de la \emph{notion de dimension} (\S{} 4) se fait par des chaînes d'idéaux premiers, équivalent arithmétique du fait que, sur les surfaces, il y a des courbes, et sur les courbes, des points. Cette formulation arithmétique, dont on démontre l'identité avec la définition par paramètres de la théorie de l'élimination, se transpose avec une légère modification à des anneaux arbitraires et y donne, conjointement avec la transposition de certaines parties des autres questions, une vue exacte de la structure des idéaux, comme je le montrerai ailleurs.

La question des \emph{zéros} est abordée chez H.-N., comme d'ordinaire, directement pour l'idéal donné, à savoir par retour à l'élimination successive, ce qui n'évite pas entièrement le caractère de vérification. À l'opposé se trouve la voie toujours suivie pour les polynômes d'une seule variable: on représente le polynôme donné comme produit de puissances de polynômes irréductibles (polynômes primaires), et l'on construit, pour chaque polynôme irréductible pris isolément, le corps des zéros comme un corps isomorphe au corps résiduel. Le degré du polynôme irréductible donne le degré de ce corps; mais le degré du polynôme primaire, dont les zéros coïncident avec ceux du polynôme irréductible, donne la dimension de l'anneau quotient par ce polynôme primaire. Si l'on passe, pour chaque polynôme irréductible, au corps de Galois, on obtient la décomposition en facteurs linéaires; et ainsi, pour le polynôme primitivement donné, la question des zéros, de la décomposition en facteurs linéaires et de la multiplicité est réglée. De manière tout à fait correspondante, on passe ici (\S{} 3) de l'anneau quotient par un idéal premier à son corps résiduel, et l'on construit un corps des zéros qui lui est isomorphe. Ce corps des zéros contient un sous-corps engendré par adjonction d'indéterminées -- par exemple $x_{i+1},\ldots,x_n$; le degré de la norme donne le degré par rapport à ce sous-corps; en même temps, le passage au corps de Galois fournit la décomposition de la forme des diviseurs élémentaires et par conséquent de la norme en facteurs linéaires. Pour l'idéal primaire associé, les zéros sont cependant les mêmes; la décomposition est également donnée avec lui, puisque la norme devient une puissance de la forme des diviseurs élémentaires de l'idéal premier (\S{} 2); le degré de la norme donne la dimension de l'anneau quotient par l'idéal primaire relativement au sous-corps issu des indéterminées.

Ainsi la question des zéros d'un idéal arbitraire est ramenée au rapport des théorèmes de décomposition pour les normes et les diviseurs élémentaires avec ceux de la théorie générale des idéaux (\S{} 5). Au fait que les zéros de l'idéal se composent à partir de ceux de ses divers idéaux premiers associés correspond la décomposition de la norme en plus grands facteurs primaires, qui deviennent égaux à des puissances de la forme des diviseurs élémentaires des divers idéaux premiers associés;\footnote{Ce parallélisme de la théorie de l'élimination avec la théorie générale des idéaux n'est pas réalisé dans la théorie de l'élimination de Kronecker; voir à ce sujet H.-N., note *).} ainsi la décomposition en facteurs linéaires est accomplie pour la norme en général. Mais la multiplicité reçoit chez H.-N. son interprétation au moyen des idéaux fondamentaux et de leurs composantes; le degré d'un plus grand facteur primaire de la norme devient égal à la dimension -- après adjonction de $x_{i+1},\ldots,x_n$ -- du quotient de l'idéal fondamental de $(i-1)$-ième étage par une composante de l'idéal fondamental de $i$-ième étage. Or l'idéal fondamental de $(i-1)$-ième étage est démontré identique à la composante isolée de la décomposition qui est définie de manière unique par tous les idéaux premiers de dimension supérieure à la $(n-i)$-ième; et une composante de l'idéal fondamental de $i$-ième étage est démontrée identique à la composante isolée de la décomposition qui est définie par l'idéal premier correspondant au facteur de la norme et par tous les idéaux premiers de dimension supérieure. Et le degré du facteur correspondant de la norme devient -- après adjonction de $x_{i+1},\ldots,x_n$ -- égal à la dimension du sous-anneau de l'anneau quotient par cette composante formé des classes modulo celle-ci qui sont représentées par des éléments de l'idéal fondamental de $(i-1)$-ième étage.

La caractérisation des idéaux premiers et des idéaux primaires propres (\S{} 6) résulte comme conséquence directe de la considération des corps résiduels correspondants. Si le domaine des coefficients initial est un corps parfait, alors les idéaux premiers ont pour norme et forme des diviseurs élémentaires des polynômes irréductibles, les idéaux primaires propres des polynômes primaires propres, et réciproquement. Si le domaine des coefficients est un corps imparfait, on ne peut donner que des conditions soit nécessaires, soit suffisantes; la norme d'un idéal premier peut devenir un polynôme primaire propre, et la forme des diviseurs élémentaires d'un idéal primaire propre peut devenir un polynôme irréductible, comme le montrent des exemples. Plus précisément, en caractéristique $p$, la norme d'un idéal premier devient la $p^g$-ième puissance ($g\geqq0$) de sa forme des diviseurs élémentaires; tandis que, pour des idéaux primaires propres dont la forme des diviseurs élémentaires est un polynôme irréductible, la norme peut devenir une puissance quelconque, comme le montrent encore des exemples. On considère enfin les idéaux premiers absolus (\S{} 7), c'est-à-dire ceux qui restent des idéaux premiers lorsque le domaine des coefficients est étendu en un corps algébriquement clos. Pour des idéaux premiers absolus à coefficients dans un corps de nombres algébriques (fini), la caractérisation conduit à la transposition d'un théorème d'abord énoncé par A. Ostrowski pour les polynômes absolument irréductibles: la propriété d'être un idéal premier est conservée modulo chaque idéal premier du corps de nombres, sauf tout au plus un nombre fini d'exceptions.

Signalons encore l'analogie de cette dernière question avec la théorie des idéaux dans les corps de nombres algébriques. Comme forme des diviseurs élémentaires d'un tel idéal, il faut prendre le plus petit entier rationnel divisible par l'idéal (note 10), qui devient donc toujours un nombre premier pour un idéal premier, tandis que, réciproquement, un idéal devient premier dès que sa norme est un nombre premier; les preuves aussi se déroulent tout à fait parallèlement. Les exemples mentionnés ci-dessus montrent donc que, dans le cas d'un corps imparfait, apparaît l'analogue des idéaux premiers de degré supérieur et des idéaux de ramification, tandis que, dans le cas d'un corps parfait, il n'y a que des idéaux premiers du premier degré et aucun idéal de ramification.

\subsection*{\S{} 1. Notions fondamentales de la théorie de l'élimination, de la théorie générale des idéaux et de la théorie des corps}

\textbf{1. L'anneau.} Soit $\bar{\mathfrak{J}}$ le domaine d'intégrité des polynômes en $y_1,\ldots,y_n$ à coefficients dans un corps $\mathsf{P}$ défini abstraitement. Que les $y$ soient soumis à une transformation ayant pour coefficients des indéterminées $u_{\mu\nu}$:\footnote{La transformation est, en vue du \S{} 3, etc., un peu plus générale que chez H.-N.; tous les résultats qui s'y trouvent sont par là conservés à plus forte raison.}
\begin{align}
 y_1&=u_{11}x_1+\cdots+u_{1n}x_n,\quad \ldots,\quad
 y_n=u_{n1}x_1+\cdots+u_{nn}x_n, \tag{1}
\end{align}
avec l'inverse
\begin{align}
 x_1&=t_{11}y_1+\cdots+t_{n1}y_n,\quad \ldots,\quad
 x_n=t_{1n}y_1+\cdots+t_{nn}y_n, \tag{1'}
\end{align}
et que les $u_{\mu\nu}$, ou -- ce qui est équivalent en raison de leur exprimabilité rationnelle mutuelle -- les $t_{\mu\nu}$, soient adjoints à $\mathsf{P}$, de sorte que naisse le domaine des coefficients $\mathsf{P}(u)=\mathsf{P}(t)$; $\mathfrak{J}$ désignera le domaine d'intégrité des polynômes en $x_1,\ldots,x_n$ à coefficients dans $\mathsf{P}(u)$. Les domaines $\mathfrak{J}$ et $\bar{\mathfrak{J}}$ sont à la base de la théorie de l'élimination; on considère des idéaux $\fraka,\frakb,\ldots$ dans $\mathfrak{J}$ et $\bar{\fraka},\bar{\frakb},\ldots$ dans $\bar{\mathfrak{J}}$. Un idéal dans un anneau quelconque est alors défini, comme d'habitude, par le fait qu'il contient, avec $\alpha$ et $\beta$, aussi $\alpha-\beta$, et, avec $\alpha$, aussi $\lambda\alpha$, où $\lambda$ désigne un élément arbitraire de l'anneau; $\frako$ désignera toujours l'idéal unité, composé de tous les éléments.

\textbf{2. Idéaux transformés.} Un idéal $\frakm$ dans $\mathfrak{J}$ s'appelle transformé s'il naît d'un idéal $\barfrakm$ dans $\bar{\mathfrak{J}}$ au moyen de (1), après adjonction des $u_{\mu\nu}$, ou des $t_{\mu\nu}$. Ainsi, lorsque $\bar f(y)$ ou $\bar g(y)$ parcourt tous les polynômes de $\barfrakm$, $\frakm$ se compose de toutes les combinaisons linéaires
\[
 f(x)=\sum U_\lambda(u)\bar f_\lambda(y)=\sum U_\lambda(u)f_\lambda(x)
\]
ou encore
\[
 g(x)=\sum T_\lambda(t)\bar g_\lambda(y)=\sum T_\lambda(t)g_\lambda(x);
\]
en particulier, tous les $f_\lambda(x)=\bar f_\lambda(y)$ y sont contenus, ou, ce qui revient au même, tous les $g_\lambda(x)=\bar g_\lambda(y)$. Comme $f(x)$, respectivement $g(x)$, ne sont déterminés qu'à des grandeurs de $\mathsf{P}(u)=\mathsf{P}(t)$ près, les $U_\lambda,T_\lambda$ peuvent être supposés, sans restriction de généralité, être des produits de puissances des $u$, respectivement des $t$. Les $f(x),g(x)$ passent de nouveau en des polynômes de $\frakm$ lorsque $U,T$ sont remplacés par des produits de puissances quelconques, donc en particulier par permutation des colonnes de (1), respectivement des lignes de (1'); ainsi $\frakm$ contient, avec n'importe quel polynôme, aussi tous ceux qui s'en déduisent par permutation des $x$, pourvu seulement que les permutations correspondantes soient effectuées dans les coefficients dépendant des $u$, respectivement des $t$. Tous les idéaux qui apparaîtront dans la suite seront -- sauf mention expresse du contraire -- supposés transformés. Les idéaux non transformés passent, par composition de (1) avec
\[
 x_i=v_{i1}z_1+\cdots+v_{in}z_n
\]
en idéaux transformés dans l'anneau de polynômes en $z$ à coefficients dans $\mathsf{P}(u,v)$, tandis que, pour les idéaux transformés, cette composition revient seulement à remplacer les $u_{\mu\nu}$ par des combinaisons bilinéaires des $u,v$.

\textbf{3. Notation.} Par $f^{(i)},g^{(i)},\ldots,a^{(i)},b^{(i)},\ldots$ on entendra constamment des polynômes de $\mathfrak{J}$ qui sont libres de $x_1,\ldots,x_{i-1}$.

\textbf{4. Idéaux fondamentaux.} À chaque idéal $\frakm$ sont associés $n$ idéaux fondamentaux $\frakg_0,\frakg_1,\ldots,\frakg_{n-1}$ par la convention suivante: l'idéal fondamental de $(i-1)$-ième étage $\frakg_{i-1}$ contient tous et seulement les polynômes $G(x)$ pour lesquels il existe un polynôme $b^{(i)}\ne 0$ -- en général variable avec $G(x)$ -- tel que $b^{(i)}G(x)\equiv 0\pmod{\frakm}$. L'existence de la base d'idéal garantit également celle d'au moins un $B^{(i)}$, fixe pour tous les $G(x)$, de sorte que
{\let\veqno\leqno\begin{equation}
 B^{(i)}\frakg_{i-1}\equiv0\pmod{\frakm},\qquad B^{(i)}\ne0 . \tag{2}
\end{equation}}
Les idéaux fondamentaux d'idéaux transformés deviennent eux-mêmes des idéaux transformés (H.-N., théorème VI). L'idéal fondamental de $i$-ième étage $\frakg_i$ est aussi défini comme l'idéal dans lequel $\frakm$ passe lorsqu'on adjoint $x_{i+1},\ldots,x_n$ à $\mathsf{P}(u)$, si l'on se restreint -- ce qui alors ne signifie aucune restriction de généralité -- aux polynômes entiers en $x_{i+1},\ldots,x_n$; $\frakg_0$ devient toujours égal à l'idéal unité $\frako$.

\textbf{5. Représentation modulaire des idéaux, diviseurs élémentaires et normes individuelles.} Si l'on conçoit chaque polynôme $f(x)$ comme une forme linéaire dans les produits de puissances $\xi$ de $x_1,\ldots,x_{i-1}$, avec des polynômes $a^{(i)}$ comme coefficients, alors l'idéal $\frakm$ passe en un module $\frakM_{i-1}$ de formes linéaires relativement au domaine des $a^{(i)}$; c'est-à-dire que $\frakM_{i-1}$ contient, avec deux formes linéaires quelconques, aussi leur différence, et, avec une forme linéaire $l(\xi)$, aussi $a^{(i)}l(\xi)$ pour tout $a^{(i)}$. On appelle module fondamental d'un tel module $\frakM$ la totalité des formes linéaires $g(\xi)$ pour lesquelles $b^{(i)}g(\xi)\equiv0(\frakM)$, $b^{(i)}\ne0$; l'idéal fondamental $\frakg_{i-1}$ passe donc, dans la représentation modulaire, au module fondamental $\frakG_{i-1}$ de $\frakM_{i-1}$. $\frakM_{i-1}$ et $\frakG_{i-1}$ dépendent d'une infinité d'indéterminées $\xi$, de telle manière que, dans chaque forme linéaire prise isolément, seulement un nombre fini de ces indéterminées interviennent. $\frakM_{i-1}$ possède la propriété caractéristique de n'avoir, après adjonction de $x_{i+1},\ldots,x_n$ à $\mathsf{P}(u)$, qu'un nombre fini de diviseurs élémentaires différents de l'unité; c'est-à-dire que $\frakG_{i-1}$ et $\frakM_{i-1}$ admettent, après cette adjonction, une représentation par base -- où $\zeta$ désigne de nouvelles indéterminées liées aux $\xi$ par une transformation linéaire réversible, entière en $x_i$, de telle manière que dans $\zeta_i=r_i(\xi),\ \xi_k=s_k(\zeta)$ n'interviennent chaque fois qu'un nombre fini de ces indéterminées --
{\csname tagsleft@true\endcsname\begin{align}
\frakG_{i-1}&=(\zeta_0,\zeta_1,\ldots,\zeta_\sigma,\zeta_{\sigma+1},\ldots,\zeta_{\sigma+\nu},\ldots); \notag\\
\frakM_{i-1}&=(E^{(i)}\zeta_0,E_1^{(i)}\zeta_1,\ldots,E_\sigma^{(i)}\zeta_\sigma,
\zeta_{\sigma+1},\ldots,\zeta_{\sigma+\nu},\ldots), \tag*{\raisebox{9.5pt}[0pt][0pt]{(3)}}
\end{align}}
où chaque $E_\mu^{(i)}$ divise le précédent (H.-N. (24) et (28)).

Le plus haut diviseur élémentaire $E^{(i)}$ est aussi défini comme le plus grand diviseur commun -- au sens polynomial -- de tous les $B^{(i)}$ qui apparaissent dans (2); il peut être supposé entier et primitif en $x_{i+1},\ldots,x_n$, et devient alors régulier en $x_i$, c'est-à-dire que $E^{(i)}$ contient le terme $x_i^\lambda$ avec coefficient non nul, s'il est de degré $\lambda$ en les $x$.

Le produit $R^{(i)}$ de tous les diviseurs élémentaires qui apparaissent dans (3) est appelé norme de $\frakG_{i-1}$ par rapport à $\frakM_{i-1}$, ou encore norme individuelle de l'idéal $\frakm$; en symboles, et en tenant compte du fait que $\frakm$ passe en $\frakg_i$ lors de l'adjonction de $x_{i+1},\ldots,x_n$:
\begin{equation}
 R^{(i)}=E^{(i)}E_1^{(i)}\cdots E_\sigma^{(i)}
 =N(\frakG_{i-1}\mid\frakM_{i-1})=N(\frakg_{i-1}\mid\frakg_i).
\end{equation}
Comme le montre (3), $R^{(i)}$ devient, après adjonction de $x_{i+1},\ldots,x_n$, égal au déterminant de la substitution de passage de $\frakG_{i-1}$ à $\frakM_{i-1}$; et le degré de $R^{(i)}$ donne la dimension, relativement à $\mathsf{P}(u,x_{i+1},\ldots,x_n)$, du quotient $\frakG_{i-1}/\frakM_{i-1}$, donc de $\frakg_{i-1}/\frakg_i$; $R^{(i)}$ peut être supposé entier et primitif en $x_{i+1},\ldots,x_n$, et devient alors régulier en $x_i$. $R^{(i)}$ se calcule comme le plus grand diviseur commun -- au sens polynomial -- de tous les déterminants à $\varrho$ lignes d'un module $\frakM^*_{i-1}$, toujours existant, de rang $\varrho$, ayant la propriété que le système quotient $\frakG^*_{i-1}\mid\frakM^*_{i-1}$ est isomorphe à $\frakG_{i-1}\mid\frakM_{i-1}$, $\frakG^*_{i-1}$ étant entendu comme le module fondamental de $\frakM^*_{i-1}$ (H.-N., théorème VII et \S{} 5).

\textbf{6. Forme des diviseurs élémentaires et norme (forme résultante) des idéaux.} Le produit $E_\frakm=E^{(1)}E^{(2)}\cdots E^{(n)}$ de tous les plus hauts diviseurs élémentaires est appelé forme des diviseurs élémentaires, et le produit $R_\frakm=R^{(1)}R^{(2)}\cdots R^{(n)}$ de toutes les normes individuelles, norme (forme résultante)\footnote{Chez H.-N. seule la désignation forme résultante est employée; la désignation norme correspond aux propriétés arithmétiques.} de $\frakm$. La forme des diviseurs élémentaires et la norme sont divisibles par l'idéal; la norme est divisible par la forme des diviseurs élémentaires, et une puissance de cette dernière par la norme. Plus généralement, on a encore:
{\csname tagsleft@true\endcsname\begin{align}
 E^{(i)}\frakg_{i-1}&\equiv0(\frakg_i),
& E^{(n)}\cdots E^{(i)}\frakg_{i-1}&\equiv0(\frakm)\quad\text{et par conséquent}\notag\\
 R^{(i)}\frakg_{i-1}&\equiv0(\frakg_i),
& R^{(n)}\cdots R^{(i)}\frakg_{i-1}&\equiv0(\frakm),\quad\text{(H.-N., théorème VIII),} \tag*{\raisebox{8.5pt}[0pt][0pt]{(4)}}
\end{align}}
et de plus: si $\frakm$ est divisible par $\frakn$, et si les normes $R_\frakm$ et $R_\frakn$ coïncident, alors les idéaux $\frakm$ et $\frakn$ coïncident aussi (H.-N., théorème IX). En remarquant que, pour un diviseur de $\frakm$, la coïncidence de $R^{(1)},R^{(2)},\ldots$ entraîne celle des idéaux fondamentaux $\frakg_1,\frakg_2,\ldots$, et que l'on a toujours $\frakg_0=\frako$, on obtient de manière tout à fait correspondante: si $\frakn$ est un diviseur propre de $\frakm$, alors la première norme individuelle de $\frakn$ qui diffère de celle correspondante de $\frakm$ devient un diviseur propre.\footnote{En revanche, des facteurs ultérieurs de $R_\frakn$ peuvent devenir des multiples des facteurs correspondants de $R_\frakm$; par exemple toujours lorsque $\frakm$ est un idéal premier et que $\frakn$ n'est pas l'idéal unité.} Si $\frakm$ contient un polynôme $G^{(i)}\ne0$, alors, par définition, $\frakg_0,\frakg_1,\ldots,\frakg_{i-1}$ deviennent égaux à l'idéal unité; et ainsi $R^{(1)},\ldots,R^{(i-1)}$, d'après la propriété des normes individuelles indiquée en 5., qui est de déterminer la dimension des quotients correspondants, deviennent égales à l'unité; par conséquent $E^{(1)},\ldots,E^{(i-1)}$ deviennent également égales à l'unité.

\textbf{7. Décomposition des normes individuelles et des diviseurs élémentaires; composantes des idéaux fondamentaux.} Par polynôme irréductible (allemand historique \emph{Primfunktion}) on entend un polynôme irréductible relativement à $\mathsf{P}(u)$; par polynôme primaire (\emph{Primärfunktion}), une puissance d'un polynôme irréductible. Un polynôme primaire propre est un polynôme primaire qui n'est pas lui-même irréductible, c'est-à-dire une puissance d'exposant strictement supérieur à un. Une décomposition d'un polynôme en plus grands facteurs primaires est une décomposition où chaque facteur est un polynôme primaire, mais où le produit de deux facteurs quelconques n'est plus primaire.\footnote{La définition de polynôme irréductible et de polynôme primaire pourrait aussi se formuler par analogie exacte avec 11., en remplaçant seulement le mot idéal par le mot polynôme. Les plus grands facteurs primaires sont l'analogue des plus grandes composantes primaires de la décomposition en 12.} À la décomposition de la norme individuelle $R^{(i)}=N(\frakg_{i-1}\mid\frakg_i)$ en plus grands facteurs primaires $Q_\lambda^{(i)}$ correspond une représentation de $\frakg_i$ comme plus petit commun multiple
\[
 \frakg_i=[\frakr_1,\frakr_2,\ldots,\frakr_s],
\]
de telle sorte que $Q_\lambda^{(i)}=N(\frakg_{i-1}\mid\frakr_\lambda)$; $\frakr_\lambda$ possède alors $\frakg_{i-1}$ comme idéal fondamental de $(i-1)$-ième étage, et coïncide avec son idéal fondamental de $i$-ième étage; les $\frakr_\lambda$ sont appelés les composantes des idéaux fondamentaux. Le plus haut diviseur élémentaire de $\frakr_\lambda$ devient égal au plus grand diviseur commun -- au sens polynomial -- de $Q_\lambda^{(i)}$ et de $E^{(i)}$, donc égal à un plus grand facteur primaire de $E^{(i)}$. Les $\frakr_\lambda$ sont déterminés de manière unique par leur comportement indiqué ci-dessus relativement aux idéaux fondamentaux et par l'exigence que la norme individuelle ou le plus haut diviseur élémentaire devienne un plus grand facteur primaire de $R^{(i)}$, respectivement de $E^{(i)}$ (H.-N., théorème X).

\textbf{8. Notion de dimension.} On attribue à l'idéal $\frakm$ la dimension maximale $n-i$ lorsque $R^{(i)}$ est la première norme individuelle différente de l'unité -- ou, ce qui est équivalent d'après 6., lorsque $\frakg_i$ est le premier idéal fondamental différent de l'idéal unité; à l'idéal unité $\frako$ on attribue la dimension $-1$. S'il n'existe qu'une seule norme individuelle $R^{(i)}$ différente de l'unité, donc si $\frakg_i=\frakg_{i+1}=\cdots=\frakm$, alors la dimension maximale est aussi appelée dimension de $\frakm$. Si $\frakm$ contient un polynôme $G^{(i)}\ne0$, il est tout au plus de dimension maximale $n-i$, comme le montre la remarque à la fin de 6. Si $\frakm$ est de dimension maximale $n-i$, et si $\frakn$ est un diviseur de $\frakm$, alors $\frakn$ est donc lui aussi tout au plus de dimension maximale $n-i$.

\textbf{9. Norme (forme résultante) et zéros.} Par zéros d'un idéal $\frakm$ on entend tous les systèmes de valeurs de $x_1,\ldots,x_i$ appartenant à un corps d'extension algébrique convenable de $\mathsf{P}(u;x_{i-1},\ldots,x_n)$ ($i=1,2,\ldots,n$), pour lesquels -- après adjonction de $x_{i+1},\ldots,x_n$ au domaine des coefficients -- tous les polynômes de $\frakm$ s'annulent. À chaque facteur $R^{(i)}$ de la norme (forme résultante) correspondent, lors de cette adjonction, autant de zéros de $\frakm$ que l'indique le degré de $R^{(i)}$; et $R^{(i)}$, écrit en combinaisons bilinéaires $w_{\mu\nu}$ des $u_{\mu\nu}$ et d'autres indéterminées $v$, admet la décomposition explicite:
\begin{equation}
 R^{(i)}=\prod\{z-(v_1\bar{x}_{1\nu}+\cdots-v_{i-1}\bar{x}_{i-1,\nu}-v_i\bar{x}_{i,\nu})\}^{\lambda_\nu},\notag
\end{equation}
où $\bar{x}_{1\nu},\ldots,\bar{x}_{i\nu}$ désignent un système de zéros associés (H.-N., théorème XIII; une nouvelle fondation en sera donnée au \S{} 3). Parmi les zéros d'un idéal de dimension maximale $n-i$, il y en a donc qui dépendent de $(n-i)$ paramètres, mais aucun qui dépende de plus de paramètres; ainsi est démontrée la concordance de la notion de dimension donnée en 8. avec la formulation usuelle.

\textbf{10. L'anneau de la théorie générale des idéaux.} Le domaine sous-jacent sera un anneau commutatif dans lequel vaut le théorème de la chaîne finie: toute chaîne d'idéaux $\fraka,\fraka_1,\ldots,\fraka_\nu,\ldots$, où chaque $\fraka_\lambda$ est un diviseur propre -- diviseur au sens des idéaux -- du précédent immédiat, s'interrompt au bout d'un nombre fini d'étapes. Cette exigence est équivalente à l'autre, selon laquelle chaque idéal possède une base d'idéal (\emph{Idealtheorie}, théorème I); elle est donc remplie en particulier pour l'anneau de polynômes. Dans les numéros suivants 11., 12., 13., on prendra pour base cet anneau général.

\textbf{11. Idéaux premiers et idéaux primaires.} Un idéal $\frakp$ s'appelle idéal premier si la divisibilité d'un produit par $\frakp$ implique celle d'au moins un facteur; $\frakq$ s'appelle idéal primaire -- ou primaire -- si la divisibilité d'un produit par $\frakq$ implique celle d'un facteur ou d'une puissance de chaque facteur. En symboles:
\[
 a\not\equiv0\pmod{\frakp},\quad b\not\equiv0\pmod{\frakp}\quad\hbox{entraîne}\quad ab\not\equiv0\pmod{\frakp};
\]
\[
 a\not\equiv0\pmod{\frakq},\quad b^x\not\equiv0\pmod{\frakq}\ \hbox{pour toute puissance }x\quad\hbox{entraîne}\quad ab\not\equiv0\pmod{\frakq}.
\]
À tout idéal primaire $\frakq$ correspond un et un seul idéal premier associé $\frakp$, qui est diviseur de $\frakq$ et dont une puissance est divisible par $\frakq$, $\frakq\equiv0(\frakp)$, $\frakp^\rho\equiv0(\frakq)$; alors $\frakp$ se compose de tous les éléments de l'anneau dont une puissance est divisible par $\frakq$. Comme le montre la définition, tout idéal premier est en même temps un idéal primaire; par idéaux primaires propres on entend ceux qui ne sont pas en même temps des idéaux premiers, pour lesquels donc $\rho>1$.

\textbf{12. Le théorème de décomposition.} Une représentation $\fraka=[\frakq_1,\ldots,\frakq_\alpha]$ d'un idéal comme plus petit commun multiple s'appelle la plus courte si aucun $\frakq$ n'est contenu dans le plus petit commun multiple des autres -- dans son complément; elle s'appelle représentation par plus grandes composantes primaires si chaque $\frakq$ est primaire, mais si le plus petit commun multiple de deux $\frakq$ quelconques n'est pas primaire. Chaque idéal admet une représentation la plus courte par un nombre fini de plus grandes composantes primaires; pour deux telles représentations différentes, le nombre des composantes et les idéaux premiers associés, tous distincts les uns des autres, coïncident. Les idéaux premiers ainsi déterminés de manière unique seront appelés les idéaux premiers, ou idéaux premiers associés, de l'idéal (\emph{Idealtheorie}, théorème IX).

\textbf{13. Composantes isolées de la décomposition.} Un idéal $\frakr=[\frakq_{i_1},\ldots,\frakq_{i_\tau}]$ est appelé composante isolée de la décomposition de $\fraka$ lorsque les $\frakq_{i_\lambda}$ interviennent dans au moins une représentation la plus courte de $\fraka$ par plus grandes composantes primaires et satisfont à la condition qu'aucun de leurs idéaux premiers associés ne soit contenu dans l'un des autres idéaux premiers de $\fraka$. Les composantes isolées de la décomposition sont déterminées de manière unique par leurs idéaux premiers; en particulier, les plus grandes composantes primaires isolées sont donc déterminées de manière unique (\emph{Idealtheorie}, théorème XIII).

\textbf{14. Caractéristique, corps parfaits et imparfaits, extension de première espèce.} Un corps $\Omega$ est dit de caractéristique zéro lorsque le corps premier qu'il contient, dérivé de l'unité, est du type du corps des nombres rationnels; il est dit de caractéristique $p$ lorsque ce corps premier est isomorphe à $\mathbb F_p$. Un corps $\Omega$ s'appelle parfait si tout polynôme irréductible à coefficients dans $\Omega$ se décompose en facteurs linéaires distincts dans un corps d'extension convenable; dans le cas contraire, il s'appelle imparfait. Tout corps de caractéristique zéro est parfait; un corps de caractéristique $p$ ne l'est que si et seulement si, avec chaque élément, il contient aussi sa racine $p$-ième. Un polynôme irréductible sur un corps imparfait est un polynôme de degré $r$ en $x^{p^f}$, et se décompose dans un corps d'extension convenable en $p^f$-ièmes puissances de $r$ facteurs linéaires distincts; $f$ est appelé exposant. Un polynôme irréductible est dit de première espèce lorsqu'il se décompose, dans un corps d'extension convenable, en facteurs linéaires distincts, l'exposant $f$ devenant nul; une extension est dite de première espèce lorsque chaque élément devient de première espèce, c'est-à-dire zéro d'un polynôme irréductible de première espèce. Toute extension qui naît par adjonction d'un élément de première espèce est de première espèce; pour les corps parfaits, il n'y a que des extensions de première espèce (Steinitz, \S{} 11 et \S{} 13).

\textbf{15. Degré réduit et exposant d'une extension finie.} Si $\mathcal K$ est une extension finie d'un corps imparfait $\Omega$, alors $\mathcal K$ contient un sous-corps $\mathfrak{L}_0$ de degré $r$, qui devient de première espèce relativement à $\Omega$ et qui se compose de tous et seulement les éléments de première espèce de $\mathcal K$. La $p^f$-ième puissance de tout élément de $\mathcal K$ appartient à $\mathcal K_0$, et il existe dans $\mathcal K$ des éléments qui sont exactement de degré $rp^f$. $r$ est appelé degré réduit, $f$ exposant de l'extension finie (Steinitz, \S{} 14, 1 et \S{} 14, 5).
\providecommand{\frakS}{\mathfrak S}
\providecommand{\fraka}{\mathfrak a}
\providecommand{\frakb}{\mathfrak b}
\providecommand{\frakc}{\mathfrak c}
\providecommand{\frakg}{\mathfrak g}
\providecommand{\frakG}{\mathfrak G}
\providecommand{\frakh}{\mathfrak h}
\providecommand{\frakm}{\mathfrak m}
\providecommand{\frakM}{\mathfrak M}
\providecommand{\frakn}{\mathfrak n}
\providecommand{\frako}{\mathfrak o}
\providecommand{\frakp}{\mathfrak p}
\providecommand{\frakq}{\mathfrak q}
\providecommand{\frakr}{\mathfrak r}
\providecommand{\barfrakm}{\overline{\mathfrak m}}

\subsection*{\S{} 2. Forme des diviseurs élémentaires et norme des idéaux premiers et primaires. Diviseurs des idéaux premiers}

À partir de maintenant, il s'agira constamment de l'anneau de polynômes défini au \S{} 1, 1. Il faut d'abord montrer que, pour les idéaux premiers et les idéaux primaires aussi, on peut se restreindre aux idéaux transformés, d'après le

\textbf{Lemme I.} L'idéal transformé $\frakp$ d'un idéal premier $\bar{\frakp}$ dans $\mathfrak{J}$ devient un idéal premier, et le transformé $\frakq$ d'un idéal primaire $\bar{\frakq}$ dans $\mathfrak{J}$ devient un idéal primaire. En d'autres termes, la propriété d'être idéal premier ou idéal primaire se conserve par adjonction des $u_{\mu\nu}$ à $\mathsf{P}$.

Supposons en effet $\fraka\not\equiv0\pmod{\frakp}$, $\frakb\not\equiv0\pmod{\frakp}$, où $\fraka$ et $\frakb$ n'ont pas nécessairement à être des idéaux transformés; soit par exemple $a(x)\not\equiv0(\frakp)$, $b(x)\not\equiv0(\frakp)$, où $a(x)$ et $b(x)$ désignent des polynômes de $\fraka$ et de $\frakb$. Par inversion de (1) suivant (1'), on obtient -- $T$ désignant, sans restriction de généralité, des produits de puissances des $T_{\mu\nu}$ --
\[
 a(x)=\sum T_j \bar a_j(y);\qquad b(x)=\sum T_j \bar b_j(y),
\]
et il faut qu'au moins un $\bar a_j(y)$ et un $\bar b_j(y)$ ne soient pas divisibles par $\bar{\frakp}$. Soient par exemple $\bar a_r(y)\not\equiv0(\bar{\frakp})$, $\bar b_s(y)\not\equiv0(\bar{\frakp})$ les termes respectivement les plus élevés pour un ordre lexicographique quelconque des $T$; alors $\bar a_r(y)\bar b_s(y)\not\equiv0(\bar{\frakp})$ aussi, et par conséquent $a(x)b(x)\not\equiv0(\frakp)$, donc $\fraka\frakb\not\equiv0(\frakp)$, ce qui démontre que $\frakp$ est un idéal premier. La preuve repose manifestement sur le fait que l'anneau quotient par un idéal premier est un anneau intègre, propriété qui se conserve par adjonction d'indéterminées.

Supposons de même $\fraka\not\equiv0\pmod{\frakq}$, $\frakb^z\not\equiv0\pmod{\frakq}$ pour toute puissance $z$; il résulte alors de l'existence de la base d'idéal qu'il doit exister un $a(x)$ de $\fraka$ et un $b(x)$ de $\frakb$ tels que $a(x)\not\equiv0\pmod{\frakq}$, $b(x)^\varkappa\not\equiv0\pmod{\frakq}$ pour toute puissance $z$; car, dans l'hypothèse contraire, une puissance de $\frakb$ serait divisible par $\frakq$.

Posons $a(x)b(x)=c(x)$, et que $\fraka^*,\frakb^*,\frakc^*$ désignent les idéaux respectivement dérivés des coefficients $\bar a_j(y),\bar b_j(y),\bar c_j(y)$ de $a(x),b(x),c(x)$. Alors $\frakb^{*z}\not\equiv0\pmod{\bar{\frakq}}$ pour toute puissance $z$, car sinon $b(x)^\varkappa$ serait divisible par $\frakq$; comme $\fraka^*\not\equiv0(\bar{\frakq})$, on doit donc aussi avoir $\fraka^*\frakb^{*\lambda}\not\equiv0(\bar{\frakq})$ pour toute puissance $\lambda$. Or, d'après le théorème de Dedekind--Mertens,\footnote{Voir H.-N., p. 63 et note 13).} il existe un exposant $h$ tel que $\fraka^*\frakb^{*h}=\frakb^{*h-1}\frakc^*$; il faut donc que $\frakc^*\not\equiv0(\bar{\frakq})$. Cela donne $c(x)\not\equiv0(\frakq)$, et par conséquent $\fraka\frakb\not\equiv0(\frakq)$, ce qui démontre que $\frakq$ est primaire.

Même dans la représentation d'un idéal comme plus petit commun multiple, on peut se restreindre aux idéaux transformés, d'après le

\textbf{Lemme II.} D'une représentation $\barfrakm=[\frakc_1,\ldots,\frakc_s]$ dans $\mathfrak{J}$ résulte une représentation $\frakm=[\frakc_1',\ldots,\frakc_s']$ dans $\frakS'$, où $\frakc_i'$ désigne le transformé de $\frakc_\lambda$. En particulier, dans toute représentation la plus courte par plus grandes composantes primaires, celles-ci peuvent donc être supposées transformées.

En effet, $\frakm$ est divisible par $[\frakc_1,\ldots,\frakc_\alpha]$, puisque chaque polynôme $f(x)$ de $\frakm$ -- éventuellement après multiplication par une puissance convenable de $U(u)$ -- est de la forme $\sum U_i(u)f_i(y)$, où chaque $f_i(y)$, comme polynôme de $\barfrakm$, est par hypothèse divisible par tous les $\frakc_\lambda$. Réciproquement, si $f(x)$ est divisible par chaque $\frakc_i'$, alors il prend la forme ci-dessus et il est donc divisible par $\frakm$. Si l'on part donc d'une décomposition en plus grandes composantes primaires dans $\mathfrak{J}$, on obtient une décomposition $\frakm=[\frakq_1',\ldots,\frakq_s']$ dans $\frakS'$; ici les $\frakq_i'$, comme transformés d'idéaux primaires, sont primaires d'après le lemme I, et il s'agit bien de plus grandes composantes primaires, puisque des idéaux premiers associés différents $\bar{\frakp}$ dans $\mathfrak{J}$ correspondent aussi à des $\frakp$ différents dans $\frakS'$.

Il s'agit maintenant d'une première caractérisation des idéaux premiers et des idéaux primaires par la forme des diviseurs élémentaires et la norme, sans séparation complète encore entre premier et primaire; et l'on a, en exacte analogie avec la théorie des idéaux dans les corps de nombres algébriques,\footnote{Comme plus haut diviseur élémentaire des idéaux dans les corps de nombres algébriques, il faut concevoir le plus petit entier rationnel divisible par l'idéal, donc en particulier le nombre premier divisible par un idéal premier, ou la puissance première divisible par un idéal primaire (puissance d'idéal premier). En fait, chaque idéal $\fraka^*$ est en même temps un module $\mathfrak{A}^*$ de formes linéaires dans les éléments $\omega_1,\ldots,\omega_n$ d'une base du corps, $(\omega_1,\ldots,\omega_n)$ étant le module fondamental de $\mathfrak{A}^*$; le plus petit entier rationnel divisible par $\fraka^*$ devient donc le plus haut diviseur élémentaire de ce module $\mathfrak{A}^*$, tandis que la norme de $\fraka^*$ devient le produit de tous les diviseurs élémentaires de $\mathfrak{A}^*$.}

\textbf{Théorème I.} La forme des diviseurs élémentaires d'un idéal premier devient égale à un polynôme irréductible, la norme à une puissance de ce polynôme irréductible. Pour les idéaux primaires, forme des diviseurs élémentaires et norme deviennent des polynômes primaires, à savoir des puissances du même polynôme irréductible, qui est lui-même la forme des diviseurs élémentaires de l'idéal premier associé. Pour les idéaux premiers et les idéaux primaires, la dimension maximale coïncide avec la dimension tout court; cette dimension est la même pour l'idéal primaire et l'idéal premier associé.

Le théorème I est manifestement satisfait pour l'idéal unité, dont la forme des diviseurs élémentaires et la norme deviennent égales à l'unité. Soit maintenant l'idéal premier $\frakp$ de dimension maximale $(n-i)$; d'après (4), \S{} 1, 6., on obtient
\[
 E^{(n)}\cdots E^{(i+1)}\frakg_i\equiv0(\frakp);\quad\hbox{mais}\quad
 E^{(n)}\cdots E^{(i+1)}\not\equiv0(\frakp),
\]
car sinon $\frakp$ serait, d'après \S{} 1, 8., tout au plus de dimension maximale $n-i-1$. Ainsi $\frakg_i\equiv0(\frakp)$, donc $\frakp$ est identique à son idéal fondamental de $i$-ième étage, et par conséquent aussi à $\frakg_{i+1},\frakg_{i+2},\ldots$. Donc $R^{(i+1)},\ldots,R^{(n)}$ deviennent égales à l'unité, et par suite aussi $E^{(i+1)},E^{(i+2)},\ldots$, comme diviseurs de $R^{(i+1)},R^{(i+2)},\ldots$; la forme des diviseurs élémentaires devient ainsi égale à $E^{(i)}$, qui doit être égal à un polynôme irréductible $P^{(i)}$. En effet, d'après \S{} 1, 5., et compte tenu du fait que $\frakg_{i-1}$ devient l'idéal unité, $E^{(i)}$ est défini comme plus grand diviseur commun -- au sens polynomial -- de tous les $B^{(i)}$ divisibles par $\frakp$; de $E^{(i)}=C_1^{(i)}C_2^{(i)}\equiv0\pmod{\frakp}$ il résulterait donc que $E^{(i)}$ doit diviser l'un de ses facteurs $C_1^{(i)}$ ou $C_2^{(i)}$. Mais comme une puissance de $E^{(i)}$ est divisible par $R^{(i)}$, $R^{(i)}$ devient une puissance -- qui peut aussi être la première -- de ce polynôme irréductible $P^{(i)}$; en même temps, $R^{(i)}$ devient la norme de $\frakp$. Comme il n'apparaît qu'une seule norme individuelle différente de l'unité, la dimension maximale de $\frakp$ coïncide enfin avec la dimension tout court.

Soit de même l'idéal primaire $\frakq$ de dimension maximale $n-i$; alors, comme ci-dessus,
\[
 E^{(n)}\cdots E^{(i+1)}\frakg_i\equiv0\pmod{\frakq},\quad
 (E^{(n)}\cdots E^{(i+1)})^\varkappa\not\equiv0\pmod{\frakq}
\]
pour toute puissance $z$; et par conséquent, comme ci-dessus, $\frakq=\frakg_i$, sa forme des diviseurs élémentaires est égale à $E^{(i)}$, sa norme à $R^{(i)}$, et sa dimension maximale à la dimension tout court. Cette dimension coïncide avec celle de l'idéal premier associé; car des relations $\frakp^\rho\equiv0(\frakq)$ et $\frakq\equiv0(\frakp)$, il résulte que la dimension $(n-i)$ de $\frakp$ est tout au plus égale à celle de $\frakq$, et celle de $\frakq$ tout au plus égale à celle de $\frakp^\rho$; de $(P^{(i)})^\rho\equiv0(\frakp^\rho)$, où $P^{(i)}$ désigne la forme des diviseurs élémentaires de $\frakp$, il résulte que cette dernière dimension maximale est tout au plus $n-i$, et ainsi $n-i$ devient la dimension de $\frakp$ et de $\frakq$. De $(P^{(i)})^\rho\equiv0(\frakq)$ on déduit en outre que $E^{(i)}$ -- comme plus grand diviseur commun de tous les $B^{(i)}$ issus de $\frakq$ -- devient une puissance de $P^{(i)}$, qui peut aussi être la première, et que le même fait vaut pour $R^{(i)}$; le théorème I est ainsi démontré dans toutes ses parties.

Une caractérisation des idéaux premiers au moyen de la notion de dimension est donnée par le

\textbf{Théorème II.} Un idéal premier ne possède aucun diviseur propre de même dimension maximale; et réciproquement, tout idéal ayant cette propriété est un idéal premier.

La preuve repose sur le

\textbf{Lemme III.} Si l'anneau quotient par un idéal premier $\frakp$ de polynômes à coefficients dans un corps $\Omega$ est de dimension finie sur $\Omega$, alors $\frakp$ ne possède aucun diviseur propre différent de l'idéal unité; en d'autres termes, cet anneau quotient est un corps.

L'hypothèse signifie qu'il existe un nombre fini $k$ tel qu'entre $(k+1)$ classes modulo $\frakp$ quelconques il existe une dépendance linéaire à coefficients dans $\Omega$ -- plus précisément dans la copie $(\Omega)$ du corps des coefficients dans le quotient --, mais qu'il existe au moins un système de $k$ classes modulo $\frakp$ linéairement indépendantes sur $\Omega$; il en résulte aussitôt que toutes les classes modulo $\frakp$ peuvent s'exprimer linéairement par un système quelconque de $k$ classes linéairement indépendantes. Soit maintenant $\fraka$ un diviseur propre de $\frakp$, et soit $a(z)\not\equiv0(\frakp)$ un polynôme de $\fraka$; la relation $c_1f_1(z)+\cdots+c_kf_k(z)\not\equiv0\pmod{\frakp}$ entraîne alors aussi
\[
 c_1a(z)f_1(z)+\cdots+c_ka(z)f_k(z)\not\equiv0\pmod{\frakp};
\]
cela signifie que, outre les classes modulo $\frakp$ de $f_1,\ldots,f_k$, celles de $af_1,\ldots,af_k$ forment aussi un système de $k$ classes linéairement indépendantes, par lesquelles en particulier la classe unité se représente linéairement. Ainsi $\fraka$ devient égal à l'idéal unité; autrement dit, l'anneau quotient est un corps, puisqu'à tout $a(z)\not\equiv0(\frakp)$ correspond toujours un $f(z)$ tel que $1\equiv a(z)f(z)(\frakp)$.

Pour démontrer la première partie du théorème II, il reste seulement à remarquer que, pour tout idéal premier de dimension $(n-i)$ -- plus généralement pour tout idéal de dimension maximale $(n-i)$ --, l'anneau quotient est de dimension finie sur $\mathsf{P}(u;x_{i+1},\ldots,x_n)$, à savoir (\S{} 1, 5.) de la dimension indiquée par le degré de $R^{(i)}$, puisque $\frakg_{i-1}$ devient ici l'idéal unité. Tout diviseur propre $\fraka$ de $\frakp$ passe donc, par adjonction de $x_{i+1},\ldots,x_n$ à $\mathsf{P}(u)$, en l'idéal unité; autrement dit, l'idéal fondamental de $i$-ième étage de $\fraka$ devient égal à l'idéal unité; cela signifie que la dimension maximale de $\fraka$ devient au plus $n-i-1$. Pour que l'énoncé vaille aussi lorsqu'un idéal non transformé $\fraka$ intervient comme diviseur, il suffit, d'après \S{} 1, 2., de passer à un nouvel anneau transformé.

Supposons maintenant réciproquement que $\frakp$ ne possède aucun diviseur propre de même dimension maximale $n-i$; et soient $\fraka\not\equiv0(\frakp)$, $\frakb\not\equiv0(\frakp)$, où $\fraka$ et $\frakb$ sont de nouveau supposés transformés. Alors, par hypothèse, comme $(\fraka,\frakp)$ et $(\frakb,\frakp)$ deviennent, en tant que plus grands communs diviseurs -- au sens des idéaux --, des diviseurs propres de $\frakp$, on obtient
\[
 G^{(i+1)}\equiv0(\fraka,\frakp);\qquad
 H^{(i+1)}\equiv0(\frakb;\frakp);\qquad
 G^{(i+1)}\neq0;\qquad
 H^{(i+1)}\neq0.
\]
Mais cela donne
\[
 G^{(i+1)}H^{(i+1)}\equiv0(\fraka\frakb,\frakp);\qquad
 G^{(i+1)}H^{(i+1)}\neq0;
\]
et l'on en déduit nécessairement que $\fraka\frakb\not\equiv0(\frakp)$, car sinon $\frakp$ serait de dimension maximale inférieure à $n-i$. Comme $\fraka,\frakb$ sont des idéaux transformés, on obtient donc que $\bar{\frakp}$, et d'après le lemme I aussi $\frakp$, est un idéal premier, ce qui démontre le théorème II.
\providecommand{\frakS}{\mathfrak S}
\providecommand{\fraka}{\mathfrak a}
\providecommand{\frakg}{\mathfrak g}
\providecommand{\frakp}{\mathfrak p}
\providecommand{\frakq}{\mathfrak q}

\subsection*{\S{} 3. Zéros des idéaux premiers et primaires}

Le passage à une fondation directe de la théorie des zéros (\S{} 1, 9.), d'abord pour les idéaux premiers, est assuré par le

\textbf{Lemme IV.} L'anneau quotient par tout idéal premier différent de l'idéal unité peut être étendu, par formation de quotients, en son corps résiduel.

Les classes modulo un idéal quelconque forment un anneau quotient, puisque somme, différence et produit sont encore bien définis, et que les lois associative, commutative et distributive se conservent lors du passage aux classes. Pour un idéal premier, cet anneau est intègre; car la définition des idéaux premiers (\S{} 1, 11.) dit que le produit de deux classes non nulles ne s'annule pas. Si l'idéal premier est différent de l'idéal unité, cet anneau contient au moins un élément non nul et peut donc être étendu à un corps par formation de quotients -- adjonction de couples d'éléments;\footnote{Steinitz, \S{} 3. L'hypothèse d'un élément non nul dans le domaine d'intégrité primitif suffit, puisque l'existence de l'unité est alors assurée par formation de quotients, et qu'ainsi le domaine primitif devient un sous-domaine du corps; Steinitz suppose l'existence de l'unité dans le domaine d'intégrité.} ainsi le corps résiduel est défini.

On placera d'abord en tête la

\textbf{Notation:} Les classes modulo l'idéal considéré seront en général désignées en mettant entre parenthèses un représentant quelconque, $(x),\ldots,(a(x)),\ldots$; de même, les corps résiduels ainsi formés seront désignés par des parenthèses. En particulier, $(P)$, $(\mathsf{P}(u))$, $(\mathsf{P}(u,x_{i+1},\ldots,x_n))$ désignent les sous-corps constitués exactement des classes représentables par des éléments de $P$, $\mathsf{P}(u)$, $P(u,x_{i+i},\ldots,x_n)$.

On placera encore en tête le fait suivant. Un corps $\mathcal R$ est dit de rang algébrique (degré de transcendance) $k$ relativement à un corps de base $\Omega$ si $\mathcal R$ contient au moins un système de $k$ grandeurs algébriquement indépendantes relativement à $\Omega$ -- des grandeurs transcendantes --, tandis que $k+1$ grandeurs quelconques de $\mathcal R$ sont algébriquement dépendantes relativement à $\Omega$. Si $\mathcal R$ peut être représenté comme corps d'extension algébrique d'un sous-corps $\Omega(z_1,\ldots,z_s)$, où les $z$ sont algébriquement indépendants -- transcendants -- relativement à $\Omega$, alors $s$ devient nécessairement égal à $k$;\footnote{Pour une preuve de ce fait connu, valable en caractéristique quelconque, voir Steinitz, \S{} 22, 9. -- Lors de l'adjonction des $u$, le rang algébrique ne peut pas croître, puisqu'il s'agit de combinaisons rationnelles des éléments primitifs à coefficients dans $\Omega(u)$; il ne peut pas décroître non plus, puisque toute relation entre éléments de $\mathcal R$ à coefficients dans $\Omega(u)$ se décompose en un nombre fini de relations à coefficients dans $\Omega$.} de même, le rang algébrique de $\mathfrak{K}(u)$ relativement à $\Omega(u)$ devient égal à $k$ lorsque $u$ désigne des indéterminées non contenues dans $\mathcal R$.

La construction du corps des zéros se fait maintenant en complète analogie avec la voie usuelle pour les polynômes irréductibles d'une indéterminée\footnote{Steinitz, \S{} 6.} selon le

\textbf{Théorème III.} Au corps résiduel $(\mathfrak{K})$ d'un idéal premier $\frakp$ de dimension $n-i$ peut être associé isomorphiquement un corps des zéros $R$ de $\frakp$, qui devient de rang algébrique $n-i$ -- dépend de $n-i$ paramètres -- et peut en particulier être conçu comme corps d'extension fini de $\mathsf{P}(u,x_{i-1},\ldots,x_n)$. L'isomorphisme entre $(\mathcal R)$ et $R$ comprend celui entre $(\mathsf{P})$ et $\mathsf{P}$, $(\mathsf{P}(u))$ et $\mathsf{P}(u)$ et, si l'on prend $x_{i+1},\ldots,x_n$ pour paramètres, aussi celui entre $(\mathsf{P}(u,x_{i+1},\ldots,x_n))$ et $\mathsf{P}(u,x_{i+1},\ldots,x_n)$. Réciproquement, tout corps des zéros de rang algébrique $n-i$ -- dépendant de $n-i$ paramètres -- est isomorphe à $(\mathcal R)$.

Le théorème III donne comme corollaire le théorème connu: si un idéal $\fraka$ s'annule en tous les zéros d'un idéal premier $\frakp$, alors $\fraka$ est divisible par $\frakp$.

Pour la preuve, il faut d'abord remarquer que le corps résiduel $(\mathcal R)$ devient de rang algébrique $n-i$ relativement à $(\mathsf{P}(u))$. En effet, les classes $(x_{i+1}),\ldots,(x_n)$ sont algébriquement indépendantes relativement à $(\mathsf{P}(u))$, puisque $\frakp$, étant de dimension $n-i$, ne contient aucun polynôme libre de $x_1,\ldots,x_i$; toute autre classe, en revanche, dépend de celles-ci. Car, d'après \S{} 1, 2., l'appartenance de la forme des diviseurs élémentaires $P^{(i)}$ à $\frakp$ entraîne aussi, pour $\lambda=1,2,\ldots,i$, l'appartenance de polynômes dépendant respectivement de $x_\lambda,x_{i+1},\ldots,x_n$. Donc $(\mathcal R)$ devient un corps d'extension fini de $(\mathsf{P}(u,x_{i+1},\ldots,x_n))$; le degré relativement à ce sous-corps devient, d'après \S{} 1, 5. -- compte tenu de ce que $\frakg_{i-1}$ est égal à l'idéal unité et que $\frakg_i$, d'après le théorème I, devient égal à $\frakp$ -- égal au degré de la norme.

Pour passer à un corps des zéros isomorphe, il faut encore remarquer que $\mathsf{P}(u,x_{i+1},\ldots,x_n)$ et $(\mathsf{P}(u,x_{i+1},\ldots,x_n))$ sont mis en rapport isomorphe par l'association, à chaque élément de $\mathsf{P}(u,x_{i+1},\ldots,x_n)$, de sa classe modulo $\frakp$. Car, tout d'abord, par cette association, les domaines d'intégrité formés des polynômes en $x_{i+1},\ldots,x_n$ et respectivement en $(x_{i+1}),\ldots,(x_n)$ se correspondent isomorphiquement, puisque $\frakp$ ne contient aucun polynôme libre de $x_1,\ldots,x_i$, et que des polynômes différents de $\mathsf{P}(u,x_{i+1},\ldots,x_n)$ correspondent donc aussi à des classes différentes; or, de domaines d'intégrité isomorphes proviennent, par formation de quotients, des corps isomorphes. Le rapport isomorphe entre $\mathsf{P}(u,x_{i+1},\ldots,x_n)$ et $(\mathsf{P}(u,x_{i+1},\ldots,x_n))$ comprend celui entre $\mathsf{P}$ et $(\mathsf{P})$, $\mathsf{P}(u)$ et $(\mathsf{P}(u))$. Pour étendre cet isomorphisme à un isomorphisme entre $(\mathcal R)$ et un corps des zéros $R$, associons isomorphiquement aux classes $(x_1),\ldots,(x_i)$ certains nouveaux éléments $\bar{x}_1,\ldots,\bar{x}_i$; c'est-à-dire qu'à chaque polynôme $(a(x))$ corresponde un polynôme $a(\bar{x}_1,\ldots,\bar{x}_i,x_{i+1},\ldots,x_n)$; à la somme et au produit correspondent la somme et le produit. Par formation de quotients -- où l'on peut se restreindre aux dénominateurs du sous-corps $(\mathsf{P}(u,x_{i+1},\ldots,x_n))$, respectivement de $\mathsf{P}(u,x_{i+1},\ldots,x_n)$ -- il correspond alors à $(\mathcal R)$ un corps $R$, extension finie de $\mathsf{P}(u,x_{i+1},\ldots,x_n)$ de degré égal à celui de la norme, que l'on peut appeler corps des zéros; car l'association dit que $f(\bar{x}_1,\ldots,\bar{x}_i,x_{i+1},\ldots,x_n)$ s'annule dès que $f(x)\equiv0(\frakp)$. À la place de $(\mathsf{P}(u,x_{i+1},\ldots,x_n))$, on pourrait aussi faire intervenir, dans tout le raisonnement, n'importe quel autre sous-corps de $(\mathcal R)$ né par adjonction de $n-i$ grandeurs algébriquement indépendantes.

Réciproquement, il est facile de montrer que tout corps des zéros $R$ de rang algébrique $n-i$ devient isomorphe à $(\mathcal R)$. Comme $R$ peut être déduit rationnellement, avec coefficients dans $\mathsf{P}(u)$, des grandeurs $\xi_1,\ldots,\xi_n$, il doit se trouver parmi ces grandeurs $n-i$ éléments algébriquement indépendants, donc transcendants relativement à $\mathsf{P}(u)$. En particulier, cela doit être vérifié pour $\xi_{i+1},\ldots,\xi_n$, puisque, à cause de $P^{(i)}\equiv0\pmod{\frakp}$, on déduit comme ci-dessus que $R$ devient un corps d'extension algébrique de $\mathsf{P}(u,\xi_{i+1},\ldots,\xi_n)$. Comme, par hypothèse, $f(\xi)$ s'annule pour tout polynôme $f(x)$ de $\frakp$, à chaque classe de l'anneau de polynômes en les $(x)$ contenu dans $(\mathcal R)$ correspond un et un seul élément de $R$, qui est un polynôme en les $\xi$; et la somme et le produit correspondent à la somme et au produit. Mais à des classes différentes correspondent aussi des éléments différents de $R$; sinon, à une classe non nulle, donc après une multiplication convenable aussi à une classe représentée par un polynôme de $(\mathsf{P}(u,x_{i+1},\ldots,x_n))$, correspondrait l'élément nul dans $R$, et il existerait, contrairement à l'hypothèse, une dépendance algébrique entre les grandeurs $\xi_{i+1},\ldots,\xi_n$. Par cette association, tous les polynômes des $\xi$ contenus dans $R$ sont épuisés; mais, par formation de quotients -- où l'on peut de nouveau se restreindre aux dénominateurs de $(\mathsf{P}(u,x_{i+1},\ldots,x_n))$, respectivement de $\mathsf{P}(u,\bar{x}_{i+1},\ldots,\bar{x}_n)$ -- ces domaines d'intégrité isomorphes donnent des corps isomorphes.

Supposons maintenant que l'idéal $\fraka$ s'annule en tous les zéros de $\frakp$, donc en particulier en tous ceux qui dépendent de $n-i$ paramètres. Si $a(x)$ est alors un polynôme quelconque de $\fraka$, $a(\xi)$ s'annule donc; à cause de l'isomorphisme entre $R$ et $(\mathcal R)$, $(a(x))$ devient par conséquent la classe nulle, et $a(x)$, donc $\fraka$, est divisible par $\frakp$. Ainsi le théorème III et son corollaire sont démontrés.

Pour passer du théorème III à la décomposition donnée au \S{} 1, 9., d'abord pour les idéaux premiers, et pour pouvoir faire des énoncés plus précis sur les exposants, il faut une étude supplémentaire de la forme des diviseurs élémentaires, selon le

\textbf{Lemme V.} Si, en caractéristique $p$, la forme des diviseurs élémentaires $E^{(i)}$ d'un idéal premier ou primaire -- plus généralement d'un idéal $\fraka$ pour lequel la dimension maximale coïncide avec la dimension tout court -- est d'exposant $f$ relativement à $x_i$, donc un polynôme en $x_i^{p^f}$, alors elle est aussi d'exposant $f$ relativement à $x_{i+1},\ldots,x_n$ et relativement aux indéterminées $u_{\mu\nu}$, respectivement $t_{\mu\nu}$, qui interviennent dans $E^{(i)}$. En particulier, lorsque l'anneau des coefficients est un corps parfait, la forme des diviseurs élémentaires $P^{(i)}$ d'un idéal premier est toujours d'exposant nul relativement à $x_i$, donc un polynôme irréductible de première espèce.

Pour cela, il faut remarquer que $\mathsf{P}(u,x_{i+1},\ldots,x_n)$ devient imparfait dès que $\mathsf{P}$ est un corps parfait de caractéristique $p$, de sorte qu'on ne peut pas conclure directement du \S{} 1, 14. que $P^{(i)}$ soit de première espèce. Supposons maintenant que $E^{(i)}$ soit d'exposant $f$ relativement à $x_i$, mais qu'elle contienne une autre indéterminée $x_{i+\lambda}$ avec l'exposant $f'<f$. Alors il existe aussi -- par permutation des $u_{\mu\nu}$, respectivement des $t_{\mu\nu}$ -- d'après \S{} 1, 2., un polynôme non nul $G^{(i)}$, divisible par l'idéal, qui devient un polynôme en $x_i^{p^{f'}}$, et qui, d'après la définition de la forme des diviseurs élémentaires, est divisible par $E^{(i)}$. Mais comme $G^{(i)}$ devient de même degré que $E^{(i)}$, il doit coïncider avec $E^{(i)}$ à un facteur de $\mathsf{P}(u)$ près, et ainsi $f'$ devient nécessairement égal à $f$. Il en résulte que $E^{(i)}$ -- que l'on peut supposer entière et primitive dans les $t_{\mu\nu}$ -- est de la forme
\[
 E^{(i)}=\sum c_\lambda(t)\,x_i^{\lambda_i p^f}\cdots x_n^{\lambda_n p^f}
       =\sum c_\lambda(t)\,g_\lambda(y^{p^f},t^{p^f})
       =\sum T_\mu(t)\,\bar h_\mu(y^{p^f}),
\]
où les $T_\mu$ désignent des produits de puissances des $t$, et où les $\bar h$ sont divisibles par $\bar{\fraka}$. On obtient donc un polynôme appartenant également à $\fraka$ si, dans les $T_\mu$, on remplace chaque exposant d'un $t$ par le plus grand multiple de $p^f$ qu'il contient; cela revient, comme le montre la représentation ci-dessus, à remplacer dans $c_\lambda(t)$ chaque exposant des $t$ par le plus grand multiple de $p^f$ qu'il contient. Par cette substitution, $E^{(i)}$ passe à un polynôme en $t,x_i,\ldots,x_n$ qui, par définition, est divisible par $E^{(i)}$, et cela aussi relativement aux $t$, en raison de la primitivité supposée; il doit donc être identique à $E^{(i)}$, puisque le degré n'a augmenté en aucun argument. Si, en particulier, $\mathsf{P}$ est un corps parfait, $E^{(i)}$, comme polynôme en $x^{p^f},t^{p^f}$, devient donc une $p^f$-ième puissance; s'il s'agit par conséquent d'un idéal premier, il faut nécessairement que $f$ soit nul; la forme des diviseurs élémentaires $P^{(i)}$ devient un polynôme irréductible de première espèce.

La décomposition est maintenant donnée par le

\textbf{Théorème IV.} La forme des diviseurs élémentaires $P^{(i)}$ d'un idéal premier admet, dans le corps de Galois dérivé du corps des zéros, la décomposition explicite
\[
 P^{(i)}=\prod (x_i-t_{1i}\bar y_{1\nu}-\cdots-t_{ni}\bar y_{n\nu})^\delta
 =\prod (t_{1i}(y_1-\bar y_{1\nu})+\cdots+t_{ni}(y_n-\bar y_{n\nu}))^\delta,
\]
où $\bar y_{1\nu},\ldots,\bar y_{n\nu}$ représentent chaque fois un système associé de zéros des indéterminées primitives $y$, indépendant de $t_{1i},\ldots,t_{ni}$, où des $\nu$ distincts correspondent à des facteurs linéaires distincts, et où l'exposant $\delta$ vaut un si $\mathsf{P}$ est parfait, et vaut $p^f$ ($f\geqq0$) si $\mathsf{P}$ est imparfait. Par là est aussi déterminée, après adjonction de nouvelles indéterminées $v$, la décomposition de la forme des diviseurs élémentaires écrite en combinaisons bilinéaires des $u$ et des $v$ à la place des $u_{\mu\nu}$:
\[
 P^{(i)}(u,v)=\prod (z-v_1\bar x_{1\nu}-\cdots-v_i\bar x_{i\nu})^\delta,
\]
où $\delta$ a la même signification que ci-dessus. De même est donnée la décomposition de la norme d'un idéal premier, ou d'un idéal primaire ayant $\frakp$ pour idéal premier associé -- comme puissance de $P^{(i)}$.

Comme corollaire du théorème IV, on obtient encore: si deux idéaux premiers coïncident dans leur forme des diviseurs élémentaires $P^{(i)}$, ils sont identiques.

La preuve du théorème IV reviendra à montrer que la forme des diviseurs élémentaires $P^{(i)}$ est identique à l'équation fondamentale du corps d'extension $(\mathsf{P}(u,x_{i+1},\ldots,x_n),(x_i))$ sur $(\mathsf{P}(u,x_{i+1},\ldots,x_n))$. On passe d'abord -- pour saisir la dépendance à l'égard des indéterminées $u_{\mu\nu}$, respectivement $t_{\mu\nu}$ -- au corps résiduel $(\bar{\mathfrak{K}})$ de $\bar{\frakp}$, qui, d'après ce qui a été dit lors de la définition du rang algébrique, doit être de rang algébrique $n-i$ relativement à $(\bar{\mathsf{P}})$, où $(\bar{\mathsf{P}})$ désigne le sous-corps représenté par les éléments de $\mathsf{P}$. Car $(\mathfrak{K})$ naît par adjonction des indéterminées $u$, respectivement $t$, à un sous-corps isomorphe à $(\bar{\mathfrak{K}})$, obtenu par formation de quotients à partir de toutes les classes, et seulement de celles, qui peuvent être représentées par des polynômes de $\bar{\mathfrak{J}}$, où $(\bar{\mathsf{P}})$ et $(\mathsf{P})$ se correspondent. Le rang algébrique de $(\mathfrak{K})$ relativement à $(\mathsf{P}(u))$ est donc égal à celui de ce sous-corps relativement à $(\mathsf{P})$, et par conséquent, en raison de l'isomorphisme, égal à celui de $(\bar{\mathfrak{K}})$ relativement à $(\bar{\mathsf{P}})$. Comme $(\bar{\mathfrak{K}})$ se déduit rationnellement, avec coefficients dans $(\bar{\mathsf{P}})$, des classes $(y_1),\ldots,(y_n)$, il doit donc y avoir parmi ces classes $(n-i)$ classes algébriquement indépendantes, et $(\bar{\mathfrak{K}})$ naît comme corps d'extension algébrique fini du sous-corps engendré par l'adjonction de ces $(n-i)$ classes à $(\bar{\mathsf{P}})$.\footnote{Cette remarque montre que le théorème III avec son corollaire aurait pu être énoncé et démontré directement pour le corps des zéros $\bar{\mathfrak{K}}$ de $\bar{\frakp}$. Mais ce n'est que par le passage à l'idéal transformé que l'on obtient que, dans les plus grandes composantes primaires de même dimension d'un idéal quelconque, les mêmes paramètres apparaissent parmi les zéros; et c'est seulement ainsi que résulte le lien avec la norme et la forme des diviseurs élémentaires d'un idéal quelconque.}

Si l'on adjoint seulement les indéterminées $t_{\sigma\tau}$ qui interviennent dans $x_{i+1},\ldots,x_n$, il naît de nouveau un corps résiduel, qui contient les classes $(y_1),\ldots,(y_n)$ et $(x_{i+1}),\ldots,(x_n)$ et qui devient un corps d'extension algébrique fini de $(\mathsf{P}(t_{\sigma\tau},x_{i+1},\ldots,x_n))$ -- il s'agit alors de corps isomorphes à des sous-corps de $(\overline{\mathcal R})$, non identiques à eux; car, pour des indéterminées différentes, les classes sont certes représentées par les mêmes éléments, mais elles ne sont pas identiques, puisqu'elles ne contiennent pas les mêmes éléments. Dans la dépendance des classes $(y)$ par rapport à $(\mathsf{P}(t_{\sigma\tau},x_{i+1},\ldots,x_n))$, seules les indéterminées $t_{\sigma\tau}$ interviennent donc; formons maintenant, par adjonction de $t_{1i},\ldots,t_{ni}$, la classe $(x_i)=(t_{1i}y_1+\cdots+t_{ni}y_n)$; et soient $(x_{i\nu})=(t_{1i}y_{1\nu}+\cdots+t_{ni}y_{n\nu})$ les $r$ valeurs conjuguées distinctes situées dans le corps de Galois relatif à $(\mathsf{P}(t_{\sigma\tau},t_{1i},\ldots,t_{ni},x_{i+1},\ldots,x_n))$. Alors -- selon qu'il s'agit ou non d'extensions de première espèce -- le produit sur les facteurs $t-(x_{i\nu})$, ou sur leurs $p^f$-ièmes puissances, devient un polynôme irréductible $(G)$ relativement à $(\mathsf{P}(t,x_{i+1},\ldots,x_n))$; car, puisqu'il existe des éléments de degré $r\cdot p^f$ (\S{} 1, 15.), $(x_i)$ doit nécessairement être un tel élément, puisque tout élément naît par spécialisation des $t_{1i},\ldots,t_{ni}$; mais $(x_i)$ devient zéro de $(G)$. Si l'on passe au corps des zéros isomorphe de $\frakp$ et au corps de Galois qui en est dérivé, $G$ admet donc une décomposition en facteurs linéaires $t-t_{1i}\bar y_{1\nu}-\cdots-t_{ni}\bar y_{n\nu}$, respectivement en leurs $p^f$-ièmes puissances. Or, comme $(G)$ s'annule pour $t=(x_i)$, $G(x_i)$ est divisible par $\frakp$ et donc par $P^{(i)}$; et, comme polynôme irréductible relativement à $\mathsf{P}(u,x_{i+1},\ldots,x_n)$ et primitif en $x_{i+1},\ldots,x_n$, $G$ doit devenir identique à $P^{(i)}$; la première décomposition du théorème IV est ainsi démontrée.

Pour passer à la seconde décomposition, il faut remarquer que la propriété d'être, ou de ne pas être, une extension de première espèce se conserve par adjonction d'indéterminées. Si l'on forme donc, après adjonction de $v_1,\ldots,v_i$ et de tous les $t_{\mu\nu}$, la classe $(z)=(v_1x_1+\cdots+v_ix_i)$ et les classes conjuguées $(z_\nu)=(v_1x_{1\nu}+\cdots+v_ix_{i\nu})$, alors le produit sur tous les $t-(z_\nu)$, respectivement sur leurs $p^f$-ièmes puissances, devient de nouveau un polynôme irréductible $(H)$ qui s'annule pour $t=(z)$. Ainsi $H$ est représentable comme produit de facteurs linéaires, comme ci-dessus, et s'identifie avec la forme des diviseurs élémentaires de l'idéal premier déduit de $\frakp$ par composition de la transformation (1) avec $z=v_1x_1+\cdots+v_ix_i$; cette forme des diviseurs élémentaires naît cependant -- à cause de la spécialisation réciproque -- de $P^{(i)}$ par remplacement des $u_{\mu\nu}$ par certaines combinaisons bilinéaires des $u$ et des $v$.\footnote{H.-N., \S{} 7.} Avec ces décompositions est également donnée la décomposition de la norme comme puissance de $P^{(i)}$,\footnote{Il s'agit, pour les extensions de première espèce, donc en particulier pour les corps parfaits, de la première puissance; pour les corps imparfaits, en revanche, de la $p^g$-ième puissance, comme on le montrera au \S{} 6, théorèmes XIII et XIV.} de même que celle de la norme et de la forme des diviseurs élémentaires d'un idéal primaire ayant $\frakp$ pour idéal premier associé.

Enfin, si deux idéaux premiers coïncident dans la forme des diviseurs élémentaires $P^{(i)}$, alors, par suite de la décomposition ci-dessus, ils coïncident dans leurs zéros; ils sont donc mutuellement divisibles l'un par l'autre et par conséquent identiques.
\providecommand{\frakS}{\mathfrak S}
\providecommand{\fraka}{\mathfrak a}
\providecommand{\frakb}{\mathfrak b}
\providecommand{\frakc}{\mathfrak c}
\providecommand{\frakm}{\mathfrak m}
\providecommand{\frako}{\mathfrak o}
\providecommand{\frakp}{\mathfrak p}
\providecommand{\frakq}{\mathfrak q}

\subsection*{\S{} 4. Formulation arithmétique de la notion de dimension.}

Une première formulation de la notion de dimension, indépendante de l'introduction des idéaux transformés, est donnée par le

\textbf{Théorème V.} \emph{La dimension d'un idéal premier est donnée par le rang algébrique du corps résiduel correspondant. La dimension maximale d'un idéal est égale au plus grand des nombres de dimension de ses idéaux premiers associés.}

La première partie du théorème V avait été démontrée dans la première remarque relative au théorème IV, où l'on avait montré que le rang algébrique du corps résiduel $(\overline{\mathcal R})$ de $\bar{\frakp}$ devient égal à celui du corps résiduel $(\mathcal R)$ de $\frakp$, donc égal à la dimension de $\frakp$, en accord avec \S{} 1, 8. Pour démontrer la seconde partie, soit $\frakm=[\frakq_1,\ldots,\frakq_\alpha]$ une représentation la plus courte par plus grandes composantes primaires; alors, tout d'abord, aucun $\frakq$, comme diviseur de $\frakm$, ne peut posséder une dimension plus élevée que $\frakm$. Mais le produit de tous les $\frakq$ est divisible par $\frakm$, et donc aussi le produit de toutes les formes des diviseurs élémentaires des $\frakq$ pris séparément; si donc $n-i$ est le plus grand des nombres de dimension des $\frakq$, alors $\frakm$ contient un polynôme $G^{(i)}\ne0$ et ne peut donc pas être de dimension plus élevée que $n-i$. Or les nombres de dimension des $\frakq$ coïncident, d'après le théorème I, avec ceux de leurs idéaux premiers associés. Si l'on définit, sans introduire d'idéaux transformés, la dimension d'un idéal premier par le rang algébrique de ce corps résiduel, la seconde partie du théorème V donne donc la définition de la dimension maximale d'un idéal quelconque.

Pour saisir d'autre part la notion de dimension par des chaînes d'idéaux premiers, ce qui devient de nouveau indépendant de l'introduction d'idéaux transformés, il faut montrer, en complément du théorème II:

\textbf{Théorème VI.} \emph{Tout idéal premier différent de l'idéal unité possède -- après adjonction éventuelle d'indéterminées -- au moins un idéal premier comme diviseur, dont la dimension a diminué exactement d'une unité.}

En effet, si $\frakp$ est de dimension $(n-i)>0$ et si $v$ désigne une nouvelle indéterminée, alors
\[
 \frakc=(\frakp,x_i-v)
\]
est un idéal ayant la propriété cherchée. Cela résulte d'une correspondance isomorphe; on a $\frakc=(\frakp^*,x_i-v)$, où $\frakp^*$ est né de $\frakp$ en remplaçant $x_i$ par l'indéterminée $v$ et en adjoignant $v$ au domaine des coefficients; de même, les classes modulo l'idéal naissent par remplacement de la classe $(x_i)$ par $(v)$. Or $\frakp$ passe en lui-même lorsque $x_i$ est adjoint à $\mathsf{P}(u)$; de
\[
 h(x_i)f(x)\equiv0\pmod{\frakp},\qquad h(x_i)\ne0
\]
on en déduit nécessairement $f(x)\equiv0\pmod{\frakp}$; car si $\frakp$ contenait un polynôme dépendant seulement de $x_i$, il en contiendrait aussi un dépendant seulement de $x_n$, et serait donc, contre l'hypothèse, tout au plus de dimension zéro. Dans la correspondance $v\sim x_i$, la classe nulle modulo $\frakc$ correspond donc nécessairement à la classe nulle modulo $\frakp$, et bien entendu la réciproque vaut aussi. Il en résulte une correspondance isomorphe entre l'anneau quotient par $\frakp$ et celui par $\frakc$; ce dernier ne possède donc pas de diviseurs de zéro, et $\frakc$ devient un idéal premier. Par cette correspondance isomorphe, la dépendance algébrique entre $(x_i)$ et $(\mathsf{P}(u,x_{i+1},\ldots,x_n))$ correspond à une relation entre $(x_{i+1}),\ldots,(x_n)$ relativement à $(\mathsf{P}(u,v))$, tandis que, pour $\lambda=1,\ldots,i-1$, les dépendances entre $(x_\lambda)$ et $(\mathsf{P}(u,x_{i+1},\ldots,x_n))$ se conservent comme telles et n'abaissent donc pas davantage le rang algébrique; le rang algébrique a diminué exactement d'une unité. Si $\frakp$ était de dimension zéro, $\frakc$ devient l'idéal unité; le résultat ci-dessus reste donc valable. Dans le domaine primitif $\bar{\mathfrak{J}}$, on obtient ainsi
\[
 \bar{\fraka}=(\bar{\frakp},v+v_1y_1+\cdots+v_ny_n),
\]
où $v_\lambda$ est posé égal à $-t_{\lambda i}$, un idéal de la nature cherchée. Si, en particulier, $\mathsf{P}$ contient une infinité d'éléments, on peut toujours spécialiser $v,v_1,\ldots,v_n$ en des éléments de $\mathsf{P}$ de telle sorte que la norme et la forme des diviseurs élémentaires, et donc la dimension de $\fraka$, deviennent égales à celles de l'idéal spécialisé $\frakb$;\footnote{Voir H.-N., dernier paragraphe du \S{} 7.} toutefois la forme des diviseurs élémentaires ne reste pas nécessairement irréductible. Mais, d'après le théorème V, $\frakb$ possède au moins un idéal premier associé de la dimension cherchée; en revenant à $\bar{\mathfrak{J}}$, le théorème VI y vaut donc sans adjonction d'indéterminées quelconques.

Les théorèmes II et VI donnent immédiatement la formulation désirée selon le

\textbf{Théorème VII.} \emph{Un idéal premier $\frakp$ est de dimension $n-i$ s'il existe -- après adjonction éventuelle d'indéterminées -- au moins une chaîne de $1+(n-i)+1$ idéaux premiers}
\[
 \frakp_0=\frakp;\quad \frakp_1,\ldots,\frakp_{n-i};\quad \frakp_{n-i+1},
\]
\emph{dont chacun est un diviseur propre de celui qui le précède immédiatement, mais s'il n'existe aucune chaîne de ce genre ayant un nombre plus grand de membres. Cette définition vaut aussi dans le domaine primitif, et cela sans introduction d'indéterminées quelconques, lorsque $\mathsf{P}$ contient une infinité d'éléments.}\footnote{Cette définition de la dimension par chaînes donne, dans le cas du corps de nombres algébriques, la dimension zéro pour les idéaux premiers ordinaires, la dimension $-1$ pour l'idéal unité et la dimension $1$ pour l'idéal nul; en effet, ce dernier satisfait aussi à la définition des idéaux premiers -- dans un corps, un produit de grandeurs non nulles est toujours non nul. Le théorème II reste valable, avec cette définition de la dimension, dans les corps de nombres algébriques.}

Il est tout d'abord clair que le dernier élément de la chaîne doit devenir égal à l'idéal unité -- lequel satisfait à la définition de l'idéal premier --, puisque autrement la chaîne pourrait être prolongée par adjonction de $\frako$; pour $\frako$ lui-même, le théorème VII donne la dimension $-1$, en accord avec \S{} 1, 8. Soit maintenant $\frakp$ différent de $\frako$ et de dimension $n-i$; alors le nombre de membres de la chaîne peut être tout au plus $1+(n-i)+1$, puisque, d'après le théorème II, deux idéaux premiers de même dimension ne peuvent pas apparaître. Mais il existe, d'après le théorème VI -- après adjonction éventuelle de nouvelles indéterminées --, au moins une chaîne de ce nombre de membres, la dimension ayant diminué exactement d'une unité à chaque membre de la chaîne. Si l'on prend le théorème VII pour définition de la dimension -- définition qui peut manifestement s'énoncer exactement de la même manière dans le domaine primitif et pour laquelle, d'après le théorème VI, l'introduction d'indéterminées devient inutile lorsque $\mathsf{P}$ contient une infinité d'éléments --, alors la définition de la dimension maximale d'un idéal quelconque est de nouveau donnée par la seconde partie du théorème V.
\providecommand{\fraka}{\mathfrak a}
\providecommand{\frakg}{\mathfrak g}
\providecommand{\frakh}{\mathfrak h}
\providecommand{\frakm}{\mathfrak m}
\providecommand{\frako}{\mathfrak o}
\providecommand{\frakp}{\mathfrak p}
\providecommand{\frakq}{\mathfrak q}
\providecommand{\frakr}{\mathfrak r}
\providecommand{\frakt}{\mathfrak t}

\subsection*{\S{} 5. Insertion des idéaux fondamentaux et de leurs composantes dans la théorie générale des idéaux. Zéros des idéaux.}

Il s'agit de caractériser les idéaux fondamentaux et leurs composantes (\S{} 1, 7.) par les composantes isolées de la décomposition (\S{} 1, 13.) au sens de la théorie générale des idéaux. Cette caractérisation montrera le parallélisme des théorèmes de décomposition des normes avec ceux de la théorie générale des idéaux, et permettra d'énoncer pour des idéaux arbitraires la théorie des zéros donnée au \S{} 3 pour les idéaux premiers et les idéaux primaires. Pour les idéaux fondamentaux vaut le

\textbf{Théorème VIII.} \emph{Les idéaux fondamentaux $\frakg_i$ de $i$-ième étage sont les composantes isolées de la décomposition qui sont déterminées de manière unique par la totalité des idéaux premiers associés de dimensions $n-1,n-2,\ldots,n-i$. Si tous les idéaux premiers associés sont de même dimension, l'idéal devient lui aussi de cette dimension; c'est-à-dire qu'un seul facteur $R^{(i)}$ de la norme apparaît. En général, la norme possède autant de facteurs $R^{(i)},R^{(i+\sigma)},R^{(i+\tau)},\ldots$ différents de l'unité qu'il apparaît de nombres de dimension différents parmi les idéaux premiers associés.}

Faisons précéder la preuve du

\textbf{Lemme VI.} \emph{Si un idéal est représenté comme plus petit commun multiple, $\frakm=[\fraka_1,\ldots,\fraka_\alpha]$, alors son idéal fondamental $\frakg_i$ est le plus petit commun multiple des idéaux fondamentaux $\frakh_i$ de $i$-ième étage des $\fraka$.}

En effet, de
\[
 b^{(i+1)}f\equiv0(\frakm),\qquad b^{(i+1)}\ne0
\]
on déduit
\[
 b^{(i+1)}f\equiv0(\fraka_\lambda),\qquad b^{(i+1)}\ne0;
\]
donc $\frakg_i$ devient divisible par $[\frakh_{i1},\ldots,\frakh_{i\alpha}]$. Réciproquement, de
\[
 c_\lambda^{(i+1)}f\equiv0(\fraka_\lambda),\qquad c^{(i+1)}\ne0
\]
on déduit aussi
\[
 c_1^{(i+1)}\cdots c_\alpha^{(i+1)}f\equiv0(\frakm),\qquad
 c_1^{(i+1)}\cdots c_\alpha^{(i+1)}\ne0,
\]
et donc la divisibilité réciproque, ce qui démontre le lemme.

Soient maintenant les idéaux premiers associés de $\frakm$ ordonnés suivant la dimension (certaines dimensions pouvant naturellement aussi manquer, $j_\lambda=0$):
\[
 \frakp_{1,1},\ldots,\frakp_{1,j_1};\quad
 \frakp_{2,1},\ldots,\frakp_{2,j_2};\quad\ldots;\quad
 \frakp_{n,1},\ldots,\frakp_{n,j_n},
\]
où, en général, $\frakp_{\lambda,\kappa}$ désigne un idéal premier de dimension $n-\lambda$; soit $\frakq_{\lambda,\kappa}$ l'idéal primaire correspondant qui intervient dans une décomposition fixe donnée de $\frakm$. Alors, par définition, l'idéal fondamental de $i$-ième étage devient égal à l'idéal unité pour tous les $\frakq_{\lambda,\kappa}$ pour lesquels $\lambda>i$; mais, pour $\lambda\leqq i$, cet idéal fondamental devient, d'après le théorème I, égal à $\frakq_{\lambda,\kappa}$. Ainsi, d'après le lemme VI, on obtient:
{\let\veqno\leqno\begin{equation}
 \frakg_i=[\frakq_{1,1},\ldots,\frakq_{1,j_1};\ldots;\frakq_{i,1},\ldots,\frakq_{i,j_i}],
 \tag{5}
\end{equation}}
et cette représentation fait reconnaître $\frakg_i$ comme composante isolée de la décomposition, puisque aucun des idéaux premiers appartenant à $\frakg_i$ ne peut être contenu dans l'un des autres idéaux premiers, tous ceux-ci étant de dimension plus basse.

Si, spécialement, il n'apparaît parmi les idéaux premiers associés que des idéaux de dimension $n-i$, alors $\frakg_0,\ldots,\frakg_{i-1}$ deviennent, d'après (5), égaux à l'idéal unité, et $\frakg_i=\frakg_{i+1}=\cdots=\frakm$; donc $R_\frakm=R^{(i)}$. Mais s'il existe des idéaux premiers de dimensions différentes, par exemple si $j_i,j_{i+\sigma},j_{i+\tau},\ldots$ sont non nuls, alors, d'après (5),
\[
 \frakg_0\ne\frakg_i\ne\frakg_{i+\sigma}\ne\frakg_{i+\tau}\ne\cdots,
\]
mais
\[
 \frakg_0=\frakg_1=\cdots=\frakg_{i-1}=\frako,
 \qquad
 \frakg_i=\frakg_{i+1}=\cdots=\frakg_{i+\sigma-1},\ldots,
\]
et donc les facteurs $R^{(i)},R^{(i+\sigma)},R^{(i+\tau)}$, et eux seuls, deviennent différents de l'unité.\footnote{Du théorème VIII, on déduit de nouveau, en tenant compte du lemme II, que les idéaux fondamentaux -- comme plus petits communs multiples d'idéaux transformés -- deviennent des idéaux transformés. Comme le théorème VIII repose seulement sur des théorèmes de H.-N. où ce fait n'est pas utilisé, on obtient ainsi une nouvelle preuve, qui en découvre en même temps la raison interne. Chez H.-N., le théorème n'est utilisé que dans la théorie des zéros, ici développée à nouveau.}

Au moyen du théorème VIII, le théorème I se précise en

\textbf{Théorème IX.} \emph{Un idéal est primaire si et seulement si sa forme des diviseurs élémentaires -- et par conséquent sa norme -- est un polynôme primaire.}

Que la condition soit remplie pour les idéaux primaires a été montré dans le théorème I. Supposons maintenant réciproquement que la forme des diviseurs élémentaires de $\frakm$ soit donnée par $Q^{(i)}=(P^{(i)})^\rho$, où $P^{(i)}$ désigne un polynôme irréductible. Alors, d'après le théorème VIII, tous les idéaux premiers de $\frakm$ sont de la même dimension $n-i$. De $\frakm=[\frakq_1,\ldots,\frakq_\alpha]$ il résulte:
\[
 Q^{(i)}\equiv0(\frakq_\lambda);\quad\hbox{donc}\quad
 (P^{(i)})^\rho\equiv0(\frakp_\lambda)\quad\hbox{et par suite}\quad
 P^{(i)}\equiv0(\frakp_\lambda).
\]
Comme tous les $\frakp_\lambda$ sont de dimension $n-i$, et que $P^{(i)}$ désigne un polynôme irréductible, $P^{(i)}$ devient la forme des diviseurs élémentaires de chaque $\frakp_\lambda$. Ainsi, d'après le corollaire du théorème IV, tous les $\frakp_\lambda$ coïncident; puisque l'on a affaire à des plus grandes composantes primaires, il ne peut donc apparaître qu'un seul $\frakp$; $\frakm$ devient primaire, le cas spécial premier n'étant pas exclu pour $\rho=1$. Les théorèmes VIII et IX conduisent à l'insertion des composantes des idéaux fondamentaux selon le

\textbf{Théorème X.} \emph{La composante $\frakr_\lambda$ de l'idéal fondamental $\frakg_i$, qui correspond au plus grand facteur primaire $Q_\lambda^{(i)}=(P_\lambda^{(i)})^{\rho_\lambda}$ de $R^{(i)}$, est donnée par la composante isolée de la décomposition qui est déterminée de manière unique par les idéaux premiers associés des dimensions $n-1,\ldots,n-i+1$ et par l'idéal premier de dimension $n-i$ dont la forme des diviseurs élémentaires devient égale à $P_\lambda^{(i)}$. Les idéaux premiers associés et les plus grands facteurs primaires de la norme se correspondent univoquement.}

D'après \S{} 1, 7., les $\frakr_\lambda$ coïncident avec leur idéal fondamental de $i$-ième étage et possèdent $\frakg_{i-1}$ comme idéal fondamental de $(i-1)$-ième étage; ils admettent donc, d'après le théorème VIII, respectivement d'après (5), une représentation la plus courte:
\[
 \frakr_\lambda=[\frakg_{i-1},\frakq_{\lambda,1},\ldots,\frakq_{\lambda,t_\lambda}],
\]
où tous les $\frakq$ sont de dimension $n-i$ et appartiennent à des idéaux premiers différents. D'après la définition de $Q_\lambda^{(i)}$ -- en tenant compte du fait que $\frakr_\lambda$ coïncide avec son idéal fondamental de $i$-ième étage --, on obtient:
\[
 Q_\lambda^{(i)}\frakg_{i-1}\equiv0\pmod{\frakq_{\lambda,\kappa}},
 \quad(\kappa=1,\ldots,t_\lambda)
\]
et par conséquent
\[
 (Q_\lambda^{(i)})^{\sigma_\lambda}\equiv0\pmod{\frakq_{\lambda,\kappa}};
 \quad\hbox{donc}\quad
 (P_\lambda^{(i)})^{\rho_\lambda\sigma_\lambda}\equiv0(\frakp_{\lambda\kappa})
\]
et par suite $P_\lambda^{(i)}\equiv0\pmod{\frakp_{\lambda,\kappa}}$. Il en résulte, comme dans la preuve du théorème IX, que, dans la représentation de $\frakr_\lambda$, il ne peut intervenir qu'un seul $\frakq_\lambda$ et un seul $\frakp_\lambda$ associé de dimension $n-i$, et que la forme des diviseurs élémentaires de $\frakp_\lambda$ est donnée par $P_\lambda^{(i)}$. Il reste à démontrer que $\frakp_\lambda$ est un idéal premier associé de $\frakm$ et à montrer que $\frakq_\lambda$ intervient effectivement dans au moins une représentation de $\frakm$ par plus grandes composantes primaires; alors $\frakr_\lambda$ est par là même reconnu comme composante isolée de la décomposition, puisqu'aucun idéal premier ne peut être absorbé dans un idéal de dimension égale ou plus basse, et précisément comme composante déterminée par $\frakp_\lambda$ et par tous les idéaux premiers associés de dimension plus élevée. Pour cette démonstration, il suffit de montrer que $\frakq_\lambda$ intervient dans une représentation la plus courte de $\frakg_i$, puisqu'une telle représentation peut toujours, d'après le théorème VIII, être complétée en une représentation la plus courte de $\frakm$; il suffit donc de montrer que la représentation
\[
 \frakg_i=[\frakr_1,\ldots,\frakr_\alpha]=[\frakg_{i-1},\frakq_1,\ldots,\frakq_\alpha]
\]
devient la plus courte. Si, par exemple, $\frakq_\lambda$ était absorbé dans son complément, c'est-à-dire si $\frakg_{i-1}\frakq_1\cdots\frakq_{\lambda-1}\frakq_{\lambda+1}\cdots\frakq_\alpha$ était divisible par $\frakq_\lambda$, il s'ensuivrait, puisque $\frakg_{i-1}\not\equiv0(\frakq_\lambda)$, qu'une puissance de $\frakq_\mu$ serait divisible par $\frakq_\lambda$ pour $\mu\ne\lambda$; donc, comme les idéaux premiers associés $\frakp_\mu$ et $\frakp_\lambda$ sont tous deux de même dimension, ces idéaux premiers coïncideraient d'après le théorème II; mais cela est impossible, puisque leurs formes des diviseurs élémentaires $P_\mu^{(i)}$ et $P_\lambda^{(i)}$ sont différentes par définition. La représentation de $\frakg_i$, à laquelle correspondent d'après le théorème VIII tous les idéaux premiers associés de dimension $n-i$, montre en outre qu'à chacun de ces idéaux premiers correspond effectivement aussi un plus grand facteur primaire de la norme; la correspondance univoque est ainsi démontrée.

Le théorème X donne immédiatement le transfert du théorème IV selon le

\textbf{Théorème XI.} \emph{Les différents plus grands facteurs primaires $Q_\lambda^{(i)}=(P_\lambda^{(i)})^{\rho_\lambda}$ de la norme admettent une représentation en produit:}
\[
 Q_\lambda^{(i)}=\prod (x_i-t_{1i}\bar y_{1\nu}-\cdots-t_{ni}\bar y_{n\nu})^{\delta_\lambda\rho_\lambda}
 =\prod (t_{1i}(y_1-\bar y_{1\nu})+\cdots+t_{ni}(y_n-\bar y_{n\nu}))^{\delta_\lambda\rho_\lambda},
\]
\emph{où $\bar y_{1\nu},\ldots,\bar y_{n\nu}$ représente chaque fois un système correspondant de zéros de l'idéal premier associé $\frakp_\lambda$ dans les indéterminées primitives $y$, indépendant de $t_{1i},\ldots,t_{ni}$, et où $\delta_\lambda$ prend la valeur un ou $p^{f_\lambda}$; ainsi la décomposition de la norme en facteurs linéaires est entièrement donnée. Le degré de $Q_\lambda^{(i)}$ donne le degré du sous-anneau de l'anneau quotient par $\frakr_\lambda$ formé des classes modulo $\frakr_\lambda$ qui sont représentées par des éléments de $\frakg_{i-1}$, relativement au sous-corps $(\mathsf{P}(u,x_{i+1},\ldots,x_n))$ obtenu par formation de quotients; pour les idéaux primaires, il s'agit donc de l'anneau quotient tout entier. Lors de l'adjonction de nouvelles indéterminées $v$, on obtient encore la décomposition:}
\[
 Q_\lambda^{(i)}(u,v)=\prod (z-v_1\bar x_{1\nu}-\cdots-v_i\bar x_{i\nu})^{\delta_\lambda\rho_\lambda}.
\]

Comme $P_\lambda^{(i)}$ a été reconnu, d'après le théorème X, comme forme des diviseurs élémentaires de $\frakp_\lambda$, il ne s'agit en réalité que d'insérer la décomposition du théorème IV; la remarque sur le degré n'est qu'une autre formulation de ce qui a été dit au \S{} 1, 5. et 7.\footnote{De même, la formulation de la multiplicité mentionnée dans la note 2) de H.-N. est une conséquence immédiate du théorème X. Car le quotient $\frakg_{i-1}\mid\frakr_\lambda$ est -- après adjonction de $x_{i+1},\ldots,x_n$ -- isomorphe à $\frakt_\lambda\mid\frakr_\lambda$, où l'on a posé $\frakt_\lambda=[\frakg_{i-1},\frakr_1,\ldots,\frakr_{\lambda-1},\frakr_{\lambda+1},\ldots,\frakr_\alpha]$, comme le montre immédiatement le retour aux modules de formes linéaires (H.-N., théorème V), compte tenu de la primalité relative des $Q_\mu^{(i)}$; $\frakt_\lambda\mid\frakr_\lambda$ devient isomorphe à $\frakt_\lambda\mid\frakq_\lambda$. Ici $\frakt_\lambda$ naît du complément de $\frakq_\lambda$ lors de l'adjonction de $x_{i+1},\ldots,x_n$.}
\providecommand{\fraka}{\mathfrak a}
\providecommand{\frakm}{\mathfrak m}
\providecommand{\frakp}{\mathfrak p}
\providecommand{\frakq}{\mathfrak q}

\subsection*{\S{} 6. Caractérisation des idéaux premiers et des idéaux primaires propres par la forme des diviseurs élémentaires et la norme.}

Dans le théorème IX, les idéaux primaires étaient caractérisés par la forme des diviseurs élémentaires et par la norme, mais seulement de telle sorte que le cas premier devait être considéré comme un cas spécial du primaire. Il s'agit maintenant de séparer les idéaux premiers des idéaux primaires propres; les raisonnements seront parallèles à ceux du \S{} 3. Des conditions qui sont soit nécessaires, soit suffisantes, s'énoncent facilement par le

\textbf{Théorème XII.} \emph{Une condition nécessaire pour un idéal premier est que la forme des diviseurs élémentaires soit un polynôme irréductible; une condition suffisante est que la norme soit un polynôme irréductible. Autrement dit, la forme des diviseurs élémentaires d'un idéal premier est toujours un polynôme irréductible, et la norme d'un idéal primaire propre est toujours un polynôme primaire propre.}

Que la forme des diviseurs élémentaires d'un idéal premier devienne un polynôme irréductible avait déjà été montré dans le théorème I. Supposons maintenant que la norme $R_\frakq$ soit égale à un polynôme irréductible; il faut montrer que, sous cette hypothèse, tout diviseur propre $\fraka$ de $\frakq$ devient de dimension maximale plus basse, ce par quoi $\frakq$ est reconnu comme idéal premier d'après le théorème II. En effet, d'après \S{} 1, 6., le premier facteur de la norme $R_\fraka$ qui diffère du facteur correspondant de $R_\frakq$ devient un diviseur propre; à cause de $R_\frakq=P^{(i)}$, le facteur $R^{(i)}$ de $R_\fraka$ doit donc être égal à un diviseur propre de $P^{(i)}$, donc à l'unité; $\fraka$ devient de dimension maximale plus basse.

Une seconde preuve simple repose sur le lemme suivant, qui est aussi à la base du théorème suivant.

\textbf{Lemme VII.} \emph{Le système des classes des polynômes $H^{(i)}$ modulo la forme des diviseurs élémentaires $E^{(i)}$ d'un idéal primaire -- plus généralement d'un idéal qui ne possède que des idéaux premiers associés de dimension $n-i$ -- est isomorphe à un sous-système de l'anneau quotient par cet idéal. Le degré de la forme des diviseurs élémentaires donne le degré de ce sous-système relativement au corps des fractions $(\mathsf{P}(u,x_{i+1},\ldots,x_n))$.}

En effet, on peut établir une correspondance en faisant se correspondre les classes représentées par les mêmes polynômes $H^{(i)}$; somme et produit correspondent alors à somme et produit. La correspondance est aussi biunivoque; car tous les polynômes $H^{(i)}$ appartenant à la classe nulle modulo l'idéal sont, par définition, divisibles par $E^{(i)}$; à la classe nulle correspond donc la classe nulle modulo $E^{(i)}$, et naturellement la réciproque vaut aussi.

Supposons maintenant que la norme d'un idéal soit un polynôme irréductible $P^{(i)}$, donc identique à la forme des diviseurs élémentaires. Par conséquent, les degrés de l'anneau quotient par la forme des diviseurs élémentaires $P^{(i)}$ et de l'anneau quotient par $\frakq$ coïncident relativement à $(\mathsf{P}(u,x_{i+1},\ldots,x_n))$; le sous-système isomorphe au premier quotient épuise donc les classes modulo $\frakq$; et, comme ce système ne peut pas posséder de diviseurs de zéro, $\frakq$ devient un idéal premier.

En général, le théorème XII ne peut pas être renforcé jusqu'à donner des conditions nécessaires et suffisantes, comme le montreront plus bas des exemples. En revanche, cela est possible lorsque le domaine des coefficients $\mathsf{P}$ est un corps parfait (\S{} 1, 14.), d'après le

\textbf{Théorème XIII.} \emph{Si le corps des coefficients $\mathsf{P}$ est parfait, alors la norme et la forme des diviseurs élémentaires d'un idéal premier deviennent des polynômes irréductibles et, par conséquent, identiques; la norme et la forme des diviseurs élémentaires d'un idéal primaire propre deviennent des polynômes primaires propres. La condition que la forme des diviseurs élémentaires soit un polynôme irréductible est donc, sur un corps parfait, nécessaire et suffisante pour un idéal premier; dans la condition, la norme peut aussi remplacer la forme des diviseurs élémentaires.}

Compte tenu du théorème XII, le théorème XIII sera démontré dès que l'on aura montré que, sur un corps parfait, l'idéal devient premier aussitôt que la forme des diviseurs élémentaires est un polynôme irréductible, et qu'alors la norme devient égale à ce polynôme. Or il a été montré dans le lemme V, \S{} 3, que, sur un corps parfait, la forme des diviseurs élémentaires devient un polynôme irréductible de première espèce dès qu'elle est irréductible. Si donc $(\mathcal R)$ désigne le corps résiduel obtenu par adjonction de $(x_i)$ à $(\mathsf{P}(u,x_{i+1},\ldots,x_n))$, c'est-à-dire $(\mathsf{P}(u,x_i,x_{i+1},\ldots,x_n))$ modulo $P^{(i)}$, alors $(x_i)$ devient zéro du polynôme irréductible de première espèce $(P^{(i)})$, et par conséquent, comme le montre par exemple la formule de Lagrange, $(y_1),\ldots,(y_n)$ sont contenus dans $(\mathcal R)$. Le sous-système de l'anneau quotient par $\frakq$ qui est isomorphe à $(\mathcal R)$ d'après le lemme VII -- où il ne s'agit que de l'apparition de dénominateurs issus de $(\mathsf{P}(u,x_{i+1},\ldots,x_n))$ -- contient donc les classes $(y_1),\ldots,(y_n)$ et épuise par conséquent l'anneau quotient par $\frakq$; mais, comme il est isomorphe à $(\mathcal R)$, il ne peut pas posséder de diviseurs de zéro, et $\frakq$ devient un idéal premier. Le degré de ce corps $(\mathcal R)$ relativement à $(\mathsf{P}(u,x_{i+1},\ldots,x_n))$ devient égal au degré de la norme de $\frakq$; d'autre part, en raison de l'isomorphisme avec $(\mathcal R)$, il est égal au degré de la forme des diviseurs élémentaires $P^{(i)}$; ainsi la norme et la forme des diviseurs élémentaires doivent coïncider.

Lorsque le corps est imparfait, le théorème XII permet encore le renforcement suivant:

\textbf{Théorème XIV.} \emph{Si le domaine des coefficients $\mathsf{P}$ est un corps imparfait de caractéristique $p$, alors la norme d'un idéal premier devient la $p^g$-ième puissance de la forme des diviseurs élémentaires; on a $0\leqq g\leqq(i-1)f$, où $f$ désigne l'exposant de la forme des diviseurs élémentaires et $n-i$ la dimension de l'idéal premier.}

Le théorème XIV revient à affirmer que le corps résiduel $(\mathcal R)$ de $\frakp$ devient de degré $p^g$ ($g\geqq0$) relativement au sous-corps $(\mathcal R')=(\mathsf{P}(u,x_i,x_{i+1},\ldots,x_n))$ isomorphe à $(\mathcal R)$. Cela résultera immédiatement des théorèmes de Steinitz indiqués au \S{} 1, 15., plus précisément du

\textbf{Lemme VIII.} \emph{Soit $\mathfrak L$ une extension finie du corps imparfait $\Omega$, de degré réduit $r$ et d'exposant $f$; soit $\mathfrak L_0$ le corps de première espèce contenu dans $\mathfrak L$, et soit $\mathfrak M$ un corps intermédiaire entre $\mathfrak L_0$ et $\mathfrak L$. Alors le degré de tout élément de $\mathfrak L$ relativement à $\mathfrak M$ devient une puissance de $p$, où $p$ désigne la caractéristique de $\Omega$.}

En effet, la $p^{f'}$-ième puissance ($f'\leqq f$) de chaque élément $z$ de $\mathfrak L$ appartient à $\mathfrak L_0$; ainsi $z$ devient zéro d'un polynôme irréductible $G(t)=(t-z)^{p^{f'}}$ à coefficients dans $\mathfrak L_0$. Soit $H(t)=(t-z)^\delta$ le polynôme irréductible dans $\mathfrak M$ ayant $z$ pour zéro; alors $H(t)$ possède aussi des coefficients dans $\mathfrak L_0(z)$, et donc les coefficients de $H(t)$ appartiennent au corps intersection de $\mathfrak M$ et de $\mathfrak L_0(z)$, donc à un corps intermédiaire entre $\mathfrak L_0$ et $\mathfrak L_0(z)$. Le degré de $z$ relativement à $\mathfrak M$ devient donc égal au degré relativement à ce corps intermédiaire, et ainsi, comme diviseur de $p^{f'}$, il est une puissance de $p$.

Pour démontrer le théorème XIV, il suffit maintenant de remarquer que le corps $(\mathfrak{R}')=(\mathsf{P}(u,x_i,x_{i+1},\ldots,x_n))$ doit être de degré $r\cdot p^f$ relativement à $(\mathsf{P}(u,x_{i+1},\ldots,x_n))$, si $r$ désigne le degré réduit et $f$ l'exposant de $(\mathfrak{K})$; et que, par conséquent, le corps de première espèce $(\mathfrak{K}_0)$ contenu dans $(\mathfrak{K})$ doit être un sous-corps de $(\mathfrak{R}')$. En effet, puisqu'il existe dans $(\mathfrak{K})$ des éléments de degré $r\cdot p^f$ (\S{} 1, 15.), il faut nécessairement que $(x_i)$ soit de ce degré, puisque tous les éléments de $(\mathfrak{K})$ naissent par spécialisation des $t_{\mu\nu}$ dans $(x_i)$; l'exposant de $(\mathfrak{K})$ coïncide donc avec l'exposant de la forme des diviseurs élémentaires, et $(x_i)^{p^f}$ devient de degré $r$ et élément de $(\mathfrak{K}_0)$, donc élément primitif de $(\mathfrak{K}_0)$. Mais comme $(\mathfrak{K})$ naît de $(\mathfrak{R}')$ par adjonction d'un nombre fini d'éléments -- par exemple $(x_1),\ldots,(x_{i-1})$ --, l'application finiment répétée du lemme VIII donne la démonstration cherchée et, en même temps, l'estimation de $g$.\footnote{Voir à ce sujet les raisonnements parallèles chez A. Ostrowski, \emph{Zur arithmetischen Theorie der algebraischen Größen}, Gött. Nachr. 1919, p. 279--298; en particulier p. 288, où le théorème correspondant est énoncé sans démonstration.}

Ajoutons maintenant les exemples qui montrent que le théorème XII ne peut pas, en général, être renforcé au-delà du théorème XIV; il s'agit d'un idéal premier dont la norme devient égale au carré de la forme des diviseurs élémentaires ($g>0$), et d'idéaux primaires propres dont les formes des diviseurs élémentaires deviennent des polynômes irréductibles; ce dernier exemple montre en même temps qu'il n'y a pas ici d'analogue du théorème XIV, mais que la norme peut devenir une puissance quelconque de la forme des diviseurs élémentaires.

Que le corps des coefficients $\mathsf{P}$ naisse par adjonction de deux indéterminées $\lambda,\mu$ au corps fini $\mathbb F_2$, et soit donc un corps imparfait de caractéristique deux. Il s'agit de l'idéal
\[
 \bar{\frakp}=(y_1^2+\lambda,\,y_2^2+\mu),
\]
qui doit être un idéal premier, puisque l'anneau quotient $(\overline{\mathcal R})$ devient isomorphe au corps $\mathsf{P}(\sqrt\lambda,\sqrt\mu)$. $(\overline{\mathcal R})$ devient de quatrième degré relativement à $(\bar{\mathsf{P}})$, de degré réduit $r=1$, d'exposant $f=2$; la forme des diviseurs élémentaires $P^{(2)}$ devient donc de second degré, la norme de quatrième degré et ainsi le carré de $P^{(2)}$, comme le laisse aussi facilement recalculer la représentation modulaire (3), \S{} 1, 5.; il vient
\[
 P^{(2)}=(u_{11}u_{22}+u_{12}u_{21})^2x_2^2+(u_{11}^2\mu+u_{21}^2\lambda)
\]
ou encore, après une multiplication convenable:
\[
 x_2^2+(t_{12}^2\lambda+t_{22}^2\mu).
\]

Que maintenant $\mathsf{P}$ naisse par adjonction de l'unique indéterminée $\lambda$ au corps fini $\mathbb F_2$, et que l'on prenne pour base l'idéal
\[
 \bar{\frakq}=(y_1^2+\lambda,\,y_2^2+\lambda)=(y_1^2+\lambda,(y_1+y_2)^2).
\]
Comme forme des diviseurs élémentaires, la spécialisation $\mu=\lambda$ dans celle qui vient d'être indiquée donne:
\[
 P^{(2)}=x_2^2+\lambda(t_{12}^2+t_{22}^2),
\]
donc un polynôme irréductible, puisque $t_{12}=(0)$, $t_{22}=(1)$ conduit au polynôme irréductible $x_2^2+\lambda$. Mais $\bar{\frakq}$ devient proprement primaire, puisqu'il contient $(y_1+y_2)^2$, mais non $(y_1+y_2)$. La norme doit devenir égale au carré, donc à la $p$-ième puissance, de $P^{(2)}$, puisque l'anneau quotient par $\bar{\frakq}$ est de dimension quatre sur $\mathsf{P}$.

Que ce dernier analogue du théorème XIV ne soit pas toujours réalisé, l'exemple suivant le montre: que $P$ naisse par adjonction de l'indéterminée $\lambda$ au corps fini $\mathbb F_3$; on prend pour base
\[
 \bar{\frakq}=(y_1^3+\lambda,(y_1-y_2)^2),
\]
de sorte qu'il s'agit, comme ci-dessus, d'un idéal primaire propre. La forme des diviseurs élémentaires devient -- en tenant compte du fait que $y_2^3+\lambda$ est aussi divisible par $\bar{\frakq}$ -- égale à
\[
 P^{(2)}=x_2^3+\lambda(t_{12}+t_{22})^3,
\]
donc à un polynôme irréductible; mais la norme devient égale au carré de $P^{(2)}$, puisque l'anneau quotient par $\bar{\frakq}$ est de dimension six. Ici n'apparaît donc pas la $p^g$-ième puissance de la forme des diviseurs élémentaires.

Le premier exemple montre que, lorsque le domaine des coefficients est un corps imparfait, peuvent apparaître, tout comme dans les corps de nombres algébriques, des idéaux premiers de degré supérieur au premier, où $p^g$ doit être désigné comme le degré de l'idéal premier. Les autres exemples sont à comprendre comme des analogues des idéaux de ramification, c'est-à-dire du fait qu'un nombre premier devient divisible par une puissance plus élevée d'un idéal premier, par un idéal primaire; car, comme il est montré dans la note 10), la forme des diviseurs élémentaires correspond au plus petit nombre rationnel entier divisible par un idéal dans un corps de nombres algébriques.
\providecommand{\frakS}{\mathfrak S}
\providecommand{\frakm}{\mathfrak m}
\providecommand{\frakp}{\mathfrak p}

\subsection*{\S{} 7. Idéaux premiers absolus.}

Un idéal premier à coefficients dans $\mathsf{P}$ est appelé \emph{idéal premier absolu} s'il reste un idéal premier dans le corps algébriquement clos $\mathsf{A}$ auquel $P$ peut être étendu.\footnote{Un corps algébriquement clos est, comme on sait, un corps dans lequel tout polynôme d'une indéterminée se décompose en facteurs linéaires. Existence et unicité essentielle du corps algébriquement clos appartenant à un corps quelconque chez Steinitz.} De manière correspondante, un polynôme irréductible $\mathsf{P}$ à coefficients dans $P$ est appelé \emph{polynôme absolument irréductible} s'il reste irréductible dans $\mathsf{A}$. La caractérisation des idéaux premiers absolus est donnée par le

\textbf{Théorème XV.} \emph{Une condition nécessaire et suffisante pour un idéal premier absolu consiste en ce que sa forme des diviseurs élémentaires -- que l'on peut supposer entière et primitive dans les $u$ -- devienne un polynôme absolument irréductible en les $x$ et les $u$. La forme des diviseurs élémentaires devient identique à la norme, de sorte que la condition peut aussi s'énoncer pour la norme.}

Pour la démonstration, remarquons d'abord qu'en général la norme et la forme des diviseurs élémentaires se conservent lors d'une extension algébrique du domaine des coefficients. En effet, les facteurs individuels $R^{(i)}$ et $E^{(i)}$ sont définis comme plus grands diviseurs communs au sens polynomial, propriété qui se conserve lors d'une extension algébrique du domaine des coefficients. Ainsi $P^{(i)}$ reste la forme des diviseurs élémentaires de $\frakp$ lors du passage de $\mathsf{P}$ à $\mathsf{A}$, donc lors du passage à un corps parfait comme domaine des coefficients; car tout corps algébriquement clos est parfait, comme le montre immédiatement la définition. D'après le théorème XIII, une condition nécessaire et suffisante pour que $\frakp$ soit un idéal premier absolu est donc que $P^{i)}$ devienne irréductible relativement à $\mathsf{A}(u)$; donc absolument irréductible comme polynôme des $x$ et des $u$, si l'on suppose $P^{(i)}$ entière et primitive dans les $u$. En même temps, d'après le théorème XIII, $P^{(i)}$ devient aussi identique à la norme de $\frakp$ dès que $\mathsf{A}$ est pris pour domaine des coefficients; il en va donc de même, d'après la remarque précédente, lorsque l'on prend $\mathsf{P}$ pour domaine de base. Le théorème XV est ainsi démontré.

Pour les polynômes absolument irréductibles $P$ à coefficients dans un corps de nombres algébriques (fini) $\mathfrak K$, on dispose maintenant du théorème\footnote{A. Ostrowski: loc. cit. (note 21)); lemme p. 296. -- Preuve plus simple chez E. Noether, \emph{Ein algebraisches Kriterium für absolute Irreduzibilität}, Math. Ann. 85 (1922), p. 26--33, no 7.} selon lequel $P$ reste irréductible, plus précisément absolument irréductible, modulo tout idéal premier de $\mathfrak K$, à l'exception tout au plus d'un nombre fini d'idéaux premiers de $\mathfrak K$. Le transfert de ce théorème aux idéaux premiers absolus s'appuie sur le

\textbf{Théorème XVI.} \emph{Si l'on prend comme domaine des coefficients, à la place d'un corps, l'anneau $\mathfrak o^*$ de tous les entiers algébriques d'un corps de nombres algébriques (fini) $\mathfrak K$, alors la norme et la forme des diviseurs élémentaires d'un idéal de polynômes se conservent comme norme et forme des diviseurs élémentaires modulo tout idéal premier de $\mathfrak K$, à l'exception tout au plus d'un nombre fini d'idéaux premiers de $\mathfrak K$.}

Pour la démonstration, il faut modifier, par rapport au \S{} 1, 1., l'anneau de polynômes qui est pris pour base. Que $\bar{\mathfrak{J}}$ soit composé de tous les polynômes en $y_1,\ldots,y_n$ à coefficients dans $\mathfrak o^*$, donc en entiers algébriques de $\mathfrak K$; que l'anneau des coefficients de $\mathfrak{J}$ soit $\mathfrak o^*[u]$, c'est-à-dire l'ensemble des polynômes en $u$ à coefficients dans $\mathfrak o^*$. Si $\bar{\frakm}$ est un idéal dans $\bar{\mathfrak{J}}$, il passe, par (1), en un idéal transformé $\frakm$ dans $\mathfrak{J}$. Il faut démontrer que, pour former la norme, il n'apparaît au dénominateur qu'un nombre fini de polynômes de $\mathfrak o^*[u]$; le théorème XVI en résultera directement.

Il faut d'abord remarquer que la représentation modulaire (3) du \S{} 1, 5., pour la formation de la norme individuelle, revient à une réduction des puissances de $x_{i-1}$ modulo un polynôme régulier en $x_{i-1}$,
\[
 C^{(i-1)}=U_{i-1}(u)x_{i-1}^{k}+\hbox{termes inférieurs};
\]
\footnote{Voir à ce sujet H.-N. \S{} 4.} tous les coefficients de $C^{(i-1)}$, et en particulier $U_{i-1}$, peuvent alors être supposés être des grandeurs de $\mathfrak o^*[u]$. Comme il s'agit de réduire successivement $x_1,\ldots,x_{i-1}$, la représentation modulaire entière dans les $x$ devient donc aussi entière dans $\mathfrak o^*[u]$, à l'exception des produits de puissances
\[
 U_1^{\lambda_1}\cdots U_{i-1}^{\lambda_{i-1}}
\]
au dénominateur. Mais le produit $U_1U_2\cdots U_n$, considéré comme polynôme en les $u$, définit par ses coefficients un idéal $\mathfrak r^*$ de $\mathfrak K$, qui n'est donc divisible que par un nombre fini d'idéaux premiers de $\mathfrak K$. Si l'on prend $\frakp^*$ différent de ceux-ci, et si l'on prend pour corps des coefficients le corps résiduel $\mathfrak o^*/\frakp^*$, alors la représentation modulaire primitive se conserve, chaque nombre de $\mathfrak o^*$ étant seulement remplacé par sa classe modulo $\frakp^*$.

En outre, comme la norme et la forme des diviseurs élémentaires ne sont déterminées qu'à des facteurs de $\mathfrak o^*[u]$ près, on peut les supposer divisibles par $\frakm$ aussi relativement à $\mathfrak o^*[u]$. Sous cette hypothèse, soit par exemple
\[
 R^{(i)}=V_i(u)x_i^{l}+\hbox{termes inférieurs},
\]
de sorte que, dans la divisibilité de tous les déterminants à $\rho$ lignes du module $\mathfrak M_{i-1}^*$ de rang $\rho$ déterminé par $C^{(i-1)}$, seules des puissances de $V_i$, en dehors de $U_1^{\lambda_1}\cdots U_{i-1}^{\lambda_{i-1}}$, apparaissent au dénominateur. Si l'on choisit donc $\frakp^*$ encore différent du nombre fini d'idéaux premiers de $\mathfrak K$ qui interviennent dans l'idéal de $\mathfrak K$ déterminé par $V_1V_2\cdots V_n$, alors $R^{(i)}$ reste un diviseur commun de ces déterminants lorsque l'on prend pour corps des coefficients le corps résiduel $\mathfrak o^*/\frakp^*$, et le rang de $\mathfrak M_{i-1}^*$ se conserve, puisque $C^{(i)}$ était un déterminant à $\rho$ lignes de $\mathfrak M_{i-1}^*$.

Enfin, il faut encore choisir $\frakp^*$ de sorte que $R^{(i)}$ reste chaque fois plus grand diviseur commun. Que $D^{(i)}$ parcoure tous les déterminants à $\rho$ lignes de $\mathfrak M_{i-1}^*$; par hypothèse, on a
\[
 U_1^{\lambda_1}\cdots U_{i-1}^{\lambda_{i-1}}V_i^{\sigma}D^{(i)}=R^{(i)}T^{(i)},
\]
où les $T^{(i)}$ ne possèdent aucun diviseur commun contenant un $x$ au sens polynomial. Par conséquent, l'idéal dérivé de tous les $T^{(i)}$ contient un polynôme
\[
 G^{(i+1)}=W_i(u)x_{i+1}^{m}+\hbox{termes inférieurs},\qquad W_i\ne0;
\]
le fait que l'on puisse supposer ici $G^{(i+1)}$ régulier en $x_{i+1}$ résulte de la composition de la transformation (1) avec une transformation de $x_{i+1},\ldots,x_n$ ayant de nouvelles indéterminées pour coefficients. Si l'on choisit donc encore $\frakp^*$ différent du nombre fini d'idéaux premiers qui interviennent dans l'idéal de $\mathfrak K$ déterminé par $W_1W_2\cdots W_n$, alors, en prenant pour corps des coefficients le corps résiduel $\mathfrak o^*/\frakp^*$, $R_\frakm$ se conserve effectivement comme norme. De façon tout à fait correspondante, on peut encore choisir $\frakp^*$ de sorte que $E_\frakm$ se conserve aussi comme forme des diviseurs élémentaires; le théorème XVI est ainsi démontré.

Des théorèmes XV et XVI résulte, comme généralisation du théorème sur les polynômes absolument irréductibles et en utilisant ce théorème, le

\textbf{Théorème XVII.} \emph{Si $\frakp$ est un idéal premier absolu à coefficients entiers dans un corps de nombres algébriques (fini) $\mathfrak K$, alors $\frakp$ reste un idéal premier, plus précisément un idéal premier absolu, modulo tout idéal premier de $\mathfrak K$, à l'exception tout au plus d'un nombre fini d'idéaux premiers de $\mathfrak K$.}

Pour la démonstration, que la norme $R_\frakp$ de $\frakp$ soit choisie comme dans le théorème XVI, de sorte qu'elle devienne aussi divisible par $\frakp$ relativement à $\mathfrak o^*[u]$; on a alors
\[
 R_\frakp=T(u)P^{(i)}(x),
\]
où l'on peut supposer $P^{(i)}(x)$ entier et primitif dans les $u$. D'après le théorème XV, $P^{(i)}$, considéré comme polynôme en $u$ et $x$ à coefficients dans $\mathfrak K$, devient donc absolument irréductible. D'après le théorème sur les polynômes absolument irréductibles, il conserve cette propriété modulo tout idéal premier de $\mathfrak K$, à l'exception tout au plus d'un nombre fini d'idéaux premiers de $\mathfrak K$. Si l'on choisit donc $\frakp^*$ différent de ce nombre fini d'idéaux premiers, puis, d'après le théorème XVI, encore différent d'un nombre fini d'autres idéaux premiers de $\mathfrak K$, alors $P^{(i)}$ -- dont les coefficients de $\mathfrak o^*$ sont remplacés par les classes correspondantes modulo $\frakp^*$ -- devient la norme de $\frakp$ lorsque l'on prend pour corps des coefficients $P=\mathfrak o^*/\frakp^*$, et cette norme devient absolument irréductible. En prenant $P$ pour corps des coefficients, $\frakp$ reste donc, d'après le théorème XV, un idéal premier absolu, ce qui démontre le théorème XVII.

\begin{center}
Göttingen, le 8 mai 1923.\\[1em]
(Reçu le 9. 3. 1923.)
\end{center}
\clearpage
\setcounter{footnote}{0}
\providecommand{\frakS}{\mathfrak S}
\providecommand{\fraka}{\mathfrak a}
\providecommand{\frakp}{\mathfrak p}
\providecommand{\frakq}{\mathfrak q}
\providecommand{\frakR}{\mathfrak R}

\section*{\NoetherNavAnchor{P25}{25. Théorie de l’élimination et théorie des idéaux}25. Théorie de l'élimination et théorie des idéaux}

\begin{center}
J. Ber. d. DMV 33 (1924), p. 116--120
\end{center}

\begin{center}
Conférence prononcée à Marbourg le 25 septembre 1923.\footnote{Exposé détaillé et indications bibliographiques dans Math. Ann. 90 et dans un travail qui paraîtra prochainement.}\\[0.5em]
Par \textsc{Emmy Noether} à Göttingen.
\end{center}

Je voudrais montrer dans ce qui suit comment la théorie de l'élimination peut se construire tout à fait parallèlement à la théorie des zéros pour les polynômes d'une indéterminée; cette construction repose sur une liaison entre la théorie de l'élimination donnée par \emph{Hentzelt} et les méthodes de la théorie des corps, telles que \emph{Steinitz} les a développées, et de la théorie des idéaux, telle que je l'ai moi-même développée.

Je veux d'abord esquisser la question dans le cas des polynômes d'une indéterminée. Comme corps des coefficients, on prend une fois pour toutes un corps fixe $P$, auquel se rapportent les notions d'irréductibilité, etc. Ici $P$ peut être supposé défini abstraitement au sens de Steinitz, mais, dans le cas particulier, il peut naturellement coïncider avec le corps des nombres rationnels ou avec le corps de tous les nombres complexes -- conformément à l'hypothèse autrefois généralement usuelle dans la théorie de l'élimination. La théorie des zéros se décompose alors, pour les polynômes d'une indéterminée, en quatre étapes:

1. \emph{Le théorème de décomposition}, qui permet de ramener la question aux polynômes irréductibles de $P[x]$. Car de
\[
 a(x)=p_1(x)^{\varrho_1}\cdots p_r(x)^{\varrho_r}=q_1(x)\cdots q_r(x)
\]
-- où les $p(x)$ désignent des polynômes irréductibles distincts, et les $q(x)$ donc les plus grands facteurs primaires -- il résulte que $a(x)$ s'annule en tous les zéros, et seulement en eux, des polynômes irréductibles $p(x)$.

2. \emph{La construction du corps des zéros pour un polynôme irréductible $p(x)$.} Un tel corps des zéros $\mathfrak K=P(\alpha)$ -- où $\alpha$ désigne un zéro de $p(x)$ -- peut toujours, d'après \emph{Steinitz}, être construit comme un corps d'extension de $P$ isomorphe au corps résiduel $P[x]/(p(x))$. Réciproquement, tout corps des zéros situé dans un corps d'extension de $P$ devient isomorphe à ce corps résiduel; les deux sont donc abstraitement identiques.

3. \emph{La décomposition de $p(x)$ en facteurs linéaires dans le corps galoisien déterminé par $p(x)$.} Par répétition de la construction donnée sous 2., on parvient au corps galoisien essentiellement uniquement déterminé $P(\alpha_1,\ldots,\alpha_m)$, qui suffit précisément pour que $p(x)$ se décompose en facteurs linéaires.

4. \emph{La signification de la multiplicité pour le polynôme initial $a(x)$.} D'après 1., 2., 3., la question des zéros et la décomposition correspondante en facteurs linéaires sont réglées pour un polynôme arbitraire $a(x)$. Comme multiplicité, on peut considérer le degré du polynôme primaire individuel $q(x)$ -- donc la dimension de $P[x]/(q(x))$ sur $P$ --; car, lorsque le corps des coefficients $P$ est algébriquement clos et que les $p(x)$ deviennent linéaires, ce nombre revient précisément à la multiplicité des zéros.

Je passe maintenant à la théorie générale de l'élimination, et je prends donc pour base l'anneau $\frakS=P[x_1,\ldots,x_n]$; la théorie de l'élimination peut alors se définir comme la \emph{théorie des zéros d'un idéal} de cet anneau de polynômes. Par un idéal $\fraka$, j'entends ici, comme d'ordinaire, un système de polynômes tel que, avec $f$ et $g$, appartienne aussi $f-g$ à $\fraka$, et qu'avec $f$ appartienne aussi $Af$ à $\fraka$, où $A$ désigne un polynôme arbitraire de $\mathfrak{J}$. Par zéro de $\fraka$, j'entends tout système de valeurs $\alpha_1,\ldots,\alpha_n$, appartenant à un corps d'extension de $P$, tel que $f(\alpha)$ s'annule pour tout polynôme $f$ appartenant à l'idéal. Si $f_1,\ldots,f_t$ désigne une base d'idéal toujours existante de $\fraka$ -- c'est-à-dire si $f=A_1f_1+\cdots+A_tf_t$ vaut pour tout $f$ de $\fraka$ --, alors les zéros de $\fraka$ sont évidemment identiques aux zéros de $f_1,\ldots,f_t$; et comme, réciproquement, tout système de polynômes en nombre fini peut être regardé comme base de l'idéal qui en est dérivé, la définition des zéros donnée ci-dessus pour la théorie de l'élimination est identique à celle qui est habituellement en usage, mais évite la difficulté qui réside dans le choix arbitraire d'une base. Pour les polynômes d'une indéterminée, où il ne s'agit que d'idéaux principaux, cette difficulté n'apparaît pas, puisque le générateur de l'idéal principal est déterminé de façon unique à une unité près -- élément de $P$ comme facteur --; toutefois la conception par les idéaux donnera aussi là un aperçu plus précis.

La théorie de l'élimination se décompose alors, conformément au cas d'une indéterminée, en quatre étapes:

I. \emph{Le théorème de décomposition} -- la représentation d'un idéal par ses plus grandes composantes primaires --, qui permet de ramener la question à une question portant sur les idéaux premiers. Un idéal s'appelle ici idéal premier lorsqu'il ne divise un produit que s'il divise au moins l'un des facteurs; il s'appelle idéal primaire lorsqu'il divise au moins l'un des facteurs ou une puissance de chaque facteur. À chaque idéal primaire $\frakq$ appartient un et un seul idéal premier associé $\frakp$, qui est diviseur de $\frakq$ et dont une puissance devient divisible par $\frakq$. On a la représentation comme plus petit commun multiple:
\[
 \mathfrak a=[\mathfrak q_1,\cdots,\mathfrak q_r];\qquad
 \frakp_1^{\sigma_1}\cdots \frakp_r^{\sigma_r}\equiv 0\pmod{\fraka}.
\]
Ici les $\frakq$ sont les plus grandes composantes primaires; c'est-à-dire que chaque $\frakq$ est primaire, mais que le plus petit commun multiple de deux quelconques des $\frakq$ n'est pas primaire. En outre, on peut supposer sans restriction de généralité qu'aucun $\frakq$ ne divise le plus petit commun multiple des autres, de sorte qu'aucun $\frakq$ ne puisse simplement être omis; les $\frakp$ sont les idéaux premiers associés aux $\frakq$. Or ce ne sont pas les $\frakq$, mais bien -- et c'est là l'analogue du cas d'une indéterminée -- leurs idéaux premiers associés $\frakp$ qui sont déterminés de façon unique par $\fraka$; et comme chaque $\frakp$ est un diviseur de $\fraka$, tandis que réciproquement un produit de puissances des $\frakp$ devient divisible par $\fraka$, les zéros de l'idéal coïncident donc avec ceux de ses idéaux premiers associés.

II. \emph{La construction du corps des zéros pour un idéal premier $\frakp$.} L'anneau quotient par un idéal premier est, par définition, intègre et contient au moins un élément non nul dès que $\frakp$ est différent de l'idéal unité; par formation de quotients, cet anneau peut donc être étendu à son corps résiduel $(\mathfrak K)$. On peut toujours construire un corps d'extension $\mathfrak K$ de $P$, isomorphe à $(\mathfrak K)$, qui devient le corps des zéros de $\frakp$; c'est-à-dire que l'on obtient $\mathfrak K=P(\alpha_1,\ldots,\alpha_n)$, où $\alpha_1,\ldots,\alpha_n$ désigne un zéro de $\frakp$. Le corps $\mathfrak K$ est alors de degré de transcendance $k$ $(0\leq k<n)$ relativement à $P$; il contient donc $k$ éléments algébriquement indépendants relativement à $P$, tandis que chaque système de $(k+1)$ éléments devient algébriquement dépendant relativement à $P$. Réciproquement, tout corps des zéros de degré de transcendance $k$, situé dans un corps d'extension de $P$, devient isomorphe à $(\mathfrak K)$; les deux sont donc abstraitement identiques.

III. \emph{Décomposition de la norme de $\frakp$ en facteurs linéaires dans le corps galoisien déterminé par $\frakp$.} Si l'on considère les indéterminées $x_1,\ldots,x_n$ comme engendrées à partir d'indéterminées primitives $y_1,\ldots,y_n$ par une transformation linéaire dont les coefficients sont des indéterminées, alors le corps des zéros $\mathfrak K$ de degré de transcendance $k$ peut être regardé comme un corps d'extension algébrique fini de $P(x_{i+1},\ldots,x_n)$ $(k=n-i)$. On peut donc passer au corps galoisien $\overline{\mathfrak K}$ relativement à $P(x_{i+1},\ldots,x_n)$, dans lequel le polynôme irréductible $P^{(i)}(x_i,x_{i+1},\ldots,x_n)$ ayant pour zéro $\alpha_i$ se décompose en facteurs linéaires. Ce polynôme irréductible $P^{(i)}$ joue, lorsque l'on revient aux indéterminées primitives $y$, le rôle de l'équation fondamentale; car $\alpha_i$ devient égale à la forme fondamentale $u_1\beta_1+\cdots+u_n\beta_n$, où $\beta$ désigne un système de zéros dans les $y$, algébriquement dépendant des paramètres $x_{i+1},\ldots,x_n$; la décomposition en facteurs linéaires donne donc le produit sur tous les zéros. Mais le polynôme irréductible $P^{(i)}$, ou une puissance de celui-ci, peut aussi être regardé comme la norme $N(\frakp)$ de $\frakp$, si l'on considère $\frakp$ comme un module de formes linéaires dans les produits de puissances de $x_1,\ldots,x_{i-1}$, si l'on prend $P(x_{i+1},\ldots,x_n)$ à la place de $P$ comme corps des coefficients, et si l'on définit la norme comme produit sur tous les diviseurs élémentaires -- dont seuls un nombre fini sont distincts de l'unité. Lorsque $P$ est un corps parfait, on obtient alors toujours $N(\frakp)=P^{(i)}$; lorsque $P$ est un corps imparfait, une puissance de $P^{(i)}$ peut apparaître comme norme.

IV. \emph{La norme de l'idéal initial et la signification de la multiplicité.} Si l'on considère de manière correspondante l'idéal $\fraka$ successivement comme module de formes linéaires dans les produits de puissances de $x_1,\ldots,x_{i-1}$ $(i=1,2,\ldots n)$, on peut montrer que ce module, en prenant pour base $P(x_{i+1},\ldots,x_n)$, ne possède chaque fois qu'un nombre fini de diviseurs élémentaires distincts de l'unité, condition de finitude qui résulte de l'existence de la base d'idéal. Le produit de ces diviseurs élémentaires chaque fois en nombre fini donne les normes individuelles; la norme $N(\fraka)$ est définie comme le produit de toutes les normes individuelles. Les plus grands facteurs primaires $Q^{(i)}(x_i,x_{i+1},\ldots,x_n)$ de $N(\fraka)$ correspondent biunivoquement aux idéaux premiers associés $\frakp$ de $\fraka$, de telle sorte que $Q^{(i)}$ devient une puissance de $P^{(i)}$, où $P^{(i)}$ (ou, le cas échéant, une puissance de $P^{(i)}$) désigne la norme $N(\frakp)$ de $\frakp$. Il vient ainsi une décomposition
\[
 N(\fraka)=N(\frakp_1)^{\varrho_1}\cdots N(\frakp_r)^{\varrho_r}=Q_1\cdots Q_r,
\]
-- où, le cas échéant, $N(\frakp)$ doit être remplacée par le plus haut diviseur élémentaire $E(\frakp)=P^{(i)}$ --, et par là, d'après I, II, III, la décomposition de la norme en facteurs linéaires est effectuée pour tous les zéros de $\fraka$. Comme multiplicité des zéros correspondant à l'idéal premier individuel $\frakp$ de $\fraka$, on peut considérer le degré de $Q^{(i)}$, qui s'interprète comme la dimension, sur $P(x_{i+1},\ldots,x_n)$, du quotient de l'idéal fondamental de $(i-1)$-ième étage par la composante de l'idéal fondamental de $i$-ième étage déterminée par $\frakp$. Cette dernière composante possède $\frakp$ et tous les idéaux premiers de dimension plus élevée comme idéaux premiers associés, tandis que la première ne possède que les idéaux premiers de dimension plus élevée.

Une pleine \emph{insertion du cas d'une indéterminée} est obtenue si l'on passe ici aussi à l'idéal principal. La décomposition donnée sous 1. se scinde alors -- selon que l'on regarde $a(x)$ comme générateur de l'idéal principal ou comme norme -- en deux décompositions séparées: d'une part, la représentation de l'idéal principal comme plus petit commun multiple de plus grandes composantes primaires selon I, ce qui revient ici à une représentation comme produit de puissances d'idéaux premiers; d'autre part, la décomposition de la norme comme produit de puissances des normes des idéaux premiers individuels selon IV. Dans la construction du corps des zéros selon 2., il s'agit en réalité du corps des zéros de l'idéal premier défini par $p(x)$, tandis que, dans 3., c'est la norme qui est décomposée en facteurs linéaires. La définition de la multiplicité selon 4. se subordonne à celle donnée en IV, puisque l'idéal fondamental de l'étage zéro est égal à l'idéal unité, tandis que, pour une indéterminée, l'idéal fondamental du premier étage est déjà égal à l'idéal lui-même.

L'\emph{insertion dans la théorie usuelle de l'élimination} est donnée par le fait que la norme de l'idéal possède la signification d'une résultante; dans cette conception, l'ancien problème de la théorie de l'élimination devient donc une partie de la théorie des idéaux de l'anneau de polynômes, et précisément la partie qui, à côté des théorèmes de décomposition des idéaux, traite aussi la décomposition de leurs normes et la connexion entre les deux décompositions.

\begin{center}
(Reçu le 19. 10. 1923.)
\end{center}
\clearpage

\section*{\NoetherNavAnchor{P26}{26. Construction abstraite de la théorie des idéaux dans un corps de nombres algébriques}26. Construction abstraite de la théorie des idéaux dans un corps de nombres algébriques}

\begin{center}
J. Ber. d. DMV 33 (1924), p. 102
\end{center}

4. E. Noether, Göttingen: Construction abstraite de la théorie des idéaux dans un corps de nombres algébriques.

On a indiqué les propriétés abstraites qui sont entièrement équivalentes à la théorie des idéaux dans un corps de nombres algébriques et qui caractérisent donc tous les anneaux possédant une telle théorie des idéaux. Ce sont les anneaux avec élément unité et sans diviseurs de zéro, dans lesquels vaut la condition de double chaîne -- toute chaîne de diviseurs formée d'idéaux s'interrompt au bout d'un nombre fini de pas, de même toute chaîne de multiples modulo tout idéal différent de l'idéal zéro -- et qui sont intégralement clos relativement à leur corps des fractions. Si l'on abandonne cette dernière hypothèse, on obtient des théorèmes sur la théorie des idéaux dans un ordre fini, qui se trouvent déjà en partie, sans démonstration, chez Dedekind (3e éd.).

Une exposition détaillée paraîtra dans les \emph{Mathematische Annalen}.
\clearpage

\section*{\NoetherNavAnchor{P27}{27. Nombres de Hilbert dans la théorie des idéaux}27. Nombres de Hilbert dans la théorie des idéaux}

\begin{center}
J. Ber. d. DMV 34 (1925), p. 101
\end{center}

E. Noether: Nombres de Hilbert dans la théorie des idéaux. Par nombres de Hilbert on entend en général des nombres qui se rapportent au dénombrement des classes modulo les idéaux. Hilbert lui-même a donné, dans sa \og{}fonction caractéristique\fg{}, une fonction qui, pour les idéaux de l'anneau de polynômes, fournit, à partir d'une certaine borne, le nombre exact de classes modulo l'idéal linéairement indépendantes. Pour les nombres de degré inférieurs, Macaulay a complété cette fonction par une fonction génératrice, qu'Ostrowski a pu évaluer explicitement. Mais Macaulay a aussi indiqué des nombres d'une autre sorte, qui se sont révélés les plus essentiels pour la théorie générale des idéaux, puisqu'ils peuvent être saisis abstraitement. Conformément aux théorèmes généraux de décomposition, il suffit de donner la définition pour les idéaux primaires $\mathfrak q$. Il s'agit du nombre de classes modulo $\mathfrak q$ linéairement indépendantes sur le corps résiduel de l'idéal premier associé $\mathfrak p$. Ces nombres se décomposent en sous-systèmes, donnés chacun par le nombre des composantes irréductibles de $\mathfrak q,\mathfrak q/\mathfrak p,\mathfrak q/\mathfrak p^2,\ldots$. Le nombre du $i$-ième sous-système est aussi caractérisé par le nombre de termes d'une série de composition qui va de $\mathfrak q/\mathfrak p^i$ à $\mathfrak q/\mathfrak p^{i-1}$. La série de composition totale et ses nombres de Hilbert forment l'équivalent de la représentation des idéaux primaires comme puissances d'idéaux premiers, représentation qui n'est plus valable dans le cas général -- par exemple dans les ordres finis d'un corps de nombres algébriques.
\clearpage

\section*{\NoetherNavAnchor{P28}{28. Caractères de groupes et théorie des idéaux}28. Caractères de groupes et théorie des idéaux}

\begin{center}
J. Ber. d. DMV 34 (1925), p. 144
\end{center}

3. E. Noether, Göttingen: Caractères de groupes et théorie des idéaux.

La théorie de Frobenius des caractères de groupes -- c'est-à-dire des représentations de groupes finis -- est conçue comme théorie des idéaux d'un anneau complètement réductible, l'anneau de groupe. On développe donc en général les théorèmes d'isomorphie pour la décomposition additive d'un anneau complètement réductible; la représentation comme somme directe d'idéaux simples directement indécomposables est unique à isomorphisme près, et les idéaux unilatères indécomposables appartenant au même idéal bilatère sont isomorphes entre eux. Si l'on réunit donc les idéaux isomorphes en une classe, il y a autant de classes d'idéaux indécomposables distinctes que d'idéaux bilatères indécomposables.

Dans la spécialisation aux systèmes complètement réductibles de grandeurs hypercomplexes -- en particulier à l'anneau de groupe -- les classes d'idéaux indécomposables et les classes de représentations irréductibles se correspondent mutuellement et de façon univoque. En outre, les idéaux bilatères indécomposables et les caractères se correspondent; et le nombre des composantes additives indécomposables d'un idéal bilatère est égal au rang des idéaux indécomposables. La théorie de Frobenius se trouve ainsi intégrée.

Un exposé détaillé doit paraître dans les Math. Ann.

\clearpage
\setcounter{footnote}{0}

\section*{\NoetherNavAnchor{P29}{29. Le théorème de finitude des invariants des groupes linéaires finis de caractéristique p}29. Le théorème de finitude des invariants des groupes linéaires finis de caractéristique $p$}

\begin{center}
Nachr. v. d. Ges. d. Wiss. zu Göttingen 1926, p. 28--35
\end{center}

\begin{center}
Par Emmy Noether à Göttingen.\\[0.5em]
Présenté par R. Courant à la séance du 30 juillet 1926.
\end{center}

Je montre dans ce qui suit que le théorème bien connu de finitude des invariants des groupes finis de substitutions linéaires\footnote{Pour une démonstration élémentaire de finitude, voir E. Noether, Der Endlichkeitssatz der Invarianten endlicher Gruppen. Math. Ann. 77 (1916), p. 89--92.} reste valable lorsque, comme domaine des coefficients, on prend au lieu d'un corps de nombres ordinaire un corps quelconque -- donc en particulier un corps de caractéristique $p$.\footnote{Pour les notions de théorie des corps, voir E. Steinitz, Algebraische Theorie der Körper. J. f. M. 137 (1910), p. 167--309. Caractéristique au \S{} 4, corps des racines au \S{} 12, corps des fractions au \S{} 3.} Mais les démonstrations de finitude connues échouent dès que l'ordre du groupe est divisible par $p$; il faut faire intervenir des faits arithmétiques plus profonds, à savoir le critère suivant, connu jusqu'ici seulement pour la caractéristique nulle:

\textbf{Critère de finitude}\footnote{Comparer E. Noether, Algebraische und Differentialinvarianten. Jahresber. d. d. Math.-Ver. 32 (1923), p. 177--184, no 3, théorème 4, ainsi que la littérature qui y est indiquée. Dans la dernière phrase du critère il y est dit par erreur ``base de $\mathfrak S$ comme $\mathfrak R$-module'' (système fondamental des grandeurs entières) au lieu de ``comme module sur un certain sous-anneau de $\mathfrak R$''.}: Un domaine d’intégrité $\mathfrak J$ formé de polynômes -- plus généralement de fonctions rationnelles -- à coefficients dans un corps est fini si et seulement si $\mathfrak J$ contient un sous-anneau fini $\mathfrak R$ tel que $\mathfrak J$ dépende algébriquement et intégralement de $\mathfrak R$. Toute base d'intégrité de $\mathfrak R$ est alors complétée en une base d'intégrité de $\mathfrak J$ par une base de $\mathfrak S$ comme module sur un certain sous-anneau de $\mathfrak R$.

Ce critère -- qui présente par lui-même un intérêt autonome -- repose sur l'existence de la base rationnelle et sur le fait que le théorème de la chaîne des diviseurs se conserve lors du passage aux grandeurs entières d'une extension finie. Ce second théorème a été démontré dans la note précédente d'Artin et van der Waerden pour une caractéristique quelconque, mais seulement en posant, pour le domaine initial, une certaine hypothèse de finitude qui est remplie dans le cas des groupes finis (l'anneau des racines représente une extension finie). Pour l'existence de la base rationnelle je donne une démonstration valable pour des corps quelconques comme domaines de coefficients; je peux ainsi démontrer le critère de finitude sous la restriction indiquée ci-dessus.

Que, pour les invariants de groupes finis, le critère de finitude soit rempli repose sur le principe des fonctions symétriques. Tandis qu'en caractéristique nulle les coefficients de la résolvante de Galois fournissent déjà la base d'intégrité, dans le cas général ils n'engendrent que le sous-anneau fini exigé dans le critère de finitude. L'existence de ce sous-anneau fini entraîne, par les mêmes raisonnements, la finitude des invariants ``relatifs'' et des invariants ``modulaires''.

Parmi les invariants considérés ici figurent, comme cas particuliers, les ``invariants projectifs mod $p$'' traités par Dickson et son école.\footnote{Voir Dickson, On Invariants and the theory of Numbers. The Madison Colloquium 1913. La littérature plus récente se trouve dans les volumes des Transactions of the American Mathematical Society parus depuis lors.} Ici le théorème de finitude était connu pour les invariants modulaires, mais non pour les invariants ``formels'' généraux.

\subsection*{\S{} 1. Le critère de finitude}

1. \textbf{Théorème de l'existence de la base rationnelle.} Première formulation: Tout système $S$ de fonctions rationnelles de $n$ indéterminées, à coefficients dans un corps $\mathsf{P}$, possède une base rationnelle finie; c'est-à-dire qu'il existe dans le système un nombre fini de fonctions telles que toute fonction du système puisse s'exprimer rationnellement, à coefficients dans $\mathsf{P}$, au moyen de ces fonctions en nombre fini.

Seconde formulation: Soit $\mathfrak K$ un corps intermédiaire entre $\mathsf{P}$ et $\mathsf{P}(x_1,\ldots,x_n)$ -- où $x_1,\ldots,x_n$ désignent des indéterminées -- de degré de transcendance $t\leqq n$; soit $y_1,\ldots,y_t$ un système irréductible\footnote{Un système est dit irréductible, d'après Steinitz, lorsqu'il est formé de fonctions algébriquement indépendantes. Un système est dit de degré de transcendance $t$ lorsqu'il contient au moins un système irréductible de $t$ éléments, mais aucun de plus d'éléments. Pour les théorèmes simples sur les systèmes irréductibles utilisés ci-dessous, voir Steinitz, \S{} 22.} de $\mathfrak K$. Alors $\mathfrak K$ devient une extension algébrique finie de $\mathsf{P}(y_1,\ldots,y_t)$.

Les deux formulations sont en fait équivalentes. Il suffit en effet de démontrer l'existence de la base rationnelle pour tous les corps intermédiaires; car le corps dérivé d'un système $S$ devient un tel corps intermédiaire $\mathfrak K$, dont la base rationnelle fournit immédiatement une base rationnelle de $S$, puisque chaque élément de $\mathfrak K$ devient une combinaison rationnelle d'un nombre fini de fonctions de $S$. Mais il résulte de la seconde formulation qu'un système irréductible arbitraire $y_1,\ldots,y_t$ de $\mathfrak K$, augmenté des fonctions en nombre fini qui caractérisent $\mathfrak K$ comme extension finie de $\mathsf{P}(y_1,\ldots,y_t)$, fournit une telle base rationnelle. Réciproquement, si, d'après la première formulation, une base rationnelle de $\mathfrak K$ est donnée, alors, par la définition du degré de transcendance, chaque élément de base doit être algébriquement dépendant de $\mathsf{P}(y_1,\ldots,y_t)$; ainsi $\mathfrak K$ se trouve caractérisé comme extension algébrique finie de $\mathsf{P}(y_1,\ldots,y_t)$.

Le théorème sera démontré dans la seconde formulation. Supposons d'abord que $\mathfrak K$ ait le même degré de transcendance $n$ que $\mathsf{P}(x_1,\ldots,x_n)$; soit $y_1,\ldots,y_n$ un système irréductible de $\mathfrak K$ et par conséquent de $\mathsf{P}(x_1,\ldots,x_n)$. À cause du degré de transcendance $n$ de $y_1,\ldots,y_n$, chaque $x_i$ devient algébrique sur $\mathsf{P}(y_1,\ldots,y_n)$; donc $\mathsf{P}(x_1,\ldots,x_n)$, et par là $\mathfrak K$ comme corps intermédiaire, devient une extension algébrique finie de $\mathsf{P}(y_1,\ldots,y_n)$.

Soit maintenant $t<n$; que la numérotation des $x$ soit choisie de telle sorte que $y_1,\ldots,y_t,x_{t+1},\ldots,x_n$ forment un système irréductible dans $\mathsf{P}(x_1,\ldots,x_n)$. Ainsi $x_{t+1},\ldots,x_n$ sont irréductibles relativement à $\mathsf{P}(y_1,\ldots,y_t)$ et donc, d'après la loi transitive de la dépendance algébrique, également irréductibles relativement à $\mathfrak K$ et à tout corps intermédiaire $\mathfrak M$ entre $\mathsf{P}(y_1,\ldots,y_t)$ et $\mathfrak K$. Cela signifie que $\overline{\mathfrak K}=\mathfrak K(x_{t+1},\ldots,x_n)$, respectivement $\overline{\mathfrak M}=\mathfrak M(x_{t+1},\ldots,x_n)$, devient une extension purement transcendante de $\mathfrak K$, respectivement de $\mathfrak M$; donc tout élément de $\overline{\mathfrak K}$ qui n'appartient pas à $\mathfrak K$ devient transcendant relativement à $\mathfrak K$, et de même tout élément de $\overline{\mathfrak M}$ qui n'appartient pas à $\mathfrak M$ devient transcendant relativement à $\mathfrak M$. De même, posons $\overline{\mathsf{P}}$ égal à l'extension purement transcendante $\mathsf{P}(x_{t+1},\ldots,x_n)$; alors $y_1,\ldots,y_t$ sont irréductibles relativement à $\overline{\mathsf{P}}$.

La démonstration va maintenant se ramener au cas déjà traité en prenant $\overline{\mathsf{P}}$ comme domaine des coefficients. De $\mathsf{P}<\mathfrak K< \mathsf{P}(x_1,\ldots,x_t;\,x_{t+1},\ldots,x_n)$\footnote{$\mathsf{P}<\mathfrak K$ signifie: $\mathsf{P}$ est contenu dans $\mathfrak K$, donc est un sous-corps de $\mathfrak K$.} il résulte en effet que $\overline{\mathsf{P}}<\overline{\mathfrak K}<\overline{\mathsf{P}}(x_1,\ldots,x_t)$, où $\overline{\mathfrak K}$ devient, relativement à $\overline{\mathsf{P}}$, du même degré de transcendance $t$ que $\overline{\mathsf{P}}(x_1,\ldots,x_t)$. Par conséquent $\overline{\mathfrak K}$ possède, relativement à $\overline{\mathsf{P}}$ comme domaine des coefficients, une base rationnelle finie, disons $y_1,\ldots,y_t; z_1,\ldots,z_s$. Les $z$ peuvent être choisis dans $\mathfrak K$, puisque $\overline{\mathfrak K}$ est le corps composé de $\mathfrak K$ et de $\overline{\mathsf{P}}$. Il faut montrer que $\mathfrak K$ est égal à $\mathfrak M=\mathsf{P}(y_1,\ldots,y_t;z_1,\ldots,z_s)$. Or $\mathfrak M$ est un corps intermédiaire entre $\mathsf{P}(y_1,\ldots,y_t)$ et $\mathfrak K$; donc tout élément de $\mathfrak K$ est algébrique relativement à $\mathfrak M$. D'autre part $\overline{\mathfrak M}$ est égal à $\overline{\mathfrak K}$ et ainsi $\mathfrak K$ est contenu dans $\overline{\mathfrak M}$. Donc $\mathfrak K$ devient un corps intermédiaire entre $\mathfrak M$ et $\overline{\mathfrak M}$; et comme tout élément de $\overline{\mathfrak M}$ qui n'appartient pas à $\mathfrak M$ est transcendant relativement à $\mathfrak M$, il s'ensuit nécessairement que $\mathfrak K$ est identique à $\mathfrak M$.

Conséquence: Si $\mathsf{P}<\mathfrak K<\mathfrak L<\mathsf{P}(x_1,\ldots,x_n)$ et si $\mathfrak L$ est algébrique relativement à $\mathfrak K$, alors $\mathfrak L$ est une extension algébrique finie de $\mathfrak K$. Car chaque élément d'une base rationnelle de $\mathfrak L$ est alors algébrique relativement à $\mathfrak K$; ainsi $\mathfrak L$ est caractérisé comme extension finie de $\mathfrak K$.

2. \textbf{Démonstration du critère de finitude.} La condition est évidemment nécessaire; car si $\mathfrak J$ est un domaine d’intégrité fini, alors $\mathfrak J$ lui-même peut être pris pour le sous-anneau fini $\mathfrak R$ qui intervient dans le critère.

La condition est suffisante sous l'hypothèse que le corps des racines $\mathsf{P}^{1/p}$ est une extension finie du corps des coefficients\footnote{Voir p. 28, note 2.} lorsque celui-ci est de caractéristique $p$; pour la caractéristique nulle, aucune hypothèse restrictive n'est nécessaire. Pour le montrer, il faut prouver que $\mathfrak J$ possède une base finie comme module sur un sous-anneau de $\mathfrak R$.

Soit $\mathfrak K$ le corps des fractions\footnote{Voir p. 28, note 2.} de $\mathfrak R$ et soit $\mathfrak L$ le corps des fractions de $\mathfrak J$. Ces corps des fractions existent, puisqu'il s'agit d'anneaux sans diviseurs de zéro; et l'on a $\mathsf{P}<\mathfrak K<\mathfrak L<\mathsf{P}(x_1,\ldots,x_n)$. Donc, d'après la conséquence sous 1., $\mathfrak L$ devient un corps d'extension algébrique fini de $\mathfrak K$; en outre, par hypothèse, chaque élément de $\mathfrak J$ est entier relativement à $\mathfrak R$, et $\mathfrak J$ devient ainsi un sous-anneau de l'anneau $\mathfrak S$ formé de tous les éléments de $\mathfrak L$ entiers sur $\mathfrak R$. Mais comme rien n'est supposé sur la clôture intégrale de $\mathfrak R$ dans $\mathfrak K$, le fait que le théorème de la chaîne des diviseurs se conserve lors d'une extension finie ne peut pas être invoqué directement. Il faut plutôt passer d'abord à un sous-anneau $\mathfrak T$ de $\mathfrak R$, lui aussi fini, pour lequel cette clôture intégrale est satisfaite.

Supposons d'abord que le domaine des coefficients $\mathsf{P}$ contienne une infinité d'éléments. Alors on peut choisir dans $\mathfrak R$ un système irréductible $g_1(x),\ldots,g_t(x)$ de telle sorte que $\mathfrak R$, et donc aussi $\mathfrak J$, dépende intégralement de $\mathfrak T=\mathsf{P}[g_1(x),\ldots,g_t(x)]$. En effet, soit $f_1(x),\ldots,f_k(x)$ la base d'intégrité de $\mathfrak R$ qui existe par hypothèse; si celle-ci ne forme pas un système irréductible, alors -- après une éventuelle renumérotation des indices -- soit $F(f_k;f_1,\ldots,f_{k-1})=0$ une relation à laquelle satisfait $f_k$. Si l'on remplace $f_1,\ldots,f_{k-1}$ par
\[
 f_1+t_1f_k,\ldots,f_{k-1}+t_{k-1}f_k,
\]
alors $f_k$, et donc aussi $f_1,\ldots,f_{k-1}$, dépend intégralement de ces nouvelles fonctions lorsque les $t_i$ sont choisis comme indéterminées; puisque $\mathsf{P}$ possède une infinité d'éléments, les $t_i$ peuvent être spécialisés en des éléments de $\mathsf{P}$ de telle sorte que cette dépendance intégrale soit conservée. Après un nombre fini d'étapes, en tenant compte de la loi transitive de la dépendance intégrale, on parvient donc à l'anneau cherché $\mathfrak T=\mathsf{P}[g_1(x),\ldots,g_t(x)]$\footnote{Voir Hilbert, Über die vollen Invariantensysteme, \S{} 1. Math. Ann. 42 (1893), p. 313--373.}. Le corps $\mathfrak L$ devient une extension finie du corps des fractions $\mathsf{P}(g_1(x),\ldots,g_t(x))$ de $\mathfrak T$, et l'anneau $\mathfrak S$ formé de tous les éléments de $\mathfrak L$ entiers sur $\mathfrak R$ devient identique à celui de tous les éléments de $\mathfrak L$ entiers sur $\mathfrak T$.

Dans $\mathfrak T$ vaut le théorème de la chaîne des diviseurs pour le système de tous les idéaux, et $\mathfrak T$ est intégralement clos dans son corps des fractions, les deux choses parce que $\mathfrak T$ est isomorphe à l'anneau de polynômes en $t$ indéterminées. De l'hypothèse que $\mathsf{P}^{1/p}$ est fini relativement à $\mathsf{P}$, il résulte en outre que $\mathfrak T^{1/p}$ est fini relativement à $\mathfrak T$.\footnote{Voir la note d'Artin--van der Waerden citée dans l'introduction.} Ainsi vaut dans $\mathfrak S$ le théorème de la chaîne des diviseurs pour le système de tous les $\mathfrak T$-modules -- en particulier $\mathfrak J$, comme sous-anneau de $\mathfrak S$ contenant $\mathfrak T$, possède une base finie comme $\mathfrak T$-module\footnotemark[\value{footnote}]. L'existence d'une telle base finie de $\mathfrak J$ comme $\mathfrak T$-module reste encore valable sans aucune hypothèse restrictive sur le domaine des coefficients $\mathsf{P}$, lorsque $\mathfrak L$ devient une extension de première espèce de $\mathsf{P}(g_1(x),\ldots,g_t(x))$, donc en particulier toujours lorsque $\mathsf{P}$ est de caractéristique nulle.\footnote{Voir par exemple E. Noether, Abstrakter Aufbau der Idealtheorie in algebraischen Zahl- und Funktionenkörpern, \S{} 3. Math. Ann. 96 (1926), p. 26--61.}

Soit $h_1(x),\ldots,h_s(x)$ une base de $\mathfrak S$ comme $\mathfrak T$-module; alors la représentation
\[
 f(x)=A_1(g)h_1(x)+\cdots+A_s(g)h_s(x)
\]
pour un $f(x)$ quelconque de $\mathfrak J$ -- où les $A$, comme éléments de $\mathfrak T$, désignent des polynômes en les $g$ -- montre que $g_1(x),\ldots,g_t(x),h_1(x),\ldots,h_s(x)$ forment la base d'intégrité cherchée.

Si le corps $\mathsf{P}$ ne contient qu'un nombre fini d'éléments, le corps $\mathsf{P}(u)$ obtenu par adjonction d'une indéterminée $u$ -- c'est-à-dire d'un élément transcendant relativement à $\mathfrak J$ -- contient cependant une infinité d'éléments. Ainsi l'anneau $\mathfrak J(u)$ possède une base d'intégrité finie relativement à $\mathsf{P}(u)$ comme domaine des coefficients; et comme $\mathfrak J(u)$ est l'anneau composé de $\mathfrak J$ et de $\mathsf{P}(u)$, cette base peut être choisie dans $\mathfrak J$ -- par décomposition suivant les puissances de $u$ -- bien que les éléments de base $g_i$ de $\mathfrak T$ ne forment plus maintenant un système irréductible. Ainsi tout polynôme $f(x)$ de $\mathfrak J$ admet une représentation $C(u)f(x)=G(u,g_i,h_i)$, où $C(u)$ désigne un polynôme non nul de $\mathsf{P}[u]$ et où $G$ devient entier en $u,g_i,h_i$. La comparaison des coefficients d'une puissance convenable de $u$ montre que les $g_i(x)$ et les $h_i(x)$ donnent une base d'intégrité de $\mathfrak J$. Le \textbf{critère de finitude est ainsi démontré en toute généralité}.

\subsection*{\S{} 2. La finitude des invariants}

1. Par \textbf{groupe linéaire fini de caractéristique $p$}, on entend un groupe $\mathfrak H$ de $h$ substitutions linéaires $A_1,\ldots,A_h$ (de déterminant non nul) à coefficients dans un corps $\mathsf{P}$ de caractéristique $p$. Ici $A_k$ représente la substitution linéaire
\[
\begin{aligned}
 x_1^{(k)}&=a_{11}^{(k)}x_1+\cdots+a_{1n}^{(k)}x_n,\\[-0.2em]
 &\ldots,\qquad
 x_n^{(k)}=a_{n1}^{(k)}x_1+\cdots+a_{nn}^{(k)}x_n,
\end{aligned}
\]
ou, en abrégé: $(x^{(k)})=A_k(x)$. Le groupe $\mathfrak H$ transforme donc la suite $(x)$ aux éléments $x_1,\ldots,x_n$ en les suites $(x^{(k)})$ aux éléments $x_1^{(k)},\ldots,x_n^{(k)}$. Comme l'identité doit être contenue parmi $A_1,\ldots,A_h$, la suite $(x)$ est elle aussi contenue parmi les suites $(x^{(k)})$.

Par \textbf{invariant rationnel entier (absolu) du groupe}, on entend une telle fonction rationnelle entière de $x_1,\ldots,x_n$ -- à coefficients dans $\mathsf{P}$\footnote{Cela ne constitue pas une restriction de la généralité, dans la mesure où tout invariant à coefficients dans un corps de caractéristique $p$ -- lequel doit posséder avec $\mathsf{P}$ un corps commun qui les contienne afin que les invariants soient définis -- peut se décomposer linéairement en un nombre fini d'invariants à coefficients dans $\mathsf{P}$. Dans le cas général, il suffit d'exprimer les coefficients par des coefficients linéairement indépendants relativement à $\mathsf{P}$; alors la propriété d'invariance est satisfaite identiquement en ces coefficients indépendants.} -- qui reste identiquement inchangée lorsqu'on applique $A_1,\ldots,A_h$. À ces invariants appartiennent donc toutes les fonctions symétriques -- à coefficients dans $\mathsf{P}$ -- des suites $(x^{(k)})$. Réciproquement, pour chaque invariant de ce genre, on a
\[
 f(x)=f(x^{(1)})=\cdots=f(x^{(h)})
 \quad\text{ou}\quad
 h\cdot f(x)=f(x^{(1)})+\cdots+f(x^{(h)}).
\]
Il s'ensuit que, dans le cas où l'ordre $h$ du groupe n'est pas divisible par la caractéristique -- donc en particulier dans le cas de la caractéristique nulle -- les fonctions symétriques des suites $(x^{(k)})$, à coefficients dans $\mathsf{P}$, épuisent aussi la totalité des invariants. Comme ces fonctions symétriques possèdent une base d'intégrité finie, la démonstration de finitude est achevée dans ce cas; en même temps, la base d'intégrité s'obtient comme le système $G_{\alpha\alpha_1\ldots\alpha_n}(x)$ des coefficients de la résolvante de Galois\footnote{Voir p. 28, note 1.}:
\[
 \Phi(z,u)=\prod_{k=1}^{h}\bigl(z-u_1x_1^{(k)}-\cdots-u_nx_n^{(k)}\bigr)
\]
\[
 =z^h+\sum G_{\alpha\alpha_1\ldots\alpha_n}(x)
 z^{\alpha}u_1^{\alpha_1}\ldots u_n^{\alpha_n}
 \quad(\alpha\ne h;\ \alpha+\alpha_1+\cdots+\alpha_n=h).
\]

2. Dans le cas où $h$ est divisible par $p$, il n'est pas nécessaire que chaque invariant du groupe soit une fonction symétrique des suites $(x^{(k)})$, puisqu'on ne peut plus conclure de $h\cdot f(x)$ à $f(x)$. Mais l'anneau $\mathfrak R$ dérivé des invariants en nombre fini $G_{\alpha\alpha_1\ldots\alpha_n}(x)$ -- les coefficients de la résolvante de Galois -- forme un sous-anneau fini du système $\mathfrak J$ de tous les invariants, et il reste à montrer que $\mathfrak J$ remplit, relativement à $\mathfrak R$, les hypothèses du critère de finitude.

On peut d'abord, sans restreindre la généralité\footnote{Voir p. 33, note.}, prendre comme domaine des coefficients le sous-corps $\overline{\mathsf{P}}$ de $\mathsf{P}$ dérivé de tous les coefficients de substitution $a_{\mu\nu}^{(k)}$. Ce corps, dérivé d'un nombre fini d'éléments, naît donc par adjonction au corps premier d'un nombre fini d'éléments algébriques ou transcendants; et, puisque le corps premier est parfait, il s'ensuit que $\overline{\mathsf{P}}^{1/p}$ est fini relativement à $\overline{\mathsf{P}}$.\footnote{Voir p. 32, note 2.}

Il faut encore montrer que $\mathfrak J$ dépend intégralement de $\mathfrak R$. D'après 1., parmi les $h$ suites $x^{(k)}$ apparaît aussi la suite initiale $x_1,\ldots,x_n$; donc la résolvante de Galois $\Phi(z,u)$ s'annule identiquement en $u$ pour $z=u_1x_1+\cdots+u_nx_n$. Si l'on spécialise chaque fois tous les $u$ à zéro sauf $u_i$, qui est posé égal à l'unité, on obtient une relation
\[
 x_i^h+\sum x_i^\alpha G_{\alpha00\ldots\alpha_i\ldots00}(x)=0
 \quad (\alpha\ne h;\ \alpha+\alpha_i=h),
\]
qui montre que chaque $x_i$ dépend intégralement de $\mathfrak R$. Il en va donc de même pour un polynôme quelconque en les $x$; en particulier $\mathfrak J$ dépend intégralement de $\mathfrak R$. Les hypothèses du critère de finitude sont ainsi remplies, et la \textbf{finitude des invariants est démontrée}.

3. Il faut remarquer que les mêmes raisonnements donnent la \textbf{finitude des invariants relatifs et modulaires}. Un polynôme $f(x)$ s'appelle \textbf{invariant relatif} du groupe si tous les $f(x^{(k)})$ ne diffèrent de $f(x)$ que par un facteur non nul de $\overline{\mathsf{P}}$. Le système des invariants absolus -- en particulier donc l'anneau $\mathfrak R$ dérivé des $G_{\alpha\alpha_1\ldots\alpha_n}(x)$ -- forme donc un sous-anneau fini de l'anneau dérivé de tous les invariants relatifs; et, comme ci-dessus, les hypothèses du critère de finitude relativement à $\mathfrak R$ sont remplies.

Un polynôme $f(x)$ s'appelle \textbf{invariant modulaire} du groupe lorsque $f(x)$ devient égal à $f(x^{(k)})$ pour tout système spécial de valeurs $x_1=\alpha_1,\ldots,x_n=\alpha_n$ de $\overline{\mathsf{P}}$. Tout invariant absolu est en même temps invariant modulaire; mais, si $\overline{\mathsf{P}}$ ne contient qu'un nombre fini d'éléments, la réciproque n'a pas besoin d'être vraie. Ici encore les hypothèses du critère de finitude sont remplies relativement au sous-anneau $\mathfrak R$ dérivé des $G_{\alpha\alpha_1\ldots\alpha_n}(x)$.
\clearpage
\setcounter{footnote}{0}

\section*{\NoetherNavAnchor{P30}{30. Construction abstraite de la théorie des idéaux dans les corps de nombres et de fonctions algébriques}30. Construction abstraite de la théorie des idéaux dans les corps de nombres et de fonctions algébriques}
\noindent\emph{Math. Ann. 96 (1927), p. 26--61.}

\begin{center}
Par\\[0.5em]
\textbf{Emmy Noether à Göttingen.}
\end{center}

Dans ce qui suit on donne une caractérisation abstraite de tous les anneaux dont la théorie des idéaux coïncide avec la théorie des idéaux de toutes les grandeurs entières du corps de nombres algébriques -- c'est-à-dire dont les idéaux se représentent de façon univoque comme produits de puissances d'idéaux premiers. À ces anneaux appartiennent encore, comme exemples les plus connus, l'anneau de toutes les grandeurs entières d'un corps de fonctions algébriques d'une indéterminée -- plus généralement de plusieurs indéterminées, dès qu'avec Kronecker on se restreint aux idéaux de dimension maximale, donc que l'on passe à une certaine localisation, le domaine fonctionnel.

Les axiomes qui sont complètement équivalents à cette théorie des idéaux sont, pour l'anneau commutatif pris pour base, les suivants:\footnote{Pour la définition des notions fondamentales de la théorie des idéaux, comparer par exemple E. Noether, \emph{Idealtheorie in Ringbereichen}, Math. Ann. 83 (1921), p. 24--66 (cité: \emph{Idealtheorie}). D'ailleurs les notions les plus essentielles sont aussi formulées brièvement dans le présent travail, surtout aux \S{}\S{}1, 4 et 5. Pour être complète, la présentation contient aussi des éléments qui ne sont pas nécessaires aux fins immédiates de ce travail.}

\begin{description}
\item[I. Condition de chaîne des diviseurs.] Toute chaîne d'idéaux dans laquelle chaque idéal est un diviseur propre du précédent s'interrompt au bout d'un nombre fini de pas -- autrement dit: pour toute chaîne de diviseurs d'idéaux, il existe un indice à partir duquel tous les idéaux deviennent égaux.
\item[II. Condition de chaîne des multiples modulo tout idéal différent de l'idéal nul.] Toute chaîne d'idéaux -- qui sont tous diviseurs d'un idéal fixe différent de l'idéal nul -- dans laquelle chaque idéal est un multiple propre du précédent s'interrompt au bout d'un nombre fini de pas.
\item[III.] Existence de l'élément unité de la multiplication.
\item[IV.] Anneau sans diviseurs de zéro.
\item[V. Clôture intégrale dans le corps des fractions.] Tout élément du corps des fractions qui est entier relativement à l'anneau appartient à l'anneau.
\end{description}

À partir de ces axiomes je développe pas à pas une théorie des idéaux de plus en plus restreinte, jusqu'à celle que l'on cherche; et je montre réciproquement que tous les axiomes y sont effectivement satisfaits.

De la condition de chaîne des diviseurs I il résulte, comme je l'ai montré (\emph{Idealtheorie} \S{}4), la représentation d'un idéal comme plus petit commun multiple d'un nombre fini d'idéaux primaires appartenant à des idéaux premiers différents. Je répète ici brièvement cette démonstration (\S{}6), parce qu'elle se simplifie en utilisant la notion de quotient d'idéaux, et parce que je veux corriger une inadvertance: la démonstration originelle utilise en effet, en un point, l'existence non supposée de l'élément unité de la multiplication.\footnote{P. Urysohn a attiré mon attention sur cette inadvertance; ma modification de la démonstration était toutefois un peu plus compliquée que celle communiquée ici, qui est due à E. Artin. Celui-ci avait déjà comblé plus tôt la lacune pour lui-même et avait aussi, indépendamment de moi, envisagé la simplification au moyen du quotient d'idéaux. -- Une autre simplification, qui évite l'introduction de la ``représentation réduite'', je la tire d'une question apparentée chez W. Krull: \emph{Algebraische Theorie der Ringe III}, Math. Ann. 92 (1924), p. 181--213, théorème I.}

L'hypothèse de la condition de chaîne des multiples a pour effet (\S{}7) que, dans l'anneau quotient par tout idéal différent de l'idéal nul, un idéal premier ne peut posséder aucun diviseur propre, à l'exception de l'idéal unité comprenant tous les éléments. Les composantes primaires de ces idéaux sont donc déterminées de façon univoque, à l'exception de celles qui appartiennent à l'idéal unité. Sous une forme un peu moins précise, ce résultat se trouve déjà chez Masazo Sono;\footnote{Masazo Sono, \emph{On the Reduction of Ideals}, Memoirs of the College of Science, Kyoto University, Series A, 7 (1924), p. 191--204.} son hypothèse de l'existence d'une série de composition dans l'anneau quotient est équivalente à la ``condition de double chaîne'' dans cet anneau, comme je le montre brièvement à la fin (\S{}10). Par condition de double chaîne j'entends ici la validité des conditions de chaîne des diviseurs et des multiples sans restriction aux idéaux différents de l'idéal nul.

Si l'axiome III -- existence de l'élément unité -- est encore rempli, l'idéal unité ne se présente pas comme idéal premier associé; les composantes primaires sont donc déterminées de façon univoque; elles deviennent étrangères deux à deux et leur plus petit commun multiple est égal au produit. Les axiomes I, II, III sont remplis pour tous les ordres finis d'un corps de nombres algébriques; ici Dedekind avait déjà énoncé (Dirichlet--Dedekind, \emph{Zahlentheorie}, 3e éd., 11e suppl., \S{}172), sans démonstration entièrement développée, la représentation univoque d'un idéal comme produit d'idéaux d'une seule espèce deux à deux étrangers (composantes primaires).

De l'hypothèse IV -- absence de diviseurs de zéro -- il résulte que les différentes puissances d'un idéal différent de l'idéal nul et de l'idéal unité sont toutes différentes, et que le corps des fractions existe. Si enfin l'axiome V de la clôture intégrale dans le corps des fractions est encore rempli, alors les idéaux primaires -- déterminés univoquement à cause de IV -- deviennent des puissances d'idéaux premiers, ce qui donne le théorème usuel de décomposition (\S{}8). Cette dernière démonstration repose sur une conséquence de la clôture intégrale tirée déjà par Dedekind (3e éd., \S{}172) dans le cas du corps de nombres: on peut représenter un élément véritablement fractionnaire comme quotient d'éléments entiers de telle sorte que même le carré du numérateur ne soit pas divisible par le dénominateur. Dans les présentations ultérieures, cette conséquence est remplacée -- également comme conséquence de la clôture intégrale -- par le théorème de Gauss généralisé, beaucoup moins transparent, ou par des théorèmes de modules correspondants, un peu plus forts; tous ces théorèmes supposent, pour leur application, des éléments de base, tandis qu'ici on travaille avec les idéaux eux-mêmes.

Dans la réciproque \hbox{\(\S 9\)}, les axiomes autres que l'axiome V s'obtiennent presque directement; celui-ci s'obtient en démontrant qu'un élément véritablement fractionnaire peut toujours se représenter de telle sorte qu'aucune puissance du numérateur ne soit divisible par le dénominateur, ce qui est donc encore un renforcement de la conséquence primitivement tirée de V. Mais cette démonstration ne réussit qu'en tirant de la représentation des idéaux les conclusions usuelles: représentation par idéaux principaux modulo un idéal fixe et théorie des idéaux fractionnaires, c'est-à-dire des quotients d'idéaux dans le corps. Le fait que la clôture intégrale résulte de la représentation par puissances d'idéaux premiers était déjà connu de Dedekind, comme le montre une remarque (3e éd., \S{}172, p. 522 en bas).

Les axiomes I à V impliquent que les quatre décompositions en général distinctes -- en plus petits idéaux étrangers deux à deux, en idéaux mutuellement premiers, en plus grandes composantes primaires et en idéaux irréductibles -- coïncident; et réciproquement, de cette coïncidence et des axiomes I--IV résulte la clôture intégrale. Pour les anneaux avec diviseurs de zéro, le fait que les autres axiomes soient remplis -- V étant à définir comme clôture intégrale dans l'anneau total des fractions -- n'implique pas nécessairement la coïncidence des décompositions, comme le montre l'exemple de l'anneau quotient par certains idéaux d'un ordre fini \hbox{\(\S 9\)}.

Au \S{}1 j'esquisse la théorie des grandeurs entières relativement à un anneau, jusqu'à la conséquence déduite par Dedekind de la clôture intégrale. Au \S{}2 je montre comment les conditions de chaîne se transfèrent de l'anneau aux modules formés de formes linéaires, avec cet anneau comme domaine de coefficients et de multiplicateurs.\footnote{Cette forme des théorèmes de modules provient de Prüfer (Neue Begründung der algebraischen Zahlentheorie, \S{}2, Math. Ann. 94 (1924), p. 198--243) et est plus transparente que le transfert usuel des théorèmes de base.} J'en tire (\S{}3) la conséquence que, lors du passage aux grandeurs entières d'un corps d'extension algébrique (de première espèce), les axiomes I à IV se conservent pour tout ordre, tandis que V vaut aussi pour l'ordre formé de toutes les grandeurs entières. Ce fait -- naturellement connu pour le corps de nombres algébriques -- permet de subordonner les corps de fonctions mentionnés au début; en particulier, par la simple démonstration que le domaine fonctionnel dérivé des polynômes entiers en plusieurs indéterminées satisfait aux axiomes, la théorie des idéaux de Kronecker est elle aussi pleinement intégrée.\footnote{La différence n'apparaît donc qu'avec les théorèmes sur les discriminants, du fait que l'anneau quotient du domaine fonctionnel par un nombre premier devient un corps imparfait, et que des extensions de seconde espèce peuvent ainsi se présenter.}

La théorie des idéaux développée ensuite ne présuppose pas ces \S{}\S{}2 et 3, qui ne servent qu'à l'intégration. Aux \S{}4 et \S{}5 sont données les bases indépendantes des axiomes I à V; ainsi ressort plus nettement que dans la fondation originelle quelles parties de la théorie sont indépendantes d'hypothèses de finitude.

\subsection*{\S~1. Théorie des grandeurs entières}

Soit donné un anneau commutatif\footnote{Un anneau est, comme on sait, défini lorsqu'une relation d'égalité est donnée dans un ensemble, et en outre deux opérations, l'addition et la multiplication, satisfaisant aux lois usuelles: l'addition est associative, commutative et réversible de façon univoque; la multiplication est associative -- et commutative lorsqu'il s'agit d'un anneau commutatif -- et distributive par rapport à l'addition. Si la définition de l'égalité que l'on prend pour base n'est pas l'identité ensembliste, alors, pour que la même définition de l'égalité soit conservée dans tout sous-système, un tel sous-système (sous-anneau, module, idéal) doit contenir, avec n'importe quel élément, tous les éléments qui lui sont égaux. De même, dans toute extension, il faut supposer que la nouvelle définition de l'égalité englobe l'ancienne. Toutes les notions telles que ``un nombre fini d'éléments'', ``correspondance univoque'', \dots{} sont entendues au sens de cette définition de l'égalité; ainsi, ``il existe un et un seul élément'' signifie par exemple qu'il n'existe qu'un seul élément à égalité près. Du reste, on peut toujours passer de l'égalité à l'identité ensembliste en réunissant les éléments égaux en une classe et en introduisant ces classes comme nouveaux éléments. Les remarques données ici sur la définition de l'égalité sont dues à R. Hölzer.} $T$ sans diviseurs de zéro, avec élément unité $e$ pour la multiplication (axiomes III et IV); dans $T$ on distingue un sous-anneau fixe $R$, contenant l'élément unité. Les notions de module, d'ordre et d'élément entier doivent partout se rapporter à cet anneau fixe $R$, que l'on omettra plus loin dans la notation par brièveté.\footnote{La théorie des grandeurs entières subsiste lorsque l'on admet des diviseurs de zéro, pourvu que l'on suppose en compensation, dans $R$, la condition de chaîne des diviseurs; d'après le \S~2, cela permet aussi la démonstration du lemme. Comme cette version de la théorie des grandeurs entières ne sera pas utilisée plus loin, on ne la signalera que dans des notes. Toutes les définitions en question subsistent en présence de diviseurs de zéro; la plupart exigent même, comme on sait et comme il est facile de le voir, des hypothèses plus faibles encore. Ce qui ne subsiste pas, c'est la formulation de Dedekind: ``Un nombre est entier lorsqu'il possède une enveloppe.'' En effet, lorsque des diviseurs de zéro apparaissent, le quotient de deux modules finis n'a pas à rester fini; c'est pourquoi l'ordre est ici défini directement, et non comme quotient de modules.} Les éléments de $R$ seront désignés par des lettres latines, ceux de $T$ en général par des lettres grecques.

\textbf{1.} Un système $M$ d'éléments de $T$ s'appelle, comme d'habitude, un $R$-module -- brièvement un module -- si, avec deux éléments $\alpha$ et $\beta$, il contient aussi leur différence, et, avec $\alpha$, aussi $r\alpha$, où $r$ désigne un élément quelconque de $R$. On dit que $M$ est divisible par $N$,
\[
        M\equiv 0 \pmod N,
\]
et que $M$ est un multiple, $N$ un diviseur, si tout élément de $M$ est aussi contenu dans $N$. Les $R$-modules dont tous les éléments appartiennent à $R$ s'appellent les idéaux de $R$.

Il est clair que l'intersection -- le plus petit commun multiple $[\ldots,M_i,\ldots]$ -- d'un nombre quelconque de modules de $T$ est encore un module; de même, $T$ est un module. Ainsi existe de façon univoque le module dérivé d'un système arbitraire $\Sigma$ d'éléments de $T$: c'est l'intersection de tous les modules qui contiennent $\Sigma$. La somme -- le plus grand commun diviseur $(\ldots,M_i,\ldots)$ -- d'un système arbitraire de modules est le module dérivé de leur réunion. Le produit $MN$ de deux modules est défini comme le module dérivé des éléments $\mu\nu$, où $\mu$ parcourt tous les éléments de $M$ et $\nu$ tous ceux de $N$. Le produit satisfait donc aux lois associative et commutative, puisque celles-ci valent pour les éléments de l'anneau. De la loi distributive dans $T$ résulte en outre, pour les modules, la loi distributive
\[
        (\ldots,M_i,\ldots)N=(\ldots,M_iN,\ldots),
\]
car, d'après cette même loi dans $T$, les deux membres sont le module dérivé de tous les éléments $\mu_i\alpha$.

Comme $R$ lui-même est un module, la condition que $R$ soit le domaine des multiplicateurs s'exprime par $RM\equiv0\pmod M$. Comme, par hypothèse, $R$ contient l'élément unité, on a aussi $M\equiv0\pmod{RM}$, et donc $M=RM$.

Un module $M$ s'appelle fini lorsqu'il existe au moins un système $\Sigma$, formé d'un nombre fini d'éléments, dont $M$ est dérivé; les éléments $\mu_1,\ldots,\mu_n$ de $\Sigma$ s'appellent une base du module
\[
        M=(\mu_1,\ldots,\mu_n).
\]
La somme et le produit de deux modules finis -- et donc d'un nombre fini de modules finis -- sont encore finis, puisqu'ils sont dérivés respectivement de la réunion ou des produits des éléments de base. Cette dernière assertion résulte de la représentation $r_1\mu_1+\ldots+r_n\mu_n$ de tous les éléments $\mu$ de $M$ par la base,\footnote{Si $R$ ne possède pas d'élément unité, la représentation par la base prend la forme $r_1\mu_1+\ldots+r_n\mu_n+n_1\mu_1+\ldots+n_n\mu_n$, où les entiers $n_i$ ne sont pas des éléments de l'anneau, mais des symboles qui désignent l'addition ou la soustraction répétée.} et de la loi commutative de la multiplication dans $T$.

\textbf{2.} Un anneau d'extension de $R$ situé dans $T$ -- c'est-à-dire un anneau qui contient tous les éléments de $R$ -- s'appelle un $R$-ordre, ou brièvement un ordre. L'intersection d'un nombre quelconque d'ordres de $T$ est encore un ordre, et $T$ lui aussi est un ordre; il existe donc de façon univoque l'ordre dérivé d'un système arbitraire $\Sigma$ d'éléments de $T$. Les sommes et les produits d'ordres se définissent de façon correspondante comme pour les modules; en particulier, $O^2=O$ pour tout ordre.

Tout ordre est en même temps un $R$-module; un ordre s'appelle fini lorsqu'il est un module fini. Puisque somme et produit se forment par formation de la somme de modules et du produit de modules, il résulte de 1. que les sommes et les produits d'ordres finis sont de nouveau des ordres finis. En outre vaut la loi transitive: si $A$ est un $R$-ordre fini et $B$ un $A$-ordre fini, alors $B$ est aussi un $R$-ordre fini. En effet, $B$ devient en même temps un $R$-ordre, donc un $R$-module. Si maintenant $\alpha_1,\ldots,\alpha_m$ forment une base de $A$ comme $R$-module, et $\beta_1,\ldots,\beta_n$ une base de $B$ comme $A$-module, les produits $\alpha_i\beta_j$ forment manifestement une base de $B$ comme $R$-module.

\textbf{3.} Un élément $\alpha$ de $T$ s'appelle entier relativement à $R$, ou brièvement entier, si l'ordre $R_\alpha$ dérivé de $\alpha$ est fini; autrement dit, si la chaîne de diviseurs des modules
\[
        U_\mu=(\alpha^0,\alpha,\ldots,\alpha^{\mu-1})
\]
s'interrompt au bout d'un nombre fini de pas, $U_n=U_{n+1}=\ldots=U_{n+\nu}=\ldots$.

La définition admet aussi la seconde forme qui sera utilisée en 4.: un élément $\alpha$ de $T$ s'appelle entier s'il existe au moins un module principal $C$ distinct du module nul -- c'est-à-dire si $C$ est dérivé d'un élément $\gamma\ne0$ -- tel que le produit de l'ordre $R_\alpha$ par $C$ soit un module fini; autrement dit, si la chaîne de modules $CU_\mu$ s'interrompt au bout d'un nombre fini de pas.\footnote{Si l'on admet des diviseurs de zéro dans $R$, il faut exiger l'existence d'un module principal régulier; c'est-à-dire que $\gamma$ doit être supposé non-diviseur de zéro, élément régulier, ce qui, pour les anneaux sans diviseurs de zéro, revient à $\gamma\ne0$.}

Enfin, la définition est aussi identique à la forme usuelle: un élément $\alpha$ de $T$ s'appelle entier s'il satisfait au moins à une équation
\[
        \alpha^n+r_1\alpha^{n-1}+\ldots+r_n=0,
\]
où les $r$ sont des éléments de $R$.

Ces différentes définitions sont effectivement identiques. En effet, comme $R_\alpha$ coïncide avec le module dérivé de toutes les puissances de $\alpha$, on peut toujours choisir, lorsque $R_\alpha$ est fini, une base du module $R_\alpha$ parmi ces puissances. Si une telle base est donnée, par exemple, par $e=\alpha^0,\alpha,\ldots,\alpha^{n-1}$, cela signifie que la chaîne des $U_\mu$ s'interrompt à $U_n$; et réciproquement, de $U_n=U_{n+\nu}=\cdots$ on obtient cette base. Mais l'interruption de la chaîne des $U_\mu$ est aussi identique à l'existence de l'équation
\[
        \alpha^n+r_1\alpha^{n-1}+\ldots+r_n=0.
\]
Avec $R_\alpha$, le module $CR_\alpha$ est lui aussi fini; la chaîne des $CU_\mu$ s'interrompt donc. Réciproquement, de la finitude de $CR_\alpha$, donc de l'interruption des $CU_\mu$, résulte aussi l'interruption des $U_\mu$, et par suite la finitude de $R_\alpha$. En effet, puisque $C$ est supposé être un module principal distinct du module nul, l'égalité $CU_n=CU_{n+\nu}$ entraîne aussi $U_n=U_{n+\nu}$.

La propriété d'anneau des grandeurs entières et la loi transitive de la dépendance entière reposent sur le

\medskip
\noindent\textbf{Lemme.} Si $O$ est un ordre fini de $T$ et si $\alpha$ est un élément quelconque de $O$, alors l'ordre $R_\alpha$ dérivé de $\alpha$ est fini; autrement dit, tous les éléments d'un ordre fini sont entiers.

La preuve résulte des raisonnements usuels. Si $\omega_1,\ldots,\omega_n$ est une base de $O$ comme $R$-module, alors, en raison de la propriété d'anneau, $\alpha\omega_i$ appartient aussi à $O$. On a donc
\[
        \alpha\omega_i=r_{i1}\omega_1+\ldots+r_{in}\omega_n,
\]
et de là
\[
        \det(r_{ij}-e_{ij}\alpha)=0,
\]
où $e_{ij}=0$ pour $i\ne j$, $e_{ii}=e$; ainsi $\alpha$ est démontré entier. Pour l'annulation du déterminant on utilise l'hypothèse que $T$, et donc $O$, est sans diviseurs de zéro.\footnote{Si la condition de chaîne des diviseurs est supposée dans $R$, mais que des diviseurs de zéro soient admis, le lemme devient une conséquence immédiate du théorème des modules du \S~2, d'après lequel la condition de chaîne se transmet aux modules provenant de $O$. Cette démonstration aussi, dans le cas du corps de nombres, est due à Dedekind (4e éd., \S~173, III), et elle représente la première application des conditions de chaîne dans la littérature. Dans les conséquences à donner sous 4., Dedekind utilise au lieu de cela le fait plus spécial que le nombre des classes modulo un idéal dans le corps de nombres algébriques est fini.}

Le système $G$ de toutes les grandeurs entières de $T$ forme un ordre. $G$ contient $R$; en effet les éléments de $R$ sont entiers, puisque $R$ forme un $R$-module fini ayant l'élément unité pour base. Mais $G$ a aussi la propriété d'anneau: car l'ordre $R_{\alpha,\beta}$ dérivé de deux grandeurs entières arbitraires $\alpha,\beta$ devient égal au produit $R_\alpha\cdot R_\beta$, donc fini. Ainsi, d'après le lemme, $\alpha+\beta$ et $\alpha\beta$, comme éléments d'un ordre fini, sont entiers.

\medskip
\noindent\textbf{Loi transitive de la dépendance entière.} Si $O$ est un ordre composé de grandeurs entières, alors tout élément de $T$ entier relativement à $O$ est aussi entier relativement à $R$. Par définition, $\alpha$ satisfait à une équation
\[
        \alpha^m+\omega_1\alpha^{m-1}+\ldots+\omega_m=0,
\]
et est donc aussi entier relativement à l'ordre $O'=R_{\omega_1,\ldots,\omega_m}$ dérivé de $\omega_1,\ldots,\omega_m$, lequel devient fini comme produit $R_{\omega_1}\ldots R_{\omega_m}$. Par la loi transitive donnée sous 2., l'ordre $O'_\alpha$ fini relativement à $O'$, dérivé de $\alpha$, devient donc aussi un $R$-ordre fini, et $\alpha$ devient entier, par le lemme, comme élément d'un ordre fini.

\textbf{4.} L'anneau $R$ s'appelle intégralement clos dans $T$ si tout élément de $T$ qui est entier relativement à $R$ appartient à $R$. D'après la seconde forme de la définition des grandeurs entières en 3., un élément $\alpha$ de $T$ appartient donc à $R$ si et seulement s'il existe au moins un module principal $C$ distinct du module nul tel que $C\cdot R_\alpha$ forme un module fini.\footnote{Si l'on admet des diviseurs de zéro dans $R$, il faut de nouveau supposer que $C$ est un module principal régulier; de même, dans la conséquence I, il faut supposer que l'élément $c$ est régulier.}

Du concept d'anneau intégralement clos, Dedekind (3e éd., \S~172) a tiré deux conséquences importantes, qui permettent de mettre ce concept à profit pour la théorie des idéaux:

\medskip
\noindent\textbf{Conséquence de Dedekind I.} Soit $R$ intégralement clos dans $T$, et supposons en outre remplie, comme hypothèse supplémentaire dans $R$, la condition de chaîne des diviseurs pour les idéaux (axiome I). Si $\beta$ est un élément non entier de $T$, et si $c\ne0$ est un élément de $R$, alors il existe un exposant $\sigma\ge0$ tel que, dans la suite des éléments
\[
        c,\;c\beta,\ldots,c\beta^r,\ldots
\]
les éléments $c,\ldots,c\beta^\sigma$ soient entiers, tandis que tous les autres ne le sont pas.

Si tous les éléments $c\beta^r$ étaient entiers, ils appartiendraient à $R$ en raison de la clôture intégrale de $R$. Mais alors la chaîne de modules $CB_1,CB_2,\ldots,CB_r,\ldots$ -- où $C$ désigne le module dérivé de $c$, et $B_r$ celui dérivé de $\beta^0,\ldots,\beta^{r-1}$ -- se transformerait en une chaîne d'idéaux $C_1,\ldots,C_r,\ldots$, avec $C_r=(c,c\beta,\ldots,c\beta^{r-1})$, qui, par hypothèse, s'interrompt au bout d'un nombre fini de pas. Mais alors, contrairement à l'hypothèse, $R_\beta$ serait fini et $\beta$ serait entier; il existe donc des éléments non entiers parmi les $c\beta^r$. En outre, avec n'importe quel élément non entier, tous les suivants sont également non entiers; car si $c\beta^\tau$ est entier, donc appartient à $R$, alors, pour $\rho<\tau$, $c\beta^\rho$ satisfait à l'équation
\[
        (c\beta^\rho)^\tau-c^{\,\tau-\rho}(c\beta^\tau)^\rho=0
\]
et il est donc aussi entier. Si donc $c\beta^{\sigma+1}$ est le premier élément non entier, alors, puisque $c$ est entier, on a $\sigma\geqq0$, et c'est l'exposant cherché.

Si l'on spécialise l'anneau d'extension $T$ au corps des fractions de $R$ -- c'est-à-dire au corps obtenu par adjonction de toutes les paires d'éléments\footnote{Dans la formulation bien connue de Steinitz, J. f. M. 137. Si l'on admet des diviseurs de zéro dans $R$, l'anneau total des fractions doit remplacer le corps des fractions; il s'obtient en inversant tous les éléments réguliers de $R$. Ici la conséquence de Dedekind II n'est plus remplie; car $n$ ne sera pas en général un élément régulier.} --, il en résulte la

\medskip
\noindent\textbf{Conséquence de Dedekind II.} Soit $R$ intégralement clos dans son corps des fractions $K$, et supposons que, dans $R$, la condition de chaîne des diviseurs pour les idéaux soit remplie. Si $\beta$ est un élément non entier de $K$, alors $\beta$ admet une représentation comme quotient d'éléments de $R$:
\[
        \beta=m/n
\]
de telle sorte que $m^2/n$ soit lui aussi non entier.

Posons $\beta=b/c$ avec $c\ne0$. Alors, dans la suite $c,c\beta=b,c\beta^2\ldots$, les deux premiers éléments sont certainement entiers, donc $\sigma\geqq1$, où $\sigma$ désigne l'exposant de la conséquence I. Si l'on pose
\[
        m=c\beta^\sigma,\qquad n=c\beta^{\sigma-1},
\]
alors $m/n=\beta$ et $m^2/n=c\beta^{\sigma+1}$ est non entier.

La loi transitive de la clôture intégrale existe sous la forme suivante: si $R$ est un anneau quelconque de $T$, et si $G$ désigne le système de toutes les grandeurs de $T$ entières relativement à $R$, alors $G$ se compose aussi de toutes les grandeurs de $T$ entières relativement à $G$. Car toute grandeur entière relativement à $G$ est aussi entière relativement à $R$, d'après la loi transitive de 3., et appartient donc à $G$.
\subsection*{\S{}2. Conditions de chaîne dans les domaines de modules finis}

Le théorème des modules, d'après lequel les conditions de chaîne se transfèrent, n'intervient dans l'application à donner ici que pour des modules d'un anneau d'extension. Mais comme la démonstration reste la même dans le cas d'un domaine de modules général, c'est un tel domaine que l'on prendra pour base.

Un système $M$ d'éléments $\alpha,\beta\ldots$ s'appelle domaine de modules relativement à un anneau $R$ lorsque, dans $M$, sont données deux opérations qui conduisent de façon univoque à des éléments de $M$: une composition additive des éléments et une multiplication des éléments par les éléments de $R$; lorsque $\mathsf{M}$ forme, par rapport à l'addition, un groupe abélien, tandis que la multiplication satisfait à la loi associative; et lorsque la loi distributive vaut sous ses deux formes.\footnote{Comparer \emph{Idealtheorie}, \S{}9.}

Pour un domaine de modules, toutes les définitions de modules données au \S{}1, 1., qui ne se rapportent pas à la multiplication, restent manifestement valables. En particulier, $\mathsf{M}$ s'appelle un domaine de modules fini lorsqu'il existe un nombre fini d'éléments $\xi_1,\ldots,\xi_k$ de $\mathsf{M}$ tels que $\mathsf{M}$ soit égal au module dérivé de $\xi_1,\ldots,\xi_k$, donc au système de toutes les formes linéaires $r_1\xi_1+\ldots+r_k\xi_k$, si $R$ possède un élément unité, tandis que, sinon, il faut encore ajouter des termes supplémentaires $n_i\xi_i$.

\medskip
\noindent\textbf{Théorème des modules.} Soit $M$ un domaine de modules fini relativement à un anneau commutatif $R$ avec élément unité, et supposons que, dans $R$, soit remplie la condition de chaîne des diviseurs, respectivement des multiples, pour les idéaux. Alors, dans $M$, est remplie la condition de chaîne des diviseurs, respectivement des multiples, pour les modules.

La démonstration repose sur le fait que l'on associe de façon univoque à chaque module de $M$ un système d'un nombre fini d'idéaux de $R$, de telle sorte qu'à un diviseur propre, respectivement à un multiple propre, du module corresponde chaque fois au moins un idéal diviseur propre, respectivement un idéal multiple propre.

Un élément de $\mathsf{M}$ est dit de longueur $i$ s'il peut se représenter d'au moins une manière comme forme linéaire en $\xi_1,\ldots,\xi_i$, tandis qu'il n'admet aucune représentation comme forme linéaire en $\xi_1,\ldots,\xi_{i-1}$. Associons maintenant à un module quelconque $W$ de $\mathsf{M}$ les $k$ idéaux $\mathfrak a_1,\ldots,\mathfrak a_k$ de $R$: on entendra par $W_i$ le module de tous les éléments de $W$ de longueur $\leqq i$, et $\mathfrak a_i$ désignera le système de tous les coefficients de $\xi_i$ dans $W_i$. Il en résulte d'abord ceci:

Si $B$ est un diviseur propre de $W$ et si $\mathfrak b_1\ldots,\mathfrak b_k$ sont les idéaux associés à $B$, alors, parmi les $\mathfrak b_i$, au moins l'un est un diviseur propre du $\mathfrak a_i$ correspondant. En effet, de $W\equiv0\pmod B$ résulte $W_i\equiv0\pmod {B_i}$, et donc $\mathfrak a_i\equiv0\pmod{\mathfrak b_i}$ pour $i=1,2,\ldots,k$. Par hypothèse, il doit y avoir dans $B$ des éléments qui ne sont pas contenus dans $W$, et donc aussi de tels éléments de plus petite longueur $\ell$. Alors $\mathfrak b_\ell$ devient un diviseur propre de $\mathfrak a_\ell$; car son coefficient $b_\ell$ satisfait
\[
        b_\ell\equiv0\pmod{\mathfrak b_\ell},
        \qquad b_\ell\not\equiv0\pmod{\mathfrak a_\ell},
\]
lorsque
\[
        \beta=b_1\xi_1+\ldots+b_\ell\xi_\ell
\]
est un tel élément de plus petite longueur dans $B$. De $b_\ell\equiv0\pmod{\mathfrak a_\ell}$ résulterait en effet l'existence d'un élément
\[
        \alpha=a_1\xi_1+\ldots+b_\ell\xi_\ell\equiv0\pmod W,
\]
et donc l'existence d'un élément $\beta-\alpha\in B$, $\beta-\alpha\notin W$, de longueur plus petite que $\ell$.

Soit maintenant donnée une chaîne de diviseurs, respectivement de multiples,
\[
        W^{(1)},W^{(2)},\ldots,W^{(\nu)},\ldots ;
\]
supposons que, dans $R$, la condition de chaîne correspondante soit remplie. Si $\mathfrak a_{1\nu},\ldots,\mathfrak a_{k\nu}$ désignent les idéaux associés à $W^{(\nu)}$, alors les suites
\[
        \mathfrak a_{i1},\mathfrak a_{i2},\ldots
\]
forment des chaînes de diviseurs, respectivement de multiples, pour $i=1,\ldots,k$; il existe donc, par hypothèse, un indice $\nu_0$ tel que l'idéal $\mathfrak a_{i\nu_0}$ devienne égal à tous les suivants pour chaque $i$. Mais alors, d'après ce qui vient d'être démontré, le module $W^{(\nu_0)}$ devient égal à tous les suivants, ce qui transfère les conditions de chaîne.

On obtient immédiatement la

\medskip
\noindent\textbf{Conséquence du théorème des modules.} Si, dans $R$, la condition de chaîne des multiples vaut modulo chaque idéal différent de l'idéal nul, alors, dans $M$, la condition de chaîne des multiples vaut modulo chaque module $C$ dont les idéaux associés $\mathfrak c_1,\ldots,\mathfrak c_k$ sont tous différents de l'idéal nul. En effet, sous ces hypothèses, les chaînes de multiples
\[
        \mathfrak a_{i1},\ldots,\mathfrak a_{i\nu},\ldots
\]
s'interrompent au bout d'un nombre fini de pas.

\medskip
\noindent\textbf{Remarque complémentaire.} Si $R$ est un anneau sans élément unité, alors se transfèrent encore la condition de chaîne des diviseurs et la condition de chaîne des multiples modulo les idéaux différents de l'idéal nul. On considère en effet, à la place de $M$, le système $M'$ de tous les éléments de la forme
\[
        a_1\xi_1+\ldots+a_k\xi_k+n_1\xi_{k+1}+\ldots+n_k\xi_{2k},
\]
qui repasse de nouveau dans $M$ lorsque l'on remplace $\xi_{k+\lambda}$ par $\xi_\lambda$. À ce $M'$ on peut associer, comme ci-dessus, les $2k$ idéaux $\mathfrak a_1,\ldots,\mathfrak a_k,\mathfrak n_1,\ldots,\mathfrak n_k$, où les $\mathfrak a_\lambda$ désignent des idéaux de $R$ et les $\mathfrak n_\lambda$ des idéaux de l'anneau des entiers rationnels. Comme, pour les $\mathfrak n_\lambda$, sont remplies la condition de chaîne des diviseurs et la condition de chaîne des multiples modulo chaque idéal différent de l'idéal nul, la démonstration de la condition de chaîne des diviseurs reste valable pour $M'$ et donc pour $M$, tandis que, pour la condition de chaîne des multiples, il faut exiger que les $2k$ idéaux associés à $C$, respectivement à $C'$, soient tous différents de l'idéal nul.
\providecommand{\mR}{\mathfrak{R}}
\providecommand{\mK}{\mathfrak{K}}
\providecommand{\mL}{\mathfrak{L}}
\providecommand{\mS}{\mathfrak{S}}
\providecommand{\mC}{\mathfrak{C}}
\providecommand{\mF}{\mathfrak{F}}
\providecommand{\mc}{\mathfrak{c}}

\subsection*{\S{}3. Passage aux corps d'extension finis}

Pour situer les corps de nombres et de fonctions, il reste à montrer comment les axiomes I à V formulés dans l'introduction se transfèrent lors du passage à des corps d'extension finis.

\textbf{1.} \emph{Supposons que, dans $\mR$, tous les axiomes I à V soient remplis, à l'exception de l'axiome II -- la condition de chaîne des multiples. Soit $\mK$ le corps des fractions de $\mR$, et soit $\mL$ un corps d'extension fini de première espèce de $\mK$.\footnote{D'après Steinitz, un corps d'extension algébrique est dit de ``première espèce'' lorsque chacun de ses éléments est racine d'un polynôme irréductible qui, dans un corps d'extension approprié, se décompose en facteurs linéaires tous distincts.} Alors, dans le système $\mS$ de toutes les grandeurs de $\mL$ entières relativement à $\mR$, les mêmes axiomes sont remplis, et, dans tout ordre de $\mS$, tous ces axiomes le sont encore, à l'exception de la clôture intégrale.}

D'après le \S{}1, 3., $\mS$ forme un anneau pour lequel les axiomes III et IV -- existence de l'élément unité et absence de diviseurs de zéro -- sont remplis. D'après la conclusion du \S{}1, 4., $\mS$ est intégralement clos dans $\mL$; en même temps, $\mL$ devient le corps des fractions de $\mS$, puisque chaque élément de $\mL$ passe, par multiplication avec un élément convenable de $\mR$, en un élément de $\mS$. L'axiome V est donc rempli.

La vérification de I résulte de raisonnements connus. Comme extension de première espèce, $\mL$ s'obtient par adjonction à $\mK$ d'un seul élément $\alpha$, que l'on peut, d'après ce qui précède, supposer entier relativement à $\mR$. Outre $\alpha$, les grandeurs conjuguées $\alpha',\alpha'',\ldots$ contenues dans le corps galoisien de $\mL$ relativement à $\mK$ sont elles aussi entières; il en va donc de même de leur produit de différences, et aussi du carré $D$ de ce produit de différences. Comme $\alpha,\alpha',\ldots$ sont tous distincts, $D$ devient une grandeur non nulle de $\mK$, donc, en raison de la clôture intégrale de $\mR$ dans $\mK$, une grandeur de $\mR$. Par les raisonnements usuels, $\mS$ est, du fait de la clôture intégrale de $\mR$, contenu dans le domaine de modules $M$, fini relativement à $\mR$, engendré par les éléments $\xi_i=\alpha^i/D$; il suffit pour cela, dans la représentation d'un élément de $\mS$ par les $\xi_i$, de passer aux éléments conjugués et de résoudre le système d'équations linéaires qui en résulte pour les coefficients de la représentation. D'après le théorème des modules du \S{}2, la condition de chaîne des diviseurs vaut pour tous les $\mR$-modules de $\mathsf{M}$, donc en particulier pour tous les idéaux de $\mS$. De même, la condition de chaîne des diviseurs vaut pour les idéaux d'un ordre quelconque de $\mS$, puisqu'il s'agit là aussi de $\mR$-modules de $\mathsf{M}$; pour tout ordre, les axiomes III et IV sont en outre remplis.

\textbf{2.} \emph{Si, dans $\mR$, tous les axiomes I à V sont remplis, il en va de même dans $\mS$. Dans tout ordre de $\mS$, tous les axiomes restent remplis, à l'exception de la clôture intégrale.}

Compte tenu de 1., il ne reste plus qu'à démontrer que l'axiome II est rempli. D'après la conséquence du théorème des modules au \S{}2, il suffit de la vérification suivante: si un idéal de $\mS$ différent de l'idéal nul est considéré comme un $\mR$-module $\mC$ dans $\mathsf{M}$, alors les idéaux $\mc_i$ de $\mR$ associés à $\mC$ sont tous différents de l'idéal nul. Or $\mC$ a le même rang linéaire fini relativement à $\mK$ que $M$; car, avec n'importe quel élément $\gamma$, $\mC$ contient aussi $\gamma\alpha^i$. Pour chaque $\xi_i$, il existe donc un élément $c_i\ne0$ de $\mR$ tel que $c_i\xi_i$ appartienne à $\mC$; par l'unicité de la représentation au moyen des $\xi_i$ -- l'indice ne devant parcourir que les valeurs inférieures au degré du corps --, $c_i$ appartient à $\mc_i$, lequel est donc différent de l'idéal nul. Ce raisonnement subsiste pour tous les ordres de même rang que $\mathsf{M}$; pour un ordre de rang plus petit, on passe au corps des fractions, qui, comme corps intermédiaire entre $\mK$ et $\mL$, est lui aussi de première espèce et dont le degré relativement à $\mK$ coïncide avec le rang linéaire de l'ordre; les raisonnements s'y conservent donc. Enfin, remarquons qu'un ordre de $\mS$ différent de $\mS$ n'est jamais intégralement clos: puisque tout ordre contient aussi $\mR$, la clôture intégrale le forcerait à contenir également $\mS$.

Pour la subordination des corps de nombres et de fonctions, il suffit donc de vérifier que les axiomes I à V sont remplis pour les domaines de base: les entiers rationnels, les polynômes en une indéterminée, et le domaine fonctionnel des polynômes en plusieurs indéterminées.

\textbf{3.} \emph{Pour l'anneau $\mR$ des entiers rationnels, respectivement des polynômes en une indéterminée à coefficients dans un corps, les axiomes I à V sont remplis.} Ici I et II résultent soit du fait que, modulo tout nombre non nul, respectivement modulo un tel polynôme en une indéterminée, il n'existe qu'un nombre fini, respectivement qu'un nombre fini de classes modulo ce polynôme, linéairement indépendantes; soit encore du fait que tout idéal devient principal. Il en résulte en effet la décomposition univoque des idéaux en produit de puissances d'idéaux premiers, et par conséquent un idéal différent de l'idéal nul n'est divisible que par un nombre fini d'idéaux. Les axiomes III et IV sont manifestement remplis, le second comme conséquence de la définition de l'égalité; la clôture intégrale V résulte de ce que, grâce à la décomposition univoque des éléments de $\mR$ en irréductibles, à des unités près -- diviseurs de l'unité --, tout élément du corps des fractions qui n'est pas contenu dans $\mR$ peut se représenter comme quotient de deux éléments premiers entre eux de $\mR$, de sorte qu'aucune puissance du numérateur n'est divisible par le dénominateur. Un tel élément ne peut donc satisfaire à aucune équation qui le caractériserait comme entier relativement à $\mR$.

\textbf{4.} \emph{Soit maintenant $\mR$ l'anneau de tous les polynômes en plusieurs indéterminées à coefficients dans un corps, ou à coefficients entiers rationnels.} Alors les axiomes III, IV, V sont remplis comme sous 3.; car ici aussi vaut la décomposition univoque des éléments en irréductibles, à des unités près. Le fait que la condition de chaîne des diviseurs I soit remplie est une conséquence immédiate du théorème de Hilbert sur l'existence d'une base d'idéal; tandis que la condition de chaîne des multiples II ne vaut pas ici.

Mais, en passant au \emph{domaine fonctionnel}, ce qui correspond à une restriction aux idéaux de dimension maximale, on peut encore obtenir la validité de l'axiome II; la validité de I peut alors se démontrer comme sous 3., sans recourir au théorème de Hilbert. Adjoignons en effet à $\mR$ une indéterminée $u$, et considérons donc l'anneau $\mR^*$ de tous les polynômes en $u$ à coefficients dans $\mR$, l'égalité étant définie comme égalité des coefficients. Comme d'ordinaire, on appelle un polynôme de $\mR^*$ primitif relativement à $\mR$ lorsque le plus grand commun diviseur de ses coefficients -- diviseur au sens d'un polynôme de $\mR$, non au sens d'un diviseur d'idéal -- devient une unité de $\mR$. Alors tout polynôme de $\mR^*$ est produit d'un élément de $\mR$ par un polynôme primitif de $\mR^*$, et le produit de polynômes primitifs de $\mR^*$ est de nouveau primitif.

\emph{La totalité des éléments du corps des fractions de $\mR^*$ dont les dénominateurs sont des polynômes primitifs de $\mR^*$ forme donc un anneau, le domaine fonctionnel $\mF$ de $\mR$.} Dans ce domaine fonctionnel $\mF$, sont unités exactement les éléments dont le numérateur et le dénominateur sont des polynômes primitifs de $\mR^*$; tout élément de $\mF$ s'écrit donc de façon univoque comme produit d'un élément de $\mR$ et d'une unité de $\mF$. Ainsi le théorème de décomposition des éléments de $\mR$ se transfère aux éléments de $\mF$; pour $\mF$, en plus des axiomes III et IV, l'axiome V est donc lui aussi rempli comme sous 3.

Dans $\mF$, tout idéal devient en outre principal. En effet, un idéal contient, avec n'importe quel élément de $\mF$, aussi l'élément de $\mR$ qui en provient par multiplication par une unité; et, avec deux éléments quelconques de $\mR$, aussi leur plus grand commun diviseur. Car, avec $f(x)=t(x)\overline f(x)$ et $g(x)=t(x)\overline g(x)$, appartient aussi à l'idéal $t(x)(\overline f(x)+u\overline g(x))$, donc aussi $t(x)$, lorsque $\overline f$ et $\overline g$ sont premiers entre eux et que leur combinaison linéaire devient ainsi une unité. Si donc on part d'un élément arbitraire $f(x)$ de l'idéal et qu'il existe dans l'idéal un élément $g(x)$ non divisible par lui, alors le plus grand commun diviseur $t(x)$ appartient aussi à l'idéal et devient un diviseur propre de $f(x)$; la répétition un nombre fini de fois conduit à un élément de base de l'idéal, de sorte que l'idéal est reconnu principal. Mais alors vaut dans $\mF$ la décomposition univoque des idéaux en produits de puissances d'idéaux premiers, correspondant à la décomposition des éléments de $\mR$; il en résulte, comme sous 3., que les axiomes I et II sont remplis. Parmi les corps d'extension du corps des fractions de $\mF$ qui entrent ici en considération, il ne s'agit que de ceux qui peuvent être engendrés par l'adjonction d'un élément de $\mS$.\footnote{Le système des grandeurs entières relativement à $\mF$ dans ces corps d'extension s'obtient aussi en passant de $\mS$ à un domaine fonctionnel correspondant, comme on passe de $\mR$ à $\mF$. Voir par exemple J. König, \emph{Algebraische Größen}, Leipzig 1903, p. 468/69.} Ainsi, compte tenu de 1. et 2., la validité des axiomes I à V est établie pour les corps de nombres et de fonctions.
\providecommand{\mR}{\mathfrak{R}}
\providecommand{\mA}{\mathfrak{A}}
\providecommand{\mB}{\mathfrak{B}}
\providecommand{\mC}{\mathfrak{C}}
\providecommand{\ma}{\mathfrak{a}}
\providecommand{\mb}{\mathfrak{b}}
\providecommand{\mc}{\mathfrak{c}}
\providecommand{\mm}{\mathfrak{m}}
\providecommand{\mo}{\mathfrak{o}}
\providecommand{\mv}{\mathfrak{v}}

\subsection*{\S{}4. Théorèmes d'isomorphisme. Sommes directes}

Dans ce qui suit, on suppose seulement la propriété d'anneau, respectivement la propriété de module, mais aucun autre axiome.

\textbf{1.} Si $\mathsf{M}$ et $\overline{\mathsf{M}}$ sont des domaines de modules relativement à $\mR$ (\S{}2), on dit que $M$ est homomorphe à $\overline M$ (plus précisément: homomorphe comme $\mR$-module), $\mathsf{M}\sim\overline{\overline{\mathsf{M}}}$, lorsque, à chaque élément de $M$, correspond un et un seul élément de $\overline M$, de telle sorte que $\overline M$ soit ainsi épuisé;\footnote{Le tout au sens de la définition de l'égalité qui prévaut dans $\overline M$; comparer la note 9.} et lorsque, dans cette correspondance, la différence et la multiplication par le même élément de $\mR$ se correspondent; c'est-à-dire lorsque de $\beta\sim\overline\beta$ et $\gamma\sim\overline\gamma$ il résulte toujours que $(\beta-\gamma)\sim(\overline\beta-\overline\gamma)$ et $r\beta\sim r\overline\beta$.

A un $\mR$-module $\mB$ de $\mathsf{M}$ correspond donc homomorphiquement un $\mR$-module $\overline{\mB}$ de $\overline{\mathsf{M}}$; et la totalité $\mC^*$ des éléments de $\mathsf{M}$ auxquels correspondent des éléments d'un $\mR$-module $\overline{\mC}$ de $\overline{\mathsf{M}}$ forme dans $M$ un $\mR$-module déterminé de façon univoque par $\overline{\mC}$, le module $\mC^*$ associé à $\overline{\mC}$, avec $\mC^*\sim\overline{\mC}$. Si, en particulier, $\mA$ est le module associé à l'élément nul de $\overline{\mathsf{M}}$, alors $\mC^*$ devient un diviseur de $\mA$ et se décompose en classes d'éléments de $M$ congrus modulo $\mA$, de telle manière que ces classes correspondent univoquement aux éléments de $\mC$. Si l'on passe de $\mB$ dans $\mathsf{M}$ à $\overline{\mB}$ dans $\overline{\mathsf{M}}$, puis de là de nouveau à $\mB^*$, on obtient $\mB^*=(\mB,\mA)$, c'est-à-dire le plus grand commun diviseur de $\mB$ et de $\mA$; ainsi $\mB^*=\mB$ lorsque $\mB$ est un diviseur de $\mA$.

Si la correspondance des éléments de $\mathsf{M}$ et de $\overline{\mathsf{M}}$ est réversible et univoque, les domaines sont dits isomorphes (plus précisément isomorphes comme $\mR$-modules), $\mathsf{M}\simeq\overline{\mathsf{M}}$.

Si $\mA$ est un $\mR$-module quelconque de $\mathsf{M}$, il naît un domaine de modules $\overline{\mathsf{M}}$ homomorphe à $M$ -- le module quotient $\mathsf{M}\mid\mA$ -- lorsque l'on prend la congruence modulo $\mA$ comme nouvelle relation d'égalité. À chaque élément de $M$ sont alors associés tous les éléments, et seulement eux, qui lui sont égaux dans $\overline{\mathsf{M}}$. Passer, dans $\overline{\mathsf{M}}$, de la définition de l'égalité à l'identité signifie que l'on réunit en une classe modulo $\mA$ tous les éléments égaux dans $\overline{\mathsf{M}}$, et que l'on prend ces classes comme nouveaux éléments de $\overline{\mathsf{M}}$.\footnote{Voir à ce sujet la note 9.}

\emph{Tout homomorphisme est produit par le passage au module quotient;} car si $\mathsf{M}\sim\overline{\overline{\mathsf{M}}}$ et si $\mA$ est le module associé à l'élément nul de $\overline{\overline{\mathsf{M}}}$, alors, comme on l'a montré ci-dessus, $\overline{\overline{\mathsf{M}}}$ devient isomorphe au module quotient $\mathsf{M}\mid\mA$.

\textbf{2. Premier théorème d'isomorphisme.} \emph{Si $\overline{\mathsf{M}}$ désigne le module quotient $\mathsf{M}\mid\mA$, et si $\mC$ est un diviseur de $\mA$, alors vaut l'isomorphisme $\overline{\mathsf{M}}\mid\overline{\mC}\simeq \mathsf{M}\mid\mC$.} En effet, la congruence modulo $\mA$ est en même temps une congruence modulo $\mC$; des éléments égaux modulo $\mA$ restent donc égaux modulo $\mC$. On peut par conséquent former le module quotient $\mathsf{M}\mid\mC$ -- c'est-à-dire l'égalité modulo $\mC$ -- en égalant d'abord les éléments modulo $\mA$, donc en passant à $\overline{\mathsf{M}}$, puis en réunissant parmi ceux-ci les éléments égaux modulo $\mC$, ce qui, dans $\overline{\mathsf{M}}$, correspond à l'égalisation modulo $\overline{\mC}$, donc à la formation de $\overline{\mathsf{M}}\mid\overline{\mC}$.

\textbf{Second théorème d'isomorphisme.} \emph{Si $\mB$ et $\mA$ sont des modules de $\mathsf{M}$, alors vaut l'isomorphisme $(\mB,\mA)\mid\mA\simeq \mB\mid[\mB,\mA]$.} En effet, d'après 1., $\mB$ devient homomorphe à $\overline{\mB}$ lorsque $(\mB,\mA)\mid\mA$ est posé égal à $\overline{\mB}$; et comme, dans cette correspondance, tous les éléments de $[\mB,\mA]$, et seulement ceux-ci, correspondent à l'élément nul de $\overline{\mB}$, l'isomorphisme ci-dessus s'obtient de nouveau d'après 1.

\textbf{3.} Si $\mR$ et $\overline{\mR}$ sont des anneaux (commutatifs), on dit que $\mR$ est homomorphe à $\overline{\mR}$ (plus précisément: homomorphe comme anneau), $\mR\sim\overline{\mR}$, lorsque, à chaque élément de $\mR$, correspond un et un seul élément de $\overline{\mR}$ de telle sorte que $\overline{\mR}$ soit ainsi épuisé, et lorsque, dans cette correspondance, différences et produits se correspondent. Si la correspondance des éléments est réversible et univoque, les anneaux sont dits isomorphes (plus précisément isomorphes comme anneaux), $\mR\simeq\overline{\mR}$.

Tous les raisonnements de 1. et 2. se conservent lorsque l'on remplace les modules par des idéaux dans $\mR$\footnote{Pour les anneaux non commutatifs, il faut prendre pour base des idéaux bilatères.}, et lorsque l'on remplace les notions d'homomorphisme de modules et d'isomorphisme de modules par celles d'homomorphisme d'anneaux et d'isomorphisme d'anneaux. En particulier, si $\mc$ est un diviseur de $\ma$, l'anneau quotient $\mc\mid\ma$ naît lorsque la congruence modulo $\ma$ est introduite comme nouvelle relation d'égalité; il faut remarquer qu'à présent les produits de classes modulo $\ma$ sont encore définis.

Alors $\mc$ est homomorphe à $\mc\mid\ma$; tout homomorphisme est produit de cette manière, et les théorèmes d'isomorphisme valent comme sous 2.:

\textbf{Premier théorème d'isomorphisme.} Si $\overline{\mR}$ désigne l'anneau quotient $\mR\mid\ma$, et si $\mc$ est un diviseur de $\ma$, alors $\overline{\mR}\mid\overline{\mc}\simeq\mR\mid\mc$ au sens de l'isomorphisme d'anneaux.

\textbf{Second théorème d'isomorphisme.} Si $\mb$ et $\ma$ sont des idéaux de $\mR$, alors vaut l'isomorphisme d'anneaux $(\mb,\ma)\mid\ma\simeq\mb\mid[\mb,\ma]$.\footnote{Ces théorèmes d'isomorphisme, connus de la théorie des groupes, se trouvent d'abord pour les modules sous une forme un peu plus spéciale chez Dedekind (comparer 3e éd., p. 484). La forme ci-dessus pour les idéaux se trouve chez Masazo Sono: \emph{On Congruences. I, II, III, IV}, Memoirs of the College of Science, Kyoto Imperial University, 2 (1917), 3 (1918, 1919), comparer 2, p. 215. Dans chaque cas, seul le second théorème d'isomorphisme est exprimé explicitement.}

\textbf{4.} Nous rassemblons maintenant les règles de calcul sur les idéaux premiers entre eux,\footnote{Et cela sous la forme qui remonte à Dedekind (4e éd., \S{}178, III à VIII). Le calcul avec l'élément unité qui apparaît partout dans les exposés récents est beaucoup plus lourd.} dont résulteront en 5. les théorèmes sur les sommes directes. Dans l'anneau commutatif $\mR$, supposons l'existence de l'élément unité $e$ de la multiplication. On a donc $\ma=\mo\ma$ pour tout idéal de $\mR$, où $\mo$ désigne l'idéal unité. Deux idéaux $\ma,\mb$ sont dits premiers entre eux si leur plus grand commun diviseur $(\ma,\mb)$ devient égal à $\mo$.

\emph{4$\alpha$. Si $(\ma,\mb)=\mo$, alors $(\ma,\mb\mc)=(\ma,\mc)$.} En effet, $(\ma,\mc)=(\ma,\mb)\cdot(\ma,\mc)=(\ma^2,\ma\mb,\ma\mc,\mb\mc)$ est divisible par $(\ma,\mb\mc)$, et réciproquement. Il en résulte:

\emph{4$\beta$. De $\mc\mb\equiv0(\ma)$ et $(\ma,\mb)=\mo$ il résulte que $\mc\equiv0(\ma)$.} En effet, $\mc\mb\equiv0(\ma)$ équivaut à $(\ma,\mb\mc)=\ma$; donc, à cause de $(\ma,\mb)=\mo$, on obtient aussi $(\ma,\mc)=\ma$, ou encore $\mc\equiv0(\ma)$. Il en résulte encore:

\emph{4$\gamma$. Si chacun des idéaux $\ma_1,\ldots,\ma_r$ est premier à chacun des idéaux $\mb_1,\ldots,\mb_s$, alors le produit des $\ma_i$ est premier au produit des $\mb_i$.} En effet, de $(\ma,\mb)=\mo$ et $(\ma,\mc)=\mo$ on obtient $(\ma,\mb\mc)=\mo$, d'où l'assertion par répétition un nombre fini de fois.

De 4$\alpha$ et 4$\gamma$ résulte:

\emph{4$\delta$. Si les idéaux $\mb_1,\ldots,\mb_r$ sont deux à deux premiers entre eux, c'est-à-dire si $(\mb_i,\mb_k)=\mo$ pour $i\ne k$, alors leur plus petit commun multiple est égal à leur produit.} Soit $(\ma,\mb)=\mo$, $\mv=[\ma,\mb]$; alors $\mv=\mo\mv=(\ma\mv,\mb\mv)\equiv0(\ma,\mb)$; comme $\ma\mb\equiv0(\mv)$, on obtient donc $\mv=\ma\mb$, d'où l'assertion par répétition un nombre fini de fois en tenant compte de 4$\gamma$.

\emph{4$\varepsilon$. Si les idéaux $\mb_1,\ldots,\mb_r$ sont deux à deux premiers entre eux, et si l'on pose $\ma_i=[\mb_1,\ldots,\mb_{i-1},\mb_{i+1},\ldots,\mb_r]=\mb_1\ldots\mb_{i-1}\mb_{i+1}\ldots\mb_r$, alors $(\ma_1,\ldots,\ma_{i-1},\ma_{i+1},\ldots,\ma_r)=\mb_i$ et donc $(\ma_1,\ldots,\ma_r)=\mo$.} En effet, d'après la loi distributive, $(\ma_1,\ma_2)=\mb_3\ldots\mb_r(\mb_2,\mb_1)=\mb_3\ldots\mb_r$; et donc $(\ma_1,\ma_2,\ma_3)=\mb_4\ldots\mb_r(\mb_3,\mb_1\mb_2)=\mb_4\ldots\mb_r$, d'où l'assertion par répétition un nombre fini de fois avec une numérotation convenable.

L'idéal $\ma_i$ est appelé le complément de $\mb_i$ dans la représentation $\mm=[\mb_1,\ldots,\mb_r]$.

\textbf{5.} Dans l'anneau commutatif $\mR$, supposons de nouveau l'existence de l'élément unité de la multiplication. L'anneau $\mR$ s'appelle somme directe des idéaux $\ma_1,\ldots,\ma_r$ -- en signes: $\mR=\ma_1+\ma_2+\ldots+\ma_r$ -- si chaque élément $c$ de $\mR$ peut se représenter d'une et d'une seule manière sous la forme $c=a_1+a_2+\ldots+a_r$, où chaque $a_i$ est un élément de $\ma_i$.

\emph{Si l'idéal nul est le plus petit commun multiple des idéaux deux à deux premiers entre eux $\mb_1,\ldots,\mb_r$, et si $\ma_i$ est le complément de $\mb_i$, alors $\mR$ devient somme directe des idéaux $\ma_1,\ldots,\ma_r$.} En effet, à cause de $\mo=(\ma_1,\ldots,\ma_r)$, tout élément de $\mR$ admet au moins une représentation additive par les $\ma_i$; à cause de $[\ma_i,\mb_i]=(0)$, cette représentation est univoque; car de $0=a_1+a_2+\ldots+a_r=a_i+b_i$ on obtient $a_i=0$. D'après le second théorème d'isomorphisme, les $\ma_i$ sont isomorphes à l'anneau quotient $\mR\mid\mb_i$. Les ``relations d'orthogonalité'' $\ma_i\ma_k=(0)$ pour $i\ne k$ et $\ma_i^2=\ma_i$ valent; la dernière parce que $\ma_i=\mo\ma_i$.

\emph{Réciproquement, s'il existe une représentation de $\mR$ comme somme directe -- $\mR=\ma_1+\ma_2+\ldots+\ma_r$ -- et si l'on pose $\mb_i=(\ma_1,\ldots,\ma_{i-1},\ma_{i+1},\ldots,\ma_r)=\ma_1+\ldots+\ma_{i-1}+\ma_{i+1}+\ldots+\ma_r$, alors les $\mb_i$ sont deux à deux premiers entre eux et leur plus petit commun multiple est égal à l'idéal nul.} En effet, de $\mo=(\ma_i,\mb_i)$ résulte $\mo=(\mb_k,\mb_i)$ pour $k\ne i$, puisque $\mb_k$ est un diviseur de $\ma_i$. Supposons en outre que $c$ soit divisible par tous les $\mb_i$; on a alors
\[
 c=a_2^{(1)}+\ldots+a_r^{(1)}=a_1^{(2)}+a_3^{(2)}+\ldots+a_r^{(2)}=\ldots=a_1^{(r)}+\ldots+a_{r-1}^{(r)},
\]
où chaque fois $a_i^{(\lambda)}\equiv0\pmod{\ma_i}$. De l'unicité de la représentation il résulte donc, puisqu'une composante devient chaque fois nulle, que $0=a_1^{(\lambda)},\ldots,0=a_r^{(\lambda)}$ pour tout $\lambda$, et par suite $c=0$.\footnote{Les théorèmes donnés sous 5. sont des cas particuliers d'un théorème général sur le rapport entre somme directe et intersection directe. Comparer H. Prüfer, Theorie der Abelschen Gruppen I, Math. Zeitschr. 20 (1924), p. 165--187, \S{}6. Les raisonnements qui y sont donnés restent valables même pour une formulation si générale du concept de groupe que les modules et les idéaux s'y subordonnent comme cas particuliers.}
\providecommand{\mR}{\mathfrak{R}}
\providecommand{\mT}{\mathfrak{T}}
\providecommand{\mA}{\mathfrak{A}}
\providecommand{\mB}{\mathfrak{B}}
\providecommand{\mC}{\mathfrak{C}}
\providecommand{\mD}{\mathfrak{D}}
\providecommand{\ma}{\mathfrak{a}}
\providecommand{\mb}{\mathfrak{b}}
\providecommand{\mc}{\mathfrak{c}}
\providecommand{\md}{\mathfrak{d}}
\providecommand{\mf}{\mathfrak{f}}
\providecommand{\mm}{\mathfrak{m}}
\providecommand{\mpideal}{\mathfrak{p}}
\providecommand{\mq}{\mathfrak{q}}

\subsection*{\S{}5. Idéaux premiers et idéaux primaires}

Soit de nouveau $\mR$ un anneau commutatif pour lequel aucun autre axiome -- pas même l'existence de l'élément unité -- n'est supposé. Nous rassemblons les faits fondamentaux sur les idéaux premiers et les idéaux primaires.\footnote{Dans \emph{Idealtheorie}, \S{}4, l'axiome I de la condition de chaîne des diviseurs est supposé; il n'y apparaît donc pas nettement quelles propriétés des idéaux premiers et primaires sont indépendantes de cet axiome.}

\textbf{1.} Un idéal $\mpideal$ de $\mR$ s'appelle \emph{idéal premier faible} lorsque l'anneau quotient $\mR\mid\mpideal$ est un anneau sans diviseurs de zéro; c'est-à-dire lorsque de $a\not\equiv0(\mpideal)$ et $b\not\equiv0(\mpideal)$ il résulte toujours que $ab\not\equiv0(\mpideal)$. Un idéal $\mpideal^*$ de $\mR$ s'appelle \emph{idéal premier fort} lorsque l'anneau quotient $\mR\mid\mpideal^*$ est un anneau sans idéaux diviseurs de zéro; c'est-à-dire lorsque de $\ma\not\equiv0(\mpideal^*)$ et $\mb\not\equiv0(\mpideal^*)$ il résulte toujours que $\ma\mb\not\equiv0(\mpideal^*)$.

\emph{Tout idéal premier fort est en même temps un idéal premier faible}, comme le montre la spécialisation de $\ma,\mb$ à des idéaux principaux. Réciproquement, \emph{tout idéal premier faible est aussi un idéal premier fort}; on peut donc parler simplement d'idéaux premiers. En effet, soit $\mpideal$ un idéal premier faible; soient $\ma\not\equiv0(\mpideal)$ et $\mb\not\equiv0(\mpideal)$; soient encore $a$ et $b$ des éléments de $\ma$ et de $\mb$ tels que $a\not\equiv0(\mpideal)$ et $b\not\equiv0(\mpideal)$. Alors on obtient $ab\not\equiv0(\mpideal)$, et donc $\ma\mb\not\equiv0(\mpideal)$; ainsi $\mpideal$ est aussi un idéal premier fort.

\textbf{2.} Un idéal $\mq$ de $\mR$ s'appelle \emph{idéal primaire faible} lorsque, dans l'anneau quotient $\mR\mid\mq$, une puissance de chaque diviseur de zéro s'annule; c'est-à-dire lorsque de $a\not\equiv0(\mq)$ et $b^\varkappa\not\equiv0(\mq)$ pour tout $\varkappa$ il résulte toujours que $ab\not\equiv0(\mq)$. Le système $\mpideal$ de tous les éléments de $\mR$ qui deviennent diviseurs de zéro dans $\mR\mid\mq$ forme un idéal, et même un idéal premier, qui devient un diviseur de $\mq$: l'idéal premier associé. En effet, avec $a$, l'élément $ra$ devient aussi diviseur de zéro; avec $a$ et $b$, il en va de même de $a-b$, puisque, avec $a^\varkappa$ et $b^\lambda$, $(a-b)^{\varkappa+\lambda}$ devient toujours divisible par $\mq$. Si en outre $a$ n'est pas diviseur de zéro, alors, par définition, aucune puissance de $a$ ne devient diviseur de zéro; de $a\not\equiv0(\mpideal)$ et $b\not\equiv0(\mpideal)$ résulte donc $a^\varkappa b^\varkappa=(ab)^\varkappa\not\equiv0(\mq)$, et par suite $ab\not\equiv0(\mpideal)$.

Un idéal $\mq^*$ de $\mR$ s'appelle \emph{idéal primaire fort} lorsque, dans l'anneau quotient $\mR\mid\mq^*$, une puissance de chaque idéal diviseur de zéro s'annule; c'est-à-dire lorsque de $\ma\not\equiv0(\mq^*)$ et $\mb^\varkappa\not\equiv0(\mq^*)$ pour tout $\varkappa$ il résulte toujours que $\ma\mb\not\equiv0(\mq^*)$. Comme pour les idéaux primaires faibles, on montre que le plus grand commun diviseur de tous les idéaux de $\mR$ qui deviennent des idéaux diviseurs de zéro dans $\mR\mid\mq^*$ forme un idéal premier, diviseur de $\mq^*$: l'idéal premier associé. Réciproquement, on dit de tous les idéaux primaires qui possèdent le même idéal premier associé $\mpideal$ qu'ils ``appartiennent à $\mpideal$''.

\emph{Tout idéal primaire fort est en même temps un idéal primaire faible}, comme le montre la spécialisation de $\ma,\mb$ à des idéaux principaux. En général, la réciproque n'est toutefois pas vraie.\footnote{C'est ce que montre l'exemple suivant, qui m'a été communiqué par R. Hölzer. Soit $\mR$ l'anneau de polynômes en une infinité dénombrable d'indéterminées $x_i$ à coefficients dans un corps; soit $\mq=(x_1^2,x_2^3,\ldots,x_\nu^{\nu+1},\ldots,x_i x_k\,[i\ne k]\ldots)$. Les diviseurs de zéro dans l'anneau quotient $\mR/\mq$ sont exactement les polynômes divisibles par $\mpideal=(x_1,x_2,\ldots,x_\nu,\ldots)$, et chaque fois une puissance de ces diviseurs de zéro devient divisible par $\mq$; ainsi $\mq$ est un idéal primaire faible ayant $\mpideal$ pour idéal premier associé. En revanche, $\mq$ n'est pas un idéal primaire fort; car si l'on pose $\ma=(x_1,x_3,x_5,\ldots,x_{2\nu+1},\ldots)$ et $\mb=(x_2,x_4,\ldots,x_{2\nu},\ldots)$, alors $\ma\mb\equiv0(\mq)$, mais aucune puissance de $\ma$ ni de $\mb$ n'est divisible par $\mq$.} \emph{Mais si, dans $\mR$, l'axiome I de la condition de chaîne des diviseurs est supposé, tout idéal primaire faible devient aussi un idéal primaire fort;} on peut donc parler simplement d'idéaux primaires. Soit alors $\mq$ un idéal primaire faible; soient $\ma\not\equiv0(\mq)$ et $\mb^\varkappa\not\equiv0(\mq)$ pour tout $\varkappa$. Il existe alors des éléments $a$ de $\ma$ et $b$ de $\mb$ tels que $a\not\equiv0(\mq)$ et $b^\varkappa\not\equiv0(\mq)$ pour tout $\varkappa$; d'où $ab\not\equiv0(\mq)$ et donc $\ma\mb\not\equiv0(\mq)$. Car si, pour chaque $b$ de $\mb$, il existait un exposant $\lambda$ tel que $b^\lambda\equiv0(\mq)$, alors on aurait $\mb^{\lambda_1+\ldots+\lambda_r}\equiv0(\mq)$, où $\lambda_1,\ldots,\lambda_r$ désignent les exposants des éléments d'une base finie.\footnote{Pour pouvoir conclure de l'axiome I à la base d'idéal finie, il faut prendre pour base un bon ordre fixé dans $\mR$ (voir \S{}6).} Ce dernier raisonnement montre encore ceci: il existe un exposant (le plus petit) $\varrho$ tel que $\mpideal^\varrho\equiv0(\mq)$, où $\mpideal$ désigne l'idéal premier associé; $\varrho$ s'appelle l'\emph{exposant} de $\mq$.

\textbf{3.} Dans ce qui suit, la condition de chaîne des diviseurs n'est de nouveau \emph{pas} supposée. Le plus petit commun multiple $[\ma_1,\ldots,\ma_r]$ d'un nombre fini d'idéaux s'appelle une \emph{représentation la plus courte} si aucun $\ma_i$ n'est contenu dans le plus petit commun multiple des autres, c'est-à-dire si aucun $\ma_i$ ne peut être omis.

\emph{Le plus petit commun multiple d'un nombre fini d'idéaux primaires faibles, respectivement forts, appartenant au même idéal premier $\mpideal$, est de nouveau un idéal primaire faible ayant $\mpideal$ pour idéal premier associé. Une représentation la plus courte par un nombre fini d'idéaux primaires faibles, respectivement forts, appartenant à des idéaux premiers différents, n'est pas un idéal primaire faible, respectivement fort.} Soit $\mf=[\mq_1,\ldots,\mq_r]$, et soit $\mpideal$ l'idéal premier associé des idéaux primaires faibles $\mq_i$; alors $\mpideal$ se compose aussi de la totalité des éléments dont une puissance devient divisible par $\mf$. De $a\not\equiv0(\mf)$ et $b^\varkappa\not\equiv0(\mf)$ pour tout $\varkappa$, il résulte donc que $b\not\equiv0(\mpideal)$ et que $a\not\equiv0(\mq_i)$ pour au moins un indice $i$; donc $ab\not\equiv0\pmod{\mq_i}$, et par suite $ab\not\equiv0\pmod{\mf}$. Si l'on remplace partout les éléments par des idéaux, on ne peut plus conclure que $\mb\not\equiv0\pmod{\mpideal}$.

Soit maintenant $\mm=[\mq_1,\ldots,\mq_r]$, avec $\mpideal_i\ne\mpideal_k$ pour $i\ne k$, une représentation la plus courte par idéaux primaires faibles, et soit $r\geqq2$. Il existe alors au moins un idéal premier associé, disons $\mpideal_1$, qui n'est contenu dans aucun des autres; car toute chaîne propre de multiples des $\mpideal$ doit s'interrompre après au plus $r$ pas. Il existe donc des éléments $a_i$ de $\mpideal_i$ tels que $a_i^{\varrho_i}\equiv0\pmod{\mq_i}$ et que $a_2\not\equiv0\pmod{\mpideal_1},\ldots,a_r\not\equiv0\pmod{\mpideal_1}$. Si, de plus, $q_1\not\equiv0\pmod{\mm}$ est un élément de $\mq_1$, alors $q_1 a_2^{\varrho_2}\ldots a_r^{\varrho_r}\equiv0\pmod{\mm}$, mais aucune puissance de $(a_2^{\varrho_2}\ldots a_r^{\varrho_r})$ n'est divisible par $\mm$, car sinon il s'ensuivrait une divisibilité par $\mpideal_1$; ainsi $\mm$ n'est pas un idéal primaire faible. Si l'on remplace partout les éléments par des idéaux, c'est-à-dire $q_1$ par $\mq_1$ et $a_i$ par $\ma_i[\not\equiv0\pmod{\mpideal_1}$, mais $\ma_i^{\varrho_i}\equiv0\pmod{\mq_i}]$, on obtient la démonstration pour les idéaux primaires forts.\footnote{Je dois à B. L. van der Waerden cette modification de ma démonstration originelle, par laquelle la condition de chaîne des diviseurs devient superflue.}

\textbf{4.} \emph{Quotient de modules et quotient d'idéaux.} Soit $\mT$ un anneau d'extension de $\mR$, et soient $\mA,\mB,\mC,\ldots$ des $\mR$-modules dans $\mT$. Le quotient $\mC=\mA:\mB$ est défini comme un module satisfaisant aux conditions suivantes: $\mC\mB\equiv0(\mA)$, et de $\mD\mB\equiv0(\mA)$ résulte $\mD\equiv0(\mC)$; ainsi $\mC$ est le module le plus petit pour lequel $\mC\mB\equiv0(\mA)$. Par là, le quotient est déterminé de façon univoque comme plus grand commun diviseur de tous les modules $\mD$ pour lesquels $\mD\mB\equiv0(\mA)$; de tels $\mD$ existent, puisque le module nul est certainement l'un d'eux.

De la définition il résulte: si $\mB$ est un multiple de $\overline{\mB}$, alors $\mA:\overline{\mB}$ devient un multiple de $\mA:\mB$. Si $\mA$ est un multiple de $\overline{\mA}$, alors $\mA:\mB$ devient un multiple de $\overline{\mA}:\mB$. On a
\[
        \mA:\mB\mC=(\mA:\mB):\mC=(\mA:\mC):\mB .
\]
Si l'anneau d'extension $\mT$ coïncide avec $\mR$, le quotient de modules passe au quotient d'idéaux $\ma:\mb$, pour lequel les règles de calcul ci-dessus valent donc également. Comme $\ma\mb\equiv0(\ma)$, il en résulte encore $\ma\equiv0(\ma:\mb)$.

\textbf{5.} \emph{Si $\ma:\mb=\ma$, alors $\mb$ s'appelle premier à $\ma$.} Un idéal $\mb$ est donc premier à $\ma$ si et seulement si de $\md\mb\equiv0(\ma)$ il résulte toujours $\md\equiv0(\ma)$. Des règles de calcul de 4. résulte: si $\mb$ et $\mc$ sont premiers à $\ma$, alors leur produit et leur plus petit commun multiple sont eux aussi premiers à $\ma$. Si $\mb$ est premier à $\ma$ et $\ma$ premier à $\mb$, alors $\ma$ et $\mb$ sont dits mutuellement premiers. D'après le \S{}4, 4$\beta$, des idéaux premiers entre eux sont aussi mutuellement premiers.

\setcounter{footnote}{25}
\providecommand{\mR}{\mathfrak{R}}
\providecommand{\ma}{\mathfrak{a}}
\providecommand{\mb}{\mathfrak{b}}
\providecommand{\mc}{\mathfrak{c}}
\providecommand{\mm}{\mathfrak{m}}
\providecommand{\mt}{\mathfrak{t}}

\subsection*{\S{}6. Théorie des idéaux sous l'hypothèse de la condition de chaîne des diviseurs}

Soit donné un anneau (commutatif) $\mR$ pour lequel l'axiome I de la condition de chaîne des diviseurs est satisfait; de plus, dans ce qui suit, soit toujours donné un bon ordre fixé de tous les éléments de $\mR$. Par là est en même temps donné aussi un bon ordre de tous les idéaux de $\mR$. En effet, les deux hypothèses entraînent aussitôt que tout idéal de $\mR$ possède une base d'idéal composée d'un nombre fini d'éléments. Si l'on passe donc du bon ordre des éléments de $\mR$ à un ordre quasi-lexicographique de toutes les sous-parties finies de $\mR$,\footnote{Il faut entendre par là ce qui suit: les éléments $a_1,\ldots,a_n$ d'une sous-partie finie $\mathfrak A_n$ soient désignés de telle sorte que l'ordre des indices concorde avec le bon ordre donné dans $\mR$. Toute sous-partie à $n$ éléments précède une telle sous-partie à $m>n$ éléments. Si deux sous-parties ont même cardinal, la première est dite antérieure lorsque le premier élément où elles diffèrent précède, dans le bon ordre, l'élément correspondant de l'autre. Ainsi un premier élément est assigné à tout système de sous-parties finies. Pour un bon ordre donné de $\mR$, aucun autre postulat du choix n'est donc nécessaire; en particulier, si $\mR$ est dénombrable, comme dans le cas du corps de nombres algébriques, aucun postulat du choix n'est en jeu.} et si l'on assigne à chaque idéal, comme base distinguée, la première, dans le bon ordre des sous-parties finies, parmi les différentes bases possibles, alors les idéaux sont bien ordonnés en même temps que les sous-parties finies.

\textbf{Théorème I.} \emph{Sous l'hypothèse de la condition de chaîne des diviseurs, tout idéal de $\mR$ admet une représentation comme plus petit commun multiple d'un nombre fini d'idéaux irréductibles, c'est-à-dire d'idéaux qui ne peuvent pas se représenter comme plus petit commun multiple de deux diviseurs propres.}

Pour la démonstration, il faut montrer ceci: si le théorème I n'est pas satisfait pour un idéal $\mm$, alors $\mm$ possède un diviseur propre pour lequel le théorème I n'est pas non plus satisfait; on peut en tirer, contrairement à la condition de chaîne des diviseurs supposée, une chaîne de diviseurs qui ne s'interrompt pas après un nombre fini de termes. En effet, $\mm$ doit être réductible, car sinon $\mm=[\mm]$ fournirait la représentation voulue. Si $\mm=[\ma,\mb]$, le théorème I ne peut pas être satisfait simultanément pour $\ma$ et pour $\mb$, car sinon il s'ensuivrait une représentation correspondante pour $\mm$. Il existe donc des diviseurs propres de $\mm$ pour lesquels le théorème I n'est pas satisfait; soit $\ma_1$ le premier d'entre eux dans le bon ordre des idéaux. En construisant de même, à partir de $\ma_1$, un diviseur propre $\ma_2$, et en général, à partir de $\ma_i$, un diviseur propre $\ma_{i+1}$, on arrive à une chaîne de diviseurs bien déterminée, qui ne s'interrompt pas après un nombre fini de termes, contre l'hypothèse.\footnote{Le théorème I vaut évidemment de même pour les domaines de modules, si la condition de chaîne des diviseurs est supposée pour le système de tous les modules. Il a un caractère purement ensembliste; il est en même temps le seul théorème dans la démonstration duquel le bon ordre des idéaux est utilisé.}

À cause de la condition de chaîne des diviseurs supposée, on peut, d'après \S{}5, 2, parler simplement d'idéaux primaires. Le lien entre primaire et irréductible est donné par:

\textbf{Théorème II.} \emph{Sous l'hypothèse de la condition de chaîne des diviseurs, tout idéal irréductible est primaire; en d'autres termes, tout idéal non primaire est réductible.}

Comme, lors du passage à l'anneau quotient $\mR/\mm$, la condition de chaîne des diviseurs est conservée par homomorphie, comme à la représentation de $\mm$ comme plus petit commun multiple correspond la représentation de l'idéal nul, et comme, à cause du premier théorème d'isomorphisme (\S{}4, 3), aux diviseurs primaires de $\mm$ correspondent de nouveau de tels idéaux dans $\mR\mid\mm$, et réciproquement, on peut se limiter à la décomposition de l'idéal nul dans l'anneau quotient. Puisque celui-ci est un anneau de même généralité, on peut d'emblée supposer que l'idéal à décomposer est l'idéal nul de $\mR$.

Soit donc l'idéal nul de $\mR$ non primaire, de sorte qu'il existe au moins un couple d'idéaux $\ma,\mb$, où $\mb$ peut encore être supposé idéal principal, tel que
\[
a\ne(0),\qquad \mb^\varkappa\ne(0)\quad\text{pour tout }\varkappa,
\qquad \ma\mb=(0).
\]
Formons la chaîne, qui par hypothèse s'interrompt après un nombre fini de termes, des quotients d'idéaux:
\[
\ma,\ \ma:\mb,\ldots,\ma:\mb^\nu,\ldots;
\]
soit par exemple $\mt=\ma:\mb^{m-1}=\ma:\mb^m\ldots$ égal à tous les idéaux suivants de la chaîne; donc $\mt=\mt:\mb$, ou bien $\mb$ est premier à $\mt$. De plus, $\mt$, comme diviseur de $\ma$, est différent de l'idéal nul, et il en va de même de $\mb^\varkappa$ pour tout exposant $\varkappa$. Pour démontrer le théorème II, il suffit donc d'établir une représentation
\[
(0)=[\mt,\mb^{m+1}].
\]
Par définition,
\[
\mb^{m-1}\mt\equiv0\pmod{\ma},\qquad\text{donc}\qquad \mb^m\mt=(0),
\]
il reste donc seulement à montrer:
\[
\mc=[\mt,\mb^{m+1}]\equiv0\pmod{\mb^m\mt}.
\]
Cela résulte des hypothèses que $\mb$ est un idéal principal et qu'il est premier à $\mt$. Tout élément $c$ de $\mc$ admet, à cause de la divisibilité par $\mb^{m+1}$, une représentation, où $n$ est compris comme symbole d'un entier,
\[
 c=k b^{m+1}+n b^{m+1}=r b^m,
\]
où $r$ désigne maintenant de nouveau un élément de $\mR$. De $c\equiv0(\mt)$ il résulte donc $r b^m\equiv0(\mt)$, et par suite $r\equiv0(\mt)$, puisque $\mt:\mb=\mt$. Mais alors $c\equiv0(\mt\mb^m)$, et le théorème II est démontré.

\textbf{Théorème III.} \emph{Sous l'hypothèse de la condition de chaîne des diviseurs, tout idéal admet une représentation la plus courte comme plus petit commun multiple d'un nombre fini de composantes primaires les plus grandes, appartenant à des idéaux premiers différents.}

Pour obtenir une telle représentation, il suffit de remplacer, chaque fois que c'est possible, la représentation par idéaux irréductibles qui existe d'après le théorème I, donc par idéaux primaires d'après le théorème II, par une représentation la plus courte. Si l'on rassemble les idéaux primaires appartenant aux mêmes idéaux premiers, on obtient d'après \S{}5, 3 la représentation cherchée. Les composantes primaires peuvent être qualifiées de plus grandes, car le plus petit commun multiple de quelconques d'entre elles n'est plus primaire.\footnote{Les idéaux premiers associés et les composantes isolées sont alors déterminés de façon univoque. Voir \emph{Idealtheorie} ou W. Krull, Math. Ann. 90 (1923), p. 55--64.}
\providecommand{\mR}{\mathfrak{R}}
\providecommand{\ma}{\mathfrak{a}}
\providecommand{\mb}{\mathfrak{b}}
\providecommand{\mm}{\mathfrak{m}}
\providecommand{\mo}{\mathfrak{o}}
\providecommand{\mpideal}{\mathfrak{p}}
\providecommand{\mq}{\mathfrak{q}}
\providecommand{\mbar}[1]{\overline{#1}}

\subsection*{\S{}7. Théorie des idéaux sous l'hypothèse de la condition de double chaîne}

Les simplifications par rapport à la théorie développée jusqu'ici reposent sur les lemmes à donner sous 1. et 2.

\textbf{1.} \emph{Si, dans un anneau commutatif sans diviseurs de zéro, la condition de chaîne des multiples est satisfaite, alors l'anneau est déjà un corps.} Il faut montrer que, pour $a\ne0$, l'équation $ax=b$ possède toujours une solution dans l'anneau. Qu'elle ne puisse pas posséder plus d'une solution résulte de l'hypothèse que l'anneau est sans diviseurs de zéro.

Soit $\ma$ l'idéal principal dérivé de $a$. Par hypothèse, la suite des puissances s'interrompt après un nombre fini de termes; soit par exemple $\ma^m$ égal à toutes les puissances suivantes. Si $\mb$ désigne l'idéal principal dérivé de $b$, on obtient donc
\[
\ma^m\mb=\ma^{m+1}\mb,
\qquad \ma^{m+1}\mb=(a^{m+1}b).
\]
Cela donne en particulier pour l'élément $a^m b$ une représentation, où $n$ est compris comme symbole d'un entier,
\[
 a^m b=k a^{m+1}b+n a^{m+1}b=a^{m+1}c,
\]
où $c$ est de nouveau un élément de $\mR$. Comme, par hypothèse, il n'existe pas de diviseurs de zéro, il en résulte $b=ac$, ce qui démontre la propriété d'être un corps.

\textbf{2.} Du point 1, on déduit immédiatement le lemme: \emph{Si, dans un anneau commutatif, la condition de chaîne des multiples est satisfaite, alors un idéal premier ne possède aucun diviseur propre différent de $\mo$.}

On en déduit de même le supplément: \emph{Si seul l'axiome II est satisfait, c'est-à-dire la condition de chaîne des multiples modulo tout idéal différent de l'idéal nul, alors tout idéal premier différent de l'idéal nul ne possède aucun diviseur propre différent de $\mo$.}

En effet, l'anneau quotient par un idéal premier $\mpideal$ devient, par définition, un anneau sans diviseurs de zéro; et comme, par homomorphie, la condition de chaîne des multiples y est également satisfaite, il est, d'après 1., un corps. Par la correspondance biunivoque entre les diviseurs de $\mpideal$ et les idéaux dans l'anneau quotient, $\mpideal$ ne possède donc aucun diviseur propre différent de $\mo$.\footnote{Il n'est pas nécessaire qu'il existe dans $\mR$, et donc dans $\mo$, l'idéal unité composé de tous les éléments de $\mR$, un élément unité, bien que, d'après 1., il existe une classe unité dans l'anneau quotient. L'exemple le plus simple de ce genre --- où il s'agit du cas plus faible du supplément --- est le système de tous les entiers pairs.}

\textbf{Théorème IV.} \emph{Si, dans un anneau commutatif, la condition de double chaîne est satisfaite, alors, dans toute représentation la plus courte d'un idéal, les composantes primaires qui n'appartiennent pas à l'idéal unité $\mo$ sont déterminées de façon univoque.}

\emph{Supplément.} Sous l'hypothèse des axiomes I et II, le théorème IV vaut pour tout idéal différent de l'idéal nul.

Soient
\[
\mm=[\mq,\mq_1,\ldots,\mq_r]
   =[\mbar\mq,\mbar\mq_1,\ldots,\mbar\mq_{\mbar r}]
\]
des représentations les plus courtes de $\mm$ par composantes primaires maximales, et soient
$\mo,\mpideal_1,\ldots,\mpideal_r$, respectivement
$\mo,\mbar\mpideal_1,\ldots,\mbar\mpideal_r$, les idéaux premiers associés. On omet $\mq$, respectivement $\mbar\mq$, lorsqu'aucun idéal primaire associé à $\mo$ n'apparaît. Comme
\[
\mo^\varrho\mpideal_1^{\varrho_1}\ldots\mpideal_r^{\varrho_r}
   \equiv0\pmod{\mm},
\]
chaque $\mbar\mpideal_i$ doit coïncider avec l'un des $\mpideal_i$ d'après le lemme. Réciproquement, chaque $\mpideal_k$ coïncide avec l'un des $\mbar\mpideal_k$; les idéaux premiers distincts de $\mo$ coïncident donc, après réindexation, $\mpideal_\lambda=\mbar\mpideal_\lambda$. En outre,
$\mq\mq_1\ldots\mq_r\equiv0(\mbar\mq_\lambda)$, d'où
$\mq_\lambda\equiv0(\mbar\mq_\lambda)$, puisque les autres composantes ne sont pas divisibles par $\mpideal_\lambda$ et sont donc premières à $\mbar\mq_\lambda$. De même,
$\mbar\mq_\lambda\equiv0(\mq_\lambda)$; l'unicité des composantes primaires non associées à $\mo$ est ainsi établie. Enfin, si $\mq$ apparaît effectivement, $\mbar\mq$ doit également apparaître, puisque les représentations sont les plus courtes, ce qui achève aussi l'unicité des idéaux premiers associés. La démonstration montre encore que toute représentation ne contenant aucune composante primaire associée à $\mo$ est déjà la plus courte.

\textbf{3.} Si l'on ajoute à l'hypothèse de la condition de double chaîne l'existence de l'élément unité, le théorème IV se renforce en théorème V.

\textbf{Théorème V.} \emph{Si, dans un anneau commutatif avec élément unité, la condition de double chaîne est satisfaite, alors tout idéal se représente de façon univoque comme produit d'un nombre fini d'idéaux primaires deux à deux premiers entre eux.}

\emph{Supplément.} Sous l'hypothèse des axiomes I, II, III, le théorème V vaut pour tout idéal différent de l'idéal nul.

À cause de l'existence de l'élément unité, on a $\mo^2=\mo$, et ainsi tout idéal primaire appartenant à $\mo$ devient égal à $\mo$; il ne peut donc pas apparaître dans une représentation la plus courte d'un idéal différent de $\mo$. Par conséquent, d'après le théorème IV, les composantes primaires sont déterminées de façon univoque. En même temps, toute représentation devient la plus courte par suppression de $\mo$. Deux idéaux premiers différents sont, d'après le lemme sous 2., toujours premiers entre eux; d'après les règles de calcul du §~4, 4., il en va donc de même des composantes primaires, et le plus petit commun multiple devient le produit.

Les hypothèses donnent encore le résultat suivant: \textbf{Lemme.} \emph{Dans un anneau commutatif avec élément unité et condition de double chaîne, les idéaux premiers qui sont premiers à un idéal $\mm$ sont, à l'exception de l'idéal unité, exactement ceux qui ne divisent pas $\mm$. Les notions de premier à et de premiers entre eux coïncident.}

\emph{Supplément.} Sous l'hypothèse des axiomes I, II et III, $\mm$ doit être supposé distinct de l'idéal nul.

Si $\mpideal$ ne divise pas $\mm$, il est distinct de tous les idéaux premiers associés à $\mm$; d'après le lemme de 2, il n'est divisible par aucun de ces idéaux premiers distincts de $\mo$. Il est donc premier à chaque composante primaire et à $\mm$; il est en même temps premier entre eux avec chacune de ces composantes et avec $\mm$.

Si $\mpideal\ne\mo$ divise $\mm$, il coïncide avec l'un des idéaux premiers associés. S'il appartient, par exemple, à la composante $\mq=\mq_1$, alors $\mq:\mpideal$ est un diviseur propre de $\mq$, car
$\mpideal^{\varrho-1}\equiv0\pmod{\mq:\mpideal}$ mais
$\mpideal^{\varrho-1}\not\equiv0\pmod{\mq}$. Comme diviseur de
$\mpideal^{\varrho-1}$, $\mq:\mpideal$ n'est divisible que par un idéal premier distinct de $\mo$ et est donc primaire. Par l'unicité de la décomposition,
$(\mq:\mpideal)\mq_2\ldots\mq_r$, qui divise $\mm:\mpideal$, est alors un diviseur propre de $\mm$; ainsi $\mpideal$ n'est pas premier à $\mm$.

Par conséquent, un idéal $\mt\ne\mo$ est premier à $\mm$ si et seulement si aucun des idéaux premiers associés à $\mt$ ne divise $\mm$; dans ce cas, $\mt$ est aussi premier entre eux avec $\mm$. Réciproquement, deux idéaux premiers entre eux sont toujours mutuellement premiers (\S{}5, 5); sous les hypothèses ci-dessus, les deux notions coïncident.
\providecommand{\mR}{\mathfrak{R}}
\providecommand{\ma}{\mathfrak{a}}
\providecommand{\mb}{\mathfrak{b}}
\providecommand{\mc}{\mathfrak{c}}
\providecommand{\mm}{\mathfrak{m}}
\providecommand{\mn}{\mathfrak{n}}
\providecommand{\mo}{\mathfrak{o}}
\providecommand{\mpideal}{\mathfrak{p}}
\providecommand{\mq}{\mathfrak{q}}
\providecommand{\mt}{\mathfrak{t}}

\subsection*{\S{}8. Théorie des idéaux sous l'hypothèse de la clôture intégrale dans le corps des fractions}

La théorie des idéaux développée jusqu'ici doit maintenant, par adjonction des deux derniers axiomes, être renforcée jusqu'à la théorie usuelle.

\textbf{1.} \emph{Lemme.} Soit $\mR$ un anneau commutatif sans diviseurs de zéro avec élément unité, et supposons dans $\mR$ la condition de chaîne des diviseurs (axiomes I, III, IV); alors de $\ma=\ma\mb$ et $\ma\ne(0)$ il résulte toujours $\mb=\mo$. On a donc
\[
        \mc^\varrho\ne\mc^{\varrho+1}
\]
pour tous les idéaux différents de l'idéal nul et de l'idéal unité.\footnote{En revanche, sous les mêmes hypothèses, on ne peut pas conclure de $\ma\mb=\ma\mc$ à $\mb=\mc$. Comme exemple, on peut prendre dans l'anneau de polynômes $K[x,y]$: $\ma=(x,y)$, $\mb=(x^2,xy,y^2)$, $\mc=(x^2,y^2)$; effectivement $\ma\mb=\ma\mc=\ma^3$, mais $\mb\ne\mc$.} La démonstration se fait comme pour le lemme \S{}1, 3, puisqu'on peut considérer $\ma\mb$ comme un module fini relativement à $\mb$. Si $\alpha_1,\ldots,\alpha_n$ désigne une base d'idéal de $\ma$ existant en vertu de I, alors $\ma=\ma\mb$ donne le système d'équations
\[
        \alpha_i=b_{i1}\alpha_1+\ldots+b_{in}\alpha_n,
        \qquad b_{ik}\equiv0\pmod{\mb},\quad i=1,\ldots,n.
\]
Comme $\mR$ est supposé sans diviseurs de zéro, il en résulte $|b_{ik}-e_{ik}|=0$, où $e_{ik}=0$ pour $i\ne k$ et $e_{ii}=e$. De l'équation
\[
        e^n+b_1e^{n-1}+\ldots+b_n=0,
        \qquad b_i\equiv0\pmod{\mb},
\]
On a alors $e\equiv0\pmod{\mb}$, et donc $\mb=\mo$.

\textbf{2.} Il reste maintenant seulement à montrer que, par adjonction de l'hypothèse de la clôture intégrale dans le corps des fractions (axiome V) aux axiomes I à IV, il s'ensuit que tout idéal primaire différent de l'idéal nul devient une puissance de son idéal premier associé. L'unicité de la représentation qui en résulte,
\[
        \mm=\mpideal_1^{\varrho_1}\ldots\mpideal_r^{\varrho_r},
\]
pour tous les idéaux différents de l'idéal nul, est déjà démontrée par le théorème V et par le lemme sous 1.; en même temps il est montré que ces idéaux premiers ne possèdent aucun diviseur propre différent de l'idéal unité. La démonstration doit être menée pour l'exposant deux en recourant à la conséquence II de Dedekind (\S{}1, 4), et, en général, par induction complète.

\emph{Lemme.} \emph{Sous l'hypothèse des axiomes I à V, il n'existe pas d'autres idéaux primaires d'exposant deux que $\mq=\mpideal^2$.}

Pour la démonstration, il faut montrer que de
\[
        \mpideal^2\equiv0\pmod{\mq},
        \qquad \mq\equiv0\pmod{\mpideal},
        \qquad \mq\not\equiv0\pmod{\mpideal^2}
\]
on déduit nécessairement $\mq=\mpideal$. Soit donc
\[
        c\equiv0\pmod{\mq},
        \qquad\text{donc }c\equiv0\pmod{\mpideal},
        \qquad\text{mais }c\not\equiv0\pmod{\mpideal^2}.
\]
Le lemme \S{}7, 3 montre alors que $\mo c:\mpideal$ devient un diviseur propre de $\mo c$. Il existe donc un élément $b$ tel que
\[
        b\mpideal\equiv0\pmod{\mo c}\equiv0\pmod{\mq},
        \qquad b\not\equiv0\pmod{\mo c}.
\]
Ainsi $\gamma=b/c$ n'est pas entier, c'est-à-dire que $\gamma$ appartient au corps des fractions, mais non à $\mR$. Il faut démontrer $b\not\equiv0(\mpideal)$; il en résulte, puisque $\mq$ est primaire et appartient à $\mpideal$, que $\mpideal\equiv0(\mq)$.

D'après la conséquence II de Dedekind, il existe des éléments $m,n$ de $\mR$ tels que
\[
        \gamma=b/c=m/n
\]
et que $m^2/n$ ne soit pas entier, donc
\[
        bn=mc,
        \qquad m\not\equiv0(\mo n),
        \qquad m^2\not\equiv0(\mo n).
\]
En multipliant $b\mpideal$ par $m$, on obtient
\[
        mb\mpideal\equiv0\pmod{\mo mc},
        \qquad\text{donc}\qquad mb\mpideal\equiv0\pmod{\mo nb},
        \qquad\text{et par suite}\qquad m\mpideal\equiv0\pmod{\mo n},
\]
ce dernier fait parce qu'il s'agit d'un anneau sans diviseurs de zéro. Il en résulte
\[
        m\not\equiv0\pmod{\mpideal},
\]
à cause de $m^2\not\equiv0\pmod{\mo n}$ et de $n\equiv0\pmod{\mpideal}$; ce dernier point découle à son tour du lemme \S{}7, 3, puisque $\mo n:\mpideal$ est un diviseur propre de $\mo n$, en raison de $m\not\equiv0(\mo n)$. Les quatre relations
\[
        m\not\equiv0(\mpideal),
        \quad n\equiv0(\mpideal),
        \quad c\equiv0(\mpideal),
        \quad {}\not\equiv0(\mpideal^2),
        \quad bn=mc
\]
donnent alors $b\not\equiv0(\mpideal)$. En effet, comme $\mpideal^2$ est primaire, on obtient d'abord $mc\not\equiv0(\mpideal^2)$; et il s'ensuit $b\not\equiv0(\mpideal)$ à cause de $bn=mc$. Le lemme est ainsi démontré.

\textbf{3.} \emph{Lemme. Sous l'hypothèse des axiomes I à V, il n'existe pas d'autres idéaux primaires d'exposant $\varrho$ que $\mq=\mpideal^\varrho$.}\footnote{La démonstration du lemme 3, sous l'hypothèse de la validité du lemme 2, se trouve chez Masazo Sono, \emph{On Congruences II}, \S{}\S{}9 et 10. Voir la note 19.}

Le lemme découle de celui de 2 sans nouvelle application des axiomes. Il faut d'abord montrer:
\[
\text{si }c\equiv0\pmod{\mpideal},\quad
c\not\equiv0\pmod{\mpideal^2},
\quad\text{alors}\quad
\mpideal^\sigma=(\mo c^\sigma,\mpideal^{\sigma+1})
\quad\text{pour tout }\sigma.
\]
D'après le lemme de 2, $\mpideal=(\mo c,\mpideal^2)$. Si l'on suppose
\[
\mpideal^{\sigma-1}=(\mo c^{\sigma-1},\mpideal^\sigma),
\]
la multiplication par $\mpideal$ et la loi distributive donnent
\[
\mpideal^\sigma
=(\mpideal c^{\sigma-1},\mpideal^{\sigma+1})
=(\mo c^\sigma,\mpideal^2c^{\sigma-1},\mpideal^{\sigma+1})
=(\mo c^\sigma,\mpideal^{\sigma+1}).
\]

Le lemme à démontrer se ramène à la seconde formulation suivante:
\[
\text{si }\mq\equiv0\pmod{\mpideal^\sigma},\quad
\mq\not\equiv0\pmod{\mpideal^{\sigma+1}},
\quad\text{et si}\quad
\mpideal^{\sigma+\lambda}\equiv0\pmod{\mq},
\quad \lambda\geqq1,
\]
alors
\[
\mpideal^{\sigma+\lambda-1}\equiv0\pmod{\mq}.
\]
Comme $\mq\equiv0\pmod{\mpideal}$, il existe pour tout idéal primaire un tel exposant $\sigma\geqq1$; en outre, la non-divisibilité
$\mq\not\equiv0(\mpideal^{\varrho+1})$ découle de
$\mpideal^\varrho\equiv0(\mq)$ et de
$\mpideal^\varrho\ne\mpideal^{\varrho+1}$, par le lemme de 1. En appliquant un nombre fini de fois la seconde formulation, on obtient
$\mpideal^\sigma\equiv0(\mq)$, donc $\mq=\mpideal^\sigma$.

Compte tenu de la représentation de $\mpideal^\sigma$, l'hypothèse de la seconde formulation donne, pour tout $q\in\mq$,
\[
q\equiv bc^\sigma\pmod{\mpideal^{\sigma+1}},
\qquad b\not\equiv0\pmod{\mpideal},
\]
ou encore
\[
bc^\sigma\equiv0\pmod{(\mq,\mpideal^{\sigma+1})}.
\]
Comme $(\mq,\mpideal^{\sigma+1})$ est primaire, il s'ensuit
\[
c^\sigma\equiv0\pmod{(\mq,\mpideal^{\sigma+1})}
\quad\text{ou}\quad
c^{\sigma+\lambda-1}\equiv0\pmod{(\mq,\mpideal^{\sigma+\lambda})},
\]
et par conséquent
\[
\mpideal^{\sigma+\lambda-1}\equiv0\pmod{\mq}.
\]

\textbf{4.} En résumé, on obtient:

\textbf{Théorème VI.} \emph{Si, dans un anneau commutatif, les axiomes I à V sont satisfaits, alors tout idéal différent de l'idéal nul et de l'idéal unité se représente de façon univoque comme produit de puissances d'un nombre fini d'idéaux premiers différents de l'idéal nul et de l'idéal unité; ces idéaux premiers ne possèdent aucun diviseur propre différent de l'idéal unité.}
\providecommand{\mR}{\mathfrak{R}}
\providecommand{\mT}{\mathfrak{T}}
\providecommand{\mK}{\mathfrak{K}}
\providecommand{\ma}{\mathfrak{a}}
\providecommand{\mb}{\mathfrak{b}}
\providecommand{\mc}{\mathfrak{c}}
\providecommand{\mm}{\mathfrak{m}}
\providecommand{\mn}{\mathfrak{n}}
\providecommand{\mo}{\mathfrak{o}}
\providecommand{\mpideal}{\mathfrak{p}}
\providecommand{\mq}{\mathfrak{q}}
\providecommand{\mt}{\mathfrak{t}}
\providecommand{\mX}{\mathfrak{X}}

\subsection*{\S{}9. Les axiomes comme conséquence de la décomposition supposée}

Afin de pouvoir, inversement, déduire ces axiomes de l'existence de la décomposition usuelle des idéaux, qui d'après le théorème VI découle des axiomes I à V, il faut supposer la décomposition sous la forme suivante.

\emph{Hypothèse.} Dans l'anneau commutatif $\mR$, tout idéal qui n'est pas engendré par l'élément nul ou par l'élément unité est représentable de façon univoque comme produit de puissances d'idéaux premiers. Ces idéaux premiers sont des idéaux simples, non engendrés par l'élément unité, c'est-à-dire qu'ils ne possèdent aucun diviseur propre différent de $\mo$; réciproquement, tous les idéaux simples sont des idéaux premiers. L'unicité vaut sous la forme stricte
\[
        \ma^\varrho\ne\ma^{\varrho+1}
\]
pour tout idéal qui n'est pas engendré par l'élément nul ou par l'élément unité.

L'existence de l'élément unité n'est pas supposée ici. S'il n'y a pas d'élément unité, la condition \og{}non engendré par l'élément unité\fg{} n'apporte aucune restriction.

\textbf{1.} \emph{Vérification de l'axiome III.} Supposons que l'axiome III ne soit pas satisfait. Il faut montrer que $\mo^2$ devient un idéal simple sans être un idéal premier. D'après l'hypothèse d'unicité stricte, on a $\mo\ne\mo^2$; donc $\mo\not\equiv0(\mo^2)$, mais $\mo^2\equiv0(\mo^2)$, et par suite $\mo^2$ n'est pas un idéal premier. Mais $\mo^2$ est simple; car $\mo^2$ ne peut être divisible par aucun idéal premier différent de $\mo$, et ne possède donc, en vertu de la représentation par produits supposée, aucun diviseur différent de $\mo$ ou de $\mo^2$.

\textbf{2.} \emph{Vérification de l'axiome IV (anneau sans diviseurs de zéro).} Si $a$ est un diviseur de zéro, donc $ab=0$ avec $a\ne0$ et $b\ne0$, la multiplication des décompositions des idéaux principaux engendrés par $a$ et $b$ donne
\[
(0)=\mpideal_1^{\varrho_1}\ldots\mpideal_k^{\varrho_k},
\]
où les $\mpideal$ sont distincts de l'idéal nul et de l'idéal unité --- cette dernière exclusion résultant de l'existence de l'élément unité. Supposons cette décomposition la plus courte, au sens où la suppression d'un facteur quelconque donne un diviseur propre de l'idéal nul --- condition qui n'est pas automatiquement satisfaite, puisqu'aucune hypothèse d'unicité de la décomposition n'a été faite pour l'idéal nul. Si un seul idéal premier intervient, alors
\[
(0)=\mpideal^\varrho=\mpideal^{\varrho+1},
\]
contrairement à l'unicité stricte. Si plusieurs idéaux premiers interviennent, on a
\[
\mpideal_1^{\varrho_1+1}
=(\mpideal_1^{\varrho_1+1},
  \mpideal_1^{\varrho_1}\ldots\mpideal_k^{\varrho_k})
=\mpideal_1^{\varrho_1}
 (\mpideal_1,\mpideal_2^{\varrho_2}\ldots\mpideal_k^{\varrho_k})
=\mpideal_1^{\varrho_1},
\]
car, grâce à l'existence de l'élément unité, $\mpideal_1$ est premier entre eux avec tous les autres idéaux premiers. On obtient encore une contradiction avec l'unicité stricte.

\textbf{3.} \emph{Vérification des axiomes I et II (conditions de chaîne).} Les hypothèses donnent, pour tout idéal distinct de l'idéal nul: \emph{de la divisibilité résulte une décomposition en produit}; c'est-à-dire que de
$\mm\equiv0(\ma)$ et $\ma\ne\mo$ résulte
$\mm=\ma\mb$ avec $\mb\not\equiv0\pmod{\mm}$. Soient
\[
\mm\equiv0\pmod{\ma},\qquad
\mm=\mpideal_1^{\varrho_1}\ldots\mpideal_r^{\varrho_r},\qquad
\ma=\overline{\mpideal}_1^{\sigma_1}\ldots
    \overline{\mpideal}_{\bar r}^{\sigma_{\bar r}}.
\]
Chaque $\overline{\mpideal}_i$ est alors identique à un $\mpideal_j$, et
$\sigma_i\leqq\varrho_i$. On obtient donc la décomposition cherchée avec
\[
\mb=\mpideal_1^{\varrho_1-\sigma_1}\ldots
    \mpideal_r^{\varrho_r-\sigma_r},
\]
certains $\sigma_i$ pouvant être nuls. Il n'existe ainsi qu'un nombre fini de diviseurs de $\mm$, correspondant aux combinaisons
$0\leqq\sigma_i\leqq\varrho_i$; l'anneau quotient par $\mm$ satisfait donc la condition de double chaîne. Les axiomes I et II sont établis; la condition de chaîne des diviseurs vaut également dans $\mR$, puisque tout diviseur propre de l'idéal nul est distinct de l'idéal nul.

\emph{La vérification de l'axiome V de la clôture intégrale dans le corps des fractions suppose les conséquences usuelles de la décomposition des idéaux; il faut donc commencer par les déduire.}

\textbf{4.} \emph{L'anneau quotient par tout idéal différent de l'idéal nul est un anneau principal.} On peut supposer que l'idéal est différent de l'idéal unité, car, pour l'anneau quotient réduit au seul élément nul, la propriété est évidemment satisfaite. Si l'idéal est primaire, $\mq=\mpideal^\varrho$, et si $c\equiv0(\mpideal)$ mais $\not\equiv0(\mpideal^2)$, le point 3 montre que $\mo c=\mpideal\ma$, avec $(\ma,\mpideal)=\mo$. Il vient donc
\[
        \mpideal^\sigma=(\mo c^\sigma,\mpideal^\varrho)
        \qquad\text{pour tout }\sigma\leqq\varrho,
\]
et, dans l'anneau quotient par un idéal primaire, tout idéal devient ainsi principal. En général, pour
\[
        \mm=\mq_1\cdot\mq_2\ldots\mq_r,
        \qquad \mq_i=\mpideal_i^{\varrho_i},
\]
et pour la représentation correspondante de l'anneau quotient comme somme directe
\[
        \mR/\mm=\overline{\ma}_1+\ldots+\overline{\ma}_r,
\]
$\overline{\ma}_i$ est isomorphe à $\mR/\mq_i$ (\S{}4, 5), et par conséquent tout idéal de $\overline{\ma}_i$ est principal. Si $\overline{\mc}$ est un idéal quelconque de l'anneau quotient, alors $\overline{\ma}_i\overline{\mc}$ devient l'idéal principal $\overline{\ma}_i\overline c_i$; il vient
\[
\overline{\mc}
=\overline{\mo}\,\overline{\mc}
=\overline{\ma}_1\overline{\mc}+\ldots+\overline{\ma}_r\overline{\mc}
=\overline{\ma}_1\overline c_1+\ldots+\overline{\ma}_r\overline c_r
=\overline{\ma}_1\overline c+\ldots+\overline{\ma}_r\overline c
=\overline{\mo}\,\overline c,
\]
où $\overline c=\overline c_1+\ldots+\overline c_r$. En effet,
$\overline{\ma}_i\overline{\ma}_k=(0)$ et
$\overline c_i\equiv0\pmod{\overline{\ma}_i}$, de sorte que les idéaux principaux
$\overline{\ma}_i\overline c_i$ et $\overline{\ma}_i\overline c$ coïncident.

Ce théorème admet encore la forme suivante: si $\mc$ est un idéal quelconque, il peut être transformé en un idéal principal par multiplication avec un idéal copremier à un idéal donné $\mb$. En effet, si l'on pose $\mm=\mb\mc$, alors $\mc$ devient principal modulo $\mm$, c'est-à-dire
\[
        \mc=(\mo c,\mm)=(\mc\ma,\mc\mb)=\mc(\ma,\mb);
\]
ainsi $\mo c=\mc\ma$ et $(\ma,\mb)=\mo$.

\textbf{5.} \emph{Théorie des idéaux fractionnaires.} On appelle idéal fractionnaire tout $\mR$-module fini dans le corps des fractions $\mK$.

\emph{Les $\mR$-modules finis de $\mK$ qui sont différents du module nul forment, pour la multiplication, un groupe abélien.}

Le produit de deux $\mR$-modules finis est de nouveau un $\mR$-module fini; la multiplication est associative et commutative; comme $\mathfrak A=\mo\mathfrak A$, l'idéal unité devient l'élément unité du système. Il reste donc seulement à montrer que l'équation
\[
        \mathfrak A\mX=\mo
\]
possède toujours une solution dans le système des $\mR$-modules finis.

\textbf{Remarque préliminaire.} Si l'équation $\mathfrak A\mT=\mathfrak B$ possède une et une seule solution, alors $\mT$ est égal au quotient de modules $\mathfrak B:\mathfrak A$ (\S{}5, 4). En effet, on a
\[
        \mT\equiv0\pmod{\mathfrak B:\mathfrak A},
        \qquad
        \mathfrak B=\mathfrak A\mT\equiv0\pmod{\mathfrak A\cdot(\mathfrak B:\mathfrak A)}\equiv0\pmod{\mathfrak B},
        \quad\text{ou encore}\quad
        \mathfrak A\cdot(\mathfrak B:\mathfrak A)=\mathfrak B.
\]
Supposons d'abord que $\mathfrak A$ soit principal, $\mathfrak A=(\alpha)$; alors $\mX=(e/\alpha)$ est solution de $\mathfrak A\mX=\mo$, et $\mX$ est de nouveau un $\mR$-module fini. Il en résulte que tout $\mR$-module fini est le quotient de modules de deux idéaux de $\mR$, ce qui justifie la dénomination d'idéal fractionnaire. En effet, soient $\tau_1,\ldots,\tau_r$ une base de $\mT$ comme $\mR$-module, soit $\tau_i=t_i/a$, et soit $\mt$ l'idéal de $\mR$ engendré par les $t_i$. Alors $\mo a\cdot\mT=\mt$; d'après la remarque préliminaire, on en déduit que $\mT=(\mt:\mo a)$, le quotient étant pris dans $\mR$.

On peut maintenant donner en général la solution de $\mathfrak A\mX=\mo$. Posons $\mathfrak A=a:\mo c$, et déterminons $\mb$ de telle sorte que $\ma\mb$ soit égal à un idéal principal $\mo a$. Alors il vient
\[
        \mathfrak A\,\mo c=\ma,
        \qquad \mathfrak A\,\mb\mo c=\mo a,
        \qquad \mathfrak A(\mb\mo c:\mo a)=\mo.
\]
Ici $\mX$, comme quotient d'un idéal par un idéal principal, est de nouveau un $\mR$-module fini. De la propriété de groupe il résulte encore que le quotient de modules de deux idéaux quelconques $\ma:\mb$, comme solution de l'équation $\mb\mX=\ma$, devient un $\mR$-module fini.

\textbf{6.} De la propriété de groupe découlent les conséquences suivantes.

6$\alpha$. On peut simplifier: la relation $\mathfrak T=\mathfrak A\mathfrak M:\mathfrak B\mathfrak M$ entraîne $\mathfrak T=\mathfrak A:\mathfrak B$, et réciproquement. En effet, l'égalité $\mathfrak T\mathfrak B\mathfrak M=\mathfrak T\mathfrak A\mathfrak M$ entraîne $\mathfrak T\mathfrak B=\mathfrak T\mathfrak A$, et réciproquement.

6$\beta$. Tout idéal fractionnaire admet une représentation $\mathfrak C=\ma:\mb$, où $\ma$ et $\mb$ sont premiers entre eux; cette condition détermine $\ma$ et $\mb$ de façon univoque. L'existence d'une telle représentation résulte de 5 et de 6$\alpha$; mais l'égalité $\mathfrak C=\ma:\mb=\overline{\ma}:\overline{\mb}$ entraîne $\ma\overline{\mb}=\mb\overline{\ma}$, et par suite l'unicité, si $\ma$ et $\mb$, de même que $\overline{\ma}$ et $\overline{\mb}$, sont supposés premiers entre eux.

6$\gamma$. Tout module principal $\mo\eta$ engendré par un élément non entier admet une représentation comme quotient d'idéaux principaux telle qu'aucune puissance du numérateur ne soit divisible par le dénominateur. Soit, dans une représentation réduite, $\mo\eta=\mb:\mc$, où $(\mb,\mc)=\mo$; soit $\mo b=\mb\ma$ et $(\ma,\mc)=\mo$. Alors
\[
        \mo\eta=\ma\mb:\ma\mc=\mo b:\ma\mc,
\]
et, comme $\ma\mc=\mo b:\mo\eta$ d'après le point 4, l'idéal $\ma\mc$ est lui aussi principal,
\[
        \mo\eta=\mo b:\mo c.
\]
Mais aucune puissance de $\ma\mb$ ne devient divisible par $\mc$, puisque $(\ma\,\mc)=\mo$ et $(\mb,\mc)=\mo$.

\textbf{7.} \emph{Vérification de l'axiome V de la clôture intégrale dans le corps des fractions.} Le point 6$\gamma$ montre que tout élément non entier de $\mK$ admet une représentation quotient
\[
        \eta=a/c
\]
telle que, dans $\mR$, aucune puissance du numérateur ne soit divisible par le dénominateur. En effet, l'égalité $\mo\eta=\mo b:\mo c$ montre que $\eta$ ne diffère du quotient $b/c$ que par un facteur unité de $\mR$. Un tel élément ne peut donc satisfaire aucune équation qui le caractérise comme entier relativement à $\mR$; car de
\[
        \eta^n+r_1\eta^{n-1}+\ldots+r_n=0
\]
On obtient $a^n\equiv0\pmod{\mo c}$ dans $\mR$.

\textbf{8.} \emph{Conséquence.} Si, dans un anneau commutatif $\mR$, les axiomes I à IV sont satisfaits et si tout idéal primaire est irréductible, alors l'axiome V de la clôture intégrale dans le corps des fractions est également satisfait. En vertu des axiomes I à III, toute puissance d'idéal premier $\mpideal^\varrho$ est en même temps un idéal primaire; cela vaut en particulier pour $\mpideal^2$, qui devient irréductible par hypothèse. Il faut en déduire
\[
        \mpideal=(\mo c,\mpideal^2);
\]
car le lemme \S{}8, 3 montre alors que tout idéal primaire devient une puissance d'idéal premier, ce qui, avec les autres hypothèses et d'après 7., entraîne la clôture intégrale.

À cause de la condition de chaîne des diviseurs, il vient
\[
        \mpideal=(\mo c_1,\ldots,\mo c_k,\mpideal^2),
\]
où seules les classes modulo $\mpideal$ entrent en considération comme multiplicateurs des $c$, et où les $c$ peuvent être supposés linéairement indépendants sur le corps résiduel $\mR/\mpideal$. Mais, pour $k>1$, cette indépendance linéaire entraîne une représentation
\[
        \mpideal^2=[\mq_1,\mq_2],
        \qquad \mq_1=(\mo c_1,\mpideal^2),
        \qquad \mq_2=(\mo c_2,\ldots,\mo c_k,\mpideal^2),
\]
de sorte que $\mq_1$ et $\mq_2$ deviennent des diviseurs propres de $\mpideal^2$. Cela contredit l'irréductibilité. Ainsi $\mpideal=(\mo c,\mpideal^2)$ est démontré.

Réciproquement, il est immédiat que, sous les axiomes I à V, tout idéal primaire est irréductible, puisqu'une puissance d'idéal premier $\mpideal^\varrho$ ne peut être le plus petit commun multiple de diviseurs propres de la forme $\mpideal^\sigma$ et $\mpideal^\tau$.

\textbf{9.} \emph{Si des diviseurs de zéro sont admis dans l'anneau, alors les axiomes I à III et la clôture intégrale dans l'anneau total des fractions n'impliquent pas que tout idéal primaire devienne irréductible.}\footnote{Voir la note 19.}

Soit $\mT$ un anneau dans lequel tous les axiomes I à IV, sauf la clôture intégrale, sont satisfaits. De tels anneaux existent: par exemple, lorsque $\mR$ satisfait aux axiomes I à V (\S{}3, 2), les ordres finis d'une extension finie du corps des fractions de $\mR$ qui sont distincts du système de tous les éléments entiers sur $\mR$ sont de tels anneaux. D'après la conséquence 8, il existe dans $\mT$ des idéaux primaires réductibles $\mq$; soit $\mpideal^\sigma$ un multiple propre de $\mq$, et soit $\overline{\mR}=\mT/\mpideal^\varrho$ l'anneau quotient. L'idéal $\overline{\mq}$ qui y correspond à $\mq$ est alors un idéal primaire réductible non nul. Dans $\overline{\mR}$, les axiomes I à III sont satisfaits, et la clôture intégrale dans l'anneau total des fractions vaut également. En effet, tout élément de $\mT$ non divisible par $\mpideal$ est copremier à $\mpideal^\sigma$; tous les éléments réguliers de $\overline{\mR}$ sont donc des unités. Ainsi $\overline{\mR}$ est identique à son anneau total des fractions.
\providecommand{\mR}{\mathfrak{R}}
\providecommand{\mA}{\mathfrak{A}}
\providecommand{\mB}{\mathfrak{B}}
\providecommand{\mE}{\mathfrak{E}}
\providecommand{\mG}{\mathfrak{G}}
\providecommand{\mH}{\mathfrak{H}}
\providecommand{\mM}{\mathfrak{M}}
\providecommand{\mU}{\mathfrak{U}}
\providecommand{\mC}{\mathfrak{C}}

\subsection*{\S{}10. Condition de double chaîne et suite de composition}

On va montrer que, pour tout domaine de modules supposé bien ordonné (\S{}2), \emph{l'hypothèse de validité de la condition de double chaîne} --- toute chaîne de diviseurs et toute chaîne de multiples de modules s'interrompt au bout d'un nombre fini de termes --- est équivalente à \emph{l'hypothèse de l'existence d'une suite de composition}.\footnote{Les modules du domaine sont des groupes abéliens pour l'addition; plus précisément, il s'agit de groupes abéliens \emph{généralisés}, puisque le domaine des multiplicateurs se compose des éléments de $\mR$. Comme les groupes sont abéliens, la notion de suite de composition coïncide avec celle de suite principale. Les théorèmes de ce paragraphe subsistent si l'on entend par \og{}modules\fg{} la totalité des sous-groupes normaux d'un groupe quelconque, non commutatif et même généralisé. La terminologie, légèrement différente de celle employée plus haut, doit rappeler la théorie des groupes.} La spécialisation du domaine de modules à un anneau commutatif donne les faits correspondants pour le système de tous les idéaux de l'anneau; dans le numéro 2, il faut alors remplacer partout l'isomorphisme de modules par l'isomorphisme d'anneaux.

Un domaine de modules $\mG$ est dit \emph{simple} s'il ne contient aucun $\mR$-module autre que $\mG$ lui-même et le module nul $\mE$. Un module $\mA$ est dit \emph{simple dans $\mG$} si le module quotient $\mG/\mA$ est simple (plus grand sous-groupe normal).

\emph{Une chaîne de multiples}
\[
        \mG,\mA_1,\ldots,\mA_r,\mE
\]
\emph{est appelée suite de composition de $\mG$ de longueur $r$} si tous les modules de la chaîne sont distincts et si chacun est \emph{simple dans le précédent}; autrement dit, si les modules quotients (groupes quotients) $\mA_{i-1}|\mA_i$, ainsi que $\mG|\mA_1$ et $\mA_r$, sont des domaines de modules simples. En formant le module quotient $\mG|\mA$, le cas où il existe une suite de composition de $\mG$ jusqu'à $\mA\ne\mE$ se ramène au cas d'une suite de composition au sens strict.

\textbf{1.} \emph{Si la condition de double chaîne vaut dans $\mG$, il existe dans $\mG$ une suite de composition}; on suppose ici $\mG\ne\mE$. L'ordre bien ordonné de $\mG$ et la condition de chaîne des diviseurs donnent, comme au \S{}6, par un ordre quasi lexicographique, un bon ordre du système de tous les modules. En vertu de ce bon ordre, la condition de chaîne des diviseurs entraîne l'existence d'au moins un module simple dans $\mG$. En effet, si $\mE_1$ est le premier diviseur propre de $\mE$ dans le bon ordre et, plus généralement, si $\mE_i$ est le premier diviseur propre de $\mE_{i-1}$, la chaîne bien déterminée
\[
        \mG,\mG_1,\ldots,\mG_i,\ldots
\]
s'interrompt au bout d'un nombre fini de termes et conduit donc nécessairement à un module $\mA_1$ simple dans $\mG$. Si $\mA_1\ne\mE$, le même procédé fournit un module $\mA_2$ simple dans $\mA_1$; en général, si $\mA_{i-1}\ne\mE$, il existe un module $\mA_i$ simple dans $\mA_{i-1}$. On obtient ainsi une chaîne de multiples bien déterminée
\[
        \mG,\mA_1,\ldots,\mA_i,\ldots,
\]
qui s'interrompt au bout d'un nombre fini de termes en vertu de la condition de chaîne des multiples et fournit donc une suite de composition de $\mG$.

\textbf{2.} \emph{Si l'on suppose qu'il existe une suite de composition dans $\mG$, alors la condition de double chaîne vaut dans $\mG$.} La démonstration procède par induction complète, sans qu'il soit nécessaire de supposer $\mG$ bien ordonné. Le passage inductif donne en même temps une démonstration du théorème de Jordan--Hölder sous la seule hypothèse de l'existence d'une suite de composition.\footnote{Voir Masazo Sono (note 19), \emph{On Congruences I}, \S{}\S{}11--13, pour le cas des anneaux. Il y traite aussi l'analogue de la suite de composition dans les groupes non commutatifs, à savoir des suites $\mR,\mA_1,\ldots,\mA_r,\mE$ dans lesquelles $\mA_i$ est chaque fois un idéal simple dans $\mA_{i-1}$, sans devoir être un idéal de $\mR$. En fait, tous les raisonnements ci-dessus subsistent pour les groupes non commutatifs si l'on entend par chaîne de multiples une suite $\mG,\mA_1,\ldots,\mA_i,\ldots$ où $\mA_i$ est chaque fois normal dans $\mA_{i-1}$, et si l'on définit les chaînes de diviseurs de manière correspondante. Voir en outre Dedekind, \emph{Über die von drei Moduln erzeugte Dualgruppe}, Math. Ann. 53 (1900), p. 371--403. Pour des domaines beaucoup plus généraux (groupes de modules), le théorème de Jordan--Hölder y est également démontré sous la seule hypothèse d'existence (\S{}6, XVI). Le second théorème d'isomorphisme y est remplacé par une relation de correspondance un peu plus faible; l'énoncé du théorème est donc lui aussi un peu plus faible, mais la démonstration est en principe identique à celle donnée ici.}

Supposons d'abord $\mG$ simple, donc $r=0$, de sorte qu'il n'existe que l'unique suite de composition $\mG,\mE$. Les assertions suivantes sont alors satisfaites:
\[
\begin{array}{ll}
\alpha) & \text{par tout module simple dans $\mG$ passe une suite de composition;}\\
\beta)  & \text{théorème de Jordan--Hölder: toutes les suites de composition ont la même}\\
         & \text{longueur et leurs systèmes de groupes quotients (modules quotients) coïncident}\\
         & \text{à l'ordre près, c'est-à-dire que les groupes quotients correspondants sont isomorphes;}\\
\gamma) & \text{par tout module quelconque différent de $\mG$ passe une suite de composition;}\\
\delta) & \text{la condition de chaîne des multiples vaut;}\\
\varepsilon) & \text{la condition de chaîne des diviseurs vaut.}
\end{array}
\]
On peut donc supposer les assertions $\alpha)$ à $\varepsilon)$ démontrées pour tout domaine de modules possédant une suite de composition de longueur $\bar r<r+1$. Il faut les établir sous \emph{l'hypothèse qu'il existe dans $\mG$ une suite de composition}
\[
        \mG,\mA,\mA_1,\ldots,\mA_r,\mE
\]
\emph{de longueur $r+1$.}

\textbf{2$\alpha$.} S'il n'existe aucun module simple dans $\mG$ distinct de $\mA$, il n'y a rien à démontrer. Soit donc $\mB$ un module simple dans $\mG$, avec $\ne\mA$; il faut montrer qu'il existe dans $\mG$ une suite de composition passant par $\mB$. Comme $\mA$ et $\mB$ sont tous deux simples dans $\mG$ et distincts, on a nécessairement $(\mA,\mB)=\mG$. Posons encore $\mM=[\mA,\mB]$; le second théorème d'isomorphisme (\S{}4, 2) donne
\[
        \mG|\mB\simeq\mA|\mM,
        \qquad
        \mG|\mA\simeq\mB|\mM.
\]
Ainsi $\mM$ est simple dans $\mA$ et dans $\mB$. Comme $\mA$ possède une suite de composition de longueur $r<r+1$, les hypothèses $\alpha)$ et $\beta)$ donnent dans $\mA$ une suite de composition de longueur $r$ passant par $\mM$, par exemple
\[
        \mA,\mM,\mM_1,\ldots,\mM_{r-1},\mE;
\]
dès lors
\[
        \mG,\mB,\mM,\mM_1,\ldots,\mM_{r-1},\mE
\]
est une suite de composition de $\mG$ passant par $\mB$.

\textbf{2$\beta$.} Pour démontrer le théorème de Jordan--Hölder, considérons les deux suites de composition de $\mG$
\[
 (1)\quad \mG,\mA,\mA_1,\ldots,\mA_r,\mE
 \qquad\text{et}\qquad
 (2)\quad \mG,\mB,\mB_1,\ldots,\mB_s,\mE,
 \quad s\geqq r.
\]
Si $\mA=\mB$, l'hypothèse valable pour $r<r+1$ règle tout. Sinon, la comparaison avec les suites construites sous 2$\alpha$,
\[
 (3)\quad \mG,\mA,\mM,\mM_1,\ldots,\mM_{r-1},\mE
 \qquad\text{et}\qquad
 (4)\quad \mG,\mB,\mM,\mM_1,\ldots,\mM_{r-1},\mE,
\]
donne ce qui suit. En vertu de l'hypothèse pour $r<r+1$, les systèmes de groupes quotients de (1) et (3) coïncident; il en va de même pour (2) et (4), puisque la suite (4), passant par $\mB$, est de longueur $r$. On a donc $s=r$. L'isomorphisme établi sous 2$\alpha$ donne en outre la concordance de (3) et (4), d'où celle de (1) et (2), et donc le théorème de Jordan--Hölder.

\textbf{2$\gamma$.} Remarque préliminaire. \emph{Soient $\mA,\mB,\mH$ des modules de $\mG$, et supposons $\mB$ simple dans $\mA$. Alors les modules $(\mB,\mH)$ et $(\mA,\mH)$ coïncident, ou bien $(\mB,\mH)$ est simple dans $(\mA,\mH)$.} Le second théorème d'isomorphisme (\S{}4, 2) donne
\[
        (\mA,\mB,\mH)/(\mB,\mH)
        \simeq \mA/[\mA,(\mB,\mH)],
\]
ou, puisque $\mB\equiv0\pmod{\mA}$,
\[
        (\mA,\mH)/(\mB,\mH)\simeq \mA/\mC,
        \qquad
        \mC=[\mA,(\mB,\mH)].
\]
On a $\mC\equiv0(\mA)$; d'autre part $\mB\equiv0(\mC)$, puisque $\mB\equiv0(\mA)$ et $\mB\equiv0((\mB,\mH))$. Ainsi $\mC$ est situé entre $\mA$ et $\mB$ et coïncide, par hypothèse, avec $\mA$ ou avec $\mB$; la remarque préliminaire en résulte.

Reprenons la suite de composition $\mG,\mA,\mA_1,\ldots,\mA_r,\mE$ de $\mG$, et soit $\mH$ un module quelconque de $\mG$, distinct de $\mG$. Pour montrer qu'une suite de composition passe par $\mH$, formons la suite de modules
\[
        \mG,(\mA,\mH),(\mA_1,\mH),\ldots,(\mA_r,\mH),\mH.
\]
D'après la remarque préliminaire, les modules distincts de cette suite, $\mG,\mC,\mC_1,\ldots,\mC_s,\mH$, forment une suite de composition de $\mG$ à $\mH$; ainsi $\mC$ --- qui, dans le cas particulier 2$\alpha$, coïncide avec $\mH$ --- est simple dans $\mG$. Il existe donc, d'après 2$\alpha$, une suite de composition de longueur $r<r+1$ dans $\mC$; par l'hypothèse d'induction, on peut y faire passer une suite de composition par $\mH$. En lui ajoutant $\mG$, on obtient dans $\mG$ une suite de composition passant par $\mH$. L'application répétée de ce raisonnement montre encore qu'une telle suite peut être choisie comme prolongement de celle construite de $\mG$ à $\mH$.

\textbf{2$\delta$.} Démonstration de la condition de chaîne des multiples. Supposons encore qu'il existe dans $\mG$ une suite de composition de longueur $r+1$, et soit
\[
        \mG,\mB,\mB_1,\ldots,\mB_\nu,\ldots
\]
une chaîne de multiples, chaque module étant un multiple propre du précédent. Supposer que la chaîne commence par $\mG$ ne restreint pas la généralité, puisqu'on peut toujours ajouter $\mG$ comme premier terme. D'après 2$\gamma$, il existe dans $\mG$ une suite de composition passant par $\mB$; ainsi $\mB$ possède une suite de composition de longueur plus courte. Par l'hypothèse d'induction, la chaîne $\mB,\mB_1,\ldots,\mB_\nu,\ldots$ s'interrompt au bout d'un nombre fini de termes; il en va donc de même après adjonction de $\mG$.

\textbf{2$\varepsilon$.} Démonstration de la condition de chaîne des diviseurs. Supposons encore qu'il existe dans $\mG$ une suite de composition de longueur $r+1$, et soit
\[
        \mE,\mC,\mC_1,\ldots,\mC_\nu,\ldots
\]
une chaîne de diviseurs, chaque module étant un diviseur propre du précédent. Supposer que la chaîne commence par $\mE$ ne restreint pas davantage la généralité. D'après 2$\gamma$, il existe dans $\mG$ une suite de composition passant par $\mC$; il existe donc dans $\mG/\mC$ une suite de composition de longueur plus courte. Par l'hypothèse d'induction, la chaîne de diviseurs
\[
        \mC/\mC,\mC_1/\mC,\ldots,\mC_\nu/\mC,\ldots
\]
s'interrompt au bout d'un nombre fini de termes dans $\mG/\mC$. En vertu de la correspondance bijective donnée par le premier théorème d'isomorphisme (\S{}4, 2), il en va de même de la chaîne initiale commençant par $\mC$, et donc aussi de celle que l'on prolonge en lui ajoutant $\mE$.

\emph{L'équivalence des deux hypothèses --- existence d'une suite de composition ou validité de la condition de double chaîne --- est ainsi démontrée.}

\emph{Reçu le 13. 8. 1925.}
\clearpage
\setcounter{footnote}{0}

\providecommand{\ma}{\mathfrak a}
\providecommand{\mb}{\mathfrak b}
\providecommand{\mc}{\mathfrak c}
\providecommand{\mo}{\mathfrak o}
\providecommand{\mO}{\mathfrak O}
\providecommand{\mP}{\mathfrak P}
\providecommand{\mR}{\mathfrak R}
\providecommand{\mT}{\mathfrak T}
\providecommand{\mX}{\mathfrak X}
\providecommand{\mZ}{\mathfrak Z}
\providecommand{\mq}{\mathfrak q}
\providecommand{\mt}{\mathfrak t}
\providecommand{\mpideal}{\mathfrak p}
\providecommand{\ba}{\overline{\mathfrak a}}
\providecommand{\bp}{\overline{\mathfrak p}}
\providecommand{\bP}{\overline P}
\providecommand{\bR}{\overline{\mathfrak R}}

\section*{\NoetherNavAnchor{P31}{31. Le théorème du discriminant pour les ordres d’un corps algébrique de nombres ou de fonctions}31. Le théorème du discriminant pour les ordres d'un corps algébrique de nombres ou de fonctions}

\begin{center}
Journal f. d. reine u. angew. Math. 157 (1927), p. 82--104

Par \emph{Emmy Noether}, à Göttingen.\\[0.4em]
\rule{9em}{0.4pt}
\end{center}

Le théorème de Dedekind sur le discriminant affirme, comme on sait: un nombre premier entre dans le discriminant du corps si et seulement s'il est divisible au moins par le carré d'un idéal premier.

Je montre dans ce qui suit que ce théorème reste valable pour tous les ordres finis d'un corps algébrique fini de nombres, c'est-à-dire pour tous les anneaux de nombres entiers qui contiennent l'anneau des entiers rationnels, sous la forme suivante: un nombre premier $p$ entre dans le discriminant d'un ordre si et seulement si, dans la décomposition de l'idéal principal $(p)$ en plus grandes composantes primaires, valable dans l'ordre,\footnote{Cf. E. Noether, \emph{Abstrakter Aufbau der Idealtheorie in algebraischen Zahl- und Funktionenkörpern}, Math. Ann. 96 (1926), p. 26--61. Pour toutes les définitions, on renvoie à ce travail, cité ci-dessous comme ``Théorie des idéaux''. Dans ce qui suit, ``ordre'' signifie toujours ``ordre fini'', et ``anneau'' un anneau commutatif.} apparaît au moins un idéal primaire propre, c'est-à-dire un idéal qui n'est pas premier.

Comme, pour l'ordre maximal, le système de tous les entiers du corps, il n'y a pas d'autres idéaux primaires que les puissances d'idéaux premiers, le théorème de Dedekind en résulte par spécialisation.

La démonstration du théorème se ramène, par passage à l'anneau quotient de l'ordre par $(p)$, à la théorie de la décomposition des anneaux qui sont de rang fini relativement à un corps contenu dans l'anneau, donc qui forment un système commutatif de grandeurs hypercomplexes avec unité principale et coefficients dans ce même corps. Ici le corps des coefficients est, en particulier, le corps résiduel des entiers rationnels modulo $p$, et à la décomposition de $(p)$ correspond, dans l'anneau quotient, la décomposition de l'idéal nul.

Si donc on appelle complètement réductible un tel anneau de rang fini lorsque son idéal nul peut être représenté comme plus petit commun multiple d'un nombre fini d'idéaux premiers, il s'agit de critères de complète réductibilité. On obtient d'abord comme tel critère que, lorsque le domaine des coefficients est un corps parfait, l'anneau est complètement réductible si et seulement si l'anneau étendu à domaine des coefficients algébriquement clos est complètement réductible. Pour cet anneau étendu, le non-évanouissement du discriminant se reconnaît très simplement comme critère nécessaire et suffisant de complète réductibilité;\footnote{Dans un cas particulier, Hilbert conclut de manière analogue: Zur Theorie der aus $n$ Haupteinheiten gebildeten komplexen Größen, Gött. Nachr. 1896.} et comme le discriminant de l'ordre, modulo chaque nombre premier, passe dans le discriminant de l'anneau quotient, le théorème du discriminant est ainsi démontré.\footnote{Des recherches apparentées, y compris dans le cas non commutatif, mais limitées à l'ordre maximal (domaine maximal) et supposant des coefficients entiers rationnels, se trouvent dans le travail tout récemment paru de A. Speiser: \emph{Allgemeine Zahlentheorie}. Naturf. Gesellschaft Zürich 71 (1926), p. 8--48. Cf. aussi la littérature américaine qui y est indiquée.}

Si l'on considère des corps de fonctions au lieu de corps de nombres, l'ordre doit maintenant contenir le domaine de base de toutes les fonctions entières rationnelles, ou bien, pour les fonctions algébriques de plusieurs indéterminées ou pour les fonctions algébriques entières, le domaine de base de toutes les fonctionnelles entières à coefficients entiers rationnels. Alors le cas d'un domaine des coefficients imparfait peut aussi se présenter dans l'anneau quotient; par exemple, dans le cas du domaine fonctionnel entier, lorsqu'il s'agit du quotient par un nombre premier. Lorsque le domaine des coefficients est un corps imparfait, il faut distinguer les idéaux premiers de première et de seconde espèce, selon que le corps résiduel associé à l'idéal premier est une extension de première ou de seconde espèce,\footnote{Dans la formulation bien connue de Steinitz, J. f. M. 137. Cf. la note relative au \S 2, no 4, de ce travail.} et, de façon correspondante, la complète réductibilité de première et de seconde espèce. Les critères indiqués plus haut se rapportent tous à la complète réductibilité de première espèce, seule à se présenter dans le cas parfait; et ainsi, dans le cas le plus général, le théorème du discriminant dit qu'un élément premier du domaine de base entre dans le discriminant d'un ordre si et seulement si, dans la décomposition idéale de $(p)$, apparaît au moins une composante proprement primaire ou un idéal premier de seconde espèce.

Ce fait a d'abord été remarqué sur des exemples par Kronecker pour l'ordre maximal, après que la Festschrift contenait encore par erreur la forme valable pour l'ordre maximal d'un corps de nombres, mais sans démonstration complète. Pour l'ordre maximal, Ostrowski\footnote{A. Ostrowski: Zur arithmetischen Theorie der algebraischen Größen. Nachr. Gött. Ges. d. Wissensch. 1919.} a donné une démonstration du théorème du discriminant ci-dessus et d'autres théorèmes de ramification qui s'y rattachent; on y trouve en même temps une bibliographie détaillée. Enfin le théorème du discriminant pour tous les ordres de corps relatifs en résulte comme conséquence directe, par localisation en chaque idéal premier du domaine de base. L'idée de ce passage, connue dans le cas particulier, a été soulignée en principe par W. Krull;\footnote{Cf. sa conférence de Dantzig: Jhrbr. d. D. Math. Ver. 34, p. 121.} une théorie systématique de ces localisés dans le cadre de recherches plus générales est donnée par H. Grell.\footnote{\emph{Beziehungen zwischen den Idealen verschiedener Ringe} (à paraître dans les Math. Annalen).}

Le traitement uniforme donné ici de tous les ordres, ainsi que des domaines de coefficients parfaits et imparfaits, repose, outre la connaissance de la théorie des idéaux, sur le fait que l'on prend d'emblée pour base la théorie des anneaux de rang fini, et donc la notion générale d'extension finie, et non la notion d'extension engendrée par un élément primitif. La base de l'ordre comme module sur le domaine de base, avec les relations entre ses éléments, tient donc lieu des puissances d'un élément primitif du corps et de l'équation définissante,\footnote{Cette conception touche de près celle que Dedekind a défendue au sujet des nombres complexes supérieurs: Zur Theorie der aus $n$ Haupteinheiten gebildeten komplexen Größen. Gött. Nachr. 1885.} respectivement des puissances de la forme fondamentale et de l'équation fondamentale.

Une telle base de l'ordre comme module sur le domaine de base passe, modulo chaque élément premier de ce domaine, en une base de l'anneau quotient, tandis que cela peut déjà ne pas être le cas pour l'ordre maximal avec une base formée des puissances d'un élément primitif; ni non plus avec la base formée des puissances de la forme fondamentale dès qu'il s'agit de fonctions algébriques entières de plusieurs indéterminées.\footnote{Cf. Ostrowski, loc. cit., ``irrégularité de première espèce''.} Autrement dit: l'anneau quotient de l'anneau de polynômes en une indéterminée, à coefficients dans le domaine de base, par une équation définissante quelconque, ou même par l'équation fondamentale, n'a pas besoin d'être isomorphe modulo $p$ à l'anneau quotient de l'ordre maximal par $(p)$; c'est dans cette cessation de l'isomorphisme que résident les difficultés principales des présentations usuelles. Comme cet anneau de polynômes auxiliaire n'intervient pas ici, ces difficultés sont par là même évitées. En même temps, la théorie additive de la décomposition de l'anneau quotient peut être comprise comme équivalent de la décomposition multiplicative de l'équation, à laquelle correspond encore ici, lors de la décomposition en facteurs premiers entre eux, une décomposition additive dans l'anneau de polynômes quotient.

Il faut remarquer que le théorème du discriminant ne peut représenter que le premier théorème d'une théorie générale de la ramification des ordres, de même que, dans l'ordre maximal, le théorème du discriminant est accompagné de la théorie plus fine de l'idéal de ramification, qui permet là une détermination plus exacte des exposants. Cette détermination des exposants, qui repose sur la représentation par puissances des facteurs primaires des équations, peut s'interpréter comme la détermination de la longueur d'une suite de composition de $\mathfrak p$ à $\mathfrak q=\mathfrak p^\varrho$ et apparaîtra probablement sous cette forme dans la théorie générale de la ramification.

\subsection*{\S 1. Anneaux d'extension de corps}

\textbf{1.} Plaçons d'abord quelques remarques sur des anneaux arbitraires. Si $T$ est un anneau, $S$ un système d'éléments de $T$, on entend par ``anneau dérivé de $S$ dans $T$'' l'intersection de tous les anneaux contenus dans $T$ qui comprennent $S$. De façon correspondante est défini l'``idéal dérivé de $S$ dans $T$'', ou le ``$P$-module dérivé de $S$ dans $T$'', lorsque $P$ est un sous-anneau de $T$.\footnote{Cf. Théorie des idéaux, \S 1.}

Si $R$ est un sous-anneau de $T$ et si $S$ est un système d'éléments de $T$, l'anneau dérivé de $R$ et $S$ dans $T$ est désigné comme l'anneau $R[S]$ né de l'adjonction annulaire de $S$ à $R$. Cet anneau se compose de tous les polynômes en les éléments de $S$ à coefficients dans $R$, les relations d'égalité et d'opération étant fixées par les relations d'égalité et d'opération dans $T$. En vertu de ces relations d'égalité et d'opération déjà fixées, il suffira, dans un cas spécial, de former seulement un sous-système de tous les polynômes, par exemple seulement toutes les formes linéaires en $S$ à coefficients dans $R$. Un tel cas spécial se présente par exemple (cf. 2.) lorsque $R$ est un anneau avec élément unité et $S$ un corps avec le même élément unité; toutes les combinaisons linéaires $\sum c_i\gamma_i$, où les $c_i$ appartiennent à $S$ et les $\gamma_i$ à $R$, forment l'anneau $R[S]$.

Mais on peut aussi, lorsque seul l'anneau $R$ est donné et qu'en outre est donné un système $S$ d'éléments non contenus dans $R$, avec des relations d'égalité ou aussi des relations d'opération mutuelles, construire l'anneau d'extension $R[S]$ né de l'adjonction de $S$ à $R$, en montrant qu'il existe toujours un anneau comprenant $R$ et $S$: à savoir celui qui est formé formellement de tous les polynômes en $S$ à coefficients dans $R$, avec fixation des règles ordinaires d'opération. Pour la définition de l'égalité, abstraction faite des règles d'opération, une pleine liberté reste possible, sous la seule restriction que les définitions d'égalité valables dans $R$ et dans $S$ soient comprises. Il faut donc nécessairement identifier seulement les polynômes qui, en vertu de l'égalité dans $R$ et dans $S$ et en vertu des opérations, peuvent être formés de la même manière à partir d'éléments égaux de $R$ et de $S$. L'anneau $T$ ainsi construit, qui comprend $R$ et $S$, devient alors identique à celui qui est dérivé dans $T$ de $R$ et de $S$.\footnote{Cf. la construction steinitzienne du corps des fractions.}

\textbf{2. Anneaux d'extension de corps.} À partir de maintenant, on supposera des anneaux $\mR$ avec élément unité, qui contiennent un corps $P$ comme sous-anneau, donc des sur-anneaux de corps. L'élément unité $e$ de $\mR$ devient en même temps élément unité de $\mathsf{P}$; et l'anneau devient un $\mathsf{P}$-module.

Si $\overline{\mathsf{P}}$ est un sur-corps de $\mathsf{P}$, l'anneau d'extension $\overline{\mathfrak R}=\mathfrak R[\overline{\mathsf{P}}]$ né de l'adjonction annulaire de $\overline{\mathsf{P}}$ à $\mR$ sera expliqué comme suit:

On suppose que l'intersection ensembliste $[\mathfrak R,\overline{\mathsf{P}}]$ de $\mR$ et de $\overline{\mathsf{P}}$ est donnée par $P$, et en outre qu'il n'existe pas de relations d'opération entre les éléments de $\mR$ et de $\overline{\mathsf{P}}$ qui n'appartiennent pas à $P$. Cette hypothèse peut toujours être réalisée en remplaçant $\overline{\mathsf{P}}$ par un corps d'extension équivalent de $P$, ou encore $\mR$ par un anneau d'extension équivalent.\footnote{Deux corps d'extension de $\mathsf{P}$ sont appelés, d'après Steinitz, ``équivalents'' lorsqu'ils peuvent être mis en isomorphisme l'un avec l'autre de telle façon que $\mathsf{P}$ se corresponde élément par élément; de façon correspondante sont à définir les anneaux d'extension équivalents de $\mathsf{P}$. Par introduction de nouveaux symboles, on peut toujours produire des extensions équivalentes; cette liberté dans l'emploi d'extensions équivalentes est essentielle pour ce qui suit.}

Comme éléments de $\bR$, on déclare toutes les combinaisons linéaires formellement constituées
\[
   \bar c_{i_1}\gamma_{i_1}+\cdots+\bar c_{i_\sigma}\gamma_{i_\sigma}
\]
où les $\bar c_i$ parcourent tous les éléments de $\overline{\mathsf{P}}$, les $\gamma_i$ tous les éléments de $\mR$.\footnote{En général, les éléments de $\mR$ ou de $\bR$ seront désignés par des lettres grecques, les éléments de $\mathsf{P}$ ou de $\overline{\mathsf{P}}$ par des lettres latines.} La définition doit être univoque: si les $\gamma$ sont remplacés par des éléments qui leur sont égaux dans $\mR$, et les $\bar c$ par des éléments qui leur sont égaux dans $\overline{\mathsf{P}}$, l'élément de $\bR$ doit aussi être déclaré égal.

La somme est définie de manière univoque, c'est-à-dire de sorte que le remplacement d'un terme par un terme égal, dans la définition d'égalité à fixer, donne le même résultat:
\[
   \sum \bar c_i\gamma_i+\sum \bar d_i\gamma_i
   =\sum(\bar c_i+\bar d_i)\gamma_i .
\]
Parmi les $\bar c$ et les $\bar d$ peuvent apparaître des zéros, si bien que les mêmes $\gamma_i$ apparaissent formellement.

Le produit est défini de manière univoque dans le même sens:
\[
   \sum \bar c_i\gamma_i\cdot \sum \bar d_k\gamma_k
   =\sum \bar c_i\bar d_k\gamma_i\gamma_k .
\]
Pour la définition d'égalité, on déclare que, si $\gamma_{i_1},\ldots,\gamma_{i_\sigma}$ sont linéairement indépendants relativement à $\mathsf{P}$, alors deux éléments
\[
   \bar c_{i_1}\gamma_{i_1}+\cdots+\bar c_{i_\sigma}\gamma_{i_\sigma}
   \quad\hbox{et}\quad
   \bar d_{i_1}\gamma_{i_1}+\cdots+\bar d_{i_\sigma}\gamma_{i_\sigma}
\]
ne doivent être égaux, et le sont toujours, que si $\bar c_{i_\lambda}$ devient égal à $\bar d_{i_\lambda}$ dans $\overline{\mathsf{P}}$. Autrement dit: les éléments de $\mR$ linéairement indépendants relativement à $\mathsf{P}$ sont déclarés linéairement indépendants relativement à $\overline{\mathsf{P}}$.

Par la définition de la somme et du produit, la définition d'égalité est ainsi donnée pour tous les éléments de $\bR$; il en résulte en effet leur exprimabilité au moyen d'éléments linéairement indépendants de $\mR$. Car, d'une relation
\[
   \mu=c_1\gamma_1+\cdots+c_s\gamma_s\quad\hbox{dans }\mR
\]
on obtient, par multiplication par $\bar b=\bar b e$ selon la définition du produit:
\[
   \bar b\mu=\bar b e\cdot(c_1\gamma_1+\cdots+c_s\gamma_s)
      =\bar b c_1\gamma_1+\cdots+\bar b c_s\gamma_s,
\]
et donc, en vertu de la définition de la somme, l'exprimabilité de tous les éléments au moyen d'éléments linéairement indépendants de $\mR$.

Lorsque tous les éléments sont exprimés au moyen d'éléments linéairement indépendants de $\mR$, la définition d'égalité devient donc l'égalité des coefficients dans $\overline{\mathsf{P}}$; elle coïncide ainsi, pour les éléments appartenant simultanément à $\mR$ et à $\bR$, ou à $\overline{\mathsf{P}}$ et à $\bR$, avec celle qui y est valable, tandis que l'égalité valable dans $\mR$ et dans $\overline{\mathsf{P}}$ ne dit rien de plus, pour les autres éléments de $\bR$, que l'univocité exigée dans la définition de ces éléments, à cause de l'hypothèse d'intersection et de la détermination supplémentaire.\footnote{Sans cette hypothèse, il y aurait au moins un élément $\bar c$ de $\overline{\mathsf{P}}$ n'appartenant pas à $P$ pour lequel $\bar c=\bar c e=\gamma$ vaudrait. Ainsi, en vertu de la définition d'égalité dans $\mR$, $e$ et $\gamma$ seraient certes linéairement indépendants relativement à $\mathsf{P}$, mais dépendants relativement à $\overline{\mathsf{P}}$. À cause des relations entre éléments de $\overline{\mathsf{P}}$ et de $\mR$, la définition d'égalité ci-dessus ne serait donc pas possible; qu'on pense par exemple à l'abaissement de degré d'un corps d'extension fini lorsque le corps de base est remplacé par un corps intermédiaire. L'extension ci-dessus est au contraire caractérisée par la possibilité d'apparition de nouveaux diviseurs de zéro; ainsi l'anneau quotient de l'anneau de polynômes à coefficients rationnels $\mathsf{P}$ par $x^2+1$ est un anneau sans diviseurs de zéro, isomorphe au corps $\mathsf{P}(\sqrt{-1})$, tandis que, lors du passage à l'anneau de polynômes à coefficients dans $\mathsf{P}(\sqrt{-1})$, apparaissent des diviseurs de zéro, selon que $(x+i)(x-i)=(x^2+1)$, mais aucun facteur n'est divisible par $x^2+1$. Il en va de même, par exemple, pour tout corps de nombres fini; c'est la conception exprimée d'abord par Dedekind pour les nombres complexes supérieurs. Cf. à ce sujet \S 6, no 4, fin.} De même, l'exigence d'univocité de la somme et du produit n'entraîne pas d'autres relations, puisque des éléments égaux peuvent être supposés égaux coefficient par coefficient et fournissent donc formellement les mêmes expressions de somme et de produit.

On vérifie maintenant aussitôt que tous les axiomes de l'opération annulaire sont satisfaits; il existe donc toujours l'anneau d'extension $\mR[\overline{\mathsf{P}}]$ engendré par adjonction annulaire.

\textbf{3. Idéaux dans $\mR$ et $\bR$.} Si $\ma$ est un idéal dans $\mR$, le produit de modules $\bR\ma$ devient un idéal $\overline{\ma}$ dans $\bR$, l'idéal étendu de $\ma$. En même temps, $\overline{\ma}$ devient l'idéal dérivé de $\ma$ dans $\bR$; il se compose de toutes les combinaisons linéaires $\sum \bar c_i a_i$, la somme portant chaque fois sur un nombre fini d'éléments, où les $a_i$ parcourent tous les éléments de $\ma$ et les $\bar c_i$ tous les éléments de $\overline{\mathsf{P}}$.

Si $\overline{\mb}$ est un idéal dans $\bR$, l'intersection $[\overline{\mb},\mR]$ devient un idéal $\mb$ dans $\mR$, l'idéal contracté de $\overline{\mb}$.\footnote{Dans la notation de H. Grell, dans le travail cité à la note 6, p. 83.} Tout idéal $\ma$ de $\mR$ est un idéal contracté, à savoir celui de son idéal étendu $\overline{\ma}$. En effet, tout élément $a$ de $[\mR,\overline{\ma}]$ est de la forme $\sum \bar c_i a_i$, où les $a_i$ peuvent être supposés linéairement indépendants relativement à $\mathsf{P}$ ou à $\overline{\mathsf{P}}$. Les éléments $a,a_i$ deviennent donc linéairement dépendants dans $\bR$, et par conséquent, d'après la définition d'égalité, aussi dans $\mR$; donc $a$ appartient à $\ma$. Comme la réciproque vaut également, $\ma$ devient égal à $[\mR,\overline{\ma}]$.

Si $\bp$ est un idéal premier dans $\bR$, alors $\mpideal=[\mR,\bp]$ devient un idéal premier dans $\mR$. En effet, d'après le second théorème d'isomorphisme,\footnote{Théorie des idéaux, \S 4, no 3. Le théorème y est énoncé seulement pour des idéaux d'un même anneau; la même considération montre toutefois sa validité sous la forme ci-dessus pour des idéaux arbitraires $\overline{\ma}$ de $\bR$. Par l'homomorphisme d'anneaux $\overline{\mathfrak R}\sim\overline{\mathfrak R}\,|\,\overline{\mathfrak a}$, l'homomorphisme d'anneaux $\mathfrak R\sim(\mathfrak R,\overline{\mathfrak a})\,|\,\overline{\mathfrak a}$ est induit pour le sous-anneau $\mR$ de $\bR$. Dans cet homomorphisme, l'élément nul correspond à tous et seulement aux éléments de $[\mR,\overline{\ma}]$; on obtient donc l'isomorphisme d'anneaux $(\mathfrak R,\overline{\mathfrak a})\,|\,\overline{\mathfrak a}\simeq \mathfrak R\,|\,[\mathfrak R,\overline{\mathfrak a}]$. Les classes peuvent alors être représentées à droite et à gauche par les mêmes éléments de $\mR$, car à chaque étape de la démonstration les éléments sont assignés aux classes qu'ils représentent.} l'anneau quotient $\mR|\mpideal$ devient isomorphe à un sous-anneau de $\overline{\mathfrak R}\,|\,\overline{\mathfrak p}$, et donc un anneau sans diviseurs de zéro. On a en effet l'isomorphisme d'anneaux:
\[
   (\mR,\bp)/\bp
     \simeq \mR/[\mR,\bp],
\]
où $(\mR,\bp)$ désigne la somme, le plus grand commun diviseur, des modules $\mR$ et $\bp$.

\subsection*{\S 2. Anneaux de rang fini. Décomposition en somme directe}

\textbf{1. Anneaux de rang fini.} Si $\mathfrak A$ est un $P$-module fini, alors $\mathfrak A$ devient de rang fini, disons $k$, relativement à $\mathsf{P}$: il y a $k$, et pas davantage, éléments linéairement indépendants relativement à $\mathsf{P}$ dans $\mathfrak A$, et tout système de $k$ éléments de $\mathfrak A$ linéairement indépendants relativement à $P$ forme donc une base de $\mathfrak A$ comme $P$-module. Si un $\mathsf{P}$-module $\mathfrak B$ de rang fini est divisible par un autre $\mathfrak A$ du même rang fini, $\mathfrak B\equiv0\pmod{\mathfrak A}$, ils sont donc identiques.

À partir de maintenant, $\mR$ sera supposé de rang fini $n$ relativement à $\mathsf{P}$. Alors tout $\mathsf{P}$-module contenu dans $\mR$ est de rang fini; dans toute chaîne propre de diviseurs ou de multiples de $P$-modules contenus dans $\mR$, le rang de chaque module doit donc être plus grand, respectivement plus petit, que celui du module immédiatement précédent; les chaînes se terminent dans le fini. En particulier, le ``théorème de la double chaîne'' vaut donc pour le système de tous les idéaux de $\mR$.

Il en résulte\footnote{Théorie des idéaux, \S 7, no 2.} qu'un idéal premier $\mpideal$ de $\mR$ ne possède aucun diviseur propre différent de l'idéal unité, donc que l'anneau quotient $\mR/\mpideal$, par un idéal premier différent de l'idéal unité, devient un corps. Si $\mq$ est un idéal primaire appartenant à $\mpideal$, donc $\mq\equiv0\pmod{\mpideal}$ et $\mpideal^{\rho}\equiv0\pmod{\mq}$, alors l'anneau quotient $\mR/\mq$ devient un anneau primaire, c'est-à-dire un anneau dans lequel une puissance (ici, en tout cas, déjà la $\rho$-ième) de tout diviseur de zéro s'annule; et $\mR/\mq$ ne contient que des diviseurs de zéro et des unités. Ce dernier point découle immédiatement du fait que tout idéal non divisible par $\mpideal$ devient premier à $\mpideal$, et donc à $\mq$.

Du théorème de la double chaîne et de l'existence de l'élément unité il résulte en outre\footnote{Théorie des idéaux, \S 7, no 3, théorème V.} que tout idéal $\ma$ de $\mR$ se représente de façon unique comme produit d'un nombre fini d'idéaux primaires deux à deux premiers entre eux. Ce produit, à cause de la coprimalité, devient égal au plus petit commun multiple et correspond à une représentation de l'anneau quotient $\mR/\ma$ comme somme directe d'anneaux primaires,\footnote{Théorie des idéaux, \S 4, no 5. Ce qui est ajouté dans le présent rassemblement se trouve souvent dispersé dans la littérature.} c'est-à-dire que $\mR/\ma$ devient le plus grand commun diviseur de ces anneaux, et la représentation de chaque élément de $\mR/\ma$ comme somme est univoque.

\textbf{2. Représentation des anneaux de rang fini comme somme directe d'anneaux primaires.} Soit donnée la représentation de l'idéal nul de $\mR$:
\[
   (0)=\mq_1\mq_2\cdots\mq_r
      =[\mq_1,\mq_2,\ldots,\mq_r];
\]
et soit la représentation correspondante comme somme directe:
\[
   \mR=\mR_1+\mR_2+\cdots+\mR_r,
   \qquad
   \mR_i\mR_k=(0)\quad\hbox{pour } i\ne k,
   \qquad
   \mR_i^2=\mR_i,
   \qquad
   \mR\mR_i=\mR_i.
\]
Alors $\mR_i$ devient un idéal de $\mR$, et plus précisément
\[
   \mR_i
   =\mq_1\cdots\mq_{i-1}\mq_{i+1}\cdots\mq_r
   =[\mq_1,\ldots,\mq_{i-1},\mq_{i+1},\ldots,\mq_r],
\]
\[
   \mathfrak q_i=\mathfrak R_1+\cdots+\mathfrak R_{i-1}+\mathfrak R_{i+1}+\cdots\mathfrak R_r,
\]
et $\mR_i$ devient isomorphe à l'anneau quotient $\mR/\mq_i$; à chaque élément $\gamma_i$ de $\mR_i$ correspond précisément la classe modulo $\mq_i$ représentée par $\gamma_i$.\footnote{Théorie des idéaux, \S 4, no 5.} Si $\ma$ est un idéal quelconque de $\mR$, la loi distributive donne:
\[
   \ma=\mR\ma
   =\mR_1\ma+\cdots+\mR_r\ma
   =\ma_1+\cdots+\ma_r;
\]
ici, comme $\mathfrak a_i=\mathfrak R_i\mathfrak a_i=\mathfrak R\mathfrak a_i$, les composantes $\ma_i$ sont des idéaux aussi bien dans $\mR_i$ que dans $\mR$. Comme idéaux, les $\ma_i$, de même que les $\mR_i$ eux-mêmes, sont des $P$-modules finis; à cause de la somme directe, leurs rangs s'additionnent pour donner le rang de $\ma$, respectivement de $\mR$.

Si la représentation de l'élément unité comme somme est donnée par
\[
   e=\varepsilon_1+\cdots+\varepsilon_r,
   \qquad
   \varepsilon_i\varepsilon_k=0\quad\hbox{pour } i\ne k,
   \qquad
   \varepsilon_i^2=\varepsilon_i,
   \qquad
   \varepsilon_i\equiv e\,(\mathfrak q_i),
\]
alors la représentation de tout élément $\gamma$ de $\mR$ s'obtient de manière univoque par
\[
   \gamma=\gamma e=\gamma\varepsilon_1+\cdots+\gamma\varepsilon_r
      =\gamma_1+\cdots+\gamma_r,
\]
où donc la composante $\gamma_i=\gamma\varepsilon_i$ est élément de $\mR_i$; en même temps il s'ensuit que $\mR_i=\mR\varepsilon_i$. Des ``relations d'orthogonalité'' vient:
\[
   \gamma\delta=(\gamma_1\varepsilon_1+\cdots\gamma_r\varepsilon_r)
      (\delta_1\varepsilon_1+\cdots+\delta_r\varepsilon_r)
   =\gamma_1\delta_1\varepsilon_1+\cdots+\gamma_r\delta_r\varepsilon_r;
\]
ainsi $\gamma_i\delta_i$ devient la $i$-ième composante du produit, et de même $\gamma_i+\delta_i$ la $i$-ième composante de la somme, faits qui suivent aussi de l'homomorphisme de $\mR$ vers $\mR_i$.

L'anneau $\mR_i$ contient un sous-corps $\mathsf{P}_i=\mathsf{P}\varepsilon_i$ isomorphe à $P$ et est de rang fini relativement à $\mathsf{P}_i$; ce rang est égal au rang que possède $\mR_i$ comme $\mathsf{P}$-module. En effet, pour tout élément $c$ de $\mathsf{P}$, on a $c\varepsilon_i\equiv ce\equiv c\,(\mathfrak q_i)$; donc $c\varepsilon_i\ne0$ dès que $c\ne0$; l'homomorphisme entre $\mathsf{P}$ et $\mathsf{P}_i$ devient un isomorphisme. De plus, toute dépendance linéaire relativement à $\mathsf{P}$ entre des éléments de $\mR_i$ passe, par multiplication par $\varepsilon_i$, en une dépendance relativement à $\mathsf{P}_i$; et réciproquement, puisque $\mathsf{P}_i=\mathsf{P}\varepsilon_i$, toute relation relativement à $\mathsf{P}_i$ représente en même temps une relation relativement à $\mathsf{P}$. Les rangs relativement à $\mathsf{P}$ et à $\mathsf{P}_i$ coïncident donc.

\textbf{3. Décomposition en somme directe des anneaux d'extension.} Si
\[
   \mR=\mR_1+\cdots+\mR_r,
\]
alors, pour $\bR$, on a une représentation
\[
   \bR=\overline{\mR}_1+\cdots+\overline{\mR}_r,
\]
où $\overline{\mR}_i$ désigne l'idéal étendu de $\mR_i$ dans $\bR$, et en même temps l'anneau dérivé de $\mR_i$ dans $\bR$. En effet, de $(0)=\mq_1\cdots\mq_r$ il résulte, par multiplication par $\bR$, la décomposition en produit $(0)=\overline{\mq}_1\cdots\overline{\mq}_r$, où les $\overline{\mq}_i$ redeviennent donc deux à deux premiers entre eux. Il en résulte une décomposition de $\bR$ en somme directe; et comme la $i$-ième composante de cette décomposition est donnée par
\[
   \overline{\mq}_1\cdots\overline{\mq}_{i-1}\,
   \overline{\mq}_{i+1}\cdots\overline{\mq}_r,
\]
elle devient égale à l'idéal étendu de $\mR_i$. Mais cet idéal étendu est, d'après \S 1, no 3, identique à l'anneau d'extension $\mR_i[\overline{\mathsf{P}}_i]$ formé à partir de $\mR_i$ relativement à $\overline{\mathsf{P}}_i=\overline{\mathsf{P}}\varepsilon_i$; il suffit donc, lors des extensions, d'étendre les composantes séparées. Alors $\overline{\mR}_i$ ne sera en général plus un anneau primaire.

\textbf{4. Idéaux premiers de première et de seconde espèce.} D'après 1., l'anneau quotient $\mR/\mpideal$, par tout idéal premier différent de l'idéal unité, devient un corps; celui-ci possède un sous-corps $(P)=(P,\mpideal)/\mpideal=P/[P,\mpideal]\simeq P$ isomorphe à $P$, puisque $\mpideal$ ne peut contenir aucun élément non nul de $P$; $(\mathsf{P})$ se compose donc de toutes, et seulement, les classes modulo $\mpideal$ qui peuvent être représentées par des éléments de $P$.\footnote{Cf. la note relative au \S 1, no 3, sur le second théorème d'isomorphisme.} L'anneau $\mR/\mpideal$ devient de rang fini relativement à $(\mathsf{P})$, donc un corps d'extension algébrique fini. Selon que cette extension est de première ou de seconde espèce,\footnote{Dans la terminologie bien connue de Steinitz: une extension algébrique s'appelle de première espèce si tout élément est racine d'un polynôme irréductible séparable, c'est-à-dire qui se décompose, dans un corps d'extension approprié, en facteurs linéaires tous distincts; sinon elle est de seconde espèce.} $\mpideal$ sera appelé idéal premier de première ou de seconde espèce.

\subsection*{\S 3. Anneaux complètement réductibles}

1. L'anneau $\mR$ s'appelle \emph{complètement réductible} lorsque, dans la décomposition en produit de son idéal nul (\S 2, no 2), toutes les composantes primaires deviennent des idéaux premiers; ou encore -- ce qui, d'après \S 2, nos 1 et 2, est identique -- lorsque $\mR$ peut se décomposer en somme directe d'un nombre fini de corps. L'anneau complètement réductible s'appelle complètement réductible de première espèce lorsque tous les idéaux premiers sont de première espèce (\S 2, no 4); dans le cas contraire, complètement réductible de seconde espèce.

Si $\mathsf{P}$ est un corps parfait, il n'y a donc que complète réductibilité de première espèce; c'est en particulier toujours le cas pour un corps algébriquement clos. Dans la suite, $\overline{\mathsf{P}}$ désignera spécialement le corps algébriquement clos dérivé de $\mathsf{P}$, et $\bR$ sera égal à $\mR[\overline{\mathsf{P}}]$.

2. \textbf{Théorème.} \emph{L'anneau $\mR$ est complètement réductible de première espèce si et seulement si l'anneau d'extension $\bR=\mR[\overline{\mathsf{P}}]$, à $\overline{\mathsf{P}}$ algébriquement clos, est complètement réductible.}

La démonstration se fait en trois étapes.

2a. De la complète réductibilité de $\bR$ -- plus généralement d'un anneau d'extension quelconque -- résulte la complète réductibilité, de première ou de seconde espèce, de $\mR$. En effet, par hypothèse l'idéal nul de $\bR$ possède une décomposition comme plus petit commun multiple d'idéaux premiers,
\[
        (\overline{0})=[\bp_1,\ldots,\bp_s].
\]
Par formation d'intersections, on obtient dans $\mR$ la représentation de l'idéal nul
\[
        (0)=[[\mR,\bp_1],\ldots,[\mR,\bp_s]],
\]
et les idéaux $[\mR,\bp_s]$ sont premiers d'après \S 1, no 3.

2b. De la complète réductibilité de première espèce de $\mR$ résulte la complète réductibilité de première espèce de $\bR$, et plus généralement celle de tout anneau intermédiaire entre $\mR$ et $\bR$.

Comme, d'après \S 2, no 3, à la représentation de $\mR$ comme somme directe correspond une telle représentation de $\bR$, de sorte que chaque composante $\bR_i$ devient l'anneau d'extension $\mR_i[\overline{\mathsf{P}}_i]$ -- où $\overline{\mathsf{P}}_i$ est un corps d'extension isomorphe à $\overline{\mathsf{P}}$ du sous-corps $\mathsf{P}_i$ de la composante $\mR_i$, isomorphe à $P$ --, il suffit de considérer la composante séparée $\mR_i$. On peut donc supposer, sans restriction de généralité, que $\mR$ est un corps, et même un corps d'extension fini de première espèce de son sous-corps $\mathsf{P}$. Ainsi $\mR$ est engendré par adjonction d'un élément primitif de première espèce à $\mathsf{P}$; autrement dit, $\mR$ devient isomorphe à l'anneau quotient $\mathsf{P}[x]/f(x)$, où $f(x)$ désigne un polynôme irréductible séparable dans l'anneau de polynômes $\mathsf{P}[x]$ en une indéterminée $x$. Si $\Omega$ est un corps d'extension de $\mathsf{P}$, alors $\Omega[x]\,|\,f(x)$ devient isomorphe à $\mR[\Omega]$; car, lors du passage de $\mathsf{P}[x]$ à $\Omega[x]$, le rang de l'anneau quotient reste égal au degré de $f(x)$, si bien qu'il s'agit exactement de la construction de l'anneau d'extension indiquée au \S 1, no 2. Dans $\Omega[x]$, $f(x)$ se décompose en polynômes irréductibles séparables deux à deux premiers entre eux, chacun engendrant un idéal premier dans $\Omega[x]$; cela signifie que $\Omega[x]\,|\,f(x)$, et donc, par l'isomorphisme, aussi $\mR[\Omega]$, devient complètement réductible de première espèce. Comme cela vaut pour tout corps d'extension $\Omega$, en particulier aussi pour le corps algébriquement clos $\overline{\mathsf{P}}$, 2b est démontré.

2c. Si $\mR$ est complètement réductible de seconde espèce, alors $\bR$ n'est pas complètement réductible.

Comme en 2b, il suffit de supposer que $\mR$ est un corps, et même un corps d'extension fini de seconde espèce de $\mathsf{P}$. Il faut montrer que, dans la représentation de $\bR$ comme somme directe d'anneaux primaires, apparaît au moins une composante proprement primaire. Pour cela, il suffit de montrer l'existence d'un élément non nul de $\bR$ dont une puissance s'annule. En effet, si $\bar\beta\ne0$, au moins une composante $\bar\beta_\lambda\ne0$; et si $\bar\beta^\sigma=0$, alors $\bar\beta_\lambda^\sigma=0$ (\S 2, no 2, homomorphisme de $\bR$ vers $\bR_i$), et donc $\bR_i$ ne devient pas un corps.

Soit $\gamma$ un élément de seconde espèce de $\mR$, d'exposant $f\geqq 1$; soit
\[
        F(\gamma)=\gamma^{kp^f}+c_1\gamma^{(k-1)p^f}+\cdots+c_k=0
\]
l'équation de plus bas degré à laquelle $\gamma$ satisfasse relativement à $\mathsf{P}$, où $p$ désigne la caractéristique de $\mathsf{P}$. Les éléments
\[
        e,\gamma,\gamma^2,\ldots,\gamma^{kp^f-1}
\]
de $\mR$ sont donc linéairement indépendants relativement à $\mathsf{P}$ et restent par conséquent, d'après la définition d'égalité, linéairement indépendants relativement à $\overline{\mathsf{P}}$ dans $\bR$. Si l'on pose alors
\[
        G(\gamma)=\gamma^{kp^{f-1}}+\sqrt[p]{c_1}\,\gamma^{(k-1)p^{f-1}}+\cdots+\sqrt[p]{c_k},
\]
alors $G(\gamma)$ devient un élément non nul de $\bR$; car, puisque $p>1$ et $f\geqq1$, $kp^{f-1}$ est toujours plus petit que $kp^f$, et il ne peut donc exister aucune dépendance entre les puissances de $\gamma$ intervenant dans $G(\gamma)$. Mais $G(\gamma)^p=F(\gamma)$ et s'annule donc; ainsi 2c est démontré.

De 2a et 2c il résulte que la complète réductibilité de première espèce de $\bR$ entraîne celle de $\mR$; de 2b résulte la réciproque. Le théorème est ainsi démontré.

3. \textbf{Conséquence.} \emph{Si $\mR$ est complètement réductible de première espèce, alors $\bR$ devient somme directe d'anneaux de rang un, donc de corps isomorphes à $\overline{\mathsf{P}}$. Une telle décomposition en anneaux de rang un existe déjà dans un anneau $\mR[\Omega]$, où $\Omega$ est un corps d'extension fini de $\mathsf{P}$.}

En effet, $\bR$ devient d'après 2b somme directe de corps qui, comme extensions finies de sous-corps isomorphes à $\overline{\mathsf{P}}$, sont identiques à ces sous-corps, à cause de la clôture algébrique de $\overline{\mathsf{P}}$. On obtient ainsi
\[
        \bR=\overline{\mathsf{P}}\bar\varepsilon_1+\cdots+\overline{\mathsf{P}}\bar\varepsilon_n,
        \qquad e=\bar\varepsilon_1+\cdots+\bar\varepsilon_n;
\]
les composantes $\bar \varepsilon_i$ de $e$ forment en même temps une base de $\bR$ comme $\overline{\mathsf{P}}$-module.

Si les $\bar e_i$ sont, d'après \S 1, no 2, représentés comme combinaisons linéaires des $\gamma$ de $\mR$ avec coefficients dans $\overline{\mathsf{P}}$, le système de tous les coefficients détermine un corps d'extension fini $\Omega$ de $\mathsf{P}$; les $\bar e_i$ appartiennent donc à $\mR[\Omega]$ et, par la propriété de cette base comme $\Omega$-module, on obtient
\[
        \mR[\Omega]=\Omega\bar e_1+\cdots+\Omega\bar e_n.
\]
Ainsi une décomposition en anneaux de rang un vaut déjà dans $\mR[\Omega]$, et dans tout anneau intermédiaire entre $\mR[\Omega]$ et $\bR$, les composantes indécomposables naissent comme anneaux d'extension de $\Omega\bar\varepsilon_\lambda$.\footnote{On peut caractériser le plus petit corps d'extension $\Omega$ de ce type comme le corps de réunion des corps galoisiens déterminés par les composantes $\mR_i=\mR\varepsilon_i$ de $\mR$. En effet, si $f_i(x)$ est choisi comme en 2b de telle sorte que $\mathsf{P}[x]/f_i(x)$ devienne isomorphe à la composante $\mR_i$, et si $\Omega^{(i)}$ est le corps d'extension galoisien contenu dans $\overline{\mathsf{P}}$ qui suffit exactement à la décomposition de $f_i(x)$ en facteurs linéaires, alors $\Omega$ est le corps de réunion de tous les $\Omega^{(i)}$. À la décomposition de $f_i(x)$ en facteurs linéaires correspond une représentation de $\Omega^{(i)}[x]\,|\,f_i(x)$ comme somme directe de corps de rang un; $\Omega^{(i)}$ est alors le plus petit corps d'extension où cela se produit. À cause de l'isomorphisme, il en va de même pour $\mR_i[\Omega^{(i)}\varepsilon_i]$ et donc pour $\mR[\Omega]$.}

\subsection*{\S 4. Représentation matricielle, traces et discriminants dans les anneaux de rang fini}

Les hypothèses sur l'anneau de base seront, dans ce paragraphe et le suivant, un peu plus générales que jusqu'ici, afin de comprendre aussi les ordres des corps de nombres et de fonctions. Il s'agit d'une formulation générale et uniforme de faits essentiellement connus.

\emph{Hypothèse.} Soit $\mR$ un anneau d'extension d'un anneau $\Sigma$ sans diviseurs de zéro; l'élément unité de $\mR$ appartient déjà à $\Sigma$. Pour tout $\Sigma$-module contenu dans $\mR$ --- donc, en particulier, pour $\mR$ lui-même --- il existe au moins une base finie de ce module, formée d'éléments linéairement indépendants relativement à $\Sigma$. À côté de $\alpha_1,\ldots,\alpha_s$, les éléments $\beta_1,\ldots,\beta_s$ forment une base linéairement indépendante du même $\Sigma$-module si et seulement si le déterminant de transformation est une unité de $\Sigma$.

Cette hypothèse est manifestement satisfaite pour les anneaux de rang fini considérés jusqu'ici, où $\Sigma$ est le corps $\mathsf{P}$. Elle l'est également pour les ordres des corps de nombres et de fonctions, où $\Sigma$ est l'anneau des entiers rationnels, respectivement le domaine fonctionnel.\footnote{Cf. \emph{Théorie des idéaux}, \S 3.}

\textbf{1. Anneaux de matrices homomorphes à $\mR$.} Soit $\alpha_1,\ldots,\alpha_s$ une base de l'idéal $\ma$ comme $\Sigma$-module, et soit $\gamma$ un élément quelconque de $\mR$. Alors $\gamma\alpha_i$ appartient encore à $\ma$, et l'on a, avec des coefficients dans $\Sigma$,\footnote{Les éléments de $\Sigma$ sont désignés par des lettres latines.}
\[
\gamma\alpha_i=c_{1i}\alpha_1+\cdots+c_{si}\alpha_s,
\qquad (i=1,\ldots,s);
\]
ou, sous forme matricielle, si $(\tau_1,\ldots,\tau_s)$ désigne chaque fois la matrice à une ligne formée de ces éléments,
\[
(\gamma\alpha_1,\ldots,\gamma\alpha_s)
=(\alpha_1,\ldots,\alpha_s)
\begin{pmatrix}
 c_{11}&\cdots&c_{1s}\\
 \vdots&&\vdots\\
 c_{s1}&\cdots&c_{ss}
\end{pmatrix}
=(\alpha_1,\ldots,\alpha_s)C.
\]

\emph{Lorsque $\gamma$ parcourt tous les éléments de $\mR$, l'ensemble des matrices $C$ qui lui sont associées au moyen de la base $\alpha_i$ forme un anneau homomorphe $R_\alpha$; l'anneau $R_\alpha$ est isomorphe à l'anneau quotient $\mR/((0):\ma)$, où $(0):\ma$ désigne le quotient de l'idéal nul par $\ma$.} En effet, à la différence $\beta-\gamma$ correspond manifestement la différence $B-C$; il en va de même du produit, car
\[
(\beta\gamma\alpha_1,\ldots,\beta\gamma\alpha_s)
=(\beta\alpha_1,\ldots,\beta\alpha_s)C
=(\alpha_1,\ldots,\alpha_s)BC.
\]\footnote{La notation est disposée de façon à rester valable pour les anneaux non commutatifs $\mR$, lorsque $\Sigma$ est supposé permuter avec tous les éléments.}
Aux éléments $c$ de $\Sigma$ correspondent en particulier les matrices diagonales $cE$, où $E$ est la matrice correspondant à l'élément unité $e$,
$\left(\begin{smallmatrix}e&&\\&\ddots&\\&&e\end{smallmatrix}\right)$.
Comme, en outre, la matrice nulle correspond à tous et seulement aux éléments $\gamma$ pour lesquels $\gamma\ma$ est l'idéal nul, on obtient l'isomorphisme annoncé.

\emph{Une base $\bar a$ de $\ma$ engendre un anneau de matrices isomorphe à $R_\ma$, et même équivalent à celui-ci: $R_{\bar a}=P^{-1}R_\ma P$.} Soit $\bar a_1,\ldots,\bar a_s$ une autre base de $\ma$ comme $\Sigma$-module; on a donc
\[
(\bar a_1,\ldots,\bar a_s)=(a_1,\ldots,a_s)P,
\]
où le déterminant de $P$ est une unité de $\Sigma$. Il vient alors
\[
(\gamma\bar\alpha_1,\ldots,\gamma\bar\alpha_s)
=(\alpha_1,\ldots,\alpha_s)CP
=(\bar \alpha_1,\ldots,\bar \alpha_s)P^{-1}CP,
\]
ou $R_{\bar \alpha}=P^{-1}R_\alpha P$.

\textbf{2. Classes d'idéaux et classes de représentation.} Deux idéaux $\ma$ et $\mb$ de $\mR$ appartiennent à la même classe d'idéaux lorsqu'ils sont isomorphes, considérés comme $\mR$-modules, c'est-à-dire lorsque leurs éléments peuvent être mis en correspondance biunivoque de telle sorte que la différence et la multiplication par un même élément de $\mR$ se correspondent; ainsi, de $\alpha\simeq\beta$ résulte toujours $\gamma\alpha\simeq\gamma\beta$, pour tout $\gamma$ de $\mR$.

Deux anneaux $R_\alpha$ et $R_{\bar\alpha}$ engendrés par la représentation matricielle de 1 appartiennent à la même classe de représentation lorsqu'ils sont équivalents, c'est-à-dire lorsqu'il existe une matrice unimodulaire $P$ telle que $R_{\bar\alpha}=P^{-1}R_\alpha P$.\footnote{La question, résolue dans certains cas particuliers, de savoir si tout anneau de matrices homomorphe peut être engendré par un idéal de $\mR$ conformément à 1 --- éventuellement en admettant des bases linéairement dépendantes --- n'est pas abordée ici. Par définition, les classes de représentation ne comprennent que les anneaux de matrices engendrés par des idéaux de $\mR$ conformément à 1.}

\emph{Les classes d'idéaux et les classes de représentation se correspondent biunivoquement.} Il a été montré en 1 que les différentes bases d'un même idéal engendrent des anneaux $R_\alpha$, $R_{\bar\alpha}$ appartenant à la même classe de représentation; mais des idéaux différents $\ma$, $\mb$ appartenant à une même classe d'idéaux engendrent eux aussi des anneaux de la même classe de représentation. Soit en effet $\beta_1,\ldots,\beta_s$ la base de $\mb$ qui correspond, par l'isomorphisme, à la base $\alpha_1,\ldots,\alpha_s$ de $\ma$; alors $\gamma\alpha_i$ et $\gamma\beta_i$ se correspondent, de même que $c_{1i}\alpha_1+\cdots+c_{si}\alpha_s$ et $c_{1i}\beta_1+\cdots+c_{si}\beta_s$; les anneaux $R_\alpha$ et $R_\beta$ sont donc identiques.

Réciproquement, si $R_\alpha$ et $R_{\bar\beta}$ appartiennent à la même classe de représentation, de sorte que $R_{\bar\beta}=P^{-1}R_\alpha P$ --- les $\alpha_i$ formant une base de $\ma$ et les $\bar\beta_i$ une base de $\mb$ ---, alors
\[
(\beta_1,\ldots,\beta_s)=(\bar\beta_1,\ldots,\bar\beta_s)P^{-1}
\]
forme aussi une base de $\mb$, et $R_\beta=PR_{\bar\beta}P^{-1}=R_\alpha$. Si donc on associe à tout élément $\alpha=c_1\alpha_1+\cdots+c_s\alpha_s$ de $\ma$ l'élément $\beta=c_1\beta_1+\cdots+c_s\beta_s$ de $\mb$, cette correspondance est biunivoque, et les sommes comme les différences se correspondent. De $R_\alpha=R_\beta$ il résulte encore que $\gamma\alpha_i$ et $\gamma\beta_i$ se correspondent, puisqu'ils s'expriment par les mêmes formes linéaires en les $\alpha_i$, respectivement les $\beta_i$; ainsi $\gamma\alpha$ et $\gamma\beta$ se correspondent en général, et les idéaux $\ma$ et $\mb$ appartiennent à la même classe d'idéaux. On appelle \emph{classe principale} la classe de l'idéal unité; toute base de $\mR$ fournit donc, comme classe de représentation, la classe principale.

\textbf{3. Trace d'un élément relativement à une classe.} Soit $C$ la matrice associée dans $R_\alpha$ à l'élément $\gamma$ de $\mR$; le passage de $R_\alpha$ à $R_{\bar\alpha}$ montre que le déterminant $|C|$, et plus généralement $|tE-C|$, où $t$ est une indéterminée, est un invariant de la classe de représentation. D'après 2, il s'agit donc aussi d'un invariant de la classe d'idéaux, bref d'un \emph{invariant de classe}. Les coefficients de $t$ dans $|tE-C|$ sont donc des invariants de classe; en particulier, on appelle le coefficient de $(-t^{s-1})$ la \emph{trace de $\gamma$ relativement à la classe}:
\[
S_{(\ma)}(\gamma)=c_{11}+c_{22}+\cdots+c_{ss}.
\]
Le coefficient indépendant de $t$, $\pm|C|$, est appelé \emph{norme de $\gamma$ relativement à la classe}.

\emph{Les traces $S_{(\ma)}(\gamma)$ relativement à une classe fixe $(\ma)$ forment un $\Sigma$-module auquel $\mR$, considéré comme $\Sigma$-module, est homomorphe.} En effet, d'après l'homomorphisme d'anneaux de $\mR$ dans $R_\alpha$ démontré en 1, on a $(\beta-\gamma)\sim(B-C)$ et $c\beta\sim cEB=cB$. Il en résulte, pour les traces,
\[
S_{(\ma)}(\beta-\gamma)=S_{(\ma)}(\beta)-S_{(\ma)}(\gamma),
\qquad
S_{(\ma)}(c\beta)=cS_{(\ma)}(\beta),
\]
d'où la propriété de module et l'homomorphisme.

\emph{Remarque. Si $\mR$ est un anneau sans diviseurs de zéro, la trace d'un élément est définie indépendamment de la classe; on peut donc parler de la trace tout court; il en va de même de la norme.} En effet, dans un anneau sans diviseurs de zéro, tout idéal différent de l'idéal nul a le rang de l'anneau; toute base d'un idéal est donc en même temps une base de l'idéal unité dans le corps des fractions. Toutes les traces et toutes les normes y sont ainsi engendrées par l'idéal unité et sont, par conséquent, indépendantes de la classe dans $\mR$.

\textbf{4. Discriminant d'un idéal relativement à une classe.} Ce n'est pas le discriminant lui-même, mais seulement l'idéal principal qui en est dérivé dans $\Sigma$, qui se révèle invariant de l'idéal et de la classe.

Soient $\varrho_1,\ldots,\varrho_r$ et $\bar\varrho_1,\ldots,\bar\varrho_r$ deux bases de l'idéal $\mathfrak r$ comme $\Sigma$-module. \emph{Les déterminants $|S_{(\ma)}(\varrho_i\varrho_k)|$ et $|S_{(\ma)}(\bar\varrho_i\bar\varrho_k)|$ ne diffèrent que par le carré d'une unité de $\Sigma$ comme facteur.} Cela vaut en effet pour les deux déterminants $|(\varrho_i\varrho_k)|$ et $|(\bar\varrho_i\bar\varrho_k)|$, qui diffèrent par le carré du déterminant de transformation entre les $\varrho$ et les $\bar\varrho$. En raison de l'homomorphisme de modules démontré en 3, les mêmes relations linéaires existent entre $S_{(\ma)}(\varrho_i\varrho_k)$ et $S_{(\ma)}(\bar\varrho_i\bar\varrho_k)$ qu'entre $\varrho_i\varrho_k$ et $\bar\varrho_i\bar\varrho_k$; sous la transformation cogrediente, le déterminant acquiert donc le même facteur.

On appelle \emph{discriminant de $\mathfrak r$ relativement à la classe $(\ma)$}, déterminé au carré d'une unité près, tout déterminant $|S_{(\ma)}(\varrho_i\varrho_k)|$; on appelle \emph{idéal discriminant} l'idéal principal qui en est dérivé dans $\Sigma$.

\emph{Remarque. Si $\mR$ est un anneau sans diviseurs de zéro, le discriminant d'un idéal est indépendant de la classe choisie.}

Cette indépendance vaut, d'après 3, pour chacun des éléments du déterminant.

\subsection*{\S 5. Sommes directes des classes}

\textbf{1. Somme directe de la classe d'idéaux ou de la classe de représentation.} Si l'idéal $\ma$ est somme directe, $\ma=\mb+\mc$, il résulte de l'isomorphisme que tout idéal $\bar\ma$ de la classe est lui aussi somme directe, $\bar\ma=\bar\mb+\bar\mc$. Les idéaux étant considérés comme $\mR$-modules, $\bar\mb$ est chaque fois isomorphe à $\mb$, et $\bar\mc$ à $\mc$; on peut donc parler de \emph{somme directe de la classe d'idéaux} et de \emph{classes composantes}.

Si un anneau $R_\alpha$ d'une classe de représentation est somme directe de sous-anneaux qui sont en même temps des idéaux de $R_\alpha$, il en va de même, par isomorphisme, de tout anneau $R_{\bar\alpha}$ de cette classe de représentation, et les composantes sont alors des anneaux équivalents. On peut donc parler de \emph{somme directe de la classe de représentation} et de \emph{classes composantes}.

\emph{À la somme directe de la classe d'idéaux correspond la somme directe de la classe de représentation; c'est en ce sens que l'on parlera de somme directe de la classe.}\footnote{La question de savoir dans quelle mesure toutes les sommes directes de classes de représentation peuvent être engendrées par des sommes directes de classes d'idéaux ne sera pas abordée ici.} Soit $\beta_1,\ldots,\beta_l$ une base de $\mb$ et $\gamma_1,\ldots,\gamma_m$ une base de $\mc$; en raison de la somme directe, les $\beta$ et les $\gamma$ réunis forment une base $a$ linéairement indépendante de $\ma$. Si $R_\alpha$ est l'anneau de matrices engendré au moyen de cette base, la matrice associée dans $R_\alpha$ à un élément quelconque $\delta$ de $\mR$ est de la forme
$\left(\begin{smallmatrix}D^{(\beta)}&0\\0&D^{(\gamma)}\end{smallmatrix}\right)$.
Définissons $R^{(\beta)}$ comme le système de toutes les matrices
$\left(\begin{smallmatrix}D^{(\beta)}&0\\0&0\end{smallmatrix}\right)$,
et définissons $R^{(\gamma)}$ de façon correspondante. Il faut montrer que $R^{(\beta)}$ et $R^{(\gamma)}$ sont des sous-anneaux de $R_\ma$ et que $R_\ma$ en est la somme directe; en même temps, $R^{(\beta)}$, respectivement $R^{(\gamma)}$, est isomorphe à l'anneau $R_\beta$, respectivement $R_\gamma$, engendré par $\mb$, respectivement $\mc$.

En effet, par l'homomorphisme, $R^{(\beta)}$ correspond à tous les éléments $\delta$ de $\mR$ pour lesquels $\delta\mc$ s'annule, c'est-à-dire aux éléments du quotient idéal $(0):\mc$; en raison de l'homomorphisme, $R^{(\beta)}$ est donc un idéal de $R_\alpha$, et il en va de même de $R^{(\gamma)}$. L'anneau $R_\alpha$ est la somme de ces idéaux, et cette somme est directe puisque le produit $R^{(\beta)}R^{(\gamma)}$ des deux anneaux s'annule. Enfin, la correspondance qui associe la matrice $D^{(\beta)}$ à la matrice de $R^{(\beta)}$ contenant $D^{(\beta)}$ et des zéros donne un isomorphisme entre $R^{(\beta)}$ et $\mR_\beta$; en ce sens, les classes composantes $(\mb)$ et $(\mc)$ de la classe d'idéaux de $\ma$ correspondent aux classes composantes de $R_\alpha$.

\textbf{2. Comportement de la trace et du discriminant pour une somme directe de classes.} Si $\ma=\mb+\mc$ et si l'on prend, pour déterminer la trace $S_{(\ma)}(\delta)$, la matrice
$\left(\begin{smallmatrix}D^{(\beta)}&0\\0&D^{(\gamma)}\end{smallmatrix}\right)$,
on lit, compte tenu de 1:

\emph{La trace d'un élément, prise relativement à une classe, est la somme des traces prises relativement aux classes composantes.}

En raison de l'annulation du produit $\mb\mc$, on lit encore:

\emph{Si $\beta$ est un élément de $\mb$, alors $S_{(\ma)}(\beta)=S_{(\mb)}(\beta)$; la trace de $\beta$, prise relativement à la classe $(\ma)$, est égale à sa trace prise relativement à la classe $(\mb)$.}

Compte tenu du \S 4, no 4, il en résulte en outre:

\emph{L'idéal discriminant de l'idéal composant $\mb$, pris relativement à la classe $(\ma)$, est égal à l'idéal discriminant de $\mb$ pris relativement à la classe composante $(\mb)$.}

Enfin:

\emph{L'idéal discriminant de $\ma$, pris relativement à la classe $(\ma)$, est égal au produit des idéaux discriminants des composantes, pris relativement aux classes composantes.} En effet, si l'on prend, tant pour former le discriminant que pour former la trace, la base des $\beta,\gamma$ correspondant à la somme directe, une base de l'idéal discriminant de $\ma$ relativement à $(\ma)$ est donnée par
\[
\left|
\begin{array}{cc}
S_{(\ma)}(\beta_i\beta_k)&0\\
0&S_{(\ma)}(\gamma_\mu\gamma_\nu)
\end{array}
\right|
=
\left|
\begin{array}{cc}
S_{(\mb)}(\beta_i\beta_k)&0\\
0&S_{(\mc)}(\gamma_\mu\gamma_\nu)
\end{array}
\right|
\quad
\left(
\begin{matrix}
 i,k=1,\ldots,l\\
\mu,\nu=1,\ldots,m
\end{matrix}
\right).
\]

\textbf{3. Passage aux anneaux d'extension.} Soit $\overline\Sigma$ un anneau d'extension sans diviseurs de zéro de $\Sigma$, tel que l'intersection de $\mR$ et de $\overline\Sigma$ soit $\Sigma$; soit l'anneau d'extension $\bR=\mR[\overline\Sigma]$, obtenu par adjonction de $\overline\Sigma$, défini comme au \S 1, no 2, en tant qu'anneau de même rang.

Si $\ma$ et $\mathfrak r$ sont des idéaux de $\mR$, et $\bar\ma$ et $\bar{\mathfrak r}$ leurs idéaux étendus dans $\bR$, l'idéal discriminant de $\bar{\mathfrak r}$, pris relativement à $(\bar\ma)$, est lui aussi l'idéal étendu, dans $\overline\Sigma$, de l'idéal discriminant de $\mathfrak r$ pris relativement à $(\ma)$. En effet, toute base de $\ma$ ou de $\mathfrak r$ comme $\Sigma$-module est aussi une base de $\bar\ma$ ou de $\bar{\mathfrak r}$ comme $\overline\Sigma$-module; on peut donc prendre, pour construire le discriminant de $\bar{\mathfrak r}$ relativement à $(\bar\ma)$, une base de $\ma$ et une base de $\mathfrak r$. En particulier, l'idéal discriminant de $\bR$ relativement à la classe principale $(\bR)$ est l'idéal étendu de celui de $\mR$ relativement à la classe principale $(\mR)$; cet idéal ainsi défini dans $\Sigma$, respectivement dans $\overline\Sigma$, sera appelé brièvement \emph{l'idéal discriminant de l'anneau}.

\subsection*{\S 6. Critère discriminantiel de la complète réductibilité de première espèce}

À partir de maintenant, $\mR$ sera de nouveau supposé anneau d'extension d'un corps $\mathsf{P}$, et $\overline{\mathsf{P}}$ sera un corps d'extension de $\mathsf{P}$. Le non-évanouissement du discriminant se révélera condition nécessaire et suffisante de complète réductibilité de première espèce. D'après le \S 5, no 3, l'idéal discriminant de $\mR$ peut être déterminé au moyen de l'anneau d'extension $\bR=\mR[\overline{\mathsf{P}}]$; d'après le \S 5, no 2, il suffit de se restreindre aux anneaux primaires dont $\bR$ se compose additivement. Nous considérerons donc d'abord les idéaux discriminants des anneaux primaires.

\textbf{1. Bases spéciales de $P$-modules.} Si $\ma$ et $\mt$ sont des idéaux de $\mR$ et si $\ma$ est divisible par $\mt$, toute base $a_1,\ldots,a_s$ de $\ma$ comme $P$-module peut être complétée, par adjonction d'éléments $\tau_{s+1},\ldots,\tau_t$, en une base de $\mt$ comme $P$-module.\footnote{La représentation matricielle d'un élément quelconque $\gamma$ fournie par cette base est alors de la forme $\left(\begin{smallmatrix}C_\ma&0\\ C_1&C_2\end{smallmatrix}\right)$; réciproquement, l'existence d'une telle représentation matricielle fournie par $\mt$ entraîne l'existence d'un multiple $\ma$ de $\mt$. Comme, dans le cas de la complète réductibilité, les composantes indécomposables ne possèdent aucun idéal différent de l'idéal nul ou de l'idéal unité, une telle représentation ne peut se présenter dans la classe de représentation de ces idéaux composantes. En y joignant la relation, donnée au \S 5, no 1, entre somme directe des classes, on retrouve la définition usuelle de la complète réductibilité dans la littérature.} C'est une conséquence immédiate du rang fini supposé relativement au corps $P$.

Supposons maintenant que $\mR$ soit un \emph{anneau primaire}; soit $\mpideal$ l'idéal premier associé à l'idéal nul, et $\varrho$ son exposant; ainsi $\mpideal^{\varrho}=(0)$ et $\mpideal^{\varrho-1}\ne(0)$. On obtient alors une base spéciale de $\mathsf{P}$-module de $\mR$ en complétant successivement une base de $\mpideal^{\varrho-1}$ en une base de $\mpideal^{\varrho-2}$, et, en général, une base de $\mpideal^\lambda$ en une base de $\mpideal^{\lambda-1}$, $\mR$ étant compté comme $\mpideal^0$. Au moyen d'une telle base on montre:

\emph{Si $\mR$ est un anneau primaire, la trace, prise relativement à la classe principale (\S 4, no 2), de tout élément divisible par l'idéal premier associé $\mpideal$ s'annule.} Soit en effet
$\pi_{0,1},\ldots,\pi_{0,i_0};\ \pi_{1,1},\ldots,\pi_{1,i_1};\ \ldots;\ \pi_{\varrho-1,1},\ldots,\pi_{\varrho-1,i_{\varrho-1}}$
une telle base spéciale de $\mathsf{P}$-module, avec
$\pi_{\lambda,j}\equiv0\pmod{\mpideal^\lambda}$ et
$\pi_{\lambda,j}\not\equiv0\pmod{\mpideal^{\lambda+1}}$.
De $\pi\equiv0\pmod{\mpideal}$ il résulte
$\pi\pi_{\lambda,j}\equiv0\pmod{\mpideal^{\lambda+1}}$; la matrice associée à $\pi$ au moyen de cette base ne contient donc que des zéros sur la diagonale, et la trace de $\pi$ s'annule.

\textbf{2. Idéaux discriminants des anneaux primaires dont le sous-anneau $\mathsf{P}$ est algébriquement clos.} Si l'idéal nul de $\mR$ est premier et si $\mathsf{P}$ est algébriquement clos, $\mR$ est de rang un, donc identique à $\mathsf{P}$ (\S 3, no 3); l'élément unité $e$ peut être pris comme base de $\mR$ comme $P$-module. Il vient alors $S_{(\mR)}(e^2)=S_{(\mR)}(e)=e$; l'idéal discriminant est égal à l'idéal unité.

On peut en outre montrer que l'idéal discriminant d'un anneau proprement primaire dont le sous-anneau $\mathsf{P}$ est algébriquement clos est l'idéal nul. Si l'on prend pour $\mR$ la base spéciale donnée en 1 --- ou si l'on complète seulement une base quelconque de $\mpideal$ en une base de $\mR$ ---, alors $i_0=1$, puisque l'anneau quotient $\mR/\mpideal$ est de rang un; on peut prendre l'élément unité $e$ pour $\pi_{0,1}$. Tous les produits d'éléments de base autres que $e^2$ sont donc divisibles par $\mpideal$, et leur trace relativement à la classe principale s'annule. Le discriminant s'annule alors lui aussi, comme déterminant d'au moins deux lignes --- $\mR$ a été supposé proprement primaire et doit donc posséder au moins un élément de base distinct de $e$ --- dans lequel tous les éléments autres que $S_{(\mR)}(e^2)$ s'annulent.

\textbf{3. Théorème.} \emph{L'anneau $\mR$ de rang fini relativement à un sous-corps $\mathsf{P}$ est complètement réductible de première espèce si et seulement si son idéal discriminant est l'idéal unité dans $\mathsf{P}$; autrement dit, si le discriminant, déterminé au carré d'une unité près, est non nul.}

La complète réductibilité de première espèce équivaut, d'après le théorème du \S 3, no 2, à la complète réductibilité de l'anneau d'extension $\bR=\mR[\overline{\mathsf{P}}]$, c'est-à-dire au fait que ses composantes indécomposables soient des corps isomorphes à $\overline{\mathsf{P}}$. D'après le \S 5, no 3, l'idéal discriminant de $\mR$ et celui de $\bR$ sont simultanément l'idéal nul, respectivement l'idéal unité; d'après le \S 5, no 2, l'idéal discriminant de $\bR$ est le produit des idéaux discriminants des composantes, pris relativement à ces composantes. Chacun de ces idéaux discriminants s'obtient à partir de ceux considérés en 2 --- où les composantes apparaissaient comme des anneaux considérés en eux-mêmes, et non comme des idéaux de $\bR$ --- en remplaçant chaque fois le corps $\overline{\mathsf{P}}_i=\overline{\mathsf{P}}\bar\varepsilon_i$, isomorphe à $\overline{\mathsf{P}}$, par $\overline{\mathsf{P}}$ lui-même. Il s'ensuit que, si $\bR$ est complètement réductible, l'idéal discriminant de $\bR$, produit d'idéaux unités, est lui-même l'idéal unité; si $\bR$ n'est pas complètement réductible, l'idéal discriminant d'au moins une composante est l'idéal nul, et l'idéal discriminant de $\bR$ est donc lui aussi l'idéal nul. Le théorème est démontré.

\textbf{4. Représentation de la trace, de la norme et du discriminant par des éléments conjugués dans le cas de la complète réductibilité de première espèce.} Soit, d'après le \S 3, no 3, lorsque $\mR$ est complètement réductible de première espèce,
\[
\mR[\Omega]=\Omega\bar\varepsilon_1+\cdots+\Omega\bar\varepsilon_n
\]
une décomposition en somme directe d'anneaux de rang un --- donc de corps isomorphes à $\Omega$ --- et soit $\Omega$ le plus petit corps d'extension dans lequel une telle décomposition soit possible. Les $n$ composantes de $\mR$ sont alors données par $K_i\bar\varepsilon_i$, où chaque $K_i$ est un sous-corps de $\Omega$ et où $\mR$ est homomorphe à $K_i$; il faut exprimer la trace, la norme et le discriminant relativement à la classe principale au moyen d'éléments des $K_i$.

Si l'on prend $\bar\varepsilon_1,\ldots,\bar\varepsilon_n$ comme base de $\mR[\Omega]$ comme $\Omega$-module, à tout élément
$\gamma=c_1\bar\varepsilon_1+\cdots+c_n\bar\varepsilon_n$
est associée la matrice
$\left(\begin{smallmatrix}c_1&0\\0&c_n\end{smallmatrix}\right)$,
où $c_\lambda$ appartient à $K_\lambda$ lorsque $\gamma$ appartient à $\mR$. On obtient ainsi
$S(\gamma)=c_1+\cdots+c_n$ et $N(\gamma)=c_1\cdots c_n$,
où $S(\gamma)$ désigne la trace et $N(\gamma)$ la norme relativement à la classe principale.

Soit en outre $\alpha_1,\ldots,\alpha_n$ une base de $\mR$ comme $P$-module; il faut exprimer le discriminant $|S(\alpha_i\alpha_k)|$ par des éléments de tous les $K_\lambda$. Écrivons
$\alpha_i=\alpha_i^{(1)}\bar\varepsilon_1+\cdots+\alpha_i^{(n)}\bar\varepsilon_n$;
alors
$\alpha_i\alpha_k=\alpha_i^{(1)}\alpha_k^{(1)}\bar\varepsilon_1+\cdots+\alpha_i^{(n)}\alpha_k^{(n)}\bar\varepsilon_n$
et
$S(\alpha_i\alpha_k)=\alpha_i^{(1)}\alpha_k^{(1)}+\cdots+\alpha_i^{(n)}\alpha_k^{(n)}$.
Il vient donc
\[
\left|S(\alpha_i\alpha_k)\right|
=
\left|\sum_{\lambda}\alpha_i^{(\lambda)}\alpha_k^{(\lambda)}\right|
=
\left|
\begin{matrix}
 \alpha_1^{(1)}&\cdots&\alpha_n^{(1)}\\
 \vdots&&\vdots\\
 \alpha_n^{(1)}&\cdots&\alpha_n^{(n)}
\end{matrix}
\right|^2.
\]

Si, en particulier, $\mR$ est lui-même un corps, donc un corps d'extension de première espèce de son sous-corps $\mathsf{P}$, l'homomorphisme de $\mR$ dans $K_i$ devient un isomorphisme qui laisse fixes les éléments de $\mathsf{P}$ --- car la composante de $\mathsf{P}$ est $\mathsf{P}\varepsilon_i$. Comme ils sont tous contenus dans le même sur-corps $\Omega$, les $K_i$ sont donc des corps conjugués relativement à $\mathsf{P}$; les $n$ isomorphismes sont tous distincts, comme le montre la représentation déterminantielle ci-dessus du discriminant non nul $|S(\alpha_i\alpha_k)|$.

\emph{Si $\mR$ est un corps, sa décomposition en somme directe dans $K_1,\ldots,K_n$ donne donc les $n$ corps d'extension conjugués, contenus dans le corps d'extension galoisien, de degré $n$ sur $\mathsf{P}$ et isomorphes à $\mR$. La représentation de la trace, de la norme et du discriminant devient la représentation bien connue par les éléments conjugués.} Rappelons encore que $\mR$ est seulement supposé équivalent aux $K_i$, et non contenu lui-même dans le corps d'extension galoisien $\Omega$.\footnote{Cf. à ce sujet la dernière note du \S 1, no 2.}

\subsection*{\S 7. Le théorème du discriminant pour les ordres d'un corps algébrique de nombres ou de fonctions}

\textbf{1.} Les corps de nombres et de fonctions considérés --- lorsqu'il s'agit de fonctions algébriques de plusieurs indéterminées ou de fonctions algébriques entières, on considère des extensions du domaine fonctionnel --- se rangent tous sous le type de corps suivant, avec fixation des éléments entiers.\footnote{\emph{Théorie des idéaux}, \S 3, nos 3 et 4.}

Soit $\mo$ un \emph{anneau principal sans diviseurs de zéro et avec élément unité} (anneau des entiers rationnels, anneau de polynômes en une indéterminée à coefficients dans un corps, domaine fonctionnel), et soit $\Omega$ son \emph{corps des fractions}; soit $K$ une \emph{extension finie de première espèce} de $\Omega$, et soit $\mO$ l'\emph{anneau de tous les éléments de $\mathfrak{K}$ entiers relativement à $\mo$}.\footnote{Pour la définition des éléments entiers relativement à un anneau, cf. \emph{Théorie des idéaux}, \S 1, no 3.} \emph{Tout sous-anneau de $\mO$ qui contient $\mo$ est appelé un ordre; $\mO$ lui-même est appelé ordre maximal.}

Tout $\mo$-module contenu dans $\mO$, en particulier tout idéal d'un ordre et l'ordre lui-même, possède une \emph{base comme $\mo$-module}.\footnote{\emph{Théorie des idéaux}, \S 3, no 1.} Si $K$ est de degré $n$ relativement à $\Omega$, il suffit de se restreindre aux ordres de rang $n$ relativement à $\mo$, car tout ordre de rang inférieur engendre, par passage au corps des fractions, un corps intermédiaire entre $\Omega$ et $K$ de ce même degré, pour lequel toutes les hypothèses restent valables.

\emph{Tout ordre $\mT$ est donc un anneau satisfaisant aux hypothèses générales des \S\S 4 et 5, et même un anneau sans diviseurs de zéro; son discriminant $D_\mT$ est ainsi défini de manière univoque, à un facteur près qui est le carré d'une unité de $\mo$.} Ce discriminant coïncide, à une unité de $\Omega$ près, avec le discriminant du corps $K$ considéré comme $\Omega$-module; il est donc \emph{non nul}, puisque $K$ a été supposé extension de première espèce (\S 6, no 3). Il admet en même temps la représentation donnée au \S 6, no 4, par le carré du déterminant d'éléments de base conjugués.

La théorie des idéaux dans $\mT$ est donnée par le théorème suivant:\footnote{\emph{Théorie des idéaux}, \S 7, no 3.} \emph{Dans $\mT$, tout idéal différent de l'idéal nul se représente de façon univoque comme produit d'un nombre fini d'idéaux primaires deux à deux premiers entre eux, dont les idéaux premiers associés ne possèdent aucun diviseur propre différent de l'idéal unité.}

\textbf{2. Passage de l'ordre aux anneaux de rang fini relativement à un corps.} Soit $p$ un élément premier de $\mo$, de sorte que l'idéal principal dérivé de $p$ dans $\mo$ soit un idéal premier $\mo p$, distinct de l'idéal nul et de l'idéal unité; désignons de même par $\mT p$ l'idéal principal dérivé de $p$ dans $\mT$.\footnote{Cette notation des idéaux principaux sera employée désormais afin d'indiquer en même temps l'anneau dans lequel ils sont considérés.} \emph{L'anneau quotient $\mR=\mT/\mT p$ est un anneau de rang $n$ relativement à un sous-corps $P$ isomorphe au corps $\mo/\mo p$; l'anneau $\mT$ est homomorphe à $\mT/\mT p$.}

L'homomorphisme $\mT\sim\mT\,|\,\mT p$ vaut en général pour tout anneau quotient; le fait que $\mR$ contienne un sous-corps $\mathsf{P}$ isomorphe à $\mo/\mo p$ résulte du second théorème d'isomorphisme,\footnote{Cf. la dernière note du \S 1.} car le sous-anneau $(\mo,\mT p)|\mT p$ est isomorphe à $\mo|[\mo,\mT p]=\mo|\mo p$.

Enfin, $\mR$ est bien de rang $n$ relativement à $P$. Soit $a_1,\ldots,a_n$ une base de $\mT$ comme $\mo$-module, et soient $(a_1),\ldots,(a_n)$ les classes modulo $\mT p$ déterminées par les $a_i$. Toute dépendance linéaire entre les $(a_i)$ relativement à $P$ provient, par l'homomorphisme, d'une congruence
$c_1\alpha_1+\cdots+c_n\alpha_n\equiv0\pmod{\mT p}$,
donc d'une relation
$c_1\alpha_1+\cdots+c_n\alpha_n=p\gamma$,
où $\gamma$ appartient à $\mT$. De la représentation de $\gamma$ par la base des $\alpha_i$ et de l'unicité de cette représentation, il résulte que chaque $c_i$ est divisible par $p$, donc que l'élément $(c_i)$ de $\mathsf{P}$ s'annule; l'indépendance linéaire des $(\alpha_i)$ relativement à $\mathsf{P}$ est ainsi démontrée.

\emph{Par l'homomorphisme, le discriminant $D_\mT$ de $\mT$ passe au discriminant de $\mR$, et la décomposition idéale de $\mT p$ passe à la décomposition de l'idéal nul de $\mR$.} En effet, en raison de l'indépendance linéaire des $(\alpha_i)$ qui vient d'être démontrée, le discriminant de $\mR$ peut être formé au moyen de cette base sous la forme $|S((\alpha_i)(\alpha_k))|$; il provient donc, par l'homomorphisme, du discriminant $|S(\alpha_i\alpha_k)|$ de $\mT$ formé au moyen de la base des $\alpha_i$. La définition de la trace repose ici sur la représentation matricielle et non sur la représentation par les éléments conjugués, valable pour $\mT$ mais pas nécessairement pour $\mR$; la trace, et par suite le discriminant, se forme donc par les opérations de l'anneau à partir de ses éléments.

Si
$\mT p=\mq_1\cdots\mq_r=[\mq_1,\ldots,\mq_r]$
est la représentation de $\mT p$ dans $\mT$, elle passe, par l'homomorphisme, à
$(0)=[(\mq_1),\ldots,(\mq_r)]$
dans $\mR$. D'après le premier théorème d'isomorphisme,\footnote{Cf. \emph{Théorie des idéaux}, \S 4, no 3, et \S 6, après le théorème II.} $\mR\,|\,(\mq_i)$ est isomorphe à $\mT\,|\,\mq_i$; ainsi $\mq_i$ et $(\mq_i)$ sont simultanément des idéaux proprement primaires ou des idéaux premiers, et, en raison de l'homomorphisme, les $(\mq_i)$ sont deux à deux premiers entre eux.

\textbf{3. Théorème du discriminant.} \emph{Un élément premier $p$ de $\mo$ entre dans le discriminant $D_\mT$ d'un ordre $\mT$ si et seulement si, dans la décomposition de l'idéal principal $\mT p$ en composantes primaires deux à deux premières entre elles dans $\mT$, apparaît au moins une composante proprement primaire ou au moins un idéal premier de seconde espèce.}\footnote{Voici un exemple d'apparition d'un idéal premier de seconde espèce. Soit $\mo$ le domaine fonctionnel dérivé de tous les polynômes en $x$ à coefficients entiers rationnels, $\Omega$ son corps des fractions, et $K$ le corps d'extension obtenu par adjonction de $\sqrt{x}$. Considérons l'ordre $\mT$ de base $1,\sqrt{x}$ comme $\mo$-module; comme le montre la substitution $x=y^2$, il s'agit de l'ordre maximal. Le discriminant est $\left|\begin{smallmatrix}1&\sqrt{x}\\1&-\sqrt{x}\end{smallmatrix}\right|^2=2^2x$; l'idéal principal $\mT2$ reste premier dans $\mT$, mais il est de seconde espèce. En effet, en raison de la base $1,\sqrt{x}$, l'ordre $\mT$ est isomorphe à l'anneau quotient $\mo[t]/(t^2-x)$, où $t$ est une nouvelle indéterminée; par suite, $\mT/\mT2=\mR$ est isomorphe à $P[t]/(t^2-x)$, où $P$ est le corps de toutes les fonctions rationnelles de $x$ et de l'indéterminée $u$ intervenant dans le domaine fonctionnel, à coefficients modulo $2$. Comme $t^2-x$ est un polynôme irréductible de $P[t]$, et même un polynôme irréductible inséparable, $\mR$ est un corps, extension de seconde espèce de $P$; ainsi $\mT2$ est un idéal premier de seconde espèce, tandis que $\mT x=(\mT\sqrt{x})^2$ est proprement primaire. Il n'est pas vrai en général que $\mR$ soit engendré par l'adjonction d'un seul élément à $P$: il suffit de remplacer l'indéterminée $x$ par $x$ et $y$ pour constater, dans l'ordre dérivé de $1,\sqrt{x},\sqrt{y},\sqrt{xy}$ --- qui est de nouveau l'ordre maximal ---, que $\mT2$ est un idéal premier de seconde espèce pour lequel $\mR$ n'est engendré qu'après adjonction de deux éléments à $P$; cf. les exemples d'Ostrowski, loc. cit.}

\emph{Si le corps résiduel $\mo/\mo p$ est parfait, tout élément premier $p$ qui entre dans $D_\mT$, et seulement un tel élément, possède donc au moins une composante proprement primaire. Si $\mo/\mo p$ est parfait et si l'on prend pour $\mT$ l'ordre maximal $\mO$, alors $p$ entre dans le discriminant si et seulement s'il est divisible au moins par le carré d'un idéal premier de $\mO$.}

En effet, en raison de la correspondance entre la décomposition idéale de l'idéal principal $\mT p$ dans $\mT$ et celle de l'idéal nul dans $\mR=\mT\,|\,\mT p$ (cf. 2), l'apparition d'une composante proprement primaire ou d'un idéal premier de seconde espèce de $\mT p$ signifie que $\mR$ n'est \emph{pas} complètement réductible de première espèce. D'après le théorème du \S 6, no 3, cela a lieu si et seulement si le discriminant de $\mR$ s'annule, ce qui, en vertu de l'homomorphisme de 2, exprime la divisibilité de $D_\mT$ par $p$.

\subsection*{\S 8. Le théorème du discriminant pour les ordres de corps relatifs}

Si, au lieu de l'anneau principal $\mo$, on prend déjà comme anneau de départ l'ordre maximal d'un corps de nombres ou de fonctions --- plus généralement, un \emph{anneau de multiplication},\footnote{J'emprunte l'expression abrégée «anneau de multiplication» à la note de W. Krull, \emph{Über Multiplikationsringe}, Heidelberger Berichte, 1925, 5e mémoire, no 3. Krull appelle toutefois «anneaux de multiplication réguliers» les anneaux sans diviseurs de zéro qui interviennent ici.} c'est-à-dire un anneau dans lequel vaut la théorie des idéaux de l'ordre maximal ---, alors l'idéal discriminant de $\mT$ relativement à $\mo$ se substitue au discriminant $D_\mT$.\footnote{Dans la littérature, pour les «corps relatifs de nombres», n'apparaît que l'idéal discriminant de l'ordre maximal, le «discriminant relatif». Cf. Hilbert, \emph{Zahlbericht}, \S 14, et Hecke, \emph{Theorie der algebraischen Zahlen}, \S 38.} Le théorème du discriminant se transfère entièrement en vertu des considérations suivantes.

\textbf{1.} Soit $\mo$ un \emph{anneau de multiplication}, donc un anneau sans diviseurs de zéro et avec élément unité, dans lequel tout idéal distinct de l'idéal nul et de l'idéal unité se représente de façon univoque comme produit de puissances d'idéaux premiers, et où ces idéaux premiers ne possèdent aucun diviseur propre différent de l'idéal unité.\footnote{Pour l'axiomatique des anneaux de multiplication, cf. \emph{Théorie des idéaux}, notamment l'introduction.} Soient en outre $\Omega,K,\mO,\mT$ définis comme au \S 7, no 1; alors, dans $\mT$, vaut la théorie des idéaux indiquée en ce lieu.\footnote{\emph{Théorie des idéaux}, \S 3, no 2, et \S 7, no 3.}

Si, en particulier, un anneau de multiplication $\mo$ ne possède qu'un seul idéal premier distinct de l'idéal nul et de l'idéal unité, alors $\mo$ est un anneau principal. Soit en effet $\mpideal$ cet idéal premier, et choisissons $p\equiv0\;(\mpideal)$, $\not\equiv0\;(\mpideal^2)$. Par hypothèse, tout idéal distinct de l'idéal nul et de l'idéal unité est une puissance de $\mpideal$; cela vaut en particulier pour l'idéal principal $\mo p$. Comme $p\not\equiv0\pmod{\mpideal^2}$, il faut donc que $\mo p=\mpideal$; toute puissance de $\mpideal$ est alors elle aussi principale, et l'idéal nul comme l'idéal unité sont également des idéaux principaux.

\textbf{2. Trace et idéal discriminant.} La trace de tout élément $\gamma$ de $\mT$ est \emph{définie comme la trace de $\gamma$ dans $K$, considéré comme $\Omega$-module}; tout système de $n$ éléments de $K$ linéairement indépendants relativement à $\Omega$ peut donc être choisi comme base d'une représentation matricielle qui définit la trace. Si $\alpha_1,\ldots,\alpha_n$ est un tel système, alors, puisque $K$ a été supposé extension de première espèce, $|S(\alpha_i\alpha_k)|$ est un élément non nul de $\Omega$, à savoir le carré du déterminant d'éléments de base conjugués (\S 6, no 4). De la clôture intégrale de $\mo$ dans $\Omega$\footnote{\emph{Théorie des idéaux}, \S 9, no 7.} il résulte que $|S(\alpha_i\alpha_k)|$ appartient à $\mo$ dès que tous les $a_i$ appartiennent à $\mO$.

\emph{Idéal discriminant.} Que $\tau_1,\ldots,\tau_n$ parcoure tous les systèmes de $n$ éléments de $\mT$; pour chacun, formons le déterminant $|S(\tau_i\tau_k)|$, qui appartient à $\mo$ et qui est non nul dès que les $\tau_i$ sont linéairement indépendants. \emph{L'idéal dérivé, dans $\mo$, de l'ensemble de ces déterminants est appelé l'idéal discriminant de l'ordre $\mT$ relativement à $\mo$;} il est donc distinct de l'idéal nul.

Soit en outre $\beta_1,\ldots,\beta_m$, avec $m\geqq n$, une base de $\mT$ comme $\mo$-module, qui existe toujours.\footnote{\emph{Théorie des idéaux}, \S 9, et \S 3, no 1.} Les déterminants, en nombre fini, correspondant chacun à $n$ éléments $\beta$, forment une base d'idéal de l'idéal discriminant. Cela résulte de l'homomorphisme de modules de $\mT$ dans le système des traces de $\mT$ (\S 4, no 3). En effet, les déterminants $|(\tau_i\tau_k)|$ formés des produits $\tau_i\tau_k$ sont des combinaisons linéaires, à coefficients dans $\mo$, des déterminants $|(\beta_i\beta_k)|$, en nombre fini, qui correspondent chacun à $n$ éléments $\beta$; il en va donc de même des $|S(\tau_i\tau_k)|$ relativement aux $|S(\beta_i\beta_k)|$. Si, en particulier, l'ordre possède une base $\alpha_1,\ldots,\alpha_n$ formée d'éléments linéairement indépendants comme $\mo$-module, le déterminant $|S(\alpha_i\alpha_k)|$ peut être pris comme base de l'idéal discriminant; la définition coïncide avec celle du \S 7, no 1.

\textbf{3. Passage aux localisés.}\footnote{Pour les démonstrations de tous les théorèmes cités dans ce numéro, cf. le travail de H. Grell cité à la note 6, p. 83.} Si $\ma$ est un idéal quelconque de $\mT$, le localisé $\mT_\ma$ est défini comme l'ensemble des éléments du corps des fractions $K$ de $\mT$ qui possèdent au dénominateur un élément de $\mT$ premier à $\ma$.\footnote{Comme les idéaux premiers de $\mT$ ne possèdent aucun diviseur propre différent de l'idéal unité, «premier à» coïncide ici avec «premier entre eux», au sens des idéaux.} En dehors de l'idéal nul et de l'idéal unité, l'anneau $\mT_\ma$ ne possède aucun idéal premier autre que les idéaux étendus des idéaux premiers associés à $\ma$. Il existe une correspondance biunivoque et un isomorphisme des anneaux quotients entre tous les idéaux de $\mT_\ma$ distincts de l'idéal nul et de l'idéal unité, et les idéaux de $\mT$ dont les idéaux premiers associés figurent parmi ceux qui sont associés à $\ma$; en outre, l'idéal nul et l'idéal unité se correspondent biunivoquement.

Les mêmes considérations valent pour l'anneau de multiplication $\mo$. Si, en particulier, $\mpideal$ est un idéal premier, $\mo_\mpideal$ ne possède qu'un seul idéal premier distinct de l'idéal nul et de l'idéal unité; en raison de l'isomorphisme des anneaux quotients, $\mo_\mpideal$ est, comme $\mo$, un anneau de multiplication et, d'après 1, il est donc principal; un élément premier $p$ de $\mo_\mpideal$ engendre l'idéal premier dérivé de $\mpideal$. L'idéal étendu d'un idéal quelconque $\mc$ de $\mo$ est engendré par la composante primaire de $\mc$ associée à $\mpideal$; il est donc égal à $(p)^\rho$ ou à l'idéal unité, selon que $\mpideal$ entre dans $\mc$ --- et précisément à la puissance $\rho$ --- ou n'y entre pas. Il en résulte encore que, si les extensions de deux idéaux $\mb$ et $\mc$ de $\mo$ coïncident dans \emph{tous} les localisés $\mo_\mpideal$, alors $\mb$ et $\mc$ sont identiques; ils sont en effet divisibles par les mêmes idéaux premiers $\mpideal$ et aux mêmes puissances.

\textbf{4. L'idéal discriminant lors du passage au localisé.} Soit $\mpideal$ un idéal premier de $\mo$ et désignons par $\mP=\mT\mpideal$ l'idéal étendu dans $\mT$. \emph{Le localisé $\mT_\mP$ possède une base de module formée d'éléments linéairement indépendants sur l'anneau principal $\mo_\mpideal$; en même temps, $\mT_\mP$ est un ordre fini sur $\mo_\mpideal$ dans le corps d'extension $K$ du corps des fractions $\Omega$ de $\mo_\mpideal$.\footnote{Pour la démonstration, cf. H. Grell.} Le théorème du discriminant (\S 7, no 3) vaut donc dans $\mT_\mP$ relativement à $\mo_\mpideal$.}

Le discriminant $D_{\mT_\mP}$ forme alors une base de l'idéal étendu, pris dans $\mo_\mpideal$, de l'idéal discriminant de $\mT$ relativement à $\mo$. Prenons en effet, ce qui est toujours possible, pour base de $\mT_\mP$ comme $\mo_\mpideal$-module des éléments $a_1,\ldots,a_n$ de $\mT$; alors $|S(a_i a_k)|$ appartient à l'idéal discriminant et, d'après 1, tout déterminant $|S(\tau_i\tau_k)|$ est divisible, dans $\mo_\mpideal$, par $|S(a_i a_k)|$.\footnote{Il en résulte, lorsque $\mo$ est l'ordre maximal d'un corps algébrique de nombres, l'accord avec le discriminant relatif défini par Hilbert, qui coïncide, comme on sait, avec celui de Hecke. Soit en effet $\omega_1,\ldots,\omega_m$ une base de l'ordre maximal du corps relatif comme $\mo$-module --- Hilbert choisit en particulier une base comme module sur l'anneau des entiers rationnels ---; on forme, suivant Hilbert, tous les déterminants de $n$ éléments de base et de leurs conjugués. Le «discriminant relatif» est l'idéal dérivé dans $\mo$ du produit deux à deux, égal ou non, de ces déterminants. Dans $\mo_\mpideal$, le discriminant relatif devient donc le carré du déterminant de la base $a_1,\ldots,a_n$ et de ses éléments conjugués, c'est-à-dire $|S(a_i a_k)|$. Les extensions du discriminant relatif et de l'idéal discriminant de l'ordre maximal relatif coïncident ainsi dans tous les localisés $\mo_\mpideal$; d'après la fin de 3, les deux idéaux sont identiques.}

\textbf{5. Théorème général du discriminant.} \emph{Un idéal premier $\mpideal$ d'un anneau de multiplication $\mo$ entre dans l'idéal discriminant d'un ordre relativement à $\mo$ si et seulement si, dans la décomposition idéale de $\mT\mpideal$ en composantes primaires deux à deux premières entre elles dans $\mT$, apparaît au moins une composante proprement primaire ou au moins un idéal premier de seconde espèce.}

Il résulte de 3 et de 4 qu'un idéal premier $\mpideal$ de $\mo$ entre dans l'idéal discriminant de $\mT$ si et seulement si l'élément premier $p$ de $\mo_\mpideal$ entre dans le discriminant $D_{\mT_\mP}$. Cela équivaut à dire que l'anneau quotient $\mT_\mP/\mT_\mP p$ n'est pas complètement réductible de première espèce; comme cet anneau est, d'après 3, isomorphe à $\mT/\mT\mpideal$, le théorème général du discriminant est démontré.

Comme spécialisation aux discriminants relatifs des corps de nombres, on obtient:

\emph{Un idéal premier $\mpideal$ de l'ordre maximal d'un corps de nombres entre dans le discriminant relatif d'un corps d'extension si et seulement si $\mpideal$ est divisible, dans l'ordre maximal du corps d'extension, par le carré d'un idéal premier.}

Göttingen, 30 mars 1926.

\begin{center}
\rule{9em}{0.4pt}
\end{center}
\clearpage
\setcounter{footnote}{0}

\providecommand{\mS}{\mathfrak S}
\providecommand{\mA}{\mathfrak A}
\providecommand{\mat}[1]{\begin{pmatrix}#1\end{pmatrix}}

\section*{\NoetherNavAnchor{P32}{32. Avec R. Brauer : sur les corps de décomposition minimaux des représentations irréductibles}32. Avec R. Brauer: Sur les corps de décomposition minimaux des représentations irréductibles}
\emph{Sitz. Ber. d. Preuß. Akad. d. Wiss. 1927, p. 221--228.}

\begin{center}
{\Large\bfseries Sur les corps de décomposition minimaux des\\
représentations irréductibles.}\\[1.0em]
Par le privat-docent Dr \textsc{Richard Brauer}\\
à Königsberg\\
et la professeure Dr \textsc{Emmy Noether}\\
à Göttingen.\\[0.8em]
\emph{(Présenté par M. \textnormal{\textsc{Schur}}.)}\\[0.5em]
\rule{4em}{0.4pt}
\end{center}

La question des corps de nombres de plus petit degré sur un corps de base $\mathrm{P}$, dans lesquels une représentation irréductible sur $\mathrm{P}$ d'un groupe fini se décompose en constituants absolument irréductibles, fut d'abord traitée par I. Schur.\footnote{Arithmetische Untersuchungen über endliche Gruppen, Sitzungsber. d. Berl. Ak. d. Wiss. 1906, p. 164. Beiträge zur Theorie der Gruppen linearer Substitutionen, Transact. of the Am. Math. Soc. Ser. 2, Vol. XV, 1909, p. 159. Dans le second travail, les résultats sont transférés aux systèmes hypercomplexes complètement réductibles. Une représentation est réunie en une classe avec toutes ses transformées (semblables).} Supposons par exemple que $\mathrm{P}$ contienne le caractère d'un tel constituant. Schur montra alors que le degré de ces corps par rapport à $\mathrm{P}$ -- l'indice -- coïncide avec le nombre de ces constituants absolument irréductibles qui appartiennent tous à la même classe; que, par conséquent, le même corps est déjà déterminé par l'exigence de séparer un constituant absolument irréductible; et encore que le degré, par rapport à $\mathrm{P}$, de tout corps dans lequel une telle décomposition a lieu devient un multiple de l'indice. Tous ces corps seront appelés corps de décomposition, même dans le cas où $\mathrm{P}$ ne contient pas le caractère. Par un exemple, Schur montra que les corps de décomposition de plus petit degré n'ont pas besoin d'être isomorphes. Si $\mathrm{P}$ ne contenait aucun caractère simple correspondant, alors les corps de décomposition doivent comprendre le corps d'un tel caractère et ses conjugués par rapport à $\mathrm{P}$; les représentations absolument irréductibles appartenant à différents caractères conjugués sont bien conjuguées, mais appartiennent évidemment à des classes différentes. Pour séparer un système de représentations absolument irréductibles d'une même classe, il suffit encore ici du corps du caractère correspondant et de l'un des corps de décomposition correspondants.

Dans ce qui suit, la question des corps de décomposition sera poursuivie. Les corps de décomposition de plus petit degré se caractérisent par des corps gauches associés; les corps de décomposition généraux, par des anneaux simples bilatères liés invariantement à ces corps gauches. Ensuite, à l'exemple du corps des quaternions, on montre que les degrés des corps de décomposition minimaux ne sont pas bornés; ici un corps de décomposition est dit minimal lorsque aucun de ses sous-corps propres n'est corps de décomposition. Cette non-bornitude résulte essentiellement d'un théorème d'existence arithmétique dont H. Hasse donnera la démonstration dans la note qui suit immédiatement. Mais l'exemple est encore traité -- plutôt à titre de vérification -- par un calcul élémentaire; et l'on montre ainsi, sans recourir à la théorie générale ni au théorème d'existence arithmétique, la non-bornitude des degrés en démontrant élémentairement un théorème d'existence moins étendu, mais suffisant.\footnote{Cet exemple est dû à E. Noether; il repose sur une modification d'un exemple provenant de R. Brauer, où était déjà montrée l'existence d'un corps de décomposition minimal dont le degré dépasse l'indice. Le traitement élémentaire de l'exemple est dû à R. Brauer. La caractérisation des corps de décomposition remonte à des recherches des auteurs, qui poursuivent par ailleurs des buts distincts. Pour E. Noether, il s'agit d'une construction de la théorie des représentations sur la base de la théorie des modules et des idéaux, sous des hypothèses générales de finitude; pour R. Brauer, de la construction des corps gauches de degré fini sur un corps de base commutatif parfait donné, à l'aide des systèmes de facteurs. Les auteurs sont arrivés indépendamment à la caractérisation des corps de décomposition de plus petit degré, R. Brauer pour un corps de base parfait, E. Noether sans cette restriction. La caractérisation des corps de décomposition généraux fut trouvée par R. Brauer à la suite d'une question de E. Noether sur la possibilité d'un plongement non commutatif; E. Noether put ici encore supprimer la restriction au domaine de base parfait.} Remarquons encore qu'il s'agit, dans cet exemple aussi, de corps de décomposition d'un groupe fini; à savoir du groupe des quaternions, qui est aussi à la base de l'exemple de Schur. Les représentations qui y sont indiquées sont en effet simultanément des représentations (isomorphes) du corps des quaternions.

\subsection*{1. Caractérisation des corps de décomposition}

Rappelons d'abord brièvement les notions fondamentales. Soit $\mS$ un système hypercomplexe relativement à un corps $\mathrm{P}_0$. Par représentation $\Gamma$ de $\mS$ -- dans $\mathrm{P}$ -- on entend un système de matrices à éléments dans $\mathrm{P}$, de telle sorte que $\Gamma$ devienne une image homomorphe de $\mS$, et que les éléments $\rho_0$ de $\mathrm{P}_0$ soient associés à des matrices diagonales $\rho_0E$; $\mathrm{P}$ doit donc comprendre $\mathrm{P}_0$. Une représentation $\Gamma$, avec toutes ses transformées $\mathit{P}^{-1}\Gamma \mathit{P}$, forme une classe de représentations. Les représentations ainsi définies comprennent les représentations des groupes finis $\mathfrak G$; le système hypercomplexe associé $\mS$, l'``anneau de groupe'', se compose de tous les ``nombres du groupe'', c'est-à-dire de toutes les combinaisons linéaires des éléments du groupe, à coefficients dans $\mathrm{P}_0$, les éléments primitifs du groupe étant considérés comme linéairement indépendants et la multiplication étant définie par la multiplication du groupe et les règles du calcul. Les notions de représentations ``réductibles'', ``irréductibles'', ``complètement réductibles'', ``absolument irréductibles'' sont définies comme d'habitude; la terminologie sera étendue aux classes de représentations. Un système $\mS$ est dit ``complètement réductible dans $\mathrm{P}$'' lorsque toutes ses classes de représentations dans $\mathrm{P}$ sont complètement réductibles.

Supposons maintenant $\mS$ complètement réductible dans $\mathrm{P}_0$; supposons $\mathrm{P}_0$ corps parfait, de sorte qu'il en soit de même pour toute extension algébrique $\mathrm{P}$ de $\mathrm{P}_0$. Les classes de représentations restent complètement réductibles lors de toute extension $\mathrm{P}$ de $\mathrm{P}_0$.\footnote{Si $\mathrm{P}_0$ n'est pas parfait, la complète réductibilité se conserve si et seulement si les composantes $\mathfrak{K}$ du centre deviennent des ``extensions de première espèce'' de $\mathrm{P}_0$. Sous cette hypothèse, toutes les conséquences ultérieures du texte restent valables.} Le centre de $\mS$ devient somme directe d'un nombre fini de corps; toute extension $P$ de $\mathrm{P}_0$ isomorphe à un tel corps $K$ devient corps d'un caractère (simple). Si l'on étend $\mathrm{P}_0$ à un $P$ isomorphe à $K$, il se détache de $\mS$ un anneau simple bilatère auquel est associée la classe de représentations absolument irréductibles du caractère. Cet anneau devient -- théorème de Wedderburn -- le produit direct d'un anneau de matrices sur $\mathrm{P}$ avec un corps gauche $\mA$, déterminé de façon unique à isomorphie près, dont le centre est donné par $\mathrm{P}$ et qui est de degré $m^2$ par rapport à $\mathrm{P}$, où $m$ désigne l'indice de Schur appartenant au caractère. Tout corps de décomposition de $\mS$ -- relativement à la représentation appartenant à ce caractère -- est en même temps un corps de décomposition de $\mA$ relativement aux représentations dans $\mathrm{P}$, et réciproquement. Les représentations conjuguées de $\mS$ s'obtiennent en partant de l'anneau simple bilatère correspondant sur $\mathrm{P}_0$; son corps associé $\mA_0$ devient isomorphe à $\mA$ et possède $\mathfrak{K}$ pour centre. Le problème des corps de décomposition est donc entièrement ramené à celui du corps gauche associé $\mA$ et de ses décompositions.\footnote{Cela correspond à la réduction de Schur au système des matrices sur $\mathrm{P}$ qui commutent avec la représentation irréductible de $\mS$ dans $\mathrm{P}$. Les transposées de ces matrices donnent une représentation de $\mA$ dans $\mathrm{P}$.} En particulier, on a le théorème suivant:

Tous les sous-corps commutatifs maximaux de $\mA$ sont de même degré $m$ par rapport à $\mathrm{P}$, où $m$ signifie l'indice de Schur. Tous ces sous-corps commutatifs maximaux sont des corps de décomposition de plus petit degré, et tout corps de décomposition de plus petit degré est contenu parmi eux. Le degré de tout corps de décomposition est un multiple de $m$. $\mA$ ne possède qu'une seule classe de représentations absolument irréductibles, et de même une seule classe de représentations irréductibles relativement à $\mathrm{P}$.

Pour caractériser les corps de décomposition généraux, désignons par $\mA_r$ le produit direct de $\mA$ avec l'anneau de matrices de degré $r^2$ relativement à $\mathrm{P}$ (système de toutes les matrices à $r$ rangées avec éléments dans $\mathrm{P}$). Alors $\mA_r$ est simple bilatère, avec $\mA$ pour corps gauche associé; et tout anneau simple bilatère, hypercomplexe relativement à $\mathrm{P}$, auquel $\mA$ est associé, s'obtient ainsi. $\mA_r$ est aussi caractérisé comme système de toutes les matrices à $r$ rangées avec éléments dans $\mA$.

Tout corps de décomposition de $\mA$ de degré $mr$ -- s'il en existe un -- est isomorphe à un sous-corps commutatif maximal de $\mA_r$. Réciproquement, tous les sous-corps commutatifs maximaux de $\mA_r$ sont des corps de décomposition, chacun d'un degré $ms$, où $s$ est un diviseur de $r$. Dans le cas régulier, tous les sous-corps commutatifs maximaux de $\mA_r$ sont de degré $mr$. On entend ici par cas régulier: le corps de base $\mathrm{P}$ a la propriété que, pour tout corps commutatif $\Sigma$ sur $\mathrm{P}$ fini sur $\mathrm{P}$, il existe encore des extensions commutatives sur $\Sigma$ de degré arbitraire.\footnote{Il n'y a donc par exemple aucun cas régulier lorsque $\mathrm{P}$ est le corps de tous les nombres réels.}

Que, parmi ces corps de décomposition pour $r>1$, il en apparaisse aussi de minimaux n'est pas encore décidé par le plongement dans $\mA_r$. Que cela puisse effectivement avoir lieu pour une infinité de $r$, c'est ce que montre l'exemple suivant.

\subsection*{2. La non-bornitude des degrés des corps de décomposition minimaux}

Pour cette non-bornitude, on prendra maintenant pour exemple le corps des quaternions, regardé comme système hypercomplexe relativement au corps $\mathrm{P}$ des nombres rationnels, avec les éléments de base $i,j,k$ outre l'unité $1$ et avec les relations connues:
\[
 i^2=j^2=k^2=-1;\qquad ij=k;\quad jk=i;\quad ki=j;\quad ji=-k;\quad kj=-i;\quad ik=-j.
\]
On montrera que tous, et seulement, les corps de nombres algébriques sont corps de décomposition en vertu de l'adjonction desquels il existe un élément ``idempotent'' $r$, pour lequel donc $r^2=r$, où $r$ est encore distinct de l'élément nul et de l'élément unité. Ces corps de nombres peuvent aussi être caractérisés comme ceux dans lesquels $(-1)$ est représentable comme somme de trois carrés, et par conséquent comme somme de deux carrés.\footnote{Que la représentabilité de $(-1)$ comme somme de trois carrés entraîne toujours une telle représentabilité comme somme de deux carrés, c'est ce que montre l'identité
\[
(c^2+d^2)(a^2+b^2+c^2+d^2)=(ac+bd)^2+(ad-bc)^2+(c^2+d^2)^2,
\]
qui résulte par exemple du théorème du produit des normes des quaternions. Au lieu d'exiger l'existence d'un élément idempotent, il aurait aussi suffi, pour la décomposition, d'exiger un élément de norme nulle.} En effet, posons
\[
        r=\alpha+\frac12\beta i+\frac12\gamma j+\frac12\delta k;
\]
donc
\[
        r^2=\left(\alpha^2-\frac14\beta^2-\frac14\gamma^2-\frac14\delta^2\right)
          +\alpha\beta i+\alpha\gamma j+\alpha\delta k,
\]
alors $r^2=r$ équivaut à:
\[
4\alpha^2-4\alpha-\beta^2-\gamma^2-\delta^2=0;\quad
2\alpha\beta-\beta=0;\quad 2\alpha\gamma-\gamma=0;\quad 2\alpha\delta-\delta=0,
\]
et donc, puisque $\beta,\gamma,\delta$ ne doivent pas s'annuler simultanément, équivaut à
\[
        (-1)=\beta^2+\gamma^2+\delta^2.
\]

Que tout tel élément idempotent produise une décomposition, et que l'existence d'un tel élément idempotent soit nécessaire pour la décomposition, résulte des faits suivants: toute représentation irréductible d'un système complètement réductible est engendrée par un idéal simple unilatère, plus exactement par la classe d'idéaux correspondante, et toute telle classe d'idéaux engendre la représentation. Ces idéaux simples unilatères apparaissent tous lors d'une décomposition en somme directe unilatère, et celle-ci détermine les éléments idempotents ``primitifs'', qui deviennent les composantes de l'unité. Tout élément idempotent donne réciproquement une décomposition en somme $\mS=r\mS+(e-r)\mS$, où $e$ désigne l'unité; les éléments ``primitifs'' détachent des idéaux simples.\footnote{Le lien entre éléments idempotents primitifs et représentations irréductibles se trouve déjà dans la dissertation de M. Herzberger: Über Systeme hyperkomplexer Größen, chapitre 4, Berlin 1923. Que les idéaux simples engendrent des représentations irréductibles est montré, dans le cas commutatif, par exemple chez E. Noether, Der Diskriminantensatz für Ordnungen, Crelle 157 (1926); dans le cas non commutatif chez W. Krull, Theorie und Anwendung der verallgemeinerten Abelschen Gruppen, Heidelberger Berichte 1926. Cf. aussi le mot final de la théorie des groupes de Speiser. Si $a_1,\cdots,a_m$ désigne une base de l'idéal relativement au corps de décomposition, et $c$ un élément arbitraire de l'anneau, de sorte que $a_i c=\gamma^i_1 a_1+\cdots+\gamma_{im} a_m$, alors $(\gamma_{ik})$ est la matrice associée à $c$; le fait que toutes les représentations irréductibles soient engendrées par des idéaux est plus profond.} Un corps gauche, considéré comme système hypercomplexe relativement à son centre, devient -- étendu à un système hypercomplexe relativement à un corps de décomposition -- somme directe de $m$ idéaux unilatères absolument simples, où $m$ est l'indice de Schur.

Dans le cas du corps des quaternions, $m$ vaut deux; et ainsi tout élément idempotent $\ne0$ et $\ne1$ produit déjà la décomposition en idéaux absolument simples à laquelle correspond la décomposition en classes de représentations absolument irréductibles. Les corps indiqués sont donc par là caractérisés comme corps de décomposition.

Il en résulte encore ceci: tous les corps cycliques $\Omega_n$ de degré $2^n$ relativement au corps des nombres rationnels, dans lesquels $(-1)$ peut s'écrire comme somme de deux carrés, sont des corps de décomposition minimaux du corps des quaternions. En effet, un corps de décomposition ne peut pas être réel, à cause de la représentabilité de $(-1)$ comme somme de carrés; d'autre part, le sous-corps de degré $2^{n-1}$ de $\Omega_n$ -- et donc tout sous-corps propre de $\Omega_n$ -- est nécessairement réel, étant l'intersection de $\Omega_n$ avec le corps de tous les nombres algébriques réels. En effet, $\Omega_n$ est un corps galoisien de degré deux relativement à cette intersection, qui doit donc nécessairement coïncider avec le sous-corps de degré $2^{n-1}$.

D'après le théorème d'existence démontré par H. Hasse,\footnote{L'existence de $\Omega_2$ se lit déjà dans la représentation paramétrique des corps cycliques de degré 4, telle qu'elle est indiquée par exemple chez Weber (Kleines Lehrbuch der Algebra, \S{} 93).} de tels corps $\Omega_n$ existent pour tout $n$; les degrés des corps de décomposition minimaux ne sont pas bornés. Ici apparaît même chaque puissance de l'indice comme degré d'un corps de décomposition minimal.\footnote{Après coup, R. Brauer put encore montrer qu'il existe même des corps de décomposition minimaux du corps des quaternions de tout degré pair, donc de tous les degrés possibles pour des corps de décomposition en général. Un tel corps $\mathrm{P}(\vartheta)$ peut être défini pour tout $n\geqq2$ par l'équation
\[
 f(\vartheta)=\vartheta^{2n}+a^2p^2\vartheta^2+b^2p^{2\beta}=0;
 \quad \text{donc}\quad -1=\left(\frac{ap}{\vartheta^{n-1}}\right)^2+
       \left(\frac{bp^\beta}{\vartheta^n}\right)^2,
\]
où l'on peut, par un choix convenable de $a,b,p,\beta$, obtenir que 1. $f(x)$ soit irréductible pour tout $n$; 2. $\mathrm{P}(\vartheta)$ ne possède que l'unique sous-corps $\mathrm{P}(\vartheta^2)$ différent de $\mathrm{P}$; 3. $\mathrm{P}(\vartheta^2)$ possède parmi ses conjugués aussi des corps réels, et ne peut donc pas être corps de décomposition. Un choix possible de $a,b,p,\beta$ est par exemple le suivant:
\[
        f(x)=x^{2n}+(25\cdot9)4x^2+9\cdot64.
\]
La démonstration repose sur une modification d'une méthode de Furtwängler pour construire des équations primitives (Math. Ann. 85, p. 34, \S{} 3), mais elle est assez pénible dans le détail.}

\subsection*{3. Traitement élémentaire des corps de décomposition du corps des quaternions}

Soit $\mathfrak{O}$ la représentation du groupe des quaternions formée des 8 matrices suivantes:\footnote{Cf. Sylvester, Math. Papers Vol. III p. 647, Vol. IV p. 122.}
\[
\pm E=\pm\mat{1&0\\0&1},\qquad
\pm I=\pm\mat{i&0\\0&-i},\qquad
\pm J=\pm\mat{0&1\\-1&0},\qquad
\pm K=\pm\mat{0&i\\ i&0}.
\]
L'anneau de groupe engendré par $\Omega$ est précisément une représentation du corps des quaternions considéré comme système hypercomplexe; nous voulons montrer élémentairement que $\Omega$ est représentable dans un corps $\mathrm{K}$ si et seulement si $-1$ est représentable dans $\mathrm{K}$ comme somme de deux carrés.

\emph{a)} Soit $\mathfrak{O}^*=T^{-1}\mathfrak{O} T$ une représentation rationnelle dans un corps $\mathrm{K}$, et soient $I^*=T^{-1}IT$, $J^*=T^{-1}JT$. Comme $J$ et $J^*$ sont deux matrices semblables rationnelles dans $\mathrm{K}$, il existe aussi une transformation $R$ rationnelle dans $K$ qui transforme $J^*$ en $J$,
\[
        R^{-1}J^*R=J.
\]
Si l'on remplace dès le début $\mathfrak{O}^*$ par la représentation également rationnelle dans $\mathrm{K}$, $R^{-1}\mathfrak{O}^*R$, on reconnaît que l'on peut supposer sans restriction $J^*=J$. Soit
\[
        I^*=\mat{a&b\\ c&d}\qquad (a,b,c,d\text{ nombres de }\mathrm{K}).
\]
De $IJ=-JI$ il résulte que $I^*J^*=-J^*I^*$, et, puisque $J^*=J$, cela donne
\[
        a=-d,
        \qquad b=c.
\]
De $I^2=-E$ il résulte que $I^{*2}=-E$, donc
\[
        \mat{-1&0\\0&-1}=\mat{a&b\\ b&-a}^2
        =\mat{a^2+b^2&0\\0&a^2+b^2},
\]
donc $a^2+b^2=-1$, et par conséquent $-1$ peut effectivement s'écrire dans $\mathrm{K}$ comme somme de deux carrés.

\emph{b)} Dans un corps $\mathrm{K}$, supposons que $-1$ soit représentable comme somme de deux carrés. Si par exemple $-1=a^2+b^2$, posons
\[
        I^*=\mat{a&b\\ b&-a},\qquad
        J^*=\mat{0&1\\-1&0},\qquad
        \mathit{K}^*=I^*J^*.
\]
Comme on le vérifie facilement, $\pm E$, $\pm I^*$, $\pm J^*$, $\pm \mathit{K}^*$ forment alors un groupe rationnel dans $K$, semblable à $\Omega$.\footnote{La caractérisation des corps de décomposition se gagne facilement aussi à l'aide des systèmes de facteurs.}

Il reste à montrer qu'il existe, pour une infinité de $n$, des corps cycliques de degré $2^n$ qui sont des corps de décomposition minimaux du groupe des quaternions.

Soit $p$ un nombre premier rationnel tel que, pour un $r>0$ convenable,
\[
        (*)\qquad 2^r+1\equiv0\pmod p
\]
et soit $2^n$ la plus haute puissance de $2$ qui divise $p-1$.

Nous montrons d'abord qu'il existe des nombres premiers $p$ pour une infinité de valeurs de $n$. Soit $k$ un entier rationnel positif arbitraire, et soit $p$ un diviseur premier de $2^{2^k}+1$. Alors la congruence exigée pour $p$ est satisfaite avec $r=2^k$. En outre, $2\pmod p$ appartient à l'exposant $2^{k+1}$, et par suite, pour la plus haute puissance $2^n$ qui divise $p-1$, le nombre $n$ est plus grand que $k$. Il existe donc en tout cas des $n$ arbitrairement grands auxquels correspondent des nombres premiers $p$.

Soit $\varepsilon$ une racine primitive $p$-ième de l'unité; nous montrons maintenant que, dans $\mathrm{P}(\varepsilon)$, $-1$ peut s'écrire comme somme de 2 carrés
\[
        -1=a^2+b^2=(a+ib)(a-ib)\qquad (a,b\text{ nombres de }\mathrm{P}(\varepsilon)).
\]
Cela équivaut à dire que, dans le corps des racines $(4p)$-ièmes de l'unité $\mathrm{P}(\varepsilon,i)=\mathrm{P}(\varepsilon\cdot i)$, il existe des nombres dont la norme relative par rapport à $\mathrm{P}(\varepsilon)$ est précisément $-1$.

Le nombre $1+i\varepsilon^\nu$ $(\nu=1,2,\ldots,p-1)$ a pour norme relative $(1+i\varepsilon^\nu)(1-i\varepsilon^\nu)=1+\varepsilon^{2\nu}$; par conséquent, tous les nombres $1+\varepsilon^\nu$ $(\nu=1,2,\ldots,p-1)$ sont des normes relatives de nombres de $\mathrm{P}(\varepsilon i)$ par rapport à $\mathrm{P}(\varepsilon)$ comme corps de base. Il en va donc de même pour
\[
\Pi=\prod_{\nu=0}^{r-1}(1+\varepsilon^{2^\nu})
=(1+\varepsilon)(1+\varepsilon^2)(1+\varepsilon^4)\ldots(1+\varepsilon^{2^{r-1}})
=\sum_{\lambda=0}^{2^r-1}\varepsilon^\lambda.
\]
Si l'on ajoute encore à $\Pi$ le terme $\varepsilon^{2^r}=\varepsilon^{-1}=\varepsilon^{p-1}$ (à cause de $(*)$), alors $\Pi+\varepsilon^{p-1}$ devient une somme de morceaux $1+\varepsilon+\varepsilon^2+\ldots+\varepsilon^{p-1}=0$, et donc
\[
        \Pi=-\varepsilon^{-1}.
\]
Or $\varepsilon$ est aussi la norme relative d'un nombre de $\mathrm{P}(\varepsilon i)$, à savoir de $\varepsilon^{\frac{p+1}{2}}$. Il en résulte de même pour
\[
        \Pi\cdot\varepsilon=-1.
\]
Ainsi $-1$ est représentable dans $\mathrm{P}(\varepsilon)$ comme somme de 2 carrés;\footnote{Pour les nombres premiers de la forme $n8\pm3$, cette considération se trouve déjà chez E. Landau, Über die Darstellung definiter Funktionen durch Quadrate, Math. Ann. t. 62, p. 272--285.} $\mathrm{P}(\varepsilon)$ est un corps de décomposition pour le groupe des quaternions.

Soit $\mathrm{K}$ le corps cyclique de degré $2^n$ contenu dans $\mathrm{P}(\varepsilon)$; nous affirmons que $\mathrm{K}$ est lui aussi corps de décomposition.\footnote{Le choix du corps de décomposition cyclique de degré $2^n$ comme sous-corps d'un corps cyclotomique est calqué sur l'exemple de Hasse.} Soit $m$ l'indice du corps des quaternions relativement à $\mathrm{K}$ comme corps de base; alors $m=1$ ou $m=2$. Mais puisqu'il existe un corps de décomposition $\mathrm{P}(\varepsilon)$ de degré relatif impair $\frac{p-1}{2^n}$ par rapport à $\mathrm{K}$, le second cas est impossible; donc $m=1$. Cela signifie précisément que $\mathrm{K}$ est corps de décomposition.

Il existe donc, pour une infinité de $n$, des corps cycliques de degré $2^n$ qui sont des corps de décomposition minimaux.

\begin{center}
\rule{7em}{0.4pt}
\end{center}
\clearpage

\section*{\NoetherNavAnchor{P33}{33. Grandeurs hypercomplexes et théorie des représentations dans une conception arithmétique}33. Grandeurs hypercomplexes et théorie des représentations dans une conception arithmétique}
\emph{Atti Congresso Bologna 2 (1928), p. 71--73.}

\setcounter{footnote}{0}
\begingroup
\renewcommand{\thefootnote}{(\arabic{footnote})}

\begin{center}
\textsc{Emmy Noether} (Göttingen -- Allemagne)\\[-0.1em]
\rule{3em}{0.4pt}\\[0.8em]
\textsc{Grandeurs hypercomplexes et théorie des représentations}\\
\textsc{dans une conception arithmétique}\footnote{Une exposition détaillée paraîtra dans la \emph{Mathematische Zeitschrift}; une suite y abordera aussi la théorie de Galois des corps gauches. Voir également, pour les \og{}corps de décomposition\fg{}, R. Brauer et E. Noether, \emph{Ber. Berl. Ak.} 1927, et, pour les théorèmes généraux de structure, la communication de G. Köthe dans les présents \emph{Atti}.}
\end{center}

Je voudrais vous montrer comment la théorie des représentations -- en particulier celle des représentations des groupes -- peut être conçue comme une théorie des classes de modules et d'idéaux, et comment celles-ci se subordonnent à leur tour à une notion généralisée de groupe, celle de groupe avec opérateurs.\footnote{Les groupes avec opérateurs remontent à W. Krull (\emph{Math. Zeitschr.} 23) et O. Schmidt (\emph{Math. Zeitschr.} 29). Le traitement donné ici du problème de la réduction d'une représentation est lui aussi dû à W. Krull, dans le cas où le corps de représentation est commutatif (\emph{Theorie und Anwendung der verallgemeinerten Abelschen Gruppen}, \emph{Heidelberger Ber.} 1926).}

Par \emph{représentation d'un groupe} dans un corps donné $P$, on entend, comme d'usage, la donnée suivante: aux éléments $a,b$ du groupe $\mathfrak G$ sont associées des matrices $A,B$ à coefficients dans $P$, de telle sorte qu'au produit $ab$ corresponde le produit $AB$; le système des matrices est donc une image homomorphe du groupe. Cette représentation fournit en même temps une représentation de l'\og{}anneau de groupe\fg{}, formé de tous les \og{}nombres du groupe\fg{}, c'est-à-dire de toutes les combinaisons linéaires finies $a\alpha+b\beta+....+c\gamma$, où $a,b,....,c$ parcourent les éléments du groupe et $\alpha,\beta,....,\gamma$ les éléments de $P$. Les éléments $a,b,....,c$ sont regardés comme linéairement indépendants, et la multiplication est définie par celle du groupe et les lois de calcul. La représentation de l'anneau de groupe est en outre un homomorphisme pour l'addition; la représentation des groupes relève ainsi du problème général de la \emph{représentation des anneaux}, où l'on exige un homomorphisme tant pour l'addition que pour la multiplication.

Deux problèmes sont ici essentiels: le premier est celui de la \emph{réduction d'une représentation donnée} et de l'unicité des constituants irréductibles qui apparaissent; le second porte sur la totalité des représentations possibles -- éventuellement seulement des représentations irréductibles -- d'un anneau donné dans un corps gauche donné.

Le premier problème est de loin le plus simple, ce qui se voit déjà au fait qu'aucune hypothèse spéciale n'est nécessaire ni sur l'anneau à représenter ni sur le corps gauche de représentation; le fait que la représentation se fasse par des matrices d'un degré fixe fournit toutes les hypothèses de finitude nécessaires. Le problème se traite complètement au moyen de la théorie des groupes avec opérateurs, sur laquelle je vais donc revenir brièvement:

Un groupe $\mathfrak{B}$ s'appelle groupe avec domaine d'opérateurs lorsqu'est donné en même temps un système de symboles $\varTheta,H$ -- les opérateurs -- tel que la liaison $\varTheta(a),H(a)$ soit définie de façon univoque pour tout $\mathrm{a}$ de $\mathfrak{B}$ et donne un élément de $\mathfrak{B}$, et que la loi distributive soit satisfaite: $\varTheta(ab)=\varTheta(a)\varTheta(b)$. Deux groupes munis du même domaine d'opérateurs sont dits isomorphes comme groupes à opérateurs lorsqu'ils sont isomorphes au sens usuel et que, de la correspondance de $a$ et $\bar a$, résulte en outre celle de $\varTheta(a)$ et $\varTheta(\bar a)$, de $H(a)$ et $H(\bar a)$, etc. Si l'on n'admet alors comme sous-groupes que ceux qui admettent eux-mêmes le domaine d'opérateurs (de sorte qu'un groupe est dit simple lorsqu'il n'a aucun diviseur normal propre admissible en dehors de l'unité), le théorème de Jordan-Hölder sur les suites de composition subsiste sous la forme suivante: si un groupe avec domaine d'opérateurs possède au moins une suite de composition, le système des groupes quotients est déterminé de façon unique, à l'ordre et à l'isomorphisme comme groupes à opérateurs près. Et, sous la même hypothèse, vaut le théorème de la décomposition en facteurs directement indécomposables, unique à isomorphisme comme groupes à opérateurs près.

Le lien avec le problème de représentation est donné par le fait que, parmi les groupes avec opérateurs, figurent en particulier les idéaux et les modules, considérés comme groupes abéliens relativement à l'addition, tandis que les opérateurs sont donnés par la multiplication par les éléments de l'anneau. Mais toute représentation est portée par un module de représentation, c'est-à-dire par un module relativement à l'anneau et au corps de représentation. Plus exactement, chaque classe de modules, c'est-à-dire chaque classe de modules de représentation isomorphes comme modules à opérateurs, détermine la même représentation. Ainsi la question de l'unicité de la réduction d'une représentation reçoit sa réponse par le théorème des suites de composition et par le théorème sur le produit direct. Appliquée au module de représentation, la suite de composition donne, pour chaque matrice, une réduction de la forme:
\[
\begin{pmatrix}
B_1&0&0\\
&B_2&0\\
A_{ik}&....&B_r
\end{pmatrix}
\]
où les matrices diagonales correspondant aux groupes quotients engendrent une représentation ``irréductible'', c'est-à-dire n'admettant plus elle-même une telle réduction.

Mais, puisque les classes de ces groupes quotients sont déterminées de façon univoque, il en va de même des représentations diagonales (univoquement au sens de l'équivalence des représentations). De manière analogue, le théorème du produit direct donne l'unicité de la ``décomposition'' en constituants ``indécomposables'':
\[
\begin{pmatrix}
C_1&&&\\
&C_2&&\\
&&....&\\
&&&C_s
\end{pmatrix},
\]
où chaque $C_i$ peut ensuite encore se réduire de façon univoque d'après le théorème des suites de composition.

Je voudrais esquisser le second problème dans le cas de la représentation des systèmes hypercomplexes, plus généralement des anneaux qui satisfont au ``théorème de double chaîne''. Le résultat est que toute classe de modules irréductibles (simples) est une classe d'idéaux, donc que les suites de composition formées d'idéaux (unilatères) fournissent déjà, par leurs groupes quotients, toutes les représentations irréductibles. Plus précisément, il suffit des suites de composition dans l'anneau quotient par le radical, ce dernier devenant un anneau ``complètement réductible''. Le problème est ainsi ramené à l'étude de structure des anneaux complètement réductibles; il se révèle donc conforme au problème de prendre comme point de départ les représentations dans les corps gauches des endomorphismes. Il faut entendre par là ceci: tous les idéaux simples (à droite) d'une composante simple bilatère d'un tel anneau appartiennent à la même classe, et possèdent donc le même anneau des endomorphismes, qui devient un corps gauche en raison de la simplicité des idéaux, et est en même temps l'anneau des endomorphismes des idéaux simples à gauche. L'anneau simple lui-même devient isomorphe au système de toutes les matrices dans le corps gauche des endomorphismes; et de cette représentation naissent toutes les représentations relativement à un sous-corps, donc en particulier relativement au domaine des coefficients du système hypercomplexe, lorsqu'on remplace les éléments du corps gauche des endomorphismes eux-mêmes par des matrices associées. Ainsi, le théorème de structure de Wedderburn connu dans le cas des systèmes hypercomplexes est réinterprété par la notion de corps gauche des endomorphismes issue de la théorie générale des groupes, et reconnu en même temps comme source de la théorie des représentations des systèmes hypercomplexes. Surgissent alors les nouvelles questions de la théorie de Galois des corps gauches, étroitement liées à la question des ``corps de décomposition''. En même temps s'ouvre une perspective sur des théorèmes de structure pour les anneaux généraux non complètement réductibles, au moyen des anneaux d'endomorphismes des idéaux indécomposables. La théorie des groupes, dans son extension aux groupes les plus généraux avec opérateurs, se montre ici encore comme principe d'organisation.

\endgroup
\clearpage

\providecommand{\mG}{\mathfrak{G}}
\providecommand{\mA}{\mathfrak{A}}
\providecommand{\mB}{\mathfrak{B}}
\providecommand{\mC}{\mathfrak{C}}
\providecommand{\mD}{\mathfrak{D}}
\providecommand{\mE}{\mathfrak{E}}
\providecommand{\mH}{\mathfrak{H}}
\providecommand{\mK}{\mathfrak{K}}
\providecommand{\mM}{\mathfrak{M}}
\providecommand{\mN}{\mathfrak{N}}
\providecommand{\mR}{\mathfrak{R}}
\providecommand{\mS}{\mathfrak{S}}
\providecommand{\mT}{\mathfrak{T}}
\providecommand{\mU}{\mathfrak{U}}
\providecommand{\mO}{\mathfrak{O}}
\providecommand{\mo}{\mathfrak{o}}
\providecommand{\OO}{\Omega}
\providecommand{\Th}{\Theta}
\providecommand{\iso}{\simeq}
\providecommand{\normlhd}{\triangleleft}

\setcounter{footnote}{0}

\section*{\NoetherNavAnchor{P34}{34. Grandeurs hypercomplexes et théorie des représentations}34. Grandeurs hypercomplexes et théorie des représentations}
\emph{Math. Zs. 30 (1929), p. 641--692.}

\begin{center}
{\Large\bfseries Grandeurs hypercomplexes et théorie des représentations\footnote{Ce qui suit est une élaboration libre, faite par B. L. van der Waerden, de mon cours du semestre d'hiver 1927/28. Nous avons préparé ensemble le texte pour l'impression. Je dois encore à B. L. van der Waerden une série de remarques critiques.}.}\\[1em]
Par\\[0.5em]
Emmy Noether à Göttingen.\\[0.8em]
\rule{3em}{0.4pt}
\end{center}

\subsection*{Introduction}

Les plus importants théorèmes généraux sur les systèmes hypercomplexes remontent à Molien (Math. Annalen, t. 41, et Rapports de Dorpat 1897). De façon essentiellement indépendante, Frobenius développa peu après, d'un point de vue unifié, la théorie des systèmes hypercomplexes et de leurs représentations -- en particulier la théorie des représentations des groupes finis. La base en est la notion, due à Dedekind, du déterminant de groupe, plus généralement du déterminant d'un système hypercomplexe quelconque. Frobenius montre qu'aux différents facteurs irréductibles de ce déterminant (le domaine des coefficients étant toujours le corps de tous les nombres complexes) correspondent les différentes classes de représentations irréductibles, et que toutes les classes de représentations irréductibles sont ainsi épuisées. Une telle classe de représentations est entièrement caractérisée par son ``système de caractères''; dans le cas des groupes finis, ce système de caractères naît de la décomposition du déterminant du système hypercomplexe commutatif dérivé des classes d'éléments conjugués du groupe. Ce déterminant se décompose en facteurs linéaires, avec les caractères comme coefficients -- généralisation directe du résultat obtenu d'abord par Dedekind, selon lequel le déterminant de groupe d'un groupe abélien fini se décompose en facteurs linéaires dont les coefficients sont les différents caractères du groupe abélien (correspondance avec Frobenius).

Ces résultats, simples et transparents du point de vue conceptuel, sont toutefois obtenus chez Frobenius par un calcul laborieux. Le développement ultérieur visa à donner de ces résultats une dérivation simplifiée, et en même temps à les étendre lorsque l'on prend pour domaine des coefficients un corps arbitraire. Ce développement ultérieur s'est accompli par des voies tout à fait séparées pour les systèmes hypercomplexes et pour la théorie des représentations. Les systèmes hypercomplexes reçurent un traitement arithmétique chez Wedderburn: tout système hypercomplexe est, de façon unique, somme directe d'anneaux bilatéralement directement indécomposables; pour les systèmes sans radical, ces anneaux indécomposables deviennent isomorphes au système de toutes les matrices à $n$ lignes dont les éléments appartiennent à un corps associé -- essentiellement déterminé de façon unique -- et non nécessairement commutatif.\footnote{Cf. les indications bibliographiques chez Dickson, \emph{Algebras and their Arithmetics} (édition allemande: \emph{Algebren und ihre Zahlentheorie}, Zürich 1927).} La théorie des représentations reçut une fondation élémentaire par Burnside et I. Schur, qui, indépendamment du système hypercomplexe, partirent directement d'une représentation donnée et opérèrent avec des théorèmes de matrices. Schur traita en particulier la question des corps de nombres de plus petit degré dans lesquels une représentation irréductible sur un corps donné se décompose absolument. Ces constituants absolument irréductibles appartiennent à un nombre fini de classes de représentations conjuguées. Le degré des corps de nombres cherchés est égal au produit du nombre de ces classes par l'indice (nombre des constituants dans chacune des classes conjuguées). L'auxiliaire principal de la théorie de Schur est le système de toutes les matrices commutant avec une représentation irréductible.\footnote{Cf. les indications bibliographiques au chapitre 10 de la théorie des groupes de Speiser.}

Dans la fondation purement arithmétique qui suit, les grandeurs hypercomplexes et la théorie des représentations apparaissent de nouveau comme un tout unifié -- comme cas particulier d'une théorie générale des anneaux non commutatifs qui ne satisfont qu'à certaines conditions de finitude. Il s'agit plus précisément d'une théorie des classes de modules et d'idéaux relativement à ces anneaux, avec pour résultat principal que les classes de modules irréductibles sont déjà épuisées par les classes d'idéaux correspondantes; en particulier, pour les anneaux sans radical, toutes les classes de modules se décomposent en classes irréductibles (deviennent complètement réductibles). C'est l'équivalent arithmétique du résultat de Frobenius selon lequel les facteurs irréductibles du déterminant régulier de groupe (respectivement de système) épuisent déjà la totalité des classes de représentations irréductibles, et selon lequel, pour les systèmes sans radical, vaut la complète réductibilité des représentations. La considération du déterminant de système et de sa décomposition peut en effet être conçue comme un passage à la norme, tandis qu'ici la décomposition directe en somme des idéaux et leurs suites de composition sont traitées directement. Quant à la théorie des représentations, elle devient, grâce à la notion de module de représentation, une théorie des classes de modules.

L'introduction des classes de modules et d'idéaux est fondée purement sur la théorie des groupes. Les modules et les idéaux sont considérés comme des groupes abéliens relativement à l'addition, avec cette restriction qu'ils admettent certaines multiplications par des éléments d'anneaux: ils forment des ``groupes avec opérateurs'' (\S{}1). Pour de tels groupes, à la place de l'isomorphisme ordinaire apparaît l'``isomorphisme d'opérateurs'' (\S{}2); les groupes isomorphes entre eux comme groupes à opérateurs (dans un domaine fixe) sont réunis en une classe -- notion de classe qui, par exemple dans le cas des idéaux d'un corps de nombres, coïncide avec la notion usuelle. La théorie des groupes avec opérateurs remonte à W. Krull et O. Schmidt (cf. note 6); elle est développée systématiquement au chapitre I. Les théorèmes qui y restent aussi valables sur les suites de composition, le produit direct (respectivement la somme directe) et les groupes complètement réductibles -- théorèmes d'unicité au sens de la division en classes mentionnée plus haut -- constituent le fondement de tout ce qui suit.

Ils donnent d'abord les théorèmes généraux d'unicité pour les constituants diagonaux irréductibles de représentations quelconques (chapitre III, \S{}16), sur la base de la correspondance biunivoque entre les classes de représentations et les classes de modules de représentation, donc les classes de groupes avec opérateurs. La question plus étendue de la totalité des représentations exige une étude plus précise de la structure des anneaux à représenter, donc, dans le cas particulier, des systèmes hypercomplexes. Au chapitre II, les résultats de Wedderburn sont obtenus de nouveau et poursuivis, selon une conception fondée sur la théorie des groupes et plaçant au premier plan les idéaux unilatères. Il apparaît alors que le ``théorème de chaîne des multiples'' pour les idéaux à droite, ou la ``condition minimale'' identique à celui-ci (dans tout ensemble d'idéaux à droite il existe au moins un élément minimal -- appartenant à cet ensemble), suffit comme condition de finitude.\footnote{Les raisonnements de Wedderburn se transposent si l'on suppose le ``théorème de double chaîne''. Cf. E. Artin, Zur Theorie der hyperkomplexen Zahlen, Hamb. Abh. 1927. Cf. aussi A. Suschkewitsch, Über die endlichen Gruppen ohne das Gesetz der eindeutigen Umkehrbarkeit, Math. Annalen 99 (1928), p. 30--50, où sont développés des théorèmes de structure parallèles pour des domaines (finis) avec une seule opération associative non inversible. La scission du ``noyau'' chez Suschkewitsch en groupes à droite et à gauche correspond à la représentation en somme d'un anneau bilatéralement simple avec élément unité par des idéaux à droite et à gauche.} On obtient l'identité des anneaux sans radical satisfaisant à la condition minimale avec les anneaux (à droite) complètement réductibles avec élément unité. Dans de tels anneaux, tous les idéaux simples à droite d'un anneau bilatéralement indécomposable appartiennent à la même classe, possèdent donc -- comme groupes avec opérateurs -- le même anneau des endomorphismes, lequel devient un corps gauche en raison de la simplicité des idéaux. C'est ce corps gauche des endomorphismes de la classe qui sert d'intermédiaire à la représentation matricielle trouvée par Wedderburn.

Cette représentation matricielle dans le corps gauche des endomorphismes résout en même temps, dans le cas complètement réductible, le problème de la représentation, dès que l'on exige une représentation relativement au corps gauche des endomorphismes ou à ses sous-corps -- donc en particulier relativement au domaine des coefficients d'un système hypercomplexe. Il n'y a -- c'est une conséquence directe du théorème qui ramène les classes de modules aux classes d'idéaux -- pas d'autres représentations irréductibles que celle de Wedderburn et celles qui en naissent lorsque les éléments du corps gauche des endomorphismes sont remplacés, de la manière naturelle, par des matrices à éléments dans un sous-corps.

Ainsi, dans le cas complètement réductible, il y a autant de classes différentes de représentations irréductibles que de classes d'idéaux; leur nombre coïncide avec le nombre des différents anneaux bilatéralement indécomposables, donc simples. Ce nombre est enfin de nouveau identique au nombre des composantes indécomposables du centre, lesquelles deviennent des corps commutatifs. Dans le cas des systèmes hypercomplexes dont le domaine de coefficients est algébriquement clos, ces derniers sont entièrement déterminés par les différents homomorphismes d'anneaux (représentations irréductibles) du centre; or ceux-ci sont, à un facteur numérique près, les caractères.

Les résultats obtenus donnent en même temps une vue d'ensemble sur les représentations irréductibles des systèmes avec radical, dans la mesure où celles-ci se ramènent aux représentations de l'anneau quotient par le radical, qui devient un système sans radical.

Comme on le montrera plus tard, pour les systèmes hypercomplexes sans radical, la théorie des ``corps de décomposition'' découle elle aussi du corps gauche des endomorphismes, c'est-à-dire la théorie des extensions commutatives du domaine des coefficients dans lesquelles les représentations se décomposent en représentations absolument irréductibles -- autrement dit, dans lesquelles a lieu une décomposition directe en somme d'idéaux à droite absolument simples. Les corps de décomposition sont identiques à ceux du corps gauche des endomorphismes, et tous les corps de décomposition de plus petit degré sont isomorphes aux sous-corps commutatifs maximaux du corps gauche des endomorphismes. Le lien avec les recherches de Schur mentionnées plus haut est établi par le fait que les transposées des matrices commutant avec la représentation fournissent précisément une représentation du corps gauche des endomorphismes.\footnote{Cf. à ce sujet R. Brauer--E. Noether, Über minimale Zerfällungskörper irreduzibler Darstellungen, Sitz.-Ber. Preuß. Akad. Wiss. 1927, p. 221.} Ces théorèmes seront développés systématiquement dans le cadre d'une théorie de Galois des corps gauches.

Les notions posées à la base -- homomorphisme d'opérateurs et anneau des endomorphismes -- donnent aussi la structure des anneaux généraux avec radical qui satisfont seulement à la condition minimale; cela sera exécuté d'un autre côté.

\section*{Chapitre I. Fondements de théorie des groupes}

\subsection*{\S{} 1. Groupes avec opérateurs}

Soit $\mG$ un groupe (fini ou infini), d'éléments $a,b,\ldots$. Par domaine d'opérateurs $\Omega$ pour le groupe $\mG$, on entend un ensemble de nouveaux symboles $\mathsf{H},\Theta,\ldots$, tel qu'à chaque $a$ de $\mG$ et à chaque $\Theta$ de $\Omega$ corresponde un $\Theta a$ de $\mG$ défini de façon univoque, et tel que la loi distributive soit satisfaite:
\[
        \Theta(ab)=\Theta a\cdot\Theta b.
\]
Ainsi chaque opérateur définit un homomorphisme du groupe dans lui-même, le groupe étant appliqué sur lui-même ou sur un sous-groupe. L'élément unité passe alors dans l'élément unité, l'inverse dans l'inverse.

Un domaine d'opérateurs est dit \emph{absolu} lorsque des opérateurs différents définissent aussi des homomorphismes différents. D'un domaine d'opérateurs arbitraire on obtient un domaine absolu en identifiant tous les éléments qui engendrent le même homomorphisme. Un domaine d'opérateurs absolu est une image biunivoque d'un sous-ensemble de l'ensemble de tous les homomorphismes du groupe dans lui-même.

Un \emph{sous-groupe admissible} $\mH$ de $\mG$ -- relativement à un domaine d'opérateurs fixe $\Omega$ -- est un sous-groupe qui admet le domaine d'opérateurs, c'est-à-dire tel que $\Theta a$ soit dans $\mH$ pour tout $a$ de $\mH$ et tout $\Theta$ de $\Omega$.

\emph{Exemples.} 1. Les opérateurs soient les automorphismes intérieurs: $\Theta a=c^{-1}ac$. Les sous-groupes admissibles sont les diviseurs normaux.

2. Les opérateurs soient tous les automorphismes. Les sous-groupes admissibles sont les ``sous-groupes caractéristiques'', qui restent en eux-mêmes sous tous les automorphismes.\footnote{Les notions expliquées ici proviennent de W. Krull, Über verallgemeinerte endliche Abelsche Gruppen, Math. Zeitschr. 23 (1925), p. 161--196, et de O. Schmidt, Über unendliche Gruppen mit endlicher Kette, Math. Zeitschr. 29 (1928), p. 34--41.}

3. Soit $\mG$ un anneau, c'est-à-dire un groupe abélien relativement à l'addition, où est aussi définie une multiplication ayant les propriétés
\[
        r(a+b)=ra+rb,\qquad (a+b)r=ar+br,\qquad ab\cdot c=a\cdot bc.
\]
Chaque élément $r$ définit à la fois deux opérateurs: les opérateurs $rx$ et $xr$. Les sous-groupes admis sont les ``idéaux'' $\mathfrak a$, à savoir: à gauche, ceux qui admettent les opérations $rx$: $r\mathfrak a\subseteqq\mathfrak a$; à droite, ceux qui admettent les opérations $xr$: $\mathfrak a r\subseteqq\mathfrak a$; bilatères, ceux qui admettent les deux opérations. Tous les idéaux deviennent triviaux $(=\{0\}\text{ ou }=\mG)$ lorsque l'anneau $\mG$ est un corps, c'est-à-dire lorsque $\mG-\{0\}$ est un groupe relativement à la multiplication.\footnote{On ne suppose donc la loi commutative de la multiplication ni pour les anneaux ni pour les corps.}

4. \emph{Modules relativement à un anneau $\mo$.} Soit $\mo$ un anneau, et $\mM$ un groupe abélien écrit additivement. On suppose donnée une multiplication $r\cdot a$ des éléments de $\mo$ par ceux de $\mM$; le produit doit appartenir de nouveau à $\mM$. On exige:
\[
\left.
\begin{aligned}
 r(a+b)&=ra+rb,\\
 rs\cdot a&=r\cdot sa,\\
 (r+s)a&=ra+sa,
\end{aligned}
\right\}\qquad r,s\in\mo,\quad a,b\in\mM.
\]
Alors $\mM$ s'appelle un $\mo$-module; plus précisément, puisque les multiplicateurs de $\mo$ s'écrivent à gauche, un $\mo$-module à gauche. L'anneau $\mo$ est en même temps domaine d'opérateurs. S'il est absolu, c'est-à-dire si des éléments différents de l'anneau donnent aussi des opérations différentes, on parle d'un \emph{domaine absolu de multiplicateurs} pour le module $\mM$.

Comme plus haut, on peut passer d'un domaine de multiplicateurs arbitraire au domaine absolu en identifiant les éléments qui engendrent la même opération. Le domaine absolu de multiplicateurs est de nouveau un anneau. Les sous-groupes admissibles d'un module $\mM$ s'appellent \emph{sous-modules}. Si, spécialement, $\mo=\mM$, on revient aux idéaux à gauche.\footnote{De la même manière, on exige pour les modules à droite:
\[
(a+b)r=ar+br,\qquad a\cdot rs=ar\cdot s,\qquad a(r+s)=ar+as.
\]}

5. \emph{Bimodules.} Si $\mM$ est à la fois un $\mo$-module à gauche et un $\mo'$-module à droite, et si en outre, pour $a$ dans $\mo$, $\alpha$ dans $\mM$, $a'$ dans $\mo'$, on a toujours
\[
        a\cdot\alpha a'=a\alpha\cdot a',
\]
alors $\mM$ s'appelle un bimodule.

6. \emph{Anneau des endomorphismes d'un groupe abélien.}\footnote{Cf. A. Châtelet, \emph{Les Groupes Abéliens finis}, Paris, Gauthier-Villars, 1925, p. 99.} Pour les homomorphismes dans lui-même d'un groupe abélien (écrit additivement), on peut définir une addition et une multiplication par les formules
\[
        (\mathsf{H}+\Theta)a=\mathsf{H}a+\Theta a,\qquad (\mathsf{H}\Theta)a=\mathsf{H}(\Theta a).
\]
Les opérations ainsi définies représentent manifestement de nouveau des homomorphismes, et appartiennent donc de nouveau au système. Le système forme un anneau, et le groupe abélien peut être considéré, d'après 4., comme un module sur cet anneau.

Par ``sous-groupes'', on entendra dorénavant toujours des sous-groupes admissibles. L'intersection $\mA\cap\mB$ de deux sous-groupes admissibles est de nouveau un sous-groupe admissible. De même le produit $\mA\mB$, pourvu qu'au moins l'un des deux sous-groupes $\mA,\mB$ soit un diviseur normal.

\providecommand{\mG}{\mathfrak{G}}
\providecommand{\mA}{\mathfrak{A}}
\providecommand{\mB}{\mathfrak{B}}
\providecommand{\mN}{\mathfrak{N}}
\providecommand{\mo}{\mathfrak{o}}

\subsection*{\S{} 2. Les théorèmes d'isomorphie}

Une application d'un groupe $\mG$ sur un groupe $\overline{\mG}$ -- où $\mG$ et $\overline{\mG}$ doivent posséder le même domaine d'opérateurs -- s'appelle un homomorphisme d'opérateurs si, premièrement, elle est un homomorphisme au sens ordinaire ($ab\rightarrow \bar a\bar b$), et si, deuxièmement, dès que $a$ est appliqué sur $\bar a$, $\Theta a$ est appliqué sur $\Theta\bar a$.\footnote{On peut étendre la notion d'homomorphisme d'opérateurs en ce sens que les groupes $\mG,\overline{\mG}$ possèdent des domaines d'opérateurs séparés, l'homomorphisme appliquant alors non seulement $\mG$ sur $\overline{\mG}$, mais aussi les domaines d'opérateurs l'un sur l'autre, de telle sorte que, si $a$ est appliqué sur $\bar a$ et $\Theta$ sur $\overline\Theta$, alors $\overline\Theta\bar a$ soit l'image de $\Theta a$. Dans cette formulation générale, la notion d'homomorphisme d'opérateurs comprend aussi celle d'homomorphisme d'anneaux; cf. \S{}8.} Notation: $\mG\sim\overline{\mG}$. Si la correspondance est biunivoque, elle s'appelle une isomorphie d'opérateurs. Notation: $\mG\simeq\overline{\mG}$. On démontre, comme pour les groupes ordinaires, le théorème d'homomorphie:

\emph{Si $\overline{\mG}$ est l'image de $\mG$ par un homomorphisme compatible avec les opérateurs, alors $\overline{\mG}$ est isomorphe, comme groupe à opérateurs, à un groupe quotient $\mG/\mN$, où $\mN$ est un diviseur normal admissible, formé de tous les éléments de $\mG$ auxquels correspond l'unité dans $\overline{\mG}$. Réciproquement, tout diviseur normal admissible $\mN$ définit un groupe quotient $\mG/\mN$ qui admet le domaine d'opérateurs de $\mG$ et qui est l'image de $\mG$ par un homomorphisme compatible avec les opérateurs.}

Un groupe s'appelle \emph{simple} lorsqu'il ne possède pas de diviseurs normaux en dehors de lui-même et du groupe unité. Tout homomorphisme d'un groupe simple est soit un isomorphisme, soit associe à chaque élément l'élément unité.

\emph{Premier théorème d'isomorphie.} \emph{Soit $\overline{\mG}$ une image homomorphe de $\mG$, soit $\overline{\mA}$ un diviseur normal de $\overline{\mG}$, et soit $\mA$ la totalité des éléments de $\mG$ auxquels sont associés des éléments de $\overline{\mA}$. Alors $\mA$ est de nouveau un diviseur normal, et}
\[
        \overline{\mG}/\overline{\mA}\simeq \mG/\mA. \tag{1}
\]
\emph{Démonstration.} $\mG\sim\overline{\mG}$, $\overline{\mG}\sim\overline{\mG}/\overline{\mA}$, donc $\mG\sim\overline{\mG}/\overline{\mA}$. Ainsi $\overline{\mG}/\overline{\mA}$ est isomorphe à un groupe quotient dans $\mG$; le diviseur normal correspondant est la totalité des éléments auxquels correspond l'unité dans $\overline{\mG}/\overline{\mA}$; or c'est précisément $\mA$.

\emph{Complément.} Si, d'après le théorème d'homomorphie, on pose $\overline{\mG}=\mG/\mN$, alors $\mA$ contiendra certainement $\mN$. À partir de $\mA$, on peut retrouver $\overline{\mA}$: $\overline{\mA}=\mA/\mN$. La correspondance $\mA\leftrightarrow\overline{\mA}$ est biunivoque. La formule (1) peut aussi s'écrire:
\[
        (\mG/\mN)/(\mA/\mN)\simeq \mG/\mA.
\]

\emph{Deuxième théorème d'isomorphie.} \emph{Soit $\mA$ un sous-groupe, $\mB$ un diviseur normal de $\mG$. Alors $\mA\cap\mB$ est un diviseur normal dans $\mA$, et}
\[
        \mA\mB/\mB\simeq \mA/\mA\cap\mB.
\]
\emph{Démonstration.} Dans l'homomorphie $\mG\sim\mG/\mB$, aux éléments $a$ de $\mA$ en particulier sont associées certaines classes $a\mB$ modulo $\mB$, qui composent ensemble le groupe $\mA\mB/\mB$. Donc $\mA\sim\mA\mB/\mB$, d'où l'on déduit l'assertion d'après le théorème d'homomorphie.

Un homomorphisme compatible avec les opérateurs d'un groupe $\mG$ dans lui-même ou dans un sous-groupe s'appelle un endomorphisme de $\mG$. Si $\mG$ est spécialement un groupe abélien (avec opérateurs), écrit additivement, et si l'on définit comme plus haut (\S{}1, exemple 6) la somme et le produit des homomorphismes, alors les endomorphismes du groupe compatibles avec les opérateurs forment un anneau, l'anneau des endomorphismes.

\emph{L'anneau des endomorphismes d'un groupe abélien simple (avec opérateurs) est un corps gauche.}

\emph{Démonstration.} Tout homomorphisme applique le groupe soit sur lui-même, soit sur le groupe nul (puisqu'il n'y a pas d'autres sous-groupes admissibles). Les homomorphismes qui n'appliquent pas tout sur le zéro sont, d'après la remarque faite plus haut sur les groupes simples, des isomorphismes, et les applications isomorphes d'un groupe sur lui-même forment un groupe. Ainsi, après omission de l'opérateur nul, l'anneau devient un groupe; l'anneau lui-même est donc un corps.

Si, spécialement, $\mG$ est un $\mo$-module à gauche, et si l'on écrit les homomorphismes d'opérateurs comme opérateurs à droite, alors $\mG$ devient un \emph{bimodule relativement à $\mo$ à gauche et à l'anneau des endomorphismes à droite}. Car le fait que les nouveaux opérateurs $\Gamma$ aient été choisis comme homomorphismes s'exprime par les formules
\[
        (a+b)\Gamma=a\Gamma+b\Gamma,
        \qquad ra\cdot\Gamma=r\cdot a\Gamma,
\]
qui (avec les autres formules déjà interprétées plus haut) caractérisent le bimodule.

Réciproquement, si un bimodule est donné, ces mêmes formules expriment que chaque élément du domaine des multiplicateurs à droite induit un homomorphisme d'opérateurs relativement aux opérateurs à gauche, et réciproquement.

\providecommand{\mG}{\mathfrak{G}}
\providecommand{\mA}{\mathfrak{A}}
\providecommand{\mB}{\mathfrak{B}}
\providecommand{\mE}{\mathfrak{E}}
\providecommand{\mH}{\mathfrak{H}}
\providecommand{\mo}{\mathfrak{o}}

\subsection*{\S{} 3. Suites de composition}

S'il existe dans $\mG$ une suite finie de sous-groupes admissibles
\[
        \mG\supset \mA_1\supset \mA_2\supset \ldots\supset \mA_r=\mE\tag{2}
\]
($\mE$ est le groupe constitué du seul élément unité $e$), telle que chaque $\mA_i$ soit un diviseur normal dans le terme précédent, et qu'il soit impossible d'insérer entre deux termes consécutifs de la suite un autre sous-groupe ayant la même propriété, alors cette suite s'appelle une suite de composition.

Les groupes quotients $\mG/\mA_1,\ \mA_1/\mA_2,\ldots,\ \mA_{r-1}/\mE$ s'appellent les facteurs de composition. Ce sont des groupes simples, c'est-à-dire qu'ils ne possèdent aucun diviseur normal admissible en dehors d'eux-mêmes et du groupe unité.

Si l'on admet parmi les opérateurs, en particulier, tous les automorphismes intérieurs, alors la suite (2) se compose uniquement de diviseurs normaux, et on l'appelle suite principale. Si l'on admet tous les automorphismes, elle s'appelle suite caractéristique. Dans les autres exemples du \S{}1, il s'agirait de suites de composition d'idéaux dans $\mo$, respectivement de suites de composition de modules ou de bimodules.

\emph{Théorème de Jordan--Hölder.} \emph{S'il existe pour un groupe $\mG$ deux suites de composition différentes:}
\[
        \mG\supset \mA\supset \mA_2\supset \ldots\supset \mA_r=\mE,
        \qquad
        \mG\supset \mB\supset \mB_2\supset \ldots\supset \mB_s=\mE,
\]
\emph{alors elles ont la même longueur: $r=s$, et les facteurs de composition}
\[
        \mG/\mA,\ \mA/\mA_2,\ldots,\mA_{r-1}/\mE
\]
\emph{sont isomorphes, dans un certain ordre, à}
\[
        \mG/\mB,\ \mB/\mB_2,\ldots,\mB_{s-1}/\mE.
\]
Pour la démonstration, voir par exemple E. Noether, Abstrakter Aufbau usw., Math. Annalen 96 (1926), \S{}10, p. 57. On y démontre en même temps:

\emph{S'il existe dans $\mG$ une suite de composition, alors on peut faire passer une suite de composition par tout diviseur normal $\mH$ de $\mG$.}

L'existence d'une suite de composition est claire pour les groupes finis, plus généralement pour les groupes dans lesquels valent une condition maximale et une condition minimale:

La \emph{condition maximale} dit que, dans tout ensemble de sous-groupes, il existe un sous-groupe maximal, c'est-à-dire un sous-groupe qui n'est plus contenu dans un autre sous-groupe de l'ensemble. Il est équivalent de dire que toute chaîne ascendante de sous-groupes $\mA_1\subset\mA_2\subset\mA_3\ldots$ s'arrête après un nombre fini de termes.

La \emph{condition minimale} dit que, dans tout ensemble de sous-groupes, il en existe un minimal, qui ne contient plus aucun autre sous-groupe de l'ensemble; ou encore que toute chaîne descendante $\mA_1\supset\mA_2\supset\mA_3\ldots$ s'arrête après un nombre fini de termes.

\emph{Si seuls les diviseurs normaux sont admis comme sous-groupes} (par exemple pour les groupes abéliens), \emph{alors réciproquement les conditions maximale et minimale résultent de l'existence d'une suite de composition.}

\emph{Démonstration.} Tout diviseur normal possède une suite de composition, dont la longueur est appelée longueur du diviseur normal. Si $\mA\subset\mB$, alors la longueur de $\mA$ est plus petite que celle de $\mB$, puisqu'on peut faire passer une suite de composition de $\mB$ par $\mA$. Ainsi tout sous-groupe de longueur la plus courte est en même temps minimal, et tout sous-groupe de plus grande longueur est en même temps maximal.

\providecommand{\mG}{\mathfrak{G}}
\providecommand{\mA}{\mathfrak{A}}
\providecommand{\mB}{\mathfrak{B}}
\providecommand{\mC}{\mathfrak{C}}
\providecommand{\mD}{\mathfrak{D}}
\providecommand{\mE}{\mathfrak{E}}
\providecommand{\mH}{\mathfrak{H}}
\providecommand{\mR}{\mathfrak{R}}
\providecommand{\mS}{\mathfrak{S}}
\providecommand{\mT}{\mathfrak{T}}

\subsection*{\S{} 4. Produits directs et intersections}

Un groupe $\mG$ s'appelle produit direct de deux facteurs, $\mG=\mA\times\mB$, si
\[
\begin{array}{ll}
1.& \mA,\mB\text{ sont des diviseurs normaux dans }\mG,\\
2.& \mA\mB=\mG,\\
3.& \mA\cap\mB=\mE.
\end{array}
\]
\emph{La définition est manifestement équivalente à ceci:} tout élément $g$ de $\mG$ se représente d'une manière univoque sous la forme $g=ab$, avec $a$ dans $\mA$, $b$ dans $\mB$, et les éléments de $\mA$ commutent avec ceux de $\mB$: $ab=ba$.

Un groupe $\mG$ s'appelle produit direct de $n$ facteurs, $\mG=\mA_1\times\ldots\times\mA_n$, si, en posant $\mB_i=\mA_1\ldots\mA_{i-1}\mA_{i+1}\ldots\mA_n$, $\mG=\mA_i\times\mB_i$ est direct pour chaque $i$.

\emph{La définition est de nouveau équivalente à ceci:} tout $g$ de $\mG$ se représente d'une manière univoque sous la forme
\[
        g=a_1a_2\ldots a_n,\qquad a_i\in\mA_i,
\]
et les éléments de $\mA_i$ commutent avec ceux de $\mA_k$.

Les faits suivants seront souvent utilisés:
\begin{enumerate}
\item Si $\mA_1\times\cdots\times\mA_n=\mR$, $\mR\times\mH=\mG$, alors $\mA_1\times\cdots\times\mA_n\times\mH=\mG$.
\item Si $\mG=\mA_1\times\ldots\times\mA_n=\mC_1\times\ldots\times\mC_n$, et $\mC_i\subseteqq\mA_i$, alors $\mC_i=\mA_i$. (Cela résulte de la représentation des éléments de $\mA_i$ au moyen des $\mC_j$.)
\item Si $\mG=\mA\times\mB$, $\mR\supseteq\mA$, alors $\mR=\mA\times(\mR\cap\mB)$. (Cela résulte de la représentation des éléments de $\mR$ sous la forme $ab$, où le second facteur appartient à la fois à $\mR$ et à $\mB$.)
\item Si $\mG=\mA\times\mB$, alors $\mA\simeq\mG/\mB$ et $\mB\simeq\mG/\mA$ (deuxième théorème d'isomorphie). Il en résulte en outre:
\item Si $\mG=\mA\times\mB$, et si $\mA$ et $\mB$ possèdent des suites de composition de longueurs $m$ et $n$, alors $\mG$ possède une suite de composition de longueur $m+n$.
\item Si $\mA\simeq\overline{\mA}$ et $\mB\simeq\overline{\mB}$, alors $\mA\times\mB\simeq\overline{\mA}\times\overline{\mB}$.
\end{enumerate}
Un groupe s'appelle directement indécomposable lorsqu'il ne peut pas être représenté comme produit direct de facteurs $\ne\mE$.

Il est clair que: \emph{tout groupe satisfaisant à la condition minimale est produit direct d'un nombre fini de groupes directement indécomposables.}\footnote{Comme W. Krull dans le cas commutatif et O. Schmidt en général l'ont montré (cf. note 6), sous l'hypothèse des conditions maximale et minimale, la représentation est unique à isomorphie près. Comme on peut adjoindre tous les automorphismes intérieurs au domaine d'opérateurs sans changer la notion d'indécomposabilité directe, il suffit d'exiger les conditions de finitude pour les diviseurs normaux ou l'existence d'une suite principale.}

Si l'on écrit additivement les groupes considérés, les notions de produit et de produit direct passent à celles de somme et de somme directe. Notation: pour une somme de groupes $(\mA_1,\mA_2,\ldots)$, pour une somme directe $\mA_1+\mA_2+\ldots$.

La notion d'``intersection directe'' ne sera pas utilisée dans la suite; cependant les théorèmes suivants sur le lien entre intersection directe et produit direct montrent comment la théorie des idéaux des chapitres suivants peut aussi se formuler avec des intersections au lieu de sommes, ce qui rétablit le lien avec les théories commutatives connues.

Un groupe $\mD$ s'appelle intersection directe de deux groupes $\mR$ et $\mS$ relativement à $\mG$, si
\[
\begin{array}{ll}
1.& \mR,\mS\text{ sont des diviseurs normaux dans }\mG,\\
2.& \mR\cap\mS=\mD,\\
3.& \mR\mS=\mG.
\end{array}
\]
De même $\mD$ s'appelle intersection directe de $n$ groupes $\mS_1,\ldots,\mS_n$ si, en posant $\mR_i=\mS_1\cap\ldots\cap\mS_{i-1}\cap\mS_{i+1}\cap\ldots\cap\mS_n$, $\mD=\mR_i\cap\mS_i$ est directe pour chaque $i$.

Il résulte de la définition que $\mD$ doit être un diviseur normal dans $\mG$. Si $\mD=\mS_1\cap\ldots\cap\mS_n$ est directe et si, dans l'homomorphisme $\mG\sim\mG/\mD$, les groupes $\mG,\mS,\mD$ passent à $\overline{\mG},\overline{\mS},\overline{\mE}$, alors $\overline{\mE}=\overline{\mS}_1\cap\ldots\cap\overline{\mS}_n$ devient directe. Réciproquement, à partir de toute représentation de cette sorte $\overline{\mE}=\overline{\mS}_1\cap\ldots\cap\overline{\mS}_n$, on revient à une représentation directe $\mD=\mS_1\cap\ldots\cap\mS_n$. Par cette correspondance, les théorèmes sur les intersections directes pour $\mD$ arbitraire sont ramenés au cas particulier $\mD=\mE$.

Dans le cas $\mD=\mE$ et pour deux facteurs, les conditions du produit direct et celles de l'intersection directe coïncident. Pour $n$ facteurs, il existe la relation biunivoque suivante entre produit direct et intersection:

\emph{I. Si $\mG=\mA_1\times\ldots\times\mA_n=\mA_i\times\mB_i$, alors $\mE=\mB_1\cap\ldots\cap\mB_n=\mB_i\cap\mC_i$ directement, et $\mC_i=\mA_i$.}

\emph{II. Si $\mE=\mS_1\cap\ldots\cap\mS_n=\mR_i\cap\mS_i$ directement, alors $\mG=\mR_1\times\ldots\times\mR_n=\mR_i\times\mT_i$, et $\mT_i=\mS_i$.}

\emph{Démonstration de I.} Soit $\mB=\mB_1\cap\ldots\cap\mB_n=\mB_i\cap\mC_i$; alors $\mC_i\supseteqq\mA_i$. Nous montrons que $\mC_i\subseteqq\mA_i$. Soit par exemple $c$ dans $\mC_1$; alors $c$ est dans $\mB_2,\ldots,\mB_n$, donc la représentation en composantes au moyen des $\mA$ donne
\[
        c=a_1a_2\ldots a_n=a_1ea_3\ldots a_n=a_1a_2e\ldots a_n=\ldots=a_1\ldots a_{n-1}e.
\]
Puisque le produit de tous les $\mA_i$ est direct, les représentations coïncident; à chaque place sauf la première se trouve une fois $e$, donc $c=a_1$, ou $\mC_1\subseteqq\mA_1$, c'est-à-dire $\mC_1=\mA_1$, et de même $\mC_i=\mA_i$. Il s'ensuit $\mB=\mB_i\cap\mC_i=\mB_i\cap\mA_i=\mE$; $\mB_i\mC_i=\mB_i\mA_i=\mG$, donc l'intersection est directe.

\emph{Démonstration de II.} Soit $\mH=\mR_1\ldots\mR_n=\mR_i\mT_i$; alors $\mT_i\subseteqq\mS_i$. Nous montrons que $\mS_i\subseteqq\mT_i$. Soit par exemple $s$ dans $\mS_1$; on obtient -- au moyen de $\mG=\mS_i\times\mR_i$ --
\[
        s=s_1e=s_2r_2=\cdots=s_nr_n.
\]
Formons $t_1=er_2\ldots r_n=t_ir_i$; alors, en tenant compte de la commutativité élément par élément de $\mR_i$ et $\mS_i$, on obtient
\[
        st_1^{-1}=s_it_i^{-1},
\]
donc un élément de tous les $\mS_i$, et par conséquent égal à $e$. Ainsi $\mS_1\subseteqq\mT_1$, ou $\mT_1=\mS_1$, et de même $\mT_i=\mS_i$. Il s'ensuit $\mH=\mR_i\mT_i=\mR_i\mS_i=\mG$ et $\mR_i\cap\mT_i=\mR_i\cap\mS_i=\mE$; ainsi $\mG$ est produit direct des $\mR_i$.

\providecommand{\mG}{\mathfrak{G}}
\providecommand{\mA}{\mathfrak{A}}
\providecommand{\mE}{\mathfrak{E}}
\providecommand{\mH}{\mathfrak{H}}

\subsection*{\S{} 5. Groupes complètement réductibles}

Un groupe s'appelle \emph{complètement réductible} lorsqu'il est le produit direct d'un nombre fini de groupes simples:
\[
        \mG=\mA_1\times\ldots\times\mA_n.
\]
Dans ce cas, les groupes
\[
        \mA_1\times\ldots\times\mA_n
        \supset \mA_1\times\ldots\times\mA_{n-1}
        \supset\ldots\supset \mA_1\supset \mE
\]
forment une suite de composition. La longueur est $n$, et les facteurs de composition sont
\[
        \sim \mA_n,\ \mA_{n-1},\ldots,\mA_1.
\]

\emph{Si $\mG$ est complètement réductible, alors tout diviseur normal est facteur direct, et l'autre facteur peut être choisi comme produit de groupes simples qui apparaissent dans une décomposition en produit donnée de $\mG$. Le diviseur normal $\mH$ est lui-même complètement réductible.}

\emph{Démonstration.} On a
\begin{equation}
        \mG=\mH\mA_1\mA_2\ldots\mA_n.\tag{3}
\end{equation}
Si l'on pose maintenant $\mH_i=\mH\mA_1\mA_2\ldots\mA_i$, alors $\mH_{i+1}=\mH_i\mA_{i+1}$. Ou bien $\mA_{i+1}\subseteqq\mH_i$, donc $\mH_{i+1}=\mH_i$; nous omettons alors dans (3) le facteur $\mA_{i+1}$. Ou bien $\mH_i\cap\mA_{i+1}$ est un diviseur normal propre de $\mA_{i+1}$, donc égal à $=\mE$, et alors $\mH_i\times\mA_{i+1}$ est direct. Après omission des facteurs superflus, (3) devient
\[
        \mG=\mH\times\mA_{i_1}\times\ldots\times\mA_{i_r},
\]
donc $\mH$ est facteur direct. En outre, on en déduit que
\[
        \mH\simeq \mG/\mA_{i_1}\times\ldots\times\mA_{i_r}
        \simeq \mA_{i_1}\times\ldots\times\mA_{j_\mu},
\]
où $\mA_{j_1},\ldots,\mA_{j_\mu}$ sont les autres $\mA_i$; donc $\mH$ est complètement réductible.

\emph{Corollaire.} \emph{Toute décomposition d'un groupe complètement réductible en facteurs directement indécomposables est une décomposition en facteurs simples, et donne par conséquent lieu à une suite de composition; il en résulte la détermination univoque des facteurs, à isomorphie près.}

\providecommand{\mG}{\mathfrak{G}}
\providecommand{\mE}{\mathfrak{E}}
\providecommand{\mH}{\mathfrak{H}}
\providecommand{\mR}{\mathfrak{R}}
\providecommand{\mo}{\mathfrak{o}}

\subsection*{\S{} 6. Modules relativement à un corps. Systèmes hypercomplexes}

Un exemple presque trivial de groupes avec opérateurs complètement réductibles est fourni par les modules ayant une base finie relativement à un corps (non nécessairement commutatif).

Soit $\mG$ un module à droite sur $\mathsf{K}$, et soit l'élément unité $e$ de $\mathsf{K}$ en même temps opérateur unité: $ae=a$ pour $a$ dans $\mG$. Tout sous-module $a\mathsf{K}$ dérivé d'un $a$ est simple, c'est-à-dire égal au module dérivé de l'un quelconque de ses éléments $\ne0$. Si donc $\mH$ est le sous-module dérivé de certains éléments, et si $a$ est un élément qui n'appartient pas à $\mH$, alors $\mH\cap a\mathsf{K}=\mE$, donc $\mH+a\mathsf{K}$ est directe. Ainsi, en partant d'un élément de base quelconque $a_1$, on peut toujours adjoindre de nouveaux éléments de base $a_2,a_3,\ldots$, et l'on obtient $\mG$ comme somme directe
\[
        \mG=a_1K+a_2K+\cdots+a_nK.
\]
\emph{Donc $\mG$ est complètement réductible.}

Le nombre $n$, longueur de la suite de composition, s'appelle le \emph{rang} de $\mG$ relativement à $\mathsf{K}$. En raison de la représentation univoque des éléments de $\mG$ (et parce que $e$ a été supposé opérateur unité), la base $a_1,\ldots,a_n$ est linéairement indépendante.

Tout sous-module $\mH$ est un sommant direct:
\[
        \mG=\mH+\mR=(h_1\mathsf{K}+\ldots+h_s\mathsf{K})+(r_1\mathsf{K}+\ldots+r_{n-s}\mathsf{K}),
\]
ou encore: \emph{on peut compléter toute base linéairement indépendante de $\mH$ en une base linéairement indépendante de $\mG$.} Les $r_i$ peuvent même être choisis parmi les éléments de base initiaux $a_i$.

Soient $(a_1,\ldots,a_n)$ et $(c_1,\ldots,c_n)$ des bases linéairement indépendantes. Alors
\[
        c_k=\sum a_i\pi_{ik},
\]
ou, écrit en matrices:
\[
        (c_1,\ldots,c_n)=(a_1,\ldots,a_n)P;
        \qquad
        P=\begin{pmatrix}
        \pi_{11}&\ldots&\pi_{1n}\\
        \vdots&&\vdots\\
        \pi_{n1}&\ldots&\pi_{nn}
        \end{pmatrix}.
\]
Réciproquement,
\[
        (a_1,\ldots,a_n)=(c_1,\ldots,c_n)Q.
\]
Il s'ensuit
\[
        (a_1,\ldots,a_n)=(a_1,\ldots,a_n)PQ,
        \qquad
        (c_1,\ldots,c_n)=(c_1,\ldots,c_n)QP,
\]
donc, puisque les $a_i$ et les $c_i$ sont linéairement indépendants,
\[
        PQ=QP=E.
\]

Les $\mathsf{K}$-modules particuliers qui sont en même temps des anneaux sont les ``systèmes hypercomplexes''.

Un anneau $\mo$, qui est en même temps un module à droite relativement à un corps commutatif $\mathsf{K}$, s'appelle \emph{système hypercomplexe relativement à $\mathsf{K}$} (également désigné dans la littérature comme ``algèbre sur $\mathsf{K}$''), si:
\begin{enumerate}
\item le rang est fini (soit $u_1,\ldots,u_n$ une base linéairement indépendante);
\item $ab\cdot\varkappa=a\cdot b\varkappa=a\varkappa\cdot b$. Autrement dit, les actions de $\mathsf{K}$ et de $\mo$ sur le module commutent;
\item l'élément unité $e$ de $\mathsf{K}$ est en même temps opérateur unité: $a\varepsilon=a$ pour $a$ dans $\mo$.
\end{enumerate}
Comme
\[
        \Bigl(\sum_i u_i\alpha_i\Bigr)\Bigl(\sum_k u_k\beta_k\Bigr)
        =\sum_{i,k}(u_i u_k)(\alpha_i\beta_k)
\]
le système est déterminé univoquement lorsque, outre $\mathsf{K}$, est donnée la \emph{table de multiplication}, qui indique comment chaque produit $u_i u_k$ s'exprime au moyen des $u_j$:
\[
        u_i u_k=\sum_j u_j\gamma^j_{ik}.
\]
Les $\gamma^j_{ik}$ doivent satisfaire aux relations connues qui résultent de la loi associative.

Si $\mo$, comme nous le supposerons souvent, possède un élément unité $e$ ($ea=ae=a$ pour tout $a$), alors on peut identifier les éléments $x$ de $\mathsf{K}$ avec $e\varkappa$, donc considérer $\mathsf{K}$ comme un sous-corps de $\mo$. Le système hypercomplexe peut alors aussi être décrit comme \emph{un anneau de rang fini relativement à un corps situé dans le centre.}

Un exemple de système hypercomplexe est l'\emph{anneau de groupe} d'un groupe fini; il apparaît lorsqu'on prend pour éléments de base $u_i$ les éléments d'un groupe fini, et pour multiplication la multiplication du groupe. Le corps $\mathsf{K}$ est ici arbitraire.

\footnotetext{La condition 2. exprime que la multiplication par un élément de $K$ signifie un homomorphisme pour $\mo$, lorsque $\mo$ est considéré aussi bien comme $\mo$-module à gauche que comme $\mo$-module à droite. En vertu de 3., ces homomorphismes sont même des isomorphismes.}

\providecommand{\mo}{\mathfrak{o}}

\subsection*{\S{} 7. Matrices}

Les matrices carrées $(\pi_{ik})$, avec $\pi_{ik}$ pris dans un anneau $\mo$, forment un anneau pour la multiplication matricielle ordinaire $\sum \pi_{ij}\varrho_{jk}=\sigma_{ik}$ et l'addition $\pi_{ik}+\varrho_{ik}=\sigma_{ik}$.

Une matrice à éléments dans un corps $\mathsf{K}$, pour laquelle il existe une inverse à droite et une inverse à gauche, s'appelle \emph{régulière}.

Nous avons vu au \S{}6: \emph{la matrice de passage entre deux bases linéairement indépendantes d'un $\mathsf{K}$-module est régulière.}

\emph{Pour la régularité, une inverse à droite suffit}; car formons un module à droite sur $\mathsf{K}$ ayant une base linéairement indépendante $(a_1,\ldots,a_n)$ (le module des formes linéaires des indéterminées $a_1,\ldots a_n$), et posons
\[
        (c_1,\ldots,c_n)=(a_1,\ldots,a_n)P.
\]
Alors
\[
        (c_1,\ldots,c_n)P^{-1}=(a_1,\ldots,a_n)PP^{-1}=(a_1,\ldots,a_n);
\]
ainsi $(c_1,\ldots,c_n)$ est de nouveau une base, donc de nouveau linéairement indépendante, et la matrice $P$ est régulière. En outre, l'inverse à droite est égale à l'inverse à gauche.

De même, une inverse à gauche suffit.

\emph{Si $P$ possède une inverse à gauche, alors $P$ n'est pas un diviseur de zéro à gauche.}
\emph{Démonstration.} La relation $PA=0$ entraîne $A=P^{-1}PA=0$.

Donc $P$ possède aussi une seule inverse à droite, laquelle, d'après \S{}6, coïncide avec l'inverse à gauche.

Si l'on entend, comme d'habitude, par $P_\varkappa$ la matrice $(\pi_{ik}\varkappa)$, et par $C_{ik}$ la matrice qui a l'élément unité à la place $ik$ et des zéros partout ailleurs, alors
\[
        (\pi_{ik})=\sum C_{ik}\pi_{ik};
\]
ainsi les matrices à éléments dans $\mathsf{K}$ forment un $\mathsf{K}$-module de rang $n^2$. Les $C_{ik}$ satisfont aux relations
\[
        C_{ij}C_{jk}=C_{ik},
        \qquad
        C_{ij}C_{\bar j k}=0,\quad j\ne\bar j.
\]
Les matrices à éléments dans $\mathsf{K}$ forment donc un $\mathsf{K}$-module fini et, si $\mathsf{K}$ est commutatif, un système hypercomplexe.

\providecommand{\ma}{\mathfrak{a}}
\providecommand{\md}{\mathfrak{d}}
\providecommand{\mo}{\mathfrak{o}}
\providecommand{\mA}{\mathfrak{A}}
\providecommand{\mB}{\mathfrak{B}}
\providecommand{\mC}{\mathfrak{C}}
\providecommand{\mM}{\mathfrak{M}}

\section*{Chapitre II. Théorie non commutative des idéaux}

\subsection*{\S{} 8. Théorème d'homomorphie pour les anneaux}

Si $\ma$ est un idéal bilatère dans $\mo$, alors l'anneau quotient $\mo/\ma$ est non seulement un $\mo$-module, mais aussi un anneau, comme on le voit sans peine. Toute image homomorphe de $\mo$ est de nouveau un anneau, et isomorphe à un anneau quotient $\mo/\ma$, où $\ma$ est un idéal bilatère. La preuve est comme pour les groupes.

En particulier, si $\mo$ est un corps, il n'y a pas d'autres idéaux que $(0)$ et $\mo$; ainsi l'homomorphisme est ici soit un isomorphisme, soit il associe à chaque élément le zéro.

Supposons en particulier que l'anneau $\mo$ soit un module à droite sur $\mathsf{K}$, où $\mathsf{K}$ doit être un anneau commutant élément par élément avec $\mo$:
\[
        ab\cdot x=a\cdot bx=ax\cdot b
\]
(par exemple lorsqu'il s'agit d'un système hypercomplexe relativement à un corps commutatif $\mathsf{K}$); alors, comme idéaux admissibles (à droite, à gauche ou bilatères), on ne considère que ceux qui sont en même temps des $\mathsf{K}$-modules, et l'on exige, en plus de l'homomorphisme d'anneaux, encore l'homomorphisme d'opérateurs relativement à la multiplication par $\mathsf{K}$. Les idéaux admissibles deviennent des bimodules relativement à $\mo$ et $\mathsf{K}$. (Dans les systèmes hypercomplexes, par ``idéaux'' tout court, on entend toujours seulement ceux qui sont admissibles relativement au corps de coefficients $\mathsf{K}$ mis à la base.) Dans cette extension aussi, le théorème d'homomorphie ci-dessus reste valable.

Un exemple d'homomorphie d'anneaux est fourni par le passage d'un domaine de multiplicateurs $\mo$ au domaine absolu de multiplicateurs (\S{}1, exemple 4). Le domaine absolu de multiplicateurs est donc isomorphe à l'anneau quotient $\mo/\md$, où $\md$ est l'idéal bilatère des éléments de $\mo$ qui annulent $\mM$.

De la validité du théorème d'homomorphie suivent, comme au \S{}2, les premier et deuxième théorèmes d'isomorphie pour l'isomorphie d'anneaux et d'opérateurs.

Tous les produits $\mA\mB$ de sous-ensembles de $\mo$ sont à comprendre comme produits de modules: $\mA\mB$ est la totalité de toutes les sommes $\sum ab$, où $a$ est dans $\mA$, $b$ dans $\mB$, et $a\mB$ est la totalité de toutes les sommes $\sum ab$, $b$ dans $\mB$. Si $\mB$ est un module à droite relativement à un anneau quelconque pris comme domaine d'opérateurs, alors le produit $\mA\mB$, respectivement $a\mB$, est lui aussi un module à droite. (En particulier: idéal à droite, idéal admissible relativement à un domaine de coefficients $\mathsf{K}$, etc.) Si $(\mA,\mB)$ désigne la somme des modules $\mA,\mB$, les règles de calcul suivantes valent:
\[
        \mA\mB\cdot\mC=\mA\cdot\mB\mC,
\]
\[
        \mA(\mB,\mC)=(\mA\mB,\mA\mC),
        \qquad
        (\mA,\mB)\mC=(\mA\mC,\mB\mC).
\]
La somme directe est notée $\mA+\mB$ (\S{}4).

Dans la suite, nous considérons souvent des anneaux qui satisfont à une condition maximale ou minimale (\S{}3) pour les idéaux à droite. Les systèmes hypercomplexes en font notamment partie, puisque, par suite de la convention ci-dessus, seuls sont admis comme idéaux des $\mathsf{K}$-modules dont le rang demeure donc borné (\S{}6).

\providecommand{\mo}{\mathfrak{o}}
\providecommand{\mr}{\mathfrak{r}}
\providecommand{\ml}{\mathfrak{l}}

\subsection*{\S{} 9. Éléments idempotents. Décomposition en somme directe en idéaux à droite}

Soit $\mo$ un anneau avec élément unité.

Si $\mr$ est un idéal à droite, alors $\mr\mo\subseteqq\mr$, et, à cause de l'élément unité, on a même $\mr\mo=\mr$. De même pour les idéaux à gauche: $\mo\ml=\ml$.

Un élément $c$ s'appelle idempotent lorsque $c^2=c$.

\emph{Si $\mo=\mr_1+\ldots+\mr_n$ est une décomposition en idéaux à droite, et si}
\[
        e=e_1+\ldots+e_n,
        \qquad e_i\in\mr_i,
\]
\emph{alors}
\[
        e_i^2=e_i,
        \qquad e_i e_k=0\quad(i\ne k)
        \quad\hbox{(``relations d'orthogonalité'')},
\]
\[
        \mr_i=e_i\mo.
\]

\emph{Démonstration.} Soit $r$ un élément de $\mr_1$. On a
\[
        r=er=e_1r+\ldots+e_nr,
        \qquad e_ir\in\mr_i,
\]
mais aussi
\[
        r=r+0+\ldots+0,
\]
donc
\[
        e_1r=r,
        \qquad e_ir=0\quad(i\ne1).
\]
La première formule donne $\mr_1\subseteqq e_1\mo$; d'autre part $e_1\mo\subseteqq\mr_1$, donc $e_1\mo=\mr_1$. En prenant $r=e_1$, on obtient
\[
        e_1^2=e_1,
        \qquad e_i e_1=0\quad(i\ne1).
\]
Il en va de même lorsque l'on remplace l'indice $1$ par $k$.

\emph{Réciproquement, si $e=\sum e_i$, $e_i e_k=0$ $(i\ne k)$, $e_i^2=e_i$, et si l'on pose $\mr_i=e_i\mo$, alors $\mo=\mr_1+\ldots+\mr_n$.}

\emph{Démonstration.} Tout élément $r$ de $\mo$ est
\[
        r=er=e_1r+\ldots+e_nr,
\]
donc
\[
        \mo=(e_1\mo,\ldots,e_n\mo).
\]
Mais la représentation est univoque; car de
\[
        0=e_1a_1+\ldots+e_na_n
\]
une multiplication par $e_i$ donne:
\[
        e_i^2a_i=e_i a_i=0.
\]

D'une représentation à droite
\[
        \mo=\mr_1+\ldots+\mr_n=e_1\mo+\ldots+e_n\mo
\]
résulte par conséquent une représentation à gauche
\[
        \mo=\ml_1+\ldots+\ml_n=\mo e_1+\ldots+\mo e_n.
\]

Si les $\mr_i$ sont directement indécomposables, alors les $\ml_i$ le sont aussi; car une décomposition, par exemple de $\ml_1$, signifierait
\[
        \ml_1=\mo e_1=\mo e_1'+\mo e_1'',
        \qquad
        \mo=\mo e_1'+\mo e_1''+\mo e_2+\ldots+\mo e_n.
\]
Il en résulterait la décomposition à droite
\[
        \mo=e_1'\mo+e_1''\mo+e_2\mo+\ldots+e_n\mo
        =\mr_1'+\mr_1''+\mr_2+\ldots+\mr_n,
\]
\[
        \mr_1\cong \mo/(\mr_2+\ldots+\mr_n)=\mr_1'+\mr_2',
\]
donc $\mr_1$ serait décomposable.

\footnotetext{Il vaut encore une autre réciproque: si $e_i e_k=0$, $e_i^2=e_i$, $c=\sum e_i$, alors $c$ est une unité à gauche pour l'anneau $e_1\mo+\ldots+e_n\mo$. Que la somme soit directe se voit tout à fait comme ci-dessus.}

\providecommand{\mo}{\mathfrak{o}}
\providecommand{\ma}{\mathfrak{a}}
\providecommand{\mb}{\mathfrak{b}}
\providecommand{\mc}{\mathfrak{c}}
\providecommand{\mr}{\mathfrak{r}}

\subsection*{\S{} 10. Décomposition en idéaux bilatères}

Soit de nouveau $\mo$ un anneau avec élément unité, et soit
\begin{equation}
        \mo=\ma_1+\ldots+\ma_n\tag{1}
\end{equation}
une décomposition de $\mo$ en idéaux bilatères. Alors les relations d'orthogonalité valent aussi pour les idéaux:
\begin{equation}
        \ma_i\ma_k=0\quad(i\ne k),
        \qquad
        \ma_i^2=\ma_i.                                             \tag{2}
\end{equation}

\emph{Démonstration.} Soit $\mb_1=\ma_2+\ldots+\ma_n$, $\mo=\ma_1+\mb_1$; alors $\ma_1\mb_1\subseteqq \ma_1\cap\mb_1=0$, donc $\ma_1\mb_1=0$, et à plus forte raison
\[
        \ma_1\ma_i=0\quad(i\ne1);
\]
\[
        \ma_1=\ma_1\mo=\ma_1(\ma_1+\mb_1)=(\ma_1^2,\ma_1\mb_1)=\ma_1^2.
\]

Réciproquement: \emph{les relations (1) et (2) montrent que les modules $\ma_i$ sont des idéaux bilatères.}
\emph{Démonstration.}
\[
        \mo\ma_i=(\ma_1+\cdots+\ma_n)\ma_i=\ma_i^2=\ma_i,
        \qquad
        \ma_i\mo=\ma_i.
\]

\emph{Les idéaux à droite dans l'anneau $\ma_i$ sont des idéaux à droite dans $\mo$.}
\emph{Démonstration.}
\[
        \mr\mo=\mr(\ma_i+\mb_i)=(\mr\ma_i,\mr\mb_i)
        \subseteq(\mr,\ma_i\mb_i)=(\mr,0)=\mr.
\]

\emph{Si $\mr$ est un idéal à droite dans $\mo$, alors $\mr$ est somme directe d'idéaux à droite $\mr\ma_i\subseteqq\ma_i$.}
\emph{Démonstration.} $\mr=\mr\mo=(\mr\ma_1,\ldots,\mr\ma_n)$, et $\mr\ma_i\mo=\mr\ma_i$, donc les $\mr\ma_i$ sont des idéaux à droite contenus dans $\mr$. En outre $\mr\ma_i\subseteqq\ma_i$, donc la somme est directe:
\[
        \mr=\mr\ma_1+\ldots+\mr\ma_n.
\]

\emph{En particulier, si $\mr$ est directement indécomposable, alors $\mr\subseteqq$ ne peut avoir qu'une seule composante; il doit donc être contenu dans un des $\ma_i$.}

Sur la base de ces théorèmes, on maîtrise la théorie des idéaux dans $\mo$ dès que l'on connaît celle des différents $\ma_i$. Nous nous restreignons donc le plus souvent aux anneaux bilatéralement indécomposables.

\emph{Deux idéaux à droite contenus dans des $\ma_i$ différents ne peuvent jamais être isomorphes au sens des opérateurs.}
\emph{Démonstration.} Un idéal situé dans $\ma_1$ est annulé par tous les autres $\ma_k$, mais non par $\ma_1$. En revanche, un idéal situé dans $\ma_2$ est annulé par $\ma_1$.

S'il est donc possible de décomposer l'anneau $\mo$ en idéaux bilatères indécomposables $\ma_1+\ldots+\ma_n$, et chaque $\ma_i$ à son tour en idéaux à droite indécomposables unilatéralement $\mr_i'+\mr_i''+\ldots$, alors les $\mr$ isomorphes au sens des opérateurs sont à chercher parmi ceux qui appartiennent au même $\ma_i$. Mais il peut arriver qu'il y ait encore différentes classes de $\mr$ -- c'est-à-dire des classes non isomorphes au sens des opérateurs -- dans le même $\ma_i$, comme le montre l'exemple suivant.

Soit $\mathsf{K}$ le corps des nombres rationnels,
\[
        \mo=e_1\mathsf{K}+e_2\mathsf{K}+u\mathsf{K}
\]
un système hypercomplexe, les $e_i$ et $u$ étant permutables avec les nombres de $\mathsf{K}$, et avec la table de multiplication
\[
\begin{array}{c|ccc}
     & e_1&e_2&u\\ \hline
 e_1 & e_1&0&u\\
 e_2 & 0&e_2&0\\
 u   & 0&u&0
\end{array}
\]
L'élément unité est $e=e_1+e_2$. Les $e_i$ satisfont aux relations d'orthogonalité. La décomposition à droite est donc
\[
        \mo=e_1\mo+e_2\mo=(e_1,u)+(e_2),
\]
et la décomposition à gauche
\[
        \mo=\mo e_1+\mo e_2=(e_1)+(e_2,u).
\]
Comme $e_2\mo$ et $\mo e_1$ sont manifestement indécomposables, $\mo e_2$ et $e_1\mo$ doivent l'être aussi.

Une décomposition bilatère est impossible; car alors une composante devrait contenir $e_1\mo$ et l'autre $e_2\mo$; la première contiendrait alors $u$, mais la seconde aussi, puisque $ue_2=u$.

Mais $e_1\mo$ et $e_2\mo$ ne sont pas isomorphes au sens des opérateurs, puisqu'ils ont des rangs différents par rapport à $\mathsf{K}$.

\emph{Soit $\mo=\ma_1+\ldots+\ma_n$, les $\ma_i$ étant bilatéralement indécomposables. Alors ils sont déterminés de manière univoque.}

\emph{Démonstration.} Supposons en outre que $\mo=\mc_1+\ldots+\mc_n$; alors
\[
        \ma_i=\ma_i\mo=(\ma_i\mc_1,\ldots,\ma_i\mc_n).
\]
Cette dernière somme est directe parce que $\ma_i\mc_k\subseteqq\mc_k$ et $\subseteqq\ma_i$; mais comme les $\ma_i$ sont indécomposables, tous les $\ma_i\mc_k$ doivent être nuls sauf un: $\ma_i\mc_{j_i}$. Ainsi
\[
        \ma_i=\ma_i\mc_{j_i}.
\]
De même:
\[
        \mc_{j_i}=\mo\mc_{j_i}=\ma_1\mc_{j_i}+\cdots+\ma_n\mc_{j_i},
\]
et comme $\ma_i\mc_{j_i}\ne0$, tous les autres termes sont $=0$,
\[
        \mc_{j_i}=\ma_i\mc_{j_i}=\ma_i,
\]
q.e.d.

\providecommand{\mo}{\mathfrak{o}}
\providecommand{\ma}{\mathfrak{a}}
\providecommand{\mA}{\mathfrak{A}}
\providecommand{\mZ}{\mathfrak{Z}}

\subsection*{\S{} 11. Le centre}

Le centre $\mZ$ d'un anneau $\mo$ est l'ensemble des $z$ qui commutent avec tous les éléments de l'anneau ($za=az$ pour tout $a$). $\mZ$ est un anneau commutatif.

Tout idéal unilatère $\ma$ dans $\mo$ possède dans $\mZ$ un ``idéal contracté'' $\ma\cap\mZ$. En effet, puisque $\ma$ aussi bien que $\mZ$ sont des $\mZ$-modules, $\ma\cap\mZ$ l'est aussi.

Tout idéal $\mA$ dans $\mZ$ possède un idéal étendu: l'idéal $\ma=(\mA,\mo\mA)$ engendré par $\mA$ dans $\mo$. Celui-ci est bilatère, car
\[
        \ma\mo=(\mo\mA,\mA)\mo=(\mo\mA\mo,\mA\mo)
        = (\mo^2\mA,\mo\mA)\subseteqq \mo\mA\subseteqq\ma.
\]
Si l'on suppose donné dans $\mo$ un élément unité, on peut écrire: $\ma=\mo\mA$. Dans ce cas vaut le théorème suivant:

\emph{À toute décomposition bilatère de $\mo$ correspond biunivoquement une décomposition du centre: de}
\begin{equation}
        \mo=\ma_1+\ldots+\ma_n,
        \qquad \mA_i=\ma_i\cap\mZ                         \tag{1}
\end{equation}
\emph{on obtient}
\begin{equation}
        \mZ=\mA_1+\ldots+\mA_n,
        \qquad \mo\mA_i=\ma_i,                              \tag{2}
\end{equation}
\emph{et réciproquement.}

\emph{Démonstration.} Soit $z$ un élément du centre, et
\begin{equation}
        z=z_1+\ldots+z_n\tag{3}
\end{equation}
sa décomposition suivant (1). Alors les $z_i$ sont des éléments du centre; car
\[
        za=z_1a+\ldots+z_na=az=az_1+\ldots+az_n;
\]
à cause du caractère direct de la somme, et puisque les $\ma_i$ sont bilatères,
\[
        z_i a=az_i.
\]
Donc $z_i$ appartient à $\ma_i\cap\mZ=\mA_i$; ainsi (3) donne la décomposition en somme affirmée. En particulier on obtient $e=\sum e_i$; donc les $e_i$ sont des éléments du centre. Enfin
\[
        z=ze_1+\ldots+ze_n,
\]
donc
\[
        \mA_i=\mZ e_i,
        \qquad
        \mo\mA_i=\mo\mZ e_i=\mo e_i=\ma_i.
\]

Réciproquement, si $\mZ=\mA_1+\ldots+\mA_n$ est une décomposition du centre, alors, d'après le \S{}9,
\[
        \mA_i=\mZ e_i,
        \qquad e_i^2=e_i,
        \qquad e_i e_k=0\quad(i\ne k).
\]
Il vient donc, à cause des relations d'orthogonalité \hbox{\(\S 9\)},
\[
\begin{aligned}
        \mo&=\mo e=\mo e_1+\ldots+\mo e_n
        =\mo\mZ e_1+\ldots+\mo\mZ e_n\\
        &=\mo\mA_1+\ldots+\mo\mA_n
        =\ma_1+\ldots+\ma_n .
\end{aligned}
\]
Si l'on pose maintenant $\ma_i\cap\mZ=\mA_1'$, alors on sait, par la première partie de l'assertion:
\[
        \mZ=\mA_1'+\ldots+\mA_n'\quad\hbox{directement}.
\]
Mais $\mA_i'\supseteqq\mA_i$, donc $\mA_i'=\mA_i$ (\S 4, 2).

\emph{Corollaire.} $\ma_i$ et $\mA_i$ sont simultanément décomposables directement en idéaux bilatères, ou ne le sont pas.

\providecommand{\mo}{\mathfrak{o}}
\providecommand{\mc}{\mathfrak{c}}
\providecommand{\mC}{\mathfrak{C}}
\providecommand{\md}{\mathfrak{d}}
\providecommand{\mr}{\mathfrak{r}}
\providecommand{\mZ}{\mathfrak{Z}}

\subsection*{\S{} 12. Idéaux nilpotents}

Un idéal $\mc$ s'appelle \emph{nilpotent} si $\mc^\varrho=0$.

\emph{Exemple.} L'idéal $(u)$ dans le \S{}10.

\emph{S'il existe dans $\mo$ un idéal à droite nilpotent $\mr$, alors il existe aussi un idéal à gauche nilpotent (même un idéal bilatère).}

\emph{Démonstration.} Soit $\mr^\varrho=0$. Alors $\mo\mr$, respectivement, si $\mo$ ne contient pas d'élément unité et que $\mo\mr$ pourrait être nul, $(\mo\mr,\mr)$, est un idéal à gauche, et
\[
        (\mo\mr)^\varrho=\mo\mr\mo\mr\ldots\mo\mr\subseteqq\mo\mr\mr\ldots\mr=\mo\mr^\varrho=0.
\]

\emph{La somme de deux idéaux à droite nilpotents est de nouveau un idéal à droite nilpotent.}

\emph{Démonstration.} Soit $\mc^\varrho=0$, $\md^\sigma=0$. Dans
\[
        (\mc,\md)^{\varrho+\sigma-1}=(\ldots,\md\mc\ldots\md\ldots\mc\ldots,\ldots)
\]
chaque terme contient soit au moins $\varrho$ facteurs $\mc$, soit au moins $\sigma$ facteurs $\md$. Dans le premier cas nous avons
\[
        \md\mc\ldots\md\ldots\mc\ldots\subseteqq\md\mc\mc\mc\ldots\mc=\md\mc^\varrho=0;
\]
dans le second cas il en va de même pour les facteurs $\md$. Ainsi
\[
        (\mc,\md)^{\varrho+\sigma-1}=0.
\]

Si maintenant la condition maximale pour les idéaux à droite est satisfaite dans $\mo$, il existe un idéal à droite nilpotent maximal. Celui-ci contient tous les autres idéaux à droite nilpotents, car sinon on pourrait former sa somme avec un tel idéal. Soit $\mc$ cet idéal. Puisque $\mo\mc$ est aussi un idéal à droite nilpotent, on a $\mo\mc\subseteqq\mc$, donc $\mc$ est un idéal à gauche. Tout idéal à gauche nilpotent $\md$ est également contenu dans $\mc$. En effet, $(\md,\md\mo)$ est un idéal à droite nilpotent, donc $\md\subseteqq\mc$. \emph{Il existe donc un idéal nilpotent maximal (bilatère) $\mc$, ou ``radical'', qui contient tous les autres idéaux (à droite et à gauche).}

Dans l'exemple du \S{}10, $(u)$ est l'idéal nilpotent maximal.

Si le radical est l'idéal nul, on parle d'un ``anneau sans radical'' (``anneau semi-simple'', ``système de Dedekind''). L'anneau quotient par le radical est toujours un anneau sans radical.

\emph{Si $\mZ$ est le centre de $\mo$, et $\mc$ le radical de $\mo$, alors $\mC=\mc\cap\mZ$ est le radical de $\mZ$.}

\emph{Démonstration.} Il est clair que $\mC$ est nilpotent. S'il existait encore un idéal nilpotent plus grand $\overline{\mC}$ dans $\mZ$, celui-ci pourrait être étendu en un idéal nilpotent $\overline{\mc}$ qui ne serait pas entièrement contenu dans $\mc$ $((\overline{\mC}\mo)^\varrho=\overline{\mC}^{\varrho}\mo^\varrho=0)$.

\emph{Conséquence.} Si $\mo$ est un anneau sans radical, alors $\mZ$ l'est aussi. Mais la réciproque n'est pas vraie, comme le montre notre exemple du \S{}10. Le centre de $\mo/\mc$ contient un anneau isomorphe à $\mZ/\mC$; celui-ci peut toutefois être un sous-anneau propre (exemple dans le \S{}10).

\providecommand{\mo}{\mathfrak{o}}
\providecommand{\ma}{\mathfrak{a}}
\providecommand{\mc}{\mathfrak{c}}
\providecommand{\md}{\mathfrak{d}}
\providecommand{\ml}{\mathfrak{l}}
\providecommand{\mr}{\mathfrak{r}}
\providecommand{\mt}{\mathfrak{t}}

\subsection*{\S{} 13. Anneaux complètement réductibles}

D'après le \S{}5, un anneau s'appelle complètement réductible à droite lorsqu'il est somme directe d'un nombre fini d'idéaux à droite simples.

\emph{Un anneau complètement réductible à droite avec élément unité ne possède aucun idéal nilpotent $\ne(0)$, donc aucun radical.}

\emph{Démonstration.} Tout idéal à droite $\mr$ est un facteur direct (\S{}5):
\[
        \mo=\mt+\mr=e_1\mo+e_2\mo.
\]
Comme $e_2^2=e_2$, on a aussi $e_2^\varrho=e_2$. Si $\mr$ était nilpotent, on aurait $e_2^\varrho=0$, donc $e_2=0$, donc $\mr=0$, q.e.d.

Nous montrons maintenant les deux réciproques. Premièrement: \emph{un anneau sans radical avec condition minimale pour les idéaux à droite est complètement réductible relativement aux idéaux à droite.}

\emph{Démonstration.} Soit $\mt$ un idéal à droite minimal $\ne0$. Nous voulons montrer que $\mt$ contient un élément idempotent.

Comme $\mt^2\subseteqq\mt$ et $\mt^2\ne0$, la minimalité donne $\mt^2=\mt$; car $\mt^2$ est un idéal à droite. Il existe donc un $a$ dans $\mt$ tel que $a\mt\ne0$. Mais alors on doit avoir $a\mt=\mt$, car $a\mt$ est un idéal à droite.

L'ensemble de tous les $b$ dans $\mt$ qui sont annulés par $a$ ($ab=0$) est un idéal à droite, et n'est pas $=\mt$, donc est $=0$. Ainsi, la relation $ab=0$ entraîne $b=0$.

Comme $\mt=a\mt$, $a$ doit pouvoir se représenter sous la forme $ac$: $a=ac$, $c\ne0$. Il s'ensuit
\[
        ac=ac^2,
        \qquad a(c-c^2)=0,
        \qquad c-c^2=0,
        \qquad c^2=c.
\]
Nous montrons ensuite que $\mt$ est un facteur direct. $c\mo$ est un idéal à droite $\subseteqq\mt$, $\ne0$, puisque $c^2=c$ appartient à $c\mo$; donc $=\mt$. Tout élément $r$ de $\mo$ admet une représentation:
\[
        r=cr+(r-cr) \qquad \hbox{(décomposition de Peirce unilatère).}
\]
Les éléments $cr$ forment l'idéal $\mt$; les éléments $r-cr$ forment un autre idéal à droite $\mr$, qui est annulé par $c$: $c\mr=0$. Donc
\[
        \mo=(\mt,\mr).
\]
Comme les éléments de $\mt$ ne sont pas annulés par $c$, la somme est directe:
\begin{equation}
        \mo=\mt+\mr.                                             \tag{1}
\end{equation}

Si l'on représente maintenant $\mo$, en vertu de la condition minimale, comme somme directe d'idéaux directement indécomposables:
\[
        \mo=\mr_1+\ldots+\mr_n,
\]
alors les $\mr_i$ doivent être simples (= minimaux); car si, par exemple, $\mr_1$ n'était pas simple et si $\mt$ était un idéal minimal dans $\mr_1$, on aurait d'après (1):
\[
        \mr_1=\mt+\mr\cap\mr_1 \qquad \hbox{\(\S 4,3\)};
\]
donc $\mr_1$ serait décomposable, ce qui est impossible. Ainsi $\mo$ est complètement réductible.

Deuxièmement: \emph{un anneau sans radical avec condition minimale possède un élément unité.}

Pour montrer d'abord l'existence d'une unité à gauche, il suffit de montrer ceci: \emph{si un idéal à droite $\ma=c_1\mo\ne\mo$ possède une unité à gauche $c_1$, alors il existe un idéal à droite $c\mo$ de longueur plus grande, qui possède également une unité à gauche $c$.} En effet, puisque nous avons déjà montré la réductibilité complète, donc le fait que toutes les longueurs sont bornées, on arrive ainsi, en un nombre fini d'étapes, à une unité à gauche pour $\mo$ lui-même.

Soit donc $\ma=c_1\mo$ et $c_1^2=c_1$. Par la formule
\[
        r=c_1r+(r-c_1r)
\]
on obtient, comme ci-dessus, une décomposition de Peirce $\mo=c_1\mo+\mr$. Soit, comme ci-dessus, $c_2$ un élément idempotent de $\mr$, donc $c_2^2=c_2$, $c_1c_2=0$ (puisque $c_1$ annule tous les éléments de $\mr$). Si l'on pose maintenant $e_1=c_1$, $e_2=c_2-c_2c_1$, alors
\[
        e_1^2=e_1,
        \quad e_1e_2=0,
        \quad e_2e_1=0,
        \quad e_2c_2=e_2,
        \quad e_2\ne0,
        \quad e_2^2=e_2.
\]
D'après la remarque faite en note, $c=e_1+e_2$ devient donc une unité à gauche pour l'anneau
\[
        c\mo=e_1\mo+e_2\mo,
\]
qui, en tant qu'idéal à droite, a une longueur plus grande que $\ma$, puisque $e_2^2=e_2\ne0$.

Pour montrer que l'unité à gauche construite $e$ est aussi une unité à droite, nous faisons la décomposition de Peirce en idéaux à gauche:
\[
        \mo=\mo e+\ml.
\]
Nous avons $\ml e=0$, $\ml=e\ml$, donc $\ml^2=\ml e\ml=0$; par conséquent, comme par hypothèse il ne peut exister aucun idéal à gauche nilpotent, $\ml=0$, donc $e$ est aussi une unité à droite et par là une unité tout court.

\emph{Théorème. De la complète réductibilité à droite et de l'existence de l'unité résulte la complète réductibilité bilatère:}
\[
        \mo=\md_1+\cdots+\md_s,
        \qquad \md_i \hbox{ simple bilatère.}
\]

Comme dans la preuve précédente, il suffit de montrer que tout idéal minimal (bilatère) $\ma$ est un facteur direct. Comme idéal à droite, $\ma$ est un facteur direct:
\[
        \mo=\ma+\mr=e_1\mo+e_2\mo.
\]
La décomposition à gauche correspondante est
\[
        \mo=\mc+\ml=\mo e_1+\mo e_2.
\]
Il s'ensuit
\[
        \mr\mc\subseteqq \mr\ma\subseteqq \mr\cap\ma=0,
        \qquad
        \ml\ma=\mo e_2e_1\mo=0,
\]
\[
        (\ma\ml)^2=\ma\ml\ma\ml=0,
        \qquad \hbox{donc } \ma\ml=0;
\]
\[
        \mr=\mr\mo=(\mr\mc,\mr\ml)=\mr\ml,
        \qquad
        \ml=\mo\ml=(\ma\ml,\mr\ml)=\mr\ml.
\]
Donc $\mr=\ml$, donc $\mr$ est bilatère; ainsi $\ma$ est un facteur direct bilatère. À partir de là, comme dans la preuve de la complète réductibilité unilatère.

Les $\md_i$ sont aussi, comme anneaux, simples bilatères, puisque tous les idéaux dans $\md_i$ sont des idéaux dans $\mo$. Nous étudions leur structure.

\providecommand{\mo}{\mathfrak{o}}
\providecommand{\ml}{\mathfrak{l}}
\providecommand{\mr}{\mathfrak{r}}
\providecommand{\mA}{\mathfrak{A}}
\providecommand{\mZ}{\mathfrak{Z}}
\providecommand{\srcfntext}[2]{%
  \begingroup
  \renewcommand{\thefootnote}{#1}%
  \footnotetext{#2}%
  \endgroup}

\subsection*{\S{} 14. Anneaux simples bilatères complètement réductibles avec élément unité}

Soit $\mo$ un tel anneau,
\[
        \mo=\mr_1+\ldots+\mr_n=e_1\mo+\ldots+e_n\mo,
        \qquad \mr_i\hbox{ simple},
\]
alors on obtient
\[
        \mo=\ml_1+\ldots+\ml_n=\mo e_1+\ldots+\mo e_n,
        \qquad \ml_i\hbox{ directement indécomposable }\hbox{\(\S 9\)}.
\]
$\ml_i\mr_i$ est un idéal bilatère $\ne0$ (puisque $e_i^2$ appartient à $\ml_i\mr_i$), donc $=\mo$.

Si l'on pose $\mA_{ik}=\mr_i\ml_k$, on obtient
\[
        \mo=(\mr_1,\ldots,\mr_n)\cdot(\ml_1,\ldots,\ml_n)
        = (\ldots,\mA_{ij},\ldots)
\]
directement; car, si $0=\sum_{i,j}a_{ij}$, le fait que $a_{ij}$ appartienne à $\mr_i$ et que $\sum\mr_i$ soit directe donne
\[
        0=\sum_j a_{ij},
\]
donc encore, puisque $a_{ij}$ est dans $\ml_j$,
\[
        0=a_{ij}.
\]

De plus, $\mA_{ij}\mr_j=\mr_i\ml_j\mr_j=\mr_i\mo=\mr_i$, donc $\mA_{ij}\ne0$. Soit $a_{ij}\ne0$ pris dans $\mA_{ij}$. Alors
\[
        a_{ij}\mr_j\subseteqq\mr_i,
        \qquad
        a_{ij}\mr_j\ne0 \hbox{ parce que } a_{ij}e_j=a_{ij},
\]
donc
\[
        a_{ij}\mr_j=\mr_i.
\]
Si, pour chaque $r_j$ de $\mr_j$, on pose
\[
        a_{ij}r_j=r_i,
\]
alors l'application $r_j\to r_i$ est un homomorphisme d'opérateurs de $\mr_j$ vers $\mr_i$. Les éléments auxquels est associée la valeur nulle forment un idéal $\subset\mr_j$, donc l'idéal nul; donc un \emph{isomorphisme}. Par suite:

\emph{Deux quelconques idéaux à droite simples $\mr_i$ et $\mr_j$ apparaissant dans une représentation de $\mo$ sont isomorphes au sens des opérateurs; les isomorphismes sont induits par la multiplication par des éléments de $\mA_{ij}$.} Comme deux idéaux à droite simples distincts $\mr,\mr'$ apparaissent toujours ensemble dans au moins une représentation de $\mo$, deux tels idéaux quelconques sont donc isomorphes.

\emph{Tous les homomorphismes de $\mr_j$ dans $\mr_i$ sont induits par la multiplication par des éléments de $\mA_{ij}$.}

\emph{Démonstration.} Si $e_j\to a_{ij}$, alors $e_j^2\to a_{ij}e_j$, donc $a_{ij}=a_{ij}e_j$ appartient à $\mr_i\ml_j=\mA_{ij}$; en outre
\[
        \mr_j=e_j\mr_j\to a_{ij}\mr_j=\mr_i.
\]

\emph{En particulier, les homomorphismes de $\mr_i$ en lui-même sont induits par les multiplications par les éléments de $\mA_{ii}$.}

Au produit de deux éléments de $\mA_{ii}$ correspond le produit des homomorphismes, et à la somme la somme. (Pour la définition, voir \S{}1, 6.) Des $a_{ii}$ différents donnent des endomorphismes différents, car $e_i$ passe en $a_{ii}e_i=a_{ii}$. Ainsi l'anneau $\mA_{ii}$ est isomorphe à l'anneau des endomorphismes de $\mr_i$.

Mais l'anneau des endomorphismes d'un idéal simple est un corps gauche (\S{}2), donc $\mA_{ii}$ est un corps gauche. Comme, de plus, tous les $\mr_i$ sont isomorphes au sens des opérateurs, leurs anneaux d'endomorphismes sont aussi isomorphes comme anneaux: $\mA_{ii}$ est déterminé de manière univoque par $\mo$, à isomorphisme d'anneaux près. Les différents $\mA_{ii}$ possibles ne sont que différentes réalisations concrètes de l'anneau des endomorphismes, défini abstraitement, des idéaux à droite simples.

\emph{Théorème principal. Tout anneau $\mo$ complètement réductible, simple bilatère, avec élément unité, est isomorphe à l'anneau des matrices de degré $n$ sur un corps gauche $\mathsf{K}$. (Le corps $\mathsf{K}$ est isomorphe au corps gauche des endomorphismes des idéaux à droite de $\mo$.)}

\emph{Démonstration\textsuperscript{13a)}. Construction des unités matricielles $c_{ik}$.} Soit $\Gamma_{11}$ l'endomorphisme identité de $\mr_1$, soit $\Gamma_{i1}$ un isomorphisme arbitraire $\Gamma_{i1}\mr_1=\mr_i$, et enfin
\[
        \Gamma_{ik}=\Gamma_{i1}\Gamma_{k1}^{-1}.
\]
On a alors, en général,
\[
        \Gamma_{ik}\Gamma_{kl}=\Gamma_{il}.
\]
Si $\Gamma_{ik}$ est induit par la multiplication par l'élément $c_{ik}$ de $\mA_{ik}$, on a de plus:
\[
        c_{ik}=e_i c_{ik}=c_{ik}e_k,
        \qquad
        c_{ik}c_{kl}e_l=c_{il}e_l,
\]
donc
\[
        c_{ik}c_{kl}=c_{il},
        \qquad
        c_{ij}c_{kl}=c_{ij}e_j e_k c_{kl}=0\quad(j\ne k).
\]
Les $c_{ik}$ sont donc des unités matricielles (\S{}7).

\srcfntext{13a)}{Récemment B. L. van der Waerden a trouvé une preuve plus simple, opérant directement avec le corps gauche abstrait des endomorphismes, qui doit être donnée dans son livre d'algèbre (Grundlehren d. Math. Wiss., Berlin: Julius Springer). [11. 6. 29.]}

\emph{Construction de $\mathsf{K}$.} La correspondance
\[
        a_{ii}=c_{i1}a_{11}c_{1i}
\]
associe à chaque $a_{11}$ de $\mA_{11}$ un ``élément conjugué'' $a_{ii}$ de $\mA_{ii}$.

Cette correspondance est un isomorphisme d'anneaux, car les isomorphismes $\Delta,\Gamma$ correspondant aux $a$ et aux $c$ sont liés par
\[
        \Delta_{ii}=\Gamma_{i1}\Delta_{11}\Gamma_{i1}^{-1},
\]
relation qui, comme on sait, représente un isomorphisme d'anneaux. Formons maintenant les éléments
\[
        \alpha=a_{11}+\ldots+a_{nn}=a_{11}+c_{21}a_{11}c_{12}+\ldots+c_{n1}a_{11}c_{1n}.
\]
À chaque $a_{11}$ est associé un $\alpha$ (biunivoquement à cause de la somme directe). L'ensemble des $\alpha$ est un corps $K\sim\mA_{11}$; car de
\[
        \alpha=a_{11}+\ldots+a_{nn},
        \qquad
        \beta=b_{11}+\ldots+b_{nn},
\]
on obtient
\[
        \alpha+\beta=(a_{11}+b_{11})+\ldots+(a_{nn}+b_{nn}),
        \qquad
        \alpha\cdot\beta=a_{11}b_{11}+\ldots+a_{nn}b_{nn};
\]
donc la correspondance $a_{11}\to\alpha$ est un isomorphisme.

En outre
\[
        c_{ik}\alpha=c_{ik}a_{kk}=c_{ik}c_{k1}a_{11}c_{1k}=c_{i1}a_{11}c_{1k},
\]
\[
        \alpha c_{ik}=a_{ii}c_{ik}=c_{i1}a_{11}c_{1i}c_{ik}=c_{i1}a_{11}c_{1k},
\]
donc chaque $\alpha$ commute avec tous les $c_{ik}$.

Enfin, les mêmes formules donnent
\[
        c_{ik}\mathsf{K}=c_{ik}\mA_{kk}=c_{ik}\mr_k\ml_k=\mr_i\ml_k=\mA_{ik},
\]
donc
\[
        \mo=\sum c_{ik}\mathsf{K}.
\]

\emph{La correspondance.} Si l'on représente chaque élément de $\mo$ sous la forme $\sum c_{ik}\alpha_{ik}$, et si l'on associe à cet élément la matrice $(\alpha_{ik})$, alors la correspondance est un isomorphisme, comme il résulte immédiatement de la définition de la multiplication des matrices. Le théorème principal est ainsi démontré.

\emph{Réciproque. L'anneau des matrices de degré $n$ sur un corps gauche $A$ est simple bilatère et complètement réductible; $A$ est isomorphe à l'anneau des endomorphismes des idéaux à droite, et $Ae$ est le $\mathsf{K}$ construit dans le théorème précédent.}

\emph{Démonstration.} Les unités matricielles (\S{}7) soient $c_{ik}$. Soit
\[
        \mr_i=\sum_k c_{ik}A.
\]
Alors manifestement $\mo=\mr_1+\ldots+\mr_n$. Les $\mr_i$ sont des idéaux à droite simples; car tout élément $\sum c_{ik}\alpha_k\ne0$ engendre tout $\mr_i$. En effet, si $\alpha_j\ne0$, alors
\[
        \Bigl(\sum c_{ik}\alpha_k\Bigr)\cdot(c_{jl}\alpha_j^{-1})=c_{il},
\]
et $\sum c_{il}\beta_l$ parcourt tous les éléments de $\mr_i$.

Les $\mr_i$ sont isomorphes au sens des opérateurs: $\mr_i=c_{ik}\mr_k$.

Il en résulte que $\mo$ est bilatéralement indécomposable, donc, d'après le théorème du \S{}13, en vertu de la complète réductibilité déjà démontrée -- et puisqu'il existe un élément unité -- simple bilatère (comme on le voit aussi au fait qu'un $\sum\sum c_{ik}\alpha_{ik}\ne0$ engendre tout l'idéal bilatère $\mo$).

En outre
\[
        \mr_i=c_{ii}\mo,
        \qquad
        \ml_i=\mo c_{ii},
\]
\[
        e=\sum c_{ii},
        \qquad c_{ii}=e_i,
\]
\[
        \mA_{ii}=\mr_i\cap\ml_i=c_{ii}\Lambda\simeq\Lambda.
\]
Si l'on pose enfin $a_{11}=c_{11}\lambda$, on obtient
\[
        \alpha=a_{11}+c_{21}a_{11}c_{12}+\ldots
        =c_{11}\lambda+c_{22}\lambda+\ldots=e\lambda,
\]
ce qui démontre tout.

\emph{Le centre de $\mo$ est le centre de $\mathsf{K}$, donc un corps.}

\emph{Démonstration.} $\mZ(\mathsf{K})$ se compose de tous les $\zeta$ qui commutent avec tous les $x$. Ceux-ci commutent aussi avec tous les $c_{ik}$, donc aussi avec les sommes $\sum c_{ik}a_{ik}$. Donc:
\[
        \mZ(\mathsf{K})\subseteqq\mZ(\mo).
\]
Réciproquement: soit $z$ dans $\mZ(\mo)$,
\[
        z=\sum_\mu\sum_\nu c_{\mu\nu}\gamma_{\mu\nu};
\]
\[
        zc_{ik}=c_{ik}z,
        \qquad
        \sum_\mu c_{\mu k}\gamma_{\mu i}=\sum_\nu c_{i\nu}\gamma_{k\nu};
\]
donc
\[
        \gamma_{\mu i}=0\ (\mu\ne i),
        \qquad
        c_{ik}\gamma_{ii}=c_{ik}\gamma_{kk},
        \qquad
        \gamma_{ii}=\gamma_{kk}=\gamma;
\]
\[
        z=\sum c_{ii}\gamma=e\gamma=\gamma.
\]
Donc $z$ appartient à $\mathsf{K}$, et commute avec tous les $x$, donc $z$ appartient à $\mZ(\mathsf{K})$.

Comme un anneau de matrices est naturellement aussi complètement réductible à gauche, et que tout anneau complètement réductible à droite avec élément unité est somme directe d'anneaux de matrices, il s'ensuit:

\emph{Tout anneau complètement réductible à droite avec élément unité l'est aussi à gauche, et réciproquement.}

Et: \emph{Le centre d'un anneau complètement réductible avec élément unité est somme directe de corps commutatifs qui correspondent aux anneaux de matrices simples bilatères.}

Il reste le théorème suivant\textsuperscript{14)}:

\emph{Deux décompositions différentes $\mo=\sum c_{ik}\mathsf{K}$, $\mo=\sum c'_{ik}\mathsf{K}'$ passent l'une dans l'autre par des automorphismes intérieurs $\alpha'\to x^{-1}\alpha x$, $c_{ik}=x^{-1}c_{ik}x$.}
\srcfntext{14)}{Cf. E. Artin, loc. cit., p. 258.}

\emph{Démonstration.} Soient $\mo=\mr_1+\ldots+\mr_n=\mr_1'+\ldots\mr_n'$, et soient $a,b$ des éléments dont les multiplications induisent deux isomorphismes réciproques des idéaux $\mr_1,\mr_1'$:
\[
        \mr_1=a\mr_1',
        \qquad
        \mr_1'=b\mr_1,
        \qquad
        ab=c_{11},
        \qquad
        ba=c'_{11}.
\]
Si l'on pose
\[
        x=\sum_i c_{i1}ac'_{1i},
        \qquad
        y=\sum_i c'_{i1}bc_{1i},
\]
on obtient
\[
        xy=e,
        \qquad
        y=x^{-1},
        \qquad
        x^{-1}c_{ik}x=c'_{ik},
        \qquad
        x^{-1}\alpha x=\alpha'.
\]

\providecommand{\mo}{\mathfrak{o}}
\providecommand{\mM}{\mathfrak{M}}

\subsection*{Chapitre III. Théorie des modules et des représentations}

\subsection*{\S{} 15. Représentations et modules de représentation\textsuperscript{15)}}
Soit $\mo$ un anneau, et $\mathsf{K}$ un anneau avec élément unité. (Dans les applications ultérieures, $\mathsf{K}$ est toujours un corps.)

Une représentation de degré $n$ de $\mo$ dans $\mathsf{K}$ est une homomorphie
\[
        \mo\sim\mathfrak D,
\]
où $\mathfrak D$ est un anneau de matrices de degré $n$ sur $\mathsf{K}$.

Par module de représentation de $\mo$ relativement à $\mathsf{K}$, on entend un bimodule $\mM$ (\S{}1, 5), qui est module à gauche sur $\mo$ et module à droite sur $\mathsf{K}$:
\[
        \mo\mM\subseteqq\mM,
        \qquad
        \mM \mathsf{K}\subseteqq\mM,
\]
qui est en outre somme directe d'un nombre fini de modules $\mathsf{K}$ monogènes:
\[
        \mM=x_1\mathsf{K}+\ldots+x_n\mathsf{K};
\]
et où l'élément unité de $\mathsf{K}$ est opérateur unité.

Tout module de représentation conduit à une représentation. Soit $c$ dans $\mo$ et
\[
        cx_k=\sum x_i\gamma_{ik},
\]
ou encore
\[
        c(x_1,\ldots,x_n)=(x_1,\ldots,x_n)C,
        \qquad C=(\gamma_{ik}).
\]
Alors les matrices $C$ forment une représentation pour $c$, car
\[
        (b+c)x_k=\sum x_i(\beta_{ik}+\gamma_{ik}),
\]
\[
%
        bcx_k
        =b\sum x_j\gamma_{jk}
          =\sum b x_j\gamma_{jk}
          =\sum\sum x_i\beta_{ij}\gamma_{jk}
        =\sum x_i(\sum\beta_{ij}\gamma_{jk})
%
\]
ou
\[
        bc(x_1,\ldots,x_n)=(x_1,\ldots,x_n)BC.
\]

Réciproquement: toute représentation $\mathfrak D$ de $\mo$ appartient à un module de représentation, et même à une base déterminée de celui-ci. On entend en effet par $\mM$ l'ensemble des formes linéaires formelles en $x_1,\ldots,x_n$,
\[
        y=\sum x_i\alpha_i.
\]
Alors $\mM$ est un module à droite sur $\mathsf{K}$. Définissons en outre, lorsque la matrice $C=(\gamma_{ik})$ est associée à l'élément $c$,
\[
        cx_k=\sum x_i\gamma_{ik},
        \qquad
        c\biggl(\sum_k x_k\alpha_k\biggr)=\sum_i x_i\sum_k\gamma_{ik}\alpha_k.
\]
Des relations d'homomorphie
\[
        c+d\to C+D,
        \qquad cd\to CD
\]
on obtient
\[
        (c+d)x_i=cx_i+dx_i,
        \qquad cdx_i=c(dx_i),
\]
et il en va de même pour les sommes $\sum x_i\alpha_i$. Les autres propriétés de bimodule
\[
        c(y+z)=cy+cz,
        \qquad c(y\alpha)=(cy)\alpha
\]
sont triviales. Ainsi $\mM$ est un module de représentation qui, par construction, appartient exactement à la représentation $\mo\sim\mathfrak D$.

Soit maintenant $(y_1,\ldots,y_n)$ une autre base du même module de représentation:
\[
        (y_1,\ldots,y_n)=(x_1,\ldots,x_n)\cdot P,
        \qquad
        (x_1,\ldots,x_n)=(y_1,\ldots,y_n)\cdot P^{-1}.
\]
Si $b\sim B$ relativement à la base $x$, on a
\[
        b(y_1,\ldots,y_n)=b(x_1,\ldots,x_n)P
        =(x_1,\ldots,x_n)BP
        =(y_1,\ldots,y_n)P^{-1}BP.
\]
On appelle les représentations $b\sim B$ et $b\sim P^{-1}BP$ équivalentes, et on les compte dans la même classe de représentations. Comme toute matrice régulière $P$ transforme de nouveau une base de $\mM$ en une base, il est démontré:

\emph{Tout module de représentation conduit de manière univoque à une classe de représentations déterminée.}\textsuperscript{15a)}

{\footnotesize\noindent 15) Les modules de représentation relativement à des corps commutatifs se trouvent d'abord chez W. Krull, \emph{Theorie und Anwendung der verallgemeinerten Abelschen Gruppen}, Sitzungsber. Heidelberger Ak. 1926, 1. Abhandl.\par
15a) Addition lors de la correction (14 avril 1929). Comme B. L. van der Waerden me le communique, on peut obtenir directement un lien invariant, indépendant du choix spécial de la base, en séparant les notions: transformation linéaire et matrice. Une transformation linéaire est un homomorphisme de deux modules de formes linéaires; une matrice est l'expression (la représentation) de cet homomorphisme pour un choix déterminé de base. I. Tout module de représentation associe univoquement au domaine des multiplicateurs $\mo$ un système homomorphe de transformations linéaires du module en lui-même. En effet, d'après la fin du \S{}2, les multiplicateurs à gauche d'un bimodule $\mM$ engendrent des homomorphismes d'opérateurs de $\mM$ en lui-même relativement aux opérateurs à droite (ici la multiplication par $\mathsf{K}$); et la correspondance est homomorphe. II. Tout système, homomorphe à $\mo$, de transformations linéaires d'un module de formes linéaires en lui-même conduit à un module de représentation. En effet, d'après encore la fin du \S{}2, ces homomorphismes, pris comme multiplicateurs à gauche, donnent un bimodule; il en va donc de même relativement au domaine des multiplicateurs $\mo$ si l'on pose $cm=\Gamma m$ pour tout $m$ de $\mM$ et tout $c$ de $\mo$, où $c\to\Gamma$ est l'homomorphisme d'anneaux initialement donné. Aux modules de représentation isomorphes comme bimodules -- c'est-à-dire par un isomorphisme compatible avec les actions à droite et à gauche -- correspondent les mêmes transformations, et aux modules qui ne sont pas isomorphes comme bimodules des transformations différentes. Si, pour un choix déterminé de base, on exprime chaque transformation par une matrice, I donne la représentation $\mo\sim\mathfrak D$, tandis que II affirme que toute représentation de cette sorte appartient à un module de représentation et à une base déterminée de celui-ci. On retrouve ainsi, sans aucun calcul, la correspondance entre modules de représentation et classes de représentations.\par}

Il est clair que deux modules de représentation isomorphes au sens des opérateurs correspondent à la même classe de représentations. Mais réciproquement aussi: si deux modules de représentation $(x_1,\ldots,x_n)$ et $(\bar x_1,\ldots,\bar x_n)$ engendrent la même représentation, alors la correspondance
\[
        x_i\to\bar x_i,
        \qquad
        \sum_i x_i\alpha_i\to\sum_i\bar x_i\alpha_i
\]
fournit un isomorphisme. Car la règle de multiplication est entièrement déterminée par les matrices $C$.

Cas particulier: si deux bases de $\mM$ engendrent la même représentation, elles se déduisent l'une de l'autre par un isomorphisme d'opérateurs de $\mM$ sur lui-même.

Ou encore: si $C=PCP^{-1}$ pour tous les $C$ et pour une matrice $P$ fixe, alors la correspondance
\[
        y_i\to x_i,
        \qquad y_i=\sum_kx_i\pi_{ik},
\]
est un isomorphisme de $\mM$ sur lui-même. Réciproquement: tout isomorphisme $y_i\to x_i$ du bimodule $\mM$ sur lui-même est réalisé par une matrice $P$ qui commute avec tous les $C$.

On peut aussi étendre la notion de représentation aux anneaux sur lesquels un domaine commutatif de multiplicateurs $\mathsf{P}$ agit centralement (\S{}8). Dans ce cas on ne prend que des anneaux de représentation $\mathsf{K}$ qui contiennent $\mathsf{P}$ dans leur centre, et l'on exige de la représentation non seulement qu'elle soit un homomorphisme d'anneaux $\mo\sim\mathfrak D$, mais même qu'elle soit un homomorphisme d'opérateurs: de $a\to A$ doit suivre
\[
        a\varrho\to A\varrho\qquad(\varrho\in P).
\]
Pour les modules de représentation, cette exigence signifie la règle de calcul supplémentaire
\[
        am\cdot\varrho=a(m\varrho)=a\varrho\cdot m.
\]
L'action de $P$ sur le module commute alors avec celle de l'anneau.

\providecommand{\mo}{\mathfrak{o}}
\providecommand{\mM}{\mathfrak{M}}
\providecommand{\mU}{\mathfrak{U}}
\providecommand{\mA}{\mathfrak{A}}
\providecommand{\mB}{\mathfrak{B}}

\subsection*{\S{} 16. Représentations réductibles}
Si $\mU$ est un sous-module du module de représentation $\mM$, et s'il est possible de choisir une base pour $\mM$ composée d'une base $z_1,\ldots,z_t$ de $\mU$, complétée par $y_1,\ldots,y_r$, donc
\[
        \mM=y_1\mathsf{K}+\ldots+y_r\mathsf{K}+z_1\mathsf{K}+\ldots+z_t\mathsf{K}\textsuperscript{16)}
\]
alors les représentations ont la forme suivante:
\[
        C=\begin{pmatrix}R&0\\ S&T\end{pmatrix},
\]
où les matrices $T$, à elles seules, forment une représentation de degré $t$ portée par $\mU$, et les matrices $R$ une représentation de degré $r$ portée par le quotient $\mM/\mU$.

\emph{Démonstration.} On a
\[
        (cy_1,\ldots,cy_r,cz_1,\ldots,cz_t)
        =(y_1,\ldots,y_r,z_1,\ldots,z_t)\cdot C.
\]
Comme les $cz_i$, en tant qu'éléments de $\mU$, s'expriment par les seuls $z_i$, le bloc supérieur droit de $C$ est nul. Si l'on nomme les autres parties de $C$, dans l'ordre ci-dessus, $R,S,T$, alors on a en particulier
\[
        (cz_1,\ldots,cz_t)=(z_1,\ldots,z_t)\cdot T;
\]
les $T$ forment donc une représentation portée par $\mU$. De plus,
\[
        (cy_1,\ldots,cy_r)\equiv (y_1,\ldots,y_r)R\pmod{\mU},
\]
tandis que les $y_i$ forment une base linéairement indépendante modulo $\mU$. Ainsi les $R$ forment une représentation engendrée par $\mM/\mU$.

Réciproquement, si une représentation ``réductible''
\[
        C=\begin{pmatrix}R&0\\ S&T\end{pmatrix}
\]
est donnée, où $R$ et $T$ sont des matrices carrées, alors, dans le module de représentation correspondant, les produits de chaque $c$ par les $t$ derniers éléments de base $z_1,\ldots,z_t$ s'expriment par ces éléments seuls, c'est-à-dire que $\mU=(z_1,\ldots,z_t)$ est un sous-module.

{\footnotesize\noindent 16) Autrement dit: si $\mU$ est facteur direct comme module $\mathsf{K}$. Si $\mathsf{K}$ est un corps, cette hypothèse est toujours satisfaite.\par}

\emph{Corollaire.} Soit
\[
        \mM\supset \mU_1\supset \mU_2\supset\cdots\supset \mU_e=0
\]
une suite de composition pour $\mM$, et soit chaque $\mU_i$ facteur direct du précédent comme module $K$, de sorte qu'on puisse choisir une base
\[
        z_{11},\ldots,z_{1r_1},\ z_{21},\ldots,z_{2r_2},\ldots,
        z_{e1},\ldots,z_{er_e}
\]
telle que les $z_{ik}$ avec $i>\nu$ forment une base de $\mU_\nu$. Alors les représentations portées par $\mM$ ont la forme suivante:
\[
 C=\begin{pmatrix}
 R_{11}&0&\ldots&0\\
 R_{12}&R_{22}&\ldots&0\\
 \vdots&\vdots&\ddots&\vdots\\
 R_{1e}&R_{2e}&\ldots&R_{ee}
 \end{pmatrix},                                      \tag{2}
\]
et les $R_{\nu\nu}$ forment des représentations portées par les facteurs de composition $\mU_{\nu-1}/\mU_\nu$.

Les représentations particulières $R_{\nu\nu}$ sont irréductibles, puisque les facteurs de composition $\mU_{\nu-1}/\mU_\nu$ sont simples. Si l'on suppose que $K$ est un corps, alors, réciproquement aussi, toute représentation réduite autant que possible sous la forme (2) conduit à une suite de composition. D'après le théorème de Jordan-Hölder, les facteurs de composition sont déterminés univoquement à isomorphie d'opérateurs près; donc les $R_{\nu\nu}$ sont déterminés univoquement à représentations équivalentes près et à l'ordre près.

Si $\mM=\mA+\mB$ est la somme directe de deux modules de représentation $\mA=(y_1,\ldots,y_r)$, $\mB=(z_1,\ldots,z_s)$, alors la représentation écrite dans la base $(y_1,\ldots,y_r,z_1,\ldots,z_s)$ a manifestement la forme
\[
        C=\begin{pmatrix}R&0\\0&S\end{pmatrix},
\]
où $R$ signifie la représentation de l'élément $c$ sur $\mA$, et $S$ sa représentation sur $\mB$; et réciproquement. Si l'on représente d'abord le module $\mM$ comme somme directe de constituants directement indécomposables, puis que l'on cherche pour ceux-ci les suites de composition comme ci-dessus, alors, pour un choix convenable de base, la représentation prend la forme
\[
\begin{array}{@{}r@{\quad}c@{}}
(3)&
C=\left(
\begin{array}{ccc}
\boxed{\begin{matrix}
R_{11}&0\\
\vdots&\ddots\\
R_{ik}&R_{rr}
\end{matrix}} & & 0\\[1.4em]
& \boxed{\begin{matrix}
S_{11}&0\\
\vdots&\ddots\\
S_{ik}&S_{ss}
\end{matrix}} & \\ & & \ddots \\
0 & & \boxed{\begin{matrix}
T_{11}&0\\
\vdots&\ddots\\
T_{ik}&T_{tt}
\end{matrix}}
\end{array}
\right)
\end{array}
\]

Si de nouveau $\mathsf{K}$ est un corps, toute représentation réduite autant que possible sous la forme (3) conduit aussi réciproquement à une représentation du module par des sommants directement indécomposables et par des facteurs de composition à l'intérieur de ceux-ci. D'après le théorème de Krull et Otto Schmidt, les classes des constituants de représentation directement indécomposables sont déterminées univoquement, à l'ordre près.

Si le module $\mM$ est complètement réductible, tous les blocs se composent d'une unique représentation irréductible, et l'on appelle la représentation complètement réductible.

\providecommand{\mo}{\mathfrak{o}}
\providecommand{\mM}{\mathfrak{M}}
\providecommand{\ma}{\mathfrak{a}}
\providecommand{\mb}{\mathfrak{b}}

\subsection*{\S{} 17. Décomposition en somme directe des modules de représentation pour les anneaux à élément unité}
Soit $\mM$ un $\mo$-module, $\mo$ étant un anneau avec unité. Alors
\[
        \mM=\mM_e+\mM_0,
\]
où, sur $\mM_e$, l'unité est opérateur unité, et, sur $\mM_0$, opérateur nul. ($\mM_e$ et $\mM_0$ sont aussi des modules à droite sur $\mathsf{K}$ si $\mM$ l'est.)

\emph{Démonstration.} Chaque $m$ de $\mM$ est représentable sous la forme
\[
        m=em+(m-em).
\]
Soit $\mM_e$ l'ensemble des $em$, et $\mM_0$ l'ensemble des $m-em$. Alors $\mM_e$ et $\mM_0$ sont évidemment des groupes additifs et des modules à droite sur $\mathsf{K}$, si $\mM$ l'est. En outre
\[
        r em=erm,
\]
donc $\mM_e$ est aussi un $\mo$-module; de même
\[
        r(m-em)=rm-rem=rm-rm,
\]
donc $\mM_0$ est un $\mo$-module. Pour les $em$, $e$ est opérateur unité; pour les $(m-em)$, opérateur nul. Donc seul le zéro appartient à la fois à $\mM_e$ et à $\mM_0$; la somme $\mM_e+\mM_0$ est directe.

S'il s'agit de modules de représentation, $\mM_0$ engendre la représentation triviale, où à chaque élément est associé le zéro. Si l'on détache le sommant $\mM_0$, il reste la représentation portée par $\mM_e$, où à l'élément unité est associée la matrice unité. Nous pouvons et voulons donc désormais nous restreindre à de telles représentations.

Soit
\[
        \mo=\ma_1+\ldots+\ma_s
\]
une somme directe d'idéaux bilatères,
\[
        e=e_1+\ldots+e_s
\]
la décomposition de $e$, et soit $\mM$ un $\mo$-module où $e$ est opérateur unité. Alors
\[
        \mM=\ma_1\mM+\ldots+\ma_s\mM
\]
est directe. Évidemment $\mM=\mo\mM=(\ma_1\mM,\ldots,\ma_s\mM)$. Posons $\mb_i=\ma_1+\cdots+\widehat{\ma_i}+\cdots+\ma_s$; alors
\[
        \mM=(\ma_i\mM,\mb_i\mM),
\]
et cette somme est directe, puisque $e_i$ est opérateur unité pour $\ma_i\mM$, et opérateur nul pour $\mb_i\mM$.

Sur le fondement de ce théorème, nous nous restreignons le plus souvent aux représentations d'anneaux bilatéralement indécomposables.

\providecommand{\mo}{\mathfrak{o}}
\providecommand{\mM}{\mathfrak{M}}
\providecommand{\ml}{\mathfrak{l}}

\subsection*{\S{} 18. Théorie des modules et des représentations des anneaux complètement réductibles}
Soit $\mo$ complètement réductible et simple bilatère, et soit $\mM$ un $\mo$-module fini pour lequel l'élément unité de $\mo$ est opérateur unité. Alors $\mM$ est complètement réductible, et les constituants simples sont isomorphes, comme modules à opérateurs, aux idéaux à gauche simples $\ml_i$.

\emph{Démonstration.} Soit
\[
        \mo=\ml_1+\ldots+\ml_n,
        \qquad
        \mM=(\mo m_1,\ldots,\mo m_k),
\]
donc
\[
        \mM=(\ldots,\ml_i m_k,\ldots).                     \tag{4}
\]
Les $\ml_i m_k$ qui sont différents de zéro sont isomorphes à $\ml_i$, car la correspondance $a\mapsto am_k$ est un isomorphisme d'opérateurs. Les $\ml_i m_k$ sont donc simples. Si l'on omet dans la représentation (4) les $\ml_i m_k$ qui sont déjà contenus dans la somme des précédents, la somme devient directe.

\emph{Conséquence.} Si, en outre, $\mM$ est directement indécomposable, alors $\mM$ est simple et isomorphe à un $\ml_i$.

Soit maintenant $\Delta$ le corps gauche des endomorphismes de ce $\mM$ simple, et $\Lambda$ celui de $\ml_i$; alors, en vertu de l'isomorphisme d'opérateurs de $\mM$ et de $\ml_i$, on a l'isomorphisme d'anneaux $\Delta\simeq\Lambda$. D'après le \S{}1, on peut considérer $\mM$ et $\ml$ comme des bimodules ayant respectivement $\Delta$ et $\Lambda$ comme domaines à droite; et ces bimodules sont aussi des modules de représentation. Les représentations portées par ces bimodules se déduisent l'une de l'autre par l'isomorphisme d'anneaux $\Delta\simeq\Lambda$. Pour l'étude de cette classe de représentations, nous pouvons donc nous restreindre à la représentation portée par $\ml_i$ dans son corps gauche des endomorphismes, car $\ml_i$, et par conséquent $\mM$, est de rang fini relativement à son corps gauche des endomorphismes (\S{}14), et l'élément unité de celui-ci est l'opérateur unité.

Supposons spécialement que, pour $\ml_1$, on ait choisi une base formée d'unités matricielles
\[
\ml_1=(c_{11},\ldots,c_{n1})
\]
et prenons comme réalisation du corps gauche des endomorphismes $\Lambda$ le sous-corps de $\mo$, désigné plus haut (\S{}14) par $\mathsf{K}$ (dans nos considérations antérieures, il faut remplacer les idéaux à droite par des idéaux à gauche et écrire en conséquence les endomorphismes à droite). Si l'on écrit
\[
a=\sum_{i,k}c_{ik}\alpha_{ik},
\]
alors
\[
a\longmapsto(\alpha_{ik})
\]
est la représentation dans $K$ portée par $\ml_1$. En effet,
\[
a\cdot c_{k1}
=\left(\sum\sum c_{ij}\alpha_{ij}\right)\cdot c_{k1}
=\sum c_{ik}\alpha_{ik}c_{k1}
=\sum c_{i1}\alpha_{ik},
\]
ou encore
\[
a(c_{11},\ldots,c_{n1})
=(c_{11},\ldots,c_{n1})(\alpha_{ik}).
\]
La représentation portée par $\ml_1$ est donc donnée par la matrice $(\alpha_{ik})$; pour $\ml_\nu=(c_{1\nu},\ldots,c_{n\nu})$, on obtient exactement la même représentation.

\[
\left\{
\begin{aligned}
&\text{Soit }\Gamma\text{ un sous-corps de }K\text{ tel que }K\text{ soit de rang fini }t\text{ sur }\Gamma:\\
&\qquad K=z_1\Gamma+\cdots+z_t\Gamma.\\
&\text{Alors }K\text{ est son propre module de représentation sur }\Gamma.\text{ Les éléments }\beta\text{ de }K\\
&\text{sont représentés par des matrices }B\text{ à coefficients dans }\Gamma,\text{ obtenues par}\quad
\beta z_j=\sum_i z_i\beta_{ij},\quad B=(\beta_{ij}).
\end{aligned}
\right.
\]
Si l'on considère maintenant $\ml_1$ comme module de représentation sur le sous-corps $\Gamma$, on obtient:

\emph{La représentation de $\mo$ dans $\Gamma$ portée par $\ml_1$ s'obtient en remplaçant, dans la représentation de $\mo$ dans $\mathsf{K}$ donnée ci-dessus,}
\[
        a\longmapsto(\alpha_{ik}),
\]
\emph{chaque élément $\alpha_{ik}$ par la matrice $A_{ik}$ qui lui correspond dans la représentation de $\mathsf{K}$ dans $\Gamma$.}

\emph{Démonstration.}
\[
        \ml_1=\sum_i c_{i1}\mathsf{K}=\sum_{i,j}c_{i1}\varkappa_j\Gamma.
\]
Relativement à la base $(\ldots,c_{i1}\varkappa_j,\ldots)$, on a
\[
c_{ik}\to
\begin{array}{c@{\quad}c}
 & \begin{array}{ccc} \boxed{\scriptstyle 1} & \boxed{\scriptstyle k} & \boxed{\scriptstyle n} \end{array}\\[-.2em]
\begin{array}{c}\boxed{\scriptstyle 1}\\ \boxed{\scriptstyle i}\\ \boxed{\scriptstyle n}\end{array}
&
\begin{pmatrix}
0&\ldots&0&\ldots&0\\
\vdots&&\vdots&&\vdots\\
0&\ldots&E_t&\ldots&0\\
\vdots&&\vdots&&\vdots\\
0&\ldots&0&\ldots&0
\end{pmatrix}
\end{array},
\]
où $E_t$ est la matrice unité et $0$ la matrice nulle de degré $t$, tandis que
\[
        \alpha\longmapsto
        \begin{pmatrix}
        A& &0\\
         &\ddots&\\
        0& &A
        \end{pmatrix}.
\]
Par conséquent,
\[
        a=\sum c_{ik}\alpha_{ik}\to
        \begin{pmatrix}
        A_{11}&\ldots&A_{1n}\\
        \vdots&&\vdots\\
        A_{n1}&\ldots&A_{nn}
        \end{pmatrix}.
\]

Pour un module $\mM$ quelconque, on ne fait de nouveau que transporter cette représentation par un isomorphisme d'anneaux. Si $\mo$ est somme directe d'anneaux simples bilatères, ces anneaux sont représentés tour à tour isomorphiquement, les autres étant représentés par zéro.

\textbf{Remarque.} À isomorphisme près, les représentations étudiées ici, dans le corps gauche des endomorphismes des idéaux à gauche et dans ses sous-corps sur lesquels il est de rang fini, sont les seules qui puissent être portées par des $\mo$-modules simples (ou par des idéaux à gauche simples). En effet, d'après une remarque antérieure (\S{}2), si l'on considère le $\mo$-module $\mM$ comme un bimodule relativement à un corps $\mathsf{K}$, les éléments de $\mathsf{K}$ induisent des homomorphismes du module; il existe donc nécessairement un homomorphisme de $\mathsf{K}$ sur un sous-corps du corps gauche des endomorphismes de ce $\mo$-module. Comme l'élément unité est opérateur unité, cet homomorphisme est un isomorphisme. En outre, le corps gauche entier des endomorphismes doit être de rang fini sur ce sous-corps; sinon le module de représentation serait lui aussi de rang infini, ce qui est exclu.

Enfin, mentionnons que les représentations irréductibles étudiées ici ne les épuisent pas toutes: ce ne sont que celles qui sont portées par des $\mo$-modules simples. Il existe en général d'autres modules de représentation --- par exemple un corps $\Omega$, considéré comme bimodule relativement à un sous-corps $\Sigma$ pris comme domaine à gauche et à $\Omega$ lui-même pris comme domaine à droite --- qui sont réductibles comme $\mo$-modules mais irréductibles comme bimodules, et qui conduisent par conséquent eux aussi à des représentations irréductibles.

\providecommand{\mo}{\mathfrak{o}}
\providecommand{\mc}{\mathfrak{c}}
\providecommand{\mM}{\mathfrak{M}}
\providecommand{\ml}{\mathfrak{l}}

\subsection*{\S{} 19. Les facteurs de composition simples dans les modules et les modules de représentation}
Soit $\mo$ un anneau satisfaisant aux conditions maximale et minimale pour les idéaux à gauche, et soit $\mc$ son idéal nilpotent maximal. Alors:

\emph{Tout $\mo$-module simple $\mM$ est soit annulé par $\mo$, soit isomorphe à un idéal à gauche simple $\ml$ de l'anneau sans radical $\mo/\mc$. Dans le second cas, le module est annulé par tous les idéaux bilatères de $\mo/\mc$ qui ne contiennent pas $\ml$, et son domaine absolu de multiplicateurs (\S{}1) est isomorphe à l'anneau simple bilatère qui contient $\ml$.}

\textbf{Démonstration.} Soit $\mc^{\varrho}=0$. On doit avoir $\mc\mM=0$, car sinon $\mc\mM=\mM$, et il s'ensuivrait
\[
        \mM=\mc\mM=\mc^2\mM=\ldots=0.
\]
On peut donc aussi considérer $\mM$ comme un $\mo/\mc$-module: tous les éléments d'une même classe modulo $\mc$ donnent le même résultat lorsqu'ils multiplient un élément de $\mM$. Comme anneau sans radical, $\mo/\mc$ possède un élément unité; $\mM$ est donc somme directe d'un module annulé par $\mo/\mc$ et d'un module sur lequel l'élément unité est opérateur unité (\S{}17). Puisque $\mM$ est simple, un seul de ces deux sommants peut apparaître. Lorsque l'élément unité est opérateur unité, on obtient, comme au \S{}18, un isomorphisme avec un idéal à gauche simple. En effet, si $m\ne0$ appartient à $\mM$, alors $\mo/\mc\cdot m\ne0$; on a donc $\ml\cdot m\ne0$ pour au moins un idéal à gauche simple $\ml$ de $\mo/\mc$, et $\ml m$ doit dès lors épuiser $\mM$. Les autres assertions sont claires pour ces idéaux à gauche et se transportent immédiatement à tous les modules qui leur sont isomorphes.

\textbf{Conséquence.} \emph{Si un $\mo$-module $\mM$ possède une suite de composition, ses facteurs de composition simples sont soit annulés par $\mo$, soit isomorphes à des idéaux à gauche simples de $\mo/\mc$.}

Soit maintenant $\mathsf{P}$ un anneau dont l'action commute avec celle de $\mo$. Tous ces résultats subsistent, à condition de considérer des bimodules à droite sur $\mathsf{P}$ et à gauche sur $\mo$ qui satisfont à
\[
        am\cdot\varrho=a\cdot m\varrho=a\varrho\cdot m\qquad(\S~15).
\]
On exige en outre que l'élément unité de $\mathsf{P}$ soit opérateur unité. Les seuls sous-modules considérés sont ceux qui sont stables par multiplication par $\mathsf{P}$, comme dans le cas des idéaux. Les sous-modules construits dans nos démonstrations, $\mc\mM,\mc^2\mM,\ldots$, satisfont à cette condition.

Si, en particulier, $\mathsf{P}$ est un corps et $\mM$ un $\mathsf{P}$-module de rang fini, alors $\mM$ est aussi un module de représentation. La représentation portée par $\mM$ est non seulement un homomorphisme d'anneaux, mais aussi un homomorphisme compatible avec l'action de $\mathsf{P}$ (\S{}15). Toutes les représentations dans $\mathsf{P}$ qui possèdent cette propriété sont fournies par de tels modules. Les représentations irréductibles correspondent aux modules simples, et réciproquement. Les théorèmes de ce paragraphe se lisent donc immédiatement comme des théorèmes sur les représentations de $\mo$ dans $\mathsf{P}$: \emph{toutes les représentations irréductibles, à l'exception des représentations nulles, sont portées par les idéaux à gauche simples de $\mo/\mc$. Si l'on réduit une représentation arbitraire, conformément au \S{}16, au moyen d'une suite de composition, les seules représentations qui apparaissent sur la diagonale, en dehors des zéros, sont équivalentes à celles que portent les idéaux à gauche simples de $\mo/\mc$.}

\textbf{Remarque.} Sous les hypothèses faites sur $\mo$ et $\mathsf{P}$, il n'est pas nécessaire d'imposer à $\mo$ une condition de finitude. En effet, toute représentation est en même temps une représentation isomorphe de son domaine absolu de multiplicateurs; celui-ci, étant isomorphe à un anneau de matrices de degré fini, est lui-même un $\mathsf{P}$-module de rang fini. Comme, par définition, seuls les $\mathsf{P}$-modules sont admis comme idéaux, les conditions maximale et minimale sont satisfaites pour ce domaine absolu de multiplicateurs.

Ce domaine absolu de multiplicateurs est un système hypercomplexe sur $\mathsf{P}$ lorsque $\mathsf{P}$ est supposé être un corps commutatif. Cette hypothèse est toujours satisfaite dès que la diagonale ne contient pas seulement des zéros. En effet, d'après la remarque du \S{}18, $\mathsf{P}$ est alors isomorphe à un sous-corps du corps gauche des endomorphismes de $\ml$; dans la réalisation par le sous-corps $\mathsf{K}$ de $\mo/\mc$ (\S{}14), ce sous-corps doit, puisque les actions commutent, être contenu dans le centre de $\mathsf{K}$.

\providecommand{\mo}{\mathfrak{o}}
\providecommand{\mc}{\mathfrak{c}}
\providecommand{\ma}{\mathfrak{a}}
\providecommand{\ml}{\mathfrak{l}}

\subsection*{Chapitre IV. Représentations de groupes et de systèmes hypercomplexes}

\subsection*{\S{} 20. Inclusion des systèmes hypercomplexes}
Nous considérons maintenant des systèmes hypercomplexes $\mo$ relativement à un corps commutatif $\mathsf{P}$. Puisque, par convention, seuls les $P$-modules entrent en considération comme idéaux dans $\mo$, les conditions maximale et minimale valent pour les idéaux à gauche et à droite (cf. le \S{}8). Toute la théorie des anneaux exposée au chapitre II est donc applicable.

Dans ce qui suit, par représentation du système hypercomplexe $\mo$, on entendra toujours une représentation dans le corps $\mathsf{P}$, et précisément une représentation telle que de $a\to A$ il résulte
\[
        a\varrho\to A\varrho\qquad(\varrho\in P),
\]
donc une homomorphie de $P$-modules. Pour les modules de représentation, cela signifie, conformément à la conclusion du \S{}15,
\[
        a\varrho\cdot m=a\cdot m\varrho.
\]
Les résultats du \S{}19 sont donc applicables. Toutes les représentations irréductibles sont portées par des idéaux à gauche simples de
\[
        \overline{\mo}=\mo/\mc,
\]
où $\mc$ est le radical. Il y a donc autant de représentations irréductibles inéquivalentes qu'il apparaît de composantes simples bilatères dans la décomposition
\[
        \overline{\mo}
        =\overline{\mo}e^{(1)}+\cdots+\overline{\mo}e^{(s)}.
\]

Pour déterminer explicitement la représentation portée par un tel idéal à gauche, on passe d'abord des éléments de $\mo$ à leurs classes résiduelles dans $\overline{\mo}$. La multiplication par un élément de $P$ est un isomorphisme pour tous les idéaux de $\overline{\mo}$; ainsi $\mathsf{P}$ est isomorphe à un sous-corps $e^{(\nu)}\mathsf{P}$ du corps gauche des endomorphismes $\mathsf{K}_\nu$ de chaque idéal à gauche simple. La représentation portée par $\ml_\nu$ sur $\mathsf{P}$ se trouve en considérant d'abord la représentation sur ce sous-corps $e^{(\nu)}\mathsf{P}$, puis en la transportant vers $P$ par l'isomorphisme $e^{(\nu)}\mathsf{P}\simeq \mathsf{P}$. Le corps gauche $\mathsf{K}_\nu$ des endomorphismes est de rang fini relativement à $\mathsf{P}$; il s'agit donc d'une représentation dans un sous-corps de degré fini du corps gauche des endomorphismes de $\ml_\nu$.

Soit
\[
        c=c_1+\ldots+c_s
\]
la décomposition d'un élément de $\overline{\mo}$ suivant ses composantes bilatères. Pour la représentation portée par $\ml_\nu$, il suffit de chercher la matrice représentante de $c_\nu$, puisque les autres $c_i$ annulent l'idéal à gauche $\ml_\nu$ et sont donc représentés par la matrice nulle. On la trouve, d'après le \S{}18, comme suit: on exprime $c_\nu$ au moyen des unités matricielles $c_{ik}^{(\nu)}$ de $\ma_\nu$:
\[
        c_\nu=\sum c_{ik}^{(\nu)}\alpha_{ik}^{(\nu)},
\]
et, dans la matrice $(\alpha_{ik}^{(\nu)})$, on remplace chaque élément $\alpha_{ik}^{(\nu)}$ par la matrice $A_{ik}^{(\nu)}$ qui lui correspond dans la représentation de $\mathsf{K}_\nu$ dans $\mathsf{P}$ portée par $\mathsf{K}_\nu$.

Comme $\mathsf{K}_\nu$ est de rang fini relativement à $\mathsf{P}$, tout élément $x$ de $\mathsf{K}_\nu$ satisfait à une équation $f(\varkappa)=0$ à coefficients dans $e^{(\nu)}\mathsf{P}$. Si $P$ est algébriquement clos, cette équation se décompose en facteurs linéaires; puisque $\mathsf{K}$ est un corps, $x$ satisfait déjà à l'un de ces facteurs linéaires et se trouve donc dans $e^{(\nu)}\mathsf{P}$. Ainsi $\mathsf{K}_\nu=e^{(\nu)}\mathsf{P}$, et les matrices $(\alpha_{ik}^{(\nu)})$ elles-mêmes sont les représentations irréductibles.

De cette remarque et de la conclusion du \S{}19 résulte ainsi le théorème de Burnside: \emph{si $\mathsf{P}$ est algébriquement clos et si son image est contenue dans le centre de $\mo$, alors toute représentation irréductible de degré $n$ de $\mo$ dans $P$, $\mo\sim\mathfrak D$, contient exactement $n^2$ matrices linéairement indépendantes.}\textsuperscript{17)}

Il en résulte aussitôt la formulation ordinaire: un système de matrices de degré $n$ sur $\mathsf{P}$, fermé pour la multiplication, est irréductible si et seulement s'il n'existe aucune relation linéaire homogène non triviale
\[
        \sum \alpha_{ik}\varrho_{ik}=0
\]
à coefficients $\varrho_{ik}$ dans $\mathsf{P}$, satisfaite par les éléments $\alpha_{ik}$ de chaque matrice du système. En effet, en ajoutant toutes les combinaisons linéaires, on obtient un anneau et un $\mathsf{P}$-module, et l'on peut considérer cet anneau comme sa propre représentation.

\textbf{Démonstration.} D'après le \S{}19, $\mathfrak D$ est isomorphe à un anneau $\overline{\mo}$ simple bilatère, complètement réductible et muni d'un élément unité. Un tel anneau est un anneau complet de matrices --- disons l'anneau de toutes les matrices de degré $m$ --- sur le corps gauche des endomorphismes de ses idéaux à gauche; l'hypothèse que $P$ est algébriquement clos implique que ce corps gauche coïncide avec $P$. Toute représentation irréductible de $\overline{\mo}$, et donc la représentation donnée, est portée par un idéal à gauche simple. Cet idéal est de rang $m$, tandis que la représentation est de degré $n$; ainsi $m=n$. Il existe donc exactement $n^2$ éléments linéairement indépendants dans $\overline{\mo}$, et par conséquent aussi dans $\mathfrak D$.

De même on obtient le théorème de Burnside généralisé\textsuperscript{18)}: sous les mêmes hypothèses, une représentation complètement réductible $\mathfrak D$ qui se décompose en composantes toutes inéquivalentes de degrés
\[
        n_1,n_2,\ldots,n_s
\]
a pour rang
\[
        n_1^2+n_2^2+\cdots+n_s^2
\]
relativement à $P$.

\textbf{Démonstration.} Soit de nouveau $\overline{\mo}$ le domaine absolu de multiplicateurs, qui est nécessairement un système hypercomplexe. Les différentes représentations inéquivalentes sont portées par des idéaux à gauche $\ml_1,\ldots,\ml_s$ appartenant à des composantes simples bilatères distinctes $\ma_i$ de $\overline{\mo}/\overline{\mc}$, où $\overline{\mc}$ est le radical. La représentation donnée est donc portée par $\ml_1+\ldots+\ml_s$ et possède $\ma_1+\ldots+\ma_s$ comme domaine absolu de multiplicateurs. Les mêmes considérations de rang que ci-dessus donnent alors l'assertion.

Si $\mo$ est un système hypercomplexe sans radical sur un corps quelconque, alors, d'après le \S{}13, $\mo$ est complètement réductible et possède un élément unité. D'après les \S{}\S{}17 et 18, toute représentation de $\mo$ est complètement réductible. Réciproquement, toute représentation complètement réductible d'un anneau $\mo$ sur lequel $P$ agit centralement a pour domaine absolu de multiplicateurs un système hypercomplexe sans radical.

\textbf{Démonstration.} Le domaine absolu de multiplicateurs $\overline{\mo}$ est isomorphe à un anneau de matrices à coefficients dans $\mathsf{P}$ qui est en même temps un $\mathsf{P}$-module; c'est donc un système hypercomplexe. D'après le \S{}19, les éléments de son radical $\overline{\mc}$ sont représentés par zéro dans chaque représentation irréductible, donc aussi dans la représentation entière. Par conséquent $\overline{\mc}=0$.

(Si la représentation ne contient que des composantes équivalentes, toutes portées par un même idéal à gauche $\ml$, alors son domaine absolu de multiplicateurs est simple bilatère.)

On appelle représentation régulière d'un système hypercomplexe $\mo$ la représentation portée par l'idéal unité $\mo$ lui-même. (Pour l'anneau de groupe, on choisit spécialement comme base les éléments du groupe $u_1,\ldots,u_h$; cf. le \S{}6.) Comme tous les idéaux à gauche de $\mo/\mc$ sont présents comme facteurs de composition dans $\mo$, toutes les représentations irréductibles se trouvent déjà dans la représentation régulière sous forme de matrices diagonales $R_{ii}$, si l'on imagine la représentation régulière transformée, conformément à la formule (2) du \S{}16, au moyen d'une suite de composition, relativement à une base appropriée.

Une représentation est connue dès que les matrices représentantes des éléments de base $u_1,\ldots,u_h$ sont connues. Si le système est un anneau de groupe, de sorte que les $u_i$ sont des éléments du groupe, toute représentation matricielle homomorphe du groupe est aussi une représentation de l'anneau de groupe. Le problème de la représentation d'un groupe est ainsi un cas particulier de la représentation des systèmes hypercomplexes.

{\footnotesize\noindent 17) On ne suppose donc pas que $\mo$ soit un système hypercomplexe; il peut être, par exemple, l'anneau de groupe constitué de toutes les combinaisons linéaires d'un nombre fini quelconque d'éléments d'un groupe infini.\par
18) G. Frobenius et I. Schur, \emph{Über die Äquivalenz der Gruppen linearer Substitutionen}, Sitzungsber. Berlin 1906.\par}

\providecommand{\mo}{\mathfrak{o}}
\providecommand{\mc}{\mathfrak{c}}
\providecommand{\mC}{\mathfrak{C}}
\providecommand{\ma}{\mathfrak{a}}
\providecommand{\mZ}{\mathfrak{Z}}

\subsection*{\S{} 21. Extension du corps de base. Les représentations du centre}
Soit
\[
        \mo=a_1P+\ldots+a_hP
\]
un système hypercomplexe, et soit $\Omega$ un corps d'extension du corps de base $\mathsf{P}$; on peut alors former
\[
        \mo\Omega=a_1\Omega+\ldots+a_h\Omega
\]
avec les anciennes règles de multiplication pour les éléments de base $a_i$.\textsuperscript{18a)} $\mo\Omega$ est de nouveau hypercomplexe relativement à $\Omega$.

Toute représentation de $\mo$ conduit à une représentation de $\mo\Omega$, puisque la représentation de tous les éléments est connue dès que les représentations des éléments de base $a_i$ sont connues.\textsuperscript{19)}

{\footnotesize\noindent 18a) Plus exactement: on introduit de nouveaux symboles $\bar a_i$ avec les anciennes règles de multiplication; $\bar a_1\Omega+\ldots+\bar a_p\Omega$ contient alors un sous-anneau isomorphe à $\mo$, de sorte que l'on peut identifier après coup les $\bar a_i$ avec les $a_i$.\par
19) Naturellement, on ne considère de nouveau que les représentations qui sont des homomorphismes de $\mathsf{P}$-modules; de $a\to A$ doit suivre $a\varrho\to A\varrho$ pour $\varrho\in P$, et de façon correspondante pour $\Omega$.\par}

Un idéal ou module de représentation, ainsi que la représentation correspondante, sont dits absolument irréductibles lorsqu'ils restent irréductibles après passage au corps algébriquement clos.

Si $\mo\Omega$ est sans radical, alors $\mo$ l'est aussi; car un idéal nilpotent $\mc$ dans $\mo$ conduirait à un idéal étendu $\mc\Omega$ dans $\mo\Omega$. La réciproque n'est pas vraie, comme nous le verrons plus bas.

Le centre de $\mo\Omega$ devient égal au centre étendu $\mZ\Omega$. Si $\mo$ est complètement réductible, alors $\mZ$ est somme directe de corps:
\[
        \mZ=\mZ_1+\cdots+\mZ_s.
\]
Si $\mo$, lors du passage à $\mo\Omega$, reste complètement réductible (c'est-à-dire si aucun idéal nilpotent ne s'ajoute), alors le nouveau centre
\[
        \mZ\Omega=\mZ_1\Omega+\cdots+\mZ_s\Omega
\]
se décompose de nouveau, dans $\Omega$, en une somme directe de corps qui, si $\Omega$ est algébriquement clos, ont, d'après une remarque du \S{}20, le rang $1$ relativement à $\Omega$ (et sont $\simeq\Omega$).

Pour les représentations du centre $\mZ$, les théorèmes suivants valent. Toute représentation irréductible d'un système hypercomplexe commutatif $\mZ$ dans le corps algébriquement clos $\Omega$ est du premier degré; ou encore: les représentations irréductibles sont identiques aux homomorphismes de $\mZ$ dans $\Omega$.

\emph{Démonstration.} Toute représentation irréductible de $\mZ$ donne une telle représentation de $\mZ_\Omega$, donc aussi une représentation de $\mZ_\Omega/\mC$, où $\mC$ est le radical de $\mZ_\Omega$. Ce quotient est un système commutatif sans radical, donc, d'après la conclusion du \S{}14, une somme directe de corps. Ceux-ci sont du premier degré, puisque $\Omega$ est supposé algébriquement clos, et ils engendrent toutes les représentations irréductibles (\S{}20). Il en résulte en même temps, puisque des représentations équivalentes du premier degré sont nécessairement égales ($\lambda^{-1}\alpha\lambda=\alpha$), que le nombre des différents homomorphismes de $\mZ$ dans $\Omega$ est égal au rang de $\mZ_\Omega/\mC$.

La dernière remarque donne, dans le cas particulier où $\mZ$ est un corps sur $\mathsf{P}$, le théorème suivant: \emph{si $\mZ$ est un corps commutatif, alors $\mZ$ est complètement réductible (sans radical) si et seulement si $\mZ$ est une extension de première espèce de $\mathsf{P}$.} En effet, pour un corps, les homomorphismes sont des isomorphismes; leur nombre est égal au rang de $\mZ_\Omega/\mC$, et il est donc égal au degré du corps ($\mZ$ étant une extension de première espèce) si et seulement si $\mC=0$. Si
\[
        \mZ=\mZ_1+\cdots+\mZ_t
\]
est complètement réductible, alors, dans chaque représentation, chaque corps $\mZ_\nu$ est aussi envoyé homomorphiquement, et, dans une représentation irréductible, l'un de ces corps est représenté isomorphiquement, les autres par zéro (\S{}19).

L'application de ces théorèmes aux systèmes hypercomplexes donne d'abord, pour les systèmes sans radical: \emph{si $\mo_\Omega$, et donc aussi $\mo$, est un système sans radical, alors, dans toute représentation absolument irréductible de $\mo$, les éléments $z$ du centre sont représentés par des matrices diagonales $E\zeta^{(\nu)}$, et $z\to\zeta^{(\nu)}$ est une représentation du premier degré de $\mZ$ dans $\Omega$. La relation ainsi établie entre les classes de représentations absolument irréductibles de $\mo$ et celles de $\mZ$ est biunivoque. Le nombre de ces classes de représentations est donc égal au rang du centre.}\textsuperscript{20)}

{\footnotesize\noindent 20) Des théorèmes correspondants valent non seulement pour les représentations dans le corps algébriquement clos $\Omega$, mais aussi pour les représentations dans les corps gauches des endomorphismes individuels, comme on l'exposera ailleurs.\par}

\textbf{Démonstration.} À la décomposition en idéaux bilatéralement indécomposables
\[
        \mo_\Omega=\ma_1+\ldots+\ma_t
\]
correspond biunivoquement une décomposition du centre en corps
\[
        \mZ_\Omega=\mZ_1+\ldots+\mZ_t,
\]
où $\mZ_\Omega$ est le centre de $\mo_\Omega$. Dans une représentation irréductible de $\mo$, l'un des $\ma_\nu$ est représenté isomorphiquement et les autres par zéro, et la classe de représentation est déterminée de manière unique par $\ma_\nu$. L'idéal $\ma_\nu$ est représenté par un anneau matriciel complet, et le centre d'un tel anneau matriciel complet se compose des matrices diagonales $E\zeta^{(\nu)}$. L'application $z\to\zeta^{(\nu}$ est un homomorphisme de $\mZ$: l'un des $\mZ_\nu$ est représenté isomorphiquement et les autres par zéro. Tout est ainsi démontré.

Pour déterminer quand un $\mo$ sans radical donne lieu à un $\mo_\Omega$ également sans radical, on dispose encore du fait suivant: si $\mZ$ possède une composante $\mZ_\nu$ qui est une extension de seconde espèce de $\mathsf{P}$, alors $\mo_\Omega$ possède sûrement un radical.\textsuperscript{21)} Car $\mZ_\Omega$ en possède déjà un dans ce cas.

{\footnotesize\noindent 21) Si $\mZ$ n'a que des composantes de première espèce, alors $\mo\Omega$ devient sans radical, comme on le démontrera également plus loin. En caractéristique nulle, I. Schur a déjà démontré le théorème équivalent selon lequel les représentations irréductibles dans $\mathsf{P}$ restent complètement réductibles lors de toute extension du domaine des coefficients (\emph{Beiträge zur Theorie der Gruppen linearer Substitutionen}, \emph{Transact. Am. Math. Soc.} 15 (1909), p.~159).\par}

\providecommand{\mG}{\mathfrak{G}}
\providecommand{\mP}{\mathfrak{P}}
\providecommand{\mZ}{\mathfrak{Z}}

\subsection*{\S{} 22. Application aux groupes abéliens}
Soit
\[
        \mG=\mG_1\times\mG_2\times\cdots\times\mG_r
\]
un groupe abélien fini, décomposé en groupes cycliques $\mathfrak{C}_i$ d'ordres $h_i$. Ses éléments sont de la forme
\[
        a=c_1^{\lambda_1}\ldots c_r^{\lambda_r}.
\]
Soit $\mathsf{P}$ un corps dont la caractéristique ne divise pas l'ordre du groupe $h=h_1h_2\ldots h_r$.

L'anneau de groupe $\mZ$ se compose de toutes les sommes
\[
        \sum a_\lambda\varrho_\lambda
        =
        \sum c_1^{\lambda_1}\ldots c_r^{\lambda_r}\varrho_\lambda
        \qquad(\varrho_\lambda\in P)
\]
et il est image homomorphe de l'anneau de polynômes $\mathsf{P}[z_1,\ldots,z_r]$ au moyen de l'application
\[
        \sum \varrho_\lambda z_1^{\lambda_1}\cdots z_r^{\lambda_r}
        \to
        c_1^{\lambda_1}\cdots c_r^{\lambda_r}\varrho_\lambda.
\]
Ainsi $\mZ\simeq \mathsf{P}[z_1,\ldots,z_r]/\mathfrak{m}$, et l'on trouve facilement:
\[
        \mathfrak{m}=(z_1^{h_1}-e,\ldots,z_r^{h_r}-e).
\]
Dans le corps d'extension $\Omega$, les $z_i^{h_i}-e$ se décomposent en facteurs tous distincts $z_i-\varepsilon_i$; l'idéal $\mathfrak{m}$ devient donc le produit, ou l'intersection, d'idéaux premiers tous distincts
\[
        \mP^{(\nu)}=(z_1-\varepsilon_1^{(\nu)},\ldots,z_r-\varepsilon_r^{(\nu)}),
\]
chacun n'ayant qu'un seul zéro $(\varepsilon_1^{(\nu)},\ldots,\varepsilon_r^{(\nu)})$. À la formation de l'intersection correspond une décomposition en somme directe (\S{}4):
\[
        \mZ_\Omega=\mZ_1+\cdots+\mZ_h,
\]
où, chaque fois,
\[
        \mZ_\nu\simeq\mZ_\Omega/\mP^{(\nu)}.
\]
Cette décomposition fournit donc la représentation
\[
        c_1^{\lambda}\ldots c_r^{\lambda_r}
        \to
        \varepsilon_1^{(\nu)\lambda_1}\ldots \varepsilon_r^{(\nu)\lambda_r},
        \qquad\text{ou}\qquad
        a\to\chi^{(r)}(a).
\]
Ces représentations sont les \emph{caractères}. Leur nombre est $h_1h_2\ldots h_r=h$, comme l'exige la théorie générale.

\providecommand{\mo}{\mathfrak{o}}
\providecommand{\mc}{\mathfrak{c}}
\providecommand{\ml}{\mathfrak{l}}

\subsection*{\S{} 23. Déterminant d'un système hypercomplexe}
Soit
\[
        \mo=a_1P+\ldots+a_hP
\]
un système hypercomplexe. J'adjoins à $\mathsf{P}$ les indéterminées $x_1,\ldots,x_h$ et je forme
\[
        \mo^*=a_1\mathsf{P}(x)+\ldots+a_h\mathsf{P}(x).
\]
Les $x$ doivent être permutables avec les $a_i$; par là les règles de calcul dans $\mo^*$ sont déterminées.

Dans $\mo^*$ se trouve l'élément général de $\mo$,
\[
        w=a_1x_1+\ldots+a_hx_h.
\]
Si, dans une représentation, $a_i\to A_i$, alors à $w$ est associée
\[
        W=A_1x_1+\ldots+A_hx_h.
\]
$W$ s'appelle la matrice du système appartenant à la représentation; spécialement, lorsque les éléments de base $a_i$ forment un groupe et que $\mo$ est donc l'anneau de groupe, elle s'appelle la matrice du groupe. S'il s'agit de la représentation régulière, on a la matrice régulière du système.

Les éléments $w_{ik}$ de $W$ sont des formes linéaires des $x$. Le déterminant du système $|W|$ est donc de degré $n$ lorsqu'il s'agit de représentations de degré $n$. En particulier, le déterminant régulier du système est de degré $h$.

Le déterminant du système ne change pas lors du passage à des représentations équivalentes, puisque $|\mathsf{P}W\mathsf{P}^{-1}|=|\mathsf{P}|\,|W|\,|\mathsf{P}^{-1}|=|W|$. Lors du passage de $(a_1,\ldots,a_h)$ à une nouvelle base $(b_1,\ldots,b_h)$ et de $w=\sum a_ix_i$ à $w=\sum b_iy_i$, on obtient les nouveaux éléments de $W$, et donc le nouveau déterminant, en effectuant dans les anciens la substitution régulière (à matrice inversible)
\[
        x_i=\sum \varrho_{ik}y_k.
\]

Si l'on dispose d'une suite de composition du module de représentation, alors, pour un choix convenable de base, la matrice $W$ a la forme
\[
        W=\begin{pmatrix}
        W_1&0&\cdots&0\\
        *&W_2&\cdots&0\\
        \vdots&\vdots&\ddots&\vdots\\
        *&*&\cdots&W_r
        \end{pmatrix}.
\]
\[
        |W|=|W_1|\cdot|W_2|\cdot\ldots\cdot|W_r|.
\]
Parmi les $|W_i|$ d'une représentation arbitraire, il n'en apparaît pas d'autres que dans la représentation régulière de $\mo$, ou même de $\mo/\mc$, où $\mc$ est le radical.

Si $\mathsf{P}$ est algébriquement clos, le déterminant $|W_i|$ appartenant à une représentation irréductible est un polynôme irréductible en les $x$, et à des représentations inéquivalentes correspondent des facteurs irréductibles distincts.

\emph{Démonstration.} Puisque toutes les représentations irréductibles de $\mo$ sont aussi des représentations de $\mo/\mc$, nous pouvons nous restreindre à l'anneau sans radical $\mo/\mc$. Dans celui-ci, nous introduisons comme base les unités matricielles $c_{ik}$; l'élément général s'écrit
\[
        w=\sum_{\nu,i,k}c_{ik}^{(\nu)}x_{ik}^{(\nu)}.
\]
Les matrices des représentations irréductibles s'écrivent
\[
        W_\nu=(x_{ik}^{(\nu)}).
\]
Les fonctions $|W_\nu|=|x_{ik}^{(\nu)}|$ sont irréductibles et distinctes entre elles.

Pour calculer les $|W_\nu|$, on peut décomposer en facteurs premiers la matrice régulière du système de $\mo/\mc$. Chaque facteur premier apparaît autant de fois que l'indique le degré de la représentation irréductible, puisque l'idéal correspondant $\ml_\nu$ apparaît autant de fois dans la suite de composition.

On peut aussi prendre pour base la matrice régulière du système de $\mo$, mais on obtient alors chaque facteur irréductible plus souvent, à savoir autant de fois que le $\ml_\nu$ correspondant apparaît comme facteur de composition parmi les idéaux à gauche. Dans cette représentation régulière, $\mo$ était considéré comme idéal à gauche; si l'on considère $\mo$ comme idéal à droite, il apparaît une seconde matrice régulière du système (la « matrice antistrophe » chez Frobenius), qui contient les mêmes facteurs irréductibles (à savoir les déterminants de système de toutes les représentations irréductibles), mais éventuellement avec d'autres exposants (voir l'exemple du \S{}10).

Le déterminant du système d'un système commutatif se décompose en facteurs linéaires, puisque toutes les représentations irréductibles sont du premier degré. Ces facteurs linéaires sont eux-mêmes les représentations irréductibles; ils donnent donc, pour les groupes abéliens, les caractères. Ce fait fut le point de départ de Dedekind dans l'étude du déterminant de groupe des groupes non abéliens.

\providecommand{\mo}{\mathfrak{o}}
\providecommand{\mc}{\mathfrak{c}}
\providecommand{\mM}{\mathfrak{M}}
\providecommand{\mR}{\mathfrak{R}}

\subsection*{\S{} 24. Traces et caractères}
Si, dans une représentation $\mo\sim\mathfrak D$ d'un système hypercomplexe $\mo$, la matrice $A$ est associée à l'élément $a$, on pose
\[
        \mathit{Sp}_{\mathfrak D}(a)=\mathit{Sp} A.
\]
Les traces sont des fonctions linéaires:
\[
        \mathit{Sp}(c+d)=\mathit{Sp}(c)+\mathit{Sp}(d),
        \qquad
        \mathit{Sp}(c\alpha)=\alpha\mathit{Sp}(c).
\]
Les représentations équivalentes ont les mêmes traces. La trace dans une représentation réductible est la somme des traces dans les représentations induites par les facteurs de composition.

\emph{Démonstration.} En effet,
\[
        \mathit{Sp}\begin{pmatrix}
        A_{11}&0&\cdots&0\\
        *&A_{22}&\cdots&0\\
        \vdots&\vdots&\ddots&\vdots\\
        *&*&\cdots&A_{ss}
        \end{pmatrix}
        =\operatorname{Sp}(A_{11})+\operatorname{Sp}(A_{22})+\cdots+\operatorname{Sp}(A_{ss}).
\]
On appelle trace principale la trace dans la représentation régulière. On appelle trace réduite la somme des traces dans les différentes représentations irréductibles.

Si $c$ est un élément de l'idéal nilpotent maximal $\mc$, alors $\mathit{Sp}(c)=0$ pour toute représentation.

\emph{Démonstration.} Il suffit d'étudier les représentations définies par les idéaux à gauche simples de $\mo/\mc$; en effet, dans toute autre représentation, seuls les facteurs de composition ainsi obtenus apparaissent, et la trace se compose additivement. Or $\mathfrak{c}$ annule tous les éléments de $\mo/\mc$; à tout $c\in\mc$ est donc associée la matrice nulle, et par conséquent $\mathit{Sp}(c)=0$.

Si $\mo$ est simple bilatère et si $\mathsf{P}$ est algébriquement clos, donc si $\mo=\sum c_{ik}\mathsf{P}$ est un anneau de matrices, alors, dans la représentation irréductible de $\mo$, à l'élément $a=\sum c_{ik}\alpha_{ik}$ est associée la matrice $(\alpha_{ik})$. Donc
\[
        \operatorname{Sp}_{\rm red}(a)=\sum_i\alpha_{ii},
        \qquad
        \operatorname{Sp}_{\rm pr}(a)=n\,\operatorname{Sp}_{\rm red}(a)=n\sum_i\alpha_{ii}.
\]
En particulier
\[
        \operatorname{Sp}_{\rm red}(c_{ik})=0\ (i\ne k),\quad
        \operatorname{Sp}_{\rm red}(c_{ii})=e,
\]
et
\[
        \operatorname{Sp}_{\rm pr}(c_{ik})=0\ (i\ne k),\quad
        \operatorname{Sp}_{\rm pr}(c_{ii})=n e.
\]

Lorsqu'on passe de $\mathsf{P}$ au corps algébriquement clos $\Omega$, la trace principale ne change pas, puisqu'elle peut être calculée à partir de la même base. Il n'en va pas de même de la trace réduite.

Les traces des éléments $a$ du système sans radical, dans les représentations absolument irréductibles, s'appellent les caractères; on les désigne par $\chi(a)$ ou, lorsqu'il faut indiquer de quelle représentation il s'agit, par $\chi^{(\nu)}(a)$.

Dans une représentation irréductible de degré $n_\nu$, les éléments du centre sont représentés, d'après le \S{}21, par des matrices diagonales $E\zeta^{(\nu)}$, où $z\to\zeta^{(\nu)}$ est une représentation du premier degré du centre (ou un homomorphisme du centre dans $\Omega$). Il en résulte pour la trace la valeur $n_\nu\theta^{(\nu)}(z)$. Donc les homomorphismes du centre sont liés aux caractères par
\[
        \chi^{(\nu)}(z)=n_\nu\theta^{(\nu)}(z)                 \tag{1}
\].
Dans le cas commutatif, $n_\nu=1$, et les caractères eux-mêmes donnent les homomorphismes (voir \S{}22). Si le corps $\Omega$ est de caractéristique zéro, ce que nous supposerons toujours dans la suite, on peut diviser (1) par $n$:
\[
        \theta^{(\nu)}(z)=\frac{\chi^{(\nu)}(z)}{n_\nu}.
\]
La propriété d'homomorphisme s'exprime par
\[
        \frac{\chi^{(\nu)}(z+z')}{n_\nu}
        =\frac{\chi^{(\nu)}(z)}{n_\nu}+\frac{\chi^{(\nu)}(z')}{n_\nu},
\]
\[
        \frac{\chi^{(\nu)}(zz')}{n_\nu}
        =\frac{\chi^{(\nu)}(z)}{n_\nu}\frac{\chi^{(\nu)}(z')}{n_\nu}
\].

Une classe de représentation est déjà déterminée de façon univoque par les seules traces des matrices. (Pour connaître les traces de toutes les matrices, il suffit naturellement de connaître les traces des éléments de base dans la représentation donnée.)

\emph{Démonstration.} Le module de représentation $\mR$ est connu lorsqu'on sait combien de fois chaque module de représentation irréductible $\mM_\nu$ apparaît comme facteur direct. Si ce nombre est $p_\nu$, la trace de $e^{(\nu)}$ est égale à $\mu_\nu n_\nu$, donc
\[
        \mu_\nu=\frac{\operatorname{Sp}(e^{(\nu)})}{n_\nu}.
\]

\providecommand{\mo}{\mathfrak{o}}
\providecommand{\mc}{\mathfrak{c}}
\providecommand{\ma}{\mathfrak{a}}

\subsection*{\S{} 25. Discriminants}
Soit
\[
        \mo=a_1P+\cdots+a_hP
\]
un système hypercomplexe. On appelle matrice discriminante la matrice dont les éléments sont les traces principales $\operatorname{Sp}(a_i a_j)$. On appelle matrice discriminante réduite la même matrice pour les traces réduites, formée dans le corps algébriquement clos. Le déterminant de la matrice s'appelle respectivement discriminant et discriminant réduit.

Le discriminant reste le même lors de l'extension du corps de base. Lors du passage à une autre base, il se multiplie par le carré du déterminant de transformation.

\emph{Démonstration.} Soit
\begin{align*}
(a_1,\ldots,a_h)&=(b_1,\ldots,b_h)P,\\
(a_i a_1,\ldots,a_i a_h)&=(a_i b_1,\ldots,a_i b_h)P,\\
\bigl(\mathit{Sp}(a_i a_1),\ldots,\mathit{Sp}(a_i a_h)\bigr)
&=\bigl(\mathit{Sp}(a_i b_1),\ldots,\mathit{Sp}(a_i b_h)\bigr)P,
\end{align*}
Ou, sous forme matricielle,
\[
        (\operatorname{Sp}(a_i a_j))=(\operatorname{Sp}(a_i b_j))P,             \tag{1}
\]
De même, si $P^t$ est la matrice transposée,
\[
\begin{pmatrix}a_1\\[-1mm]\vdots\\[-1mm]a_h\end{pmatrix}
=P^t
\begin{pmatrix}b_1\\[-1mm]\vdots\\[-1mm]b_h\end{pmatrix},
\]
et, comme précédemment,
\[
        (\operatorname{Sp}(a_i b_j))=P^t(\operatorname{Sp}(b_i b_j)).           \tag{1a}
\]
De (1) et (1a), on déduit
\[
        (\operatorname{Sp}(a_i a_j))=P^t(\operatorname{Sp}(b_i b_j))P.
\]
Le passage aux déterminants fournit l'assertion. Le déterminant n'est donc déterminé qu'à un carré près dans $P$; son annulation est une propriété invariante. On déduit encore de (1) que l'on peut aussi déterminer le discriminant, à un facteur $\ne0$ près, à partir de deux bases différentes, comme $|\operatorname{Sp}(a_i b_j)|$.

Le discriminant s'annule si $\mo$ possède un idéal nilpotent $\mc$.

\emph{Démonstration.} Comme éléments de base, choisissons
\[
(c_1,\ldots,c_t,d_{t+1},\ldots,d_h),
\]
où $(c_1,\ldots,c_t)$ est une base de $\mc$. Le discriminant est
\[
\left|\begin{array}{cc}
\mathit{Sp}(c_i c_k)&\mathit{Sp}(c_i d_k)\\
\mathit{Sp}(d_i c_k)&\mathit{Sp}(d_i d_k)
\end{array}\right|
=
\left|\begin{array}{cc}
0&0\\
0&\mathit{Sp}(d_i d_k)
\end{array}\right|=0.
\]
La partie supérieure devient donc nulle; cela vaut aussi pour le discriminant réduit.

\emph{Conséquence.} Le discriminant s'annule également lorsqu'un idéal nilpotent apparaît après passage au corps algébriquement clos.

Soit
\[
        \mo=\ma_1+\ldots+\ma_s
\]
une somme directe bilatère, et soient $M,M_1,\ldots,M_s$ les matrices discriminantes des anneaux $\mo,\ma_1,\ldots,\ma_s$. Alors
\[
        M=\begin{pmatrix}
        M_1&0&\cdots&0\\
        0&M_2&\cdots&0\\
        \vdots&\vdots&\ddots&\vdots\\
        0&0&\cdots&M_s
        \end{pmatrix},
\]
donc $|M|=|M_1|\cdots|M_s|$.

\emph{Démonstration.} Choisissons pour $\mo$ une base composée des bases des $\ma_i$. Si $a\in\ma_i$, alors
\[
        \mathit{Sp}_{\mo}(\ma)=\mathit{Sp}_{\ma_i}(a).
\]
Les deux membres désignent ici des traces principales, prises respectivement dans les anneaux $\mo$ et $\ma_i$. En outre, $a_i a_k=0$ si $a_i\in\ma_i$, $a_k\in\ma_k$ et $i\ne k$; donc aussi $\mathit{Sp}(a_i a_k)=0$. D'où l'assertion, qui vaut de même pour le discriminant réduit.

Pour un anneau de matrices $\sum c_{ik}\mathsf{P}$, choisissons deux bases différentes, à savoir les $c_{ik}$ ordonnés une fois d'après les premiers indices, une fois d'après les seconds indices. La table de multiplication a alors cette forme:
\[
\begin{array}{c|c|c|c}
& \overbrace{c_{11}\;\ldots\;c_{1n}}^{\mathfrak r_1}
& \overbrace{c_{21}\;\ldots\;c_{2n}}^{\mathfrak r_2}
& \ldots \\ \hline
\mathfrak l_1\left\{\begin{array}{c}c_{11}\\[-1mm]\vdots\\[-1mm]c_{n1}\end{array}\right.
& (c_{ik})&0&0\\[1.3em]\hline
\mathfrak l_2\left\{\begin{array}{c}c_{12}\\[-1mm]\vdots\\[-1mm]c_{n2}\end{array}\right.
&0&(c_{ik})&0\\[1.3em]\hline
\vdots&0&0&(c_{ik})
\end{array}
\]
Le discriminant devient
\[
        D=\left|\begin{array}{cccc}
        \mathit{Sp}(c_{11})& & &0\\
        &\mathit{Sp}(c_{22})& & \\
        & &\ddots& \\
        0& & &\mathit{Sp}(c_{nn})
        \end{array}\right|^{n}
        =\begin{cases}
        e, & \hbox{pour les traces réduites},\\
        n^{n^2}\cdot e&\hbox{pour les traces principales}.
        \end{cases}
\]
\emph{Conséquence.} Si $\mo$ se décompose, dans le corps d'extension algébriquement clos, en anneaux de matrices de degrés $n_1,\ldots,n_s$, alors
\[
        D=n_1^{n_1^2}n_2^{n_2^2}\ldots n_s^{n_s^2}\cdot e,
        \qquad
        D_{\rm red}=e.
\]
Le discriminant réduit est donc toujours $\ne0$ pour les systèmes sans radical; l'autre ne l'est que si aucun $n_i$ n'est divisible par la caractéristique du corps. Par conséquent: le non-évanouissement du discriminant réduit est nécessaire et suffisant pour les systèmes sans radical après clôture algébrique du corps $\mathsf{P}$. Dans le cas de la caractéristique zéro, le non-évanouissement du discriminant est nécessaire et suffisant pour la non-apparition d'un radical après clôture algébrique.\textsuperscript{22)}

{\footnotesize\noindent 22) Cf., pour les systèmes commutatifs, E. Noether, \emph{Diskriminantensatz für Ordnungen ...}, J. f. M. 157 (1927), p. 82--104, \S{}\S{}4--6. La méthode de preuve y est la même, mais plus compliquée dans le détail; le passage d'une base à une autre (\S{}4, 4.) doit être remplacé par celui qui est donné ici (au début du paragraphe), puisque le déterminant $\left|(\varrho_i\varrho_k)\right|$ s'annule.\par}

\subsection*{\S{} 26. Inclusion de l'anneau de groupe}
Si $a_1,\ldots,a_h$ sont les éléments d'un groupe fini et que l'on forme l'anneau de groupe sur un corps dont la caractéristique ne divise pas $h$, alors le discriminant $D\ne0$; donc l'anneau de groupe est un anneau sans radical.

\emph{Démonstration.} Montrons d'abord, en convenant que $\mathit{Sp}$ désigne désormais la trace principale:
\[
        \operatorname{Sp}(e)=he,
        \qquad
        \operatorname{Sp}(a_i)=0\quad(a_i\ne e).
\]
Dans la représentation régulière, en effet,
\[
        e\mapsto\begin{pmatrix}e&0&\cdots&0\\0&e&\cdots&0\\\vdots&\vdots&\ddots&\vdots\\0&0&\cdots&e\end{pmatrix},
        \qquad \operatorname{Sp}(e)=he.
\]
Pour $a_i\ne e$, on a
\[
a_i a_k=a_\lambda,\qquad \lambda\ne k.
\]
La matrice qui exprime les produits $a_i a_1,\ldots,a_i a_h$ au moyen de $a_1,\ldots,a_h$ n'a donc que des zéros sur sa diagonale principale; ainsi $\mathit{Sp}(a_i)=0$.

Nous prenons maintenant les bases $a_1,\ldots,a_h$ et
$a_1^{-1},\ldots,a_h^{-1}$. La matrice
$\bigl(\mathit{Sp}(a_i a_k^{-1})\bigr)$ est
\[
\begin{pmatrix}
he&&&0\\
&he&&\\
&&\ddots&\\
0&&&he
\end{pmatrix}.
\]
Par conséquent,
\[
D=h^h e\ne0,\hfill\text{q. e. d.}
\]

On appelle classe d'un élément de groupe la somme formée dans l'anneau de groupe
\[
        K_a=\sum s^{-1}as,\tag{1}
\]
où l'on ne somme que sur les différents $s^{-1}as$. Les classes $K_a$ sont des éléments du centre, puisqu'elles commutent avec tous les éléments du groupe. Elles engendrent le centre; car si un élément
$\sum a_i\varrho_i$ de l'anneau de groupe commute avec tous les $s$, alors
\[
\begin{aligned}
\sum_i a_i\varrho_i
&=s^{-1}\biggl(\sum_i a_i\varrho_i\biggr)s
 =\sum_i(s^{-1}a_i s)\varrho_i\\
&=\frac1h\sum_i\biggl(\sum_s s^{-1}a_i s\biggr)\varrho_i
 =\frac1h\sum_i\biggl(\frac h{h_i}K_i\biggr)\varrho_i,
\end{aligned}
\]
où $h$ est le nombre des éléments du groupe et $h_i$ le nombre des éléments de la classe $K_i$.

Par conséquent, le rang du centre est égal au nombre des classes d'éléments de groupe conjugués, et donc aussi le nombre des représentations absolument irréductibles est égal à ce nombre de classes. Les matrices $A$ et $S^{-1}AS$ ont la même trace; donc les éléments de groupe $a$ et $s^{-1}as$ ont la même trace. Si l'on prend la trace des deux côtés de (1), il s'ensuit
\[
        \chi(K_i)=h_i\chi(a_i),
\]
où $h_i$ est le nombre des éléments de la classe $K_i$. En mots: le caractère d'un élément de groupe est égal au caractère de la classe, divisé par le nombre des éléments de la classe.

\begin{center}
(Reçu le 12 août 1928.)
\end{center}

\clearpage
\section*{\NoetherNavAnchor{P35}{35. Sur les domaines maximaux de fonctions à coefficients entiers}35. Sur les domaines maximaux de fonctions à coefficients entiers}

\begin{center}
{\bfseries Emmy Noether (Göttingen).}\par
\emph{Rec. Soc. Math. Moscou \textbf{36} (1929), p. 65--72}
\end{center}

Les recherches qui suivent touchent de près au problème de Hilbert des fonctions relativement entières\footnote{D. Hilbert, \emph{Mathematische Probleme} (Gött. Nachr. 1900, p. 253--297), problème 14.}; elles remontent à une question occasionnelle d'E. Hecke, qui avait pu régler directement par le calcul un cas particulier dont il avait besoin comme lemme auxiliaire\footnote{E. Hecke, \og{} Über die Konstruktion relativ-Abelscher Zahlkörper durch Modulfunktionen von zwei Variabeln \fg{} (Math. Ann. 74 (1913), p. 465--510), \S\,2. Chez Hecke, le domaine d'intégrité fini $\mathfrak{I}$ est engendré par trois quotients de séries entières en deux variables, entre lesquels il existe une relation algébrique.}.

La question est la suivante. Par \emph{domaine maximal} de fonctions à coefficients entiers en $x_1,\ldots,x_n$ -- polynômes, séries entières ou quotients de telles fonctions -- à l'intérieur d'un domaine donné $\mathfrak{Z}$, j'entends un domaine que l'on ne peut plus élargir par adjonction de dénominateurs entiers : l'appartenance de $a\cdot g(x)$ -- où $g(x)$ appartient à $\mathfrak{Z}$ et où $a$ désigne un entier rationnel non nul -- entraîne toujours celle de $g(x)$. Chaque domaine d'intégrité $\mathfrak{I}$ dans $\mathfrak{Z}$ engendre de façon univoque un domaine maximal $\mathfrak{M}(\mathfrak{I})$, composé de toutes les fonctions $g(x)$ de $\mathfrak{Z}$ pour lesquelles il existe au moins un $a \neq 0$ tel que $a\cdot g(x)$ appartienne à $\mathfrak{I}$.

Il s'agit de savoir ce que l'on peut dire des dénominateurs $a$ de ce domaine maximal $\mathfrak{M}(\mathfrak{I})$ lorsque le domaine d'intégrité initial $\mathfrak{I}$ était fini, c'est-à-dire lorsqu'il se composait de tous les polynômes à coefficients entiers en un nombre fini de fonctions $f_1(x),\ldots,f_r(x)$ de $\mathfrak{Z}$. Pour $\mathfrak{Z}$ on prendra alors tous les polynômes ou séries entières à coefficients entiers en $x_1,\ldots,x_n$, respectivement les quotients qui ne contiennent pas d'autres dénominateurs que ceux qui interviennent dans les $f(x)$ et leurs puissances.

Je montre -- sous l'hypothèse additionnelle, dans le cas des séries entières ou de leurs quotients, qu'il n'y ait entre les $f(x)$ que des dépendances algébriques -- que les $a$ peuvent toujours être choisis de telle sorte que seul un nombre fini de nombres premiers différents les divise, ceux-ci pouvant toutefois apparaître en général à toute puissance.

Ce dernier point est immédiatement clair : si $a\cdot g(x)$ appartient à $\mathfrak{I}$, mais si $a^{\varrho-1}g(x)^{\varrho}$ n'y appartient pour aucun exposant $\varrho$, alors toutes les puissances $a^{\varrho}$ apparaissent; par exemple, pour $\mathfrak{I}=[2x]$, on a $x^{\varrho}=(2x)^{\varrho}/2^{\varrho}$.

La question de la bornitude ne peut se formuler que de la manière suivante : peut-on donner un système de générateurs, en général infini, de $\mathfrak{M}(\mathfrak{I})$ tel que, dans les dénominateurs correspondants, les exposants des nombres premiers demeurent bornés? Puisqu'il n'y a qu'un nombre fini de nombres premiers en jeu, cette question est, par des raisonnements connus\footnote{Cf. Hilbert, loc. cit., ou l'exposé plus détaillé dans E. Noether, \emph{Die Endlichkeit des Systems der ganzzahligen Invarianten binärer Formen} (Gött. Nachr. 1919, p. 138--156), \S\,1, IV.}, identique à celle de la finitude de $\mathfrak{M}(\mathfrak{I})$, donc au problème hilbertien des fonctions relativement entières dans sa formulation entière; c'est une question que je ne puis décider par les méthodes employées ici\footnote{Ainsi, par exemple, dans le cas des invariants entiers (projectifs) d'un système fondamental de formes, on montre qu'ils peuvent être représentés au moyen d'un système complet d'invariants au sens usuel, avec seulement un nombre fini de nombres premiers différents au dénominateur -- résultat qui ne découle pas de la méthode habituelle du procédé $\varOmega$ --; mais on n'affirme rien sur la finitude entière elle-même, connue jusqu'ici seulement dans le cas des formes fondamentales binaires. Cf. la note citée sous 3).}.

La preuve du fait que seulement un nombre fini de nombres premiers distincts $p$ interviennent dans les dénominateurs $a$ repose sur ceci : à chacun de ces $p$ correspond au moins une relation modulo $p$ entre les fonctions $f(x)$ qui engendrent $\mathfrak{I}$, relation à coefficients entiers qui n'est pas identiquement satisfaite en $x$. Autrement dit, si $\mathfrak{o}$ désigne l'anneau de polynômes à coefficients entiers en les indéterminées $u_1,\ldots,u_r$, alors $\mathfrak{I}$ est, par la correspondance $u_i \to f_i(x)$, une image homomorphe de $\mathfrak{o}$; cet homomorphisme a pour noyau un idéal premier $\mathfrak{a}$ de $\mathfrak{o}$; ainsi $\mathfrak{I}$ est isomorphe à l'anneau quotient $\mathfrak{o}/\mathfrak{a}$. De même, $\mathfrak{I}_p$ est image homomorphe de $\mathfrak{o}_p$ -- l'indice désigne le passage aux classes résiduelles modulo $p$, possible sauf pour les nombres premiers, en nombre fini, qui apparaissent dans les dénominateurs de $\mathfrak{I}$. Cet homomorphisme a de nouveau pour noyau un idéal premier $\mathfrak{c}_p$ de $\mathfrak{o}_p$; cet idéal contient $\mathfrak{a}_p$. Alors $p$ intervient dans le dénominateur de $\mathfrak{M}(\mathfrak{I})$ si et seulement si l'inclusion $\mathfrak{a}_p\subsetneq\mathfrak{c}_p$ est stricte (autrement dit, $\mathfrak{c}_p$ est un diviseur propre de $\mathfrak{a}_p$ au sens des idéaux). Cela ne peut se produire que si, ou bien le degré de transcendance de $\mathfrak{I}_p$ diminue par rapport à celui de $\mathfrak{I}$, ou bien $\mathfrak{a}_p$ perd sa propriété d'être premier\footnote{Par \og{} degré de transcendance \fg{} d'un système on entend, comme on sait, le nombre -- invariant -- de fonctions algébriquement indépendantes dont le système dépend algébriquement; un système formé de fonctions algébriquement indépendantes est appelé irréductible. Par degré de transcendance (dimension) d'un idéal premier on entend celui de son corps résiduel. Un idéal premier ne possède pas de diviseur premier propre du même degré de transcendance; cf., par exemple, B. L. v. d. Waerden, \og{} Zur Nullstellentheorie der Polynomideale \fg{}, Math. Ann. 96 (1926), p. 183--208.}. Que cela n'arrive que pour un nombre fini de $p$ résulte de ce que, modulo $p$ aussi, une dépendance algébrique entraîne l'annulation du déterminant fonctionnel (\S\,1), et de ce que $\mathfrak{a}$ se révèle en outre être un idéal premier absolu, si bien qu'il ne peut perdre sa propriété de premier que modulo un nombre fini de $p$\footnote{E. Noether, \og{} Eliminationstheorie und allgemeine Idealtheorie \fg{}, Math. Ann. 90 (1923), p. 229--261, \S\,7. Le théorème résulte, par la théorie de l'élimination\textsuperscript{*)}, du théorème plus simple selon lequel un polynôme absolument irréductible ne perd cette propriété que modulo un nombre fini de nombres premiers (A. Ostrowski, \og{} Zur arithmetischen Theorie der algebraischen Größen \fg{}, Gött. Nachr. 1919, p. 279--298, là p. 296. Preuve plus simple chez E. Noether, \og{} Ein algebraisches Kriterium für absolute Irreduzibilität \fg{}, Math. Ann. 85 (1922), p. 26--33, \S\,7). Si $\mathfrak{I}$ possède une \emph{base d'intégrité} --- c'est-à-dire ici une famille finie de générateurs comme $\mathbf Z$-algèbre --- qui contient exactement une fonction de plus que son degré de transcendance -- comme dans le cas considéré par Hecke (note 2) --, de sorte que $\mathfrak{a}$ devienne un idéal principal, il suffit de faire appel à ce théorème plus simple, en évitant la théorie de l'élimination. Partout, les nombres premiers peuvent être remplacés par des idéaux premiers d'un corps de nombres algébriques fini. Le théorème vaut manifestement aussi lorsque $\mathfrak{a}$ est l'idéal nul.}.

Remarquons encore que tous les raisonnements restent valables, avec de légères modifications, si l'on remplace les entiers rationnels par les entiers d'un corps de nombres algébriques fini et les nombres premiers par ses idéaux premiers (\S\,4).

\begin{quote}
\small
\noindent\textsuperscript{*)} Plus précisément, puisque $\mathfrak{a}$ a ses coefficients en caractéristique zéro, la conclusion tirée de la théorie de l'élimination est indépendante du lemme auxiliaire V de ce travail. La démonstration de ce lemme n'est en effet pas correcte (avec la substitution indiquée, tout pourrait s'annuler identiquement); une démonstration correcte ne se trouve que chez B. L. v. d. Waerden, \emph{Eine Verallgemeinerung des Bézoutschen Theorems}, Math. Ann. 99 (1928), p. 497--541; cf. la note 37.

Une autre correction, facile à voir, est nécessaire à la p. 257 du travail (le théorème XIV n'est pas non plus utilisé au \S\,7). Il y est affirmé que tous les éléments de $(\mathfrak{R})$ s'obtiennent par spécialisation -- en éléments de $(\mathfrak{R}')$ -- des $t_{\mu,\nu}$ dans $(x_i)$; en réalité, ils s'obtiennent par spécialisation d'une forme linéaire plus générale
\[
\sum s_{\lambda_1\ldots\lambda_n}(y_1)^{\lambda_1}\cdots(y_n)^{\lambda_n},
\]
dont l'exposant est toutefois le même que celui de $(x_i)$.
\end{quote}

\subsection*{\S\,1. Déterminants fonctionnels sur des domaines de coefficients de caractéristique $p$}

\smallskip
\noindent\textbf{1.}
Soient $f_1(x_1,\ldots,x_n),\ldots,f_t(x_1,\ldots,x_n)$ des fonctions rationnelles ou, plus généralement, des séries entières formellement définies ou des quotients de telles séries, à coefficients dans un corps $\mathrm{P}$ de caractéristique $p$. Les dérivées $\partial f_i/\partial x_k$ seront elles aussi supposées formellement définies, toutes les règles usuelles du calcul étant conservées. Par $\varOmega$ on entendra un corps d'extension algébriquement clos de $\mathrm{P}$.

Alors, comme en caractéristique zéro, on a :

\noindent\textbf{Théorème.}
\textit{S'il existe entre les $f(x)$ une dépendance algébrique à coefficients dans $\varOmega$, alors le rang de la matrice fonctionnelle $(\partial f_i/\partial x_k)$ ($i=1,\ldots,t$; $k=1,\ldots,n$) est inférieur à $t$}\footnote{Dans la version originale, j'avais supposé que le domaine des coefficients $P$ était un corps parfait; je dois à B. L. v. d. Waerden l'extension à des corps arbitraires, même non parfaits, par passage de $P$ à $\varOmega$. Que la réciproque en caractéristique $p$ ne soit pas valable même pour $P$ parfait est montré par l'exemple $f=x^p$, $df/dx=0$.}.

\noindent\textit{Démonstration.}
Dans l'anneau de polynômes $\varOmega[u_1,\ldots,u_t]$, soit $\mathfrak{m}$ l'idéal premier de tous les polynômes $G(u)$ pour lesquels $G(f)=0$ identiquement en $x$. Par hypothèse, $\mathfrak{m}$ est différent de l'idéal nul; soit $G(u)\ne0$ (et donc non constant) un polynôme de degré minimal dans $\mathfrak{m}$, ce qui implique en particulier que $G(u)$ est irréductible. On obtient alors, comme d'habitude :
\[
\begin{aligned}
\left(\frac{\partial G}{\partial u_1}\right)_{f}\frac{\partial f_1}{\partial x_1}
 +\cdots+
\left(\partial G/\partial u_t\right)_f\partial f_t/\partial x_1&=0,\\[-1mm]
&\vdots\\[-1mm]
\left(\frac{\partial G}{\partial u_1}\right)_{f}\frac{\partial f_1}{\partial x_n}
 +\cdots+
\left(\partial G/\partial u_t\right)_f\partial f_t/\partial x_n&=0.
\end{aligned}
\]
On en déduit que, soit le rang de la matrice fonctionnelle est inférieur à $t$, soit toutes les $\left(\partial G/\partial u_i\right)_{f}$ s'annulent. Mais puisque $G(u)$ avait été choisi de degré minimal, il s'ensuit l'annulation de toutes les $\partial G/\partial u_i$.

Ainsi l'on aurait $G(u)=H(u_1^p,\ldots,u_t^p)$; et puisque le domaine des coefficients $\varOmega$ est algébriquement clos, donc parfait, on aurait en outre $G(u)=H(u_i^p)=(T(u))^p$, contrairement au choix de $G(u)$. La preuve vaut naturellement aussi en caractéristique zéro, le dernier pas étant alors omis.

\medskip
\noindent\textbf{2. Corollaire.}

Soient $F_1(x_1,\ldots,x_n),\ldots,F_t(x_1,\ldots,x_n)$ des fonctions à coefficients entiers (polynômes, séries entières ou quotients de telles fonctions) dont la matrice fonctionnelle est exactement de rang $t$. Choisissons le nombre premier $p$ de façon qu'il ne divise ni les dénominateurs des $F_i(x)$ ni tous les déterminants à $t$ lignes de la matrice fonctionnelle.

Alors $F_1(x),\ldots,F_t(x)$ restent algébriquement indépendantes modulo $p$.

\noindent\textit{Démonstration.}
Soit $\mathrm{P}$ le corps résiduel modulo $p$, et soient $f_1(x),\ldots,f_t(x)$ les fonctions sur $\mathrm{P}$ obtenues à partir de $F_1(x),\ldots,F_t(x)$ en remplaçant les coefficients entiers par leurs classes résiduelles modulo $p$. Les $f_i(x)$ sont définies, puisque $p$ ne divisait pas les dénominateurs des $F_i(x)$; leur matrice fonctionnelle $(\partial f_i/\partial x_k)$ s'obtient également à partir de $(\partial F_i/\partial x_k)$ en remplaçant les coefficients entiers par leurs classes résiduelles modulo $p$ -- dans $(\partial F_i/\partial x_k)$ n'apparaissent pas d'autres dénominateurs que ceux des $F_i(x)$. Par le choix de $p$, $(\partial f_i/\partial x_k)$ conserve le rang $t$; les $f_i(x)$ sont algébriquement indépendantes dans $\varOmega$ et, à plus forte raison, dans $\mathrm{P}$.

\subsection*{\S\,2. Formulation du théorème}

\smallskip
\noindent\textbf{1.}
Comme dans l'introduction, soit $\mathfrak{I}$ un domaine d'intégrité (anneau sans diviseurs de zéro) de fonctions à coefficients entiers en $x_1,\ldots,x_n$ : polynômes ou séries entières à coefficients entiers rationnels, ou quotients de tels objets. Dans le cas des séries entières, il peut s'agir de séries formellement définies ou de séries convergeant dans un certain domaine; on utilise seulement le fait que les fonctions considérées forment un domaine d'intégrité.

Soit $\mathfrak{Q}$ le corps des fractions de $\mathfrak{I}$, et soit $\mathfrak{Z}$ un domaine d'intégrité compris entre $\mathfrak{I}$ et $\mathfrak{Q}$. Comme dans l'introduction, on appelle $\mathfrak{M}$ \emph{maximal} dans $\mathfrak{Z}$ si l'appartenance de $a\cdot g(x)$ -- $a$ entier non nul, $g(x)$ dans $\mathfrak{Z}$ -- à $\mathfrak{M}$ entraîne toujours celle de $g(x)$. À $\mathfrak{I}$ correspond de manière univoque le domaine maximal $\mathfrak{M}(\mathfrak{I})$ engendré par $\mathfrak{I}$ dans $\mathfrak{Z}$, comme intersection de tous les domaines maximaux de $\mathfrak{Z}$ qui contiennent $\mathfrak{I}$ : en effet, $\mathfrak{Z}$ lui-même en est un, et avec un nombre quelconque de tels domaines leur intersection est encore maximale et contient $\mathfrak{I}$. $\mathfrak{M}(\mathfrak{I})$ se compose de toutes les fonctions $g(x)$ de $\mathfrak{Z}$ pour lesquelles il existe au moins un $a \neq 0$ tel que $a\cdot g(x)$ appartienne à $\mathfrak{I}$.

En effet, ces $g(x)$ forment un domaine maximal $\mathfrak{N}$ dans $\mathfrak{Z}$ qui contient $\mathfrak{I}$ : si $b\cdot h(x)$ appartient à $\mathfrak{N}$, avec $b \neq 0$ et $h(x)$ dans $\mathfrak{Z}$, alors il existe un $a\ne0$ tel que $ab\cdot h(x)$ appartienne à $\mathfrak{I}$, et comme $ab \neq 0$, $h(x)$ appartient à $\mathfrak{N}$. Tout domaine maximal $\mathfrak{M}$ de $\mathfrak{Z}$ qui contient $\mathfrak{I}$ contient donc $\mathfrak{N}$; car si $a\cdot g(x)$ appartient à $\mathfrak{I}$, alors $a\cdot g(x)$ appartient à $\mathfrak{M}$, puisque $\mathfrak{I}$ est contenu dans $\mathfrak{M}$, et ainsi $g(x)$ appartient à $\mathfrak{M}$.

\medskip
\noindent\textbf{2.}
Supposons maintenant que $\mathfrak{I}$ soit un domaine d'intégrité fini avec élément unité, ayant $f_1(x),\ldots,f_r(x)$ comme \emph{base d'intégrité}; autrement dit, cette famille engendre $\mathfrak{I}$ comme $\mathbf Z$-algèbre, et $\mathfrak{I}$ se compose de tous les polynômes à coefficients entiers en les $f_i(x)$. De plus, $\mathfrak{Z}$ se compose de tous les polynômes ou séries entières à coefficients entiers en $x$, respectivement des quotients qui ne contiennent pas d'autres dénominateurs que ceux qui apparaissent dans les fonctions $f_i(x)$, en nombre fini, et naturellement avec toutes leurs puissances.

Supposons encore que dans au moins une des $f_i(x)$ -- en général dans toutes -- les $x$ interviennent réellement; sinon $\mathfrak{I}$ et $\mathfrak{Z}$ ne se composent simplement que de tous les polynômes à coefficients entiers en un nombre rationnel fractionnaire $1/c$, et il n'y a rien à prouver.

S'il s'agit de polynômes ou de fonctions rationnelles, le corps des fractions $\mathfrak{Q}$ est alors de degré de transcendance $t$, avec $1 \leqq t \leqq r$, relativement au corps $\mathbf{K}$ des rationnels; et si, par exemple, $f_1(x),\ldots,f_t(x)$ forment un système irréductible (voir note 5), alors $\mathfrak{Q}$ est une extension algébrique finie du corps purement transcendant $\mathrm{K}(f_1(x),\ldots,f_t(x))$. En même temps, la matrice fonctionnelle de $f_1(x),\ldots,f_r(x)$, et de même celle de $f_1(x),\ldots,f_t(x)$, est exactement de rang $t$\footnote{Pour une démonstration purement algébrique de ce théorème sur les déterminants fonctionnels -- il s'agit de la caractéristique zéro --, cf. B. L. v. d. Waerden, \og{} Algebraische Theorie der Differentiation \fg{}, \emph{Nieuw Archief v. Wiskunde} (2) 15 (1926), p. 111--120.}.

Pour que cette structure de $\mathfrak{Q}$ se conserve aussi dans le cas des séries entières ou de leurs quotients, il faut poser une hypothèse additionnelle : si $t$ est le rang de la matrice fonctionnelle des $f_i(x)$ et, par exemple, aussi le rang de la matrice fonctionnelle de $f_1(x),\ldots,f_t(x)$ -- de sorte que celles-ci forment, par \S\,1, 1., un système irréductible --, alors $\mathfrak{Q}$ doit de nouveau être une extension algébrique finie du corps purement transcendant $\mathbf{K}(f_1(x),\ldots,f_t(x))$; c'est-à-dire que $f_{t+1}(x),\ldots,f_r(x)$ doivent être algébriques relativement à $\mathbf{K}(f_1(x),\ldots,f_t(x))$\footnote{Cette hypothèse additionnelle est satisfaite lorsque -- comme dans le cas considéré par Hecke (note 2) -- il s'agit de fonctions analytiques et que $f_1(x),\ldots,f_t(x)$ sont analytiquement indépendantes, tandis que $f_{t+1}(x),\ldots,f_r(x)$ dépendent algébriquement d'elles. Mais l'hypothèse est intentionnellement formulée en termes formels -- rang d'une matrice fonctionnelle -- afin qu'elle ait également un sens clair dans le cas de séries de puissances définies formellement. On obtient ainsi que toutes les considérations demeurent dans le cadre purement formel et algébrique, et que les théorèmes sur les fonctions analytiques ne sont nécessaires que pour vérifier, dans le cas particulier, que l'hypothèse est satisfaite.}. On a de nouveau ici $t \geqq 1$, car -- caractéristique zéro! -- les $\partial f_i/\partial x_k$ ne s'annulent pas toutes.

La condition vaut en même temps pour tout système de $t$ fonctions de la base d'intégrité dont la matrice fonctionnelle est exactement de rang $t$; car celles-ci forment un système irréductible dont les autres dépendent algébriquement.

\medskip
\noindent\textbf{3.}
Soit maintenant $\mathfrak{o}$ l'anneau de polynômes à coefficients entiers en les indéterminées $u_1,\ldots,u_r$. Alors $\mathfrak{I}$ est image homomorphe d'anneaux de $\mathfrak{o}$, comme le montre immédiatement la correspondance $u_i \to f_i(x)$. Ainsi $\mathfrak{I}$ est isomorphe comme anneau à $\mathfrak{o}/\mathfrak{a}$, où $\mathfrak{a}$ est un idéal premier, puisque $\mathfrak{I}$ est un anneau sans diviseurs de zéro; $\mathfrak{a}$ se compose de tous les polynômes $G(u)$ pour lesquels $G(f)$ s'annule identiquement en $x$.

Comme on l'a déjà indiqué dans l'introduction, il s'agit du théorème suivant :

\medskip
\noindent\textbf{Théorème.}
\textit{Sous l'hypothèse additionnelle de 2., soient $\mathfrak{o}_p$, $\mathfrak{a}_p$, $\mathfrak{I}_p$ les domaines obtenus à partir de $\mathfrak{o}$, $\mathfrak{a}$, $\mathfrak{I}$ en remplaçant les coefficients entiers par leurs classes résiduelles modulo $p$. Alors, sauf pour au plus un nombre fini de nombres premiers -- parmi lesquels figurent ceux qui apparaissent dans les dénominateurs des $f_i(x)$, afin que $\mathfrak{I}_p$ soit défini --, l'homomorphisme de $\mathfrak{o}_p$ vers $\mathfrak{I}_p$ a pour noyau $\mathfrak{a}_p$; ainsi $\mathfrak{I}_p$ est isomorphe comme anneau à $\mathfrak{o}_p/\mathfrak{a}_p$, et donc à $\mathfrak{o}/(\mathfrak{a},p\mathfrak{o})$.}

Du théorème résulte le fait que seuls un nombre fini de nombres premiers apparaissent dans les dénominateurs $a$ de $\mathfrak{M}(\mathfrak{I})$, grâce au corollaire suivant.

\medskip
\noindent\textbf{Corollaire.}
\textit{Si $p$ est choisi de telle sorte que $\mathfrak{I}_p$ soit isomorphe à $\mathfrak{o}_p/\mathfrak{a}_p$, alors, dès que $p\cdot g(x)$ appartient à $\mathfrak{I}$ et que $g(x)$ appartient à $\mathfrak{Z}$, on a toujours $g(x)$ dans $\mathfrak{I}$.}

Autrement dit, $\mathfrak{I}$ est maximal dans $\mathfrak{Z}$ relativement à tous les nombres premiers distincts des nombres premiers exceptionnels, en nombre fini.

On a $\mathfrak{o}_p/\mathfrak{a}_p\simeq\mathfrak{o}/(\mathfrak{a},p\mathfrak{o})$. En effet, par définition $\mathfrak{o}_p$ est l'anneau quotient $\mathfrak{o}/p\mathfrak{o}$, et de même $\mathfrak{a}_p$ est $(\mathfrak{a},p\mathfrak{o})/p\mathfrak{o}$. Donc -- premier théorème d'isomorphie -- $\mathfrak{o}_p/\mathfrak{a}_p$, et avec lui aussi $\mathfrak{I}_p$ par hypothèse, est isomorphe à $\mathfrak{o}/(\mathfrak{a},p\mathfrak{o})$.

Cela signifie que
\[
H(f_i(x))\equiv0\ (p)\quad\hbox{dans }\mathfrak{I}
\quad\Longrightarrow\quad
H(u)\equiv0\ (\mathfrak{a},p\mathfrak{o})\quad\hbox{dans }\mathfrak{o}.
\]

On peut reformuler cela ainsi : si $H(f_i(x)) = pg(x)$ [$H(f_i(x))$ dans $\mathfrak{I}$; $g(x)$ dans $\mathfrak{Z}$], alors $H(u_i)-pF(u_i)\equiv0\ (\mathfrak{a})$, donc aussi $H(f_i(x))=pF(f_i(x))$ et par suite $g(x)=F(f_i(x))$; ainsi $g(x)$ appartient à $\mathfrak{I}$.

Une répétition finie donne enfin : si $a$ n'est divisible par aucun des nombres premiers exceptionnels, alors l'appartenance de $a\cdot g(x)$ à $\mathfrak{I}$ et celle de $g(x)$ à $\mathfrak{Z}$ entraînent que $g(x)$ appartient à $\mathfrak{I}$; le corollaire est ainsi démontré.

\subsection*{\S\,3. Démonstration du théorème}

\smallskip
\noindent\textbf{1. Lemme auxiliaire.}
\textit{L'idéal premier $\mathfrak{a}$, noyau de l'homomorphisme de $\mathfrak{o}$ sur $\mathfrak{I}$ (\S\,2, 3.), est un idéal premier absolu.}

Soit $\mathfrak{o}^*$ l'anneau de polynômes en $u_1,\ldots,u_r$ à coefficients algébriques arbitraires -- même fractionnaires --, et soit $\mathfrak{a}^*$ l'idéal de $\mathfrak{o}^*$ engendré par $\mathfrak{a}$; ainsi $\mathfrak{a}^*$ se compose de toutes les combinaisons linéaires $a_1G_1(u)+\cdots+a_sG_s(u)$, où les $a_i$ sont des nombres algébriques arbitraires et les $G_i(u)$ des polynômes de $\mathfrak{a}$. Il faut prouver que $a^*$ est un idéal premier.

Cela sera prouvé si l'on montre que $a^*$ est le noyau de l'homomorphisme de $\mathfrak{o}^*$ vers $\mathfrak{I}^*$ -- où $\mathfrak{I}^*$ désigne le domaine d'intégrité formé de tous les polynômes en les $f_i(x)$ à coefficients algébriques arbitraires --, c'est-à-dire que $H(u)$ appartient à $\mathfrak{a}^*$ ($H(u)$ venant de $\mathfrak{o}^*$) lorsque $H(f_i(x))$ s'annule.

Soit $\omega_1,\ldots,\omega_s$ une base linéairement indépendante du corps de nombres fini engendré par les coefficients, en nombre fini, de $H(u)$. On obtient alors :
\[
a\cdot H(u)=\omega_1H_1(u)+\cdots+\omega_sH_s(u),
\]
où les $H_1(u),\ldots,H_s(u)$ appartiennent à $\mathfrak{o}$ et où $a \neq 0$ est un entier rationnel. L'égalité $H(f)=0$ entraîne alors $\omega_1H_1(f)+\cdots+\omega_sH_s(f)=0$, et donc $H_1(f)=0,\ldots,H_s(f)=0$. Ainsi $H_1(u),\ldots,H_s(u)$ appartiennent à $\mathfrak{a}$, et par conséquent $H(u)$ appartient à $\mathfrak{a}^*$.

\medskip
\noindent\textbf{2.}
La preuve du théorème peut maintenant se donner sous la forme suivante.

Si $p$ est choisi de telle sorte que
\begin{itemize}
\item[I.] $\mathfrak{I}_p$ soit défini;
\item[II.] $\mathfrak{a}_p$ reste un idéal premier, et même absolu;
\item[III.] le degré de transcendance de $\mathfrak{I}_p$ n'ait pas diminué par rapport à celui de $\mathfrak{I}$,
\end{itemize}
alors l'homomorphisme de $\mathfrak{o}_p$ vers $\mathfrak{I}_p$ a pour noyau $\mathfrak{a}_p$. Ces conditions sont satisfaites pour tous les nombres premiers sauf, au plus, un nombre fini de nombres premiers exceptionnels.

Que, pour les conditions I, II et III, il ne s'agisse que d'un nombre fini de nombres premiers exceptionnels résulte immédiatement de ce qui précède.

I est clair : les nombres premiers exceptionnels ne sont que ceux, en nombre fini, qui divisent les dénominateurs des $f_i(x)$; II résulte de ce que $\mathfrak{a}$ est, par le lemme auxiliaire, un idéal premier absolu, et conserve donc cette propriété modulo tous les nombres premiers sauf au plus un nombre fini (cf. note\textsuperscript{6}). Enfin III est démontré dans le corollaire de \S\,1, 2., puisque la matrice fonctionnelle de $f_1(x),\ldots,f_t(x)$ est, par \S\,2, 2., exactement de rang $t$.

Soit maintenant $p$ choisi distinct de ces nombres premiers exceptionnels. L'homomorphisme de $\mathfrak{o}_p$ vers $\mathfrak{I}_p$ a pour noyau un idéal $\mathfrak{c}_p$, qui est premier -- puisque $\mathfrak{I}_p$ aussi est un anneau sans diviseurs de zéro -- et dont le degré de transcendance est donné par celui de $\mathfrak{I}_p$, donc est au moins $t$. En effet, puisque $\mathfrak{I}_p$ est isomorphe à $\mathfrak{o}_p/\mathfrak{c}_p$, le corps résiduel de $\mathfrak{c}_p$ est isomorphe au corps des fractions de $\mathfrak{I}_p$.

Comme $\mathfrak{a}_p$ est lui aussi un idéal premier et que $\mathfrak{a}_p\subseteq\mathfrak{c}_p$, ces deux idéaux coïncident si et seulement si leurs degrés de transcendance coïncident\footnote{Cf., par exemple, B. L. v. d. Waerden, \og{} Zur Nullstellentheorie der Polynomideale \fg{}, Math. Ann. 96 (1926), p. 183--208, \S\,3.}. Compte tenu de cette inclusion, le degré de transcendance de $\mathfrak{a}_p$ est également au moins $t$; plus précisément, les classes modulo $\mathfrak{a}_p$ de $u_1,\ldots,u_t$ forment un système irréductible. Il faut montrer que ce degré de transcendance est exactement égal à $t$, ce qui achèvera la preuve.

C'est ici que l'on utilise, dans le cas des séries entières ou de leurs quotients, l'hypothèse additionnelle sur la structure du corps des fractions $\mathfrak{Q}$ de $\mathfrak{I}$ (\S\,2, 2.). D'après elle, l'idéal premier $\mathfrak{a}$ était dans tous les cas de degré de transcendance $t$; la numérotation était choisie de telle sorte que les classes modulo $\mathfrak{a}$ de $u_1,\ldots,u_t$ forment un système irréductible, et que les classes de $u_{t+1},\ldots u_r$ dépendent algébriquement d'elles. Dans $\mathfrak{a}$ étaient donc contenus des polynômes -- à coefficients entiers et primitifs --
\[
G_1(u_{t+1};u_1,\ldots,u_t),\ldots,G_{r-t}(u_r;u_1,\ldots,u_t)\tag{1}
\]
qui passent à des polynômes -- non nuls par primitivité --
\[
G^*_1(u_{t+1};u_1,\ldots,u_t),\ldots,G^*_{r-t}(u_r;u_1,\ldots,u_t)\tag{2}
\]
de $\mathfrak{a}_p$. Mais, dans chaque $G^*_{\lambda}$, $u_{t+\lambda}$ doit intervenir effectivement; sinon -- contrairement au choix de $p$ -- il y aurait une dépendance algébrique entre les classes modulo $\mathfrak{a}_p$ de $u_1,\ldots,u_t$. Ainsi il est démontré que $\mathfrak{a}_p$ est de degré de transcendance $t$; il est donc identique à $\mathfrak{c}_p$ et constitue le noyau de l'homomorphisme, ce qui prouve tout.

\subsection*{\S\,4. Extension aux coefficients entiers algébriques}

\smallskip
Comme on l'a déjà remarqué à la fin de l'introduction, tous les raisonnements restent valables, avec de légères modifications, si, au lieu des entiers rationnels, on prend les entiers d'un corps de nombres algébriques fini $\mathbf{K}$, et si, au lieu des nombres premiers $p$, on prend les idéaux premiers $\mathfrak{p}$ de ce corps. Le \S\,1 reste entièrement inchangé, de même que les raisonnements qui conduisent dans \S\,2 à la formulation du théorème; dans la preuve (\S\,3, 2.) et dans le corollaire de \S\,2, 3., quelques ajouts sont nécessaires.

\medskip
\noindent\textbf{1.}
Dans \S\,3, 2., il faut ajouter, comme idéaux premiers exceptionnels, ceux, en nombre fini, qui divisent les polynômes $G_1(u),\ldots,G_{r-t}(u)$, \S\,3, 2., (1), puisque ceux-ci ne peuvent plus être supposés primitifs. Ainsi est garantie la non-annulation des polynômes $G^*_1(u),\ldots,G^*_{r-t}(u)$ (\S\,3, 2., (2)); cela représente un renforcement de la condition III. La condition II demeure, car le théorème sur l'irréductibilité absolue ne change pas; de même, naturellement, I demeure.

\medskip
\noindent\textbf{2.}
Le corollaire qui donne le résultat final doit alors être formulé, conformément à la forme finale donnée là, de la manière suivante :

\noindent\textbf{Corollaire.}
\textit{Si l'entier $\alpha$ de $\mathbf{K}$ n'est divisible par aucun des idéaux premiers exceptionnels, en nombre fini, alors l'appartenance de $\alpha\cdot g(x)$ à $\mathfrak{I}$ et celle de $g(x)$ à $\mathfrak{Z}$ entraînent que $g(x)$ appartient à $\mathfrak{I}$.}

Si les entiers $\lambda_1,\ldots,\lambda_s$ de $\mathbf{K}$ sont choisis de telle sorte que chacun soit divisible par un idéal premier exceptionnel, mais non par son carré ni par les autres, alors il suffit, comme dénominateurs $\alpha$, de prendre les produits de puissances de ces $\lambda$.

\noindent\textit{Démonstration.}
Soit
\[
(\alpha)=\mathfrak{p}_1^{\varrho_1}\ldots\mathfrak{p}_g^{\varrho_g}
\]
la décomposition de $\alpha$ en idéaux premiers, et soient $\mathfrak{p}_1,\ldots,\mathfrak{p}_g$ distincts des idéaux premiers exceptionnels; ainsi $\mathfrak{I}_p$ est isomorphe à $\mathfrak{o}_p/\mathfrak{a}_p$ pour chacun des $\mathfrak{p}_i$. De l'anneau de tous les entiers de $\mathbf{K}$, on passe à la localisation $\mathfrak{R}$ définie par $(\alpha)$, en inversant tous et seulement les entiers premiers avec $(\alpha)$, c'est-à-dire ceux qui ne sont divisibles par aucun des idéaux premiers $\mathfrak{p}_1,\ldots,\mathfrak{p}_g$. Dans $\mathfrak{R}$, les idéaux premiers $\mathfrak{p}_1,\ldots,\mathfrak{p}_g$ restent premiers, mais deviennent principaux; tous les autres deviennent l'idéal unité, et la décomposition de $(\alpha)$ demeure la même\footnote{Cf., par exemple, H. Grell, \og{} Zur Theorie der Ordnungen in algebraischen Zahl- und Funktionenkörpern \fg{}, Math. Ann. 97 (1927), p. 524--558, \S\,2, 1.}. Par conséquent, lors du passage à $\mathfrak{R}$ comme domaine de coefficients, le corollaire et sa preuve restent exactement sous la forme de \S\,2, 3.

Mais cela signifie : si $\alpha\cdot g(x)$ appartient à $\mathfrak{I}$, alors $g(x)=F(f_i(x))$, où le polynôme $F(u)$ a ses coefficients dans $\mathfrak{R}$. Si $\delta$ est un dénominateur commun de ces coefficients -- donc $\delta$ premier avec $\alpha$ --, cela admet aussi la formulation suivante : si $\alpha\cdot g(x)$ appartient à $\mathfrak{I}$, alors $\delta\cdot g(x)$ appartient à $\mathfrak{I}$ avec $\delta$ premier à $\alpha$. Ainsi $(\delta,\alpha)$ est l'idéal unité, et l'on obtient donc : si $\alpha\cdot g(x)$ appartient à $\mathfrak{I}$, alors $g(x)$ appartient à $\mathfrak{I}$. La première assertion du corollaire ci-dessus est démontrée.

Pour prouver la seconde assertion, soit $\mathfrak{R}^*$ la localisation dans $\mathrm{K}$ définie par les idéaux premiers exceptionnels, en nombre fini, et soient $\lambda_1,\ldots,\lambda_s$ des éléments de base de ces idéaux premiers dans $\mathfrak{R}^*$, que l'on peut supposer entiers. Dans $\mathfrak{R}^*$, $\alpha$ s'écrit, à une unité de $\mathfrak{R}^*$ près, comme un produit de puissances des $\lambda$. En revenant aux entiers, on obtient donc
\[
\gamma\alpha=\beta\lambda_1^{\varrho_1}\cdots\lambda_s^{\varrho_s},
\]
où $\gamma$ et $\beta$ ne sont divisibles par aucun des idéaux premiers exceptionnels. Si $\alpha\cdot g(x)$ appartient à $\mathfrak{I}$, alors $\beta \cdot \lambda_1^{\sigma_1} \ldots \lambda_s^{\sigma_s} g(x)$ appartient à $\mathfrak{I}$, et la première assertion du corollaire entraîne que $\lambda_1^{\sigma_1} \ldots \lambda_s^{\sigma_s} g(x)$ appartient aussi à $\mathfrak{I}$.

\bigskip
\noindent\rule{\textwidth}{0.4pt}

\medskip
\begingroup
\begin{center}
{\fontspec{Noto Serif}\Large\bfseries О максимальных областях целочисленных функций.}\par
\vspace{0.4em}
{\itshape Sur les domaines maximaux de fonctions à coefficients entiers.}\par
\vspace{0.6em}
{\bfseries Emmy Noether (Göttingen).}\par
\vspace{0.4em}
\emph{Résumé (traduit du russe).}
\end{center}

\smallskip
\textit{Tout domaine d'intégrité $\mathfrak{I}$ (Integritätsbereich, anneau sans diviseurs de zéro) définit un certain domaine maximal $\mathfrak{M}(\mathfrak{I})$ du même type, composé de toutes les fonctions $g(x)$ pour lesquelles il existe au moins un $a\ne0$ tel que $a\cdot g(x)$ soit contenu dans $\mathfrak{I}$.}

\textit{Dans ce travail on étudie la question des \og{} dénominateurs \fg{} $a$ de ce domaine maximal $\mathfrak{M}(\mathfrak{I})$, sous l'hypothèse que le domaine $\mathfrak{I}$ était fini. On y démontre que ces $a$ peuvent toujours être choisis de telle sorte qu'en eux n'intervienne, au total, qu'un nombre fini de nombres premiers comme facteurs.}

\textit{En même temps, dans le cas où les $g(x)$ sont des séries entières ou des quotients de telles séries, on suppose que les fonctions $f_1(x),\ldots,f_r(x)$ qui définissent le domaine $\mathfrak{I}$ ne peuvent être liées entre elles que par des relations algébriques.}

\textit{La preuve repose sur la réduction du problème à certaines questions de la théorie générale des idéaux. Cette réduction s'accomplit en considérant l'apparition du nombre premier $p$ dans le dénominateur $a$ comme une relation modulo $p$ entre les fonctions $f_i$.}

\begin{center}
(Rec. Math. XXXVI:1; 1929).
\end{center}
\endgroup

\bigskip
\noindent\rule{\textwidth}{0.4pt}

\setcounter{footnote}{0}
\clearpage
\section*{\NoetherNavAnchor{P36}{36. Différentiation des idéaux et différente}36. Différentiation des idéaux et différente}

\begin{center}
\emph{J. Ber. d. DMV 39 (1930), p. 17}
\end{center}

2. E. Noether, Göttingen : différentiation des idéaux et différente.

On montre comment la différente d'un corps de nombres algébriques peut être conçue comme le \emph{quotient différentiel} d'un \emph{idéal} définissant. Que cette définition coïncide avec la définition usuelle résulte de théorèmes de structure intéressants en eux-mêmes, qui constituent, en un sens analogue, la généralisation de la formule d'interpolation de Lagrange.

Un exposé détaillé doit paraître dans les \emph{Mathematische Annalen}.

\setcounter{footnote}{0}
\clearpage
\section*{\NoetherNavAnchor{P37}{37. Base normale dans les corps sans ramification supérieure}37. Base normale dans les corps sans ramification supérieure}

\begin{center}
\emph{Journal f. d. reine u. angew. Math. 167 (1932), p. 147--152}
\end{center}

\begin{center}
Par \emph{Emmy Noether} à Göttingen.
\end{center}

Par base normale d'un corps de nombres galoisien \(K/k\), on entend une base de l'anneau des entiers \(\mathfrak{O}\) de \(K\) sur l'anneau des entiers \(\mathfrak{o}\) de \(k\), composée d'éléments conjugués. Une condition nécessaire à l'existence d'une telle base est, comme on le sait\footnote{Voir A. Speiser, Gruppendeterminante und Körperdiskriminante, \S 6. Math. Ann. 77 (1916), p. 546--562.}, qu'il n'apparaisse dans \(K/k\) aucun corps de ramification supérieure. Cela vaut encore lorsque l'on se restreint à une place, c'est-à-dire lorsque l'on passe à l'extension \(\mathfrak{p}\)-adique \(k_{\mathfrak{p}}\) de \(k\) et à l'extension correspondante \(K_{\mathfrak{p}}\) de \(K\). Je montre dans ce qui suit la validité de la réciproque:

\emph{Pour toutes les places \(\mathfrak{p}\) où \(\mathfrak{p}\) ne divise pas le degré de \(K/k\), il existe une base normale. En particulier donc: si \(K/k\) ne possède pas de ramification supérieure, il existe à chaque place ramifiée une base normale; le discriminant y devient le carré d'un déterminant de groupe.}

À cette dernière assertion il faut ajouter que, contrairement à une conjecture d'E. Artin, le passage à ses conducteurs\footnote{E. Artin, Über die gruppentheoretische Struktur der Diskriminante. J. f. Math. 164 (1931), p. 1--11.} exige encore d'autres considérations: la décomposition du déterminant de groupe suivant les caractères irréductibles donne des résolvantes de Lagrange généralisées; un regroupement convenable donne une décomposition dans le corps de base élargi par les caractères. Dans les cas les plus simples, je peux identifier ces nouveaux facteurs aux conducteurs d'Artin, au moyen de la décomposition des résolvantes de Lagrange du point de vue de la théorie des idéaux.\footnote{[6. 10. 31.] Que, même en général, des considérations relevant de la théorie des idéaux soient encore nécessaires est montré par l'exemple suivant, communiqué par M. Deuring, et susceptible de variations arbitraires, d'un ordre non maximal dont le discriminant peut se représenter comme le carré du déterminant de groupe, tandis qu'à des caractères conjugués correspondent des facteurs idéaux différents: soit \(k\) le corps des racines cinquièmes de l'unité, \(K=k(\sqrt[5]{2})\). On considère l'ordre \([1,2\sqrt[5]{2},\sqrt[5]{2^2},\sqrt[5]{2^3},\sqrt[5]{2^4}]\), et cela à la place \((2)\), où, d'après les considérations de cette note, il possède une base normale. La décomposition suivant les caractères donne, pour le déterminant de groupe correspondant, précisément les éléments de base ci-dessus comme facteurs. Les ``conducteurs'' deviennent donc, excepté \(1\): \(2\sqrt[5]{2}\cdot\sqrt[5]{2^4}\), \(\sqrt[5]{2^2}\cdot\sqrt[5]{2^3}\), \(\sqrt[5]{2^3}\cdot\sqrt[5]{2^2}\), \(\sqrt[5]{2^4}\cdot 2\sqrt[5]{2}\), c'est-à-dire \(4,2,2,4\); ils sont donc différents pour des caractères conjugués.} Je voudrais encore remarquer que, même dans le cas général où il n'existe pas de base normale, une scission en résolvantes de Lagrange généralisées est toujours possible, au moyen d'une ``série de composition suivant les modules de Galois''; et qu'il en résulte, de façon analogue, une décomposition dans le corps de base élargi; ici recommencent des questions de théorie des idéaux.

La preuve de l'existence de la base normale sous les hypothèses ci-dessus repose sur le fait que \(\mathfrak{O}/\mathfrak{o}\), considéré comme module de Galois, c'est-à-dire comme \(\mathfrak{o}\)-module ayant les substitutions du groupe de Galois pour domaine d'opérateurs, est isomorphe, comme module, à un idéal unilatère de l'anneau de groupe entier étendu par \(\mathfrak{o}\). Celui-ci est un ordre maximal à toutes les places de ramification ordinaires; or les idéaux dans les ordres maximaux d'algèbres semi-simples \(\mathfrak{p}\)-adiques sont des idéaux principaux\footnote{Voir H. Hasse, Über \(\mathfrak{p}\)-adische Schiefkörper und ihre Bedeutung für die Arithmetik hyperkomplexer Zahlsysteme, théorèmes 46 et 47. Math. Ann. 104 (1931), p. 495--534. Pour les algèbres commutatives, la propriété d'idéal principal vaut déjà après localisation de \(k\) en \(\mathfrak{p}\); on peut donc, lorsqu'il s'agit de groupes et de corps abéliens, se restreindre à cette extension plus faible pour définir la place.}; ils sont donc ici engendrés par un élément au moyen des substitutions du groupe, ou, respectivement, de l'anneau de groupe. En revenant au module de Galois, cela signifie exactement l'existence de la base normale. Aux places de ramification supérieure, l'anneau de groupe entier devient un ordre non maximal, et l'idéal associé à \(\mathfrak{O}/\mathfrak{o}\) devient un idéal non principal.

Toutes ces considérations subsistent manifestement lorsqu'il s'agit de corps de fonctions d'une variable, pourvu que la caractéristique du corps des constantes ne divise pas le degré de \(K/k\).

\subsection*{\S1. L'anneau de groupe entier étendu \(\mathfrak{p}\)-adiquement}

\(\mathfrak{o}\), respectivement \(\mathfrak{o}_{\mathfrak{p}}\), désigneront les anneaux des entiers d'un corps de nombres algébrique \(k\) et de son extension \(\mathfrak{p}\)-adique \(k_{\mathfrak{p}}\), où \(\mathfrak{p}\) est un idéal premier de \(k\). De plus, \((\mathfrak{G})\) désignera l'algèbre de groupe à coefficients rationnels et \([\mathfrak{G}]\) l'anneau de groupe à coefficients entiers d'un groupe \(\mathfrak{G}\) à \(n\) éléments. Par \((\mathfrak{G})_k\), respectivement \((\mathfrak{G})_{k_{\mathfrak{p}}}\), on entend l'algèbre de groupe à coefficients dans \(k\), respectivement dans \(k_{\mathfrak{p}}\); de même, par \([\mathfrak{G}]_{\mathfrak{o}}\), respectivement \([\mathfrak{G}]_{\mathfrak{o}_{\mathfrak{p}}}\), l'extension de l'anneau de groupe entier à un tel anneau à coefficients dans \(\mathfrak{o}\), respectivement dans \(\mathfrak{o}_{\mathfrak{p}}\).

Ainsi \((\mathfrak{G})_k\) et \((\mathfrak{G})_{k_{\mathfrak{p}}}\) sont des algèbres sans radical (des algèbres semi-simples), et \([\mathfrak{G}]_{\mathfrak{o}}\), respectivement \([\mathfrak{G}]_{\mathfrak{o}_{\mathfrak{p}}}\), des ordres de \((\mathfrak{G})_k\), respectivement de \((\mathfrak{G})_{k_{\mathfrak{p}}}\).

\paragraph{Théorème 1.}
\emph{Si l'idéal premier \(\mathfrak{p}\) ne divise pas le nombre d'éléments \(n\) de \(\mathfrak{G}\), alors \([\mathfrak{G}]_{\mathfrak{o}_{\mathfrak{p}}}\) est un ordre maximal de \((\mathfrak{G})_{k_{\mathfrak{p}}}\); si \(\mathfrak{p}\) divise \(n\), alors \([\mathfrak{G}]_{\mathfrak{o}_{\mathfrak{p}}}\) n'est pas maximal.}

Le discriminant de l'anneau de groupe entier, donc aussi celui de \([\mathfrak{G}]_{\mathfrak{o}}\) et de \([\mathfrak{G}]_{\mathfrak{o}_{\mathfrak{p}}}\), est égal à une puissance du nombre d'éléments \(n\).\footnote{Voir par exemple E. Noether, Hyperkomplexe Größen und Darstellungstheorie, \S 26. Math. Ztschr. 30 (1929), p. 641--692.} Si donc \(\mathfrak{p}\) ne divise pas \(n\), le discriminant de \([\mathfrak{G}]_{\mathfrak{o}_{\mathfrak{p}}}\) est une unité. Mais cela signifie que \([\mathfrak{G}]_{\mathfrak{o}_{\mathfrak{p}}}\), en maintenant fixe \(\mathfrak{o}_{\mathfrak{p}}\), ne peut être élargi à un ordre plus grand; car un tel élargissement retirerait du discriminant un facteur carré non inversible. Or \(\mathfrak{o}_{\mathfrak{p}}\) était déjà supposé être l'anneau des entiers, donc maximal dans \(k_{\mathfrak{p}}\). Ainsi est démontrée la première partie du théorème 1, seule utilisée dans la suite.

Si au contraire \(\mathfrak{p}\) divise \(n\), la somme directe correspondant à la représentation triviale
\[
        \mathfrak{C}=E^{(1)}[\mathfrak{G}]_{\mathfrak{o}_{\mathfrak{p}}}+(1-E^{(1)})[\mathfrak{G}]_{\mathfrak{o}_{\mathfrak{p}}},
        \qquad
        E^{(1)}=\frac1n\sum_{S\in\mathfrak{G}} S,
\]
est une extension propre de \([\mathfrak{G}]_{\mathfrak{o}_{\mathfrak{p}}}\). En effet, puisque \(E^{(1)}=\frac1n\sum S\), \(\mathfrak{C}\) est une extension propre; \(\mathfrak{C}\) est un ordre avec les générateurs dépendants \(E^{(1)}\), \((1-E^{(1)})S\), \(S\in\mathfrak{G}\) (\(E^{(1)}S=E^{(1)}\)).

\paragraph{Théorème 2.}
\emph{Si \(\mathfrak{p}\) ne divise pas \(n\), tout idéal entier ou fractionnaire de \([\mathfrak{G}]_{\mathfrak{o}_{\mathfrak{p}}}\) est principal.}

En effet, d'après le théorème 1, \([\mathfrak{G}_{\mathfrak o_\mathfrak p}]\) est un ordre maximal d'une algèbre semi-simple \(\mathfrak{p}\)-adique; les idéaux sont donc des idéaux principaux.\footnote{Hasse, loc. cit. La propriété d'idéal principal n'y est démontrée que pour les algèbres simples; elle en découle pourtant aussi pour les algèbres semi-simples par des raisonnements connus. En effet, les composantes d'un ordre maximal sont encore maximales; les composantes de l'idéal considéré sont donc des idéaux principaux, par exemple de base \(a_i\). Alors \(a=\sum a_i\) devient une base de l'idéal initial. -- Pour les groupes abéliens, voir encore la note 3.}

\subsection*{\S2. Modules de Galois, isomorphismes de modules, base normale}

Soit \(K/k\) galoisien, de groupe \(\mathfrak{G}\); \(z^S\) désignera l'élément de \(K\) obtenu à partir de \(z\) par la substitution \(S\) de \(\mathfrak{G}\). Les substitutions de \(\mathfrak{G}\) induisent sur \(K/k\), considéré comme \(k\)-module, c'est-à-dire comme groupe abélien additif ayant \(k\) pour domaine d'opérateurs, des automorphismes de module. Par conséquent, \(K/k\) peut être défini comme module sur \((\mathfrak{G})_k\) par la convention:
\[
        z\Bigl(\sum_i S_i c_i\Bigr)=\sum_i z^{S_i}c_i,
        \qquad z\in K,\; S_i\in\mathfrak{G},\; c_i\in k.
\]

\paragraph{Définition.}
\emph{Tout \((\mathfrak{G})_k\)-module contenu dans \(K/k\) est appelé module de Galois rationnel. De même, les \([\mathfrak{G}]_{\mathfrak{o}}\)-modules contenus dans \(K/k\) sont appelés modules de Galois entiers; en particulier, \(\mathfrak{O}/\mathfrak{o}\) est donc un module de Galois entier.}

\paragraph{Théorème 3.}
\emph{Comme module de Galois rationnel, \(K/k\) est isomorphe à \((\mathfrak{G})_k\) comme module.}

Cela résulte du fait que \(K/k\) possède toujours une base normale de corps, c'est-à-dire une \(k\)-base composée d'éléments conjugués \(z^S\).\footnote{Car, si \(a_1,\ldots,a_n\) est une base de \(K/k\) et si les \(u_i\) désignent des indéterminées, les \(\sum u_i a_i^S\) sont linéairement indépendants lorsque \(S\) parcourt les éléments de \(\mathfrak{G}\); par conséquent, les \(u_i\) peuvent aussi être spécialisés en éléments de \(k\) de telle sorte que l'indépendance linéaire soit conservée.} La correspondance
\[
        S\to z^S,
        \qquad
        \sum S_i c_i\to\sum z^{S_i}c_i
        \quad(c_i\in k),
        \quad\hbox{donc } ST\to z^{ST},
\]
donne un homomorphisme de modules qui, en raison du même rang sur \(k\), est un isomorphisme.

\paragraph{Théorème 4.}
\emph{Comme module de Galois entier, \(\mathfrak{O}/\mathfrak{o}\) est isomorphe, comme module, à un \([\mathfrak{G}]_{\mathfrak{o}}\)-module de \((\mathfrak{G})_k\), de rang maximal \(n\). De même, \(\mathfrak{O}_{\mathfrak{p}}/\mathfrak{o}_{\mathfrak{p}}\) est isomorphe, comme module, à un \([\mathfrak{G}]_{\mathfrak{o}_{\mathfrak{p}}}\)-module de rang \(n\) dans \((\mathfrak{G})_{k_{\mathfrak{p}}}\).}

L'isomorphisme de \((\mathfrak{G})_k\)-modules décrit au théorème 3, de \(K/k\) vers \((\mathfrak{G})_k\), associe au \([\mathfrak{G}]_{\mathfrak{o}}\)-module \(\mathfrak{O}/\mathfrak{o}\) un \([\mathfrak{G}]_{\mathfrak{o}}\)-module de rang \(n\) dans \((\mathfrak{G})_k\). De plus, l'isomorphisme de \((\mathfrak{G})_k\)-modules s'étend à un isomorphisme de \((\mathfrak{G})_{k_{\mathfrak{p}}}\)-modules de \(K_{\mathfrak{p}}/k_{\mathfrak{p}}\) sur \((\mathfrak{G})_{k_{\mathfrak{p}}}\), en laissant simplement les \(c_i\) de la correspondance du théorème 3 parcourir tous les éléments de \(k_{\mathfrak{p}}\). Cet isomorphisme induit alors un isomorphisme de \([\mathfrak{G}]_{\mathfrak{o}_{\mathfrak{p}}}\)-modules de \(\mathfrak{O}_{\mathfrak{p}}/\mathfrak{o}_{\mathfrak{p}}\). Ici, \(K_{\mathfrak{p}}/k_{\mathfrak{p}}\) est à considérer comme une algèbre hypercomplexe sur \(k_{\mathfrak{p}}\). \(K_{\mathfrak{p}}/k_{\mathfrak{p}}\) naît de \(K/k\) par extension des coefficients; il devient en général une algèbre à diviseurs de zéro, plus précisément une somme de corps isomorphes, donc une algèbre semi-simple.

\paragraph{Théorème 5.}
\emph{À toute place \(\mathfrak{p}\) qui ne divise pas le degré \(n\) de \(K/k\), \(\mathfrak{O}_{\mathfrak{p}}/\mathfrak{o}_{\mathfrak{p}}\) possède une base normale.}\footnote{Dans le cas des corps abéliens, la place peut alors aussi être définie, d'après les notes 3 et 5, par localisation en \(\mathfrak{p}\).}

En effet, d'après le théorème 1, \([\mathfrak{G}]_{\mathfrak{o}_{\mathfrak{p}}}\) y est un ordre maximal; les \([\mathfrak{G}]_{\mathfrak{o}_{\mathfrak{p}}}\)-modules de rang \(n\) dans \((\mathfrak{G})_{k_{\mathfrak{p}}}\) deviennent donc des idéaux entiers ou fractionnaires, et même des idéaux principaux d'après le théorème 2. Si donc l'idéal associé à \(\mathfrak{O}_{\mathfrak{p}}/\mathfrak{o}_{\mathfrak{p}}\) est égal à \(W[\mathfrak{G}]_{\mathfrak{o}_{\mathfrak{p}}}\), il possède la \(\mathfrak{o}_{\mathfrak{p}}\)-base \(WS_i\); l'isomorphisme de modules montre que \(w^{S_i}\) est une \(\mathfrak{o}_{\mathfrak{p}}\)-base de \(\mathfrak{O}_{\mathfrak{p}}/\mathfrak{o}_{\mathfrak{p}}\), lorsque \(w\) désigne l'élément de \(\mathfrak{O}_{\mathfrak{p}}/\mathfrak{o}_{\mathfrak{p}}\) associé à \(W\). Les \(n\) éléments conjugués \(w^{S_i}\) forment donc une base normale de \(\mathfrak{O}_{\mathfrak{p}}/\mathfrak{o}_{\mathfrak{p}}\).

Si \(\mathfrak{p}\) divise \(n\), le module associé à \(\mathfrak{O}_{\mathfrak{p}}/\mathfrak{o}_{\mathfrak{p}}\) ne peut pas être un module principal, puisqu'il ne peut alors exister aucune base normale.

\paragraph{Remarque supplémentaire au théorème 3.}
De l'isomorphisme de modules entre \((\mathfrak{G})_k\) et \(K/k\) établi au théorème 3\footnote{La preuve suppose manifestement seulement que \(K\) soit galoisien de première espèce sur \(k\), et que \(k\) possède une infinité d'éléments. Cette dernière restriction est inutile: M. Deuring a trouvé une preuve du théorème 3 valable aussi pour les corps finis, dont il résulte réciproquement l'existence de la base normale de corps. En même temps, on obtient ainsi une preuve de l'existence de l'élément primitif, qui vaut uniformément pour les corps finis et infinis.} découlent les résultats de Speiser sur le problème des formes de Klein,\footnote{Loc. cit. La notion de module de Galois en est abstraite. -- Si la caractéristique de \(k\) divise \(n\), il s'agit de séries de composition suivant les modules de Galois et de leurs facteurs irréductibles.} que l'on peut formuler ainsi:

\emph{Toute représentation irréductible définie sur \(k\) du groupe de Galois \(\mathfrak{G}\) de \(K/k\) est engendrée par un module de Galois rationnel irréductible contenu dans \(K/k\); toute représentation absolument irréductible l'est par un module de Galois contenu dans \(K_{\mathsf{Z}}/\mathsf{Z}\).}

Ici, \(\mathsf{Z}\) désigne un corps de décomposition contenant \(k\), relatif à la composante correspondante de l'anneau de groupe; \(K_{\mathsf{Z}}/\mathsf{Z}\) est à considérer comme une algèbre hypercomplexe sur \(\mathsf{Z}\), obtenue par extension des coefficients à partir de \(K/k\). En particulier, \(K_{\mathsf{Z}}\) reste un corps si \(\mathsf{Z}\) peut être choisi linéairement disjoint de \(K\) sur \(k\), de sorte que \(K_{\mathsf{Z}}/\mathsf{Z}\) devienne une extension accessoire (le cas traité par Speiser).

L'assertion elle-même résulte, en vertu de cet isomorphisme de modules, de l'assertion correspondante pour \((\mathfrak{G})_k\), respectivement \((\mathfrak{G})_Z\), où les idéaux irréductibles (relativement à \((\mathfrak G)_k\), respectivement à \((\mathfrak G)_Z\)) engendrent la représentation.

Si \(v_1,\ldots,v_t\) est une base d'un tel module de Galois, et si \(S\to \bar S\) est la représentation correspondante, on a donc:
\[
        \begin{pmatrix}\vdots\\ v_i^S\\ \vdots\end{pmatrix}
        =\bar S\begin{pmatrix}\vdots\\ v_i\\ \vdots\end{pmatrix}.
\]
C'est la formulation de Klein et Speiser.

\subsection*{\S3. Le discriminant comme déterminant de groupe dans les corps sans ramification supérieure}

Pour décomposer le discriminant, il suffit, comme on sait, de considérer la décomposition aux différentes places ramifiées; dans les corps sans ramification supérieure, il ne s'agit donc que de places possédant une base normale.

\paragraph{Théorème 6.}
\emph{Dans les corps galoisiens \(K/k\) sans ramification supérieure, le discriminant devient, à toute place ramifiée, égal au carré du déterminant de groupe du groupe de Galois, les indéterminées étant remplacées par la base normale. Aux représentations irréductibles correspondent des décompositions du déterminant de groupe en résolvantes de Lagrange généralisées; le discriminant se décompose en idéaux du corps de base élargi par les caractères, les mêmes idéaux correspondant aux caractères complexes conjugués.}

En effet, avec \(w^S\) comme base normale, le discriminant \(\Delta\) devient le carré de
\[
        D=\bigl|w^{ST^{-1}}\bigr|;\qquad S,T\in\mathfrak{G}.
\]
Or \(D\) est le déterminant de groupe avec les \(w^S\) à la place des indéterminées \(x_S\).

La décomposition en résolvantes de Lagrange résulte de raisonnements de Speiser, loc. cit. Soit
\[
        D=D_1^{f_1}\cdots D_t^{f_t},
        \qquad
        M=f_1M_1+\cdots+f_tM_t
\]
la décomposition du déterminant de groupe, respectivement de la matrice de groupe, correspondant aux représentations absolument irréductibles. Si \(S\to\bar S\) désigne la représentation correspondant à \(M_\lambda\), les \(\bar S\) appartenant au corps de base élargi, il vient donc:
\[
        M_\lambda=\sum_S w^S\bar S,
        \qquad
        M_\lambda^{T^{-1}}=\sum_S w^{ST^{-1}}\bar S
          =\sum_R w^R\bar R\bar T=M_\lambda\bar T,
\]
et, en passant au déterminant,
\[
        D_\lambda^{T^{-1}}=D_\lambda|\bar T|=D_\lambda\varepsilon_T.
\]
Ainsi, les \(D_\lambda\) sont les résolvantes de Lagrange généralisées; \(\varepsilon_T\), déterminant de \(\bar T\), est une racine de l'unité (représentation irréductible du groupe quotient de \(\mathfrak{G}\) par son groupe des commutateurs). Dans le cas des groupes cycliques, les \(D_\lambda\) deviennent simplement les résolvantes de Lagrange \(\sum w^S\chi_\lambda(S)\).

En général, les \(D_\lambda\) appartiennent à l'algèbre hypercomplexe obtenue à partir de \(K_{\mathfrak{p}}/k_{\mathfrak{p}}\) par extension du domaine des coefficients \(k_{\mathfrak{p}}\) au corps du caractère correspondant; et il s'agit là d'\emph{éléments entiers} de cette algèbre. En effet, le déterminant formé avec les indéterminées \(x_S\) possède comme coefficients des entiers du corps du caractère; les \(w^S\), eux, sont des éléments entiers de \(K_{\mathfrak{p}}\).

Pour passer de la décomposition du déterminant de groupe à celle du discriminant, on regroupe chaque fois une représentation avec son adjointe. Si \(\lambda,\bar\lambda\) désignent des représentations adjointes, on a donc simultanément:
\[
        D_\lambda^{T^{-1}}=D_\lambda\varepsilon_T,
        \qquad
        D_{\bar\lambda}^{T^{-1}}=D_{\bar\lambda}\varepsilon_T^{-1};
\]
en même temps, les déterminants formés pour des indéterminées \(x_S\) et appartenant à \(\lambda,\bar\lambda\) sont conjugués complexes, tandis qu'en général, à des caractères conjugués correspondent, pour les indéterminées \(x_S\), des déterminants conjugués.

Si l'on pose donc \(\Delta_\lambda=D_\lambda D_{\bar\lambda}\), il suffit d'adjoindre le sous-corps réel du \(\lambda\)-ième caractère, et l'on a
\[
        \Delta_\lambda^T=\Delta_\lambda
        \qquad\hbox{pour tout }T\hbox{ de }\mathfrak{G};
\]
et par suite\footnote{Comme il s'agit de l'extension hypercomplexe des coefficients d'un corps, le raisonnement usuel de la théorie de Galois doit être modifié. Soient \(\mathsf{P}\) le corps d'extension en question de \(k_{\mathfrak{p}}\) et
\[
        (K_{\mathfrak{p}})_{\mathsf{P}}=(K_{\mathfrak{p}})_{\mathsf{P}} e^{S_1}+\cdots+(K_{\mathfrak{p}})_{\mathsf{P}} e^{S_r}
\]
la décomposition en somme directe de corps. Alors l'extension \( (K_{\mathfrak{p}})_{\mathsf{P}} e^{S_i}/\mathsf{P}e^{S_i}\) est galoisienne (le groupe étant un sous-groupe du groupe de décomposition d'un facteur premier de \(\mathfrak{p}\)) et les \(e^{S_i}\) sont conjugués relativement aux classes latérales. De \(\Delta_\lambda^T=\Delta_\lambda\), il résulte d'abord que la \(i\)-ième composante de \(\Delta_\lambda\) appartient à \(\mathsf{P}e^{S_i}\), donc s'écrit \(\gamma_i e^{S_i}\), avec \(\gamma_i\) dans \(\mathsf{P}\). Il en résulte encore, à cause de la conjugaison des \(e^{S_i}\), que tous les \(\gamma_i\) sont égaux, donc que, effectivement, \(\Delta_\lambda=\gamma\sum e^{S_i}=\gamma\) appartient à \(\mathsf{P}\).}: \(\Delta_\lambda\) appartient au corps de base élargi par le sous-corps réel du \(\lambda\)-ième caractère.

\emph{La décomposition du discriminant}
\[
        \Delta=\Delta_1^{f_1}\cdots\Delta_t^{f_t}
\]
\emph{est donc une décomposition dans le corps de base élargi par les sous-corps réels des caractères, les mêmes facteurs correspondant aux caractères complexes conjugués. Pour les autres caractères conjugués, on ne peut rien conclure sans argument supplémentaire [cf. 2a)].}

Le théorème 6 est ainsi démontré dans toutes ses parties. Si l'on peut montrer que les idéaux engendrés par les \(\Delta_\lambda\) appartiennent en réalité déjà au corps de base \(k\) et sont égaux pour les caractères conjugués, ils coïncident avec les conducteurs d'Artin, dès que l'on associe aux caractères composés les produits de puissances correspondants des \(\Delta_\lambda\). En effet, on a aussi bien l'équation fonctionnelle que le fait, découlant directement de la base normale (voir Speiser, loc. cit.), que les représentations triviales des sous-groupes correspondent aux discriminants des sous-corps associés. De la coïncidence à chaque place résulte alors la coïncidence des idéaux globaux. Si l'on se restreint aux caractères rationnellement irréductibles, les faits indiqués donnent cette coïncidence. Pour les caractères absolument irréductibles, la coïncidence doit encore être vérifiée directement dans le cas le plus simple:

\paragraph{Théorème 7.}
\emph{Si \(k\) est le corps des nombres rationnels, et si \(K/k\) est cyclique de degré premier \(l\), ce degré étant premier au discriminant, donc sans ramification supérieure, alors les idéaux engendrés par les \(\Delta_\lambda\), pour les représentations non triviales, sont égaux entre eux et coïncident avec le conducteur de \(K/k\).}

Cela résulte de la décomposition connue des résolvantes de Lagrange.\footnote{Hilbert, Zahlbericht, \S 108. Comme, d'après le théorème 5, l'existence de la base normale est établie -- la place peut, d'après la note 7, être définie simplement par localisation en \(\mathfrak{p}\) --, il n'est plus nécessaire de recourir au fait, utilisé dans le Zahlbericht, que les corps absolument cycliques sont des corps cyclotomiques.} À une place ramifiée \(\mathfrak{p}\) de \(K/k\), on a en effet, pour \(k_\varepsilon\), le corps des racines \(l\)-ièmes de l'unité, et pour \(K_\varepsilon\) (\(K_\varepsilon\) reste un corps, puisque \(k_\varepsilon\) est linéairement disjoint de \(K\) et que l'extension \(\mathfrak{p}\)-adique est ici superflue d'après la note 7):
\[
        (p)=\mathfrak{p}_1\cdots\mathfrak{p}_{l-1},
        \qquad
        \mathfrak{p}_i=\mathfrak{P}_i^{\,l},
\]
avec \(\mathfrak{p}_i\) dans \(k_\varepsilon\), \(\mathfrak P\) dans \(K_\varepsilon\). De plus, pour les résolvantes de Lagrange, on a:
\[
        (\Omega_\lambda)=\mathfrak{P}_1^{r_1}\cdots\mathfrak{P}_{l-1}^{r_{l-1}},
        \qquad
        (\Omega_{\bar\lambda})=
        \mathfrak{P}_1^{l-r_1}\cdots\mathfrak{P}_{l-1}^{l-r_{l-1}},
\]
où \(r_1,\ldots,r_{l-1}\) est une permutation des nombres \(1,\ldots,l-1\). Or les résolvantes de Lagrange \(\Omega_\lambda\) sont ici identiques aux facteurs \(D_\lambda\) du déterminant de groupe; il vient donc:
\[
        (\Delta_\lambda)=(D_\lambda D_{\bar\lambda})
        =(\Omega_\lambda\Omega_{\bar\lambda})
        =\mathfrak{P}_1^{l}\cdots\mathfrak{P}_{l-1}^{l}
        =\mathfrak{p}_1\cdots\mathfrak{p}_{l-1}=(p);
\]
donc coïncidence avec le conducteur de \(K/k\).

\begin{center}
Reçu le 24 août 1931.
\end{center}

\setcounter{footnote}{0}
\clearpage
\section*{\NoetherNavAnchor{P38}{38. Avec R. Brauer et H. Hasse : preuve d’un théorème principal dans la théorie des algèbres}38. Avec R. Brauer et H. Hasse: preuve d'un théorème principal dans la théorie des algèbres}

\begin{center}
\emph{Journal f. d. reine u. angew. Math. 167 (1932), p. 399--404}
\end{center}

\begin{center}
\textbf{Preuve d'un théorème principal dans la théorie des algèbres.}\\[0.5em]
Par \emph{R. Brauer} à Koenigsberg, \emph{H. Hasse} à Marbourg et \emph{E. Noether} à Goettingen\footnote{La rédaction de cette note fut prise en charge par H. Hasse.}.
\end{center}

Enfin, nos efforts conjoints ont permis d'établir le théorème suivant, fondamental pour la théorie structurale des algèbres sur les corps de nombres algébriques et au-delà:

\paragraph{Théorème principal.}
\emph{Toute algèbre à division normale sur un corps de nombres algébrique est cyclique (ou, comme on dit aussi, du type de Dickson).}

Nous avons le plaisir particulier de dédier ce résultat --- succès dû pour l'essentiel à la méthode \(p\)-adique --- à M. Kurt Hensel, fondateur de cette méthode, à l'occasion de son soixante-dixième anniversaire.

Notre preuve comporte trois réductions, dont chacun de nous a fourni une.\footnote{Elles sont reproduites dans l'ordre de leur genèse, qui est inverse de l'ordre systématique.}

\paragraph{1.}
La première réduction fut donnée par H. Hasse sur la base de la théorie des algèbres cycliques sur les corps de nombres algébriques qu'il avait récemment développée.\footnote{H. Hasse, \emph{Theorie der zyklischen Algebren ueber einem algebraischen Zahlkoerper}, Goett. Nachr. 1931. Un exposé détaillé des preuves, qui développe notamment la théorie des corps de décomposition et des produits croisés élaborée par E. Noether dans un cours, fondamentale là comme ici, paraîtra prochainement dans les Trans. Amer. Math. Soc. sous le titre \emph{Theory of cyclic algebras over an algebraic number field}. Ce travail sera cité dans la suite par H. H; 1--6 constituent, à la différence de la langue près, la première note.}

\paragraph{Réduction 1.}
Il suffit, pour démontrer le théorème principal, d'établir que:
\[
\text{I. Toute algèbre sur }\Omega\text{ qui se décompose partout est } \sim \Omega .
\]
Pour abréger l'expression, ici comme dans la suite, ``algèbre \(A\) sur \(\Omega\)'' signifiera toujours ``algèbre simple normale \(A\) sur le corps de nombres algébrique \(\Omega\)'', c'est-à-dire un système hypercomplexe simple de centre \(\Omega\). De plus, \(A\sim\Omega\) signifie que \(A\) est une algèbre complète de matrices sur \(\Omega\), donc l'appartenance de \(A\) à la classe spéciale, au sens de la similitude -- c'est-à-dire l'égalité des algèbres à division associées sur \(\Omega\) --, déterminée par \(\Omega\) lui-même comme algèbre à division sur \(\Omega\) [H, 5]. Enfin, par algèbre sur \(\Omega\) se décomposant partout, on entend une algèbre telle que, pour toute place première \(\mathfrak p\) de \(\Omega\), l'extension \(\mathfrak p\)-adique, c'est-à-dire l'algèbre obtenue par extension du corps des coefficients \(\Omega\) au corps \(\mathfrak p\)-adique correspondant \(\Omega_{\mathfrak p}\), se décompose:
\[
A_{\mathfrak p}\sim \Omega_{\mathfrak p}.
\]

\paragraph{Preuve.}
Soit \(D\) une algèbre à division sur \(\Omega\). Si \(\mathsf{Z}\) est alors un corps cyclique sur \(\Omega\) tel que, pour toute place première \(\mathfrak p\) de \(\Omega\), le \(\mathfrak p\)-degré de \(\mathsf{Z}\) soit un multiple du \(\mathfrak p\)-indice de \(D\) [cf. la note relative à H, 17 Bb], alors, pour les diviseurs premiers correspondants \(\mathfrak P\) dans \(\mathsf{Z}\), le corps \(\mathfrak P\)-adique \(\mathsf{Z}_{\mathfrak P}\) est chaque fois corps de décomposition de \(D_{\mathfrak p}\) [H, 18.1]. Il en résulte que l'algèbre \(D_{\mathsf{Z}}\), obtenue par extension du corps des coefficients de \(D\) à \(Z\), se décompose partout sur \(\mathsf{Z}\). Si I est déjà acquis, pour tout \(\Omega\), il s'ensuit donc que \(D_{\mathsf{Z}}\sim \mathsf{Z}\). Cela signifie que \(D\) possède le corps de décomposition cyclique \(\mathsf{Z}\), donc qu'elle est représentable cycliquement [H, 5]. Alors \(D\) possède aussi un tel corps de décomposition cyclique dont le degré coïncide avec le degré de \(D\) lui-même; autrement dit, \(D\) est cyclique [H, 6, théorème 6]. Sous l'hypothèse I, le théorème principal est donc démontré.

\paragraph{2.}
Pour la deuxième réduction, R. Brauer donna l'impulsion décisive en communiquant par lettre à H. Hasse comment la question voisine de la valeur exacte de l'exposant d'une algèbre à division (question qui est d'ailleurs résolue en même temps par le théorème principal, voir ci-dessous) peut être ramenée, au moyen du théorème des groupes de Sylow, au cas d'un corps de décomposition résoluble. Cette idée put alors être utilisée aussi pour une réduction correspondante de la question de cyclicité.

\paragraph{Réduction 2.}
L'assertion I est démontrée s'il est montré que:
\[
\text{II. Toute algèbre partout décomposée ayant un corps de décomposition résoluble sur }\Omega\text{ est }\sim\Omega .
\]

\paragraph{Preuve.}
Soit \(A\) une algèbre partout décomposée sur \(\Omega\). Si \(\mathsf{K}\) est un corps de décomposition galoisien pour \(A\), si de plus \(p\) est un nombre premier quelconque et \(\Sigma\) l'un des corps de Sylow de \(K\) sur \(\Omega\) relatifs à \(p\) (c'est-à-dire le corps des invariants d'un groupe de Sylow relatif à \(p\) du groupe de Galois), alors \(A_{\Sigma}\) est une algèbre également partout décomposée sur \(\Sigma\), qui possède le corps de décomposition résoluble \(\mathsf{K}\). Si II est déjà acquis (pour tout \(\Omega\)), il s'ensuit donc que \(A_{\Sigma}\sim \Sigma\). Cela signifie que \(A\) possède le corps de décomposition \(\Sigma\), de degré premier à \(p\). L'indice de \(A\), étant diviseur de ce degré, est donc premier à \(p\). Comme cela vaut pour tout nombre premier \(p\), il est donc égal à 1, c'est-à-dire \(A\sim\Omega\). Sous l'hypothèse II, l'assertion I est donc démontrée.\footnote{L'idée de la réduction à un corps de décomposition résoluble au moyen du théorème des groupes de Sylow avait déjà été employée auparavant par R. Brauer, pour montrer que tout diviseur premier de l'indice apparaît aussi dans l'exposant (Ueber den Zusammenhang von arithmetischen und invariantentheoretischen Eigenschaften von Gruppen linearer Substitutionen, Berl. Akad.-Ber. 1926). Plus récemment A. A. Albert a développé, pour cette idée et plus généralement pour une série de théorèmes généraux de la théorie de R. Brauer et E. Noether, des preuves simples indépendantes de la théorie des représentations (1. \emph{On direct products, cyclic division algebras, and pure Riemann matrices}; 2. \emph{On direct products}; tous deux dans Trans. Amer. Math. Soc. 33 (1931); pour la réduction ici en question, voir surtout le théorème 23 dans 2.).

\emph{Addition pendant l'impression.} En outre A. A. Albert, ayant en sa possession la communication épistolaire de H. Hasse selon laquelle le théorème principal avait été démontré par lui pour les algèbres abéliennes (voir plus bas dans le texte), en a immédiatement déduit, indépendamment de nous, les faits suivants:

a) le théorème principal pour les degrés de la forme \(2^e\),

b) le théorème 1 ci-dessous (exposant = indice),

c) outre l'idée fondamentale de la réduction 2, celle aussi de la réduction 3 suivante, naturellement sans relation avec la réduction 1, et par conséquent avec le résultat: pour les algèbres à division \(D\) de degré \(p^e\), puissance d'un nombre premier, sur \(\Omega\), il existe un corps d'extension \(\Omega'\) de degré premier à \(p\) sur \(\Omega\), tel que \(D_{\Omega'}\) soit cyclique.

Ces trois résultats sont naturellement dépassés par notre démonstration désormais achevée du théorème principal. Ils montrent cependant qu'A. A. Albert a lui aussi contribué indépendamment à sa preuve. Enfin, après avoir pris connaissance de notre démonstration, A. A. Albert a remarqué que notre théorème central I découle en quelques lignes des théorèmes 13, 10 et 9 d'un travail de lui en cours d'impression (Bull. Amer. Math. Soc. 37 (1931)). La preuve de ces théorèmes repose essentiellement sur les mêmes raisonnements que nos réductions 2 et 3.}

\paragraph{3.}
La troisième réduction --- dont H. Hasse reconnut, une fois en possession des deux premières, qu'elle conduisait à la preuve définitive --- fut donnée par E. Noether à la suite d'une communication de H. Hasse. Celle-ci démontrait le théorème principal dans le cas particulier d'un corps de décomposition abélien, au moyen de la réduction générale 1 et d'une réduction explicite du système de facteurs correspondant; E. Noether en dégagea précisément le noyau sous la forme de la troisième réduction. Indépendamment d'E. Noether, R. Brauer avait d'ailleurs déjà envisagé cette troisième réduction.

\paragraph{Réduction 3.}
L'assertion II est démontrée s'il est montré que:
\[
\begin{gathered}\text{III. Toute algèbre partout décomposée ayant un corps de décomposition cyclique}\\\text{de degré premier sur }\Omega\text{ est }\sim\Omega .\end{gathered}
\]

\paragraph{Preuve.}
Soit \(A\) une algèbre partout décomposée ayant un corps de décomposition résoluble \(\mathsf{K}\) sur \(\Omega\). Soit
\[
K=\Lambda_0>\Lambda_1>\cdots>\Lambda_r=\Omega
\]
une chaîne de corps entre \(\mathsf{K}\) et \(\Omega\), telle que \(\Lambda_i\) soit toujours cyclique de degré premier sur \(\Lambda_{i+1}\). Alors \(A_{\Lambda_1}\) est d'abord une algèbre également partout décomposée sur \(\Lambda_1\), qui possède le corps de décomposition cyclique de degré premier \(\mathsf{K}=\Lambda_0\). Si III est déjà acquis (pour tout \(\Omega\)), il s'ensuit donc que \(A_{\Lambda_1}\sim \Lambda_1\). Cela signifie que \(\Lambda_1\) est corps de décomposition pour \(A\). On démontre maintenant de la même façon, en partant de \(\Lambda_1\) au lieu de \(\Lambda_0\), que \(A_{\Lambda_2}\sim\Lambda_2\), et ainsi de suite, jusqu'à obtenir finalement \(A_{\Lambda_r}\sim\Lambda_r\), c'est-à-dire \(A\sim\Omega\). Sous l'hypothèse III, l'assertion II est donc démontrée.

\paragraph{4.}
Or III est acquis d'après les résultats de H. Hasse [H, 24]. Par conséquent, en remontant les réductions 3, 2, 1, on obtient l'exactitude du théorème principal.

E. Noether indique encore à ce propos ce qui suit: l'exactitude de III, plus généralement pour des corps de décomposition cycliques de degré quelconque, revient au théorème de norme de Hasse\footnote{\label{fn:p38-hasse-normensatz}Cité dans H, 3.11; démontré dans: H. Hasse, \emph{Beweis eines Satzes und Widerlegung einer Vermutung ueber das allgemeine Normenrestsymbol}, Goett. Nachr. 1931.}. Pour la réduction 3, seul le cas spécial du degré premier est cependant nécessaire, donc le théorème de norme originel de Hilbert-Furtwaengler\footnote{Voir H. Hasse, Bericht ueber neuere Untersuchungen und Probleme in der Theorie der algebraischen Zahlkoerper II, Jahresber. der D. M.-V., Erg.-Bd. 6 (1930), \S 8, ainsi que la littérature qui y est citée.}. Ainsi la réduction III fournit une nouvelle preuve simple du théorème de norme de Hasse. Cette preuve se déroulerait ainsi:

Soit \(Z\) un corps cyclique sur \(\Omega\) et \(\alpha\) un nombre de \(\Omega\), pour lequel le symbole de reste normique
\[
\left(\frac{\alpha,Z}{\mathfrak p}\right)=1
\]
vaut pour toute place première \(\mathfrak p\) de \(\Omega\). Alors toute algèbre cyclique
\[
A=(\alpha,Z)
\]
\footnote{Dans la notation de H, 1. -- L'indication de l'automorphisme \(S\) de \(Z\) lié à \(\alpha\) a été omise ici comme inessentielle.} se décompose partout [H, 17.7]. D'après la réduction 3 il s'ensuit donc, en utilisant seulement le théorème de norme de Hilbert-Furtwaengler, que \(A\sim\Omega\). Mais alors \(\alpha\) est la norme d'un élément de \(Z\) [H, 15.4].

En rapport avec le théorème de norme, nous remarquons encore que l'assertion I doit être regardée comme sa véritable généralisation aux cas supérieurs (même non abéliens), tandis que la généralisation littérale, comme H. Hasse l'a montré\textsuperscript{\ref{fn:p38-hasse-normensatz}}, n'est pas vraie en général.

\begin{center}
\textbf{Conséquences (H. Hasse).}
\end{center}

\paragraph{5.}
Comme conséquence la plus immédiate du théorème principal, signalons que la théorie des algèbres cycliques sur les corps de nombres algébriques développée par H. Hasse a désormais acquis la portée d'une théorie générale de la structure et des invariants des algèbres simples normales (en particulier des algèbres à division normales) sur les corps de nombres algébriques. En particulier, le résultat suivant est maintenant démontré en toute généralité:

\paragraph{Théorème 1.}
L'exposant d'une algèbre simple normale sur un corps de nombres algébrique est égal à son indice.

De plus, l'assertion I donne le fait remarquable suivant:

\paragraph{Théorème 2.}
L'idéal fondamental d'une vraie algèbre à division normale sur \(\Omega\) est divisible par au moins une place première de \(\Omega\), et même par au moins deux.

Ici l'idéal fondamental est défini par exemple comme la norme (réduite) de la différente (réduite).\footnote{Dans le cas spécial des algèbres de quaternions rationnelles, cela revient à la notion de ``nombre fondamental'' introduite par H. Brandt (Idealtheorie in Quaternionenalgebren, Math. Ann. 99 (1928)). Comme il ne s'agit pas ici d'un nombre, mais d'un idéal, il fallait dire ``idéal fondamental''; on ne doit pas confondre cela avec les dénominations de Dedekind idéal fondamental et nombre fondamental pour la différente et le discriminant, qui justement pour la raison indiquée se révèlent inutilisables pour les corps relatifs. -- Pour la définition de la différente (réduite), voir la note suivante. La norme (réduite) d'un idéal est également définie par E. Noether ``place première par place première'', plus précisément, conformément au fait que, pour les places premières individuelles, tout idéal est principal, tout simplement par formation des normes de nombres (réduites) des nombres de base d'idéaux principaux pour les places premières individuelles.} En outre, une place première infinie est incluse dans l'idéal fondamental si et seulement si, pour elle, l'algèbre se réduit à l'algèbre de quaternions (et non au corps des nombres réels ou complexes).

\paragraph{Preuve.}
D'après les résultats de H. Hasse\footnote{\label{fn:p38-hasse-padisch}H. Hasse, \emph{Ueber \(\mathfrak p\)-adische Schiefkoerper und ihre Bedeutung fuer die Arithmetik hyperkomplexer Zahlsysteme}, Math. Ann. 104 (1931); voir là les théorèmes 42, 59.}, une place première \(\mathfrak p\) de \(\Omega\) divise l'idéal fondamental si et seulement si le \(\mathfrak p\)-indice est différent de \(1\). Si tous les \(\mathfrak p\)-indices étaient égaux à \(1\), l'algèbre se décomposerait partout, et d'après l'assertion I il n'y aurait alors pas de vraie algèbre à division. (Le fait qu'il ne puisse pas non plus y avoir un seul \(\mathfrak p\)-indice différent de \(1\) résulte de ce que les symboles de reste normique, dont les ordres sont les \(\mathfrak p\)-indices [H, 17.7], satisfont au théorème du produit (loi de réciprocité) [H, 3.8].)

\paragraph{6.}
En outre, l'assertion I permet de déterminer (caractérisation arithmétique) tous les corps de décomposition \(\mathsf{K}\) appartenant à une algèbre \(A\) sur \(\Omega\). Dans une généralisation exacte de l'état de choses déjà établi par H. Hasse dans le cas spécial où \(A,\mathsf{K}\) sont cycliques [H, 6, théorème 2], on obtient:

\paragraph{Théorème 3.}
Pour une algèbre \(A\) sur \(\Omega\), un corps de nombres algébrique \(\mathsf{K}\) sur \(\Omega\) est corps de décomposition si et seulement si, pour tous les diviseurs premiers \(\mathfrak P_i\) dans \(\mathsf{K}\) de toutes les places premières \(\mathfrak p\) de \(\Omega\), le \(\mathfrak P_i\)-degré \(n_{\mathfrak P_i}\) de \(\mathsf{K}\) est chaque fois un multiple du \(\mathfrak p\)-indice \(m_{\mathfrak p}\) de \(A\).

Ici le \(\mathfrak P_i\)-degré de \(\mathsf{K}\) est le degré du corps \(\mathfrak P_i\)-adique \(\mathsf{K}_{\mathfrak P_i}\), correspondant à \(\mathfrak P_i\), sur le corps \(\mathfrak p\)-adique \(\Omega_{\mathfrak p}\), c'est-à-dire, si
\[
\mathfrak p=\prod_i \mathfrak P_i^{\,e_{\mathfrak P_i}},
        \qquad \mathsf{N}_{\mathsf{K}\Omega}(\mathfrak P_i)=\mathfrak p^{f_{\mathfrak P_i}},
\]
est la décomposition de \(\mathfrak p\) dans \(\mathsf{K}\), le produit du degré \(f_{\mathfrak P_i}\) par l'ordre de ramification \(e_{\mathfrak P_i}\) de \(\mathfrak P_i\) relativement à \(\mathfrak p\), \(n_{\mathfrak P_i}=f_{\mathfrak P_i}e_{\mathfrak P_i}\) [H, 4]. Et le \(\mathfrak p\)-indice de \(A\) est l'indice \(m_{\mathfrak p}\) de l'extension \(\mathfrak p\)-adique \(A_{\mathfrak p}\) [H, 6].

\paragraph{Preuve.}
Dire que \(\mathsf{K}\) est corps de décomposition pour \(A\), c'est-à-dire \(A_{\mathsf{K}}\sim \mathsf{K}\), équivaut d'après l'assertion I à dire que \(A_{\mathsf{K}_{\mathfrak P_i}}\sim \mathsf{K}_{\mathfrak P_i}\) pour tout \(\mathfrak P_i\).

Pour une place première finie \(\mathfrak p\), soit
\[
        A_{\mathfrak p}\sim(\pi,W_{\mathfrak p})
\]
la représentation cyclique arithmétiquement distinguée de \(A_{\mathfrak p}\) [H, 16; en outre loc. cit.\textsuperscript{\ref{fn:p38-hasse-padisch}}, théorème 38], donc \(W_{\mathfrak p}\) le corps non ramifié de degré \(m_{\mathfrak p}\) sur \(\Omega_{\mathfrak p}\) et \(\pi\) un nombre de \(\Omega_{\mathfrak p}\) divisible exactement par \(\mathfrak p^1\). Soit de plus \(\mathit{K}_{\mathfrak P_i}\) le composé et
\[
        \Delta_{\mathfrak P_i}=K_{\mathfrak P_i}\cap W_{\mathfrak p}
\]
l'intersection de \(W_{\mathfrak p}\) et \(\mathsf{K}_{\mathfrak P_i}\). Comme le sous-corps non ramifié maximal contenu dans \(\mathsf{K}_{\mathfrak P_i}\) est de degré \(f_{\mathfrak P_i}\) sur \(\Omega_{\mathfrak p}\), et comme il n'existe, pour chaque degré, qu'un seul corps non ramifié sur \(\Omega_{\mathfrak p}\), \(\Delta_{\mathfrak P_i}\) est
\[
        \text{de degré }d_{\mathfrak P_i}=(m_{\mathfrak p},f_{\mathfrak P_i})\text{ sur }\Omega_{\mathfrak p}.
\]
Par conséquent \(\mathit{K}_{\mathfrak P_i}\) est de degré \(\frac{m_{\mathfrak p}}{d_{\mathfrak P_i}}\) sur \(\mathsf{K}_{\mathfrak P_i}\), et non ramifié.

Or
\[
        A_{K_{\mathfrak P_i}}=(A_{\mathfrak p})_{K_{\mathfrak P_i}}\sim(\pi,K_{\mathfrak P_i}W_{\mathfrak p})
\]
[H, 15.5]. Par suite, \(A_{\mathsf{K}_{\mathfrak P_i}}\sim \mathsf{K}_{\mathfrak P_i}\) équivaut à ce que \(\pi\) soit norme de \(\mathit{K}_{\mathfrak P_i}\) relativement à \(\mathsf{K}_{\mathfrak P_i}\). Mais, en raison du caractère non ramifié de \(\mathit{K}_{\mathfrak P_i}\) sur \(\mathsf{K}_{\mathfrak P_i}\), cela est le cas si et seulement si l'ordre \(e_{\mathfrak P_i}\) de \(\pi\) en \(\mathfrak P_i\) est un multiple du degré \(\frac{m_{\mathfrak p}}{d_{\mathfrak P_i}}\) de \(\mathit{K}_{\mathfrak P_i}\) sur \(\mathsf{K}_{\mathfrak P_i}\), ou encore, ce qui revient au même à cause de
\[
\left(\frac{m_{\mathfrak p}}{d_{\mathfrak P_i}},\frac{f_{\mathfrak P_i}}{d_{\mathfrak P_i}}\right)=1,
\]
si
\[
\frac{f_{\mathfrak P_i}}{d_{\mathfrak P_i}}e_{\mathfrak P_i}=\frac{n_{\mathfrak P_i}}{d_{\mathfrak P_i}}
\]
est un multiple de \(\frac{m_{\mathfrak p}}{d_{\mathfrak P_i}}\), c'est-à-dire si \(n_{\mathfrak P_i}\) est un multiple de \(m_{\mathfrak p}\), comme il était affirmé.

Dans le cas où \(\mathfrak p\) est une place première infinie et où l'on n'est pas dans le cas trivial \(m_{\mathfrak p}=1\), on a \(m_{\mathfrak p}=2\) et \(A_{\mathfrak p}\) est l'algèbre de quaternions sur le corps réel \(\Omega_{\mathfrak p}\). Dire que \(A_{\mathsf{K}_{\mathfrak P_i}}=(A_{\mathfrak p})_{\mathsf{K}_{\mathfrak P_i}}\sim \mathsf{K}_{\mathfrak P_i}\) équivaut alors à dire que \(\mathsf{K}_{\mathfrak P_i}\) est le corps des nombres complexes, c'est-à-dire à \(n_{\mathfrak P_i}=f_{\mathfrak P_i}=2=m_{\mathfrak p}\) (\(e_{\mathfrak P_i}\) n'intervient pas ici), ce qui donne ici encore l'assertion (\(n_{\mathfrak P_i}\), comme \(m_{\mathfrak p}\), ne peut prendre que les valeurs 1 ou 2).

\paragraph{7.}
On arrive à des résultats significatifs si l'on inverse la question, résolue dans le théorème 3, de tous les corps de décomposition \(\mathsf{K}\) d'une algèbre fixe \(A\): pour un corps algébrique fixe \(\mathsf{K}\) sur \(\Omega\), on considère l'ensemble de toutes les algèbres \(A\) sur \(\Omega\) décomposées par \(K\). Ces algèbres \(A\) forment un sous-groupe \(\mathfrak K\), déterminé par \(K\), du groupe de R. Brauer \(\mathfrak A\) de toutes les algèbres (plus exactement de toutes les classes d'algèbres semblables) sur \(\Omega\).\footnote{R. Brauer, \emph{Ueber Systeme hyperkomplexer Groessen}, Jahresber. d. D. M.-V. 38 (1929), p. 47/48. -- Voir aussi H, 13.1.} Si l'on suppose \(\mathsf{K}\) galoisien sur \(\Omega\), on parvient ainsi à des théorèmes que l'on peut regarder comme des généralisations de théorèmes principaux de la théorie des corps de classes (théorie des corps de nombres relativement abéliens) aux corps de nombres relativement galoisiens généraux:

\paragraph{Théorème 4.}
\emph{(Théorème de décomposition).} Le degré relatif \(f\) des diviseurs premiers dans \(\mathsf{K}\) d'un idéal premier \(\mathfrak p\) de \(\Omega\) ne divisant pas le discriminant relatif de \(K\) est égal au plus petit exposant pour lequel
\[
        A_{\mathfrak p}^{\,f}\sim\Omega_{\mathfrak p}
\]
vaut pour toutes les algèbres \(A\) du groupe \(\mathfrak K\) associé à \(\mathsf{K}\).

\paragraph{Preuve.}
Comme \(f\) est le \(\mathfrak p\)-degré de \(\mathsf{K}\) (le même pour tous les diviseurs premiers de \(\mathfrak p\) dans \(\mathsf{K}\)), tandis que, d'autre part, le \(\mathfrak p\)-indice de \(A\) est égal à l'indice, donc à l'exposant, de \(A_{\mathfrak p}\), il résulte du théorème 3 que \(f\) est en tout cas un multiple de ce plus petit exposant; il reste donc seulement, pour la preuve, à montrer qu'il existe dans le groupe \(\mathfrak K\) des algèbres \(A\) ayant exactement le \(\mathfrak p\)-indice \(f\).

D'après le théorème de densité de Frobenius, il existe sûrement encore un autre idéal premier \(\mathfrak p'\) de \(\Omega\), ne divisant pas le discriminant relatif de \(K\), dont les diviseurs premiers dans \(\mathsf{K}\) ont le degré relatif \(f\). Soit alors \(Z\) un corps cyclique sur \(\Omega\) dont le \(\mathfrak p\)-degré et le \(\mathfrak p'\)-degré aient la valeur \(f\). Soit de plus \(\alpha\) un nombre de \(\Omega\), pour lequel les symboles de reste normique
\[
\left(\frac{\alpha,Z}{\mathfrak p}\right)\quad\text{et}\quad
\left(\frac{\alpha,Z}{\mathfrak p'}\right)
\]
aient des valeurs réciproques d'ordre \(f\), tandis que pour tous les autres diviseurs \(\mathfrak q\) du conducteur de \(Z\) on a
\[
\left(\frac{\alpha,Z}{\mathfrak q}\right)=1.
\]
D'après le théorème généralisé de la progression arithmétique, \(\alpha\) peut de plus être choisi de façon qu'en dehors éventuellement de \(\mathfrak p,\mathfrak p'\) et des \(\mathfrak q\), il ne contienne plus qu'un seul idéal premier \(\mathfrak r\) de \(\Omega\), exactement à la première puissance. Pour celui-ci, d'après le théorème du produit pour le symbole de reste normique (loi de réciprocité), on a également
\[
\left(\frac{\alpha,Z}{\mathfrak r}\right)=1
\]
[H, 3.8--10]. Toute algèbre
\[
        (\alpha,Z)=A
\]
a alors le \(\mathfrak p\)-indice et le \(\mathfrak p'\)-indice \(f\), mais pour toutes les autres places premières de \(\Omega\) l'indice \(1\) [H, 17.7]. D'après le théorème 3, il en résulte, en tenant compte de l'hypothèse sur \(\mathfrak p\) et du choix de \(\mathfrak p'\), que \(\mathsf{K}\) est corps de décomposition pour \(A\). Ainsi sont effectivement mises en évidence dans le groupe \(\mathfrak K\) des algèbres \(A\) de \(\mathfrak p\)-indice \(f\).

\paragraph{Théorème 5.}
\emph{(Théorème d'unicité et d'ordre.)} Si \(\mathsf{K}\leqq \mathsf{K}'\), alors \(\mathfrak K\leqq\mathfrak K'\), et réciproquement.

En particulier, l'association des groupes d'algèbres \(\mathfrak K\) aux corps galoisiens \(\mathsf{K}\) est donc biunivoque.

\paragraph{Preuve.}
a) De \(\mathsf{K}\leqq \mathsf{K}'\), il résulte trivialement que \(\mathfrak K\leqq\mathfrak K'\): toute algèbre \(A\) décomposée par \(K\) est, a fortiori, décomposée par \(\mathsf{K}'\).

b) Soit réciproquement \(\mathfrak K\leqq\mathfrak K'\). Alors, d'après le théorème de décomposition, on a \(f_{\mathfrak p}\mid f'_{\mathfrak p}\) pour tout non-diviseur \(\mathfrak p\) du discriminant relatif. En particulier, on a donc \(f_{\mathfrak p}=1\) pour presque tous les \(\mathfrak p\) tels que \(f'_{\mathfrak p}=1\). Mais il en résulte, par le procédé analytique usuel\footnote{Voir à ce sujet par exemple H. Hasse [loc. cit. note 6], \S 25, III.}, que \(\mathsf{K}\leqq \mathsf{K}'\).

\paragraph{8.}
Enfin, le théorème principal apporte un progrès essentiel à la question, étudiée par I. Schur\footnote{I. Schur, \emph{Arithmetische Untersuchungen ueber endliche Gruppen linearer Substitutionen}, Berl. Akad.-Ber. 1906.}, des corps de nombres dans lesquels les représentations absolument irréductibles d'un groupe fini sont réalisables:

\paragraph{Théorème 6.}
Les représentations absolument irréductibles d'un groupe fini \(\mathfrak G\) sont toutes possibles dans des corps cyclotomiques, par exemple toujours dans le corps des racines \(n^h\)-ièmes de l'unité, si \(n\) est l'ordre de \(\mathfrak G\) et si \(h\) est suffisamment grand.

\paragraph{Preuve.}
Si l'on passe à l'algèbre de groupe \(G=\mathbb Q[\mathfrak G]\), les représentations absolument irréductibles \(\Gamma_i\) de \(\mathfrak G\) deviennent les représentations absolument irréductibles des constituants simples \(G_i\) de l'algèbre semi-simple \(G\), et leurs centres sont chaque fois les corps \(\Omega_i\) des caractères correspondants,\footnote{Voir à ce sujet: a) R. Brauer et E. Noether, \emph{Ueber minimale Zerfaellungskoerper irreduzibler Darstellungen}, Berl. Akad.-Ber. 1927; \S 1. b) R. Brauer, \emph{Ueber Systeme hyperkomplexer Zahlen}, Math. Zeitschr. 30 (1929); théorème 3. c) E. Noether, \emph{Hyperkomplexe Groessen und Darstellungstheorie}, Math. Zeitschr. 30 (1929); \S\S 21, 24, 26.} donc, en raison de la finitude de \(\mathfrak G\), des corps cyclotomiques, et même certainement des sous-corps du corps des racines \(n\)-ièmes de l'unité.

D'après le théorème 3, un corps cyclique \(\mathsf{Z}_i\) sur \(\Omega_i\) est maintenant corps de décomposition pour \(G_i\), si, pour toute place première \(\mathfrak p\) de \(\Omega_i\), son \(\mathfrak p\)-degré \(n_{i\mathfrak p}\) est un multiple du \(\mathfrak p\)-indice \(m_{i\mathfrak p}\) de \(G_i\). \(m_{i\mathfrak p}\) n'est différent de 1 que pour les diviseurs premiers de l'idéal fondamental de \(G_i\) relativement à \(\Omega_i\), donc certainement seulement pour les diviseurs premiers du discriminant absolu de \(G\); comme celui-ci divise \(n^n\) -- \(n^n\) est le discriminant d'un ordre (non maximal) dans \(G\)\footnote{Voir E. Noether [loc. cit. note 12c], \S 26.} --, donc tout au plus pour les diviseurs premiers \(\mathfrak p\) de \(n\). Pour faire de \(n_{i\mathfrak p}\), pour ces \(\mathfrak p\), un multiple de \(m_{i\mathfrak p}\), il suffit de prescrire que les \(\mathfrak p\) en question soient ramifiés dans \(\mathsf{Z}_i\) avec un ordre chaque fois divisible par \(m_{i\mathfrak p}\). Or c'est ce que réalise, comme on le voit facilement, le corps des racines \(n^h\)-ièmes de l'unité pour \(h\) suffisamment grand.

Dans la dernière phrase du travail cité, I. Schur constate que, dans tous les cas connus jusqu'ici, le corps des racines \(n\)-ièmes de l'unité suffit déjà. La question de savoir si cela vaut toujours, et si les méthodes développées ici suffisent pour trancher cette question, reste réservée à des recherches ultérieures.

\begin{center}
Reçu le 11 novembre 1931.
\end{center}

\setcounter{footnote}{0}
\clearpage
\section*{\NoetherNavAnchor{P39}{39. Systèmes hypercomplexes dans leurs rapports avec l’algèbre commutative et la théorie des nombres}39. Systèmes hypercomplexes dans leurs rapports avec l'algèbre commutative et la théorie des nombres}

\begin{center}
\emph{Verhandl. Intern. Math.-Kongreß Zürich 1 (1932), p. 189--194}
\end{center}

\begin{center}
\textsc{Par Emmy Noether, Göttingen}
\end{center}

\paragraph{1.}
La théorie des systèmes hypercomplexes, des algèbres, a connu ces dernières années un fort essor; mais ce n'est que tout récemment que l'importance de cette théorie pour les questions commutatives est devenue claire. Je voudrais parler aujourd'hui de cette importance du non-commutatif pour le commutatif et l'examiner en détail à propos de deux questions classiques remontant à Gauss: le théorème du genre principal et le théorème de la norme, qui lui est étroitement lié. Ces questions ont sans cesse changé de formulation au cours du temps: chez Gauss, elles apparaissent comme l'achèvement de sa théorie des formes quadratiques; puis elles jouent un rôle essentiel dans la caractérisation des corps de nombres relativement cycliques et abéliens par la théorie du corps de classes; enfin elles peuvent s'énoncer comme des théorèmes sur les automorphismes et sur la décomposition des algèbres, et cette dernière formulation donne en même temps un transfert des théorèmes aux corps de nombres galoisiens arbitraires sur leur corps de base.

Par cette esquisse, que je développerai plus loin, je voudrais en même temps expliquer le \emph{principe} de l'application du non-commutatif au commutatif: \emph{on cherche, au moyen de la théorie des algèbres, à obtenir, pour des faits connus sur les formes quadratiques ou les corps cycliques, des formulations invariantes et simples, c'est-à-dire des formulations qui ne dépendent que des propriétés structurales des algèbres. Une fois ces formulations invariantes démontrées -- et c'est le cas dans les exemples indiqués ci-dessus -- on obtient par là même le transfert de ces faits aux corps galoisiens arbitraires.}

\paragraph{2.}
Avant d'entrer dans le détail, je voudrais encore donner un aperçu général des différentes méthodes et des autres résultats. Il faut d'abord remarquer que la difficulté principale réside dans l'obtention de la formulation pour les corps galoisiens généraux, pour laquelle il n'existait aucun point d'appui sans la méthode hypercomplexe; dans les exemples cités, le passage sous-jacent au non-commutatif est obtenu par la \emph{considération simultanée du corps et du groupe} au moyen du ``produit croisé'' et de ses constantes de multiplication, les ``systèmes de facteurs'' (voir 3). On obtient ainsi une algèbre simple normale sur le corps de base, et toute telle algèbre peut essentiellement être engendrée ainsi. De tels produits croisés ont d'abord été considérés par Dickson,\footnote{Voir son livre \emph{Algebren und ihre Zahlentheorie}, Zürich 1927, \S 34.} tandis que la théorie des systèmes de facteurs, par Speiser, Schur et R. Brauer,\footnote{Voir par exemple R. Brauer, \emph{Untersuchungen über die arithmetischen Eigenschaften von Gruppen linearer Substitutionen}, et la littérature indiquée à la note 2 de cet article, Math. Z. 28 (1928).} s'était développée à partir d'un tout autre point de départ, celui de la question des représentations absolument irréductibles. Ce n'est que par la fusion des deux théories que l'on a pu obtenir une construction assez simple et assez étendue pour attaquer ainsi des questions commutatives.\footnote{J'ai développé cette construction pour la première fois dans un cours d'hiver 1929/30, reproduit au chap. 2 de H. Hasse, \emph{Theory of cyclic algebras}, Trans. Amer. Math. Soc. 34 (1932). \emph{Note éditoriale: l'édition allemande diplomatique ED0014 conserve la lecture imprimée \og{}Transac, 134\fg{}; la présente traduction donne la référence bibliographique contrôlée, vol.~34.} Un exposé de M. Deuring sur les nombres hypercomplexes et les applications arithmétiques, à paraître dans la collection \emph{Ergebnisse der Mathematik}, donne une vue d'ensemble du domaine traité dans cette conférence.}

Il en résulte en même temps, pour des faits connus, des preuves nouvelles et transparentes: je voudrais signaler ici une preuve hypercomplexe de la loi de réciprocité pour les corps cycliques, donnée par H. Hasse et devant paraître prochainement dans les Math. Ann.,\footnote{H. Hasse, \emph{Die Struktur der R. Brauerschen Algebrenklassengruppe über einem algebraischen Zahlkörper (insbesondere Normenrestsymbol und Reziprozitätsgesetz)}. Math. Ann. 107 (1932/33).} au moyen d'une formulation invariante de son symbole de reste normique, reposant sur la théorie du produit croisé. Puis une démonstration, par la méthode hypercomplexe, de la théorie locale du corps de classes, reposant sur la même base, a été donnée récemment par C. Chevalley; il fallut toutefois développer à cette fin de nouveaux théorèmes algébriques sur les systèmes de facteurs.\footnote{C. Chevalley, \emph{Sur la théorie du symbole de restes normiques}; paraîtra dans J. f. Math. 169.} Je dois toutefois formuler une réserve: la seule méthode des produits croisés ne donne vraisemblablement \emph{pas la théorie complète des corps de nombres galoisiens}. Cela résulte de nouveaux résultats encore non publiés d'Artin, qui se rattachent à la preuve de Hasse ci-dessus dans l'esprit du principe indiqué, mais qui ne donnent que des égalités numériques au lieu de théorèmes d'isomorphie complets.

Des méthodes qui donnent un isomorphisme complet -- il est vrai, un isomorphisme de modules -- existent déjà en algèbre. Il s'agit du prolongement d'idées dues à A. Speiser,\footnote{A. Speiser, \emph{Gruppendeterminante und Körperdiskriminante}. Math. Ann. 77 (1916).} plus précisément de la conception du corps galoisien comme ``module de Galois'', c'est-à-dire comme module sur le corps de base admettant comme opérateurs les substitutions du groupe galoisien. Et il existe un isomorphisme de modules entre le corps et l'anneau de groupe (algèbre de groupe), en ce sens qu'il y a une correspondance biunivoque entre les éléments telle que les formes linéaires sur le corps de base se correspondent, et qu'aux substitutions du groupe galoisien dans le corps est associée la multiplication dans l'anneau de groupe. Ce théorème, que j'avais formulé,\footnote{E. Noether, \emph{Normalbasis bei Körpern ohne höhere Verzweigung}, théorème 3, J. f. Math. 167 (1932) (la preuve contient une lacune).} a été démontré par M. Deuring,\footnote{M. Deuring, \emph{Galoissche Theorie und Darstellungstheorie}. Math. Ann. 107 (1932).} qui y a fondé une preuve de la théorie de Galois, où l'isomorphisme de modules réalise la correspondance entre groupe et corps. Des théorèmes de structure plus poussés -- également dus à Deuring -- sont parallèles à des faits formels des séries \(L\) d'Artin et donnent un accès structurel aux conducteurs d'Artin. Ces séries \(L\) et conducteurs d'Artin,\footnote{E. Artin, \emph{Über eine neue Art von \(L\)-Reihen}, Math. Sem. Hamburg 3 (1924). \emph{Zur Theorie der \(L\)-Reihen mit allgemeinen Gruppencharakteren}, ibid., 8 (1931). \emph{Die gruppentheoretische Struktur der Diskriminanten algebraischer Zahlkörper}, J. f. M. 164 (1931).} formés avec des caractères de groupes généraux, constituent, avec Speiser\textsuperscript{6)}, le premier lien entre théorie des nombres et théorie des représentations, une première avancée au-delà des corps abéliens. Ils ont donné une forte impulsion à tout le développement; en particulier, la théorie des modules de Galois s'est orientée d'après eux.

\paragraph{3.}
Je voudrais maintenant examiner en détail les problèmes annoncés d'emblée: le théorème de la norme et le théorème du genre principal. D'abord la \emph{définition du produit croisé}: soit \(K/k\) un corps galoisien de degré \(n\), \(\mathfrak{G}\) son groupe. Le produit croisé signifie un plongement simultané de \(K\) et de \(\mathfrak{G}\) dans une algèbre \(A\), de telle sorte que les automorphismes de \(K\) deviennent intérieurs. Soient \(u_{S_1},\ldots,u_{S_n}\) des symboles correspondant aux \(n\) éléments du groupe; on pose alors d'abord \(A\) comme module de formes linéaires de rang \(n\) sur \(K\):
\[
\begin{gathered}
(1)\qquad A=u_{S_1}K+\cdots+u_{S_n}K\\
\text{\normalfont(c'est-à-dire que \(A\) se compose de toutes les formes linéaires}\\
\text{\normalfont\(u_{S_1}a_1+\cdots+u_{S_n}a_n\), les \(a_i\) étant arbitraires dans \(K\)).}
\end{gathered}
\]
Par l'exigence de l'automorphisme intérieur -- engendré par les \(u_S\), plus généralement par \(u_SK^\ast\)\footnote{\(K^\ast\) s'obtient de \(K\) en en retranchant le zéro; cette notation est employée de façon générale.} -- \(A\) devient un anneau, donc une algèbre de rang \(n^2\) sur \(k\). Cette exigence s'exprime en effet par
\[
(2)\qquad u_S^{-1}zu_S=z^S,\qquad \text{ou}\qquad zu_S=u_Sz^S
\]
pour tout \(z\in K\),\footnote{\(z^S\) désigne comme d'habitude l'élément provenant de \(z\) par la substitution \(S\).} et par
\[
(3)\qquad u_Su_T=u_{ST}a_{S,T}\quad \text{avec }a_{S,T}\text{ dans }K^\ast
\]
\[
\begin{gathered}
(4)\qquad a_{ST,R}\,a_{S,T}^{\,R}=a_{S,TR}\,a_{T,R}\\
\text{\normalfont(loi associative provenant de \([u_Su_T]u_R=u_S[u_Tu_R]\)).}
\end{gathered}
\]
\(A\) s'appelle le produit croisé de \(K\) avec \(\mathfrak{G}\) pour le système de facteurs \(a_{S,T}\); on démontre que \(A\) devient une algèbre simple normale sur \(k\), donc un anneau de matrices de degré \(r\), \(D_r\), sur l'algèbre à division associée \(D\), et que \(K\) devient un sous-corps commutatif maximal, donc un corps de décomposition (c'est-à-dire que l'extension du domaine des coefficients \(k\) par un corps isomorphe à \(K\) donne une algèbre décomposée, une algèbre matricielle pleine sur le centre). Réciproquement, pour une algèbre à division donnée \(D\), il existe toujours des anneaux de matrices \(D_r\) qui peuvent être engendrés de la manière indiquée comme produits croisés.

Si l'on passe de \(u_S\) à \(v_S=u_Sc_S\), avec \(c_S\) dans \(K^\ast\), ce qui engendre le même automorphisme, il apparaît des systèmes de facteurs ``associés''.
\[
(5)\qquad \bar a_{S,T}=a_{S,T}\frac{c_S^{\,T}c_T}{c_{ST}}
\]
Les systèmes de facteurs associés sont regroupés en une classe \((a)\); de même, on regroupe toutes les algèbres semblables à \(A\), c'est-à-dire tous les \(D_r\) avec \(r=1,2,\ldots\), en une classe \(\mathfrak A\). \emph{Les classes \(\mathfrak A\) et \((a)\) se correspondent biunivoquement: les classes à corps de décomposition fixé \(K\) forment, pour le produit direct, un groupe abélien isomorphe au produit terme à terme des classes de systèmes de facteurs. L'élément neutre est la classe des algèbres décomposées et, du côté des systèmes de facteurs, le système de toutes les grandeurs de transformation}
\[
        \frac{c_S^{\,T}c_T}{c_{ST}}.
\]
Il s'agit du groupe des classes d'algèbres mis en évidence par R. Brauer.

\paragraph{4.}
Je vais maintenant, par \emph{spécialisation au cas des corps de décomposition cycliques}, établir le lien avec la \emph{notion de norme} et parvenir ainsi à la formulation du \emph{théorème de la norme} généralisé selon le principe exposé au début. Si \(Z\) est cyclique, \(S\) une substitution génératrice de son groupe -- l'algèbre correspondante est alors dite cyclique -- on peut faire correspondre aux puissances de \(S\) les puissances de \(u\), donc
\[
\begin{aligned}
(1')\qquad &A=Z+uZ+\cdots+u^{n-1}Z\\
(2')\qquad &zu=uz^S\\
(3')\qquad &u^n=a\\
(4')\qquad &a\text{ appartient au corps de base }k^\ast\\
(5')\qquad &\bar a=a\cdot N(c),\quad \text{si l'on pose }v=uc.
\end{aligned}
\]
Chaque système de facteurs consiste donc ici en un seul élément \(a\) situé dans le corps de base -- notation \(A=(\alpha,Z)\) --; la classe neutre des systèmes de facteurs est donnée par les normes de \(Z^\ast\), et le groupe des classes d'algèbres devient isomorphe au groupe quotient
\[
        k^\ast/N(Z^\ast).
\]
Une algèbre cyclique \((\alpha,Z)\) est donc décomposée si et seulement si \(a\) est la norme d'un élément de \(Z\).

Ce lien entre norme et décomposition donne la formulation du ``théorème de la norme'', à savoir le \emph{théorème sur les algèbres décomposées: si une algèbre est décomposée à chaque place, elle est décomposée sur le corps de base.} Ici, la ``place'' est définie, comme d'habitude en théorie des nombres, en remplaçant le corps de base \(k\) par son extension \(\mathfrak p\)-adique \(k_{\mathfrak p}\), où \(\mathfrak p\) est un idéal premier de \(k\) (respectivement, pour les places infinies, en étendant \(k\) et ses conjugués par le corps des nombres réels).

En effet, le théorème de la norme pour les corps cycliques y est contenu. Car, pour les algèbres cycliques \((\alpha,Z)\), le théorème signifie, d'après ce qui vient d'être dit: si \(a\) est une norme \(\mathfrak p\)-adique en toute place (finie ou infinie), alors \(a\) est la norme d'un nombre de \(Z\), ou, sans passer au \(\mathfrak p\)-adique: \emph{si \(a\) est un résidu normique suivant tout idéal premier \(\mathfrak p\) de \(k\) (et satisfait certaines conditions de signe), alors \(a\) est la norme d'un nombre de \(Z\).} Cette dernière forme est le théorème de la norme démontré dans la théorie du corps de classes à l'aide des moyens analytiques connus. Et la preuve du théorème général sur les algèbres décomposées s'obtient à partir de ce cas spécial cyclique par des considérations purement algébriques-arithmétiques.\footnote{R. Brauer, H. Hasse, E. Noether, \emph{Beweis eines Hauptsatzes in der Theorie der Algebren}. J. f. M. 167 (1932).} Une première conséquence importante en a été tirée par Hasse: \emph{toute algèbre simple normale sur un corps de nombres algébrique est cyclique.} C'est dans la recherche d'une preuve de ce fait longtemps conjecturé que naquit la formulation générale.

\paragraph{5.}
Une deuxième conséquence du théorème sur les algèbres décomposées -- de nouveau par des raisonnements purement algébriques-arithmétiques -- est le théorème du genre principal placé en tête.\footnote{La preuve doit paraître dans les Math. Ann.} Sa formulation invariante repose sur le fait que les relations (2) à (5) définissant le produit croisé sont purement multiplicatives, donc gardent un sens si \(K^\ast\) est remplacé par un groupe abélien \(\mathfrak{J}\), soumis seulement à la condition que son groupe d'automorphismes contienne un sous-groupe isomorphe à \(\mathfrak{G}\). À la place de (1) apparaît alors l'``extension de \(\mathfrak{G}\) par \(\mathfrak{J}\)'' au sens de la théorie des groupes. Si l'on prend pour \(\mathfrak{J}\) le groupe de tous les idéaux de \(K\), les systèmes de facteurs deviennent donc des systèmes d'idéaux; une partition en classes dans \(\mathfrak{J}\) induit une partition en classes pour les systèmes de facteurs, et celle-ci sera en général une partition en classes d'idéaux plus fine que l'initiale. Car l'exigence que la multiplication des \(u_S\) avec des classes d'idéaux (absolues) soit univoque dit seulement que les grandeurs de transformation \(c_S^{\,T}c_T/c_{ST}\) -- la classe neutre des systèmes de facteurs d'éléments -- appartiennent à la classe neutre des systèmes de facteurs d'idéaux. En fait, comme le montre la spécialisation aux cas connus, une partition un peu moins fine suffit déjà. Je définis: \emph{dans la classe neutre des systèmes de facteurs, on range les éléments \(a_{S,T}\) qui, à toutes les places ramifiées (finies et infinies) de \(K\), engendrent des algèbres décomposées.}

L'extension de \(\mathfrak{G}\) qui en résulte sera notée \(\mathfrak{G}^\ast\); \((\mathfrak c)\) désignera la classe idéale absolue de \(\mathfrak c\). Alors:

\emph{Formulation invariante du théorème du genre principal: si, pour la partition en classes ainsi définie, la substitution \(v_S=u_S(\mathfrak c_S)\) donne un automorphisme de \(\mathfrak{G}^\ast\) -- tous ces \((\mathfrak c_S)\) formant le genre principal --, alors l'automorphisme est intérieur et est engendré par une classe idéale \((\mathfrak b)\).}

Les cas particuliers connus résultent d'une formulation équivalente, un peu plus explicite: \emph{si les grandeurs de transformation \((\mathfrak c_S^{\,T})(\mathfrak c_T)/(\mathfrak c_{ST})\), formées à partir des classes d'idéaux \((\mathfrak c_S)\), appartiennent à la classe neutre des systèmes de facteurs d'idéaux, alors il existe une classe idéale \((\mathfrak b)\) telle que les \((\mathfrak c_S)\) deviennent des puissances symboliques d'exposant \((1-S)\): \((\mathfrak c_S)=(\mathfrak b)/(\mathfrak b^{S})\) pour tout \(S\) de \(\mathfrak{G}\).} En effet, l'hypothèse exprime exactement la propriété d'automorphisme; le fait que celui-ci soit intérieur s'exprime par \(v_S=(\mathfrak b)^{-1}u_S(\mathfrak b)=u_S(\mathfrak b)^{1-S}=u_S(\mathfrak c_S)\).

La spécialisation au cas cyclique donne donc (en tenant compte de la normalisation): si \(N([\mathfrak c])\) appartient à la classe neutre des systèmes de facteurs d'idéaux, alors \((\mathfrak c)\) devient une puissance symbolique d'exposant \((1-S)\): \((\mathfrak c)=(\mathfrak b)^{1-S}\).

\paragraph{6.}
Pour parvenir de là à la forme connue pour les corps cycliques et les formes quadratiques, je remarque d'abord que le théorème est énoncé pour l'ensemble complet des classes d'idéaux, mais que son contenu reste le même lorsque, comme d'habitude, on se restreint aux idéaux premiers aux places ramifiées.

Alors le genre principal défini ici devient, pour les corps quadratiques, celui de Gauss; car aux classes d'idéaux correspondent les formes quadratiques, et aux normes des classes les nombres représentables par les formes. Dire que la classe neutre engendre aux places ramifiées de \(K\) des algèbres décomposées signifie donc que ces nombres représentables deviennent des résidus quadratiques aux places ramifiées; les formes correspondantes possèdent ainsi le caractère total de la forme principale, et forment donc le genre principal de Gauss. Il est connu que la puissance symbolique d'exposant \((1-S)\) devient la duplication.

Pour les corps cycliques, la formulation devient la suivante: le genre principal se compose de toutes les classes d'idéaux dont les normes deviennent des résidus normiques aux places ramifiées (finies et infinies). Mais cela est, comme on sait, équivalent à ``résidu normique suivant le conducteur'', et ainsi le théorème usuel apparaît comme spécialisation.

Du reste, dans le cas général de corps galoisiens quelconques, on peut aussi introduire un conducteur composé uniquement des places ramifiées, de telle sorte que, après normalisation des systèmes de facteurs pour chacune de ces places, le rayon ne contienne que des éléments de la classe neutre.

Ici apparaît la question du rapport avec les conducteurs d'Artin mentionnés dans l'aperçu 2, qui sont composés des mêmes idéaux premiers; et ainsi la question du rapport avec la théorie des modules de Galois, la seconde méthode hypercomplexe. Jusqu'où ces deux méthodes pourront porter, seul l'avenir le montrera.

\setcounter{footnote}{0}
\clearpage
\providecommand{\frA}{\mathfrak A}
\providecommand{\frB}{\mathfrak B}
\providecommand{\frC}{\mathfrak C}
\providecommand{\frR}{\mathfrak R}
\providecommand{\frS}{\mathfrak S}
\providecommand{\frT}{\mathfrak T}
\providecommand{\frM}{\mathfrak M}
\providecommand{\frG}{\mathfrak G}
\providecommand{\frako}{\mathfrak o}

\section*{\NoetherNavAnchor{P40}{40. Algèbre non commutative}40. Algèbre non commutative}

\begin{center}
\emph{Math. Zs. 37 (1933), p. 514--541}
\end{center}

\begin{center}
\textbf{Algèbre non commutative.}\\[0.5em]
Par\\[0.25em]
Emmy Noether à Goettingen.
\end{center}

Les théorèmes principaux de l'algèbre commutative sont, comme on sait, contenus dans la théorie de Galois; celle-ci est précédée par la théorie des corps de rupture et des corps de décomposition, c'est-à-dire des corps qui suffisent pour qu'un polynôme donné fasse apparaître un facteur linéaire, respectivement se décompose entièrement en facteurs linéaires.

Je développe ici les parties correspondantes de l'algèbre dans le non-commutatif, en particulier dans l'hypercomplexe. Et je travaille alors par principe avec des méthodes non commutatives, avec la représentation dans des corps gauches. À la fin je montre comment les théorèmes mentionnés ci-dessus de l'algèbre commutative peuvent se fonder d'une manière tout à fait parallèle par représentation dans des corps commutatifs.

La théorie des représentations sur laquelle repose ce travail --- précédée d'une brève théorie des endomorphismes (\S 1) --- prolonge celle qui est fondée sur les modules de représentation.\footnote{E. Noether, Hyperkomplexe Grossen und Darstellungstheorie, Math. Zeitschr. 30 (1929), p. 641--692, cité ci-dessous comme Darstellungstheorie. Voir la reproduction chez van der Waerden, Moderne Algebra II.} En particulier, je considère, à côté de la représentation directe, la représentation réciproque --- chacune pouvant se ramener à l'autre ---, désormais produite au moyen du module de représentation réciproque. L'avantage est que ce module --- qui est donc un bimodule --- peut aussi être regardé comme un module unilatéral sur un anneau d'extension (\S 2). La représentation dans des corps gauches se ramène ainsi à la théorie des idéaux de l'anneau d'extension; les classes de représentations irréductibles correspondent aux classes d'idéaux irréductibles de cet anneau, en exacte analogie avec les faits qui valent pour le système lui-même lorsqu'on représente un système hypercomplexe sur son corps commutatif de coefficients (\S 3 et 4).

On en déduit les théorèmes structuraux sur les anneaux de matrices sur des corps gauches, grâce à la remarque (\S 5) selon laquelle \emph{tout sous-anneau fournit une représentation par ce corps gauche, ou une représentation réciproque par le corps gauche anti-isomorphe}. Ce corps gauche anti-isomorphe constitue un premier analogue du corps de scission minimal et du corps de décomposition dans le cas commutatif: il réalise toutes les représentations réciproques de degré 1 et, lorsque le rang sur le centre --- qui n'avait pas été supposé fini jusque-là --- est fini, une décomposition complète en facteurs directs de rang 1. On obtient ainsi la théorie de Galois pour les corps gauches (\S 6), ainsi que le fait (\S 7) que des algèbres à division anti-isomorphes (corps gauches de rang fini sur leur centre) déterminent des classes inverses dans le groupe de R. Brauer des classes d'algèbres.

Une seconde démonstration de la théorie de Galois, valable plus généralement pour les systèmes simples (et présentée plus tôt dans le texte), découle presque immédiatement d'un théorème sur les sous-anneaux commutant entre eux, lui-même rattaché presque directement à l'étude préliminaire des endomorphismes.

Jusqu'ici, les considérations sont entièrement non commutatives. La question des corps de décomposition \emph{commutatifs} est néanmoins elle aussi traitée par des méthodes non commutatives, au moyen de la représentation du corps de décomposition par l'algèbre à division, conformément à la remarque préparée au \S 5 (\S 7). L'article se conclut par le passage au cas commutatif mentionné plus haut (\S 8) et par la théorie des corps de décomposition des systèmes arbitraires (\S 9); on obtient en même temps le lien avec la théorie usuelle des représentations sur un corps commutatif.

R. Brauer a fondé la théorie des corps de décomposition sur cette théorie de la représentation dans le commutatif --- et sur les systèmes de facteurs ``irrationnels''.\footnote{Voir notre note commune: Uber minimale Zerfallungskorper irreduzibler Darstellungen. Sitz. Ber. d. Preuss. Ak. d. Wiss. 1927, p. 221--228. (La remarque de la p. 222 contient un ``rapport justificatif''. Les théorèmes principaux sont nés indépendamment et presque simultanément.) Voir aussi R. Brauer, Uber Systeme hyperkomplexer Zahlen, \S 3. Math. Zeitschr. 30 (1929), p. 79--107. A. A. Albert a retrouvé plus tard ces théorèmes de façon indépendante.} À cause de cette fondation commutative il doit --- tout comme Albert plus tard --- supposer que le centre est un corps parfait, restriction inutile d'après la fondation non commutative. Avec la même restriction, et également avec la théorie de la représentation dans le commutatif, R. Brauer et K. Shoda ont --- après avoir pris connaissance de ma théorie de Galois pour les corps gauches --- développé davantage la théorie: R. Brauer donna le théorème mentionné ci-dessus sur les sous-anneaux permutables, K. Shoda indépendamment la théorie complète aussi pour les systèmes semi-simples.\footnote{R. Brauer, Uber die algebraische Struktur von Schiefkorpern, \S 2. Journ. f. Math. 166 (1932), p. 241--252. K. Shoda, Uber die galoissche Theorie der halbeinfachen hyperkomplexen Systeme. Math. Ann. 107 (1932), p. 252--258. Ces résultats avaient aussi été développés par J. Levitzki.}

Enfin il faut mentionner que, pour le cas des corps gauches, on trouve une brève présentation chez van der Waerden II (\S 128), à la suite de mon cours de l'été 1928, où j'avais parlé pour la première fois de ces questions. La présentation de van der Waerden apporte une série de simplifications par rapport à ce cours; je les ai reprises et prolongées en partie dans un second cours (hiver 1929/30),\footnote{Le premier cours a été élaboré par G. Kothe, le second par M. Deuring.} en partie seulement ici. En particulier, le transfert de la méthode de preuve hypercomplexe de la théorie de Galois au commutatif (\S 8) ne date que du second cours, où la preuve non commutative (\S 6, 3.) avait été essentiellement simplifiée par rapport au premier cours. Le théorème sur les sous-anneaux permutables (\S 5) et la seconde méthode de preuve de la théorie de Galois non commutative qui s'appuie sur lui (\S 6, 1. et 2.) n'ont été ajoutés qu'après le second cours, après connaissance de R. Brauer (note 3) et de van der Waerden \S 128; dans \S 5, 3., les hypothèses de finitude sont toutefois beaucoup plus faibles que là.

Certains raisonnements du premier cours --- méthode de formation d'intersections d'idéaux de différences --- ont, bien que plus compliqués, l'avantage de se transposer aux systèmes de rang infini et aux systèmes entiers.\footnote{La méthode est reproduite pour l'essentiel chez G. Kothe, Schiefkorper unendlichen Ranges uber dem Zentrum, \S 5, Math. Ann. 105 (1931), p. 15--39. Voir aussi la note 10. Cette méthode d'intersection est maintenant remplacée par le fait presque trivial (\S 4, 1.) --- remarqué pour la première fois par van der Waerden --- que les idéaux bilatères appartiennent aux modules invariants. Tout peut ainsi se fonder par la décomposition plus simple en sommes directes, ce qui échoue cependant dans le cas du rang infini et dans le cas entier.} Pour les systèmes commutatifs entiers, on parvient ainsi au lien entre différentiation idéale et différente,\footnote{Voir la conférence de Prague, Jahresber. d. Deutsch. Math. Ver. 39 (1930), p. 17 (pagination oblique).} sur lequel je reviendrai à l'occasion plus exactement.

La théorie des corps de décomposition trouve son application principale dans la théorie des produits croisés et de leurs systèmes de facteurs, qui constituent à leur tour le fondement d'applications arithmétiques; il ne sera plus question ici de cela.\footnote{La théorie des produits croisés est développée dans le second cours; elle est reproduite --- avec de petites modifications, pour ne pas devoir supposer la théorie donnée ici --- chez H. Hasse: Theory of cyclic algebras over an algebraic numberfield, chap. II. Transactions of the Amer. Math. Soc. 34 (1932), p. 171--214. Une présentation plus étroitement rattachée au cours paraîtra dans un rapport de M. Deuring dans les ``Ergebnisse der Mathematik''.}

\subsection*{\S 1. Endomorphismes, modules et bimodules}

La théorie de la représentation fondée sur les modules de représentation repose sur la théorie de l'anneau des endomorphismes des groupes abéliens, avec ou sans opérateurs. Les raisonnements implicitement sous-jacents, donc les relations entre application et lois de calcul, seront formulés dans quelques propositions simples afin d'éviter les répétitions (voir aussi \S 2).

\paragraph{1. Application multiplicative. Loi associative.}
Soit d'abord \(\frG\) un groupe sans opérateurs, et soit \(\frA\) son domaine absolu des endomorphismes, c'est-à-dire le système de tous les homomorphismes de \(\frG\) dans lui-même. \(\frA\) est fermé multiplicativement, puisque le produit \( \sigma\tau \) est défini par
\[
        g(\sigma\tau)=(g\sigma)\tau,
        \qquad g\in\frG,\quad \sigma,\tau\in\frA                         \tag{1}
\]
et satisfait évidemment à la loi associative.

\(\frG\) devient, comme on sait, un groupe avec opérateurs lorsqu'est donnée une famille \(\frB\) de symboles \(\Theta,H,\ldots\), de telle sorte que les combinaisons \(g\Theta,gH,\ldots\) fournissent des éléments de \(\frG\) définis de façon univoque et engendrent des endomorphismes (homomorphismes de \(\frG\) dans lui-même):
\[
        (g\cdot h)\Theta=g\Theta\cdot h\Theta .
\]
Il existe donc une application univoque (en général non biunivoque) de \(\frB\) sur un sous-ensemble \(\overline{\frB}\) du domaine absolu des endomorphismes \(\frA\). Si le domaine d'opérateurs \(\frB\) est fermé multiplicativement, l'application univoque de \(\frB\) sur la partie correspondante du domaine des endomorphismes est multiplicativement homomorphe si et seulement si la relation associative
\[
        g(\Theta H)=(g\Theta)H, \qquad g\in\frG,\ \Theta,H\in\frB,                       \tag{1a}
\]
est satisfaite. De façon correspondante, l'antihomomorphisme est défini lorsqu'il s'agit d'opérateurs à gauche.

En effet, l'application de \(\frB\) sur \(\overline{\frB}\) est donnée par \(g\Theta=g\sigma\) pour tout \(g\) de \(\frG\), donc aussi par
\[
        (g\Theta)H=(g\sigma)\tau=g(\sigma\tau),
\]
de sorte que (1a) devient nécessaire et suffisante pour l'homomorphie multiplicative. Le fait que les opérateurs à gauche engendrent un antihomomorphisme vient de ce que les endomorphismes écrits à gauche \(\sigma^\ast g,\tau^\ast\sigma^\ast g\) doivent se lire de droite à gauche: \(\tau^\ast\sigma^\ast g=g\sigma\tau\), tandis que les opérateurs se lisent toujours de gauche à droite.

\paragraph{2. Application homomorphe relativement aux opérateurs.}
Si \(\frG\) est un groupe avec opérateurs, l'homomorphie relativement aux opérateurs est définie par
\[
        (g\Theta)\sigma=(g\sigma)\Theta,                                          \tag{2}
\]
respectivement, pour les opérateurs à gauche, par
\[
        (\Theta g)\sigma=\Theta(g\sigma).                                          \tag{2*}
\]
c'est-à-dire par la commutation des deux actions ou, dans le cas des opérateurs à gauche, par la loi d'associativité mixte. En appliquant, comme en 1., le raisonnement qui identifie un domaine d'opérateurs à son image dans le domaine des endomorphismes, on obtient pour deux domaines d'opérateurs distincts:

Si deux domaines d'opérateurs \(\frB\) et \(\frC\) sont donnés, alors les éléments de \(\frC\) engendrent des \(\frB\)-endomorphismes, et simultanément ceux de \(\frB\) des \(\frC\)-endomorphismes, si et seulement si les actions des deux domaines commutent sur \(\frG\) ou, respectivement, si la loi d'associativité mixte est satisfaite:
\[
        (g\Theta)\bar H=(g\bar H)\Theta,                                                    \tag{2a}
\]
respectivement
\[
        (\Theta g)\bar H=\Theta(g\bar H).                               \tag{2a*}
\]

\paragraph{3. Modules et bimodules sur des anneaux.}
Un module à droite \(\frM\) sur un anneau \(\frR\) est un groupe abélien additif avec les éléments de \(\frR\) comme opérateurs à droite; outre la relation associative (1a), les relations distributives y sont satisfaites:
\[
        (g+h)\rho=g\rho+h\rho,\qquad
        g(\rho+\sigma)=g\rho+g\sigma.                                  \tag{3a}
\]
Il en va de même pour les modules à gauche.

Parmi les bimodules il faut distinguer deux espèces. Les modules à droite sur deux anneaux \(\frR\) et \(\frS\) s'appellent bimodules lorsque, outre les combinaisons valables séparément pour \(\frR\) et \(\frS\), on a encore
\[
        (m\rho)\sigma=(m\sigma)\rho
\]
Un groupe qui est un \(\frR\)-module à gauche et un \(\frS\)-module à droite est appelé bimodule lorsque s'ajoute la loi d'associativité mixte
\[
        \rho(m\sigma)=(\rho m)\sigma .
\]

La signification de ces combinaisons vient du fait que, dans le cas des groupes abéliens, le domaine absolu des endomorphismes forme un anneau, l'anneau absolu d'endomorphismes\footnote{Pour les groupes non abéliens il s'agit d'un anneau ``généralisé''; voir à ce sujet un travail de H. Fitting à paraître dans les Math. Annalen. [Math. Annalen 107, p. 514--542.]}. Dans le détail, le raisonnement par application donne:

Si le domaine d'opérateurs \(\frR\) d'un groupe abélien (écrit additivement) \(\frM\) est un anneau, l'application canonique de \(\frR\) dans l'anneau absolu d'endomorphismes, d'image \(\overline{\frR}\), est un homomorphisme d'anneaux si et seulement si \(\frM\) est un module à droite sur \(\frR\) -- donc si (1a) et (3a) sont satisfaites --, et un antihomomorphisme d'anneaux pour les modules à gauche.

Si \(\frM\) est un module à droite sur les anneaux \(\frR\) et \(\frS\), alors \(\frM\) est un bimodule si et seulement si l'application de \(\frR\) dans l'anneau des \(\frS\)-endomorphismes définie par son action est un homomorphisme d'anneaux et, simultanément, si l'application correspondante de \(\frS\) dans l'anneau des \(\frR\)-endomorphismes en est un. Si \(\frM\) est un module à gauche sur \(\frR\) et un module à droite sur \(\frS\), la condition de bimodule (2a*) signifie que l'application de \(\frS\) dans l'anneau des \(\frR\)-endomorphismes est un homomorphisme d'anneaux, tandis que celle de \(\frR\) dans l'anneau des \(\frS\)-endomorphismes est un antihomomorphisme d'anneaux.

Il en résulte en particulier:

1) Si \(\frR\) est un anneau unitaire, alors \(\frR\) est isomorphe à son anneau des endomorphismes comme \(\frR\)-module à gauche et anti-isomorphe à son anneau des endomorphismes comme \(R\)-module à droite. En effet, la loi d'associativité mixte (2a*) est satisfaite, donc l'assertion d'homomorphie s'applique. L'existence de l'élément unité donne tout endomorphisme sous la forme \(e\to a\), avec \(a\in\frR\); on a donc un isomorphisme, et \(\frR\) épuise l'anneau des \(\frR\)-endomorphismes.

2) Si \(\frR\) est anneau complet de matrices sur un corps gauche \(\mathsf{A}\), en général non commutatif,
\[
        \frR=\sum c_{ik}\mathsf{A}=\sum \mathsf{A} c_{ik},
\]
alors \(\mathsf{A}\) est anti-isomorphe au corps gauche des endomorphismes des idéaux à droite simples, et isomorphe au corps gauche des endomorphismes des idéaux à gauche simples. En effet, tout idéal à droite simple est isomorphe, comme module à opérateurs, à un \(\mathsf{A}\)-module à gauche et \(\frR\)-module à droite, et \(\mathsf{A}\) épuise les \(\frR\)-endomorphismes; de même pour les idéaux à gauche.

\paragraph{4. Passage d'un bimodule à droite à un module unilatéral sur un anneau-produit.}
L'intérêt des bimodules à droite repose essentiellement sur ce passage.

\paragraph{Théorème.}
Si \(\frM\) est un module à droite sur un anneau \(\frT\) contenant deux sous-anneaux \(\frR\) et \(\frS\) qui commutent élément par élément, alors \(\frM\) peut aussi être regardé comme un bimodule sur \(\frR\) et \(\frS\). Réciproquement, si l'on se donne un bimodule sur \(\frR\) et \(\frS\) et qu'il existe un anneau-produit \(\frT\) dans lequel \(\frR\) et \(\frS\) commutent élément par élément, alors \(\frM\) peut aussi être regardé comme un module à droite sur \(\frT\), l'action de \(\frT\) prolongeant les actions données de \(\frR\) et de \(\frS\).

La première assertion vient de ce que, par définition, l'action de \(\frT\) définit un homomorphisme de \(\frT\) sur un sous-anneau \(\overline{\frT}\) de l'anneau absolu d'endomorphismes, tandis que les images \(\overline{\frR}\) et \(\overline{\frS}\) commutent élément par élément. C'est précisément la relation (2a) qui caractérise le bimodule.

Réciproquement, si \(\frM\) est donné comme \(\frR,\frS\)-bimodule (module à droite), les actions de \(\frR\) et de \(\frS\) définissent des homomorphismes vers des sous-anneaux \(\overline{\frR},\overline{\frS}\) de l'anneau absolu d'endomorphismes, dont les éléments commutent deux à deux. Dans l'anneau absolu d'endomorphismes existe, pour \(\overline{\frR},\overline{\frS}\), l'anneau-produit\footnote{Anneau-produit signifie donc seulement: l'anneau engendré par deux anneaux ayant un sur-anneau commun; le produit n'a pas besoin d'être direct. -- Le fait qu'un anneau-produit n'ait pas toujours à exister pour une intersection donnée est montré par l'exemple suivant: soit \(\frako\) l'anneau des entiers; posons \(\frR=\frako[x]\), \(\frS=\frako[y]\) avec \(2x-1=0\), \(2y-1=0\), de sorte que \(\frako\) soit l'intersection de \(\frR\) et \(\frS\), si aucune combinaison n'est fixée entre \(x\) et \(y\). Mais si \(\frR\) et \(\frS\) se trouvent dans un sur-anneau commun, alors on obtient: \((2x-1)y-x(2y-1)=0\), donc \(x=y\). Il n'existe donc pas d'anneau-produit ayant pour intersection \(\frako\).} \(\overline{\frT}\), lequel, en raison de la permutabilité élément par élément, consiste en tous les éléments de la forme
\[
        \sum \bar r_i\bar s_i+\bar r+\bar s
\]
(si \(\overline{\frR},\overline{\frS}\) possèdent des éléments unités, les termes supplémentaires \(\bar r,\bar s\) disparaissent). S'il existe maintenant aussi un anneau-produit \(\frT\), dans lequel \(\frR\) et \(\frS\) sont permutables élément par élément, celui-ci consiste de façon correspondante en tous les éléments de la forme \(\sum r_i s_i+r+s\). Les homomorphismes
\[
        \frR\to\overline{\frR},\qquad \frS\to\overline{\frS}
\]
se prolongent donc, par la correspondance
\[
        \sum r_i s_i+r+s
        \to
        \sum \bar r_i\bar s_i+\bar r+\bar s
\]
en un homomorphisme \(\frT\to\overline{\frT}\); ainsi la combinaison
\[
        m\Bigl(\sum_i r_i s_i+r+s\Bigr)
        =\sum_i (m r_i)s_i+m r+m s
\]
est univoque et définit, d'après \S 1, 3., un \(\frT\)-module qui englobe le \(\frR,\frS\)-module donné.

\clearpage
\providecommand{\frA}{\mathfrak A}
\providecommand{\frB}{\mathfrak B}
\providecommand{\frR}{\mathfrak R}
\providecommand{\frS}{\mathfrak S}
\providecommand{\frT}{\mathfrak T}
\providecommand{\frM}{\mathfrak M}
\providecommand{\frN}{\mathfrak N}
\providecommand{\mA}{\mathfrak A}
\providecommand{\mbarA}{\overline{\mathfrak A}}

\subsection*{\S 2. Représentation réciproque et directe}

\paragraph{1. Modules de formes linéaires.}
Le passage de la théorie développée au \S 1 à la théorie de la représentation repose sur la spécialisation des modules qui y figurent en modules de formes linéaires et sur la considération de leurs anneaux d'endomorphismes.

Un \(\frS\)-module à droite \(\frM\) -- où \(\frS\) est un anneau avec élément unité -- s'appelle module de formes linéaires dans \(\frS\) si \(\frM\) est somme directe de \(n\) \(\frS\)-modules monogènes:
\[
        \frM=m_1\frS+\ldots+m_n\frS,
\]
de telle sorte que \(m_i\frS\) soit isomorphe à \(\frS\) comme module à opérateurs, donc que \(m_i s=0\) entraîne toujours \(s=0\). Il en résulte que l'action de \(\frS\) définit un isomorphisme de \(\frS\) sur un sous-anneau de l'anneau absolu d'endomorphismes.

\paragraph{Théorème 1.}
Si \(\frM\) est un module de formes linéaires dans \(\frS\), alors l'anneau des \(\frS\)-endomorphismes \(\mA\) de \(\frM\) devient anti-isomorphe -- et non pas seulement antihomomorphe -- à l'anneau \(\mbarA\) de toutes les matrices à \(n\) lignes dans \(\frS\).

Un \(\frS\)-endomorphisme arbitraire \(\alpha\) de \(\frM\) est complètement décrit par la correspondance
\[
        m_i\to \overline m_i=m_i\alpha;
\]
car il s'ensuit, par définition,
\[
        \sum m_i s_i\to
        \sum \overline m_i s_i
        =\sum (m_i\alpha)s_i
        =\sum (m_i s_i)\alpha.
\]
Si maintenant, en écriture matricielle,
\[
        (m_1\alpha,\ldots,m_n\alpha)=(m_1,\ldots,m_n)A,
\]
alors -- lorsque \(\alpha\) parcourt tous les endomorphismes -- la correspondance \(\alpha\to A\) donne une relation biunivoque entre \(\mA\) et \(\mbarA\). En effet, puisque \(\frM\) a été supposé module de formes linéaires, \(A\) est déterminé univoquement par \(\alpha\), et toute matrice \(A\) du \(n\)-ième degré dans \(\frS\) engendre un endomorphisme; il s'ensuit en outre:
\[
        \alpha+\beta\to A+B,\qquad
        \alpha\beta\to BA.
\]
Ce dernier point d'après
\[
        (m_1\alpha\beta,\ldots,m_n\alpha\beta)
        =(m_1,\ldots,m_n)A\beta
        =(m_1,\ldots,m_n)\beta A
        =(m_1,\ldots,m_n)BA
\]
à cause de \(ms\beta=m\beta s\).

\paragraph{2. Représentation et module de représentation.}
Si l'on entend par représentation réciproque, respectivement directe, de \(n\)-ième degré de l'anneau \(\frR\) dans \(\frS\) un antihomomorphisme, respectivement un homomorphisme, d'anneaux de \(\frR\) dans un sous-anneau de l'anneau de toutes les matrices à \(n\) lignes dans \(\frS\), alors le théorème 1 admet aussi la forme suivante:

\paragraph{Théorème 1'.}
L'anneau des endomorphismes \(\mA\) d'un module de formes linéaires de \(n\)-ième degré dans \(\frS\) admet une représentation réciproque isomorphe (fidèle) par l'anneau complet de matrices \(\mbarA\) de toutes les matrices à \(n\) lignes dans \(\frS\).

De là résultent, pour les représentations réciproques, les théorèmes usuels pour les représentations directes.

\paragraph{Définition.}
Un module de formes linéaires \(\frM\) dans \(\frS\) s'appelle module de représentation réciproque de \(\frR\) dans \(\frS\) si \(\frM\) est bimodule sur \(\frR\) et \(\frS\) et, en même temps, module à droite sur \(\frR\) et sur \(\frS\). Il s'appelle module de représentation directe si \(\frM\) est bimodule et, en même temps, \(\frR\)-module à gauche et \(\frS\)-module à droite.

\paragraph{Théorème 2.}
Tout module de représentation réciproque, respectivement directe, engendre une classe de représentations réciproques, respectivement directes, équivalentes de \(\frR\) dans \(\frS\), et toutes les représentations réciproques, respectivement directes, sont engendrées ainsi.

Soit d'abord \(\frM\) un module de représentation réciproque. D'après \S 1, 3., l'action de \(\frR\) définit alors un homomorphisme (et non un antihomomorphisme) de \(\frR\) sur un sous-anneau \(\overline{\frR}\) de l'anneau des \(\frS\)-endomorphismes de \(\frM\); d'après \S 2, théorème 1', \(\overline{\frR}\) admet une représentation réciproque (fidèle), qui devient donc aussi représentation réciproque de \(\frR\). Réciproquement, si une représentation réciproque \(\frR\to\frR^*\) est donnée, sa composition avec l'anti-isomorphisme de \(\frR^*\) à un sous-anneau \(\overline{\frR}\) de l'anneau des \(\frS\)-endomorphismes donne un homomorphisme de \(\frR\) sur \(\overline{\frR}\); ainsi \(\frM\) devient un \(\frR\)-module, et même un bimodule, donc un module de représentation d'après \S 1, 3. Le passage à d'autres bases de \(\frS\) donne la classe des représentations équivalentes.

S'il s'agit d'un module de représentation directe, alors l'action de \(\frR\) définit un antihomomorphisme sur \(\overline{\frR}\). La composition avec la représentation réciproque de \(\overline{\frR}\) donne donc une représentation directe de \(\frR\). Inversement, d'une représentation directe résulte un antihomomorphisme sur \(\overline{\frR}\); ensuite \(\frM\) devient un module de représentation directe d'après \S 1, 3.

\paragraph{Remarque.}
Naturellement les deux sortes de représentations peuvent se ramener l'une à l'autre: une représentation directe de \(\frR\) est une représentation réciproque de l'anneau opposé à \(\frR\), et inversement. Une autre réduction consiste à passer de \(\frS\) à un anneau opposé et à remplacer les matrices par les transposées. Cela correspond au passage d'un \(\frR\)-module à gauche et \(\frS\)-module à droite à un module à gauche sur \(\frR\) et sur \(\frS\).

\paragraph{3. Passage du module de représentation réciproque au module sur un anneau d'extension.}
L'avantage du module de représentation réciproque repose avant tout sur le passage à l'anneau d'extension, possible d'après \S 1, 4. Pour le cas spécial présent -- \(\frM\) module de formes linéaires sur \(\frS\) -- il s'agit donc du théorème de passage pour les modules de représentation:

\emph{Si \(\frM\) est un module à droite sur un anneau \(\frT\), qui contient deux sous-anneaux \(\frR\) et \(\frS\) permutables élément par élément de telle sorte que \(\frM\) devienne module de formes linéaires sur \(\frS\), alors \(\frM\) peut aussi être conçu comme module de représentation réciproque de \(\frR\) dans \(\frS\).}

\emph{Réciproquement, si un module de représentation réciproque de \(\frR\) dans \(\frS\) est donné, et s'il existe un anneau-produit \(\frT\) dans lequel \(\frR\) et \(\frS\) sont permutables élément par élément, alors \(\frM\) peut aussi être conçu comme \(\frT\)-module, de telle sorte que la combinaison avec \(\frT\) soit prolongement des combinaisons données avec \(\frR\) et \(\frS\).}

Dans la théorie de la représentation on utilise essentiellement la

\paragraph{Conséquence du théorème de passage pour les modules de représentation.}
Si \(\frM\) est un module de représentation réciproque de \(\frR\) dans \(\frS\) et en même temps un \(\frT\)-module, avec \(\frT\) comme anneau-produit, alors l'isomorphisme de \(\frM\), comme \(\frR,\frS\)-module, avec un module \(\frN\) équivaut à l'isomorphisme comme \(\frT\)-module. \emph{À chaque classe de \(\frT\)-modules isomorphes correspond donc une classe de représentations réciproques de \(\frR\) dans \(\frS\), et inversement.}

La liaison de ces théorèmes avec \S 1, 2. et 3. donne le rapport fondamental pour la suite entre les matrices permutables avec une représentation et les \(\frR\)-, \(\frS\)-endomorphismes.

\paragraph{Théorème sur les matrices permutables.}
\emph{Si \(\frR\to\frR^*\) est une représentation réciproque de \(n\)-ième degré de \(\frR\) dans \(\frS\), et si \(\frB^*\) désigne l'anneau de toutes les matrices (de \(n\)-ième degré dans \(\frS\)) permutables élément par élément avec \(\frR^*\), alors \(\frB^*\) donne une représentation anti-isomorphe des \(\frR\)-, \(\frS\)-endomorphismes -- donc de ces \(\frS\)-endomorphismes qui sont en même temps \(\frR\)-endomorphismes -- du module de représentation engendré, ou encore, si l'anneau-produit \(\frT\) existe, des \(\frT\)-endomorphismes.}

Soit \(\frB\) l'anneau de tous les \(\frR\)-, \(\frS\)-endomorphismes de \(\frM\). Comme sous-anneau de l'anneau \(\mA\) des \(\frS\)-endomorphismes, \(\frB\) admet une représentation anti-isomorphe dans \(\frS\). D'après \S 1, 2. et 3., \(\frB\) se compose de tous les endomorphismes de \(\mA\) qui commutent avec l'action de \(\frR\), donc qui satisfont à la relation (2a). Par conséquent, \(\frB\) est formé de tous les endomorphismes qui commutent élément par élément avec ceux de \(\overline{\frR}\), où \(\overline{\frR}\) désigne l'image de \(\frR\) dans \(\mA\). Comme les représentations de \(\frB\) et de \(\overline{\frR}\) sont des anti-isomorphismes, et non pas de simples antihomomorphismes, c'est bien le résultat à démontrer.

\paragraph{Remarque.}
Les \(\frR\)-, \(\frS\)-endomorphismes signifient une correspondance \(m_i\to m_i'\) telle qu'au moyen des \(m_i\) naisse la même représentation \(\frR^*\); les \(m_i'\) n'ont pas nécessairement à fournir une \(\frS\)-base de \(\frM\), conformément au fait que les matrices de \(\frB^*\) ne sont pas nécessairement inversibles. Si la somme \(\sum m_i'\frS\) n'est plus directe, il faut entendre ``représentation'' dans un sens transposé:
\[
        (m_1',m_2',\ldots m_n')a=(m_1',\ldots m_n')A
\]
reste juste, mais \(A\) n'est plus déterminé univoquement par \(a\).

Une autre remarque concernant le théorème de passage est la suivante: les matrices diagonales \(E\cdot s\) fournissent une représentation \emph{isomorphe (et non anti-isomorphe)} de \(\frS\), la représentation ``identique'' de \(\frS\). Le fait qu'il s'agisse ici d'un isomorphisme \emph{et non d'un anti-isomorphisme} résulte de ce que l'anneau des \(\frS\)-endomorphismes de \(\frM\) possède un sous-anneau anti-isomorphe à \(\frS\), dont il s'agit de prendre la représentation réciproque. Car chaque sommant monogène \(m_i\frS\) de \(\frM\) est isomorphe à \(\frS\) comme module à droite avec opérateurs; son anneau des \(\frS\)-endomorphismes est donc anti-isomorphe à \(\frS\) (\S 1, 3., conséquence 1).

La correspondance de \(\frT\) à des matrices dans \(\frS\), définie univoquement par \(\frR\) et \(\frS\), n'est donc \emph{pas une représentation de l'anneau \(\frT\), mais l'extension de la représentation réciproque de \(\frR\) à une représentation homomorphe relativement aux opérateurs de \(\frS\):}
\[
        \sum r_i s_i\to \sum R_i s_i.
\]
En particulier, la représentation réciproque de \(\frR\) est donc en même temps \emph{homomorphe relativement aux opérateurs de l'intersection \([\frR,\frS]\), située dans le centre de \(\frR\), de \(\frR\) et \(\frS\).}

\clearpage
\providecommand{\frG}{\mathfrak G}
\providecommand{\frL}{\mathfrak L}
\providecommand{\frM}{\mathfrak M}
\providecommand{\frN}{\mathfrak N}
\providecommand{\frT}{\mathfrak T}

\subsection*{\S 3. Modules relativement à un corps gauche}

Dans ce qui suit il ne s'agira que de représentations dans des corps gauches en général non commutatifs. On rassemblera donc quelques théorèmes simples sur les modules de formes linéaires sur un corps gauche.

\paragraph{1. Base normale d'un sous-module d'après une base donnée du module total.}
Soit \(A\) un corps gauche et soit
\[
        \frN=x_1A+\ldots+x_nA
\]
un module de formes linéaires de rang \(n\); soit en outre
\[
        \frL=z_1A+\ldots+z_\ell A
\]
un sous-module de rang \(\ell\le n\). Les \(z_i\) s'appellent base normale d'après \(x\) s'ils sont de la forme
\[
        z_i=x_i-(x_{l+1}\alpha_{i,l+1}+\ldots+x_n\alpha_{i,n}),
        \qquad (i=1,\ldots,l),
\]
Tout module \(\frL\) possède, après numérotation convenable des \(x_i\), une base normale. En effet, après une numérotation convenable on obtient
\[
        \frN=\frL+x_{l+1}\mathsf{A}+\ldots+x_n\mathsf{A},
\]
donc
\[
        x_i\equiv x_{l+1}\alpha_{i,l+1}+\ldots+x_n\alpha_{i,n}
        \pmod{\frL},\qquad i=1,\ldots,\ell. \tag{2}
\]
et par conséquent les
\[
        z_i=x_i-(x_{l+1}\alpha_{i,l+1}+\ldots+x_n\alpha_{i,n})
\]
correspondants appartiennent à \(\frL\) et épuisent \(\frL\), puisqu'ils forment, avec \(x_{l+1},\ldots,x_n\), une base de \(\frN\).

\paragraph{2. Module d'extension.}
Soit de nouveau \(\mathsf{A}\) un corps gauche, et \(\mathsf{P}\) un sous-corps gauche. Un \(\mathsf{A}\)-module \(\frN\) de rang \(n\) s'appelle module d'extension d'un \(\mathsf{P}\)-module \(\frM\); \(\frN=\frM_{\mathsf{A}}\), si \(\frM\) est sous-module de même rang de \(\frN\). Si
\[
        \frM=y_1P+\ldots+y_nP,
\]
alors
\[
        \frN=\frM_A=y_1A+\ldots+y_nA.
\]
Réciproquement, à chaque \(\mathsf{P}\)-module \(\frM=y_1\mathsf{P}+\ldots+y_n\mathsf{P}\) correspond, à isomorphie de modules près, un module d'extension \(\frM_{\mathsf{A}}\) déterminé de façon univoque. Car le module \(\overline y_1\mathsf{A}+\ldots+\overline y_n\mathsf{A}\) contient un sous-module \(\overline{\frM}\) isomorphe à \(\frM\), qui peut être identifié à \(\frM\).

\paragraph{Conséquence a).}
Si \(z_1,\ldots,z_t\) sont des éléments de \(\frM\) linéairement indépendants relativement à \(P\), alors ils sont aussi linéairement indépendants sur \(\mathsf{A}\) dans \(\frN=\frM_{\mathsf{A}}\) (car ils peuvent être complétés en une base de \(\frM\), et donc de \(\frN\)). Tout \(\mathsf{P}\)-module
\[
        \frT=z_1P+\ldots+z_tP
\]
contenu dans \(\frM\) engendre donc un module d'extension
\[
        \frT_A=z_1A+\ldots+z_tA\subset \frM_A.
\]

\paragraph{Conséquence b).}
Tout \(\mathsf{P}\)-module \(\frT=z_1\mathsf{P}+\ldots+z_t\mathsf{P}\) contenu dans \(\frM\) est module de contraction, c'est-à-dire l'intersection de son module d'extension \(\frT_{\mathsf{A}}\) avec \(\frM\):
\[
        \frT=\frT_A\cap\frM .
\]
Car \(\frT\leqq \frT_{\mathsf{A}}\cap\frM\); si, réciproquement, \(a\) appartient à \(\frT_{\mathsf{A}}\cap\frM\), on a
\[
        a=\sum z_i\alpha_i,
\]
donc les éléments \(a,z_1,\ldots,z_t\) de \(\frM\) sont dépendants sur \(\mathsf{A}\), donc d'après la conséquence a) aussi sur \(\mathsf{P}\); par suite \(a\) appartient à \(\frT\).

\paragraph{3. Théorème sur les modules invariants.}
Soit de nouveau \(\frN=\frM_{\mathsf{A}}\) module d'extension d'un \(\mathsf{P}\)-module \(\frM\) de rang \(n\), et soit en outre \(\mathsf{P}\) le corps gauche complet des invariants relativement à un groupe \(\frG\) d'automorphismes d'anneaux de \(\mathsf{A}\). Le groupe \(\frG\) est défini comme domaine d'opérateurs de \(\frN=\frM_{\mathsf{A}}\) par les prescriptions
\[
        G\cdot m=m\quad\text{pour }m\text{ dans }\frM\text{ et }G\text{ dans }\frG
\]
et
\[
        G\cdot\sum m_i\alpha_i=\sum m_i\,G(\alpha_i)
        \quad\text{pour }m\text{ dans }\frM\text{ et }\alpha\text{ dans }\mathsf{A}.
\]

\paragraph{Lemme.}
D'après ces prescriptions, \(\frM\) est constitué par la totalité des éléments de \(\frN\) qui restent fixes sous \(\frG\). En effet, si \(x_1,\ldots,x_n\) est une \(\mathsf{P}\)-base de \(\frM\), donc \(w=x_1\alpha_1+\ldots+x_n\alpha_n\), alors \(G(w)=w\) pour tout \(G\) de \(\frG\) donne aussi \(G(\alpha_i)=\alpha_i\), et par conséquent, d'après l'hypothèse, \(\alpha_i\) appartient à \(\mathsf{P}\).

\paragraph{Théorème.}
Les modules d'extension \(\frL_{\mathsf{A}}\) des sous-modules \(\frL\) de \(\frM\), et eux seuls, sont des sous-modules admissibles relativement à \(\frG\) comme domaine d'opérateurs.

La première moitié est claire. Soit réciproquement \(\frL\) admissible relativement à \(\frG\), et soit \(z_1,\ldots,z_l\) une base normale de \(\frL\) (voir 1.) d'après \(x_1,\ldots,x_n\), où \(x_1,\ldots,x_n\) est choisi comme \(\mathsf{P}\)-base de \(\frM\), donc est admis élément par élément par \(\frG\). Par hypothèse,
\[
        G(z_i)=x_i-\bigl(x_{l+1}G(\alpha_{i,l+1})
              +\ldots+x_nG(\alpha_{i,n})\bigr)
\]
est élément de \(\frL\) pour tout \(G\) de \(\frG\), donc, pour \(G\) fixe, exprimable linéairement par les \(z\); ainsi, comme il résulte de la comparaison des coefficients en \(x_1,\ldots,x_l\), \(G(z_i)=z_i\) pour tout \(G\) de \(\frG\). Les \(z\) appartiennent donc -- d'après le lemme -- à \(\frM\), et \(\frL\) devient extension du \(\mathsf{P}\)-module
\[
        \frT=z_1P+\ldots+z_\ell P
\]
contenu dans \(\frM\), avec \(\frT=\frL\cap\frM\) (d'après 2).\footnote{Les théorèmes de ce paragraphe restent valables, au moyen de raisonnements simples de bon ordre, même lorsque le rang de \(\frM\) sur \(\mathsf{P}\) devient infini (voir G. Kothe: Ein Beitrag zur Theorie der kommutativen Ringe ohne Endlichkeitsvoraussetzung \S 1, Gott. Nachr. 1931, p. 195--207).}

\clearpage
\providecommand{\frR}{\mathfrak R}
\providecommand{\frS}{\mathfrak S}
\providecommand{\frG}{\mathfrak G}

\subsection*{\S 4. Systèmes hypercomplexes et leurs classes de représentations}

\paragraph{1. Théorème d'extension.}
Nous relions maintenant les théorèmes de représentation du \S 2 aux théorèmes sur les modules du \S 3. À la place des anneaux généraux \(\frR\), nous considérons des systèmes hypercomplexes avec élément unité -- \(S,T\ldots\) -- sur un corps commutatif \(\mathsf{P}\), donc
\[
        S=x_1\mathsf{P}+\ldots+x_n\mathsf{P};
\]
à la place des anneaux de représentation arbitraires \(\frS\), seulement des corps \(\mathsf{A},\mathsf{B}\ldots\), et, sauf remarque contraire, de tels corps avec centre \(\mathsf{P}\). Dans ce cas il existe toujours l'anneau-produit dans lequel \(S\) et \(\mathsf{A}\) deviennent permutables élément par élément, avec intersection \(\mathsf{P}\); il s'agit du produit direct \(S\times \mathsf{A}\), qui devient en même temps le module d'extension \(S_{\mathsf{A}}\) de \(S\) au sens du \S 3,
\[
        S_{\mathsf{A}}=x_1\mathsf{A}+\ldots+x_n\mathsf{A},
\]
et qui est donc aussi désigné comme anneau d'extension.

Dans cette spécialisation, le théorème sur les modules invariants (\S 3, 3.) devient le

\paragraph{Théorème d'extension.}
Tout idéal bilatéral de \(S_{\mathsf{A}}\) est idéal d'extension d'un idéal de \(S\), à savoir extension de son intersection avec \(S\).

En effet, si le groupe \(\frG\) pris pour base est celui des automorphismes intérieurs de \(A\), alors \(P\), centre de \(A\), devient corps complet des invariants, tandis que la permutabilité élément par élément de \(A\) avec \(S\) dit que \(\frG\), d'après \S 3, 3., est étendu à \(S_{\mathsf{A}}\) de telle sorte que \(S\) devienne domaine complet des invariants. Tout idéal bilatéral \(\mathfrak a\) de \(S_{\mathsf{A}}\) est un sous-groupe admissible -- à cause de \(\alpha^{-1}\mathfrak a\alpha\leqq\mathfrak a\) -- donc, d'après \S 3, 3., extension de son module-intersection \(\mathfrak a\cap S\), qui devient un idéal de \(S\).

Addendum: Si \(S\) est commutatif, alors \(S\) devient le centre de \(S_{\mathsf{A}}\). Car \(S\) est contenu dans le centre et est domaine complet des invariants par rapport aux automorphismes intérieurs engendrés par \(\mathsf{A}\). Plus généralement, le centre de \(S\) devient aussi le centre de \(S_{\mathsf{A}}\).

Notation: Par \(\mathsf{A}_r,\mathsf{B}_n\ldots\) on entendra toujours les anneaux matriciels du degré indiqué sur \(\mathsf{A},\mathsf{B}\ldots\), donc
\[
        A_r=\sum A c_{ik}=\sum c_{ik}A.
\]

Dans ce qui suit on utilisera essentiellement une conséquence, le théorème d'extension pour les systèmes simples: si \(S\) est un système simple\footnote{Système simple signifie, comme d'habitude, ``simple bilatéral''.}, alors \(S_{\mathsf{A}}\) devient lui aussi simple:
\[
        S_A=\sum B c_{ik}=\sum c_{ik}B=B_t
\]
avec corps \(\mathsf{B}\) associé, dont le centre contient \(\mathsf{P}\). Ici \(\mathsf{A}\), et donc aussi \(\mathsf{B}\), peut être de rang infini sur \(\mathsf{P}\); seul \(S\) est supposé hypercomplexe. Si \(S\) et \(\mathsf{A}\) possèdent le centre commun \(\mathsf{P}\), alors \(\mathsf{P}\) devient aussi centre de \(S_{\mathsf{A}}\).

Plus généralement encore: si \(S\) est un système hypercomplexe simple, alors le produit direct \(S\times \mathsf{A}_r\) devient lui aussi simple. Car \(S\times \mathsf{A}_r\) coïncide avec \(S_{\mathsf{A}}\times \mathsf{P}_r\), et devient donc égal à \(\mathsf{B}_{tr}\), lorsque \(S_{\mathsf{A}}=\mathsf{B}_t\).

\paragraph{2. Classes de représentations.}
Par représentation -- réciproque ou directe -- d'un système hypercomplexe, on entendra, comme d'habitude, une représentation homomorphe relativement aux opérateurs du domaine de coefficients \(\mathsf{P}\) (\S 2, fin) et distincte de la représentation nulle. La liaison du théorème de passage et de sa conséquence (\S 2, 3.) avec le théorème d'extension de 1. donne le

\paragraph{Théorème sur les classes de représentations.}
Si \(S\) est hypercomplexe avec élément unité, avec \(\mathsf{P}\) comme domaine de coefficients, et si \(\mathsf{A}\) est un corps de centre \(\mathsf{P}\), alors \(S\) possède autant de classes irréductibles différentes de représentations réciproques de \(S\) dans \(\mathsf{A}\) qu'il y a de classes d'idéaux à droite simples dans l'anneau quotient de \(S_{\mathsf{A}}\) par son radical. -- Si \(S\) est un système simple, alors il possède exactement une classe irréductible de représentations réciproques et une classe irréductible de représentations directes dans \(\mathsf{A}\).\footnote{Voir Darstellungstheorie, note 20), pour le cas où \(\mathsf{A}\) est le corps gauche des endomorphismes associé de \(S\).}

Car, d'après la conséquence du théorème de passage (\S 2, 3.), les classes simples de représentations réciproques et les classes simples de modules sur \(S_{\mathsf{A}}\) se correspondent biunivoquement. Puisque \(S_{\mathsf{A}}\) est de rang fini sur le corps \(\mathsf{A}\), il satisfait aux conditions maximale et minimale pour les idéaux unilatères; ainsi la première partie de l'énoncé découle directement de Darstellungstheorie, \S 19. Si \(S\) est simple, alors \(S_{\mathsf{A}}\) l'est aussi d'après le théorème d'extension. À cause des conditions maximale et minimale, il est en même temps complètement réductible, donc ne possède qu'une seule classe d'idéaux unilatères simples, c'est-à-dire que \(S\) ne possède qu'une seule classe irréductible de représentations réciproques dans \(\mathsf{A}\). Avec \(S\), le système qui lui est anti-isomorphe est lui aussi simple et ne possède qu'une seule classe irréductible de représentations réciproques; donc \(S\) ne possède aussi qu'une seule classe simple de représentations directes.

\paragraph{3. Relation de rang.}
Le fait que, si \(S\) est simple,
\[
        S_A=\sum Bc_{ik}
\]
est de rang fini aussi bien sur \(A\) que sur \(B\), donne la relation de rang:
\[
        n=rt\quad\text{avec}\quad n=(S:P),\quad t^2=(S_A:B),
        \quad r=\text{degré de la représentation réciproque irréductible}.      \tag{1}
\]

En effet, \(S_{\mathsf{A}}\) est lui-même module de représentation réciproque pour \(S\) dans \(\mathsf{A}\) (la représentation se trouve déjà dans \(\mathsf{P}\)). Comme anneau matriciel sur \(\mathsf{B}\), \(S_{\mathsf{A}}\) se décompose en \(t\) idéaux à droite simples, donc en \(S_{\mathsf{A}}\)-modules, dont le rang sur \(\mathsf{A}\) est égal au degré \(r\) de la représentation réciproque irréductible. Puisque \(n\) est en même temps le rang de \(S_{\mathsf{A}}\) sur \(\mathsf{A}\), la relation de rang en résulte.

\paragraph{4. Renforcement de la relation de rang, lorsque \(\mathsf{A}\) est de rang fini sur \(\mathsf{P}\).}
Dans ce cas les trois relations valent:
\[
        (S:P)=n=rt,                                                     \tag{1}
\]
\[
        (B:P)\,t=(A:P)\,r,                                               \tag{2}
\]
\[
        (S:P)(B:P)=(A:P)\,r^2.                                           \tag{3}
\]
Ici (2) signifie le rang sur \(\mathsf{P}\) d'un idéal à droite simple de \(S_{\mathsf{A}}\), exprimé d'après 3. par le rang sur \(\mathsf{B}\), respectivement sur \(\mathsf{A}\), tandis que (3) signifie le rang d'un idéal à droite simple dans un anneau matriciel \(\mathsf{B}_n\) et résulte de (2) par multiplication par \(r\), en vertu de (1).

\clearpage
\providecommand{\frM}{\mathfrak M}

\subsection*{\S 5. Application des théorèmes de représentation aux anneaux de matrices}

La représentation réciproque de \(S\) dans \(\mathsf{A}\), considérée au \S 4, peut aussi être comprise comme homomorphisme direct de \(S\) dans un sous-anneau de \(A_r\), où \(A\) désigne le corps anti-isomorphe à \(\mathsf{A}\). Le principe des paragraphes suivants consiste, inversement, à ramener l'étude des anneaux de matrices \(A_r\) à la théorie de la représentation dans le corps anti-isomorphe \(\mathsf{A}\).

\paragraph{1. Plongement réductible et irréductible dans \(A_r\).}
Soient \(\mathsf{A}\) et \(A\) des corps anti-isomorphes, de rang fini ou infini sur leur centre \(\mathsf{P}\).\footnote{En général, les corps ou systèmes réciproques seront désignés par des lettres grecques et latines correspondantes.} Le système simple \(S\) sur \(\mathsf{P}\) s'appelle plongeable irréductiblement, respectivement réductiblement, dans \(A_r\), lorsque \(S\) possède une représentation réciproque irréductible, respectivement réductible, de degré \(r\) dans \(\mathsf{A}\). Si \(S\) est plongeable irréductiblement dans \(A_r\), alors il est plongeable réductiblement dans tous les anneaux matriciels \(\mathsf{A}_{rs}\) avec \(s>1\), puisqu'il possède des représentations réductibles de ces degrés et d'eux seuls (\S 4, 2.). Ainsi \(S\) est isomorphe à certains sous-anneaux de \(A_r\) et de \(A_{rs}\). L'isomorphisme devient prolongement de l'identique de \(\mathsf{P}\);\footnote{Comme déjà pour les représentations, les isomorphismes qui interviennent dans ce qui suit sont toujours des prolongements de l'identique de \(\mathsf{P}\) -- isomorphes sur \(\mathsf{P}\) --, sans que cela soit chaque fois mentionné spécialement.} d'après \S 4, 3., \(r\) est alors un diviseur du rang \(n\) de \(S\) sur \(\mathsf{P}\).

De cette manière on obtient tous les sous-anneaux simples, contenant \(\mathsf{P}\) et de rang fini sur \(\mathsf{P}\), des différents \(A_f\) (\(f=1,2\ldots\)); car tout sous-anneau de ce genre est anti-isomorphe à un sous-anneau de \(\mathsf{A}_f\), et admet donc une représentation réciproque réductible ou irréductible dans \(\mathsf{A}\).

\paragraph{2. Théorème sur les automorphismes intérieurs.}
Si \(S^{(1)}\) et \(S^{(2)}\) sont deux sous-anneaux simples de \(A_f\), contenant \(P\), de rang fini sur \(\mathsf{P}\), et si \(S^{(1)}\) et \(S^{(2)}\) sont isomorphes sur \(\mathsf{P}\), alors cet isomorphisme est engendré par un automorphisme intérieur de \(A_f\).

En effet, \(S^{(1)}\) et \(S^{(2)}\) signifient des plongements -- en général réductibles -- d'un système simple \(S\) dans \(A_f\), donc des représentations directes de degré \(f\) dans \(A\); ils appartiennent ainsi à la même classe de représentations directes réductibles dans \(A\); la matrice qui effectue le passage engendre l'automorphisme.

Si \(A\) lui-même est de rang fini sur \(\mathsf{P}\), donc hypercomplexe, le théorème devient le théorème connu: deux sous-systèmes simples de \(A_f\), contenant le centre \(P\) et isomorphes sur \(\mathsf{P}\), se transforment l'un dans l'autre par un automorphisme intérieur de \(A_f\). En particulier tout automorphisme de \(A_f\) est intérieur.\footnote{Cela n'est plus vrai lorsque \(A\) est de rang infini sur \(\mathsf{P}\) (voir Kothe, dans le travail cité à la note 5).}

\paragraph{3. Théorème sur les sous-anneaux permutables élément par élément.}
Ce théorème découle de nouveau, d'après le principe mentionné au début du paragraphe, du théorème sur les matrices permutables du \S 2, 3.:

Si \(S\) est un système simple contenant \(\mathsf{P}\), dans \(A_f\), et si \(R\) désigne la totalité des éléments de \(A_f\) qui sont permutables élément par élément avec \(S\), alors \(R\) devient lui aussi un anneau de matrices: \(R=B_s\). Les corps associés \(\mathsf{B}\) de \(S_{\mathsf{A}}\) et \(B\) de \(R\) deviennent anti-isomorphes. L'intersection de \(R\) et de \(S\) devient le centre de \(S\). \(R\) devient corps si et seulement si \(S\) est plongé irréductiblement dans \(A_f\).

Car d'après \S 2, 3., \(R\) est isomorphe à l'anneau des endomorphismes du module de représentation réciproque \(\frM\) de \(S\) dans \(\mathsf{A}\); mais c'est l'anneau des endomorphismes de \(\frM\), conçu comme \(S_{\mathsf{A}}\)-module. Si \(\frM\) se décompose en \(s\) \(S_{\mathsf{A}}\)-modules simples et si, comme au \S 4, 1., on pose \(S_{\mathsf{A}}=\sum \mathsf{B}c_{ik}\), alors \(R\) devient isomorphe à \(\sum Bc_{ik}\), où \(\mathsf{B}\) et \(B\) sont anti-isomorphes (car \(\mathsf{B}\) est anti-isomorphe au corps gauche des endomorphismes des idéaux à droite simples de \(S_{\mathsf{A}}\), donc des \(S_{\mathsf{A}}\)-modules simples, d'après \S 1, 3., fin). \(R\) devient corps -- \(R=B\) -- si et seulement si \(s\) est égal à un, donc si \(S\) est plongé irréductiblement. L'affirmation selon laquelle l'intersection de \(R\) et de \(S\) est le centre de \(S\) résulte immédiatement de la définition de \(R\).

\paragraph{4. Théorème de commutation renforcé dans le cas d'un anneau matriciel hypercomplexe.}
Dans ce cas, les relations de rang du \S 4, 4. donnent le renforcement suivant:

Les sous-systèmes simples contenant \(\mathsf{P}\) dans \(A_f\) -- où \(A\) est de rang fini sur son centre \(\mathsf{P}\) -- se décomposent en paires \(S,\overline S\) telles que \(\overline S\) soit constitué de la totalité des éléments permutables élément par élément avec \(S\), et réciproquement. Les corps associés de \(S_{\mathsf{A}}\) et de \(\overline S\) sont anti-isomorphes, de même ceux de \(S\) et de \(\overline S_{\mathsf{A}}\). L'intersection de \(S\) et de \(\overline S\) devient égale au centre commun. L'un des sous-anneaux est corps si et seulement si l'autre est plongé irréductiblement. Le produit des nombres de rang sur \(\mathsf{P}\) de \(S\) et de \(\overline S\) est égal au rang de \(A_f\). Si \(f=rs=\overline r\,\overline s\) -- où \(S\) est plongeable irréductiblement dans \(A_r\) et \(\overline S\) dans \(A_{\overline r}\) --, alors les corps associés de \(S\) et de \(\overline S\) deviennent isomorphes à une paire de sous-anneaux permutables de \(A_g\), avec \(f=g s\overline s\).

Il ne reste essentiellement qu'à démontrer l'assertion sur le produit des nombres de rang. Si d'abord \(S\) est plongé irréductiblement, donc si \(\overline S\) est corps et anti-isomorphe à \(\mathsf{B}\) -- où \(S_{\mathsf{A}}=\mathsf{B}_t\) --, alors la relation de rang (3) du \S 4, 4. donne l'assertion sur le produit, puisque maintenant \(f=r\), et que le rang de \(\overline S\) est égal à celui de \(\mathsf{B}\). Si \(S\) est plongé réductiblement, donc si \(\overline S=B_s\), alors \(f=rs\), et la multiplication de (3) par \(s^2\) donne l'assertion. Il en résulte immédiatement que \(S\) est constitué de la totalité des éléments permutables élément par élément avec \(\overline S\), et par suite tout le reste de 3., sauf la dernière assertion. Celle-ci résulte par passage deux fois à un plongement irréductible. Dans \(A_r\) -- avec \(f=rs\) -- \(S\) est plongé irréductiblement, et \(S\) et \(R\) deviennent des systèmes deux à deux permutables, où \(R\) est isomorphe au corps \(B\). Mais de plus \(S\), à cause de la réciprocité de \(S\) et de \(\overline S\), est égal à \(\overline B_{\overline s}\) -- où \(\overline B\) a, pour \(\overline S\), la même signification que \(B\) pour \(S\) --, donc \(R\) est plongé réductiblement dans \(A_r\), irréductiblement dans \(A_g\) avec \(r=g\overline s\). Ici les systèmes deux à deux permutables deviennent donc isomorphes à \(B\) et \(\overline B\).

\clearpage
\providecommand{\Gcal}{\mathcal G}

\subsection*{\S 6. Théorie galoisienne des systèmes simples}

Le théorème de commutation renforcé du \S 5, 4. exprime la théorie galoisienne des systèmes simples relativement au centre comme domaine fondamental. Nous considérons d'abord les corps gauches, puis les systèmes simples en général.

\paragraph{1. Théorie galoisienne des corps gauches de rang fini sur le centre.}
Le groupe galoisien \(\Gcal\) de \(A\) est défini comme le groupe de tous les automorphismes qui prolongent l'identique du centre \(\mathsf{P}\), donc, d'après \S 5, 2., comme le groupe de tous les automorphismes intérieurs. On a donc
\[
        \Gcal \simeq A^*/P^*,
\]
où \(A^*\), respectivement \(\mathsf{P}^*\), désigne le groupe multiplicatif des éléments non nuls de \(A\), respectivement de \(\mathsf{P}\). Aux sous-groupes \(\mathfrak H\) de \(\Gcal\) correspondent donc biunivoquement les sous-groupes \(H^*\) de \(A^*\) contenant \(\mathsf{P}^*\), par
\[
        \mathfrak H \sim H^*/P^* .
\]
Un sous-groupe \(\mathfrak H\) de \(\Gcal\) s'appelle fermé si \(H^*\), après adjonction du zéro, devient un corps gauche \(H\). Le fait que \(C\) admette le groupe \(\mathfrak H\) équivaut au fait que \(C\) soit permutable élément par élément avec \(H\). Ainsi le théorème de commutation \S 5, 4. devient le

\paragraph{Théorème principal de la théorie galoisienne des corps gauches.}
Les corps gauches \(C\) compris entre \(\mathsf{P}\) et \(A\) et les sous-groupes fermés \(\mathfrak H\) de \(\Gcal\) peuvent être coordonnés biunivoquement de telle sorte que \(C\) devienne le corps gauche complet des invariants de \(\mathfrak H\), et \(\mathfrak H\) le groupe complet des invariants de \(C\).

\paragraph{2. Théorie galoisienne des systèmes simples.}
Si \(A_f\) est un système simple de centre \(\mathsf{P}\), alors le groupe galoisien \(\Gcal\) est de nouveau le groupe des automorphismes intérieurs, donc isomorphe à
\[
        A_f^*/P^*,
\]
où maintenant \(A_f^*\) désigne le groupe multiplicatif des éléments réguliers de \(A_f\). Un sous-groupe \(\mathfrak H\simeq H^*/\mathsf{P}^*\) de \(\Gcal\) s'appelle simplement fermé si le sous-anneau \(H\) engendré par \(H^*\) est un système simple, et si \(H^*\) est constitué par la totalité des éléments réguliers de \(H\). Le passage du théorème de commutation à la théorie galoisienne est donné ici par le

\paragraph{Lemme.}
Tout système simple de \(A_f\) contenant \(P\) est engendré par le groupe \(H^*\) de ses éléments réguliers. Cela est clair lorsque \(\mathsf{P}\) possède une infinité d'éléments; car l'``élément général'' formé avec des indéterminées est régulier, et par spécialisation convenable on peut former des éléments de base réguliers sur \(\mathsf{P}\). Mais le lemme vaut en général d'après Shoda.\footnote{Shoda, travail cité à la note 3). L'introduction des indéterminées y est évitée.}

En vertu du lemme, la permutabilité de \(S\) avec un système simple \(\overline S\) équivaut de nouveau au fait que \(S\) admette les automorphismes induits par les éléments réguliers de \(\overline S\). Ainsi le théorème de commutation \S 5, 4. devient ici encore le

\paragraph{Théorème principal de la théorie galoisienne des systèmes simples.}
Les systèmes simples \(S\) de \(A_f\) contenant \(\mathsf{P}\) et les sous-groupes simplement fermés \(\mathfrak H\) de \(\Gcal\) peuvent être coordonnés biunivoquement de telle sorte que \(S\) devienne le domaine complet des invariants de \(\mathfrak H\), et \(\mathfrak H\) le groupe complet des invariants de \(S\).

\paragraph{3. Preuve au moyen du principe de prolongement.}
Je donne encore, pour le cas d'un corps gauche, une seconde preuve du théorème principal; elle repose sur le principe du prolongement des isomorphismes, c'est-à-dire des représentations du premier degré, et, puisqu'elle se passe du dénombrement des rangs, impose moins d'hypothèses de finitude. En particulier on voit ainsi dans quelle direction un transfert à des corps gauches \(S\) de rang infini sera possible. La preuve n'utilise pas non plus le fait qu'il n'existe qu'une seule classe de représentations réciproques irréductibles; elle vaut donc exactement de la même façon dans le cas commutatif; voir \S 8, 2. Je fais précéder quelques lemmes.

\paragraph{Hypothèses.}
Soit un système simple \(S\) sur \(\mathsf{P}\) admettant une représentation réciproque du premier degré dans \(\mathsf{A}\), donc un corps gauche; ici \(\mathsf{A}\) peut être de rang fini ou infini sur son centre \(\mathsf{P}\). Soit donnée, d'après \S 4, 3., une décomposition de \(S_{\mathsf{A}}\) en \(n\) idéaux à droite simples, isomorphes comme modules à opérateurs:
\[
        S_{\mathsf{A}}=\mathfrak r_1+\ldots+\mathfrak r_n
            =e_1S_{\mathsf{A}}+\ldots+e_nS_{\mathsf{A}}
            =e_1\mathsf{A}+\ldots+e_n\mathsf{A} ;
\]
la représentation réciproque engendrée par \(e_i\) est donc définie par
\[
        e_i s=e_i\sigma .
\]

\paragraph{Lemme 1.}
Les \(n\) représentations réciproques engendrées par \(e_1,\ldots,e_n\) sont distinctes, donc sont des représentations différentes de la même classe de représentations. Ainsi \(S\) possède au moins autant de représentations différentes que son rang.

Pour la preuve, nous étendons les représentations de \(S\), d'après la fin du \S 2, à des correspondances de \(S_{\mathsf{A}}\), qui deviennent homomorphes relativement aux opérateurs de \(\mathsf{A}\). Ces correspondances sont ici définies par
\[
        e_i w=e_i\omega
\]
avec \(\omega\in A\) pour chaque \(w\in S_A\). Si maintenant \(e_i\) et \(e_j\), \((i\ne j)\), engendraient la même représentation, on aurait donc \(e_iw=e_jw\) pour chaque \(w\in S_A\). Cela est impossible à cause de \(e_ie_i=e_i\) et de \(e_je_i=0\).

\paragraph{Lemme 2.}
Si \(T\) est un corps gauche compris entre \(\mathsf{P}\) et \(S\), et si \(s\) est le rang de \(S\) sur \(T\), alors tout anti-isomorphisme de \(T\) dans \(\mathsf{A}\) admet au moins \(s\) prolongements différents.

Soit l'anti-isomorphisme, donc la représentation réciproque de \(T\) dans \(\mathsf{A}\), fourni par une décomposition
\[
        T_{\mathsf{A}}=E_1T_{\mathsf{A}}+\ldots+E_hT_{\mathsf{A}}
            =E_1A+\ldots+E_hA,
        \qquad E_it=E_i\tau ,
\]
où les \(E_i\) sont des idempotents. Ce n'est pas une restriction de la généralité, puisqu'une représentation engendrée par un élément de base \(E_i\alpha\) est aussi engendrée par \(\alpha^{-1}\cdot E_i\alpha\), donc par un idempotent. La multiplication à droite par \(S\) donne
\[
        S_A=E_1S_A+\ldots+E_hS_A .
\]
Ici les \(E_iS_{\mathsf{A}}\) ont tous le rang \(s\) sur \(\mathsf{A}\). En effet, le rang de \(E_iS_{\mathsf{A}}\) est au plus \(s\), comme le montre l'insertion d'une \(T\)-base de \(S\) en vertu de la permutabilité de \(S\) et de \(\mathsf{A}\):
\[
        S=Tw_1+\ldots+Tw_s,
\]
\[
        E_iS_{\mathsf{A}}=E_iT_{\mathsf{A}}w_1+\ldots+E_iT_{\mathsf{A}}w_s
              =E_iw_1\mathsf{A}+\ldots+E_iw_s\mathsf{A} .
\]
Comme la somme \(S_{\mathsf{A}}\) est de rang \(n=hs\), chaque \(E_iS_{\mathsf{A}}\) a donc exactement le rang \(s\). Ainsi
\[
        E_iS_A=r_{i1}+\ldots+r_{is}
              =e_{i1}A+\ldots+e_{is}A,
\]
où les \(s\) représentations engendrées par \(e_{i1},\ldots,e_{is}\) sont toutes distinctes d'après le lemme 1. Elles sont toutes des prolongements de la représentation de \(T\) engendrée par \(E_i\); en effet, la relation \(E_it=E_i\tau\) donne
\[
        (e_{i1}+\ldots+e_{is})t=(e_{i1}+\ldots+e_{is})\tau
\]
et donc \(e_{il}t=e_{il}\tau\) à cause de la décomposition en somme directe.

\paragraph{Définition.}
La division en classes \(\mathfrak I_T\) des anti-isomorphismes \(\mathfrak I\) de \(S\) sur \(\mathsf{A}\), induite par \(T\), est définie ainsi: tous et seulement les anti-isomorphismes de \(\mathfrak I\) qui induisent le même anti-isomorphisme sur \(T\) sont considérés comme équivalents.

\paragraph{Théorème principal.}
Si \(\mathfrak I_T\) est la division en classes induite par \(T\), alors \(T\) est sous-corps gauche maximal relativement à cette division en classes.

En effet, si \(L\) est un corps gauche intermédiaire entre \(T\) et \(S\), de degré \(l\) sur \(T\), alors, d'après le lemme 2, chaque classe de \(\mathfrak I_T\) se décompose en au moins \(l\) classes de \(\mathfrak I_L\). Le passage de cette version du théorème principal à la version usuelle se fait en transformant l'ensemble d'anti-isomorphismes \(\mathfrak I\) en groupe d'automorphismes par multiplication de ses éléments avec l'inverse d'un élément fixe quelconque; alors une classe de \(\mathfrak I_T\) devient le groupe d'invariants de \(T\). L'affirmation que les sous-groupes fermés sont des groupes complets d'invariants y est déjà contenue; il suffit, pour \(\mathfrak H\simeq H^*/\mathsf{P}^*\), de prendre le corps gauche \(H\) comme corps gauche intermédiaire \(T\).

\clearpage
\subsection*{\S 7. La classe d'algèbres semblables. Corps de décomposition}

À partir de maintenant il ne s'agira plus que de corps gauches \(\mathsf{A},A\) et d'anneaux de matrices \(A_r\) qui sont de rang fini sur leur centre \(\mathsf{P}\); donc d'algèbres simples normales sur \(\mathsf{P}\), qui seront appelées brièvement algèbres sur \(\mathsf{P}\), ou même seulement algèbres. \(\mathsf{A},A\) sont alors des algèbres à division. Dans une classe d'algèbres semblables --- \(A_r\sim A_s\) ---, on rassemble tous les \(A_r\) ayant le même \(A\) associé. Les \(A\) isomorphes y sont identifiés.

\paragraph{1. Le groupe des classes d'algèbres.}
Si \(A_r,B_s\) sont des algèbres sur \(\mathsf{P}\), il en va de même pour leur produit, le produit direct sur \(\mathsf{P}\); car d'après la fin du \S 3 on obtient
\[
        A_r\times_P B_s=C_t,
\]
où \(C\) est une algèbre à division de centre \(\mathsf{P}\), déterminée univoquement à isomorphie près. La multiplication de \(A_r,B_s\) par des anneaux de matrices sur \(\mathsf{P}\) montre que cette multiplication est déterminée univoquement pour les classes, lesquelles forment ainsi un système commutatif fermé multiplicativement. En outre vaut le

\paragraph{Théorème de R. Brauer.}
Les classes d'algèbres semblables forment, par rapport à la formation du produit direct, un groupe abélien.

En effet, le système possède une classe unité, constituée par les anneaux de matrices sur \(\mathsf{P}\), donc par les algèbres semblables à \(\mathsf{P}\). Il possède en outre, pour chaque classe \((A)\), l'inverse \((A)^{-1}\), constitué par les algèbres semblables à \(\mathsf{A}\), où \(\mathsf{A}\) est anti-isomorphe à \(A\). Cela résulte du théorème de commutation \S 5, 3., par la spécialisation \(S=A\). Comme \(A\) admet un plongement irréductible dans lui-même, on obtient \(B=\mathsf{P}\) et \(A_{\mathsf{A}}=\mathsf{P}_t\).\footnote{Les théorèmes plus fins sur le groupe des classes d'algèbres utilisent la théorie des systèmes de facteurs; il n'en sera pas question ici. On remarquera que jusqu'ici aucune considération sur les corps commutatifs n'intervient; par exemple, le fait que le rang \((A:\mathsf{P})\) soit un carré n'a été ni utilisé ni démontré.}

\paragraph{2. Corps de décomposition d'une classe d'algèbres.}
La théorie des corps de décomposition repose de nouveau entièrement sur le principe énoncé au début du \S 5; les corps d'extension commutatifs sont représentés dans \(\mathsf{A}\), respectivement dans \(A\). Faisons précéder le

\paragraph{Lemme.}
Si \(A\) est un corps gauche de centre \(\mathsf{P}\), et si \(\mathsf{Z}\) est un corps d'extension commutatif fini de \(\mathsf{P}\), alors \(A_{\mathsf{Z}}\) devient simple de centre \(\mathsf{Z}\). Car cela vaut, d'après \S 4, 1., pour le système \(Z_{\mathsf{A}}\), anti-isomorphe à \(A_{\mathsf{Z}}\). Pour chaque classe \((A)\) d'algèbres semblables à \(A\) sur \(\mathsf{P}\), \(\mathsf{Z}\) engendre donc une classe d'extension \((A)_{\mathsf{Z}}\) d'algèbres simples sur \(Z\), semblables à \(A_{\mathsf{Z}}\).

L'algèbre à division \(D\) associée à une classe d'extension résulte aussitôt du théorème de commutation du \S 5, 3.:

\paragraph{Théorème sur l'algèbre à division associée à une classe d'extension.}
Soit \((A)_{\mathsf{Z}}\) une classe d'extension; que \(Z\) signifie un plongement irréductible, donc une représentation, de \(\mathsf{Z}\) dans \(A_r\). Alors
\[
        (A)_Z=(D),
\]
où \(D\) est constitué par la totalité des éléments de \(A_r\) permutables élément par élément avec \(Z\). Il suffit de remplacer, dans le \S 5, 3., le système simple \(S\) qui y figure par \(Z\), en tenant compte du fait que \(Z_{\mathsf{A}}\) et \(A_{\mathsf{Z}}\) sont anti-isomorphes. S'il s'agit d'un plongement réductible de \(\mathsf{Z}\), alors la totalité des éléments permutables élément par élément avec \(Z\) donne un anneau de matrices sur \(D\).

Un corps d'extension commutatif \(\mathsf{Z}\) de \(\mathsf{P}\) s'appelle, comme on sait, corps de décomposition de la classe \((A)\), lorsque la classe d'extension \((A)_{\mathsf{Z}}\) se décompose complètement, c'est-à-dire devient égale à la classe unité sur \(\mathsf{Z}\), donc à la classe de tous les anneaux de matrices sur \(\mathsf{Z}\).

\paragraph{Théorème sur la caractérisation des corps de décomposition.}
Un corps d'extension commutatif fini \(\mathsf{Z}\) de \(\mathsf{P}\) est corps de décomposition de la classe \((A)\) si et seulement si son plongement irréductible dans \(A_r\) donne un sous-corps commutatif maximal de \(A_r\).

En effet, la propriété pour \(\mathsf{Z}\) d'être corps de décomposition de \((A)\) équivaut à \((D)=(\mathsf{Z})\). Il faut donc, d'après le théorème sur l'algèbre à division associée, que le plongement irréductible \(Z\) soit un sous-corps maximal commutatif. Si cette condition est satisfaite, alors nécessairement \(D=Z\), puisque chaque \(a\) de \(D\), adjoint à \(Z\), engendre un corps d'extension commutatif.

\paragraph{Remarques.}
1. On en déduit en même temps qu'un sous-corps maximal commutatif \(Z\), dans un plongement irréductible, engendre un sous-anneau maximal commutatif. Car la totalité des éléments permutables élément par élément avec \(Z\) forme un corps.

2. Le théorème donne en même temps l'existence des corps de décomposition; car l'algèbre à division associée \(A\) possède certainement des sous-corps maximaux commutatifs.

\paragraph{3. Relations de rang, indice, extensions commutatives infinies.}
L'assertion de rang du \S 5, 4. donne d'abord le fait connu que le rang d'une algèbre simple sur son centre est un carré. En effet, si \(Z\) est maximal commutatif dans \(A\), alors, puisque dans \(A\) tout plongement est irréductible, on obtient
\[
        (Z:P)\cdot(Z:P)=(A:P),
\]
donc
\[
        (A:P)=m^2,\qquad (Z:P)=m,
\]
où \(m\) est l'indice de Schur, qui signifie donc en même temps le degré de tous les corps de décomposition plongeables dans l'algèbre à division \(A\) elle-même. En outre, pour le degré \(n\) d'un corps de décomposition général, on obtient
\[
        n=m\cdot r\text{\footnotemark},
\]
\footnotetext{En général, il existe aussi des corps de décomposition « minimaux », c'est-à-dire de tels corps pour lesquels aucun sous-corps propre n'est corps de décomposition, et cela pour tous les degrés. Voir Brauer--Noether, le travail cité à la note 2).}
par exemple à cause de \((A_r:\mathsf{P})=m^2\cdot r^2\), ou encore d'après la relation de rang \S 4, 3.; ce dernier point parce que l'indice \(m\) est aussi égal au nombre des composantes simples de \(A_{\mathsf{Z}}\), qui devient anneau de matrices sur \(\mathsf{Z}\) de rang \(m^2\). Plus généralement, si \(A_\Lambda\sim D\), si \(L\) est l'image d'un plongement irréductible de \(\Lambda\) dans \(A_r\), et si \(d^2=(D:L)\) est le rang de l'algèbre à division \(D\) sur son centre \(L\), alors on obtient
\[
        l\cdot\,d=m\cdot r .
\]
Car, d'après \S 5, 4. et le théorème sur l'algèbre à division associée de 2., on a
\[
        (A_r:P)=(L:P)\cdot(D:P),
\]
donc
\[
        m^2r^2=l^2d^2 .
\]

De l'existence des corps de décomposition résultent, pour les corps d'extension commutatifs infinis \(\Omega\) de \(\mathsf{P}\), algébriques ou transcendants, les faits suivants: la classe d'extension \((A)_\Omega\) est une classe d'algèbres simples de centre \(\Omega\). En particulier, si \(\Omega\) est algébriquement clos, alors \(A_\Omega\) devient un anneau de matrices de degré \(m\). L'indice est donc égal au ``nombre absolu de composantes''. En effet, si \(\mathsf{Z}\) est un corps de décomposition de \(A\) et si \(\overline\Omega\) est le composé de \(\Omega\) et de \(Z\), alors \(A_{\overline\Omega}\) devient anneau de matrices sur \(\overline\Omega\), donc sans radical et de centre \(\overline\Omega\). Ainsi \(A_\Omega\) est lui aussi sans radical et son centre est \(\Omega\), car un élément de \(A_\Omega\) distinct de \(\Omega\) qui serait central se trouverait aussi dans le centre de \(A_{\overline\Omega}\), donc dans un corps. \(A_\Omega\) est donc une algèbre simple sur \(\Omega\), et certainement un anneau de matrices sur \(\Omega\) lorsque \(\Omega\) est algébriquement clos.

\paragraph{4. Existence de corps de décomposition séparables.}
Chaque classe \((A)\) possède des corps de décomposition séparables, même de tels corps plongeables dans \(A\) elle-même.\footnote{Voir G. Köthe, Über Schiefkörper mit Unterkörpern zweiter Art über dem Zentrum, Journ. f. Math. 166 (1932), p. 182--184. La présente preuve, beaucoup plus simple, remonte à une remarque de M. Zorn.}

Soit le corps de base \(\mathsf{P}\) de caractéristique \(p\), et soit l'indice
\[
        m=s p^f
\]
avec \(s\) premier à \(p\). Si alors \(\mathsf{Z}\) est un corps de décomposition de \(A\) de degré \(m\), et \(\mathsf{Z}_0\) l'extension de première espèce qu'il contient, donc
\[
        (Z:Z_0)=p^f,
\]
alors l'indice de \(A_{\mathsf{Z}_0}\sim D\) devient, d'après 3., égal à \(p^f\). Il faut montrer l'existence d'un corps d'extension séparable \(\mathsf{Z}_1\) de \(\mathsf{Z}_0\) tel que l'indice de \(D_{\mathsf{Z}_1}\) devienne un diviseur propre de \(p^f\). Or \(D_\Omega\), avec \(\Omega\) algébriquement clos, est sans radical d'après 3.; le discriminant réduit de \(D\) ne s'annule donc pas. Il existe donc au moins un élément \(d\) de \(D\) avec trace réduite non nulle; mais \(d\) ne peut pas déjà appartenir au corps de base \(\mathsf{Z}_0\), puisque la trace de tout élément \(\alpha\) de \(\mathsf{Z}_0\) devient \(p^f\alpha\), donc s'annule. Ainsi \(\mathsf{Z}_1=\mathsf{Z}_0(d)\) devient un corps d'extension propre, et même séparable; l'indice de \(D_{\mathsf{Z}_1}\) devient un diviseur propre de \(p^f\). La répétition un nombre fini de fois conduit donc à un corps de décomposition séparable de \((A)\), dont le degré \(m\) coïncide avec l'indice de \(A\).

\clearpage
\providecommand{\Gcal}{\mathcal G}
\providecommand{\Sp}{\operatorname{Sp}}

\subsection*{\S 8. Corps de décomposition, corps de scission et théorie galoisienne dans le commutatif}

Pour passer des corps de décomposition des systèmes simples sur leur centre à ceux de systèmes simples arbitraires, il reste encore à développer la théorie de décomposition du centre; on y ajoutera une théorie galoisienne parallèle au \S 6, 3. Tous les corps et systèmes qui interviennent dans ce paragraphe sont commutatifs.

\paragraph{1. Corps de décomposition et corps de scission des systèmes simples commutatifs.}
Que le système commutatif \(Z\) soit simple sur \(\mathsf{P}\), donc un corps, et que \(\Omega\) soit un corps d'extension algébriquement clos de \(\mathsf{P}\). Alors on sait (cf. \emph{Darstellungstheorie}, \S 21): \(Z_\Omega\) reste système sans radical si et seulement si \(Z\) est séparable, donc extension de première espèce, sur \(\mathsf{P}\). Car alors et seulement alors il possède autant d'applications isomorphes du premier degré dans \(\Omega\) que son rang sur \(P\), donc aussi autant de représentations différentes, qui coïncident ici avec les classes de représentations.

Le fait qu'un radical puisse apparaître dans \(Z_\Omega\), donc qu'il n'y ait pas nécessairement décomposition en somme directe de composantes absolument simples, conduit aux définitions suivantes: un corps d'extension \(\Lambda\) de \(\mathsf{P}\) s'appelle corps de décomposition si \(Z_\Lambda\) se décompose en facteurs de composition du premier degré; autrement dit, si toutes les représentations absolument irréductibles se trouvent déjà dans \(\Lambda\). Un corps d'extension \(\mathsf{T}\) de \(\mathsf{P}\) s'appelle corps de scission si \(Z_{\mathsf{T}}\) scinde au moins un facteur de composition du premier degré; autrement dit, s'il existe au moins une représentation absolument irréductible de \(Z\) dans \(\mathsf{T}\).

Les corps de décomposition et de scission se caractérisent par le

\paragraph{Théorème.}
Un corps \(\mathsf{T}\) est corps de scission si et seulement s'il contient un sous-corps \(Z'\) isomorphe à \(\mathsf{Z}\); un corps \(\Lambda\) est corps de décomposition si et seulement s'il contient un sous-corps \(\Gamma\) isomorphe au corps galoisien appartenant à \(\mathsf{Z}\).

Cela résulte directement de la définition des corps de décomposition et de scission au moyen des représentations absolument irréductibles. \paragraph{Conséquence.}
On en déduit en particulier, par analogie avec le non-commutatif, qu'un corps \(Z'\) est corps de scission minimal, c'est-à-dire qu'aucun sous-corps propre n'est corps de scission, si et seulement s'il est isomorphe à \(Z\) (autrement dit, si un plongement commutatif irréductible de \(Z'\) dans \(Z\) en fait un sous-corps commutatif maximal de \(Z\)). Un corps de scission minimal est en même temps corps de décomposition si et seulement si \(Z\) est normal, donc corps galoisien sur \(\mathsf{P}\).

C'est la raison pour laquelle on appelle une algèbre simple normale sur son centre. Le fait qu'il n'y ait, dans un \(\Omega\) algébriquement clos fixé sur \(\mathsf{P}\), qu'un seul corps de décomposition minimal, le corps galoisien appartenant à \(\mathsf{Z}\), par opposition aux cas non commutatifs où il y en a en général une infinité non isomorphes,\footnote{R. Brauer--E. Noether, loc. cit.} s'explique par l'unicité de la décomposition en somme directe dans le commutatif, par opposition à l'unicité seulement à isomorphie près comme modules à opérateurs dans le non-commutatif; ou, ce qui revient au même, par le nombre fini de représentations absolument irréductibles différentes, au lieu des représentations différentes en nombre infini du non-commutatif, lesquelles appartiennent pourtant à la même classe (cf. 2.).

\paragraph{2. Fondation hypercomplexe de la théorie galoisienne dans le cas des corps séparables \(Z/\mathsf{P}\).}
En analogie exacte avec le \S 6, 3., la théorie des isomorphismes vaut ici pour des \(Z\) arbitraires séparables sur \(P\); dans le cas d'un \(Z\) galoisien, on peut alors passer au groupe usuel d'automorphismes.

\paragraph{Hypothèses.}
Soit le système simple \(Z\) sur \(\mathsf{P}\) un corps d'extension fini séparable; soit \(\Omega\) un corps de décomposition fini ou infini sur \(\mathsf{P}\), \(\mathsf{Z}'=\mathsf{Z}^{(1)}\) un sous-corps isomorphe à \(Z\), et \(\Gamma\) le corps galoisien correspondant. D'après 1., on a donc la décomposition en somme directe
\[
        Z_\Omega=\mathfrak r^{(1)}+\ldots+\mathfrak r^{(n)}
        =e^{(1)}Z_\Omega+\ldots+e^{(n)}Z_\Omega
        =e^{(1)}\Omega+\ldots+e^{(n)}\Omega .
\]
La représentation \(Z\to\mathsf{Z}^{(i)}\), avec \(Z^{(i)}\leqq\Gamma\), engendrée par \(e^{(i)}\), est donc définie par
\[
        e^{(i)}z=e^{(i)}\zeta^{(i)}.
\]

\paragraph{Lemme 1.}
Les \(n\) représentations engendrées par \(e^{(1)},\ldots,e^{(n)}\) sont toutes distinctes. La preuve est comme au \S 6, 3., ou encore d'après 1.; car elles appartiennent à des classes de représentations différentes.

\paragraph{Lemme 2.}
Si \(T\) est un corps intermédiaire entre \(\mathsf{P}\) et \(Z\), et si \(s\) est le rang de \(Z\) sur \(T\), alors tout isomorphisme de \(T\) dans \(\Omega\) admet au moins, et par conséquent exactement, \(s\) prolongements.

La preuve est comme au \S 6, 3. Que le ``au moins'' soit ici renforcé en ``exactement'' résulte de la décomposition univoque de \(Z_\Omega\), d'après laquelle \(Z\) ne possède pas au moins, mais exactement, \(n\) isomorphismes dans \(\Omega\) (cf. à ce sujet les remarques à la fin de 1.).

Comme au \S 6, 3., on définit maintenant la division en classes \(\mathfrak I_T\), induite par \(T\), de l'ensemble fini \(\mathfrak I\) des isomorphismes de \(Z\): tous et seulement les isomorphismes de \(\mathfrak I\) qui induisent le même isomorphisme sur \(T\) sont considérés comme équivalents dans \(\mathfrak I_T\). On démontre le

\paragraph{Théorème principal.}
Si \(\mathfrak I_T\) est la division en classes induite par \(T\), alors \(T\) est sous-corps maximal relativement à cette division en classes.

À partir de là, dans le cas d'un \(Z\) galoisien, on peut passer exactement comme au \S 6, 3., par composition avec l'inverse d'un isomorphisme, au groupe d'automorphismes et à la formulation usuelle. L'assertion correspondante pour les sous-groupes résulte de la considération des idempotents, qui est intéressante en elle-même.

\paragraph{3. Les idempotents de \(Z_\Omega\).}
Soit \(S\) l'une des \(n\) substitutions qui transportent \(\mathsf{Z}\) dans les divers \(\mathsf{Z}^{(i)}\), c'est-à-dire la composée des isomorphismes \(\mathsf{Z}\to Z\) et \(Z\to\mathsf{Z}^{(i)}\), le premier désignant l'isomorphisme réciproque de \(Z\to\mathsf{Z}\). Il n'est pas nécessaire que \(Z\) soit galoisien. Pour les éléments \(a\in Z_Z\), on définit les substitutions \(S\) comme opérateurs en prescrivant qu'elles induisent l'identité sur \(Z\); pour chaque \(a\), le conjugué \(a^S\in Z_{\mathsf{Z}^{(i)}}\) est ainsi défini. Avec ces prescriptions vaut le

\paragraph{Théorème.}
Les \(n\) idempotents \(e^{(i)}\) de \(Z_\Omega\) sont conjugués:
\[
        e^{(i)}=e^S,\qquad e=e^{(1)};
\]
ils se trouvent dans les \(n\) systèmes d'extension conjugués \(Z_{\mathsf{Z}^{(i)}}\).

En effet, l'application de \(S\) transporte les relations de définition
\[
        ez=e\zeta^{(1)},\qquad e^2=e
\]
en
\[
        e^S z=e^S\zeta^{(i)},\qquad (e^S)^2=e^S .
\]
À cause de l'unicité de la décomposition, les idempotents sont eux aussi déterminés univoquement; donc \(e^S=e^{(i)}\). De plus, comme \(Z\) est un corps de scission et porte déjà la représentation \(ez=e\zeta\), on a déjà \(e\in Z_Z\), et par conséquent \(e^{(i)}\in Z_{Z^{(i)}}\).

Si \(Z\) est spécialement galoisien, il en résulte immédiatement la partie du théorème principal relative aux sous-groupes; car si \(\mathfrak H\) est un sous-groupe du groupe galoisien \(\Gcal\), alors, à cause de la détermination univoque des composantes d'une somme directe, l'élément
\[
        \sum_{H\in\mathfrak H} e^H
\]
admet toutes les substitutions de \(\mathfrak H\), mais aucune autre. Puisque le groupe \(\Gcal\) était défini pour \(\mathsf{Z}\), il s'agit de la théorie galoisienne de \(\mathsf{Z}\); par l'isomorphie, celle de \(Z'\) est ainsi traitée en même temps.

\paragraph{4. Rapport des idempotents avec les bases complémentaires.}
On ne suppose de nouveau pas nécessairement que \(Z\) soit galoisien.

\paragraph{Théorème.}
Si \(a_1,\ldots,a_n\) est une base de \(Z\) sur \(\mathsf{P}\), et si l'idempotent
\[
        e=a_1\beta_1+\ldots+a_n\beta_n
\]
est posé, donc
\[
        e^S=a_1\beta_1^S+\ldots+a_n\beta_n^S,
\]
alors
\[
        \beta_1^S,\ldots,\beta_n^S
\]
signifie l'image de la base complémentaire \(b_1,\ldots,b_n\) de \(a_1,\ldots,a_n\) dans l'application fournie par \(e^S\).

En effet, pour l'application fournie par \(e^S\), \(e^Sz=e^S\zeta^S\), en particulier \(e^Sa_i=e^S\alpha_i^S\); à cause de
\[
        e^S\cdot e^S=e^S,\qquad e^S\cdot e^R=0\quad(S\ne R)
\]
on a les correspondances \(e^S\to1\), \(e^R\to0\). Il vient donc
\[
        1=\alpha_1^S\beta_1^S+\ldots+\alpha_n^S\beta_n^S;\qquad
        0=\alpha_1^S\beta_1^R+\ldots+\alpha_n^S\beta_n^R\quad(S\ne R).
\]
Écrit en matrices, c'est l'équation de définition des bases complémentaires:
\[
\begin{pmatrix}
\alpha_1 & \ldots & \alpha_n\\
\ldots & \ldots & \ldots\\
\alpha_1^S & \ldots & \alpha_n^S\\
\ldots & \ldots & \ldots
\end{pmatrix}
\cdot
\begin{pmatrix}
\beta_1 & \ldots & \beta_1^R & \ldots\\
\ldots & \ldots & \ldots & \ldots\\ \ldots & \ldots & \ldots & \ldots\\
\beta_n & \ldots & \beta_n^R & \ldots
\end{pmatrix}
=
E.
\]
Le passage à la définition par traces résulte du fait que, pour
\[
        a_i=\sum_S e^S\alpha_i^S,\qquad b_i=\sum_S e^S\beta_i^S
\]
la relation matricielle, après permutation des matrices, devient
\[
        \mathit{Sp}(a_i b_i)=1;\qquad \mathit{Sp}(a_i b_k)=0\quad\text{pour }i\ne k\text{\srcfnmark{21)}}.
\]
\srcfntext{21)}{Voir \emph{Darstellungstheorie}, \S 25.}
Les \(a_i\) et \(b_i\) appartiennent à \(Z\), puisqu'ils admettent toutes les substitutions \(S\) (lemme du \S 3, 3.).\srcfn{22)}{Des relations entre les bases complémentaires et les idempotents se trouvent pour la première fois chez Dedekind, \emph{Zur Theorie der aus \(n\) Haupteinheiten gebildeten komplexen Größen}; voir le t. II des \emph{Gesammelten Werke}, p. 4--5.}

La contragrédience des bases complémentaires résulte ici du fait que les idempotents \(e^S\) sont des invariants de \(Z\), indépendants de la base sous-jacente, donc du fait que tous les
\[
        a_1\beta_1^S+\ldots+a_n\beta_n^S
\]
se transforment en eux-mêmes lors du passage à une autre base
\[
        \bar a_1,\ldots,\bar a_n
\]
de \(\mathsf{Z}/\mathsf{P}\).

\clearpage
\subsection*{\S 9. Corps de décomposition et corps de scission de systèmes arbitraires}

\paragraph{1. Définitions et réduction au cas simple.}
Aux définitions données au \S 8 correspondent, dans le cas général, les définitions suivantes: que \(S\) soit hypercomplexe sur \(\mathsf{P}\). Un corps d'extension commutatif \(\Lambda\) de \(P\) s'appelle corps de décomposition de \(S\) si, dans une suite de composition par idéaux unilatères, \(S_\Lambda\) se décompose en facteurs de composition absolument simples; autrement dit, si toutes les représentations irréductibles de \(S\) dans \(\Lambda\) sont déjà absolument irréductibles. Un corps d'extension \(\mathsf{T}\) de \(\mathsf{P}\) s'appelle corps de scission si \(S_{\mathsf{T}}\) scinde au moins un facteur de composition absolument simple, donc si au moins une représentation irréductible dans \(\mathsf{T}\) est déjà absolument irréductible.

Ces définitions contiennent manifestement celles qui ont été données au \S 8 pour le commutatif, et au \S 7 pour les algèbres simples; pour ces dernières, parce que \(A_\Lambda\), pour toute extension commutative du centre, est complètement réductible, de sorte que les facteurs de composition et les composantes de la décomposition en somme directe coïncident. Les anneaux complets de matrices sur le centre, et eux seuls, donnent ici des composantes absolument simples, donc aussi des représentations absolument irréductibles. Comme en outre \(A_\Lambda\) est simple bilatère, donc se décompose en composantes isomorphes comme modules à opérateurs, la scission d'une composante absolument simple donne déjà la décomposition complète: corps de scission et corps de décomposition coïncident.

Des définitions résultent immédiatement les faits suivants: les corps de décomposition de \(S\) sont donnés comme corps composés d'un corps de décomposition pour chacune des représentations irréductibles dans \(\mathsf{P}\); est corps de scission de \(S\) tout corps de scission d'une représentation irréductible dans \(\mathsf{P}\). En raison de la correspondance biunivoque entre les systèmes simples, composantes de l'anneau quotient par le radical, et les représentations irréductibles dans \(\mathsf{P}\), il suffit donc de considérer les systèmes simples. Ici les résultats s'obtiendront par combinaison des \S 7 et \S 8.

\paragraph{2. Corps de décomposition et corps de scission des systèmes simples.}
Nous considérons d'abord le cas d'un système simple \(A\) dont le centre \(Z\) est séparable sur \(\mathsf{P}\). Il résulte des \S 8, 1. et \S 8, 3. que la décomposition bilatère de \(A_\Omega\) correspondant à la décomposition en somme directe de \(Z_\Omega\), où \(\Omega\) désigne un corps de décomposition de \(Z\), donne \(n\) composantes conjuguées, isomorphes à \(A\),
\[
        e^{(i)}A
\]
de \(A_\Omega\), qui deviennent des algèbres simples sur leur centre
\[
        e^{(i)}Z^{(i)}=e^{(i)}Z
\]
Ici les substitutions \(S\) sont définies comme opérateurs pour \(A_{\mathsf{Z}}\) par la stipulation qu'elles induisent l'identité sur \(A\).

En effet, d'après le \S 8, 3., les idempotents sont conjugués; il en va donc de même des composantes \(e^{(i)}A\), par définition des substitutions pour \(A\). Comme \(A\) est simple bilatère, \(e^{(i)}A\) est isomorphe à \(A\) comme anneau. De plus on a \(e^{(i)}A=e^{(i)}A_{\mathsf{Z}^{(i)}}\), puisque \(e^{(i)}Z=e^{(i)}\mathsf{Z}^{(i)}=e^{(i)}Z_{\mathsf{Z}^{(i)}}\); le centre coïncide donc avec le domaine des coefficients, et les composantes sont des algèbres normales.

Les résultats du \S 7, 2. donnent alors encore ceci: si \(A\) est un système simple sur \(\mathsf{P}\) dont le centre \(Z\) est séparable sur \(\mathsf{P}\), alors \(A_\Omega\) est aussi sans radical, donc complètement réductible; ici \(\Omega\) peut désigner un corps d'extension quelconque, fini ou infini. Car, d'après le \S 7, 2., \(e^{(i)}A_{\mathsf{Z}^{(i)}}\) reste simple lors de toute extension des coefficients, donc sans radical. Il en va donc de même de la somme directe; or celle-ci est \(A_\Omega\), ou bien elle provient de \(A_\Omega\), si \(\Omega\) n'est pas un corps de décomposition de \(Z\), par extension des coefficients.

Le \S 7, 2 donne également la

\paragraph{Caractérisation des corps de scission et des corps de décomposition.}
Si \(A\) est une algèbre à division de centre séparable \(Z\), alors les plongements irréductibles des corps de scission sont donnés par tous, et seulement par, les sous-corps commutatifs maximaux de \(A_f\) qui contiennent \(Z\); les corps de décomposition sont tous, et seulement, les corps composés d'un corps de scission avec un corps de décomposition du centre, en particulier avec le corps galoisien appartenant au centre. Un corps de scission minimal peut être corps de décomposition, même sans que le centre soit galoisien.

Car un corps de scission doit contenir un corps isomorphe à \(Z\); l'assertion de plongement sur les corps de scission résulte maintenant des faits précédents et du \S 7, 2. Un corps de décomposition doit en outre contenir un corps de décomposition \(\Omega\) de \(Z\). Mais tout corps de scission au-dessus de \(\Omega\) est en même temps corps de décomposition, puisque les composantes \(e^{(i)}A_\Omega\) sont simples et possèdent un centre isomorphe à \(\Omega\). Si \(Z\) est galoisien, les corps minimaux de scission et de décomposition coïncident donc. Mais même dans le cas non galoisien il peut y avoir, contrairement au cas commutatif, des corps de scission minimaux qui sont en même temps corps de décomposition, à savoir dès qu'un tel corps de scission minimal contient le corps galoisien appartenant à \(Z\). Exemple: un corps de quaternions ayant un centre isomorphe à \(\mathsf{P}(\sqrt[3]{2})\), \(\mathsf{P}\) corps des nombres rationnels; le corps galoisien appartenant à \(\mathsf{P}(\sqrt[3]{2})\) est corps minimal de scission et de décomposition.

S'il s'agit finalement d'un centre inséparable, alors \(Z_\Omega\), et par suite \(A_\Omega\), devient un système avec radical. Soient par exemple \(\mathfrak C\) le radical de \(Z_\Omega\) et \(\mathfrak c\) l'idéal d'extension dans \(A_\Omega\). Alors, pour \(Z_\Omega/\mathfrak C\), on a une décomposition en somme directe en composantes conjuguées, en analogie exacte avec le \S 8, 2.; cela donne exactement comme ci-dessus la décomposition en somme directe de \(A_\Omega/\mathfrak c\), de telle sorte que les composantes \(\overline e^{(i)}A\) sont isomorphes à \(A\), avec centre \(\overline e^{(i)}\mathsf{Z}^{(i)}\). Ainsi les théorèmes sur les corps de scission et de décomposition restent valables. Le dénombrement des rangs montre en outre que les facteurs de composition correspondant à une suite de composition de \(\mathfrak C\) deviennent isomorphes à \(A_\Omega/\mathfrak c\) comme modules à opérateurs, et non seulement homomorphes.

Formulés en termes de représentations irréductibles, ces résultats donnent donc le résumé suivant: si l'on a une représentation irréductible dans \(\mathsf{P}\) d'un système hypercomplexe, alors cette représentation est absolument complètement réductible si et seulement si le centre correspondant est séparable sur \(\mathsf{P}\). Dans tous les cas, les corps de décomposition et de scission de la représentation peuvent se caractériser par plongement dans le système simple correspondant, comme ci-dessus. En particulier, le plus petit degré d'un corps de scission est égal au produit du degré du centre par l'indice de la classe d'algèbres correspondante sur le centre; le degré de tout corps de scission en est un multiple. La représentation se décompose, au sens des suites de composition, en autant de classes de représentations absolument irréductibles distinctes mais conjuguées que vaut le degré du sous-corps séparable maximal du centre. Chaque classe apparaît autant de fois que vaut le produit de l'indice par le degré du centre relativement à cette extension séparable.

Pour un corps de base parfait, où il n'y a que des corps d'extension séparables, nous retrouvons ainsi des résultats connus de I. Schur, avec l'ajout de la caractérisation par plongement, qui se trouve déjà chez R. Brauer.

\begin{center}
(Reçu le 8 juin 1932.)
\end{center}

\clearpage

\setcounter{footnote}{0}
\section*{\NoetherNavAnchor{P41}{41. Le théorème du genre principal pour les corps de nombres relativement galoisiens}41. Le théorème du genre principal pour les corps de nombres relativement galoisiens}

\begin{center}
Par\\[0.4ex]
\textbf{Emmy Noether à Göttingen.}\\[0.8ex]
\emph{Math. Ann. 108 (1933), p. 411--419}
\end{center}

Dans la période récente, il est apparu que les méthodes non commutatives, en particulier la théorie des algèbres, permettent de formuler et de démontrer des théorèmes connus sur les corps de nombres relativement cycliques et relativement abéliens pour des corps relativement galoisiens arbitraires.\footnote{Voir à ce sujet mon exposé de Zurich, \emph{Hyperkomplexe Systeme in ihren Beziehungen zur kommutativen Algebra und Zahlentheorie}, dans les \emph{Verhandlungen des Internationalen Mathematiker-Kongresses Zürich 1932}, vol. I, où sont notamment discutées en détail les formulations hypercomplexes du théorème de norme et du théorème du genre principal; on y trouvera également les références bibliographiques. Pour le théorème sur les algèbres décomposées, voir en outre H. Hasse, \emph{Die Struktur der R. Brauerschen Algebrenklassengruppe} (6.1), Math. Ann. 107 (1932).} Je rappelle surtout le théorème de norme, qui s'énonce généralement comme un théorème sur les algèbres décomposées. Je montre ci-dessous que, de façon correspondante, le théorème du genre principal peut lui aussi se formuler, en général, de manière hypercomplexe: comme un théorème sur les automorphismes intérieurs ou encore sur les représentations croisées. La démonstration se déduit du théorème mentionné sur les algèbres décomposées par des considérations purement algébrico-arithmétiques, tandis que, dans le premier théorème, le théorème de norme dans le cas cyclique --- le degré premier suffit --- demeure le noyau transcendant.

Pour mieux faire comprendre la formulation, je commence par rappeler le théorème du genre principal au sens minimal, qui ne sera pas utilisé dans la suite,\footnote{Je dis ``au sens minimal'', parce qu'il s'agit de l'analogue algébrique du théorème du genre principal dans des corps arbitraires, tandis que ``au petit sens'' a déjà la signification du passage au corps de base \(p\)-adique.} un théorème algébrique élémentaire,\footnote{Le théorème se trouve chez A. Speiser, \emph{Zahlentheoretische Sätze aus der Gruppentheorie}, p.~3, Math. Z. 5 (1919), p. 1--6; il y est déjà énoncé comme théorème sur les représentations croisées.} qui, dans le cas cyclique particulier, devient le théorème 90 bien connu du \emph{Zahlbericht} de Hilbert: \(\Norm(a)\) est alors égal à un si et seulement si \(a=b^{1-S}\). J'énonce déjà ce théorème, lui aussi, comme un théorème sur les automorphismes intérieurs ou sur les représentations croisées, et je le démontre au moyen des théorèmes généraux sur les automorphismes intérieurs et sur la représentation des algèbres comme produits croisés.

Dans le théorème du genre principal proprement dit, la dernière construction est remplacée par le produit croisé du groupe des classes d'idéaux et du groupe de Galois, mais avec une division plus fine, induite, des systèmes de facteurs en classes d'idéaux; on saisit ainsi un analogue de la division en classes de rayon de la théorie du corps de classes. Comme on ne dispose pas encore ici de théorèmes généraux sur les automorphismes et les représentations, le théorème du genre principal peut être considéré comme un premier pas dans cette direction. Je donne donc, comme je viens de le dire, une réduction au moyen du théorème sur les algèbres décomposées. Le théorème est ainsi ramené au théorème correspondant pour les idéaux, au lieu des classes d'idéaux; ici les théorèmes de Brandt sur le rapport entre les ordres maximaux d'une algèbre\footnote{La théorie de Brandt des ordres maximaux est exposée et refondée chez H. Hasse, \emph{Über \(p\)-adische Schiefkörper und ihre Bedeutung für die Arithmetik hyperkomplexer Zahlsysteme}, Math. Ann. 104 (1931).} peuvent être regardés comme l'équivalent des théorèmes sur les automorphismes. Avec les considérations sur les ordres maximaux de produits croisés,\footnote{Une note de C. Chevalley (\emph{Hamburger Berichte}), une de H. Hasse et une de E. Noether ont été annoncées; les deux dernières doivent paraître dans un volume à la mémoire de Herbrand.} ils donnent une démonstration essentiellement parallèle à la démonstration au sens minimal, quoique plus compliquée parce qu'elle revient aux places individuelles. Je préfère donc la démonstration presque triviale indiquée par E. Artin, qui fournit ici un analogue encore plus simple de la démonstration de Speiser\textsuperscript{3)}.

\subsection*{\S{}1. Théorème du genre principal au sens minimal}

Dans ce paragraphe, \(k\) désigne un corps commutatif quelconque, pas nécessairement un corps de nombres; \(K/k\) est une extension séparable galoisienne de degré \(n\), et \(\Gg\) est le groupe de Galois correspondant.

\noindent\textbf{1.} Rassemblons d'abord les faits connus sur les produits croisés.\footnote{Cette théorie, à la suite d'un de mes cours (1929/30), est exposée chez H. Hasse, \emph{Theory of cyclic algebras}, chap.~2, Trans. Amer. Math. Soc. 34 (1932). Voir en outre un rapport de M. Deuring sur les nombres hypercomplexes, qui doit paraître dans les \emph{Ergebnisse der Mathematik}. Le présent résumé est tiré de mon exposé de Zurich cité à la note~1.} Le produit croisé signifie un plongement simultané de \(K\) et de \(\Gg\) dans une algèbre \(A\), de telle sorte que les automorphismes de \(K\) deviennent intérieurs. Si \(u_{S_1},\ldots,u_{S_n}\) sont des symboles associés aux \(n\) éléments du groupe, on pose d'abord \(A\), comme module de formes linéaires de rang \(n\) sur \(K\), sous la forme
\[
        A=u_{S_1}K+\ldots+u_{S_n}K.\tag{1}
\]
Par l'exigence que l'automorphisme intérieur soit engendré par \(u_S\), plus généralement par \(u_S K^*\),\footnote{\(K^*\) s'obtient de \(K\) en en retranchant zéro; cette notation sera employée partout.} \(A\) devient un anneau, donc une algèbre de rang \(n^2\) sur \(k\). Cette exigence s'exprime par
\[
        u_S^{-1}zu_S=z^S,\qquad\text{ou}\qquad
        z u_S=u_S z^S,                                      \tag{2}
\]
pour tout \(z\in K\), ainsi que par
\[
        u_Su_T=u_{ST}a_{S,T},\qquad a_{S,T}\in K^*,           \tag{3}
\]
et par la loi associative
\[
        a_{ST,R}a_{S,T}^{\,R}=a_{S,TR}a_{T,R}.                \tag{4}
\]
La loi associative équivaut ici à \([u_Su_T]u_R=u_S[u_Tu_R]\).
Si l'on définit le produit de deux éléments quelconques par
\[
 \Bigl(\sum_S u_S b_S\Bigr)\Bigl(\sum_T u_T c_T\Bigr)
   =\sum_{S,T}u_Su_T\,b_S^{\,T}c_T ,
\]
alors \(A\) devient le produit croisé de \(K\) avec \(\Gg\) pour le système de facteurs \(a_{S,T}\); on écrit
\[
        A=(a_{S,T},K,\Gg).
\]
On démontre que \(A\) est une algèbre simple normale sur \(k\), donc un anneau de matrices \(D_r\) de degré \(r\) sur l'algèbre à division associée \(D\), et que \(K\) est un sous-corps commutatif maximal, donc un corps de décomposition. Réciproquement, pour toute algèbre à division normale de cette sorte \(D\), il existe des anneaux de matrices \(D_r\) qui peuvent être engendrés de cette manière comme produits croisés.

Si l'on passe de \(u_S\) à \(v_S=u_Sc_S\), avec \(c_S\in K^*\), ce qui engendre le même automorphisme, on obtient des systèmes de facteurs associés
\[
        \bar a_{S,T}
        =a_{S,T}c_S^Tc_T/c_{ST} .                 \tag{5}
\]
En particulier, d'après (5), les systèmes de facteurs
\[
        \frac{c_S^{\,T}c_T}{c_{ST}}
\]
sont associés à un; ces grandeurs s'appellent grandeurs de transformation. Les systèmes de facteurs associés sont regroupés dans une classe \((a)\); les grandeurs de transformation \(c_S^{\,T}c_T/c_{ST}\) appartiennent donc en particulier à la classe \((1)\). De même, on regroupe dans une classe \(\mathfrak A\) toutes les algèbres semblables à \(A\), c'est-à-dire tous les \(D_r\), pour \(r=1,2,\ldots\). Les classes \(\mathfrak A\) et \((a)\) se correspondent biunivoquement. Les classes ayant pour corps de décomposition fixé \(K\) forment, pour la formation du produit direct, un groupe abélien isomorphe au produit terme à terme des classes de systèmes de facteurs. L'élément unité est respectivement la classe des algèbres décomposées et le système de toutes les grandeurs de transformation \(c_S^{\,T}c_T/c_{ST}\). Il s'agit du groupe de R. Brauer des classes d'algèbres.

Pour les algèbres cycliques, on obtient le cas particulier suivant: si \(Z/k\) est cyclique et si \(S\) est une substitution génératrice, les puissances de \(u\) correspondent aux puissances de \(S\), et l'on a
\[
        A=Z+uZ+\ldots+u^{n-1}Z,                              \tag{1'}
\]
\[
        zu=uz^S.                                               \tag{2'}
\]
\[
        u^n=\alpha.                                            \tag{3'}
\]
\[
        \alpha\in k^*.                                       \tag{4'}
\]
\[
        \bar\alpha=\alpha\,\Norm(c)
        \quad\text{lorsque l'on pose }v=uc\alpha                  \tag{5'}
\]
Chaque système de facteurs consiste donc ici en un seul élément du corps de base, avec la notation \(A=(\alpha,Z,S)\); la classe unité est donnée par les normes venues de \(Z^*\), et le groupe des classes d'algèbres est isomorphe à \(k^*/\Norm(Z^*)\).

\noindent\textbf{2.} La formulation du théorème du genre principal au sens minimal utilise le fait que les relations (2) à (5) sont purement multiplicatives et définissent donc en même temps le groupe d'extension \(\Gg^*\) de \(\Gg\), constitué des éléments des \(n\) complexes \(u_S K^*\):
\[
        \Gg^*=\{\,u_{S_1}K^*,\ldots,u_{S_n}K^*\,\}.            \tag{6}
\]
\(\Gg^*\) possède \(K^*\) comme sous-groupe normal abélien, et \(\Gg^*/K^*\simeq\Gg\).\footnote{Les relations (2) à (5) sont précisément les relations de définition d'une extension de groupes. Une telle extension existe dès qu'il existe un sous-groupe du groupe complet des automorphismes de \(K^*\), considéré seulement comme groupe multiplicatif, qui soit une image homomorphe de \(\Gg\). Dans le produit croisé, cet homomorphisme est un isomorphisme, et le sous-groupe en question est exactement l'ensemble des automorphismes à la fois multiplicatifs et additifs, c'est-à-dire le groupe de Galois de \(K\); on obtient ainsi, en plus de l'extension de groupes, l'extension d'anneaux.} \(\Gg^*\) est formé de la totalité des éléments réguliers \(g^*\) qui transforment \(K\) en lui-même, donc tels que \(g^{*-1}Kg^*=K\). Ainsi les automorphismes d'anneau de \(K\) sont engendrés par tous les éléments de \(\Gg^*\), et seulement par eux. En effet, chacun de ces \(g^*\) engendre un automorphisme d'anneau de \(K\). Réciproquement, tout automorphisme de cette sorte est engendré par la transformation au moyen d'un élément \(v=u_Sw\), où \(w\) induit l'identité sur \(K\). Comme \(w\) commute alors avec tous les éléments de \(K\), il appartient à \(K\), donc à \(K^*\), puisque \(K\) est un sous-anneau commutatif maximal.

\paragraph{Théorème du genre principal au sens minimal.}
Première formulation: tout automorphisme de groupe de \(\Gg^*\) qui prolonge l'automorphisme identique de \(K^*\) est intérieur et est engendré par un élément de \(K^*\).

Deuxième formulation: si les grandeurs de transformation
\[
        \frac{c_S^{\,T}c_T}{c_{ST}}
\]
ont toutes la valeur un, alors les \(c_S\) sont des puissances symboliques \((1-S)\)-ièmes; autrement dit, il existe un \(b\in K^*\) tel que
\[
        c_S=b^{1-S}=\frac{b}{b^S}\qquad(S\in\Gg).
\]
Troisième formulation: le groupe \(\Gg\) ne possède qu'une seule classe de représentations croisées du premier degré dans \(K^*\) appartenant au système de facteurs un.

Une représentation \(u_S\to C_S\) s'appelle croisée de système de facteurs \(a_{S,T}\) si
\[
        C_S^{\,T}C_T=C_{ST}a_{S,T}
\]
est vérifié. Deux représentations appartiennent à la même classe si
\[
        C_S=B^{-S}D_SB .
\]

Je démontre la première formulation, puis je montre qu'elle coïncide avec la deuxième et la troisième. La démonstration repose sur le fait que tout automorphisme de groupe du type indiqué se prolonge en un automorphisme d'anneau de \(A\), lequel est, comme on le sait, un automorphisme intérieur. Soient \(v_S\) les images des \(u_S\) par cet automorphisme. Alors \(\sum v_SK\) est de nouveau le produit croisé de \(\Gg\) avec \(K\) pour le même système de facteurs; en effet, par hypothèse, toutes les relations (2) à (4) restent valables. La correspondance
\[
        \sum u_Sb_S\to \sum v_Sb_S
\]
définit un automorphisme d'anneau qui fixe les éléments de \(K\). Il est donc engendré par un élément de \(K^*\).

Le passage à la deuxième formulation repose sur ceci: d'après (2), les \(v_S\) doivent être de la forme \(v_S=u_Sc_S\); puisque les systèmes de facteurs sont eux aussi conservés, les grandeurs de transformation correspondantes doivent, d'après (5), devenir égales à un. Réciproquement, (2) à (5) montrent que la substitution \(v_S=u_Sc_S\) engendre un automorphisme du type indiqué. La première formulation donne donc
\[
        v_S=b^{-1}u_Sb=u_S\,\frac{b}{b^S},
        \qquad c_S=b^{1-S}.
\]
Dans le cas cyclique, l'hypothèse signifie en particulier \(\Norm(c)=1\), et \(c=b^{1-S}\) est le théorème connu.\footnote{La démonstration donnée ici vaut en particulier lorsque \(k\) n'a qu'un nombre fini d'éléments, cas où la démonstration de Hilbert échoue, tandis que celle de Speiser reste valable.} La troisième formulation est immédiatement équivalente à la deuxième, car l'hypothèse dit que \(u_S\to c_S\) forme une représentation croisée du premier degré à système de facteurs un; \(c_S=b^{1-S}\) signifie que la représentation est équivalente à la représentation unité, \(c_S\sim1\).

Il faut toutefois remarquer que cette troisième formulation n'est qu'un cas particulier du théorème beaucoup plus général suivant: pour un système de facteurs donné \(a_{S,T}\), il n'existe qu'une seule classe de représentations croisées irréductibles; elle est de degré \(m\), où \(m\) est l'indice de Schur de l'algèbre correspondante, c'est-à-dire le degré de l'algèbre à division associée.\footnote{I. Schur, \emph{Einige Bemerkungen zu der vorstehenden Arbeit des Herrn A. Speiser}, Math. Z. 5 (1919), p. 7--10. Pour une démonstration fondée sur les modules de représentation croisée, issus de mon cours cité à la note~6, voir le rapport de M. Deuring également cité à la note~6.}

\subsection*{\S{}2. Le théorème du genre principal}

\paragraph{1. Remarques préliminaires sur la division en classes.}
Dans le théorème du genre principal de la théorie du corps de classes relatif aux corps relativement cycliques, deux divisions en classes différentes interviennent, comme on le sait: les classes d'idéaux absolues dans le corps supérieur, et les classes de rayon dans le corps inférieur. Dans le cas général, cela correspond au fait qu'une division en classes des idéaux du corps supérieur induit, pour les systèmes de facteurs, une division plus fine en classes d'idéaux.

Notons d'abord \(\Ii\) le groupe de tous les idéaux de \(K\). Comme \(\Gg\) engendre des automorphismes de \(\Ii\), l'extension par \(\Gg\) est définie par les relations (2) à (5) et se compose des \(n\) complexes \(\{\,\cdots u_S\Ii\cdots\,\}\). Si, dans \(\Ii\), on passe par équivalence au groupe \(\bar\Ii\) des classes d'idéaux absolues, les \(n\) complexes \(u_S\bar\Ii\) sont encore définis. En effet, \(\Gg\) engendre aussi des automorphismes du groupe des classes, puisque la classe principale se transforme en elle-même.\footnote{Plus généralement, \(\{\,\ldots,u_S\bar\Ii,\ldots\,\}\) est défini dès que l'on prend pour classe principale un sous-groupe de \(\Ii\) invariant par \(\Gg\); les grandeurs de transformation provenant des idéaux de ce sous-groupe remplacent alors celles qui proviennent des idéaux principaux.} Ainsi la relation d'équivalence qui mène de \(\Ii\) à \(\bar\Ii\) se transporte de \(\{\,\cdots u_S\Ii\cdots\,\}\) à \(\{\,\cdots u_S\bar\Ii\cdots\,\}\): aux idéaux équivalents \(\mathfrak a\) et \(\mathfrak b\) correspondent les produits équivalents \(u_S\mathfrak a\) et \(u_S\mathfrak b\). En particulier les \(u_S\) deviennent équivalents à tous les produits \(u_S(c_S)\), où \((c_S)\) parcourt les idéaux principaux de \(\Ii\).

Si l'on exige maintenant que les relations de produit (3) soient univoques au sens de cette équivalence, cela veut dire que toutes les grandeurs de transformation provenant d'idéaux principaux,
\[
        \frac{(c_S)^{\,T}(c_T)}{(c_{ST})},
\]
deviennent équivalentes au système de facteurs unité. Autrement dit:

Dans le système des \(n\) complexes \(\{\,\cdots u_S\bar\Ii\cdots\,\}\), où \(\bar\Ii\) est le groupe des classes d'idéaux absolues de \(K\) et où \(H\) désigne en particulier la classe principale, les relations de produit (3) pour les \(u_SH\) sont définies de façon univoque si et seulement si, dans la division induite en classes d'idéaux des systèmes de facteurs, la classe principale contient toutes les grandeurs de transformation provenant d'idéaux principaux, c'est-à-dire dont les générateurs \((c_S)\) appartiennent à \(H\).

À la suite de ce fait, je donne, comme généralisation de la division en classes de rayon, la définition suivante.

\paragraph{Définition.}
La division induite en classes d'idéaux des systèmes de facteurs est fixée comme suit: la classe principale se compose de tous les idéaux principaux \((a_{S,T})\) pour lesquels, avec au moins un choix d'éléments de base \(a_{S,T}\), l'algèbre
\[
        (a_{S,T},K,\Gg)
\]
se décompose en toutes les places ramifiées, finies et infinies, de \(K\).

Ces idéaux principaux forment un groupe pour la composition des systèmes de facteurs; ils contiennent les grandeurs de transformation
\[
        \frac{(c_S)^{\,T}(c_T)}{(c_{ST})},
\]
car les systèmes de facteurs \(c_S^{\,T}c_T/c_{ST}\) engendrent des algèbres décomposées partout.

\paragraph{2. Formulation du théorème du genre principal.}
Le théorème est énoncé, conformément au théorème au sens minimal, sous trois formulations équivalentes.

Première formulation: si, en prenant pour base la division induite en classes d'idéaux des systèmes de facteurs, la substitution
\[
        v_S=u_S \bar c_S\text{\footnotemark}
\]
\footnotetext{\(\bar{\mathfrak c}_S\) désigne la classe d'idéaux de \(\mathfrak c_S\).}
donne un automorphisme des \(n\) complexes \(\{\,\cdots u_S\bar\Ii\cdots\,\}\), c'est-à-dire si, pour les \(v_S\), les mêmes relations (3) subsistent au sens de cette division en classes, alors cet automorphisme est intérieur et est engendré par une classe d'idéaux \(\bar{\mathfrak b}\).

Deuxième formulation: si les grandeurs de transformation formées à partir des classes d'idéaux \(\bar{\mathfrak c}_S\),
\[
        \frac{\bar c_S^{\,T}\bar c_T}{\bar c_{ST}}
\]
appartiennent toutes à la classe principale de la division induite en classes des systèmes de facteurs --- l'ensemble de ces vecteurs \(\{\ldots\bar{\mathfrak c}_S\ldots\}\) de classes d'idéaux forme le genre principal ---, alors les classes \(\bar{\mathfrak c}_S\) sont des puissances symboliques \((1-S)\)-ièmes; c'est-à-dire qu'il existe une classe d'idéaux \(\bar{\mathfrak b}\) telle que
\[
        \bar c_S=\bar b^{1-S}=\frac{\bar b}{\bar b^S}
        \qquad(S\in\Gg).
\]

Troisième formulation: en prenant pour base la division induite en classes d'idéaux des systèmes de facteurs, le groupe \(\Gg\) ne possède qu'une seule classe de représentations croisées du premier degré dans \(\bar\Ii\) qui appartienne à la classe unité des systèmes de facteurs.

L'équivalence de ces formulations se déduit comme dans le \S{}1. Le passage de la première à la deuxième formulation est même plus simple ici, puisque l'automorphisme est déjà supposé engendré par \(v_S=u_S\bar{\mathfrak c}_S\). Cette hypothèse est nécessaire, car \(\bar\Ii\) n'a plus à être un sous-groupe commutatif maximal: il se peut que toutes les classes de \(\bar\Ii\) soient ambiguës. Pour la troisième formulation, il faut remarquer que, dans \(\bar\Ii\), seule la composition multiplicative est définie; les matrices de représentation admises sont donc uniquement celles qui n'ont qu'un seul élément non nul dans chaque ligne et dans chaque colonne. Pour les représentations du premier degré, cela n'impose aucune restriction.

\paragraph{3. Démonstration du théorème du genre principal.}
La démonstration de la deuxième formulation repose sur deux lemmes.

\paragraph{Lemme 1 (théorème du genre principal pour les idéaux).}
Si les grandeurs de transformation formées à partir des idéaux \(\mathfrak c_S\),
\[
        \frac{\mathfrak c_S^{\,T}\mathfrak c_T}{\mathfrak c_{ST}}
\]
donnent l'idéal unité, alors les \(\mathfrak c_S\) sont des puissances symboliques \((1-S)\)-ièmes:
\[
        \mathfrak c_S=\mathfrak b^{1-S}\qquad(S\in\Gg).
\]
Cela équivaut de nouveau aux première et troisième formulations; dans la première, comme pour les classes d'idéaux, il faut supposer que l'automorphisme est engendré par \(v_S=u_S\mathfrak c_S\). Démonstration:\footnote{Voir à ce sujet la fin de l'introduction.} \(\mathfrak b=\sum_S\mathfrak c_S\) a la propriété voulue, la somme étant la somme de modules, donc le plus grand diviseur commun. En effet,
\[
        \mathfrak b=\sum_S \mathfrak c_S,\qquad
        \mathfrak b^T=\sum_S \mathfrak c_S^{\,T},\qquad
        \mathfrak b^T\mathfrak c_T
        =\sum_S \mathfrak c_S^{\,T}\mathfrak c_T
        =\sum_S \mathfrak c_{ST}=\mathfrak b,
\]
d'où \(\mathfrak c_T=\mathfrak b^{1-T}\).

\paragraph{Lemme 2 (autre formulation du théorème sur les algèbres décomposées).}
Si l'algèbre
\[
        A=(a_{S,T},K,\Gg)
\]
se décompose en toutes les places ramifiées finies et infinies de \(K\), et si les idéaux principaux \((a_{S,T})\) forment des grandeurs de transformation d'idéaux,
\[
        (a_{S,T})
        =\frac{\mathfrak c_S^{\,T}\mathfrak c_T}{\mathfrak c_{ST}},
\]
alors \(A\) se décompose partout, donc absolument. Les \((a_{S,T})\) deviennent alors des grandeurs de transformation d'idéaux principaux,
\[
        (a_{S,T})=\frac{(d_S^{\,T})(d_T)}{(d_{ST})}\text{\footnotemark}.
\]
\footnotetext{Plus précisément, on a \(a_{S,T}=d_S^{\,T}d_T/d_{ST}\). Pour la démonstration du théorème du genre principal, on n'utilise cependant que l'énoncé affaibli portant sur les idéaux principaux.}
Réciproquement, toute algèbre décomposée partout satisfait cette condition.

La démonstration repose sur le comportement connu de \(A\) lors d'une extension \(\mathfrak p\)-adique. Soit \(k_{\mathfrak p}\) l'extension \(\mathfrak p\)-adique de \(k\), et soit \(K_\PP\) l'extension \(\PP\)-adique correspondante de \(K\), où \(\PP\) désigne un idéal premier au-dessus de \(\mathfrak p\); soit \(A_{\mathfrak p}\) l'extension de \(A\) obtenue par passage de \(k\) à \(k_{\mathfrak p}\). Soit encore \(\mathfrak Z\) le groupe de décomposition de \(\PP\), et \(a_{\mathfrak Z}\) la partie correspondante du système de facteurs, c'est-à-dire l'ensemble des \(a_{S,T}\) pour lesquels \(S,T\in\mathfrak Z\). Alors
\[
        A_{\mathfrak p}\sim
        (a_{\mathfrak Z},K_\PP/k_{\mathfrak p},\mathfrak Z)\text{\footnotemark}.
\]
\footnotetext{Voir Hasse, cité à la note~6, §~14. Une démonstration un peu plus simple est la suivante. Soit \(Z\) le corps de décomposition correspondant à \(\PP\); alors \(k_{\mathfrak p}\) contient un corps isomorphe à \(Z\), et \(A_Z\sim(a_{\mathfrak Z},K/Z,\mathfrak Z)\). En effet, \(A_Z\) est semblable à l'ensemble des éléments de \(A\) qui commutent élément par élément avec \(Z\); voir par exemple B. L. van der Waerden, \emph{Moderne Algebra}, vol.~II, p.~210, ou mon travail sur l'algèbre non commutative à paraître dans la Math. Z. L'extension de \(A_Z\) à \(A_{\mathfrak p}\) donne le résultat. Le choix de l'un des idéaux premiers conjugués \(\PP\) est indifférent: les algèbres obtenues sont isomorphes et se déduisent les unes des autres par transformation avec \(u_T\), lorsque \(\Gg=\sum\mathfrak ZT\).}
Si \(\mathfrak p\) est en particulier non ramifié, donc si \(\mathfrak Z\) est cyclique, il s'agit de la représentation distinguée de \(A_{\mathfrak p}\) au moyen du corps d'inertie non ramifié \(K_\PP/k_{\mathfrak p}\).\footnote{Voir, par exemple, le travail de H. Hasse cité à la note~4, §~3.}

Il reste à montrer que, sous l'hypothèse du lemme 2, \(A_{\mathfrak p}\) se décompose aussi dans ce cas non ramifié. Pour les places infinies \(\mathfrak p\), on a ici \(K_\PP=k_{\mathfrak p}\), donc \(A_{\mathfrak p}\) se décompose. Pour \(\mathfrak p\) fini, tous les idéaux \(\mathfrak c_S\) deviennent principaux, disons \((c_S)\), dans \(K_\PP\). L'hypothèse sur \((a_{S,T})\) signifie donc
\[
        a_{S,T}=e_{S,T}\frac{c_S^{\,T}c_T}{c_{ST}},
\]
où les \(e_{S,T}\) sont des unités; ainsi le système de facteurs \(a_{\mathfrak Z}\) est associé au système de facteurs d'unités \(e_{S,T}\). Cette propriété subsiste lors du passage à la représentation cyclique normalisée: \(\alpha\) est donc une unité dans \(A_{\mathfrak p}=(\alpha,K_\PP/k_{\mathfrak p},R)\), où \(R\) engendre \(\mathfrak Z\). En effet, le fait que les \(e_{S,T}\) forment un système de facteurs s'exprime par
\[
        u_{R^i}u_{R^j}=u_{R^{i+j}}e_{R^i,R^j}.
\]
Si l'on pose \(u_R=u\), on obtient
\[
        u^i=u_{R^i}e_{R,R^{i-1}}e_{R,R^{i-2}}\ldots e_{R,R},
\]
et donc, puisque \(u_E=e_{E,E}\), si \(f\) est l'ordre de \(\mathfrak Z\),
\[
        u^f=\alpha=e_{E,E}e_{R,R^{f-1}}\ldots e_{R,R}.
\]
Comme, dans une extension non ramifiée \(K_\PP/k_{\mathfrak p}\), toutes les unités sont des normes\textsuperscript{16)}, il s'ensuit que \(A_{\mathfrak p}\) se décompose. Ainsi \(A\) se décompose partout, donc, d'après le théorème sur les algèbres décomposées, absolument.

Autrement dit, l'hypothèse implique que les \(a_{S,T}\) sont des grandeurs de transformation d'éléments,
\[
        a_{S,T}=\frac{d_S^{\,T}d_T}{d_{ST}},
        \qquad d_S\in K^*.
\]
Sous une forme affaiblie, cela signifie que les \((a_{S,T})\) deviennent des grandeurs de transformation d'idéaux principaux. Réciproquement, si \(A\) est décomposée partout, alors les idéaux principaux \((a_{S,T})\) deviennent eux aussi des grandeurs de transformation d'idéaux, puisque cela vaut à chaque place, et \(A\) se décompose en particulier aux places ramifiées. Il s'agit en fait d'une autre formulation de l'hypothèse du théorème sur les algèbres décomposées.

Les deux lemmes donnent immédiatement la démonstration du théorème du genre principal. Si \(\mathfrak c_S\) désigne un idéal dans la classe \(\bar{\mathfrak c}_S\), l'hypothèse sur les classes, compte tenu de la définition de la division induite en classes des systèmes de facteurs, dit précisément que les grandeurs de transformation des \(\mathfrak c_S\) deviennent des idéaux principaux \((a_{S,T})\) satisfaisant aux hypothèses du lemme 2. Il existe donc des idéaux principaux \((d_S)\) tels que les idéaux
\[
        \mathfrak d_S=\frac{\mathfrak c_S}{(d_S)},
\]
qui appartiennent eux aussi aux classes \(\bar{\mathfrak c}_S\), satisfassent à l'hypothèse du lemme 1: leurs grandeurs de transformation sont égales à l'idéal unité. Il existe donc, d'après le lemme 1, un idéal \(\mathfrak b\) tel que
\[
        \mathfrak d_S=\mathfrak b^{1-S}.
\]
Cela signifie que les classes \(\bar{\mathfrak c}_S\) sont des puissances symboliques \((1-S)\)-ièmes de la classe \(\bar{\mathfrak b}\).

Dans le cas cyclique, le théorème du genre principal devient le théorème connu dès que, en adoptant la normalisation, la classe principale de la division induite en classes est le groupe des idéaux principaux \((\alpha)\) pour lesquels \(\alpha\) devient un reste normique en toutes les places ramifiées.

\begin{center}
(Reçu le 27. 10. 1932.)
\end{center}

\clearpage

\setcounter{footnote}{0}
\section*{\NoetherNavAnchor{P42}{42. Produits croisés décomposés et leurs ordres maximaux}42. Produits croisés décomposés et leurs ordres maximaux}

\begin{center}
\emph{Actualités scientifiques et industrielles 148 (1934), p. 5--15}
\end{center}

Dans le cas le plus simple des produits croisés décomposés,\footnote{La théorie des produits croisés est exposée, à la suite d'un de mes cours, dans la section II de H. Hasse, \emph{Theory of cyclic algebras over an algebraic number field}, Trans. Amer. Math. Soc. 34 (1932). Je n'utilise ici que les premières notions fondamentales. La théorie de Brandt des ordres maximaux est exposée et entièrement refondée, dans le cadre de résultats plus étendus, chez H. Hasse, \emph{Ueber \(p\)-adische Schiefkörper und ihre Bedeutung für die Arithmetik hyperkomplexer Zahlsysteme}, Math. Ann. 104 (1931) (cité H.).} c'est-à-dire de ceux qui donnent des algèbres décomposées et possèdent donc des systèmes de facteurs associés à un, les relations peuvent être suivies de manière largement explicite. Cela vaut encore pour des corps de décomposition non galoisiens, c'est-à-dire pour des sous-corps commutatifs maximaux. Les faits algébriques simples qui relèvent de cette situation sont développés au \S{}1. Au \S{}2, en prenant pour corps de base un corps de nombres algébrique, je donne des représentations explicites de tous les ordres maximaux et de leurs idéaux au moyen des modules et des modules complémentaires du corps de décomposition sous-jacent \(k\).\footnote{Ces représentations explicites m'étaient connues depuis longtemps; l'impulsion des considérations suivantes est venue de la communication du théorème de H. Hasse cité à la note 4.}

Je donne en outre une division en domaines: on range dans un même domaine tous les ordres maximaux qui ont la même intersection avec \(k\). Ainsi les domaines correspondent bijectivement aux ordres de \(k\); à l'anneau des entiers correspond le domaine principal. Les ordres maximaux d'un domaine se transforment les uns dans les autres par des idéaux formés à partir des modules de l'ordre correspondant; en particulier, dans le domaine principal, il s'agit des extensions des idéaux de \(k\), c'est-à-dire des modules sur l'anneau des entiers. Au \S{}3, je passe aux algèbres simples arbitraires. En affinant la division en domaines, de telle sorte que non seulement l'intersection, mais aussi les composantes \(p\)-adiques des ordres maximaux aux places ramifiées, doivent coïncider, on peut transférer les théorèmes obtenus dans le cas décomposé. En particulier, les ordres maximaux d'un domaine principal se transforment alors eux aussi les uns dans les autres par transformation au moyen d'extensions des idéaux de \(k\).

Les théorèmes sur le domaine principal se trouvent chez Chevalley\footnote{C. Chevalley, \emph{Sur certains idéaux d'une algèbre simple}, Abh. Math. Sem. Hamburg, 1934.} et Hasse;\footnote{H. Hasse, \emph{Ueber gewisse Ideale in einer einfachen Algebra}, dans le travail qui précède celui-ci; I de la présente collection.} Chevalley montre aussi que, sans la division affinée relative aux places ramifiées, les théorèmes ne valent plus en général.

\subsection*{\S{}I. Unités matricielles des produits croisés décomposés et bases complémentaires des corps de décomposition}

\paragraph{1.}
Dans ce paragraphe, \(\Omega\) désigne un corps de base quelconque. Soit \(k/\Omega\) une extension séparable galoisienne de degré \(n\), soit \(\Gg\) son groupe de Galois avec les éléments \(\mathrm{S},\mathrm{T},\ldots\), et soient \(u_s,u_t,\ldots\) les opérateurs correspondants. On considère les produits croisés à système de facteurs un, donc les algèbres décomposées
\[
        K=\Gg\ast k=u_{S_1}k+\ldots+u_{S_n}k,\qquad
        u_Su_T=u_{ST},
\]
avec
\[
        z u_s=u_s z^s \qquad(z\in k).
\]
Si l'on pose
\[
        E_\Gg=\sum_{S\in\Gg}u_S,
\]
on obtient, à partir du système de facteurs un, les relations\footnote{Si la caractéristique de \(k\) ne divise pas \(n\), alors \(\frac1n\mathrm{E}_{\mathfrak g}\) est l'idempotent de la représentation unité. La forme normale de \(\mathrm{K}\) reste toutefois valable lorsque la caractéristique divise \(n\); l'unique hypothèse est que \(k/\Omega\) soit séparable.}
\[
        \mathrm{E}_{\mathfrak g}u_s=u_s\mathrm{E}_{\mathfrak g}=\mathrm{E}_{\mathfrak g}\tag{1}
\]
pour tout \(S\in\Gg\), c'est-à-dire que \(\mathrm{E}_{\mathfrak g}\) engendre la représentation unité de \(\Gg\), ainsi que
\[
        z\mathrm{E}_{\mathfrak g}=\sum_su_sz^s,\tag{2a}
\]
et
\[
        E_\Gg z E_\Gg
        =E_\Gg\sum_Su_Sz^S
        =E_\Gg\,\Sp(z),                                      \tag{2b}
\]
où \(\Sp(z)\) désigne la trace de \(z\) comme élément de \(k/\Omega\).

\paragraph{Théorème.}
Si \(a_1,\ldots,a_n\) est une base de \(k/\Omega\) et si \(\widetilde a_1,\ldots,\widetilde a_n\) est la base complémentaire correspondante, alors
\[
        c_{ik}=\widetilde a_iE_\Gg a_k
\]
donne un système d'unités matricielles de \(K\). \(K\) se représente comme produit de modules sous la forme
\[
        K=kE_\Gg k
          =k\Bigl(\sum_Su_S\Bigr)k
          =\sum_{i,k}(\widetilde a_i\mathrm{E}_{\mathfrak g}a_k)\,\Omega.
\]
En effet, d'après (2b), on obtient
\[
        c_{ij}c_{lk}
        =\widetilde a_i\mathrm{E}_{\mathfrak g}a_j\cdot\widetilde a_l\mathrm{E}_{\mathfrak g}a_k
        =\widetilde a_i\mathrm{E}_{\mathfrak g}a_k\cdot\Sp(a_j\widetilde a_l)
        =c_{ik}\cdot\Sp(a_j\widetilde a_l).
\]
C'est précisément la propriété des unités matricielles, en raison des équations de définition
\[
        \Sp(a_j\widetilde a_l)=\delta_{jl}
\]
des bases complémentaires.

\paragraph{Remarque 1.}
Si \(\mathrm{Z}\) est un corps d'extension commutatif quelconque de \(\Omega\), tout reste valable lorsque \(a_i\) désigne une base de \(k_{\mathrm Z}/\mathrm{Z}\) et que \(k_{\mathrm Z}\) devient une somme directe de corps. Les substitutions du groupe sont alors définies par la condition qu'elles induisent l'identité sur \(\mathrm{Z}\), conformément au fait que chaque \(u_s\) commute élément par élément avec \(Z\) dans \(K_Z\).

\paragraph{Remarque 2.}
Le fait que \(\mathrm{K}=k\mathrm{E}_{\mathfrak g}k\) résulte aussi immédiatement de ce que le membre de droite est un idéal non nul du système simple \(\mathrm{K}\).

\paragraph{2.}
Supposons maintenant que \(\Omega\) soit de caractéristique nulle. La représentation \(\mathrm{K}=k\mathrm{E}_{\mathfrak g}k\) de l'algèbre décomposée reste valable même lorsque \(k\) est un corps de décomposition non galoisien de \(\mathrm{K}/\Omega\). On le montre en passant au corps de décomposition galoisien. Soit \(\bar k\) le corps galoisien associé à \(k\), soit \(\bar\Gg\) son groupe, soient \(\bar u_s\) les opérateurs correspondants, et soit \(\bar{\mathrm K}=\bar k\mathrm{E}_{\bar{\mathfrak g}}\bar k\) un produit croisé décomposé. Le sous-corps \(k\) correspond au sous-groupe \(\Hh\) d'ordre \(\mathrm{h}\). On pose
\[
        \bar\Gg=\Hh\mathrm{S}_1+\ldots+\Hh\mathrm{S}_n
              =\mathrm{S}_1^{-1}\Hh+\ldots+\mathrm{S}_n^{-1}\Hh,
        \qquad
        \mathrm{E}_{\mathfrak h}=\frac1h\sum\bar u_{\mathfrak h},
\]
et
\[
        \mathrm{E}_{\mathfrak g}=\frac1h \mathrm{E}_{\bar{\mathfrak g}}
              =\sum_i u_{s_i},\qquad
        u_{s_i}=\mathrm{E}_{\mathfrak h}\bar u_{s_i}.
\]
Alors l'idempotent \(\mathrm{E}_{\mathfrak h}\) engendre en particulier la représentation unité de \(\Hh\), et \(\mathrm{E}_{\mathfrak g}\), comme \(\mathrm{E}_{\bar{\mathfrak g}}\), engendre la représentation unité de \(\bar\Gg\). Pour les \(\mathrm{E}_{\mathfrak g}\) et les \(u_{S_i}\) ainsi définis, les relations (1) valent d'un côté, et (2) entièrement. Ici \(z\) parcourt tous les éléments de \(k\), de sorte que \(z^s\) parcourt ceux du corps conjugué \(k^s\), lequel est, comme sous-corps de \(\bar k\), également un sous-corps de \(\bar{\mathrm K}\). On a d'abord
\[
        \mathrm{E}_{\mathfrak g}=\mathrm{E}_{\mathfrak g}\mathrm{E}_{\mathfrak h}=\mathrm{E}_{\mathfrak h}\mathrm{E}_{\mathfrak g}=\mathrm{E}_{\mathfrak h}\mathrm{E}_{\mathfrak g}\mathrm{E}_{\mathfrak h},\tag{3}
\]
\[
        z\mathrm{E}_{\mathfrak h}=\mathrm{E}_{\mathfrak h} z,                                      \tag{4}
\]
et par conséquent
\[
        zu_s=u_sz^s.\tag{5}
\]
En effet, (3) découle de (1), selon la définition de \(\mathrm{E}_{\mathfrak h}\), pour \(\mathrm{E}_{\bar{\mathfrak g}}\), donc aussi pour \(\mathrm{E}_{\mathfrak g}\); et (4) exprime que \(z\), comme élément de \(k\), est invariant sous les substitutions de \(\Hh\): \(zu_h=u_hz\). La relation (5) découle de (4) par multiplication à droite par \(\bar u_{s_i}\). On obtient maintenant (1) d'un côté sous la forme
\[
        E_\Gg u_{S_i}=E_\Gg,
        \qquad\text{car}\qquad
        E_\Gg u_{S_i}=E_\Gg E_\Hh\bar u_{S_i}
        =E_\Gg\bar u_{S_i}=E_\Gg .                         \tag{1'}
\]
La relation (2a) se déduit de (5) par sommation, puis (2b) découle de (1'). La démonstration de \(\mathrm{K}=k\mathrm{E}_{\mathfrak g}k\) se fait alors exactement comme au numéro 1: la preuve que les \(\widetilde a_i\mathrm{E}_{\mathfrak g}a_k\) sont des unités matricielles n'utilisait que (2b). Comme anneau complet de matrices sur \(\Omega\), \(\mathrm{K}\) contient tout corps de degré \(n\) comme sous-corps, donc en particulier aussi \(k\).

\paragraph{3.}
En vue du \S{}3, remarquons encore que le passage de \(\mathrm{K}\) à \(\bar{\mathrm K}\) est possible même lorsque \(\mathrm{K}\) ne se décompose pas.\footnote{On en déduit l'analogue de la représentation en produit croisé dans le cas non galoisien; ce point devait vraisemblablement être développé dans une dissertation.} À côté de \(k\), \(\bar k\) est alors lui aussi un corps de décomposition; \(\mathrm{K}\) est donc semblable à \(\bar{\mathrm K}\), le produit croisé de \(\bar\Gg\) avec \(\bar k\). Comme \(k\) est corps de décomposition, le sous-groupe \(\Hh\) correspondant à \(k\) a pour système de facteurs le système un; l'idempotent \(\mathrm{E}_{\mathfrak h}\) est donc défini comme ci-dessus. Il s'ensuit, en sens inverse, que \(\mathrm{K}\) est isomorphe à \(\mathrm{E}_{\mathfrak h}\bar{\mathrm K}\mathrm{E}_{\mathfrak h}\). Cette algèbre est isomorphe à l'anneau des endomorphismes de l'idéal à gauche \(\bar{\mathrm K}\mathrm{E}_{\mathfrak h}\), donc semblable à \(\bar{\mathrm K}\). Son rang coïncide avec celui de \(\mathrm{K}\): en effet, \(\bar{\mathrm K}\mathrm{E}_{\mathfrak h}\) est de rang \(n\) sur \(\bar k\), puisque les \(\bar u_{s_i}^{-1}\mathrm{E}_{\mathfrak h}\) sont indépendants sur \(\bar k\), tandis que \(\bar{\mathrm K}\) est de rang \(nh\). Cela signifie que \(\bar{\mathrm K}\) se décompose en \(h\) idéaux isomorphes à \(\bar{\mathrm K}\mathrm{E}_{\mathfrak h}\), ce qui, pour l'anneau des endomorphismes, revient à multiplier par l'algèbre matricielle complète de degré \(h\).

Dans le cas où \(\mathrm{K}\), et donc \(\bar{\mathrm K}\), se décompose, on retrouve ainsi le résultat du numéro 2:
\[
\begin{aligned}
\mathrm{E}_{\mathfrak h}\bar{\mathrm K}\mathrm{E}_{\mathfrak h}
 =\mathrm{E}_{\mathfrak h}\bar k\mathrm{E}_{\bar{\mathfrak g}}\bar k\mathrm{E}_{\mathfrak h}
 =(\mathrm{E}_{\mathfrak h}\bar k\mathrm{E}_{\mathfrak h})\mathrm{E}_{\bar{\mathfrak g}}(\mathrm{E}_{\mathfrak h}\bar k\mathrm{E}_{\mathfrak h})
 &=\Sp(\bar k/k)\mathrm{E}_{\mathfrak h}\mathrm{E}_{\bar{\mathfrak g}}\mathrm{E}_{\mathfrak h}\Sp(\bar k/k)\\
 &=k\mathrm{E}_{\mathfrak g}k.
\end{aligned}
\]

\subsection*{\S{}II. Les ordres maximaux des produits croisés décomposés et leur division en domaines}

Désormais \(\Omega\) est un corps de nombres algébrique fini,\footnote{Les considérations de ce paragraphe restent valables lorsque \(\Omega\) est un corps de fonctions en une indéterminée.} \(\mathrm{K}\) une algèbre décomposée de degré \(n\) sur \(\Omega\), et \(k\) un sous-corps commutatif maximal galoisien ou non galoisien, de sorte que \(\mathrm{K}=k\mathrm{E}_{\mathfrak g}k\). Soit \(\mathfrak a,\widetilde{\mathfrak a}\) un couple de \(o\)-modules finis complémentaires de \(k\), où \(o\) est l'anneau des entiers de \(\Omega\) et \(\oo\) l'anneau des entiers de \(k\). La structure des ordres maximaux de \(\mathrm{K}\), relativement à \(k\), est décrite par les théorèmes suivants.

\paragraph{Théorème 1.}
Le produit de modules
\[
        \mathfrak O_a=\widetilde a E_\Gg a
\]
est un ordre maximal de \(K\); en particulier
\[
        \mathfrak O=\mathfrak O_o=\widetilde oE_\Gg o
\]
est un ordre maximal.

\paragraph{Théorème 2.}
La transformation de \(\mathfrak O\) en \(\mathfrak O_{\mathfrak a}\) se fait au moyen de l'idéal à gauche
\[
        \mathfrak L=\widetilde oE_\Gg a
\]
et de l'idéal à droite réciproque
\[
        \mathfrak R=\widetilde a E_\Gg o
\]
relativement à \(\mathfrak O\).

\paragraph{Remarque.}
La rencontre des deux modules complémentaires \(\mathfrak a,\widetilde{\mathfrak a}\) dans le produit \(\mathfrak L\mathfrak R\) conduit précisément, par opposition au simple produit \(\mathfrak a\widetilde{\mathfrak a}\), à la relation de réciprocité, grâce à la formation de trace effectuée par \(\mathrm{E}_{\mathfrak g}\).

\paragraph{Théorème 3.}
Tous les idéaux à gauche et à droite relativement à \(\mathfrak O\) sont épuisés par ceux indiqués dans le théorème 2. Les ordres maximaux \(\mathfrak O_{\mathfrak a}\) épuisent donc l'ensemble de tous les ordres maximaux de \(\mathrm{K}\). Deux ordres maximaux quelconques \(\mathfrak O_{\mathfrak a}\) et \(\mathfrak O_{\mathfrak b}\) se transforment l'un dans l'autre par
\[
        \mathfrak L_a=\widetilde aE_\Gg b,
        \qquad
        \mathfrak R_a=\widetilde bE_\Gg a
\]
et ces \(\mathfrak L_{\mathfrak a}\) et \(\mathfrak R_{\mathfrak a}\) épuisent les idéaux à gauche et à droite relativement à \(\mathfrak O_{\mathfrak a}\).

\paragraph{Théorème 4.}
Lorsque \(\mathfrak O_{\mathfrak a}\) parcourt tous les ordres maximaux de \(\mathrm{K}\), l'intersection
\[
        [\mathfrak O_a,k]
\]
parcourt tous les ordres de \(k\); cette intersection est chaque fois égale à l'ordre de \(\mathfrak a\). Si l'on rassemble en un domaine tous les ordres maximaux appartenant au même ordre de \(k\), alors la transformation à l'intérieur d'un domaine se fait par des idéaux à gauche \(\widetilde{\mathfrak a}\mathrm{E}_{\mathfrak g}\mathfrak b\) relativement à \(\mathfrak O_{\mathfrak a}\) et par des idéaux à droite réciproques, où \(\mathfrak a\) et \(\mathfrak b\) sont des modules de l'ordre correspondant dans \(k\).

\paragraph{Théorème 4a.}
En particulier, la transformation des ordres maximaux du domaine principal, c'est-à-dire du domaine de tous les ordres maximaux contenant l'anneau des entiers \(\mathfrak o\) de \(k\), se fait par des idéaux
\[
        \widetilde aE_\Gg b,
\]
où \(\mathfrak a\) et \(\mathfrak b\) sont des idéaux de \(k\). Autrement dit: si \(\mathfrak L\) est un idéal à gauche d'un ordre maximal du domaine principal et si \(\mathfrak o\) est contenu dans l'ordre à droite de \(\mathfrak L\), alors \(\mathfrak L\) est l'extension d'un \(\mathfrak o\)-idéal.

Les démonstrations reposent sur le passage aux places individuelles. Il suffit chaque fois de passer à la localisation en \(p\), où \(p\) parcourt tous les idéaux premiers de \(\Omega\). On note \(o\) l'anneau des entiers de \(\Omega\), et \(o_p\) sa localisation en \(p\), c'est-à-dire l'anneau principal de tous les éléments \(p\)-entiers, représentables par des fractions dont le dénominateur est premier à \(p\). Si \(\mathfrak C\) est un \(o\)-module quelconque de \(\mathrm{K}\), on désigne comme composante \(\mathfrak C_p\) le module étendu \(\mathfrak C o_p\). Alors:

\paragraph{Lemme.}
\(\mathfrak C\) est l'intersection de tous les modules étendus \(\mathfrak C_p\). En effet, si \(w\) appartient à \(\mathfrak C_p\), il existe un \(\sigma\in\oo\), non divisible par \(p\), tel que \(\sigma w\in\mathfrak C\). Si \(w\) appartient à l'intersection de tous les \(\mathfrak C_p\), et si \(g\) désigne l'idéal de tous les \(\sigma\in\oo\) tels que \(\sigma w\in\mathfrak C\), alors \(g\) ne peut être divisible par aucun \(p\); donc \(g=o\).\footnote{Si l'on définit les places par extension \(p\)-adique, le lemme correspondant s'énonce ainsi: l'intersection de toutes les composantes \(p\)-adiques \(\mathfrak C^{(p)}\) avec \(\mathrm{K}\) est égale à \(\mathfrak C\). Cela résulte de \([\mathfrak C^{(p)},\mathrm{K}]=\mathfrak C_p\), comme on le voit à l'aide d'une \(o_p\)-base de \(\mathfrak C_p\) qui est aussi une base de \(\mathfrak C^{(p)}\). Dans H., théorème 66, ce lemme n'est énoncé que pour les idéaux, c'est-à-dire pour les modules de rang maximal.} Notons que le lemme ne suppose pas que le module ait rang maximal.

\paragraph{Conséquences.}
Chaque module est déterminé par ses composantes. Si \(\mathfrak D\) est un surmodule propre de \(\mathfrak C\), alors au moins un \(\mathfrak D_p\) est un surmodule propre de \(\mathfrak C_p\). En particulier, si \(\mathfrak M\) est un ordre de \(\mathrm{K}\) et si toutes ses composantes \(\mathfrak M_p\) sont des ordres maximaux sur \(o_p\), alors \(\mathfrak M\) est lui aussi un ordre maximal.

\paragraph{Démonstration du théorème 1.}
\(\mathfrak O_{\mathfrak a}\) est un ordre, car \(\mathfrak O_{\mathfrak a}\) est un \(o\)-module fini, puisque \(\mathfrak a\) et \(\widetilde{\mathfrak a}\) le sont; il est en outre de rang maximal. De plus
\[
        \mathfrak O_a^2
        =\widetilde aE_\Gg a\,\widetilde aE_\Gg a
        =\widetilde aE_\Gg a\,\Sp(\widetilde a a)\text{\srcfnmark{9)}}
        \subseteq \widetilde aE_\Gg a=\mathfrak O_a,
\]
\srcfntext{9)}{Pour tout module \(\mathfrak c\), \(\Sp(\mathfrak c)\) désigne l'idéal de \(o\) formé des traces de tous les éléments de \(\mathfrak c\).}
car \(\Sp(\widetilde a a)\subseteq\oo\). Pour montrer la maximalité de \(\mathfrak O_{\mathfrak a}\), et donc qu'il contient \(o\), il suffit, d'après la conséquence, de montrer la maximalité de chaque composante. Comme \(o_p\) est un anneau principal, \(\mathfrak a_p\) et \(\widetilde{\mathfrak a}_p\) possèdent des bases complémentaires indépendantes \(a_1,\ldots,a_n\) et \(\widetilde a_1,\ldots,\widetilde a_n\). Les produits \(\widetilde a_i\mathrm{E}_{\mathfrak g}a_k\) forment, d'après le \S{}1, des unités matricielles; ainsi
\[
        \widetilde a_pE_\Gg a_p
        =\sum_{i,k}(\widetilde a_iE_\Gg a_k)\oo_p
\]
est maximal.

\paragraph{Remarque.}
Puisque \(\mathfrak O_{\mathfrak a}\) est maximal, donc \(\mathfrak O_{\mathfrak a}^2=\mathfrak O_{\mathfrak a}\), on obtient encore \(\Sp(\widetilde{\mathfrak a}\mathfrak a)=o\). Cela découle aussi directement du fait que cette relation vaut en chaque place, comme le montrent les bases complémentaires, donc partout par le lemme.

\paragraph{Démonstration du théorème 2.}
La règle de trace donne:
\[
\begin{aligned}
\mathfrak O\mathfrak L
 &=\widetilde{\mathfrak o}\mathrm{E}_{\mathfrak g}\mathfrak o\,\widetilde{\mathfrak o}\mathrm{E}_{\mathfrak g}\mathfrak a
 =\widetilde{\mathfrak o}\mathrm{E}_{\mathfrak g}\mathfrak a=\mathfrak L,
&\qquad
\mathfrak R\mathfrak O
 &=\widetilde{\mathfrak a}\mathrm{E}_{\mathfrak g}\mathfrak o\,\widetilde{\mathfrak o}\mathrm{E}_{\mathfrak g}\mathfrak o
 =\widetilde{\mathfrak a}\mathrm{E}_{\mathfrak g}\mathfrak o=\mathfrak R,\\
\mathfrak L\mathfrak R
 &=\widetilde{\mathfrak o}\mathrm{E}_{\mathfrak g}\mathfrak a\,\widetilde{\mathfrak a}\mathrm{E}_{\mathfrak g}\mathfrak o
 =\widetilde{\mathfrak o}\mathrm{E}_{\mathfrak g}\mathfrak o=\mathfrak O,
&
\mathfrak R\mathfrak L
 &=\widetilde{\mathfrak a}\mathrm{E}_{\mathfrak g}\mathfrak o\,\widetilde{\mathfrak o}\mathrm{E}_{\mathfrak g}\mathfrak a
 =\widetilde{\mathfrak a}\mathrm{E}_{\mathfrak g}\mathfrak a=\mathfrak O\,.
\end{aligned}
\]

\paragraph{Démonstration du théorème 3.}
Soit \(\mathfrak L\) un idéal à gauche relativement à \(\mathfrak O\). Alors, en particulier,
\[
        \mathfrak L=\widetilde oE_\Gg\mathfrak L
        =(\widetilde oE_\Gg l_1,\ldots,\widetilde oE_\Gg l_t)
\]
avec \(l_1,\ldots,l_t\), par exemple, une \(o\)-base de \(\mathfrak L\). Si l'on pose \(l_\mu=\sum\limits_i u_{s_i}z_{\mu,i}\), où \(z_{\mu,i}\in k\),\footnote{Lorsque \(k\) n'est pas galoisien, il faut remplacer \(z\) par \(\bar z\mathrm{E}_{\mathfrak h}\), et par conséquent \(a\) par \(\bar a\mathrm{E}_{\mathfrak h}\). On obtient alors \(\displaystyle E_\Gg l=E_\Gg\bar aE_\Hh=E_\Gg a\), avec \(a\in k\), car les relations (3), (4) et (5) donnent \(\displaystyle kE_\Gg k=kE_\Gg kE_\Hh=\sum_i u_{S_i}\{k^{S_i},k\}E_\Hh\).} on obtient
\[
        E_\Gg l_\mu=E_\Gg\sum_i z_{\mu,i}=E_\Gg a_\mu,
\]
donc \(\mathfrak L=\widetilde{\mathfrak o}\mathrm{E}_{\mathfrak g}\mathfrak a\). Lorsque \(\mathfrak L\) est régulier, \(a\) doit avoir le rang maximal \(n\); alors \(\widetilde aE_\Gg o\) est l'idéal à droite réciproque. Par composition, \(\widetilde{\mathfrak a}\mathrm{E}_{\mathfrak g}\mathfrak b\) parcourt tous les idéaux à gauche relativement à \(\mathfrak O_{\mathfrak a}\), et \(\widetilde bE_\Gg a\) les idéaux à droite réciproques.

\paragraph{Démonstration du théorème 4.}
L'intersection \(t=[k,\mathfrak O_a]\) est un ordre, puisque \(\mathfrak O_{\mathfrak a}\) ne contient que des éléments entiers. Comme sous-ensemble de \(\mathfrak O_{\mathfrak a}\), \(t\) devient un domaine de multiplicateurs à droite et il est le plus grand ordre de \(o\) qui satisfasse cette condition; il contient donc certainement l'ordre \(o_a\) de \(a\). Pour montrer que \(t=o_a\), il suffit de prouver l'égalité des composantes. Soit \(t\) un élément de \(t_p\). Alors, pour tout \(c_{ik}=\widetilde a_i\mathrm{E}_{\mathfrak g} a_k\), le produit à droite \(c_{ik}t\) est une combinaison linéaire des \(c_{ij}\) à coefficients entiers, c'est-à-dire dans \(o_p\). D'autre part, \(a_kt\) est une combinaison linéaire des \(a_j\) à coefficients dans \(\Omega_p\). Par l'unicité de la représentation au moyen des \(c_{ij}\), ces coefficients appartiennent nécessairement à \(o_p\); ainsi \(t\) appartient à l'ordre de \(\mathfrak a_p\). Donc \(t\) est égal à l'ordre de \(a\). Que tout ordre de \(k\) puisse apparaître comme intersection résulte du fait que le module \(a\) parcourt en particulier tous les ordres. La transformation découle alors immédiatement du théorème 3.

\paragraph{Démonstration du théorème 4a.}
Le théorème est la spécialisation du théorème 4 au domaine principal. Si l'on part d'abord de \(\mathfrak O=\widetilde oE_\Gg o\) et si \(o\) est aussi contenu dans l'ordre à droite de \(\mathfrak L=\widetilde{\mathfrak o}\mathrm{E}_{\mathfrak g}\mathfrak a\), alors \(\widetilde aE_\Gg a\) appartient au domaine principal; donc \(a\) est un idéal de \(k\), et \(\mathfrak L=\mathfrak O a\). Si l'on part de \(\mathfrak O_{\mathfrak a}\), on obtient de même \(\mathfrak L=\widetilde aE_\Gg b\); \(b\), et donc aussi \(a^{-1}b\), est un idéal de \(k\). Il vient
\[
        \mathfrak L=\widetilde aE_\Gg aa^{-1}b
        =\mathfrak O_a a^{-1}b.
\]

\subsection*{\S{}III. Ordres maximaux de produits croisés arbitraires}

Les théorèmes déduits pour les produits croisés décomposés subsistent aussi dans le cas général, dès que l'on affine la division en domaines. Les faits relatifs au domaine principal valent alors exactement; les autres doivent être formulés comme des assertions sur les places individuelles. Ici, la place doit être comprise comme extension \(p\)-adique.

\paragraph{Définition.}
Soit \(K\) une algèbre simple normale sur \(\Omega\), et soit \(k\) un sous-corps commutatif maximal. On range dans un domaine relatif à \(k\) tous les ordres maximaux de \(K\) qui ont la même intersection avec \(k\) et qui possèdent en outre, aux places ramifiées de \(K\), les mêmes composantes \(p\)-adiques. Si l'intersection avec \(k\) est l'anneau des entiers, il s'agit d'un domaine principal.

Pour les domaines principaux, on a alors:

\paragraph{Théorème 1.}
Les ordres maximaux d'un domaine principal se transforment les uns dans les autres par transformation au moyen d'idéaux d'extension des idéaux de \(k\). Autrement dit: les idéaux à gauche, premiers à la différente, des ordres maximaux d'un domaine principal, qui contiennent \(o\) dans leur ordre à droite, sont des extensions d'idéaux de \(k\).

\paragraph{Théorème 2.}
Si \(K\) est de degré premier, il n'existe qu'un seul domaine principal; celui-ci se transforme en lui-même par des idéaux d'extension des idéaux de \(k\). Les idéaux correspondants d'un ordre maximal fixé \(\mathfrak O\) sont les idéaux de \(k\) étendus avec les idéaux bilatères de \(\mathfrak O\).

\paragraph{Démonstration.}
Deux ordres maximaux d'un domaine principal ne diffèrent qu'en un nombre fini de places; par définition, ce sont des places où \(K\) se décompose. Mais là \(K\) a la forme considérée dans les \S{}\S{}I--II. Cela est clair si \(k\) est galoisien: le système de facteurs, associé à un, peut alors être remplacé par le système un. D'après le \S{}I, nos 2 et 3, cela vaut aussi dans le cas non galoisien. Il s'agit donc, en chacune de ces places et par suite globalement,\footnote{Voir H., \S{}8, ou le travail de Hasse qui précède, cité à la note 4.} d'idéaux d'extension.

Si \(K\) est de degré premier, alors \(K\) est soit décomposée, soit une algèbre à division. Dans le second cas, elle reste une algèbre à division à toutes les places ramifiées et ne possède donc qu'un seul domaine principal, si bien que la transformation par idéaux d'extension vaut de nouveau. La transformation la plus générale s'obtient par composition des idéaux d'extension avec les idéaux qui transforment l'ordre maximal en lui-même, donc avec les idéaux bilatères; ceux-ci, comme on le sait, ne diffèrent des idéaux centraux qu'aux places ramifiées, donc des idéaux d'extension. Les autres assertions en découlent.

On a encore:

\paragraph{Théorème 3.}
Les ordres maximaux d'un domaine quelconque se transforment les uns dans les autres par des idéaux premiers à la différente qui, en chaque place, sont composés de modules de l'ordre correspondant, comme indiqué au \S{}II, théorème 3. Les composantes de ces modules, pour un idéal fixé, ne diffèrent de l'ordre qu'en un nombre fini de places.

Cela découle immédiatement du \S{}II; il faut seulement remarquer que les opérateurs \(u_s\) doivent être remplacés, aux différentes places, par des opérateurs équivalents correspondants.

La question de savoir s'il existe, pour tout ordre de \(k\), des ordres maximaux ayant cet ordre comme intersection demande ici une étude plus précise aux places ramifiées. Un domaine principal existe toujours, comme on le déduit de la représentation en produit croisé, éventuellement généralisée, qui conduit à un ordre contenant \(o\), dès que les systèmes de facteurs sont supposés entiers.\footnote{Une autre démonstration est donnée par Hasse et Chevalley, aux références des notes 4 et 3.}

Une autre question est de savoir dans quelle mesure il est possible de caractériser les corps de décomposition algébriques au moyen de la division arithmétique en domaines. Autrement dit: des corps de décomposition différents engendrent-ils toujours des divisions en domaines différentes, engendrent-ils déjà des domaines principaux différents? Les domaines principaux de tous les corps de décomposition épuisent-ils peut-être l'ensemble de tous les ordres maximaux?

Que, par un seul ordre maximal, une telle caractérisation ne soit certainement pas possible, les exemples les plus simples le montrent. Ainsi l'ordre maximal
\[
        \left[\,1,i,j,\frac{1+i+j+k}{2}\,\right]
\]
du corps de quaternions contient, à côté de l'anneau des entiers \([1,i]\), aussi l'anneau des entiers
\[
        \left[\,1,\frac{1+i+j+k}{2}\,\right]
\]
du corps des racines troisièmes de l'unité.

\begin{flushleft}
Göttingen, août 1932.
\end{flushleft}
\clearpage

\setcounter{footnote}{0}
\section*{\NoetherNavAnchor{P43}{43. Différentiation des idéaux et différente}43. Différentiation des idéaux et différente}

\begin{center}
\emph{Journal f. d. reine u. angew. Math. 188 (1950), p. 1--21}\\[0.75ex]
Par \emph{Emmy Noether} \(\dagger\)\srcfnmark{1)}.
\end{center}

\srcfntext{1)}{Le présent travail, laissé par E. Noether et rédigé durant l'hiver 1927/28, a été aimablement mis à la disposition de l'éditeur par H. Grell. Il donne un exposé détaillé de la conférence de même titre présentée au congrès de la D. M. V. à Prague en 1929; voir \emph{Jahresbericht D. M. V.} 39 (1930), p. 17, en italique. À partir du \S{} 6, 3, le manuscrit devait manifestement faire encore l'objet d'une révision plus précise. Le lecteur averti comprendra toutefois sans difficulté l'exposé esquissé, laissé ici inchangé à quelques retouches mineures près.}

\subsection*{Introduction}

Le théorème principal de la théorie de la ramification des corps de nombres algébriques affirme, comme on sait, que la différente (l'idéal de ramification) est divisible au moins par la \((\varrho-1)\)-ième puissance d'un idéal premier qui entre à la \(\varrho\)-ième puissance dans le nombre premier \(p\), exactement par la \((\varrho-1)\)-ième puissance lorsque \(\varrho\) n'est pas divisible par \(p\), et par une puissance plus élevée lorsque \(\varrho\) est divisible par \(p\).

Ce théorème est l'analogue du fait que le quotient différentiel \(f'(x)\) d'un polynôme \(f(x)\) est divisible au moins par la \((\varrho-1)\)-ième puissance d'un facteur linéaire qui entre à la \(\varrho\)-ième puissance dans \(f(x)\), exactement par la \((\varrho-1)\)-ième puissance lorsque \(\varrho\) n'est pas divisible par la caractéristique du domaine des coefficients, et par une puissance plus élevée lorsque \(\varrho\) est divisible par cette caractéristique.

Je montre dans ce qui suit qu'il s'agit ici de plus que d'une analogie formelle. \emph{La différente peut être conçue comme le quotient différentiel d'un idéal définissant du corps de nombres \(\mathsf{K}\), dans un anneau de polynômes entiers associé en \(x_1,\ldots,x_n\), le quotient différentiel étant pris au point \(x=\omega\), donc pour \(x_1=\omega_1,\ldots,x_n=\omega_n\), où \(\omega_1,\ldots,\omega_n\) est une base de l'anneau des entiers \(\mathfrak o\) de \(K\) comme \(\mathbb Z\)-module.} L'idéal définissant \(\mathfrak M\) correspond à l'ensemble des relations entre les \(\omega\); il se compose donc de tous les polynômes entiers \(f(x)\) pour lesquels \(f(\omega)\) s'annule. Le quotient différentiel \(\mathfrak M'[x\to\omega]\), quant à lui, est défini comme quotient des différences, pris au point \(x=\omega\).

En effet, puisque \(f(\omega)=0\), considérons l'idéal \(\mathfrak M\) comme formé de toutes les différences \(f(x)-f(\omega)\); formons ensuite l'idéal des différences \(\mathfrak P=(x_1-\omega_1,\ldots,x_n-\omega_n)\) et le quotient des différences \(\mathfrak A=\mathfrak M:\mathfrak P\), la formation du quotient désignant le quotient d'idéaux usuel, pris dans l'anneau de polynômes en les \(x\) dont les coefficients sont les entiers de \(\mathsf{K}\), donc les éléments de \(\mathfrak o\). La différente \(\mathfrak D\) est alors définie par
\[
\mathfrak D=\mathfrak A[x\to\omega].
\]
L'extension du domaine des coefficients est nécessaire pour que l'idéal des différences et le quotient des différences soient seulement définis. Il s'agit donc de la généralisation directe du quotient différentiel d'un polynôme en une indéterminée \(x\),
\[
f'(\xi)=(f(x)-f(\xi)):(x-\xi),
\]
pris au point \(x=\xi\); ici encore, l'anneau de polynômes doit être étendu par adjonction de \(\xi\) pour que les différences soient définies. En fait, le quotient différentiel d'idéal devient lui aussi le quotient différentiel polynomial dès que la base des \(\omega\) se compose des puissances d'un seul élément.

La définition de la différente donnée ici se rattache ainsi directement à des faits connus; elle introduit cependant des termes intermédiaires non invariants, en ce que l'idéal \(\mathfrak M\) dépend de la base choisie. On obtient une définition équivalente et \emph{entièrement invariante} en introduisant, au lieu de l'anneau de polynômes entiers, un anneau \(\mathfrak O\) isomorphe à l'ordre maximal \(\mathfrak o\) --- et donc isomorphe à l'anneau quotient par \(\mathfrak M\) --- puis en étendant son domaine des coefficients par \(\mathfrak o\) pour obtenir \(\mathfrak O_{\mathfrak o}\).\srcfn{2)}{L'extension du domaine des coefficients --- formation du produit direct --- est traitée au \S{} 1 aussi pour les anneaux non commutatifs, dans une généralité supérieure à celle qui est nécessaire aux fins du présent travail. Le \S{} 2, spécialisé aux anneaux commutatifs et aux bases finies, suffit à celles-ci.} L'idéal des différences \(\mathfrak P\) est alors engendré par toutes les différences \((x-\xi)\), où \(x\) et \(\xi\) se correspondent par l'isomorphisme de \(\mathfrak O\) sur \(\mathfrak o\); le quotient des différences \(\mathfrak A\) est le quotient de l'idéal nul par \(\mathfrak P\); la différente s'obtient à partir de \(\mathfrak A\) en remplaçant chaque \(x\) par le \(\xi\) qui lui est isomorphe.

Cette définition invariante sera placée en tête de ce qui suit: elle permet de définir directement la différente d'un anneau commutatif quelconque relativement à un sous-anneau, pourvu que soient satisfaites certaines conditions relatives à la possibilité d'étendre le domaine des coefficients. C'est par exemple toujours le cas lorsqu'il existe une base du module considéré, éventuellement après passage à certaines localisations. Ainsi se trouvent définies en particulier la différente d'un ordre quelconque d'un corps de nombres, la différente de l'anneau quotient d'un ordre, par exemple modulo une puissance d'un nombre premier, la différente relative et la différente dans le cas des fonctions algébriques de plusieurs indéterminées. La concordance de la définition différentielle de la différente avec la définition usuelle résulte d'une étude structurelle précise, fondée sur la décomposition en somme directe, de l'anneau d'extension galoisien associé au corps de nombres par formation d'un produit direct\srcfn{3)}{Il s'agit ici de poursuivre les considérations qui sous-tendent mon travail sur le discriminant, dont la connaissance n'est toutefois pas supposée. Voir E. Noether, \emph{Der Diskriminantensatz für die Ordnungen eines algebraischen Zahl- und Funktionenkörpers}, J. Math. 157 (1927), p. 82--104.} --- étude qui, dans le cas particulier de l'anneau quotient par un polynôme en une indéterminée, revient à analyser la formule d'interpolation de Lagrange. Plus précisément, voici ce que j'entends par là:

Introduisons un anneau \(\mathfrak K\), isomorphe au corps de nombres \(K\), qui contient \(\mathfrak O\), puis étendons son domaine des coefficients --- les nombres rationnels --- par le corps galoisien \(F\) de \(\mathsf{K}\). Le système étendu \(\mathfrak K_\Gamma\) devient la somme directe de \(n\) composantes simples correspondant aux \(n\) isomorphismes de \(\mathsf{K}\) (passage aux corps conjugués); ces composantes sont des extensions du quotient des différences \(\mathfrak A\) et de ses conjugués. Simultanément, l'idéal nul de \(\mathfrak K_\Gamma\) devient l'intersection des \(n\) idéaux déduits de l'idéal des différences et de ses conjugués. Les composantes simples sont de la forme \(\Gamma e^{(i)}\), où \(e^{(i)}\) désigne la composante de l'unité; puisque \(e^{(1)}\) passe à l'unité pour \(x=\xi\), le quotient des différences devient \(\mathfrak d e^{(1)}\), où \(\mathfrak d\) désigne la différente de \(\mathfrak o\). La différente \(\mathfrak d\) se compose de tous les éléments du corps qui, multipliés par \(e^{(1)}\), appartiennent à \(\mathfrak O_{\mathfrak o}\), produit direct de l'ordre avec lui-même. Il en va de même pour les conjugués.

Pour en déduire la définition au moyen du module complémentaire, il faut remarquer que \(e^{(i)}\) devient une forme linéaire dans la base \(x_1,\ldots,x_n\) de \(\mathfrak O\), respectivement de \(\mathfrak K\), qui correspond aux \(\omega\); or les coefficients des \(x\) dans les \(e^{(i)}\) donnent précisément une base du module complémentaire de \(\mathfrak o\), et de ses conjugués. Comme le quotient des différences \(\mathfrak A\) a ses coefficients dans \(\mathfrak o\), on en déduit aussitôt que la différente est le quotient de \(\mathfrak o\) par son module complémentaire, ce qui donne la définition de Dedekind. Les considérations qui précèdent ne sont pas limitées à l'ordre maximal; elles caractérisent la différente de tout ordre comme le quotient de cet ordre par son module complémentaire.

Dans le cas de l'ordre maximal, la concordance de la définition différentielle avec la définition au moyen de l'équation fondamentale résulte du fait mentionné plus haut: le quotient différentiel d'idéal devient le quotient différentiel polynomial dès que la base se compose des puissances d'un élément, ici de la forme fondamentale. Le lien conceptuel entre cette dernière définition et celle de Dedekind est ainsi également établi.

En prenant pour base l'équation fondamentale, on obtient directement, comme on sait, le théorème principal mentionné au début. Pour déterminer exactement les exposants, on considère la différente \(\mathfrak d[p^t]\) de l'anneau quotient de l'ordre maximal modulo \(p^t\), avec \(t\) suffisamment grand --- \(t=n\) suffit, ce qui est aussi nécessaire pour les corps de degré deux ---; il est essentiel que la différente \(\mathfrak d\), considérée modulo \(p^t\), passe exactement dans \(\mathfrak d[p^t]\). Puisque la somme directe des anneaux correspond à la somme directe des différentes, on peut de nouveau se restreindre à l'anneau quotient modulo \(\mathfrak p^{t\varrho}\). Or ce dernier est isomorphe à l'anneau quotient par un polynôme en une indéterminée à coefficients modulo \(p^t\); la différentiation de ce polynôme donne --- comme l'a montré Ö. Ore, qui est arrivé par une autre voie à un tel polynôme\srcfn{4)}{Note de l'éditeur: cette référence n'est pas complétée dans le manuscrit. Voir les indications du \S{} 7, qui n'est donné que sous forme d'esquisse.} --- une vue exacte de tous les exposants possibles à l'aide des nombres supplémentaires.

Enfin, l'étude structurelle donne encore le théorème connu selon lequel la différente d'un sur-corps de \(\mathsf{K}\) est égale au produit de la différente relative par la différente de \(\mathsf{K}\), comme conséquence du fait que la même propriété multiplicative est satisfaite par les unités de la décomposition en somme directe, et donc par les modules complémentaires.

Tous les faits indiqués valent de même pour les différentes relatives, par passage à certaines localisations, exactement comme je l'ai exposé dans le travail sur le discriminant pour les discriminants relatifs. Par passage au domaine fonctionnel, tout subsiste encore pour les corps de fonctions algébriques dont le corps des coefficients est de caractéristique zéro; en caractéristique \(p\), en revanche, et pour les fonctions algébriques entières, les propriétés structurelles demeurent valables --- pour l'essentiel même sans passage au domaine fonctionnel ---, mais la théorie de la ramification change par suite de l'apparition de corps d'extension inséparables, ce qui ne sera pas développé ici. Comme résultat principal, je voudrais désigner les théorèmes de structure entièrement invariants esquissés ci-dessus, valables aussi pour les fonctions algébriques de plusieurs indéterminées (\S{} 5 et \S{} 6); le passage au module complémentaire constitue déjà une résolution en identités entre coefficients et n'est effectué que pour rattacher ces résultats à la théorie connue.

\subsection*{\S{} 1. Extension du domaine des coefficients d'un anneau, ou formation du produit direct. Idéal définissant\srcfnmark{5)}}
\srcfntext{5)}{Je dois à B. L. van der Waerden des remarques critiques sur le présent paragraphe.}

Dans ce paragraphe, on considère des anneaux généraux, pour lesquels on ne suppose pas la loi commutative de la multiplication. En revanche, pour simplifier, on suppose l'existence d'un élément unité et ce, dans tout le travail, sauf mention expresse du contraire.

\paragraph{1. Définition du produit direct.}
Un anneau \(\mathfrak O\times\mathfrak o\), ou \(\mathfrak O_{\mathfrak o}\), est appelé \emph{produit direct} des anneaux \(\mathfrak O\) et \(\mathfrak o\) relativement à \(\mathfrak h\) --- ou encore anneau \emph{obtenu à partir de \(\mathfrak O\) par extension du domaine des coefficients \(\mathfrak h\) à \(\mathfrak o\)}\srcfn{6)}{La seconde expression, qui sera le plus souvent employée dans la suite, doit rappeler l'extension du domaine des coefficients; la première est la désignation usuelle du produit direct.} --- lorsque les conditions suivantes sont satisfaites:

\begin{enumerate}
\item[(1)] \(\mathfrak O_{\mathfrak o}\) contient \(\mathfrak O\) et \(\mathfrak o\) comme sous-anneaux et est égal au produit de ces anneaux (donc à l'intersection de tous les anneaux de \(\mathfrak O_{\mathfrak o}\) qui contiennent \(\mathfrak O\) et \(\mathfrak o\)).

\item[(2)] L'intersection de \(\mathfrak O\) et de \(\mathfrak o\) est \(\mathfrak h\), où \(\mathfrak h\) contient l'élément unité de \(\mathfrak O_{\mathfrak o}\).

\item[(3)] \(\mathfrak O\) et \(\mathfrak o\) sont permutables élément par élément; ainsi \(\mathfrak O_{\mathfrak o}\) se compose de l'ensemble des combinaisons bilinéaires
\[
x_{i_1}y_{i_1}+\cdots+x_{i_r}y_{i_r},
\]
où les \(x\), respectivement les \(y\), parcourent tous les systèmes finis d'éléments de \(\mathfrak O\), respectivement de \(\mathfrak o\); \(\mathfrak h\) est situé dans le centre de \(\mathfrak O\) et de \(\mathfrak o\), donc aussi dans celui de \(\mathfrak O_{\mathfrak o}\).

\item[(4)] Deux éléments de \(\mathfrak O_{\mathfrak o}\) sont égaux seulement si --- et naturellement toujours si --- ils deviennent formellement égaux d'au moins une manière en vertu des relations qui existent dans \(\mathfrak O\), d'une part, et dans \(\mathfrak o\), d'autre part. Il faut entendre par là ce qui suit. Si
\[
x_1y_1+\cdots+x_ny_n=0
\quad\hbox{dans }\mathfrak O_{\mathfrak o},
\]
cette expression peut, par l'adjonction de sommes de combinaisons formellement nulles \((\gamma h)\delta-\gamma(h\delta)\), où \(h\) appartient à \(\mathfrak h\), \(\gamma\) à \(\mathfrak O\) et \(\delta\) à \(\mathfrak o\), être transformée en
\[
(z_1\alpha_1+\cdots+z_r\alpha_r)+(t_1\beta_1+\cdots+t_s\beta_s)
\]
avec \(z_1=0,\ldots,z_r=0\) et \(\beta_1=0,\ldots,\beta_s=0\), où les \(\alpha,\beta\) sont dans \(\mathfrak o\) et les \(z,t\) dans \(\mathfrak O\). Ainsi les relations \(z=0\) sont des relations entre les \(x,\gamma,h\) dans \(\mathfrak O\), et les relations \(\beta=0\) des relations entre les \(y,\delta,h\) dans \(\mathfrak o\).\srcfn{7)}{Si l'on ne suppose pas l'existence de l'élément unité, des termes linéaires supplémentaires apparaissent encore dans (3) et (4).}
\end{enumerate}

Deux éléments de \(\mathfrak O_{\mathfrak o}\) sont donc égaux seulement si leur différence s'annule de la manière indiquée.

La condition de produit direct fixe ainsi de manière univoque les relations d'égalité et les lois de composition (addition et multiplication) à partir de celles des facteurs. Il en résulte ceci: si \(\overline{\mathfrak O},\overline{\mathfrak o}\) et leur intersection \(\overline{\mathfrak h}\) sont isomorphes à \(\mathfrak O,\mathfrak o\) et \(\mathfrak h\), et si le produit direct \(\mathfrak O\times\mathfrak o\) existe, alors \(\overline{\mathfrak O}\times\overline{\mathfrak o}\) existe également et est isomorphe à \(\mathfrak O\times\mathfrak o\).

Tout anneau \(\mathfrak S\) qui est égal au produit de \(\overline{\mathfrak O}\) et de \(\overline{\mathfrak o}\), de telle sorte que \(\overline{\mathfrak O}\) et \(\overline{\mathfrak o}\) soient permutables élément par élément, est une image homomorphe de \(\mathfrak O\times\mathfrak o\), l'application de \(\overline{\mathfrak O}\) sur \(\mathfrak O\), respectivement de \(\overline{\mathfrak o}\) sur \(\mathfrak o\), restant isomorphe. En effet, l'égalité dans \(\mathfrak O\times\mathfrak o\) est comprise dans celle de \(\mathfrak S\), et coïncide, sur \(\overline{\mathfrak O}\), respectivement sur \(\overline{\mathfrak o}\), avec celle de \(\mathfrak O\), respectivement de \(\mathfrak o\), dans \(\mathfrak O\times\mathfrak o\).\srcfn{8)}{Voir aussi \S{} 1, 4.}

\paragraph{2. Conditions nécessaires et suffisantes d'existence du produit direct relativement à un sous-anneau.}
Soient \(\mathfrak O\) et \(\mathfrak o\) des anneaux dont l'intersection est l'anneau \(\mathfrak h\), situé dans le centre de chacun d'eux. Le produit direct relativement à \(\mathfrak h\) ne doit pas nécessairement exister, comme le montre l'exemple suivant. Soit \(\mathfrak h\) l'anneau des entiers rationnels, \(\mathfrak O=\mathfrak h[x]\) avec \(1-2x=0\), et \(\mathfrak o=\mathfrak h[y]\) avec \(1-2y=0\), où \(y\) désigne un nouveau symbole. Ainsi \(\mathfrak O\) et \(\mathfrak o\) sont des anneaux équivalents relativement à \(\mathfrak h\), dont l'intersection ensembliste est \(\mathfrak h\), tandis qu'aucune loi de composition n'est donnée entre les symboles \(x\) et \(y\). Il n'existe aucun anneau \(R\) contenant \(\mathfrak O\) et \(\mathfrak o\) tel que, dans \(R\) --- où des lois de composition existent donc maintenant entre \(x\) et \(y\) ---, leur intersection soit \(\mathfrak h\). En effet, on aurait dans \(R\)
\[
0=(1-2x)y-x(1-2y)=y-x,
\]
sans même avoir utilisé la permutabilité de \(x\) et de \(y\).\srcfn{9)}{La localisation d'un anneau commutatif dans un anneau fixé --- ici par adjonction de \(\frac12\) --- est précisément unique. C'est l'analogue du fait que deux extensions galoisiennes équivalentes et distinctes ne peuvent être plongées dans un même sur-corps. Il existe toutefois ici, par formation du produit direct, un plongement dans un anneau à diviseurs de zéro.}

La \emph{condition nécessaire et suffisante d'existence du produit direct} est la suivante:

\emph{Il doit exister au moins un anneau \(R\) contenant \(\mathfrak O\) et \(\mathfrak o\), tel que l'intersection de \(\mathfrak O\) et de \(\mathfrak o\) soit \(\mathfrak h\) et que \(\mathfrak O\) et \(\mathfrak o\) deviennent permutables élément par élément.}

Si le produit direct existe, il est lui-même un tel anneau \(R\). Réciproquement, pour construire le produit direct, on peut supposer --- et cette hypothèse sera toujours faite dans la suite --- qu'il n'existe aucune loi de composition entre les éléments de \(\mathfrak O\) et de \(\mathfrak o\) qui n'appartiennent pas à \(\mathfrak h\), car on peut toujours y parvenir en passant à une extension équivalente de \(\mathfrak h\), c'est-à-dire en introduisant d'autres symboles.\srcfn{9a)}{Voir \S{} 2, 1 et la note 12).} Supposons en outre que \(\mathfrak h\) contienne l'élément unité \(e\) de \(\mathfrak O\) et de \(\mathfrak o\) (voir toutefois la note 7)). Définissons maintenant \(\mathfrak T\) comme l'ensemble de toutes les combinaisons bilinéaires
\[
x_{i_1}y_{i_1}+\cdots+x_{i_r}y_{i_r}
=y_{i_1}x_{i_1}+\cdots+y_{i_r}x_{i_r},
\]
où les \(x\) et les \(y\) parcourent tous les systèmes finis d'éléments de \(\mathfrak O\), respectivement de \(\mathfrak o\), et où les \(x\) et les \(y\) deviennent donc permutables dans \(\mathfrak T\). Deux de ces combinaisons sont posées égales si, et seulement si, elles deviennent formellement égales en vertu des relations dans \(\mathfrak O\) et dans \(\mathfrak o\), au sens précisé en 1, (4). Définissons en outre la somme et le produit de
\[
\alpha_{i_1}x_{i_1}+\cdots+\alpha_{i_r}x_{i_r}
\quad\hbox{et}\quad
\beta_{i_1}x_{i_1}+\cdots+\beta_{i_r}x_{i_r},
\]
où des zéros peuvent aussi figurer parmi les \(\alpha,\beta\), par
\[
(\alpha_{i_1}+\beta_{i_1})x_{i_1}+\cdots+(\alpha_{i_r}+\beta_{i_r})x_{i_r}
\quad\hbox{et}\quad
\sum\alpha_{i_\lambda}\beta_{i_\mu}x_{i_\lambda}x_{i_\mu}.
\]
Cette définition est univoque au sens de l'égalité, car toute relation nulle (au sens de 1, (4)), multipliée à droite ou à gauche par un élément de \(\mathfrak o\) ou de \(\mathfrak O\), est encore une telle relation; il en va de même de la somme et de la différence de ces relations nulles. Le système \(\mathfrak T\) forme un anneau contenant \(\mathfrak O\) et \(\mathfrak o\), qui satisfait aux conditions du produit direct quant à l'égalité, à l'addition et à la multiplication. \(\mathfrak T\) satisfait aussi à la condition d'intersection, par suite de l'hypothèse générale de plongement. En effet, soit \(\mathfrak S\) le produit de \(\mathfrak O\) et de \(\mathfrak o\) dans \(R\); d'après 1, \(\mathfrak S\) est une image homomorphe de \(\mathfrak T\), isomorphe relativement à \(\mathfrak O\) et à \(\mathfrak o\). Une relation \(x=y\) dans \(\mathfrak T\), avec \(x\) dans \(\mathfrak O\) et \(y\) dans \(\mathfrak o\), entraîne donc la même relation dans \(\mathfrak S\); par hypothèse, les deux éléments appartiennent alors à \(\mathfrak h\) dans \(\mathfrak S\), et il en est de même dans \(\mathfrak T\) en raison de l'isomorphisme sur \(\mathfrak O\) et \(\mathfrak o\) pris séparément. L'hypothèse de plongement est donc nécessaire et suffisante.

\paragraph{3. Idéal définissant d'un anneau relativement à un sous-anneau.}
Un système \(S\) d'éléments \(\ldots,z,\ldots\) de \(\mathfrak O\) est appelé \emph{système générateur} de \(\mathfrak O\) relativement à \(\mathfrak h\) si \(\mathfrak O\) est égal à l'anneau engendré dans \(\mathfrak O\) par \(\mathfrak h\) et \(S\), donc à l'intersection de tous les anneaux de \(\mathfrak O\) qui contiennent \(\mathfrak h\) et \(S\); en symboles, \(\mathfrak O=\mathfrak h[S]\). Supposons de nouveau que \(\mathfrak h\) soit un sous-anneau du centre de \(\mathfrak O\). Alors \(\mathfrak O\) est une image homomorphe de l'anneau de polynômes non commutatifs \(\mathfrak h[..Z..]\), où les \(Z\) sont des indéterminées et où le système des indéterminées a la même cardinalité que \(S\). On entend ici par anneau de polynômes non commutatifs le système de toutes les sommes finies \(\sum h_iP_i=\sum P_ih_i\), où les \(P_i\) sont des produits de puissances, y compris l'unité correspondant à l'exposant zéro:
\[
P_i=Z_{i_1}^{\varrho_1}Z_{i_2}^{\varrho_2}\cdots Z_{i_n}^{\varrho_n}
\]
(les mêmes indices pouvant naturellement se répéter), et où l'égalité est définie comme l'égalité des coefficients relativement aux divers produits de puissances. Il résulte de l'homomorphisme que \(\mathfrak O\) est isomorphe à l'anneau quotient \(\mathfrak h[..Z..]/\mathfrak M\), où l'idéal bilatère \(\mathfrak M\) se compose de tous les polynômes \(F(Z)\) (c'est-à-dire des éléments de \(\mathfrak h[..Z..]\)) tels que \(F(z)\) s'annule dans \(\mathfrak O\), les \(Z\) et les \(z\) se correspondant par l'homomorphisme. \emph{L'idéal \(\mathfrak M\) est appelé idéal définissant relativement à \(\mathfrak h\), plus précisément idéal correspondant au système générateur des \(z\).} La définition montre que des anneaux isomorphes, relativement à des sous-anneaux \(\mathfrak h,\overline{\mathfrak h}\) qui se correspondent isomorphiquement et à des systèmes générateurs \(S,\overline S\), possèdent aussi des idéaux définissants isomorphes; lorsque \(\mathfrak h\) coïncide avec \(\overline{\mathfrak h}\), on peut en particulier les supposer identiques.

Puisque \(\mathfrak h\) est supposé contenu dans le centre de \(\mathfrak O\), l'anneau \(\mathfrak O\) est certainement commutatif s'il possède un système générateur \(S\) formé d'un seul élément \(z\). Si \(\mathfrak M\) est en outre un idéal principal, l'équation \(G(z)=0\), correspondant au polynôme générateur \(G(Z)\) de \(\mathfrak M\), est appelée équation définissante de \(\mathfrak O\) relativement à \(\mathfrak h\).

\paragraph{4. Idéal définissant d'un produit direct.}
Supposons, comme en 3, que \(\mathfrak O\) soit isomorphe à \(\mathfrak h[..Z..]/\mathfrak M\), où \(\mathfrak M\) est donc un idéal définissant. À côté de \(\mathfrak h[..Z..]\), considérons l'anneau de polynômes \(\mathfrak o[..Z..]\), c'est-à-dire le système de toutes les sommes finies
\[
\sum\gamma_iP_i=\sum P_i\gamma_i
\qquad\hbox{avec }\gamma_i\hbox{ dans }\mathfrak o,
\]
l'égalité étant définie comme égalité des coefficients. \(\mathfrak o[Z]\) contient \(\mathfrak h[Z]\) et est le produit direct de \(\mathfrak h[..Z..]\) et de \(\mathfrak o\) relativement à \(\mathfrak h\); autrement dit, il provient de \(\mathfrak h[..Z..]\) par extension du domaine des coefficients.\srcfn{10)}{En effet, les conditions d'égalité et de composition sont satisfaites, et l'égalité \(\gamma=h+h_1P_1+\cdots+h_rP_r\) dans \(\mathfrak o[..Z..]\) entraîne \(\gamma=h\). Voir à ce sujet \S{} 2, 2. Le fait que \(\mathfrak o[..Z..]\) soit un produit direct n'est pas utilisé dans ce numéro.} Notons \(\mathfrak M_{\mathfrak o}\) l'idéal étendu de \(\mathfrak M\) relativement à \(\mathfrak o[..Z..]\) --- l'intersection de tous les idéaux bilatères de \(\mathfrak o[..Z..]\) qui contiennent \(\mathfrak M\) ---; c'est donc le système de toutes les combinaisons linéaires d'éléments de \(\mathfrak M\) à coefficients dans \(\mathfrak o\), qui, puisque \(\mathfrak o\) commute avec les \(Z\), forme déjà un idéal bilatère. \emph{Si le produit direct \(\mathfrak O_{\mathfrak o}\) relativement à \(\mathfrak h\) existe, alors \(\mathfrak O_{\mathfrak o}\) est isomorphe à \(\mathfrak o[..Z..]/\mathfrak M_{\mathfrak o}\).} En effet, \(\mathfrak O_{\mathfrak o}\) est une image homomorphe de \(\mathfrak o[..Z..]\), donc isomorphe à l'anneau quotient par l'idéal définissant de \(\mathfrak O_{\mathfrak o}\) relativement à \(\mathfrak o\), correspondant au système générateur des \(z\). Or, d'après la définition de l'égalité en 1, (4), les relations entre les \(z\) dans \(\mathfrak O_{\mathfrak o}\) s'écrivent comme des combinaisons linéaires, à coefficients dans \(\mathfrak o\), des relations qui existent dans \(\mathfrak O\); ainsi \(\mathfrak M_{\mathfrak o}\) est l'idéal définissant.

\paragraph{Remarque.}
Le même raisonnement donne encore une représentation symétrique, mais plus lourde, du produit direct et de l'idéal définissant correspondant. Supposons \(\mathfrak o\) isomorphe à \(\mathfrak h[..Y..]/\mathfrak N\), où les \(Y\) sont des symboles distincts des \(Z\), de sorte que l'anneau de polynômes \(\mathfrak h[..Y..Z..]\) soit défini. Alors \(\mathfrak O_{\mathfrak o}\) est isomorphe à
\[
\mathfrak h[..Y..Z..]/(\mathfrak M,\mathfrak N,..YZ-ZY..),
\]
où l'idéal définissant du membre de droite est engendré par \(\mathfrak M,\mathfrak N\) et toutes les combinaisons \(YZ-ZY\). (Ces dernières expriment la permutabilité des éléments de \(\mathfrak o\) et de \(\mathfrak O\); les premières, le fait que les relations dans \(\mathfrak O_{\mathfrak o}\) sont des conséquences de celles qui existent séparément dans \(\mathfrak O\) et dans \(\mathfrak o\).)

\subsection*{\S{} 2. Anneaux possédant une base indépendante comme $\mathfrak h$-module. Construction du produit direct}

\paragraph{1. Construction du produit direct.}
Un système \(T\) d'éléments \(\ldots,t,\ldots\) de \(\mathfrak O\) est appelé base du \(\mathfrak h\)-module \(\mathfrak O\) si \(\mathfrak O\) est égal au \(\mathfrak h\)-module engendré dans \(\mathfrak O\) par \(T\). Puisque l'unité appartient à \(\mathfrak h\), l'anneau \(\mathfrak O\) se compose alors de toutes les combinaisons linéaires \(h_{i_1}t_{i_1}+\cdots+h_{i_r}t_{i_r}\). Cette base du \(\mathfrak h\)-module est dite indépendante lorsque
\[
h_{i_1}t_{i_1}+\cdots+h_{i_r}t_{i_r}=0
\quad\hbox{entraîne toujours}\quad
h_{i_1}=0,\ldots,h_{i_r}=0.
\]
\emph{Si \(\mathfrak O\) possède une base indépendante comme \(\mathfrak h\)-module contenant l'élément unité, le produit direct \(\mathfrak O_{\mathfrak o}\) existe pour tout anneau d'extension \(\mathfrak o\) de \(\mathfrak h\), pourvu seulement que \(\mathfrak h\) soit un sous-anneau du centre et que l'intersection ensembliste \([\mathfrak O,\mathfrak o]\) soit \(\mathfrak h\).}\srcfn{11)}{On pourrait naturellement, de manière symétrique, imposer plutôt l'hypothèse de base à \(\mathfrak o\). D'après \S{} 2, 2, on suppose --- sans restreindre la généralité; voir E. Noether, \emph{Der Diskriminantensatz für die Ordnungen eines algebraischen Zahl- und Funktionenkörpers}, ce journal 157 (1927), p. 82--104, \S{}\S{} 1, 2 --- qu'il n'existe aucune loi de composition entre les éléments de \(\mathfrak O\) et de \(\mathfrak o\) qui n'appartiennent pas à \(\mathfrak h\).} D'après \S{} 1, 2, il suffit de construire un anneau de plongement \(\mathfrak R\) correspondant, lequel se révélera déjà être ici le produit direct. Définissons \(\mathfrak R\) comme l'ensemble de toutes les combinaisons linéaires
\[
\gamma_{i_1}t_{i_1}+\cdots+\gamma_{i_r}t_{i_r}
=t_{i_1}\gamma_{i_1}+\cdots+t_{i_r}\gamma_{i_r},
\]
où les \(t_{i_r}\) parcourent tous les systèmes finis d'éléments de la base indépendante \(T\), et les \(\gamma\) sont les éléments correspondants de \(\mathfrak o\). Définissons l'égalité comme égalité des coefficients relativement aux \(t_{i_\nu}\), et la somme et le produit suivant les règles de calcul usuelles, comme au \S{} 1, 2. Par la propriété de base, le produit de deux quelconques des \(t_{i_\nu}\) est une combinaison linéaire finie à coefficients dans \(\mathfrak h\); ainsi \(\mathfrak R\) est un anneau. Il contient \(\mathfrak O\), puisque \(\mathfrak h\) est un sous-anneau de \(\mathfrak o\), et il contient \(\mathfrak o\), puisque \(T\) contient l'élément unité. Puisque les \(\gamma\) sont en outre permutables élément par élément tant avec les \(h_{i_\nu}\) qu'avec les \(t_{i_\nu}\), \(\mathfrak O\) est permutable élément par élément avec \(\mathfrak o\) dans \(\mathfrak R\). Enfin, l'intersection de \(\mathfrak O\) et de \(\mathfrak o\) dans \(\mathfrak R\) est \(\mathfrak h\), car une relation
\[
\gamma=\gamma e=h_1t_1+\cdots+h_rt_r
\]
entraîne, d'après la définition de l'égalité --- puisque \(e\) appartient par hypothèse à la base ---, qu'un des \(t\) est égal à \(e\), que les autres s'annulent et, par conséquent, que \(\gamma=h\). Le produit direct existe donc; il est égal à \(\mathfrak R\), lequel est le produit de \(\mathfrak O\) et de \(\mathfrak o\) et satisfait aux conditions de permutabilité et d'égalité (les éléments de la base étant indépendants, cette dernière condition se réduit à l'égalité des coefficients). Cette base \(T\) de \(\mathfrak O\) devient simultanément une \(\mathfrak o\)-base de \(\mathfrak O_{\mathfrak o}\).

\paragraph{2. Modules d'extension et de restriction.}
Si \(\mathfrak B\) est un \(\mathfrak h\)-module de \(\mathfrak O\), on entend par son extension \(\mathfrak B_{\mathfrak o}\) le \(\mathfrak o\)-module de \(\mathfrak O_{\mathfrak o}\) qui est l'intersection de tous les \(\mathfrak o\)-modules de \(\mathfrak O_{\mathfrak o}\) contenant \(\mathfrak B\). Si \(\mathfrak C\) est un \(\mathfrak o\)-module de \(\mathfrak O_{\mathfrak o}\), son module de restriction \(\mathfrak C_{\mathfrak h}\) dans \(\mathfrak O\) est défini comme l'intersection \([\mathfrak C,\mathfrak O]\). Si \(\mathfrak B\) est en même temps un idéal de \(\mathfrak O\), alors \(\mathfrak B_{\mathfrak o}\) est un idéal de \(\mathfrak O_{\mathfrak o}\); si \(\mathfrak C\) est un idéal de \(\mathfrak O_{\mathfrak o}\), alors \([\mathfrak C,\mathfrak O]\) est un idéal de \(\mathfrak O\).

\paragraph{Théorème.}
\emph{Si un module \(\mathfrak B\) de \(\mathfrak O\) possède une base de \(\mathfrak h\)-module indépendante qui peut être complétée en une base de \(\mathfrak h\)-module indépendante de \(\mathfrak O\), alors \(\mathfrak B\) est la restriction à \(\mathfrak O\) de son module étendu \(\mathfrak B_{\mathfrak o}\).}

\paragraph{Démonstration.}
Soit \(T\) une base de \(\mathfrak h\)-module indépendante de \(\mathfrak B\), qui peut par hypothèse être complétée, au moyen d'un système \(\overline T\) d'éléments \(\overline t\) de \(\mathfrak O\), en une base de \(\mathfrak h\)-module indépendante de \(\mathfrak O\). Alors \(T\) est aussi une \(\mathfrak o\)-base indépendante de \(\mathfrak B_{\mathfrak o}\), et la réunion de \(T\) et de \(\overline T\) une telle base de \(\mathfrak O_{\mathfrak o}\).

Le module \(\mathfrak B\) est contenu dans l'intersection \([\mathfrak B_{\mathfrak o},\mathfrak O]\). Réciproquement, soit \(c\) un élément de ce module de restriction. Comme élément de \(\mathfrak B_{\mathfrak o}\), \(c\) admet une expression au moyen des éléments de \(T\), laquelle est simultanément une expression relativement à la \(\mathfrak o\)-base \(T,\overline T\) de \(\mathfrak O_{\mathfrak o}\):
\[
c=\gamma_1t_1+\cdots+\gamma_rt_r
\qquad\hbox{avec }\gamma\hbox{ dans }\mathfrak o.
\]
Comme élément de \(\mathfrak O\), \(c\) admet une expression relativement à la \(\mathfrak h\)-base \(T,\overline T\) de \(\mathfrak O\), soit
\[
c=h_{i_1}t_{i_1}+\cdots+h_{i_r}t_{i_r}
+h_{j_1}\overline t_{j_1}+\cdots+h_{j_s}\overline t_{j_s},
\]
qui, puisque \(\mathfrak O\) est un sous-anneau de \(\mathfrak O_{\mathfrak o}\) et \(\mathfrak h\) un sous-anneau de \(\mathfrak o\), est elle aussi une expression relativement à la \(\mathfrak o\)-base \(T,\overline T\) de \(\mathfrak O_{\mathfrak o}\). Les deux expressions doivent donc avoir les mêmes coefficients; les coefficients des \(\overline t\) s'annulent et
\[
c=h_{i_1}t_{i_1}+\cdots+h_{i_r}t_{i_r}
\]
montre que \(c\) appartient à \(\mathfrak B\). Ainsi \(\mathfrak B\) est bien le module de restriction.

\paragraph{Remarque complémentaire.}
La démonstration n'utilise que des hypothèses un peu plus faibles que celles formulées dans le théorème et prouve les autres. On a donc la forme renforcée suivante:

\emph{Le module \(\mathfrak B\) est le module de restriction de son extension \(\mathfrak B_{\mathfrak o}\) s'il existe dans \(\mathfrak O\) une base de \(\mathfrak h\)-module indépendante \(T,\overline T\), où \(T\) est contenu dans \(\mathfrak B\), telle que \(T\) soit en même temps une base (indépendante) du \(\mathfrak o\)-module \(\mathfrak B_{\mathfrak o}\). Alors \(T\) se révèle aussi être une base indépendante du \(\mathfrak h\)-module \(\mathfrak B\).}

\paragraph{3. Critère suffisant d'existence du produit direct.}
Il peut arriver (par exemple pour la différente relative) qu'une base indépendante comme \(\mathfrak h\)-module n'existe qu'après extension du domaine des coefficients \(\mathfrak h\). L'existence du produit direct peut alors se déduire du critère suffisant suivant:

\emph{S'il existe un anneau d'extension \(\mathfrak f\) de \(\mathfrak h\) tel que les produits directs \(\mathfrak O_{\mathfrak f}\) et \(\mathfrak o_{\mathfrak f}\) relativement à \(\mathfrak h\), puis \(\mathfrak O_{\mathfrak f}\times\mathfrak o_{\mathfrak f}\) relativement à \(\mathfrak f\), existent, alors \(\mathfrak O\times\mathfrak o\) relativement à \(\mathfrak h\) existe.}

Il faut montrer que \(\mathfrak R=\mathfrak O_{\mathfrak f}\times\mathfrak o_{\mathfrak f}\) est un anneau de plongement au sens du \S{} 1, 2. En effet, \(\mathfrak O\) et \(\mathfrak o\), comme sous-anneaux de \(\mathfrak O_{\mathfrak f}\) et de \(\mathfrak o_{\mathfrak f}\), deviennent permutables élément par élément; de plus,
\[
[\mathfrak O,\mathfrak o]\subseteqq[\mathfrak O_{\mathfrak f},\mathfrak o_{\mathfrak f}]=\mathfrak f.
\]
On a aussi \([\mathfrak O,\mathfrak o]\subseteqq\mathfrak O\); par conséquent \([\mathfrak O,\mathfrak o]\subseteqq[\mathfrak O,\mathfrak f]=\mathfrak h\), donc \([\mathfrak O,\mathfrak o]=\mathfrak h\), ce qui prouve que \(\mathfrak R\) est un anneau de plongement.

\paragraph{Remarque.}
La démonstration n'utilise l'hypothèse d'intersection que pour \(\mathfrak O_{\mathfrak f}\), et non pour \(\mathfrak o_{\mathfrak f}\); ce dernier anneau peut donc être simplement le produit de \(\mathfrak o\) et de \(\mathfrak f\).

\subsection*{\S{} 3. Différente d'un anneau relativement à un sous-anneau. Quotient différentiel d'un idéal définissant}

À partir de maintenant, tous les anneaux considérés sont supposés commutatifs.\srcfn{12)}{Sans cette hypothèse, il faudrait définir une différente à droite et une différente à gauche, puis rechercher dans quelle mesure elles coïncident entre elles et, dans le cas des ordres de systèmes semi-simples, dans quelle mesure elles coïncident avec les constructions connues (comme généralisation des recherches du \S{} 4).} On suppose en outre toujours que les anneaux \(\mathfrak O\) et \(\mathfrak o\) sont des extensions isomorphes, plus précisément équivalentes, du sous-anneau \(\mathfrak h\) (c'est-à-dire que l'application de \(\mathfrak O\) sur \(\mathfrak o\) prolonge l'application identique de \(\mathfrak h\)), et que le produit direct \(\mathfrak O_{\mathfrak o}\) relativement à \(\mathfrak h\) existe.

\paragraph{1. Définition de la différente.}
Supposons que le produit direct \(\mathfrak O_{\mathfrak o}\) relativement à \(\mathfrak h\) existe. Que \(x\) parcoure tous les éléments de \(\mathfrak O\), et que \(\xi\) désigne chaque fois l'élément correspondant dans \(\mathfrak o\) par l'isomorphisme; notons \(\mathfrak P\) l'\emph{idéal des différences} bilatère de \(\mathfrak O_{\mathfrak o}\), engendré par toutes les différences \(x-\xi\):
\[
\mathfrak P=\{..(x-\xi)..\}\srcfnmark{13)}
\]
\srcfntext{13)}{Cette notation sera toujours employée pour les idéaux engendrés par des ensembles.}
Le \emph{quotient des différences} \(\mathfrak A\) est défini comme le quotient de l'idéal nul par \(\mathfrak P\); il se compose donc de tous les éléments \(a\) de \(\mathfrak O_{\mathfrak o}\) tels que \(a\mathfrak P\) s'annule:
\[
\mathfrak A=(0):\mathfrak P.
\]
La \emph{différente} \(\mathfrak d\) de \(\mathfrak o\) relativement à \(\mathfrak h\) est définie comme \(\mathfrak A[x\to\xi]\); elle s'obtient donc à partir de \(\mathfrak A\) en remplaçant les \(x\), dans chaque \(a=\alpha_{i_1}x_{i_1}+\cdots+\alpha_{i_r}x_{i_r}\) de \(\mathfrak A\), par les \(\xi\) qui leur correspondent isomorphiquement. De même, la différente \(\mathfrak D\) de \(\mathfrak O\) relativement à \(\mathfrak h\) est définie comme \(\mathfrak A[\xi\to x]\). Puisque \(\mathfrak A\) est défini symétriquement en \(x\) et \(\xi\), les différentes \(\mathfrak d\) et \(\mathfrak D\), qui se déduisent l'une de l'autre en permutant \(\xi\) et \(x\), se correspondent par l'isomorphisme de \(\mathfrak o\) sur \(\mathfrak O\). Puisque \(\mathfrak A\) est un idéal de \(\mathfrak O_{\mathfrak o}\), donc en particulier un \(\mathfrak o\)-module et un \(\mathfrak O\)-module, \(\mathfrak d\) et \(\mathfrak D\) sont des idéaux de \(\mathfrak o\) et de \(\mathfrak O\).

\paragraph{2. Quotient différentiel d'un idéal définissant.}
Soit, comme au \S{} 1, 3, \(\mathfrak M\) l'idéal définissant de \(\mathfrak O\) relativement à \(\mathfrak h\), correspondant au système générateur des \(z\). D'après le \S{} 1, 3, \(\mathfrak M\) est aussi l'idéal définissant de \(\mathfrak o\) relativement à \(\mathfrak h\), correspondant au système générateur isomorphe des \(\xi\); et, d'après le \S{} 1, 4, \(\mathfrak O_{\mathfrak o}\) est isomorphe à \(\mathfrak o[..Z..]/\mathfrak M_{\mathfrak o}\). Dans l'anneau de polynômes \(\mathfrak o[..Z..]\), l'idéal des différences et le quotient des différences sont maintenant définis par
\[
\mathfrak Q=\{..(Z-\xi)..\}
\]
et
\[
\mathfrak C=\mathfrak M_{\mathfrak o}:\mathfrak Q
=\{..(F(Z)-F(\xi))..\}:\{..(Z-\xi)..\}.
\]
Le \emph{quotient différentiel} \(\mathfrak M'\) est défini comme \(\mathfrak C[\xi\to Z]\). Il faut remarquer que cette définition est univoque, indépendamment de la \emph{représentation} des éléments de \(\mathfrak o\) par les générateurs \(\xi\). En effet, puisque \(\mathfrak C\) est un quotient d'idéaux, il contient \(\mathfrak M_{\mathfrak o}\); deux représentations d'un élément de \(\mathfrak o\) diffèrent d'un polynôme \(F(\xi)\), donc, après substitution de \(Z\), d'un \(F(Z)\) appartenant à \(\mathfrak M\), et par suite déjà à \(\mathfrak C\). L'idéal \(\mathfrak M'\) est un idéal de \(\mathfrak h[..Z..]\) contenant \(\mathfrak M\). \emph{Pour \(Z\to\xi\), le quotient différentiel \(\mathfrak M'\) passe dans la différente \(\mathfrak d\) de \(\mathfrak o\).} En effet, par une application homomorphe, les quotients se correspondent en même temps que les idéaux. Ainsi \(\mathfrak A\) est l'image homomorphe de \(\mathfrak C\) par l'homomorphisme entre \(\mathfrak o[..Z..]\) et \(\mathfrak O_{\mathfrak o}\), c'est-à-dire par la correspondance entre \(Z\) et \(z\), puisque la même chose vaut pour \(\mathfrak P\) et l'idéal nul relativement à \(\mathfrak Q\) et \(\mathfrak M_{\mathfrak o}\).\srcfn{13a)}{Note de l'éditeur: il est essentiel que \(\mathfrak M_{\mathfrak o}\) soit l'image réciproque de \((0)\), afin de pouvoir conclure à la correspondance des quotients.} Il en résulte que, par l'homomorphisme, la substitution \(\xi\to Z\) dans \(\mathfrak C\) correspond à la substitution \(\xi\to z\) dans \(\mathfrak A\); la différente \(\mathfrak D\) de \(\mathfrak O\) est donc l'image homomorphe de \(\mathfrak M'\) par la correspondance entre \(Z\) et \(z\). Ainsi
\[
\mathfrak M'[Z\to z]=\mathfrak D
\qquad\hbox{et}\qquad
\mathfrak M'[Z\to\xi]=\mathfrak D[z\to\xi]=\mathfrak d.
\]

\paragraph{3. Différente dans le cas d'une équation définissante.}
Supposons que \(\mathfrak O\) possède une équation définissante \(G(z)=0\) (\S{} 1, 3, fin), et supposons en outre que le coefficient du terme de plus haut degré soit l'unité, donc
\[
G(Z)=Z^n+h_1Z^{n-1}+\cdots+h_n,
\]
avec des coefficients \(h_i\) dans \(\mathfrak h\). Le lien avec le quotient différentiel polynomial est alors établi par le théorème suivant.

\paragraph{Théorème.}
\emph{La différente de \(\mathfrak o\), respectivement de \(\mathfrak O\), relativement à \(\mathfrak h\), est définie (c'est-à-dire que le produit direct \(\mathfrak O_{\mathfrak o}\) relativement à \(\mathfrak h\) existe) et est l'idéal principal ayant pour élément de base \(G'(\xi)\), respectivement \(G'(z)\), où \(G'(Z)\) désigne le quotient différentiel polynomial. Le quotient différentiel d'idéal est \(\mathfrak M'=(G'(Z),G(Z))\).}

\paragraph{Démonstration.}
L'hypothèse implique d'abord que les éléments \(e,z,\ldots,z^{n-1}\) forment une base indépendante de \(\mathfrak O\) comme \(\mathfrak h\)-module. En effet, puisque le coefficient dominant de \(G(Z)\) est l'unité, ces puissances forment effectivement une base du \(\mathfrak h\)-module; celle-ci est indépendante parce que \(\mathfrak M\) ne peut contenir aucun polynôme de degré inférieur à \(n\). En effet, le degré d'un polynôme \(K(Z)G(Z)\) est égal à la somme des degrés des facteurs, de nouveau en vertu de l'hypothèse sur le coefficient dominant. L'existence de cette base du module entraîne donc celle du produit direct \(\mathfrak O_{\mathfrak o}\) relativement à \(\mathfrak h\), et la différente est ainsi définie. Le reste résulte de la considération du quotient des différences \(\mathfrak C\) dans \(\mathfrak o[..Z..]\) (voir 2). Ce quotient des différences est l'idéal principal ayant pour polynôme générateur
\[
%
C(Z,\xi)=(Z^{n-1}+Z^{n-2}\xi+\cdots+\xi^{n-1})
+h_1(Z^{n-2}+\cdots+\xi^{n-2})+\cdots+h_{n-1}.
%
\]
En effet, \(C\) s'obtient par division formelle de \(G(Z)-G(\xi)\) par \(Z-\xi\) et appartient donc à \(\mathfrak C\), puisque \(Z-\xi\) est un élément de base de \(\mathfrak Q\). Modulo \(C\), un polynôme quelconque se réduit à un polynôme \(D\) dont le degré en \(Z\) est strictement inférieur à \(n-1\); le produit \(D(Z-\xi)\) est donc de degré strictement inférieur à \(n\) et ne peut appartenir à \(\mathfrak M\) que s'il s'annule. Mais \(D(Z-\xi)=0\) entraîne \(D=0\), puisque le coefficient de \(Z\) dans \(Z-\xi\) est l'unité. Ainsi \(C\) est l'élément de base de \(\mathfrak C\). Or la substitution \(\xi\to Z\) dans \(C\) et dans ses multiples (donc dans \(G(Z)=C(Z-\xi)\) et ses multiples) donne
\[
\mathfrak M'=(G'(Z),G(Z)),
\]
et, par suite,
\[
\mathfrak d=(G'(\xi))
\qquad\hbox{et}\qquad
\mathfrak D=(G'(z)).
\]
\subsection*{\S{} 4. Différente d'une somme directe.}

On fait sur $\mathfrak O$ les hypothèses suivantes :

\begin{enumerate}
\item[(1)] $\mathfrak O$ est une somme directe d'idéaux
\[
\mathfrak O=\mathfrak R_1+\cdots+\mathfrak R_r
=\mathfrak R_i+\mathfrak S_i
=\mathfrak Oe_1+\cdots+\mathfrak Oe_r,
\]
où les $e_i$ désignent les composantes de l'unité.

\item[(2)] Chaque $\mathfrak R_i$ possède une base indépendante $T_i$ comme $\mathfrak h_i$-module, où $\mathfrak h_i$ est la composante de $\mathfrak h$ dans $\mathfrak R_i$, donc $\mathfrak h_i=\mathfrak h e_i$. Cette base $T_i$ contient, comme élément, l'unité $e_i$ de $\mathfrak R_i$.

\item[(3)] Toute intersection $[\mathfrak h,\mathfrak S_i]$ est réduite à l'élément nul ($\mathfrak S_i$ ne contient pas de constantes).
\end{enumerate}

En vertu de l'isomorphisme, les mêmes hypothèses valent pour $\mathfrak o$.

\paragraph{Théorème.}
\emph{Sous les hypothèses indiquées, la différente de $\mathfrak O$ relativement à $\mathfrak h$ est la somme directe des différentes des $\mathfrak R_i$ relativement aux $\mathfrak h_i$.}

La démonstration repose sur l'établissement d'une suite de faits particuliers.

\paragraph{I.}
Il s'ensuit d'abord que $\mathfrak O$ possède lui aussi une base indépendante comme $\mathfrak h$-module et contenant $e$, de sorte que la différente de $\mathfrak O$, respectivement de $\mathfrak o$, relativement à $\mathfrak h$ est définie ; cette base est constituée par la réunion $T$ des bases $T_i$ des $\mathfrak R_i$. En effet, de
\[
x=xe_1+\cdots+xe_r
\]
et de
\[
xe_i=(h_1e_i)t_{1i}e_i+\cdots+(h_\lambda e_i)t_{\lambda i}e_i
=h_1(t_{1i}e_i)+\cdots+h_\lambda(t_{\lambda i}e_i)
\]
il résulte que $T$ est bien une $\mathfrak h$-base de $\mathfrak O$. Celle-ci est aussi indépendante : une relation entre les éléments de $T$ entraîne, puisque la somme est directe, des relations entre les éléments de chaque $T_i$, et par suite $h e_i=0$ pour tous les coefficients. Mais cela donne $h=0$ d'après l'hypothèse (3) ; en effet, puisque
\[
h\equiv h e_i\ (\mathfrak S_i)
\]
la relation $h e_i=0$ implique que $h$ appartient à $[\mathfrak h,\mathfrak S_i]$. Comme $T$ contient en outre, comme éléments, les unités $e_i$ des $\mathfrak R_i$, dont la somme est $e$, on peut remplacer dans $T$ l'un des $e_i$ par $e$ ; on obtient ainsi une $\mathfrak h$-base $\overline T$ satisfaisant à toutes les conditions. Il en résulte l'existence de $\mathfrak O_{\mathfrak o}$ et de la différente de $\mathfrak O$ relativement à $\mathfrak h$.

\paragraph{II.}
La condition d'intersection (3) entraîne maintenant la condition plus forte
\[
[\mathfrak o,\mathfrak S_{i(\mathfrak o)}]=0,
\]
où $\mathfrak S_{i(\mathfrak o)}$ désigne l'idéal de $\mathfrak O_{\mathfrak o}$ déduit de $\mathfrak S_i$. Soit
\[
\alpha=\alpha e=\beta_1t_1+\cdots+\beta_\lambda t_\lambda,
\]
où $t_1,\ldots,t_\lambda$ sont les éléments de la base obtenue en retranchant $T_i$ de $T$, et où $\alpha$ et les $\beta$ appartiennent à $\mathfrak o$. Si l'on forme $\overline T$ en remplaçant $e_i$ par $e$, alors $e,t_1,\ldots,t_\lambda$ sont des éléments distincts de la $\mathfrak h$-base $\overline T$ de $\mathfrak O$. Or l'égalité dans $\mathfrak O_{\mathfrak o}$ devient l'égalité des coefficients dans toute base de cette sorte ; la relation
\[
\alpha=\alpha e=\beta_1t_1+\cdots+\beta_\lambda t_\lambda,
\]
donne donc $\alpha=0$, ce qui \emph{démontre la propriété d'intersection renforcée}. Il en résulte en particulier (comme plus haut) que
\[
\alpha\text{ élément de }\mathfrak o\quad\text{et}\quad\alpha\ne0
\]
implique $\alpha e_i\ne0$.

\paragraph{III.}
Une conséquence immédiate de cette condition d'intersection est que $\mathfrak o$ est \emph{isomorphe comme anneau} à ses \emph{composantes} $\mathfrak o e_i$ correspondant à la décomposition de $\mathfrak O$ ; cet isomorphisme comprend l'isomorphisme
\[
\mathfrak h\simeq\mathfrak h_i=\mathfrak h e_i.
\]
En effet, puisque $e_i^2=e_i$ --- et puisque $e_i$ commute avec tous les éléments de $\mathfrak o$ ---, $\mathfrak o e_i$ est une image de $\mathfrak o$ isomorphe comme anneau ; la correspondance est biunivoque, car $\alpha\ne0$ entraîne $\alpha e_i\ne0$. Au sous-anneau $\mathfrak h$ correspond alors la composante $\mathfrak h_i=\mathfrak h e_i$.

\paragraph{IV.}
D'\emph{autres isomorphismes} résultent si l'on fait encore intervenir la décomposition correspondante de $\mathfrak o$ :
\[
\mathfrak o=\mathfrak r_1+\cdots+\mathfrak r_r
=\mathfrak r_i+\mathfrak s_i
=\mathfrak o\varepsilon_1+\cdots+\mathfrak o\varepsilon_r.
\]
On obtient alors
\[
\mathfrak o\simeq\mathfrak o e_i,
\quad \mathfrak O\simeq\mathfrak O\varepsilon_i,
\quad \mathfrak r_\lambda\simeq\mathfrak r_\lambda e_i,
\quad \mathfrak R_\lambda\simeq\mathfrak R_\lambda\varepsilon_i;
\quad \mathfrak h\simeq\mathfrak h e_i\simeq\mathfrak h\varepsilon_i\simeq\mathfrak h e_i\varepsilon_i.
\]
La relation $\mathfrak O\simeq\mathfrak O\varepsilon_i$ correspond exactement à $\mathfrak o\simeq\mathfrak o e_i$ lorsqu'on part, dans le raisonnement précédent, de $\mathfrak o$ au lieu de $\mathfrak O$. De même, les deux relations suivantes se correspondent ; la première résulte de ce que, dans $\mathfrak o\simeq\mathfrak o e_i$, le sous-anneau $\mathfrak r_\lambda$ est lui aussi mis en correspondance avec $\mathfrak r_\lambda e_i$. On obtient de même
\[
\mathfrak h\simeq\mathfrak h e_i\simeq\mathfrak h\varepsilon_i.
\]
Mais puisque $\mathfrak h\varepsilon_i$ est un sous-anneau de $\mathfrak o$ (et de $\mathfrak r_i$), on a enfin
\[
\mathfrak h\varepsilon_i\simeq\mathfrak h\varepsilon_i e_i.
\]

\paragraph{V.}
De la représentation de $\mathfrak O$ comme somme directe, il résulte ceci : \emph{si $\mathfrak C=(0):\mathfrak B$ est le quotient de l'idéal nul par l'idéal $\mathfrak B$ dans $\mathfrak O$, alors toute composante $\mathfrak C_i$ est égale au quotient de l'idéal nul de $\mathfrak R_i$ par la composante $\mathfrak B_i$.} De
\[
\mathfrak B\mathfrak C=(\mathfrak B e_1)(\mathfrak C e_1)+\cdots+(\mathfrak B e_r)(\mathfrak C e_r)=0
\]
il résulte, puisque la somme est \emph{directe}, que, séparément,
\[
(\mathfrak B e_i)(\mathfrak C e_i)=0.
\]
Réciproquement, si un élément $r e_i$ de $\mathfrak R_i$ satisfait à la condition $(\mathfrak B e_i)(r e_i)=0$, on a aussi $(\mathfrak B e_k)(r e_i)=0$ pour $i\ne k$.

Ainsi $r e_i$ est un élément de $\mathfrak C$ et en même temps de $\mathfrak R_i$ ; il appartient donc à la composante $\mathfrak C_i=\mathfrak C e_i$. On reconnaît ainsi $\mathfrak C_i=\mathfrak C e_i$ comme le quotient de l'idéal nul par $\mathfrak B_i$, donc aussi comme le quotient de l'idéal nul de $\mathfrak R_i$ par $\mathfrak B_i$, car de
\[
r e_i\mathfrak B_i\equiv0\quad(\mathfrak S_i)
\]
il résulte toujours que $r e_i\mathfrak B_i=0$.\srcfn{14)}{Le théorème résulte aussi directement du théorème employé au no 2, selon lequel, dans une application homomorphe, les idéaux aussi bien que les quotients se correspondent.} Puisque, avec $\mathfrak O$, l'anneau
\[
\mathfrak O_{\mathfrak o}=\mathfrak O_{\mathfrak o}e_1+\cdots+\mathfrak O_{\mathfrak o}e_r
\]
est lui aussi une somme directe, \emph{le théorème vaut donc également pour les idéaux de $\mathfrak O_{\mathfrak o}$.}

\paragraph{VI. Décomposition de la différente en composantes.}
Par définition (voir \S{} 3, 1), on avait $\mathfrak D=\mathfrak A[\xi\to x]$ ; on obtient donc
\[
\mathfrak D=\mathfrak D e_1+\cdots+\mathfrak D e_r
=\mathfrak A[\xi\to x]e_1+\cdots+\mathfrak A[\xi\to x]e_r.
\]
Mais puisque $e_i$ se correspond à lui-même dans $\mathfrak O$, puisque de plus l'isomorphisme envoie $\varepsilon_i$ sur $e_i$, et puisque $e_i^n=e_i$, on a
\[
\mathfrak A[\xi\to x]e_i=\mathfrak A e_i[\xi\to x]
=\mathfrak A e_i\varepsilon_i[\xi\to x].
\]
D'après V, $\mathfrak A e_i\varepsilon_i$ est égal au quotient $(0):\mathfrak P e_i\varepsilon_i$ --- et également au quotient de l'idéal nul de $\mathfrak R_i\mathfrak r_i$ par $\mathfrak P e_i\varepsilon_i$. En effet, les décompositions en sommes directes de $\mathfrak O$ et $\mathfrak o$ engendrent, par multiplication, la décomposition
\[
\mathfrak O_{\mathfrak o}=\sum\mathfrak O_{\mathfrak o}e_i\varepsilon_k
=\sum\mathfrak R_i\mathfrak r_k,
\]
qui est directe, puisque $\sum e_i\varepsilon_k=e$ et que les relations d'orthogonalité sont satisfaites pour les unités.

\paragraph{VII.}
\emph{Montrons que $\mathfrak D e_i=\mathfrak A e_i\varepsilon_i[\xi\to x]=(0):\mathfrak P e_i\varepsilon_i$ est la différente de $\mathfrak R_i$ relativement à $\mathfrak h_i$.} Cette vérification repose sur les faits suivants : $\mathfrak R_i\mathfrak r_i$ est égal à $\mathfrak R_i\varepsilon_i\times\mathfrak r_i e_i$ relativement à $\mathfrak h e_i\varepsilon_i$ ; son idéal des différences est donné par $\mathfrak P e_i\varepsilon_i$, son quotient des différences par $\mathfrak A e_i\varepsilon_i$, tandis que $\mathfrak D e_i\varepsilon_i$ est la différente de $\mathfrak R_i\varepsilon_i$ et, de manière correspondante, $\mathfrak d\varepsilon_i e_i$ celle de $\mathfrak r_i e_i$. Les isomorphismes
\[
\mathfrak R_i\varepsilon_i\simeq\mathfrak R_i
\qquad\text{et}\qquad
\mathfrak r_i\simeq\mathfrak r_i e_i
\]
donnent alors respectivement $\mathfrak D e_i$ et $\mathfrak d\varepsilon_i$ comme différentes de $\mathfrak R_i$ relativement à $\mathfrak h e_i$ et de $\mathfrak r_i$ relativement à $\mathfrak h\varepsilon_i$.

\paragraph{VIII.}
\emph{$\mathfrak R_i\mathfrak r_i$ est égal au produit direct $\mathfrak R_i\varepsilon_i\times\mathfrak r_i e_i$ relativement à $\mathfrak h e_i\varepsilon_i$.} Comme $\mathfrak R_i$ possède la base $T_i$ comme $\mathfrak h_i$-module, l'isomorphisme (IV) $\mathfrak R_i\simeq\mathfrak R_i\varepsilon_i$ montre que $T_i\varepsilon_i$ est une base de $\mathfrak R_i\varepsilon_i$ comme $\mathfrak h e_i\varepsilon_i$-module. Or, d'après IV, $\mathfrak r_i e_i$ est isomorphe à $\mathfrak R_i\varepsilon_i$ --- puisque $\mathfrak r_i\simeq\mathfrak R_i$ en vertu de $\mathfrak o\simeq\mathfrak O$ --- et est un anneau d'extension de $\mathfrak h e_i\varepsilon_i$. Le produit direct est donc constitué de toutes les formes linéaires des éléments de $T_i\varepsilon_i$ à coefficients dans $\mathfrak r_i e_i$, l'égalité étant définie comme égalité des coefficients. Mais, en vertu de $\mathfrak O\times\mathfrak o$, c'est précisément $\mathfrak R_i\mathfrak r_i$, puisque $T_i=T_i e_i$, $\mathfrak r_i=\mathfrak r_i\varepsilon_i$ et que $T_i$ appartient à une base de $\mathfrak O$ comme module.

\paragraph{IX.}
\emph{L'idéal des différences de $\mathfrak R_i\varepsilon_i\times\mathfrak r_i e_i$ est donné par $\mathfrak P e_i\varepsilon_i$, et le quotient des différences par $\mathfrak A e_i\varepsilon_i$.} Pour le voir, écrivons $\mathfrak P$ sous la forme
\[
\mathfrak P=\{\ldots,xe_1-\xi\varepsilon_1,\ldots,xe_r-\xi\varepsilon_r,\ldots\}.
\]
Alors
\[
\mathfrak P e_i\varepsilon_i
=\{\ldots,xe_i\varepsilon_i-\xi\varepsilon_i e_i,\ldots\},
\]
et cet idéal est donc bien l'idéal des différences de $\mathfrak R_i\varepsilon_i\times\mathfrak r_i e_i$. Mais puisque (d'après VI) $\mathfrak A e\varepsilon_i$ est égal au quotient de l'idéal nul de cet anneau par $\mathfrak P e_i\varepsilon_i$, on l'a identifié au quotient des différences.

\paragraph{X.}
\emph{La différente de $\mathfrak R_i\varepsilon_i$ est égale à $\mathfrak D e_i\varepsilon_i$.} La différente de $\mathfrak R_i\varepsilon_i$ s'obtient en remplaçant, dans $\mathfrak A e_i\varepsilon_i$, chaque élément de $\mathfrak r_i e_i$ par son correspondant isomorphe ; on remplace donc chaque $(\xi\varepsilon_i)e_i$ par $(xe_i)\varepsilon_i$. Autrement dit, puisque $\xi\varepsilon_i$ est encore un $\xi$ de $\mathfrak o$ qui passe à $xe_i$, on effectue dans $\mathfrak A e_i\varepsilon_i$ la substitution $\xi\to x$ --- qui transforme $(\xi\varepsilon_i)e_i$ en $(xe_i)e_i=xe_i$ --- puis on multiplie par $\varepsilon_i$. Or on avait (d'après VI) $\mathfrak A e_i\varepsilon_i[\xi\to x]=\mathfrak D e_i$ ; la différente de $\mathfrak R_i\varepsilon_i$ relativement à $\mathfrak h e_i\varepsilon_i$ est donc $\mathfrak D e_i\varepsilon_i$, et, symétriquement, la différente de $\mathfrak r_i e_i$ est $\mathfrak d e_i\varepsilon_i$.

\paragraph{XI.}
\emph{La différente de $\mathfrak R_i$ est égale à $\mathfrak D e_i$.} En effet, par l'isomorphisme, la différente de $\mathfrak R_i$ passe à celle de $\mathfrak R_i\varepsilon_i$ par multiplication par $\varepsilon_i$ ; puisque cette dernière est égale à $\mathfrak d e_i\varepsilon_i$, la première est donc égale à $\mathfrak D e_i$. Ainsi
\[
\mathfrak D=\mathfrak D e_i+\cdots+\mathfrak D e_r,
\]
\emph{et $\mathfrak D e_i$ est la différente de $\mathfrak R_i$} ; symétriquement,
\[
\mathfrak d=\mathfrak d\varepsilon_1+\cdots+\mathfrak d\varepsilon_r,
\]
\emph{et $\mathfrak d\varepsilon_i$ est la différente de $\mathfrak r_i$,} ce qui démontre le théorème.

\subsection*{\S{} 5. Anneau d'extension galoisien d'un corps. Théorèmes de structure.}

\paragraph{1. Les anneaux considérés.}
Soit $\mathsf{K}$ un \emph{corps de degré $n$ sur un corps de base $\mathsf{P}$}, et supposons $\mathsf{K}$ séparable sur $P$. Soit $\Gamma$ un corps galoisien correspondant --- déterminé à isomorphisme près --- qui contient donc les $n$ conjugués $\mathsf{K}^{(i)}$ et en est le compositum, avec $\mathsf{K}=\mathsf{K}^{(1)}$. On note $\mathfrak K$ une extension de $P$, isomorphe à $K$, et par suite à tous les $\mathsf{K}^{(i)}$, et équivalente relativement à $\mathsf{P}$, telle que l'intersection $[\mathfrak K,\Gamma]$ soit égale à $\mathsf{P}$ (c'est-à-dire qu'il n'existe pas de corps commun contenant proprement $\Gamma$ et $\mathfrak K$).\srcfn{15)}{On suppose naturellement --- ce qui ne restreint pas la généralité --- qu'il n'existe aucune relation entre les éléments de $K$ et de $\mathfrak K$ qui n'appartiennent pas à $P$ (voir \S{} 1, 2, \S{} 2, 1 et la note 11).} Puisque $\mathfrak K$ possède une base indépendante, finie et contenant l'unité comme $\mathsf{P}$-module, les produits directs $\mathfrak K_\Gamma$, $\mathfrak K_{\mathsf{K}}$, $\mathfrak K_{\mathsf{K}^{(i)}}$ existent d'après \S{} 2, 1 ; $\mathfrak K_\Gamma$ est le produit de ses $n$ sous-anneaux $\mathfrak K_{\mathsf{K}^{(i)}}$. On appelle $\mathfrak K_\Gamma$ l'\emph{anneau d'extension galoisien} de $\mathfrak K$. En même temps, $\mathfrak K_\Gamma$ est le produit direct de $\mathfrak K_{\mathsf{K}}$ et de $\Gamma$ relativement à $\mathsf{K}$, car toute base indépendante de $\mathfrak K$ sur $P$ devient aussi une base indépendante de $\mathfrak K_{\mathsf{K}}$ sur $K$, et de même pour les conjugués.\srcfn{16)}{Les théorèmes qui valent uniformément pour tous les conjugués seront le plus souvent énoncés seulement pour $\mathfrak K_{\mathsf{K}}$, en omettant donc les indices.} Les relations entre les idéaux de ces anneaux sont régies par le théorème suivant.

\paragraph{Théorème.}
\emph{Tout idéal de $\mathfrak K$ est l'idéal contracté de son extension à $\mathfrak K_{\mathsf{K}}$ ou à $\mathfrak K_\Gamma$ ; tout idéal de $\mathfrak K_{\mathsf{K}}$ est l'idéal contracté de son extension à $\mathfrak K_\Gamma$.}

\paragraph{Démonstration.}
Comme $\mathsf{P}$ et $\mathsf{K}$ sont des corps, tout $\mathsf{P}$-module de $\mathfrak K$, respectivement tout $\mathsf{K}$-module de $\mathfrak K_{\mathsf{K}}$, possède sur $\mathsf{P}$, respectivement sur $\mathsf{K}$, une base indépendante et finie que l'on peut compléter en une telle base de $\mathfrak K$, respectivement de $\mathfrak K_{\mathsf{K}}$. Le théorème résulte alors du théorème du \S{} 2, 2.

Notons $\mathfrak P_{\mathsf{K}^{(i)}}^{(i)}$ l'idéal des différences de $\mathfrak K_{\mathsf{K}^{(i)}}$, c'est-à-dire
\[
\mathfrak P_{\mathsf{K}^{(i)}}^{(i)}=\{\ldots,x-\xi^{(i)},\ldots\},
\]
et en particulier
\[
\mathfrak P_{\mathsf{K}}=\{\ldots,x-\xi,\ldots\}.
\]
Puisque $\xi$, respectivement $\xi^{(i)}$, est l'élément de $\mathsf{K}$, respectivement de $\mathsf{K}^{(i)}$, qui correspond à $x$ par l'isomorphisme, les $n$ idéaux des différences sont conjugués entre eux.

Notons $\mathfrak A_{\mathsf{K}}$, respectivement $\mathfrak A_{\mathsf{K}^{(i)}}$, les quotients des différences $(0):\mathfrak P_{\mathsf{K}}$ dans $\mathfrak K_{\mathsf{K}}$ et leurs conjugués. Leurs idéaux étendus à $\mathfrak K_\Gamma$ seront notés $\mathfrak P_\Gamma$ et $\mathfrak A_\Gamma$. Du théorème de ce numéro sur les idéaux étendus et contractés résulte en particulier ceci :

\emph{L'idéal des différences $\mathfrak P_{\mathsf{K}}$ et le quotient des différences $\mathfrak A_{\mathsf{K}}$ sont les idéaux contractés de leurs extensions $\mathfrak P_\Gamma$ et $\mathfrak A_\Gamma$.}

\paragraph{2. Décomposition de l'anneau d'extension galoisien en somme directe.}
Le lien entre l'idéal des différences, le quotient des différences et la structure de l'anneau d'extension galoisien est donné par le

\paragraph{Théorème de structure.}
\emph{L'idéal nul de l'anneau d'extension galoisien $\mathfrak K_\Gamma$ est l'intersection des extensions des $n$ idéaux des différences conjugués ; $\mathfrak K_\Gamma$ lui-même est la somme directe des extensions des $n$ quotients des différences conjugués. $\mathfrak K_{\mathsf{K}^{(i)}}$ est une somme directe, et son idéal nul est l'intersection de l'idéal des différences et du quotient des différences correspondants.}

\paragraph{Démonstration.}
L'anneau quotient par chaque $\mathfrak P_\Gamma^{(i)}$ est de rang un relativement à $\Gamma$, car tout élément est congru modulo $\mathfrak P_\Gamma^{(i)}$ à un élément de $\Gamma$, tandis que $\mathfrak P_\Gamma^{(i)}$ ne peut contenir aucun élément $\ne0$ de $\Gamma$ (puisque $\mathfrak P_\Gamma^{(i)}$ s'annule pour $x\to\xi^{(i)}$). Comme, en raison de la séparabilité, les $\mathfrak P_\Gamma^{(i)}$ sont tous distincts, ils sont premiers entre eux deux à deux (la somme de deux quelconques d'entre eux est égale à $\mathfrak K_\Gamma$). Par conséquent, le rang des idéaux
\[
\mathfrak P_\Gamma^{(1)},
[\mathfrak P_\Gamma^{(1)},\mathfrak P_\Gamma^{(2)}],\ldots,
[\mathfrak P_\Gamma^{(1)},\ldots,\mathfrak P_\Gamma^{(n)}]
\]
diminue d'au moins une unité à chaque étape ; l'intersection de tous ces idéaux est l'idéal nul, tandis que l'intersection de moins de tous les conjugués doit être distincte de l'idéal nul, puisqu'ils sont premiers entre eux deux à deux (le rang a donc diminué exactement d'une unité à chaque étape). Des théorèmes connus donnent alors la représentation suivante, uniquement déterminée, de $\mathfrak K_\Gamma$ comme somme directe :\srcfn{17)}{Voir, par exemple, E. Noether, \emph{Abstrakter Aufbau der Idealtheorie in algebraischen Zahl- und Funktionenkörpern}, Math. Ann. 96 (1927), 26--61, \S{} 4, 5.}
\begin{align}
\mathfrak K_\Gamma
&=\mathfrak C_\Gamma^{(1)}+\cdots+\mathfrak C_\Gamma^{(n)}
=\Gamma e^{(1)}+\cdots+\Gamma e^{(n)}
=\mathfrak C_\Gamma^{(i)}+\mathfrak P_\Gamma^{(i)}, \tag{1}\\
e&=e^{(1)}+\cdots+e^{(n)},\qquad
\mathfrak C_\Gamma^{(i)}=
[\mathfrak P_\Gamma^{(1)},\ldots,\mathfrak P_\Gamma^{(i-1)},
\mathfrak P_\Gamma^{(i+1)},\ldots].\notag
\end{align}
Il faut montrer que $\mathfrak C_\Gamma$ est égal à $\mathfrak A_\Gamma$. Cela résultera de l'égalité correspondante dans $\mathfrak K_{\mathsf{K}}$. On a ici la décomposition
\begin{equation}
\mathfrak K_{\mathsf{K}}=\mathfrak C_{\mathsf{K}}+\mathfrak P_{\mathsf{K}}
=\mathfrak K_{\mathsf{K}} e^{(1)}+\mathfrak K_{\mathsf{K}}(e-e^{(1)}).\tag{2}
\end{equation}
En effet, $e^{(1)}$ appartient à $\mathfrak K_{\mathsf{K}}$ : $\mathfrak C_\Gamma^{(1)}$ est invariant par toutes les permutations de $\mathfrak P_\Gamma^{(2)},\ldots,\mathfrak P_\Gamma^{(n)}$, et il en va donc de même de $e^{(1)}$, composante de l'unité dans $\mathfrak C_\Gamma^{(1)}$. Si l'on exprime alors $e^{(1)}$ dans une base indépendante de $\mathfrak K_{\mathsf{K}}$ sur $K$ --- ce qui est possible puisque celle-ci est aussi une base indépendante de $\mathfrak K_\Gamma$ sur $\Gamma$ ---, les coefficients sont chacun invariants par toutes les permutations de $\mathsf{K}^{(2)},\ldots,\mathsf{K}^{(n)}$ et appartiennent donc à $\mathsf{K}$. Comme $(e^{(1)})^2=e^{(1)}$, la représentation
\[
\mathfrak K_{\mathsf{K}}=\mathfrak K_{\mathsf{K}} e^{(1)}+\mathfrak K_{\mathsf{K}}(e-e^{(1)})
\]
est une somme directe ; son extension à $\mathfrak K_\Gamma$ reste directe et conduit donc, d'après (1), à $\mathfrak K_\Gamma=\mathfrak C_\Gamma^{(1)}+\mathfrak P_\Gamma^{(1)}$. Ainsi, d'après la fin du \S{} 2, 2,
\[
\mathfrak K_{\mathsf{K}} e^{(1)}=[\mathfrak C_\Gamma,\mathfrak K_{\mathsf{K}}]
=[\Gamma e^{(1)},\mathfrak K_{\mathsf{K}}]=\mathsf{K}e^{(1)}=\mathfrak C_{\mathsf{K}};
\]
\[
\mathfrak K_{\mathsf{K}}(e-e^{(1)})=[\mathfrak P_\Gamma,\mathfrak K_{\mathsf{K}}]=\mathfrak P_{\mathsf{K}},
\]
ce qui démontre (2).

Mais (2) entraîne $\mathfrak C_{\mathsf{K}}=\mathfrak A_{\mathsf{K}}$, et par suite $\mathfrak C_\Gamma=\mathfrak A_\Gamma$. En effet, $e^{(1)}(e-e^{(1)})=0$ donne $\mathfrak C_{\mathsf{K}}\mathfrak P_{\mathsf{K}}=0$ ; réciproquement, $b\mathfrak P_{\mathsf{K}}=0$ donne aussi $b(e-e^{(1)})=0$, donc $b=be=be^{(1)}$, et par conséquent
\[
\mathfrak C_{\mathsf{K}}=(0):\mathfrak P_{\mathsf{K}}=\mathfrak A_{\mathsf{K}}
\]
puis, par extension, $\mathfrak C_\Gamma=\mathfrak A_\Gamma$. On a donc
\begin{align}
\mathfrak K_\Gamma
&=\mathfrak A_\Gamma^{(1)}+\cdots+\mathfrak A_\Gamma^{(n)}
=\Gamma e^{(1)}+\cdots+\Gamma e^{(n)};
&0&=[\mathfrak P_\Gamma^{(1)},\ldots,\mathfrak P_\Gamma^{(n)}]; \tag{3}\\
\mathfrak K_{\mathsf{K}^{(i)}}
&=\mathfrak A_{\mathsf{K}^{(i)}}+\mathfrak P_{\mathsf{K}^{(i)}}^{(i)}
=\mathsf{K}^{(i)}e^{(i)}+\mathfrak P_{\mathsf{K}^{(i)}}^{(i)};
&0&=[\mathfrak A_{\mathsf{K}^{(i)}},\mathfrak P_{\mathsf{K}^{(i)}}^{(i)}].\notag
\end{align}
Le théorème est ainsi démontré dans toutes ses parties. Signalons encore qu'il s'agit d'un \emph{théorème de structure entièrement invariant}, puisque toutes les constructions qui y interviennent ($\mathfrak K_{\mathsf{K}}$, $\mathfrak K_\Gamma$, $\mathfrak P_{\mathsf{K}}$, $\mathfrak A_{\mathsf{K}}$) sont définies sans choix de base.\srcfn{18)}{Une formulation au moyen d'une équation définissante et de la décomposition du polynôme correspondant en facteurs linéaires ne serait pas invariante.}

\paragraph{3. Représentation de $\mathfrak K$ par des composantes conjuguées.}
Aux idéaux des différences et aux quotients des différences conjugués correspond une représentation des éléments de $\mathfrak K$ par des composantes conjuguées.

\paragraph{Théorème.}
\emph{Si
\[
x=\xi^{(1)}e^{(1)}+\cdots+\xi^{(n)}e^{(n)}
\]
est la représentation par composantes d'un élément $x$ de $\mathfrak K$, alors les $n$ composantes sont conjuguées entre elles. En particulier, $\xi^{(1)},\ldots,\xi^{(n)}$ sont les $n$ valeurs conjuguées qui correspondent à $x$ par les isomorphismes de $\mathfrak K$ sur $\mathsf{K}^{(1)},\ldots,\mathsf{K}^{(n)}$.}

Les relations (3) du no 2 montrent que $e^{(1)},\ldots,e^{(n)}$ sont conjugués entre eux : les automorphismes de $\Gamma$ qui envoient $\mathfrak P_{\mathsf{K}}$ sur $\mathfrak P_{\mathsf{K}^{(i)}}^{(i)}$ envoient en même temps $e^{(1)}$ sur $e^{(i)}$, ce dernier étant bien un élément de $\mathfrak K_{\mathsf{K}^{(i)}}$. Les mêmes relations donnent $xe^{(i)}=\xi^{(i)}e^{(i)}$, puisque $\mathfrak P_{\mathsf{K}^{(i)}}^{(i)}e^{(i)}=0$, et par suite
\[
(x-\xi^{(i)})e^{(i)}=0
\]
(on peut aussi le déduire du fait que $\mathfrak K$, étant un corps, est isomorphe à sa composante, laquelle est donc donnée par $\mathsf{K}^{(i)}e^{(i)}$, avec correspondance isomorphe entre $x$ et $\xi^{(i)}e^{(i)}$).

Le théorème entraîne le

\paragraph{Corollaire.}
\emph{Si $\mathfrak B$ est un module absolu dans $\mathfrak K$ (c'est-à-dire que, avec deux éléments quelconques, $\mathfrak B$ contient aussi leur différence), et si $\mathfrak B^{(1)},\ldots,\mathfrak B^{(n)}$ sont ses composantes dans $\mathfrak K_\Gamma$, alors $\mathfrak B$ est égal à l'intersection $[\mathfrak K,\mathfrak B^{(1)}+\cdots+\mathfrak B^{(n)}]$.}

En effet, $\mathfrak B$ est contenu dans cette intersection ; réciproquement, pour qu'un élément de $\mathfrak B^{(1)}+\cdots+\mathfrak B^{(n)}$ appartienne à $\mathfrak K$, il doit, d'après le théorème, être une somme de conjugués et est donc un élément de $\mathfrak B$.

\paragraph{4. Bases complémentaires de $\mathfrak K$ et leur relation avec les composantes de l'unité.}
Lorsque l'on passe de la structure invariante à une représentation au moyen d'une base, cette structure permet de regrouper les bases par paires complémentaires.

\paragraph{Théorème.}
\emph{À toute $\mathsf{P}$-base $t_1,\ldots,t_n$ de $\mathfrak K$ est associée une base complémentaire $T_1,\ldots,T_n$ de $\mathfrak K$, dont la complémentaire est à son tour la base initiale des $t$. Les bases $\mathsf{A}_1^{(i)},\ldots,\mathsf{A}_n^{(i)}$ de $\mathsf{K}^{(i)}$, isomorphes à la base complémentaire $T_1,\ldots,T_n$, sont définies comme les coefficients des composantes $e^{(i)}$ de l'unité dans la représentation, induite par les $t$,
\[
e^{(i)}=\mathsf{A}_1^{(i)}t_1+\cdots+\mathsf{A}_n^{(i)}t_n.
\]
Si $\alpha_1^{(i)},\ldots,\alpha_n^{(i)}$ sont les bases correspondant aux $t_i$, les matrices des $\alpha$ et des $\mathsf{A}$ sont inverses. Si deux bases $(s)$ et $(t)$ de $\mathfrak K$ sont reliées par une transformation
\[
(s)=(t)P
\]
leurs bases complémentaires sont reliées de manière contragrédiente par
\[
(S)=P^{-1}(T).
\]}

\paragraph{Démonstration.}
Soient $e^{(i)}=A_1^{(i)}t_1+\cdots+A_n^{(i)}t_n$ les représentations, dans la base, des $n$ composantes $e^{(i)}$ de l'unité, lorsqu'on prend pour point de départ une $P$-base $t_1,\ldots,t_n$ de $\mathfrak K$, qui est donc une $\Gamma$-base de $\mathfrak K_\Gamma$. D'après 2 et 3, $A_1^{(i)},\ldots,A_n^{(i)}$ appartiennent à $K^{(i)}$, et des systèmes de coefficients $A$ conjugués correspondent aux différents $e^{(i)}$. Si $\alpha_1^{(i)},\ldots,\alpha_n^{(i)}$ désigne la base de $K^{(i)}$ isomorphiquement associée à $t_1,\ldots,t_n$, on a encore les relations
\begin{equation}
e^{(i)}[t\to\alpha^{(i)}]=e;
\qquad e^{(i)}[t\to\alpha^{(j)}]=0\quad\text{pour }i\ne j. \tag{1}
\end{equation}
En effet, $\mathfrak P_{\mathsf{K}^{(i)}}^{(i)}$ contient les éléments $t_1-\alpha_1^{(i)},\ldots,t_n-\alpha_n^{(i)}$ --- qui forment une base dépendante de $\mathfrak P_{\mathsf{K}^{(i)}}^{(i)}$ comme $P$-module --- et s'annule donc sous la substitution $t=\alpha^{(i)}$. La première relation (1) résulte par conséquent de
\[
e\equiv e^{(i)}\;(\mathfrak P_{\mathsf{K}^{(i)}}^{(i)})
\]
tandis que la seconde résulte de ce que $e^{(i)}$ est un élément de tous les $\mathfrak P_{\mathsf{K}^{(j)}}^{(j)}$ pour $i\ne j$. Développées coefficient par coefficient, les relations (1) s'écrivent sous forme matricielle
\begin{equation}
\begin{pmatrix}
\alpha_1^{(1)}&\ldots&\alpha_n^{(1)}\\
\vdots&&\vdots\\
\alpha_1^{(n)}&\ldots&\alpha_n^{(n)}
\end{pmatrix}
\begin{pmatrix}
A_1^{(1)}&\cdots&A_1^{(n)}\\
\vdots&&\vdots\\
A_n^{(1)}&\cdots&A_n^{(n)}
\end{pmatrix}=E_n. \tag{2}
\end{equation}
où $E_n$ désigne la matrice unité ; les $\mathsf{A}$ sont donc déjà déterminés de manière unique par les $\alpha$. La relation (2) montre que le déterminant des deux matrices ne s'annule pas. Elle montre en outre que $\mathsf{A}_1^{(i)},\ldots,\mathsf{A}_n^{(i)}$ sont linéairement indépendants relativement à $\mathsf{P}$ : une relation de dépendance linéaire serait aussi satisfaite par les systèmes conjugués et diminuerait donc le rang de la matrice des $\mathsf{A}$, contrairement au fait que le rang de $E_n$ est égal à $n$. Ainsi, $\mathsf{A}_1^{(i)},\ldots,\mathsf{A}_n^{(i)}$ forment pour chaque $i$ une base de $\mathsf{K}^{(i)}$ ; comme ces bases sont conjuguées, elles sont toutes isomorphes à une même base $T_1,\ldots,T_n$ de $\mathfrak K$, où, d'après 3, $T_\lambda$ est déterminé par
\[
T_\lambda=\mathsf{A}_\lambda^{(1)}e^{(1)}+\cdots+\mathsf{A}_\lambda^{(n)}e^{(n)}
\]
La base $T_1,\ldots,T_n$ est dite complémentaire de $t_1,\ldots,t_n$ ; de même, $\mathsf{A}_1^{(i)},\ldots,\mathsf{A}_n^{(i)}$ est dite complémentaire de $\alpha_1^{(i)},\ldots,\alpha_n^{(i)}$. Comme les $\mathsf{A}$ sont déjà déterminés de manière unique par (2), et comme cette relation --- par passage à la matrice transposée --- détermine les $\alpha$ à partir de la base des $T$, donc des $A$, la base initiale des $t$ est la complémentaire de celle des $T$. Par la correspondance induite par les composantes de l'unité, toutes les bases se répartissent ainsi en paires de bases complémentaires. Soit $s,S$ une autre paire de bases complémentaires, avec
\[
e^{(i)}=\mathsf{B}_1^{(i)}s_1+\cdots+\mathsf{B}_n^{(i)}s_n
\qquad\text{et}\qquad (s)=(t)P,
\]
De l'identité
\[
\mathsf{A}_1t_1+\cdots+\mathsf{A}_nt_n=\mathsf{B}_1s_1+\cdots+\mathsf{B}_ns_n
\]
--- considérée identiquement en $t$ après la substitution $(s)=(t)P$ ---, il résulte que $\mathsf{B},\mathsf{A}$ se transforment de manière contragrédiente à $s,t$. En raison de l'isomorphisme, il en va donc de même de $S$ et $T$. Le théorème est ainsi démontré dans toutes ses parties.

\paragraph{Remarque.}
Si $c=C_1t_1+\cdots+C_nt_n$ est la représentation d'un élément de $\mathfrak K$ au moyen d'une $\mathsf{P}$-base et $c=c_1T_1+\cdots+c_nT_n$ la représentation du même élément dans la base complémentaire, alors
\[
C_\lambda=\operatorname{Sp}(cT_\lambda),
\qquad c_\lambda=\operatorname{Sp}(ct_\lambda),
\]
où $\operatorname{Sp}(x)$ désigne, comme d'ordinaire, la somme des conjugués :
\[
\operatorname{Sp}(x)=\xi^{(1)}+\xi^{(2)}+\cdots+\xi^{(n)}
\]
pour tout $x$ de $\mathfrak K$. En effet,
\[
c=\gamma^{(1)}e^{(1)}+\cdots+\gamma^{(n)}e^{(n)}
=\sum\limits_\lambda\gamma^{(1)}\mathsf{A}_\lambda^{(1)}t_\lambda+
\cdots+\sum\limits_\lambda\gamma^{(n)}\mathsf{A}_\lambda^{(n)}t_\lambda,
\]
et la comparaison des coefficients donne $C_\lambda=\operatorname{Sp}(cT_\lambda)$. La seconde relation s'obtient de même en échangeant $t$ et $T$.\srcfn{19)}{Les considérations de ce numéro montrent que les théorèmes de structure invariants sont essentiellement plus simples que les théorèmes non invariants relatifs aux bases. Ces derniers n'ont été employés que pour établir le lien avec les résultats connus. En fait, la décomposition de l'anneau d'extension galoisien et l'introduction des composantes de l'unité fournissent tout ce que donnent les bases complémentaires (et le module complémentaire au paragraphe suivant). Là encore, le lien entre la différente et le module complémentaire n'est donné qu'à des fins de classement. Il en va de même pour l'étude de la formule d'interpolation de Lagrange au no 5.}

\paragraph{5. Choix d'une équation définissante. Formule d'interpolation de Lagrange.}
[Note des éditeurs : ce numéro est resté non rédigé dans le manuscrit.]

\subsection*{\S{} 6. La différente d'un ordre.}

Les théorèmes de structure qui ne concernent encore que les corps et leurs anneaux d'extension galoisiens vont maintenant être renforcés dans le sens de l'intégralité : on distinguera certains sous-anneaux de $\mathsf{P}$ ainsi que les éléments de $\mathfrak K$ et de $\mathsf{K}$ entiers relativement à ceux-ci, c'est-à-dire les ordres. Toutes les notations du \S{} 5 sont conservées.\srcfn{20)}{Pour les notions d'élément entier relativement à un anneau et d'anneau intégralement fermé, voir le \S{} 1 de l'article cité à la note 17.}

\paragraph{1. Les anneaux considérés.}
Soit $\mathfrak h$ un sous-anneau de $\mathsf{P}$ tel que $\mathsf{P}$ soit le corps des fractions de $\mathfrak h$ ; supposons en outre $\mathfrak h$ \emph{intégralement fermé dans $\mathsf{P}$} (tout élément de $\mathsf{P}$ entier relativement à $\mathfrak h$ appartient à $\mathfrak h$). Soit encore $\mathfrak O$ un $\mathfrak h$-ordre de $\mathfrak K$, c'est-à-dire un sous-anneau de $\mathfrak K$ contenant $\mathfrak h$, dont le corps des fractions est $\mathfrak K$ et qui possède une base finie, non nécessairement indépendante, comme $\mathfrak h$-module. Tous les éléments de $\mathfrak O$ sont alors entiers (relativement à $\mathfrak h$) ; la base de $\mathfrak O$ comme $\mathfrak h$-module est en même temps une base, en général dépendante, de $\mathfrak K$ comme $P$-module. Notons $\mathfrak o$ l'ordre de $\mathsf{K}$ isomorphe à $\mathfrak O$, et $\mathfrak o^{(i)}$ ses conjugués dans $\mathsf{K}^{(i)}$. Soit $\mathfrak g$ l'anneau de $\Gamma$ engendré par les $n$ ordres conjugués ; $\mathfrak g$ est lui aussi un ordre. En effet, $\mathfrak g$ contient $\mathfrak h$. Les produits de tous les éléments des bases des $\mathfrak o^{(i)}$ comme $\mathfrak h$-modules forment une base de $\mathfrak g$ comme $\mathfrak h$-module, et en même temps une base de $\Gamma$ comme $P$-module.

\emph{Les produits directs $\mathfrak O_{\mathfrak o}$ et $\mathfrak O_{\mathfrak g}$ existent,} comme le montre le critère suffisant du \S{} 2, 3 (il faut y remplacer l'anneau $\mathfrak f$ par $\mathsf{P}$). Dans $\mathfrak K_\Gamma$ existe d'abord le produit direct $\mathfrak O_{\mathsf{P}}=\mathfrak K$ relativement à $\mathfrak h$. En effet, $\mathfrak h$ est contenu dans l'intersection $[\mathfrak O,\mathsf{P}]$ ; réciproquement, tout élément de l'intersection est entier parce qu'il appartient à $\mathfrak O$, donc appartient à $\mathfrak h$ parce qu'il est un élément de $\mathsf{P}$ et que $\mathfrak h$ est intégralement fermé. Les combinaisons linéaires à coefficients dans $\mathsf{P}$ de $n$ éléments linéairement indépendants quelconques de $\mathfrak O$ donnent alors $\mathfrak K$ comme produit direct. De même,
\[
\mathfrak o_{\mathsf{P}}=\mathsf{K}
\qquad\text{et}\qquad
\mathfrak g_{\mathsf{P}}=\Gamma;
\]
du produit direct $\mathfrak K_{\mathsf{K}}=\mathfrak O_{\mathsf{P}}\times\mathfrak o_{\mathsf{P}}$ relativement à $\mathfrak h$, le critère du \S{} 2, 3 donne l'existence de $\mathfrak O_{\mathfrak o}$ relativement à $\mathfrak h$ ; de même, $\mathfrak K_\Gamma=\mathfrak O_{\mathsf{P}}\times\mathfrak g_{\mathsf{P}}$ donne l'existence de $\mathfrak O_{\mathfrak g}$. Avec $\mathfrak O_{\mathfrak o}$ existent aussi les anneaux conjugués $\mathfrak O_{\mathfrak o^{(i)}}$. On appelle $\mathfrak O_{\mathfrak g}$ l'\emph{anneau d'extension galoisien de l'ordre $\mathfrak O$}.

La différente de $\mathfrak O$, respectivement de $\mathfrak o$, relativement à $\mathfrak h$ est donc définie ; il en va de même des différentes conjuguées. Notons
\[
\mathfrak P=\{\ldots,y-\eta,\ldots\}
\]
l'idéal des différences de $\mathfrak O_{\mathfrak o}$, et
\[
\mathfrak P^{(i)}=\{\ldots,y-\eta^{(i)},\ldots\}
\]
les idéaux conjugués dans $\mathfrak O_{\mathfrak o^{(i)}}$ ; notons de même
\[
\mathfrak A^{(i)}=(0):\mathfrak P^{(i)}
\quad\text{--- le quotient étant pris dans $\mathfrak O_{\mathfrak o^{(i)}}$ ---}
\]
les \emph{quotients des différences} conjugués, et en particulier
\[
\mathfrak A=(0):\mathfrak P\quad\text{dans }\mathfrak O_{\mathfrak o}.
\]
Notons encore
\[
\mathfrak D=\mathfrak A[\eta\to y]
\]
la différente de $\mathfrak O$, et
\[
\mathfrak d^{(i)}=\mathfrak A[y\to\eta^{(i)}]
\]
les différentes des $n$ ordres conjugués $\mathfrak o^{(i)}$, en particulier
\[
\mathfrak d=\mathfrak A[y\to\eta]
\]
la différente de $\mathfrak o$.\srcfn{21)}{L'hypothèse que $\mathfrak h$ est intégralement fermé dans $\mathsf{P}$ est satisfaite, par exemple, lorsque $\mathfrak h$ est l'anneau des entiers d'un corps de nombres, et en particulier lorsque $\mathfrak h$ est l'anneau des entiers rationnels ; la différente et la différente relative d'ordres arbitraires, y compris dans des corps relatifs, sont donc définies. On peut aussi prendre pour $\mathfrak h$ l'anneau de polynômes en $n$ indéterminées, ce qui donne la différente dans les corps de fonctions algébriques en $n$ indéterminées.}

\paragraph{2. Théorèmes de structure.}
Au théorème de structure pour les anneaux d'extension galoisiens de corps (\S{} 5, 2) correspond le

\paragraph{Théorème de structure pour l'anneau d'extension galoisien des ordres.}
\emph{L'idéal des différences et le quotient des différences relatifs à l'ordre et au corps sont liés comme suit : $\mathfrak P_{\mathsf{K}}$ est l'extension de l'idéal des différences $\mathfrak P$ de $\mathfrak O_{\mathfrak o}$ ; de même, $\mathfrak A_{\mathsf{K}}$ est l'extension du quotient des différences $\mathfrak A$ de $\mathfrak O_{\mathfrak o}$, et $\mathfrak A$ est l'idéal contracté de $\mathfrak A_{\mathsf{K}}$. Entre le quotient des différences, la différente et la composante de l'unité, on a les relations}
\[
\mathfrak A=\mathfrak d e^{(1)},
\qquad
\mathfrak A^{(i)}=\mathfrak d^{(i)}e^{(i)},
\]
\emph{et par conséquent}
\[
\mathfrak d=[\mathfrak o,\mathfrak A^{(1)}+\cdots+\mathfrak A^{(n)}].
\]
\emph{La différente $\mathfrak d$ de $\mathfrak o$ est donc constituée de tous les éléments $x$ de $\mathsf{K}$ tels que $\varkappa e^{(1)}$ appartienne à $\mathfrak O_{\mathfrak o}$ ; elle est distincte de l'idéal nul.}

Le fait que $\mathfrak P_{\mathsf{K}}$ soit l'extension de $\mathfrak P$ résulte de $\mathfrak K=\mathfrak O_P$ et de $K=\mathfrak o_{\mathsf{P}}$ : chaque $x-\xi$ de $\mathfrak P_{\mathsf{K}}$ s'exprime donc linéairement, à coefficients dans $\mathsf{K}$, au moyen des $y-\eta$ de $\mathfrak P$. Il en résulte en outre que $\mathfrak A_{\mathsf{K}}=(0):\mathfrak P_{\mathsf{K}}$ est aussi défini comme le quotient $(0):\mathfrak P$ dans $\mathfrak K_{\mathsf{K}}$. Ainsi, $\mathfrak A=(0):\mathfrak P$ dans $\mathfrak O_{\mathfrak o}$ est constitué de tous les éléments de $\mathfrak O_{\mathfrak o}$ qui appartiennent à $\mathfrak A_{\mathsf{K}}$ ; on a $\mathfrak A=[\mathfrak O_{\mathfrak o},\mathfrak A_{\mathsf{K}}]$, de sorte que $\mathfrak A$ est l'idéal contracté de $\mathfrak A_{\mathsf{K}}$. Comme $\mathfrak A_{\mathsf{K}}=\mathsf{K}e^{(1)}$, on a donc $\mathfrak A=\mathfrak z e^{(1)}$, où $\mathfrak z$ est un $\mathfrak o$-module dans $\mathsf{K}$, puisque $\mathfrak A$ est lui aussi un $\mathfrak o$-module. Mais, par définition de la différente,
\[
\mathfrak d=\mathfrak A[x\to\xi]
=(\mathfrak z e^{(1)})[x\to\xi]
=\mathfrak z\bigl(e^{(1)}[x\to\xi]\bigr)=\mathfrak z e=\mathfrak z
\]
(voir \S{} 5, 4, (1)). Ainsi $\mathfrak z=\mathfrak d$ ; le quotient des différences $\mathfrak A$ est la composante dans $\mathsf{K}$ de la différente $\mathfrak d$ de $\mathfrak o$, et de même pour les conjugués. Le corollaire du \S{} 5, 3 entraîne alors la relation ci-dessus pour $\mathfrak d$. Et
\[
\mathfrak A=\mathfrak d e^{(1)}=[\mathsf{K}e^{(1)},\mathfrak O_{\mathfrak o}]
\]
montre que $\mathfrak d$ est constitué de tous les éléments $x$ de $\mathsf{K}$ pour lesquels $\varkappa e^{(1)}$ appartient à $\mathfrak O_{\mathfrak o}$. Ainsi $\mathfrak d$ est distincte de l'idéal nul : comme élément de $\mathfrak K_{\mathsf{K}}=\mathfrak O_{\mathsf{K}}$, $e^{(1)}$ est une combinaison linéaire d'un nombre fini d'éléments de $\mathfrak O$ à coefficients dans $\mathsf{K}$ ; c'est donc le quotient d'un élément de $\mathfrak O_{\mathfrak o}$ par un élément de $\mathfrak o$ (même par un élément de $\mathfrak h$), et ce dénominateur donne un élément non nul de $\mathfrak d$. Par conséquent, $\mathfrak d$ est de rang $n$ relativement à $\mathsf{P}$, ou $\mathfrak d_{\mathsf{P}}=\mathsf{K}$, de sorte que $\mathfrak A_{\mathsf{K}}$ est l'extension de $\mathfrak A$.

\begin{center}
\textbf{À partir d'ici, esquisse.}
\end{center}

\paragraph{3. Relation avec le module complémentaire lorsqu'il existe une base indépendante.}
Si $t_1,\ldots,t_n$ est une base indépendante de $\mathfrak O$ comme $\mathfrak h$-module, alors toute base $(s)=(t)P$ de $\mathfrak K$ l'est aussi dès que $P$ est à coefficients dans $\mathfrak h$ et unimodulaire. Mais $(S)=P^{-1}(T)$ montre alors qu'avec $T$, $S$ est aussi une base d'un même $\mathfrak h$-module, appelé module complémentaire de $\mathfrak O$.

La dernière partie du théorème de structure du no 2 affirme alors que $\mathfrak d$ est égale au quotient
\[
\mathfrak o:\mathfrak e
\]
dans $\mathsf{K}$, où $\mathfrak e$ désigne le module complémentaire. En effet, soit
\[
e^{(1)}=\mathsf{A}_1t_1+\cdots+\mathsf{A}_nt_n,
\]
où les $\mathsf{A}_\lambda$ forment une base de $\mathfrak e$ ; alors $\mathfrak d$ est constitué de tous les éléments $\delta$ tels que $\delta e^{(1)}$ appartienne à $\mathfrak O_{\mathfrak o}$, c'est-à-dire tels que $\delta \mathsf{A}_\lambda$ appartienne à $\mathfrak o$, donc $\mathfrak d=\mathfrak o:\mathfrak e$.

\paragraph{Cas particulier.}
Ordre régulier ayant pour base comme $\mathfrak h$-module $e,z,\ldots,z^{n-1}$ :
\[
\mathfrak d=\{f'(z)\}.
\]

\paragraph{Remarque.}
La remarque du \S{} 5, 4 donne l'autre définition connue du module complémentaire $\mathfrak E$ de $\mathfrak O$. \emph{$\mathfrak E$ est constitué de tous les éléments $c$ de $\mathfrak K$ pour lesquels les traces $\operatorname{Sp}(ct)$, c'est-à-dire les traces $\operatorname{Sp}(co)$ pour tout $o$ de $\mathfrak O$, sont entières, c'est-à-dire appartiennent à $\mathfrak h$} ; il en va de même de $\mathfrak e$ relativement à $\mathfrak o$ et à $\mathsf{K}$.

En effet, soit $c=c_1T_1+\cdots+c_nT_n$, avec $c_i$ dans $\mathfrak h$, une représentation d'un élément $c$ de $\mathfrak E$ dans la base de $\mathfrak K$ complémentaire de $t$. Alors $c_\lambda=\operatorname{Sp}(ct_\lambda)$ appartient à $\mathfrak h$ ; réciproquement, si $\operatorname{Sp}(ct_\lambda)$ appartient à $\mathfrak h$, alors $c_\lambda$ appartient à $\mathfrak h$, et donc $c$ à $\mathfrak E$.

\paragraph{Complément.}
\emph{S'il existe dans l'ordre une base indépendante comme $\mathfrak h$-module, alors le quotient des différences $\mathfrak A^{(i)}$ est l'intersection de tous les idéaux des différences distincts de $\mathfrak P^{(i)}$ ; plus précisément,}
\[
\mathfrak A^{(i)}=
[\mathfrak O_{\mathfrak o^{(i)}},
\mathfrak P_{\mathfrak g}^{(1)},\ldots,
\mathfrak P_{\mathfrak g}^{(i-1)},
\mathfrak P_{\mathfrak g}^{(i+1)},\ldots,
\mathfrak P_{\mathfrak g}^{(n)}].
\]
En effet, d'après la forme renforcée du \S{} 2, 2, $\mathfrak P$ est alors aussi l'idéal contracté de $\mathfrak P_{\mathsf{K}}$ : si $e,t_2,\ldots,t_n$ est une base de $\mathfrak O$ comme module, donc aussi de $\mathfrak O_{\mathfrak o}$, alors
\[
e,t_2-\alpha_2,\ldots,t_n-\alpha_n;
\]
en est également une ; mais $t_2-\alpha_2,\ldots,t_n-\alpha_n$ est en même temps une base de $\mathfrak P_{\mathsf{K}}$ comme $\mathsf{K}$-module. Les hypothèses du \S{} 2, 2 (complément) sont donc satisfaites. Or on avait
\[
\mathfrak A_\Gamma=[\mathfrak P_\Gamma^{(2)},\ldots,\mathfrak P_\Gamma^{(n)}];
\qquad
\mathfrak A_\mathfrak g=[\mathfrak A_\Gamma,\mathfrak O_\mathfrak g]
=[\ldots,[\mathfrak P_\Gamma^{(\lambda)},\mathfrak O_\mathfrak g],\ldots]
=[\mathfrak P_\mathfrak g^{(2)},\ldots,\mathfrak P_\mathfrak g^{(n)}].
\]
Mais $\mathfrak A=[\mathfrak A_\mathfrak g,\mathfrak O_{\mathfrak o}]$ ; par conséquent,
\[
\mathfrak A=[\mathfrak O_{\mathfrak o},\mathfrak P_\mathfrak g^{(2)},\ldots,\mathfrak P_\mathfrak g^{(n)}].
\]

\paragraph{4. Relation avec le quotient différentiel d'une équation fondamentale, lorsqu'une telle équation existe.}
Si l'on étend l'ordre $\mathfrak O$ par $n$ indéterminées $u_1,\ldots,u_n$, c'est-à-dire si l'on forme le produit direct $\mathfrak O_{\mathfrak U}$, où $\mathfrak U$ désigne l'anneau de polynômes $\mathfrak h[u_1,\ldots,u_n]$, donc $\mathfrak O_{\mathfrak U}=\mathfrak O[u_1,\ldots,u_n]$, l'extension $\mathfrak C_{\mathfrak U}$ de tout idéal $\mathfrak C$ de $\mathfrak O$ est constituée de tous les polynômes en $u$ à coefficients dans $\mathfrak C$ ; réciproquement, $\mathfrak C$ est caractérisé comme le système de tous les coefficients de $\mathfrak C_{\mathfrak U}$.

Supposons maintenant qu'il existe une équation fondamentale, c'est-à-dire que $e,U,\ldots,U^{n-1}$ forme une base comme $\mathfrak h$-module après passage à la localisation obtenue en inversant un ou plusieurs polynômes primitifs $H(u)$ de $\mathfrak h[u]$. Les idéaux $\mathfrak C_{\mathfrak U}$ deviennent alors des systèmes obtenus à partir de $\mathfrak C_{\mathfrak U}$ par multiplication par des unités ($H(u)$ et ses puissances négatives et positives). Pour retrouver $\mathfrak C$, on peut donc d'abord, par multiplication par ces unités, revenir à $\mathfrak C_{\mathfrak U}$ (entier en $u$), puis de celui-ci à $\mathfrak C$. En particulier, soit $\mathfrak D_U$ l'idéal auquel passe la différente $\mathfrak D$ de $\mathfrak O$ ; $\mathfrak D_{\mathfrak U}$ est donc la différente de $\mathfrak O_{\mathfrak U}$. Mais, dans $\mathfrak O_{\mathfrak U}$ --- après localisation par les $H(u)$ primitifs ---, il existe une équation définissante $G(U)=0$. Ainsi $G'(U)$ est un polynôme générateur de la différente, c'est-à-dire que celle-ci se déduit des coefficients de $G'(U)$. (Il faut encore montrer plus précisément que $(G'(U))$ ne contient réellement que $\mathfrak D_{\mathfrak U}$, et rien de plus ; peut-être cela n'est-il vrai que lorsque $\mathfrak o$ est intégralement fermé dans $\mathsf{K}$. La multiplication par des diviseurs de $H(u)$ pourrait en effet encore faire apparaître, comme $\mathfrak O_{\mathfrak U}$, d'autres éléments appartenant à $\mathfrak D_{\mathfrak U}$. Dans le cas de la fermeture intégrale (5 axiomes), cela résulte du théorème selon lequel les idéaux des coefficients, ou contenus, se multiplient. Il devient alors possible d'obtenir une théorie de la ramification en considérant l'équation fondamentale modulo $p$.)

\subsection*{\S{} 7. Théorie de la ramification de l'anneau des entiers modulo $p^t$.}

Plaçons-nous dans un corps de nombres et soit $\mathfrak h$ son anneau d'entiers ; on peut aussi se placer dans un corps relatif sur la localisation de l'anneau des entiers en un idéal premier, auquel cas $p$ désigne un générateur de l'idéal premier dans cet anneau localisé. Adjoignons les indéterminées $u_1,\ldots,u_n$ ; l'anneau des entiers aussi bien que l'anneau quotient par $p^t$ possèdent alors la base $e,U,\ldots,U^{n-1}$ comme $\mathfrak h$-module ; dans les deux cas, la différente est donc donnée par $G'(u)$. \emph{La différente de l'anneau des entiers passe modulo $p^t$ à celle de l'anneau quotient par $p^t$} (cela vaut-il pour des ordres arbitraires, ou non ?). Mais si
\[
p=\mathfrak p_1^{\varrho_1}\cdots\mathfrak p_r^{\varrho_r},
\]
alors l'anneau quotient $\mathfrak R$ par $p^t$ est une somme directe dont les composantes sont isomorphes aux anneaux quotients par $\mathfrak p_i^{t\varrho_i}$ ; ainsi
\[
\mathfrak R=\mathfrak R e_1+\cdots+\mathfrak R e_r;
\]
Pour $i\ne k$, $e_k$ est divisible par $\mathfrak p_i^{t\varrho_i}$ ; par conséquent, si
\[
\mathfrak d[p^t]=\mathfrak d[p^t]e_1+\cdots+\mathfrak d[p^t]e_r
=\mathfrak d_1+\cdots+\mathfrak d_r
\]
et si $t$ est assez grand, $\mathfrak p_i$ divise $\mathfrak d_i$ à la même puissance que $\mathfrak d[p^t]$, donc que $\mathfrak d$. Or, d'après \S{} 4, $\mathfrak d_i$ est la différente de $\mathfrak R_i$ ; il suffit donc de considérer la différente de l'anneau quotient par $\mathfrak p^{\varrho t}$. Choisissons $\xi$ (\emph{Zahlbericht}, théorèmes 29 et 30) de manière que $\xi$ soit un entier et que
\[
\varphi(\xi)\equiv0(\mathfrak p),
\quad \not\equiv0(\mathfrak p^2),
\]
où $\varphi(x)$ est un polynôme irréductible modulo $p$ de degré $f$. Si, pour $t\geqq2$, on remplace ce $\xi$ par $\xi^*=\xi e_1$, où $e_1$ est le premier idempotent qui intervient dans la décomposition de l'anneau quotient $\mathfrak R$ par $p^t$, alors, en plus de $\varphi(\xi^*)\equiv0(\mathfrak p)$, $\not\equiv0(\mathfrak p^2)$, on a même $\mathfrak p_1\equiv(p,\varphi(\xi^*))$. Une base relativement aux entiers modulo $p^t$ est alors constituée par les $\xi^\mu\varphi(\xi)^\nu$, avec $\mu=0,1,\ldots,f-1$ et $\nu=0,1,\ldots,t-1$. Une base indépendante sur $\mathfrak h$ (classes modulo $p^t$) est donnée par
\begin{equation}
\xi^{\mu_0}\varphi(\xi)^{\nu_0}
\quad\text{avec}\quad
\mu_0=0,1,\ldots,f-1;
\quad \nu_0=0,1,\ldots,\varrho-1. \tag{1}
\end{equation}
En effet, $(\varphi(\xi),\mathfrak p^{\varrho t})=\mathfrak p$ ; comme $(p)=\mathfrak p^\varrho$ dans l'anneau quotient, le passage des idéaux principaux aux éléments donne $\varphi(\xi)=pM(\xi)$, où $M(\xi)$ est une unité de l'anneau quotient, c'est-à-dire $M(\xi)\not\equiv0(\mathfrak p)$. On peut donc (puisque $p^t=0$) exprimer successivement $\xi^{\varrho t}$ à l'aide de puissances inférieures, d'abord dans $M(\xi)$, puis dans l'anneau quotient ; (1) est donc bien une base. L'équation définissante relativement à cette base est
\[
F(x)=\varphi(x)^\varrho+pM(x),
\]
où $M(x)$ n'est pas divisible modulo $p$ par $\varphi(x)$ (Ø. Ore, Math. Ann. 96, p. 348), et par conséquent
\[
M(\xi)\not\equiv0(\mathfrak p).
\]
Mais la base est aussi indépendante. En effet, de
\[
a_0(\xi)+a_1(\xi)\varphi(\xi)+\cdots+a_{\varrho-1}(\xi)\varphi(\xi)^{\varrho-1}
\equiv0(\mathfrak p^{t\varrho})
\]
il résulte $\equiv0(\mathfrak p)$ ; donc $a_0(\xi)\equiv0(\mathfrak p)$, et par suite
\[
a_0(x)\equiv0(p);
\]
En divisant par $\mathfrak p$, on obtient successivement
\[
a_0(x)\equiv0(p),\ldots,a_{\varrho-1}(x)\equiv0(p)
\]
puis, en divisant encore par $p$,
\[
a_0(\alpha)\equiv0(p^2),\ldots,a_{\varrho-1}(\alpha)\equiv0(p^2)
\]
et ainsi de suite jusqu'à
\[
a_0(\alpha)\equiv0(p^t),\ldots,a_{\varrho-1}(\alpha)\equiv0(p^t).
\]
La différente de l'anneau quotient est donc donnée par le quotient différentiel polynomial $F'(\xi)$ --- tout est ainsi ramené à Ore, p. 345, \S{} 3. Le théorème 10 d'Ore (voir Dedekind--Hensel) montre que $p^t=p^n$ suffit pour tout $\mathfrak p$ ; les nombres supplémentaires résultent en outre des travaux d'Ore.

Tout ce qui précède vaut en même temps --- par passage à une localisation --- pour les différentes relatives ; ainsi, Ore, Math. Ann. 97, p. 594, \S{} 4, résulte immédiatement d'Ore, Math. Ann. 96. Il suffit d'ajouter les énoncés sur l'expression de la différente du corps supérieur comme produit de la différente du corps inférieur et de la différente relative (voir l'introduction du présent article).

\begin{center}
Reçu le 25 octobre 1949.
\end{center}

\clearpage
\begin{center}
\vspace*{3.5em}
{\LARGE\bfseries \pdfbookmark[0]{Algèbre des grandeurs hypercomplexes}{noether.group.BOOK}\NoetherNavAnchor{BOOK}{Algèbre des grandeurs hypercomplexes — cours de la professeure E. Noether, semestre d’hiver 1929/30, rédigé par le professeur M. Deuring}Algèbre des grandeurs hypercomplexes\par}
\vspace{3.5em}
{\large Cours de la professeure E.\ Noether\par}
{\large Semestre d'hiver 1929/30\par}
\vspace{1.5em}
{\large Rédigé par le professeur M.\ Deuring\par}
\vspace{4em}
{\Large\bfseries Table des matières\par}
\end{center}
\vspace{1em}
\begingroup
\small
\setlength{\parindent}{0pt}
\newcommand{\frtocline}[2]{\noindent #1\dotfill #2\par}
\newcommand{\frtocsec}[3]{\noindent\hspace*{1.2em}\S\,#1.\quad #2\dotfill #3\par}
\frtocline{\textbf{Introduction}}{3}
\medskip
\frtocline{\textbf{Chapitre I : Représentations}}{3}
\frtocsec{1}{Définition de la représentation directe --- définition de la représentation réciproque}{3}
\frtocsec{2}{Classes de représentations}{4}
\frtocsec{3}{Modules de représentation}{4}
\frtocsec{4}{Le rapport entre modules de représentation et représentations}{5}
\medskip
\frtocline{\textbf{Chapitre II : Théorie de Galois des corps commutatifs}}{7}
\frtocsec{5}{Extension du corps des coefficients des systèmes hypercomplexes}{7}
\frtocsec{6}{Les représentations irréductibles des systèmes commutatifs}{8}
\frtocsec{7}{Les isomorphismes d'un corps}{8}
\frtocsec{8}{Le corps de décomposition}{9}
\frtocsec{9}{Les isomorphismes de $\mathsf{Z}_1$ sur les $\mathsf{Z}_i$}{9}
\frtocsec{10}{Le groupe de Galois}{9}
\frtocsec{11}{Le théorème fondamental de la théorie de Galois}{9}
\frtocsec{12}{La signification formelle des composantes $e_i$}{11}
\frtocsec{13}{Le théorème général de prolongement}{12}
\medskip
\frtocline{\textbf{Chapitre III : Groupes abéliens}}{13}
\frtocsec{14}{L'anneau de groupe}{13}
\frtocsec{15}{Les anneaux de groupe des groupes abéliens}{14}
\frtocsec{16}{Les relations entre caractères}{15}
\frtocsec{17}{La théorie de Galois des groupes abéliens}{15}
\medskip
\frtocline{\textbf{Chapitre IV : Anneaux simples bilatères}}{18}
\frtocsec{18}{Un lemme auxiliaire}{18}
\frtocsec{19}{Représentations des systèmes hypercomplexes simples bilatères dans des corps d'extension non commutatifs de leurs corps de coefficients}{19}
\frtocsec{20}{Corps non commutatifs}{22}
\frtocsec{21}{La théorie de Galois des corps non commutatifs}{26}
\clearpage
\frtocline{\textbf{Chapitre V : Systèmes de facteurs}}{29}
\frtocsec{22}{Le groupe des corps de centre donné}{29}
\frtocsec{23}{Systèmes de facteurs}{30}
\frtocsec{24}{Multiplication des systèmes de facteurs}{33}
\frtocsec{25}{Représentation normale de $\mathfrak K_r$ avec des sous-corps commutatifs maximaux galoisiens}{39}
\frtocsec{26}{Multiplication des représentations croisées}{43}
\medskip
\frtocline{\textbf{Chapitre VI : Théorie des produits croisés}}{44}
\frtocsec{27}{Représentation des $\mathfrak K_r$ comme produits croisés}{44}
\frtocsec{28}{Théorème du produit pour les systèmes de facteurs}{47}
\frtocsec{29}{Théorème du genre principal dans le cas minimal}{50}
\frtocsec{30}{Spécialisation aux corps de décomposition cycliques}{50}
\frtocsec{31}{Applications du cas particulier cyclique}{52}
\endgroup
\clearpage

\section*{\NoetherNavAnchor{INTRO}{Introduction}Introduction}\label{fr-r823-introduction}

Ce cours traite de la théorie générale des représentations des anneaux, issue, d'une part, de la théorie des représentations des groupes finis et, d'autre part, de la théorie algébrique et arithmétique des systèmes de nombres hypercomplexes. Par rapport à la conception antérieure de la théorie des représentations (Frobenius, Schur), cette conception plus générale présente les avantages suivants.

Puisque l'on considère non pas le groupe, mais l'anneau de groupe, et, plus généralement, un système hypercomplexe ou un anneau quelconque, on peut employer la théorie des anneaux et des idéaux, ce qui libère la théorie de nombreux calculs peu transparents. L'emploi de la théorie générale des anneaux et des idéaux ne se borne pas à simplifier la théorie : il la fait aussi progresser --- par exemple dans la question des rapports entre un groupe et ses sous-groupes --- et lui ouvre de nouveaux domaines d'application. À l'aide de la théorie générale des représentations, on peut donner un fondement nouveau et très élégant à la théorie de Galois et --- ce qui est peut-être plus important encore --- édifier une théorie de Galois des corps gauches.

Les notions fondamentales de la théorie des groupes, des anneaux et des idéaux, ainsi que les théorèmes généraux de la théorie des représentations, se trouvent dans le travail de Mlle Noether, \og Grandeurs hypercomplexes et théorie des représentations\fg{}, \emph{Mathematische Zeitschrift} 30, p.\ 641. Ce travail, cité ci-après sous l'abréviation M.Z., doit être considéré comme un complément au présent exposé. On trouve également la même matière que dans M.Z. dans la rédaction du cours de Mlle Noether sur les grandeurs hypercomplexes et les caractères de groupes, semestre d'hiver 1927/28.

Les notions de représentation et de module de représentation sont brièvement rappelées ici afin d'introduire le module de représentation réciproque et les représentations réciproques, dont la nouveauté est purement formelle.

\section*{Chapitre I. Représentations}\label{fr-r823-chapitre-i-representations}

Soit $\mathfrak{o}$ un anneau --- l'anneau que l'on représente --- et $\mathsf T$ un autre anneau --- celui dans lequel on le représente.

\subsection*{\NoetherNavAnchor{S01}{§ 1. Définition de la représentation directe — définition de la représentation réciproque}\S{} 1. Définition de la représentation directe}\label{fr-r823-definition-representation-directe}

Une \emph{représentation directe} de degré $n$ de $\mathfrak{o}$ dans $\mathsf T$ est un sous-anneau $\mathfrak D$ de l'anneau $\mathsf T_n$ des matrices de degré $n$ à éléments dans $\mathsf T$, sur lequel $\mathfrak{o}$ est envoyé par un homomorphisme d'anneaux. En symboles,
\[
\mathfrak{o}\to\mathfrak D\subseteqq\mathsf T_n.
\]

\subsection*{Définition de la représentation réciproque}\label{fr-r823-definition-representation-reciproque}

Une \emph{représentation réciproque} de degré $n$ de $\mathfrak{o}$ dans $\mathsf T^*$ est un sous-anneau $\mathfrak D^*$ de l'anneau $\mathsf T_n^*$ des matrices de degré $n$ à éléments dans $\mathsf T^*$, sur lequel $\mathfrak{o}$ est envoyé par un antihomomorphisme d'anneaux. Dire qu'il s'agit d'un antihomomorphisme signifie ceci : si les éléments $c,d$ de $\mathfrak{o}$ correspondent aux éléments $c^*,d^*$ de $\mathfrak D^*$, alors l'élément $cd$ de $\mathfrak{o}$ doit correspondre à l'élément $d^*\cdot c^*$ de $\mathfrak D^*$, et l'élément $c+d$ de $\mathfrak{o}$ à l'élément $c^*+d^*$ de $\mathfrak D^*$. En symboles,
\[
\mathfrak{o}\nrightarrow\mathfrak D^*\subseteqq T_n^*.
\]

Voici un exemple de la notion d'antihomomorphisme : l'application qui associe à chaque matrice carrée d'ordre $n$ sur un corps commutatif sa transposée est un anti-isomorphisme d'anneaux.

Plus généralement, pour tout anneau $\mathsf T$, on peut construire un anneau $\mathsf T^*$ qui lui est anti-isomorphe de la manière suivante. À chaque élément $c$ de $\mathsf T$, on associe un symbole $c^*$; on définit l'addition de ces symboles par $c^*+d^*=(c+d)^*$ et leur multiplication par
\[
c^*\cdot d^*=(dc)^*.
\]
L'ensemble des symboles $c^*$ est un anneau $\mathsf T^*$ anti-isomorphe à $\mathsf T$.

\subsection*{\NoetherNavAnchor{S02}{§ 2. Classes de représentations}\S{} 2. Classes de représentations}\label{fr-r823-classes-representations}

On peut transformer une représentation directe $\mathfrak D$ de $\mathfrak{o}$ dans $\mathsf T$ en une représentation isomorphe $\mathfrak D'$ en la transformant élément par élément au moyen d'une matrice régulière $P$ : à un élément $c$ de $\mathfrak{o}$, on associe, au lieu de $C\subseteqq\mathfrak D$, la matrice $P^{-1}CP\subseteqq P^{-1}\mathfrak D P$.

Deux représentations que l'on peut ainsi transformer l'une en l'autre sont dites \emph{équivalentes}.

Toutes les représentations équivalentes à une représentation donnée forment une \emph{classe de représentations directes} (c.\ r.\ d.).

On définit de même l'équivalence des représentations réciproques : deux représentations réciproques sont équivalentes si l'une s'obtient à partir de l'autre par transformation au moyen d'une matrice régulière. Toutes les représentations équivalentes à une représentation réciproque fixée forment une \emph{classe de représentations réciproques}.

\subsection*{\NoetherNavAnchor{S03}{§ 3. Modules de représentation}\S{} 3. Modules de représentation}\label{fr-r823-modules-representation}

Définition du \emph{module de représentation direct} (m.\ r.\ d.).

Un module de représentation direct de $\mathfrak{o}$ relativement à $\mathsf T$ est tout bimodule $\mathfrak{o}$--$\mathsf T$, $\mathfrak M$ (c'est-à-dire un groupe abélien, écrit additivement, qui est un module à gauche sur $\mathfrak{o}$, un module à droite sur $\mathsf T$, et satisfait à la loi d'associativité mixte : pour $c\in\mathfrak{o}$, $m\in\mathfrak M$, $\tau\in\mathsf T$, on a $(cm)\tau=c(m\tau)$), qui peut s'écrire comme somme directe d'un nombre fini de $\mathsf T$-modules monogènes --- avec la condition que $x_i\tau=0$ entraîne $\tau=0$ :
\[
\mathfrak M=x_1\cdot\mathsf T+x_2\cdot\mathsf T+\cdots+x_n\cdot\mathsf T.
\]

Définition du \emph{module de représentation réciproque} (m.\ r.\ r.).

Un module de représentation réciproque de $\mathfrak{o}$ relativement à $\mathsf T^*$ est tout groupe $\mathfrak M^*$ qui est à la fois un $\mathfrak{o}$-module à gauche et un $\mathsf T^*$-module à gauche, satisfait à la loi de commutation (si $c\in\mathfrak{o}$, $m\in\mathfrak M^*$ et $\tau\in\mathsf T^*$, alors $c(\tau^*m^*)=\tau^*(cm)$), et peut s'écrire comme somme directe d'un nombre fini de $\mathsf T^*$-modules simples --- avec la condition que $\tau^*\cdot x_i^*=0$ entraîne $\tau^*=0$ :
\[
\mathfrak M^*=\mathsf T^*\cdot x_1+\cdots+\mathsf T^*\cdot x_n.
\]

\subsection*{\NoetherNavAnchor{S04}{§ 4. Le rapport entre modules de représentation et représentations}\S{} 4. Le rapport entre modules de représentation et représentations}\label{fr-r823-rapport-modules-representations}

\textbf{1.} \emph{Chaque module de représentation détermine de façon unique une classe de représentations (un m.\ r.\ d. détermine une c.\ r.\ d.; un m.\ r.\ r., une c.\ r.\ r.).}

Voici la construction.

La multiplication du m.\ r.\ d. $\mathfrak M$ par un élément $c$ de $\mathfrak{o}$ définit un endomorphisme de $\mathfrak M$ qui, en vertu de la loi $(cm)\tau=c(m\tau)$, admet les éléments de $\mathsf T$ comme opérateurs. C'est donc une \emph{transformation linéaire} de $\mathfrak M$ en lui-même : la correspondance $m\mapsto cm=m'$ possède la propriété
\[
\text{si }y_i\mapsto y_i'(=cy_i),\text{ alors aussi }
\sum y_i\tau_i\mapsto\sum y_i'\tau_i\left(=c\sum y_i\tau_i\right).
\]

Définissons la somme $\mathfrak T_1+\mathfrak T_2$ de deux transformations linéaires
\[
\mathfrak T_1:m\mapsto m'(=c_1m),\qquad
\mathfrak T_2:m\mapsto m''(=c_2m)
\]
comme la transformation $m\mapsto m'+m''(=(c_1+c_2)m)$, et leur produit $\mathfrak T_1\cdot\mathfrak T_2$ comme la transformation composée $m\mapsto(m'')'=c_1c_2m$. L'ensemble des différentes transformations linéaires de $\mathfrak M$ engendrées par les éléments de $\mathfrak{o}$ est alors une image homomorphe d'anneau $\overline{\mathfrak D}$ de $\mathfrak{o}$.

Choisissons maintenant une base quelconque $x_1,\ldots,x_n$ de $\mathfrak M$. La transformation linéaire de $\mathfrak M$ engendrée par $c$ est entièrement déterminée dès que l'on sait comment elle transforme les éléments de base $x_i$, c'est-à-dire dès que l'on connaît la matrice $(\gamma_{ij})=C$ du système
\[
x'_j=cx_j=\sum x_i\gamma_{ij},
\]
car, pour un autre élément $x=\sum x_j\tau_j$, on a simplement $x'=\sum x'_j\tau_j$. Les matrices $C$ correspondent bijectivement aux différentes transformations linéaires (chacune détermine l'autre) et forment, comme on le voit immédiatement, un anneau isomorphe à $\overline{\mathfrak D}$; elles constituent donc une image homomorphe de $\mathfrak{o}$, c'est-à-dire une représentation de degré $n$ de $\mathfrak{o}$.

L'isomorphisme entre $\overline{\mathfrak D}$ et la représentation s'exprime par les formules suivantes. Si la transformation linéaire engendrée par $c\in\mathfrak{o}$ correspond à la matrice $C$, de sorte que
\[
c(x_1,\ldots,x_n)=(x_1,\ldots,x_n)C,
\]
et si, de même, $d\in\mathfrak{o}$ correspond à $D\in\mathsf T_n$, de sorte que
\[
d(x_1,\ldots,x_n)=(x_1,\ldots,x_n)D,
\]
alors, manifestement,
\[
\begin{aligned}
(c+d)(x_1,\ldots,x_n)
&=(cx_1+dx_1,\ldots,cx_n+dx_n)\\
&=(x_1,\ldots,x_n)(C+D)
\end{aligned}
\]
et
\[
\begin{aligned}
cd(x_1,\ldots,x_n)
&=c(dx_1,\ldots,dx_n)=c((x_1,\ldots,x_n)D)\\
&=(cx_1,\ldots,cx_n)D=(x_1,\ldots,x_n)CD.
\end{aligned}
\]

Ainsi, lorsqu'on les réalise au moyen d'une base déterminée $x_\nu$ de $\mathfrak M$, les transformations linéaires de $\mathfrak M$ en lui-même engendrées par $\mathfrak{o}$ donnent une représentation de degré $n$ de $\mathfrak{o}$ dans $\mathsf T$.

Si $(y_1,\ldots,y_n)$ est une autre base de $\mathfrak M$, elle se déduit de $x_1,\ldots,x_n$ par multiplication par une matrice régulière :
\[
(y_1,\ldots,y_n)=(x_1,\ldots,x_n)P.
\]
Si une transformation linéaire de $\mathfrak M$ est réalisée, dans la base $x_\nu$, par la matrice $C$, elle est réalisée, dans la base $y_\nu$, par $P^{-1}CP$. En effet, de
\[
c(x_1,\ldots,x_n)=(x_1,\ldots,x_n)C
\]
on déduit
\[
\begin{aligned}
c(y_1,\ldots,y_n)
&=c(x_1,\ldots,x_n)P=((x,\ldots,x_n)C)P\\
&=(x_1,\ldots,x_n)CP=((y_p,\ldots,y_n)P^{-1})CP\\
&=(y_1,\ldots,y_n)P^{-1}CP.
\end{aligned}
\]
Il en résulte immédiatement que $\mathfrak M$ fournit toutes les représentations d'une classe de représentations lorsque l'on emploie successivement toutes les bases de $\mathfrak M$ pour réaliser les transformations linéaires de $\mathfrak M$ induites par $\mathfrak{o}$.

\textbf{2.} \emph{Toute classe de représentations est engendrée par un module de représentation de la manière indiquée en 1.} Nous nous limitons de nouveau aux représentations directes.

\emph{Démonstration.} Choisissons une représentation déterminée de la classe, donnée par la correspondance homomorphe $c\mapsto C$. Si $n$ est son degré, formons, avec $n$ indéterminées $x_1,\ldots,x_n$, un $\mathsf T$-module à droite
\[
\mathfrak M=x_1\mathsf T+\dots+x_n\mathsf T.
\]
On lui applique les règles de calcul usuelles
\[
\sum x_i\tau_i+\sum x_i\overline{\tau}_i=\sum x_i(\tau_i+\overline{\tau}_i),
\]
$\sum x\tau_i=\sum x_i\overline{\tau}_i$ si et seulement si $\tau_i=\overline{\tau}_i$ pour tout $i$, et
\[
\left(\sum x_i\tau_i\right)\varrho=\sum x_i\tau_i\varrho.
\]
On en fait un $\mathfrak{o}$-module à gauche par la convention
\[
c\sum x_i\tau_i=\sum x_i'\tau_i,
\]
lorsque $(x_1',\ldots,x_n')=(x_1,\ldots,x_n)C$, où $C$ désigne la matrice de la représentation associée à l'élément $c$ de $\mathfrak{o}$.

Cette convention fait de $\mathfrak M$ un bimodule $\mathfrak{o}$--$\mathsf T$. En effet, les relations d'homomorphisme
\[
c+d\mapsto C+D,\qquad cd\mapsto CD
\]
entraînent
\[
(c+d)x_i=cx_i+dx_i,
\]
ainsi que
\[
cd\cdot x_i=c\cdot dx_i,
\qquad
c\left(\sum x_i\tau_i+\sum x_i\overline{\tau}_i\right)
=c\sum x_i\tau_i+c\sum x_i\overline{\tau}_i;
\]
ce sont les propriétés d'un module à gauche. D'autre part, on a manifestement
\[
c\left(\sum x_i\tau_i\varrho\right)
=\left(c\sum x_i\tau_i\right)\cdot\varrho,
\]
ce qui est la loi d'associativité mixte.

Ainsi $\mathfrak M$ est un module de représentation, et il engendre la représentation choisie de la manière indiquée en 1 lorsque l'on emploie la base $x_1,\ldots,x_n$ pour réaliser ses transformations linéaires. Les autres bases donnent les autres représentations de la classe.

Tout se passe de la même manière pour les représentations et modules de représentation réciproques. On obtient que le module
\[
\mathfrak M^*=\mathsf T^*x_1^*+\cdots+\mathsf T^*x_n^*
\]
engendre une représentation réciproque de $\mathfrak D$ dans $\mathsf T^*$ en associant à l'élément $c$ de $\mathfrak{o}$ la matrice $C$ de l'équation
\[
c\begin{pmatrix}x_1\\ \vdots\\ x_n\end{pmatrix}
=\begin{pmatrix}cx_1\\ \vdots\\ cx_n\end{pmatrix}
=C\begin{pmatrix}x_1\\ \vdots\\ x_n\end{pmatrix}.
\]
Les différentes bases donnent de nouveau les différentes représentations d'une même classe.

\section*{Chapitre II. Théorie de Galois des corps commutatifs}\label{kapitel-ii.-galoissche-theorie-kommutativer-koerper}

La théorie des représentations développée jusqu'ici suffit à fonder la théorie de Galois; elle présente sur la théorie usuelle l'avantage de ne pas devoir recourir à des systèmes particuliers d'éléments générateurs ni à des théorèmes sur les polynômes en des indéterminées.\srcfn{1}{Pour ce qui suit, comparer M.Z. \S{} 15--20. M.Z. \S{} 21 contient les premiers théorèmes de la section suivante sur la théorie de Galois. Pour la définition des systèmes hypercomplexes, voir M.Z. \S{} 6.}

\subsection*{\NoetherNavAnchor{S05}{§ 5. Extension du corps des coefficients des systèmes hypercomplexes}\S{} 5. Extension du corps des coefficients des systèmes hypercomplexes}\label{koeffizientenerweiterung-hyperkomplexer-systeme}

Nous énonçons le théorème général suivant, dont nous laisserons la démonstration de côté, puisqu'elle est presque immédiate. Soit \(\mathfrak{o}=c_1\mathsf P+\cdots+c_n\mathsf P\) un système hypercomplexe sur le corps commutatif \(\mathsf P\). Soit \(\Omega\) un anneau d'extension de \(\mathsf P\), dont le centre contient \(\mathsf P\). Nous définissons le \emph{système étendu} \(\mathfrak{o}_{\Omega}\) comme l'ensemble de toutes les sommes formelles \(\sum c_i\omega_i\), \(\omega_i\in\Omega\); les conventions

\[
\begin{array}{l}
\sum c_i\omega_i=\sum c_i\omega_i' \quad\text{si et seulement si } \omega_i=\omega_i' \text{ pour tout } i;\\
\sum c_i\omega_i+\sum c_i\omega_i'=\sum c_i(\omega_i+\omega_i');\\
\omega\sum c_i\omega_i=\sum c_i\omega\omega_i;\quad
\sum c_i\omega_i\cdot\omega=\sum c_i\omega_i\omega;\\
\sum c_i\omega_i=\sum \omega_i c_i
\end{array}
\]

en font d'abord un \(\Omega\)-module où les actions gauche et droite de \(\Omega\) coïncident, puis la convention supplémentaire

\[
\sum c_i\omega_i\cdot \sum c_i\omega_i
=\sum_i\sum_j c_ic_j\omega_i\omega_j
\]

--- où \(c_ic_j\) est déjà défini comme élément de \(\mathfrak{o}\) --- en fait un anneau d'extension de \(\mathfrak{o}\). \(\mathfrak{o}_{\Omega}\) est de rang \(n\) sur \(\Omega\) et peut, sous la forme

\[
\mathfrak{o}_{\Omega}=c_1\Omega+\cdots+c_n\Omega
\]

être écrit.

La définition de \(\mathfrak{o}_{\Omega}\) est indépendante de la base \(c_1,\ldots,c_n\). Si \(\Omega\) est un corps commutatif, \(\mathfrak{o}_{\Omega}\) est un système hypercomplexe sur \(\Omega\).

\subsection*{\NoetherNavAnchor{S06}{§ 6. Les représentations irréductibles des systèmes commutatifs}\S{} 6. Les représentations irréductibles des systèmes commutatifs}\label{die-irreduziblen-darstellungen-kommutativer-systeme}

Soit \(\mathfrak Z=z_1\mathsf P+\cdots+z_n\mathsf P\) un système commutatif sur \(\mathsf P\). Nous voulons déterminer toutes les représentations irréductibles de \(\mathfrak Z\) dans un corps d'extension algébriquement clos \(\Omega\) de \(\mathsf P\). Toute représentation de \(\mathfrak Z\) dans \(\Omega\) donne une représentation de \(\mathfrak Z_\Omega=z_1\Omega+\cdots+z_n\Omega\) --- car les représentations des éléments de base déjà contenus dans \(\mathsf P\) suffisent à déterminer la représentation entière --- et, réciproquement, toute représentation de \(\mathfrak Z_\Omega\) induit par restriction une représentation de \(\mathfrak Z\); nous pouvons donc chercher aussi les représentations de \(\mathfrak Z_\Omega\). D'après M.Z. \S{} 20, celles-ci sont contenues dans la représentation régulière de \(\mathfrak Z_\Omega\), c'est-à-dire dans la représentation de \(\mathfrak Z_\Omega\) obtenue en employant \(\mathfrak Z_\Omega\) lui-même comme module de représentation. Il suffit même de se limiter à \(\mathfrak Z_\Omega/\mathfrak C\) comme module de représentation, où \(\mathfrak C\) désigne le radical de \(\mathfrak Z_\Omega\) (voir M.Z. \S{} 19).

Or \(\mathfrak Z_\Omega/\mathfrak C\) est un système complètement réductible, donc une somme directe de corps de degré fini sur \(\Omega\). Puisque \(\Omega\) est algébriquement clos, ces composantes sont isomorphes à \(\Omega\).

\[
\mathfrak Z_\Omega/\mathfrak C=e_1\Omega+\cdots+e_t\Omega.
\]

Les \(e_\nu\) sont les composantes de l'unité. On a donc :

\emph{\(\mathfrak Z_\Omega\) possède exactement \(t\) classes distinctes de représentations irréductibles dans \(\Omega\), toutes de degré un. Ici \(t\) est le rang de \(\mathfrak Z_\Omega/\mathfrak C\), donc au plus égal à \(n\).}

Chacune de ces classes ne contient qu'une seule représentation, car transformer une représentation de degré un dans un corps commutatif ne change pas la représentation : \(\pi^{-1}\cdot\alpha\cdot\pi=\alpha\).

Il existe par conséquent exactement \(t\) homomorphismes d'anneaux distincts \(\Theta_1,\ldots,\Theta_t\) de \(\mathfrak Z_\Omega\) sur des sous-anneaux de \(\Omega\).

\subsection*{\NoetherNavAnchor{S07}{§ 7. Les isomorphismes d’un corps}\S{} 7. Les isomorphismes d'un corps}\label{die-isomorphismen-eines-koerpers}

Supposons maintenant que \(\mathfrak Z\) soit un corps. Tout homomorphisme \(\Theta_i\) de \(\mathfrak Z_\Omega\) sur un sous-anneau de \(\Omega\) doit envoyer \(\mathfrak Z\) isomorphiquement sur un sous-corps \(\mathsf{Z}_i\) de \(\Omega\). En effet, ou bien \(\mathsf{Z}_i=0\) --- ce qui n'est pas le cas ici, puisque l'unité de \(\mathfrak Z\) correspond à l'unité de \(\Omega\) --- ou bien \(\mathsf{Z}_i\) est une image isomorphe de \(\mathfrak Z\). Par conséquent :

\emph{Si \(\mathfrak Z\) est une extension finie de degré \(n\) du corps \(\mathsf P\), le nombre des corps distincts \(\mathsf{Z}_i\) contenus dans \(\Omega\) et isomorphes à \(\mathfrak Z\) est égal au rang \(t\) de \(\mathfrak Z_\Omega/\mathfrak C\). Le nombre des \(\mathsf{Z}_i\) est donc au plus égal à \(n\).}

On dit que \(\mathfrak Z\) est de première espèce sur \(\mathsf P\) lorsque le nombre des \(\mathsf{Z}_i\) est égal à \(n\), et de seconde espèce lorsque ce nombre est inférieur à \(n\).

\(\mathfrak Z\) est de première espèce sur \(\mathsf P\) si et seulement si \(\mathfrak Z_\Omega\) est complètement réductible.

On voit ici l'avantage de cette façon de fonder la théorie de Galois sur la fondation usuelle. Dans la théorie de Galois usuelle, on considère les isomorphismes de l'un des \(\mathsf{Z}_i\) sur les autres. Cette dissymétrie est ici évitée. Pour montrer qu'un corps n'a pas plus de corps conjugués que ne l'indique son degré, on emploie ordinairement un élément primitif. Ici, ce fait résulte simplement de ce que toutes les représentations irréductibles de \(\mathfrak Z_\Omega\) sont contenues dans la représentation régulière.

\subsection*{\NoetherNavAnchor{S08}{§ 8. Le corps de décomposition}\S{} 8. Le corps de décomposition}\label{der-zerfaellungskoerper}

Les \(\mathsf{Z}_i\) sont les représentations irréductibles de \(\mathfrak Z\) dans \(\Omega\). Si donc un corps \(\Gamma\) compris entre \(\mathsf P\) et \(\Omega\), \(\mathsf P\subseteqq\Gamma\subseteqq\Omega\), contient tous les \(\mathsf{Z}_i\), alors \(\mathfrak Z_\Gamma/\mathfrak C\) doit déjà se décomposer en \(t\) corps de degré un sur \(\Gamma\). Réciproquement, pour que \(\mathfrak Z_\Gamma/\mathfrak C\) se décompose en \(t\) corps, qui sont alors les représentations irréductibles de degré un de \(\mathfrak Z_\Gamma\), les \(\mathsf{Z}_i\) doivent être contenus dans \(\Gamma\). Si l'on appelle corps de décomposition de \(\mathfrak Z\) un corps \(\Gamma\) tel que \(\mathfrak Z_\Gamma/\mathfrak C\) devienne une somme directe de \(t\) corps, nous avons donc :

\emph{Le corps galoisien des \(\mathsf{Z}_i\) isomorphes à \(\mathfrak Z\), c'est-à-dire leur composé, est le corps de décomposition minimal de \(\mathfrak Z\).}

\subsection*{\NoetherNavAnchor{S09}{§ 9. Les isomorphismes de Z1 sur les Zi}\S{} 9. Les isomorphismes de \(\mathsf{Z}_1\) sur les \(\mathsf{Z}_i\)}\label{die-isomorphismen-von-z1-auf-die-zi}

Restreint à \(\mathfrak Z\), \(\Theta_i\) est un isomorphisme de \(\mathfrak Z\) sur \(\mathsf{Z}_i\). Ainsi \(S_i=\Theta_i\Theta_1^{-1}\) est un isomorphisme de \(\mathsf{Z}_1\) sur \(\mathsf{Z}_i\). Puisque les \(\Theta_i\) sont distincts, les \(S_i\) le sont aussi. Les \(S_i\) sont les isomorphismes d'un corps \(\mathsf{Z}_1\) sur ses conjugués que l'on considère ordinairement dans la théorie de Galois. Il n'en existe pas d'autres, puisqu'il n'existe pas d'autres \(\Theta_i\).

\subsection*{\NoetherNavAnchor{S10}{§ 10. Le groupe de Galois}\S{} 10. Le groupe de Galois}\label{die-galoissche-gruppe}

Si \(\mathfrak Z\) est un corps galoisien, c'est-à-dire si les \(\mathsf{Z}_i\) coïncident comme ensembles d'éléments, les \(S_i\) forment le groupe de tous les isomorphismes de \(Z\) sur lui-même qui fixent \(\mathsf P\). En effet, tout produit \(S_iS_j\) est un automorphisme de \(\mathsf{Z}\), et puisqu'il n'en existe pas d'autres que les \(S_i\), \(S_iS_j\) est égal à un \(S_k\). Les \(S_i\) forment donc le groupe de tous les automorphismes de \(Z\) qui fixent élément par élément \(\mathsf P\); c'est le groupe de Galois de \(\mathsf{Z}\) relativement à \(\mathsf P\).

\subsection*{\NoetherNavAnchor{S11}{§ 11. Le théorème fondamental de la théorie de Galois}\S{} 11. Le théorème fondamental de la théorie de Galois}\label{der-hauptsatz-der-galoisschen-theorie}

Supposons maintenant que \(\mathfrak Z\) soit une extension galoisienne de première espèce de \(\mathsf P\). Nous démontrons le théorème fondamental de la théorie de Galois :

\emph{Les corps \(\Gamma\) compris entre \(\mathsf P\) et \(\mathfrak Z\) peuvent être mis en correspondance bijective avec les sous-groupes \(\mathfrak H\) du groupe de Galois \(\mathfrak G\) de \(\mathsf{Z}\) relativement à \(\mathsf P\), de telle sorte que le groupe \(\mathfrak H\) associé au corps \(\Gamma\) soit constitué de tous les automorphismes de \(\mathsf{Z}\) qui fixent \(\Gamma\) élément par élément, et que \(\Gamma\) soit l'ensemble de tous les éléments de \(\mathsf{Z}\) qui restent fixes sous tous les automorphismes contenus dans \(\mathfrak H\).}

Nous l'énonçons brièvement comme suit :

\begin{enumerate}
\item Si \(\mathfrak H\) est le groupe d'invariance de \(\Gamma\), alors \(\Gamma\) est le corps des invariants de \(\mathfrak H\).
\item Si \(\Gamma\) est le corps des invariants de \(\mathfrak H\), alors \(\mathfrak H\) est le groupe d'invariance de \(\Gamma\).
\end{enumerate}

\paragraph{1. se ramène au lemme auxiliaire :}\label{laesst-sich-zurueckfuehren-auf-den-hilfssatz}

\emph{Si \(\mathsf P \subseteqq \Sigma \subseteqq \mathsf T \subseteqq \mathsf Z\) et si \(\mathsf T\) est de degré \(t\) sur \(\Sigma\), tout isomorphisme de \(\Sigma\) sur un corps conjugué (nécessairement contenu dans \(\mathsf Z\)) se prolonge en \(t\) isomorphismes distincts de \(\mathsf T\) sur des corps conjugués.}

Supposons en effet que \(\mathfrak H\) soit le groupe d'invariance de \(\Sigma\) et que \(\mathsf T\) soit le corps des invariants de \(\mathfrak H\). Nous pouvons alors raisonner comme suit. Les éléments de \(\mathfrak H\) sont les prolongements à tout \(\mathsf Z\) de l'isomorphisme identique de \(\Sigma\). Si l'on applique ces prolongements seulement à \(\mathsf T\), le lemme auxiliaire montre qu'ils se répartissent en \(t\) classes, de telle sorte que les isomorphismes d'une même classe coïncident sur \(\mathsf T\), tandis que les isomorphismes de classes différentes sont distincts sur \(\mathsf T\). Mais puisque \(\mathsf T\) est le corps des invariants de \(\mathfrak H\), on doit avoir \(t=1\), \(\mathsf T=\Sigma\), c.q.f.d.

Pour démontrer le lemme auxiliaire, revenons aux corps \(\mathfrak S\), \(\mathfrak T\) compris entre \(\mathsf P\) et \(\mathsf Z\), auxquels \(\Sigma\) et \(\mathsf T\) sont associés par \(\Theta_1\). Nous devons alors démontrer :

\emph{Tout isomorphisme \(\mathsf{H}_i\) de \(\mathfrak S\) sur un \(\Sigma_i\) se prolonge, d'exactement \(t\) manières distinctes, en des isomorphismes de \(\mathfrak T\) sur des sous-corps \(\mathsf T_i\) de \(\mathsf Z\).}

Puisque, par hypothèse, \(\mathfrak Z\) est de première espèce sur \(\mathsf P\), \(\mathfrak Z_\Omega/\mathfrak C\) s'identifie à \(\mathfrak Z_\Omega\); ainsi \(\mathfrak S_\Omega\) se décompose en autant de corps que l'indique son degré \(s\) :
\[
\mathfrak S_\Omega = E_1\Omega+\cdots+E_s\Omega .
\]
Les \(E_\nu\Omega\), considérés comme modules de représentation, engendrent les \(s\) isomorphismes de \(\mathfrak S\) sur \(s\) sous-corps \(\Sigma_1,\ldots,\Sigma_s\). Pour trouver les isomorphismes de \(\mathfrak T\), nous avons besoin de \(\mathfrak T_\Omega\). Puisque les \(E_\nu\Omega\) sont des idéaux de \(\mathfrak S_\Omega\), on a \(E_\nu\Omega=\mathfrak S_\Omega E_\nu\). En outre, il est manifeste que
\[
\mathfrak T_\Omega=\mathfrak T\cdot\Omega=\mathfrak T\cdot\mathfrak S_\Omega .
\]
Par conséquent,
\[
\mathfrak T_\Omega
=E_1\mathfrak S_\Omega\cdot\mathfrak T+\cdots+E_s\mathfrak S_\Omega\cdot\mathfrak T
=E_1\mathfrak T_\Omega+\cdots+E_s\mathfrak T_\Omega .
\]
C'est une décomposition de \(\mathfrak T_\Omega\) en idéaux. \(E_\nu\cdot\mathfrak T_\Omega\) se décompose en \(t\) idéaux simples. En effet, on a les relations de degrés
\[
[\mathfrak T_\Omega\cdot E_\nu:\mathfrak S_\Omega E_\nu]
\leqq [\mathfrak T_\Omega:\mathfrak S_\Omega]=[\mathfrak T:\mathfrak S]=t,
\]
donc, puisque \(\mathfrak S_\Omega E_\nu=\Omega E_\nu\cong\Omega\),
\[
[\mathfrak T_\Omega\cdot E_\nu:\Omega]
=[\mathfrak T_\Omega E_\nu:\mathfrak S_\Omega E_\nu]\leqq t,
\]
et, puisque la somme de tous les degrés \([\mathfrak T_\Omega E_\nu:\Omega]\) doit être égale à \(s\cdot t\),
\[
[\mathfrak T_\Omega\cdot E_\nu:\Omega]=t,\qquad
E_\nu\mathfrak T_\Omega=e_\nu^{(1)}\cdot\Omega+\cdots+e_\nu^{(t)}\Omega .
\]
Chaque \(e_\nu^{(i)}\Omega\) donne une représentation irréductible de \(\mathfrak T_\Omega\). Il nous reste à montrer que toutes ces représentations, lorsqu'on les applique seulement à \(\mathfrak S_\Omega\), coïncident avec la représentation engendrée par \(E_\nu\cdot\Omega\). Cela suffit pour conclure. Or c'est clair, car
\[
s\cdot E_\nu=E_\nu\sigma_\nu\quad
(s\in\mathfrak S_\Omega\text{ est représenté par }\sigma_\nu)
\]
entraîne
\[
\sum s\cdot e_\nu^{(i)}=\sum e_\nu^{(i)}\sigma_\nu,
\]
et donc, en raison de l'unicité de l'écriture comme somme,
\[
s\cdot e_\nu^{(i)}=e_\nu^{(i)}\sigma_\nu,
\]
les nouvelles représentations associent donc elles aussi l'élément \(\sigma_\nu\) à l'élément \(s\).

\paragraph{2. \(\mathsf T\) est le corps des invariants de \(\mathfrak H\).}
Nous voulons montrer que \(\mathfrak H\) est le groupe d'invariance de \(\mathsf T\); il suffit pour cela d'exhiber un élément de \(\mathsf T\) qui n'est pas fixe sous un automorphisme n'appartenant pas à \(\mathfrak H\).

Les éléments \(S_i\) du groupe de Galois \(\mathfrak G\) de \(Z\) sont les automorphismes de \(\mathsf Z\) qui fixent \(\mathsf P\). Nous les faisons agir aussi sur \(\mathfrak Z_{\mathsf Z}\) par la convention
\[
S_i z=z,\quad \text{si } z \text{ appartient à } \mathfrak Z.
\]
L'application de \(S_i\) à
\[
\mathfrak Z_{\mathsf Z}=e_1\mathsf Z+\cdots+e_n\mathsf Z
\]
doit permuter les composantes directes \(e_\nu\mathsf Z\), puisqu'elles sont déterminées de façon unique; on doit en particulier avoir
\[
S_i e_\nu=e_{\nu'},\qquad\text{et}\qquad S_i e_\nu\ne S_k e_\nu
\quad\text{pour } i\ne k
\]
Si \(S_1,\ldots,S_h\) sont les éléments de \(\mathfrak H\), choisissons un quelconque \(e_\nu\), disons \(e'\), et considérons tous les \(e_\nu\) qui s'en déduisent au moyen des \(S_i\), \(i\leqq h\) :
\[
e_1=S_1e',\ldots,e_h=S_he'.
\]
Puisque, pour \(i\leqq h, j\leqq h\),
\[
S_i e_j=S_iS_j e'=S_ke'=e_k \quad \text{avec}\quad k\leqq h
\]
chaque \(S_i\) de \(\mathfrak H\) permute entre eux les \(e_1,\ldots,e_h\), et par conséquent
\[
E_1=e_1+\cdots+e_h
\]
est un élément de \(\mathfrak Z\), invariant sous tous les \(S_i\in\mathfrak H\). Si l'on exprime \(E_1\) au moyen d'une base de \(\mathfrak Z\), ses coefficients doivent être des éléments de \(\mathsf Z\) invariants sous \(\mathfrak H\); ils appartiennent donc à \(\mathsf T\).

\(E_1\) n'est invariant sous aucun \(S_i\) n'appartenant pas à \(\mathfrak H\). En effet, pour un tel \(S_i\),
\[
S_i e_1=e_i+\cdots
\]
contient la composante \(e_i\), qui ne figure pas parmi \(e_1,\ldots,e_h\), tandis que la représentation de \(E_1\) comme somme est unique.

\subsection*{\NoetherNavAnchor{S12}{§ 12. La signification formelle des composantes e\_i}\S{} 12. La signification formelle des composantes \(e_i\)}\label{die-formale-bedeutung-der-komponenten-ei}

Fixons une base \(z_1,\ldots,z_n\) de \(\mathfrak Z\).
\[
\mathfrak Z=z_1\mathsf P+\cdots+z_n\mathsf P
\]
(Nous supposons maintenant seulement que \(\mathfrak Z\) est de première espèce.) L'isomorphisme \(\Theta_i\) envoie le système \(z_1,\ldots,z_n\) sur une base \(\zeta_1^{(i)},\ldots,\zeta_n^{(i)}\) de \(\mathsf Z_i\), selon les relations
\[
z_\nu e_i=e_i\zeta_\nu^{(i)} .
\]

Réciproquement, la composante \(e_j\) de l'unité doit pouvoir s'exprimer au moyen des \(z_\nu\), avec des coefficients dans \(\mathsf Z\) :
\[
e_j=\sum_\nu \varrho_\nu^{(j)}z_\nu .
\]
On a donc
\[
e_j=e_j^2=\sum \varrho_\nu^{(j)}z_\nu e_j
=\sum \varrho_\nu^{(j)}\cdot\zeta_\nu^{(j)}\cdot e_j
\]
et
\[
0=e_je_k=\sum \varrho_\nu^{(j)}\cdot z_\nu\cdot e_k
=\sum \varrho_\nu^{(j)}\cdot\zeta_\nu^{(k)}\cdot e_k,
\]
d'où :
\[
\sum_\nu \varrho_\nu^{(j)}\cdot\zeta_\nu^{(j)}=1;
\qquad
\sum_\nu \varrho_\nu^{(j)}\cdot\zeta_\nu^{(k)}=0,\quad
\text{si } j\ne k.
\]
Cela signifie que les matrices
\[
\begin{pmatrix}
\zeta_1^{(1)} & \cdots & \zeta_n^{(1)}\\
\cdots & \cdots & \cdots\\
\cdots & \cdots & \cdots\\
\zeta_1^{(n)} & \cdots & \zeta_n^{(n)}
\end{pmatrix}
\quad\text{et}\quad
\begin{pmatrix}
\varrho_1^{(1)} & \cdots & \varrho_1^{(n)}\\
\cdots & \cdots & \cdots\\
\cdots & \cdots & \cdots\\
\varrho_n^{(1)} & \cdots & \varrho_n^{(n)}
\end{pmatrix}
\]
sont inverses l'une de l'autre. On peut aussi l'exprimer ainsi : \(\varrho_1^{(i)},\ldots,\varrho_n^{(i)}\) est la base complémentaire de \(\zeta_1^{(i)},\ldots,\zeta_n^{(i)}\). Elle intervient dans la théorie de la différente d'un corps de nombres.

\subsection*{\NoetherNavAnchor{S13}{§ 13. Le théorème général de prolongement}\S{} 13. Le théorème général de prolongement}\label{der-allgemeine-fortsetzungssatz}

Le théorème sur le prolongement des \(\mathsf{H}_i\), démontré à la page 10, peut être généralisé.

Soit \(\mathfrak Z=z_1\mathsf P+\cdots+z_n\mathsf P\) un système commutatif, hypercomplexe et complètement réductible. Soient \(\mathfrak S\) et \(\mathfrak T\) des sous-systèmes de \(\mathfrak Z\) --- nécessairement eux aussi complètement réductibles dans cette situation --- tels que
\[
\mathsf P\subseteqq\mathfrak S\subseteqq\mathfrak T\subseteqq\mathfrak Z .
\]
D'après M.Z. \S{} 13, chacun des trois systèmes possède une unité. Ces trois unités ne sont toutefois nullement nécessairement égales. (Elles ne le sont certainement pas lorsque \(\mathfrak S\) et \(\mathfrak T\) sont des idéaux de \(\mathfrak Z\) distincts de \(\mathfrak Z\).) Soit \(E\) l'unité de \(\mathfrak S\). Le système \(\mathfrak T\cdot E\) est un \(\mathfrak S\)-module de rang fini :
\[
E\mathfrak T=a_1\mathfrak S+\cdots+a_t\mathfrak S
\]
(ce n'est généralement pas le cas de \(\mathfrak T\)!).

\emph{Démonstration.} \(\mathfrak S\) est une somme directe de corps : \(\mathfrak S=\mathfrak K_1+\cdots+\mathfrak K_\gamma\). Il lui correspond une décomposition de \(\mathfrak T E\) en somme directe. Chaque facteur direct de \(\mathfrak T E\) doit contenir l'un des corps \(\mathfrak K_i\), faute de quoi il serait annulé par \(E\); mais \(E\) est l'unité de \(\mathfrak T E\). Par conséquent, \(\mathfrak T E\) est une somme directe d'anneaux d'extension des \(\mathfrak K_i\), c'est-à-dire une somme directe de la forme
\[
\mathfrak T E=a_{11}\mathfrak K_1+\cdots+a_{1p_1}\mathfrak K_1+\cdots+a_{rp_r}\mathfrak K_r,
\]
ou bien \(E\mathfrak T=a_{11}\mathfrak S+\cdots+a_{rp_r}\mathfrak S\), c.q.f.d.

\begin{center}
Théorème général de prolongement.
\end{center}

\emph{Toute représentation irréductible de \(\mathfrak S\) dans \(\Omega\), donnée par un homomorphisme \(\mathsf{H}_i\), se prolonge, d'exactement \(t\) manières distinctes, en des représentations de \(\mathfrak T\).}

La démonstration est analogue à celle des pages 10--11. On a
\[
\mathfrak S_\Omega=E_1\Omega+\cdots+E_s\Omega,\qquad
E_\nu\Omega=\mathfrak S_\Omega E_\nu,
\]
\[
\mathfrak T_\Omega\cdot E=\mathfrak T\cdot E\Omega
=\mathfrak T\mathfrak S\Omega=\mathfrak T\mathfrak S_\Omega,
\]
\[
\mathfrak T_\Omega\cdot E
=\mathfrak T(E_1\mathfrak S_\Omega+\cdots+E_s\mathfrak S_\Omega)
=E_1\mathfrak T_\Omega+\cdots+E_s\mathfrak T_\Omega,
\]
\[
[\mathfrak T_\Omega E_\nu:\Omega]
=[\mathfrak T_\Omega E_\nu:\Omega E_\nu]
=[\mathfrak T_\Omega E_\nu:E_\nu\mathfrak S_\Omega]
\leqq[\mathfrak T_\Omega\cdot E:\mathfrak S_\Omega]
=[\mathfrak T\cdot E:\mathfrak S]=t,
\]
\[
\sum[\mathfrak T_\Omega\cdot E_\nu:\Omega]=st
\]
et donc
\[
[\mathfrak T_\Omega E_\nu:\Omega]=t.
\]
Ainsi \(\mathfrak T_\Omega E_\nu\) se décompose en \(t\) idéaux, dont chacun, en tant que composante de \(\mathfrak T_\Omega E\), et donc aussi de \(\mathfrak T_\Omega\), fournit une représentation de \(\mathfrak T_\Omega\) qui, comme aux pages 10--11, se révèle être un prolongement de celle donnée par \(\mathsf{H}_\nu\). Les autres représentations irréductibles de \(\mathfrak T_\Omega\) --- engendrées par les composantes de \(\mathfrak T_\Omega\) qui ne sont pas contenues dans \(\mathfrak T_\Omega\cdot E_\nu\) --- ne peuvent pas être des prolongements des représentations de \(\mathfrak S_\Omega\) engendrées par les \(\mathsf{H}_\nu\), parce que ces idéaux sont annulés par \(\mathfrak T_\Omega\cdot E_\nu\), et à plus forte raison par \(\mathfrak S_\Omega\).

\section*{Chapitre III. Groupes abéliens}\label{fr-r823-chapitre-iii-groupes-abeliens}

\subsection*{\NoetherNavAnchor{S14}{§ 14. L’anneau de groupe}\S{} 14. L'anneau de groupe}\label{fr-r823-anneau-de-groupe}

Soit \(\mathfrak G\) un groupe fini d'éléments \(a_1,\ldots,a_h\), et soit \(\mathsf P\) un corps commutatif. En prenant pour multiplication celle du groupe, nous formons un système hypercomplexe
\[
\mathfrak o[\mathfrak G]=a_1\mathsf P+\cdots+a_h\mathsf P,
\]
l'anneau de groupe de \(\mathfrak G\) sur \(\mathsf P\). (M.Z. \S{} 6; Speiser, \emph{Theorie d. Gruppen end. Ord.}, \S{} 48.)

Nous cherchons les représentations de \(\mathfrak o[\mathfrak G]\) dans \(\mathsf P\), ou dans un corps algébriquement clos \(\Omega\) contenant \(\mathsf P\). Comme il a déjà été mentionné, les représentations irréductibles sont engendrées par les idéaux à gauche simples de \(\mathfrak o[\mathfrak G]/\mathfrak C\). La question de savoir quand \(\mathfrak o[\mathfrak G]\) possède un radical présente donc le même intérêt qu'en théorie de Galois. On a à ce sujet le théorème général suivant :

\begin{center}
\emph{\(\mathfrak o[\mathfrak G]\) est complètement réductible si et seulement si l'ordre \(h\) de \(\mathfrak G\) n'est pas divisible par la caractéristique \(p\) de \(\mathsf P\).}
\end{center}

Nous démontrons d'abord :
\[
\text{De } h\equiv0(p) \text{ il résulte que } \mathfrak C\ne 0.
\]
\[
\mathfrak a=\mathsf P(a_1+\cdots+a_h)
\]
est un idéal bilatère de \(\mathfrak o\). En effet,
\[
\mathsf P\cdot(a_1+\cdots+a_h)a
=\mathsf P\cdot(a_1a+\cdots+a_ha)
=\mathsf P(a_1+\cdots+a_h)=\mathfrak a
\]
et
\[
a\cdot\mathsf P(a_1+\cdots+a_h)
=\mathsf P(aa_1+\cdots+aa_h)
=\mathsf P(a_1+\cdots+a_h)=\mathfrak a,
\]

en vertu de la propriété de groupe des \(a_i\). Naturellement, \(\mathfrak a\) n'est pas nul, mais il est nilpotent :
\[
\left(\sum_i a_i\right)^2=\sum_i\sum_j a_i a_j
=h\sum_j a_j=0.
\]

La réciproque,
\[
\text{de } h\not\equiv0(p) \text{ il résulte que } \mathfrak C=0,
\]
n'est démontrée ici que pour les groupes abéliens. (Dans M.Z. \S{} 26, cette partie du théorème est démontrée en général.)

\subsection*{\NoetherNavAnchor{S15}{§ 15. Les anneaux de groupe des groupes abéliens}\S{} 15. Les anneaux de groupe des groupes abéliens}\label{fr-r823-anneaux-de-groupe-groupes-abeliens}

Nous désignons par \(\mathfrak A=\{a_1,\ldots,a_h\}\) un groupe abélien. Alors \(\mathfrak o\) est commutatif, de sorte que nous pouvons appliquer les théorèmes généraux sur les systèmes hypercomplexes commutatifs.

\emph{Si \(h\) n'est pas divisible par la caractéristique \(p\) de \(\mathsf P\), alors \(\mathfrak o[\mathfrak A]\) possède exactement \(h\) représentations irréductibles distinctes dans \(\Omega\). (Elles sont de degré un, puisque \(\mathfrak o\) est commutatif.)}

Il s'ensuit que \(\mathfrak o[\mathfrak A]/\mathfrak C\) a le même rang que \(\mathfrak o\), donc que \(\mathfrak C=0\); c'est la seconde partie, pour les groupes abéliens, du théorème laissé sans démonstration ci-dessus.

\emph{Démonstration.} \(\mathfrak A\) est le produit direct de groupes cycliques \(\mathfrak Z_i\) d'ordres \(h_i\) :
\[
\mathfrak A=\mathfrak Z_1\times\cdots\times\mathfrak Z_t,
\qquad
h=h_1h_2\ldots h_t .
\]
Si \(z_i\) est un élément générateur de \(\mathfrak Z_i\), l'élément qui lui correspond dans une représentation irréductible (de degré un) doit être une racine \(h_i\)-ième de l'unité \(\varepsilon_\nu^{(h_i)}\). Réciproquement, si l'on associe à chaque \(z_i\) une racine \(h_i\)-ième de l'unité \(\varepsilon_\nu^{(h_i)}\), on obtient une représentation irréductible de \(\mathfrak o\) dans \(\Omega\). Or, puisque \(h\not\equiv0(p)\), on a a fortiori \(h_i\not\equiv0(p)\); il existe donc exactement \(h_1h_2\ldots h_t=h\) représentations distinctes de \(\mathfrak o\), c.q.f.d.

(On voit aussi aisément que, si \(h\equiv0(p)\), donc si \(h_i\equiv0(p)\) pour au moins un \(h_i\), il existe au plus \(h_i/p\) racines \(\varepsilon_\nu^{(h_i)}\) distinctes, et par conséquent moins de \(h\) représentations; \(\mathfrak o_\Omega\) doit donc contenir un radical non nul. Mais cela n'entraîne immédiatement rien quant au radical de \(\mathfrak o\).)

Nous noterons \(\Theta_1,\ldots,\Theta_h\) les \(h\) applications homomorphes distinctes de \(\mathfrak o_\Omega\) sur des sous-anneaux de \(\Omega\); à chaque \(\Theta_i\) correspond un facteur direct \(e_i\Omega\) de \(\mathfrak o_\Omega\), considéré comme module de représentation. Il nous arrivera de n'appliquer \(\Theta_i\) qu'à \(\mathfrak A\), et de parler alors d'une représentation de \(\mathfrak A\). Ces applications de \(\mathfrak A\) sur des groupes de racines de l'unité sont les \(h\) caractères de \(\mathfrak A\). Un caractère particulier est le caractère principal, qui associe \(1\) à chaque élément de \(\mathfrak A\). Notons \(\Theta_1\) l'application correspondante.

La représentation par le caractère principal est possible pour tout groupe, et non seulement pour les groupes abéliens. Comme cette représentation est déjà définie dans \(\mathsf P\), si \(\mathfrak o\) est complètement réductible, il doit s'en détacher un facteur direct simple \(E\cdot\mathfrak o\) qui fournisse cette représentation. Nous pouvons calculer \(E\) sans difficulté. Il faut que \(E=\sum x_i a_i\). Premièrement,
\[
aE=E\cdot 1\quad\text{ou}\quad \sum x_i a a_i=\sum x_i a_i,
\]
d'où, pour \(a=a_j a_i^{-1}\),
\[
x_1a_ja_i^{-1}a_1+\cdots+x_i a_j+\cdots+x_na_ja_i^{-1}a_n
=x_1a_1+\cdots+x_na_n,
\]
et par conséquent
\[
x_i=x_j.
\]

Deuxièmement,

\[
E=E^2=x^2\left(\sum a_i\right)^2=x^2h\sum a_i=xhE.
\]

Dans le cas \(h\equiv0(p)\), il en résulte que \(E=0\), c'est-à-dire que \(\mathfrak o\) ne contient pas le module de représentation du caractère principal et ne pouvait donc pas être complètement réductible. Ainsi, une moitié du théorème de la page 13 est de nouveau démontrée.

Mais si \(h\not\equiv0(p)\), et donc si \(\mathfrak o\) est complètement réductible, on obtient immédiatement

\[
x=\frac{1}{h},\qquad E=\frac{1}{h}(a_1+\cdots+a_h).
\]

\subsection*{\NoetherNavAnchor{S16}{§ 16. Les relations entre caractères}\S{} 16. Les relations entre caractères}\label{fr-r823-relations-entre-caracteres}

Le fait que \(E=(1/h)\cdot\sum a_i\) fournisse le caractère principal conduit immédiatement aux relations entre caractères. Il faut avoir

\[a \cdot e_i = e_i \cdot \Theta_i a\]

et donc

\[e_i \cdot \Theta_i e_j = \begin{cases} e_i & i = j \\ 0 & i \neq j \end{cases}\]

ou encore

\[\Theta_i e_j = \begin{cases} 1 & i = j \\ 0 & i \neq j \end{cases}\]

En particulier,

\[\Theta_i e_1 = \begin{cases} 1 & i = 1 \\ 0 & i \neq 1 \end{cases}\]

ou, puisque \(e_1 = (1/h) \sum a_i\),

\[
\sum_\nu \Theta_i a_\nu=\begin{cases} h & i=1\\ 0 & i\ne 1 \end{cases}
\]

Ce sont les relations connues entre caractères :

La somme des valeurs d'un caractère sur tous les éléments du groupe est nulle, sauf pour le caractère principal, pour lequel elle vaut \(h\).

\subsection*{\NoetherNavAnchor{S17}{§ 17. La théorie de Galois des groupes abéliens}\S{} 17. La théorie de Galois des groupes abéliens}\label{fr-r823-theorie-galois-groupes-abeliens}

Pour les anneaux de groupe des groupes abéliens, on peut établir des théorèmes analogues à ceux que l'on démontre dans la théorie de Galois des corps commutatifs. Dans cette théorie de Galois, les \(\Theta_i\Theta_1^{-1}\) formaient un groupe \(\mathfrak G\). Le théorème principal exprime une correspondance bijective, fondée sur certaines propriétés d'invariance, entre les sous-groupes de \(\mathfrak G\) et les sous-corps de \(\mathsf{Z}\).

Les homomorphismes \(\Theta_i\) de \(\mathfrak o[\mathfrak A]\) peuvent être réunis en un groupe \(\mathsf A\), dont les sous-groupes entretiennent avec les sous-groupes de \(\mathfrak A\) une relation analogue.

Si l'on définit le produit \(\Theta_i \Theta_k = \Theta\) par

\[\Theta a = \Theta_i a \cdot \Theta_k a,\]

l'ensemble des \(\Theta_i\) est un groupe \(\mathsf A\) isomorphe à \(\mathfrak A\).

\emph{Démonstration.} En vertu de

\[\begin{aligned}
\Theta a \cdot \Theta b &= \Theta_i a \cdot \Theta_k a \cdot \Theta_i b \cdot \Theta_k b \\
&= \Theta_i a \cdot \Theta_i b \cdot \Theta_k a \cdot \Theta_k b \\
&= \Theta_i a b \cdot \Theta_k a b = \Theta a b,
\end{aligned}\]

\(\Theta = \Theta_i \cdot \Theta_k\) est un homomorphisme de \(\mathfrak A\) sur \(\mathsf A\); comme il n'existe pas d'homomorphismes de \(\mathfrak A\) autres que les \(\Theta_i\), \(\Theta\) est égal à un certain \(\Theta_j\). Les \(\Theta_i\) forment donc bien un groupe de \(h\) éléments.

Désignons maintenant par \(\varepsilon_{h_\nu}\) une racine primitive \(h_\nu\)-ième de l'unité. Alors \(\Theta_i\) envoie l'élément générateur \(z_\nu\) du facteur cyclique \(\mathfrak Z_\nu\) de \(\mathfrak A\) sur une puissance de \(\varepsilon_{h_\nu}\) :

\[
\Theta_i z_\nu=\varepsilon_{h_\nu}^{\varrho_\nu^{(i)}}.
\]

À l'homomorphisme \(\Theta_i\), associons l'élément \(\prod_\nu z_\nu^{\varrho_\nu^{(i)}}\) de \(\mathfrak A\). Il s'agit d'une application bijective de \(\mathsf A\) sur \(\mathfrak A\), car \(\Theta_i\) étant déterminé par les \(\varrho_\nu^{(i)}\), à des \(\Theta_i\) distincts correspondent des \(\prod_\nu z_\nu^{\varrho_\nu^{(i)}}\) distincts. L'application est isomorphe. En effet, d'une part,

\[
\Theta_i\Theta_j z_\nu=\Theta_i z_\nu\cdot \Theta_j z_\nu
=\varepsilon_{h_\nu}^{\varrho_\nu^{(i)}+\varrho_\nu^{(j)}},
\]

et, d'autre part,

\[
\prod_\nu z_\nu^{\varrho_\nu^{(i)}}\cdot
\prod_\nu z_\nu^{\varrho_\nu^{(j)}}=
\prod_\nu z_\nu^{\varrho_\nu^{(i)}+\varrho_\nu^{(j)}},
\quad\text{c.q.f.d.}
\]

Nous définissons maintenant le

\emph{domaine des invariants d'un sous-groupe \(\mathsf B\) de \(\mathsf A\)} :

le sous-groupe \(\mathfrak B\) des éléments de \(\mathfrak A\) que tout \(\Theta\in\mathsf B\) envoie sur \(1\);

et le

\emph{groupe d'invariance d'un sous-groupe \(\mathfrak B\) de \(\mathfrak A\)} :

le sous-groupe \(\mathsf B\) des éléments de \(\mathsf A\) qui envoient tout \(a\in\mathfrak B\) sur \(1\).

Nous pouvons alors formuler le théorème principal :

Les sous-groupes \(\mathfrak B\) de \(\mathfrak A\) et les sous-groupes \(\mathsf B\) de \(\mathsf A\) peuvent être mis en correspondance bijective de telle sorte que, si \(\mathfrak B\) et \(\mathsf B\) se correspondent, \(\mathsf B\) soit le groupe d'invariance de \(\mathfrak B\), et \(\mathfrak B\) le domaine des invariants de \(\mathsf B\).

Nous devons donc montrer :

\begin{enumerate}
\item Si \(\mathsf B\) est le groupe d'invariance de \(\mathfrak B\), alors \(\mathfrak B\) est le domaine des invariants de \(\mathsf B\).
\item Si \(\mathfrak B\) est le domaine des invariants de \(\mathsf B\), alors \(\mathsf B\) est le groupe d'invariance de \(\mathfrak B\).
\end{enumerate}

\paragraph{1.}
Cela découle, comme à la page 10, du théorème général de prolongement. Supposons que \(\mathsf B\) soit le groupe d'invariance de \(\mathfrak B\), et que \(\mathfrak B'\) soit le domaine des invariants de \(\mathsf B\); \(\mathfrak B'\) contient \(\mathfrak B\). Les éléments de \(\mathsf B\) sont les prolongements à tout \(\mathfrak A\) de l'unique homomorphisme de \(\mathfrak B\) qui envoie \(\mathfrak B\) sur \(1\). Si l'on restreint ces prolongements à \(\mathfrak B'\), il résulte du théorème général de prolongement qu'ils se répartissent en \(t\) classes, de sorte que les homomorphismes d'une même classe soient identiques sur \(\mathfrak B'\), et ceux de classes distinctes, distincts sur \(\mathfrak B'\); ici \(t\) est le degré de \(\mathfrak o[\mathfrak B']\) sur \(\mathfrak o[\mathfrak B]\), donc l'indice de \(\mathfrak B\) dans \(\mathfrak B'\), car \(\mathfrak o[\mathfrak B]\) et \(\mathfrak o[\mathfrak B']\) ont le même élément unité, à savoir l'élément unité de tout \(\mathfrak A\). Or, puisque \(\mathfrak B'\) est le domaine des invariants de \(\mathsf B\), il faut que \(t=1\), donc que \(\mathfrak B'=\mathfrak B\), c.q.f.d.

\paragraph{2. Nous ramenons cette assertion à la première.}

Considérons les éléments \(a\) de \(\mathfrak A\) comme des homomorphismes de \(\mathsf A\). Nous posons en effet \(a\Theta\) égal à la racine de l'unité \(\Theta a\), en abrégé \(a\Theta=\Theta a\). Ainsi, \(a\) est bien défini comme homomorphisme de \(\mathsf A\), car
\[
a\Theta\cdot a\Theta'=\Theta a\cdot\Theta' a=\Theta\Theta' a=a\Theta\Theta'.
\]

Le produit de deux homomorphismes \(a,b\) est le produit de groupe ordinaire :
\[
\begin{array}{cccccc}
ab\Theta&=&a\Theta\cdot b\Theta&=&\Theta a\cdot\Theta b=\Theta ab&=ab\Theta\\
&&\text{produit d'homomorphismes}&&\text{produit de groupe}&
\end{array}
\]

Des éléments \(a,b\) distincts comme éléments du groupe sont aussi distincts comme homomorphismes : si \(a\Theta=b\Theta\) pour tout \(\Theta\), alors \(\Theta a=\Theta b\) pour tout \(\Theta\), ou \(\Theta ab^{-1}=1\) pour tout \(\Theta\). Les \(\Theta\) sont donc les prolongements à tout \(\mathfrak A\) de l'homomorphisme identique du sous-groupe cyclique \(\mathfrak Z\) de \(\mathfrak A\) engendré par \(ab^{-1}\). D'après le théorème de prolongement, l'indice de \(\mathfrak Z\) dans \(\mathfrak A\) est donc égal à \(h\), c'est-à-dire \(\mathfrak Z=e\), et \(a=b\).

Traduisons maintenant dans l'ancienne interprétation les assertions exprimées dans la nouvelle :
\[
\begin{array}{ll}
\text{a)}& \mathfrak B\ \text{est le groupe d'invariance de }\mathsf B.\\
\text{b)}& \mathsf B\ \text{est le domaine des invariants de }\mathfrak B.
\end{array}
\]

Pour a), \(\mathfrak B\) se compose de tous les \(a\in\mathfrak A\) tels que \(a\Theta=1\) pour tout \(\Theta\in\mathsf B\), c'est-à-dire de tous les \(a\in\mathfrak A\) tels que \(\Theta\cdot a=1\) pour tout \(\Theta\in\mathsf B\). La traduction de a) est donc
\[
\text{a')} \quad \mathfrak B\ \text{est le domaine des invariants de }\mathsf B.
\]

De même, la traduction de b) est
\[
\text{b')} \quad \mathsf B\ \text{est le groupe d'invariance de }\mathfrak B.
\]

Ainsi, la traduction de l'assertion déjà démontrée en 1,

si \(\mathfrak B\) est le groupe d'invariance de \(\mathsf B\), alors \(\mathsf B\) est le domaine des invariants de \(\mathfrak B\),

est précisément l'assertion qu'il fallait démontrer :

si \(\mathfrak B\) est le domaine des invariants de \(\mathsf B\), alors \(\mathsf B\) est le groupe d'invariance de \(\mathfrak B\).

Si l'on écrit les relations entre caractères de la page 15 pour \(\mathsf A\) au lieu de \(\mathfrak A\), on obtient les équations
\[
\sum_{i=1}^{h} \Theta_i a_\nu =
\begin{cases}
h & \nu=1\\
0 & \nu\ne 1.
\end{cases}
\]

La somme de tous les caractères en un élément du groupe est nulle, sauf pour l'élément unité, pour lequel elle vaut \(h\).

\section*{Chapitre IV. Anneaux simples bilatères}\label{fr-r823-chapitre-iv-anneaux-simples-bilateres}

\subsection*{\NoetherNavAnchor{S18}{§ 18. Un lemme}\S{} 18. Un lemme}\label{fr-r823-un-lemme}

Nous désignons par \(\mathsf K\) un corps gauche, et par \(\mathfrak G\) un groupe d'automorphismes (d'applications isomorphes de \(\mathsf K\) sur lui-même) de \(\mathsf K\). L'ensemble des éléments de \(\mathsf K\) qui restent inchangés par tout automorphisme appartenant à \(\mathfrak G\) est un sous-corps \(\mathsf P\) de \(\mathsf K\), le \emph{corps des invariants de} \(\mathfrak G\).

Si \(\tau\ne 0\) est un élément quelconque de \(\mathsf K\), on obtient un automorphisme de \(\mathsf K\) en associant à chaque élément \(\alpha\) de \(\mathsf K\) l'élément \(\tau^{-1}\alpha\tau\). Par analogie avec la terminologie correspondante de la théorie des groupes, nous appelons un tel automorphisme un \emph{automorphisme intérieur}. Le corps des invariants \(\mathsf P\) du groupe de tous les automorphismes intérieurs d'un corps gauche \(\mathsf K\) est son centre.

Soient maintenant donnés un corps gauche \(\mathsf K\), un groupe d'automorphismes \(\mathfrak G\) de \(\mathsf K\), et son corps des invariants \(\mathsf P\).

\[
\mathfrak M=x_1\mathsf P+\cdots+x_n\mathsf P
\]

est un \(\mathsf P\)-module de rang \(n\). Nous définissons l'action d'un élément \(\gamma\) de \(\mathfrak G\) sur les éléments de \(\mathfrak M\) par

\[
\gamma\cdot x_i=x_i,
\qquad\text{donc}\qquad
\gamma\sum \varrho_i x_i=\sum \varrho_i x_i.
\]

Il en résulte que \(\mathfrak G\) est devenu un groupe d'automorphismes du module étendu

\[
\mathfrak M_{\mathsf K}=x_1\mathsf K+\cdots+x_n\mathsf K;
\]

le domaine des invariants de \(\mathfrak G\) dans \(\mathfrak M_{\mathsf K}\) est alors \(\mathfrak M\). Sous ces hypothèses, on a le théorème suivant :

Si
\[
\mathfrak C=z_1\mathsf K+\cdots+z_r\mathsf K
\]

est un sous-module de \(\mathfrak M_{\mathsf K}\) que tout élément \(\gamma\) de \(\mathfrak G\) envoie sur lui-même (mais pas nécessairement élément par élément), alors \(\mathfrak C\) est le module étendu, formé avec \(\mathsf K\), d'un sous-module \(\mathfrak c\) de \(\mathfrak M\) :

\[
\mathfrak C=\mathfrak c_{\mathsf K}.
\]

\emph{Démonstration.} 1. Remarque préliminaire : il faut avoir \(\mathfrak c=\mathfrak C\cap\mathfrak M\). En effet, on a en général
\[
\mathfrak c=\mathfrak c_{\mathsf K}\cap\mathfrak M.
\]

L'inclusion \(\mathfrak c\subseteqq\mathfrak c_{\mathsf K}\cap\mathfrak M\) est claire. Pour démontrer \(\mathfrak c_{\mathsf K}\cap\mathfrak M\subseteqq\mathfrak c\), prenons une base \(y_1,\ldots,y_r\) de \(\mathfrak c\) :

\[
\mathfrak c=y_1\mathsf P+\cdots+y_r\mathsf P.
\]

Alors

\[
\mathfrak c_{\mathsf K}=y_1\mathsf K+\cdots+y_r\mathsf K.
\]

Tout élément \(c\) de \(\mathfrak c_{\mathsf K}\cap\mathfrak M\) admet une expression par rapport à la base comme élément de \(\mathfrak c_{\mathsf K}\),
\[
c=\sum_{i=1}^{r}y_i\varkappa_i,
\]
et une expression par rapport à la base comme élément de \(\mathfrak M\),

\[
c=\sum_{i=1}^{r}y_i\varrho_i+
\sum_{j=r+1}^{n}y_j\varrho_j,
\]

où \(y_{r+1},\ldots,y_n\) désigne un système d'éléments de \(\mathfrak M\) qui complète \(y_1,\ldots,y_r\) en une base de \(\mathfrak M\) (M.Z. \S{} 6).

Puisque les deux expressions peuvent être considérées comme des expressions par rapport à une base dans \(\mathfrak M_{\mathsf K}\), elles sont identiques : \(\varkappa_i=\varrho_i\ (i=1,\ldots,r)\), \(\varrho_j=0\ (j=r+1,\ldots,n)\). Par conséquent,

\[
c=\sum_{i=1}^{r}y_i\varrho_i,
\qquad\text{c.q.f.d.}
\]

2. Pour démontrer le théorème en question, construisons une base \(z_1,\ldots,z_r\) de \(\mathfrak C\) dont les éléments \(\mathsf{Z}_i\) appartiennent à \(\mathfrak M\). Alors, en posant

\[
\mathfrak c=z_1\mathsf P+\cdots+z_r\mathsf P,
\]

on aura bien

\[
\mathfrak C=\mathfrak c_{\mathsf K}.
\]

Prenons une base quelconque \(y_1,\ldots,y_r\) de \(\mathfrak C\), et complétons-la en une base de \(\mathfrak M_{\mathsf K}\) par des éléments convenables \(x_{r+1},\ldots,x_n\) d'une base \(x_1,\ldots,x_n\) de \(\mathfrak M\), qui est donc aussi une base de \(\mathfrak M_{\mathsf K}\). Cela est possible d'après M.Z. \S{} 5.

\[
\mathfrak M_{\mathsf K}=\mathfrak C+x_{r+1}\mathsf K+\cdots+x_n\mathsf K.
\]

Il en résulte, pour les \(r\) premiers éléments \(x_1,\ldots,x_r\) de la base de \(\mathfrak M\), des équations

\[
x_i=z_i-\sum_{j=r+1}^{n}x_j\varkappa_j^{(i)},\qquad,
\qquad
z_i\in\mathfrak C,
\quad
\varkappa_j^{(i)}\in\mathsf K.
\]

Comme les \(z_i\), joints à \(x_{r+1},\ldots,x_n\), forment manifestement une base de \(\mathfrak M_{\mathsf K}\), ils sont linéairement indépendants; \(z_1,\ldots,z_r\) est donc une base de \(\mathfrak C\). C'est une base de la forme cherchée : les \(z_i\) appartiennent à \(\mathfrak M\).

Soit \(\gamma\) un élément de \(\mathfrak G\). D'une part,

\[
\begin{array}{rcl}
(1)\qquad \gamma z_i
&=&\gamma x_i+\displaystyle\sum_{j=r+1}^{n}\gamma x_j\cdot \gamma\varkappa_j^{(i)}
=x_i+\displaystyle\sum_{j=r+1}^{n}x_j\gamma\varkappa_j^{(i)}.
\end{array}
\]

D'autre part,
\[
\begin{array}{rcl}
(2)\qquad \gamma z_i&=&\displaystyle\sum_{\mu=1}^{r}z_\mu\sigma_\mu^{(i)},
\qquad \sigma_\mu^{(i)}\in\mathsf K,
\end{array}
\]

puisque \(\mathfrak C\) est stable par tout \(\gamma\in\mathfrak G\). On peut encore écrire (2) sous la forme

\[
\begin{array}{rcl}
(3)\qquad \gamma z_i
&=&\displaystyle\sum_{\mu=1}^{r}x_\mu\sigma_\mu^{(i)}
+\displaystyle\sum_{j=r+1}^{n}x_j\sum_{\mu=1}^{r}\varkappa_j^{(\mu)}\sigma_\mu^{(i)}
\end{array}
\]

La comparaison des coefficients dans (1) et (3) donne

\[
\sigma_j^{(i)}=
\begin{cases}
1 & i=j,\\
0 & i\ne j.
\end{cases}
\]

Par conséquent,

\[
\gamma\cdot z_i=z_i.
\]

Les \(z_i\) appartiennent donc au domaine des invariants de \(\mathfrak G\), c'est-à-dire à \(\mathfrak M\), c.q.f.d.

\subsection*{\NoetherNavAnchor{S19}{§ 19. Représentations des systèmes hypercomplexes simples bilatères dans des corps non commutatifs étendant leurs corps de coefficients}\S{} 19. Représentations des systèmes hypercomplexes simples bilatères dans des corps non commutatifs étendant leurs corps de coefficients}\label{fr-r823-representations-systemes-simples-bilateres}

\[
\mathfrak S=x_1\mathsf P+\cdots+x_n\mathsf P
\]

est un système hypercomplexe simple bilatère, donc complètement réductible, et \(\mathsf K\) un corps gauche dont le centre est \(\mathsf P\).

Théorème 1. \(\mathfrak S_{\mathsf K}\) est simple bilatère (donc aussi complètement réductible).

\emph{Démonstration.} 1. Tout idéal bilatère \(\mathfrak a\) de \(\mathfrak S_{\mathsf K}\) est invariant par le groupe \(\mathfrak G\) des automorphismes intérieurs de \(\mathsf K\). En effet, si, \(\varkappa\) désignant un élément non nul de \(\mathsf K\), nous posons \(\varkappa^{-1}\mathfrak a\varkappa=\overline{\mathfrak a}\), alors \(\overline{\mathfrak a}\subseteqq\mathfrak a\) en vertu de la propriété d'idéal de \(\mathfrak a\); mais, considéré comme \(\mathsf K\)-module, \(\overline{\mathfrak a}\) a le même rang sur \(\mathsf K\) que \(\mathfrak a\), de sorte qu'il faut avoir \(\overline{\mathfrak a}=\mathfrak a\).

2. D'après le lemme que nous venons d'établir, tout idéal bilatère \(\mathfrak a\) de \(\mathfrak S_{\mathsf K}\) doit donc être l'idéal étendu d'un sous-module \(\mathfrak A\) de \(\mathfrak S\), soit \(\mathfrak a=\mathfrak A_{\mathsf K}\). \(\mathfrak A\) est un idéal bilatère de \(\mathfrak S\), car \(\mathfrak A=\mathfrak S\cap\mathfrak a\) donne
\[
\mathfrak S\mathfrak A\subseteqq\mathfrak S\mathfrak S\cap\mathfrak S\mathfrak a
\subseteqq\mathfrak S\cap\mathfrak a=\mathfrak A
\]
et
\[
\mathfrak A\mathfrak S\subseteqq\mathfrak S\mathfrak S\cap\mathfrak a\mathfrak S
\subseteqq\mathfrak S\cap\mathfrak a=\mathfrak A.
\]

Or \(\mathfrak S\) ne contient que les deux idéaux bilatères \(0\) et \(\mathfrak S\); il n'existe donc, dans \(\mathfrak S_{\mathsf K}\), que les deux idéaux bilatères \(0\) et \(\mathfrak S_{\mathsf K}\), c.q.f.d.

Nous passons maintenant aux représentations irréductibles de \(\mathfrak S\) dans \(\mathsf K\); nous établirons plus précisément les représentations réciproques, ce qui offre un avantage formel. En définitive, il est indifférent de prendre des représentations directes ou réciproques, puisqu'elles se correspondent bijectivement.

Soit \(\mathfrak M\) un module de représentation réciproque de \(\mathfrak S\) dans \(\mathsf K\), c'est-à-dire un \(\mathfrak S\)-module à gauche et un \(\mathsf K\)-module à gauche de rang fini,
\[
\mathfrak M=\mathsf K\cdot y_1+\cdots+\mathsf K\cdot y_r,
\]
satisfaisant à la loi de commutation
\[
s(\varkappa m)=\varkappa (sm).
\]

Nous allons montrer que \(\mathfrak M\) peut aussi être considéré comme un \(\mathfrak S_{\mathsf K}\)-module à gauche. Il faut d'abord définir le produit d'un élément de \(\mathfrak S_{\mathsf K}\) par un élément de \(\mathfrak M\). Si \(s\in\mathfrak S\), \(x\in\mathsf K\), \(m\in\mathfrak M\), nous posons
\[
s\varkappa \cdot m=s\cdot(\varkappa \cdot m),
\]
et, par conséquent, pour un élément général \(\sum s_i\varkappa _i\) de \(\mathfrak S_{\mathsf K}\),
\[
\sum s_i\varkappa _i\cdot m=\sum s_i\cdot(\varkappa _i\cdot m).
\]

On vérifie que cette définition est univoque, c'est-à-dire qu'elle ne dépend pas de la manière dont on représente un élément de \(\mathfrak S_{\mathsf K}\) comme somme de produits \(s_i\varkappa _i\).

Il s'agit donc de déduire de
\[
\begin{array}{rcl}
(1)\qquad \displaystyle\sum_{i=1}^{p}s_i\varkappa _i&=&0
\end{array}
\]
que
\[
\sum_{i=1}^{p}s_i(\varkappa _i m)=0.
\]
Supposons que \(s_1,\ldots,s_q\) soient linéairement indépendants et que \(s_{q+1},\ldots,s_p\) s'expriment au moyen de \(s_1,\ldots,s_q\) :
\[
s_j=\sum_{i=1}^{q}s_i\varrho_i^{(j)},
\qquad j=q+1,\ldots,p.
\]

Les coefficients doivent appartenir à \(\mathsf P\), car \(s_j\) ne peut s'exprimer que d'une seule manière, et l'expression est déjà possible dans \(\mathfrak S\). L'équation (1) donne

\[
\varkappa_i+\sum_{j=q+1}^{p}\varrho_i^{(j)}\varkappa_j=0,
\qquad i=1,\ldots,q.
\]

Par conséquent,

\[
\begin{aligned}
\sum_{i=1}^{p}s_i(\varkappa _i m)
&=s_1(\varkappa _1m)+\cdots+s_q(\varkappa _qm)
  +\left(\sum_i s_i\varrho_i^{(q+1)}\right)(\varkappa _{q+1}m)\\
&\quad+\cdots+\left(\sum_i s_i\varrho_i^{(p)}\right)(\varkappa _pm)\\
&=s_1(\varkappa _1m)+\cdots+s_q(\varkappa _qm)\\
&\quad+s_1\varrho_1^{(q+1)}(\varkappa _{q+1}m)+\cdots+s_q\varrho_q^{(q+1)}(\varkappa _{q+1}m)
  \qquad\text{par la propriété de module}\\
&\quad+\cdots\\
&\quad\cdots\\
&\quad+s_1\varrho_1^{(p)}(\varkappa _pm)+\cdots+s_q\varrho_q^{(p)}(\varkappa _pm)\\
&=s_1(\varkappa _1m)+\cdots+s_q(\varkappa _qm)\\
&\quad+s_1(\varrho_1^{(q+1)}\varkappa _{q+1}m)+\cdots+s_q(\varrho_q^{(q+1)}\varkappa _{q+1}m)\\
&\quad+\cdots\\
&\quad\cdots\\
&\quad+s_1(\varrho_1^{(p)}\varkappa _pm)+\cdots+s_q(\varrho_q^{(p)}\varkappa _pm)
  \qquad\begin{gathered}\text{par les lois}\\ \text{d'associativité }s'(sm)\\ =s's\cdot m\text{ et}\\ \varkappa '(\varkappa m)=\varkappa '\varkappa \cdot m\end{gathered}\\
&=s_1\left(\varkappa _1m+\sum\varrho_1^{(q+1)}\varkappa _{q+1}\cdot m\right)+\cdots\\
&=s_1\left(\left(\varkappa _1+\sum\varrho_1^{q+1}\varkappa _{q+1}\right)m\right)+\cdots\\
&=0.
\end{aligned}
\]

Parmi les propriétés de \(\mathfrak M\) comme \(\mathfrak S_{\mathsf K}\)-module, les lois distributives \((S+S')m=Sm+S'm\), \(S(m+m')=Sm+Sm'\) sont claires. La loi d'associativité \(S(S'm)=SS'\cdot m\) demande encore quelque attention; puisque les lois distributives sont vérifiées, nous pouvons nous borner au cas

\[
s\varkappa \cdot(s'\varkappa ' m)=s\varkappa (s'\varkappa '\cdot m).
\]

D'une part, d'après la définition de \(s\varkappa \cdot m\),

\[
s\varkappa \cdot(s'\varkappa '\cdot m)=s\cdot(\varkappa \cdot(s'\cdot(\varkappa '\cdot m)));
\]

d'autre part,

\[
s\varkappa  s'\varkappa '\cdot m=ss'\varkappa \varkappa '\cdot m=ss'\cdot(\varkappa \varkappa '\cdot m)
=s\cdot(s'\cdot(\varkappa \cdot(\varkappa '\cdot m))),
\]

ce qui, en vertu des deux relations de commutation, est égal à \(=s\cdot(\varkappa \cdot(s'\cdot(\varkappa '\cdot m)))\) :

\[
\varkappa \cdot s'=s'\cdot \varkappa,
\]
\[
s'\cdot(\varkappa \cdot m)=\varkappa \cdot(s'\cdot m).
\]

Ainsi, \(\mathfrak M\) est bien un \(\mathfrak S_{\mathsf K}\)-module à gauche, évidemment fini. Réciproquement, tout \(\mathfrak S_{\mathsf K}\)-module à gauche fini est un module de représentation réciproque de \(\mathfrak S\) dans \(\mathsf K\), car il est manifestement un \(\mathsf K\)-module fini et satisfait aussi à la condition de commutation, puisque

\[
s(\varkappa \cdot m)=s\varkappa \cdot m=\varkappa s\cdot m=\varkappa (s\cdot m).
\]

(Voir aussi la démonstration plus simple du \S{} 24.)

Or, d'après M.Z. \S{} 18, tout \(\mathfrak S_{\mathsf K}\)-module fini est somme de modules simples qui sont, comme modules à opérateurs sur \(\mathfrak S_{\mathsf K}\), isomorphes aux idéaux à gauche simples de \(\mathfrak S_{\mathsf K}\). Réciproquement, tout idéal à gauche de \(\mathfrak S_{\mathsf K}\) est un module de représentation de \(\mathfrak S\) dans \(\mathsf K\); comme tous les idéaux à gauche de \(\mathfrak S_{\mathsf K}\) sont isomorphes comme modules à opérateurs, on obtient le

\textbf{Théorème 2.} \emph{L'unique classe de représentations réciproques irréductibles de \(\mathfrak S\) dans \(\mathsf K\) est celle qui est engendrée par un idéal à gauche simple de \(\mathfrak S_{\mathsf K}\).}

Si cette représentation irréductible est de degré \(r\), le degré de toute représentation est un multiple de \(r\). En juxtaposant des représentations irréductibles, on peut naturellement construire, pour tout multiple donné \(rs\) de \(r\), une représentation de degré \(rs\); cela correspond à la formation d'un module de représentation qui est somme directe de \(s\) modules simples.

Nous connaissons déjà une représentation de \(\mathfrak S\) dans \(\mathsf K\), à savoir celle qui est engendrée par \(\mathfrak S\) lui-même comme module de représentation (représentation principale, M.Z. \S{} 20); elle est déjà définie dans \(\mathsf P\). Son degré est \(n\), donc \(n\) est divisible par \(r\) :

\[
n=rt.
\]

Le quotient \(t\) est facile à déterminer : puisque \(r\) est le rang d'un seul idéal à gauche sur \(\mathsf K\), et \(n\) le rang de tout \(\mathfrak S_{\mathsf K}\), \(t\) est le nombre d'idéaux à gauche de \(\mathfrak S_{\mathsf K}\); ainsi \(t^2\) est le rang de \(\mathfrak S_{\mathsf K}\) sur le corps gauche des endomorphismes de ses idéaux simples.

Spécialisons nos considérations au cas où \(\mathfrak S\) est un corps commutatif.

\textbf{Théorème 3.} \emph{Si \(\mathfrak S\) est un corps commutatif, alors \(\mathfrak S\) est le centre de \(\mathfrak S_{\mathsf K}\).}

\emph{Démonstration.} Chaque élément de \(\mathfrak S\) commutant avec tout élément de \(\mathfrak S\) et avec tout élément de \(\mathsf K\), donc avec tout élément de \(\mathfrak S_{\mathsf K}\), \(\mathfrak S\) appartient au centre de \(\mathfrak S_{\mathsf K}\). Un élément de \(\mathfrak S_{\mathsf K}\) ne commute avec tout élément de \(\mathfrak S_{\mathsf K}\) que s'il commute au moins avec tout élément de \(\mathsf K\), donc si ses coefficients appartiennent au centre de \(\mathsf K\), c'est-à-dire à \(\mathsf P\), autrement dit s'il appartient à \(\mathfrak S\), c.q.f.d.

Puisque, pour les anneaux commutatifs, il n'y a pas de différence entre anti-isomorphie et isomorphie directe, toute représentation réciproque peut aussi être considérée comme directe. La représentation irréductible de \(\mathfrak S\) dans \(\mathsf K\) doit donc également être fournie par un module de représentation direct. Ce module de représentation direct est un idéal à droite simple de \(\mathfrak S_{\mathsf K}\). En effet, \(\mathfrak r\) est un \(\mathsf K\)-module à droite fini de rang \(r\), et, puisque \(\mathfrak r\mathfrak S=\mathfrak S\mathfrak r\), également un \(\mathfrak S\)-module à gauche (avec, bien entendu, la loi d'associativité mixte); c'est donc un module de représentation direct de rang \(r\), et, puisqu'il n'existe qu'une classe de représentations irréductibles, il doit engendrer la représentation donnée.

\subsection*{\NoetherNavAnchor{S20}{§ 20. Corps non commutatifs}\S{} 20. Corps non commutatifs}\label{fr-r823-corps-non-commutatifs}

Nous désignons maintenant par

\[
\mathfrak K=y_1\mathsf P+\cdots+y_m\mathsf P
\]

un corps gauche dont le centre est \(\mathsf P\). Nous allons établir pour \(\mathfrak K\) une série de théorèmes analogues à ceux de la théorie des extensions algébriques finies de corps commutatifs.

Si \(\mathsf{Z}=\xi_1\mathsf P+\cdots+\xi_n\mathsf P\) désigne un système hypercomplexe sur \(\mathsf P\), les deux systèmes étendus \(\mathfrak K_{\mathsf{Z}}\) et \(\mathsf{Z}_{\mathfrak K}\) sont identiques. En effet, nous pouvons écrire pour \(\mathfrak K_{\mathsf{Z}}\)

\[
\mathfrak K_{\mathsf{Z}}=y_1\xi_1\mathsf P+\cdots+y_m\xi_n\mathsf P,
\]

et pour \(\mathsf{Z}_{\mathfrak K}\)

\[
\mathsf{Z}_{\mathfrak K}=\xi_1y_1\mathsf P+\cdots+\xi_ny_m\mathsf P;
\]

or les deux expressions sont identiques, puisque \(y_i\xi_j=\xi_jy_i\).

Dans ce qui suit, \(\mathsf{Z}\) sera toujours un corps commutatif; les théorèmes du paragraphe 19 sont donc applicables. D'après le théorème 1 du paragraphe 19, \(\mathfrak K_{\mathsf{Z}}\) est toujours simple bilatère.

\textbf{Théorème 1.} \emph{Soit \(\Omega\) un corps algébriquement clos contenant \(\mathsf P\). Alors le centre de \(\mathfrak K_\Omega\) est \(\Omega\).}

\emph{Démonstration.} Nous nous appuyons sur le théorème 3 correspondant du \S{} 19. Il est clair que \(\Omega\) est contenu dans le centre de \(\mathfrak K_\Omega\). Si \(\alpha\) est un élément quelconque du centre de \(\mathfrak K_\Omega\), ses coefficients appartiennent à une extension algébrique finie \(\mathsf{Z}\) de \(\mathsf P\), à savoir celle qu'ils engendrent. Évidemment, \(\alpha\) doit aussi appartenir au centre de \(\mathfrak K_{\mathsf{Z}}\). Or ce centre est \(\mathsf{Z}\); ainsi \(\alpha\) appartient à \(\mathsf{Z}\), et par conséquent aussi à \(\Omega\), c.q.f.d.

Nous démontrons de manière tout à fait analogue le

\textbf{Théorème 2.} \emph{\(\mathfrak K_\Omega\) est simple bilatère.}

\emph{Démonstration.} Soit \(\mathfrak a=z_1\Omega+\cdots+z_r\Omega\) un idéal bilatère. Les coefficients de \(z_1,\ldots,z_r\) appartiennent à une extension finie \(\mathsf{Z}\) de \(\mathsf P\); dans \(\mathfrak K_{\mathsf{Z}}\), \(\mathfrak a'=z_1\mathsf{Z}+\cdots+z_r\mathsf{Z}\) est un idéal bilatère. Comme \(\mathfrak K_{\mathsf{Z}}\) est simple bilatère, \(\mathfrak a'\) est soit nul, soit égal à \(\mathfrak K_{\mathsf{Z}}\); par conséquent, \(\mathfrak a\) est soit nul, soit égal à \(\mathfrak K_\Omega\). Ainsi \(\mathfrak K_\Omega\) est simple bilatère, c.q.f.d.

D'après ce théorème, \(\mathfrak K_\Omega\) est un anneau de matrices sur un corps gauche d'extension de son centre \(\Omega\). Ce corps, étant un sous-anneau de \(\mathfrak K_\Omega\) tout entier, doit être de degré fini sur \(\Omega\), donc égal à \(\Omega\) :

\textbf{Théorème 3.} \emph{\(\mathfrak K_\Omega\) est un anneau de matrices sur \(\Omega\),}

\[
\mathfrak K_\Omega=\sum\Omega c_{ik}
\]

\emph{et, par conséquent, le degré de \(\mathfrak K\) sur son centre \(\mathsf P\) est un carré, \(m=t^2\).}

Le nombre \(t\) des idéaux à droite (ou à gauche) en lesquels se décompose \(\mathfrak K_\Omega\) sera appelé, en abrégé, le nombre absolu de composantes de \(\mathfrak K\).

La propriété de \(\Omega\) selon laquelle \(\mathfrak K_\Omega\) se décompose en \(t\) idéaux simples appartient naturellement déjà à certaines extensions finies \(\mathsf{Z}\) de \(\mathsf P\), par exemple à celles que l'on obtient en adjoignant à \(\mathsf P\) les coefficients des \(c_{ik}\). Nous allons maintenant étudier ces corps de décomposition.

\textbf{Définition.} Une extension algébrique finie \(\mathsf{Z}\) de \(\mathsf P\) est appelée \emph{corps de décomposition} de \(\mathfrak K\) si le nombre d'idéaux à droite simples en lesquels se décompose \(\mathfrak K_{\mathsf{Z}}\) est déjà égal au nombre absolu de composantes \(t\)\srcfn{1}{Comparer la notion de corps de décomposition à la page 9.}.

\textbf{Théorème 4.} \emph{Si \(\mathsf{Z}\) est un corps de décomposition de \(\mathfrak K\), alors \(\mathsf{Z}\) est le corps gauche des endomorphismes des idéaux à droite de \(\mathfrak K_{\mathsf{Z}}\). Réciproquement, si le corps gauche des endomorphismes des idéaux à droite de \(\mathfrak K_{\mathsf{Z}}\) est \(\mathsf{Z}\), alors \(\mathsf{Z}\) est un corps de décomposition.}

\emph{Démonstration.} 1. Soit \(\mathsf{Z}\) un corps de décomposition, et soit \(\mathfrak T\) le corps gauche des endomorphismes des idéaux à droite de \(\mathfrak K_{\mathsf{Z}}\), de sorte que

\[
\mathfrak K_{\mathsf{Z}}=\sum\mathfrak T c_{ik}.
\]

Il en résulte que

\[
\mathfrak K_\Omega=\sum\mathfrak T_\Omega c_{ik}.
\]

Le rang de \(\mathfrak K_\Omega\) sur \(\Omega\) est donc le produit de \(t^2\) par le rang de \(\mathfrak T_\Omega\); d'autre part, il est égal à \(t^2\). Le rang de \(\mathfrak T_\Omega\), et donc aussi celui de \(\mathfrak T\), est par conséquent égal à \(1\); ainsi \(\mathfrak T=\mathsf{Z}\), c.q.f.d.

2. Si \(\mathsf{Z}\) est le corps gauche des endomorphismes des idéaux simples de \(\mathfrak K_{\mathsf{Z}}\), alors

\[
\mathfrak K_{\mathsf{Z}}=\sum\mathsf Zc_{ik}.
\]

Donc

\[
\mathfrak K_\Omega=\sum\Omega c_{ik}.
\]

Le nombre des \(c_{ik}\) doit par conséquent être égal à \(t^2\), c.q.f.d.

\textbf{Théorème 5.} \emph{La représentation irréductible \(\mathfrak Z\) d'un corps de décomposition \(\mathsf{Z}\) de \(\mathfrak K\) dans \(\mathfrak K\) --- il n'existe, d'après le théorème 2 de la page 22, qu'une seule classe de représentations irréductibles --- est un sous-corps commutatif maximal de l'anneau de matrices \(\mathfrak K_r\) (où \(r\) est le degré de la représentation irréductible de \(\mathsf{Z}\)).}

\emph{Réciproquement, si la représentation irréductible \(\mathfrak Z\) d'une extension commutative \(\mathsf{Z}\) de \(\mathsf P\) est un sous-corps commutatif maximal de \(\mathfrak K_r\), alors \(\mathsf{Z}\) est un corps de décomposition.}

\emph{Démonstration.} 1. Si \(n\) est le degré du corps de décomposition \(\mathsf{Z}\) sur \(\mathsf P\), le degré \(r\) de la représentation irréductible \(\mathfrak Z\) de \(\mathsf{Z}\) est égal à \(n/t\) (page 22). Soit \(\mathfrak Z\subseteqq\mathfrak Z^*\subseteqq\mathfrak K_r\), où \(\mathfrak Z^*\) est un corps commutatif maximal de degré \(n^*\geqq n\) sur \(\mathsf P\). Nous devons démontrer que \(\mathfrak Z^*=\mathfrak Z\). Il suffit pour cela de montrer que \(n^*=n\). Si \(r^*\) désigne le degré de la représentation irréductible d'une extension \(\mathsf{Z}^*\) de \(Z\), isomorphe à \(\mathfrak Z^*\), alors on a de nouveau \(n^*=r^*t\). D'autre part, \(r\) est divisible par \(r^*\), puisqu'il existe dans \(\mathfrak Z^*\) une représentation de degré \(r\) de \(\mathsf{Z}^*\); ainsi \(r^*\leqq r\). De

\[
n\leqq n^*=r^*t\leqq rt=n
\qquad\text{il résulte que}\qquad n^*=n.
\]

2. Soit \(\mathfrak Z\), la représentation irréductible de degré \(r\) de l'extension commutative \(\mathsf{Z}\) de \(\mathsf P\) dans \(\mathfrak K\), un sous-corps commutatif maximal de \(\mathfrak K_r\). Soit \(\mathfrak T\) le corps gauche des endomorphismes des idéaux simples de \(\mathfrak K_{\mathsf{Z}}\), de sorte que

\[
\mathfrak K_{\mathsf{Z}}=\sum_{1}^{t'}\mathfrak T c_{ik}.
\]

\(\mathfrak T\) contient \(\mathsf{Z}\). Nous devons montrer que \(\mathfrak T=\mathsf{Z}\). Or un idéal à droite simple \(\mathfrak r\) de \(\mathfrak K_r\) n'est pas seulement un module de représentation de \(\mathsf{Z}\) dans \(\mathfrak K\), mais aussi un module de représentation de \(\mathfrak T\), car

\[
\mathfrak r=\mathfrak T c_{i1}+\cdots+\mathfrak T c_{ik}.
\]

Ainsi \(\mathfrak T\mathfrak r=\mathfrak r\), \(\mathfrak r\) est un \(\mathfrak T\)-module à gauche et, naturellement, un \(\mathfrak K\)-module à droite.

Il existe donc dans \(\mathfrak K_r\) un sous-corps gauche \(\overline{\mathfrak T}\), isomorphe à \(\mathfrak T\), qui contient \(\mathfrak Z\). Si \(\mathfrak T\) était distinct de \(\mathsf{Z}\), donc \(\overline{\mathfrak T}\) de \(\mathfrak Z\), l'adjonction à \(\mathfrak Z\) d'un élément \(\alpha\) de \(\mathfrak T\) n'appartenant pas à \(\mathfrak Z\) donnerait dans \(\mathfrak K_r\) une extension commutative propre de \(\mathfrak Z\), contrairement à la maximalité de \(\mathfrak Z\), c.q.f.d.

Du théorème 5 découle immédiatement le

\textbf{Théorème 6.} \emph{Les sous-corps commutatifs maximaux de \(\mathfrak K\) ont le degré \(t\) sur \(\mathsf P\) et sont des corps de décomposition\srcfn{1}{La conjecture selon laquelle tout corps de décomposition contiendrait alors aussi un sous-corps commutatif maximal de \(\mathfrak K\) est fausse; il peut exister des corps de décomposition qui ne possèdent pas cette propriété, et le \(r\) correspondant est nécessairement supérieur à 1. Voir R. Brauer et E. Noether, \emph{Über minimale Zerfällungskörper irreduzibler Darstellungen}; \emph{Sitzungsberichte d. preuss. Akademie der Wisssch.} 1927, XXXII, p. 221.}.}

Nous établissons maintenant un théorème sur l'isomorphie de sous-anneaux d'un \(\mathfrak K_r\), sur lequel reposent une connaissance plus précise de la structure du groupe de Galois d'un \(\mathfrak K\) et les faits qui s'y rattachent.

\textbf{Théorème 7.} \emph{Soient \(\mathfrak S_1\) et \(\mathfrak S_2\) des sous-anneaux simples bilatères de l'anneau de matrices \(\mathfrak K_r\) sur \(\mathfrak K\), contenant \(\mathsf P\), et supposons qu'il existe entre eux un isomorphisme \(\mathfrak S_1\cong\mathfrak S_2\) qui fixe chaque élément de \(\mathsf P\). Alors cet isomorphisme est la transformation par un élément inversible \(\tau\) de \(\mathfrak K_r\).}

\emph{Si \(\sigma_1\in\mathfrak S_1\) et \(\sigma_2\in\mathfrak S_2\) se correspondent, alors}

\[
\tau^{-1}\sigma_1\tau=\sigma_2.
\]

\emph{Démonstration.} Construisons un anneau \(\Sigma\) anti-isomorphe à \(\mathfrak S_1\) et à \(\mathfrak S_2\), et identifions à \(\mathsf P\) son sous-corps isomorphe à \(\mathsf P\). Alors \(\mathfrak S_1\) et \(\mathfrak S_2\) sont deux représentations réciproques de même degré \(r\) de \(\Sigma\) dans \(\mathfrak K\). Or, d'après le théorème 2 de la page 22, il n'existe qu'une seule classe de représentations réciproques irréductibles, et par conséquent une seule classe de représentations de chaque degré donné de \(\Sigma\) dans \(\mathfrak K\) (multiple de la classe irréductible). Ainsi \(\mathfrak S_1\) et \(\mathfrak S_2\) doivent appartenir à la même classe; elles se déduisent donc l'une de l'autre par transformation, c.q.f.d.

Du théorème 7 découle immédiatement le

\textbf{Théorème 8.} \emph{Le groupe des automorphismes de \(\mathfrak K\) qui fixent chaque élément de \(\mathsf P\) --- c'est-à-dire le groupe de Galois \(\mathfrak G\) de \(\mathfrak K\) relativement à son centre --- est le groupe des automorphismes intérieurs de \(\mathfrak K\).}

Il suffit de poser, dans le théorème 7, \(r=1\) et \(\mathfrak S_1=\mathfrak S_2=\mathfrak K\).

Les deux théorèmes suivants montrent l'utilité de la théorie générale développée jusqu'ici pour traiter des questions particulières.

\textbf{Théorème 9.} \emph{Le seul corps gauche de degré fini dont le centre est le corps \(\mathsf{R}\) des nombres réels est le corps des quaternions.}

\emph{Démonstration.} Soit \(\mathfrak K\) un corps gauche de centre \(\mathsf{R}\) (et de rang fini sur \(\mathsf{R}\)). D'après le théorème 6, le degré de \(\mathfrak K\) sur \(\mathsf{R}\) est le carré du degré \(t\) d'un sous-corps commutatif maximal de \(\mathfrak K\); comme celui-ci ne peut être qu'un corps isomorphe au corps des nombres complexes, on a \(t=2\), et \(\mathfrak K\) est de degré \(4\) sur \(\mathsf{R}\). Soit \(\mathfrak Z_1\) un sous-corps commutatif maximal de \(\mathfrak K\); nous pouvons poser \(\mathfrak Z_1=\mathsf{R}(i)\), où \(i^2=-1\). Alors \(\mathfrak Z_2=\mathsf{R}(-i)\) est aussi un corps isomorphe à \(\mathfrak Z_1\), qui coïncide avec \(\mathfrak Z_1\) comme ensemble; l'isomorphisme est donné par \(i\mapsto-i\). D'après le théorème 7, il existe donc un élément \(j^*\in\mathfrak K\) tel que

\[
j^{*-1}ij^*=-1.
\]

Le rang de \(\mathfrak Z_1(j^*)\) sur \(\mathfrak Z_1\) est au moins égal à \(2\), puisque \(j^*\) ne peut appartenir à \(\mathfrak Z_1\). Ainsi \(\mathfrak Z_1(j^*)\) est de rang \(4\) sur \(\mathsf{R}\), donc égal à \(\mathfrak K\), et nous pouvons écrire

\[
\mathfrak K=\mathsf R+\mathsf Ri+\mathsf Rj^*+\mathsf Rij^*.
\]

Comme

\[
\begin{aligned}
j^{*-2}1j^{*2}&=1\\
j^{*-2}ij^{*2}&=i\\
j^{*-2}jj^{*2}&=j^*\\
j^{*-2}ij^*j^{*2}&=ij^*
\end{aligned}
\]

\(j^{*2}\) commute avec tous les éléments de \(\mathfrak K\); il appartient donc au centre \(\mathsf{R}\), et \(j^{*2}\) est un nombre réel \(r\). Le nombre \(r\) ne peut pas être positif, car le polynôme \(x^2-r\) aurait alors, dans le corps commutatif \(\mathsf{R}(j^*)\), les quatre racines \(+\sqrt r,-\sqrt r,+j^*\) et \(-j^*\). Ainsi \(j^{*2}=-m^2\), avec \(m\in R\). Posons \(j=j^*m^{-1}\); alors \(j^2=-1\), et \(1,i,j,k=ij\) sont les unités quaternioniennes, c.q.f.d.

La démonstration vaut manifestement si l'on remplace \(\mathsf{R}\) par un corps réel clos quelconque.

\textbf{Théorème 10 (théorème de Wedderburn).} \emph{Tout corps gauche fini est commutatif.}

\emph{Démonstration.} Soit \(\mathfrak K\) un corps gauche ayant un nombre fini d'éléments, et soit \(\mathsf P\) son centre. Soit \(t^2\) le degré de \(\mathfrak K\) sur \(\mathsf P\). Tous les sous-corps commutatifs maximaux de \(\mathfrak K\) sont de degré \(t\) sur \(\mathsf P\); ils ont donc le même nombre d'éléments et sont par conséquent isomorphes d'après un théorème connu de la théorie des corps de Galois (théorème de Moore). Comme tout corps de Galois est une extension galoisienne de son corps premier, ces isomorphismes peuvent être choisis de manière que chaque élément de \(\mathsf P\) se corresponde à lui-même. D'après le théorème 7, tous les sous-corps commutatifs maximaux de \(\mathfrak K\) se déduisent donc d'un même \(\mathfrak Z\) par transformation. Or \(\mathfrak K\) est la réunion de tous ses sous-corps commutatifs maximaux. (L'adjonction d'un élément à \(\mathsf P\) engendre un corps commutatif.) Considérons maintenant seulement le groupe multiplicatif \(\mathfrak K^*\) de \(\mathfrak K\) (c'est-à-dire \(\mathfrak K\) privé de zéro). Il résulte de ce qui précède que \(\mathfrak K^*\) est la réunion des sous-groupes conjugués au sous-groupe \(\mathfrak Z^*\). Mais, si \(t>1\), cela conduit à une contradiction. En effet, si \(\mathfrak K^*\) contient \(M\) éléments, \(\mathfrak Z^*\) en contient \(N\), et si \(L\) est l'indice de \(\mathfrak Z^*\) dans \(\mathfrak K^*\), alors \(M=NL\), et le nombre de sous-groupes conjugués distincts est au plus \(L\). Même si ce nombre était égal à \(L\), il faudrait encore, pour que tous les sous-groupes conjugués épuisent ensemble le groupe tout entier, que deux quelconques d'entre eux n'aient aucun élément commun. Or ils ont tous l'élément unité en commun.

\subsection*{\NoetherNavAnchor{S21}{§ 21. La théorie de Galois des corps non commutatifs}\S{} 21. La théorie de Galois des corps non commutatifs}\label{fr-r823-theorie-galois-corps-non-commutatifs}

Nous désignons par \(\mathsf{G}\) le groupe de Galois de \(\mathfrak K\) relativement à son centre \(\mathsf P\). D'après le théorème 8, \(\mathsf{G}\) est isomorphe à \(\overline{\mathsf G}=\mathfrak K^*/\mathsf P^*\); nous regardons cet isomorphisme comme fixé une fois pour toutes. Nous appelons \emph{fermé} un sous-groupe \(\mathsf{H}\) de \(\mathsf{G}\) lorsque le sous-groupe \(\mathfrak H^*\) de \(\mathfrak K^*\), qui lui correspond par l'isomorphisme \(G\cong\mathfrak K^*/\mathsf P^*\), devient un corps gauche après adjonction de zéro.

\textbf{Théorème 11. Théorème principal de la théorie de Galois.} \emph{Les sous-groupes fermés \(\mathsf{H}\) de \(\mathsf{G}\) et les corps gauches \(\mathfrak S\) compris entre \(\mathsf P\) et \(\mathfrak K\) peuvent être mis en correspondance bijective de telle sorte que, si \(\mathsf{H}\) et \(\mathfrak S\) se correspondent, \(\mathsf{H}\) soit le groupe d'invariance de \(\mathfrak S\), et \(\mathfrak S\) le corps gauche des invariants de \(\mathsf{H}\).}

\emph{Démonstration.} Décomposons notre théorème en trois assertions :
\begin{enumerate}
\item Le groupe d'invariance \(\mathsf{H}\) d'un corps gauche \(\mathfrak S\) est fermé.
\item Si \(\mathfrak S\) a pour groupe d'invariance \(\mathsf{H}\), alors \(\mathfrak S\) est le corps gauche des invariants de \(\mathsf{H}\).
\item Si \(\mathsf{H}\) a pour corps gauche des invariants \(\mathfrak S\), alors \(\mathsf{H}\) est le groupe d'invariance de \(\mathfrak S\).
\end{enumerate}

\begin{enumerate}
\item Par l'isomorphisme \(G\cong\mathfrak K^*/\mathsf P^*\), au groupe d'invariance \(H\) d'un corps gauche \(\mathfrak S\) correspond l'ensemble \(\mathfrak H^*\) de tous les éléments non nuls de \(\mathfrak K\) qui commutent élément par élément avec \(\mathfrak S\). Avec \(0\), ils forment naturellement un corps gauche; ainsi \(\mathsf{H}\) est fermé.

\item Nous ramenons la troisième assertion à la deuxième. Au corps gauche \(\mathfrak S\) est associé un sous-groupe \(\mathsf{S}\) de \(\mathsf{G}\), à savoir le sous-groupe de \(G\) qui correspond à \(\mathfrak S^*/\mathsf P^*\). De même, au groupe \(\mathsf{H}\) correspond un sous-corps gauche \(\mathfrak H\) de \(\mathfrak K\), et l'on a \(\mathsf{H}\cong \mathfrak H^*/\mathsf P^*\). Si \(\mathfrak S\) est le corps gauche des invariants de \(\mathsf{H}\), alors \(\mathfrak S\) se compose de tous les éléments de \(\mathfrak K\) qui commutent élément par élément avec \(\mathfrak H\); on peut donc dire aussi que \(\mathsf{S}\) est le groupe d'invariance de \(\mathfrak H\). De façon correspondante, si \(\mathsf{H}\) est le groupe d'invariance de \(\mathfrak S\), alors \(\mathfrak H\) est le corps gauche des invariants de \(\mathsf{S}\). La troisième assertion peut donc aussi se formuler ainsi :

\item[\(3'\).] Si \(\mathfrak H\) a pour groupe d'invariance \(\mathsf{S}\), alors \(\mathfrak H\) est le corps gauche des invariants de \(\mathsf{S}\); c'est précisément la deuxième assertion, qu'il reste maintenant à démontrer.

\item[3.] Cette démonstration se déroule, à quelques petites différences près, comme dans le cas commutatif :
\end{enumerate}

Construisons un corps gauche \(K\) anti-isomorphe à \(\mathfrak K\), dont nous identifions le centre à \(\mathsf P\). Au lieu des automorphismes de \(\mathfrak K\), nous pouvons considérer les anti-isomorphismes de \(\mathfrak S\) sur des sous-corps gauches de \(\mathsf{K}\).

Nous allons de nouveau démontrer un lemme analogue au théorème de prolongement employé dans le cas commutatif.

Soient \(\mathfrak S\) et \(\mathfrak T\) des corps gauches compris entre \(\mathsf P\) et \(\mathfrak K\),
\[
\mathsf P \subseteqq \mathfrak S \subseteqq \mathfrak T \subseteqq \mathfrak K.
\]
Nous affirmons que tout anti-isomorphisme de \(\mathfrak S\) sur un sous-corps gauche de \(\mathsf{K}\) (représentation réciproque de degré un de \(\mathfrak S\) dans \(\mathsf{K}\)) peut être prolongé en des anti-isomorphismes de \(\mathfrak T\) d'au moins autant de manières que l'indique le degré \(s\) de \(\mathfrak T\) sur \(\mathfrak S\).

On étend par \(\mathsf{K}\) les corps gauches \(\mathsf P\), \(\mathfrak S\), \(\mathfrak T\), \(\mathfrak K\) :
\[
\mathsf{K}\subseteqq \mathfrak S_{\mathsf{K}}\subseteqq \mathfrak T_{\mathsf{K}}\subseteqq \mathfrak K_{\mathsf{K}}.
\]
Chaque \(\mathfrak S_{\mathsf{K}}\) est, d'après le théorème 1 de la page 20, simple bilatère et complètement réductible; ses idéaux à gauche simples --- isomorphes entre eux comme modules à opérateurs --- sont des modules de représentation pour l'unique classe de représentations irréductibles de \(\mathfrak S\) dans \(\mathsf{K}\). Comme cette classe de représentations est de degré \(1\) --- car, par construction de \(\mathsf{K}\), il existe dans \(\mathsf{K}\) des corps gauches anti-isomorphes à \(\mathfrak S\) ---, tous les idéaux à gauche simples sont de rang \(1\). Autrement dit, \(\mathfrak S_{\mathsf{K}}\) devient somme d'autant d'idéaux à gauche que l'indique son degré sur \(\mathsf P\) :
\[
\mathfrak S_K=K E_1+\cdots+K E_p,
\qquad K E_\nu=\mathfrak S_K E_\nu.
\]

On a alors
\[
\begin{aligned}
\mathfrak T_{\mathsf{K}}&=\mathfrak T\cdot \mathsf{K}=\mathfrak T\cdot \mathfrak S_{\mathsf{K}},\\
\mathfrak T_{\mathsf{K}}&=\mathfrak T\mathfrak S_{\mathsf{K}} E_1+\cdots+\mathfrak T\mathfrak S_{\mathsf{K}} E_p
              =\mathfrak T_{\mathsf{K}} E_1+\cdots+\mathfrak T_{\mathsf{K}} E_p.
\end{aligned}
\]
Chaque idéal \(\mathfrak T_{\mathsf{K}}\cdot E_\nu\) se décompose en \(s\) idéaux simples, car
\[
[\mathfrak T_K E_\nu:K]=[\mathfrak T_K E_\nu:K E_\nu]
=[\mathfrak T_K E_\nu:\mathfrak S_K E_\nu]
\leqq[\mathfrak T_K:\mathfrak S_K]=[\mathfrak T:\mathfrak S]=s
\]
et
\[
\sum [\mathfrak T_{\mathsf{K}} E_\nu:\mathsf{K}]=sp;
\]
donc
\[
[\mathfrak T_{\mathsf{K}} E_\nu:\mathsf{K}]=s,
\qquad
\mathfrak T_{\mathsf{K}} E_\nu=\mathsf K e_1^{(\nu)}+\cdots+\mathsf K e_s^{(\nu)}.
\]

La représentation de \(\mathfrak T\) fournie par \(\mathsf K e_i^{(\nu)}\) prolonge la représentation de \(\mathfrak S\) donnée par \(\mathsf K E_\nu\), car, si \(\alpha\in\mathfrak S\),
\[
\alpha E_\nu=\sigma E_\nu,
\qquad
\sum \alpha e_i^{(\nu)}=\sigma\sum e_i^{(\nu)},
\qquad
\alpha e_i^{(\nu)}=\sigma e_i^{(\nu)}.
\]

Il ne reste plus qu'à démontrer que les représentations fournies par les \(\mathsf K e_i^{(\nu)}\) sont toutes distinctes. Dans le cas commutatif, cela était clair, puisqu'elles appartenaient à des classes de représentations distinctes. Ici, au contraire, elles appartiennent toutes à la même classe de représentations; nous devons donc examiner les représentations elles-mêmes. Considérons à cet effet les applications \(\mathsf{K}\)-homomorphes de tout \(\mathfrak T_k\) fournies par les \(\mathsf K e_i^{(\nu)}\). Elles sont déjà entièrement déterminées dès que l'on connaît les représentations de \(\mathfrak T\), car, pour construire l'application tout entière, il suffit de savoir comment les éléments de base sont représentés. Nous pouvons donc nous borner à montrer que les applications de \(\mathfrak T_K\) ainsi fournies par les \(\mathsf K e_i^{(\nu)}\) sont toutes distinctes. Si nous représentons \(e_i^{(\nu)}\), une première fois au moyen du module de représentation \(\mathsf K e_i^{(\nu)}\), une seconde fois au moyen de \(\mathsf K e_j^{(\nu)}\), nous trouvons, en vertu de
\[
e_i^{(\nu)2}=e_i^{(\nu)},
\]
la première fois \(1\), et, en vertu de
\[
e_i^{(\nu)}e_j^{(\nu)}=0,
\]
la seconde fois \(0\), comme élément représentant \(e_i^{(\nu)}\). Les deux représentations sont donc bien distinctes. (On peut appliquer exactement les mêmes raisonnements dans le cas commutatif, mais ils y sont superflus.)

Désormais, nous concluons exactement comme aux pages 10 et 16. Soit \(\mathsf{H}\) le groupe d'invariance de \(\mathfrak S\), et soit \(\mathfrak T\) le corps gauche des invariants de \(H\). Les éléments de \(H\) sont tous les prolongements à tout \(\mathfrak K\) de l'automorphisme identique de \(\mathfrak S\). Comme le lemme qui vient d'être démontré assure qu'il en existe parmi eux au moins \(s\) qui sont distincts sur \(\mathfrak T\), il faut que \(s=1\), donc que \(\mathfrak T=\mathfrak S\), c.q.f.d.

\section*{Chapitre V. Systèmes de facteurs}\label{fr-r823-chapitre-v-systemes-facteurs}

\subsection*{\NoetherNavAnchor{S22}{§ 22. Le groupe des corps gauches de centre donné}\S{} 22. Le groupe des corps gauches de centre donné}\label{fr-r823-groupe-corps-centre-donne}

Nous prenons pour base de nos recherches un corps commutatif fixe \(\mathsf P\). L'ensemble de tous les anneaux de matrices \(\mathfrak K_r\), dont les corps gauches des endomorphismes \(\mathfrak K\) sont des corps gauches de rang fini sur \(\mathsf P\) ayant \(\mathsf P\) lui-même pour centre, forme un système multiplicativement fermé pour la \og formation du produit direct\fg{}. Le produit direct de deux tels anneaux \(\mathfrak K_r\) et \(\mathfrak L_s\) est simplement l'anneau \(\mathfrak K_r\cdot\mathfrak L_s\), obtenu à partir de \(\mathfrak K_r\) par extension du corps de base à \(\mathfrak L_s\) :
\[
\mathfrak K_r\times\mathfrak L_s=\mathfrak K_r\cdot\mathfrak L_s.
\]
D'après la remarque de la page 22, cette formation du produit direct est commutative :
\[
\mathfrak K_r\times\mathfrak L_s=\mathfrak L_s\times\mathfrak K_r.
\]
Dire que le système de tous les \(\mathfrak K_r\) est multiplicativement fermé signifie donc seulement que \(\mathfrak K_r\times\mathfrak L_s=\mathfrak K_r\mathfrak L_s\) est toujours encore un anneau de matrices dont le corps gauche des endomorphismes \(\mathfrak M\) a pour centre \(\mathsf P\) --- le rang fini relativement à \(\mathsf P\) va de soi. D'après le théorème 1, page 20, \(\mathfrak K_r\times\mathfrak L_s\) est simple bilatère et complètement réductible; c'est donc un anneau de matrices sur un corps gauche d'extension \(\mathfrak M\) de \(\mathsf P\). Que le centre de \(\mathfrak M\) soit égal à \(\mathsf P\) résulte immédiatement de ce que, premièrement, \(\mathsf P\) est assurément contenu dans le centre et, deuxièmement, celui-ci doit être contenu dans le centre \(\mathsf P\) de \(\mathfrak K\) (ou de \(\mathfrak L\)). La multiplication \(\mathfrak K_r\times\mathfrak L_s\) est manifestement associative. Les \(\mathfrak K_r\) forment un système multiplicativement fermé, mais non un groupe : en effet, si \(\mathfrak K_r\times\mathfrak L_s=\mathfrak M_m\), alors assurément \(m\geqq rs\); on ne pourra donc trouver aucun \(\mathfrak L_s\) satisfaisant à l'équation \(\mathfrak K_r\times\mathfrak L_s=\mathfrak M_m\) lorsque \(r>m\).

Réunissons maintenant tous les \(\mathfrak K_r\) ayant le même \(\mathfrak K\) en une classe \(\{\mathfrak K\}\). On a alors ceci : si \(\mathfrak K_r\) et \(\mathfrak K_{r'}\) appartiennent à la même classe \(\{\mathfrak K\}\), et \(\mathfrak L_s\) et \(\mathfrak L_{s'}\) à la même classe \(\{\mathfrak L\}\), alors \(\mathfrak K_r\times\mathfrak L_s\) et \(\mathfrak K_{r'}\times\mathfrak L_{s'}\) appartiennent eux aussi à la même classe. Il suffit de montrer que \(\mathfrak K_r\times\mathfrak L_s\) appartient à la même classe que \(\mathfrak K\times\mathfrak L\), autrement dit que, si \(\mathfrak K\times\mathfrak L=\mathfrak M_m\), alors \(\mathfrak K_r\times\mathfrak L_s\) est également un anneau de matrices sur \(\mathfrak M\). Cela se voit aisément :
\[
\mathfrak K_r\times\mathfrak L
=\left(\sum_1^r\mathfrak K c_{ik}\right)\mathfrak L
=\sum_1^r\mathfrak K_{\mathfrak L}c_{ik}
=\sum_1^r\mathfrak M_m c_{ik}.
\]
C'est un anneau de matrices de degré \(r\) sur \(\mathfrak M_m\), donc naturellement un anneau de matrices de degré \(mr\) sur \(\mathfrak M\). On conclut alors exactement de même :
\[
\mathfrak K_r\times\mathfrak L_s
=\left(\sum_1^s\mathfrak L d_{ik}\right)_{\mathfrak K_r}
=\sum_1^s\mathfrak L_{\mathfrak K_r}d_{ik}
=\sum_1^s\mathfrak M_{mr}d_{ik}
=\mathfrak M_{mrs},
\]
ce qui démontre déjà notre assertion. Retenons que \(\mathfrak L\times\mathfrak K=\mathfrak M_m\) entraîne \(\mathfrak K_r\times\mathfrak L_s=\mathfrak M_{mrs}\).

Il résulte de ce qui vient d'être démontré que, si l'on définit le produit de deux classes \(\{\mathfrak K\}\) et \(\{\mathfrak L\}\) comme la classe du produit d'un représentant de \(\{\mathfrak K\}\) avec un représentant de \(\{\mathfrak L\}\), cette définition est indépendante du choix des représentants.

L'application \(\mathfrak K_r\mapsto\{\mathfrak K_r\}\) est donc un homomorphisme du système de tous les \(\mathfrak K_r\) sur l'ensemble de toutes les classes \(\{\mathfrak K\}\); l'ensemble \(\mathscr K\) des classes \(\{\mathfrak K\}\) est par conséquent un système multiplicativement fermé. Nous allons montrer que \(\mathscr K\) est un groupe. Son élément unité est, comme on le voit immédiatement, la classe \(\{\mathsf P\}\), puisque \(\mathfrak K_r\times\mathsf P=\mathfrak K_r\).

Il ne reste donc qu'à démontrer, pour chaque classe \(\{\mathfrak K\}\), l'existence d'une classe \(\{\mathfrak K\}^{-1}\). On constate qu'il faut poser \(\{\mathfrak K\}^{-1}=\{\overline{\mathfrak K}\}\), où \(\{\overline{\mathfrak K}\}\) désigne le corps gauche anti-isomorphe à \(\mathfrak K\). En d'autres termes, \(\mathfrak K\times\overline{\mathfrak K}\) doit être un anneau de matrices sur \(\mathsf P\); ou encore, si \(\mathfrak K\) a le rang \(t^2\), et donc \(\mathfrak K\times\overline{\mathfrak K}\) le rang \(t^4\) sur \(\mathsf P\), alors \(\mathfrak K\times\overline{\mathfrak K}\) doit se décomposer en \(t^2\) idéaux à gauche simples. Cela résulte de ce qu'un idéal à gauche simple de \(\mathfrak K\times\overline{\mathfrak K}\), considéré comme module de représentation d'une représentation réciproque irréductible de \(\overline{\mathfrak K}\) dans \(\mathfrak K\) --- laquelle est de degré 1 ---, a le rang 1 sur \(\mathfrak K\), donc le rang \(t^2\) sur \(\mathsf P\).

Nous dirons que \(\mathfrak K\) est un \emph{corps gauche simple} s'il n'existe entre \(\mathsf P\) et \(\mathfrak K\) aucun corps gauche dont le centre soit également \(\mathsf P\).

\textbf{Théorème.} \emph{Tout corps gauche \(\mathfrak K\) est le produit direct de sous-corps gauches simples \(\mathfrak K^{(i)}\),}
\[
\mathfrak K=\mathfrak K^{(1)}\times\cdots\times\mathfrak K^{(p)}.
\]

\emph{Démonstration.} Il suffit de montrer que, si \(\mathfrak K^{(1)}\) est un sous-corps gauche de \(\mathfrak K\) ayant \(\mathsf P\) pour centre, il existe un sous-corps gauche \(\mathfrak S\) de \(\mathfrak K\), de centre \(\mathsf P\), qui satisfasse à l'équation
\[
\mathfrak K^{(1)}\times\mathfrak S=\mathfrak K.
\]
Il existe en tout cas une classe \(\{\mathfrak S\}\) ayant la propriété
\[
\{\mathfrak K^{(1)}\}\times\{\mathfrak S\}=\{\mathfrak K\},
\]
à savoir
\[
\{\mathfrak S\}=\{\overline{\mathfrak K^{(1)}}\}\times\{\mathfrak K\}.
\qquad
\text{Soit}\qquad
\overline{\mathfrak K^{(1)}}\times\mathfrak K=\mathfrak S_\lambda.
\]
Alors
\[
\mathfrak K^{(1)}\times\mathfrak S_\lambda
=\mathfrak K^{(1)}\times\overline{\mathfrak K^{(1)}}\times\mathfrak K
=\mathsf P_{t^2}\times\mathfrak K
=\mathfrak K_{t^2},
\qquad
t^2=\text{rang de }\mathfrak K^{(1)}.
\]
Mais \(\mathfrak K^{(1)}\times\mathfrak S_\lambda\) est un anneau de matrices d'un degré au moins égal à \(\lambda\), donc \(\lambda\leqq t^2\). \(\mathfrak S_\lambda=\overline{\mathfrak K^{(1)}}\times\mathfrak K\) doit se décomposer en idéaux simples de rang 1 sur \(\mathfrak K\), parce que \(\overline{\mathfrak K^{(1)}}\) admet des représentations réciproques du premier degré dans \(\mathfrak K\); par suite \(\lambda\geqq t^2\), donc \(\lambda=t^2\), et
\[
\mathfrak K_{t^2}=\mathfrak K^{(1)}\times\mathfrak S_{t^2}.
\]
Ainsi \(\mathfrak K^{(1)}\times\mathfrak S\) doit lui-même être un corps gauche, et
\[
\mathfrak K^{(1)}\times\mathfrak S_{t^2}
=\bigl(\mathfrak K^{(1)}\times\mathfrak S\bigr)_{t^2}
=\mathfrak K_{t^2}.
\]
\(\mathfrak K^{(1)}\times\mathfrak S\) est le corps gauche des endomorphismes des idéaux à gauche simples dans une décomposition de \(\mathfrak K_{t^2}\), tandis que \(\mathfrak K\) est le corps gauche des endomorphismes des idéaux à gauche simples dans une autre décomposition de \(\mathfrak K_{t^2}\). Mais, puisque deux idéaux à gauche distincts quelconques de \(\mathfrak K_{t^2}\) peuvent être placés dans une même décomposition, \(\mathfrak K\) et \(\mathfrak K^{(1)}\times\mathfrak S\) sont isomorphes, et cet isomorphisme peut être choisi de manière à laisser fixe chaque élément de \(\mathsf P\). D'après le théorème 7, page 25, cet isomorphisme est donc engendré par un automorphisme intérieur de \(\mathfrak K_{t^2}\) :
\[
\mathfrak K
=\lambda^{-1}\bigl(\mathfrak K^{(1)}\times\mathfrak S\bigr)\lambda
=\lambda^{-1}\mathfrak K^{(1)}\lambda\times\lambda^{-1}\mathfrak S\lambda,
\qquad
\lambda\in\mathfrak K_{t^2}.
\]
\(\lambda^{-1}\mathfrak K^{(1)}\lambda\) est un sous-corps gauche de \(\mathfrak K\) isomorphe à \(\mathfrak K^{(1)}\); d'après le théorème 7, page 25, il existe donc un \(\mu\) convenable de \(\mathfrak K\) tel que
\[
\mu^{-1}\lambda^{-1}\mathfrak K^{(1)}\lambda\mu=\mathfrak K^{(1)},
\]
et par conséquent
\[
\mu^{-1}\mathfrak K\mu
=\mathfrak K
=\mathfrak K^{(1)}\times\mu^{-1}\lambda^{-1}\mathfrak S\lambda\mu,
\]
ce qui démontre l'assertion.

\subsection*{\NoetherNavAnchor{S23}{§ 23. Systèmes de facteurs}\S{} 23. Systèmes de facteurs}\label{fr-r823-systemes-facteurs}

Dans les recherches qui suivent, nous supposons que le corps de base \(\mathsf P\) est parfait. On peut remplacer cette hypothèse par l'hypothèse plus générale selon laquelle on ne considère que les corps gauches \(\mathfrak K\) de centre \(\mathsf P\) ayant la propriété que chaque \(\mathfrak K_r\) contient un sous-corps commutatif maximal de première espèce sur \(\mathsf P\). (Comme Köthe l'a montré depuis, cette propriété est satisfaite pour tout \(\mathfrak K\).)

Par \(\mathsf{Z}\), nous entendons un corps de décomposition du corps gauche \(\mathfrak K\) de centre \(\mathsf P\). Supposons que la représentation irréductible \(\mathfrak Z\) de \(\mathsf{Z}\) dans \(\mathfrak K\) soit de degré \(r\); elle est alors un sous-corps commutatif maximal de \(\mathfrak K_r\). Soit \(\Gamma\) le corps galoisien de \(\mathsf{Z}\); on a donc
\[
\mathfrak Z_\Gamma=e_1\Gamma+\cdots+e_n\Gamma
\qquad\text{(page 9)},
\]
où \(n\) est le degré de \(\mathsf{Z}\) sur \(\mathsf P\). On obtient alors aussi, pour l'anneau \(\mathfrak K_r\) étendu par \(\Gamma\), une décomposition
\[
\mathfrak K_{r\Gamma}=\mathfrak K_{r\Gamma}e_1+\cdots+\mathfrak K_{r\Gamma}e_n
\]
\[
\mathfrak K_{r\Gamma}=\mathfrak l_1+\cdots+\mathfrak l_n,\qquad
\mathfrak l_i=\mathfrak K_{r\Gamma}e_i.
\]
Les \(\mathfrak l_i\) sont manifestement des idéaux à gauche. Ce sont même des idéaux à gauche simples, car, \(\mathsf{Z}\) étant un corps de décomposition, \(\mathfrak K_{r\Gamma}\) est un anneau de matrices de degré \(r\cdot t=n\) sur \(\Gamma\). (Théorème 4, page 23; \(n=rt\), page 22.) La décomposition ainsi obtenue de \(\mathfrak K_{r\Gamma}\) en idéaux à gauche simples est particulièrement adaptée aux propriétés algébriques de \(\mathfrak K\) et de \(\mathsf{Z}\). En effet, si l'on définit, comme à la page 18, l'action d'une substitution \(S\) du groupe de Galois \(\mathfrak G\) de \(\Gamma/\mathsf P\) sur tous les éléments de \(\mathfrak K_{r\Gamma}\) en posant, pour tout élément \(z\) de \(\mathfrak K_r\),
\[
S(z)=z,
\]
alors les idéaux \(\mathfrak l_i\) sont conjugués : l'application d'un \(S\) à un \(\mathfrak l_j\) donne un autre idéal,
\[
S(\mathfrak l_j)=\mathfrak l_i,
\]
car toute composante indécomposable \(e_j\) de l'unité doit passer à une autre \(e_i\). (Voir aussi page 11.)

Le module de représentation \(e_i\Gamma\) envoie \(\mathfrak Z\) sur un sous-corps \(\mathsf{Z}_i\) de \(\Gamma\). L'élément \(e_i\) appartient déjà à \(\mathfrak Z_{\mathsf{Z}_i}\), puisque la représentation \(\mathsf{Z}_i\) de \(\mathfrak Z\) est déjà contenue dans \(\mathsf{Z}_i\), de sorte que \(\mathfrak Z_{\mathsf{Z}_i}\) doit détacher le module de représentation \(e_i\mathsf{Z}_i\). Soit \(\{\mathsf{Z}_i,\mathsf{Z}_k\}\) le compositum de \(\mathsf{Z}_i\) et de \(\mathsf{Z}_k\). L'anneau \(\mathfrak K_{r\{\mathsf{Z}_i,\mathsf{Z}_k\}}\) détache les deux idéaux à gauche
\[
\bar{\mathfrak l}_i=\mathfrak K_{i\{\mathsf{Z}_i,\mathsf{Z}_k\}}e_i,\qquad
\bar{\mathfrak l}_k=\mathfrak K_{r\{\mathsf{Z}_i,\mathsf{Z}_k\}}e_k,
\]
qui sont simples, puisque \(\bar{\mathfrak l}_i=\mathfrak K_{r\Gamma}\bar{\mathfrak l}_i\). On a donc
\[
\bar{\mathfrak l}_i=\sum_{\lambda=1}^m\{\mathsf{Z}_i,\mathsf{Z}_k\}c_{\lambda i}^{(ik)},\qquad
\bar{\mathfrak l}_k=\sum_{\lambda}\{\mathsf{Z}_i,\mathsf{Z}_k\}c_{\lambda k}^{(ik)}.
\]
\(c_{\lambda\varrho}^{(ik)}\) est un système d'unités matricielles de \(\mathfrak K_{r\{\mathsf{Z}_i,\mathsf{Z}_k\}}\).

Tout isomorphisme \(a_i\mapsto a_k\) de \(\bar{\mathfrak l}_i\) sur \(\bar{\mathfrak l}_k\), ayant \(\mathfrak K_{r\{\mathsf{Z}_i,\mathsf{Z}_k\}}\) pour domaine d'opérateurs, est caractérisé par l'élément \(p_{ik}\) tel que \(e_i\mapsto p_{ik}\), puisque l'on a alors \(a_i=a_ie_i\mapsto a_ip_{ik}\), donc \(\mathfrak l_i\cdot p_{ik}=\mathfrak l_k\). Tous les \(p_{ik}\) possibles sont manifestement de la forme
\[
p_{ik}=c_{ik}^{(ik)}\gamma_{ik},\qquad \gamma_{ik}\ne0\quad\text{dans}\quad \{\mathsf{Z}_i,\mathsf{Z}_k\}.
\]
L'application \(a_i\mapsto a_ip_{ik}\) définit alors aussi un isomorphisme d'opérateurs de \(\mathfrak l_i\) sur \(\mathfrak l_k\). Nous choisissons maintenant certains systèmes de tels \(p_{ik}\) en leur imposant une condition de conjugaison.

Définissons l'action d'un \(S\) sur un indice \(i\) par
\[
S(i)=k,\qquad\text{si}\qquad S(\mathsf{Z}_i)=\mathsf{Z}_k.
\]
Pour toute paire d'indices \((\nu,\mu)\), il existe au moins une paire conjuguée \(S(\nu,\mu)\) de la forme \(S(\nu,\mu)=(1,k)\). Dans chaque classe de paires d'indices conjuguées, on choisit une paire fixe \((1,k)\), puis, pour cette paire \((1,k)\), on pose
\[
p_{1k}=c_{1k}^{(1k)}\gamma_{1k},\qquad \gamma_{1k}\ne0\quad\text{dans}\quad \{\mathsf{Z}_1,\mathsf{Z}_k\},
\]
de façon arbitraire par ailleurs; ensuite, si \(S(1,k)=(\nu,\mu)\), on pose
\[
p_{\nu\mu}=S(p_{1k}).
\]
Alors \(p_{\nu\mu}\) a effectivement la forme
\[
p_{\nu\mu}=c_{\nu\mu}^{(\nu\mu)}\gamma_{\nu\mu},\qquad
\gamma_{\nu\mu}\ne0\quad\text{dans}\quad \{\mathsf{Z}_\nu,\mathsf{Z}_\mu\}.
\]
En effet, l'action de \(S\) sur l'isomorphisme
\[
a_1\mapsto a_1p_{1k},\qquad\text{en particulier}\qquad e_1\mapsto p_{1k}
\quad\text{de }\bar{\mathfrak l}_1\text{ sur }\bar{\mathfrak l}_k
\]
donne un isomorphisme
\[
a_\nu\mapsto a_\nu p_{\nu\mu},\qquad\text{en particulier}\qquad e_\nu\mapsto p_{\nu\mu}
\quad\text{de }\bar{\mathfrak l}_\nu\text{ sur }\bar{\mathfrak l}_\mu.
\]

Si \(c_{\lambda\varrho}\) désigne un système d'unités matricielles de \(\mathfrak K_{r\Gamma}\), il est clair que \(p_{\nu\mu}\) peut aussi s'écrire sous la forme
\[
p_{\nu\mu}=c_{\nu\mu}\bar\gamma_{\nu\mu}.
\]
Les relations
\[
c_{ij}c_{1k}=0,\quad\text{si}\quad j\ne1,
\qquad
c_{ij}c_{jk}=c_{ik}
\]
entraînent donc les relations
\[
\begin{aligned}
p_{ij}p_{1k}&=0,\quad &&\text{si }j\ne1,\\
p_{ij}p_{jk}&=\alpha_{ik}^{(j)}p_{ik},\quad
&&\alpha_{ik}^{(j)}\text{ dans }\Gamma\text{ (plus précisément dans }\{\mathsf{Z}_i,\mathsf{Z}_j,\mathsf{Z}_k\}\text{)}
\end{aligned}
\]
Les \(n^3\) grandeurs \(\alpha\) forment le \emph{système de facteurs de \(\mathfrak K\) relativement au choix de \(\mathsf{Z}\)}, associé au système \(p_{ik}\) de \og pseudo-unités matricielles\fg{}.

Les \(\alpha\) sont conjugués : de \(S(i,k,j)=(\nu,\mu,\tau)\), il résulte que
\[
S\bigl(\alpha_{ik}^{(j)}\bigr)=\alpha_{\nu\mu}^{\tau}.
\]

Tout système \(p_{ik}^*\) de pseudo-unités matricielles s'obtient à partir d'un système \(p_{ik}\) en multipliant les \(p_{ik}\) initiaux par des \(\delta_{ik}\ne0\) arbitraires de \(\{\mathfrak l_1,\mathfrak l_k\}\) :
\[
p_{1k}^*=p_{1k}\delta_{1k},
\]
puis en posant de nouveau, lorsque \(S(1,k)=(\nu,\mu)\),
\[
p_{\nu\mu}^*=S(p_{1k}^*).
\]
On a donc
\[
p_{\nu\mu}^*=p_{\nu\mu}\delta_{\nu\mu},
\]
où \(\delta_{\nu\mu}\) appartient à \(\{\mathsf{Z}_\nu,\mathsf{Z}_\mu\}\). Les \(\delta_{\nu\mu}\) sont conjugués, et de \(S(i,k)=(\nu,\mu)\), il résulte que \(S(\delta_{ik})=\delta_{\nu\mu}\). Si \(\alpha^*\) est le système de facteurs associé à \(p_{ik}^*\), on a
\[
\alpha_{ik}^{(i)*}=\frac{\delta_{ij}\delta_{jk}}{\delta_{ik}}\alpha_{ik}^{(j)}.
\]

Les systèmes de facteurs \(\alpha^*\) et \(\alpha\) sont dits associés (de même \(p_{ik}^*\) et \(p_{ik}\)); tous les systèmes de facteurs de \(\mathfrak K\) relativement au choix de \(\mathsf{Z}\) forment une classe \(\{\alpha\}\) de systèmes de facteurs associés.

La classe \(\{\alpha\}\) ne dépend pas de la représentation \(\mathfrak Z\) de \(\mathsf{Z}\) dans \(\mathfrak K_r\) que l'on prend pour base de sa définition. En effet, le passage d'un \(\mathfrak Z\) à un autre pouvant s'effectuer par transformation au moyen d'un élément \(\lambda\) de \(\mathfrak K_r\), lequel est donc invariant par \(\mathfrak G\), un système \(p_{ik}\) associé à \(\mathfrak Z\) donne un système
\[
p'_{ik}=\lambda^{-1}p_{ik}\lambda,
\]
associé à \(\mathfrak Z'\); mais les systèmes de facteurs de \(p_{ik}\) et de \(p'_{ik}\) sont identiques.

Un système de facteurs n'est rien d'autre qu'une table de multiplication de \(\mathfrak K_{n\Gamma}\), mais une table adaptée aux propriétés algébriques particulières de \(\mathfrak K\).

\subsection*{\NoetherNavAnchor{S24}{§ 24. Multiplication des systèmes de facteurs}\S{} 24. Multiplication des systèmes de facteurs}\label{fr-r823-multiplication-systemes-facteurs}

Au lieu de parler d'un corps de décomposition d'un corps gauche \(\mathfrak K\), nous parlons aussi d'un corps de décomposition de \(\{\mathfrak K\}\). Cela a un sens puisque, si \(\mathfrak K_{\mathsf{Z}}\) se décompose en composantes absolument irréductibles, il en va de même de tout anneau \(\mathfrak K_{r\mathsf{Z}}\), et réciproquement.

Deux corps gauches quelconques \(\mathfrak K\) et \(\mathfrak S\) de centre \(\mathsf P\) ont des corps de décomposition communs, par exemple le compositum d'un corps de décomposition de \(\mathfrak K\) et d'un corps de décomposition de \(\mathfrak S\).

\textbf{Théorème 1.} \emph{Les \(\{\mathfrak K\}\) qui ont un corps \(\mathsf{Z}\) pour corps de décomposition forment un sous-groupe \(\mathscr K_{\mathsf{Z}}\) du groupe \(\mathscr K\) de tous les \(\{\mathfrak K\}\).}

\emph{Démonstration.} Si \(\mathfrak K\) et \(\mathfrak S\) ont \(\mathsf{Z}\) pour corps de décomposition, \(\mathfrak K_{\mathsf{Z}}\) et \(\mathfrak S_{\mathsf{Z}}\), et donc aussi \((\mathfrak K\times\mathfrak S)_{\mathsf{Z}}\), se décomposent en composantes absolument irréductibles; ainsi \(\mathsf{Z}\) est un corps de décomposition de \(\{\mathfrak K\times\mathfrak S\}\) (théorème 4, page 23).

Des corps gauches anti-isomorphes ont tous leurs corps de décomposition en commun.

Considération auxiliaire. Soient
\[
\mathfrak K_r=\mathfrak L_1+\cdots+\mathfrak L_r,\qquad
\mathsf P_m=\mathfrak z_1+\cdots+\mathfrak z_m
\]
des décompositions de \(\mathfrak K_r\) et de \(\mathsf P_m\) en idéaux à gauche simples. Alors
\[
\mathfrak K_{rm}=\mathfrak K_r\times\mathsf P_m
=\mathfrak L_1\mathfrak z_1+\cdots+\mathfrak L_i\mathfrak z_j+\cdots+\mathfrak L_r\mathfrak z_m
\]
est une décomposition de \(\mathfrak K_{rm}\) en \(rm\) idéaux à gauche non nuls, donc simples. Par conséquent, tout idéal à gauche \(\mathsf{A}\) à \(r\) composantes de \(\mathfrak K_{rm}\) est relié, par un automorphisme intérieur de \(\mathfrak K_r\), à l'idéal à gauche à \(r\) composantes
\[
\mathfrak L_1\mathfrak z_1+\cdots+\mathfrak L_r\mathfrak z_1=\mathfrak K_r\mathfrak z_1,
\]
à savoir
\[
\mathsf{A}=\tau^{-1}\mathfrak K_r\tau\cdot\tau^{-1}\mathfrak z_1\tau
=\mathfrak K'_r\cdot\mathfrak z'_1.
\]

Soit \(\mathfrak K_{r\Gamma}=\mathfrak l_1+\cdots+\mathfrak l_n\) une décomposition de \(\mathfrak K_{r\Gamma}\) en idéaux à gauche conjugués, associée à un \(\mathfrak Z\) dans \(\mathfrak K_r\), avec \(\mathfrak l_i=\mathfrak K_re_i\). Puisque
\[
S(\tau^{-1}\mathfrak l_i\tau\cdot\mathfrak z'_1)
=S(\tau^{-1}\mathfrak l_i\tau)\,\mathfrak z'_1
=\tau^{-1}(S\mathfrak l_i)\tau\cdot\mathfrak z'_1
=\tau^{-1}\mathfrak l_k\tau\cdot\mathfrak z'_1,
\]
on obtient, dans
\[
\tag{1}
\mathsf{A}_\Gamma=\mathfrak l'_1\mathfrak z'_1+\cdots+\mathfrak l'_n\mathfrak z'_1,
\]
une décomposition de \(\mathsf{A}_\Gamma\) en idéaux à gauche conjugués. On a
\[
\mathfrak l'_i\mathfrak z'_1=\mathfrak K_{rm\Gamma}\cdot e'_iE',
\]
si l'on pose \(\tau^{-1}e_i\tau=e'_i\) et si \(E'\) désigne l'unité de \(\mathfrak z_1\).

Toute autre décomposition
\[
\tag{2}
\mathsf{A}_\Gamma=\mathfrak q_1+\cdots+\mathfrak q_n,\qquad
S\mathfrak q_i=\mathfrak q_k\quad\text{si}\quad
S\mathfrak l'_i\mathfrak z'_1=\mathfrak l'_k\mathfrak z'_1
\]
en idéaux à gauche conjugués s'obtient à partir de (1) par transformation au moyen d'un élément \(\lambda\) de \(\mathfrak K_{rm}\) :
\[
\mathfrak q_i=\lambda^{-1}\mathfrak l'_i\mathfrak z'_1\lambda,
\qquad
\text{donc, si \(e_i^*\) est l'unité de \(\mathfrak q_i\),}\quad
e'_i=\lambda^{-1}e'_iE'\lambda.
\]

La démonstration résulte immédiatement du lemme suivant.

\textbf{Lemme.} \emph{Soit \(\mathsf{A}\) un idéal à gauche de \(\mathfrak K_n\), et soit}
\[
\mathsf{A}=\mathfrak l_1+\cdots+\mathfrak l_n
=\mathfrak K_{n\Gamma}e_1+\cdots+\mathfrak K_{n\Gamma}e_n
\]
\emph{une décomposition --- dont l'existence est supposée --- en idéaux à gauche conjugués; on peut, sans restreindre la généralité, supposer que les \(l_i\) sont conjugués.}

\emph{Soit \(\mathfrak z\) isomorphe à \(\mathsf{A}\) comme module à opérateurs, et supposons satisfaites les mêmes conditions :}
\[
\mathfrak z=\mathfrak q_1+\cdots+\mathfrak q_n
=\mathfrak K_{n\Gamma}e_1^*+\cdots+\mathfrak K_{n\Gamma}e_n^*.
\]
\emph{(Ainsi les \(\mathfrak q_i\), de même que les \(e_i^*\), sont conjugués.) Alors il existe un élément régulier \(\lambda\) de \(\mathfrak K_n\) tel que \(\mathfrak q_i=\lambda^{-1}\mathfrak l_i\lambda\).}

\emph{Démonstration.} 1. Voici pourquoi on peut supposer les \(e_i\) conjugués. Si \(e_1\) est un idempotent générateur quelconque de \(\mathfrak l_1\), de sorte que \(\mathfrak l_1=\mathfrak K_{n\Gamma}e_1\), l'application d'un automorphisme à \(\Gamma\) donne \(\mathfrak l_i=\mathfrak K_{n\Gamma}e_i\), où \(e_i\) est donc conjugué à \(e_1\) et est de nouveau un idempotent générateur.

2. Tout \(\mathfrak l\) est isomorphe à tout \(\mathfrak q\). En effet, les \(\mathfrak l\), étant conjugués, ont tous la même longueur (absolue); le nombre des \(\mathfrak l\) et celui des \(\mathfrak q\) coïncident, et \(\mathfrak z\) et \(\mathsf{A}\) sont supposés isomorphes.

3. Fixons, par la correspondance
\[
\tag{3}
e_1\mapsto s_1=e_1s_1;\qquad
e_1^*\mapsto t_1=e_1^*t_1;\qquad
s_1t_1=e_1;\qquad
t_1s_1=e_1^*,
\]
un isomorphisme entre \(\mathfrak l_1\) et \(\mathfrak q_1\), ainsi que son inverse. L'application de tous les automorphismes transforme (3) en
\[
e_i\mapsto s_i=e_is_i;\qquad
e_i^*\mapsto t_i=e_i^*t_i;\qquad
s_it_i=e_i;\qquad
t_is_i=e_i^*,
\]
ce qui fixe les isomorphismes \(\mathfrak q_i=\mathfrak l_is_i\) et \(\mathfrak l_i=\mathfrak q_it_i\).

4. Posons maintenant
\[
s=s_1+\cdots+s_n;\qquad t=t_1+\cdots+t_n.
\]
Il vient
\[
\mathfrak q_i=\mathfrak l_is;\qquad
\mathfrak l_i=\mathfrak q_it;\qquad
st=e_1+\cdots+e_n=\xi;\qquad
ts=e_1^*+\cdots+e_u^*=\xi^*.
\]
Ici \(\xi\) et \(\xi^*\) appartiennent à \(\mathfrak K_n\), et l'on a
\[
\mathsf{A}\xi=\mathsf{A},\qquad \mathfrak z\xi^*=\mathfrak z\quad
(\text{puisque }\mathfrak l_i\xi=\mathfrak l_il_i=\mathfrak l_i),
\]
\[
\xi^2=\xi,\qquad \xi^{*2}=\xi^*.
\]
Posons maintenant \(e=\overline E+E=\overline E^*+E^*\), et par suite
\[
\mathfrak K_n=\overline{\mathsf A}+\mathsf{A}
=\mathfrak K_n\overline\xi+\mathfrak K_n\xi
=\overline{\mathfrak z}+\mathfrak z
=\mathfrak K_n\overline\xi^*+\mathfrak K_n\xi^*.
\]

Alors \(\overline{\mathsf A}\) et \(\overline{\mathfrak z}\) sont encore isomorphes. Si donc
\[
\overline E\mapsto\overline s=\overline E\,\overline s
\qquad\text{et}\qquad
\overline E^*\mapsto\overline t=\overline E^*\overline t
\]
sont des isomorphismes réciproques, et si l'on pose
\[
\lambda=s+\overline s,\qquad \mu=t+\overline t,
\]
il vient
\[
\lambda\mu=\mu\lambda=e,\qquad
\mathfrak l_i\lambda=\mathfrak q_i,\qquad
\lambda^{-1}\mathfrak q_i=\mathfrak q_i
\quad
(\lambda^{-1}\mathfrak q_i\subseteqq\mathfrak q_i
=\lambda\cdot\lambda^{-1}\mathfrak q_i\leqq\lambda^{-1}\mathfrak q_i),
\]
et par conséquent
\[
\lambda^{-1}\mathfrak l_i\lambda=\mathfrak q_i,
\qquad\text{ce qu'il fallait démontrer.}
\]

Pour la décomposition (1), on peut de nouveau déterminer des éléments
\[
r_{ik}=c_{ik}^{\prime(ik)}\gamma_{ik},\qquad \gamma_{ik}c\{\mathsf{Z}_i\mathsf{Z}_k\},
\]
qui, par \(e'_iE'\mapsto r_{ik}\), réalisent des isomorphismes d'opérateurs entre les idéaux \(\mathfrak l'_i\mathfrak z'_1\) et \(\mathfrak l'_k\mathfrak z'_1\), et qui satisfont aussi à la condition de conjugaison \(S(r_{ik})=r_{\nu\mu}\) lorsque \(S(i,k)=(\nu,\mu)\). Ils déterminent, par
\[
r_{ij}r_{jk}=\alpha_{ik}^{(j)}r_{ik},
\]
un système de facteurs \(\alpha\) de \(\mathsf{A}_\Gamma\). Il importe maintenant que l'on puisse aussi employer la décomposition arbitraire (2) pour déterminer les systèmes de facteurs de \(\mathsf{A}_\Gamma\). En effet, si l'on remplace un système \(r_{ik}\) associé à (1) par \(r'_{ik}=\lambda^{-1}r_{ik}\lambda\), alors \(e_i^*\mapsto r'_{ik}\) donne un isomorphisme d'opérateurs de \(\mathfrak q_i\) sur \(\mathfrak q_k\); les \(r'_{ik}\) satisfont à la condition de conjugaison, puisque \(\lambda\) est invariant par \(\mathfrak q\), et ils ont le même système de facteurs que les \(r_{ik}\).

Si l'on prend un système de pseudo-unités matricielles \(p_{ik}\) de \(\mathfrak K\), alors
\[
r_{ik}=\tau^{-1}p_{ik}\tau\cdot E'
\]
est un système \(r_{ik}\) de \(\mathsf{A}_\Gamma\) ayant le même système de facteurs. Nous obtenons donc le lemme suivant.

\textbf{Théorème.} \emph{Les systèmes de facteurs de \(\mathsf{A}_\Gamma\) formés à l'aide d'une décomposition quelconque}
\[
\mathsf{A}_\Gamma=\mathfrak q_1+\cdots+\mathfrak q_n
\]
\emph{de \(\mathsf{A}_\Gamma\) en idéaux à gauche conjugués sont les mêmes que ceux de \(\mathfrak K_r\).}

Le théorème suivant montre que, si \(\alpha_{ik}^{(j)}\) et \(\beta_{ik}^{(j)}\) sont des systèmes de facteurs de \(\{\mathfrak K\}\) et de \(\{\mathfrak S\}\), relativement au choix de \(\mathsf{Z}\), alors \(\varepsilon_{ik}^{(j)}=\alpha_{ik}^{(j)}\cdot\beta_{ik}^{(j)}\) est un système de facteurs de \(\{\mathfrak K\times\mathfrak S\}\), relativement au choix de \(\mathsf{Z}\). Il est donc légitime de définir le produit \(\{\alpha\}\times\{\beta\}\) de deux classes \(\{\alpha\}\), \(\{\beta\}\) de systèmes de facteurs associés comme la classe \(\{\alpha_{ik}^{(j)}\beta_{ik}^{(j)}\}\).

\textbf{Théorème 2.} \emph{Supposons que \(\{\mathfrak K\}\) possède, relativement au choix de \(\mathsf{Z}\), la classe \(\{\alpha\}\) de systèmes de facteurs associés. Les \(\{\alpha\}\) de tous les \(\{\mathfrak K\}\) ayant \(\mathsf{Z}\) pour corps de décomposition forment un groupe que l'application \(\{\mathfrak K\}\mapsto\{\alpha\}\) met isomorphiquement en correspondance avec \(\mathscr K_{\mathsf{Z}}\).}

Pour le démontrer, il faut établir deux points.

1. Si \(\{\mathfrak K\}\) a pour système de facteurs \(\alpha_{ik}^{(j)}\), et \(\{\mathfrak S\}\) le système \(\beta_{ik}^{(j)}\), alors \(\varepsilon_{ik}^{(j)}=\alpha_{ik}^{(j)}\beta_{ik}^{(j)}\) est un système de facteurs de \(\{\mathfrak K\times\mathfrak S\}\).

2. Si \(\{\mathfrak K\}\) a pour système de facteurs \(\alpha_{ik}^{(j)}=1\), alors \(\mathfrak K=\mathsf P\).

1. Soient \(r,s\) les degrés des représentations irréductibles de \(\mathsf{Z}\) dans \(\mathfrak K\), \(\mathfrak S\), et soient
\[
\mathfrak K_{r\Gamma}=\mathfrak l_1+\cdots+\mathfrak l_n,\qquad
\mathfrak S_{s\Gamma}=\mathfrak m_1+\cdots+\mathfrak m_n
\]
des décompositions en idéaux conjugués, où \(\mathfrak l_i\) et \(\mathfrak m_i\) sont supposés correspondants (ils sont dérivés de composantes correspondantes de \(\mathfrak Z_\Gamma\) et de \(\mathfrak Z'_\Gamma\)). Alors
\[
(\mathfrak K_r\times\mathfrak S_s)_\Gamma
=\mathfrak l_1\mathfrak m_1+\cdots+\mathfrak l_i\mathfrak m_j+\cdots+\mathfrak l_n\mathfrak m_n
\]
est une décomposition en \(n^2\) idéaux non nuls, donc simples. Ainsi
\[
\mathsf{A}_1=\mathfrak l_1\mathfrak m_1+\mathfrak l_2\mathfrak m_2+\cdots+\mathfrak l_n\mathfrak m_n
\]
est un idéal de longueur absolue \(n\). \(\mathsf{A}_1\) est invariant par tous les \(S\), puisque les \(\mathfrak m_i\) se transforment comme les \(\mathfrak l_i\); il existe donc un idéal \(\mathsf{A}\) de \(\mathfrak K_r\times\mathfrak S_s\) tel que
\[
\mathsf{A}_1=\mathsf{A}_\Gamma
\]
(lemme, page 18).

Si \(u\) est le degré de la représentation irréductible de \(\mathfrak Z\) dans le corps gauche des endomorphismes \(\mathfrak M\) des idéaux simples de \(\mathfrak K\times\mathfrak S\), alors \(\mathsf{A}\) se décompose en \(u\) idéaux simples (puisque sa longueur absolue est \(n\)). Notre lemme montre donc que les systèmes de facteurs de \(\mathsf{A}\) sont également ceux de \(\mathfrak M\).

Or une décomposition de \(\mathsf{A}_\Gamma\) en idéaux conjugués est précisément celle qui a servi à la définition :
\[
\tag{3}
\mathsf{A}_\Gamma=\mathfrak l_1\mathfrak m_1+\mathfrak l_2\mathfrak m_2+\cdots+\mathfrak l_n\mathfrak m_n.
\]
Si \(\alpha_{ik}^{(j)}\) et \(\beta_{ik}^{(j)}\) correspondent aux pseudo-unités matricielles \(p_{ik}\) et \(q_{ik}\) de \(\mathfrak K\) et de \(\mathfrak S\), alors \(r_{ik}=p_{ik}q_{ik}\) est manifestement un système de pseudo-unités matricielles de \(\mathsf{A}_\Gamma\) pour la décomposition (3); le système de facteurs correspondant \(\varepsilon_{ik}^{(j)}=\alpha_{ik}^{(j)}\beta_{ik}^{(j)}\) est donc aussi un système de facteurs de \(\mathfrak M\).

2. Supposons que \(\{\mathfrak K\}\) ait pour système de facteurs \(\alpha_{ik}^{(j)}=1\), avec les pseudo-unités matricielles correspondantes \(p_{ik}\) et la décomposition
\[
\mathfrak K_{r\Gamma}=\mathfrak l_1+\cdots+\mathfrak l_n
\]
en idéaux conjugués. Nous devons démontrer \(\mathfrak K=\mathsf P\), ou \(\mathfrak K_r=\mathsf P_r\), ou encore l'existence d'un idéal à gauche \(\mathfrak l\) à \(n\) composantes de \(\mathfrak K_r\) --- car un idéal à gauche simple de \(\mathfrak K_r\) est de rang \(t^2r\), de sorte qu'il faut avoir \(n=t^2rh\), avec \(h\) entier, ou bien \(t=1\).

L'idéal de \(\mathfrak K_{r\Gamma}\)
\[
\mathfrak L=\mathfrak l_1(p_{11}+\cdots+p_{1n})
\]
est non nul --- \(e_1\sum_i p_{1i}\ne0\) --- et simple, donc de rang \(n\). Il est invariant par \(\mathfrak G\) :
\[
S\mathfrak L
=\mathfrak L\mathfrak l_1\sum_\nu S p_{1\nu}
=\mathfrak l_i\sum_\mu p_{i\mu}
=\mathfrak l_i\sum_\mu p_{i1}p_{1\mu},
\]
parce que \(\alpha_{ik}^{(j)}=1\). Par conséquent
\[
S\mathfrak L=\mathfrak l_ip_{i1}\sum_\mu p_{1\mu}
=\mathfrak l_1\sum_\mu p_{1\mu}
=\mathfrak L.
\]
Ainsi \(\mathfrak L=\mathfrak l_\Gamma\), où \(\mathfrak l\) désigne un idéal de \(\mathfrak K_r\) de rang \(n\).

\textbf{Théorème 3.} \emph{Tout élément \(\{\mathfrak K\}\) du groupe \(\mathscr K\) a un ordre fini \(\lambda\), qui divise le nombre absolu de composantes \(t\) de \(\mathfrak K\).}

\emph{Première démonstration (Brauer).} Considérons deux décompositions de \(\mathfrak K_{r\Gamma}\) en idéaux à gauche conjugués :
\[
\mathfrak K_{r\Gamma}=\mathfrak l_1+\cdots+\mathfrak l_n=\mathfrak z_1+\cdots+\mathfrak z_n,\qquad
\mathfrak l_i=\mathfrak K_{r\Gamma}e_i,\quad
\mathfrak z_i=\mathfrak K_{r\Gamma}e_i^*.
\]
Que \(a_i\mapsto a_ir_{ik}\), en particulier \(e_i\mapsto r_{ik}\), définisse un isomorphisme d'opérateurs de \(\mathfrak l_i\) sur \(\mathfrak z_k\). Dans l'isomorphisme inverse, \(e_k^*\) passe à un \(s_{ki}\) de \(\mathfrak l_i\), avec \(s_{ki}\cdot r_{ik}=e_k^*\). Si nous choisissons les \(r_{ik}\) conjugués,
\[
S(r_{ik})=r_{\nu\mu}\quad\text{si}\quad S(i,k)=(\nu,\mu),
\]
alors les \(s_{ik}\) sont également conjugués :
\[
S(s_{ki})S(r_{ik})=S(s_{ki}r_{ik})=S(e_k^*)=e_\mu^*
=s_{\mu\nu}r_{\nu\mu}=s_{\mu\nu}S(r_{ik}),
\]
donc \(S(s_{ki})=s_{\mu\nu}\). Comme \(e_i\mapsto e_ir_{ij}s_{jk}\) définit un isomorphisme d'opérateurs de \(\mathfrak l_i\) sur \(\mathfrak l_k\), on doit avoir
\[
r_{ij}s_{jk}=p_{ik}\xi_{ik}^{(j)},\qquad
\xi_{ik}^{(j)}\ne0\in\Gamma.
\]

Puisque les \(r\), les \(s\) et les \(p\) sont conjugués, les \(\xi\) le sont aussi : de \(S(i,j,k)=(\nu,\tau,\mu)\), il résulte que
\[
S(\xi_{ik}^{(j)})=\xi_{\nu\mu}^{\tau}.
\]

De
\[
r_{i\lambda}s_{\lambda j}r_{j\lambda}\varepsilon_{\lambda k}
=r_{i\lambda}e_\lambda^*s_{\lambda k}
=r_{i\lambda}s_{\lambda k}
\]
il résulte que
\[
p_{ij}\xi_{ij}^{(\lambda)}\cdot p_{jk}\xi_{ik}^{(\lambda)}
=p_{ik}\xi_{ik}^{(\lambda)},
\]
ou, puisque \(p_{ij}p_{jk}=p_{ik}\alpha_{ik}^{(j)}\),
\[
\alpha_{ik}^{(j)}
=\frac{\xi_{ik}^{(\lambda)}}{\xi_{ij}^{(\lambda)}\xi_{jk}^{(\lambda)}}.
\]
Posons \(\delta_{\nu\mu}=\prod_{\lambda}\xi_{\nu\mu}^{(\lambda)}\). Alors \(\delta_{\nu\mu}\) est invariant par tous les \(S\) qui fixent \(\nu\) et \(\mu\), donc appartient à \(\{\mathsf{Z}_\nu,\mathsf{Z}_\mu\}\); en outre les \(\delta_{\nu\mu}\) sont conjugués, puisque les \(\xi_{\nu\mu}^{(j)}\) le sont. Comme
\[
\alpha_{ik}^{(j)n}
=\frac{\prod_{\lambda}\xi_{ik}^{(\lambda)}}
       {\prod_{\lambda}\xi_{ij}^{(\lambda)}\prod_{\lambda}\xi_{jk}^{(\lambda)}}
=\frac{\delta_{ik}}{\delta_{ij}\delta_{jk}},
\]
la formule des systèmes de facteurs associés (page 32) montre que \(\alpha_{ik}^{(j)n}\) est associé au système \(\beta_{ik}^{(j)}=1\). Puisque l'ordre de \(\{\mathfrak K\}\) est égal à celui de \(\{\alpha\}\), l'ordre de \(\{\mathfrak K\}\) divise donc \(n\). Si l'on choisit en particulier \(\mathfrak Z\) comme sous-corps commutatif maximal de \(\mathfrak K\), on a \(n=t\), et \(\lambda\) divise donc \(t\).

\emph{Deuxième démonstration (Schur).} Cette démonstration repose sur une représentation des \(p_{ik}\) par des matrices régulières \(P_{ik}\) d'ordre \(n\) sur \(\Gamma\), satisfaisant aux mêmes conditions de conjugaison que les \(p_{ik}\). La matrice \(P_{ik}\) appartient à \(\{\mathsf{Z}_i,\mathsf{Z}_k\}\); si \(p_{ik}\) correspond à \(P_{ik}\), alors \(S(p_{ik})\) correspond à \(S(P_{ik})\). Cette représentation est homomorphe : si \(P_{ik}\) correspond à \(p_{ik}\) et \(P_{kj}\) à \(p_{kj}\), alors \(P_{ik}P_{kj}\) correspond à \(p_{ik}p_{kj}\), c'est-à-dire
\[
P_{ik}P_{kj}=\alpha_{ij}^{(k)}P_{ij}
\]
\((P_{ik_1}P_{k_2j}=0\) ne peut se produire lorsque \(k_1\neq k_2\)). Une fois cette représentation des \(p_{ik}\) obtenue, on conclut comme suit. Le passage aux déterminants donne
\[
|P_{ik}|\cdot |P_{kj}|=\alpha_{ij}^{(k)n}\cdot |P_{ij}|.
\]
Comme les matrices \(P_{ik}\) sont régulières, \(|P_{\nu\mu}|\ne0\), et par conséquent
\[
\alpha_{ij}^{(k)n}=\frac{\delta_{ik}\delta_{kj}}{\delta_{ik}},
\]
où l'on a posé \(|P_{\nu\mu}|=\delta_{\nu\mu}\). On poursuit alors comme dans la démonstration de Brauer.

\emph{Représentation des \(p_{ik}\) par les matrices \(P_{ik}\).}

Les idéaux \(\mathfrak l_i\) de la décomposition
\[
\mathfrak K_{r\Gamma}=\mathfrak l_1+\cdots+\mathfrak l_n
\]
sont des modules de représentation réciproque de \(\mathfrak K_r\) dans \(\Gamma\) (théorème 2, page 32). Si l'on fixe une base \(k_1,\ldots,k_n\) de \(\mathfrak K_r\) relativement à \(\mathfrak Z\), alors
\[
\mathfrak K_{r\Gamma}=k_1\mathfrak Z_\Gamma+\cdots+k_n\mathfrak Z_\Gamma,
\]
et donc
\[
\mathfrak l_i=\mathfrak K_{r\Gamma}e_i
=k_1\mathfrak Z_\Gamma e_i+\cdots+k_n\mathfrak Z_\Gamma e_i
=k_1e_i\Gamma+\cdots+k_ne_i\Gamma.
\]
Il en résulte
\[
\mathfrak l_k=\mathfrak l_i p_{ik}
=k_1p_{ik}\Gamma+\cdots+k_np_{ik}\Gamma.
\]
Les deux bases \((k_1e_i,\ldots,k_ne_i)\) et \((k_1p_{ik},\ldots,k_np_{ik})\) de \(\mathfrak l_k\) peuvent être transformées l'une dans l'autre par multiplication par une matrice régulière \(P_{ik}\) :
\[
P_{ik}\cdot
\begin{pmatrix}
k_1e_i\\ \vdots\\ k_ne_i
\end{pmatrix}
=
\begin{pmatrix}
k_1p_{ik}\\ \vdots\\ k_np_{ik}
\end{pmatrix},
\qquad
\text{brièvement}\quad
P_{ik}(k_\varrho e_i)=(k_\varrho p_{ik}).
\]
On en déduit la conjugaison des \(P_{ik}\) :
\[
S(P_{ik})(k_\varrho S(e_k))=(k_\varrho S(p_{ik}))
\quad\text{ou}\quad
S(P_{ik})(k_\varrho e_\mu)=(k_\varrho p_{\nu\mu})
=P_{\nu\mu}(k_\varrho e_\mu),
\]
c'est-à-dire
\[
S(P_{ik})=P_{\nu\mu}\quad\text{si}\quad S(i,k)=(\nu,\mu).
\]
Il s'ensuit à son tour que \(P_{ik}\) appartient à \(\{\mathsf{Z}_i,\mathsf{Z}_k\}\), qui est précisément le corps des invariants des \(S\) fixant \(i\) et \(k\). L'homomorphie se voit tout aussi aisément :
\[
\begin{split}
P_{ik}P_{kj}(k_\varrho e_j)
&=P_{ik}(k_\varrho p_{kj})
=(k_\varrho p_{ik}p_{kj})\\
&=\alpha_{ij}^{(k)}(k_\varrho p_{ij})
=\alpha_{ij}P_{ij}(k_\varrho e_j),
\quad\text{donc}\\
&\cdot\, P_{ik}P_{kj}=\alpha_{ij}^{(k)}P_{ij},
\quad\text{ce qu'il fallait démontrer.}
\end{split}
\]

Le théorème 3 peut être renforcé.

\textbf{Théorème 4.} \emph{Si \(p\) est un facteur premier du nombre absolu de composantes \(t\) de \(\mathfrak K\), alors \(p\) divise aussi l'ordre \(\lambda\) de \(\{\mathfrak K\}\). (Brauer)}

\emph{Démonstration.} Soit \(p^\sigma\) la plus grande puissance de \(p\) contenue dans le degré \(n\) du corps de décomposition \(\mathsf{Z}\), soit \(\mathfrak H\) un sous-groupe d'ordre \(p^\sigma\) de \(\mathfrak G\) (il existe d'après le \emph{théorème de Sylow}, Speiser, 1\textsuperscript{re} éd., \S{} 19, page 42, théorème 43), et soit \(\mathsf{T}\) le corps des invariants de \(\mathfrak H\).

Alors \([\Gamma:\mathsf{T}]=p^\sigma\), tandis que \([\mathsf{T}:\mathsf P]=n/p^\sigma\) n'est pas divisible par \(p\), et a fortiori pas par \(t\). Ainsi \(\mathsf{T}\) n'est pas un corps de décomposition, et le corps gauche des endomorphismes \(\mathfrak S\) de \(\mathfrak K_{\mathsf{T}}=\mathfrak S_m\) est différent de \(\mathsf P\). Mais \(\mathfrak S\) a pour corps de décomposition \(\Gamma\), dont le degré sur le centre \(\mathsf{T}\) de \(\mathfrak S\) est \(p^\sigma\); le nombre absolu de composantes de \(\mathfrak S\) est donc une puissance \(p^\beta>1\). Par le théorème 3, l'ordre de \(\{\mathfrak S\}\) est par conséquent une puissance \(p^\alpha>1\) de \(p\). Or
\[
\{\mathfrak K\}^{\lambda}=\{\mathsf P\},
\qquad
\{\mathfrak K_{\mathsf{T}}\}^{\lambda}=\{\mathsf{T}\}
\]
entraîne \(\lambda\equiv0\pmod {p^\alpha}\), ce qu'il fallait démontrer. (Brauer a donné des exemples montrant que l'on n'a pas nécessairement \(\lambda=t\)\srcfn{1}{Brauer: Untersuchungen über die arithmetischen Eigenschaften von Gruppen linearer Substitutionen II. Math. Zeitschr. 31, S. 733 \S{} 5, (1930).}).

\textbf{Théorème 5.} \emph{Soit \(t\), décomposé en facteurs premiers, \(t=\prod_i p_i^{\varrho_i}\), avec \(\varrho_i\geqq1\) et \(p_i\ne p_j\) si \(i\ne j\). Alors \(\mathfrak K=\mathfrak K^{(1)}\times\mathfrak K^{(2)}\times\cdots\), où \(\mathfrak K^{(i)}\) est un corps gauche de nombre absolu de composantes \(p_i^{\varrho_i}\).}

\emph{Démonstration.} D'après les théorèmes 3 et 4, l'ordre \(\lambda\) de \(\{\mathfrak K\}\) est
\[
\lambda=\prod_i p_i^{\alpha_i},\qquad 1\leqq\alpha_i\leqq\varrho_i.
\]
Le sous-groupe cyclique \(\mathscr \mathsf{Z}\) de \(\mathscr K\) engendré par \(\{\mathfrak K\}\) est le produit direct de sous-groupes \(\mathscr \mathsf{Z}_i\) d'ordres \(p_i^{\alpha_i}\); par conséquent
\[
\{\mathfrak K\}=\prod_i\{\mathfrak K^{(i)}\}.
\]
Il en résulte que tout corps de décomposition \(\mathsf{Z}\) de \(\mathfrak K\) est aussi un corps de décomposition de \(\mathfrak K^{(i)}\). Le nombre absolu de composantes \(p_i^{\sigma_i}\) de \(\mathfrak K^{(i)}\) divise donc \(t\), et \(\sigma_i\leqq\varrho_i\). Soit \(\prod_i \mathfrak K^{(i)}=\mathfrak K_m\); alors
\[
\prod p_i^{\varrho_i}=\prod p_i^{\varrho_i}\cdot m,
\qquad
\sigma_i\geqq\varrho_i,
\]
d'où
\[
\sigma_i=\varrho_i,\qquad m=1,\qquad
\mathfrak K=\prod \mathfrak K^{(j)},
\quad\text{ce qu'il fallait démontrer.}
\]

\subsection*{\NoetherNavAnchor{S25}{§ 25. Représentation normale de K\_r avec sous-corps commutatif maximal galoisien}\S{} 25. Représentation normale de \(\mathfrak K_r\) avec sous-corps commutatif maximal galoisien\protect\footnotemark}\label{fr-r823-representation-normale-sous-corps-galoisien}
\footnotetext{Une justification plus simple est donnée au \S{} 27.}

Supposons galoisien le sous-corps commutatif maximal \(\mathfrak Z\) de \(\mathfrak K_r\). D'après le théorème 7, page 25, les \(n\) automorphismes \(\mathsf{H}_i\) de \(\mathfrak Z/\mathsf P\) sont des automorphismes intérieurs de \(\mathfrak K_r\); il existe donc des éléments \(u_i\) tels que
\[
\mathsf{H}_i z=u_i^{-1}zu_i\quad\text{pour tout }z\text{ de }\mathfrak Z.
\]
Chaque \(u_i\) est déterminé à un facteur non nul \(y\) de \(\mathfrak Z\) près. En effet, si \(u^{-1}zu=u'^{-1}zu'\) pour tout \(z\) de \(\mathfrak Z\), alors \((u^{-1}u')^{-1}z(u^{-1}u')=z\) pour tout \(z\in\mathfrak Z\). L'adjonction de \(u^{-1}u'\) à \(\mathfrak Z\) donne donc un sous-corps commutatif de \(\mathfrak K_r\), lequel, par maximalité de \(\mathfrak Z\), coïncide avec \(\mathfrak Z\). Ainsi \(u^{-1}u'\) appartient à \(\mathfrak Z\).

Nous allons démontrer le théorème suivant :
\[
\mathfrak K_r=u_1\mathfrak Z+\cdots+u_n\mathfrak Z.
\]
Il suffit manifestement d'établir l'indépendance linéaire des \(u_i\) relativement à \(\mathfrak Z\). La démonstration procède par construction explicite des \(u_i\). Soient
\[
\mathfrak Z_{\mathsf{Z}}=e_1\mathsf{Z}+\cdots+e_n\mathsf{Z},\qquad
\mathfrak K_{r\mathsf{Z}}=\mathfrak l_1+\cdots+\mathfrak l_n,\quad
\mathfrak l_i=\mathfrak K_{r\mathsf{Z}}e_i,
\]
et soit \(p_{ik}\) un système de pseudo-unités matricielles associé à cette décomposition. À la page 8, nous avons désigné par \(\Theta_i\) l'isomorphisme de \(\mathfrak Z\) sur \(\mathsf{Z}\) fourni par le module de représentation \(e_i\mathsf{Z}\), de sorte que
\[
S_i=\Theta_i\Theta_1^{-1}
\]
sont les \(n\) automorphismes de \(\mathsf{Z}\) (mais seulement lorsqu'ils sont appliqués à \(\mathsf{Z}\) lui-même, et non plus lorsqu'il s'agit de l'action, toujours sous-entendue ici, des \(S_i\) sur tout \(\mathfrak K_{r\mathsf{Z}}\)!). Pour simplifier, nous écrirons \(\Theta_{S_i}\) à la place de \(\Theta_i\) dans \(S_i=\Theta_i\Theta_1^{-1}\), puis nous remplacerons aussi, dans les \(p_{ik}\), les indices \(i,k\) par les automorphismes correspondants \(S\) :
\[
p_{ik}=p_{S_iS_k}.
\]

Les \(n\) automorphismes de \(\mathfrak Z\) sont (\(E\) désignant la substitution identique)
\[
\mathsf{H}=\Theta_E^{-1}\Theta.
\]
Nous choisissons maintenant la numérotation des \(\mathsf{H}\) de la manière suivante :
\[
\mathsf{H}_S=\Theta_E^{-1}\Theta_{S^{-1}}.
\]

Dans ces conditions, nous pouvons poser
\[
u_T=\sum_S p_{S,ST}=\sum_S S(p_{E,T}).
\]
En effet :

1. \(u_T^{-1}\) existe :
\[
u_T^{-1}=\sum_S p_{ST,S}
\frac{1}{\alpha_{S,S}^{ST}\overline{\gamma}_{S,ST}\overline{\gamma}_{ST,S}}.
\]
(Pour la signification de \(\bar{\gamma}\), voir page 32.)

2. \(u_T\) est un élément de \(\mathfrak K_r\), parce qu'il est invariant par \(\mathfrak G\) :
\[
R(u_T)=\sum_S RS(p_{E,T})=\sum_{S'} S'(p_{E,T})=u_T.
\]

3. \(u_T\) engendre \(\mathsf{H}_T\), c'est-à-dire que, si \(z\in\mathfrak Z\), alors \(u_T^{-1}zu_T=\mathsf{H}_T(z)\), ou encore
\[
z u_T=u_T \mathsf{H}_T(z).
\]

En effet,
\[
z=\sum_S\Theta_S(z)e_S,\qquad
u_T=\sum_S S(e_Ep_{E,T})=\sum_S e_S S(p_{E,T}),
\]
donc
\[
zu_T=\sum_{S,S'}\Theta_S(z)e_Se_{S'}S(p_{E,T})
=\sum_S\Theta_S(z)S(p_{E,T}).
\]
D'autre part,
\[
u_T=\sum_S S(p_{E,T}e_T)=\sum_S S(p_{E,T})e_{ST},
\]
et l'égalité suivante vaut pour l'application aux éléments de \(\mathfrak Z\) :
\[
\Theta_S \mathsf{H}_T=\Theta_S\Theta_E^{-1}\Theta_{T^{-1}}
=S\Theta_{T^{-1}}=ST^{-1}\Theta_E=\Theta_{ST^{-1}}.
\]
Ainsi
\[
\mathsf{H}_T(z)=\sum_S e_S\Theta_S \mathsf{H}_T(z)=\sum_S e_S\Theta_{ST^{-1}}(z),
\]
c'est-à-dire
\[
\begin{aligned}
u_T\mathsf{H}_T(z)
&=\sum_{S,S'} S(p_{E,T})e_{ST}\cdot e_{S'}\Theta_{S'T^{-1}}(z)\\
&=\sum_S S(p_{E,T})\Theta_S(z)
\end{aligned}
\]
et par conséquent
\[
z\cdot u_T=u_T\cdot \mathsf{H}_T(z).
\]

4. Les \(u_S\) sont linéairement indépendants relativement à \(\mathfrak Z\).

De
\[
\sum_T z_Tu_T=0
\]
et de \(z=\sum_S\Theta_S(z)e_S\), il résulte que
\[
\sum_{T,S,L}\Theta_S(z_T)e_Sp_{L,LT}
=\sum_{S,T}\Theta_S(z_T)p_{S,ST}\quad\text{donc}\quad
\Theta_S(z_T)=0,
\]
et par conséquent \(z_T=0\).

Désormais, nous écrirons \(z^S\), \og puissance symbolique\fg{}, à la place de \(\mathsf{H}_S(z)\). Les \(u_S\) doivent satisfaire des relations de la forme
\[
u_Ru_T=u_{RT}a_{R,T},\qquad a_{R,T}\ne0\quad\text{dans }\mathfrak Z.
\]

Comme
\[
\begin{aligned}
u_Ru_T
&=\sum_{S,S'} S(p_{E,R})S'(p_{E,T})
=\sum_S p_{S,SR}p_{SR,SRT}\\
&=\sum_S p_{S,SRT}\alpha_{S,SRT}^{(SR)}
=\sum_S p_{S,SRT}\alpha_{S,SRT}^{(SR)}e_{SRT}
\end{aligned}
\]
et comme
\[
a_{R,\,T} = \sum\limits_{S} e_{SRT}\, \Theta_{SRT} \left(a_{R,\,T}\right),
\]
donc
\[
u_{RT} a_{R,T} = \sum\limits_{S} p_{S,SRT} \Theta_{SRT} (a_{R,T}),
\]
il vient
\[
\Theta_{SRT}(a_{R,\,T}) = \alpha_{S,SRT}^{(SR)}, ~~ a_{R,\,T} = \sum\limits_{S} e_{SRT} \alpha_{S,SRT}^{(SR)} = \sum\limits_{L} e_{L} \alpha_{LT^{-1}R^{-1},L}^{(LT^{-1})}.
\]

Nous appellerons le système des \(a_{R,\,T}\) un \emph{petit système de facteurs} de \(\mathfrak K_r\). Le grand système de facteurs \(\alpha_{BC}^{(A)}\) et le petit système de facteurs \(a_{R,\,T}\) se déterminent réciproquement. Au produit \(\varepsilon_{BC}^{(A)} = \alpha_{BC}^{(A)} \beta_{BC}^{(A)}\) de deux grands systèmes de facteurs correspond le produit \(f_{R,\,T} = a_{R,\,T}b_{R,\,T}\) des petits systèmes de facteurs \(a_{R,\,T}\), \(b_{R,\,T}\) associés à \(\alpha_{BC}^{(A)}\), \(\beta_{BC}^{(A)}\); cela résulte immédiatement des formules précédentes. Au grand système de facteurs unité \(\alpha_{BC}^{(A)} = 1\) correspond le petit système de facteurs unité \(a_{R,\,T} = 1\), et réciproquement.

On peut remplacer \(u_S\) par \(u_S'=u_Sr_S\), où \(r_S\neq 0\) est choisi arbitrairement dans \(\mathfrak Z\). Aux \(u_S'\) correspond un petit système de facteurs \(a_{R,T}'\) :
\[
\text{(1)}\qquad u_R'u_T'=u_{RT}'a_{R,T}'.
\]
Nous dirons que \(a_{R,T}^{\prime}\) est associé à \(a_{R,T}\). On a
\[
u_R r_R u_T r_T = u_{RT} r_{RT} a'_{R,T} = u_R u_T r_R^T r_T = u_{RT} a_{R,T} r_R^T r_T,
\]
d'où
\[
\text{(2)}\qquad a'_{R,T}=a_{R,T}\frac{r_R^T r_T}{r_{RT}}.
\]

À \(u_S' = u_S r_S\) correspond le système suivant de pseudo-unités matricielles :
\[
p_{S,ST}^{\prime}=p_{S,ST}r_{T}.
\]
Le petit système de facteurs \(a'_{R,\,T}\) correspond donc à un grand système de facteurs \(\alpha^{(A)\prime}_{BC}\) associé au système \(\alpha^{(A)}_{BC}\) qui correspond à \(a_{R,\,T}\). Les petits systèmes de facteurs associés correspondent aux grands systèmes de facteurs associés, et réciproquement. La correspondance entre les classes \(\{\alpha\}\) de grands systèmes de facteurs associés et les classes correspondantes \(\{a\}\) de petits systèmes de facteurs associés est un isomorphisme de groupes.

Nous arrivons maintenant à une représentation des \(u_S\) par des matrices d'ordre \(n\) sur \(\mathfrak Z\), représentation qui diffère essentiellement de celles considérées jusqu'ici.

Si l'on multiplie à droite par \(u_S\) la base sur \(\mathfrak Z\)
\[
\text{(I)}\qquad k_{S_1}, \ldots, k_{S_n}
\]
de \(\mathfrak K_r\), on obtient
\[
\text{(II)}\qquad k_{S_1}u_S, \ldots, k_{S_1}u_S.
\]
Nous appellerons \emph{matrice représentant \(u_S\)} la matrice régulière \(A_S\) qui transforme (I) en (II) :
\[
\text{(III)}\qquad (k_{S_1}u_S, \ldots, k_{S_n}u_S) = (k_{S_1}, \ldots, k_{S_n})A_S.
\]
La représentation doit porter sur tous les éléments de \(\mathfrak K_r\) qui, comme les \(u_S\), engendrent un automorphisme de \(\mathfrak Z/\mathsf{P}\), c'est-à-dire sur tous les éléments \(u_Sr_S\), où \(r_S \neq 0\) est arbitraire dans \(\mathfrak Z\). Tous ces éléments forment un groupe \(\mathfrak G^*\), qui contient comme sous-groupe normal le groupe multiplicatif \(\mathfrak Z^*\) de \(\mathfrak Z\). Le groupe quotient \(\mathfrak G^*/\mathfrak Z^*\) est isomorphe au groupe de Galois \(\mathfrak G\) de \(\mathfrak Z/\mathsf{P}\) par la correspondance qui associe à chaque élément d'une classe résiduelle de \(\mathfrak G^*\) modulo \(\mathfrak Z^*\) l'automorphisme de \(\mathfrak Z\) qu'il engendre. Le groupe \(\mathfrak G^*\) est représenté de la manière indiquée ci-dessus :

\(u_S \mapsto A_S\)

Mais cette représentation n'est ni homomorphe, ni même antihomomorphe. De
\(u_S \mapsto A_S\), \(u_T \mapsto A_T\)
cela n'implique pas que \(u_Su_T \mapsto A_SA_T\), mais, puisque
\[
\begin{aligned}
(k_{S_1}, \ldots, k_{S_n})u_Su_T
&=(k_{S_1}, \ldots, k_{S_n})A_Su_T
=(k_{S_1}, \ldots, k_{S_n})u_TA_S^T\\
&=(k_{S_1}, \ldots, k_{S_n})A_TA_S^T:\ u_Su_T\mapsto A_TA_S^T.
\end{aligned}
\]
on obtient la dernière correspondance. Nous appellerons cette représentation de \(\mathfrak G^*\) la \emph{représentation croisée}. Si \(h_{S_1}, \ldots, h_{S_n}\) est une autre base de \(\mathfrak K_r\) relativement à \(\mathfrak Z\), et si
\[
\text{(IV)}\qquad (h_{S_1}, \ldots, h_{S_n}) = (k_{S_1}, \ldots, k_S)M,
\]
alors (III) et (IV) donnent
\[
(h_{S_1},\ldots,h_{S_n})\,u_S=(h_{S_1},\ldots,h_{S_n})\,M^{-1}A_S\,M^S.
\]
Autrement dit, si \(u_S\) est représenté par \(A_S\) dans la base \(k\), il est représenté par \(M^{-1}A_SM^S\) dans la base \(h\), où \(M\) est la matrice qui transforme la base \(k\) en la base \(h\).

La formule (III), qui définit la représentation croisée, montre que celle-ci est à l'anneau \(\mathfrak K_r\) dans une relation analogue à celle d'une représentation réciproque ordinaire --- écrite toutefois à droite plutôt qu'à gauche --- avec son module de représentation. La différence provient de ce qu'ici, à la place de la relation de commutation
\[
(m^*\tau^*)c = (m^*c)\tau^*
\]
(voir page 5) du module de représentation réciproque, intervient la relation
\[
(kz)u_S=(ku_S)z^S\qquad k\text{ dans }\mathfrak K_r,\quad z\text{ dans }\mathfrak Z.
\]

On peut aussi former une représentation croisée de \(\mathfrak G^*\) au moyen d'une base d'un idéal à droite \(\mathfrak r\) de \(\mathfrak K_r\). Tout idéal à droite de \(\mathfrak K_r\) engendre ainsi une classe de représentations équivalentes; les représentations \(u_S \mapsto A_S, u_S \mapsto B_S\) sont dites équivalentes s'il existe une matrice \(M\) telle que \(B_S = M^{-1}A_SM^S\).

Au lieu de nous en tenir aux représentations de \(\mathfrak G^*\) dans \(\mathfrak Z\) engendrées par les idéaux à droite de \(\mathfrak K_r\), nous pouvons maintenant définir en général une représentation croisée de degré \(m\) de \(\mathfrak G^*\) dans \(\mathfrak Z\) comme suit :

\begin{enumerate}
\item À chaque \(u_S\) correspond une matrice \(A_S\) d'ordre \(m\) sur \(\mathfrak Z\).
\item De \(u_S \mapsto A_S\) et \(u_T \mapsto A_T\), il résulte que \(u_Su_T \mapsto A_TA_S^T\).
\end{enumerate}

Prenons un module à droite sur \(\mathfrak Z\) de rang \(m\),
\[
\mathfrak{M}=x_1\mathfrak{Z}+\cdots+x_m\mathfrak{Z},
\]
et faisons-en, au moyen de
\[
(x_1u_S,\ldots,x_mu_S)=(x_1,\ldots,x_m)A_S,
\]
un module à droite sur \(\mathfrak G^*\), avec la relation de commutation
\[
(mz) u_S = (mu_S) z^S,
\]
exactement comme dans la construction, page 5, du module de représentation associé à une représentation donnée. Nous montrons maintenant, comme à la page 20, que \(\mathfrak M\) devient un module à droite fini sur \(\mathfrak K_r\) en posant
\[
m \cdot z u_S = mz \cdot u_S = mu_S \cdot z^S,
\]
et, en général,
\[
m \cdot \sum z_S u_S = \sum m z_S \cdot u_S = \sum m u_S \cdot z_S^S.
\]
Ici, aucun calcul analogue à celui de la page 20 n'est nécessaire, car notre définition s'est servie d'une base déterminée.

Or, d'après M.Z. \S{} 18, tout module simple fini \(\mathfrak M\) sur \(\mathfrak K_r\) est isomorphe, comme module à opérateurs, à un idéal à droite simple de \(\mathfrak K_r\). On conclut donc, comme au théorème 2, page 22, qu'il n'existe qu'une seule classe irréductible de représentations croisées de \(\mathfrak G^*\) dans \(\mathfrak Z\) : celle qui correspond aux idéaux à droite simples.

Le degré de cette représentation irréductible est \(t\), où \(t\) est l'indice (exposant absolu) de \(\mathfrak K\). Cela résulte d'un simple dénombrement des rangs. En effet, ce degré est égal au rang \((\mathfrak r:\mathfrak Z)\), sur \(\mathfrak Z\), d'un idéal à droite simple. Or, page 22, on a
\[
(\mathfrak Z:\mathsf{P})=n=rt;\qquad (\mathfrak K_r:\mathsf{P})\equiv n^2;\qquad (\mathfrak r:\mathsf{P})\equiv n\cdot t,
\]
la dernière égalité résultant de \(\mathfrak K_r = \mathfrak r_1 + \cdots + \mathfrak r_r\).

De
\[
(\mathfrak{r}:P)=(\mathfrak{r}:\mathfrak Z)\cdot(\mathfrak Z:P),
\]
ou \(n\cdot t=(\mathfrak r:\mathfrak Z)\cdot n\), il s'ensuit donc que \((\mathfrak r:\mathfrak Z)=t\).

\subsection*{\NoetherNavAnchor{S26}{§ 26. Multiplication des représentations croisées}\S{} 26. Multiplication des représentations croisées}\label{fr-r823-multiplication-representations-croisees}

Jusqu'ici, nous avons toujours considéré la représentation croisée comme une représentation de \(\mathfrak G^*\). Mais nous pouvons aussi la concevoir comme une représentation du groupe de Galois \(\mathfrak G\) de \(\mathfrak Z/\mathsf{P}\), de telle sorte qu'à tout élément \(S\) de \(\mathfrak G\) soit associée la matrice \(A_S\) qui représente un élément \(u_S\) de \(\mathfrak K_r\) engendrant, par transformation --- \(u_S^{-1}zu_S=z^S\) ---, l'automorphisme \(S\).

Pour caractériser une telle représentation \(S\mapsto A_S\)\srcfn{1}{Une telle représentation (indépendamment des systèmes hypercomplexes)
\[
S\mapsto A_S;\quad T\mapsto A_T;\quad ST\mapsto A_T\mathrm{A}_s^T=A_{ST}\cdot a_{S,T^{-1}}
\]
(ou sa transposée) a servi de point de départ à Speiser (Math. Zschr. 5) pour définir les systèmes de facteurs \(a_{S,T}\). Ces systèmes de facteurs se révèlent donc identiques à ceux qui sont introduits ici à partir d'un autre fondement.}, il faut préciser dans quel \(\mathfrak K_r\) elle est prise, quels \(u_S\) et quels \(k_{S_1},\ldots,k_{S_n}\) de \(\mathfrak K\) sur \(\mathfrak Z\) sont employés. Un changement de base donne une représentation équivalente; un changement des \(u\), des systèmes de facteurs associés \(a_{S,\,T}\). Les objets essentiels sont donc les \(\mathfrak K_r\) dans lesquels la représentation est prise, c'est-à-dire la classe \(\{a_{S,\,T}\}\) de systèmes de facteurs associés.

Soient maintenant données deux représentations de cette sorte de \(\mathfrak G\) :
\[
\begin{array}{lll}
S\mapsto u_S\mapsto A_S & \text{petit système de facteurs correspondant} & a_{R,T},\\
S\mapsto v_S\mapsto B_S & \text{petit système de facteurs correspondant} & b_{R,T}.
\end{array}
\]

Par \emph{produit de Kronecker} \(A\times B\) de deux matrices \(A=(a_{ik})\), \(B=(b_{\mu\nu})\), on entend la matrice \(A\times B=(a_{ik}b_{\mu\nu})\).

\noindent\textbf{Assertion.} \emph{L'application \(S\to A_S\times B_S\) est elle aussi une représentation croisée de \(\mathfrak G\); plus précisément, elle appartient au produit \(\mathfrak K_r\times\mathfrak L_s\) lorsque \(A_S\) provient de \(\mathfrak K_r\) et \(B_S\) de \(\mathfrak L_s\).}

\noindent\emph{Démonstration.} La représentation \(S\mapsto A_S\times B_S\) est une représentation croisée. Soient
\[
S\mapsto A_S\times B_S,\qquad T\mapsto A_T\times B_T,
\]
et
\[
ST\mapsto A_{ST}\times B_{ST}
=(A_TA_S^T\times B_TB_S^T)a_{S,T}^{-1}b_{S,T}^{-1}.
\]
Il faut donc démontrer
\[
(A_T\times B_T)(A_S\times B_S)^T a_{S,T}^{-1}b_{S,T}^{-1}
=(A_TA_S^T\times B_TB_S^T)a_{S,T}^{-1}b_{S,T}^{-1}.
\]
Cela résulte de la considération suivante.

Si l'on représente les matrices \(A\) et \(B\) au moyen des unités matricielles \(c_{ik}\) et \(d_{\mu\nu}\) (qui n'ont rien à voir avec les unités matricielles de \(\mathfrak K_r\) ou de \(\mathfrak L_s\)),
\[
A=\sum c_{ik}\alpha_{ik};\qquad
B=\sum d_{\mu\nu}\beta_{\mu\nu},\qquad\text{on a donc}
\]
\[
A\times B=B\times A
=\sum c_{ik}d_{\mu\nu}\cdot\alpha_{ik}\cdot\beta_{\mu\nu}.
\]
C'est-à-dire qu'il s'agit d'une matrice dans les nouvelles unités matricielles \(c_{ik}d_{\mu\nu}=d_{\mu\nu}c_{ik}\). La multiplication des matrices \(A\) et \(B\) correspond donc bijectivement à la multiplication des éléments dans le produit direct des deux anneaux de matrices :
\[
\sum c_{ik}\mathfrak Z\times\sum d_{\mu\nu}\mathfrak Z.
\]
Par conséquent, toute matrice \(A\) commute avec toute matrice \(B\) (mais les matrices \(A\) ne commutent pas entre elles, ni les matrices \(B\) entre elles); et cette commutation de toute matrice \(A\) avec toute matrice \(B\) donne précisément --- après multiplication par \(a_{S,T}^{-1}b_{S,T}^{-1}\) --- la relation à démontrer.

\section*{Chapitre VI. Théorie des produits croisés}\label{fr-r823-chapitre-vi-theorie-produits-croises}

\subsection*{\NoetherNavAnchor{S27}{§ 27. Représentation des K\_r comme produits croisés}\S{} 27. Représentation des \(\mathfrak K_r\) comme produits croisés}\label{fr-r823-representation-produits-croises}

Nous allons maintenant donner un nouveau fondement, indépendant, à la théorie des produits croisés développée au \S{} 24 à partir des systèmes de facteurs. L'importance de ces produits croisés tient à ce qu'un corps galoisien et son groupe y sont décrits \emph{simultanément} dans \(\mathfrak K_r\).

\noindent\textbf{Théorème 1.} \emph{Si \(\mathfrak Z\) est un corps de décomposition galoisien de \(\mathfrak K_r\), donc un sous-corps commutatif maximal, et si \(\mathfrak G\) est son groupe de Galois, alors \(\mathfrak K_r\) est égal au produit croisé \(\mathfrak Z\cdot\mathfrak G\). Réciproquement, tout produit croisé \(\mathfrak Z\cdot\mathfrak G\) engendre un \(\mathfrak K_r\) de cette sorte.}

\noindent\textbf{Démonstration et définitions. Définition 1.} Soit \(\mathfrak G=\{E,S,\ldots,T\}\) le groupe de Galois de \(\mathfrak Z\). Dire que \(\mathfrak K_r\) est égal au produit croisé signifie
\[
\mathfrak K_r=\mathfrak Zu_E+\cdots+\mathfrak Zu_T,
\]
avec \(zu_S=u_Sz^S\) et \(u_Su_T=a_{S,T}u_{ST}\) (les \(a_{S,T}\), dans \(\mathfrak Z\), forment un petit système de facteurs).

\noindent\textbf{Définition 2.} Le groupe générateur \(\mathfrak G^*\) est défini comme l'ensemble des éléments réguliers \(g\) de \(\mathfrak K_r\) qui, par transformation, envoient \(\mathfrak Z\) sur lui-même dans son ensemble. On a \(g^{-1}\mathfrak Zg=\mathfrak Z\), donc \(g^{-1}zg=z^S\), avec \(S\) dans \(\mathfrak G\). Le groupe \(\mathfrak G\) est ainsi une image homomorphe du groupe \(\mathfrak G^*\), puisque tout automorphisme de \(\mathfrak Z\) est engendré par un \(g\), \(\mathfrak G^*\to\mathfrak G\); de plus,
\[
\mathfrak G\cong\mathfrak G^*/\mathfrak T^*
\quad\text{avec}\quad
\mathfrak T^*\supseteqq\mathfrak Z^*.
\]
On montrera plus loin que \(\mathfrak T^*=\mathfrak Z^*\) --- il n'est pas encore démontré que \(\mathfrak Z\) est également un sous-anneau commutatif maximal; cela peut aussi se démontrer indépendamment.

\noindent\textbf{1.} \emph{Démonstration de la représentation comme produit croisé.} Soient \(u_E,u_S,\ldots,u_T\) des éléments de \(\mathfrak K_r\) qui engendrent les automorphismes \(E,\ldots,T\), donc tels que \(zu_S=u_Sz^S\); en particulier \(u_E=e\), d'après 2. Il faut montrer que la somme
\[
\mathfrak M=\{\mathfrak Zu_E,\ldots,\mathfrak Zu_T\}
\]
est directe et égale à \(\mathfrak Zu_E+\cdots+\mathfrak Zu_T\); elle est alors égale à \(\mathfrak K_r\), puisque les rangs coïncident. Cela se montre comme suit (Schwarz).

\(\mathfrak M\) est un bimodule sur \(\mathfrak Z\), donc un module de représentation de \(\mathfrak Z\) dans \(\mathfrak Z\); et \(\mathfrak M\) engendre les \(n\) représentations distinctes. Comme \(\mathfrak Z\) est commutatif de première espèce, \(\mathfrak M\) est donc au moins de rang \(n\); le reste de la démonstration est exactement le même.

\noindent\textbf{2.} \(\mathfrak G\cong\mathfrak G^*/\mathfrak Z^*\), c'est-à-dire \(\mathfrak T^*=\mathfrak Z^*\) : seuls les éléments de \(\mathfrak Z\) induisent l'identité sur \(\mathfrak Z\) par transformation.

Supposons en effet que \(tz=zt\) pour tout \(z\) de \(\mathfrak Z\), et posons
\[
t=z_0+z_1u_S+\cdots+z_{n-1}u_T.
\]
Alors
\[
zt=zz_0+\cdots+zz_{n-1}u_T,
\]
et
\[
tz=z_0z+\cdots+z_{n-1}u_Tz
=zz_0+z^{S^{-1}}z_1u_S+\cdots+z^{T^{-1}}z_{n-1}u_T.
\]
La somme étant directe, il s'ensuit que
\[
(z-z^{S^{-1}})z_1u_S=0,\ldots,
(z-z^{T^{-1}})z_{n-1}u_T=0,
\]
et, puisque les éléments \(u_S,\ldots\) sont réguliers,
\[
(z-z^{S^{-1}})z_1=0,\ldots,
(z-z^{T^{-1}})z_{n-1}=0,
\qquad\text{pour tout }z.
\]
(On peut aussi conclure du fait que le rang est \(n\).) Or, pour chacun de \(S,\ldots\), tous distincts de \(E\), il existe par définition au moins un \(z\) tel que \(z-z^{S^{-1}}\neq0\), et un \(\bar z\) tel que \(\bar z-\bar z^{T^{-1}}\neq0\). On obtient donc \(z_i=0\), \(i=1,\ldots,(n-1)\), et ainsi \(t=z_0\) appartient à \(\mathfrak Z^*\), ce qu'il fallait démontrer.

\noindent\textbf{3.} Les \(u\) satisfont aux lois de multiplication \(u_Su_T=a_{S,T}u_{ST}\), où les \(a_{S,T}\) satisfont aux lois d'associativité croisées
\[
a_{R,S}\cdot a_{RS,T}=a_{S,T}^{R^{-1}}\cdot a_{R,ST}
\]
(les \(a_{S,T}\) forment un petit système de facteurs).

En effet, d'après 2,
\[
\mathfrak Z^*u_S\cdot\mathfrak Z^*u_T=\mathfrak Z^*u_{ST};
\]
ainsi
\[
u_Su_T=a_{S,T}u_{ST},
\]
avec \(a_{S,T}\) dans \(\mathfrak Z^*\), donc dans \(\mathfrak Z\), et \(a_{S,T}\neq0\). De \(u_R\cdot u_Su_T=u_Ru_S\cdot u_T\), on déduit ensuite
\[
\begin{aligned}
u_Ra_{S,T}u_{ST}
&=a_{S,T}^{R^{-1}}\cdot u_Ru_{ST}
 =a_{S,T}^{R^{-1}}a_{R,ST}u_{RST}\\
&=a_{R,S}u_{RS}\cdot u_T
 =a_{R,S}a_{RS,T}u_{RST},
\end{aligned}
\]
ce qui donne les lois de multiplication indiquées.

\noindent\textbf{4. Réciproque :} produit croisé \(=\mathfrak K_r\). Déclarons que la somme directe
\[
\mathfrak M=\mathfrak Zu_E+\cdots+\mathfrak Zu_T
\]
est un anneau au moyen des relations \(zu_S=u_Sz^S\), \(u_Su_T=a_{S,T}u_{ST}\), où les \(a_{S,T}\) satisfont aux relations d'associativité croisées, de sorte que \(u_Ru_Su_T\) est défini sans ambiguïté :
\[
\begin{array}{ll}
1.&z_1\cdot z_2u_S=z_1z_2\cdot u_S;\\[2pt]
2.&zu_S\cdot u_T=z\cdot u_Su_T
\end{array}
\]
Ainsi \(\mathfrak M\) devient un anneau associatif.

\noindent\textbf{Théorème réciproque 2.} \emph{Le produit croisé \(\mathfrak Z\cdot\mathfrak G\) est simple bilatère et a \(\mathsf{P}\) pour centre (d'après des idées de R. Brauer).} La démonstration repose sur les lemmes suivants.

\noindent\textbf{Lemme 1.} \emph{Tout élément qui commute élément par élément avec \(\mathfrak Z\) appartient à \(\mathfrak Z\).}

\noindent\emph{Démonstration.} Soit \(w=\sum_S c_{\lambda_S}u_S\), où les \(c_{\lambda_S}\) appartiennent à \(\mathfrak Z\), et supposons que \(wz=zw\) pour tout \(z\) de \(\mathfrak Z\). Alors
\[
\sum_S zc_{\lambda_S}u_S
 =\sum_S c_{\lambda_S}u_Sz
 =\sum_S c_{\lambda_S}z^{S^{-1}}u_S.
\]
Ainsi \(zc_{\lambda_S}=c_{\lambda_S}z^{S^{-1}}\), ou \(c_{\lambda_S}(z-z^{S^{-1}})=0\), pour tout \(z\) et tout \(S\), en vertu de l'indépendance linéaire des \(u\). Il s'ensuit que \(c_{\lambda_S}=0\) pour \(S\neq E\), puisqu'il existe pour chaque tel \(S\) au moins un \(z\) tel que \(z-z^{S^{-1}}\neq0\). Ainsi \(w\) appartient à \(\mathfrak Z\).

\noindent\textbf{Lemme 2.} \emph{Tout bimodule \(\mathfrak A\) sur \(\mathfrak Z\) qui est isomorphe à \(\mathfrak Zu_S=u_S\mathfrak Z\) comme bimodule (\(\mathfrak Z\) agissant à gauche et à droite) est identique à \(\mathfrak Zu_S\).}

\noindent\emph{Démonstration de 2.} Le module à gauche \(\mathfrak A\) est isomorphe à \(\mathfrak Zu_S\). Si donc \(u_S\mapsto a\) avec \(a\) dans \(\mathfrak A\), alors \(\mathfrak Zu_S\mapsto\mathfrak Za\), et par suite \(\mathfrak A=\mathfrak Za\). De l'isomorphisme de bimodules découle (1) \(za=az^S\), puisque \(u_Sz\mapsto az\); et il en résulte que \(\mathfrak A=\mathfrak Zu_S\). En effet, si l'on pose \(a=u_Su_S^{-1}a=u_Sw\), alors \(za=zu_Sw=u_Sz^Sw\); mais, d'après (1), \(za=u_Swz^S\). Après multiplication par \(u_S\), on obtient donc \(z^Sw=wz^S\) pour tout \(z\) de \(\mathfrak Z\). Comme \(z^S\) parcourt lui aussi tout \(\mathfrak Z\), \(w\) commute élément par élément avec \(\mathfrak Z\) et appartient donc à \(\mathfrak Z\) d'après 1.

\noindent\textbf{Lemme 3.} \emph{Tout bimodule sur \(\mathfrak Z\) contenu dans \(\mathfrak Z\cdot\mathfrak G\) est de la forme \(\mathfrak Zu_{S_1}+\cdots+\mathfrak Zu_{S_k}\), où \(u_{S_i}\) désigne n'importe quels \(k\) éléments distincts parmi les \(u_S,\ldots,u_T\).}

\noindent\emph{Démonstration.} Le produit croisé \(\mathfrak Z\cdot\mathfrak G\) est complètement réductible comme bimodule sur \(\mathfrak Z\), puisque chaque facteur direct \(\mathfrak Zu_S\) est de rang 1. De plus, comme bimodules, tous les facteurs appartiennent à des classes différentes et ne sont donc pas isomorphes entre eux. Tout bimodule \(\mathfrak M\) sur \(\mathfrak Z\) est ainsi, par complète réductibilité, un facteur direct et est lui-même complètement réductible. On a \(\mathfrak M=\bar z\mathfrak A_{S_i}\), où \(\mathfrak A_{S_i}\) est isomorphe à \(\mathfrak Zu_{S_i}\); d'après 2, il lui est donc identique. (Tous les \(\mathfrak A_{S_i}\) appartiennent à des classes différentes.)

\noindent\textbf{Lemme 4.} \emph{Tout anneau contenu dans \(\mathfrak Z\cdot\mathfrak G\) et contenant \(\mathfrak Z\), en particulier \(\mathfrak Z\cdot\mathfrak G\) lui-même, est simple bilatère. --- Les sous-anneaux qui contiennent \(\mathfrak Z\) correspondent bijectivement aux sous-groupes de \(\mathfrak G\).}

\noindent\emph{Démonstration.} Tout anneau \(\mathfrak o\) contenant \(\mathfrak Z\) est un bimodule sur \(\mathfrak Z\), donc, d'après 3, de la forme \(\sum\mathfrak Z\cdot u_{S_i}\). Puisque \(\mathfrak o\) est un anneau, le produit de deux \(u_{S_i}\) appartient encore à \(\mathfrak o\), de même que le module sur \(\mathfrak Z\) correspondant. Les \(u_{S_i}\) forment donc, aux systèmes de facteurs près, un groupe, puisqu'ils sont en nombre fini. Ainsi \(\mathfrak o=\mathfrak Z\cdot\mathfrak H\), où \(\mathfrak H\) est un sous-groupe de \(\mathfrak G\). (Plus précisément, les \(\mathfrak Z\cdot u_{S_i}\) forment un système multiplicatif.) Soit maintenant \(\mathfrak a\) un idéal bilatère non nul quelconque de \(\mathfrak o\). C'est un bimodule sur \(\mathfrak o\), donc aussi un bimodule sur \(\mathfrak Z\) et un sous-module de \(\mathfrak o\). Avec un \(\mathfrak Zu_{S_i}\), il contient donc tout le groupe \(\mathfrak H\), et devient ainsi \(\mathfrak o=\mathfrak Z\cdot\mathfrak H\).

\noindent\textbf{Lemme 5.} \emph{\(\mathsf{P}\) est le centre de \(\mathfrak Z\cdot\mathfrak G\).}

\noindent\emph{Démonstration.} Si \(w\) commute élément par élément avec \(\mathfrak Z\cdot\mathfrak G\), il commute en particulier avec \(\mathfrak Z\), et appartient donc à \(\mathfrak Z\) d'après 1. Comme \(wu_S=u_Sw^S\), la commutation entraîne alors \(w=w^S\). Ainsi \(w\) appartient à \(\mathsf{P}\).

\noindent\emph{Remarque 1.} La démonstration montre que le théorème demeure valable lorsque \(\mathfrak Z\) est une extension algébrique galoisienne infinie de \(\mathsf{P}\). En effet, les théorèmes de complète réductibilité subsistent (Krull, Zahlringe, Math. Ann. 99). Dans ce cas, \(\mathfrak Z\cdot\mathfrak G\) est le système de toutes les sommes finies \(\sum_{S_\lambda}u_{S_\lambda}\cdot c_{S_\lambda}\).

Seulement, en 4, il faut supposer pour \(\mathfrak o\) que \(\mathfrak o=\mathfrak Z\cdot\mathfrak H\), ce qui est bien le cas du produit complet \(\mathfrak Z\cdot\mathfrak G\). En effet, lorsque \(\mathfrak o\) est infini sur \(\mathfrak Z\), l'appartenance de \(u_{S_i}u_{S_j}\) ne permet plus de conclure que le système possède la propriété de groupe.

\noindent\emph{Remarque 2.} Dans cette construction directe, il faut raccorder ici la théorie de la représentation matricielle croisée développée aux \S\S{} 25/26. On emploiera ci-dessous, à la fin du \S{} 28, le fait que la représentation croisée irréductible est de degré \(t\), où \(t\) est le nombre absolu de composantes (l'indice) de \(\mathfrak K\). Soulignons encore que cette théorie des représentations établit l'identité entre les systèmes de facteurs introduits ici comme constantes de multiplication et la notion usuelle dans la littérature.

\subsection*{\NoetherNavAnchor{S28}{§ 28. Théorème du produit pour les systèmes de facteurs}\S{} 28. Théorème du produit pour les systèmes de facteurs}\label{fr-r823-theoreme-produit-systemes-facteurs}

Soient
\[
\mathfrak K_r=\mathfrak Z\cdot\mathfrak G
=\mathfrak Zu_E+\cdots+\mathfrak Zu_T,
\qquad u_Su_T=u_{ST}a_{S,T},
\]
\[
\overline{\mathfrak K}_r
=\overline{\mathfrak Z}\,\overline{\mathfrak G}
=\overline{\mathfrak Z}\bar u_E+\cdots+\overline{\mathfrak Z}\bar u_T,
\qquad
\bar u_S\bar u_T=\bar u_{ST}\bar b_{S,T}.
\]
Ici \(\mathfrak Z\) et \(\overline{\mathfrak Z}\) sont des corps galoisiens isomorphes, et \(\mathfrak G\), \(\overline{\mathfrak G}\) leurs groupes. On a alors
\[
\mathfrak K_r\times\overline{\mathfrak K}_r
=\widetilde{\mathfrak Z}\,\widetilde{\mathfrak G}\times \mathsf{P}_n
=\mathfrak A_f\times \mathsf{P}_n
=\mathfrak A_j,
\]
avec \(\tilde u_S\tilde u_T=\tilde u_{ST}\tilde c_{S,T}\), où \(\widetilde{\mathfrak Z}\) est isomorphe à \(\mathfrak Z\) et à \(\overline{\mathfrak Z}\). Cela résulte du dénombrement des rangs et du fait que \(\mathfrak Z\) est aussi un corps de décomposition.

\noindent\textbf{Assertion.} \emph{On a \(\tilde c_{S,T}=\tilde a_{S,T}\tilde b_{S,T}\), où \(\tilde a,\tilde b\) sont les éléments de \(\widetilde{\mathfrak Z}\) isomorphes respectivement à \(a\) et à \(\bar b\).}

La démonstration repose sur le fait que l'anneau des endomorphismes de tout idéal à gauche de \(\mathfrak A_j\) de longueur absolue \(n\) est isomorphe à \(\mathfrak A_f=\widetilde{\mathfrak Z}\widetilde{\mathfrak G}\), et que cet anneau des endomorphismes se détermine directement pour un idéal à gauche convenable de \(\mathfrak A_j\). Plus généralement, avec \(s\) au lieu de \(n\), on a le théorème suivant, énoncé pour \(\mathfrak K_r\) par commodité.

\noindent\textbf{1. Théorème 1.} \emph{L'anneau des endomorphismes d'un idéal à gauche \(\mathfrak l\) de longueur absolue \(n\) dans \(\mathfrak K_{rs}=\mathfrak K_r\times \mathsf{P}_s\) est isomorphe à \(\mathfrak K_r\).}

La démonstration résulte des lemmes suivants.

\noindent\textbf{Lemme 1.} \emph{Tous les idéaux à gauche de longueur absolue \(n\) de \(\mathfrak K_{rs}\) sont isomorphes comme modules à opérateurs. Leurs anneaux d'endomorphismes sont donc isomorphes comme anneaux. En particulier, \(\mathfrak l=\mathfrak K_r\times\mathfrak b\) est de longueur absolue \(n\), où \(\mathfrak b\) désigne un idéal simple de \(\mathsf{P}_s\).}

\noindent\emph{Démonstration.} Le fait que les idéaux à gauche aient la longueur absolue \(n\), c'est-à-dire qu'ils se décomposent en \(n\) idéaux simples absolument indécomposables, entraîne qu'ils ont la même longueur dans \(\mathfrak K_{rs}\), autrement dit qu'ils comportent le même nombre de facteurs directs simples dans \(\mathfrak K_{rs}\). Il en résulte qu'ils sont isomorphes comme modules à opérateurs, et donc que leurs anneaux d'endomorphismes sont isomorphes comme anneaux. En outre, \(\mathfrak K_r=\mathfrak Z\cdot\mathfrak G\) est de longueur absolue \(n\), c'est-à-dire de rang \(n^2\); ainsi \(\mathfrak K_r\times \mathsf{P}_s\) est de longueur absolue \(n\cdot s\), et \(\mathfrak K_r\times\mathfrak b\) de longueur absolue \(n\).

\noindent\textbf{Lemme 2.} \emph{Si \(\mathfrak o\) est un anneau unitaire, si \(\mathfrak l=\mathfrak o e_1\) est un idéal à gauche et un facteur direct, donc si \(e_1\) est idempotent, alors l'anneau des endomorphismes de \(\mathfrak l\) est isomorphe à \(e_1\mathfrak o e_1\).}

\noindent\emph{Démonstration.} La multiplication par un élément \(e_1re_1\) définit un endomorphisme d'opérateurs de \(\mathfrak l\). Des éléments distincts \(e_1re_1\) définissent des homomorphismes distincts, car ils associent à \(e_1\) des éléments distincts, et tout homomorphisme est ainsi engendré : \(e_1\mapsto re_1\), avec \(re_1=e_1re_1\) puisque \(e_1e_1=e_1\). Or la somme et le produit dans \(e_1\mathfrak o e_1\) correspondent à la somme et au produit des homomorphismes associés.

\noindent\textbf{Lemme 3.} \emph{Si l'on pose \(\mathsf{P}_s=\sum_1^s c_{ik}\mathsf{P}\), avec les unités matricielles \(c_{ik}\), et \(\mathfrak b=c_{11}\mathsf{P}+\cdots+c_{s1}\mathsf{P}\), alors l'anneau des endomorphismes de \(\mathfrak K_r\times\mathfrak b\) est l'anneau \(\mathfrak K_rc_n\), isomorphe à \(\mathfrak K_r\).}

\noindent\emph{Démonstration.} En effet, \(c_{11}\) commute élément par élément avec \(\mathfrak K_r\), et aucun élément de \(\mathfrak K_r\) n'est annulé, car \(c_{11}\) est un idempotent générateur de \(\mathfrak K_r\times \mathsf{P}_s\); de plus,
\[
c_{11}\cdot\mathfrak K_r\times \mathsf{P}_s\cdot c_{11}
=\mathfrak K_r\times c_{11}\mathsf Pc_{11}
=\mathfrak K_rc_{11}.
\]
Les trois lemmes donnent immédiatement le théorème 1.

\noindent\textbf{2. Détermination et décomposition de \(\mathfrak K_r\times\overline{\mathfrak K}_r\) en idéaux à gauche de longueur absolue \(n\).} Par définition,
\[
\mathfrak A_f=\mathfrak K_r\times\overline{\mathfrak K}_r
=\mathfrak Z\cdot\mathfrak G\times\overline{\mathfrak Z}\cdot\overline{\mathfrak G}.
\]
Les éléments surlignés et non surlignés commutent élément par élément, donc
\[
\mathfrak A_j
=(\mathfrak G\times\overline{\mathfrak G})\cdot(\mathfrak Z\times\overline{\mathfrak Z})
=\mathfrak G\times\overline{\mathfrak G}\cdot(\mathfrak Ze_1+\cdots+\mathfrak Ze_n),
\]
avec \(\mathfrak Ze_i=\overline{\mathfrak Z}e_i\), par multiplication d'un corps galoisien par lui-même. Tous les facteurs directs ayant le même rang relativement à \(\mathfrak Z\), ils ont aussi la même longueur absolue; or celle de \(\mathfrak A_f\) est égale à \(=n^2\). Chaque facteur direct est donc de longueur absolue \(n\). Nous prendrons désormais pour base
\[
\mathfrak l=(\mathfrak G\times\overline{\mathfrak G})\mathfrak Ze_1,
\]
où \(ze_1=e_1\bar z\) fixe un isomorphisme déterminé entre \(\mathfrak Z\) et \(\overline{\mathfrak Z}\).

\noindent\textbf{3. Représentation croisée de \(e_1\mathfrak A_fe_1\), et par suite démonstration du théorème du produit.} --- D'après 1 et 2, les systèmes de facteurs du produit \(\mathfrak K_r\times\overline{\mathfrak K}_r\) sont isomorphes à ceux d'une représentation croisée de \(e_1\mathfrak A_je_1\equiv\mathfrak A\).

\noindent\textbf{Théorème 2.} \emph{On a \(\mathfrak A=\widetilde{\mathfrak Z}\cdot\widetilde{\mathfrak G}\), avec \(\widetilde{\mathfrak Z}=\mathfrak Ze_1=\overline{\mathfrak Z}e_1\), \(\overline{\mathfrak G}=\{-,e_1u_S\bar u_Se_1,-\}\), où}
\[
e_1u_S\bar u_Se_1\cdot e_1u_T\bar u_Te_1
=e_1u_{ST}\bar u_{ST}e_1\cdot a_{S,T}e_1\bar b_{S,T}e_1.
\]
\emph{Ainsi, en posant \(\tilde a_{S,T}=a_{S,T}e_1\) et \(\tilde b_{S,T}=\bar b_{S,T}e_1\), le théorème du produit est démontré; \(\tilde a\) et \(\tilde b\) appartiennent au même corps \(\widetilde{\mathfrak Z}\).}

\noindent\emph{Démonstration.} L'élément \(e_1\) devient l'unité de l'anneau. Les conjugués de \(\tilde z=ze_1\) sont donc donnés par \(\tilde z^Se_1\). Il faut montrer que les éléments \(u_S\bar u_Se_1\) induisent ces substitutions; la représentation croisée sera alors établie.

\noindent\textbf{Lemme 1.} \(e_1u_S\bar u_S=u_S\bar u_Se_1\), donc \(=e_1u_S\bar u_Se_1\).

\noindent\emph{Démonstration.} Posons \(u_S^{-1}e_1u_S=e^S\)\srcfn{1}{Cette convention donne \(e_1u_S=u_Se^S\).}, ou \(e_1=e^E\). Les \(e^S\) parcourent alors les conjugués \(e_1,e_2,\ldots,e_n\). On a (1)\srcfn{2}{Car \(e_1\bar u_R=\bar u_Re^{R^{-1}}\).}
\[
e_1\bar u_R=\bar u_Re^{R^{-1}}.
\]
En effet, \(e_1=\sum a_i\bar A_i=\sum a_i^R\bar A_i^R\), où \(a_i\) parcourt une base quelconque de \(\mathfrak Z\) et \(\bar A_i\) la base complémentaire correspondante de \(\overline{\mathfrak Z}\); la même relation vaut donc pour \(a_i^R\) et \(\bar A_i^R\). Ainsi \(e_1\bar u_R=\bar u_Re^{R^{-1}}\), ce qui démontre (1). Il en résulte
\[
e_1u_S\bar u_S=u_Se^S\bar u_S=u_S\bar u_S\cdot e^{SS^{-1}}=u_S\bar u_Se_1,
\]
ce qui démontre le lemme 1.

\noindent\textbf{Lemme 2.} \(ze_1\cdot e_1u_S\bar u_Se_1=e_1u_S\bar u_Se_1\cdot z^Se_1\).

\noindent\emph{Démonstration.} En effet,
\[
ze_1\cdot e_1u_S\bar u_Se_1
=zu_S\bar u_Se_1\quad(\text{lemme 1})
=u_Sz^S\bar u_Se_1
=u_S\bar u_Se_1\cdot z^Se_1
\equiv b=e_1b,
\]
car \(z\) commute avec \(\bar u\) et avec \(e_1\). Posons \(u'=u_S\bar u_Se_1\). Le lemme 2 affirme que \(u'\) indexe les substitutions de \(\mathfrak Ze_1\), de sorte que la représentation croisée indiquée dans le théorème est bien obtenue.

\noindent\textbf{Lemme 3.} \emph{Le système de facteurs est égal à \(a_{S,T}e_1\cdot\bar b_{S,T}e_1\).}

\noindent\emph{Démonstration.}
\[
\begin{aligned}
e_1u_S\bar u_Se_1\cdot e_1u_T\bar u_Te_1
&=u_S\bar u_S\cdot u_T\bar u_Te_1
 =u_{ST}a_{S,T}\cdot\bar u_{ST}\bar b_{S,T}e_1\\
&=u_{ST}\bar u_{ST}\cdot a_{S,T}\bar b_{S,T}e_1
 =e_1u_{ST}\bar u_{ST}e_1\cdot a_{S,T}e_1\cdot\bar b_{S,T}e_1.
\end{aligned}
\]
Le théorème de 2, et donc le théorème du produit pour les systèmes de facteurs, est ainsi démontré.

\noindent\textbf{4. Théorème du produit pour les classes de systèmes de facteurs associés :}
\[
\{a_{S,T}\}\cdot\{b_{S,T}\}=\{a_{S,T}\cdot b_{S,T}\}.
\]
Cela résulte directement de 3, car le passage à de nouveaux générateurs \(v\) et \(\bar v\) entraîne aussi un tel passage à \(v\bar ve_1\). En effet, \(u=vc\) et \(\bar u=\bar v\bar d\) donnent
\[
u\bar ue_1=vc\cdot\bar v\bar de_1=v\bar v\cdot c\bar de_1=v\bar ve_1\cdot c\bar de_1.
\]

\noindent\textbf{5.} Si \(\{a\}\) est le système de facteurs de \(\{\mathfrak K\}\), et \(t\) l'indice de \(\{\mathfrak K\}\), alors \(\{a\}^t\) est égal à la classe unité. Cela résulte de ce que la représentation croisée irréductible de \(\mathfrak G\) est de degré \(t\), exactement comme dans la démonstration de Schur qui suit la représentation des \(p_{ik}\) (\S{} 24, p. 37).

\noindent\textbf{6.} Isomorphisme entre le groupe \(\mathscr K_{\mathfrak Z}\) des \(\{\mathfrak K\}\) associés à \(\mathfrak Z\), c'est-à-dire ayant \(\mathfrak Z\) pour corps de décomposition, et le groupe des classes \(\{a_{S,T}\}\) de systèmes de facteurs associés. D'après la construction du \S{} 27, il existe une correspondance bijective entre \(\{\mathfrak K\}\) et \(\{a_{S,T}\}\); d'après 4, cette correspondance est homomorphe, et c'est donc un isomorphisme. Il en résulte en particulier que les classes unités se correspondent :

\noindent\emph{Le produit croisé est égal à \(\mathsf{P}\), c'est-à-dire qu'il est un anneau de matrices sur son centre, si et seulement si le système de facteurs est associé au système unité.}

\subsection*{\NoetherNavAnchor{S29}{§ 29. Théorème du genre principal dans le cas minimal}\S{} 29. Théorème du genre principal dans le cas minimal}\label{fr-r823-theoreme-genre-principal-cas-minimal}

Soit \(\mathfrak K_r=\mathfrak Z\cdot\mathfrak G\) un produit croisé, et soit \(\mathfrak G^*\) le groupe de tous les éléments qui, par des automorphismes intérieurs, envoient \(\mathfrak Z\) sur lui-même dans son ensemble (écrit en classes latérales : \(\mathfrak Z^*u_E,\mathfrak Z^*u_S\)).

\noindent\textbf{Définition.} Le \og \emph{genre principal dans le cas minimal}\fg{} \(\mathscr S\) est formé de tous les automorphismes du groupe \(\mathfrak G^*\) qui prolongent l'identité de \(\mathfrak Z^*\). La raison de cette dénomination ressort de la spécialisation aux corps cycliques; voir \S{} 30.

\noindent\textbf{Théorème.} \emph{Tout automorphisme du genre principal est de la forme \(u_S\mapsto u_S\cdot b^Sb^{-1}=u_Sb^{S-1}\), où \(b\) appartient à \(\mathfrak Z\); réciproquement, toute correspondance de cette forme engendre un automorphisme du genre principal.}

\noindent\emph{Démonstration.} Supposons que l'automorphisme de groupe soit donné par \(u_S\mapsto v_S\), \(z\mapsto z\), pour tout \(z\) de \(\mathfrak Z\), et qu'il respecte les produits. On peut alors le prolonger en un automorphisme d'anneau de \(\mathfrak K_r\) par la correspondance
\[
\sum u_Sc_{\lambda_S}\mapsto\sum v_Sc_{\lambda_S},
\]
où les \(c_{\lambda_S}\) appartiennent à \(\mathfrak Z\). Or tout automorphisme de cette sorte est, comme on le sait, intérieur.

Soit \(b\) un élément qui l'engendre; on a donc \(v_S=bu_Sb^{-1}\), \(z=bzb^{-1}\). Comme \(b\) commute avec tous les \(z\), il appartient à \(\mathfrak Z\), et par conséquent
\[
v_S=bu_Sb^{-1}=u_Sb^Sb^{-1}.
\]
Réciproquement, si l'on pose \(v_S=u_Sb^Sb^{-1}\), avec \(b\) dans \(\mathfrak Z\), alors \(v_S=bu_Sb^{-1}\). Il s'agit donc d'un automorphisme de \(\mathfrak K_r\) qui envoie \(\mathfrak Z\) sur lui-même élément par élément et, par définition, \(\mathfrak G^*\) sur lui-même dans son ensemble; c'est donc un automorphisme du genre principal.

\subsection*{\NoetherNavAnchor{S30}{§ 30. Spécialisation aux corps de décomposition cycliques}\S{} 30. Spécialisation aux corps de décomposition cycliques}\label{fr-r823-corps-decomposition-cycliques}

Lorsque l'on forme le produit croisé spécialement avec un corps cyclique \(\mathfrak Z\), on obtient une interprétation de théorèmes connus sur les normes. On a d'abord le résultat suivant.

\noindent\textbf{Théorème 1.} \emph{Le groupe \(\mathscr K_{\mathfrak Z}\) des \(\{\mathfrak K\}\) ayant pour corps de décomposition cyclique fixe \(\mathfrak Z\) sur \(\mathsf{P}\) est isomorphe au groupe \(\mathsf{P}^*/N\mathfrak Z^*\), où \(N\mathfrak Z^*\) est le groupe formé de toutes les normes \(Nz\) des éléments \(z\neq0\) de \(\mathfrak Z\).}

\noindent\emph{Démonstration.} La démonstration repose sur une normalisation de la représentation croisée et sur quelques lemmes à son sujet.

\noindent\textbf{Définition.} La représentation croisée cyclique est dite \emph{normalisée} si elle est de la forme
\[
\mathfrak K_r=\mathfrak Z+\mathfrak Zu+\cdots+\mathfrak Zu^{n-1},
\]
avec \(zu=uz^S\), \(u^n=a\), où \(S\) est un élément générateur de \(\mathfrak G\) (autrement dit, si l'on pose \(u_{S^r}=u^r\), \(r=0,\ldots,n-1\), au lieu de poser plus généralement \(u_{S^r}=b_ru^r\), avec \(b_r\) dans \(\mathfrak Z\)).

\noindent\textbf{Lemme 1.} \emph{La représentation croisée cyclique peut toujours être prise normalisée. Alors \(a\) appartient à \(\mathsf{P}\), et tout \(a\) de \(\mathsf{P}\) donne une telle représentation.}

\noindent\emph{Démonstration.} La première assertion est claire : de \(u^{-1}zu=z^S\), il résulte que \(u^{-r}zu^r=z^{S^r}\), et donc que \(u^r\) engendre la substitution \(z\mapsto z^{S^r}\). De plus, \(u^n\) induisant l'identité \(z\mapsto z^{S^n}\), on a \(u^n=a\), avec \(a\) dans \(\mathfrak Z\). Mais \(u^r\) commute avec tous les \(u^r\); il en va donc de même de \(a\), et par suite \(a^{S^r}=a\). Ainsi \(a\) appartient à \(\mathsf{P}\).

La seule constante de multiplication essentielle est donc \(a\). Les conditions d'associativité deviennent ainsi dépourvues de contenu; tout \(a\) de \(\mathsf{P}\) donne un système associatif et, par conséquent, un \(\mathfrak K_r\).

\noindent\textbf{Lemme 2.} \emph{Au lieu de considérer la classe entière \(\{a_{S,T}\}\) des systèmes de facteurs associés, il suffit de considérer les systèmes de facteurs normalisés contenus dans la classe, c'est-à-dire ceux qui proviennent de représentations normalisées. Nous appellerons cette sous-partie de la classe la classe normalisée. Le groupe des classes entières \(\{a_{S,T}\}\) est isomorphe au groupe des classes normalisées.}

\noindent\emph{Démonstration.} D'après le lemme 1, aucune classe normalisée n'est vide; la correspondance est donc bijective. La loi de composition reste l'ancienne, puisqu'elle est définie au moyen d'un représentant quelconque.

\noindent\textbf{Lemme 3.} \emph{Le groupe des systèmes de facteurs normalisés est isomorphe au groupe multiplicatif \(\mathsf{P}^*\), et le groupe des classes normalisées est isomorphe à \(\mathsf{P}^*/N\mathfrak Z^*\).}

\noindent\emph{Démonstration.} Il résulte du lemme 1 que tout système de facteurs normalisé est de la forme \([1,1,\ldots,a,a,\ldots]\), puisque
\[
u^\nu\cdot u^\mu=u^{\nu+\mu}\quad(\nu+\mu<n),
\qquad
u^\nu\cdot u^\mu=au^\varrho\quad(\tau+\mu=n+\varrho,\ 0\leqq\varrho\leqq n-2).
\]
Or la correspondance
\[
[1,\ldots,a,a,\ldots]\mapsto a,
\qquad
[1,1,\ldots,b,b,\ldots]\mapsto b,
\]
\[
[1,\ldots,ab,\ldots]\mapsto ab,
\]
est un isomorphisme de groupes sur \(\mathsf{P}^*\), puisque, de nouveau d'après le lemme 1, \(a\) parcourt tous les éléments de \(\mathsf{P}^*\). Deux systèmes de facteurs normalisés associés s'obtiennent l'un de l'autre par une transformation \(v=uc\), avec \(c\) dans \(\mathfrak Z\). Or cette transformation donne
\[
v=uc,\qquad v^2=u^2\cdot c^S\cdot c;
\]
\[
v^r=u^r c^{S^{r-1}}c^{S^{r-2}}\ldots c;
\qquad
v^n=u^nN(c).
\]
De \(u^n=a\), il résulte donc que \(v^n=aN(c)\). Une classe normalisée est ainsi formée des systèmes de facteurs \([1,\ldots,aN(c),aN(c),\ldots]\), où \(c\) parcourt tous les éléments non nuls de \(\mathfrak Z\). Cela démontre la seconde partie du lemme.

Les lemmes 2 et 3 donnent le théorème 1 ci-dessus, compte tenu de l'isomorphisme entre les \(\{\mathfrak K\}\) qui y figurent et les classes \(\{a_{S,T}\}\).

\noindent\textbf{Théorème 2.} \emph{Le groupe du genre principal dans le cas minimal est isomorphe au groupe des éléments \(c\) de \(\mathfrak Z\) tels que \(N(c)=1\). Ainsi \(N(c)=1\) entraîne l'existence d'un \(b\neq0\) dans \(\mathfrak Z\) tel que \(c=b^S\cdot b^{-1}=b^{S-1}\) (puissance symbolique \(S-1\)-ième).}

\noindent\emph{Démonstration.} La démonstration du lemme 3 a montré que \(N(c)=1\) équivaut à ce que la transformation \(v=uc\) définisse un automorphisme du genre principal. À des \(c\) distincts correspondent des automorphismes distincts, et réciproquement; au produit correspond le produit.

\noindent\textbf{Remarque.} Le théorème 2 justifie le nom de genre principal. On peut, conformément à son sens, appeler le groupe quotient
\[
\mathfrak Z^*/\mathscr H
\quad(\mathscr H=\text{genre principal})
\cong\mathfrak Z^*/\mathfrak Z^{*S-1}
\]
le groupe des genres dans le cas minimal. Celui-ci est donc isomorphe à \(N(\mathfrak Z^*)\); en combinaison avec le théorème 1, on obtient
\[
\begin{aligned}
\text{groupe des genres dans le cas minimal}
&=\mathfrak Z^*/\mathscr H
 \cong\mathfrak Z^*/\mathfrak Z^{*S-1}
 \cong N(\mathfrak Z^*)\\
&=\text{groupe des }\{\mathfrak K\}\text{ associés à }\mathfrak Z
 =\mathscr K_{\mathfrak Z}
 \cong \mathsf{P}^*/N(\mathfrak Z^*).
\end{aligned}
\]

\subsection*{\NoetherNavAnchor{S31}{§ 31. Applications du cas particulier cyclique}\S{} 31. Applications du cas particulier cyclique}\label{fr-r823-applications-cas-cyclique}

La représentation croisée cyclique fournit de nouvelles démonstrations des théorèmes établis au \S{} 20 : tout corps gauche fini est commutatif, et il n'existe, hormis les quaternions, aucun véritable corps gauche fini sur le corps des nombres réels. Ce sont des démonstrations que l'on pourrait qualifier de démonstrations de la théorie des corps de classes dans le cas minimal.

\noindent\textbf{1. Tout corps gauche fini est commutatif (théorème 10, \S{} 20).}

La démonstration repose sur les deux faits connus suivants relatifs aux corps commutatifs finis (corps de Galois).

\noindent\textbf{1.} \emph{Les corps de Galois sont cycliques sur chacun de leurs sous-corps.}

\noindent\textbf{2.} \emph{Si \(\mathfrak Z\) est un corps de Galois, donc si \(\mathfrak Z\) est cyclique sur \(\mathsf{P}\), tout élément de \(\mathsf{P}\) est la norme d'un élément de \(\mathfrak Z\). On a \(\mathsf{P}^*=N(\mathfrak Z^*)\).}

Soit maintenant \(\mathfrak K\) un corps gauche fini, \(\mathsf{P}\) son centre, et \(\mathfrak Z\) un corps de décomposition cyclique. Alors, d'après la fin du \S{} 30 et le point 1, le groupe \(\mathscr K_{\mathfrak Z}\) des \(\{\mathfrak K\}\) associés à \(\mathfrak Z\) est isomorphe à \(\mathsf{P}^*/N\mathfrak Z^*\); d'après 2, il se réduit donc à l'unité \(\{\mathsf{P}\}\). Cela signifie qu'il n'existe aucun corps gauche différent de \(\mathsf{P}\) ayant \(\mathsf{P}\) pour centre; ainsi \(\mathfrak K=\mathsf{P}\).

\noindent\textbf{2. Les quaternions sont le seul véritable corps gauche fini sur le corps \(\mathsf{P}\) des nombres réels (théorème 9, \S{} 20).}

\noindent\emph{Démonstration.} Il faut que \(\mathsf{P}\) soit le centre, puisque la seule extension commutative de \(\mathsf{P}\) est le corps \(\mathfrak Z\) des nombres complexes, et que, \(\mathfrak Z\) étant algébriquement clos, il n'existe aucun corps gauche fini sur \(\mathfrak Z\) qui ait \(\mathfrak Z\) pour centre. En même temps, seul \(\mathfrak Z\) peut être corps de décomposition, et \(\mathfrak Z\) est cyclique sur \(\mathsf{P}\). De plus, tout nombre positif est une norme; le groupe \(\mathsf{P}^*/N\mathfrak Z^*\) est donc de degré 2. Hormis \(\{\mathsf{P}\}\), il n'existe par conséquent qu'une seule classe \(\{\mathfrak K\}\), celle des quaternions. On obtient en même temps la forme normale en choisissant \(-1\) comme système de facteurs normalisé dans la classe associée. Il vient alors \(\mathfrak K=\mathfrak Z+\mathfrak Zu\), avec \(u^2=-1\). Ainsi \(\mathfrak K=\mathsf P+\mathsf Pi+(\mathsf P+\mathsf Pi)u\), avec \(1^2=-1\), \(u^2=-1\).

\noindent\textbf{Remarque relative à 1.} On peut également démontrer le fait connu 2 à partir de la fin du \S{} 30. On a \(\mathscr S=\mathfrak Z^{*S-1}\cong\mathfrak Z^*/\mathsf{P}^*\), puisque les éléments de \(\mathsf{P}^*\), et eux seuls, engendrent l'automorphisme identique du genre principal, ou encore sont envoyés sur l'unité par l'application de \(S-1\). Lorsque \(\mathfrak Z\) est fini, le nombre d'éléments de \(\mathscr S\) est donc égal à l'indice de \(\mathsf{P}^*\) dans \(\mathfrak Z^*\). Le nombre d'éléments de \(\mathfrak Z^*/\mathscr S\) est par suite égal à celui de \(\mathsf{P}^*\); il en va donc de même de \(N(\mathfrak Z^*)=\mathfrak Z^*/\mathscr S\), et ainsi \(\mathsf{P}^*=N(\mathfrak Z^*)\). (Comparer : nombre des genres = nombre des classes ambiguës. Mais il faut partout ajouter \og dans le cas minimal\fg{} à cette terminologie de théorie des nombres.)

\noindent\textbf{3.} \emph{Si \(\mathsf{P}\) est un corps de base \(p\)-adique, et si \(\mathfrak Z\) est cyclique de degré \(n\) sur \(\mathsf{P}\), alors le groupe \(\mathscr K_{\mathfrak Z}\) des \(\{\mathfrak K\}\) associés à \(\mathfrak Z\) est isomorphe au groupe de Galois de \(\mathfrak Z\).}

D'après la fin du \S{} 30, c'est une conséquence immédiate de la \og théorie locale des corps de classes\fg{}, qui affirme
\[
P^*/N\mathfrak Z^*\cong\mathfrak G.
\]

\noindent\textbf{Remarque relative à 3.} Il en résulte que l'ordre, dans le groupe, des éléments de \(\mathscr K_{\mathfrak Z}\), c'est-à-dire l'exposant de \(\{\mathfrak K\}\), est égal à leur indice. En effet, l'exposant d'un élément générateur de \(\mathscr K_{\mathfrak Z}\) est égal à \(n\), et il en va donc de même de l'indice, qui divise \(n\) et est multiple de l'exposant. Le groupe \(\mathscr K_{\mathfrak Z}\) ne contient ainsi que des classes dont l'indice divise \(n\) et pour lesquelles \(\mathfrak Z\) est un corps de décomposition. H. Hasse a montré en outre, Math. Ann. 104, que \(\mathscr K_{\mathfrak Z}\) est formé de toutes les classes \(\{\mathfrak K\}\) dont l'indice divise \(n\). Tout corps de décomposition cyclique \(\mathfrak Z\) de degré \(n\) sur un corps \(p\)-adique \(\mathsf{P}\) conduit donc au même groupe \(\mathscr K_{\mathfrak Z}\) de classes de corps gauches \(\{\mathfrak K\}\).

Les deux faits suivants :

\noindent\textbf{1.} le groupe de tous les \(\{\mathfrak K\}\) sur \(\mathsf{P}\) dont l'indice divise \(n\) est cyclique de degré \(n\);

\noindent\textbf{2.} tout corps de décomposition cyclique de degré \(n\) sur \(\mathsf{P}\) engendre le groupe entier des \(\{\mathfrak K\}\),

peuvent être considérés comme le contenu essentiel du théorème réciproque de la théorie locale des corps de classes.

\clearpage

\clearpage
\subsection*{\pdfbookmark[0]{Complément, bibliographie et listes}{noether.group.KAPFERER}\NoetherNavAnchor{KAPFERER}{Conditions de multiplicité nécessaires et suffisantes pour le théorème fondamental de Noether sur les fonctions algébriques — par H. Kapferer}Conditions de multiplicité nécessaires et suffisantes pour le théorème fondamental de Noether sur les fonctions algébriques : par H. Kapferer}
\label{fr-r823-conditions-multiplicite-kapferer-noether}

\begin{center}
(Avec un complément, en commun avec E. Noether.)

\emph{Math. Ann.} 97 (1927), p. 559--567
\end{center}

Ce qui suit constitue un renforcement du théorème fondamental de Noether, sous la forme de conditions de multiplicité nécessaires et suffisantes\srcfn{1)}{Pour un exposé détaillé, je renvoie à mon travail de même titre, paru dans le numéro spécial des \emph{Heidelberger Akademie-Berichte} de 1927, \og Beiträge zur Algebra\fg{}, préparé à l'initiative d'A. Loewy, à Fribourg-en-Brisgau, et qui contient des travaux de différents auteurs.}. Je réunis l'ensemble des faits qui constituent le théorème fondamental de Noether --- en suivant la conception du théorème fondée sur la théorie des idéaux\srcfn{2)}{Voir, par exemple, B. L. van der Waerden, \og Zur Nullstellentheorie der Polynomideale\fg{}, \emph{Math. Annalen} 96 (1926), p. 203.}, mais sans toutefois faire usage, ici ni dans la suite, de notions ou de théorèmes de la théorie des idéaux --- dans les deux énoncés distincts suivants\srcfn{3)}{Au \S{} 1 du mémoire cité, j'ai démontré séparément et directement les deux énoncés, c'est-à-dire sans recourir à la théorie des idéaux.}.

\textbf{Théorème Ia.} Pour tout zéro commun $x=\alpha$, $y=\beta$ de deux polynômes premiers entre eux $\varphi,\psi$ en $x,y$ --- à coefficients complexes arbitraires ---, il existe un plus petit entier rationnel positif $\varrho$ tel que le module $(\varphi,\psi,\mathfrak{p}^\varrho)$ désigne exactement le même ensemble de polynômes que le module $(\varphi,\psi,\mathfrak{p}^\sigma)$, où $\mathfrak{p}=(x-\alpha,y-\beta)$ et où $\sigma$ désigne un entier rationnel arbitraire supérieur à $\varrho$\srcfn{4)}{Les lettres $\mathfrak{p}$ et $\mathfrak{q}$ ont été choisies pour rappeler l'accord de contenu entre les objets ainsi désignés et les notions d'idéal premier $\mathfrak{p}$ et d'idéal primaire $\mathfrak{q}$, usuelles dans la littérature récente.}.

\textbf{Théorème Ib.} Si l'on désigne par $\mathfrak{q}$ cet ensemble de polynômes, déterminé de manière conceptuellement univoque, ou par $\mathfrak{q}_i$ lorsqu'il s'agit du $i$-ième des $s$ zéros distincts, $x=\alpha_i$, $y=\beta_i$, alors :

La condition nécessaire et suffisante pour qu'un polynôme donné $K$ en $x,y$ possède la propriété de divisibilité $K\equiv0(\varphi,\psi)$ est que les $s$ congruences suivantes soient simultanément satisfaites :
\[
K\equiv0(\mathfrak{q}_1),\quad K\equiv0(\mathfrak{q}_2),\quad\ldots,\quad K\equiv0(\mathfrak{q}_s).
\]

Historiquement --- que l'on compare l'introduction du mémoire fondamental de Max Noether sur ce sujet\srcfn{5)}{Max Noether, \og Über einen Satz aus der Theorie der algebraischen Funktionen\fg{}, \emph{Math. Annalen} 6 (1872).} ---, le théorème fondamental est né des efforts consacrés à la résolution de la tâche fondamentale suivante : \emph{pour tout polynôme donné $K$ en $x,y$, on doit pouvoir décider s'il est divisible ou non par $(\varphi,\psi)$}, c'est-à-dire aussi s'il existe deux autres polynômes $\lambda$ et $\mu$ en $x,y$ tels que l'équation
\[
K=\lambda\cdot\varphi+\mu\cdot\psi
\]
soit une identité en $x,y$.

Cette tâche n'est pas résolue par le théorème de Noether; celui-ci doit plutôt être compris comme une \emph{réduction}, certes fondamentale\srcfn{6)}{La réduction en question peut se comparer à celle que chacun connaît en théorie élémentaire des nombres : la divisibilité d'un entier naturel $A$ par un entier naturel $B$ équivaut à la divisibilité de $A$ par $p_1^{a_1}$, par $p_2^{a_2}$, $\ldots$, par $p_s^{a_s}$, où $p_1,p_2,\ldots,p_s$ désignent les différents diviseurs premiers de $B=p_1^{a_1}p_2^{a_2}\ldots p_s^{a_s}$. C'est précisément la notion d'\og idéal primaire\fg{} --- désignée dans le texte par $\mathfrak{q}$ --- qu'il faut concevoir comme analogue d'une puissance de nombre premier. [Voir E. Noether, \og Idealtheorie in Ringbereichen\fg{}, \emph{Math. Annalen} 83 (1921), introduction.]}, de la tâche posée à une tâche plus simple, à savoir la recherche des conditions nécessaires et suffisantes pour que chacune des congruences $K\equiv0(\mathfrak{q})$ soit satisfaite. Quand a-t-on donc $K\equiv0(\mathfrak{q})$? Ce n'est que lorsque cette dernière question est résolue que le problème initial de M. Noether l'est également. Tel est le but du présent mémoire; ce but est effectivement atteint par l'indication d'un nombre fini de \emph{conditions de multiplicité, nécessaires et suffisantes pour $K\equiv0(\mathfrak{q})$, et dont la vérification ne présente pas non plus de difficulté pratique}. Une condition suffisante pour $K\equiv0(\mathfrak{q})$, mais qui n'est pas également nécessaire, est connue et souvent employée et citée dans la littérature, par exemple sous la forme suivante :

Pour que la congruence $K\equiv0(\varphi,\psi)$ soit satisfaite, il \emph{suffit} que chaque point d'intersection de $\varphi=0$, $\psi=0$ apparaisse dans le polynôme $K$ avec une multiplicité \og suffisamment\fg{} élevée.

Cet énoncé est contenu comme cas particulier dans le théorème I; en effet, qu'un point soit de multiplicité $\varrho$ dans $K$ équivaut à $K\equiv0(\mathfrak{p}^\varrho)$; on a alors, a fortiori, $K\equiv0(\varphi,\psi,\mathfrak{p}^\varrho)$, c'est-à-dire aussi $K\equiv0(\mathfrak{q})$. Mais la réciproque n'est pas vraie.

La solution annoncée du problème peut se caractériser brièvement comme suit. À chaque $\mathfrak{q}$ on peut associer un système de $t$ entiers rationnels \emph{positifs} $(g_1),(g_2),\ldots,(g_t)$, et au module $(K,\mathfrak{q})$ un même nombre d'entiers rationnels \emph{positifs ou négatifs} $(K_1),(K_2),\ldots,(K_t)$, de telle sorte que soit valable le \emph{théorème principal} suivant :

\textbf{Théorème II.} \emph{La condition nécessaire et suffisante pour $K\equiv0(\mathfrak{q})$ est que les $t$ conditions}
\[
(K_1)\geqq(g_1),\quad (K_2)\geqq(g_2),\quad \ldots,\quad (K_t)\geqq(g_t)
\]
\emph{soient simultanément satisfaites. On a en outre}
\[
(g_1)+(g_2)+\ldots+(g_t)=\mu,
\]
\emph{où $\mu$ désigne la multiplicité d'intersection correspondant à $\mathfrak{q}$ et définie par le résultant.}

Pour la démonstration, supposons --- sans restreindre la généralité --- que $x=0$, $y=0$ soit le zéro de $(\varphi,\psi)$ correspondant à $\mathfrak{q}$. En outre, si $A(x,y)$ est un polynôme quelconque\srcfn{7)}{Si l'on admet aussi pour $A(x,y)$ des fonctions rationnelles fractionnaires $C(x,y):D(x,y)$, on pose $(A)=(C)-(D)$; dans ce dernier cas, $(A)$ peut donc aussi prendre des valeurs négatives.} en $x,y$, le symbole $(A)$ désignera le nombre qui indique combien de fois $y$ intervient comme facteur dans $A(0,y)$; si toutefois $A(0,y)$ s'annule identiquement, on posera $(A)=\infty$. Les nombres $(g_1),(g_2),\ldots,(g_t)$ associés à $\mathfrak{q}$ sont définis par le système d'équations suivant, qui constitue la base méthodique du procédé :
\begin{equation}
\tag{1}
\begin{aligned}
f_1\cdot g_1(0,y):y^{(g_1)} - g_1\cdot f_1(0,y):y^{(g_1)} &= x\cdot h_1\\
f_2\cdot g_2(0,y):y^{(g_2)} - g_2\cdot f_2(0,y):y^{(g_2)} &= x\cdot h_2\\
&\vdots \qquad\qquad\vdots\qquad\qquad\vdots\\
f_r\cdot g_r(0,y):y^{(g_r)} - g_r\cdot f_r(0,y):y^{(g_r)} &= x\cdot h_r
\end{aligned}
\end{equation}

Ce système se construit à partir du seul couple de polynômes donné $\varphi,\psi$; les polynômes $f_1,g_1$ sont en effet définis comme identiques à $\varphi,\psi$, à l'ordre près éventuellement; cet ordre doit être choisi de telle sorte que $(f_1)\geqq(g_1)$. Par cette convention, et seulement alors, $h_1$ est rationnel entier.

Les polynômes $f_2,g_2$ sont définis comme identiques à $h_1,g_1$, sous la condition auxiliaire $(f_2)\geqq(g_2)$. En général, le couple de polynômes $f_r,g_r$ est défini comme identique à $h_{r-1},g_{r-1}$, sous la condition auxiliaire $(f_r)\geqq(g_r)$.

Le système d'équations (1) peut certes être poursuivi indéfiniment de la manière indiquée; mais le théorème principal n'exige que la formation des $t$ premières équations du système, où $t$ est déterminé par le fait suivant :

\textbf{Lemme auxiliaire I.} Il existe un entier naturel $t$ [et même $t\leqq\mu$, où $\mu$ désigne la multiplicité avec laquelle le point $x=0$, $y=0$ intervient comme point d'intersection de $\varphi=0$, $\psi=0$] tel que les $t$ premiers nombres définis par (1), à savoir $(g_1),(g_2),\ldots,(g_t)$, soient tous \emph{strictement positifs}, tandis que $(g_{t+1})$ et tous les suivants sont simultanément \emph{nuls}.

En effet, compte tenu du théorème du produit pour les résultants, la multiplicité d'intersection au point $x=0$, $y=0$ de $f_r=0$, $g_r=0$ est égale à
\[
\mu-(g_1)-(g_2)-\ldots-(g_r),
\]
d'où résulte l'\emph{arrêt après au plus $\mu$ étapes}. Il en résulte en même temps l'égalité
\begin{equation}
\tag{2}
\mu=(g_1)+(g_2)+\ldots+(g_t),
\end{equation}
formule remarquable parce qu'elle fournit, tout à fait généralement, une méthode rationnelle et d'exécution pratique facile pour déterminer la multiplicité d'intersection, dès que le zéro considéré est connu.

Une fois connus les nombres $(g_i)$, les nombres $(K_i)$ sont déterminés par le système d'équations suivant, correspondant à (1) :
\begin{equation}
\tag{3}
\begin{aligned}
K_1\cdot g_1(0,y):y^{(g_1)} - g_1\cdot K_1(0,y):y^{(g_1)} &= x\cdot K_1\\
K_2\cdot g_2(0,y):y^{(g_2)} - g_2\cdot K_2(0,y):y^{(g_2)} &= x\cdot K_2\\
&\vdots \qquad\qquad\vdots\qquad\qquad\vdots\\
K_r\cdot g_r(0,y):y^{(g_r)} - g_r\cdot K_r(0,y):y^{(g_r)} &= x\cdot K_{r+1}
\end{aligned}
\end{equation}

$K_1$ est défini comme identique à $K$. Il est manifeste que $K_2$ est rationnel entier si et seulement si $(K_1)\geqq(g_1)$. Si cette dernière condition est satisfaite, il s'ensuit ensuite que $K_3$ est rationnel entier si et seulement si $(K_2)\geqq(g_2)$. En général --- si $K_r$ est déjà rationnel entier --- $K_{r+1}$ est rationnel entier si et seulement si $(K_r)\geqq(g_r)$.

La démonstration du théorème principal (théorème II) repose maintenant sur le second lemme suivant.

\textbf{Lemme auxiliaire II.} Si l'on pose $(f_i,g_i,\mathfrak{p}^\varrho)=q^{(i)}$, alors, pour chaque $i$ de la suite infinie des entiers $1,2,\ldots$, on a :

Pour la congruence
\[
K_i\equiv0(\mathfrak{q}^{(i)}),
\]
\emph{il est nécessaire et suffisant que les deux conditions}
\[
(K_i)\geqq(g_i)\quad \text{et}\quad K_{i+1}\equiv0(\mathfrak{q}^{(i+1)})
\]
\emph{soient simultanément satisfaites.}

La démonstration\srcfn{8)}{Voir aussi la démonstration dans la remarque I du \og complément\fg{}.} du lemme auxiliaire s'obtient en combinant les systèmes d'équations (1) et (3). Compte tenu de ce que $g_{t+1}$, conformément au lemme auxiliaire I, ne s'annule plus à l'origine, de sorte que le module $\mathfrak{q}^{(t+1)}$ est constitué de \emph{tous} les polynômes (c'est l'idéal unité), et que, par conséquent, $K_{t+1}\equiv0(\mathfrak{q}^{(t+1)})$ devient trivial, l'application $t$ fois répétée du lemme auxiliaire II donne la démonstration du théorème principal.

\begin{center}
\rule{4cm}{0.4pt}
\end{center}

Sans démonstrations --- on les trouvera dans mon mémoire cité plus haut ---, j'ajoute encore deux compléments au théorème principal :

\textbf{Complément I.} Le nombre $t$ du théorème principal est identique au plus petit\srcfn{9)}{On obtient donc ici les exposants \emph{exacts}, contrairement aux estimations de Mlle Grete Hermann [\og Die Frage der endlich vielen Schritte in der Theorie der Polynomideale\fg{} (avec utilisation de théorèmes posthumes de K. Hentzelt), \emph{Math. Annalen} 95 (1926)], où, pour une problématique certes beaucoup plus générale, seules des bornes supérieures sont indiquées.} entier positif $\nu$ tel que
\[
x^v\equiv0(\mathfrak{q}).
\]

\textbf{Complément II.} Les systèmes d'équations (1) et (3), ainsi que ceux qui en résultent par permutation de $x$ et $y$, fournissent, lorsqu'on les combine convenablement, une méthode rationnelle permettant de déterminer numériquement et exactement l'exposant $\varrho$ \og appartenant\fg{} à $\mathfrak{q}$, c'est-à-dire l'exposant minimal $\varrho$ du théorème Ia.

E. Bertini\srcfn{10)}{E. Bertini, \emph{Math. Annalen} 34 (1889).} a donné une formule générale pour $\varrho$. Le fait que cette formule ne donne pas toujours la valeur minimale de $\varrho$ est montré dans le mémoire cité sur l'exemple $\varphi=x^3+y^4$; $\psi=x^2+y^3$. Remarquons enfin que notre théorème principal fournit lui-même une nouvelle manière de justifier la formule de Bertini.

\begin{center}
\rule{4cm}{0.4pt}
\end{center}

Une \emph{spécialisation du théorème principal} digne d'intérêt résulte de l'hypothèse restrictive selon laquelle le zéro du $\mathfrak{q}=(\varphi,\psi,\mathfrak{p}^\varrho)$ considéré est un point \emph{simple} d'au moins l'une des deux courbes $\varphi=0$, $\psi=0$. Si, par exemple, ce zéro est un point simple de $\psi$, on peut montrer que le nombre $t$ du théorème principal est \emph{exactement} égal à $\mu$ et que les $t$ inégalités du théorème principal peuvent être remplacées par une \emph{unique} condition de multiplicité. On obtient ainsi le théorème général suivant :

\textbf{Théorème III.} Si \emph{tous} les points d'intersection de $\varphi=0$, $\psi=0$ sont des points \emph{simples} de $\psi$, alors --- en supposant $K$ premier à $\psi$ --- la condition nécessaire et suffisante pour que la congruence $K\equiv0(\varphi,\psi)$ soit satisfaite est que $K$ s'annule en tous ces points, mais de telle sorte que la multiplicité d'intersection dans le couple de courbes $K=0$, $\psi=0$ soit, en chacun d'eux, \emph{au moins} aussi grande que dans le couple de courbes $\varphi=0$, $\psi=0$.

Ce théorème III est encore remarquable pour une raison particulière : la condition de multiplicité qu'il exige est telle que l'on peut décider, par un nombre fini de \emph{différentiations} et indépendamment de la formule (2) fondée sur (1), si elle est satisfaite ou non. Cela apparaît immédiatement si je cite le \og théorème\fg{} suivant, que j'ai formulé en 1923\srcfn{11)}{H. Kapferer, \og Über die Multiplizität der Schnittpunkte von zwei algebraischen Kurven\fg{}, \emph{Jahresber. der D. M.-V.} 32 (1923).} :

Si $A$ et $B$ sont deux polynômes quelconques premiers entre eux en $x,y$, et si $P$ est un point commun, \emph{qui soit un point simple de $B$}, alors la multiplicité de ce point comme point d'intersection est égale à un nombre $\mu$ si et seulement si la multiplicité du même point dans le couple de courbes $A'=0$, $B=0$ est \emph{exactement} $\mu-1$; ici
\[
A'=\frac{\partial A}{\partial x}\cdot\frac{\partial B}{\partial y}
   -\frac{\partial A}{\partial y}\cdot\frac{\partial B}{\partial x}.
\]

\begin{center}
\rule{4cm}{0.4pt}
\end{center}

Mentionnons enfin que le théorème principal (théorème II), moyennant une légère modification, se transpose également aux polynômes homogènes en trois variables $x,y,z$. Cette transposition est particulièrement élégante dans le cas particulier ci-dessus. On peut en effet montrer :

\textbf{Complément au théorème III.} Le théorème III demeure entièrement inchangé, jusque dans son libellé, si l'on entend par $\varphi,\psi,K$ des polynômes \emph{homogènes} en \emph{trois} variables\srcfn{12)}{La question des conditions pour que soit satisfaite une congruence homogène en trois variables $K(x,y,z)\equiv0(\varphi(x,y,z),\psi(x,y,z))$ constitue aussi --- quoique sous des hypothèses très spécialisées --- la question centrale d'une étude d'A. Ostrowski, \og Über die Maxwellsche Erzeugung der Kugelfunktionen\fg{}, \emph{Jahresb. der D. M.-V.} 33 (1925). Dans cette note, on suppose que $\varphi$ se décompose entièrement en facteurs linéaires, que $\varphi$ et $K$ sont de même degré et que $\psi=x^2+y^2+z^2$. Puisque $\psi$ y est donc entièrement dépourvu de singularités, le théorème principal de la note d'Ostrowski peut être considéré comme un cas particulier de notre théorème III.}.

\subsection*{\NoetherNavAnchor{SUPPLEMENT}{Complément, en commun avec E. Noether}Complément, en commun avec E. Noether\srcfn{13)}{L'interprétation issue de la théorie des idéaux est due à E. Noether; la démonstration de (5) et (6) sur laquelle elle repose est due à H. Kapferer.}}
\label{fr-r823-complement-kapferer-noether}

Le système d'équations (1) fournit en même temps une méthode rationnelle de \emph{détermination d'un système complet de restes modulo $\mathfrak{q}$}; on entend par là un système complet de représentants des classes résiduelles, en nombre \emph{fini}, \emph{linéairement indépendantes} modulo $q$ relativement au domaine des coefficients $\mathsf{P}$. Les différentes étapes de la démonstration du théorème principal deviennent ainsi très transparentes; on obtient en même temps, comme résultat annexe, que la \emph{multiplicité résultantielle $\mu$ coïncide avec la longueur de l'idéal $\mathfrak{q}$};

c'est-à-dire avec le rang de l'anneau quotient, ou avec le nombre de classes résiduelles linéairement indépendantes relativement à $\mathsf{P}$\srcfn{14)}{On supposera --- sans restreindre la généralité --- que le corps $\mathsf{P}$ est algébriquement clos; dans cette mesure, tous les théorèmes précédents sont effectivement valables, à l'exception de la condition de différentiation. La longueur de $\mathfrak{q}$ est aussi définie comme la longueur d'une suite de composition d'idéaux allant de $\mathfrak{p}$ à $\mathfrak{q}$. Voir E. Noether, \emph{Jahresber. d. d. Math.-Ver.} 34 (1925), p. 101 (pagination oblique); plus en détail, H. Grell, \emph{Math. Annalen} 97 (1927), p. 529. Les démonstrations publiées de la coïncidence des deux nombres de multiplicité (pour $n$ variables) sont très compliquées; la plus simple est la démonstration posthume, encore inédite, de Hentzelt.}.

Avec les notations de (1) et du lemme auxiliaire II, le système complet de restes est caractérisé par le théorème suivant.

\textbf{Théorème IV.} \emph{Un système complet de restes modulo $\mathfrak{q}$ est donné par les polynômes}
\[
h_t,\ h_ty,\ldots,h_ty^{(g_t)-1};\ \ldots;\ h_i,\ h_iy,\ldots,h_iy^{(g_i)-1};\ \ldots;\ h_1,\ h_1y,\ldots.\ h_1y^{(g_1)-1}.
\]
\emph{Pour chaque $i$, $h_i,h_iy,\ldots,h_iy^{(g_i)-1}$ forment un système complet de restes de $\mathfrak{q}^{(i+1)}$ modulo $\mathfrak{q}^{(i)}$.}

Il résulte immédiatement de la définition que $\mathfrak{q}^{(i+1)}=(\mathfrak{q}^{(i)},h_i)$, donc que $\mathfrak{q}^{(i)}\equiv0(\mathfrak{q}^{(i+1)})$. Le théorème IV sera donc démontré dès que l'on aura établi
\begin{equation}
\tag{4}
h_i x\equiv0(\mathfrak{q}^{(i)});\qquad h_i y^{(g_i)}\equiv0(\mathfrak{q}^{(i)})
\end{equation}
et
\begin{equation}
\tag{5}
\text{de } h_iF(y)\equiv0(\mathfrak{q}^{(i)}) \text{ on déduit que } F(y)\equiv0(y^{(g_i)}).
\end{equation}

Les relations (4) expriment en effet que les combinaisons linéaires de $h_i,h_iy,\ldots,h_iy^{(g_i)-1}$ épuisent les classes résiduelles de $\mathfrak{q}^{(i+1)}$ modulo $\mathfrak{q}^{(i)}$, tandis que (5) exprime leur indépendance linéaire. Comme $\mathfrak{q}^{(t+1)}$ devient l'idéal unité --- ce qui résulte maintenant de la longueur finie de $\mathfrak{q}$ ---, on en déduit la démonstration du théorème IV et, en même temps, le fait que la longueur de $\mathfrak{q}$ est égale à $(g_1)+\ldots+(g_t)$, et coïncide donc, d'après (2), avec la multiplicité résultantielle.

La première des relations (4) se lit dans (1); la seconde s'obtient comme suit.

On a $h_ig_i\equiv0(\mathfrak{q}^{(i)})$; donc, en vertu de la première relation (4), aussi $h_ig_i(0,y)\equiv0(\mathfrak{q}^{(i)})$, ou
$h_iy^{(g_i)}\cdot(g_i(0,y):y^{(g_i)})\equiv0(\mathfrak{q}^{(i)})$; puisque $g_i(0,y):y^{(g_i)}\not\equiv0(y)$, il s'ensuit que $h_iy^{(g_i)}\equiv0(\mathfrak{q}^{(i)})$.

Avant de démontrer (5), faisons la \emph{remarque} suivante, qui se lit immédiatement dans (1) : $f_i$ et $g_i$ ne peuvent avoir pour diviseur commun qu'un polynôme $d(y)$ en $y$ qui n'est pas divisible par $y$ :
$f_i=d_i(y)\bar f_i$, $g_i=d_i(y)\bar g_i$, où $\bar f_i$ et $\bar g_i$ sont premiers entre eux. Puisque $d(y)\not\equiv0(\mathfrak{p})$, on a donc aussi $q^{(i)}=(\bar f_i,\bar g_i,\mathfrak{p}^\varrho)$; cet idéal est égal à $(\bar f_i,\bar g_i,\mathfrak{p}^\sigma)$ pour tout $\sigma\geqq\varrho$, en tant que diviseur de $\mathfrak{q}$ pour tout tel $\sigma$; ainsi $\mathfrak{q}^{(i)}$ est la composante primaire de $(\bar f_i,\bar g_i)$ correspondant à l'origine.

Soient $\gamma_1,\delta_1;\ldots;\gamma_r,\delta_r$ les zéros de $(\bar f_i,\bar g_i)$ distincts de $x=0$, $y=0$, et soient $\tau_1,\ldots,\tau_r$ les exposants des composantes primaires correspondantes de $(\bar f_i,\bar g_i)$; dans le produit
\[
P=((x-\gamma_1)+u(y-\delta_1))^{\tau_1}\ldots
((x-\gamma_r)+u(y-\delta_r))^{\tau_r},
\]
spécialisons l'indéterminée $u$ en un élément de $\mathsf{P}$ de telle sorte que $P\not\equiv0(\mathfrak{p})$. Si l'on pose $G(x,y)$ égal au produit du $P$ spécialisé par $d(y)$, on a donc $G(x,y)\not\equiv0(\mathfrak{p})$ et $\mathfrak{q}^{(i)}\cdot G(x,y)\equiv0(f_i,g_i)$.

De l'hypothèse $h_iF(y)\equiv0(\mathfrak{q}^{(i)})$, il résulte donc que $h_iF(y)G(x,y)\equiv0(f_i,g_i)$, ou --- conjointement avec (1) --- sous forme d'identités :
\[
h_iF(y)G=Af_i+Bg_i;\qquad
h_ix=(g_i(0,y):y^{(g_i)})\cdot f_i-(f(0,y):y^{(g_i)})g_i;
\]
on en déduit
\[
\bar f_i(xA-F(y)H)\equiv0(\bar g_i);\qquad
H=G(x,y)\cdot g_i(0,y):y^{(g_i)}\not\equiv0(\mathfrak{p}).
\]

Puisque $\bar f_i$ et $\bar g_i$ sont premiers entre eux, il s'ensuit encore que
\[
xA-F(y)H\equiv0(\bar g_i);
\]
donc $F(y)H\equiv0(\bar g_i,x)\equiv0(\bar g_i,x,\mathfrak{p}^\sigma)=(y^{(g_i)},x)$, et par conséquent $F(y)\equiv0(y^{(g_i)})$. Ainsi (5), et donc le théorème IV, sont démontrés.

\textbf{Remarque I.} Les relations (4) et (5), conjointement avec le lemme auxiliaire II, peuvent s'écrire sous forme de relations réciproques entre quotients d'idéaux\srcfn{15)}{Par le quotient $\mathfrak a:\mathfrak b$, on entend, comme on sait, l'ensemble des polynômes $c(x,y)$ tels que $\mathfrak b c\equiv0(\mathfrak a)$.}, la démonstration de la seconde relation fournissant en même temps une démonstration du lemme auxiliaire II :
\begin{equation}
\tag{6}
\mathfrak{q}^{(i)}:\mathfrak{q}^{(i+1)}=(\mathfrak{q}^{(i)},x)=(y^{(g_i)},x);\qquad
\mathfrak{q}^{(i)}:(\mathfrak{q}^{(i)},x)=\mathfrak{q}^{(i+1)}.\srcfnmark{16)}
\end{equation}
\srcfntext{16)}{Par la théorie générale des idéaux irréductibles, chacune de ces deux relations réciproques entraîne l'autre; voir Macaulay, \emph{Modular Systems}, nos 72 et 73, \emph{Cambridge Tracts} 19 (1916).}

La première relation résume (4) et (5). Quant à la seconde, remarquons que $x\cdot \mathfrak{q}^{(i+1)}\equiv0(\mathfrak{q}^{(i)})$ a été démontré dans (4); la réciproque s'obtient de manière analogue à (5). De $xM\equiv0(\mathfrak{q}^{(i)})$, on déduit, comme ci-dessus,
\[
xMG\equiv0(f_i,g_i);\quad\text{donc}\quad
xMG\cdot g_i(0,y):y^{(g_i)}\equiv0(xh_i,g_i).
\]

Puisque, par définition, $g_i\not\equiv0(x)$ --- $(f_i)\geqq(g_i)$ ---, on en tire
\[
M\cdot(G\cdot g_i(0,y):y^{(g_i)})\equiv0(h_i,g_i)\equiv0(\mathfrak{q}^{(i+1)}),
\quad\text{donc}\quad M\equiv0(\mathfrak{q}^{(i+1)}).
\]

La nécessité de la condition du lemme auxiliaire II s'établit maintenant ainsi : $K_i\equiv0(\mathfrak{q}_i)$ entraîne $K_i(0,y)\equiv0(y^{(g_i)},x)$, donc $(K_i)\geqq(g_i)$ et, par (3),
\[
xK_{i+1}\equiv0(\mathfrak{q}^{(i)});\quad\text{donc, par la seconde relation (6),}\quad
K_{i+1}\equiv0(\mathfrak{q}^{(i+1)}).
\]

La suffisance de la condition se lit immédiatement; l'application itérée un nombre fini de fois du lemme auxiliaire II donne alors le théorème principal. De même, l'itération de la seconde relation (6) donne le complément I, à savoir $\mathfrak{q}:x^t=\mathfrak{q}^{(t+1)}$, donc l'idéal unité; $t$ est bien la plus petite puissance ayant cette propriété, puisque $\mathfrak{q}:x^{t-1}=\mathfrak{q}^{(t)}$.

\textbf{Remarque II.} Pour un polynôme $K$ qui n'est pas divisible par $\mathfrak{q}$, on peut encore affirmer ceci : si $K\equiv0(\mathfrak{q}^{(i+1)})$, $\not\equiv0(\mathfrak{q}^{(i)})$, alors $(K,\mathfrak{q}^{(i)})=(h_iy^{(K)},\mathfrak{q}^{(i)})$. En effet, d'après le théorème IV, $K\equiv h_iy^\lambda c(y)$ modulo $(\mathfrak{q}^{(i)})$, avec $c(y)\not\equiv0(y)$; il vient donc
\[
(K,\mathfrak{q}^{(i)})=(h_iy^\lambda,\mathfrak{q}^{(i)}).
\]

On a alors $y^{(g_i)-\lambda}K\equiv0(\mathfrak{q}^{(i)})$, et, d'après (5), il s'agit de la plus petite puissance ayant cette propriété. Or, d'après le lemme auxiliaire II, $(g_i)-(K)$ est cette plus petite puissance; ainsi $\lambda=(K)$.

\begin{center}
(Reçu le 24 octobre 1926.)
\end{center}

\clearpage

\clearpage
\begin{center}
{\Large\bfseries \NoetherNavAnchor{BIBLIOGRAPHIE}{Bibliographie}Bibliographie}
\end{center}
\label{fr-r823-bibliographie}

\begin{enumerate}
\setlength{\itemsep}{0.35em}
\item Sur la formation du système de formes de la forme biquadratique ternaire.\\
\emph{Sitz. Ber. d. Physikal.-mediz. Sozietät in Erlangen} 39 (1907), p. 176--179
\item Sur la formation du système de formes de la forme biquadratique ternaire.\\
\emph{J. Reine Angew. Math.} 134 (1908), p. 23--90 et deux tableaux
\item Sur la théorie des invariants des formes à $n$ variables.\\
\emph{J. Ber. d. DMV} 19 (1910), p. 101--104
\item Sur la théorie des invariants des formes à $n$ variables.\\
\emph{J. Reine Angew. Math.} 139 (1911), p. 118--154
\item Corps de fonctions rationnelles.\\
\emph{J. Ber. d. DMV} 22 (1913), p. 316--319
\item Corps et systèmes de fonctions rationnelles.\\
\emph{Math. Ann.} 76 (1915), p. 161--196
\item Le théorème de finitude des invariants des groupes finis.\\
\emph{Math. Ann.} 77 (1916), p. 89--92
\item Sur la représentation rationnelle entière des invariants d'un système comprenant un nombre arbitraire de formes fondamentales.\\
\emph{Math. Ann.} 77 (1916), p. 93--102. (Voir aussi le no 16.)
\item Les domaines les plus généraux de nombres transcendants entiers.\\
\emph{Math. Ann.} 77 (1916), p. 103--128. (Voir aussi le no 16.)
\item Les équations fonctionnelles de l'application isomorphe.\\
\emph{Math. Ann.} 77 (1916), p. 536--545
\item Équations à groupe prescrit.\\
\emph{Math. Ann.} 78 (1918), p. 221--229. (Voir aussi le no 16.)
\item Invariants d'expressions différentielles arbitraires.\\
\emph{Nachr. v. d. Ges. d. Wiss. zu Göttingen} 1918, p. 37--44
\item Problèmes variationnels invariants.\\
\emph{Nachr. v. d. Ges. d. Wiss. zu Göttingen} 1918, p. 235--257
\item La théorie arithmétique des fonctions algébriques d'une variable, dans ses rapports avec les autres théories et avec la théorie des corps de nombres.\\
\emph{J. Ber. d. DMV} 28 (1919), p. 182--203
\item La finitude du système des invariants entiers des formes binaires.\\
\emph{Nachr. v. d. Ges. d. Wiss. zu Göttingen} 1919, p. 138--156
\item Sur le développement en série dans la théorie des formes.\\
\emph{Math. Ann.} 81 (1920), p. 25--30
\item Avec W. Schmeidler : Modules dans des domaines non commutatifs, notamment ceux formés d'expressions différentielles et aux différences.\\
\emph{Math. Z.} 8 (1920), p. 1--35
\item Sur un travail de K. Hentzelt, tombé à la guerre, consacré à la théorie de l'élimination.\\
\emph{J. Ber. d. DMV} 30 (1921), p. 101
\item Théorie des idéaux dans les anneaux.\\
\emph{Math. Ann.} 83 (1921), p. 24--66
\item Un critère algébrique d'irréductibilité absolue.\\
\emph{Math. Ann.} 85 (1922), p. 26--33
\item Calcul formel des variations et invariants différentiels.\\
\emph{Encyklopädie d. math. Wiss.} III, 3 (1922), p. 68--71 (dans R. Weitzenböck, \emph{Differentialinvarianten})
\item Édition de K. Hentzelt $\dagger$ : Sur la théorie des idéaux de polynômes et des résultants.\\
\emph{Math. Ann.} 88 (1923), p. 53--79
\item Invariants algébriques et différentiels.\\
\emph{J. Ber. d. DMV} 32 (1923), p. 177--184
\item Théorie de l'élimination et théorie générale des idéaux.\\
\emph{Math. Ann.} 90 (1923), p. 229--261
\item Théorie de l'élimination et théorie des idéaux.\\
\emph{J. Ber. d. DMV} 33 (1924), p. 116--120
\item Construction abstraite de la théorie des idéaux dans un corps de nombres algébriques.\\
\emph{J. Ber. d. DMV} 33 (1924), p. 102
\item Nombres de Hilbert dans la théorie des idéaux.\\
\emph{J. Ber. d. DMV} 34 (1925), p. 101
\item Caractères de groupes et théorie des idéaux.\\
\emph{J. Ber. d. DMV} 34 (1925), p. 144
\item Le théorème de finitude des invariants des groupes linéaires finis de caractéristique $p$.\\
\emph{Nachr. v. d. Ges. d. Wiss. zu Göttingen} 1926, p. 28--35
\item Construction abstraite de la théorie des idéaux dans les corps algébriques de nombres et de fonctions.\\
\emph{Math. Ann.} 96 (1927), p. 26--61
\item Le théorème du discriminant pour les ordres d'un corps algébrique de nombres ou de fonctions.\\
\emph{J. Reine Angew. Math.} 157 (1927), p. 82--104
\item Avec R. Brauer : Sur les corps de décomposition minimaux des représentations irréductibles.\\
\emph{Sitz. Ber. d. Preuß. Akad. d. Wiss.} 1927, p. 221--228
\item Grandeurs hypercomplexes et théorie des représentations dans une conception arithmétique.\\
\emph{Atti Congresso Bologna} 2 (1928), p. 71--73
\item Grandeurs hypercomplexes et théorie des représentations.\\
\emph{Math. Z.} 30 (1929), p. 641--692
\item Sur les domaines maximaux de fonctions entières.\\
\emph{Rec. Soc. Math. Moscou} 36 (1929), p. 65--72
\item Différentiation des idéaux et différente.\\
\emph{J. Ber. d. DMV} 39 (1930), p. 17
\item Base normale dans les corps sans ramification supérieure.\\
\emph{J. Reine Angew. Math.} 167 (1932), p. 147--152
\item Avec R. Brauer et H. Hasse : Démonstration d'un théorème principal de la théorie des algèbres.\\
\emph{J. Reine Angew. Math.} 167 (1932), p. 399--404
\item Systèmes hypercomplexes dans leurs rapports avec l'algèbre commutative et la théorie des nombres.\\
\emph{Verhandl. Intern. Math.-Kongreß Zürich} 1 (1932), p. 189--194
\item Algèbre non commutative.\\
\emph{Math. Z.} 37 (1933), p. 514--541
\item Le théorème du genre principal pour les corps de nombres relativement galoisiens.\\
\emph{Math. Ann.} 108 (1933), p. 411--419
\item Produits croisés scindés et leurs ordres maximaux.\\
\emph{Actualités scientifiques et industrielles} 148 (1934), 15 p.
\item Différentiation des idéaux et différente.\\
\emph{J. Reine Angew. Math.} 188 (1950), p. 1--21
\end{enumerate}

\noindent Conditions de multiplicité nécessaires et suffisantes pour le théorème fondamental de Noether sur les fonctions algébriques, par H. Kapferer\\
(avec un complément, en commun avec E. Noether).\\
\emph{Math. Ann.} 97 (1927), p. 559--567

\clearpage
\begin{center}
{\Large\bfseries Liste des communications brèves et des comptes rendus d'Emmy Noether}
\end{center}
\label{fr-r823-liste-communications-comptes-rendus}

\begin{center}
\emph{\NoetherNavAnchor{NOTICES}{Communications brèves parues dans les Jahresberichte der Deutschen Mathematiker-Vereinigung}Communications brèves parues dans les Jahresberichte der Deutschen Mathematiker-Vereinigung}
\end{center}

\begingroup
\setlength{\parindent}{0pt}
\setlength{\parskip}{0.15em}
Sur les invariants d'expressions différentielles arbitraires, 27, partie II (1918), p. 28\par
Problèmes variationnels invariants, 27, partie II (1918), p. 47\par
Finitude des invariants binaires entiers, 28, partie II (1919), p. 29\par
Sur les polynômes entiers et les séries de puissances, 28, partie II (1919), p. 29--30\par
Théorèmes de décomposition de la théorie des modules, 29, partie II (1920), p. 54\par
Composantes élémentaires et théorie générale des idéaux, 30, partie II (1921), p. 32\par
Remarque sur une formule de Lipschitz, 30, partie II (1921), p. 48\par
L'analogue d'un théorème de Hilbert dans la théorie générale des idéaux, 32, partie II (1923), p. 20--21\par
Remarques de théorie des groupes sur le problème de l'axiomatique, 32, partie II (1923), p. 21\par
Théorie de Galois dans des corps arbitraires, 33, partie II (1925), p. 120\par
Nombres de Hilbert dans la théorie des idéaux, 34, partie II (1926), p. 101\par
Quotients différentiels d'idéaux et théorie de la ramification, 38, partie II (1929), p. 81\par
\endgroup

\begin{center}
\emph{\NoetherNavAnchor{REVIEWS}{Comptes rendus parus dans les Jahresberichte der Deutschen Mathematiker-Vereinigung}Comptes rendus parus dans les Jahresberichte der Deutschen Mathematiker-Vereinigung}
\end{center}

\begingroup
\setlength{\parindent}{0pt}
\setlength{\parskip}{0.35em}
E. Study, \emph{Einleitung in die Theorie der Invarianten linearer Transformationen auf Grund der Vektorenrechnung}. Première partie. (\emph{Die Wissenschaft}, vol. 71.) Brunswick, 1923, Vieweg.\\
36, partie II (1927), p. 71\par
\emph{Leopold Kroneckers Werke}, vol. 4. Édité par K. Hensel. VI et 509 p. Leipzig et Berlin, 1929, B. G. Teubner.\\
40, partie II (1931), p. 11--12\par
\emph{Leopold Kroneckers Werke}, vol. 5. Édité par K. Hensel. 527 p. Leipzig et Berlin, 1930, B. G. Teubner.\\
41, partie II (1932), p. 17\par
\emph{Leopold Kroneckers Werke}, vol. III.2. Édité par K. Hensel. 215 p. Leipzig et Berlin, 1931, B. G. Teubner.\\
42, partie II (1933), p. 38--39\par
D. E. Rutherford, \emph{Modular Invariants}. \emph{Cambridge Tracts} no 27. Cambridge, 1932, VII et 84 p.\\
43, partie II (1934), p. 94--95\par
\endgroup

\clearpage
\begin{center}
\emph{\NoetherNavAnchor{LIVRES}{Livres réalisés avec la participation d’Emmy Noether}Livres réalisés avec la participation d'Emmy Noether}
\end{center}
\label{fr-r823-livres-participation-emmy-noether}

\begingroup
\setlength{\parindent}{0pt}
\setlength{\parskip}{0.5em}
R. Dedekind, \emph{Gesammelte mathematische Werke}, I--III. Édité par R. Fricke, Emmy Noether et Øystein Ore. Vieweg, Brunswick, 1930--1932\\
(vol. I : 1930; II : 1931; III : 1932)\par
G. Cantor, \emph{Briefwechsel mit R. Dedekind}. Édité par Emmy Noether et J. Cavaillès, 60 p. \emph{Actualités scientifiques et industrielles} 581, Paris, 1937\par
\endgroup
\end{document}
